\documentclass[11pt]{article}

\usepackage[T1]{fontenc}
\usepackage[utf8]{inputenc}
\usepackage{lmodern}
\usepackage{microtype}
\usepackage[margin=1.05in]{geometry}

\usepackage{amsmath,amssymb,amsthm,mathtools,aliascnt}
\usepackage{bm,mathrsfs}
\usepackage{booktabs,longtable,array}
\usepackage{enumitem}
\usepackage{tocloft}
\usepackage{xspace}
\usepackage{xcolor}
\usepackage{tikz}
\usetikzlibrary{arrows.meta,positioning,fit,calc,shapes.geometric}
\usepackage{pgfplots}
\usepgfplotslibrary{fillbetween}
\pgfplotsset{compat=1.18}
\usepackage[numbers,sort&compress]{natbib}
\usepackage[hidelinks,unicode]{hyperref}

\allowdisplaybreaks
\emergencystretch=2em
\setcounter{secnumdepth}{3}
\setcounter{tocdepth}{3}
\setlist{nosep,leftmargin=*}
\setlength{\cftsecnumwidth}{4em}
\setlength{\cftsubsecnumwidth}{5em}
\setlength{\cftsubsubsecnumwidth}{6em}
\providecommand{\tightlist}{%
  \setlength{\itemsep}{0pt}\setlength{\parskip}{0pt}}

\numberwithin{section}{part}

\theoremstyle{plain}
\newtheorem{theorem}{Theorem}[section]
\newaliascnt{proposition}{theorem}
\newtheorem{proposition}[proposition]{Proposition}
\aliascntresetthe{proposition}
\newaliascnt{lemma}{theorem}
\newtheorem{lemma}[lemma]{Lemma}
\aliascntresetthe{lemma}
\newaliascnt{corollary}{theorem}
\newtheorem{corollary}[corollary]{Corollary}
\aliascntresetthe{corollary}

\theoremstyle{definition}
\newaliascnt{definition}{theorem}
\newtheorem{definition}[definition]{Definition}
\aliascntresetthe{definition}
\newaliascnt{assumption}{theorem}
\newtheorem{assumption}[assumption]{Assumption}
\aliascntresetthe{assumption}
\newaliascnt{convention}{theorem}
\newtheorem{convention}[convention]{Convention}
\aliascntresetthe{convention}
\newaliascnt{example}{theorem}
\newtheorem{example}[example]{Example}
\aliascntresetthe{example}

\theoremstyle{remark}
\newaliascnt{remark}{theorem}
\newtheorem{remark}[remark]{Remark}
\aliascntresetthe{remark}

\usepackage[nameinlink,capitalise]{cleveref}
\crefname{part}{Part}{Parts}
\crefname{theorem}{Theorem}{Theorems}
\crefname{proposition}{Proposition}{Propositions}
\crefname{lemma}{Lemma}{Lemmas}
\crefname{corollary}{Corollary}{Corollaries}
\crefname{definition}{Definition}{Definitions}
\crefname{assumption}{Assumption}{Assumptions}
\crefname{convention}{Convention}{Conventions}
\crefname{example}{Example}{Examples}
\crefname{remark}{Remark}{Remarks}

\newcommand{\given}{\,\middle|\,}
\newcommand{\sem}[1]{\lbrack\!\lbrack #1\rbrack\!\rbrack}
\newcommand{\satcost}[1]{\operatorname{Cost}_{\mathrm{sat},n}(#1)}
\newcommand{\satprofile}[2]{G_{#1,n}^{\mathrm{sat}}(#2)}
\newcommand{\visprofile}[2]{m_{#1,n}^{\mathrm{sat}}(#2)}
\newcommand{\obsalg}{\mathcal F}
\newcommand{\obsalgpoly}[1]{\mathcal F_{#1}^{\mathrm{poly}}}
\newcommand{\Geom}{\mathsf{Geom}}
\newcommand{\Caus}{\mathsf{Caus}}
\newcommand{\Abs}{\mathsf{Abs}}
\newcommand{\Lift}{\mathsf{Lift}}
\newcommand{\Bdry}{\mathsf{Bdry}}
\newcommand{\Res}{J_{\mathrm{res}}}
\newcommand{\PartI}{\hyperref[part:task-spaces]{Part~I}\xspace}
\newcommand{\PartII}{\hyperref[part:dynamics]{Part~II}\xspace}
\newcommand{\PartIII}{\hyperref[part:spectral-complexity]{Part~III}\xspace}
\newcommand{\PartIV}{\hyperref[part:learning]{Part~IV}\xspace}
\newcommand{\PartV}{\hyperref[part:dynamical-systems]{Part~V}\xspace}
\newcommand{\PartVI}{\hyperref[part:attention]{Part~VI}\xspace}
\newcommand{\PartVII}{\hyperref[part:structural-metacomplexity]{Part~VII}\xspace}
\newcommand{\PartVIII}{\hyperref[part:lpn]{Part~VIII}\xspace}
\newcommand{\PartsItoII}{\hyperref[part:task-spaces]{Parts~I}--\hyperref[part:dynamics]{II}\xspace}
\newcommand{\PartsItoIII}{\hyperref[part:task-spaces]{Parts~I}--\hyperref[part:spectral-complexity]{III}\xspace}
\DeclareMathOperator{\Var}{Var}
\DeclareMathOperator{\Cov}{Cov}
\DeclareMathOperator{\supp}{supp}
\DeclareMathOperator{\rank}{rank}
\DeclareMathOperator{\diag}{diag}

\definecolor{fwInk}{HTML}{2F3E46}
\definecolor{fwBlue}{HTML}{4C78A8}
\definecolor{fwTeal}{HTML}{4C9F92}
\definecolor{fwAmber}{HTML}{C58B2A}
\definecolor{fwCoral}{HTML}{B9655A}
\definecolor{fwViolet}{HTML}{8064A2}
\definecolor{fwGray}{HTML}{E9ECEF}

\tikzset{
  fw diagram/.style={>=Stealth,font=\small,line width=0.75pt},
  fw box/.style={draw=fwInk,rounded corners=2pt,align=center,
    minimum height=7mm,inner xsep=5pt,inner ysep=3pt,fill=white},
  fw source/.style={fw box,draw=fwBlue,fill=fwBlue!13},
  fw geom/.style={fw box,draw=fwTeal,fill=fwTeal!14},
  fw use/.style={fw box,draw=fwAmber,fill=fwAmber!14},
  fw provenance/.style={fw box,draw=fwViolet,fill=fwViolet!12},
  fw terminal/.style={fw box,draw=fwCoral,fill=fwCoral!12},
  fw residual/.style={fw box,draw=fwInk!65,fill=fwGray},
  fw arrow/.style={-{Stealth[length=2.2mm,width=1.5mm]},draw=fwInk},
  fw dashed arrow/.style={fw arrow,dashed},
  fw group/.style={draw=fwInk!55,dashed,rounded corners=3pt,inner sep=5pt}
}

\pgfplotsset{
  fw axis/.style={
    width=.92\linewidth,
    height=6.2cm,
    axis line style={draw=fwInk},
    tick style={draw=fwInk},
    tick label style={font=\scriptsize},
    label style={font=\small},
    legend style={draw=none,fill=none,font=\scriptsize,cells={anchor=west}},
    grid=major,
    grid style={draw=fwGray}
  }
}

\title{Presentation-Relative Structural Complexity:\\
Quantitative Limits for General Intelligent Computation}
\author{Guillem Duran Ballester\\
\href{mailto:guillem@fragile.tech}{\texttt{guillem@fragile.tech}}}
\date{}

\hypersetup{
  pdftitle={Presentation-Relative Structural Complexity: Quantitative Limits for General Intelligent Computation},
  pdfauthor={Guillem Duran Ballester}
}

\begin{document}

\maketitle

\begin{abstract}
Can the capabilities of an arbitrary intelligent system be bounded without
fixing its architecture?  We develop a presentation-relative theory that
replaces algorithm-specific analysis with a saturated structural profile: the
strongest target-relevant structure that any resource-bounded computation can
construct from the public observables, constructors, terminal data, and cost
model of a task.  Universal accountability theorems convert this profile into
quantitative performance bounds that hold simultaneously for every admissible
algorithm, including architectures and runtime-generated representations not
anticipated in advance.

The framework unifies represented geometry, abstraction, learned features,
information acquisition, memory, Bayesian inference, control, reinforcement
learning, and attention within a common cost-, provenance-, temporal-,
source-, and noise-aware ledger.  It separates prospective sources from
retrospective explanations and charges generated representations at their
complete saturated cost.  Source-resolved Rademacher, sequential Rademacher,
PAC--Bayes, and effective-dimension bounds keep statistical capacity,
structural ancestry, noise contraction, construction cost, and deployment
attached to the same computation.  The resulting statistical, computational,
and information-theoretic noise thresholds quantify transitions among
learnable, resource-limited, and information-limited regimes.

Universal quantification over the declared resource class gives precise
architecture-independent bounds on learning, reasoning, control, and
target-directed performance for increasingly general intelligent systems,
including artificial general intelligence.  As a proof-producing stress
test, the LPN specialization isolates one compact
residual-ancestry certification target and proves that a deterministically
accepted finite-to-asymptotic certificate discharges it.  The instantiated
certificate analysis is not executed here, so the current LPN hardness
consequences remain conditional on this internal certificate; verifying it
would make them unconditional relative to the declared presentation.
\end{abstract}

\begin{center}
\fbox{%
  \begin{minipage}{0.91\linewidth}
  \centering
  \textbf{Universal capability bound.}
  For every declared task presentation \(\mathfrak M\), scale \(n\), and
  resource budget \(B\), the framework has the schematic form
  \[
  \operatorname{Performance}_{\mathfrak M,n}(B)
  \;\le\;
  \operatorname{Overhead}(B,n)\,
  \operatorname{SaturatedProfile}_{\mathfrak M,n}
  \bigl(\Psi(B,n)\bigr).
  \]
  Here the left-hand side is optimized over every admissible algorithm in the
  declared resource class.  The right-hand side depends only on the task
  presentation and the cost-filtered closure of all structure derivable from
  it.  Consequently, the bound applies to arbitrary architectures without
  enumerating them or anticipating how they represent, search, learn, remember,
  attend, or control.  Advice, target information unavailable through the
  declared presentation, structures above the transformed budget
  \(\Psi(B,n)\), and enriched interfaces are assigned to distinct accounting
  classes rather than included in the saturated profile without charge.
  \end{minipage}%
}
\end{center}

\tableofcontents

\clearpage
\part{Task Spaces and Observable Structure}\label{part:task-spaces}\setcounter{section}{-1}

\section{Introduction: Static Task
Presentations}\label{sec-static-introduction}
\label{introduction-the-task-before-the-solution}

\paragraph{Part contract.}
\emph{Inputs:} only the declared data of a task presentation.
\emph{Output:} the static observable and target-relevance structures used in
\cref{part:dynamics,part:spectral-complexity}.
\emph{First pass:} the examples and audit procedures may be skimmed after the
definitions they illustrate.

Most performance analyses begin with a chosen computational method.  They
inspect its states, representations, and transitions, then derive a bound for
that method:
\[
\text{choose an algorithm}
\;\longrightarrow\;
\text{analyze its computation}
\;\longrightarrow\;
\text{bound its performance}.
\]
Such a bound is tied to the analyzed design.  A different architecture,
representation, search rule, or controller requires a new analysis unless the
argument identifies an invariant shared by the entire admissible resource
class.

This paper begins instead with the information and operations supplied by the
task.  Across theoretical computer science, optimization, machine learning,
and control, these data comprise a state space, a public observable interface,
admissible constructors, resource accounting, and terminal criteria
distinguishing success from failure.  The analysis follows the reverse order:
\[
\begin{gathered}
\text{specify the public task presentation}
\;\longrightarrow\;
\text{close all affordable structure}\\[2pt]
\longrightarrow\;
\text{measure the saturated target-relevant profile}
\;\longrightarrow\;
\text{bound every admissible computation}.
\end{gathered}
\]
The resulting bound is architecture-independent within the declared
presentation and resource class.  Its defining closure property is summarized
in the following budget-level statement.

\begin{center}
\fbox{%
  \begin{minipage}{0.91\linewidth}
  \raggedright
  \textbf{Saturation in one paragraph.}
  At budget \(B\), saturation closes the public presentation under every
  admissible oracle-free constructor and represented computation whose
  complete derivation is admissible at \(B\); it does not select a preferred
  algorithm.  If an algorithm generates a quotient, feature map, potential,
  kernel, controller, value representation, or target-aligned observable at
  runtime, the complete generated record re-enters the same closure with its
  derivation and admissible budget set.  Under scalar cost accounting, it
  first appears at its least saturated representation cost and remains outside
  the budget-\(B\) slice while that cost exceeds \(B\).  For a partially
  ordered resource budget, it contributes precisely at the budgets in its
  represented admissible budget set.  The resulting profile is therefore
  intrinsic to the presentation and budget, rather than to the method that
  discovered the representation.
  \end{minipage}%
}
\end{center}

The exact eight-state calculation in
\cref{ex-eight-state-affordable-algebra} instantiates this closure at a
finite budget: it lists the entire affordable quotient slice, computes every
conditional-algebra increment, and turns their conserved charge into a bound
for every admissible transition rule in the example.

\subsection{What this paper introduces}\label{subsec-overview-general}

Stated in classical terms, the manuscript develops a single quantitative
theory around three objects: a canonical decomposition of sub-Markov
generators, a cost filtration on an algebra of observables, and the
hitting-time analysis of a distinguished target set. No familiarity with
the internal terminology of the framework is needed to read this
subsection; the remainder of the paper makes each ingredient precise.

\emph{The presented object.} A computational or learning problem is
recorded as a measurable state space together with a distinguished
unital algebra of bounded observables, a slot for state-to-state
structure, and a pair of terminal events (success and failure). This is
the common measurable core of decision problems, optimization
landscapes, statistical models, and controlled systems; it fixes what is
public about a problem before any algorithm is chosen.

\emph{The generator decomposition.} Any search procedure, algorithm, or
stochastic dynamics acting on such a space and respecting a stated
resource budget induces a positive sub-Markov operator, hence, by the
positive maximum principle, a generator. The first structural theorem is
an exhaustive decomposition of every such generator into five
components with distinct classical meanings---a reversible (Dirichlet-form)
part, a drift part directed by a declared potential, a
lumping part in the sense of Markov-chain lumpability (quotients of the
state space), a state-augmentation part (lifts to enlarged carriers), and a
killing/boundary part---plus a residual term that provably carries no
structure attributable to the five listed sources. Normal forms are
proved for each component, and the decomposition is stable under
composition and under valid lifts.

\emph{The cost filtration.} The observable algebra carries a filtration
indexed by a nonnegative cost: an object lies at level \(B\) when it
admits a public derivation of total cost at most \(B\) from the declared
data. The value of this filtration at the working budget---the
\emph{saturated} cost---measures how much structure is legally available
to any procedure operating within that budget. All later upper and lower
bounds are stated against this filtration rather than against a machine
model.

\emph{Hitting times and hardness.} Complexity is then measured by the
discounted hitting functional of the target set, equivalently by the
target-projected resolvent of the killed generator. The central
inequality bounds this functional by the structure extractable at the
saturated budget: each unit of discounted target mass is charged, via an
exact first-hit decomposition, to a specific affordable component of the
generator. Hardness statements take the form ``the saturated profile is
negligible, hence the hitting functional is negligible,'' and are
presentation-relative: they quantify over every procedure representable
within the declared budget, not over a fixed algorithm.

\emph{Consequences.} Because the saturated filtration is a
reproducing-kernel object, the same ledger yields effective-dimension
and two-sided generalization bounds for kernel and neural-network
learners, including under label noise and dependent (mixing or adaptive)
sampling. For controlled Markov systems and partially observed models,
the decomposition of the selected generator produces a coordinatewise
certificate covering committors, reach--avoid probabilities, Bellman
error, and learned policies. Softmax attention is identified with a
Gibbs-measure selector whose normalization satisfies exact
Kullback--Leibler and Fisher--Rao identities, so trained attention
layers are audited by the same accounting. A meta-complexity part
inverts the picture: strict Pareto frontiers record the least resource
vectors at which a certificate of prescribed quality exists, are exact
generalized inverses of the saturated profiles, and specialize to
circuit and universal-program minimization. Finally, the framework is
instantiated on Learning Parity with Noise: after all geometric,
quotient, and provenance sources are exhausted, the surviving
coordinate of the saturated profile is isolated explicitly.  At fixed
positive noise, every materialized located spectral block has exponentially
small alignment, while compact non-lumpable residual-ancestry records form
one internal certification target.  The paper proves a deterministic
finite-to-asymptotic verification theorem: exact source coverage, certified
capacity and attenuation, outward-rounded error bounds, symbolic uniform
scale laws, a crossover check, and a finite-prefix check together imply the
required bound \((1-\eta)^n\).  The present manuscript does not execute that
final certificate analysis.  Consequently, its fixed-noise hardness theorem
and the resulting \(P\ne NP\) consequence remain conditional on acceptance of
this explicit proof certificate.  An accepted certificate would discharge
the sole remaining premise internally; no external LPN or
complexity-theoretic hardness assumption is used.

\Cref{fig-intro-channel-dictionary} summarizes the generator
decomposition together with its classical reading; the internal channel
names used throughout the paper appear in the left column.

\begin{figure}[tbp]
\centering
\begin{tikzpicture}[fw diagram,node distance=3.2mm and 9mm]
  \node[fw box,minimum width=3.1cm,minimum height=9mm,font=\small\bfseries]
    (gen) {represented\\generator \(L\)};
  \node[fw geom,right=18mm of gen,yshift=27.5mm,text width=5.9cm]
    (geom) {\(\Geom\): reversible part\\
      {\scriptsize symmetric Dirichlet form; detailed balance}};
  \node[fw use,below=of geom,text width=5.9cm]
    (caus) {\(\Caus\): directed drift\\
      {\scriptsize declared potential; authorized current}};
  \node[fw source,below=of caus,text width=5.9cm]
    (abs) {\(\Abs\): quotient part\\
      {\scriptsize Markov lumpability; coarse-graining}};
  \node[fw provenance,below=of abs,text width=5.9cm]
    (lift) {\(\Lift\): lifted-state part\\
      {\scriptsize state augmentation on enlarged carriers}};
  \node[fw terminal,below=of lift,text width=5.9cm]
    (bdry) {\(\Bdry\): boundary part\\
      {\scriptsize killing, absorption, terminal readout}};
  \node[fw box,draw=fwInk!65,fill=fwGray,below=of bdry,text width=5.9cm]
    (res) {\(\Res\): residual\\
      {\scriptsize transport without structural provenance}};
  \foreach \t in {geom,caus,abs,lift,bdry,res}
    \draw[fw arrow] (gen.east) -- (\t.west);
\end{tikzpicture}
\caption{Dictionary for the exhaustive generator decomposition
(\cref{thm-channel-exhaustiveness}). Every budget-admissible positive
sub-Markov operator splits into the five structural components and a
residual; the small print gives the classical counterpart of each
component. The intrinsic sources of reusable structure are \(\Geom\)
and \(\Abs\); the channels \(\Caus\), \(\Lift\), and \(\Bdry\) orient,
reorganize, and read out that structure, and \(\Res\) is transport that
provably admits no such attribution.}
\label{fig-intro-channel-dictionary}
\end{figure}

The labels in \cref{fig-intro-channel-dictionary} name typed framework
objects, not only their classical operator cores.  A row of
\cref{tab-intro-channel-meanings} applies to a complete represented record;
the classical ingredient in its middle column supplies an algebraic
coordinate, while the final column states the additional conditions under
which that coordinate has framework status.

\begin{table}[tbp]
\centering
\small
\setlength{\tabcolsep}{4pt}
\renewcommand{\arraystretch}{1.16}
\begin{tabular}{@{}
  >{\raggedright\arraybackslash}p{0.105\linewidth}
  >{\raggedright\arraybackslash}p{0.255\linewidth}
  >{\raggedright\arraybackslash}p{0.535\linewidth}@{}}
\toprule
Framework term & Classical ingredient & Additional framework meaning \\
\midrule
\(\Geom\) &
reversible or symmetric operator; Dirichlet form &
represented intrinsic geometry with public provenance, qualification,
complete representation cost, and target-visibility accounting \\

\(\Caus\) &
directed current or drift &
authorized target-oriented use of an exposed source, using the represented
projection of the actual selected current under a declared authority envelope \\

\(\Abs\) &
quotient, sufficient statistic, or lumping &
represented abstraction with explicit fibers, terminal compatibility,
qualification, temporal availability, and saturated representation cost \\

\(\Lift\) &
state augmentation &
represented enlargement with a charged interface that preserves target
fractions and inherits the provenance and qualification of its base source \\

\(\Bdry\) &
killing, absorption, or terminal value &
represented terminal readout that evaluates previously exposed structure and
creates no independent interior transport source \\

\(\Res\) &
orthogonal remainder &
post-saturation transport remaining after all affordable
channel-qualified target-relevant representatives have been reclassified and
charged; it carries no prospective structural source charge \\
\bottomrule
\end{tabular}
\caption{What the channel names mean.  The middle column records the
classical operator ingredient; the final column records the additional typed
requirements imposed by the framework.  The precise gates are the budgeted
representative condition of \cref{def-admissible-channel-representatives}, the
authority condition of \cref{def-caus-authority-envelope}, the quotient and
lift conditions of \cref{def-abs-lift-split,def-valid-lift}, the boundary
readout condition of \cref{def-boundary-value-channel}, and the saturated
residual condition of
\cref{cor-no-post-saturation-residual-source-alternative}.}
\label{tab-intro-channel-meanings}
\end{table}

\PartI formalizes this structure independently of any solution method.
A task presentation consists of complete states, observables, a declared
slot for state-to-state structure, and terminal data. Search procedures,
stochastic processes, generators, transition kernels, algorithms, and
complexity bounds are introduced only in subsequent parts. The symbol
\(\operatorname{Dyn}\) therefore denotes a static relation slot in this
part; its dynamical interpretation begins in \PartII.

The static spine is

\[
\text{states}
\longrightarrow
\text{observables}
\longrightarrow
\text{observable event algebra}
\longrightarrow
\text{terminal sectors}
\longrightarrow
\text{the presented task}.
\]

\subsection{Contributions}\label{subsec-contributions}

The contributions form three levels.  The first consists of the framework's
new mathematical objects.  The second contains the theorems made possible by
those objects.  The third records the established analytic ingredients used
inside the proofs.

\paragraph{A. New foundational objects.}

\begin{itemize}
\item \emph{Task presentations as machine-independent computational objects.}
A task-space signature separates complete states, public observables,
state-to-state structure, admissible constructors, resource accounting, and
terminal data
(\cref{def-task-space-signature,def-observable-algebra}).

\item \emph{Budgeted derivability and saturated represented closure.}
Budgeted derivability closes the declared presentation under oracle-free
constructors and represented computations, while saturated representation
cost records the least complete budget of each derived object
(\cref{def-budgeted-derivability-closure,def-saturated-representation-degree}).

\item \emph{Typed structural channels.}
The channels \(\Geom,\Caus,\Abs,\Lift,\Bdry\), and \(\Res\) carry cost,
provenance, qualification, target-use, and residual semantics in addition to
their operator-theoretic coordinates
(\cref{def-admissible-channel-representatives,def-hodge-channel-split,def-abs-lift-split}).

\item \emph{Affordable conditional-algebra frontiers.}
An affordable frontier records the represented backward slice, its nested
conditional algebras, complete carrier cost, and contextual target charges
(\cref{def-affordable-conditional-algebra-frontier}).

\item \emph{Prospective source and retrospective accounting types.}
The framework distinguishes structure available when a transition is chosen
from terminal-compatible records constructed after the trajectory is
complete
(\cref{def-pre-hit-temporally-admissible-source,cor-first-hit-accounting-not-source}).

\item \emph{Algorithm-independent saturated capability profiles.}
The master and channelwise profiles optimize over every qualified represented
structure generated within budget, including structures constructed at
runtime
(\cref{def-extraction-functional,def-saturated-representation-degree}).

\item \emph{Structural meta-frontiers.}
Strict Pareto frontiers are generalized inverses of saturated profiles over
complete resource vectors and retain their temporal and source types
(\cref{def-strict-structural-meta-frontier,def-temporally-typed-structural-meta-frontiers}).

\item \emph{Source-resolved saturated learning ledgers.}
Learning records retain the exact structural source cell, statistical
capacity, information, noise, selection, dynamics, and deployment costs
(\cref{def-complete-source-resolved-learning-ledger}).
\end{itemize}

\paragraph{B. New theorems enabled by these objects.}

\begin{itemize}
\item \emph{Exhaustive channels and exact structural-charge conservation.}
Every represented positive sub-Markov operator has the exhaustive channel
decomposition and explicit normal forms, while every affordable frontier
satisfies an exact Pythagorean charge identity
(\cref{thm-channel-exhaustiveness,thm-exhaustive-geom-normal-form,thm-exhaustive-lift-normal-form,thm-exhaustive-bdry-normal-form,thm-exhaustive-residual-normal-form,thm-exact-affordable-algebra-structural-charging}).

\item \emph{Universal algorithm accountability.}
The quantitative backward cut and exact first-hit ledger bound every
budget-admissible computation by the corresponding saturated profile
(\cref{thm-quantitative-backward-source-factorization,lem:exact-affordable-first-hit,thm-universal-saturated-profile-accountability,thm-prospective-selector-accountability}).

\item \emph{No-retrospective-source and no-uncharged-value principles.}
Completed first-hit labels enter prospective profiles only through an
independent pre-action representation at complete cost, and represented value
cannot amplify prospective target gain without a charged source
(\cref{cor-no-retrospective-first-hit-source-loophole,thm-no-uncharged-prospective-value-amplification}).

\item \emph{Source-cell exhaustion.}
The five-family provenance normal form and the exact
state--coefficient classification exhaust the prospective source cells and
their meta-frontiers
(\cref{thm-five-family-abs-provenance,thm-universal-state-coefficient-source-decomposition,thm-source-cell-meta-complexity-exhaustion}).

\item \emph{Source-resolved Rademacher and PAC--Bayes compilation.}
The learning ledger gives cellwise Rademacher union and PAC--Bayes mixture
bounds, together with operator, effective-dimension, label-noise, and
dependent-sampling consequences
(\cref{thm-source-resolved-saturated-learning-compilation,thm-fixed-operator-rademacher-bound,thm-krr-effective-dimension-bound}).

\item \emph{Controlled-system, attention, and compact-belief compilation.}
Selected generators, Bellman and reach--avoid quantities, attention records,
and compact latent beliefs compile into the same source- and cost-resolved
capability ledger
(\cref{thm-controlled-policy-channel-decomposition,thm-universal-dynamical-certificate,thm-attention-structural-source-compilation,thm-attention-target-gain-envelope,thm-compact-latent-information-sufficiency,thm-joint-compact-latent-attention-acquisition}).

\item \emph{Explicit statistical, computational, and information noise
thresholds.}
The source-resolved ledger separates the sample, represented-computation, and
information limits and compares their phase transitions
(\cref{thm-three-noise-threshold-comparison,thm-dynamical-noise-three-thresholds}).

\item \emph{Proof-producing finite-to-asymptotic certificate machinery.}
For the LPN application, exact quotient and geometric exhaustion isolates one
residual-ancestry certification target.  Exact enumeration, deterministic
verification, symbolic inequalities, outward-rounded interval bounds,
finite-prefix checks, and uniform scale propagation form a sufficient
certificate whose deterministic acceptance discharges that target.  The
paper proves this discharge implication and specifies every verification
field; it does not execute the final instantiated certificate analysis
(\cref{thm-lpn-constant-noise-materialized-spectral-closure,lem-lpn-residual-parity-spectral-fibers,cor-empirical-discovery-deterministic-universal-closure,thm-lpn-proof-grade-empirical-residual-discharge}).
\end{itemize}

\paragraph{C. Standard ingredients used inside the proofs.}

\begin{itemize}
\item Markov and sub-Markov kernels, generators, and killed resolvents;
\item Dirichlet forms, reversible operators, and Hodge projection;
\item conditional expectation and projection Pythagoras;
\item lumpability, quotient operators, and state augmentation;
\item Rademacher complexity, PAC--Bayes inequalities, effective dimension,
and concentration bounds;
\item Bellman equations, Bayesian identities, and Gibbs selectors; and
\item interval arithmetic, finite enumeration, and induction.
\end{itemize}

These established ingredients serve inside a presentation-relative
architecture that assigns complete cost, provenance, temporal availability,
source type, and deployment semantics before compiling them into universal
capability bounds.

\subsection{How to read this paper}\label{subsec-reading-guide}

The manuscript has a three-part structural spine followed by learning,
controlled-systems, attention, structural meta-complexity, and LPN parts. The beginning of each part states its inputs, output, and a suggested
first pass. The complete organization is as follows.

\begin{description}[style=nextline]
\item[\PartI: presented structure.]
The \hyperref[the-task-space-signature]{task-space signature} fixes the state,
observable, relation, and terminal interfaces. The part then constructs the
\hyperref[the-observable-algebra]{observable algebra}, formalizes static target
relevance, and closes with
\hyperref[admissible-presentation-equivalence]{admissible presentation
equivalence}. Its output is the complete static object on which all later
dynamics act.

\item[\PartII: admissible dynamics.]
Budget-admissible movement is written in generator form and decomposed into
its structural channels. The
\hyperref[the-exhaustiveness-theorem]{channel-exhaustiveness theorem} proves
that the list is complete, while the
\hyperref[the-substrate-bridge-algorithms-as-markov-kernels]{operational
embedding} places ordinary computations on the same transcript carrier. Its
output is an algorithm-independent channel decomposition with the exhaustive
channel normal forms of
\cref{thm-exhaustive-geom-normal-form,thm-exhaustive-lift-normal-form,thm-exhaustive-bdry-normal-form,thm-exhaustive-residual-normal-form}
and explicit resource provenance.

\item[\PartIII: saturated extraction cost.]
The \hyperref[the-resolvent-stopwatch]{killed resolvent} measures target
arrival. The
\hyperref[the-combined-geomabs-extraction-functional]{combined extraction
functional} then assembles the saturated ledger, quantitative backward cut,
and structural charging identity. Exact first-hit accounting culminates in
the \hyperref[thm-universal-saturated-profile-accountability]{universal
accountability theorem}. The
\hyperref[the-saturated-cost-invariant]{saturated cost invariant} and the
\hyperref[the-barrier-theorems-and-the-operator-invariant]{presentation-extension
tests} complete the framework core.

\item[\PartIV: learning.]
The saturated structural objects are read as kernels and RKHS features. This
part derives effective-dimension and
\hyperref[the-tightened-two-sided-learning-bound]{two-sided learning bounds},
treats \hyperref[neural-networks-the-learned-kernel-realization]{learned neural
kernels}, transports these capacity certificates through
\hyperref[learning-under-symmetric-label-corruption]{symmetric label
corruption}, develops computable evaluation and deployment certificates, and
extends structural Rademacher and PAC--Bayes control to independent episodes,
mixing trajectories, and predictable adaptive paths.  It ends with the
noise-indexed statistical, computational, and information threshold comparison.
It uses \PartsItoIII.  Its statistical results feed the quantitative LPN
phase comparison, while the structural LPN reduction uses \PartsItoIII
directly.

\item[\PartV: controlled dynamical systems.]
Represented policies select cost-tagged generators on finite or clocked
carriers.  This part combines the existing reversible geometry, authorized
Caus current, committor accounting, reach--avoid capacity, Bellman
contraction, the audited functional-inequality profile, and dependent-data
learning certificates in the
\hyperref[thm-universal-dynamical-certificate]{universal dynamical-systems
theorem} and the
\hyperref[thm-master-structural-dynamical-learning-certificate]{master
learned-control certificate}.  Continuous and stochastic models enter through the
\hyperref[thm-controlled-continuum-lift-transfer]{represented Lift transfer}.
The optimal structural Bayesian section also defines compact latent beliefs,
their information and sufficiency coordinates, and the exact joint
acquisition--encoding ceiling in
\cref{thm-compact-latent-information-sufficiency,thm-optimal-compact-latent-acquisition}.
The trained-model chapters give complementary white-box and black-box
certificates: \cref{sec-certified-neural-implementation} audits a complete
training record, while \cref{sec-black-box-structural-audit} audits an oracle
and represented test law at finite query cost.

\item[\PartVI: attention.]
Attention rows are treated as represented Gibbs selectors.  Their score
geometry, directed current, fibers, values, masks, and residual transport are
compiled into the existing source-resolved ledger in
\cref{thm-attention-structural-source-compilation};
\cref{thm-attention-target-gain-envelope} gives the target-gain envelope.
The log-likelihood specialization identifies attention with the compact latent
Bayesian update, and the joint theorem equates optimal compact latent attention
with optimal acquisition on its saturated transfer class
(\cref{cor-bayesian-posterior-attention-flow,thm-joint-compact-latent-attention-acquisition}).

\item[\PartVII: structural meta-complexity.]
The complete prospective records are indexed by the Pareto cost required to
cross a target-gain threshold.  The strict structural meta-frontier is the
exact generalized inverse of the saturated profile in
\cref{thm-strict-meta-profile-generalized-inverse}.  The part separates
representation, extraction, decision, search, certification, and prospective
deployment; recovers circuit and program minimization as typed
specializations; and derives the meta-accountability implication in
\cref{thm-structural-meta-accountability}.

\item[\PartVIII: LPN.]
The \hyperref[the-lpn-task-presentation]{public LPN presentation} is inserted
into the framework.  Its noise--resource phase diagram compares the
polynomial, statistical, and information thresholds.  Public verification is
separated from extraction; the constant-noise spectral ledger evaluates
materialized located spectra and the exact residual-parity fiber gate; the
geometric and five abstraction-provenance sources are core-refined; and the
resulting saturated-profile reduction yields the
\hyperref[sec-lpn-saturated-accounting-criterion]{exact obstruction
criterion}.  Fixed-noise search-LPN hardness implies \(P\ne NP\) by
\cref{cor-lpn-p-not-equal-np}.

\item[Index of notation.]
The appendix collects the load-bearing symbols, their definitions, and their
first substantive uses.
\end{description}

Four reading routes cover the usual uses of the manuscript.

\begin{itemize}
\item For the framework itself, read \PartsItoIII in order; the shortest spine
is the task signature and observable algebra, channel and residual exhaustiveness,
the four channel normal forms, budgeted derivability closure, saturated representation cost, backward cut,
structural charging, and universal accountability.
\item For the LPN theorem, read
\cref{def-task-space-signature,def-observable-algebra},
\cref{thm-channel-exhaustiveness,thm-exhaustive-geom-normal-form,thm-exhaustive-lift-normal-form,thm-exhaustive-bdry-normal-form,thm-exhaustive-residual-normal-form},
\cref{def-budgeted-derivability-closure,def-saturated-representation-degree},
\cref{thm-five-family-abs-provenance,thm-exact-affordable-algebra-structural-charging},
\cref{lem:exact-affordable-first-hit,thm-universal-saturated-profile-accountability},
and then \cref{part:lpn}.
For the quantitative noise thresholds within that part, read
\cref{sec-lpn-noise-resource-phase,cor-lpn-quantitative-noise-phase-localization}.
For the constant-noise closure perimeter, read
\cref{sec-lpn-constant-noise-spectral-closure,thm-lpn-exact-constant-noise-spectral-perimeter}.
\item For the learning results, read the static presentation and observable
algebra in \PartI, the channel decomposition in \PartII, and the saturated
profile and accountability results in \PartIII before entering \PartIV.
\item For controlled dynamics, read the generator and channel decomposition in
\PartII, the killed-resolvent and temporal-accountability results in \PartIII,
and the geometry-learning results in \PartIV before entering \PartV.
\item For a trained black-box classifier, read
\cref{def-complete-black-box-predictor-audit,%
thm-finite-query-simultaneous-black-box-audit,%
thm-black-box-training-noise-identifiability,%
thm-end-to-end-black-box-structural-certificate}.  The resulting certificate
is test-law relative; \cref{thm-black-box-audit-deployment-transfer} supplies
the represented deployment-law conversion.
\item For attention and active acquisition, read the source-resolved learning
ledger in \PartIV and the optimal active-sampling certificate in \PartV before
entering \PartVI.  The unified compact-belief route is
\cref{def-complete-compact-latent-belief-record,thm-compact-latent-information-sufficiency,thm-optimal-compact-latent-acquisition,cor-bayesian-posterior-attention-flow,thm-joint-compact-latent-attention-acquisition}.
\item For structural meta-complexity, read
\cref{def-structural-meta-profile-datum,def-strict-structural-meta-frontier,thm-strict-meta-profile-generalized-inverse,thm-structural-meta-accountability}.
The classical representation-minimization specializations are
\cref{thm-circuit-minimization-specialization,thm-program-minimization-specialization};
the source-resolved and magnification interfaces are
\cref{thm-source-cell-meta-complexity-exhaustion,thm-represented-profile-magnification}.
\end{itemize}

\Cref{fig-framework-dependency} gives both the dependency structure and the
contents of each part. Horizontal arrows show the progression within a part;
the vertical arrows show which parts supply the next stage.

\begin{figure}[tbp]
\centering
\begin{tikzpicture}[
  fw diagram,
  node distance=4.5mm and 5mm,
  part node/.style={fw box,minimum width=2.4cm,font=\small\bfseries},
  contents node/.style={fw box,text width=10.5cm,minimum height=8mm,
    font=\scriptsize,align=center}
]
  \node[part node] (p1) {\PartI};
  \node[contents node,right=of p1] (c1)
    {\hyperref[the-task-space-signature]{task presentation}
     $\longrightarrow$ \hyperref[the-observable-algebra]{observable algebra}
     $\longrightarrow$ \hyperref[what-makes-an-observable-useful]{target relevance}
     $\longrightarrow$ \hyperref[admissible-presentation-equivalence]{presentation equivalence}};

  \node[part node,below=of p1] (p2) {\PartII};
  \node[contents node,right=of p2] (c2)
    {\hyperref[budget-admissible-movement-and-the-search-operator]{admissible movement}
     $\longrightarrow$ \hyperref[the-fixed-space-decomposition]{generator decomposition}
     $\longrightarrow$ \hyperref[the-exhaustiveness-theorem]{exhaustive channels}
     $\longrightarrow$ \hyperref[the-substrate-bridge-algorithms-as-markov-kernels]{computation embedding}};

  \node[part node,draw=fwBlue,fill=fwBlue!13,below=of p2] (p3) {\PartIII};
  \node[contents node,draw=fwBlue,fill=fwBlue!13,right=of p3] (c3)
    {\hyperref[the-resolvent-stopwatch]{killed resolvent}
     $\longrightarrow$ \hyperref[def-saturated-representation-degree]{saturated ledger}
     $\longrightarrow$ \hyperref[thm-quantitative-backward-source-factorization]{backward cut and charging}
     $\longrightarrow$ \hyperref[thm-universal-saturated-profile-accountability]{first hit and accountability}};

  \node[part node,draw=fwAmber,fill=fwAmber!14,below=of p3] (p4) {\PartIV};
  \node[contents node,draw=fwAmber,fill=fwAmber!14,right=of p4] (c4)
    {\hyperref[the-observable-algebra-is-a-reproducing-kernel-hilbert-space]{RKHS realization}
     $\longrightarrow$ \hyperref[saturated-cost-and-effective-dimension]{effective dimension}
     $\longrightarrow$ \hyperref[the-tightened-two-sided-learning-bound]{learning bounds}
     $\longrightarrow$ \hyperref[computable-recipes]{evaluation and certification}};

  \node[part node,draw=fwTeal,fill=fwTeal!14,below=of p4] (p5) {\PartV};
  \node[contents node,draw=fwTeal,fill=fwTeal!14,right=of p5] (c5)
    {\hyperref[def-controlled-dynamical-presentation]{controlled presentation}
     $\longrightarrow$ \hyperref[thm-controlled-policy-channel-decomposition]{selected-generator channels}
     $\longrightarrow$ \hyperref[thm-controlled-committor-bellman]{Bellman and committors}
     $\longrightarrow$ \hyperref[thm-universal-dynamical-certificate]{coordinatewise certificate}};

  \node[part node,draw=fwViolet,fill=fwViolet!12,below=of p5] (p6) {\PartVI};
  \node[contents node,draw=fwViolet,fill=fwViolet!12,right=of p6] (c6)
    {\hyperref[def-complete-attention-record]{attention record}
     $\longrightarrow$ \hyperref[thm-attention-gibbs-gradient-flow]{Gibbs geometry}
     $\longrightarrow$ \hyperref[thm-attention-structural-source-compilation]{source compilation}
     $\longrightarrow$ \hyperref[thm-attention-target-gain-envelope]{target-gain envelope}};

  \node[part node,draw=fwBlue,fill=fwBlue!13,below=of p6] (p7) {\PartVII};
  \node[contents node,draw=fwBlue,fill=fwBlue!13,right=of p7] (c7)
    {\hyperref[def-strict-structural-meta-frontier]{certificate frontier}
     $\longrightarrow$ \hyperref[thm-strict-meta-profile-generalized-inverse]{profile inverse}
     $\longrightarrow$ \hyperref[thm-source-cell-meta-complexity-exhaustion]{source cells}
     $\longrightarrow$ \hyperref[thm-structural-meta-accountability]{meta-accountability}};

  \node[part node,draw=fwCoral,fill=fwCoral!12,below=of p7] (p8) {\PartVIII};
  \node[contents node,draw=fwCoral,fill=fwCoral!12,right=of p8] (c8)
    {\hyperref[the-lpn-task-presentation]{LPN presentation}
     $\longrightarrow$ \hyperref[sec-lpn-noise-resource-phase]{noise phases}
     $\longrightarrow$ \hyperref[sec-lpn-constant-noise-spectral-closure]{spectral perimeter}
     $\longrightarrow$ \hyperref[cor-lpn-abs-source-exhaustion]{exact-\(\Abs\) exhaustion}
     $\longrightarrow$ \hyperref[thm-lpn-one-way]{fixed-noise hardness}
     $\longrightarrow$ \hyperref[cor-lpn-p-not-equal-np]{\(P\ne NP\)}};

  \node[part node,draw=fwInk!65,fill=fwGray,below=of p8] (pa)
    {\hyperref[index-of-notation]{Appendix}};
  \node[contents node,draw=fwInk!65,fill=fwGray,right=of pa] (ca)
    {notation $\longrightarrow$ defining result $\longrightarrow$ first substantive use};

  \draw[fw arrow] (p1) -- (p2);
  \draw[fw arrow] (p2) -- (p3);
  \draw[fw arrow] (p3) -- (p4);
  \draw[fw arrow] (p4) -- (p5);
  \draw[fw arrow] (p5) -- (p6);
  \draw[fw arrow] (p6) -- (p7);
  \draw[fw arrow] (p7) -- (p8);
  \draw[fw dashed arrow] (p8) -- (pa);
\end{tikzpicture}
\caption{Map of the manuscript. \PartsItoIII form the structural spine.
\PartIV derives learning bounds, \PartV assembles controlled dynamical-system
certificates, \PartVI compiles attention, \PartVII takes the certificate-cost
inverse of the saturated profile, and \PartVIII applies the framework to LPN.
The \hyperref[index-of-notation]{notation
appendix} is a navigation aid rather than a further proof stage.}
\label{fig-framework-dependency}
\end{figure}
\clearpage

A score is an observable on the state space rather than a relation
between states. Likewise, the success set is one component of the
terminal interface. Keeping these objects distinct provides the static
data required by the dynamical construction in \PartII.

\begin{example}[SAT before any solver]\label{ex-intro-sat-static}

A SAT formula \(\Phi\) presents the finite state set
\(X_\Phi=\{0,1\}^n\), whose elements are complete truth assignments.
Clause readouts, variable projections, and the violated-clause count are
observables on this set. The satisfying assignments form the positive
terminal sector. This presentation specifies the available information,
but no rule for selecting subsequent assignments.

\end{example}

\section{The Task-Space Signature}\label{sec-task-space-signature}
\label{the-task-space-signature}

A task-space signature collects the data specified independently of a
solution procedure. Its four components are the state space, observable
interface, dynamic relation slot, and terminal interface. Only the first,
second, and fourth components receive substantive structure in \PartI.

The task-space signature packages the four kinds of data declared before any solution procedure is chosen.
\begin{definition}[Task-space signature]\label{def-task-space-signature}

For an instance \(I\), a task-space signature is a tuple

\[
\mathcal T_I=(X_I,\operatorname{Obs}_I,\operatorname{Dyn}_I,\operatorname{End}_I).
\]

The components are defined as follows.

\medskip\noindent\textbf{State space.}
\(X_I\) is a nonempty set of complete states. Completeness means that
every primitive observable and every state-dependent terminal predicate
is defined on every \(x\in X_I\).

\medskip\noindent\textbf{Observable interface.}
There is an index set \(A_I\). For each \(a\in A_I\), let
\((Y_{I,a},\mathcal B_{I,a})\) be a measurable value space. A primitive
observable is a map

\[
O_{I,a}:X_I\to Y_{I,a}.
\]

The observation interface is the indexed family

\[
\operatorname{Obs}_I=(O_{I,a})_{a\in A_I}.
\]

The observable event algebra is the initial \(\sigma\)-algebra

\[
\mathcal F_I=\mathcal F_{\operatorname{Obs}_I}
=
\sigma\{O_{I,a}^{-1}(B):a\in A_I,\ B\in\mathcal B_{I,a}\}.
\]

Thus \(\mathcal F_I\) is the smallest \(\sigma\)-algebra on \(X_I\) for which
all primitive observables are measurable.

All observable value spaces used for atom and fiber statements are
measurably separated on their realized ranges: distinct realized values are
separated by a measurable set, and every realized singleton is measurable.
Standard Borel and finite discrete value spaces satisfy this convention.  It
ensures that equality of realized observable values agrees with
indistinguishability by measurable preimages.

\medskip\noindent\textbf{Relation slot.}
The dynamic slot is declared as an
\(\mathcal F_I\otimes\mathcal F_I\)-measurable relation

\[
\operatorname{Dyn}_I=R_I\subseteq X_I\times X_I,
\qquad
R_I\in\mathcal F_I\otimes\mathcal F_I.
\]

No process, kernel, generator, algorithm, or traversal rule is assigned
to \(R_I\) at this stage. It records the relation slot whose dynamical
interpretation is introduced in \PartII.

The fourth component \(\operatorname{End}_I\) is the terminal interface of
\cref{def-task-terminal-interface}.

\end{definition}

The terminal interface separates state-terminal sectors from a run-dependent
stopping predicate.
\begin{definition}[Terminal interface]\label{def-task-terminal-interface}

The terminal interface is

\[
\operatorname{End}_I=(\operatorname{End}_I^+,\operatorname{End}_I^-,\operatorname{End}_I^0).
\]

The first two components are measurable state predicates, equivalently
disjoint measurable sectors

\[
E_I^+=\{x:\operatorname{End}_I^+(x)=1\},
\qquad
E_I^-=\{x:\operatorname{End}_I^-(x)=1\},
\qquad
E_I^+\cap E_I^-=\varnothing.
\]

The state-terminal sector is \(S_I=E_I^+\sqcup E_I^-\). When terminal
status is a function only of the present state, failure is the
complement of success inside \(S_I\):

\[
E_I^-=S_I\setminus E_I^+.
\]

Thus failure is defined within the state-terminal sector rather than as
\(X_I\setminus E_I^+\), which may contain nonterminal states. The
component \(\operatorname{End}_I^0\) is a run-dependent stopping
predicate. For a clock or resource space
\((C_I,\mathcal C_I)\), one may formally write

\[
\operatorname{End}_I^0:
\bigsqcup_{m\ge0}(X_I\times C_I)^{m+1}\to\{0,1\},
\]

measurable for the corresponding disjoint-union product \(\sigma\)-algebra.
Its use requires runs and is deferred with dynamics.

\end{definition}

\Cref{fig-task-presentation} records the four declared components
of the task-space signature before any solution procedure is chosen.

\begin{figure}[tbp]
\centering
\begin{tikzpicture}[fw diagram,node distance=5mm and 5mm]
  \node[fw source,minimum width=2.7cm] (state) {states\\\(X_I\)};
  \node[fw source,minimum width=2.7cm,right=of state] (obs)
    {observables\\\(\operatorname{Obs}_I\)};
  \node[fw use,minimum width=2.7cm,right=of obs] (dyn)
    {relation slot\\\(\operatorname{Dyn}_I\)};
  \node[fw terminal,minimum width=2.7cm,right=of dyn] (end)
    {terminal data\\\(\operatorname{End}_I\)};
  \node[fw box,minimum width=4.2cm,below=11mm of $(obs)!0.5!(dyn)$] (task)
    {presented task \(\mathcal T_I\)};
  \draw[fw arrow] (state) -- (task);
  \draw[fw arrow] (obs) -- (task);
  \draw[fw arrow] (dyn) -- (task);
  \draw[fw arrow] (end) -- (task);
\end{tikzpicture}
\caption{The task presentation fixes states, public observables, the
relation slot, and terminal data independently of an algorithm; see
\cref{def-task-space-signature}.}
\label{fig-task-presentation}
\end{figure}

Accordingly, \(X_I\) is the space of complete candidate states,
\(\mathcal F_I\) is the observable \(\sigma\)-algebra generated by the
primitive readouts, and \(\operatorname{End}_I\) separates successful,
failed, and run-dependent terminal conditions. The relation \(R_I\)
remains undeveloped until \PartII.

\begin{example}[Running example: the LPN presentation]
\label{ex-lpn-running-part1}

For the public LPN data \((A,b,\eta)\), take
\(X_I=\mathbb F_2^n\). Candidate coordinates and the residual map
\(x\mapsto Ax+b\) generate the public observable interface. The relation
slot remains undeveloped in \PartI, while the target for the inversion task is
\(T_I=\{s\}\). The secret \(s\) and noise vector are not primitive public
observables. \Cref{def-lpn-task-object} gives the complete specialization.

\end{example}

\begin{proposition}[Complete states are scorable and testable]\label{prop-complete-states-testable}

Let \(D_I:X_I\to\mathbb R_{\ge0}\) be a defect score included in
\(\operatorname{Obs}_I\), and suppose success is zero defect. Then
\(D_I\) is \(\mathcal F_I\)-measurable and

\[
E_I^+=D_I^{-1}(\{0\})\in\mathcal F_I.
\]

Thus every complete state is both scorable by \(D_I\) and testable for
success.

\end{proposition}

\begin{proof}\phantomsection\label{proof-complete-states-testable}

Because \(D_I\) is a declared observable, it is measurable by
construction of \(\mathcal F_I\). Since \(\{0\}\) is Borel in
\(\mathbb R_{\ge0}\), the preimage \(D_I^{-1}(\{0\})\) belongs to
\(\mathcal F_I\). Completeness of \(X_I\) gives a value \(D_I(x)\) and a
terminal verdict for every \(x\in X_I\).

\end{proof}

\begin{example}[SAT signature]\label{ex-signature-sat}

For a CNF formula \(\Phi=C_1\wedge\cdots\wedge C_m\) on \(n\) variables,
take

\[
X_\Phi=\{0,1\}^n,
\qquad
O_{\Phi,j}(x)=\mathbf 1\{C_j(x)=1\},
\qquad
D_\Phi(x)=\sum_{j=1}^m(1-O_{\Phi,j}(x)).
\]

The success sector is \(E_\Phi^+=D_\Phi^{-1}(\{0\})\). If the
presentation has an explicit terminal-unsatisfied label, that sector is
\(E_\Phi^-\); otherwise unsatisfied assignments are non-success states,
not automatically terminal failures. The relation
\(R_\Phi\subseteq X_\Phi\times X_\Phi\) may be declared in the
signature, but \PartI does not use it.

\end{example}

\section{The Observable Algebra}\label{sec-observable-algebra}
\label{the-observable-algebra}

\subsection{Observable partitions and measurable readouts}
\label{subsec-observable-partitions-readouts}

The observable interface determines the measurable subsets of the state
space. Its atoms are the minimal distinguishable cells, whereas the
level sets of a single observable may form a coarser partition. A score
is one observable within this algebra.

The observable algebra records exactly the distinctions generated by the primitive readouts.
\begin{definition}[Observable algebra, atoms, and level sets]\label{def-observable-algebra}

Fix an instance \(I\). The observable algebra is

\[
\mathcal F_I=
\sigma\{O_{I,a}^{-1}(B):a\in A_I,\ B\in\mathcal B_{I,a}\}.
\]

Two states are observationally equivalent, written
\(x\equiv_{\operatorname{Obs}_I}y\), when

\[
O_{I,a}(x)=O_{I,a}(y)
\qquad\text{for every }a\in A_I.
\]

Under the measurable-separation convention of
\cref{def-task-space-signature}, when \(X_I\) is finite the atoms of
\(\mathcal F_I\) are the equivalence classes of
\(\equiv_{\operatorname{Obs}_I}\). More generally, when the displayed
intersection is measurable and is an atom, the atom containing \(x\) is

\[
[x]_{\mathcal F_I}=\bigcap\{A\in\mathcal F_I:x\in A\}.
\]

If \(\phi:(X_I,\mathcal F_I)\to(Y,\mathcal B_Y)\) is a derived
observable whose realized singletons are measurable, its level set over
\(y\in Y\) is

\[
L_y^\phi=\phi^{-1}(\{y\})=\{x\in X_I:\phi(x)=y\}.
\]

The subalgebra generated by \(\phi\) is
\(\sigma(\phi)=\{\phi^{-1}(B):B\in\mathcal B_Y\}\subseteq\mathcal F_I\).

\end{definition}

An atom is a maximal observational-equivalence class: no event or
function generated by the primitive observables separates two states in
the same atom. A level set may be coarser because it depends on a single
observable rather than the full observable family.

\begin{proposition}[Observable events and functions are atom-constant]\label{prop-atoms-characterize-events}

Assume \(X_I\) is finite and the primitive value spaces obey the
measurable-separation convention of \cref{def-task-space-signature}. A set
\(A\subseteq X_I\) belongs to
\(\mathcal F_I\) if and only if it is a union of atoms. A function
\(g:X_I\to Z\) into a measurably separated space is
\(\mathcal F_I\)-measurable if and only if it is constant on every atom
of \(\mathcal F_I\).

\end{proposition}

\begin{proof}\phantomsection\label{proof-atoms-characterize-events}

Each generator \(O_{I,a}^{-1}(B)\) is a union of observational
equivalence classes. Complements and countable unions of such unions have
the same property, so every event in \(\mathcal F_I\) is a union of
equivalence classes. Conversely, fix one class \(C\). For each
\(y\notin C\), measurable separation supplies an observable index and a
measurable value event that contains the value at a fixed point of \(C\)
and excludes the value at \(y\). Because \(X_I\) is finite, intersecting
these finitely many preimages produces \(C\). Thus every class, and hence
every union of classes, is measurable. For functions, measurable
separation in \(Z\) implies that unequal values are separated by a
measurable preimage, contradicting membership in one atom. If a function
is atom-constant, every preimage is a union of atoms and is measurable.

\end{proof}

\begin{proposition}[Measurability and budget admissibility]\label{prop-measurability-is-not-affordability}

For a finite task presentation, \(\mathcal F_I\)-measurability is
equivalent to atom constancy. Budget admissibility imposes the stronger
representation-cost condition that the observable admit a representative
whose constructor charge lies below the ceiling \(b(n)\) of a budget structure
\(\mathcal B=(\mathcal R,\operatorname{cost},b,\mathfrak C)\), or lie in
a valid saturated representation generated within that budget. A
measurable observable can have high degree or high cost in one public
filtration. If a valid quotient, basis change, or lift generated from
the observable algebra represents the same target-relevant distinction
at lower cost, the later saturated invariant uses that lower
representation cost. The polynomial clause is the special budget
\(\mathcal B_{\mathrm{poly}}\) with \(b_{\mathrm{poly}}(n)=n^{O(1)}\).

\end{proposition}

\begin{proof}\phantomsection\label{proof-measurability-is-not-affordability}

Atom constancy bounds neither the number of distinguished atoms nor the
dimension or constructor charge of the smallest representing coordinate
band. Measurability therefore determines provenance but not budget
admissibility. A represented quotient or valid lifted coordinate may
represent the same atom union at smaller charge, in which case the cost
of that representation applies. Degree is one coordinate of the
admissible representation cost rather than its defining quantity.

\end{proof}

\subsection{Budgeted generation and refinement}
\label{subsec-budgeted-generation-refinement}

The budget structure records how resource bounds are ordered and combined.
\begin{definition}[Budget structure]\label{def-budget-structure}

A budget structure for a task family \((I_n)\) is
\(\mathcal B=(\mathcal R,\operatorname{cost},b,\mathfrak C)\).
\medskip\noindent\textbf{Resource monoid and transfer system.}
\(\mathcal R=(C,\preceq,\oplus,\mathbf 0)\) is an ordered commutative
resource monoid with monotone accumulation \(\oplus\), with
\(\mathbf0\preceq c\) for every \(c\in C\) (hence
\(c\preceq c\oplus d\));
\(\operatorname{cost}\) assigns a resource charge to each admissible
constructor representative; \(b(n)\in C\) is the size-indexed capacity
ceiling; and \(\mathfrak C\) is a typed transfer system

\[
\mathfrak C=
(\mathfrak C_{\mathrm{res}},\mathfrak C_{\mathrm{num}}).
\]

The resource part \(\mathfrak C_{\mathrm{res}}\) consists of declared
monotone overhead maps \(\theta:C\to C\), closed under the finite
compositions used below. Two resource ceilings are class-equivalent
when each is bounded by the image of the other under such a map. The
numeric part \(\mathfrak C_{\mathrm{num}}\) consists of functions
\(\beta:\mathbb N\to[1,\infty)\), contains the constant one, and is
closed under finite products, maxima, and the declared
class-preserving compositions. Thus an expression involving a resource
ceiling uses \(\mathfrak C_{\mathrm{res}}\), whereas notation such as
\(\beta\in\mathfrak C\), \(\beta^{-1}\), or a negligibility margin
always means \(\beta\in\mathfrak C_{\mathrm{num}}\). No inverse or
numeric product is applied to an element of the resource monoid.

\medskip\noindent\textbf{Time capacity.}
Whenever a stopwatch or finite horizon is used, the budget structure
also declares a monotone numerical time-capacity map
\(\mathsf t:C\to[1,\infty]\). The number \(\mathsf t(B)\) is the
elementary-transition capacity of resource ceiling \(B\); it does not
identify \(B\) with a scalar.
Compatibility means explicitly that for every resource overhead
\(\theta\in\mathfrak C_{\mathrm{res}}\) used on a size-indexed ceiling
there is \(\beta\in\mathfrak C_{\mathrm{num}}\) such that

\[
\mathsf t(\theta(B(n)))
\le \beta(n)\,\mathsf t(B(n))
\]

eventually. The same requirement applies to any additional declared
numeric capacity map used by a later specialization. This is the
operative content of static/dynamic compatibility axiom B4.

\end{definition}

Repeated use and constructor composition are charged inside the resource
monoid rather than by implicit scalar multiplication.
\begin{definition}[Resource accumulation and constructor accounting]
\label{def-budget-constructor-accounting}

\medskip\noindent\textbf{Repeated resource use.}
For \(k\in\mathbb N\) and \(c\in C\), write

\[
k\odot c=\underbrace{c\oplus\cdots\oplus c}_{k\text{ terms}},
\qquad 0\odot c=\mathbf 0.
\]

This is resource accumulation under \(k\) repeated uses; it is not
numeric multiplication unless the resource monoid is numeric.

\medskip\noindent\textbf{Constructor accounting.}
The cost charge is sub-accumulative: a constructor assembled from
\(\gamma_1,\ldots,\gamma_k\) by an admissible grammar operation has
charge bounded by
\(\operatorname{cost}(\gamma_1)\oplus\cdots\oplus\operatorname{cost}(\gamma_k)\oplus\operatorname{cost}_{\mathrm{op}}\).
A declared bounded-iteration rule is required to provide a resource
overhead \(\theta\in\mathfrak C_{\mathrm{res}}\) such that the actual
finite repeated sum of all per-step and retained charges is
\(\preceq\theta(B)\) whenever the input record and iteration count lie
within ceiling \(B\). This is the finite-sum domination used by every
later ledger; it is not inferred from notation such as \(n^{O(1)}\).
A representative is \(\mathcal B\)-admissible when some equivalent
representative has charge \(\preceq b(n)\). Both parts of the transfer
system are closed
under the bounded compositions allowed by the resource envelope, so a
bounded composition of admissible representatives remains admissible
after replacing \(b\) by a class-equivalent ceiling. Budget admissibility
is representative-independent, and the static ceiling \(b(n)\) is
compatible with the dynamic stopping level up to class-preserving
overhead.

\end{definition}

The following convention names the axioms and fixes the polynomial
specialization used throughout the manuscript.
\begin{convention}[Budget axioms and polynomial specialization]
\label{conv-budget-axioms-polynomial}

\medskip\noindent\textbf{Budget axioms.}
These are the budget axioms used below: \textbf{B1} monotone ordered
accumulation of constructor charges; \textbf{B2} closure of the transfer
class under the bounded compositions used by the admissible grammar;
\textbf{B3} representative independence of budget admissibility; and
\textbf{B4} compatibility between the static ceiling \(b(n)\) and
dynamic stopping levels. No later argument may use a special property of
polynomials until the budget structure is instantiated as
\(\mathcal B_{\mathrm{poly}}\).

\medskip\noindent\textbf{Polynomial specialization.}
The polynomial budget \(\mathcal B_{\mathrm{poly}}\) is the
transfer-equivalence class of ceilings \(n^C\) with fixed
\(C<\infty\), with numeric resource, ordinary addition, and total encoded
execution-and-representation cost, including time, peak native state,
outcome generation, transcript, evaluation, and retained-output charges.
One representative ceiling is
\(b_{\mathrm{poly},C_0}(n)=n^{C_0}\); admissibility permits a
class-equivalent polynomial enlargement with another fixed exponent. Its
time-capacity map is the identity, and both parts of its transfer system
are polynomially bounded.

\end{convention}

\Cref{fig-budget-typed-ledger} displays the separation between resource
accumulation and the numerical capacity quantities derived from it.

\begin{figure}[tbp]
\centering
\begin{tikzpicture}[fw diagram]
  \node[fw source,minimum width=3.0cm] (dag) at (-5.0,2.1)
    {constructor DAG\\and retained data};
  \node[fw provenance,minimum width=3.2cm] (charge) at (-1.5,2.1)
    {resource charges\\\((C,\oplus,\preceq)\)};
  \node[fw terminal,minimum width=2.7cm] (ceiling) at (2.1,2.1)
    {resource ceiling\\\(B\)};
  \node[fw use,minimum width=3.0cm] (slice) at (5.2,2.1)
    {represented slice\\cost \(\preceq B\)};
  \draw[fw arrow] (dag) -- (charge);
  \draw[fw arrow] (charge) -- (ceiling);
  \draw[fw arrow] (ceiling) -- (slice);
  \node[fw box,minimum width=2.8cm] (time) at (-1.5,0)
    {time capacity\\\(\mathsf t(B)\)};
  \node[fw box,minimum width=2.8cm] (margin) at (2.1,0)
    {numeric factors\\\(\beta(n)\)};
  \node[fw geom,minimum width=3.0cm] (bound) at (5.2,0)
    {numeric bounds\\and margins};
  \draw[fw dashed arrow] (ceiling) -- node[left,font=\scriptsize] {typed map} (time);
  \draw[fw dashed arrow] (ceiling) -- node[right,font=\scriptsize] {B4/\(\theta\)} (margin);
  \draw[fw arrow] (time) -- (bound);
  \draw[fw arrow] (margin) -- (bound);
\end{tikzpicture}
\caption{Typed budget accounting. Constructor costs accumulate in the
resource monoid and are compared with a resource ceiling. Numerical time
capacities and transfer factors arise only through declared typed maps;
resource elements are never multiplied or inverted as scalars.}
\label{fig-budget-typed-ledger}
\end{figure}

Budget filtering restricts represented observables by their charged construction cost.
\begin{definition}[Budget-filtered represented observable algebra]\label{def-budgeted-observable-closure-static}

The static algebra \(\mathcal F_I\) is the ambient measurable algebra
of the presentation. Affordability is carried by represented
constructions rather than by membership in this ambient algebra. A
static construction record \(\rho\) is a finite directed acyclic graph
whose leaves are declared public data and whose internal vertices are
declared constructors. Its denotation on \(I\) is written
\(\sem{\rho}_I\), and its charge includes all leaves,
constructor descriptions, intermediate represented data, and output
operations. A family of records is uniform when one oracle-free
construction scheme produces the record from \(\operatorname{enc}(I)\)
for every instance in the family. Instance-dependent constants may
occur only as outputs of that scheme; inserting them into the record is
nonuniform advice.

For a resource ceiling \(B\in C\), define the represented budget slice

\[
\mathfrak A^{\mathrm{rep}}_{I,\le B}
=
\left\{
\sem{\rho}_I:
\rho\text{ is family-uniform and }
\operatorname{Cost}_I(\rho)\preceq B
\right\}.
\]

These slices form a filtered represented algebra. The word
``family-uniform'' binds one derivation scheme across all sizes and
instances; it is not re-chosen on each fixed size slice. If \(f_i\) has
a record of cost \(B_i\), an admissible \(k\)-ary constructor produces
\(h(f_1,\ldots,f_k)\) at accumulated charge

\[
B_1\oplus\cdots\oplus B_k
\oplus \operatorname{Cost}(h)
\oplus \operatorname{Cost}_{\mathrm{out}}.
\]

Consequently a fixed slice need not be an algebra: applying a Boolean,
linear, nonlinear, quotient, or fiber constructor can move the result
to a larger slice. The measurable envelope

\[
\mathcal F^{\mathrm{env}}_{I,\le B}
=\sigma\bigl(\mathfrak A^{\mathrm{rep}}_{I,\le B}\bigr)
\subseteq\mathcal F_I
\]

is used only for conditional expectation and atom calculations.
Membership in \(\mathcal F^{\mathrm{env}}_{I,\le B}\) does not assign
cost at most \(B\) to an event or observable; that conclusion requires
its own represented construction record. For a budget structure
\(\mathcal B\), write
\(\mathfrak A^{\mathcal B}_{I_n}
=\mathfrak A^{\mathrm{rep}}_{I_n,\le b(n)}\). A polynomial represented
family is a family \(f=(f_I)\) for which there are one fixed exponent
\(C<\infty\) and one uniform derivation \(\rho\) satisfying

\[
\sem{\rho}_I=f_I,
\qquad
\operatorname{Cost}_I(\rho)\le (1+n)^C
\quad
(n\ge1,\ I\in\mathfrak I_n).
\]

The notation \(\mathfrak A^{\mathrm{poly}}_{I_n}\) denotes the
size-\(n\) slice of these uniformly polynomial families.  The exponent
cannot vary with \(n\). \PartIII extends these static records to the full
derivability closure, including legal lifted computation.

\end{definition}

This definition formalizes the enlargement obtained by adjoining represented observables.
\begin{definition}[Refinement by adding observables]\label{def-observable-refinement}

If \(\operatorname{Obs}_I\subseteq\operatorname{Obs}'_I\) and both
families are read on the same \(X_I\), then

\[
\mathcal F_{\operatorname{Obs}_I}\subseteq\mathcal F_{\operatorname{Obs}'_I}.
\]

The atomic partition \(\mathcal A_{\operatorname{Obs}'_I}\) refines
\(\mathcal A_{\operatorname{Obs}_I}\): every new atom is contained in an
old atom.

\end{definition}

\begin{proof}\phantomsection\label{proof-observable-refinement}

The generators for \(\mathcal F_{\operatorname{Obs}_I}\) are included
among the generators for \(\mathcal F_{\operatorname{Obs}'_I}\), so the
generated \(\sigma\)-algebras are nested. In the finite case, a larger
\(\sigma\)-algebra has more measurable sets and therefore can only split old
atoms into smaller nonempty measurable pieces.

\end{proof}

\Cref{fig-observable-refinement-tree} summarizes the algebra and atom
refinement in \cref{def-observable-algebra,def-observable-refinement}.

\begin{figure}[tbp]
\centering
\begin{tikzpicture}[fw diagram]
  \node[fw source,minimum width=2.8cm] (obs) at (-3.5,2.2)
    {primitive readouts\\\(\operatorname{Obs}_I\)};
  \node[fw source,minimum width=2.8cm] (obsp) at (3.5,2.2)
    {adjoined readouts\\\(\operatorname{Obs}'_I\)};
  \node[fw box,minimum width=2.8cm] (alg) at (-3.5,0.9)
    {\(\mathcal F_{\operatorname{Obs}_I}\)};
  \node[fw box,minimum width=2.8cm] (algp) at (3.5,0.9)
    {\(\mathcal F_{\operatorname{Obs}'_I}\)};
  \node[fw use,minimum width=2.4cm] (coarse) at (-3.5,-0.6)
    {atom \(C\)};
  \node[fw geom,minimum width=1.8cm] (c1) at (1.2,-0.6) {\(C_1\)};
  \node[fw geom,minimum width=1.8cm] (c2) at (3.5,-0.6) {\(C_2\)};
  \node[fw geom,minimum width=1.8cm] (c3) at (5.8,-0.6) {\(C_3\)};
  \draw[fw arrow] (obs) -- (alg);
  \draw[fw arrow] (obsp) -- (algp);
  \draw[fw arrow] (obs) -- node[above,font=\scriptsize] {adjoin} (obsp);
  \draw[fw arrow] (alg) -- node[above,font=\scriptsize] {include} (algp);
  \draw[fw arrow] (alg) -- (coarse);
  \foreach \n in {c1,c2,c3}
    \draw[fw arrow] (algp) -- (\n);
  \node[font=\scriptsize,fwInk] at (1.15,-1.25)
    {\(C=C_1\sqcup C_2\sqcup C_3\)};
\end{tikzpicture}
\caption{Adjoining represented observables enlarges the generated algebra
and refines its atomic partition, as in
\cref{def-observable-algebra,def-observable-refinement}.}
\label{fig-observable-refinement-tree}
\end{figure}

The score observable places a scalar readout inside the same represented algebra as the task data.
\begin{definition}[Score observable]\label{def-score-as-observable}

A score is a real derived observable

\[
D_I:(X_I,\mathcal F_I)\to(\mathbb R,\mathcal B_{\mathbb R}).
\]

Its level sets \(D_I^{-1}(c)\) partition states by score value. The
score has no distinguished status in the algebra except that a terminal
predicate may be defined from one of its level sets.

\end{definition}

\begin{example}[SAT atoms and clause readouts]\label{ex-observable-algebra-sat}

Let \(\Phi=C_1\wedge\cdots\wedge C_m\) and expose only the
clause-satisfaction vector

\[
O_\Phi(x)=(\mathbf 1\{C_1(x)=1\},\ldots,\mathbf 1\{C_m(x)=1\})\in\{0,1\}^m.
\]

Two assignments lie in the same atom exactly when they satisfy the same
clauses. The defect score \(D_\Phi(x)=m-\sum_jO_{\Phi,j}(x)\) is a
coarser observable: its level set \(D_\Phi^{-1}(c)\) contains all atoms
whose readout has exactly \(m-c\) satisfied clauses. Adding variable
projections \(x\mapsto x_i\) refines the atoms, possibly down to
singleton assignments.

\end{example}

\section{Families and Instances}\label{sec-families-instances}
\label{families-and-instances}

A task family specifies a collection of instances, each with a concrete
signature. Static family uniformity requires a common public construction
rule across the family, excluding observables fitted to an individual
solution.

A task family separates the uniform public construction from its individual presented instances.
\begin{definition}[Task family and instance]\label{def-task-family}

A task family is a tuple

\[
\mathfrak T=(\mathfrak I,n,(\mathcal T_I)_{I\in\mathfrak I}),
\]

where \(\mathfrak I\) is an instance set, \(n:\mathfrak I\to\mathbb N\)
is a size map, and each \(I\in\mathfrak I\) has a task-space signature
\(\mathcal T_I=(X_I,\operatorname{Obs}_I,\operatorname{Dyn}_I,\operatorname{End}_I)\).
The size-\(n\) slice is

\[
\mathfrak I_n=\{I\in\mathfrak I:n(I)=n\},
\qquad
\mathfrak I=\bigsqcup_{n\ge1}\mathfrak I_n.
\]

Unless a probability distribution is explicitly specified, a statement
over \(\mathfrak I_n\) is interpreted uniformly over all instances in
the slice.

\end{definition}

Static family uniformity requires one public construction across the task family.
\begin{definition}[Static family-uniform observable structure]\label{def-static-family-uniformity}

A family-indexed observable \(\phi=(\phi_I)_{I\in\mathfrak I}\), with
\(\phi_I:(X_I,\mathcal F_I)\to\mathbb R\), is static family-uniform if
there is one construction rule \(\mathsf\Phi\), depending only on the
public presentation of \(I\), such that

\[
\phi_I=\mathsf\Phi(I)
\qquad\text{for every }I\in\mathfrak I.
\]

The rule may use \(X_I\), \(\operatorname{Obs}_I\),
\(\operatorname{End}_I\), the declared static slot
\(\operatorname{Dyn}_I\), and the public encoding of \(I\). It may not
use per-instance advice absent from the presentation.

\end{definition}

Family uniformity distinguishes persistent structure from regularities
specific to a single finite instance. It requires the same construction
of an observable, quotient, or exposed relation throughout the task
family.

\begin{proposition}[Static family uniformity is independent of algorithmic certification]\label{prop-family-uniformity-not-algorithm-certificate}

Static family uniformity is a property of the observable structures in
the task presentation, independent of membership in a subsequent
budgeted runtime class. A solver, adversary, or inverter need not output
an observable, quotient, invariant, proof trace, or certificate. A
uniform algorithm supplies a public transition rule under a resource
envelope; \PartII embeds each bounded slice of this rule in its clocked
transcript carrier and
decomposes the induced operator.

\end{proposition}

\begin{proof}\phantomsection\label{proof-family-uniformity-not-algorithm-certificate}

\PartI defines family-uniform observables through constructions from the
public presentation before dynamics are introduced. By contrast, a
budget-admissible algorithm is a uniform transition rule on complete
configurations whose instance-dependent inputs lie in the budgeted
observable closure. Family uniformity therefore constrains static
structural claims without imposing an output-certificate requirement on
the algorithm. When the induced operator uses family-uniform structure,
the decomposition in \PartII identifies it. Movement without such
budget-exposed provenance is assigned to the residual component; data
outside the declared observable closure are presentation-inaccessible,
oracle-supplied, or over budget. The polynomial case is obtained by
setting \(\mathcal B=\mathcal B_{\mathrm{poly}}\).

Uniform code and instance-dependent data have distinct status. Fixed
public procedures may include arithmetic, Gaussian elimination,
branching, memory updates, and finite outcome generation. Instance-dependent inputs,
including a modulus, quotient basis, planted label, target pointer, or
value table, must be declared in the presentation, generated by an
admissible construction from the budgeted observable closure, or charged
as advice, oracle data, or presentation enrichment. The criterion is
extensional: an observable has low cost precisely when it admits a
low-cost representation in the observable grammar. Consequently, a
first-hitting-time table, committor, occupation law, or future-success
partition of a subsequent solver is a static family observable only if
it also has an independent representation in the public task algebra.

\end{proof}

The transfer-class margin states the asymptotic separation required after admissible losses are charged.
\begin{definition}[Transfer-class asymptotic margin]\label{def-uniform-margin}

Let \(Q_I\ge0\) be a nonnegative quantity assigned to each instance and
fix the transfer class \(\mathfrak C\) of the declared budget structure.
It has a \(\mathfrak C\)-lower margin on \(\mathfrak T\) when there are
\(\beta\in\mathfrak C\) and \(n_0\in\mathbb N\) such that

\[
\inf_{I\in\mathfrak I_n}Q_I\ge\frac1{\beta(n)}
\qquad\text{for every }n\ge n_0.
\]

It is negligible relative to \(\mathfrak C\) when for every
\(\beta\in\mathfrak C\), \(\sup_{I\in\mathfrak I_n}Q_I\le1/\beta(n)\)
for all sufficiently large \(n\). For \(\mathcal B_{\mathrm{poly}}\),
these are the usual inverse-polynomial lower margin and negligible upper
margin.

\end{definition}

\begin{proposition}[Family structure requires family-uniform
observables]\label{prop-one-instance-not-family}

Let \(\phi_I\) be an observable built for one instance
\(I\in\mathfrak I_n\). If there is no static family-uniform rule
\(\mathsf\Phi\) producing observables with the same claimed margin over
all instances in \(\mathfrak I_n\), then \(\phi_I\) is not structural
evidence for the family.

\end{proposition}

\begin{proof}\phantomsection\label{proof-one-instance-not-family}

A family-level claim is quantified over \(\mathfrak I_n\). An observable
constructed for a single instance provides no corresponding object on
the remaining instances, and an arbitrary extension need not preserve
the claimed margin. The infimum over the slice therefore requires a
public rule that constructs the observable for every instance.

\end{proof}

\begin{example}[Same SAT name, different exposed structure]\label{ex-family-sat}

The family name SAT may include two formulas \(\Phi_1,\Phi_2\) with the
same number of variables but different clause sets. If
\(\operatorname{Obs}_{\Phi_i}\) exposes only clause readouts, then the
atoms are clause-signature cells. If another presentation also exposes
public XOR rows, those rows refine the atoms and can define a quotient.
Both are SAT instances, but their static exposed structures differ. A
planted satisfying assignment is family-uniform only if it is produced
by a public rule in the presentation; otherwise it is instance advice.

\end{example}

\section{Static Target Relevance of
Observables}\label{sec-static-target-relevance}
\label{what-makes-an-observable-useful}

Static target relevance is determined jointly by an observable, the
terminal sectors, the observable algebra, and the task family. The
following conditions characterize whether an observable represents
target structure independently of a solution method.

\subsection{Target-relevance criteria}\label{subsec-static-target-criteria}

We first fix the data against which static target relevance is evaluated.

\begin{definition}[Static target-relevance setting]\label{def-usefulness-environment}

Fix an instance \(I\). Let

\[
T_I=E_I^+=\{x:\operatorname{End}_I^+(x)=1\},
\qquad
F_I=E_I^-=\{x:\operatorname{End}_I^-(x)=1\},
\qquad
S_I=T_I\sqcup F_I.
\]

The sets \(T_I,F_I\in\mathcal F_I\) are disjoint state-terminal sectors.
Failure is \(S_I\setminus T_I\), not \(X_I\setminus T_I\). The
run-dependent stop component \(\operatorname{End}_I^0\) is not part of
this state-terminal split.

Let \(\mu_I\) be a reference probability measure on
\((X_I,\mathcal F_I)\) with \(0<\mu_I(T_I)\) and \(0<\mu_I(F_I)\). When
a degree coordinate is used, let \((\mathcal A_I^{\le d})_{d\ge0}\) be a
finite-dimensional observable-degree filtration with
\(\mathcal A_I^{\le d}\subseteq L^0(X_I,\mathcal F_I;\mathbb R)\). This
filtration is an optional cost coordinate; budget admissibility is
determined by the declared budget structure.

\end{definition}

The complement defining failure is taken within \(S_I\), thereby
preserving the distinction between failed and nonterminal states.

\begin{definition}[Five static target-relevance conditions]\label{def-static-usefulness-clauses}

Let \(\phi=(\phi_I)_{I\in\mathfrak I}\) be an observable family, with
each \(\phi_I:(X_I,\mathcal F_I)\to\mathbb R\) of finite variance. Its
instance clauses are evaluated at \(I\), while affordability and
uniformity use the single family record below.

\textbf{Alignment.} The standardized terminal separation is

\[
\operatorname{Sep}_I(\phi)=
\frac{\mathbb E_{T_I,\mu_I}[\phi]-\mathbb E_{F_I,\mu_I}[\phi]}
{\sqrt{\operatorname{Var}_{\mu_I}(\phi)}}
\]

when the variance is nonzero, and is \(0\) otherwise. Alignment requires
\(|\operatorname{Sep}_I(\phi)|>0\). The sign
\(s_I(\phi)=\operatorname{sign}(\operatorname{Sep}_I(\phi))\) orients
the observable.

\textbf{Extremum-consistency.} With \(\psi=s_I(\phi)\phi\), there is a
threshold \(\theta\in\mathbb R\) such that

\[
T_I=\{x\in S_I:\psi(x)\ge\theta\}
\quad\text{up to }\mu_I\text{-null sets},
\qquad
\operatorname*{argmax}_{x\in X_I}\psi(x)\subseteq T_I.
\]

\textbf{Budget admissibility.} For the chosen budget structure
\(\mathcal B=(\mathcal R,\operatorname{cost},b,\mathfrak C)\),

\[
\operatorname{Constr}_{\mathcal B}(\phi)=1
\quad\Longleftrightarrow\quad
\begin{gathered}
\text{there is one family-uniform derivation }\rho\text{ such that}\\
\sem{\rho}_I=\phi_I,qquad
\operatorname{Cost}_I(\rho)\preceq b(n)
\quad(I\in\mathfrak I_n)
\end{gathered}
\]

on every relevant size slice. Cheap representatives chosen separately
for each instance do not satisfy this clause.

For \(\mathcal B_{\mathrm{poly}}\), a representative chosen from the
optional degree filtration is charged for its band dimension and
evaluation cost. Raw degree is only one possible cost coordinate: a
succinct quotient or other saturated representation need not have low
raw degree.

\textbf{Family-uniformity.} The family \((\phi_I)_{I\in\mathfrak I}\) is
produced by one public rule and its clause margins have a
\(\mathfrak C\)-lower bound over each large size slice.

\textbf{Monotone terminal filtration.} With \(\psi=s_I(\phi)\phi\),
define superlevel sets \(H_t=\{x\in X_I:\psi(x)\ge t\}\). The terminal
purity function

\[
\pi_I(t)=
\frac{\mu_I(T_I\cap H_t)}{\mu_I(S_I\cap H_t)}
\]

where the denominator is positive, is nondecreasing in \(t\), and some
threshold \(\theta\) satisfies \(S_I\cap H_\theta=T_I\) up to null sets.

\end{definition}

\Cref{fig-static-target-relevance-gates} displays the five simultaneous
conditions in \cref{def-static-usefulness-clauses}.

\begin{figure}[tbp]
\centering
\begin{tikzpicture}[fw diagram,node distance=5mm]
  \node[fw source,minimum width=2.25cm] (align) {alignment};
  \node[fw geom,minimum width=2.25cm,right=of align] (ext) {target extremum};
  \node[fw use,minimum width=2.25cm,right=of ext] (cost) {budget};
  \node[fw provenance,minimum width=2.25cm,right=of cost] (unif) {uniform rule};
  \node[fw terminal,minimum width=2.25cm,right=of unif] (filt) {filtration};
  \node[fw box,minimum width=4.0cm,below=11mm of cost] (rel)
    {static target relevance\\\(U^{\mathrm{stat}}_{\mathcal B,n}(\phi)\)};
  \foreach \x in {align,ext,cost,unif,filt}
    \draw[fw arrow] (\x) -- (rel);
\end{tikzpicture}
\caption{Static target relevance is a same-observable conjunction.  Target
separation and extremum consistency are combined with a charged construction,
one family rule, and a monotone terminal filtration; failure of any gate sets
the product in \cref{def-static-usefulness-core} to zero.}
\label{fig-static-target-relevance-gates}
\end{figure}

Alignment measures separation between successful and failed terminal
states. Extremum consistency requires the selected extremum to coincide
with the success set. Budget admissibility places the observable in the
budgeted subalgebra, and family uniformity requires a common public
construction. The final condition orders the observable's level sets as
a nested terminal filtration toward success, independently of any
gradient, flow, or traversal.

The static target-relevance functional combines alignment, consistency, affordability, uniformity, and filtration into one coordinate.
\begin{definition}[Static target-relevance functional under \(\mathcal B\)]\label{def-static-usefulness-core}

Let \(\operatorname{ExtCons}_I(\phi)\),
\(\operatorname{Constr}_{\mathcal B}(\phi)\), and
\(\operatorname{Fil}_I(\phi)\) be the indicators of
extremum-consistency, budget admissibility, and monotone terminal
filtration. The instance-level static target-relevance functional is

\[
U^{\mathrm{stat}}_{I,\mathcal B}(\phi)=
|\operatorname{Sep}_I(\phi)|
\operatorname{ExtCons}_I(\phi)
\operatorname{Constr}_{\mathcal B}(\phi)
\operatorname{Fil}_I(\phi).
\]

For an observable family \(\phi=(\phi_I)\), let
\(\operatorname{Unif}_{\mathfrak T}(\phi)\) be one when the family is
produced by a single public construction rule and zero otherwise. Set

\[
U^{\mathrm{stat}}_{\mathcal B,n}(\phi)=
\operatorname{Unif}_{\mathfrak T}(\phi)
\inf_{I\in\mathfrak I_n}U^{\mathrm{stat}}_{I,\mathcal B}(\phi_I).
\]

The observable is statically target-relevant under \(\mathcal B\) on the
size-\(n\) slice if
\(U^{\mathrm{stat}}_{\mathcal B,n}(\phi)\ge 1/\beta_{\mathfrak C}(n)\) for
the transfer-class margin of the budget. For
\(\mathcal B_{\mathrm{poly}}\),
\(\beta_{\mathfrak C}(n)=\operatorname{poly}(n)\), giving the standard
inverse-polynomial condition.

\end{definition}

In the instance-level display,
\(\operatorname{Constr}_{\mathcal B}(\phi)\) is the single family
indicator just defined; it is not a new pointwise existential
quantifier. The other terms carry the implicit subscript \(I\).

\begin{lemma}[Necessity of the static target-relevance conditions]\label{thm-static-clauses-necessary}

An observable family has transfer-class static target relevance under
\(\mathcal B\) only if it is family-uniform, its alignment has a
transfer-class lower margin, and the extremum-consistency, budget, and
filtration indicators equal one on every sufficiently large size
slice. If family uniformity fails, one of the three indicators vanishes
on some instance in every large slice, or the alignment supremum on the
slice is negligible relative to \(\mathfrak C\), then
\(U^{\mathrm{stat}}_{\mathcal B,n}(\phi)\) is respectively zero or
negligible.

\end{lemma}

\begin{proof}\phantomsection\label{proof-static-clauses-necessary}

By definition,

\[
U^{\mathrm{stat}}_{\mathcal B,n}(\phi)
=
\operatorname{Unif}_{\mathfrak T}(\phi)
\inf_{I\in\mathfrak I_n}
\left(
|\operatorname{Sep}_I(\phi_I)|
\operatorname{ExtCons}_I(\phi_I)
\operatorname{Constr}_{\mathcal B}(\phi_I)
\operatorname{Fil}_I(\phi_I)
\right).
\]

The first factor is zero when family uniformity fails. Each of the last
three factors inside the infimum is an indicator, so its failure on an
instance makes the infimum zero. If the alignment supremum is bounded
by a negligible function on the size slice, then the entire expression
is bounded by the same function because all indicators lie in
\([0,1]\). Conversely, a positive transfer-class lower margin for the
product forces the uniformity factor and all three indicators to be one
and forces the alignment infimum to have that margin. These are exactly
the five stated conditions.

\end{proof}

\subsection{Projection ceilings for generated algebras}
\label{subsec-target-projection-ceilings}

The target projection mass quantifies how much terminal information is present
in a generated observable algebra.

\begin{definition}[Target projection mass of a generated observable algebra]\label{def-generated-algebra-target-projection}

Let \(\mathfrak A_I\subseteq\mathcal F_I\) be a sub-\(\sigma\)-algebra
generated by declared observables, represented quotients, or any
specified family-uniform observable constructors. Let
\(P_{\mathfrak A_I}\) denote conditional expectation onto
\(L^2(X_I,\mathfrak A_I,\mu_I)\). For a target sector \(T_I\) with
\(0<\mu_I(T_I)<1\), define the centered target contrast and the target
projection mass by

\[
y_{T_I}=1_{T_I}-\mu_I(T_I),
\qquad
\operatorname{ProjMass}_{\mathfrak A_I}(T_I)
=
\frac{\|P_{\mathfrak A_I}y_{T_I}\|_{L^2(\mu_I)}^2}{\|y_{T_I}\|_{L^2(\mu_I)}^2}.
\]

Equivalently, if the atoms or fibers of \(\mathfrak A_I\) are denoted by
\(C\), then \(P_{\mathfrak A_I}1_{T_I}\) is the conditional target
density \(\mu_I(T_I\mid C)\). The projection mass is the normalized
variance of this conditional target density.

\end{definition}

\begin{lemma}[Generated algebra projection ceiling]\label{thm-generated-algebra-projection-ceiling}

Fix a task instance \(I\), a target sector \(T_I\), and a generated
observable algebra \(\mathfrak A_I\subseteq\mathcal F_I\). For every
square-integrable observable \(g\) measurable with respect to
\(\mathfrak A_I\),

\[
\frac{|\langle g-\mathbb E_{\mu_I}g,\,y_{T_I}\rangle_{L^2(\mu_I)}|^2}
{\|g-\mathbb E_{\mu_I}g\|_{L^2(\mu_I)}^2\,\|y_{T_I}\|_{L^2(\mu_I)}^2}
\le
\operatorname{ProjMass}_{\mathfrak A_I}(T_I),
\]

with the left side defined as \(0\) when \(g\) is constant. Hence every
postprocessing, nonlinear combination, quotient readout, threshold,
kernel feature, learned readout, or candidate monotone potential
measurable in \(\mathfrak A_I\) has target correlation bounded by the
target projection mass of \(\mathfrak A_I\).

If \(\mu_I(S_I)=1\), write \(p_I=\mu_I(T_I)\).  Then the normalized
correlation and the separation coordinate of
\cref{def-static-usefulness-clauses} satisfy the exact identity

\[
\frac{|\langle g-\mathbb E_{\mu_I}g,y_{T_I}\rangle|^2}
{\|g-\mathbb E_{\mu_I}g\|_2^2\|y_{T_I}\|_2^2}
=p_I(1-p_I)|\operatorname{Sep}_I(g)|^2.
\]

Consequently,

\[
|\operatorname{Sep}_I(g)|^2
\le
\frac{\operatorname{ProjMass}_{\mathfrak A_I}(T_I)}
{p_I(1-p_I)}.
\]

If \(\mu_I(S_I)=1\) throughout a family of generated algebras
\((\mathfrak A_I)_{I\in\mathfrak I_n}\) satisfies
\[
\sup_{I\in\mathfrak I_n}
\frac{\operatorname{ProjMass}_{\mathfrak A_I}(T_I)}
{p_I(1-p_I)}
=\operatorname{negl}_{\mathfrak C}(n),
\]
then every family-uniform observable measurable in those algebras has
squared separation bounded by the same negligible function and therefore
cannot have a transfer-class lower margin.  In particular, it is enough to
combine negligible projection mass with a transfer-class lower bound on
\(p_I(1-p_I)\) strong enough to preserve negligibility after division.
Extremum consistency and monotone terminal filtration restrict this
class further and preserve the same projection ceiling.

\end{lemma}

\begin{proof}\phantomsection\label{proof-generated-algebra-projection-ceiling}

Let \(\bar g=g-\mathbb E_{\mu_I}g\). Since \(g\) is
\(\mathfrak A_I\)-measurable,
\(\bar g\in L^2(X_I,\mathfrak A_I,\mu_I)\). Orthogonal projection gives

\[
\langle \bar g,y_{T_I}\rangle
=
\langle \bar g,P_{\mathfrak A_I}y_{T_I}\rangle.
\]

Cauchy's inequality gives
\(|\langle \bar g,P_{\mathfrak A_I}y_{T_I}\rangle|^2\le\|\bar g\|_2^2\|P_{\mathfrak A_I}y_{T_I}\|_2^2\).
Dividing by \(\|\bar g\|_2^2\|y_{T_I}\|_2^2\) proves the displayed
bound. Any deterministic or learned postprocessing of the generating
observables remains \(\mathfrak A_I\)-measurable, so it satisfies the
same bound.  When \(\mu_I(S_I)=1\),
\(\|y_{T_I}\|_2^2=p_I(1-p_I)\) and
\[
\langle \bar g,y_{T_I}\rangle
=p_I(1-p_I)
\bigl(\mathbb E_{T_I,\mu_I}g-\mathbb E_{F_I,\mu_I}g\bigr),
\]
the two normalization identities follow by direct substitution. The
family statement follows by taking the supremum over the size slice only
after retaining the factor \(p_I(1-p_I)\); this is the required
rare-target normalization. Extremum-consistency and monotone terminal filtration are
additional predicates on the same measurable readouts; they select a
subclass and therefore cannot enlarge the Hilbert projection onto
\(L^2(\mathfrak A_I)\).

\end{proof}

\begin{lemma}[Affordable composition and target-projection monotonicity]
\label{thm-affordable-composition-no-creation}

Let \(f_i:X_I\to Y_i\) be represented observables and let

\[
J=(f_1,\ldots,f_k):X_I\longrightarrow Y_1\times\cdots\times Y_k
\]

be their represented joint observable. Let \(h\) be a represented
postprocessing on \(\operatorname{im}J\), and set \(q=h\circ J\).
If \(P_J\) and \(P_q\) are conditional expectation onto
\(L^2(\sigma(J))\) and \(L^2(\sigma(q))\), respectively, then

\[
P_qP_J=P_JP_q=P_q
\]

and, for the centered target contrast \(y_{T_I}\),

\begin{equation}
\label{eq-affordable-composition-pythagoras}
\|P_Jy_{T_I}\|_2^2
=
\|P_qy_{T_I}\|_2^2
+
\|(P_J-P_q)y_{T_I}\|_2^2.
\end{equation}

In particular,

\begin{equation}
\label{eq-affordable-composition-monotonicity}
\frac{\|P_qy_{T_I}\|_2^2}{\|y_{T_I}\|_2^2}
\le
\frac{\|P_Jy_{T_I}\|_2^2}{\|y_{T_I}\|_2^2}.
\end{equation}

Suppose records \(\rho_i\) represent the \(f_i\), a record
\(\rho_{\mathrm{join}}\) represents their joint fibers, and a record
\(\rho_h\) represents the evaluation and fiber operations of \(h\).
Then \(q\) has the composite record
\(\rho_q=\rho_h\circ\rho_{\mathrm{join}}(\rho_1,\ldots,\rho_k)\), with

\begin{equation}
\label{eq-affordable-composition-cost}
\operatorname{Cost}_I(\rho_q)
\preceq
\bigoplus_{i=1}^k\operatorname{Cost}_I(\rho_i)
\oplus\operatorname{Cost}_I(\rho_{\mathrm{join}})
\oplus\operatorname{Cost}_I(\rho_h).
\end{equation}

Thus the postprocessing can select or compress a target-aligned
component of the joint observable, while its availability at a budget
is determined by the complete represented derivation
\eqref{eq-affordable-composition-cost}.

\end{lemma}

\begin{proof}\phantomsection\label{proof-affordable-composition-no-creation}
Since \(q=h\circ J\), one has
\(\sigma(q)\subseteq\sigma(J)\). Conditional expectations onto the nested
closed subspaces \(L^2(\sigma(q))\subseteq L^2(\sigma(J))\) satisfy
\(P_qP_J=P_JP_q=P_q\). Hence \(P_qy_{T_I}\) and
\((P_J-P_q)y_{T_I}\) are orthogonal and sum to \(P_Jy_{T_I}\), which
proves \eqref{eq-affordable-composition-pythagoras} and
\eqref{eq-affordable-composition-monotonicity}. The composite record is
the directed acyclic graph obtained by adjoining the join and
postprocessing vertices to the input records. The budget axioms give
\eqref{eq-affordable-composition-cost}.
\end{proof}

\begin{corollary}[Terminal projection ceiling for an expanded derivation]
\label{cor-expanded-derivation-projection-ceiling}

Let \(\widehat\rho\) be an expanded derivation whose terminal vertex is
a represented observable or quotient \(q_{\widehat\rho}\). Let
\(J_{\widehat\rho}\) be the joint observable formed from all represented
inputs to that terminal vertex, including retained lifted coordinates.
Then a represented map \(h_{\widehat\rho}\) satisfies

\[
q_{\widehat\rho}
=h_{\widehat\rho}\circ J_{\widehat\rho},
\qquad
\frac{\|P_{q_{\widehat\rho}}y_{T_I}\|_2^2}
     {\|y_{T_I}\|_2^2}
\le
\frac{\|P_{J_{\widehat\rho}}y_{T_I}\|_2^2}
     {\|y_{T_I}\|_2^2}.
\]

The records for \(J_{\widehat\rho}\),
\(h_{\widehat\rho}\), and the terminal fiber interface are subrecords
of \(\widehat\rho\) and retain their complete charges. Thus a terminal
postprocessing may select target-aligned mass already present in its
joint represented inputs, but cannot increase their target projection.

\end{corollary}

\begin{proof}\phantomsection\label{proof-expanded-derivation-projection-ceiling}

The incoming edges of the terminal vertex give the joint map
\(J_{\widehat\rho}\), and the vertex evaluation rule gives
\(h_{\widehat\rho}\). Apply
\Cref{thm-affordable-composition-no-creation}. The expanded
record retains the incoming records and the terminal constructor charge
by \Cref{def-expanded-derivation-normal-form}.

\end{proof}

\begin{example}[Affordable XOR synergy]
\label{ex-affordable-xor-synergy}

Two binary observables \(f\) and \(g\) may each have zero target
projection while \(J=(f,g)\) determines the target through
\(q=f\oplus g\). In that case the target component belongs to
\(L^2(\sigma(f,g))\), and the XOR record gives it a constant-overhead
readout. An arbitrary table on \(\operatorname{im}(f,g)\) has the same
measurability relation but enters a budget slice only through a uniform
record that constructs, stores, evaluates, and represents its fibers.
An instance-dependent table inserted without such a record is
nonuniform advice. A succinct family-uniform formula is an admissible
representation and must be included at its actual cost, independently
of how that formula was discovered.

\end{example}

The tuple map and terminal readout are distinct components of the
presentation. Observables \(\Phi=(\phi_1,\ldots,\phi_r)\) may be exposed
without identifying their successful fibers. Subsequent parts account
for any derivation, search, or learning required to construct this
readout.

\section{Target-Discrimination Certificates and Budgeted Extraction}
\label{sec-target-discrimination-extraction}

We next isolate the represented evidence required for a composite observable to
separate the target.

\begin{definition}[Target-discrimination certificate]\label{def-target-discrimination-certificate}

Let \(\mathfrak A_I\subseteq\mathcal F_I\) be a generated observable
algebra and let \(T_I\subseteq S_I\) be the success sector. A
target-discrimination certificate for \((\mathfrak A_I,T_I)\) is
represented data in the instance presentation, or in a declared
family-uniform budget-admissible generated class, that gives one of the
following terminal readouts:

\begin{enumerate}

\item

an exact fiber readout: \(T_I\) is a union of atoms or fibers of
\(\mathfrak A_I\), together with a represented map
\(r_I:X_I/\mathfrak A_I\to\{0,1\}\) whose pullback is \(1_{T_I}\) on
terminals;

\item

a threshold or extremum readout: an \(\mathfrak A_I\)-measurable
observable and threshold satisfying extremum-consistency;

\item

a monotone terminal filtration with certified transfer-class purity
margins;

\item

an enrichment certificate: a represented family of cells or fibers with
transfer-class separation in conditional target density.

\end{enumerate}

For a tuple \(\Phi=(\phi_1,\ldots,\phi_r)\), the tuple map alone is an
uncalibrated observable. It defines a target-discriminating quotient only
when the certificate is represented or generated. A learned certificate
incurs the sample, capacity, and selection costs treated in \cref{part:learning}.
An instance-specific certificate is nonuniform advice; one outside the
declared observable closure is presentation-inaccessible rather than
exposed structure.

The certificate represents the terminal interpretation of the fibers;
it is not a proof that the observable is aligned and is not an output
required from an algorithm. Target alignment is the objective
projection quantity above. When the terminal verifier is public, its
readout supplies terminal compatibility, while the projection and later
dynamic utility are evaluated by the framework.

\end{definition}

\begin{lemma}[Observable combinations require certified target discrimination]\label{thm-raw-combinations-require-discrimination}

Fix an instance \(I\), a generated observable algebra \(\mathfrak A_I\),
and a target sector \(T_I\). A family-uniform observable \(g\)
measurable in \(\mathfrak A_I\) can contribute static target-relevant structure
only through the target projection mass of \(\mathfrak A_I\) together
with a target-discrimination certificate satisfying the static
target-relevance conditions.

Consequently, if \(\operatorname{ProjMass}_{\mathfrak A_I}(T_I)\) is
negligible on a family size slice, every \(\mathfrak A_I\)-measurable
combination has negligible normalized target extraction.  A conclusion for
the separation coordinate additionally uses the terminal-balance normalization
of \cref{thm-generated-algebra-projection-ceiling}. If the projection mass is
non-negligible but no target-discrimination certificate is represented
or generated, then the tuple or algebra remains uncalibrated. Subsequent
use of the absent readout is accounted for as learning and selection
cost, residual structure, advice, oracle access, or presentation
enrichment. It does not define an \(\Abs\) source, a
\(\Caus\) potential, or a target \(\Bdry\) readout in
\PartI.

If extremum-consistency or monotone terminal filtration fails, the
observable has zero static target relevance for extremal or monotone
uses. It may still contribute through a separately certified nonmonotone
quotient or readout, with that readout included in the represented
abstraction.

\end{lemma}

\begin{proof}\phantomsection\label{proof-raw-combinations-require-discrimination}

The projection ceiling bounds the target correlation of every
\(\mathfrak A_I\)-measurable postprocessing by
\(\operatorname{ProjMass}_{\mathfrak A_I}(T_I)\). Static target relevance
also requires the five stated conditions. A tuple map specifies fibers
but not their terminal labels; the target-discrimination certificate
supplies this additional represented readout. It makes no assertion
about the size of the projection and need not be produced by the
algorithm. Without it, selecting successful
fibers requires data outside the exposed observable algebra and is
classified as nonuniform, presentation-inaccessible, or residual
information. Extremum consistency and monotone terminal filtration
further restrict the measurable readouts without increasing projection
mass or assigning terminal meaning in the absence of a certificate.

\end{proof}

\Cref{fig-target-extraction-audit} collects the projection and certification
stages used by the budgeted extraction ledger.

\begin{figure}[tbp]
\centering
\begin{tikzpicture}[fw diagram]
  \node[fw source,minimum width=3.6cm] (slice) at (0,2.4)
    {represented budget slice\\\(\mathfrak A^{\mathrm{rep}}_{I,\le B}\)};
  \node[fw geom,minimum width=3.2cm] (env) at (-3.4,0.9)
    {measurable envelope\\\(\mathfrak A^{\mathrm{env}}_{I,B}\)};
  \node[fw provenance,minimum width=3.2cm] (cert) at (3.4,0.9)
    {target-discrimination\\certificate};
  \node[fw geom,minimum width=3.2cm] (proj) at (-3.4,-0.7)
    {\(\operatorname{ProjMass}\)};
  \node[fw use,minimum width=3.2cm] (readout) at (3.4,-0.7)
    {certified readouts\\\(\mathcal R^{\mathrm{cert}}_{I,B}\)};
  \node[fw terminal,minimum width=4.2cm] (act) at (0,-2.2)
    {\(\operatorname{ActMass}_{I,B}\le\operatorname{ProjMass}\)};
  \draw[fw arrow] (slice) -- (env);
  \draw[fw arrow] (slice) -- (cert);
  \draw[fw arrow] (env) -- (proj);
  \draw[fw arrow] (cert) -- (readout);
  \draw[fw arrow] (proj) -- (act);
  \draw[fw arrow] (readout) -- (act);
\end{tikzpicture}
\caption{The generated algebra bounds the available target projection,
while the represented target-discrimination certificate selects the
affordable actionable part.  The objects are those of
\cref{def-generated-algebra-target-projection,def-target-discrimination-certificate,def-budgeted-extraction-ledger}.}
\label{fig-target-extraction-audit}
\end{figure}

\subsection{Budgeted extraction ledgers and ceilings}
\label{subsec-budgeted-extraction-ledgers}

The following ledger records target mass together with the cost of its complete
represented derivation.

\begin{definition}[Budgeted extraction accounting]\label{def-budgeted-extraction-ledger}

Fix an instance \(I\), target sector \(T_I\), and public budget \(B\).
Let \(\mathfrak A^{\mathrm{rep}}_{I,\le B}\) be the represented budget
slice of Definition
\hyperref[def-budgeted-observable-closure-static]{Budget-filtered
represented observable algebra}, and let
\(\mathfrak A^{\mathrm{env}}_{I,B}\) denote its measurable envelope.
For a square-integrable readout \(g\), define its normalized target
extraction by

\[
\operatorname{Ext}_I(g;T_I)
=
\frac{|\langle g-\mathbb E_{\mu_I}g,\,y_{T_I}\rangle_{L^2(\mu_I)}|^2}
{\|g-\mathbb E_{\mu_I}g\|_{L^2(\mu_I)}^2\,\|y_{T_I}\|_{L^2(\mu_I)}^2},
\]

with value \(0\) when \(g\) is constant. Let
\(\mathcal R^{\mathrm{cert}}_{I,B}(T_I)\) be the readouts
\(g\in\mathfrak A^{\mathrm{rep}}_{I,\le B}\) whose own uniform record,
represented fibers, and target-discrimination readout have total cost at
most \(B\) and satisfy the static target-relevance conditions. The
actionable target mass at budget \(B\) is

\[
\operatorname{ActMass}_{I,B}(T_I)
=
\sup_{g\in\mathcal R^{\mathrm{cert}}_{I,B}(T_I)}
\operatorname{Ext}_I(g;T_I),
\]

where the supremum is \(0\) if the certified class is empty. Thus
\(\operatorname{ProjMass}_{\mathfrak A^{\mathrm{env}}_{I,B}}(T_I)\)
measures the target signal present in the measurable envelope, while
\(\operatorname{ActMass}_{I,B}(T_I)\) measures the part of that signal
that has a represented or generated terminal readout at total cost at
most \(B\). The latter supremum does not range over unrepresented
members of the envelope.

\end{definition}

These ceilings isolate the target mass attributable to each audited extraction channel.
\begin{definition}[Audit ceilings for extracted target mass]\label{def-extraction-audit-ceilings}

For a certified budget-\(B\) readout class
\(\mathcal R\subseteq\mathcal R^{\mathrm{cert}}_{I,B}(T_I)\), the
empirical, effective-dimension, and residual-order ceilings are the bounds of
\cref{def-empirical-extraction-audit,def-effective-dimension-extraction-audit,def-residual-order-extraction-audit}.
Each ceiling enters an extraction bound only when its hypotheses and complete
certificate hold for the same readout class and probability space.

\end{definition}

\begin{definition}[Empirical extraction-audit ceiling]
\label{def-empirical-extraction-audit}

A learned or searched class
\(\mathcal H_B\) has a one-sided audit penalty
\(\mathfrak P_{\mathcal H_B}(M,\delta)\) when there is an empirical
audit statistic \(\widehat{\operatorname{Aud}}_M(h;T_I)\) such that,
from \(M\) audit samples or independent family instances, with
probability at least \(1-\delta\),

\[
\sup_{h\in\mathcal H_B}
\operatorname{Ext}_I(h;T_I)
\le
\sup_{h\in\mathcal H_B}
\widehat{\operatorname{Aud}}_{M}(h;T_I)
+
\mathfrak P_{\mathcal H_B}(M,\delta).
\]

\end{definition}

\begin{definition}[Effective-dimension extraction ceiling]
\label{def-effective-dimension-extraction-audit}

A kernel or spectral band
\((K,\lambda)\) has captured target mass \(m_K(\lambda)\) when every
readout in the band lies in the spectral subspace or ridge-effective
span counted by \(N_{\mathrm{eff}}(K,\lambda)\) and has target
extraction at most \(m_K(\lambda)\). The quantity
\(N_{\mathrm{eff}}(K,\lambda)\) is charged as statistical mode count in
the learning budget.

\end{definition}

\begin{definition}[Residual-order extraction ceiling]
\label{def-residual-order-extraction-audit}

A residual-noise certificate of order
\(q\) and flip rate \(\eta\in(0,1/2)\) means that the target-relevant
correlation of every nonconstant readout \(g\) in the class admits a
factorization

\[
\frac{\langle g-\mathbb E g,y_{T_I}\rangle}
{\|g-\mathbb E g\|_2\|y_{T_I}\|_2}
=\beta^{r_g}a_g,
\qquad r_g\ge q,
\qquad |a_g|\le1,
\]

where \(\beta=1-2\eta\). Such a factorization is supplied, for example,
when the target-relevant covariance passes through \(r_g\) independent
mean-\(\beta\) residual signs and the remaining normalized correlation
has modulus at most one. The resulting normalized target extraction is
bounded by \(\beta^{2q}\).

\end{definition}

\begin{lemma}[Budgeted extraction bound]\label{thm-budgeted-extraction-bound}

Fix \(I\), \(T_I\), budget \(B\), and a certified readout class
\(\mathcal R\subseteq\mathcal R^{\mathrm{cert}}_{I,B}(T_I)\). Then

\[
\sup_{g\in\mathcal R}\operatorname{Ext}_I(g;T_I)
\le
\operatorname{ActMass}_{I,B}(T_I)
\le
\operatorname{ProjMass}_{\mathfrak A^{\mathrm{env}}_{I,B}}(T_I).
\]

If \(\mathcal R\subseteq\mathcal H_B\) and the empirical audit ceiling
holds, then with probability at least \(1-\delta\),

\[
\sup_{g\in\mathcal R}\operatorname{Ext}_I(g;T_I)
\le
\sup_{h\in\mathcal H_B}\widehat{\operatorname{Aud}}_{M}(h;T_I)
+
\mathfrak P_{\mathcal H_B}(M,\delta).
\]

If \(\mathcal R\) lies in a budget-admissible kernel band
\((K,\lambda)\), then

\[
\sup_{g\in\mathcal R}\operatorname{Ext}_I(g;T_I)
\le
m_K(\lambda).
\]

If every \(g\in\mathcal R\) has residual-noise order at least \(q\) at
flip rate \(\eta\), then

\[
\sup_{g\in\mathcal R}\operatorname{Ext}_I(g;T_I)
\le
(1-2\eta)^{2q}.
\]

Consequently, when all displayed audits apply,

\[
\begin{aligned}
\sup_{g\in\mathcal R}\operatorname{Ext}_I(g;T_I)
\le \min\Bigl\{&
\operatorname{ProjMass}_{\mathfrak A^{\mathrm{env}}_{I,B}}(T_I),
\operatorname{ActMass}_{I,B}(T_I),\\
&\sup_{h\in\mathcal H_B}\widehat{\operatorname{Aud}}_{M}(h;T_I)
  +\mathfrak P_{\mathcal H_B}(M,\delta),\\
&m_K(\lambda),
(1-2\eta)^{2q}
\Bigr\}.
\end{aligned}
\]

A term enters the bound when its certificate belongs to the declared
presentation or is generated by the budgeted constructor grammar.
An uncertified term is not a member of the certified supremum at that
budget. If every applicable certified ceiling in the displayed minimum is
negligible, the certified extraction of the readout class is likewise
negligible.

\end{lemma}

\begin{proof}\phantomsection\label{proof-budgeted-extraction-bound}

The first inequality is the definition of
\(\operatorname{ActMass}_{I,B}\), since
\(\mathcal R\subseteq\mathcal R^{\mathrm{cert}}_{I,B}(T_I)\). The second
is the generated-algebra projection ceiling applied to the algebra
\(\mathfrak A^{\mathrm{env}}_{I,B}\). The empirical audit inequality is exactly the
certified one-sided uniform concentration statement for
\(\mathcal H_B\), restricted to the subclass \(\mathcal R\). The kernel
inequality follows because the band projection contains only the target
spectral mass \(m_K(\lambda)\); every readout in the band is a vector in
that subspace or ridge-effective span. For the residual-order
inequality, the certificate gives a normalized correlation
\(\beta^{r_g}a_g\) with \(r_g\ge q\) and \(|a_g|\le1\). Squaring and
using \(0<\beta<1\) gives
\(\operatorname{Ext}_I(g;T_I)\le\beta^{2r_g}\le\beta^{2q}\).
Since each displayed quantity is an upper bound for
the same nonnegative supremum, the supremum is bounded by their minimum.

\end{proof}

\begin{corollary}[Vanishing ceilings give vanishing extraction]\label{cor-vanishing-ceiling-vanishing-extraction}

For a family \((I_n,T_{I_n})\), let \(B(n)\preceq b(n)\) be a declared
admissible budget level and let \(\mathcal R_n\) be certified readout
classes at that budget. If any applicable ceiling in Theorem
\hyperref[thm-budgeted-extraction-bound]{Budgeted extraction bound} is
negligible relative to the transfer class uniformly over the size slice,
then

\[
\sup_{g\in\mathcal R_n}\operatorname{Ext}_{I_n}(g;T_{I_n})
=
\operatorname{negl}_{\mathfrak C}(n).
\]

In particular, combinations from an algebra with negligible target
projection mass cannot produce a transfer-class target readout. An
uncertified projection remains outside the statically actionable class.
When learning or search constructs the readout, its empirical audit
penalty, effective dimension, selection cost, and residual-noise order
enter the accounting before the extracted mass is included.

\end{corollary}

\begin{proof}\phantomsection\label{proof-vanishing-ceiling-vanishing-extraction}

The theorem bounds the nonnegative extraction supremum by each
applicable ceiling. A uniformly negligible ceiling therefore forces the
supremum to be negligible. The final statements are the projection,
certification, and learning-charge cases of the same bound.

\end{proof}

\section{Finite Computable Target Audits}
\label{sec-finite-target-audits}

The abstract extraction ledger admits a finite matrix audit on each declared
observable algebra.

\begin{definition}[Computable observable audit]\label{def-computable-observable-audit}

For a finite presented instance \(I\), let
\(S=\{x_1,\ldots,x_m\}=\operatorname{supp}\mu_I\), set
\(t_i=1_{T_I}(x_i)\), and use the exact reference weights
\(w_i=\mu_I(x_i)>0\).  Given an exposed feature matrix
\(\Phi\in\mathbb R^{m\times r}\), optional quotient fibers \(c:S\to C\),
optional kernel Gram matrix \(K\succeq0\), and declared audit parameters
\((M,\delta,\eta,q,\lambda)\), the computable observable audit is the
tuple

\[
\mathsf{Audit}(I,B)
=
\bigl(
\widehat{\operatorname{ProjMass}},
\widehat{\operatorname{ActMass}},
\widehat{\operatorname{InertMass}},
\widehat{\operatorname{ThrMargin}},
\widehat{\operatorname{MonFil}},
\widehat N_{\mathrm{eff}},
\widehat m_K,
\widehat{\mathfrak P},
\widehat R_q,
\widehat Z_{\mathrm{null}},
\widehat U_{\mathrm{cert}}
\bigr).
\]

Here \(\widehat{\operatorname{ProjMass}}\) is the weighted least-squares
projection of the centered target contrast onto the exposed feature span
or the conditional-density projection onto quotient fibers;
\(\widehat{\operatorname{ActMass}}\) is the largest extraction among
represented readouts passing the target-discrimination certificate;
\(\widehat{\operatorname{InertMass}}=\max\{0,\widehat{\operatorname{ProjMass}}-\widehat{\operatorname{ActMass}}\}\);
\(\widehat{\operatorname{ThrMargin}}\) and
\(\widehat{\operatorname{MonFil}}\) measure threshold separation and
monotone terminal filtration for the audited scores;
\(\widehat N_{\mathrm{eff}}=\operatorname{Tr}(K(K+\lambda I)^{-1})\);
\(\widehat m_K\) is the target spectral mass captured by the
ridge-effective kernel modes; \(\widehat{\mathfrak P}\) is an empirical
Rademacher or covering penalty for the searched class;
\(\widehat R_q=(1-2\eta)^{2q}\) is the residual-order ceiling;
\(\widehat Z_{\mathrm{null}}\) is the random-label or permutation null
score; and \(\widehat U_{\mathrm{cert}}\) is the minimum of all
applicable certified ceilings.

The equality \(S=\operatorname{supp}\mu_I\) and the exact weights are part
of this definition.  Replacing them by observations from a sample produces
the sampled audit of \cref{def-sampled-computable-observable-audit}, not an
exact population projection.

\end{definition}

\begin{definition}[Sampled computable observable audit]
\label{def-sampled-computable-observable-audit}

Let \(Z_1,\ldots,Z_N\) be an audit sample from \(\mu_I\), independent of
any training data used to construct the audited class.  Let
\(\widehat C_1,\ldots,\widehat C_J\) be the empirical versions of the
applicable projection, actionable-mass, threshold, filtration, kernel, or
residual-order ceilings.  A simultaneous sampling certificate consists of
nonnegative radii \(\varepsilon_j(N,\delta)\), proved uniformly over every
readout and every output-selection rule admitted by the audited class, such
that the single event

\[
\mathcal E_{N,\delta}
=
\bigcap_{j=1}^J
\left\{C_j\le \widehat C_j+\varepsilon_j(N,\delta)\right\}
\]

has probability at least \(1-\delta\).  Here each \(C_j\) is the
corresponding population ceiling on the same probability space.  Define

\[
\widehat U_{\mathrm{samp}}
=
\min_{1\le j\le J}
\min\{1,\widehat C_j+\varepsilon_j(N,\delta)\}.
\]

If a coordinate has no proved simultaneous radius, it is omitted from the
minimum.  In particular, a pointwise concentration bound for a fixed readout
does not certify a class chosen using the audit sample.

\end{definition}

\begin{lemma}[Soundness of the computable observable audit]\label{thm-computable-audit-soundness}

Assume the exact finite audit data of
\cref{def-computable-observable-audit} are part of the declared presentation
or are generated by the budget-\(B\) observable constructor grammar. Let
\(\mathcal R_{I,B}^{\mathrm{aud}}\) be the audited readout class: all
represented readouts measurable in the exposed feature span, quotient
algebra, or kernel band that pass the stated target-discrimination
certificate and whose search procedure is covered by the declared
empirical audit. Then every quantity in \(\mathsf{Audit}(I,B)\) is
computable from the finite matrices, labels, certificates, and audit
parameters. Moreover, with the confidence attached to the empirical
audit term,

\[
\sup_{g\in\mathcal R_{I,B}^{\mathrm{aud}}}
\operatorname{Ext}_I(g;T_I)
\le
\widehat U_{\mathrm{cert}}.
\]

Any budget-admissible computation whose terminal readout is measurable
in the audited algebra and whose model selection lies in the audited
search class satisfies the same bound. A computation that outputs a
readout outside the audited algebra is a new claimed observable
representative; it must supply provenance, budget, terminal readout, and
learning or selection certificate before it enters a later audit.

For the sampled audit of
\cref{def-sampled-computable-observable-audit}, the corresponding population
conclusion is

\[
\sup_{g\in\mathcal R_{I,B}^{\mathrm{aud}}}
\operatorname{Ext}_I(g;T_I)
\le \widehat U_{\mathrm{samp}}
\qquad\text{on }\mathcal E_{N,\delta},
\]

and therefore holds with probability at least \(1-\delta\).  No population
claim follows from the sampled matrices alone when the simultaneous
certificate is absent.

\end{lemma}

\begin{proof}\phantomsection\label{proof-computable-audit-soundness}

Because the exact list is \(\operatorname{supp}\mu_I\) and
\(w_i=\mu_I(x_i)\), weighted least squares computes the population
orthogonal projection of the centered target contrast onto the finite
feature span. Grouping by
quotient fibers computes conditional target densities and therefore the
quotient projection mass. Threshold margins and monotone filtrations are
finite scans over sorted score values. The effective dimension and
captured kernel mass are obtained from the eigendecomposition of the
finite Gram matrix. The residual-order ceiling is the scalar
\((1-2\eta)^{2q}\). The null score is computed from either an analytic
random-label variance or from the declared permutation sample. The
empirical penalty is supplied by the Rademacher or covering theorem for
the audited search class.

The projection ceiling bounds every measurable postprocessing of the
exposed algebra. The target-discrimination certificate restricts the
projection to actionable readouts. The budgeted extraction theorem adds
the empirical, effective-dimension, and residual-order ceilings. Taking
their minimum gives \(\widehat U_{\mathrm{cert}}\). A terminal readout
produced by any computation inside the audited algebra is one member of
\(\mathcal R_{I,B}^{\mathrm{aud}}\), so the same supremum bound applies.
A readout outside the audited algebra defines a different represented
observable class and therefore lies outside the stated audit.

In the sampled case, each event in the intersection
\(\mathcal E_{N,\delta}\) converts one empirical coordinate into a population
upper ceiling for the entire audited, selection-closed class.  On their
intersection, every applicable quantity
\(\widehat C_j+\varepsilon_j\) bounds the same extraction supremum, so their
minimum does as well.  The defining simultaneous certificate gives
\(\Pr(\mathcal E_{N,\delta})\ge1-\delta\).

\end{proof}

\Cref{fig-finite-audit-regimes} distinguishes the exact and sampled routes to
a population certificate.

\begin{figure}[tbp]
\centering
\begin{tikzpicture}[fw diagram]
  \node[fw box,minimum width=3.2cm] (data) at (0,2.8) {finite audit input};
  \node[fw source,minimum width=3.8cm] (full) at (-3.7,1.3)
    {full support\\\(S=\operatorname{supp}\mu\), \(w=\mu\)};
  \node[fw provenance,minimum width=3.6cm] (sample) at (3.7,1.3)
    {independent sample\\\(Z_{1:N}\sim\mu\)};
  \node[fw geom,minimum width=3.8cm] (exact) at (-3.7,-0.25)
    {exact population\\matrix quantities};
  \node[fw use,minimum width=3.8cm] (emp) at (3.7,-0.25)
    {empirical quantities\\+ simultaneous radii};
  \node[fw terminal,minimum width=4.4cm] (cert) at (0,-1.9)
    {population extraction certificate};
  \draw[fw arrow] (data) -- (full);
  \draw[fw arrow] (data) -- (sample);
  \draw[fw arrow] (full) -- (exact);
  \draw[fw arrow] (sample) -- (emp);
  \draw[fw arrow] (exact) -- (cert);
  \draw[fw arrow] (emp) -- node[right,font=\scriptsize]
    {on \(\mathcal E_{N,\delta}\)} (cert);
\end{tikzpicture}
\caption{The two finite-audit regimes. Full-support enumeration with exact
weights computes population projections directly. A sample produces population
bounds only after a simultaneous concentration certificate covering the
audited class and its output selection.}
\label{fig-finite-audit-regimes}
\end{figure}

\begin{example}[SAT defect as a statically target-relevant observable]\label{ex-usefulness-sat}

For SAT, the defect \(D_\Phi(x)\) counts violated clauses. The oriented
observable \(\psi=-D_\Phi\) has success sector
\(T_\Phi=\{x:D_\Phi(x)=0\}\) as its maximum level set. Its superlevel
sets are assignments with at most a specified number of violated
clauses. Its family-level target relevance requires separation of
successful and failed terminal states, recovery of the correct terminal
threshold, budget-admissible degree in the clause algebra, and family
uniformity. These properties do not define a state-to-state relation.

\end{example}

\section{Pointwise Observables and Geometric
Regularity}\label{sec-pointwise-geometric-regularity}
\label{pointwise-observables-and-geometric-regularity}

Pointwise observables provide forward evaluation without a represented
fiber decomposition. Geometric regularity is therefore imposed only for
uses based on local pointwise comparisons; invertibility alone does not
require it.

The pointwise class contains observables supplied by forward evaluation without represented fibers.
\begin{definition}[Pointwise observable]\label{def-pointwise-observable}

An \(\mathcal F_I\)-measurable observable \(\phi:X_I\to Y\) is pointwise
in the presentation when the interface supplies forward evaluation

\[
x\longmapsto\phi(x)
\]

but does not supply compact representations of the fibers
\(\phi^{-1}(y)\). The level sets exist as subsets of \(X_I\), but become
represented cells only through enumeration or additional structure.

\end{definition}

The comparability relation identifies static records connected by an admissible represented transfer.
\begin{definition}[Static comparability relation]\label{def-static-comparability}

A geometric presentation augments the observable data by a symmetric
measurable relation

\[
E_I\subseteq X_I\times X_I,
\qquad
E_I=E_I^{\mathsf T},
\qquad
E_I\in\mathcal F_I\otimes\mathcal F_I,
\qquad
(x,x)\notin E_I.
\]

The relation \(E_I\) specifies the comparable state pairs. It is static
presentation data rather than an algorithmic transition rule.

In the finite weighted case, let \(c_I:X_I\times X_I\to[0,\infty)\) be
a symmetric measurable edge conductance, supported on \(E_I\) and zero on
the diagonal.  If the primitive datum is instead a jump-rate table
\(q_I\), its conductance is \(c_I(x,y)=\mu_I(x)q_I(x,y)\) and symmetry
requires detailed balance.

\end{definition}

Local regularity of a pointwise observable is relative to the declared
comparison structure: the same scalar function may have different
regularity under different comparability relations. The following
functional measures this variation statically.

The regularity functional measures variation across the declared static comparison relation.
\begin{definition}[Static geometric regularity functional]\label{def-static-regularity-functional}

For a real pointwise observable \(\phi:X_I\to\mathbb R\) on a finite
weighted comparability presentation, define the static quadratic form

\[
\mathcal E_{E_I}(f,g)=
\frac12\sum_{x,y\in X_I}c_I(x,y)(f(x)-f(y))(g(x)-g(y)).
\]

The defect-regularity ratio of \(\phi\) is

\[
\rho_{E_I}(\phi)=
\frac{\mathcal E_{E_I}(\phi,\phi)}{\operatorname{Var}_{\mu_I}(\phi)}
\]

when the variance is nonzero, and \(+\infty\) otherwise. This quantity
measures how much the observable varies across declared-comparable pairs
per unit variance.

\end{definition}

This quadratic form has the algebraic form of a Dirichlet form
\citep{LevinPeres2017}, but here it serves only as a regularity
functional over a static relation and introduces no generator.

\begin{proposition}[Exact Fourier--Dirichlet identity on the binary cube]
\label{prop-hypercube-fourier-dirichlet-identity}

Let \(X=\mathbb F_2^n\) carry the uniform law, let
\(\chi_a(x)=(-1)^{\langle a,x\rangle}\), and take Hamming comparability with
conductances
\[
c(x,y)=2^{-n}\mathbf 1_{\{|x-y|_1=1\}}.
\]
For \(f=\sum_{a\in\mathbb F_2^n}\widehat f(a)\chi_a\),
\begin{align}
\label{eq-hypercube-fourier-dirichlet-diagonalization}
\mathcal E_E(\chi_a,\chi_b)
&=2|a|\,\mathbf1_{\{a=b\}},\\
\label{eq-hypercube-fourier-dirichlet-energy}
\mathcal E_E(f,f)
&=2\sum_{a\ne0}|a|\,|\widehat f(a)|^2.
\end{align}
Hence, when \(f\) is nonconstant,
\begin{equation}
\label{eq-hypercube-regularity-mean-fourier-weight}
\rho_E(f)
=2\,
\frac{\sum_{a\ne0}|a|\,|\widehat f(a)|^2}
     {\sum_{a\ne0}|\widehat f(a)|^2}.
\end{equation}
Thus the static regularity ratio is exactly twice the mean Hamming weight
under the normalized spectral-energy measure of \(f-\mathbb E f\).

Equation \eqref{eq-hypercube-regularity-mean-fourier-weight} is an identity
for the complete represented observable.  Target qualification remains the
separate comparison required by
\cref{def-pointwise-geometric-regularity,def-pre-hit-temporally-admissible-source}.
In particular, neither low mean Fourier weight nor a diffuse formal spectrum
alone supplies a target orientation, a coefficient locator, or a prospective
transition.

\end{proposition}

\begin{proof}
For a coordinate flip,
\[
\chi_a(x)-\chi_a(x+e_i)
=\chi_a(x)\bigl(1-(-1)^{a_i}\bigr).
\]
Substitution into \cref{def-static-regularity-functional} gives
\[
\mathcal E_E(\chi_a,\chi_b)
=\frac12\sum_{i=1}^n
\bigl(1-(-1)^{a_i}\bigr)
\bigl(1-(-1)^{b_i}\bigr)
\mathbb E_x[\chi_{a+b}(x)].
\]
Character orthogonality makes this zero for \(a\ne b\); for \(a=b\),
exactly \(|a|\) summands equal \(4/2\).  This proves
\cref{eq-hypercube-fourier-dirichlet-diagonalization}.  Bilinearity and
Parseval give \cref{eq-hypercube-fourier-dirichlet-energy}; division by
\(\operatorname{Var}(f)=\sum_{a\ne0}|\widehat f(a)|^2\) gives
\cref{eq-hypercube-regularity-mean-fourier-weight}.
\end{proof}

\begin{definition}[Geometric regularity of a pointwise observable]\label{def-pointwise-geometric-regularity}

Let \(\psi=s_I(\phi)\phi\) be the oriented observable. The static
off-target local-obstruction set is

\[
\operatorname{ObsMax}_{E_I}(\psi)=
\{x\in X_I\setminus T_I:
\psi(y)\le\psi(x)\text{ for every }y\text{ with }(x,y)\in E_I\}.
\]

The pointwise geometric regularity factor is

\[
\operatorname{Reg}^{\mathrm{Geom}}_{I}(\phi)=
\mathbf 1\{\operatorname{ObsMax}_{E_I}(\psi)=\varnothing\}
\cdot
\frac1{1+\rho_{E_I}(\psi)}.
\]

The observable is geometrically regular when this factor has a
transfer-class lower margin over the family.

\end{definition}

\begin{proposition}[Non-identifiability of fiber quotients from pointwise access]\label{prop-pointwise-no-fiber-quotient}

If \(\phi\) is pointwise, two evaluations decide the pairwise predicate
\(x\sim_\phi y\Longleftrightarrow\phi(x)=\phi(y)\).  Pointwise access alone
does not supply a represented quotient object: a uniform cell encoding,
construction of the represented fiber indexed by a value, and the declared
finite Boolean fiber operations require their own constructor record.  Such a
quotient is available precisely when the presentation supplies that record or
the budgeted grammar derives it.

\end{proposition}

\begin{proof}\phantomsection\label{proof-pointwise-no-fiber-quotient}

Pointwise access returns \(\phi(x)\), and hence decides equality of the
values at two supplied states.  The definition of a fiber-represented
observable requires more: a uniform representation of the cell associated
with a realized value, together with its construction and manipulation
costs.  These operations are not constructors in the pointwise interface.
They enter the represented algebra only through an additional declared or
derived record, at which point \cref{def-invertible-observable} applies.

\end{proof}

\begin{example}[SAT defect over Hamming comparability]\label{ex-pointwise-sat}

For SAT, take \(E_\Phi\) to relate assignments at Hamming distance one
\citep{Hamming1950}.
This is a static comparability relation on \(\{0,1\}^n\). The defect
score \(D_\Phi\) is pointwise if the presentation lets us evaluate the
number of violated clauses at a given assignment but does not give
compact descriptions of all assignments with the same defect. The ratio

\[
\rho_{E_\Phi}(D_\Phi)=
\frac{\frac12\sum_{x,y}c_\Phi(x,y)(D_\Phi(x)-D_\Phi(y))^2}
{\operatorname{Var}_{\mu_\Phi}(D_\Phi)}
\]

is a static measurement of defect variation across one-bit comparable
assignments, independent of any bit-flip transition rule.

\end{example}

\section{Fiber-Represented Observables, Lumpability, and Algebraic
Quotients}\label{sec-fiber-observables-quotients}
\label{invertible-observables-lumpability-and-algebraic-quotients}

A fiber-represented observable supplements pointwise evaluation with a
representation of its level sets, thereby inducing a static quotient of
the state space. Compatibility of this quotient with a search process is
treated separately in \PartII.

Fiber representation adds an affordable description of an observable's level sets to its pointwise evaluation rule.
\begin{definition}[Fiber-represented observable]\label{def-invertible-observable}

An observable \(\phi:X_I\to Y\) is fiber-represented in
the task presentation when the presentation supplies compact
descriptions of every nonempty fiber

\[
\phi^{-1}(y)=\{x\in X_I:\phi(x)=y\}.
\]

Compact representation requires a uniform value-indexed fiber constructor,
fiber membership, construction of the compact fiber description, and finite
Boolean operations on represented fibers to have constructor charge below the
declared budget ceiling.  It
does not require selecting or finding an element of a nonempty fiber.
Under the
polynomial budget, this is the standard polynomial bound on description
length and manipulation cost.

\end{definition}

\Cref{fig-pointwise-fiber-quotient-hierarchy} separates pointwise evaluation,
represented fibers, algebraic quotients, and the independent geometric
comparability interface.

\begin{figure}[tbp]
\centering
\begin{tikzpicture}[fw diagram]
  \node[fw source,minimum width=3.1cm] (point) at (-4.8,2.0)
    {pointwise evaluation\\\(x\mapsto\phi(x)\)};
  \node[fw box,minimum width=3.1cm] (equal) at (-1.4,2.0)
    {pairwise value test\\\(\phi(x)=\phi(y)\)};
  \node[fw provenance,minimum width=3.3cm] (fiber) at (2.2,2.0)
    {represented fiber\\constructor};
  \node[fw use,minimum width=3.0cm] (quot) at (5.6,2.0)
    {algebraic quotient\\\(q_\phi:X\to Q_\phi\)};
  \draw[fw arrow] (point) -- (equal);
  \draw[fw dashed arrow] (equal) -- (fiber);
  \draw[fw arrow] (fiber) -- (quot);
  \node[fw terminal,minimum width=3.4cm] (lump) at (5.6,0)
    {terminal-lumpability\\\(T=q^{-1}(T_Q)\)};
  \draw[fw arrow] (quot) -- (lump);
  \node[fw geom,minimum width=3.4cm] (comp) at (-1.4,0)
    {static comparability\\\(E_I,c_I\)};
  \node[fw geom,minimum width=3.4cm] (geom) at (2.2,0)
    {geometric regularity\\\(\mathcal E_{E_I}\)};
  \draw[fw arrow] (comp) -- (geom);
\end{tikzpicture}
\caption{Observable interfaces have distinct types. Pointwise evaluation
supplies value comparisons; a represented quotient additionally requires a
charged fiber constructor. Terminal lumpability is a property of that quotient,
whereas comparability and geometric regularity form an independent static
interface.}
\label{fig-pointwise-fiber-quotient-hierarchy}
\end{figure}

A pointwise observable supplies state values, whereas a fiber-represented
observable also represents the associated fibers. This additional data
names measurable cells without specifying transitions between them.

The induced algebraic quotient names the represented fibers without assigning dynamics to them.
\begin{definition}[Algebraic quotient induced by a fiber-represented observable]\label{def-algebraic-quotient}

Let \(\phi\) be fiber-represented. Define an equivalence relation

\[
x\sim_\phi y
\quad\Longleftrightarrow\quad
\phi(x)=\phi(y).
\]

The quotient set is

\[
Q_\phi=X_I/{\sim_\phi},
\qquad
q_\phi:X_I\to Q_\phi,
\qquad
q_\phi(x)=[x]_\phi.
\]

Because the fibers are represented, the quotient cells
\(q_\phi^{-1}(q)\) are represented observable objects.

\end{definition}

The terminology below is adapted from finite-state Markov-chain
lumpability \citep{KemenySnell1976}.

\begin{definition}[Static lumpability]\label{def-static-lumpability}

A represented quotient \(q:X_I\to Q\) is statically terminal-lumpable
when the terminal sectors are unions of quotient cells:

\[
T_I=q^{-1}(T_Q),
\qquad
F_I=q^{-1}(F_Q)
\]

for some disjoint subsets \(T_Q,F_Q\subseteq Q\). Equivalently, if
\(q(x)=q(y)\), then \(x\) and \(y\) have the same success/failure
terminal status whenever either is state-terminal.

\end{definition}

Static lumpability requires each quotient cell to have consistent
terminal status. It is a compatibility condition between the quotient
partition and terminal structure; dynamical lumpability is introduced
in \PartII.

\begin{proposition}[Quotient events are exactly unions of represented fibers]\label{prop-quotient-algebra-level-sets}

Let \(\phi:X_I\to Y\) be fiber-represented with quotient map \(q_\phi\). In the
finite case, assume the realized singletons of \(Y\) are measurable as in
\cref{def-task-space-signature}. A set \(A\subseteq X_I\) is measurable with respect to
\(\sigma(\phi)\) if and only if it is a union of quotient cells.

\end{proposition}

\begin{proof}\phantomsection\label{proof-quotient-algebra-level-sets}

The atoms of \(\sigma(\phi)\) are the nonempty fibers \(\phi^{-1}(y)\),
which are exactly the quotient cells. By the atom characterization for
finite observable algebras, the measurable sets in \(\sigma(\phi)\) are
precisely unions of those atoms.

\end{proof}

\begin{example}[SAT parity quotient]\label{ex-invertible-sat}

Suppose a SAT presentation declares a public parity row
\(\ell(x)=b\cdot x\pmod 2\). The observable
\(\ell:X_\Phi\to\mathbb F_2\) is fiber-represented when the affine fibers
\(\ell^{-1}(0)\) and \(\ell^{-1}(1)\) are represented by the public
linear equation. The quotient \(X_\Phi/{\sim_\ell}\) has two cells. It
is statically terminal-lumpable only if the satisfying assignments and
terminal failures are unions of those parity cells. If a parity cell
mixes satisfying and terminal-failing assignments, the parity quotient
is not a valid static success quotient.

\end{example}

\section{Observable Versus Hidden
Structure}\label{sec-observable-hidden-structure}
\label{observable-versus-hidden-structure}

A relation may hold on an ambient instance without being generated by
its presentation. The distinction is measurable: the observable algebra
separates exposed structure from information outside the primitive
presentation.

Exposure distinguishes relations generated by the presentation from relations available only in the ambient instance.
\begin{definition}[Exposed and hidden structure]\label{def-exposed-hidden}

A set \(A\subseteq X_I\) is exposed when \(A\in\mathcal F_I\). A
\(k\)-ary relation \(R\subseteq X_I^k\) is exposed when

\[
R\in\mathcal F_I^{\otimes k}.
\]

A function \(g:X_I\to Z\) is exposed when it is
\(\mathcal F_I\)-measurable. A structure is hidden relative to the
presentation when it may be true in an ambient description of the
instance but is not exposed in this sense.

\end{definition}

This definition concerns whether the declared observation interface
already generates the structure. Formally hidden structure may hold in
the ambient instance while remaining presentation-inaccessible. \PartII
classifies movement based on such structure as residual relative to the
public presentation unless a budget-admissible computation generates it
from exposed observables in a lifted state space.

The subsequent source hierarchy distinguishes exposed local geometry,
defined by the comparability relation and regularity functional, from
represented abstraction, defined by fibers and quotient cells. The
directed/drift channel \(\Caus\) orients these sources, the
lifted-state channel \(\Lift\) reorganizes exposed or generated
coordinates, and the boundary/killing channel \(\Bdry\) evaluates
terminal observables. Each operation therefore acts on geometry or
abstraction already represented in the observable algebra or generated
in an admissible lifted configuration.

\begin{proposition}[Finite membership test for exposed relations]\label{prop-product-atom-test}

Assume \(X_I\) is finite with atomic partition \(\mathcal A_I\). A
relation \(R\subseteq X_I^k\) is exposed if and only if its indicator is
constant on every product atom \(A_1\times\cdots\times A_k\) with
\(A_j\in\mathcal A_I\).

\end{proposition}

\begin{proof}\phantomsection\label{proof-product-atom-test}

The product \(\sigma\)-algebra \(\mathcal F_I^{\otimes k}\) is generated by
rectangles whose sides are observable events. In the finite case its
atoms are the products of atoms of \(\mathcal F_I\). By the atom
characterization, a subset of \(X_I^k\) is measurable in the product
algebra exactly when it is a union of these product atoms. This is
equivalent to constancy of its indicator on each product atom.

\end{proof}

\begin{example}[Hidden modulus]\label{ex-hidden-modulus}

Let \(X=\{0,1\}^n\) and suppose the ambient instance has a residue
relation \(x\equiv_M y\) for a modulus \(M\). If \(M\) and the residue
map are not generated by \(\operatorname{Obs}_I\), then the relation

\[
R_M=\{(x,y):x\equiv_M y\}
\]

is hidden precisely when \(R_M\notin\mathcal F_I\otimes\mathcal F_I\).
In a finite presentation this means that some product atom contains both
residue-equivalent and non-residue-equivalent pairs.

\end{example}

\begin{example}[Hidden planted assignment in SAT]\label{ex-hidden-sat}

A random planted SAT generator may choose a satisfying assignment
\(x^\star\) and then sample clauses consistent with it. If the final
task presentation exposes only the clauses, the singleton
\(\{x^\star\}\) is exposed only when it is a union of atoms of the
clause-readout algebra. If another assignment has the same public
readout but is not \(x^\star\), the planted label is hidden even though
it was used by the external generator that produced the instance.

\end{example}

\section{Scores and State-to-State
Relations}\label{sec-scores-state-relations}
\label{why-scores-do-not-move}

A score is a function on states, whereas a state-to-state relation is a
subset of \(X_I\times X_I\). These data have different mathematical
types.

Score coloring records the state partition induced by a scalar score.
\begin{definition}[Score coloring]\label{def-score-coloring}

A score observable is a measurable map
\(D_I:(X_I,\mathcal F_I)\to(\mathbb R,\mathcal B_{\mathbb R})\). It
induces the coloring

\[
\kappa_D:X_I\to\mathbb R,
\qquad
\kappa_D(x)=D_I(x),
\]

and the level-set equivalence relation
\(x\sim_D y\Longleftrightarrow D_I(x)=D_I(y)\). It does not by itself
specify any relation of admissible pairs between different score levels.

\end{definition}

A monotone transformation changes score values while preserving their
level-set partition. The score alone canonically induces only equality
of score values, not a relation between distinct fibers.

\begin{proposition}[Non-identifiability of transition relations from scores]\label{prop-score-no-state-relation}

Let \(X\) be a set with \(|X|\ge2\), and let \(D:X\to\mathbb R\). The
data of \(D\) determine the partition into fibers \(D^{-1}(c)\). They do
not determine a distinguished relation \(R\subseteq X\times X\) between
distinct fibers.

\end{proposition}

\begin{proof}\phantomsection\label{proof-score-no-state-relation}

Every construction from \(D\) is invariant under score-preserving
permutations of \(X\). It can test equality and order of score values,
but cannot distinguish two state pairs with the same ordered pair of
scores. For distinct fibers \(A\) and \(B\), both the empty and complete
relations on \(A\times B\) are compatible with \(D\). The score therefore
does not determine a distinguished state-to-state relation.

\end{proof}

\begin{example}[SAT defect partitions assignments]\label{ex-score-sat}

The SAT defect \(D_\Phi(x)\) partitions assignments by their numbers of
violated clauses; the zero-defect fiber is the success set. This
partition specifies neither state pairs nor a bit-flip, clause-repair,
or quotient-comparison rule. Such relations require additional
structure.

\end{example}

\section{Examples of Static Task Presentations}\label{sec-static-examples}
\label{worked-instruments}

The SAT and subset-sum examples illustrate state spaces, observable
readouts, atoms, level sets, and the distinction between exposed and
hidden relations.

\subsection{SAT as a Task Space}\label{subsec-static-sat-example}
\label{sat-as-a-task-space}

SAT provides the first example \citep{Cook1971}. A formula
\(\Phi=C_1\wedge\cdots\wedge C_m\)
on \(n\) variables has state space \(X_\Phi=\{0,1\}^n\). Clause
predicates and the defect score are observables. The success sector is
the zero-defect level set.

The SAT presentation instantiates the static signature with assignments, clause readouts, and the satisfying sector.
\begin{definition}[SAT task presentation]\label{def-sat-task-space}

For a CNF formula \(\Phi\), define

\[
X_\Phi=\{0,1\}^n,
\qquad
O_{\Phi,j}(x)=\mathbf 1\{C_j(x)=1\},
\qquad
D_\Phi(x)=\#\{j:C_j(x)=0\}.
\]

The observable algebra is generated by the declared clause readouts and
any additional public observables. The positive terminal sector is

\[
T_\Phi=D_\Phi^{-1}(\{0\}).
\]

\end{definition}

\subsection{Subset Sum and Exposed
Arithmetic}\label{subsec-static-subset-sum-example}
\label{subset-sum-and-exposed-arithmetic}

Subset sum provides a second instance of the same static construction
\citep{Karp1972}. The
state set can be the binary inclusion cube \(X=\{0,1\}^n\). The public
input consists of weights and a target,

\[
a_1,\ldots,a_n,
\qquad
t.
\]

The natural observable is the partial-sum defect

\[
D(x)=\left|\sum_{i=1}^n a_i x_i-t\right|.
\]

If the presentation also exposes residues modulo a public modulus, then
residue classes are observable quotient cells. If the modulus is not in
the observable algebra, the same residue relation may be true in an
ambient description but hidden from the task presentation.

\begin{example}[Observable and hidden residue structure]\label{ex-subset-sum-static}

Let \(r_M(x)=\sum_i a_ix_i\pmod M\). If \(M\) is public and \(r_M\) is
declared or derived from the observables, then the residue fibers
\(r_M^{-1}(b)\) are exposed. If \(M\) is absent from
\(\operatorname{Obs}_I\), then the relation \(r_M(x)=r_M(y)\) is hidden
unless it is still measurable in \(\mathcal F_I\otimes\mathcal F_I\).

\end{example}

The same pattern appears in SAT. Public XOR clauses expose affine
fibers. A planted assignment or secret parity used only by an instance
generator is hidden unless the final presentation declares it.

\section{Admissible Presentation
Equivalence}\label{sec-admissible-presentation-equivalence}
\label{admissible-presentation-equivalence}

Because the invariant is presentation-relative, changing the observable
interface may change its value. Presentations that encode the same
observable information at comparable representation cost must,
nevertheless, yield equivalent structural conclusions.

Admissible equivalence identifies presentations that preserve observable information and its charged representation cost.
\begin{definition}[Admissible presentation equivalence]\label{def-admissible-presentation-equivalence}

Two task-family presentations
\(\mathcal T=(X,\mathcal F,T,E,\mathcal R_{\le d})\) and
\(\mathcal T'=(X',\mathcal F',T',E',\mathcal R'_{\le d})\) are
admissibly equivalent with degree distortions \(a(d)\) and \(b(d)\) when
there are actual mutually inverse measurable bijections
\(\phi:X\to X'\) and \(\psi:X'\to X\).  Their pointwise pullbacks

\[
\Phi^*:L^\infty(X',\mathcal F')\to L^\infty(X,\mathcal F),
\qquad
\Psi^*:L^\infty(X,\mathcal F)\to L^\infty(X',\mathcal F')
\]

are inverse Boolean-algebra isomorphisms on the represented observable
events and observables, including their values on null states, and
satisfy the following conditions.  The induced maps on equivalence
classes in \(L^\infty\) are denoted by the same symbols.

\textbf{Terminals.} \(\Phi^*1_{T'}=1_T\) and \(\Psi^*1_T=1_{T'}\), and
the same holds for all named terminal sectors.

\textbf{Measures and Hilbert structures.} The pullbacks extend to
unitary maps between the declared state and edge Hilbert spaces. In
particular, the reference measure is transported exactly and the edge
inner products used for geometric, Hodge, quotient, lift, and boundary
projections are preserved. A merely bounded nonconstant
Radon--Nikodym density gives equivalent norms but need not preserve a
correlation, angle, zero projection, or orthogonal channel. Such a
change is not an admissible equivalence for the exact invariance theorem
unless separate cross-pairing and projection-comparison estimates are
declared and charged.

\textbf{Static structure.} Observable comparability relations,
represented fibers, quotient cells, and static terminal-lumpability are
transported by \(\Phi^*\) and \(\Psi^*\).

\textbf{Cost filtration.} Every represented observable, quotient,
comparability relation, or static constructor admissible for
\(\mathcal B'\) pulls back to one admissible for a class-equivalent
enlargement of \(\mathcal B\), and conversely. The distortion functions
preserve the transfer class: polynomial-dimensional bands map to
polynomial-dimensional bands only in the special case
\(\mathcal B=\mathcal B_{\mathrm{poly}}\).

\end{definition}

\begin{lemma}[Static invariance under admissible equivalence]\label{thm-static-presentation-invariance}

Under admissible presentation equivalence, exposed versus hidden
structure, terminal membership, represented quotient structure, static
terminal-lumpability, comparability geometry, and the five static
target-relevance conditions are preserved up to the stated budget/cost distortion
and with their Hilbert-space projection and correlation values
transported exactly.

\end{lemma}

\begin{proof}\phantomsection\label{proof-static-presentation-invariance}

A pointwise Boolean-algebra isomorphism sends atoms to atoms, measurable
sets to measurable sets, observable values to the corresponding values,
and nonmeasurable distinctions to nonmeasurable distinctions relative
to the transported presentation. In particular, pointwise maximizers
and essential maximizers are both transported exactly; no change on a
null state is discarded. Terminal
preservation gives the same success and failure sectors. Transport of
represented fibers sends quotient cells to quotient cells, so unions of
terminal cells and static lumpability are invariant. Transport of
comparability relations and the unitary edge transport send the static
Dirichlet quadratic form to the corresponding transported quadratic
form. Unitarity preserves normalized projections and cross-pairings.
Finally, the
cost-filtration condition sends each budget-admissible static observable or
quotient on one side to an admissible representative on the other after
the declared transfer-class distortion. Thus the static target-relevance
conditions are preserved under the declared budget distortion.

\end{proof}

A presentation containing an additional quotient, modulus, planted
label, or target pointer changes the observable algebra and is therefore
inequivalent. A recoding of the same observable algebra with comparable
constructor cost is equivalent, and the subsequent invariants agree up
to the prescribed cost distortion.

\section{Summary of the Static Task
Presentation}\label{sec-static-summary}
\label{summary-the-closed-static-object}

\PartI represents each problem instance as the static object

\[
\mathcal T_I=(X_I,\operatorname{Obs}_I,\operatorname{Dyn}_I,\operatorname{End}_I),
\]
where \(X_I\), \(\operatorname{Obs}_I\), and \(\operatorname{End}_I\)
carry the developed structure and \(\operatorname{Dyn}_I\) remains a
static relation slot. The observable algebra, its atoms and level sets,
and its pointwise and represented-fiber observables determine the
exposed static structure. Target projection mass bounds the target
signal obtainable from any generated subalgebra, while a
target-discrimination certificate assigns terminal meaning to its
fibers. Budgeted extraction accounting then bounds the actionable mass
after certification, learning, spectral, and residual-noise audits. For
finite presentations, the computable audit evaluates these quantities
from matrices, labels, quotient fibers, kernels, and declared audit
parameters. The intrinsic sources used in \PartII are exposed local
geometry and represented abstraction; the directed/drift, lifted-state,
and boundary/killing channels operate on these sources.

\paragraph{Dependency trail.}
The static interface is fixed by
\cref{def-observable-refinement,def-score-as-observable,def-static-family-uniformity,def-uniform-margin}.
Target relevance and its audits use
\cref{def-static-usefulness-clauses,def-static-usefulness-core,def-target-discrimination-certificate,def-budgeted-extraction-ledger,def-extraction-audit-ceilings,def-computable-observable-audit}.
The geometric and quotient vocabulary is
\cref{def-pointwise-observable,def-static-comparability,def-static-regularity-functional,def-pointwise-geometric-regularity,def-invertible-observable,def-algebraic-quotient,def-static-lumpability,def-exposed-hidden,def-score-coloring,def-sat-task-space}.
The signature and observable-algebra consequences are
\cref{prop-complete-states-testable,prop-atoms-characterize-events,prop-measurability-is-not-affordability,prop-family-uniformity-not-algorithm-certificate,prop-one-instance-not-family}.
Static target relevance is developed in
\cref{thm-static-clauses-necessary,thm-generated-algebra-projection-ceiling,cor-expanded-derivation-projection-ceiling,thm-raw-combinations-require-discrimination,cor-vanishing-ceiling-vanishing-extraction,thm-computable-audit-soundness}.
The pointwise and quotient distinctions are
\cref{prop-pointwise-no-fiber-quotient,prop-quotient-algebra-level-sets,prop-product-atom-test,prop-score-no-state-relation}.
Presentation invariance and the resulting static closure are
\cref{thm-static-presentation-invariance,thm-static-closure}.
The running static examples are collected in
\cref{ex-intro-sat-static,ex-signature-sat,ex-observable-algebra-sat,ex-family-sat,ex-affordable-xor-synergy,ex-usefulness-sat,ex-pointwise-sat,ex-invertible-sat,ex-hidden-sat,ex-score-sat,ex-subset-sum-static,ex-summary-sat}.

\begin{lemma}[Static closure of \PartI]\label{thm-static-closure}

Every object developed in \PartI is determined by the state set,
observable algebra, terminal sectors, family presentation, represented
fibers, target-discrimination certificates, budgeted extraction accounts,
and static comparability relations. The target signal available to any
generated subalgebra is bounded by the Hilbert projection of the target
contrast onto that subalgebra, while terminal relevance requires a
represented or generated target-discrimination certificate. Certified
extraction is further bounded by the applicable empirical,
effective-dimension, and residual-order audits. The only intrinsic
reusable sources for later structured movement are the local geometric
source from exposed comparability and the abstraction source from
represented fibers. None requires a run, transition kernel, generator,
search process, or algorithm.

\end{lemma}

\begin{proof}\phantomsection\label{proof-static-closure}

The task signature is defined from sets, maps, measurable structures,
and terminal predicates. Atoms, level sets, refinements, exposed
relations, hidden relations, quotient cells, static lumpability, score
partitions, target projection masses, target-discrimination certificates,
and budgeted extraction accounts are all set-theoretic or measurable
constructions from those data. The comparability relation \(E_I\) is
declared as a symmetric measurable relation, and the regularity
functional is a quadratic form on functions over that relation. No
definition quantifies over runs or assigns probabilities to
state-to-state evolution. Therefore the development is closed at the
static level.

\end{proof}

The forward correspondences are now named and ready for \PartII:

\begin{longtable}[]{@{}
  >{\raggedright\arraybackslash}p{(\columnwidth - 2\tabcolsep) * \real{0.5000}}
  >{\raggedright\arraybackslash}p{(\columnwidth - 2\tabcolsep) * \real{0.5000}}@{}}
\toprule\noalign{}
\begin{minipage}[b]{\linewidth}\raggedright
Static object in \PartI
\end{minipage} & \begin{minipage}[b]{\linewidth}\raggedright
Later role in \PartII
\end{minipage} \\
\midrule\noalign{}
\endhead
\bottomrule\noalign{}
\endlastfoot
declared slot \(\operatorname{Dyn}_I\) & relation slot for search and
state-to-state dynamics \\
comparability relation \(E_I\) & intrinsic local geometry to be
activated dynamically as \(\Geom\) \\
regularity functional \(\mathcal E_{E_I}\) & quantitative form of the
intrinsic local geometry \\
fiber-represented observable \(\phi\) & intrinsic represented abstraction to be
activated as \(\Abs\) \\
static terminal-lumpability & compatibility condition for the later
quotient/lumping channel \\
terminal sectors \(\operatorname{End}_I\) & boundary/value readout through
\(\Bdry\) \\
score observable \(D_I\) & state partition without a state-to-state
transition relation \\
target projection mass of a generated algebra & ceiling on target signal
available to any combination of its observables \\
target-discrimination certificate & equips observable fibers with a
terminal-relevant quotient readout; without it, the tuple remains
uncalibrated \\
budgeted extraction accounting & quantitative upper bound on certified
extraction after projection, certification, learning, spectral, and
residual-noise audits \\
computable observable audit & finite matrix implementation of the accounting
for a declared exposed algebra and audit class \\
\end{longtable}

\begin{example}[SAT static closure]\label{ex-summary-sat}

For SAT, \PartI gives the assignment cube, clause observables, defect
score, satisfying sector, observable atoms, level sets, possible exposed
parity quotients, possible hidden planted labels, and the static Hamming
comparability relation when declared. The selection of subsequent
assignments belongs to the dynamics of \PartII.

\end{example}

\clearpage
\part{Budget-Admissible Dynamics and Structural Decomposition}\label{part:dynamics}\setcounter{section}{-1}

\section{Introduction: Budget-Admissible
Dynamics}\label{sec-dynamics-introduction}
\label{introduction-from-the-presented-task-to-budget-admissible-movement}

\paragraph{Part contract.}
\emph{Inputs:} the task presentation, observable algebra, terminal sectors,
comparability data, and represented quotients from \cref{part:task-spaces}.
\emph{Output:} the exhaustive structural-channel decomposition in
\cref{thm-channel-exhaustiveness}, the explicit channel normal forms in
\cref{thm-exhaustive-geom-normal-form,thm-exhaustive-lift-normal-form,thm-exhaustive-bdry-normal-form,thm-exhaustive-residual-normal-form},
and their stability under represented composition.
\emph{First pass:} the post-exhaustion comparison band and worked examples may
be skimmed.

\PartII equips the static task presentation with dynamics and develops a
decomposition of every budget-admissible mass-transport operator.

\PartI defined the presented task

\[
\mathcal T_I=(X_I,\operatorname{Obs}_I,\operatorname{Dyn}_I,\operatorname{End}_I),
\]

while leaving \(\operatorname{Dyn}_I\) without an assigned process,
kernel, generator, traversal rule, or algorithm. The state space,
observable algebra, terminal sectors, static comparability relation, and
represented quotients now provide the data for the dynamic construction.

The proposed algebraic decomposition contains five structural channel
slots and one residual component; at a fixed budget, only charged
represented projections are exposed:

\[
\Geom,\qquad
\Caus,\qquad
\Abs,\qquad
\Lift,\qquad
\Bdry,\qquad
\Res,
\]

Every budget-admissible represented computation without oracle access
induces such an operator on its full observable transcript carrier. The
results below establish exhaustiveness of the
decomposition.

The geometric/reversible channel \(\Geom\) and the
quotient/lumping channel \(\Abs\) represent the intrinsic
sources: exposed local comparability and represented abstraction. The
directed/drift channel \(\Caus\) orients these sources, the
lifted-state channel \(\Lift\) reorganizes their coordinates,
and the boundary/killing channel \(\Bdry\) evaluates terminal
observables. The orthogonal remainder is denoted by
\(\Res\).

The directed channel is relative to a declared authority envelope.  A
problem may permit an affordable control to select the skew part of its
generator, may impose that part as dynamic data, or may allow a hybrid
of imposed and controlled motion.  Subsequent projections and lifts
transport the selected operator and its current; they do not replace
them by a more favorable target-aligned flow.

For a declared budget structure
\(\mathcal B=(\mathcal R,\operatorname{cost},b,\mathfrak C)\),
admissibility is filtered by the capacity functional \(b(n)\).
Instance-dependent information is restricted to the budgeted saturated
observable closure generated by the presented observable algebra,
public parameters, admissible lifted constructors, and finite-horizon
computations whose transcripts have charge below \(b(n)\). Relations,
quotients, potentials, lift interfaces, boundary values, transition
data, trajectory statistics, and finite-horizon target-contact
observables generated in this way enter the closure at their saturated
computation cost. Objects outside the closure are associated with an
enriched presentation, oracle or advice data, over-budget recovery, or
the residual component. The polynomial model
is the special case \(\mathcal B_{\mathrm{poly}}\) with
\(b_{\mathrm{poly}}(n)=n^{O(1)}\).

The analysis proceeds through five closure results. The positive maximum
principle yields nonnegative transition and killing rates. A fixed local
relation then separates local symmetric, local directed, nonlocal, and
boundary terms. Hodge and quotient projections resolve these terms into
the five structural channels and the residual component. Any
family-uniform budget-admissible reduction, quotient, or target alignment
is reclassified from the post-exhaustion residual band into the
corresponding structural channel. \PartIII treats the saturated
representation cost of these components. Finally, finite composition
and valid lifting preserve the classification.

The principal decomposition results are stated for finite state spaces,
including the SAT and subset-sum examples and the finite transcript
carriers of resource-bounded computations. \PartIII gives a conditional
sectorial-form and carrier interface for continuum realizations;
specific applications still require their own domain, closure,
capacity, and regularity hypotheses. Arrival times, resolvents, saturated cost
thresholds, and quantitative bounds are deferred to \PartIII.

\section{Budget-Admissible Movement and the Search
Operator}\label{sec-budget-admissible-movement}
\label{budget-admissible-movement-and-the-search-operator}

A frontier is a nonnegative measure on the state space of \PartI. In the
finite setting it is a vector in \(\mathbb R^{X_I}_{\ge0}\) representing
the mass currently assigned to each state before channel decomposition.

The frontier formalizes the current mass distribution on which a search operator acts.
\begin{definition}[Frontier and search operator]\label{def-frontier-search-operator}

Fix an instance \(I\) with finite state space \(X_I\) and observable
algebra \(\mathcal F_I\). Let \(\mathcal M(X_I)\) be the real vector
space of signed measures on \(X_I\), and let \(\mathcal M_+(X_I)\) be
its positive cone. A frontier is an element \(\mu\in\mathcal M_+(X_I)\).

A discrete search operator is a linear map
\(P:\mathcal M(X_I)\to\mathcal M(X_I)\) such that
\(P\mathcal M_+(X_I)\subseteq\mathcal M_+(X_I)\). With row-vector
frontiers this means

\[
(\mu P)(y)=\sum_{x\in X_I}\mu(x)P_{xy},
\qquad P_{xy}\ge0.
\]

The dual operator on observable functions \(f\in\mathbb R^{X_I}\) is

\[
(Kf)(x)=\sum_{y\in X_I}P_{xy}f(y),
\]

and it is positive in the sense that

\[
f\ge 0 \quad\Longrightarrow\quad Kf\ge 0.
\]

A continuous-time search operator is a linear operator
\(L:\mathbb R^{X_I}\to\mathbb R^{X_I}\) whose semigroup is positive
and sub-Markov:

\[
e^{tL}f\ge0\quad(f\ge0),
\qquad
e^{tL}\mathbf1\le\mathbf1
\quad(t\ge0).
\]

It is Markov when equality holds in the second relation.  On a finite
state space this condition is equivalent to the positive maximum
principle below. Evolution is written on observables as
\(\partial_t u_t=Lu_t\), or dually on frontiers as
\(\partial_t\mu_t=L^*\mu_t\).

\end{definition}

Budget admissibility is inherited from the closure defined in \PartI.
A dynamic edge, kernel, or generator may use only distinctions exposed
at the declared budget. The available finite data are generated by the
observable algebra, its product algebra on edges, and the coordinates of
an admissible lifted configuration; after lifting, the same condition is
imposed with respect to the lifted observable algebra.

When a budget-admissible computation generates additional data from this
closure, including through finite-horizon iteration and postprocessing,
those data enter the lifted observable coordinates at their accumulated
cost and the decomposition is applied on the lifted space. The
polynomial case sets \(\mathcal B=\mathcal B_{\mathrm{poly}}\). Data
outside the closure and any admissible budget-generated lift are
presentation-inaccessible, oracle-supplied, or over budget. Movement
without exposed structural provenance is assigned to the residual
component.

Budget-admissible movement restricts transition rules to the declared saturated closure and its charged lifts.
\begin{definition}[Budget-admissible movement]\label{def-admissible-operator}

Let \(\mathcal A_I\) be the atomic partition of the finite observable
algebra \(\mathcal F_I\). A kernel \(P=(P_{xy})\) is budget-admissible
at the declared budget when it is Markov or sub-Markov, when its
positive support is contained in an exposed relation represented within
the budgeted saturated closure,

\[
\operatorname{supp}_{\mathrm{off}}P
:=\{(x,y):x\ne y,\ P_{xy}>0\}
\subseteq R_I,
\qquad
R_I\in\mathcal F_I\otimes\mathcal F_I.
\]

Diagonal holding probabilities are unrestricted by this support condition;
equivalently, \(\operatorname{supp}P\subseteq R_I\cup\Delta_{X_I}\).
The rate table must also be product-measurable and represented at charged
cost: if \(x,x'\) lie in the same atom of \(\mathcal F_I\) and \(y,y'\)
lie in the same atom, then \(P_{xy}=P_{x'y'}\) whenever both pairs are
in the same product atom. For a generator with rates \(q(x,y)\), the
same condition is imposed on \(q\). Equivalently,
\(q:X_I\times X_I\to[0,\infty)\) is
\(\mathcal F_I\otimes\mathcal F_I\)-measurable and its representation
lies inside the active budgeted closure.

\end{definition}

This definition extends the exposed-versus-hidden distinction of \PartI
to dynamics. A hidden modulus, planted assignment, secret parity, or
target pointer may exist in the ambient instance, but it supports a
budget-admissible transition only when generated from exposed data in an
admissible lifted configuration. Otherwise its use is oracle-supplied,
presentation-inaccessible, over budget, or residual relative to the
public presentation.

\section{The Generator Form and the Positive Maximum
Principle}\label{sec-generator-pmp}
\label{the-generator-form-and-the-positive-maximum-principle}

A budget-admissible infinitesimal frontier update preserves
nonnegativity. Dually, the positive maximum principle requires the
generator to be nonpositive at a nonnegative maximum. On finite state
spaces this condition determines the matrix form
\citep{Courrege1965,Norris1997}.

The positive maximum principle fixes the finite-generator condition that preserves nonnegative frontier evolution.
\begin{definition}[Positive maximum principle]\label{def-pmp}

A linear operator \(L\) on observable functions satisfies the positive
maximum principle when

\[
f(x_0)=\max_{x\in X_I} f(x)\ge 0
\quad\Longrightarrow\quad
(Lf)(x_0)\le 0.
\]

This is the local infinitesimal form of positive sub-Markov preservation
for the frontier semigroup.

\end{definition}

\begin{lemma}[Finite generator form]\label{thm-finite-generator-form}

Let \(X\) be finite and let \(L\) be a linear operator on functions
\(f:X\to\mathbb R\). If \(L\) satisfies the positive maximum principle,
then there are nonnegative rates \(q(x,y)\ge0\) for \(x\ne y\) and a
nonnegative killing rate \(\kappa(x)\ge0\) such that

\[
(Lf)(x)
=
\sum_{y\ne x}q(x,y)\bigl(f(y)-f(x)\bigr)
-\kappa(x)f(x).
\]

Conversely, every operator of this form satisfies the positive maximum
principle.

\end{lemma}

\begin{proof}\phantomsection\label{proof-finite-generator-form}

Write \(L_{xy}=(L\mathbf 1_y)(x)\), so \((Lf)(x)=\sum_y L_{xy}f(y)\).
Fix \(x\ne y\). Let \(f(y)=-1\) and \(f(z)=0\) for \(z\ne y\). Then
\(f(x)=0=\max_z f(z)\), so the positive maximum principle at \(x\) gives
\((Lf)(x)=-L_{xy}\le0\). Hence \(L_{xy}\ge0\). Put \(q(x,y)=L_{xy}\) for
\(x\ne y\).

Apply the principle to the constant function \(1\). Since every point is
a nonnegative maximum, \((L1)(x)\le0\). Define
\(\kappa(x)=-(L1)(x)\ge0\). Then

\[
L_{xx}
=-\sum_{y\ne x}q(x,y)-\kappa(x),
\]

because \((L1)(x)=L_{xx}+\sum_{y\ne x}L_{xy}\). This gives the displayed
representation. Conversely, if \(f\) has a nonnegative maximum at \(x\),
then \(f(y)-f(x)\le0\) for every \(y\), and \(-\kappa(x)f(x)\le0\). Thus
\((Lf)(x)\le0\).

\end{proof}

Thus every positive finite-state generator consists of off-diagonal
transition rates and diagonal loss. With a represented terminal
observable space, diagonal loss is the zero-boundary-value specialization
of boundary readout. These terms exhaust the matrix entries.

\section{The Fixed-Space
Decomposition}\label{sec-fixed-space-decomposition}
\label{the-fixed-space-decomposition}

The finite generator theorem yields a rate table, while the static
comparability relation \(E_I\) from \PartI identifies local state pairs.
Relative to this declared relation, every rate has a unique local or
nonlocal location.

The positive weight \(m\) is a reference measure used to express local
rates as conductance and current. Any positive weight provides algebraic
coordinates for the same finite rate table. Quantitative statements,
however, use only the admissible reference measures of
\cref{def-reference-measure-legality}. A
task-dependent weight supplied by the presentation or a public lifted
initialization is charged at that level; otherwise it is
presentation-inaccessible, oracle-supplied, or over budget. Changing an
admissible weight changes the conductance-current coordinates but not the
underlying transport operator.

Reference-measure legality determines which conductance and current coordinates may enter quantitative certificates.
\begin{definition}[Reference measure legality]\label{def-reference-measure-legality}

A base reference measure is admissible when it is part of the task
presentation or is transported from an admissibly equivalent
presentation by its unitary measure transport. A lifted reference
measure on \(\Omega\) is admissible when it is fixed by the lifted
presentation before the measured operator is used. Concretely, it is
one of the following: a public product gauge on finitely encoded lifted
coordinates, including clocked transcript coordinates; the public initialization lift
\(\Lambda\mu\); or a measure absolutely continuous with respect to one
of these public gauges whose density is
\(\mathcal G_{\mathrm{obs}}\)-measurable, generated from the public
presentation and lifted initialization data, and bounded above and below
by transfer-class factors on the family. For
\(\mathcal B_{\mathrm{poly}}\), these are polynomial factors.

A measure used only for a static integral may have smaller support. A
measure used to realize \(L\) on \(L^2\), define an orthogonal Hodge or
channel projection, or compress to a quotient must have full support on
the represented carrier. Equivalently, one may first restrict the
carrier to its support, provided that support is invariant under the
operator and all quotient and terminal maps. Without full support or
this invariant-support reduction, a transition into a null state need
not define an operator on \(L^2\) equivalence classes.

A public product gauge may be used for the exact clocked-transcript and
first-hit accounting of an arbitrary finite computation.  It supports a
structural \(\Lift\) comparison only when the separate
projection and target-fraction conditions of
\Cref{def-valid-lift} hold.  Thus the configuration gauge
does not create a target-mass transfer identity.

Trajectory-defined laws, including occupation laws
\(\sum_{t\le p(n)}\mu_0P^t\), killed or hitting-conditioned occupation
laws, and measures obtained from an invariant, Green, or resolvent
problem for the measured operator, are admissible only when they also
possess an independent representation satisfying the preceding
conditions. Their dynamic definition alone does not provide a reference
gauge for the presentation.

\end{definition}

\Cref{fig-reference-measure-legality} summarizes the admissible routes to a
reference law and the additional representation gate for trajectory-defined
laws.

\begin{figure}[tbp]
\centering
\begin{tikzpicture}[fw diagram]
  \node[fw box,minimum width=3.4cm] (candidate) at (0,3.0)
    {candidate reference law};
  \node[fw source,minimum width=3.3cm] (presented) at (-4.7,1.55)
    {presented or unitarily\\transported base law};
  \node[fw source,minimum width=3.3cm] (lifted) at (0,1.55)
    {public lifted gauge\\or initialization lift};
  \node[fw residual,minimum width=3.6cm] (trajectory) at (4.7,1.55)
    {occupation, invariant,\\hitting, Green, resolvent};
  \node[fw use,minimum width=3.7cm] (density) at (2.4,0)
    {independent represented density\\with controlled bounds};
  \node[fw terminal,minimum width=3.8cm] (admit) at (-1.7,-1.55)
    {admissible full-support\\reference law};
  \node[fw residual,minimum width=3.5cm] (reject) at (4.1,-1.55)
    {inadmissible as an\\uncharged gauge};
  \draw[fw arrow] (candidate) -- (presented);
  \draw[fw arrow] (candidate) -- (lifted);
  \draw[fw arrow] (candidate) -- (trajectory);
  \draw[fw arrow] (presented) -- (admit);
  \draw[fw arrow] (lifted) -- (admit);
  \draw[fw dashed arrow] (trajectory) -- (density);
  \draw[fw arrow] (density) -- (admit);
  \draw[fw arrow] (trajectory) -- (reject);
\end{tikzpicture}
\caption{Reference-measure legality. Presented and public lifted gauges are
admissible subject to the stated support conditions. A trajectory-defined law
enters only through an independent charged representation with controlled
density; its operator-dependent definition alone is not a reference gauge.}
\label{fig-reference-measure-legality}
\end{figure}

\begin{proposition}[Exclusion of operator-dependent trajectory measures]\label{prop-trajectory-measures-not-gauges}

Let \(P\) or \(L\) be the induced operator whose channels are being
decomposed. A reference law constructed from the trajectory, occupation,
hitting-conditioned, invariant, Green, or resolvent behavior of that
same operator is inadmissible as an uncharged lifted reference measure in the
channel decomposition or in the \PartIII quantitative projections.

\end{proposition}

\begin{proof}\phantomsection\label{proof-trajectory-measures-not-gauges}

Such a law is operator-dependent because it is defined from \(P\),
\(L\), or their killed or conditioned dynamics rather than from the
presented observable algebra, its budgeted closure, and the public
initialization lift. It also differs from a gauge change: a gauge
modifies \(L^2\) coordinates by a controlled observable density, whereas
an occupation law may concentrate on a trajectory or reachable band and
alter the Hilbert space and channel projections. Finally, it is not an
initialization lift, since in general
\(\pi_*\sum_{t\le p(n)}\mu_0P^t\) is the time-occupation law of the
projected run rather than the base frontier \(\mu\) required by
\(\pi_*\Lambda\mu=\mu\). The law therefore supplies dynamical information
rather than reference coordinates. An independent public representation
with controlled density restores admissibility at the cost of that
representation.

\end{proof}

Consequently, a clock, program-counter, or depth coordinate aligned with
the target under an admissible public lifted reference law belongs to
the directed/drift or geometric/reversible channel. Alignment arising
only under the solver's trajectory or occupation law is supplied by an
inadmissible reference measure.

The fixed-space split decomposes one represented rate table into symmetric, skew, killing, and residual terms.
\begin{definition}[Fixed-space rate split]\label{def-fixed-space-split}

Let \(L\) have rates \(q(x,y)\) and killing \(\kappa(x)\). Fix a
symmetric local relation \(E_I\subseteq X_I\times X_I\setminus\Delta\)
and a positive reference weight \(m(x)>0\). Put
\(q_E(x,y)=q(x,y)\mathbf 1_{E_I}(x,y)\) and write the local oriented
edge mass

\[
r_E(x,y)=m(x)q_E(x,y).
\]

For \(x\ne y\), split the local edge mass into symmetric conductance and
antisymmetric current:

\[
c_E(x,y)=\frac12\bigl(r_E(x,y)+r_E(y,x)\bigr),
\qquad
j_E(x,y)=\frac12\bigl(r_E(x,y)-r_E(y,x)\bigr).
\]

The four fixed-space operators, in zero-boundary reduction, are

\[
\begin{aligned}
(L^{\mathsf{Sym}}f)(x)
&=\frac1{m(x)}\sum_{y:(x,y)\in E_I}c_E(x,y)(f(y)-f(x)),\\
(L^{\mathsf{Anti}}f)(x)
&=\frac1{m(x)}\sum_{y:(x,y)\in E_I}j_E(x,y)(f(y)-f(x)),\\
(L^{\mathsf{Jump}}f)(x)
&=\sum_{y\ne x:(x,y)\notin E_I}q(x,y)(f(y)-f(x)),\\
(L^{\Bdry}f)(x)
&=-\kappa(x)f(x).
\end{aligned}
\]

Here \(c_E(x,y)=c_E(y,x)\ge0\) and \(j_E(x,y)=-j_E(y,x)\). The
antisymmetric term is a signed current operator; it is not itself
required to be a positive generator. The displayed boundary term is the
special case in which removed mass has boundary value zero.

\end{definition}

The boundary-value extension records the destination and value of mass removed
from the interior carrier.
\begin{definition}[Boundary-value channel extension]
\label{def-boundary-value-channel}

For terminal readout, let \(\partial X_I\) be a finite boundary/value
space, let \(b(x,\zeta)\ge0\) be the rate of transfer from \(x\in X_I\)
to \(\zeta\in\partial X_I\), and put
\(\kappa(x)=\sum_{\zeta}b(x,\zeta)\). For an interior observable
\(f:X_I\to\mathbb R\) and a boundary readout
\(v:\partial X_I\to\mathbb R\), the same boundary channel acts on the
extended observable pair by

\[
(L^{\Bdry}(f,v))(x)
=
\sum_{\zeta\in\partial X_I} b(x,\zeta)\bigl(v(\zeta)-f(x)\bigr)
=
\kappa(x)\bigl(Bv(x)-f(x)\bigr),
\]

where \(Bv(x)=\kappa(x)^{-1}\sum_{\zeta}b(x,\zeta)v(\zeta)\) when
\(\kappa(x)>0\). Setting \(v\equiv0\) gives the zero-value killing term
\(-\kappa(x)f(x)\). Thus killing/removing support is the zero-value
specialization of boundary value processing.

A boundary value is budget-admissible as \(\Bdry\) when the
terminal observable and value readout are declared by the presentation
or generated from budget-admissible observable structure. A committor or
hitting-time table of an induced kernel becomes a boundary readout when
it admits an observable representation, including one produced by a
budget-admissible finite-horizon computation. An uncomputed resolvent or
infinite-horizon hitting process does not constitute an admissible
boundary constructor.

\end{definition}

\begin{lemma}[Fixed-space split]\label{thm-fixed-space-split}

For a finite budget-admissible generator, the exposed local relation
\(E_I\), and a positive reference weight \(m\) used only to coordinatize
local edge mass, the decomposition

\[
L=L^{\mathsf{Sym}}+L^{\mathsf{Anti}}+L^{\mathsf{Jump}}+L^{\Bdry}
\]

is canonical relative to these coordinates. The four terms are the
local symmetric, local directed, nonlocal, and boundary/killing
components. Dependence on \(m\) is confined to the conductance-current
coordinates.

\end{lemma}

\begin{proof}\phantomsection\label{proof-fixed-space-split}

The generator form determines all off-diagonal rates and diagonal loss.
The relation \(E_I\) partitions off-diagonal pairs into local and
nonlocal pairs. On local pairs, the ordered edge mass \(r_E(x,y)\) has
the unique decomposition \(r_E=c_E+j_E\) into its symmetric and
antisymmetric parts. Substituting
\(q_E(x,y)=m(x)^{-1}(c_E(x,y)+j_E(x,y))\) into the local part of \(L\)
gives \(L^{\mathsf{Sym}}+L^{\mathsf{Anti}}\). The nonlocal rates give
\(L^{\mathsf{Jump}}\). The diagonal loss, equivalently transfer to an
observable boundary value space, gives \(L^{\Bdry}\). The
supports are disjoint, and the symmetric/antisymmetric split is unique
in the chosen edge-mass coordinates. Since \(m\) only defines those
coordinates, it leaves the underlying movement channels unchanged.

\end{proof}

\begin{proposition}[Exhaustiveness of the local rate decomposition]\label{prop-no-third-local-primitive}

Let \(A\) be any linear operator satisfying the positive maximum
principle whose off-diagonal support is contained in the observable
local relation \(E_I\). Then \(A\) has only local pair rates and
diagonal loss, equivalently zero-value boundary readout. After the
symmetric/current split above, no third local term remains.

\end{proposition}

\begin{proof}\phantomsection\label{proof-no-third-local-primitive}

Apply the finite generator theorem to \(A\). Its off-diagonal
coefficients are nonnegative rates \(a(x,y)\ge0\), and the support
hypothesis gives \(a(x,y)=0\) when \((x,y)\notin E_I\). Its diagonal
part is exactly \(-\sum_{y\ne x}a(x,y)-\kappa_A(x)\) with
\(\kappa_A(x)\ge0\). Thus every local positive generator is already a
sum of pair-difference terms and boundary loss. If terminal values are
carried, that loss is the boundary value term \(\kappa_A(Bv-f)\); if the
terminal value is zero, it is killing. Splitting the local edge mass
into \(c_E+j_E\) accounts for all local pair terms, so a further local
primitive would have zero coefficients in every matrix entry.

\end{proof}

This algebraic split locates the rates but does not yet identify their
structural provenance.

\section{From the Algebraic Split to Structural
Channels}\label{sec-local-structural-channels}
\label{the-channel-bridge-from-syntactic-split-to-structural-channels}

The fixed-space decomposition separates the rate table algebraically.
The structural decomposition further classifies each component by its
observable provenance.

The relevant subspaces are generated by the observable algebra, its
budgeted saturated closure, and admissible coordinates in a lifted
configuration. A potential, quotient, or lifted interface outside this
closure can enter the operator only through an admissible lifted
derivation. Otherwise it is presentation-inaccessible, oracle-supplied,
over budget, or part of an enriched presentation.  Its operator term is
retained in the budget-filtered residual ledger, but it has no affordable
target-aligned source status for the original presentation.

\begin{longtable}[]{@{}
  >{\raggedright\arraybackslash}p{(\columnwidth - 4\tabcolsep) * \real{0.3333}}
  >{\raggedright\arraybackslash}p{(\columnwidth - 4\tabcolsep) * \real{0.3333}}
  >{\raggedright\arraybackslash}p{(\columnwidth - 4\tabcolsep) * \real{0.3333}}@{}}
\toprule\noalign{}
\begin{minipage}[b]{\linewidth}\raggedright
Syntactic part
\end{minipage} & \begin{minipage}[b]{\linewidth}\raggedright
Location
\end{minipage} & \begin{minipage}[b]{\linewidth}\raggedright
Structural channel
\end{minipage} \\
\midrule\noalign{}
\endhead
\bottomrule\noalign{}
\endlastfoot
\(\mathsf{Sym}\) & paired local rates inside \(E_I\) &
\(\Geom\) (geometric/reversible) \\
local and nonlocal skew flow & oriented imbalance &
\(\Caus\) (directed/drift) \(+\Res\) \\
nonlocal jump & outside \(E_I\) & \(\Abs\) (quotient/lumping)
\(+\Lift\) (lifted-state) when
supplied by observable quotient or valid lifted-interface data;
otherwise \(\Res\) \\
\(\Bdry\) & boundary readout, value processing, zero-value
killing & \(\Bdry\) (boundary/killing) \\
\end{longtable}

Provenance determines the nonlocal classification. A directed component
belongs to \(\Caus\) when it is the gradient of a represented
potential constructed from local geometry, quotient coordinates, or
lifted-state coordinates. A nonlocal transition has algebraic quotient or
lift provenance when represented by an observable quotient or valid lifted
interface, respectively. The quotient projection contributes a qualified
\(\Abs\) source only through \cref{def-paper1-abs-usefulness-gate}; the lift
retains the qualification of its represented base source. The remainder is
residual transport. These are properties of the presented operator,
independent of whether a representation has been discovered.

\subsection{Projected channels and causal authority}
\label{subsec-projected-channels-caus-authority}

We first define the represented channel projections and the authority envelope
for their directed component.

\begin{definition}[Budgeted channel representatives]\label{def-admissible-channel-representatives}

The representative spaces for \(\Caus\), \(\Abs\), valid
\(\Lift\), and \(\Bdry\) consist of objects in
the presented observable algebra, the public initialization lift, and
the budgeted saturated closure of lifted observable coordinates. At a
declared budget this includes budget-admissible finite-horizon
computations performed in the lifted configuration and charged at their
saturated computation cost. Stopped dynamics, powers, occupations, Green
operators, resolvents, committors, first-hitting maps of the same
operator, and ambient relations enter these representative spaces when
the denoted finite object is generated by a budget-admissible computation
from the closure, or otherwise has a budget-admissible
saturated representative at the claimed budget. A solver-indexed
future-hitting observable such as

\[
D_L(\omega)=\inf\{t\ge0:P^t\omega\in T\}
\]

for a deterministic transition rule, or the analogous stopped hitting
distribution for a randomized kernel, has the cost of an admissible
representation rather than cost zero by definition. A computed
finite-horizon version belongs to the budgeted saturated closure at the
cost of its simulation. An uncomputed future-hitting functional is a
semantic property of the induced operator rather than a representative. Thus
the criterion is extensional and cost-sensitive: any budget-admissible
saturated representation of the same denotation is sufficient.

\end{definition}

The Hodge split separates the directed-flow components used by the operator ledger.
\begin{definition}[Hodge split of directed flow]\label{def-hodge-channel-split}

Fix the represented generator rates \(q(x,y)\) and a legal full-support
reference weight \(m\).  On every ordered off-diagonal pair put

\[
r(x,y)=m(x)q(x,y),
\qquad
j(x,y)=\frac12\bigl(r(x,y)-r(y,x)\bigr).
\]

Let \(\vec{\mathcal E}\) contain one declared orientation of every unordered
pair on which \(j\) may be nonzero, including local and nonlocal pairs, and
define the actual skew-current vector by
\(F_J(e)=j(x,y)\) for \(e=(x,y)\).  Fix a represented positive edge metric
\(w=(w_e)_{e\in\vec{\mathcal E}}\) and set
\(\mathcal H_{\mathrm{edge}}=\mathbb R^{\vec{\mathcal E}}\) with

\[
\langle a,b\rangle_w
=\sum_{e\in\vec{\mathcal E}}w_ea_eb_e.
\]

The orientation, reference weight, and edge metric are declared as part of
the channel-coordinate record before the projection is evaluated.  Their
representations and all metric arithmetic are charged.  Thus the projection
below is canonical relative to a named represented metric; an arbitrary
post-hoc choice of \(w\) is not an exposed \(\Caus\) certificate.

Let \(\mathcal P_{\mathrm{str}}\le\mathbb R^{X_I}\) be the observable
finite-dimensional space of exposed structural potentials: geometric
coordinates, quotient coordinates from represented fibers, and lifted
interface coordinates satisfying Definition
\hyperref[def-admissible-channel-representatives]{Budgeted channel
representatives}. Define the gradient map

\[
G:\mathcal P_{\mathrm{str}}\to\mathcal H_{\mathrm{edge}},
\qquad
(G\Phi)(x,y)=\Phi(y)-\Phi(x).
\]

Let \(\mathcal C_{\mathrm{str}}=\operatorname{im}G\). Since
\(\mathcal H_{\mathrm{edge}}\) is finite dimensional, the orthogonal
projection \(\Pi_{\mathcal C_{\mathrm{str}}}\) is well-defined. This is
the finite-dimensional Hodge projection used below
\citep{Eckmann1944,JiangEtAl2011}. Define

\[
F_{\mathrm{Caus}}
=
\Pi_{\mathcal C_{\mathrm{str}}}F_J,
\qquad
F_{\mathrm{res}}
=
F_J-F_{\mathrm{Caus}}.
\]

The corresponding skew matrices are denoted \(J_{\mathrm{Caus}}\) and
\(\Res^{\mathrm{skew}}\). This unique orthogonal
decomposition separates the structural-gradient projection from the
directed flow orthogonal to every exposed structural gradient. A
Poisson solution, committor, first-hitting-time function, or other
future functional enlarges the potential space only when it has an
admissible representation. A finite-horizon potential produced by
budget-admissible computation is a lifted observable charged at its
saturated computation cost.

The displayed projection is an unconditional finite-dimensional
algebraic projection relative to the declared potential subspace and edge
metric. At a
budget \(B\), it is called an exposed \(\Caus\) projection only
when one represented record of cost \(\preceq B\) supplies a finite
basis for \(\mathcal P_{\mathrm{str}}\), its gradients, the orientation,
reference weight, edge metric, and the projection coefficients, or a uniform procedure that
computes them, with all arithmetic and output costs charged. The set of
objects having individual cost at most \(B\) is not silently replaced
by its free linear span. An algebraic gradient component lacking this
record remains in the budget-filtered residual ledger, although the
unfiltered orthogonal identity remains valid.

\end{definition}

The authority envelope specifies which causal currents the presentation permits the certificate to use.
\begin{definition}[Caus authority envelope]
\label{def-caus-authority-envelope}

Let a realized budget-admissible generator on \(L^2(X_I,\mu_I)\) have
the operator decomposition

\[
L=L^{\mathsf s}+A+L^{\mathsf b},
\qquad
(L^{\mathsf s})^*=L^{\mathsf s},
\qquad
A^*=-A,
\]

where the diagonal killed or boundary contribution may be included in
\(L^{\mathsf s}\) when convenient.  A Caus authority envelope at budget
\(B\) is an affine represented family

\[
\mathfrak A^{\Caus}_{I,\le B}
=
A_I^{\mathrm{imp}}+
\mathcal A^{\mathrm{ctrl}}_{I,\le B}.
\]

The imposed operator \(A_I^{\mathrm{imp}}\) is part of the declared
dynamic realization.  Each
\(A_I^{\mathrm{ctrl}}\in\mathcal A^{\mathrm{ctrl}}_{I,\le B}\) is a
skew represented control constructed from the public presentation with
its saturated cost charged.  A choice is admissible only when the total
operator

\[
L^{\mathsf s}+A_I^{\mathrm{imp}}+A_I^{\mathrm{ctrl}}
+L^{\mathsf b}
\]

satisfies the positive maximum principle, the declared support
relation, and the terminal conditions.  The three authority regimes
are

\[
\begin{array}{c|c|c}
\text{regime}&A_I^{\mathrm{imp}}&
\mathcal A^{\mathrm{ctrl}}_{I,\le B}\\ \hline
\text{free}&0&\text{all budget-admissible represented controls},\\
\text{imposed}&A_I^{\mathrm{imp}}&\{0\},\\
\text{hybrid}&A_I^{\mathrm{imp}}&
\text{a declared budget-admissible control family}.
\end{array}
\]

Thus free authority removes an external restriction on the directed
operator; it does not remove the derivation or selection cost of the
chosen control.

\end{definition}

The current record distinguishes authorized operator transport from auxiliary geometric descent.
\begin{definition}[Authorized operator current and geometric descent current]
\label{def-authorized-caus-use-current}

Fix a generator selected from
\(\mathfrak A^{\Caus}_{I,\le B}\). Its authorized
\(\Caus\) operator current is the structural-gradient
projection of the actual edge current induced by
the selected skew operator
\(A_I^{\mathrm{imp}}+A_I^{\mathrm{ctrl}}\).

Separately, a represented symmetric conductance \(c_\Phi\) and a
represented target-oriented potential \(D\) define the auxiliary
geometric descent current

\[
J_{\Phi,D}^{\mathrm{grad}}(x,z)
=-c_\Phi(x,z)\bigl(D(z)-D(x)\bigr),
\]

It may orient a \(\Geom\) drift or descent certificate when the
same full generator supplies the stated local-consistency, drift, or
carrier comparison. It is not the \(\Caus\) operator component
of that generator unless the certificate includes an explicit
represented identity equating it with the structural-gradient
projection of the actual selected current.

The potential is target-normalized: \(D\ge0\) and \(D=0\) on the
positive terminal sector. A provenance record contains the authority
envelope, selected control, the actual-current identity or auxiliary
geometric-current identity as appropriate, every quotient or coordinate
transport, and the comparison gate. All entries belong to one charged
saturated derivation.

\end{definition}

\begin{lemma}[Caus authority and no substitution]
\label{thm-caus-authority-no-substitution}

For a fixed dynamic realization and authority envelope, the exposed
Caus component is the Hodge projection of the actual edge current of
the selected skew operator. A quotient, coordinate change, or valid
lift may transport this current through a represented dynamic interface,
but may not replace it by another target-aligned current. An auxiliary
geometric descent current may be used only in its Geom certificate role,
unless an explicit represented equality identifies it with the actual
operator current. A different represented current contributes only as
a newly charged control in the free or hybrid envelope, or as residual
directed motion relative to the imposed envelope.

\end{lemma}

\begin{proof}\phantomsection\label{proof-caus-authority-no-substitution}

The Hodge projection in
\Cref{def-hodge-channel-split} is unique for the selected
generator.  Its input edge current is fixed by the imposed operator and
the selected represented control. The auxiliary geometric current is
fixed by its represented conductance and target-normalized potential,
but a same-generator comparison does not make it equal to the skew
current. A legal transport applies the represented interface to the
appropriate data and retains their derivation provenance. A different
operator current is not the image of the selected generator. It must
consequently be supplied as a new control, if the
authority envelope permits it, or remains in the orthogonal residual
ledger.  This exhausts the alternatives in the budgeted channel
representative space.

\end{proof}

\begin{theorem}[Reversible geometric carrier and exhaustive directed
classification]
\label{thm-reversible-geom-exhaustive-directed-classification}

Fix a finite presented task, a budget \(B\), and a represented generator
selected from the authority envelope of
\cref{def-caus-authority-envelope}.  Write its existing operator split as

\begin{equation}
\label{eq-reversible-geom-directed-operator-split}
L=L^{\mathsf s}+A+L^{\mathsf b},
\qquad
(L^{\mathsf s})^*=L^{\mathsf s},
\qquad
A^*=-A.
\end{equation}

The local conductance represented by \(L^{\mathsf s}\) is the
\(\Geom\) carrier.  In particular, the geometric carrier is reversible with
respect to the declared reference measure.  Let \(F_A\) be the actual edge
current of \(A\), let
\(\mathcal C_{\mathrm{str}}=\operatorname{im}G\) be the exposed structural
gradient space of \cref{def-hodge-channel-split}, and put

\begin{equation}
\label{eq-reversible-geom-caus-residual-current-split}
F_A
=
\Pi_{\mathcal C_{\mathrm{str}}}F_A
+
\bigl(I-\Pi_{\mathcal C_{\mathrm{str}}}\bigr)F_A.
\end{equation}

The first summand is the unique exposed \(\Caus\) current of the selected
generator whenever its complete projection and authority record has cost at
most \(B\).  The second summand is the directed \(\Res\) coordinate.  For
every exposed structural potential \(D\in\mathcal P_{\mathrm{str}}\), one has
the exact identities

\begin{align}
\label{eq-directed-target-pairing-factors-through-caus}
\left\langle F_A,-GD\right\rangle_w
&=
\left\langle
\Pi_{\mathcal C_{\mathrm{str}}}F_A,-GD
\right\rangle_w,\\
\label{eq-directed-residual-zero-structural-gradient-pairing}
\left\langle
\bigl(I-\Pi_{\mathcal C_{\mathrm{str}}}\bigr)F_A,-GD
\right\rangle_w
&=0.
\end{align}

Consequently, every target-oriented use of the selected directed current that
is visible at budget \(B\) is evaluated through the existing authorized
\(\Caus\) alignment factor of the reversible \(\Geom\) carrier.  A directed
term whose charged projection record is unavailable at budget \(B\), or whose
current is orthogonal to the exposed structural gradient space, remains in
the budget-filtered \(\Res\) coordinate.  If a later affordable represented
record supplies a structural potential, projection, quotient, or lift that
qualifies the term, saturated reclassification assigns it to the corresponding
existing channel at the complete cost of that record.

Thus the channel decomposition contains no independent nonreversible
geometric source.  Nonreversibility of the full selected generator is carried
by the authorized \(\Caus\) current and its directed residual complement.  An
auxiliary geometric descent current remains a derived use of the reversible
carrier unless the represented equality required by
\cref{def-authorized-caus-use-current} identifies it with the actual selected
operator current.

The general-resolvent stability factor
\(S_{q,T}^{\mathrm{nr}}(\lambda)\) of
\cref{def-normalized-resolvent-defect} does not add a source channel.  It is a
charged target-transfer comparison for a quotient of the full selected
\(\Geom+\Caus\) generator.  Its active geometric slot retains the reversible
carrier, and its causal slot retains the transported current and authority
record prescribed by the master certificate.  Failure of that provenance
sets the existing alignment factor to zero; it does not create a
nonreversible \(\Geom\) alternative.

\end{theorem}

\begin{proof}\phantomsection\label{proof-reversible-geom-exhaustive-directed-classification}

Self-adjointness of \(L^{\mathsf s}\) in the declared
\(L^2(X_I,\mu_I)\) realization is detailed balance for its off-diagonal
conductance table.  By the fixed-space and channel decompositions in
\cref{thm-fixed-space-split,thm-channel-exhaustiveness}, this symmetric local
table is exactly the \(\Geom\) coordinate, whereas the actual skew current of
\(A\) is the input to the Hodge projection.  Orthogonal projection onto
\(\mathcal C_{\mathrm{str}}\) gives
\cref{eq-reversible-geom-caus-residual-current-split}; uniqueness follows from
the orthogonal direct sum
\[
\mathcal H_{\mathrm{edge}}
=
\mathcal C_{\mathrm{str}}
\oplus
\mathcal C_{\mathrm{str}}^\perp.
\]

For \(D\in\mathcal P_{\mathrm{str}}\), the vector \(GD\) belongs to
\(\mathcal C_{\mathrm{str}}\).  The second summand of
\cref{eq-reversible-geom-caus-residual-current-split} is orthogonal to that
space, which proves
\cref{eq-directed-residual-zero-structural-gradient-pairing}; adding this zero
term proves
\cref{eq-directed-target-pairing-factors-through-caus}.  The charged-projection
criterion and the assignment of an unavailable projection to filtered
residual are part of \cref{def-hodge-channel-split}.  The authority and
same-current requirements are
\cref{def-authorized-caus-use-current,thm-caus-authority-no-substitution}.
Therefore every represented directed pairing used in a geometric certificate
is already the authorized \(\Caus\) factor, while the orthogonal or
unrepresented remainder has no exposed structural-gradient pairing at budget
\(B\).  Saturated reclassification is
\cref{prop-family-uniform-residual-reclassification,thm-no-uniform-residual-correlation}.

Finally, \cref{def-normalized-resolvent-defect} defines
\(S_{q,T}^{\mathrm{nr}}(\lambda)\) as a scalar comparison of the quotient and
full-process target functionals.  The causal-provenance and dynamic-stability
slots of \cref{def-master-compensated-quotient-geometry-certificate} retain,
respectively, the actual selected current and that comparison factor.
\Cref{prop-derived-channels-create-no-structure} prevents either derived slot
from becoming an additional primitive source.  Hence a non-self-adjoint
full-generator comparison changes neither the reversible meaning of
\(\Geom\) nor the exhaustive \(\Caus/\Res\) classification of its skew
current.

\end{proof}

\Cref{fig-caus-authority-routing} shows the no-substitution routing enforced by
the authority envelope.

\begin{figure}[tbp]
\centering
\begin{tikzpicture}[fw diagram]
  \node[fw source,minimum width=3.7cm] (selected) at (-4.6,2.7)
    {selected skew operator\\\(A^{\mathrm{imp}}+A^{\mathrm{ctrl}}\)};
  \node[fw provenance,minimum width=3.2cm] (current) at (-1.0,2.7)
    {actual current\\\(F_A\)};
  \node[fw box,minimum width=3.4cm] (hodge) at (2.6,2.7)
    {represented Hodge\\projection};
  \node[fw use,minimum width=2.7cm] (caus) at (1.0,0.9)
    {authorized \(\Caus\)};
  \node[fw residual,minimum width=2.9cm] (res) at (4.3,0.9)
    {directed \(\Res\)};
  \draw[fw arrow] (selected) -- (current);
  \draw[fw arrow] (current) -- (hodge);
  \draw[fw arrow] (hodge) -- (caus);
  \draw[fw arrow] (hodge) -- (res);
  \node[fw geom,minimum width=3.7cm] (aux) at (-4.6,-0.8)
    {auxiliary descent\\\((c_\Phi,D)\mapsto J^{\mathrm{grad}}\)};
  \node[fw terminal,minimum width=3.0cm] (gate) at (-0.8,-0.8)
    {represented equality\\\(J^{\mathrm{grad}}=F_A\)};
  \node[fw geom,minimum width=3.0cm] (geomuse) at (-4.6,-2.35)
    {Geom comparison use};
  \draw[fw arrow] (aux) -- (geomuse);
  \draw[fw dashed arrow] (aux) -- (gate);
  \draw[fw dashed arrow] (gate) -- (caus);
\end{tikzpicture}
\caption{Causal-authority routing. The \(\Caus\) current is the represented
Hodge projection of the actual selected operator current. An auxiliary
geometric descent current remains a Geom comparison unless a charged equality
record identifies it with that actual current.}
\label{fig-caus-authority-routing}
\end{figure}

The target sector does not enter this projection. Target alignment
without an exposed structural gradient is classified as an
oracle-supplied edge, a presentation-inaccessible edge, or residual
directed motion.

\section{Quotient, Lifted-State, and Operational Channels}
\label{sec-quotient-lift-operational-channels}

The nonlocal structural term now separates into represented quotient motion and
valid lifted-state transport.

\begin{definition}[Structured nonlocal quotient and lift channels]\label{def-abs-lift-split}

Let \(\mathcal H_{\mathrm{nl}}\) be the finite-dimensional Hilbert space
of all nonlocal off-diagonal operator tables remaining after the
declared \(\Geom\) and \(\Caus\) projections, with the
same weighted edge inner-product convention. Its elements are signed
operator components; an orthogonal projection need not itself be a
positive generator. An observable represented quotient \(q:X_I\to Q\)
contributes an abstraction subspace
\(\mathcal S_{\mathrm{Abs}}(q)\le\mathcal H_{\mathrm{nl}}\) spanned by
represented nonlocal operator tables whose quotient-constant subspace is
invariant and whose restriction descends to a quotient operator.
With lift \(U:L^2(Q)\to L^2(X_I)\), exact descent means

\[
L_XU=UL_Q.
\]

The abstraction subspace is the span of these observable quotient
subspaces, including symmetric or directed nonlocal tables after the
earlier projections. An observable lifted interface contributes a subspace
\(\mathcal S_{\mathrm{Lift}}\le\mathcal H_{\mathrm{nl}}\) generated by
nonlocal moves on lifted configuration coordinates: memory,
stack state, proof fragments, subproblem interfaces, or product factors.
To make the structured split unique, set

\[
\mathcal S_{\mathrm{struct}}
=\mathcal S_{\mathrm{Abs}}\oplus
\bigl(\mathcal S_{\mathrm{Lift}}\cap\mathcal S_{\mathrm{Abs}}^\perp\bigr).
\]

Projection onto \(\mathcal S_{\mathrm{Abs}}\) defines the quotient/lumping
channel coordinate \(\Abs\) at the algebraic-ledger stage, while projection onto
\(\mathcal S_{\mathrm{Lift}}\cap\mathcal S_{\mathrm{Abs}}^\perp\) defines
the lifted-state component \(\Lift\). The remaining nonlocal
edge mass has neither quotient nor lift provenance and is included in
\(\Res\). A single observed transition does not determine
symmetry of the full rate table.

The algebraic channel coordinate records quotient provenance. A represented
fiber partition contributes to the target-bearing \(\Abs\) structural source
profile precisely when its complete record satisfies the quotient-utility gate
of \cref{def-paper1-abs-usefulness-gate}, the uniform saturated-cost condition
of \cref{def-abs-extraction}, and, for prospective source use, the temporal
condition of \cref{def-pre-hit-temporally-admissible-source}. Exact quotient
records used only to preserve a completed hitting law remain in the accounting
profile of \cref{def-source-accounting-profile-separation}. Thus quotient
provenance, accounting admissibility, and qualified structural-source status
are separate stages of the same ledger.

Likewise, the spans and orthogonal projections in this definition are
unconditional algebraic coordinates. Their budget-\(B\) exposed
versions require one charged represented record containing finite bases
for the quotient and lift table spaces, the represented interfaces, and
the projection procedure and output. Individual low-cost quotient or
lift records do not make their unrestricted linear span or its
orthogonal projection free. Any algebraic portion without such a record
is retained in \(J_{\mathrm{res},B}\) for the filtered accounting.

\end{definition}

Validity identifies the lifts that preserve the declared observable and target interfaces.
\begin{definition}[Valid Lift]\label{def-valid-lift}

Write \(\mathcal F_{\mathrm{obs}}=\mathcal F_I\) for the observable
algebra of the presented task. Let \(T_I\in\mathcal F_{\mathrm{obs}}\)
be any target or terminal sector named by the task. A finite lift is
valid when it consists of a finite lifted state space \(\Omega\), a
lifted observable algebra \(\mathcal G_{\mathrm{obs}}\), a measurable
projection \(\pi:\Omega\to X_I\), and a positive lifting map
\(\Lambda:\mathcal M(X_I)\to\mathcal M(\Omega)\).  The map \(\Lambda\)
is linear, sends \(\mathcal M_+(X_I)\) into
\(\mathcal M_+(\Omega)\), and satisfies

\[
\pi^{-1}\mathcal F_{\mathrm{obs}}\subseteq\mathcal G_{\mathrm{obs}},
\qquad
\pi_*\Lambda\mu=\mu
\quad\text{for every }\mu\in\mathcal M_+(X_I).
\]

In particular, \(\Lambda\) preserves total mass and is a Markov kernel from
base states to lifted fibers.  Linearity is part of Valid Lift; it is what
permits the fiber disintegration and composition identities below.

The structural-Lift target sector is \(T_\Omega=\pi^{-1}(T_I)\). If the
instance presentation or a budget-generated computation specifies a
different terminal sector, that sector is a legal computation interface
but is a structural Valid-Lift target only when the same
target-fraction and average identities below are verified separately.
For the pullback target, every nonzero frontier satisfies the
target-fraction identity

\[
\frac{(\Lambda\mu)(T_\Omega)}{(\Lambda\mu)(\Omega)}
=
\frac{\mu(T_I)}{\mu(X_I)}.
\]

Equivalently, for every bounded \(\mathcal F_{\mathrm{obs}}\)-measurable
target-relevant observable \(h\), the lifted average of \(h\circ\pi\)
under \(\Lambda\mu\) equals the original average of \(h\) under \(\mu\).
A budget-controlled lift, including the polynomial-budget lifts used in
later examples, must satisfy these measurability and
fraction-preservation conditions. The lifted-state channel may therefore
transform, organize, and expose information in the observable algebra or
explicitly charged lifted coordinates while preserving target mass.

\end{definition}

\subsection{Resource envelopes and operational transcripts}
\label{subsec-resource-envelopes-transcripts}

Resource envelopes place the quotient and lift constructions on the complete
operational transcript carrier.

\begin{definition}[Budget-compatible resource envelope]\label{def-resource-envelope}

For a task family \((I_n)\) and budget structure
\(\mathcal B=(\mathcal R,\operatorname{cost},b,\mathfrak C)\), a
budget-compatible resource envelope bounds update steps, native
internal-state capacity, live work bits or cells, stack depth, stored
tables, learned records, query transcripts, preprocessing output,
outcome records, processor registers, and total parallel work. Data
within the envelope are
generated by a uniform rule without oracle access and are stored,
queried, or updated with accumulated charge \(\preceq b(n)\). Remaining
data are advice, oracle-supplied, or over budget. For
\(\mathcal B_{\mathrm{poly}}\), this recovers the usual polynomial
resource envelope.

Parallelism is charged by total represented work. A parallel
configuration records the finite list of active branch states, branch
weights or scheduler state, and shared memory/interface records. The
envelope bounds the total number of represented branch states and cells;
unrecorded exponential branching is therefore inadmissible.
An internal operational state is charged by its native capacity and
update rule, rather than by the cardinality of an extensional coordinate
expansion. Thus this clause charges actual replicated work and does not
replace one represented internal state by a list of hypothetical
branches.

\end{definition}

This definition charges an operational transition rule as a represented object of the presentation.
\begin{definition}[Represented operational transition rule]
\label{def-represented-operational-transition-rule}

A represented operational transition rule consists of a complete native
operational state \(z\), a finite represented outcome interface
\(A(z)\), a subprobability law \(p(\,\cdot\mid z)\) on that interface,
a represented successor map \((z,a)\mapsto z\cdot a\), and finite
represented public-readout and terminal-record interfaces. Completeness
means that the conditional law of the next native state and every
declared record is determined by the current native state. The native
transition law is

\[
P(z,B)
=
\sum_{a\in A(z)}p(a\mid z)\,
\mathbf 1_B(z\cdot a).
\]

The native state contains every internal coordinate needed
by the update semantics, together with the finite control, work, and
history records used by later transitions. It need not be a member of a
finite enumerated set. The represented initializer produces \(z_0\);
any nondeterministic initialization is incorporated into the first
finite outcome interface. Starting from that initializer, a
finite outcome prefix \(h=(a_1,\ldots,a_t)\) determines the current
native state recursively by

\[
z(\varnothing)=z_0,
\qquad
z(ha)=z(h)\cdot a.
\]

The observable algebra contains the finite outcome and terminal records
and the declared public effects of \(z(h)\); semantic internal
coordinates are not automatically adjoined as free readouts. For a
finite horizon, the represented records

\[
\xi_t=(t,a_1,\ldots,a_t,v_0,\ldots,v_t),
\]

where each \(v_j\) is the declared finite public terminal/readout record
at time \(j\), form a finite padded transcript carrier. Execution keeps
only the current native state and appends the represented records. It
does not retain the sequence \(z_0,\ldots,z_t\). Representation cost
charges the update implementation, precision, outcome and terminal
records, accumulated transcript and total work, together with the peak
native-state capacity. It does not require an extensional native-state
vector, transition table, or enumeration of native states. The
represented row or effect at an isolated prefix \(h\) may reconstruct
\(z(h)\) by replaying that prefix through the successor map, using one
live native state and charging at most \(t\) updates. Sequential
execution already carries \(z(h)\) and needs no such reconstruction.
Neither evaluation retains the intermediate native states. The
definition uses only the conditional transition law on complete native
states and makes no classification by implementation substrate.

\end{definition}

The full transcript representation places every stored operational field on one charged carrier.
\begin{definition}[Full operational transcript representation]\label{def-full-configuration-lift}

Let \(A_n\) be a represented operational transition rule under a finite
resource envelope \(\mathcal B\), and let \(H\) be the step bound supplied
by that envelope. Its full operational transcript representation is the
clocked transcript carrier \(\Omega_{A,n}^{[H]}\) of padded outcome,
public-readout, and terminal records from \Cref{def-represented-operational-transition-rule}. A legal transcript
prefix \(h\) determines the current complete native state \(z(h)\) by
the represented initializer and successor recursion. The live state may
contain the candidate state, finite control, work or memory contents,
head positions, stacks, preprocessed data, stored tables, learned
records, priority queues, active parallel branch states, scheduler
state, branch weights, and output registers. These are charged through
the peak native-state capacity and their update operations; their
successive values are not copied into the finite carrier. Any outcome,
query, public-history, or terminal record retained by the computation is
instead appended to the transcript and charged there. Malformed or
unused syntactic transcripts are completed by a fixed inert failure
transition, so the transcript update is total without changing the
legal computation.

The representation ledger distinguishes sequential execution from
isolated carrier access. Sequential execution carries the current native
state. Evaluating a kernel row or declared native-state effect from an
isolated transcript prefix replays that prefix through the outcome
successor interface and charges every repeated update, while retaining
only one native state.

The lifted observable algebra \(\mathcal G_{A,n}^{[H]}\) is generated by
the transcript coordinates and declared represented effects on the
current native state \(z(h)\). The complete native state determines the
update but is exposed only through these declared effects. The transition
representation itself requires neither a target nor a task projection. When the rule
is coupled to a presented task, this algebra also contains the pullback
of the base observable algebra, and the task interface adjoins an
exposed total readout \(\pi_A:\Omega_{A,n}^{[H]}\to X_n\), a public
initialization, and a terminal verifier. In that case the lifted target
may be \(T_{A,n}=\pi_A^{-1}(T_n)\), with any additional terminal-output
convention exposed in \(\mathcal G_{A,n}^{[H]}\). The initialization lift
\(\Lambda_A\) attaches the public initialization distribution over
auxiliary records to a base frontier.

When these task-interface data are present, the transcript
representation is additionally a valid structural Lift when
\(\pi_A^{-1}\mathcal F_n\subseteq\mathcal G_{A,n}^{[H]}\),
\(\pi_{A*}\Lambda_A\mu=\mu\), and the target-fraction identity of
\cref{def-valid-lift} holds. Preprocessing is
part of the initialized native state and is charged to \(\mathcal B\).
Finite outcomes are either recorded in the transcript or marginalized
into its Markov kernel. Persistent adaptive memory, tables, and actual
parallel branches remain in the current native state unless a declared
represented effect exposes them; retained public records are transcript
coordinates.

The finite transcript representation exists without the displayed
valid-Lift identities. Those identities are required only when the
auxiliary coordinates are used as a structural \(\Lift\) source
or when a target-fraction comparison is transported to the base task.
The Markov embedding and exact first-hit construction act directly on
\(\Omega_{A,n}^{[H]}\) and require no structural-Lift classification.

\end{definition}

\subsection{Admissibility and transcript completeness}
\label{subsec-lift-admissibility-transcript-completeness}

The following conditions identify valid lifted interfaces and establish that
resource-bounded computations have complete transcript representations.

\begin{definition}[Structural admissibility conditions for Lift]\label{def-structural-lift-gate}

A lifted component is a budget-admissible \(\Lift\) contribution at
saturated public budget \(B\) precisely when the task presentation and
the induced generator on the full represented carrier satisfy all of the
following objective conditions:

\textbf{Generation.} The lifted coordinates are generated by the public
observable algebra and the uniform transition rule without oracle access within
the resource envelope.

\textbf{Measurability.} The base observables pull back into the lifted
algebra:
\(\pi^{-1}\mathcal F_{\mathrm{obs}}\subseteq\mathcal G_{\mathrm{obs}}\).

\textbf{Projection.} The lift preserves the base frontier under
projection: \(\pi_*\Lambda\mu=\mu\) for every base frontier under
consideration.

\textbf{Reference law.} Every lifted \(L^2\) reference measure used to
prove cost, regularity, projection, or reach satisfies Definition
\hyperref[def-reference-measure-legality]{Reference measure legality}.
In particular, occupation, hitting-conditioned, invariant, Green, and
resolvent measures of the induced operator are not admissible reference
laws unless they have an independent public representation with
controlled density.

\textbf{Target-fraction preservation.} For every named target or
terminal sector \(T\), the lifted target is the pullback \(\pi^{-1}T\)
unless an alternative terminal convention is itself exposed, and the
identity \((\Lambda\mu)(\pi^{-1}T)/(\Lambda\mu)(\Omega)=\mu(T)/\mu(X)\)
holds.

\textbf{Cost.} The description length, represented state dimension,
lifted interfaces, and comparison data used by the charge have saturated
public cost at most \(B\).

These conditions are predicates of the presented or lifted operator,
rather than outputs required from an algorithm, inverter, or adversary.
When they hold at budget \(B\), the Lift component exists independently
of whether its representation has been identified. When they fail, the
component is not exposed at budget \(B\) and is classified as another
represented channel, presentation-inaccessible or oracle-supplied data,
presentation enrichment, or \(\Res\), according to the
provenance theorem.

\end{definition}

\Cref{fig-operational-to-structural-lift-gates} separates unconditional
operational embedding from the additional gates required for a structural
\(\Lift\) certificate.

\begin{figure}[tbp]
\centering
\begin{tikzpicture}[fw diagram]
  \node[fw source,minimum width=3.8cm] (rule) at (-5.2,1.5)
    {bounded operational\\transition rule};
  \node[fw box,minimum width=3.8cm] (transcript) at (-1.2,1.5)
    {finite clocked\\transcript kernel};
  \node[fw provenance,minimum width=3.5cm] (task) at (2.8,1.5)
    {represented task\\projection \(\pi\)};
  \node[fw terminal,minimum width=3.7cm] (gates) at (2.8,-0.4)
    {measurability + frontier\\+ target-fraction gates};
  \node[fw use,minimum width=3.3cm] (lift) at (6.9,-0.4)
    {charged structural\\\(\Lift\)};
  \node[fw residual,minimum width=3.8cm] (embed) at (-1.2,-0.4)
    {operational embedding\\without Lift claim};
  \draw[fw arrow] (rule) -- (transcript);
  \draw[fw arrow] (transcript) -- (task);
  \draw[fw arrow] (task) -- (gates);
  \draw[fw arrow] (gates) -- (lift);
  \draw[fw arrow] (transcript) -- (embed);
\end{tikzpicture}
\caption{Operational embedding and structural lifting are distinct. Every
bounded represented transition rule has a finite clocked transcript kernel.
It becomes a structural \(\Lift\) only after a represented task projection,
measurability, frontier-preservation, target-fraction, reference-law, and cost
gates are verified.}
\label{fig-operational-to-structural-lift-gates}
\end{figure}

\begin{proposition}[Channel admissibility conditions are operator properties]\label{prop-gate-predicates-objective}

The channel conditions used by later bounds, including quotient lumpability,
valid-lift conditions, target-fraction preservation, saturated-cost
bounds, boundary readout, and any mixing or regularity comparison, are
objective predicates of the presented task object, the lifted
configuration, and the induced operator. They are not membership
conditions for a runtime class and are not obligations imposed on the
algorithm, inverter, or adversary.

Absent a proof of a condition, the corresponding quantitative theorem
cannot assign the component a budget-admissible charge. The induced
operator nevertheless remains in the \PartII decomposition. At a fixed
budget, a component without the required representative is classified as
another represented channel, presentation-inaccessible or oracle-supplied
information, presentation enrichment, or \(\Res\).

\end{proposition}

\begin{proof}\phantomsection\label{proof-gate-predicates-objective}

The generator form, fixed-space split, Hodge projection, quotient
projection, valid-lift projection, and boundary split apply directly to
the induced operator. The admissibility conditions enter only when a
subsequent quantitative theorem assigns a cost to a component. Their
truth is independent of the availability of a proof. Without a proof
from the presentation and induced operator, that quantitative
application is unavailable, while the structural decomposition remains
valid. Verification is therefore an application hypothesis rather than
a condition for budget-admissible execution.

\end{proof}

\begin{proposition}[Full transcript representation of resource-bounded computations]\label{prop-full-configuration-lift-complete}

Every represented operational transition rule under a finite resource
envelope has a full finite transcript representation once every external
input on which its transition depends is included in the initialized
native state or in the represented interface. Every datum used by a
future transition is either retained in the current native state,
recorded in the outcome/terminal transcript, marginalized into the
induced transcript kernel, or supplied as an external input. The ledger
charges the current native state at peak capacity and the transcript at
its accumulated represented size; it does not store or charge a list of
intermediate native states. Relative to a fixed presented task, an external
instance-dependent datum not declared or generated within budget is
advice, oracle-supplied data, or over-budget information; this affects
affordability for that presentation, not existence of the transition
representation.
When the projection and initialization also satisfy
\Cref{def-valid-lift}, this representation is a valid
structural \(\Lift\).

\end{proposition}

\begin{proof}\phantomsection\label{proof-full-configuration-lift-complete}

The resource envelope bounds the horizon and each finite outcome,
public-readout, and terminal-record alphabet. Their padded transcript
carrier is therefore finite. For a legal prefix \(h\), the initializer
and successor recursion determine \(z(h)\), and the represented
conditional outcome law appends one record. Extend this rule by the
declared inert failure update on malformed transcripts. This gives a
Markov or sub-Markov kernel on the finite transcript carrier without
enumerating the native-state space. Preprocessing belongs to the
initialized native state. Adaptive memory, learned records, tables,
replicated branches, and scheduler state are updated inside the single
live native state and are charged through its peak capacity; declared
public records are appended to the transcript. A transition depending on an unrecorded datum is a
transition system with an external input and is represented after that
input is adjoined to its represented interface. For a fixed task presentation, using
such a datum without declaring or deriving it is advice, oracle
information, or data over budget for \(\mathcal B\). This proves the
unconditional finite transcript representation. The task projection
and target-fraction identities are additional Valid Lift conditions
invoked only for the corresponding structural channel.

\end{proof}

\begin{proposition}[The directed, lifted-state, and boundary channels are derived]\label{prop-derived-channels-create-no-structure}

The reusable sources of observable structure are exposed local geometry
and represented abstraction. The channels \(\Caus\),
\(\Lift\), and \(\Bdry\) process these sources. More
precisely, \(\Caus\) is the
gradient projection of directed flow onto potentials represented in the
observable or lifted algebra; \(\Lift\) is valid only when the
lifted coordinates are observable or explicitly charged and the target
fraction is preserved by projection; and \(\Bdry\) processes
exposed terminal observables into boundary value functions represented
by the observable grammar. Removing support is the zero-boundary-value
case. A committor or hitting-time value of the induced generator is
counted as \(\Bdry\) when it is represented by
exposed terminal data together with already budget-admissible
Geom/Abs/Caus/Lift structure, or when a finite-horizon approximation of
the same value is generated by budget-admissible computation and
charged as a lifted observable.

\end{proposition}

\begin{proof}\phantomsection\label{proof-derived-channels-create-no-structure}

For \(\Caus\), the definition is
\(F_{\mathrm{Caus}}=\Pi_{\mathcal C_{\mathrm{str}}}F_J\) with
\(\mathcal C_{\mathrm{str}}=\operatorname{im}G\). If the observable
structural potential space is zero, then \(\mathcal C_{\mathrm{str}}=0\)
and the causal component is zero. A directed edge that points at the
target without such a potential is residual or oracle-supplied rather
than a directed/drift structural component.

For \(\Lift\), validity gives \(\pi_*\Lambda\mu=\mu\) and
\(T_\Omega=\pi^{-1}(T_I)\). Therefore
\((\Lambda\mu)(T_\Omega)=\mu(T_I)\) and
\((\Lambda\mu)(\Omega)=\mu(X_I)\). The displayed target-fraction
identity follows. Any lifted coordinate not generated by the observable
algebra must be generated as part of a valid lifted observable state
before use. Otherwise it is presentation-inaccessible, oracle-supplied,
or over budget, and its associated movement remains residual relative to
the public presentation.

For \(\Bdry\), the general boundary form is
\((L^{\Bdry}(f,v))(x)=\kappa(x)(Bv(x)-f(x))\). With an interior
transport operator \(A\) from the other channels, the associated
boundary value function \(V\) is equivalently characterized by the
finite Dirichlet value equation

\[
(AV)(x)+\kappa(x)\bigl(Bv(x)-V(x)\bigr)=0
\qquad (x\in X_I).
\]

Thus \(\Bdry\) maps observable terminal readouts to interior
value functions, with support removal given by the special case
\(Bv\equiv0\). Because this channel has no interior off-diagonal transport
edge, it evaluates or removes mass according to exposed terminal data and
budget-admissible transport structure. The map
\(L\mapsto(\lambda I-L_N)^{-1}k_T\), or the corresponding
infinite-horizon hitting table of the same induced kernel, requires an
admissible representation to enter the boundary channel. A finite-horizon
table generated by budget-admissible simulation is charged at its
saturated computation cost. These three derived channels therefore act
only on exposed geometry, represented abstraction, or explicitly charged
lifted data.

\end{proof}

\begin{theorem}[Exhaustive boundary normal form]
\label{thm-exhaustive-bdry-normal-form}

Fix the represented boundary space \(\partial X_I\), transfer rates
\(b(x,\zeta)\ge0\), and boundary value input \(v\) from
\cref{def-boundary-value-channel}.  Put
\(\kappa(x)=\sum_{\zeta\in\partial X_I}b(x,\zeta)\), and define the two
boundary blocks

\begin{align}
\bigl(L_{\mathrm{kill}}^{\Bdry}(f,v)\bigr)(x)
&=-\kappa(x)f(x),\label{eq-bdry-killing-block}\\
\bigl(L_{\mathrm{val}}^{\Bdry}(f,v)\bigr)(x)
&=\sum_{\zeta\in\partial X_I}b(x,\zeta)v(\zeta)
=\kappa(x)Bv(x).\label{eq-bdry-value-block}
\end{align}

Every finite boundary component has the unique block decomposition

\begin{equation}
\label{eq-exhaustive-bdry-normal-form}
L^{\Bdry}
=
L_{\mathrm{kill}}^{\Bdry}
+
L_{\mathrm{val}}^{\Bdry}.
\end{equation}

The value block has the unique boundary-atom expansion

\begin{equation}
\label{eq-bdry-atom-expansion}
L_{\mathrm{val}}^{\Bdry}
=
\sum_{\zeta\in\partial X_I}
b_\zeta\,\operatorname{ev}_\zeta,
\qquad
b_\zeta(x)=b(x,\zeta),
\qquad
\operatorname{ev}_\zeta(v)=v(\zeta).
\end{equation}

Thus the algebraic \(\Bdry\) channel has exactly two blocks: homogeneous
boundary loss and inhomogeneous represented boundary-value readout.  Killing
and support removal are the specialization \(v=0\).  Terminal indicators,
support colorings, committors, hitting-time values, and other scalar terminal
tables belong to the value block precisely when the same denotation is a
declared boundary observable or has a budget-admissible represented
derivation.  A finite-horizon value generated by an admissible computation is
in the latter provenance class at its saturated computation cost.

At budget \(B\), the complete record consists of the boundary space, the rate
vectors \(b_\zeta\), the represented value input, and the evaluation
procedure.  A record exposed at cost at most \(B\) is charged to
\(\Bdry\).  The filtered-residual convention of
\cref{thm-exhaustive-residual-normal-form} retains a structurally typed
boundary block whose complete record is unavailable at that budget.  The
channel contains no interior off-diagonal transport block.

\end{theorem}

\begin{proof}\phantomsection\label{proof-exhaustive-bdry-normal-form}

The finite generator form assigns every diagonal loss not balanced by an
interior off-diagonal rate to transfer into the represented boundary space.
For the extended input pair \((f,v)\), the boundary term is

\[
\sum_{\zeta\in\partial X_I}
b(x,\zeta)\bigl(v(\zeta)-f(x)\bigr)
=
-\kappa(x)f(x)
+
\sum_{\zeta\in\partial X_I}b(x,\zeta)v(\zeta),
\]

which proves \eqref{eq-exhaustive-bdry-normal-form}.  Restriction to inputs
\((f,0)\) recovers \(L_{\mathrm{kill}}^{\Bdry}\), while restriction to
\((0,v)\) recovers \(L_{\mathrm{val}}^{\Bdry}\); hence the two-block
decomposition is unique.  The coordinate evaluations
\(\{\operatorname{ev}_\zeta\}_{\zeta\in\partial X_I}\) form the canonical
basis of \((\mathbb R^{\partial X_I})^*\).  The coefficients of the value
block in that basis are exactly the rate vectors \(b_\zeta\), proving
\eqref{eq-bdry-atom-expansion} and its uniqueness.

\Cref{prop-derived-channels-create-no-structure} identifies a represented
terminal observable or a budget-generated finite-horizon value as the
admissible input to this block.  The charged projection convention gives the
budget statement.  Both blocks map from the interior/boundary input pair to
the interior value coordinate, and neither introduces an interior
off-diagonal rate.  Therefore the two blocks exhaust \(\Bdry\).

\end{proof}

\Cref{fig-bdry-normal-form} displays the algebraic and provenance layers of
\cref{thm-exhaustive-bdry-normal-form}.

\begin{figure}[tbp]
\centering
\begin{tikzpicture}[fw diagram]
  \node[fw terminal,minimum width=3.4cm] (root) at (0,2.7)
    {boundary block \(L^{\Bdry}\)};
  \node[fw residual,minimum width=3.2cm] (kill) at (-2.8,1.25)
    {homogeneous loss\\\(-\kappa f\)};
  \node[fw terminal,minimum width=3.2cm] (value) at (2.8,1.25)
    {value readout\\\(\sum_\zeta b_\zeta v(\zeta)\)};
  \node[fw source,text width=2.6cm] (declared) at (0.7,-0.15)
    {declared\\terminal data};
  \node[fw use,text width=3.2cm] (generated) at (4.8,-0.15)
    {represented\\derived value};
  \draw[fw arrow] (root) -- (kill);
  \draw[fw arrow] (root) -- (value);
  \draw[fw arrow] (value) -- (declared);
  \draw[fw arrow] (value) -- (generated);
\end{tikzpicture}
\caption{The exhaustive boundary normal form of
\cref{thm-exhaustive-bdry-normal-form}: homogeneous killing and inhomogeneous
value readout are the two algebraic blocks; the value is declared or supplied
by a charged represented derivation.}
\label{fig-bdry-normal-form}
\end{figure}

\section{Partial Lumpability and Full-Operator Accounting}
\label{sec-partial-lumpability-accounting}

Partial lumpability is measured on the restricted block while its contribution
remains charged to the full represented operator.

\begin{definition}[Partial lumpability defect]\label{def-partial-lumpability-defect}

Let \(q:X_I\to Q\) be a represented quotient. Let
\(U:\mathbb R^Q\to\mathbb R^{X_I}\) be the lift \((Ug)(x)=g(q(x))\), and
let \(P:\mathbb R^{X_I}\to\mathbb R^Q\) be a declared projection with
\(PU=I_Q\). For a finite generator \(L_X\), the canonical quotient
operator and defect are

\[
L_Q=PL_XU,
\qquad
E_q=L_XU-UL_Q=(I-UP)L_XU.
\]

The quotient is exactly lumpable when \(E_q=0\). It is partially or
approximately lumpable when \(E_q\ne0\), with the size of \(E_q\)
measured relative to the quotient and within-fiber dynamics. This is the
standard distinction between exact and approximate aggregation of a
finite Markov chain \citep{KemenySnell1976,Buchholz1994}. The quotient
operator \(L_Q\) is a valid quotient generator only when it satisfies
the finite generator form on \(Q\); otherwise the part failing that form
is counted as defect rather than as \(\Abs\).

\end{definition}

\begin{proposition}[Full-operator accounting of partial lumpability]\label{prop-partial-lumpability-no-escape}

A partially lumpable quotient decomposes the action of \(L_X\) on
quotient observables, but that restricted identity is not by itself a
decomposition of the full operator. With \(\Pi=UP\), the exact
full-domain identity is

\begin{equation}
\label{eq-full-operator-partial-lumpability}
L_X=UL_QP+E_qP+L_X(I-\Pi).
\end{equation}

The channel projections are applied to all three terms. The represented
quotient-descending portion of \(UL_QP\) that lies in
\(\mathcal S_{\mathrm{Abs}}\) is an \(\Abs\) contribution;
the remaining local, causal, lifted, boundary, and residual portions of
the complete identity are assigned by the same channel decomposition.

\end{proposition}

\begin{proof}\phantomsection\label{proof-partial-lumpability-no-escape}

The restricted identity

\[
L_XU=UL_Q+E_q
\]

is an algebraic decomposition of the image of quotient observables.
Because \(PU=I_Q\), \(\Pi=UP\) is idempotent, and

\[
L_X=L_X\Pi+L_X(I-\Pi)
=UL_QP+E_qP+L_X(I-\Pi),
\]

which proves \eqref{eq-full-operator-partial-lumpability}. The block
\(UL_QP\) is a full-domain quotient-intertwined block; it is not
automatically an entire channel component. Apply the declared Hilbert
projections to all three full-domain terms. This is a linear-algebraic
decomposition and does not assert that \(E_qP\), \(L_X(I-\Pi)\), or
any projected piece is a positive Markov generator. A piece with a
represented geometric, causal, quotient, lift, or boundary witness is
charged to that channel, and the orthogonal complement is charged to
\(\Res\). Any hitting or resolvent statement still uses the
full positive generator \(L_X\).

\end{proof}

\Cref{fig-partial-lumpability-accounting} displays the intertwining defect
and the full-domain identity of
\cref{def-partial-lumpability-defect,prop-partial-lumpability-no-escape}.

\begin{figure}[tbp]
\centering
\begin{tikzpicture}[fw diagram]
  \node[fw source,minimum width=2.0cm] (q0) at (-3.0,2.0) {\(Q\)};
  \node[fw box,minimum width=2.0cm] (x0) at (3.0,2.0) {\(X_I\)};
  \node[fw source,minimum width=2.0cm] (q1) at (-3.0,0) {\(Q\)};
  \node[fw box,minimum width=2.0cm] (x1) at (3.0,0) {\(X_I\)};
  \draw[fw arrow] (q0) -- node[above] {\(U\)} (x0);
  \draw[fw arrow] (q0) -- node[left] {\(L_Q\)} (q1);
  \draw[fw arrow] (x0) -- node[right] {\(L_X\)} (x1);
  \draw[fw arrow] (q1) -- node[below] {\(U\)} (x1);
  \node[fw residual,minimum width=3.0cm] (defect) at (0,1.0)
    {\(E_q=L_XU-UL_Q\)};
  \draw[draw=fwInk,dashed,-{Stealth[length=2.2mm,width=1.5mm]}]
    (defect) -- (x1);
  \node[fw terminal,minimum width=8.6cm] (full) at (0,-1.7)
    {\(L_X=UL_QP+E_qP+L_X(I-UP)\)};
  \draw[fw arrow] (q1) -- (full);
  \draw[fw arrow] (x1) -- (full);
\end{tikzpicture}
\caption{Exact lumpability is the commuting case \(E_q=0\).  Partial
lumpability retains the defect and the within-fiber complement in the full
operator identity of \cref{prop-partial-lumpability-no-escape}.}
\label{fig-partial-lumpability-accounting}
\end{figure}

The residual analysis later treats iteration of approximate quotients
across a residual degree threshold. The structural conclusion used here
is that approximate intertwining supplies the full identity
\eqref{eq-full-operator-partial-lumpability}, whose three terms are
covered by the same channel projections.

Accordingly, the comparability relation and regularity functional induce
the geometric/reversible channel, represented fibers and static
lumpability induce the quotient/lumping channel, level-set filtrations
and exposed potentials induce the directed/drift channel, lifted task
states induce the lifted-state channel, and terminal readouts induce the
boundary/killing channel. Nonlocal transitions without these provenance
witnesses belong to \(\Res\).

\section{Exhaustiveness of the Structural Decomposition}
\label{sec-structural-decomposition-exhaustiveness}
\label{the-exhaustiveness-theorem}

The generator form reduces the operator to transition rates and diagonal
loss; after adjoining the boundary-value space, the latter is boundary
readout with killing as its zero-value specialization. The fixed-space
decomposition separates local symmetric, local directed, nonlocal, and
boundary terms. The structural projections below exhaust these
components.

\subsection{Exhaustive channel decomposition}
\label{subsec-exhaustive-channel-decomposition}

The first result gives the exact operator identity and its budget-filtered
interpretation.

\begin{theorem}[Component exhaustiveness]\label{thm-channel-exhaustiveness}

Every finite search generator satisfying the positive maximum principle
has, relative to any declared finite channel subspaces on
\((X_I,\mathcal F_I)\) or on a full lifted observable configuration, the
unconditional algebraic decomposition

\[
L
=
L^{\Geom}
+L^{\Caus}
+L^{\Abs}
+L^{\Lift}
+L^{\Bdry}
+L^{\Res}
\]

in the direct-sum operator Hilbert space determined by the fixed-space,
Hodge, quotient, lift, and boundary projections. In particular,

\begin{equation}
\label{eq-channel-pythagorean-accounting}
\|L\|_{\mathscr H_{\mathrm{ch}}}^2
=
\sum_{j\in\{\Geom,\Caus,\Abs,
\Lift,\Bdry\}}
\|L^j\|_{\mathscr H_{\mathrm{ch}}}^2
+\|L^{\Res}\|_{\mathscr H_{\mathrm{ch}}}^2,
\end{equation}

where \(\|\cdot\|_{\mathscr H_{\mathrm{ch}}}\) denotes the
direct-sum Hilbert norm used to define the channel projections. It is
not the induced operator norm.

At budget \(B\), a displayed component is an exposed structural channel
only when its subspace, projection, and output have the charged
representation required above. Algebraic portions lacking that record
are included in the filtered residual \(L^{\Res,B}\). This
budget convention changes no full-operator identity or Pythagorean
identity for the corresponding declared algebraic subspaces.

Here the local relation, structural potentials, quotient subspaces, and
valid lift subspaces are exactly those generated by the observable
algebra of the presented or lifted task. The reference weight \(m\) only
fixes conductance/current coordinates. The exposed structural
components are the geometric/reversible \(\Geom\),
directed/drift \(\Caus\), quotient/lumping \(\Abs\),
lifted-state \(\Lift\), and boundary/killing
\(\Bdry\) channels. The term \(L^{\Res}\) is the
orthogonal remainder left after the declared projections. It includes
directed local cycle/current components orthogonal to the represented
gradient space as well as unclassified nonlocal transport. These terms
exhaust the algebraic operator decomposition together with the five
structural projections; their budget-exposed versions use the charged
projection convention above.

\end{theorem}

\begin{proof}\phantomsection\label{proof-channel-exhaustiveness}

By the finite generator theorem, the operator consists only of
nonnegative off-diagonal rates and nonnegative diagonal loss. In
zero-boundary reduction this loss is killing. With terminal values
represented, it is the boundary value operator \(\kappa(Bv-f)\). This is
\(\Bdry\).

Fix the exposed local relation \(E_I\) and reference weight \(m\).
The off-diagonal rate table splits uniquely into local symmetric, local
antisymmetric, and nonlocal entries in these coordinates. Local
symmetric conductance is \(\Geom\). The weight \(m\) provides
conductance-current coordinates for the same rate table and does not
alter transitions or target data.

Apply the declared Hodge projection to the directed current. This gives
the gradient component \(\Caus\) and an orthogonal directed
remainder, which may contain local cycle current and is assigned to
\(\Res\). After setting aside that remainder, the remaining
nonlocal table lies in the finite Hilbert space
\(\mathcal H_{\mathrm{nl}}\). Orthogonal
projection onto the observable
quotient subspace \(\mathcal S_{\mathrm{Abs}}\), then onto the valid
observable lifted subspace
\(\mathcal S_{\mathrm{Lift}}\cap\mathcal S_{\mathrm{Abs}}^\perp\), gives
\(\Abs\) and \(\Lift\). Partial lumpability gives no
additional case: the full identity
\eqref{eq-full-operator-partial-lumpability} is projected by this same
component calculus, and no restricted block is silently identified with
the full \(\Abs\) component. Any remaining nonlocal component
is also assigned to \(\Res\). These projections establish
the algebraic identity. Restricting exposed terms to charged projection
records and moving every unrepresented projected denotation into the
filtered residual gives the budget-indexed statement. Hence the five
structural projections and their residual complement exhaust the rate
table.

\end{proof}

The residual complement has its own canonical normal form within the same
projection calculus.
\begin{theorem}[Exhaustive residual normal form]
\label{thm-exhaustive-residual-normal-form}

Fix the data and projection order of
\cref{def-fixed-space-split,def-hodge-channel-split,def-abs-lift-split}, and
let \(L\) be a finite search generator covered by
\cref{thm-channel-exhaustiveness}.  Let \(F_J\) be the complete directed-flow
vector passed to the Hodge projection.  Define the directed residual
coordinate \(L_{\mathrm{dir}}^{\Res}\) to be the operator represented by

\begin{equation}
\label{eq-residual-directed-normal-form}
F_{\mathrm{res}}
=
\bigl(I-\Pi_{\mathcal C_{\mathrm{str}}}\bigr)F_J,
\end{equation}

After setting aside \(L_{\mathrm{dir}}^{\Res}\), let
\(L_{\mathrm{nl}}\in\mathcal H_{\mathrm{nl}}\) be the nonlocal table
remaining after the declared \(\Geom\) and \(\Caus\) projections, in the
projection order of \cref{thm-channel-exhaustiveness}.  Define the nonlocal
residual coordinate by

\begin{equation}
\label{eq-residual-nonlocal-normal-form}
L_{\mathrm{nl}}^{\Res}
=
\bigl(I-\Pi_{\mathcal S_{\mathrm{struct}}}\bigr)L_{\mathrm{nl}},
\qquad
\mathcal S_{\mathrm{struct}}
=
\mathcal S_{\mathrm{Abs}}
\oplus
\bigl(\mathcal S_{\mathrm{Lift}}
\cap\mathcal S_{\mathrm{Abs}}^\perp\bigr).
\end{equation}

Then the algebraic residual has the unique orthogonal decomposition

\begin{equation}
\label{eq-exhaustive-residual-normal-form}
L^{\Res}
=
L_{\mathrm{dir}}^{\Res}
+
L_{\mathrm{nl}}^{\Res},
\qquad
\bigl\langle L_{\mathrm{dir}}^{\Res},
L_{\mathrm{nl}}^{\Res}\bigr\rangle_{\mathscr H_{\mathrm{ch}}}=0.
\end{equation}

Consequently,

\begin{equation}
\label{eq-residual-pythagorean-accounting}
\|L^{\Res}\|_{\mathscr H_{\mathrm{ch}}}^2
=
\|L_{\mathrm{dir}}^{\Res}\|_{\mathscr H_{\mathrm{ch}}}^2
+
\|L_{\mathrm{nl}}^{\Res}\|_{\mathscr H_{\mathrm{ch}}}^2.
\end{equation}

The first summand is precisely the directed flow orthogonal to every exposed
structural gradient.  It includes cycle current and, when
\(\mathcal P_{\mathrm{str}}\) is a proper subspace of the full potential
space, the orthogonal components of gradients whose potentials have no
represented structural realization.  The second summand is precisely the
remaining nonlocal table orthogonal to both the represented quotient
subspace and the orthogonalized valid-lift subspace.  These two summands
exhaust the unfiltered residual.

At a budget \(B\), every term retained in the filtered residual
\(L^{\Res,B}\) has exactly one of the following ledger provenances:

\begin{enumerate}[label=\textup{(\roman*)}]
\item the directed algebraic remainder \(L_{\mathrm{dir}}^{\Res}\);
\item the nonlocal algebraic remainder \(L_{\mathrm{nl}}^{\Res}\);
\item an algebraic component in one of the five ordered channel slots
\(\Geom,\Caus,\Abs,\Lift,\Bdry\) whose subspace, projection procedure, or
projected output lacks the complete charged representation required at
budget \(B\).
\end{enumerate}

The third case retains the unique algebraic channel tag supplied by the
ordered orthogonal decomposition; it is a budget-ledger assignment rather
than an additional algebraic summand.  Whenever its complete charged record
is exposed at a budget, the charged projection convention assigns the
component to its algebraic channel, and orthogonality removes it from the
filtered residual.
Every partial-lumpability defect and within-fiber complement is passed through
the same decomposition by
\cref{prop-partial-lumpability-no-escape}.  Thus partial quotients, finite
composition, lifted transcripts, memory, and adaptive branching create no
further residual class.

\end{theorem}

\begin{proof}\phantomsection\label{proof-exhaustive-residual-normal-form}

By \cref{thm-fixed-space-split}, the fixed-space rate table is the unique sum
of local symmetric, local directed, nonlocal, and boundary terms.  The local
symmetric term is \(\Geom\), and the boundary term is \(\Bdry\).
\Cref{prop-no-third-local-primitive} excludes any further local primitive.

Apply the Hodge projection of \cref{def-hodge-channel-split} to the complete
directed-flow vector.  Orthogonal projection gives

\[
F_J
=
\Pi_{\mathcal C_{\mathrm{str}}}F_J
+
\bigl(I-\Pi_{\mathcal C_{\mathrm{str}}}\bigr)F_J.
\]

The first term is the \(\Caus\) coordinate and the second is
\eqref{eq-residual-directed-normal-form}.  After this directed remainder is
set aside, \cref{def-abs-lift-split} places the remaining nonlocal table in
\(\mathcal H_{\mathrm{nl}}\).  Its structural subspace is the orthogonal sum

\[
\mathcal S_{\mathrm{Abs}}
\oplus
\bigl(\mathcal S_{\mathrm{Lift}}
\cap\mathcal S_{\mathrm{Abs}}^\perp\bigr).
\]

Projection onto the two displayed summands gives \(\Abs\) and \(\Lift\),
respectively.  Projection onto their orthogonal complement gives exactly
\eqref{eq-residual-nonlocal-normal-form}.  The fixed-space blocks and the
successive structural projections are the direct-sum coordinates of
\(\mathscr H_{\mathrm{ch}}\).  Hence the two residual coordinates are
orthogonal and give
\eqref{eq-exhaustive-residual-normal-form} and
\eqref{eq-residual-pythagorean-accounting}.  The symmetric and boundary
blocks have already been exhausted, so no third algebraic residual coordinate
remains.

The charged projection convention in
\cref{thm-channel-exhaustiveness,def-hodge-channel-split,def-abs-lift-split}
retains in \(L^{\Res,B}\) every algebraic projected denotation whose complete
projection record is unavailable at budget \(B\).  The unfiltered ordered
orthogonal decomposition supplies its unique channel tag, proving the third
ledger case without replacing the budgeted objects by a free linear span.
The charged projection convention removes every channel-qualified term once
its complete representation record is exposed.  Finally,
\cref{prop-partial-lumpability-no-escape} sends all three
terms of the full partial-lumpability identity back through these same
projections.  By
\cref{prop-full-configuration-lift-complete}, a resource-bounded operational
rule is represented on its complete finite lifted transcript carrier.
\Cref{thm-channel-exhaustiveness} applies the identical calculation on that
carrier, including memory, branching, and retained outcomes.  This exhausts
the algebraic and budget-filtered alternatives.

\end{proof}

\begin{corollary}[Completeness of the residual provenance normal form]
\label{cor-no-residual-provenance-outside-normal-form}

For every presentation-relative finite generator on its complete represented
carrier, each admitted residual operator term is either
Hodge-orthogonal directed flow or quotient/lift-orthogonal nonlocal transport.
At a fixed budget it may instead be a structurally typed algebraic component
retained by the filtered-residual convention until its complete charged
record is exposed.  These alternatives exhaust the residual provenance of
the admitted generator.

\end{corollary}

\begin{proof}\phantomsection\label{proof-no-residual-provenance-outside-normal-form}

Apply \cref{thm-exhaustive-residual-normal-form} on the complete represented
carrier supplied by \cref{prop-full-configuration-lift-complete}.

\end{proof}

A transition rule using data outside the admitted presentation is assigned by
\cref{prop-oracle-accounting} to the corresponding enriched presentation.
Target arrival produced by arbitrary interaction of the two algebraic
residual coordinates with the five channel coordinates is included in the
full-process bounds of
\cref{prop-residual-excess-reclassification,cor-residual-family-baseline}.

\Cref{fig-channel-decomposition} displays the structural classification
given by \cref{thm-channel-exhaustiveness}, together with the Abs provenance
classification of \cref{thm-five-family-abs-provenance}.
\Cref{fig-residual-normal-form} displays the exhaustive residual normal form
of \cref{thm-exhaustive-residual-normal-form}.

\begin{figure}[tbp]
\centering
\begin{tikzpicture}[fw diagram]
  \node[fw box,minimum width=4.0cm] (root) at (0,3.6)
    {structural part of represented \(L\)};
  \coordinate (channelSplit) at (0,2.9);
  \node[fw geom,minimum width=1.75cm] (geom) at (-5.2,2.0) {\(\Geom\)};
  \node[fw use,minimum width=1.75cm] (caus) at (-2.6,2.0) {\(\Caus\)};
  \node[fw source,minimum width=1.75cm] (abs) at (0,2.0) {\(\Abs\)};
  \node[fw use,minimum width=1.75cm] (lift) at (2.6,2.0) {\(\Lift\)};
  \node[fw terminal,minimum width=1.75cm] (bdry) at (5.2,2.0) {\(\Bdry\)};
  \draw[fw arrow] (root) -- (channelSplit);
  \foreach \n in {geom,caus,abs,lift,bdry}
    \draw[fw arrow] (channelSplit) -| (\n);

  \coordinate (absSplit) at (0,1.15);
  \node[fw provenance,minimum width=1.75cm] (add) at (-4.0,0.2) {\(\mathsf{Add}\)};
  \node[fw provenance,minimum width=1.75cm] (mult) at (-2.0,0.2) {\(\mathsf{Mult}\)};
  \node[fw provenance,minimum width=1.75cm] (ord) at (0,0.2) {\(\mathsf{Ord}\)};
  \node[fw provenance,minimum width=1.75cm] (sym) at (2.0,0.2) {\(\mathsf{Sym}\)};
  \node[fw provenance,minimum width=1.75cm] (filt) at (4.0,0.2) {\(\mathsf{Filt}\)};
  \draw[fw arrow] (abs) -- node[right,font=\scriptsize] {provenance} (absSplit);
  \foreach \n in {add,mult,ord,sym,filt}
    \draw[fw arrow] (absSplit) -| (\n);
\end{tikzpicture}
\caption{The five structural channels of
\cref{thm-channel-exhaustiveness}.  The Abs channel has the five exhaustive,
possibly overlapping provenance families of
\cref{thm-five-family-abs-provenance}.  The residual complement is classified
separately in \cref{fig-residual-normal-form}.}
\label{fig-channel-decomposition}
\end{figure}

\begin{figure}[tbp]
\centering
\begin{tikzpicture}[fw diagram]
  \node[fw box,minimum width=3.5cm] (root) at (0,3.2)
    {filtered residual \(L^{\Res,B}\)};
  \coordinate (split) at (0,2.45);
  \node[fw residual,minimum width=3.6cm] (dir) at (-4.3,1.2)
    {Hodge-orthogonal\\directed remainder};
  \node[fw residual,minimum width=3.8cm] (nl) at (0,1.2)
    {quotient/lift-orthogonal\\nonlocal remainder};
  \node[fw use,minimum width=3.6cm] (unexposed) at (4.3,1.2)
    {unexposed structural\\portion at budget \(B\)};
  \node[fw provenance,minimum width=5.0cm] (slots) at (4.3,-0.25)
    {\(\Geom\quad\Caus\quad\Abs\quad\Lift\quad\Bdry\)};
  \draw[fw arrow] (root) -- (split);
  \foreach \n in {dir,nl,unexposed}
    \draw[fw arrow] (split) -| (\n);
  \draw[fw arrow] (unexposed) -- (slots);
  \node[font=\scriptsize,fwInk] at (-2.15,-0.15)
    {orthogonal sum \(L^{\Res}\)};
\end{tikzpicture}
\caption{Exhaustive residual classification from
\cref{thm-exhaustive-residual-normal-form}.  The two gray leaves form the
orthogonal unfiltered residual.  The budget-filtered ledger also retains a
structurally typed portion until its complete charged projection record is
exposed.}
\label{fig-residual-normal-form}
\end{figure}

\begin{example}[Running example: the LPN transcript channels]
\label{ex-lpn-running-part2}

Apply \cref{thm-channel-exhaustiveness} to any represented clocked LPN
transcript generator built from the presentation in
\cref{ex-lpn-running-part1}. Its positive operator decomposes into
\(\Geom,\Caus,\Abs,\Lift,\Bdry\), and \(\Res\) on that transcript
carrier. This is the complete channel ledger; the LPN source analysis in
\cref{part:lpn} determines which target-aligned source terms remain.

\end{example}

\begin{remark}[Operator components and executable dynamics]
\label{rem-channel-components-need-not-be-generators}

The decomposition is an exact operator identity. An orthogonal
component or a partial sum of components need not preserve positivity
and need not be a Markov generator. Hitting probabilities, resolvents,
and first-arrival laws are therefore evaluated for the full positive
generator \(L\). Component terms provide structural accounting; they
are used as independent dynamics only when positivity and the finite
generator form are verified separately. In particular,
\(L^{\Res}\) is an orthogonal remainder, not an executable
residual process by default.

\end{remark}

\begin{proposition}[Channel decomposition is invariant under admissible presentation
equivalence]\label{prop-channel-functoriality-presentation-equivalence}

Let two presentations be admissibly equivalent in the sense of \PartI,
and suppose the equivalence transports the dynamic data: admissible edge
algebras, local relations, represented quotient subspaces, valid lift
interfaces, terminal readouts, and the state and edge Hilbert structures
by the unitary pullbacks required in \Cref{def-admissible-presentation-equivalence}. If \(L'\) is the transported generator
of \(L\), then each component transports to the corresponding component:

\[
\begin{aligned}
L^{\Geom}&\leftrightarrow L'^{\Geom}, &
L^{\Caus}&\leftrightarrow L'^{\Caus}, &
L^{\Abs}&\leftrightarrow L'^{\Abs},\\
L^{\Lift}&\leftrightarrow L'^{\Lift}, &
L^{\Bdry}&\leftrightarrow L'^{\Bdry}, &
L^{\Res}&\leftrightarrow L'^{\Res}.
\end{aligned}
\]

Thus an admissible recoding preserves each structural channel and may
change only the coordinate cost of its transported representation.
A bounded-density but nonunitary reweighting is not covered by this
exact componentwise assertion.

\end{proposition}

\begin{proof}\phantomsection\label{proof-channel-functoriality-presentation-equivalence}

The generator form is invariant under conjugating a finite Markov
generator by a positive observable-algebra isomorphism: off-diagonal
rates remain rates and diagonal loss remains diagonal loss. The
transported local relation sends local symmetric conductance to local
symmetric conductance and local oriented current to local oriented
current. The Hodge projection is the orthogonal projection onto the
transported structural-gradient subspace, so causal and orthogonal
residual parts correspond. Represented quotient subspaces and valid lift
subspaces are transported by hypothesis, so the Abs and Lift orthogonal
projections correspond. Terminal readouts are transported by terminal
preservation, so Bdry corresponds. The residual is the orthogonal
complement left after these transported subspaces are removed, so it
corresponds as well.

\end{proof}

\begin{corollary}[Completeness of the structural channel classification]\label{cor-no-cleverness}

A budget-admissible operator obtains observable target-relevant
structure through the five exposed structural channels. The intrinsic
sources are exposed local geometry
\(\Geom\) and represented abstraction \(\Abs\). The
other exposed channels are uses of those sources: \(\Caus\)
orients exposed potentials, valid \(\Lift\) reorganizes
observable or generated coordinates while preserving target fractions,
and \(\Bdry\) reads terminal observables as value functions. A
budget-admissible method without oracle access either represents one of
these objects on its clocked transcript carrier or through a declared
native-state effect, or contributes to
the residual operator component, which includes local cycle current and
unclassified nonlocal transport.  After saturated reclassification, this
residual is an operator remainder and contains no affordable target-aligned
source.

\end{corollary}

\begin{proof}\phantomsection\label{proof-no-cleverness}

\Cref{thm-information-provenance-exhaustion} partitions every target-relevant edge effect,
after the full transcript and declared native-state effects are included, into public structural
data, lifted/generated data, and residual operator transport. The public
structural part is then decomposed by
\cref{thm-channel-exhaustiveness} into the
five exposed channels. The classification of generative sources identifies
the only intrinsic reusable sources among those channels as valid
\(\Geom\), exact \(\Abs\), and their coupled
quasi-lumpable case. \(\Caus\), \(\Lift\), and
\(\Bdry\) retain the source they orient, reorganize, or read in
their complete carrier-cost record. Their conditional target charge is
transferred to that source only when a represented comparison proves the
transfer inequality. By definition, every effect outside these projections belongs to
the orthogonal residual \(\Res\), completing the
classification.

\end{proof}

The same conclusion admits an information-provenance formulation.
Public-algebra factorization, channel decomposition,
partial-lumpability defect accounting, and finite-composition closure
exhaust the origins of target-relevant information. Later low-degree and
enumeration results quantify residual correlation but are not required
for this structural partition.

\section{Information Provenance and the Residual Band}
\label{sec-information-provenance-residual-band}

Having decomposed the operator, we classify the represented origin of each
state-dependent edge rule.

\begin{definition}[Information provenance of an edge rule]\label{def-information-provenance}

Fix the finite observable algebra \(\mathcal F_I\) with atom map
\(\alpha:X_I\to\mathcal A_I\). A provenance witness for an edge rule
\(q:X_I\times X_I\to[0,\infty)\) is a finite list of auxiliary variables
\(D=(D_1,\ldots,D_r)\) such that \(q\) is measurable with respect to the
finite algebra generated by \((\alpha(x),\alpha(y),D)\). The witness is
a mathematical object associated with the presentation or induced
operator, independently of its algorithmic discovery. Relative to the
task presentation, each auxiliary variable has one of three provenances:

\textbf{Public.} It is generated by \(\mathcal F_I\), or by the declared
product algebra \(\mathcal F_I\otimes\mathcal F_I\) for edge data.

\textbf{Lifted/generated.} It is generated beyond the base observable
algebra by a budget-admissible computation without oracle access from the
public presentation and is present as a coordinate in the lifted
observable configuration.

\textbf{Residual.} After the public and lifted/generated data have been
included, the edge flow lies outside the observable structural
subspaces. It is recorded as \(\Res\). This includes
directed local cycle current orthogonal to every exposed structural
gradient, as well as observed nonlocal jumps whose full
symmetric/reversible rate table is not part of the observable structure.
This provenance tag records operator transport only.  Any affordable
represented target-aligned use is reclassified and charged before the
post-saturation residual is formed.

\end{definition}

\begin{lemma}[Exhaustion of information provenance]\label{thm-information-provenance-exhaustion}

Let \(L\) be a finite budget-admissible generator on the full clocked
transcript carrier of a budget-admissible operational rule without
oracle access. Assume that the transcript contains every retained finite
outcome, query, public-history, and terminal record, and that every
declared effect of the current native state is represented and charged.
Then every target-relevant edge effect has one of the
following provenances:

\[
\begin{gathered}
\text{public structural}
\quad\sqcup\quad
\text{lifted/generated data}
\quad\sqcup\quad
\text{residual operator transport},\\
\text{where the residual includes local cycle current and unclassified
nonlocal transport}.
\end{gathered}
\]

After budget-generated data are lifted into the observable state, the
combined represented public and lifted/generated part is decomposed by
the five exposed channel projections
\[
\Geom,\quad
\Caus,\quad
\Abs,\quad
\Lift,\quad
\Bdry.
\]
The term \(\Res\) records the movement left after these
projections and has the exhaustive directed/nonlocal normal form of
\cref{thm-exhaustive-residual-normal-form}.

\end{lemma}

\begin{proof}\phantomsection\label{proof-information-provenance-exhaustion}

\textbf{Step 1: public factorization.} A finite function or relation is
\(\mathcal F_I\)-measurable precisely when it is constant on atoms.
Hence any admissible rate table \(q(x,y)\) factors through the product
atom map \((x,y)\mapsto(\alpha(x),\alpha(y))\) whenever that table is
public-\(\mathcal F_I\otimes\mathcal F_I\)-measurable. This accounts for
the public part. Tables depending on generated coordinates are handled
after adjoining those coordinates in Step 3.

\textbf{Step 2: structural partition.} The positive maximum principle
gives rates and diagonal loss. The fixed-space split partitions the
rates into local symmetric, local directed, and nonlocal parts, and
records the diagonal loss as boundary readout. The Hodge projection
partitions directed flow into the structural-gradient subspace and its
residual. The quotient/valid-lift projections extract the nonlocal flow
with represented abstraction or target-fraction-preserving lift
provenance. A partially lumpable quotient supplies the full-domain
identity
\(L=UL_QP+E_qP+L(I-UP)\); all three terms are passed through the channel
projections, and only the quotient-descending projection of the first
term is \(\Abs\). The boundary component processes terminal observables
into value functions; if the terminal value is zero, it is ordinary
killing. Every remaining directed or nonlocal edge effect has no observable
structural provenance and is assigned to \(\Res\).

\textbf{Step 3: lifted/generated data become observable coordinates.} If
an edge rule depends on data outside \(\mathcal F_I\), then those data are
either generated by a budget-admissible computation from public history
or are presentation-inaccessible, oracle-supplied, or over budget. In
the generated case replace \(X_I\) by the lifted space
\(\Omega=X_I\times D\), where \(D\) is the finite generated transcript
or declared-readout component. The datum is now a coordinate observable on
\(\Omega\), so Step 2 applies on the lifted space. If \(D\) contains
target information that is neither primitive public data nor generated
by the charged derivation just described, then its use is
presentation-inaccessible, oracle-supplied, nonuniform, or over budget.
Target information actually derived from public data by that
budget-admissible computation remains in the lifted/generated case and
is charged at the cost of the derivation.

\textbf{Step 4: conditional outcomes give positive mixing.} At a
complete operational state, the induced kernel is the nonnegative sum
of its represented outcome transitions. The positive cone and the
finite structural subspaces are closed under finite nonnegative sums.
Mixing typed edges therefore preserves the channel subspaces and their
residual projection.

\textbf{Step 5: finite composition is lifted state evolution.} A finite
computation \(g_{t+1}=F_t(g_t,\phi_1,\ldots,\phi_k,D,\xi_t)\) is represented
by its outcome/terminal transcript together with the current native
field \(g_t\), whose peak capacity is charged. Intermediate native
fields are not automatically observable coordinates. If the final field supplies an exposed
budget-admissible structural potential, only the authorized Hodge
projection of the actual selected current is a \(\Caus\)
operator component; an auxiliary potential may instead support a Geom
descent certificate. If the final field supplies represented level
sets, they are a candidate quotient. They enter \(\Abs\) only
after terminal compatibility, exact descent, and nonzero utility are
verified; otherwise the partial-lumpability and residual calculus
applies. If the computation uses a family, stack,
table, or interface, that organization is \(\Lift\) only when
the lift is observable, charged, and target-fraction preserving. A
stopped-future functional produced by budget-admissible
finite-horizon computation is charged as such a lifted observable when
it satisfies the channel conditions; an uncomputed semantic future of
the same operator has no such representation. If the computation supplies none of these
structural coordinates, or if the proposed lift enriches the target
fraction without being charged as operator information, its nonlocal
movement is projected to \(\Res\) or classified as
inadmissible oracle data. These alternatives exhaust the finite
transcript coordinates and declared native-state effects.

\end{proof}

\begin{proposition}[Reclassification of family-uniform residual alignment]\label{prop-family-uniform-residual-reclassification}

Let \((I_n)\) be a task family and let \(J_{\mathrm{res},I}\) be the
residual component of a budget-admissible generator after the five
exposed structural projections. Suppose a family-uniform
budget-admissible relation, potential, quotient, lifted interface,
terminal readout, or finite-horizon trajectory statistic is generated
from the public presentations, including through presentation-admissible
oracle evaluators, and is verified to
satisfy the defining predicate of one of the five channel subspaces.
Then its projection onto that subspace is reclassified into one of
\(\Geom\), \(\Caus\), \(\Abs\),
\(\Lift\), or \(\Bdry\). An uncomputed committor,
first-hitting-time partition, occupation law, or stopped-future
functional of the residual operator is a witness only when it also has a
budget-admissible representation, including one computed at a
budget-admissible finite horizon.
Consequently a component that remains \(\Res\) has no
family-uniform budget-admissible one-step structural representative.
If the full process has positive excess over the declared neutral baseline,
\cref{thm-prospective-selector-accountability} exposes the represented pre-hit
sampler carrying that excess. At saturated cost, an affordable target-aligned
source in its ancestry is reclassified into an exposed channel; otherwise the
purported source is over budget, nonuniform advice, presentation-inaccessible,
or an inadmissible oracle enrichment. The completed first-hit quotient remains
an accounting record of the full process.

\end{proposition}

\begin{proof}\phantomsection\label{proof-family-uniform-residual-reclassification}

By \PartI's family-uniformity clause, a target-aligned construction that
works across the family by one public rule is an observable family
invariant, provided the construction has a
budgeted channel representative in the sense of Definition
\hyperref[def-admissible-channel-representatives]{Budgeted channel
representatives}. If the rule is produced by a budget-admissible
computation, after expanding every presentation-admissible oracle call by its
represented evaluator, its records are coordinates in the lifted
observable configuration. Thus the aligned edge predicate, potential,
quotient, lifted interface, or terminal readout is measurable in the
presented or lifted observable algebra. The additional hypothesis that
it satisfies a channel predicate, rather than measurability alone,
makes it a structural representative.

The information-provenance theorem applies on that lifted observable
space. An aligned edge effect with a verified budgeted channel
representative has public or lifted/generated provenance, so the channel
exhaustiveness theorem assigns it to the appropriate exposed component:
local geometry, causal gradient, represented abstraction, valid lift, or
boundary readout. The residual operator's own committor,
first-hitting-time partition, occupation law, or stopped-future
functional requires an admissible representation. If the same
finite object is generated by budget-admissible computation and its
induced operator effect satisfies a channel predicate, that structural
projection is reclassified. Any orthogonal remainder stays residual.
If no budget-admissible channel-qualified representative exists, any semantic
target pairing has no admissible source representation at that budget. The
operator term remains one of the two algebraic coordinates in
\cref{thm-exhaustive-residual-normal-form}; it supplies no affordable
target-aligned source in the saturated ledger. The prospective hazard identity
accounts for neutral contact and exposes every affordable excess, while the
completed first-hit quotient records the resulting arrival retrospectively.

\end{proof}

\begin{lemma}[No represented structural channel witness remains residual]\label{thm-no-uniform-residual-correlation}

After the five exposed structural projections and the family-uniform
reclassification step have been performed, \(\Res\) contains
no family-uniform budget-admissible target-aligned \emph{channel
representative} lying in one of the five projected structural subspaces
while remaining in \(\Res\). Equivalently, if a represented
edge effect, quotient, potential, structural Valid Lift, or boundary
readout satisfies the defining predicate of an exposed channel, its
orthogonal projection to that channel is removed from the residual.

The algebraic projection and the saturated affordability statement are
distinct. A semantic pairing of the algebraic remainder with the target has no
algorithmic force at budget \(B\) without a family-uniform admissible
representation. Once such a representation is generated affordably, the
family-uniform reclassification and post-exhaustion irreducibility results
assign its target-aligned component to an exposed channel. Thus the saturated
residual contains no affordable target-aligned source. Neutral target contact
is retained in the declared enumeration baseline, and every affordable excess
is carried by the prospective sampler of
\cref{thm-prospective-selector-accountability}.

\end{lemma}

\begin{proof}\phantomsection\label{proof-no-uniform-residual-correlation}

Let \(S_1,\ldots,S_5\) be the five declared closed channel subspaces in
the finite operator Hilbert space, orthogonalized in the order fixed by
the channel definition, and let \(R\) be their orthogonal complement.
The residual is \(\Pi_RL\). If a represented structural witness gives a
component \(K\in S_j\), then \(K\perp R\), so its \(S_j\)-projection is
part of the exposed channel and cannot simultaneously be a nonzero
component of \(\Pi_RL\). This is precisely the asserted
reclassification. This proves the algebraic assertion. The affordability
assertion follows by applying
\cref{prop-family-uniform-residual-reclassification,thm-irreducible-band-closure}
to every family-uniform represented target-aligned use at its least saturated
cost. A denotation lacking such a representation is inactive at the declared
budget.

\end{proof}

\begin{corollary}[No post-saturation residual source alternative]
\label{cor-no-post-saturation-residual-source-alternative}

Fix a budget \(B\) and form the saturated closure before taking the residual
coordinate.  Every family-uniform, budget-admissible represented mechanism
invoked to explain target progress passes through the following exhaustive
ordered accounting:

\begin{enumerate}[label=\textup{(\roman*)}]
\item every represented projection satisfying an exposed channel predicate is
      assigned to \(\Geom\), \(\Caus\), \(\Abs\), \(\Lift\), or \(\Bdry\)
      at its complete saturated cost;
\item a \(\Caus\), valid-\(\Lift\), or \(\Bdry\) use inherits the contextual
      charge of the represented \(\Geom\), \(\Abs\), or coupled source that it
      orients, transports, or reads;
\item the post-projection orthogonal component lies in \(\Res\) and supplies no
      prospective structural source charge.
\end{enumerate}

This ordered accounting is exhaustive.  In particular, an affordable represented
target-aligned use of an operator component cannot remain a residual source:
its least-cost channel-qualified representative is removed from \(\Res\) and
charged before the post-saturation residual is formed.  Neutral target contact
belongs to the declared enumeration baseline.  Every affordable excess over
that baseline is carried by the represented pre-hit selector and its charged
source ancestry.  A completed first-hit record remains an accounting record
and creates no additional source alternative.

\end{corollary}

\begin{proof}\phantomsection\label{proof-no-post-saturation-residual-source-alternative}

Channel exhaustiveness and the full residual normal form are
\cref{thm-channel-exhaustiveness,thm-exhaustive-residual-normal-form}.  The
no-creation identity for the three derived channels is
\cref{prop-derived-channels-create-no-structure}.  Apply
\cref{prop-family-uniform-residual-reclassification} at the least saturated
cost of every affordable represented target-aligned use.  Its structural
projection is assigned to the corresponding exposed channel, while the
orthogonal complement remains one of the two residual operator coordinates.
\Cref{thm-no-uniform-residual-correlation} then gives zero prospective source
status to that post-reclassification remainder.  Finally,
\cref{thm-prospective-selector-accountability,cor-first-hit-accounting-not-source}
assign neutral contact to the baseline, every nonbaseline term to a represented
pre-hit selector, and the completed first-hit record to accounting.  These
operations cover every represented mechanism in the saturated closure.

\end{proof}

\begin{proposition}[Distinction between provenance and degree]\label{prop-provenance-is-not-degree}

The \PartII decomposition assigns each target-relevant edge effect a
provenance and a transport type. Budget admissibility at degree \(d\)
additionally requires a valid saturated
representation of degree or effective cost at most \(d\) in the \PartI
filtration or in a lifted observable configuration generated from it.

Consequently, \(\mathcal F_I\otimes\mathcal F_I\)-measurability records
only extensional observability. Exposed provenance and placement in a
budget-admissible cost band require a uniform represented derivation.
Valid Lift may lower, preserve, or raise
the raw coordinate degree of a target-relevant distinction; \PartIII
charges the distinction to the least saturated representation budget at
which it is represented. Below that budget the effect is not an exposed
channel for the filtered problem. For a declared budget structure, a
distinction requiring recovery, representation, or application above
\(b(n)\) is residual from the budgeted observable presentation,
presentation-inaccessible, oracle-supplied, or structure for an enriched presentation. In the
polynomial instance, this is the usual superpolynomial-recovery
statement.

\end{proposition}

\begin{proof}\phantomsection\label{proof-provenance-is-not-degree}

The information-provenance theorem is a finite represented decomposition:
public data, uniformly lifted or generated data, and residual transport
exhaust the possible sources of edge information. Degree is a separate
filtration on represented observables and generated interfaces. Since a
finite measurable function may distinguish many atoms, ambient
measurability gives neither provenance nor a dimension bound.
Conversely, a generated quotient or lifted interface
can give a compact representation of a high raw-degree coordinate. Thus
the channel assignment is exact at the level of source, while budget
admissibility is determined later by the minimal saturated
representation cost of the assigned source.

\end{proof}

\subsection{Post-Exhaustion Band: Irreducibility, Enumeration Complexity, and
Composition Closure}\label{subsec-post-exhaustion-residual-band}
\label{post-exhaustion-band-irreducibility-unscannability-and-composition-closure}

The residual band records represented operator effects remaining after all declared structural projections.
\begin{definition}[Post-exhaustion residual band]\label{def-post-exhaustion-residual-band}

Let \((I_n)\) be a task family equipped with the \PartI observable
filtration and the \PartIII saturated public cost budget \(B\). For a
budget-admissible generator \(L_n\), first remove all components with
exposed geometry, exposed structural gradients, represented quotients,
valid lifted interfaces, and boundary readouts. Also remove the exact
quotient part of every partially lumpable quotient and send its defect
back through the same channel decomposition. The remaining component is
\(L_n^{\Res}\).

The residual cost band above \(B\), denoted \(\mathcal R_{n,>B}(L)\), is
the part of \(L_n^{\Res}\) whose target-aligned structural factor or terminal
readout has no saturated admissible representation of public cost at most
\(B\): it does not factor through a base observable band, through a represented
quotient, or through a valid lifted coordinate whose charged saturated cost is
at most \(B\). The raw transition implementation may be represented; what is
absent below \(B\) is an admissible target-aligned structural representative.
This band is defined after accounting
for quotient, lifted-state, and boundary components. Membership is
determined by saturated representation cost rather than measurability
alone.

\end{definition}

\begin{lemma}[Properties of the post-exhaustion residual band]\label{thm-irreducible-band-closure}

Fix a budget-compatible resource envelope and a budgeted task family.
After the five exposed structural channels have been removed,
\(\mathcal R_{n,>B}(L)\) has the following three properties.

\textbf{Irreducibility.} It admits no family-uniform budget-admissible
observable reduction, including one using presentation-admissible oracle
evaluators, that produces target alignment at the declared threshold. Any such reduction is a represented quotient, a
structural potential, a valid lifted coordinate, a finite-horizon
trajectory statistic generated by budget-admissible computation,
or a boundary readout, and is therefore reclassified into
\(\Abs\), \(\Caus\), \(\Lift\), or
\(\Bdry\) before the residual band is formed. A first-hitting
partition or stopped-future value of the residual operator enters this
classification precisely when it has a budget-admissible observable
representative.

\textbf{Direct-enumeration lower bound.} Let \(B_{n,d}\) be an
independent degree-\(d\) coordinate band of dimension \(N_{n,d}\). A
direct coordinate scan exhausts \(B_{n,d}\) only if the inspected
coordinates span \(B_{n,d}\); hence it requires at least \(N_{n,d}\)
coordinate tests. In the Boolean Walsh filtration,
\(N_{n,d}=\binom{n}{d}\) for the exact degree-\(d\) band. If
\(\binom{n}{d}\) is superpolynomial, no polynomial-cardinality direct
scan exhausts that band.

\textbf{Closure under valid composition.} A finite composition of
classified channel components and valid lifts either exposes a quotient,
potential, lifted interface, or terminal value function and is charged
to one of the five structural channels, or it leaves residual operator
movement---nonlocal transport or directed cycles---in
\(\Res\). Target information outside the budgeted
saturated observable closure is presentation-inaccessible,
oracle-supplied, over budget, or part of an enriched presentation.

\end{lemma}

\begin{proof}\phantomsection\label{proof-irreducible-band-closure}

For irreducibility, suppose a residual subcomponent above a cost level
\(B\) has a family-uniform budget-admissible observable reduction
\(r_n\), after expanding every presentation-admissible oracle call by its
represented evaluator, with target alignment at the declared threshold and
with saturated representation cost at most \(B\). If \(r_n\) is a map
\(X_n\to Q_n\) whose fibers and quotient operations are generated and
manipulated within the declared budget, then it is a represented
observable-quotient candidate once its fibers are generated by the public
rule. Its exact intertwined projection has algebraic \(\Abs\) provenance.
It contributes to the \(\Abs\) structural source profile only through the
qualification conditions of
\cref{def-paper1-abs-usefulness-gate,def-abs-extraction}. If the quotient is only
partial, the identity \(L_XU=UL_Q+E_q\) separates the exact
\(\Abs\) part from a defect that is decomposed again. If the
rule produces a potential, it lies in the budgeted structural-gradient
space and is \(\Caus\). If it exposes memory, tables, products,
proof objects, interfaces, or finite-horizon trajectory statistics as
represented output or retained public records, those readouts are part
of the transcript observable algebra. Internal work state remains
charged at peak native capacity and is not a free observable. The exposed
readouts are \(\Lift\) only when the valid-lift conditions and
target-fraction preservation hold. If it produces terminal values, it is
\(\Bdry\). The residual operator's own first-hitting partition,
committor, or stopped-future table is included when generated by
budget-admissible finite-horizon computation or represented by another
admissible observable. These are precisely the cases in the
information-provenance theorem and the family-uniform
residual-reclassification proposition. Thus a
family-uniform aligned reduction of cost at most \(d\) cannot remain in
\(\mathcal R_{n,>d}(L)\). If the least saturated representation cost is
larger than \(d\), the effect is not a channel below \(d\); at that
filtered level it remains in the residual band, is
presentation-inaccessible, oracle-supplied, or enriched-presentation
data, or is charged only when the budget is raised
to its represented cost.

For direct enumeration, let \(W\le B_{n,d}\) be the span of all
coordinates inspected by the scan. Each coordinate test adds at most one
dimension, so \(\dim W\le |S_n|\), where \(S_n\) is the inspected set.
If \(|S_n|<N_{n,d}=\dim B_{n,d}\), then \(W\ne B_{n,d}\), and there
exists a nonzero coordinate direction in \(B_{n,d}\cap W^\perp\) unseen
by the scan. Therefore exhaustive scanning requires
\(|S_n|\ge N_{n,d}\). For Boolean Walsh characters the exact
degree-\(d\) band consists of one independent character for each
\(d\)-subset of variables, so \(N_{n,d}=\binom{n}{d}\). Superpolynomial
\(\binom{n}{d}\) therefore exceeds every polynomial-cardinality scan.
This establishes a dimension lower bound for direct enumeration.

For composition closure, use the finite insertion identity and interface
expansion proved in the
\hyperref[composition-closure-lift-creates-nothing]{composition-closure
lemma}. Those identities decompose every composed edge into the same
typed pieces: local conductance, directed structural gradient, quotient
or valid-lift jump, boundary readout, and residual transport. If the
composition creates a family-uniform aligned coordinate, the first
paragraph reclassifies it into one of the five structural channels. If
it creates no such coordinate, it remains residual. Hence valid
composition remains within the five structural channels and preserves
irreducibility of the post-exhaustion residual band.

\end{proof}

\section{Represented Quotient Geometry and Geom/Abs Sources}
\label{sec-represented-quotient-geometry}

The residual-band analysis leaves the represented fiber algebra and the
conditions for a geometric source certificate.

\begin{lemma}[Fiber algebra of a represented quotient]
\label{thm-abs-fiber-algebra}

Let \(q:X_I\to Q\) be a represented quotient, and let \(Q_q\subseteq Q\)
be its set of nonempty fibers. Its quotient-observable algebra is

\[
\mathscr A(q)
:=
\{a\circ q:a:Q_q\to\mathbb R\}
\subseteq\mathbb R^{X_I},
\]

with pointwise addition and multiplication. If \(k=|Q_q|\), then

\[
                       \mathscr A(q)\cong\mathbb R^k
\]

as a unital commutative real algebra. Its primitive idempotents are the
indicators of the nonempty quotient fibers. In particular,
\(\mathscr A(q)\) is reduced, and both its nilradical and its Jacobson
radical are zero.

Consequently, the abstract isomorphism type of the pointwise quotient
algebra records the number of quotient cells but not the represented
derivation that produced them. Construction provenance is carried by
the quotient's uniform derivation record, represented fiber operations,
and saturated cost rather than by an Artinian decomposition of
\(\mathscr A(q)\).

\end{lemma}

\begin{proof}\phantomsection\label{proof-abs-fiber-algebra}

Enumerate the nonempty fibers as \(C_i=q^{-1}(u_i)\),
\(1\le i\le k\). The map

\[
\mathscr A(q)\longrightarrow\mathbb R^k,
\qquad
a\circ q\longmapsto
\bigl(a(u_1),\ldots,a(u_k)\bigr),
\]

is a well-defined unital algebra isomorphism. Under this map the fiber
indicators \(\mathbf 1_{C_i}\) are the standard coordinate idempotents.
If \(f^r=0\) for some \(r\ge1\), every coordinate of the image of \(f\)
has zero \(r\)-th power in \(\mathbb R\), and hence \(f=0\). Thus the
algebra is reduced and its nilradical is zero. A finite product of
fields also has zero Jacobson radical. The remaining assertions follow.
The displayed isomorphism depends only on the fiber partition; the
derivation and its cost are additional represented data.

\end{proof}

\begin{lemma}[Characterization of transfer-class geometric certificates]\label{thm-valid-geom-source-characterization}

A pointwise observable \(\phi:X_I\to\mathbb R\), together with its
declared conductance and comparison data, is a transfer-class
\(\Geom\) certificate precisely when the following five
conditions hold:

\begin{enumerate}

\item

\textbf{Represented provenance.} \(\phi\) is generated in the budgeted
observable or lifted closure at the stated budget.

\item

\textbf{Exposed locality.} The differences \(\phi(y)-\phi(x)\) used by
the source are supported on the exposed comparability relation \(E_I\).

\item

\textbf{Symmetric conductance.} The local conductance
\(c_E(x,y)=c_E(y,x)\) is exposed, so \(\phi\) defines the symmetric
Dirichlet form \(\sum_{(x,y)\in E_I}c_E(x,y)(\phi(y)-\phi(x))^2\).

\item

\textbf{Positive killed gap.} The killed local form has a represented
positive Poincare/Cheeger gap toward the terminal sector at the stated
budget \citep{LawlerSokal1988,LevinPeres2017}.

\item

\textbf{Target reach.} The level-set algebra of \(\phi\) has
non-negligible projection on the centered target contrast.

\end{enumerate}

Conditions (1) and (2) specify provenance and locality. Condition (3)
identifies the symmetric local source and separates \(\Geom\)
from \(\Caus\). Conditions (4) and (5) impose the rate and
alignment requirements.

\end{lemma}

\begin{proof}\phantomsection\label{proof-valid-geom-source-characterization}

The symmetric local part of a positive generator is represented by an
exposed conductance on \(E_I\) and pointwise potential differences. A
source outside the represented closure lacks budget-admissible
provenance; one using unexposed comparisons is outside the local
\(\Geom\) channel. Nonsymmetric orientation belongs to
\(\Caus\). Without a positive killed-gap comparison the record
does not supply the required transfer-class rate certificate; without
target reach it does not supply the required alignment certificate.
Conversely, conditions (1)--(3) construct the represented symmetric
local channel, while conditions (4) and (5) provide exactly the rate
and alignment margins required by the definition of a transfer-class
certificate. This proves the characterization of certificate records.
It does not assert that every symmetric local generator has a positive
killed gap or non-negligible target reach.

\end{proof}

\subsection{The five Geom/Abs source configurations}
\label{subsec-geom-abs-source-configurations}

The final phase combines exact quotient structure with ambient or quotient
geometry through the compensated certificate.

\begin{definition}[Geometrically compensated quotient]\label{def-coupled-abs-geom-source}

Let \(q:\Omega\to Q\) be a represented terminal-compatible quotient.
Write

\[
U:H_Q\longrightarrow H_\Omega,
\qquad
P:H_\Omega\longrightarrow H_Q,
\qquad
PU=I_Q,
\]

and let \(L_Q=PL_\Omega U\). The lumpability defect is the operator

\[
E_q=L_\Omega U-UL_Q:H_Q\longrightarrow H_\Omega.
\]

A geometrically compensated quotient consists of \(q\), represented
Dirichlet data for the quotient modes and the within-fiber modes, and a
charged comparison certificate showing that the coupling \(E_q\) is
small relative to both forms. In the reversible case the form
comparison is controlled by the Schur-complement parameter
\(\theta_\lambda(q)<1\) of
\Cref{def-reversible-compensated-defect}; target visibility
additionally uses the target-functional margin
\(\vartheta_\lambda^T(q)<1\) of
\Cref{cor-target-functional-reversible-transfer}. In the
general case it uses the target-functional normalized resolvent loss
\(\Delta_\lambda^T(q)<1\) of
\Cref{def-normalized-resolvent-defect}. The reciprocal of the
resulting positive stability factor \(S_q(\lambda)\) must lie in the
declared transfer class.

When the quotient is used for directed target progress, the same
certificate also contains a represented potential and its
\(\Caus\) current. Their quantitative alignment and drift margin
are those of \Cref{def-quantitative-caus-alignment} and
\Cref{thm-caus-drift-survives-defect}. Thus geometric
compensation controls an operator between quotient and ambient spaces;
it does not identify that operator with a single edge gradient.

\end{definition}

The source classification enumerates the represented configurations available to the two structural channels.
\begin{theorem}[Classification of represented Geom/Abs source configurations]\label{thm-complete-generic-source-theory}

For any presented finite task family, every dynamically qualified
source record in the declared channel calculus whose reusable
target-bearing data are represented by \(\Geom\),
\(\Abs\), or their composition has one of the following five
configurations of the compensated quotient--geometry form:

\begin{enumerate}
\item represented ambient \(\Geom\), oriented by an
authority-compatible \(\Caus\) current and represented
potential when it is used for directed progress;
\item exact represented \(\Abs\), corresponding to a null geometry
slot;
\item represented \(\Geom\) on an exact \(\Abs\) quotient,
with authority-compatible \(\Caus\) providing target
orientation;
\item represented \(\Geom\) on a quasi-lumpable \(\Abs\)
quotient whose defect satisfies a charged reversible
Schur-complement compensation bound, again with \(\Caus\)
providing target orientation under the declared authority envelope; and
\item represented \(\Geom\) on a quasi-lumpable \(\Abs\)
quotient whose defect satisfies a charged general resolvent
compensation bound.
\end{enumerate}

The last three cases combine the two source coordinates; they do not
define an additional coordinate. Within this calculus the remaining
channels are derived uses:
\(\Caus\) supplies only the orientation authorized by the
selected imposed/control generator, \(\Lift\)
reorganizes represented coordinates, and \(\Bdry\) supplies a
terminal readout. A quotient whose defect has no transfer-class
compensation certificate supplies no quotient-to-ambient comparison at
that budget. Its defect may instead be resolved by refining the
quotient, representing a different geometry, or retaining the
unclassified operator part in \(\Res\).  The theorem's domain is the
represented, budget-admissible channel-certificate class of
\cref{def-coupled-abs-geom-source,thm-valid-geom-source-characterization};
affordability is the charged derivation condition of that class.

\end{theorem}

\begin{proof}\phantomsection\label{proof-complete-generic-source-theory}

Component exhaustiveness leaves local symmetric, local oriented,
nonlocal quotient, lifted, boundary, and residual pieces. The local
symmetric piece is \(\Geom\), characterized by
\Cref{thm-valid-geom-source-characterization}. The represented
quotient piece is \(\Abs\). Its pointwise fiber algebra is
described by \Cref{thm-abs-fiber-algebra}, while its construction
provenance is the resource-filtered derivation established in
\Cref{thm-affordable-exact-quotient-normal-form}. If no
nonidentity quotient is used, identity support gives the first case. A
nonidentity exact quotient with a null geometry slot gives the second
case. If a represented geometry is constructed on a quotient, the intertwining defect
\(E_q=L_\Omega U-UL_Q\) separates the exact and quasi-lumpable cases.
Exact intertwining gives the third case. For nonzero defect, the
comparison is valid precisely when the defect is controlled relative to
the quotient and within-fiber forms, as quantified in
\Cref{quotient-defect-and-geom-form-comparison}. The reversible
Schur-complement and general resolvent records give the fourth and fifth
cases. The transformed-gradient decomposition then places the
target-directed part in \(\Caus\), and
\Cref{thm-caus-authority-no-substitution} restricts it to an
authorized current of the selected generator. Valid Lift and Bdry remain
charged uses of these sources, while every unrepresented remainder lies
in \(\Res\). These alternatives exhaust the dynamically
qualified represented Geom/Abs source records specified in the
statement.

\end{proof}

\begin{theorem}[Exhaustive geometric normal form]
\label{thm-exhaustive-geom-normal-form}

Fix the symmetric local relation \(E_I\), reference weight \(m\), and
conductance \(c_E\) of \cref{def-fixed-space-split}.  Let

\[
\mathcal E_I^{\mathrm{un}}
=
\bigl\{\{x,y\}:x\ne y,\ (x,y)\in E_I\bigr\}
\]

be the associated unordered edge set.  For \(e=\{x,y\}\), define the
weighted edge Laplacian

\[
\bigl(\Delta_e^{(m)}f\bigr)(z)
=
\begin{cases}
m(x)^{-1}\bigl(f(y)-f(x)\bigr),&z=x,\\
m(y)^{-1}\bigl(f(x)-f(y)\bigr),&z=y,\\
0,&z\notin\{x,y\}.
\end{cases}
\]

Then every algebraic geometric component has the unique edge-conductance
normal form

\begin{equation}
\label{eq-exhaustive-geom-edge-normal-form}
L^{\Geom}
=
\sum_{e\in\mathcal E_I^{\mathrm{un}}}
c_e\Delta_e^{(m)},
\qquad
c_e=c_E(x,y)=c_E(y,x)\ge0.
\end{equation}

A target-bearing transfer-class certificate for this component satisfies the
five necessary and sufficient conditions of
\cref{thm-valid-geom-source-characterization}.  After composing its
represented quotient maps and retaining the complete derivation cost, its
active geometry has one or more of the following exhaustive provenance
configurations:

\begin{enumerate}[label=\textup{(\roman*)}]
\item ambient represented geometry, including ambient geometry on a complete
valid lifted carrier, with identity quotient;
\item represented geometry on an exact \(\Abs\) quotient;
\item represented geometry on a quasi-lumpable \(\Abs\) quotient with a
charged reversible Schur-complement compensation record;
\item represented geometry on a quasi-lumpable \(\Abs\) quotient with a
charged general resolvent-compensation record.
\end{enumerate}

The last two provenance families may overlap when both represented
compensation records are available.  Selecting and charging one record gives
the disjoint certificate tag used by
\cref{thm-affordable-geom-abs-normal-form}.  Authorized \(\Caus\) orientation,
valid \(\Lift\), and \(\Bdry\) readout retain this geometric source and its
carrier cost; they introduce no additional geometric configuration.  At a
budget \(B\), a geometric record is exposed precisely when its relation,
conductance, potential or comparison data, target-reach data, projection
procedure, and output are represented at total saturated cost at most \(B\).
Otherwise its algebraic Geom tag is retained in the budget-filtered residual
ledger of \cref{thm-exhaustive-residual-normal-form}.

\end{theorem}

\begin{proof}\phantomsection\label{proof-exhaustive-geom-normal-form}

The symmetric local block in \cref{thm-fixed-space-split} is

\[
(L^{\mathsf{Sym}}f)(x)
=
\frac1{m(x)}
\sum_{y:(x,y)\in E_I}
c_E(x,y)\bigl(f(y)-f(x)\bigr).
\]

Grouping the two orientations of each unordered pair gives
\eqref{eq-exhaustive-geom-edge-normal-form}.  Conversely, the off-diagonal
matrix entries recover
\(c_e=m(x)(L^{\Geom})_{xy}=m(y)(L^{\Geom})_{yx}\), so the coefficients and
the edge sum are unique.  By
\cref{prop-no-third-local-primitive,thm-channel-exhaustiveness}, every
symmetric local operator term is in this block and no further local
geometric primitive remains.

\Cref{thm-valid-geom-source-characterization} gives the necessary and
sufficient represented-provenance, locality, conductance, killed-gap, and
target-reach conditions.  Apply
\cref{thm-complete-generic-source-theory} to any record satisfying them.  Its
five source configurations consist of ambient Geom, null-geometry exact Abs,
Geom on exact Abs, and the reversible and general compensated
quasi-lumpable cases.  Removing the null-geometry configuration, which has
zero active Geom slot, leaves exactly the four configurations in the
statement.  The derived-channel proposition
\cref{prop-derived-channels-create-no-structure} and the exact Lift identity
\cref{lem-exact-lift-identity} preserve the source assignment under
orientation, lifting, and terminal readout.  The charged projection
convention proves the final budget clause.

\end{proof}

\Cref{fig-geom-normal-form} displays the four active-geometry configurations
of \cref{thm-exhaustive-geom-normal-form}.

\begin{figure}[tbp]
\centering
\begin{tikzpicture}[fw diagram]
  \node[fw geom,minimum width=4.0cm] (root) at (0,2.8)
    {qualified \(\Geom\)\\\(\sum_e c_e\Delta_e^{(m)}\)};
  \coordinate (split) at (0,2.05);
  \node[fw geom,minimum width=2.5cm] (ambient) at (-4.5,0.65)
    {ambient\\\(q=\mathrm{id}\)};
  \node[fw source,minimum width=2.5cm] (exact) at (-1.5,0.65)
    {exact quotient\\\(E_q=0\)};
  \node[fw use,minimum width=2.7cm] (rev) at (1.7,0.65)
    {quasi quotient\\reversible};
  \node[fw provenance,minimum width=2.7cm] (general) at (4.8,0.65)
    {quasi quotient\\general resolvent};
  \draw[fw arrow] (root) -- (split);
  \foreach \n in {ambient,exact,rev,general}
    \draw[fw arrow] (split) -| (\n);
\end{tikzpicture}
\caption{The four exhaustive, possibly overlapping active-Geom provenance
configurations of \cref{thm-exhaustive-geom-normal-form}.  Every configuration
uses the same unique local edge-conductance normal form.}
\label{fig-geom-normal-form}
\end{figure}

\Cref{fig-geom-abs-configurations} compresses the five
source configurations into their quotient and geometry slots.

\begin{figure}[tbp]
\centering
\begin{tikzpicture}[fw diagram]
  \node[fw box,minimum width=3.7cm] (record) at (0,1.65)
    {represented \(\Geom/\Abs\) record};
  \coordinate (split) at (0,0.95);
  \node[fw geom,minimum width=2.35cm] (g) at (-5.2,0)
    {\(q=\mathrm{id}\)\\ambient \(\Geom\)};
  \node[fw source,minimum width=2.35cm] (a) at (-2.6,0)
    {\(q\ne\mathrm{id}\)\\exact \(\Abs\)};
  \node[fw geom,minimum width=2.35cm] (ag) at (0,0)
    {\(\Abs\to\Geom\)\\\(E_q=0\)};
  \node[fw use,minimum width=2.35cm] (rev) at (2.6,0)
    {\(\Abs\dashrightarrow\Geom\)\\reversible};
  \node[fw provenance,minimum width=2.35cm] (res) at (5.2,0)
    {\(\Abs\dashrightarrow\Geom\)\\resolvent};
  \draw[fw arrow] (record) -- (split);
  \foreach \n in {g,a,ag,rev,res}
    \draw[fw arrow] (split) -| (\n);
\end{tikzpicture}
\caption{The compensated quotient--geometry certificate has five
alternative configurations, classified by
\cref{thm-complete-generic-source-theory}.}
\label{fig-geom-abs-configurations}
\end{figure}

Consequently, the five exposed structural channels together with
\(\Res\) provide the complete source decomposition used in
the subsequent quantitative analysis.

\section{Stability Under Composition and Valid
Lifting}\label{sec-composition-valid-lifting}
\label{composition-closure-lift-creates-nothing}

Algorithms may augment the state space with stacks, products, proof
trees, dynamic-programming tables, learned records, or subproblem
interfaces. When this augmentation defines a valid lift of the
observable task, the resulting generator remains within the same
infinitesimal decomposition.

\begin{proposition}[Defect calculus for composed quotients]
\label{prop-composed-quotient-defect}

Let

\[
X\xrightarrow{q_1}Q_1\xrightarrow{q_2}Q_2
\]

be represented quotients with lifts \(U_1,U_2\), projections
\(P_1,P_2\), and \(P_iU_i=I\). Starting from \(L_0=L_X\), define

\[
L_1=P_1L_0U_1,
\qquad
L_2=P_2L_1U_2,
\]

and

\[
E_1=L_0U_1-U_1L_1,
\qquad
E_2=L_1U_2-U_2L_2.
\]

For \(\sharp\in\{\mathsf s,\mathsf a\}\), let
\(L_i^\sharp\) be obtained by compressing the corresponding component
at each stage, and set

\[
E_1^\sharp=L_0^\sharp U_1-U_1L_1^\sharp,
\qquad
E_2^\sharp=L_1^\sharp U_2-U_2L_2^\sharp.
\]

The composite lift is \(U_{12}=U_1U_2\), and its defect satisfies the
exact identity

\begin{equation}
\label{eq-composed-quotient-defect}
E_{q_2\circ q_1}=E_1U_2+U_1E_2.
\end{equation}

Componentwise,

\begin{equation}
\label{eq-composed-causal-defect}
E_{q_2\circ q_1}^{\sharp}
=E_1^\sharp U_2+U_1E_2^\sharp,
\qquad
\sharp\in\{\mathsf s,\mathsf a\}.
\end{equation}

Consequently,

\[
\|E_{q_2\circ q_1}\|
\le
\|E_1\|\,\|U_2\|+\|U_1\|\,\|E_2\|.
\]

For conditional-expectation projections the lifts are isometries, and
the right-hand side is \(\|E_1\|+\|E_2\|\). Exact lumpability therefore
composes.  The same estimate holds separately for the symmetric and
causal defects.  The saturated cost of the composite includes both quotient
derivations, both represented-fiber interfaces, and the composition
constructor.

\end{proposition}

\begin{proof}\phantomsection\label{proof-composed-quotient-defect}

Insert the two defect identities:

\[
\begin{aligned}
L_0U_1U_2-U_1U_2L_2
&=(L_0U_1-U_1L_1)U_2
  +U_1(L_1U_2-U_2L_2)\\
&=E_1U_2+U_1E_2.
\end{aligned}
\]

The norm estimate follows from the triangle inequality and
submultiplicativity. The cost statement is the join-and-composition
clause of the budgeted derivability closure.  Replacing each \(L_i\) by
\(L_i^{\mathsf s}\) or \(L_i^{\mathsf a}\) gives
\eqref{eq-composed-causal-defect} by the same calculation.

\end{proof}

The lift closure collects the represented observables generated by valid lift interfaces.
\begin{definition}[Lift closure]\label{def-lift-closure}

Let \(\mathcal G_0\) be the class of budget-admissible fixed-space
generators decomposed into the five exposed structural channels plus
residual accounting. The valid lift closure is the least class
containing \(\mathcal G_0\) and closed under finite direct sums, tensor
products, product semigroups, bounded iteration, split/base/merge
decompositions, and finite interface contraction along observable
interfaces, subject at every stage to the valid-lift conditions
\(\pi^{-1}\mathcal F_{\mathrm{obs}}\subseteq\mathcal G_{\mathrm{obs}}\),
\(\pi_*\Lambda\mu=\mu\), and target-fraction preservation.

\end{definition}

\begin{lemma}[Exact Lift identity]\label{lem-exact-lift-identity}

Let \(\pi:\widetilde X\to X\) be a valid lift with lifted initialization
\(\Lambda\mu\).  Set the lifted measure \(\widetilde\mu:=\Lambda\mu\), and
let the target be
\(T\subseteq X\). Then

\[
\widetilde\mu(\pi^{-1}T)=\mu(T),
\qquad
\int_{\widetilde X}h\circ\pi\,d\widetilde\mu
=
\int_Xh\,d\mu
\]

for every base-measurable target-relevant observable \(h\). These
identities compose exactly under finite products, staged lifts, bounded
iteration, and finite interface contractions that satisfy the valid-lift
conditions. Hence Lift has zero target-fraction defect: it reorganizes
represented source mass but does not create a new target source.

\end{lemma}

\begin{proof}\phantomsection\label{proof-exact-lift-identity}

The identity \(\pi_*\widetilde\mu=\mu\) is the valid-lift projection
condition. Applying the definition of pushforward to \(1_T\) gives
\(\widetilde\mu(\pi^{-1}T)=\mu(T)\). Applying it to a bounded base
observable \(h\) gives the average identity. For a composition of valid
lifts, pushforwards compose; induction gives the same identities for the
composed projection. Interface contraction is valid only when the
contracted interface has an observable positive representative
preserving the same projection identities. Therefore finite lifted
composition preserves target fractions exactly and carries provenance
back to the represented base source.

\end{proof}

\begin{lemma}[Fiber disintegration of a finite valid lift]
\label{lem-valid-lift-fiber-disintegration}
Let \((\Omega,\mathcal G_{\mathrm{obs}},\pi,\Lambda)\) be a finite valid lift
of \(X_I\) in the sense of \cref{def-valid-lift}.  For each \(x\in X_I\),
define the positive measure

\[
\kappa_x=\Lambda\delta_x.
\]

Then \(\kappa_x\) is a probability measure supported on the fiber
\(\pi^{-1}(x)\), and every positive frontier has the unique fiber-kernel
normal form

\begin{equation}
\label{eq-exhaustive-lift-fiber-disintegration}
\Lambda\mu
=
\sum_{x\in X_I}\mu(x)\kappa_x,
\qquad
\operatorname{supp}\kappa_x\subseteq\pi^{-1}(x).
\end{equation}

Let \(U:\mathbb R^{X_I}\to\mathbb R^\Omega\) be pullback and define

\[
(P_\kappa h)(x)
=
\sum_{\omega\in\pi^{-1}(x)}
\kappa_x(\omega)h(\omega).
\]

Then \(P_\kappa U=I_{X_I}\).

\end{lemma}

\begin{proof}\phantomsection\label{proof-valid-lift-fiber-disintegration}

The valid-lift identity gives
\(\pi_*\Lambda\delta_x=\delta_x\). Positivity implies that
\(\Lambda\delta_x\) vanishes outside \(\pi^{-1}(x)\), and preservation of
total mass gives \(\kappa_x(\Omega)=1\). Linearity of \(\Lambda\) on the
point-mass basis gives
\eqref{eq-exhaustive-lift-fiber-disintegration}, and evaluation on each
\(\delta_x\) proves uniqueness. Finally,
\[
(P_\kappa Ug)(x)
=\sum_{\omega\in\pi^{-1}(x)}\kappa_x(\omega)g(\pi(\omega))
=g(x),
\]
so \(P_\kappa U=I_{X_I}\).

\end{proof}

\begin{theorem}[Exhaustive finite-Lift normal form]
\label{thm-exhaustive-lift-normal-form}

Let \((\Omega,\mathcal G_{\mathrm{obs}},\pi,\Lambda)\) be a finite valid lift
of \(X_I\), and let \(\kappa_x\), \(U\), and \(P_\kappa\) be the unique
fiber data of \cref{lem-valid-lift-fiber-disintegration}.

For every finite lifted generator
\(L_\Omega\), put

\[
L_X=P_\kappa L_\Omega U,
\qquad
E_\pi=L_\Omega U-UL_X,
\qquad
\Pi_\pi=UP_\kappa.
\]

The lifted operator has the exact full-domain decomposition

\begin{equation}
\label{eq-exhaustive-lift-operator-normal-form}
L_\Omega
=
UL_XP_\kappa
+
E_\pi P_\kappa
+
L_\Omega(I-\Pi_\pi).
\end{equation}

Apply the ordered channel projections to all three terms.  The represented
quotient-descending projection of \(UL_XP_\kappa\) is \(\Abs\).  After the
declared \(\Geom\), \(\Caus\), and \(\Abs\) projections, the represented
fiber/interface projection

\begin{equation}
\label{eq-exhaustive-lift-channel-projection}
L^{\Lift}
=
\Pi_{\mathcal S_{\mathrm{Lift}}\cap
\mathcal S_{\mathrm{Abs}}^\perp}L_{\mathrm{nl}}
\end{equation}

is the complete \(\Lift\) component.  Every part of
\(E_\pi P_\kappa+L_\Omega(I-\Pi_\pi)\) outside this subspace is assigned by
the same \(\Geom,\Caus,\Abs,\Bdry,\Res\) decomposition.  Consequently every
finite structural Lift has exactly the following normal-form data:

\begin{enumerate}[label=\textup{(\roman*)}]
\item the unique family of target-neutral fiber laws
\((\kappa_x)_{x\in X_I}\);
\item the represented lifted interface and its orthogonalized
\(\Lift\)-channel projection \eqref{eq-exhaustive-lift-channel-projection};
\item the quotient defect and within-fiber complement of
\eqref{eq-exhaustive-lift-operator-normal-form}, recursively assigned by the
existing channel decomposition.
\end{enumerate}

Finite products, staged lifts, bounded iteration, split/base/merge
constructions, and represented interface contractions produce new fiber
kernels of the same form and preserve the exact target-fraction identity.
At budget \(B\), the fiber laws, interface, projection, update rule, and
output are included in the charged record.  A structurally typed Lift whose
complete record is unavailable at \(B\) is retained by the filtered-residual
convention of \cref{thm-exhaustive-residual-normal-form}.

\end{theorem}

\begin{proof}\phantomsection\label{proof-exhaustive-lift-normal-form}

By \cref{lem-valid-lift-fiber-disintegration},
\(P_\kappa U=I_{X_I}\), so \(\Pi_\pi=UP_\kappa\) is idempotent. Therefore

\[
\begin{aligned}
L_\Omega
&=L_\Omega\Pi_\pi+L_\Omega(I-\Pi_\pi)\\
&=L_\Omega UP_\kappa+L_\Omega(I-\Pi_\pi)\\
&=UL_XP_\kappa+E_\pi P_\kappa+L_\Omega(I-\Pi_\pi),
\end{aligned}
\]

which is \eqref{eq-exhaustive-lift-operator-normal-form}.  This is
\cref{prop-partial-lumpability-no-escape} for the represented projection
\(\pi\).  The structured nonlocal split of \cref{def-abs-lift-split} first
projects onto \(\mathcal S_{\mathrm{Abs}}\) and then onto
\(\mathcal S_{\mathrm{Lift}}\cap\mathcal S_{\mathrm{Abs}}^\perp\), proving
\eqref{eq-exhaustive-lift-channel-projection}.  Component exhaustiveness
assigns the orthogonal remainder, the defect, and the within-fiber complement
without adding another Lift class.

For every named target \(T\subseteq X_I\), fiber support gives

\[
\kappa_x(\pi^{-1}T)=\mathbf1_T(x),
\qquad
(\Lambda\mu)(\pi^{-1}T)=\mu(T).
\]

This is the exact identity of \cref{lem-exact-lift-identity}.  Pushforwards
and fiber kernels compose under the constructors in
\cref{def-lift-closure}, so the same normal form and target identity persist
under finite valid composition.  The structural Lift gate
\cref{def-structural-lift-gate} supplies the represented-generation,
reference-law, and cost conditions, and the charged projection convention
gives the budget statement.

\end{proof}

\Cref{fig-lift-normal-form} displays the fiber and operator layers of
\cref{thm-exhaustive-lift-normal-form}.

\begin{figure}[tbp]
\centering
\begin{tikzpicture}[fw diagram]
  \node[fw use,minimum width=3.5cm] (root) at (0,3.0)
    {valid finite \(\Lift\)};
  \node[fw source,minimum width=3.4cm] (fiber) at (-3.5,1.65)
    {fiber laws\\\(\kappa_x\subseteq\pi^{-1}(x)\)};
  \node[fw box,minimum width=5.0cm] (operator) at (3.0,1.65)
    {\(UL_XP_\kappa+E_\pi P_\kappa+L_\Omega(I-UP_\kappa)\)};
  \node[fw source,text width=2.4cm] (abs) at (-0.2,0.15)
    {descending\\\(\Abs\)};
  \node[fw use,text width=3.0cm] (lift) at (3.1,0.15)
    {represented interface\\\(\Lift\)};
  \node[fw residual,text width=3.0cm] (defect) at (6.6,0.15)
    {defect/complement\\reproject};
  \draw[fw arrow] (root) -- (fiber);
  \draw[fw arrow] (root) -- (operator);
  \draw[fw arrow] (operator) -- (abs);
  \draw[fw arrow] (operator) -- (lift);
  \draw[fw arrow] (operator) -- (defect);
\end{tikzpicture}
\caption{The exhaustive finite-Lift normal form of
\cref{thm-exhaustive-lift-normal-form}.  A valid initialization is a unique
family of fiber laws; the lifted operator separates into its descending
quotient block, represented Lift interface, and recursively decomposed
defect/complement.}
\label{fig-lift-normal-form}
\end{figure}

\begin{lemma}[Closure of the channel decomposition under composition]\label{lem-gaplessness-composition}

If the observable target-directed structure of each factor lies in the
five exposed structural channels and every lift is valid, then the same
holds for the composed generator. Nonlocal movement without observable
structural provenance remains residual unless the composition generates
an exposed quotient or valid-lift coordinate in the enlarged observable
space.

\end{lemma}

\begin{proof}\phantomsection\label{proof-gaplessness-composition}

For product spaces, infinitesimal composition has the insertion form

\[
L_{X_1\otimes\cdots\otimes X_k}
=
\sum_{i=1}^k
I_{X_1}\otimes\cdots\otimes L_{X_i}\otimes\cdots\otimes I_{X_k}.
\]

Each summand is the same generator type acting on one factor and
identity on the others. Local conductance remains local conductance on
its factor, directed current remains directed current, quotient
structure remains quotient structure on a represented macro-coordinate,
valid lifted organization is recorded as \(\Lift\), boundary
readout remains boundary readout, and non-structural nonlocal movement
remains residual. The projection maps of valid lifts compose, so the
identities \(\pi_*\Lambda\mu=\mu\) and
\((\Lambda\mu)(\pi^{-1}T_I)/(\Lambda\mu)(\Omega)=\mu(T_I)/\mu(X_I)\) are
preserved through finite products and finite staged lifts.

For a finite interface \(X=A\sqcup I\), eliminating the interface
through a declared positive interface inverse produces

\[
L_{\mathrm{eff}}(\lambda)
=
L_{AA}+L_{AI}(\lambda I-L_{II})^{-1}L_{IA}.
\]

The positive inverse and the resulting Schur-complement matrix are
treated as one represented interface-contraction constructor, with the
inverse, precision, evaluation, and output table fully charged. A finite
interface does not make its Neumann/path expansion finite when cycles
are present. After the contraction is formed, apply the finite channel
decomposition directly to its resulting rate table: local conductance,
directed current, quotient or valid lifted jump, boundary readout, and
residual operator transport are projected in the usual way. An interface
containing an undeclared solver, presentation-inaccessible target
pointer, or target-enriching lift that violates fraction preservation
introduces oracle-supplied or over-budget information. Otherwise
composition remains within the same component decomposition.

\end{proof}

\begin{corollary}[Target-mass preservation under valid lifting]\label{cor-lift-creates-nothing}

A valid Lift composes, exposes, routes, and extracts structure represented
in its factors or observable interfaces while preserving target-relevant
mass fractions under projection. Target enrichment outside those factors
is presentation-inaccessible, oracle-supplied, over budget, or residual
relative to the public presentation.

\end{corollary}

\begin{proof}\phantomsection\label{proof-lift-creates-nothing}

For a valid lift \(\pi:\Omega\to X_I\), the lifted target is
\(T_\Omega=\pi^{-1}(T_I)\) and the initialization lift satisfies
\(\pi_*\Lambda\mu=\mu\). Therefore \((\Lambda\mu)(T_\Omega)=\mu(T_I)\)
and, for every base-measurable observable \(h\), the lifted average of
\(h\circ\pi\) equals the base average of \(h\). These identities persist
through finite compositions of valid lifted interfaces. Generated
coordinates may route represented \(\Geom\) and \(\Abs\)
sources, but target fractions and target-aligned averages remain those
of the pre-lift source. A lift changing the target fraction is invalid
for the original presentation.

\end{proof}

\section{Operational Representation: Computations as Markov
Kernels}\label{sec-operational-computation-embedding}
\label{the-substrate-bridge-algorithms-as-markov-kernels}

\subsection{Finite transcript kernels}
\label{subsec-finite-transcript-kernels}

The operator formulation applies to represented computations through
their Markov representation. Under the resource envelope of Definition
\hyperref[def-resource-envelope]{Budget-compatible resource envelope},
the clocked outcome/terminal transcript carrier of Definition
\hyperref[def-full-configuration-lift]{Full operational transcript representation}
is finite. Each transcript prefix determines the current native state,
and one operational step defines a Markov or sub-Markov kernel on the
finite transcript carrier.

The finite-envelope statement below is a budget-slice theorem, not a
finiteness axiom for computation.  \Cref{cor-all-turing-machines-represented} represents an arbitrary
unbounded or nonhalting DTM as one countable process; its finite kernels
are exact stopped restrictions used for finite-budget accounting.

\begin{lemma}[Operational computation embedding]\label{thm-algorithm-kernel-embedding}

Fix an input instance and a finite resource envelope with horizon \(H\).
Any represented operational computation embeds as a finite Markov or
sub-Markov kernel \(P^{[H]}\) on the clocked transcript carrier
\(\Omega^{[H]}\) of \Cref{def-full-configuration-lift}; its
legal transcript prefixes form an invariant subset from the public
initialization. For any
scalar \(\rho>0\),

\[
L^{[H]}=\rho(P^{[H]}-I)
\]

is a positive continuous-time generator on \(\Omega^{[H]}\).  It is a
represented generator when \(\rho\) is represented and charged; the
canonical choice \(\rho=1\) is always available.  The embedding depends
only on the complete operational transition law. Its execution ledger
charges the accumulated finite transcript and one peak native state; it
does not enumerate or store the intervening native states.

\end{lemma}

\begin{proof}\phantomsection\label{proof-algorithm-kernel-embedding}

The bounded horizon and padded finite record alphabets form a finite
represented transcript carrier. For a legal prefix \(h\), let \(z(h)\)
be the current native state obtained by the initializer and successor
recursion. The transition from \(h\) appends an outcome \(a\) and its
declared terminal/readout record with conditional weight
\(p(a\mid z(h))\). During execution the update is applied to the single
live state \(z(h)\); earlier native states are neither retained nor
exposed. An isolated row evaluation reconstructs \(z(h)\) by the same
prefix recursion and charges its repeated updates; this affects the
representation ledger, not the finite-carrier identity. Extend malformed and unused transcripts by the fixed inert
failure rule from \Cref{def-full-configuration-lift}. This
gives a total Markov or sub-Markov kernel \(P^{[H]}\) without enumerating
the native-state space, the reachable transcript subset, or an
extensional transition table.

The matrix \(L^{[H]}=\rho(P^{[H]}-I)\) has off-diagonal entries
\(L^{[H]}_{\xi\xi'}=\rho P^{[H]}_{\xi\xi'}\ge0\) for \(\xi\ne\xi'\). Its generator-form rates and
killing are

\[
q(\xi,\xi')=\rho P^{[H]}_{\xi\xi'}\quad (\xi\ne\xi'),
\qquad
\kappa(\xi)=\rho\Bigl(1-\sum_{\xi'} P^{[H]}_{\xi\xi'}\Bigr)\ge0.
\]

Then
\(L^{[H]}_{\xi\xi}=\rho(P^{[H]}_{\xi\xi}-1)
=-\sum_{\xi'\ne\xi}q(\xi,\xi')-\kappa(\xi)\).
Hence \(L^{[H]}\) satisfies the generator form and the positive maximum
principle. The continuous-time process generated by
\(\rho(P^{[H]}-I)\) performs the same clocked transcript transition
chain at Poisson-distributed jump counts, so it represents the same
finite-horizon step relation without adding a movement primitive.

\end{proof}

\Cref{fig-operational-transcript} shows the finite clocked carrier
used to represent an arbitrary computation at horizon \(H\).

\begin{figure}[tbp]
\centering
\begin{tikzpicture}[fw diagram,node distance=7mm]
  \node[fw source] (input) {public input};
  \node[fw box,right=of input] (s0) {\(s_0\)};
  \node[fw box,right=of s0] (s1) {\(s_1\)};
  \node[right=of s1] (dots) {\(\cdots\)};
  \node[fw box,right=of dots] (sh) {\(s_H\)};
  \node[fw terminal,right=of sh] (out) {readout};
  \draw[fw arrow] (input) -- (s0);
  \draw[fw arrow] (s0) -- (s1);
  \draw[fw arrow] (s1) -- (dots);
  \draw[fw arrow] (dots) -- (sh);
  \draw[fw arrow] (sh) -- (out);
  \node[fw group,fit=(s0)(s1)(dots)(sh),
    label={[font=\scriptsize,fwInk]above:\(\Omega^{[H]}\)}] (carrier) {};
  \node[fw provenance,below=12mm of carrier] (kernel)
    {\(P^{[H]}\quad\longmapsto\quad L^{[H]}=\rho(P^{[H]}-I)\)};
  \draw[fw arrow] (carrier) -- (kernel);
\end{tikzpicture}
\caption{A bounded computation becomes a represented clocked transcript
kernel and hence a positive generator by
\cref{thm-algorithm-kernel-embedding}.}
\label{fig-operational-transcript}
\end{figure}

\subsection{Countable processes and finite stopping}
\label{subsec-countable-processes-finite-stopping}

The countable-process record connects an unbounded description to finite charged stopping filtrations.
\begin{definition}[Countably represented process and finite stopping filtration]
\label{def-countably-represented-filtered-process}

A countably represented process is a tuple

\[
\mathfrak P=(\Omega,\mathcal P(\Omega),P,\mu_0),
\]

where \(\Omega\) is a countable set of uniformly finite-encoded states,
\(P\) is a Markov or sub-Markov kernel specified by one represented
transition rule, and \(\mu_0\) is a represented initial law.  Put

\[
R_H(\mu_0)=\bigcup_{t=0}^H\operatorname{supp}(\mu_0P^t).
\]

The process is \emph{finitely stoppable} when every \(R_H(\mu_0)\) is
finite and, uniformly from the codes of \(P,\mu_0\), and \(H\), the
finite disjoint-time carrier

\[
\Omega^{[H]}
=
\left(\bigsqcup_{t=0}^H R_t(\mu_0)\times\{t\}\right)
\sqcup\{\dagger_H\}
\]

and its stopped kernel are represented.  For \(t<H\), the transition
from \((x,t)\) is the pushforward of \(P(x,\cdot)\) under
\(y\mapsto(y,t+1)\); states at time \(H\) enter the absorbing timeout
state \(\dagger_H\).  A sub-Markov deficit is sent to the same state.
The sets \(R_H\) are increasing.  The clocked kernels need not be
nested, but their path laws agree with the original process on every
cylinder event determined before their deadlines.  The countable tuple
together with all of these stopped kernels is its \emph{filtered process
object}.

This definition is a representation layer, not a finiteness assumption
on the underlying computation.  Finite-state profile and decomposition
theorems may be applied to the stopped kernels; a direct infinite-state
analytic theorem requires its own operator-domain hypotheses.

\end{definition}

\begin{corollary}[All Turing machines are represented]
\label{cor-all-turing-machines-represented}

Let \(M\) be an arbitrary deterministic Turing machine, with no
halting, time, space, reversibility, locality, lumpability, or regularity
assumption beyond its finite transition table.  Let \(\Omega_M\) be the
countable set of all finite configurations of \(M\), with the discrete
sigma-algebra.  Extend the transition map \(F_M\) by fixing halting
configurations, and define

\[
P_M(c,A)=\mathbf 1_A(F_M(c)).
\]

For every finite input \(I\), with initial configuration \(c_I\),

\[
\mathfrak P_{M,I}
=
(\Omega_M,\mathcal P(\Omega_M),P_M,\delta_{c_I})
\]

is a single countably represented, finitely stoppable process object.
It represents the complete run, whether or not that run halts.  Its
finite stopped slice at horizon \(H\) reproduces the first \(H\) Turing
transitions exactly and is a deterministic finite Markov kernel whose
generator is the point-mass specialization of
\Cref{thm-algorithm-kernel-embedding}.

Let \(G_{M,I,H}\) instead denote the partial, unclocked transition graph
containing the configurations and transition edges encountered before
time \(H\), without replacing its boundary edge by timeout. Then
\(G_{M,I,H}\subseteq G_{M,I,H+1}\), and their union is the complete
reachable computation graph of \(M\) on \(I\). Every finite prefix and
every halting run is represented exactly in a sufficiently large
slice. Taking the same exhaustion over all finite input encodings (or
all finite starting configurations) recovers the ordinary countable
configuration graph of \(M\). The clocked timeout kernels themselves
need not be nested, and no result below assumes that they are.

Consequently every DTM, including every unbounded or nonhalting DTM, is
represented by the framework as one countable filtered object.  No
structural Valid-Lift, target, geometric, reversibility, or lumpability
hypothesis is used for this representation.  Point-mass kernels are a
proper subclass of the permitted stochastic and sub-Markov kernels;
continuous-time kernels and general represented positive generators add
further operator objects.  Thus the represented operator class strictly
contains deterministic Turing evolution.

\end{corollary}

\begin{proof}\phantomsection\label{proof-all-turing-machines-represented}

Every Turing configuration has a finite encoding, so \(\Omega_M\) is
countable.  The finite transition table computes \(F_M(c)\) uniformly
from \(c\), hence \(P_M\) is a represented point-mass kernel.  From the
point-mass initial law, at most one configuration is reached at each
time, so \(|R_H(\delta_{c_I})|\le H+1\).  \Cref{def-countably-represented-filtered-process} therefore supplies a
finite stopped kernel for every \(H\), and its transitions before the
deadline are exactly those of \(M\).

Removing the artificial deadline edge gives the partial graph
\(G_{M,I,H}\). Increasing \(H\) only appends configurations and
Turing-transition edges, so these partial graphs are nested. Their union
contains exactly every configuration and edge reached after some finite
number of steps on \(I\). Every ordinary configuration has a finite
encoding, and ranging also over finite starting configurations recovers
the full countable Turing configuration graph. This proves both the
finite-slice representation and its exhaustion statement.

Finally, on any represented carrier containing two states, a kernel
having a row with two positive transition probabilities is permitted by
the framework but is not a deterministic point-mass kernel.  Hence the
operator inclusion is proper; sub-Markov killing and general represented
positive generators provide further permitted cases.

\end{proof}

\begin{remark}[Operator domain does not restrict Turing representation]
\label{rem-countable-dtm-operator-domain}

On \(\ell^\infty(\Omega_M)\),

\[
(P_Mf)(c)=f(F_M(c)),
\qquad
L_M=P_M-I
\]

defines, respectively, a positive contraction \(P_M\) and a bounded
generator \(L_M\) with
\(\|L_M\|\le2\).  Thus the complete countable computation is itself an
operator object without an additional analytic assumption.  The
finite-dimensional \(L^2\) projections, Hodge decompositions, spectral
gaps, Schur complements, and matrix resolvents developed below apply
unconditionally to every finite stopped slice.  Applying a specifically
Hilbert-space or spectral statement directly on the countable carrier
requires a declared reference measure and the usual domain, closure, and
limit hypotheses.  Those hypotheses govern that analytic extension;
they are not conditions for representing the Turing machine.

\end{remark}

The operational representation retains every finite outcome, query,
public-history, and terminal record in the transcript carrier. Memory,
stacks, learned databases, priority queues, preprocessing state, actual
parallel branches, and decomposition work remain in the live native
state and are charged at peak capacity; a declared represented effect
exposes any of them that is used as an observable. The channel
decomposition applies to the induced transcript kernel rather than only
to candidate assignments.

\subsection{Extraction and structural admissibility}
\label{subsec-extraction-structural-admissibility}

\Cref{prop-operator-extraction-not-algorithm-proof-witness} separates
extraction of the represented operator from certification of its structure.
\begin{proposition}[Operator extraction is independent of algorithmic certification]\label{prop-operator-extraction-not-algorithm-proof-witness}

Budget-admissible membership is a statement about a uniform operational
transition rule without oracle access under the declared resource
envelope. In the polynomial instance, the procedure supplies its code,
native state bound, and resource bound. No quotient,
invariant, proof of correctness, or explanation of target alignment is
required. Once the transition rule and a finite horizon are fixed, the
outcome/terminal transcript determines a Markov kernel \(P^{[H]}\) and
hence a generator \(L^{[H]}=\rho(P^{[H]}-I)\). The channel decomposition
is then applied to \(L^{[H]}\). Any structural
witness used in later bounds is a budget-admissible
observable-structural representative of the induced operator, not
part of the algorithm's membership in the declared budgeted class. A
future-hitting statistic generated by budget-admissible finite-horizon
computation is counted at that cost; an unrepresented semantic future of
the operator is not a channel representative. In the polynomial
instance, this is membership in the declared uniform polynomial
operational class.

\end{proposition}

\begin{proof}\phantomsection\label{proof-operator-extraction-not-algorithm-proof-witness}

A represented operational transition rule gives the conditional outcome
law and native successor. On the finite clocked transcript carrier this
yields the kernel \(P^{[H]}\) in the embedding theorem. The finite
generator theorem and channel decomposition depend on \(L^{[H]}\) and
the observable algebra of that carrier, independently of an explanatory
annotation of the code. After embedding, every family-uniform target-aligned structure has
a provenance component. It is task structure at the declared budget
precisely when represented in the budgeted saturated observable closure
below the ceiling. Otherwise it belongs to the residual component,
presentation-inaccessible or oracle-supplied information, over-budget
recovery, or an enriched presentation containing the missing datum. The
polynomial case is the instance
\(\mathcal B_{\mathrm{poly}}\).

\end{proof}

\begin{proposition}[Operational embedding and structural budget admissibility]\label{prop-substrate-structural-affordability-split}

Every uniform resource-bounded computation without oracle access
determines a finite generator on each bounded clocked transcript carrier.
This embedding does not assign zero structural representation
cost to the generator, transition graph, committor, first-hitting
partition, or hitting-time function. A finite-horizon object generated
by budget-admissible simulation is instead a lifted observable charged
at the saturated cost of that computation.

A target-relevant component of the induced generator becomes budget-admissible
structure when the component is represented as one of the existing
channel objects at the relevant saturated public cost: a declared
conductance or potential for \(\Geom/\Caus\),
represented lumpable fibers for \(\Abs\), a
target-fraction-preserving valid interface for \(\Lift\), or an
observable terminal value readout for \(\Bdry\). A
budget-admissible finite-horizon simulation is one way to produce such a
representative, and the resulting denotation is counted at the cost of
that simulation. Naming an uncomputed infinite future, occupation law,
or resolvent of the same generator does not define a constructor. A
denotation available only through such an operation is outside the budget-admissible
algebra of the original task unless the same object has a
budget-admissible saturated representative.

\end{proposition}

\begin{proof}\phantomsection\label{proof-substrate-structural-affordability-split}

The embedding theorem records the complete transition rule in the finite
transcript kernel \(P^{[H]}\) and generator \(L^{[H]}\). This construction uses only the
next-step dynamics and the resource envelope, after the admissible
instance-dependent data have been fixed by the budgeted saturated
observable closure. It charges one peak native state and the accumulated
finite transcript, rather than a sequence of native states. Saturated representation cost is the minimal
description cost in that observable structural grammar for the component
extracted from \(L^{[H]}\), including the cost of any budget-admissible finite-horizon
simulation used to generate it. \PartIII records provenance, execution,
and structural representation as distinct readings inside the same
budgeted closure. The embedding theorem establishes the induced operator
and leaves the representation costs of associated denotations unchanged.
A committor, first-hitting partition, or hitting-time table obtained by
solving a global boundary problem for \(L^{[H]}\) is an operator property. The
same finite object is a budget-admissible datum exactly when generated by
budget-admissible computation or otherwise represented by a
budget-admissible quotient,
potential, geometry, lifted interface, or boundary readout.

\end{proof}

\begin{proposition}[Oracle accounting]\label{prop-oracle-accounting}

An oracle or advice mechanism is external to the presented
task. If an edge or observable supplied by such a mechanism is already
\(\mathcal F_I\otimes\mathcal F_I\)-measurable and structurally
explained, it belongs to one of the exposed structural channels. Extra
target information outside the presentation is oracle-supplied and
therefore outside the original budgeted presentation. Directed motion
without exposed structural provenance lies in \(\Res\).

\end{proposition}

\begin{proof}\phantomsection\label{proof-oracle-accounting}

Admissibility is relative to the primitive presentation and its budgeted
saturated closure. If an oracle-supplied relation or edge table is
already measurable in the public product algebra and has a represented
\(\Geom\), \(\Caus\), \(\Abs\), \(\Lift\),
or \(\Bdry\) provenance, the \PartII projections place
it in that channel. If the oracle supplies data not generated by the
presentation, the transition rule uses oracle access and is therefore
outside the budget-admissible rules for the original task; it is an
oracle/advice-relative rule or a changed signature. Finally, if the
supplied motion has no exposed structural provenance after the
projections, the residual complement assigns it to \(\Res\).
The information-provenance theorem exhausts these alternatives.

\end{proof}

Thus every presentation-relative budget-admissible algorithm without
oracle access induces a budget-admissible generator on each finite
clocked transcript carrier. Presentation-relative means that
the instance-dependent data read by the transition rule lie in the
budgeted saturated observable closure of the original task, and that any
finite-horizon trajectory statistic used as structure is generated by
budget-admissible computation and charged there. Every such
generator decomposes into five algebraic structural channel slots and a
residual.  Their charged projections are the exposed channels. Remaining
directed cycle current and nonlocal movement lacking exposed provenance
are assigned to \(\Res\). The decomposition derives provenance from the
induced operator without requiring an algorithmic proof witness. \PartIII
determines the saturated public cost at which each extracted component
appears.

\section{Operator-Level
Universality}\label{sec-operator-universality}
\label{why-this-is-algorithm-independent}

The decomposition is a property of operators on the observable algebra.
An algorithm enters the analysis through the generator induced on its
clocked transcript carrier.

\begin{proposition}[Universality of the operator decomposition]\label{prop-structural-universality}

Universality follows from object exhaustiveness: every budget-admissible
positive generator has exactly five algebraic channel slots and a
residual. Only charged represented projections are exposed at a budget;
remaining directed cycle current and nonlocal movement lacking observable
structural provenance are recorded as \(\Res\). Every
resource-bounded represented computation without oracle access embeds as
such a generator on each finite clocked transcript carrier.
For a declared budget structure, budget-admissible means generated from the
budgeted saturated observable closure of the presented task. Therefore
every algorithmic use of budget-admissible task structure is represented by the
channel decomposition. In particular, an algorithm uses the
intrinsic problem geometry supplied by \(\Geom\) and
\(\Abs\), through causal orientation, valid lifting, and
boundary/value readout. Structure outside the budgeted saturated closure
lies outside the admissible data of the original presentation.

\end{proposition}

\begin{proof}\phantomsection\label{proof-structural-universality}

By the operational embedding theorem, every resource-bounded represented
computation without oracle access becomes a positive finite generator on
its clocked outcome/terminal transcript carrier. The ledger separately
charges the live internal state and memory at peak native capacity. The channel
exhaustiveness theorem applies to every such generator and decomposes
its movement into the five algebraic structural slots plus its residual
component.  At a budget, only charged represented projections are
exposed. The classification of generative sources
identifies \(\Geom\) and \(\Abs\) as the intrinsic
sources, with exact and compensated quotient geometry as their joint
modes. The channels \(\Caus\), \(\Lift\), and
\(\Bdry\) retain the sources they orient, reorganize, or read in
the complete carrier ledger; target charge moves between contextual
extensions only by a represented transfer. Because the
embedding includes every retained admissible instance-dependent record
and every declared native-state effect in the observable space, every induced effect either changes the primitive
signature, uses oracle or advice information, requires over-budget
recovery, or appears in this decomposition. Universality therefore
concerns the operator class rather than an enumeration of algorithms.

\end{proof}

Once the observable algebra, budget-admissibility condition, positivity
principle, and clocked transcript carrier are fixed, every algorithm in
the stated class induces an operator already covered by the
decomposition. A budget-admissible method without oracle access may change
the exposed transcript/readout space, generate an exposed edge relation,
compute a finite-horizon trajectory statistic, or select a different
kernel using admissible instance data. Each case still produces a
positive generator on the budgeted observable algebra. Intrinsic problem
structure appears through \(\Geom\) or \(\Abs\), its
derived uses through \(\Caus\), valid \(\Lift\), or
\(\Bdry\), and the remainder through \(\Res\). A
semantic value function on the configuration graph is budget-admissible
only when it has an admissible representation, including one produced by
finite-horizon computation from the observable closure.

\section{Examples: Dynamic Channels for SAT and Subset
Sum}\label{sec-dynamic-examples}
\label{worked-the-channels-on-sat-and-subset-sum}

The static SAT and subset-sum presentations from \PartI now illustrate
the dynamic channel decomposition.

\subsection{SAT}\label{subsec-dynamic-sat-example}\label{sat}

For a CNF formula \(\Phi\), the state space is \(X_\Phi=\{0,1\}^n\). A
one-bit Hamming relation is an exposed local relation when the
presentation declares bit flips. Symmetric one-bit rates belong to the
geometric/reversible channel \(\Geom\). A directed local imbalance
belongs to \(\Caus\) when it descends an exposed structural potential, such as a computable
clause-defect or propagation potential. If the directed imbalance has no
such provenance, it is \(\Res\).

XOR-SAT supplies the clean abstraction witness. Public parity equations
define a represented quotient

\[
\pi(x)=Ax,
\qquad
\pi^{-1}(b)=\{x:Ax=b\}.
\]

A jump or walk that descends to the quotient and satisfies the
intertwining relation \(L_XU=UL_Q\) belongs to \(\Abs\). Row
reduction operates on the represented quotient.

Random planted 3-SAT illustrates the opposite case. The satisfying
states exist, but a planted assignment outside the final formula's
observable closure does not define a budget-admissible operator edge. A
directed search path without a represented potential or quotient
coordinate contributes to \(\Res\).

\subsection{Subset Sum}\label{subsec-dynamic-subset-sum-example}\label{subset-sum}

For subset sum, states are inclusion vectors \(x\in\{0,1\}^n\), and the
public defect observable is

\[
D(x)=\left|\sum_i a_i x_i-t\right|.
\]

One-bit inclusion moves belong to the geometric/reversible channel when declared local. A
public residue map

\[
r_M(x)=\sum_i a_i x_i\pmod M
\]

gives a quotient/lumping channel when the modulus and residue fibers are
exposed by the task presentation. A hidden modulus may hold in the
ambient arithmetic, but it is budget-inadmissible unless made observable
by the presentation or generated within budget in the lifted
configuration. Otherwise it is presentation-inaccessible,
oracle-supplied, over budget, or residual relative to the public
presentation.

\section{Summary of the Dynamic
Decomposition}\label{sec-dynamic-summary}
\label{summary-the-exhaustive-dynamical-object}

\PartII represents dynamics as a search operator on frontiers. Under
positive sub-Markov evolution, equivalently the positive maximum
principle in finite dimension, this operator has generator form; under budget admissibility
it is measurable with respect to the observable algebra and represented
in the active budgeted closure. Relative to the local relation and
structural projections, it decomposes as

\[
\Geom,\qquad
\Caus,\qquad
\Abs,\qquad
\Lift,\qquad
\Bdry,\qquad
\Res.
\]

\paragraph{Dependency trail.}
The dynamic objects and structural subspaces are
\cref{def-frontier-search-operator,def-admissible-operator,def-fixed-space-split,def-structural-lift-gate,def-partial-lumpability-defect}.
Their information, residual, coupled-source, and lifting records are
\cref{def-information-provenance,def-post-exhaustion-residual-band,def-coupled-abs-geom-source,def-lift-closure}.
Generator and fixed-space accounting are established in
\cref{prop-trajectory-measures-not-gauges,thm-fixed-space-split,prop-no-third-local-primitive}.
The structural gates and complete lifted carrier are treated in
\cref{prop-gate-predicates-objective,prop-full-configuration-lift-complete,prop-derived-channels-create-no-structure,prop-partial-lumpability-no-escape,prop-channel-functoriality-presentation-equivalence}.
The provenance and residual consequences are
\cref{thm-exhaustive-residual-normal-form,cor-no-residual-provenance-outside-normal-form,cor-no-cleverness,prop-family-uniform-residual-reclassification,thm-no-uniform-residual-correlation,prop-provenance-is-not-degree,thm-irreducible-band-closure}.
Composition, lifting, and the operational bridge are summarized by
\cref{lem-gaplessness-composition,cor-lift-creates-nothing,prop-substrate-structural-affordability-split,prop-structural-universality,thm-dynamical-closure}.
The interpretive qualifications concerning channel components and the
countable DTM carrier are recorded in
\cref{rem-channel-components-need-not-be-generators,rem-countable-dtm-operator-domain}.

\begin{lemma}[Dynamical closure of \PartII]\label{thm-dynamical-closure}

Every budget-admissible positive search operator, including every
represented computation without oracle access on its finite clocked
transcript carrier, has exactly five algebraic
structural channel slots: \(\Geom\), \(\Caus\),
\(\Abs\), \(\Lift\), and \(\Bdry\). After those
five projections are removed, the remaining movement is the residual
component \(\Res\). At a budget \(B\), a slot is exposed
only through a charged represented projection record; an unrepresented
algebraic portion is retained in the filtered residual. The
intrinsic sources of reusable problem geometry are \(\Geom\) and
\(\Abs\). Exact quotient structure is \(\Abs\); partial
quotient structure is exact \(\Abs\) plus a defect operator that
is decomposed again. A quotient-level geometric comparison transfers to
the ambient process only through the quantitative compensation criterion
of \PartIII. The channels \(\Caus\), valid
\(\Lift\), and \(\Bdry\) orient, reorganize, or read out
those intrinsic sources.  The \(\Caus\) projection is the one
authorized by the selected imposed/control generator, and legal lifts
preserve its current identity. Any local cycle current orthogonal to the
exposed gradient space, or nonlocal movement without quotient, valid
lift, or boundary-readout provenance, is residual transport. Composition
and valid lifting preserve this classification.

\end{lemma}

\begin{proof}\phantomsection\label{proof-dynamical-closure}

The positive maximum principle gives rates and diagonal loss. The
fixed-space split separates local symmetric, local directed, and
nonlocal rates, together with boundary processing. The structural
projections identify local symmetric rates as geometric/reversible,
directed flow as a structural gradient plus residual, and nonlocal
structure as quotient/lumping or lifted-state when the observable algebra and
target-fraction preservation supply that provenance, partial-lumpability
defects by recursive component decomposition, and terminal observables
as boundary value functions. Every remaining directed or nonlocal
operator portion is \(\Res\) residual accounting.  Its directed and
nonlocal coordinates are exhausted by
\cref{thm-exhaustive-residual-normal-form}.  The budget-filtered assignment
uses the charged projection convention of
\cref{thm-channel-exhaustiveness,thm-exhaustive-residual-normal-form}. The Hodge projection is
taken from the selected imposed/control generator, so its uniqueness
also gives the no-substitution property. The operational embedding theorem places
algorithms into this same generator class, and the composition lemma
prevents lifted constructions from adding another structural primitive.
Hence the displayed decomposition is complete.

\end{proof}

Structured movement is generated by exposed geometry and represented
abstraction. The directed/drift channel orients exposed potentials, the
lifted-state channel reorganizes observable coordinates while preserving
target fractions, and the boundary/killing channel maps exposed terminal
observables to value functions, with termination or absorption as the
zero-value case. Nonlocal motion without one of these provenance
witnesses belongs to \(\Res\). In particular, a restart
without declared reversible structure is an ordered nonlocal jump in the
residual component.

\PartIII introduces the quantitative arrival-time analysis and the
saturated public cost at which each exposed structural channel becomes
budget-admissible. Its structural quantities are defined on the five
exposed channels, with \(\Res\) treated separately as the
component without exposed structural provenance.

\clearpage
\part{Spectral Complexity and Saturated Extraction Cost}\label{part:spectral-complexity}\setcounter{section}{-1}

\section{Introduction: Discounted Hitting Functionals and Spectral
Cost}\label{introduction-from-exhausted-movement-to-measured-time}
\label{sec:spectral-complexity-introduction}

\paragraph{Part contract.}
\emph{Inputs:} the static presentation of \cref{part:task-spaces} and the
channel decomposition of \cref{part:dynamics}.
\emph{Output:} the saturated structural profile, exact charging identities,
first-hit accounting, and universal accountability used in
\cref{part:learning,part:lpn}.
\emph{First pass:} the Walsh specialization, continuum interface, and examples
may be skimmed after the general profile results.

\PartI specifies the static task, and \PartII classifies its
budget-admissible dynamics. This part introduces the corresponding
quantitative theory of target-hitting times.

\PartI defined the presented object

\[
\mathcal T_I=(X_I,\operatorname{Obs}_I,\operatorname{Dyn}_I,\operatorname{End}_I)
\]

and developed the observable algebra, terminal sectors, pointwise and
represented-fiber observables, measurable and presentation-inaccessible
structure, and static comparability. \PartII proves that every
budget-admissible evolution decomposes into five structural components
and a residual term:

\[
\Geom,\qquad
\Caus,\qquad
\Abs,\qquad
\Lift,\qquad
\Bdry,\qquad
\Res.
\]

These components have distinct analytic roles. The intrinsic structure
is encoded by the geometric or reversible component \(\Geom\)
and the quotient or lumping component \(\Abs\). The directed or
drift component \(\Caus\), the lifted-state component
\(\Lift\), and the boundary or killing component
\(\Bdry\) orient, reorganize, and evaluate that intrinsic
structure. In particular, \(\Bdry\) consists of terminal or
support indicators and boundary values; it has no off-diagonal support
and therefore induces no transport by itself. The residual term records
transport without a representative in the exposed structural
decomposition.

For a budget-admissible operator, the central quantity is the discounted
mass reaching the positive terminal sector. The analysis is organized
by the comparison

\[
\boxed{
\text{complexity}
=
\text{arrival time}
=
\text{target resolvent}
\le
\text{saturated extraction at budget }B
}
\]

and its associated invariant is the least saturated public cost at which
an exposed structural component produces non-negligible target
transport.

The interaction of Geom and Abs is expressed by one compensated
quotient--geometry certificate. Its five parameter regimes are ambient
geometry, exact quotient structure with a null geometry slot, geometry
on an exact quotient, reversible compensation of a nonzero quotient
defect, and general resolvent compensation of a nonzero quotient defect.
The identity quotient gives pure ambient Geom, while zero defect and a
null slot give pure Abs. Caus supplies the represented descending
current whenever the geometry slot is active. The saturated profile
charges the complete certificate, including its quotient fibers,
geometry, potential, comparison window, and defect loss.

\section{Global Notation and the Saturated Closure}
\label{sec:global-saturated-closure}

We begin by fixing the common finite-state notation and the budgeted closure in
which all subsequent certificates are represented.

\begin{definition}[Global notation for \PartIII]\label{def-paper3-global-notation}

A family is indexed by a size parameter \(n\). The finite state space is
\(X_n\), the reference measure is \(\mu_n\), and the positive terminal
sector is

\[
T_n=E_{I_n}^+\subseteq X_n,
\qquad
N_n=X_n\setminus T_n.
\]

The centered target contrast is

\[
y_n=\mathbf 1_{T_n}-\mu_n(T_n),
\qquad
\bar y_n=\frac{y_n}{\|y_n\|_{L^2(\mu_n)}}
\quad \text{when }0<\mu_n(T_n)<1.
\]

The ambient resource model is a budget structure
\(\mathcal B=(\mathcal R,\operatorname{cost},b,\mathfrak C)\): a
resource space, a constructor cost, a size-indexed capacity functional
\(b(n)\), and a transfer class. A numeric margin has transfer-class size
when it is bounded below by \(1/\beta(n)\) for some
\(\beta\in\mathfrak C\), and it is negligible relative to
\(\mathfrak C\) when it is bounded above by \(1/\beta(n)\) for every
\(\beta\in\mathfrak C\) eventually. For \(\mathcal B_{\mathrm{poly}}\),
these specialize to inverse-polynomial and negligible. A representative
lies within \(\mathcal B\) when its saturated description,
construction, fiber-use, lift, comparison, boundary-readout, and
finite-postprocessing charges lie below \(b(n)\). The optional degree
symbol \(d\) denotes one cost coordinate, often the log-cost level
\(d=\lceil\log_2 B\rceil\) for a numeric ceiling \(B\), not raw Walsh
degree.

The polynomial model is the instance \(\mathcal B_{\mathrm{poly}}\) with
\(C=\mathbb N\), \(\preceq=\le\), \(\oplus=+\), and
\(b_{\mathrm{poly}}(n)=n^{C_0}\), with the polynomial transfer class
taking the union over constants \(C_0\).

The dense Boolean Walsh band is one charged coordinate inside this
profile. If \(r\) is raw Walsh degree, the dense scan/projection cost is

\[
\operatorname{Cost}_W(n,r)=\dim V_{n,r}^{W}=\sum_{k\le r}\binom{n}{k}.
\]

Thus, under \(\mathcal B_{\mathrm{poly}}\), a dense Walsh representation
lies within budget precisely at fixed raw \(r\) for each fixed exponent
in the polynomial budget instance, while a log-cost parameter satisfies
the polynomial budget when \(d=O(\log n)\). Under another budget, the
admissible Walsh cutoff is the largest \(r\) with
\(\operatorname{Cost}_W(n,r)\preceq b(n)\). These are different
coordinates in the same saturated cost profile.

\end{definition}

The saturated represented closure collects all observables obtainable within the declared budget.
\begin{definition}[Budgeted saturated represented closure]\label{def-budgeted-saturated-observable-closure}

For a presented task family \((I)_{I\in\mathfrak I_n}\), let
\(\mathcal R^{\mathrm{sat}}_{n,\le B}\) be the denotations produced by
one family-uniform oracle-free represented construction of saturated
public cost at most \(B\). The available constructors include declared
public observables, represented quotients and fibers, admissible public
basis changes, budget-generated clocked transcript carriers and their
declared readouts, structural Valid-Lift interfaces, boundary
readouts, forward observable transport data, and finite represented
postprocessings. Every denotation retains its complete construction
record.

The ambient measurable envelope

\[
\mathcal F^{\mathrm{env}}_{I,\le B}
=
\sigma\left\{
R_I:R\in\mathcal R^{\mathrm{sat}}_{n,\le B}
\right\}
\]

supports conditional expectations and finite atom calculations. It is
not the budget slice: an element of this envelope belongs to
\(\mathcal R^{\mathrm{sat}}_{n,\le B}\) only when it has its own
family-uniform construction record of total cost at most \(B\).
Constructor application therefore satisfies a filtered closure rule,
with the input, operation, output, fiber, and postprocessing costs added
to determine the output slice.

For a budget structure
\(\mathcal B=(\mathcal R,\operatorname{cost},b,\mathfrak C)\), the
within-budget represented closure is
\(\mathcal R^{\mathcal B}_n
=\mathcal R^{\mathrm{sat}}_{n,\le b(n)}\). It is the source of legal
instance-dependent data for computation under \(\mathcal B\). Legal
finite-horizon iteration extends this closure only through a represented
computation whose transcript, lifted coordinates, precision, and
postprocessing are included in its accumulated charge. An ambient
relation or semantic future object requiring over-budget construction,
oracle access, nonuniform advice, an uncomputed resolvent, or an
undeclared primitive has no within-budget representative for the
original presentation.

The polynomial represented closure consists of uniform families
\(R=(R_I)\) for which one fixed exponent \(C<\infty\) and one uniform
derivation \(\rho\) satisfy

\[
\sem{\rho}_I=R_I,
\qquad
\operatorname{Cost}_I(\rho)\le (1+n)^C
\quad
(n\ge1,\ I\in\mathfrak I_n).
\]

Write \(\mathcal R^{\mathrm{poly}}_n\) for the size-\(n\) slice of
these families. The exponent is uniform in \(n\) and in the instance.

\end{definition}

\Cref{fig-saturated-closure} distinguishes the represented budget
slice from the larger ambient collection of semantic objects.

\begin{figure}[tbp]
\centering
\begin{tikzpicture}[fw diagram,node distance=12mm]
  \node[fw source] (public) {public primitives};
  \node[fw provenance,right=of public] (construct) {represented\\constructors};
  \node[fw box,right=of construct,minimum width=3.2cm] (closure)
    {\(\mathcal R^{\mathrm{sat}}_{n,\le B}\)};
  \draw[fw arrow] (public) -- (construct);
  \draw[fw arrow] (construct) -- (closure);
  \node[fw group,fit=(public)(construct)(closure),
    label={[font=\scriptsize,fwInk]above:complete cost \(\le B\)}] (budget) {};
  \node[fw residual,below=13mm of construct] (outside)
    {oracle, advice, or over-budget object};
  \draw[fw dashed arrow] (outside.east) -| node[pos=0.25,below,font=\scriptsize]
    {not represented} ($(budget.south east)+(0,-1mm)$);
\end{tikzpicture}
\caption{The saturated closure contains exactly the uniformly represented
objects whose complete construction records fit the declared budget of
\cref{def-budgeted-saturated-observable-closure}.}
\label{fig-saturated-closure}
\end{figure}

The derivability closure records the constructors and provenance available at a fixed budget.
\begin{definition}[Budgeted derivability closure: represented grammar]\label{def-budgeted-derivability-closure}

The budgeted derivability closure consists of the represented grammar below,
the uniform budget filtration of
\cref{def:uniform-derivation-budget-filtration}, and the admissibility
criterion of \cref{def:derivation-admissibility-criterion}.  The present
definition fixes the grammar and therefore retains the public label for the
combined construction.

For the same presentation, budget structure, and resource envelope, let
\(\mathsf{Der}_{\mathfrak I}\) be the smallest class of
family-uniform derivation schemes generated by the following clauses.
For \(\rho\in\mathsf{Der}_{\mathfrak I}\), write
\(\sem{\rho}_I\) for its denotation on \(I\) and
\(\operatorname{Cost}_I(\rho)\) for its complete charged execution and
representation cost.

\medskip\noindent\textbf{(i) Local coding.}
The derivation coding is locally finite: the description alphabet is
finite, the charged cost dominates encoded description length, and
there are only finitely many uniform derivation descriptions of cost at
most \(B\) for every finite \(B\).  Encoded coefficients, tables, and
precision parameters contribute their full description lengths.  This
is a convention on the resource coding, independent of the presented
task family.

\medskip\noindent\textbf{(ii) Primitive records.} Declared public observables, public parameters,
terminal predicates, budget-generated transcript coordinates,
finite outcome-generation interfaces,
and the constructor signatures declared by the presentation are leaves
of a derivation, with their public description and access costs.

\medskip\noindent\textbf{(iii) Constructors.} Derivations are closed under the admissible
oracle-free constructors of \PartsItoII: finite postprocessing,
represented joins, quotient and fiber formation, public basis changes, valid
lifts with \(\pi_*\Lambda\mu=\mu\) and target-fraction preservation,
boundary readouts through already derived structure, conductance and
potential formation, forward observable transport data, channel
projection data, and the finite comparison data required by the
discounted hitting estimates. The charged cost of a constructor output is the sum of
input costs plus the declared constructor charge for its arity, output
description length, evaluation rule, represented fiber operations,
lifted-interface size, retained coefficients or tables, and displayed
comparison constants. For \(\mathcal B_{\mathrm{poly}}\), admissible
constructor records have polynomially bounded total charge.

\medskip\noindent\textbf{(iv) Computation closure.} If one uniform oracle-free procedure has
all instance-dependent inputs represented by prior derivations and
halts within the declared budget-compatible resource envelope, its
finite output has the composite derivation obtained by adjoining the
represented outcome/terminal transcript and terminal output record. Its
full padded transcript carrier, initializer, operational update, terminal interface, and any public
full-support product gauge are legal computation records independently
of the structural Valid-Lift identities. The charged cost includes input access,
time, peak native-state capacity, outcome generation, accumulated transcript, retained data, output
representation, and postprocessing. This clause includes finite-horizon trajectory statistics,
first-hit quotients, truncated value functions, represented propagation
components, and target-contact statistics computed by legal simulation
on the clocked transcript carrier with the current native state charged
at peak capacity.

\medskip\noindent\textbf{(v) Evaluator normalization.} Every finite state-dependent
primitive and evaluator in an oracle-free derivation has an explicit
charged table or a uniform charged implementation over a fixed finite
operational basis of encoded coordinate access, affine and product
Boolean gates, comparisons and selection, represented group-action
instructions, and clocked state or memory updates. The implementation
record is part of the derivation. State-dependent access supplied
without either representation belongs to the
\hyperref[def-oracle-extended-closure]{oracle-extended derivability
closure}.  This clause normalizes an evaluator already present in a
derivation; it does not add a semantic function merely because the
operational basis is functionally complete.  An explicit table retains its
full encoded length, storage, and access cost.  A compact circuit or other
factorization is admitted only through its own represented derivation and is
charged at its actual description, evaluation, storage, and retained-state
cost.

\end{definition}

\begin{definition}[Uniform derivations and budget filtration]
\label{def:uniform-derivation-budget-filtration}

Fix the presentation, resource envelope, and represented grammar of
\cref{def-budgeted-derivability-closure}.

\medskip\noindent\textbf{(vi) Uniform provenance.} The program producing the derivation is
fixed across the family. Instance-dependent strings, coefficients,
tables, or fiber labels must be generated from
\(\operatorname{enc}(I)\) by prior vertices of the same derivation.
Changing the program or inserting a string separately for each instance
is nonuniform advice rather than a public derivation. A fixed succinct
uniform formula remains admissible at its actual description and
execution cost; the framework does not charge the process by which that
formula was discovered.

\medskip\noindent\textbf{(vii) Minimality and budget slices.} No other derivations belong to
\(\mathsf{Der}_{\mathfrak I}\). For a denoted family \(R=(R_I)\),
define its upward-closed set of admissible uniform budgets by

\[
\mathfrak B_{\mathrm{sat},n}(R)
:=
\left\{
B:\ \exists\rho\in\mathsf{Der}_{\mathfrak I}
\quad
\sem{\rho}_I=R_I,
\quad
\operatorname{Cost}_I(\rho)\preceq B
\quad( I\in\mathfrak I_n)
\right\}.
\]

The primary budget slice is therefore

\[
\mathcal D_{n,\le B}
:=
\{R:B\in\mathfrak B_{\mathrm{sat},n}(R)\}.
\]

When the resource order is partial, the saturated cost of \(R\) is the
Pareto-minimal frontier of
\(\mathfrak B_{\mathrm{sat},n}(R)\), not a putative least vector. When a
declared admissible scalar cost \(\kappa:C\to[0,\infty]\) has been
fixed, define

\[
\operatorname{Cost}_{\mathrm{sat},n}^{\kappa}(R)
=
\inf_{\rho:\,\sem{\rho}_I=R_I\ \forall I\in\mathfrak I_n}
\ \sup_{I\in\mathfrak I_n}
\kappa\!\left(\operatorname{Cost}_I(\rho)\right),
\]

and omit the superscript only in scalar-cost statements. Thus no
infimum or supremum in the resource monoid is used to define a budget
slice. The saturated represented class used below is
identified with this closure:

\[
\mathcal R^{\mathrm{sat}}_{n,\le B}=\mathcal D_{n,\le B}.
\]

Here admissibility of \(\kappa\) means that it is monotone, dominates
encoded description length up to a fixed class-preserving factor, and
has finite description sublevels.  These conditions prevent a
scalarization from discarding an unbounded charged coordinate.  A
budget-indexed expression such as
\(\operatorname{Cost}_{\mathrm{sat}}(R)\preceq B\) always means
\(B\in\mathfrak B_{\mathrm{sat},n}(R)\), witnessed by one actual uniform
derivation. The typography \(\le B\) is used only after a numeric
scalarization has been declared.

\end{definition}

\begin{definition}[Derivation admissibility criterion]
\label{def:derivation-admissibility-criterion}

Use the grammar and budget filtration of
\cref{def-budgeted-derivability-closure,%
def:uniform-derivation-budget-filtration}.

\medskip\noindent\textbf{(viii) Admissibility criterion.}
An object is admissible within the resource envelope precisely when it
is the denotation of one uniform derivation whose cost is charged to
that envelope. Ambient measurability and membership in a measurable
envelope do not supply such a derivation. A datum without one is not admissible for the presentation under
\(\mathcal B\); its use requires presentation-inaccessible information,
nonuniform advice, oracle access, over-budget recovery, or an enlarged
primitive signature.

\end{definition}

\begin{lemma}[Description length is a charged resource]
\label{lem-description-length-charged-resource}

Fix the locally finite derivation coding of
\cref{def-budgeted-derivability-closure} and an admissible scalarization
\(\kappa\) from
\cref{def:uniform-derivation-budget-filtration}.  There is a coding constant
\(c_{\rm desc}\ge1\), independent of the presented family and size slice,
such that every represented derivation \(\rho\) satisfies
\begin{equation}
\label{eq-description-length-charged-resource}
\bigl|\operatorname{enc}(\rho)\bigr|
\le
c_{\rm desc}\!\left(
1+\sup_{I\in\mathfrak I_n}
\kappa\bigl(\operatorname{Cost}_I(\rho)\bigr)
\right).
\end{equation}
Consequently, a derivation below a polynomial scalar ceiling has polynomial
encoded description length.  A single family-uniform derivation has one fixed
finite program code across all sizes.  A size-dependent table or coefficient
list is admissible only when that fixed program constructs it from the public
presentation and its output, construction, storage, and evaluation costs fit
the current ceiling.  Inserting such data separately for each instance is
nonuniform advice and belongs to an enriched presentation.

\end{lemma}

\begin{proof}
Local coding in
\cref{def-budgeted-derivability-closure} makes encoded description length a
coordinate dominated by complete cost.  Admissibility of \(\kappa\) in
\cref{def:uniform-derivation-budget-filtration} preserves that domination up
to one fixed class-preserving coding factor, which gives
\cref{eq-description-length-charged-resource}.  The same two definitions
charge generated tables, retained coefficients, output length, storage, and
evaluation.  Their uniform-provenance clause requires every
instance-dependent string to be produced by an earlier vertex of the same
fixed derivation.  A separately inserted string therefore has no derivation
in the original presentation and is classified by
\cref{def:derivation-admissibility-criterion,prop-oracle-accounting}.
\end{proof}

\begin{theorem}[Complete charging of explicit and implicit spectral aggregation]
\label{thm-complete-spectral-aggregation-charging}

Let \(I\) be a presented instance, let \(X_I\) be a finite abelian carrier,
and let \(h_I:X_I\to\mathbb C\) be the denotation of a represented
family-uniform derivation \(\rho_h\).  An orthogonal expansion
\[
h_I=\sum_{\gamma\in\widehat X_I}\widehat h_I(\gamma)\chi_\gamma
\]
is an analytic identity.  It creates no derivation vertex, coefficient
oracle, locator, or prospective source.  At every declared charged ceiling,
the following alternatives exhaust the represented uses of this expansion.

\begin{enumerate}[label=\textup{(\roman*)}]
\item If a derivation materializes a finite set \(L_I\subseteq\widehat X_I\)
      of indices or coefficients, the complete record contains the rule
      constructing \(L_I\), every coefficient evaluator, the represented
      output, its storage and precision, and every subsequent access.  Under
      an admissible scalarization, the output and evaluation coordinates of
      its complete cost dominate \(|L_I|\) up to the fixed coding constant.

\item If \(h_I\) is evaluated without materializing its formal expansion,
      then \(\rho_h\) is the compression record.  Its circuit,
      factorization, recurrence, tensor rule, shared subexpressions,
      generated parameters, memory, precision, and evaluation cost are
      charged exactly as represented.  The cardinality of the formal
      spectral support is not charged term by term unless those terms are
      individually constructed or accessed.

\item If a compact evaluator is used prospectively, the complete derivation
      additionally contains the represented comparison, index-selection,
      aggregation, potential, transition, attention, or active-acquisition
      rule that converts its output into the claimed target-directed use.
      Every instance-dependent field of that rule is generated from the
      declared public presentation by an ancestral derivation.  A selection
      predicate involving a nonpublic target is unavailable unless an
      independent represented derivation supplies the same comparison, in
      which case that derivation and its complete cost are retained.

\item The complete prospective record in item~\textup{(iii)} is assigned by
      the existing structural gates: qualified nonidentity fibers enter an
      exact or compensated quotient mode; identity-supported qualified
      transport enters the ambient \(\Geom/\Caus\) mode; and a readout with
      no represented authorized intermediate use is terminal accounting.
      The implicit spectrum is not an additional source modality.
\end{enumerate}

Consequently, an exponentially supported formal spectrum has two distinct
budget interpretations.  An explicitly accessed expansion pays its
materialized cardinality.  A compact implicit aggregation pays the complete
cost of the structural compression rule and of its prospective consumer.
Neither representation supplies free coefficient location or free target
alignment.

\end{theorem}

\begin{proof}
Evaluator normalization in
\cref{def-budgeted-derivability-closure} admits either a charged table or a
charged uniform implementation and retains the implementation record.
Local coding and
\cref{lem-description-length-charged-resource} charge every materialized
index, coefficient, output coordinate, precision field, and stored table.
This proves item~\textup{(i)}.  The same evaluator-normalization clause
charges a compact circuit or factorization at its actual representation and
execution cost, rather than at the support size of an analytic basis
expansion; this proves item~\textup{(ii)}.

Computation closure and the dependency principle retain every ancestor of a
prospective use.  Uniform provenance requires all instance-dependent fields
to be generated by the same family-uniform derivation.  Temporal
availability and target qualification are
\cref{def-pre-hit-temporally-admissible-source,thm-no-free-succinct-transition-rule};
they prove item~\textup{(iii)}.  Finally,
\cref{thm-affordable-geom-abs-normal-form,cor-prospective-nonlinear-readout-routing,cor-terminal-verifier-not-prospective-structural-activity}
give the exhaustive assignment in item~\textup{(iv)}.
\end{proof}

\Cref{fig-derivation-grammar-budget-gate} records the distinction between a
semantic object and a denotation admitted by the charged derivation grammar.

\begin{figure}[tbp]
\centering
\resizebox{\linewidth}{!}{%
\begin{tikzpicture}[fw diagram,node distance=7mm and 10mm]
  \node[fw source,text width=2.5cm] (prim) {declared public\\primitives};
  \node[fw provenance,text width=2.8cm,right=of prim] (ctor)
    {represented joins, quotients,\\lifts, and bases};
  \node[fw use,text width=2.8cm,right=of ctor] (comp)
    {computation and\\evaluator normalization};
  \node[fw box,text width=2.7cm,right=of comp] (uniform)
    {one family-uniform\\provenance record};
  \node[fw terminal,text width=2.7cm,right=of uniform] (slice)
    {complete cost \(\preceq B\)\\\(\mathcal D_{n,\le B}\)};
  \draw[fw arrow] (prim) -- (ctor);
  \draw[fw arrow] (ctor) -- (comp);
  \draw[fw arrow] (comp) -- (uniform);
  \draw[fw arrow] (uniform) -- (slice);
  \node[fw residual,text width=3.2cm,below=11mm of comp] (outside)
    {oracle, advice, semantic-only,\\or over-budget datum};
  \draw[fw dashed arrow] (outside) -- node[right,font=\scriptsize]
    {no admitted derivation} (uniform);
\end{tikzpicture}%
}
\caption{The derivability filtration is generated from declared primitives
by represented constructors and charged computations.  Family uniformity and
the complete resource record are independent gates: semantic measurability
without both gates does not place a denotation in the budget slice.}
\label{fig-derivation-grammar-budget-gate}
\end{figure}

\subsection{Resource-filtered derivation normal form}
\label{subsec:resource-filtered-derivation-normal-form}

The derivation normal form exposes every constructor and retained datum in a
budget-admissible representation.

\begin{definition}[Expanded resource-filtered derivation normal form]
\label{def-expanded-derivation-normal-form}

Let \(\rho\in\mathsf{Der}_{\mathfrak I}\). Its expanded normal form
\(\operatorname{Exp}(\rho)\) is the same uniform typed derivation graph
with every computation-closure vertex displayed through the finite
outcome/terminal transcript
already charged by \Cref{def-budgeted-derivability-closure}.
Every vertex has one of the following presentation-independent forms:

\begin{enumerate}

\item a declared primitive or public parameter;

\item a finite represented composition, join, quotient/fiber
constructor, or public change of coordinates;

\item a budget-generated transcript coordinate, elementary
update, finite outcome-interface update, retained terminal record, or a
separately certified structural Valid-Lift interface;

\item a represented conductance, potential, forward-transport,
channel-projection, or comparison-data constructor; or

\item a boundary or terminal readout through previously derived data.

\end{enumerate}

All instance-dependent inputs to a vertex occur at earlier vertices of
the same graph. The normal-form cost is inherited from the original
record:

\[
\operatorname{Cost}_I^{\mathrm{nf}}
   (\operatorname{Exp}(\rho))
:=
\operatorname{Cost}_I(\rho).
\]

Displaying a transcript therefore does not charge its data a second time.
Let \(\mathsf{Der}^{\mathrm{nf}}_{\mathfrak I}\) be the family of
expanded records obtained in this way.

\end{definition}

\begin{lemma}[Expanded derivation equivalence]
\label{thm-expanded-derivation-equivalence}

Expansion and contraction preserve uniform denotation and total
saturated cost. More precisely, for every
\(\rho\in\mathsf{Der}_{\mathfrak I}\),

\[
\sem{\operatorname{Exp}(\rho)}_I=\sem{\rho}_I,
\qquad
\operatorname{Cost}_I^{\mathrm{nf}}
   (\operatorname{Exp}(\rho))
=\operatorname{Cost}_I(\rho)
\]

on every instance. Conversely, contracting the displayed elementary
transcript vertices of any
\(\widehat\rho\in\mathsf{Der}^{\mathrm{nf}}_{\mathfrak I}\) produces
a derivation in \(\mathsf{Der}_{\mathfrak I}\) with the same denotation
and cost. Hence the budget slices of the two derivation systems have
identical denotations.

\end{lemma}

\begin{proof}

Proceed by structural induction on the derivation graph. Primitive
vertices agree. Each constructor vertex retains its inputs, output,
evaluation rule, represented interface, and declared charge, so the
induction hypothesis passes through every constructor clause. For a
computation-closure vertex, \Cref{def-budgeted-derivability-closure}
already includes the finite outcome/terminal transcript, peak native and
work-state capacity, retained data, and postprocessing in the record and its charge.
Expansion displays these existing entries without changing the induced
map or adding a charge. Contraction reverses this bookkeeping operation.
The uniform program and every instance-dependent dependency are
unchanged, which preserves family uniformity and proves both directions.

\end{proof}

\section{Five-Family Provenance for Represented Evaluators}
\label{sec:five-family-evaluator-provenance}

Every normalized state-dependent evaluator instruction receives exhaustive
constructor provenance.  Exact quotient and active geometric records are the
two structural applications used below.

\begin{definition}[Five-family evaluator provenance]
\label{def-five-family-abs-provenance}

Let

\[
\mathfrak J_{\Abs}^5
=
\{\mathsf{Add},\mathsf{Mult},\mathsf{Ord},
  \mathsf{Sym},\mathsf{Filt}\}.
\]

The five represented provenance families for state-dependent evaluator
instructions are as follows.

\begin{enumerate}

\item \(\mathsf{Add}\) contains encoded coordinate access, affine and
additive combinations, characters, module operations, parity maps, and
additive convolution.

\item \(\mathsf{Mult}\) contains Boolean conjunction, products,
composition gates, polynomial and ring operations, equality tests, and
represented fiber formation.

\item \(\mathsf{Ord}\) contains comparisons, branching and selection,
weights, metric comparisons, quadratic comparisons, and ordered
reduction rules.

\item \(\mathsf{Sym}\) contains represented group actions, orbit maps,
equivariant operations, averaging operators, and isotypic projections.
The authenticated action record represents the action on the active
public object and verifies equivariance of the state-dependent primitive
interface used by the orbit or isotypic constructor. A represented
permutation without this equivariance record retains its operational
implementation provenance but does not acquire a \(\mathsf{Sym}\) tag.

\item \(\mathsf{Filt}\) contains clocks, memory updates, recurrence,
bounded iteration, refinement, outcome histories, transcript
coordinates, flags, and kernel towers. Nilpotent update rules are
filtration provenance in this sense. This tag records the operational
carrier and does not, by itself, authenticate an independent generative
source. Its target charge is governed by
\hyperref[thm-filtration-provenance-charging]{filtration provenance
charging} below.

\end{enumerate}

The predicates may overlap. A symmetry or order instruction may retain
its additive and multiplicative implementation tags. Pure tuple,
relabeling, and interface instructions expose their state-dependent
arguments and inherit the union of their argument provenance.  The same
opcode assignment applies to the represented derivation of an
instance-dependent coefficient when that coefficient is used by a structural
source record.  Public constants independent of both the represented state
and the presented instance remain untagged coefficients.  Thus a coefficient
computed from public instance data retains the provenance and complete cost
of its computation; declaring its output state-independent does not erase
that ancestry.

For a normalized state-dependent vertex, or for a normalized
instance-dependent coefficient vertex used by a structural source, write

\[
\operatorname{Prov}_5(v)
\subseteq\mathfrak J_{\Abs}^5
\]

for its nonempty set of authenticated opcode tags. For a pure interface
vertex, \(\operatorname{Prov}_5(v)\) is the union of the tags of its
exposed state-dependent arguments.

For an expanded derivation \(\widehat\rho\), its five-family expansion
\(\operatorname{Exp}_5(\widehat\rho)\) replaces every finite
state-dependent evaluator by the charged operational implementation
provided by evaluator normalization. Boolean gates use the affine and
product basis over \(\mathbb F_2\); finite-precision arithmetic,
comparison, lookup, and memory access use their charged Boolean
implementations. Represented symmetry instructions retain their
authenticated action data. Bounded stateful composition retains its
clock, work state, outcome record, and transcript under
\(\mathsf{Filt}\).

\end{definition}

The five-family normal form exhausts the represented provenance of structural
source evaluators.
\begin{theorem}[Budgeted five-family evaluator-provenance normal form]
\label{thm-five-family-abs-provenance}

For every finite task presentation under the budgeted derivability
closure and every oracle-free derivation
\(\rho\in\mathsf{Der}_{\mathfrak I}\), the record
\(\operatorname{Exp}_5(\operatorname{Exp}(\rho))\) is a uniform
derivation with the same denotation. Its cost is at most the
class-preserving overhead fixed by the encoded operational model:

\[
\operatorname{Cost}_I
  \bigl(\operatorname{Exp}_5(\operatorname{Exp}(\rho))\bigr)
\preceq
\Psi_5\bigl(\operatorname{Cost}_I(\rho)\bigr)
\]

on every instance.

Every state-dependent instruction in this normal form has provenance in
at least one member of \(\mathfrak J_{\Abs}^5\), or is a pure
interface exposing earlier tagged arguments. In particular, if the
terminal denotation is a represented quotient \(q\), its backward slice
admits a represented joint frontier \(J_5\) and a charged postprocessing
\(h_5\) such that

\[
q=h_5\circ J_5,
\qquad
\sigma(q)\subseteq\sigma(J_5).
\]

Consequently, for every centered target contrast \(y\),

\[
\|P_qy\|_2^2\le\|P_{J_5}y\|_2^2.
\]

Every coordinate of \(J_5\) is a declared public state-dependent
primitive or the output of a prior charged represented constructor.
Thus the quotient visibility is carried by the public derivation ancestry.
The represented joint frontier retains the complete state-dependent and
instance-coefficient ancestry with every applicable provenance tag.

The same charged normalization applies to budget-generated transcript and
exact first-hit quotients already present in the saturated execution class.
It neither constructs a new evaluator nor certifies target alignment.

\end{theorem}

\begin{proof}

Fix the represented derivation \(\rho\).  Evaluator normalization supplies a
charged finite implementation of each state-dependent primitive and
evaluator occurring in that derivation.  Affine gates and multiplication
over \(\mathbb F_2\) form a complete Boolean gate basis because conjunction
is multiplication and negation is addition of the constant one.  This is a
coding fact about the implementation of \(\rho\), not an admission rule for
other Boolean maps.  The finite encoded arithmetic, comparison, selection,
or table lookup retained by \(\rho\) compiles to the basis with its complete
implementation cost; in particular, an explicit table retains its encoded
length and storage cost.  The resulting gates have \(\mathsf{Add}\),
\(\mathsf{Mult}\), or \(\mathsf{Ord}\) provenance.  Authenticated
group-action instructions also carry \(\mathsf{Sym}\), and every stateful
execution record carries \(\mathsf{Filt}\).  The encoded implementation and
evaluator-normalization clauses give the class-preserving cost map
\(\Psi_5\).

Expand interface instructions until their state-dependent arguments are
exposed. The terminal backward slice is a finite derivation DAG, so its
tagged state-dependent vertices form a finite represented joint
frontier. Topological contraction of the remaining downstream record
gives \(h_5\) and the factorization \(q=h_5\circ J_5\). The sigma-algebra
inclusion follows. The target-projection inequality is
the \hyperref[thm-affordable-composition-no-creation]{affordable
composition and target-projection theorem}. Computation
closure and the exact first-hit construction provide derivation records
of the same form for finite-horizon algorithms.

\end{proof}

\Cref{fig-five-family-provenance} shows the task-independent charged
provenance frontier retained by a budgeted exact quotient evaluator.

\begin{figure}[tbp]
\centering
\begin{tikzpicture}[fw diagram,node distance=5mm]
  \node[fw provenance,minimum width=1.75cm] (add) {\(\mathsf{Add}\)};
  \node[fw provenance,minimum width=1.75cm,right=of add] (mult) {\(\mathsf{Mult}\)};
  \node[fw provenance,minimum width=1.75cm,right=of mult] (ord) {\(\mathsf{Ord}\)};
  \node[fw provenance,minimum width=1.75cm,right=of ord] (sym) {\(\mathsf{Sym}\)};
  \node[fw provenance,minimum width=1.75cm,right=of sym] (filt) {\(\mathsf{Filt}\)};
  \node[fw box,below=12mm of ord,minimum width=2.2cm] (joint) {joint frontier \(J_5\)};
  \node[fw use,right=10mm of joint] (post) {\(h_5\)};
  \node[fw source,right=10mm of post] (quot) {quotient \(q\)};
  \foreach \n in {add,mult,ord,sym,filt}
    \draw[fw arrow] (\n) -- (joint);
  \draw[fw arrow] (joint) -- (post);
  \draw[fw arrow] (post) -- (quot);
\end{tikzpicture}
\caption{Every exact quotient factors through a charged joint frontier
whose state-dependent instructions have the provenance classified by
\cref{thm-five-family-abs-provenance}.}
\label{fig-five-family-provenance}
\end{figure}

\subsection{Exhaustive source provenance for valid geometry}
\label{subsec-exhaustive-geom-source-provenance}
\label{subsec:exhaustive-geom-source-provenance}

The geometric normal form classifies the carrier on which a valid geometry
acts.  The normalized derivation DAG classifies every represented source from
which its active geometric data are generated.

\begin{definition}[Active Geom ancestry, charged cut, and provenance cells]
\label{def-active-geom-source-frontier}

Let
\[
\mathfrak S_{\Geom}^{4}
=
\{\mathsf X,\mathsf Q_{\mathrm{ex}},
  \mathsf Q_{\mathrm{rev}},\mathsf Q_{\mathrm{nr}}\}
\]
denote the four active-geometry support tags of
\cref{thm-exhaustive-geom-normal-form}: ambient identity support, exact
quotient support, reversible compensated quotient support, and
general-resolvent compensated quotient support.

Let \(\mathscr C\) be a dynamically qualified represented source record in
one of the four configurations of
\cref{thm-exhaustive-geom-normal-form}, with active geometry and complete
saturated cost at most \(B\).  Its \emph{active Geom record} is the
represented tuple containing:

\begin{enumerate}[label=\textup{(\roman*)}]
\item the normalized kernel, selected generator, reference gauge, and
      initialization in \(\mathfrak N\);
\item the exposed relation and conductance of \(\mathfrak G\);
\item every potential, local-difference evaluator, target-reach evaluator,
      killed-gap or mixing comparison, regularity comparison, and
      same-generator local-consistency record used by its visibility;
\item the selected current and its complete authority record when directed
      progress is invoked; and
\item for nonidentity support, the quotient, represented fibers, quotient
      dynamics, defect ledger, target-functional comparison, and selected
      exact, reversible, or general stability record.
\item every represented qualification indicator used by the source.  Such an
      indicator is either identically one for an already-qualified record or
      is retained as an active pre-hit gate together with its complete
      represented derivation and evaluation cost.
\end{enumerate}

Write this complete tuple as
\(\operatorname{Act}_{\Geom}(\mathscr C)\).  Its entries are conjunctive
validity data: a qualified geometry must pass every applicable gate.  The
classification below does not replace this joint record by any single gate.

First expand every presentation-admissible oracle interface into its charged
family-uniform represented evaluator.  Then apply
\(\operatorname{Exp}_5(\operatorname{Exp}(\rho_{\mathscr C}))\) to one
complete derivation of this tuple.  Adjoin a pure tuple-interface vertex whose
arguments are precisely the active fields above, and take the backward slice
of that vertex.  After pure interfaces expose their arguments, let
\(\mathcal A_{\Geom,5}(\mathscr C)\) be the encoded-order list of every
tagged non-interface vertex in this finite backward slice: every
state-dependent vertex and every instance-dependent coefficient vertex whose
value is an ancestor of an active field or gate.  Thus the list contains both
primitive leaves and the tagged internal constructors that combine them.  A
public constant independent of the instance is omitted.  This is the
\emph{complete active Geom ancestral set}.

Let \(t_{\mathscr C}\) be the adjoined active-tuple interface vertex and set
\[
\Gamma_{\Geom,5}(\mathscr C)=(t_{\mathscr C}),
\qquad
J_{\Geom,5}(\mathscr C)=\operatorname{Act}_{\Geom}(\mathscr C).
\]
This singleton is an antichain and meets every path terminating at the
active-tuple vertex.  It is the active-tuple analogue of the permitted
singleton terminal cut in \cref{def-charged-backward-source-cut}.  The cut
supplies factorization and cost accounting; the complete ancestral set
supplies provenance.  Keeping these two objects separate prevents an
internal constructor from disappearing from the signature while preserving
the antichain and separation properties required of a cut.

The \emph{five-family source signature} of \(\mathscr C\) is
\[
\operatorname{Prov}_{\Geom,5}(\mathscr C)
=
\bigcup_{v\in\mathcal A_{\Geom,5}(\mathscr C)}
\operatorname{Prov}_5(v)
\subseteq\mathfrak J_{\Abs}^{5}.
\]
It is nonempty.  Mixed source derivations retain every applicable tag.

For \(\sigma\in\mathfrak S_{\Geom}^{4}\) and
\(c\in\mathfrak J_{\Abs}^{5}\), let
\[
G_{\sigma,c,n}^{\mathrm{src,sat}}(B)
\]
be the restriction of the existing prospective Geom profile to records whose
selected master support tag is \(\sigma\) and whose source signature contains
\(c\), with supremum \(0\) for an empty class.  These quantities are
provenance restrictions of \(G_n^{\mathrm{src,sat}}\); they introduce no new
structural channel.

Let
\[
\mathfrak P_{\Geom,5}
=
2^{\mathfrak J_{\Abs}^{5}}\setminus\{\varnothing\}
\]
be the thirty-one nonempty exact provenance signatures.  For
\(S\in\mathfrak P_{\Geom,5}\), let
\[
G_{\sigma,S,n}^{\mathrm{src,sat},=}(B)
\]
be the restriction to records with selected support tag \(\sigma\) and
canonical source signature exactly \(S\).  Exact-signature restrictions are
pairwise disjoint.  Their union is the complete normalized prospective Geom
record family.  The superscript \(=\) distinguishes exact signatures from the
overlapping single-tag restrictions above.

\end{definition}

\begin{theorem}[Exhaustive five-family source classification of valid Geom]
\label{thm-exhaustive-five-family-geom-source-classification}

For every presented finite task family, every temporally admissible,
dynamically qualified Geom source, including one using
presentation-admissible oracle interfaces, has:

\begin{enumerate}[label=\textup{(\roman*)}]
\item exactly one selected support tag
      \(\sigma\in\mathfrak S_{\Geom}^{4}\);
\item a nonempty source signature
      \[
      \operatorname{Prov}_{\Geom,5}(\mathscr C)
      \subseteq
      \{\mathsf{Add},\mathsf{Mult},\mathsf{Ord},
        \mathsf{Sym},\mathsf{Filt}\};
      \]
\item a charged factorization of its complete active tuple through its joint
      source frontier,
      \begin{equation}
      \label{eq-geom-active-frontier-factorization}
      \operatorname{Act}_{\Geom}(\mathscr C)
      =
      h_{\mathscr C}\circ J_{\Geom,5}(\mathscr C),
      \end{equation}
      where the downstream record \(h_{\mathscr C}\) retains public constants
      independent of the instance, all active gate evaluations, and every
      comparison listed in
      \cref{def-active-geom-source-frontier}; and
\item a class-preserving complete-cost bound
      \begin{equation}
      \label{eq-geom-active-frontier-cost}
      \operatorname{Cost}_I
      \bigl(\rho_{\Gamma_{\Geom,5}(\mathscr C)}\bigr)
      \preceq
      \Psi_{\Geom,5}(B),
      \end{equation}
      where \(\Psi_{\Geom,5}\) is the five-family expansion cost together
      with the charged represented join interface.
\end{enumerate}

Consequently,

\begin{equation}
\label{eq-geom-five-family-envelope-cover}
G_n^{\mathrm{src,sat}}(B)
\le
\max_{\substack{\sigma\in\mathfrak S_{\Geom}^{4}\\
                 c\in\mathfrak J_{\Abs}^{5}}}
G_{\sigma,c,n}^{\mathrm{src,sat}}
  \bigl(\Psi_{\Geom,5}(B)\bigr)
\le
G_n^{\mathrm{src,sat}}
  \bigl(\Psi_{\Geom,5}(B)\bigr).
\end{equation}

Thus every nonnegligible valid Geom source at a budget in the transfer class
produces a nonnegligible record in at least one of the twenty
support--provenance restrictions
\[
\mathfrak S_{\Geom}^{4}\times\mathfrak J_{\Abs}^{5}.
\]
The support coordinate is unique.  The provenance restrictions may overlap,
because a mixed derivation retains every applicable five-family tag.
Negligibility of these twenty provenance-restricted envelopes throughout a
class-preserving budget closure implies negligibility of the complete
prospective Geom profile.

Authorized \(\Caus\) orientation does not add a source cell: its current and
authority comparison occur in the active tuple and retain the source
signature of that tuple.  Valid \(\Lift\) and \(\Bdry\) records retain the
same source assignment.  A represented record failing the Geom qualification
gates contributes zero to every cell, while its post-saturation
gradient-orthogonal complement belongs to \(\Res\).

\end{theorem}

\begin{proof}

\Cref{thm-exhaustive-geom-normal-form} gives the four active-geometry support
configurations.  Selecting one exact, reversible, or general comparison
record gives the disjoint master support tag of
\cref{thm-affordable-geom-abs-normal-form}; hence \(\sigma\) is unique.

\Cref{thm-oracle-admissibility-presentation-separation} expands every
class-preserving oracle interface into its represented evaluator without
changing the presented task family.  Evaluator normalization and
\cref{thm-expanded-derivation-equivalence,thm-five-family-abs-provenance}
expand every finite state-dependent evaluator in the complete active record
without changing its denotation.  Each exposed state-dependent vertex has at
least one tag in \(\mathfrak J_{\Abs}^{5}\).  The active tuple contains a
state-dependent potential, conductance, kernel, current, quotient, or
comparison whenever its visibility is nonzero, so its source signature is
nonempty.

The active backward slice is a finite derivation DAG.  Its terminal
active-tuple vertex is a singleton separating antichain, and its represented
output is \(J_{\Geom,5}(\mathscr C)=\operatorname{Act}_{\Geom}(\mathscr C)\).
Taking \(h_{\mathscr C}\) to be the represented identity on that tuple proves
\cref{eq-geom-active-frontier-factorization}.  The complete ancestral set
\(\mathcal A_{\Geom,5}(\mathscr C)\), rather than the cut alone, retains the
construction of every tagged internal vertex for provenance.  The ancestral
records and the downstream active tuple are subrecords of the original
expanded derivation.  Only the represented join interface adds a constructor
charge.  Evaluator normalization and computation closure therefore give
\cref{eq-geom-active-frontier-cost}, with a class-preserving
\(\Psi_{\Geom,5}\).  This is the same singleton-terminal cost argument as in
\cref{thm-quantitative-backward-source-factorization}, applied to the
represented active tuple rather than a quotient-map output.

Every qualified prospective Geom record of cost at most \(B\) is consequently
present, with unchanged visibility, in at least one restricted cell at the
enlarged cost.  Taking its support tag and any member of its nonempty source
signature proves the first inequality in
\cref{eq-geom-five-family-envelope-cover}.  Each restricted envelope is a
subfamily of the complete Geom profile at the same enlarged budget, proving
the second inequality.  The index set has \(4\cdot5=20\) elements, so the
negligibility conclusion follows from closure of the transfer class under
finite maxima.

Finally,
\cref{prop-derived-channels-create-no-structure,thm-caus-authority-no-substitution,lem-exact-lift-identity}
preserve the source assignment under authorized orientation, valid lifting,
and boundary readout.  Component exhaustiveness and the residual normal form
\cref{thm-channel-exhaustiveness,thm-exhaustive-residual-normal-form} assign
the remaining gradient-orthogonal component to \(\Res\).

\end{proof}

\begin{theorem}[Canonical exact-signature decomposition of valid Geom]
\label{thm-canonical-exact-signature-geom-decomposition}

For every presented finite task family and saturated budget \(B\), the
canonical coefficient-complete ancestral set of
\cref{def-active-geom-source-frontier} assigns every temporally admissible,
dynamically qualified Geom source to exactly one cell
\[
(\sigma,S)
\in
\mathfrak S_{\Geom}^{4}\times\mathfrak P_{\Geom,5}.
\]
The support tag \(\sigma\) records ambient identity, exact quotient,
reversible compensated quotient, or general-resolvent compensated quotient
support.  The exact signature \(S\) records all opcode families occurring in
the active state-dependent inputs and in the represented derivations of the
instance-dependent coefficients used by the active tuple.  In particular,
computing a center, weight, ordering, conductance, kernel parameter, current
parameter, or authority comparison from public instance data cannot be
hidden in an untagged coefficient.

At the class-preserving normalized ceiling
\(\widehat B=\Psi_{\Geom,5}(B)\),
\begin{equation}
\label{eq-exact-signature-geom-decomposition}
G_n^{\mathrm{src,sat}}(B)
\le
\max_{\substack{\sigma\in\mathfrak S_{\Geom}^{4}\\
                 S\in\mathfrak P_{\Geom,5}}}
G_{\sigma,S,n}^{\mathrm{src,sat},=}(\widehat B)
\le
G_n^{\mathrm{src,sat}}(\widehat B).
\end{equation}
For the normalized record family at a fixed ceiling, the first comparison is
an equality.  Thus the Geom source family is indexed by
\(4(2^5-1)=124\) support--provenance normal forms, with empty cells assigned
value zero and at most 124 cells nonempty.  Authorized
\(\Caus\) is the gradient projection inside the selected Geom cell; its
gradient-orthogonal complement is \(\Res\) and is absent from all 124 cells.

\end{theorem}

\begin{proof}

The four support alternatives are disjoint by
\cref{thm-exhaustive-geom-normal-form,thm-affordable-geom-abs-normal-form}.
Evaluator normalization assigns a nonempty subset of
\(\mathfrak J_{\Abs}^{5}\) to every tagged non-interface vertex of the
coefficient-complete active ancestral DAG.  Taking their union gives one and
only one exact signature \(S\).
The active backward slice includes the represented derivation of every
instance-dependent coefficient used in an active field or gate, so
contraction leaves only instance-independent constants untagged.  This proves
existence and uniqueness of \((\sigma,S)\).

Expansion preserves denotation and has cost at most \(\widehat B\) by
\cref{thm-expanded-derivation-equivalence,thm-five-family-abs-provenance,eq-geom-active-frontier-cost}.
The exact-signature cells are disjoint and cover the normalized record
family.  A supremum over a finite disjoint union equals the maximum of the
cell suprema, proving
\cref{eq-exact-signature-geom-decomposition}.  The final current assignment
is
\cref{prop-derived-channels-create-no-structure,thm-exhaustive-residual-normal-form}.

\end{proof}

\begin{theorem}[No unclassified Geom source remains after quotient
assignment]
\label{thm-no-unclassified-geom-source-after-quotient-assignment}

Fix a presented finite task family and a saturated ceiling \(B\), and put
\(\widehat B=\Psi_{\Geom,5}(B)\).  At the normalized ceiling, the complete
prospective Geom profile is the exact finite maximum
\begin{equation}
\label{eq-no-unclassified-geom-source-finite-maximum}
G_n^{\mathrm{src,sat}}(\widehat B)
=
\max_{\substack{\sigma\in\mathfrak S_{\Geom}^{4}\\
                 S\in\mathfrak P_{\Geom,5}}}
G_{\sigma,S,n}^{\mathrm{src,sat},=}(\widehat B).
\end{equation}
Consequently its quotient-free remainder is exactly
\begin{equation}
\label{eq-quotient-free-geom-is-exact-ambient-maximum}
G_{X,n}^{\mathrm{src,sat}}(\widehat B)
=
\max_{S\in\mathfrak P_{\Geom,5}}
G_{\mathsf X,S,n}^{\mathrm{src,sat},=}(\widehat B).
\end{equation}

Every record on the right-hand side has identity support, a nonempty exact
signature drawn from
\[
\{\mathsf{Add},\mathsf{Mult},\mathsf{Ord},
  \mathsf{Sym},\mathsf{Filt}\},
\]
and one coefficient-complete represented ancestral DAG containing the
potential, geometry, transition rule, current, authority comparison, target
readout, every instance-dependent coefficient, and their complete costs.
Authorized \(\Caus\), valid \(\Lift\), and \(\Bdry\) operations inherit this
same source cell and create no additional one.  The gradient-orthogonal
component is \(\Res\) and is absent from the prospective Geom family.

Thus a represented computation contributes prospective geometric structure
only through a quotient-supported exact or compensated record, or through
one of the ambient identity-support cells in
\cref{eq-quotient-free-geom-is-exact-ambient-maximum}.  There is no additional
algorithmic, coefficient, mixed-operation, orientation, spectral, or
postprocessing Geom source outside this decomposition.

\end{theorem}

\begin{proof}

\Cref{thm-affordable-geom-abs-normal-form} gives the four pairwise disjoint
support tags: ambient identity support, exact quotient support, reversible
compensated quotient support, and general-resolvent compensated quotient
support.  \Cref{thm-canonical-exact-signature-geom-decomposition} assigns
every normalized qualified record exactly one nonempty coefficient-complete
five-family signature \(S\), including the derivations of all
instance-dependent coefficients.  Their product indexing gives the disjoint
record classification whose profile form is
\cref{eq-no-unclassified-geom-source-finite-maximum}.  Restricting that
classification to the ambient tag \(\mathsf X\) gives
\cref{eq-quotient-free-geom-is-exact-ambient-maximum}.

Source inheritance under authorized orientation, lifting, and boundary
readout is
\cref{prop-derived-channels-create-no-structure,thm-caus-authority-no-substitution,%
lem-exact-lift-identity}.  The orthogonal assignment is
\cref{thm-exhaustive-residual-normal-form}.  These results exhaust support,
coefficient provenance, temporal source status, and derived-channel status,
which proves the final assertion.

\end{proof}

\section{State--Coefficient Indexing of the Saturated Source Profile}
\label{subsec-universal-state-coefficient-source-provenance}
\label{sec:state-coefficient-source-profile}

The master source normal form admits an exact finite reindexing by the
provenance already present in each complete represented record.  The indices
below add neither a source channel nor an admissibility predicate or
numerical visibility factor.

\subsection{State--coefficient source coordinates}
\label{subsec:state-coefficient-source-coordinates}

\begin{definition}[State--coefficient provenance and restricted source-profile
notation]
\label{def-state-coefficient-source-provenance}

This definition fixes the ancestral signatures and their profile cells.  The
same-record interaction coordinate attached to those cells is defined
separately in \cref{def:state-coefficient-structural-interaction}.

Let \(\mathscr C\) be a coefficient-complete normalized record of one of two
temporal types: a prospective source in one of the five source configurations
of \cref{thm-complete-generic-source-theory}, subsequently formalized as the
five master families in
\cref{def-master-saturated-quotient-geometry-profile}, or a completed
terminal-compatible exact quotient retained only for accounting.  Its
\emph{complete ancestral set} (with the historical notation)
\(\mathcal A_{\mathrm{src},5}(\mathscr C)\) is defined as follows.

\begin{enumerate}[label=\textup{(\roman*)}]
\item For a pure exact-\(\Abs\) source record or a completed exact-quotient
      accounting record, expand every
      presentation-admissible oracle interface and then apply
      \(\operatorname{Exp}_5\circ\operatorname{Exp}\) to the complete
      backward slice of the terminal quotient.  After pure interfaces expose
      their arguments, define
      \(\mathcal A_{\mathrm{src},5}(\mathscr C)\) to be the encoded-order
      list of every tagged non-interface vertex in the complete ancestral DAG
      of that quotient, separated into state-dependent vertices and
      instance-dependent coefficient vertices.
\item For any of the four active-\(\Geom\) records, set
      \[
      \mathcal A_{\mathrm{src},5}(\mathscr C)
      =\mathcal A_{\Geom,5}(\mathscr C),
      \]
      with \(\mathcal A_{\Geom,5}\) as in
      \cref{def-active-geom-source-frontier}.
\end{enumerate}

The same syntactic ancestral-set construction applies to a completed
terminal-compatible exact quotient retained only for accounting.  In that
use the notation \(\mathcal A_{\mathrm{src},5}\) is historical: it records
normalized ancestry and does not assign prospective-source status.  The
record's source/accounting temporal type is retained separately throughout.

In both cases, an instance-independent public constant is omitted, whereas
the complete represented derivation of every instance-dependent coefficient
is retained.  The encoded order and the applicable master family make this
ancestral set canonical for the normalized record.  The applicable terminal
singleton---the quotient-map vertex in the pure exact-\(\Abs\) mode and the
active-tuple vertex in an active-\(\Geom\) mode---is the associated canonical
charged cut.  The terminal cut is used for factorization, whereas the
complete ancestral set is used for provenance.

Split \(\mathcal A_{\mathrm{src},5}(\mathscr C)\) into its represented
state-dependent vertices
\(\mathcal A_{\mathrm{src},5}^{\mathrm{st}}(\mathscr C)\) and its represented
instance-dependent coefficient vertices
\(\mathcal A_{\mathrm{src},5}^{\mathrm{cf}}(\mathscr C)\).  Define
\begin{align}
\label{eq-state-source-exact-signature}
S_{\mathrm{st}}(\mathscr C)
&=
\bigcup_{v\in\mathcal A_{\mathrm{src},5}^{\mathrm{st}}(\mathscr C)}
\operatorname{Prov}_5(v),\\
\label{eq-coefficient-source-exact-signature}
S_{\mathrm{cf}}(\mathscr C)
&=
\bigcup_{v\in\mathcal A_{\mathrm{src},5}^{\mathrm{cf}}(\mathscr C)}
\operatorname{Prov}_5(v).
\end{align}
The state signature is nonempty because a prospective structural-source
record contains a state-dependent quotient or active geometric field, and
that tagged vertex belongs to the complete ancestral DAG.  The coefficient
signature may be empty.  In the active-\(\Geom\) modes their union is the exact signature of
\cref{thm-canonical-exact-signature-geom-decomposition}.  For
\[
\varnothing\ne S_{\mathrm{st}}\subseteq\mathfrak J_{\Abs}^{5},
\qquad
S_{\mathrm{cf}}\subseteq\mathfrak J_{\Abs}^{5},
\]
write
\[
G_{\sigma,S_{\mathrm{st}},S_{\mathrm{cf}},n}
 ^{\mathrm{src,sat},=}(B)
\]
for the restriction to records having precisely these two signatures and
support tag \(\sigma\).  When a family functional \(\mathfrak M_n\) is
declared, write
\[
G_{\sigma,S_{\mathrm{st}},S_{\mathrm{cf}},\mathfrak M,n}
 ^{\mathrm{src,sat},=}(B)
\]
for the supremum of
\(\mathfrak M_n(I\mapsto Q_I V_{\mathscr C,I})\) over the same refined
Geom cell, with the qualification gate included in the complete source
record as specified below.

\end{definition}

\begin{definition}[State--coefficient structural-interaction coordinate]
\label{def:state-coefficient-structural-interaction}

Fix a budget \(B\), a nonempty state signature
\(S_{\mathrm{st}}\subseteq\mathfrak J_{\Abs}^5\), a coefficient signature
\(S_{\mathrm{cf}}\subseteq\mathfrak J_{\Abs}^5\), and a monotone sublinear
positively homogeneous family functional \(\mathfrak M_n\).  The following
symbol abbreviates the restriction of the
existing prospective saturated source profile to one exact signature pair:
\begin{equation}
\label{eq-affordable-coefficient-interaction-coordinate}
\begin{aligned}
&\operatorname{StructAct}_{\mathfrak M,
S_{\mathrm{st}},S_{\mathrm{cf}},n}^{\mathrm{sat}}(B)\\
&\quad:=
\sup_{(\rho,\mathscr C,Q)}
\mathfrak M_n\!\left(
I\longmapsto Q_I V_{\mathscr C,I}
\right).
\end{aligned}
\end{equation}
Here \(\mathscr C\) ranges over the five temporally admissible prospective
source modes, \(Q_I\in\{0,1\}\) is its represented qualification indicator,
and \(\rho\) is the complete family-uniform represented derivation of the
qualified pair \((\mathscr C,Q)\), of saturated cost at most \(B\).  Either
\(Q\equiv1\), or \(Q\), its complete pre-hit derivation, storage, and
evaluation are an active gate of \(\mathscr C\).  In the latter case every
tagged ancestor of \(Q\) is adjoined to
\(\mathcal A_{\mathrm{src},5}(\mathscr C)\) before the exact signatures are
formed.  Thus no qualification predicate is external to the charged source
record.  The complete normalized ancestral DAG of the qualified pair has
exact state--coefficient signature
\((S_{\mathrm{st}},S_{\mathrm{cf}})\), and \(V_{\mathscr C,I}\) is the
complete source visibility, including every dynamic, target-utility,
compatibility, regularity, and temporal-source gate applicable to its master
mode.  The record \(\rho\) contains the derivation, storage, and evaluation
of every instance-dependent coefficient used by that complete source.  The
supremum is zero when this represented class is empty.  Thus
\(\operatorname{StructAct}\) is profile-restriction notation, not a new
structural coordinate.

This coordinate retains the complete same-record source product.  A public
target verifier or target-discriminating readout alone therefore contributes
no structural activity: it enters this coordinate only as part of a
prospective source record that passes all applicable gates.  A coefficient
supplied without a family-uniform derivation belongs to advice or to an
enlarged presentation.  An over-budget coefficient is absent from the budget
slice.

\end{definition}

\begin{theorem}[Exact state--coefficient decomposition and coefficient
accountability for Geom]
\label{thm-exact-state-coefficient-geom-accountability}

For every presented finite task family and saturated budget \(B\), every
temporally admissible dynamically qualified Geom source has exactly one
structural index
\[
(\sigma,S_{\mathrm{st}},S_{\mathrm{cf}}),
\qquad
\sigma\in\mathfrak S_{\Geom}^{4},
\qquad
\varnothing\ne S_{\mathrm{st}}\subseteq\mathfrak J_{\Abs}^{5},
\qquad
S_{\mathrm{cf}}\subseteq\mathfrak J_{\Abs}^{5}.
\]
At the coefficient-complete normalized ceiling
\(\widehat B=\Psi_{\Geom,5}(B)\),
\begin{equation}
\label{eq-exact-state-coefficient-geom-decomposition}
G_n^{\mathrm{src,sat}}(B)
\le
\max_{\substack{\sigma,S_{\mathrm{st}},S_{\mathrm{cf}}\\
S_{\mathrm{st}}\ne\varnothing}}
G_{\sigma,S_{\mathrm{st}},S_{\mathrm{cf}},n}
 ^{\mathrm{src,sat},=}(\widehat B)
\le
G_n^{\mathrm{src,sat}}(\widehat B).
\end{equation}
Let
\(G_{n,\operatorname{Norm}(B)}^{\mathrm{src,sat}}(\widehat B)\)
denote the supremum over the exact image of the cost-
\(B\) family under coefficient-complete normalization; a semicolon
\(\operatorname{Norm}(B)\) on a cell denotes the same restriction.
Disjointness gives
the literal equality
\[
G_{n,\operatorname{Norm}(B)}^{\mathrm{src,sat}}(\widehat B)
=
\max_{\substack{\sigma,S_{\mathrm{st}},S_{\mathrm{cf}}\\
S_{\mathrm{st}}\ne\varnothing}}
G_{\sigma,S_{\mathrm{st}},S_{\mathrm{cf}},n;
\operatorname{Norm}(B)}^{\mathrm{src,sat},=}(\widehat B).
\]
The inequalities in
\cref{eq-exact-state-coefficient-geom-decomposition} compare this exact
normalized image with the enclosing ambient budget slices.  Hence the
normalized Geom family is indexed by
\[
4(2^5-1)2^5=3968
\]
support--state-provenance--coefficient-provenance cells, with every empty
cell assigned value zero; therefore at most 3968 cells are nonempty.  Its coarser
124-cell decomposition is recovered exactly from
\begin{equation}
\label{eq-coarse-signature-from-state-coefficient-signatures}
G_{\sigma,S,n}^{\mathrm{src,sat},=}(B)
=
\max_{\substack{
\varnothing\ne S_{\mathrm{st}}\subseteq\mathfrak J_{\Abs}^{5}\\
S_{\mathrm{cf}}\subseteq\mathfrak J_{\Abs}^{5}\\
S_{\mathrm{st}}\cup S_{\mathrm{cf}}=S}}
G_{\sigma,S_{\mathrm{st}},S_{\mathrm{cf}},n}
 ^{\mathrm{src,sat},=}(B).
\end{equation}

Let \(\mathfrak M_n\) evaluate the same family of records.  There is a
class-preserving normalization overhead \(\Psi_D\) such that every refined cell
satisfies
\begin{equation}
\label{eq-geom-cell-dominated-by-coefficient-interaction}
G_{\sigma,S_{\mathrm{st}},S_{\mathrm{cf}},\mathfrak M,n}
 ^{\mathrm{src,sat},=}(B)
\le
\operatorname{StructAct}_{\mathfrak M,
S_{\mathrm{st}},S_{\mathrm{cf}},n}^{\mathrm{sat}}
\bigl(\Psi_D(B,n)\bigr).
\end{equation}
Thus an instance-dependent coefficient enters a Geom source only through a
represented affordable interaction inside the complete qualified source
record.  Its provenance, qualification event, full gate product, and complete
cost remain in the right-hand coordinate.

\end{theorem}

\begin{proof}

The coefficient-complete ancestral set of
\cref{def-active-geom-source-frontier} gives a disjoint partition of its
vertices into state-dependent and instance-dependent coefficient vertices.
Taking the union of the authenticated five-family tags in each part gives
\cref{eq-state-source-exact-signature,eq-coefficient-source-exact-signature}.
Both sets are uniquely determined by the normalized backward slice, and the
state signature is nonempty.  Together with the unique support tag from
\cref{thm-canonical-exact-signature-geom-decomposition}, this proves the
existence and uniqueness of the displayed structural index.

Expansion preserves denotation and has cost at most \(\widehat B\).
The refined indexed cells are disjoint and cover the normalized record
family; an empty indexed cell has supremum zero.  Hence a supremum over the
family is the maximum of their suprema.  This proves
\cref{eq-exact-state-coefficient-geom-decomposition}.  For a fixed coarse
signature \(S\), the refined cells with
\(S_{\mathrm{st}}\cup S_{\mathrm{cf}}=S\) are again a finite disjoint cover,
which proves
\cref{eq-coarse-signature-from-state-coefficient-signatures}.  There are
\(2^5-1\) choices for the nonempty state signature and \(2^5\) choices for
the coefficient signature; multiplication by the four support tags gives
3968 cells.

Let \(\mathscr C=(\mathfrak N,\Phi,D,J,\mathfrak K)\) be a record in one
refined cell and let \(Q_I\) be its qualification indicator.  The normalized
complete source record, including every active gate and every ancestral
coefficient derivation, is one of the records in
\cref{eq-affordable-coefficient-interaction-coordinate} at the
class-preserving normalization overhead \(\Psi_D(B,n)\).  Its exact
state--coefficient signature is unchanged.  Apply the monotone family
functional and take the supremum over the refined Geom cell to obtain
\cref{eq-geom-cell-dominated-by-coefficient-interaction}.

\end{proof}

\subsection{Universal source decomposition}
\label{subsec:universal-state-coefficient-decomposition}

\begin{theorem}[Universal state--coefficient decomposition of affordable
structural sources]
\label{thm-universal-state-coefficient-source-decomposition}

Let
\[
\mathfrak S_{\mathrm{src}}^{5}
=
\{\mathsf A,\mathsf X,\mathsf Q_{\mathrm{ex}},
  \mathsf Q_{\mathrm{rev}},\mathsf Q_{\mathrm{nr}}\}
\]
denote, respectively, pure exact \(\Abs\), ambient \(\Geom\), \(\Geom\) on
an exact quotient, reversible compensated quotient--geometry, and
general-resolvent compensated quotient--geometry.  These are the five
certificate families of
\cref{def-master-saturated-quotient-geometry-profile}.

There are class-preserving normalization and complete-source overheads
\(\Psi_{\mathrm{src},5}\), \(\Psi_{\mathrm{read}}\), and
\(\Psi_{\mathrm{all}}\), fixed by the encoded operational model.
The first expands admissible oracle interfaces, performs the five-family
evaluator normalization, and constructs the canonical ancestral set and its
charged cut.  The second
retains the complete qualified source product at the normalized ceiling, and
the third is a common class-preserving majorant of the finite compositions
used below.  Normalize
every temporally admissible prospective structural source at
\(\widehat B=\Psi_{\mathrm{src},5}(B)\).  For
\(\chi\in\mathfrak S_{\mathrm{src}}^{5}\) and
\(\varnothing\ne S_{\mathrm{st}}\subseteq\mathfrak J_{\Abs}^{5}\) and
\(S_{\mathrm{cf}}\subseteq\mathfrak J_{\Abs}^{5}\), let
\[
\mathfrak U_{\chi,S_{\mathrm{st}},S_{\mathrm{cf}},n}
 ^{\mathrm{src,sat},=}(\widehat B)
\]
be the restriction of the master source coordinate \(\chi\) to records with
the exact state--coefficient signature of
\cref{def-state-coefficient-source-provenance}.  When a family functional
\(\mathfrak M_n\) is declared, write
\(\mathfrak U_{\chi,S_{\mathrm{st}},S_{\mathrm{cf}},\mathfrak M,n}
 ^{\mathrm{src,sat},=}(B)\) for the supremum of
\(\mathfrak M_n(I\mapsto Q_I V_{\mathscr C,I})\) over that same cell, where
the qualification gate belongs to the complete charged pair as in
\cref{def:state-coefficient-structural-interaction}, and use
\(\mathfrak U_{\mathfrak M,n}^{\mathrm{src,sat}}\) and
\(I_{\mathfrak M,n}^{\mathrm{src,sat}}\) for the corresponding master and
combined evaluations.  Then
\begin{equation}
\label{eq-universal-state-coefficient-source-maximum}
\mathfrak U_n^{\mathrm{src,sat}}(B)
\le
\max_{\substack{
\chi\in\mathfrak S_{\mathrm{src}}^{5}\\
\varnothing\ne S_{\mathrm{st}}\subseteq\mathfrak J_{\Abs}^{5}\\
S_{\mathrm{cf}}\subseteq\mathfrak J_{\Abs}^{5}}}
\mathfrak U_{\chi,S_{\mathrm{st}},S_{\mathrm{cf}},n}
 ^{\mathrm{src,sat},=}(\widehat B)
\le
\mathfrak U_n^{\mathrm{src,sat}}(\widehat B).
\end{equation}
Let
\(\mathfrak U_{n,\operatorname{Norm}(B)}^{\mathrm{src,sat}}(\widehat B)\)
denote the supremum over the exact normalized image of the cost-
\(B\) record family; a semicolon \(\operatorname{Norm}(B)\) on a cell
denotes the same restriction.  The cell partition gives
\[
\mathfrak U_{n,\operatorname{Norm}(B)}^{\mathrm{src,sat}}(\widehat B)
=
\max_{\substack{\chi,S_{\mathrm{st}},S_{\mathrm{cf}}\\
S_{\mathrm{st}}\ne\varnothing}}
\mathfrak U_{\chi,S_{\mathrm{st}},S_{\mathrm{cf}},n;
\operatorname{Norm}(B)}^{\mathrm{src,sat},=}(\widehat B).
\]
The inequalities in
\cref{eq-universal-state-coefficient-source-maximum} compare this exact
image with its enclosing ambient slices.  Thus that normalized prospective
source family is
indexed by
\[
5(2^5-1)2^5=4960
\]
source-mode--state-provenance--coefficient-provenance cells, with empty
cells assigned value zero; hence at most 4960 cells are nonempty.

For every monotone sublinear positively homogeneous family functional
\(\mathfrak M_n\), every cell satisfies
\begin{equation}
\label{eq-universal-source-cell-structural-interaction-bound}
\mathfrak U_{\chi,S_{\mathrm{st}},S_{\mathrm{cf}},\mathfrak M,n}
 ^{\mathrm{src,sat},=}(B)
\le
\operatorname{StructAct}_{\mathfrak M,
S_{\mathrm{st}},S_{\mathrm{cf}},n}^{\mathrm{sat}}
\bigl(\Psi_{\mathrm{read}}(B,n)\bigr).
\end{equation}
Consequently,
\begin{align}
\label{eq-universal-master-structural-interaction-bound}
\mathfrak U_{\mathfrak M,n}^{\mathrm{src,sat}}(B)
&\le
\max_{\substack{
\varnothing\ne S_{\mathrm{st}}\subseteq\mathfrak J_{\Abs}^{5}\\
S_{\mathrm{cf}}\subseteq\mathfrak J_{\Abs}^{5}}}
\operatorname{StructAct}_{\mathfrak M,
S_{\mathrm{st}},S_{\mathrm{cf}},n}^{\mathrm{sat}}
\bigl(\Psi_{\mathrm{all}}(B,n)\bigr),\\
\label{eq-universal-combined-structural-interaction-bound}
I_{\mathfrak M,n}^{\mathrm{src,sat}}(B)
&\le
2
\max_{\substack{
\varnothing\ne S_{\mathrm{st}}\subseteq\mathfrak J_{\Abs}^{5}\\
S_{\mathrm{cf}}\subseteq\mathfrak J_{\Abs}^{5}}}
\operatorname{StructAct}_{\mathfrak M,
S_{\mathrm{st}},S_{\mathrm{cf}},n}^{\mathrm{sat}}
\bigl(\Psi_{\mathrm{all}}(B,n)\bigr).
\end{align}

\(\Caus\), valid \(\Lift\), and \(\Bdry\) retain the cell of the source
record that they use.  Their represented currents, lifted coordinates, and
boundary readouts remain in its complete derivation.  The post-saturation
\(\Res\) coordinate contains no affordable target-bearing source and
therefore contributes no cell to
\cref{eq-universal-state-coefficient-source-maximum}.

\end{theorem}

\begin{proof}

The five master certificate families form an exhaustive disjoint partition by
\cref{thm-complete-generic-source-theory,thm-affordable-geom-abs-normal-form}
and the selected certificate record in
\cref{def-master-saturated-quotient-geometry-profile}.  In particular, the
classification is made by the displayed source carrier and stability record,
not by the numerical value or eventual success of the certificate.
In the pure exact-\(\Abs\) family,
\cref{thm-oracle-admissibility-presentation-separation} first expands every
presentation-admissible oracle interface, and the coefficient-complete
evaluator expansion of \cref{thm-five-family-abs-provenance} then applies to
the target-bearing quotient readout.  The complete-ancestry expansion of
\cref{thm-exact-state-coefficient-geom-accountability} applies to the other
four families.  The encoded backward slice and the state--coefficient type of
each retained vertex are fixed by
\cref{def-state-coefficient-source-provenance}.  Hence every normalized
record has one unique state signature and one unique coefficient signature.
The terminal quotient in the pure exact-\(\Abs\) mode, and the complete active
tuple in each active-\(\Geom\) mode, contains a state-dependent source vertex.
By \cref{thm-five-family-abs-provenance}, that state-dependent vertex has at
least one five-family tag, so \(S_{\mathrm{st}}\ne\varnothing\).
Intersecting the five master families with these disjoint signature pairs,
and assigning value zero to every empty intersection, proves
\cref{eq-universal-state-coefficient-source-maximum}.  The number of cells is
the product of five source modes, \(2^5-1\) nonempty state signatures, and
\(2^5\) coefficient signatures.  The oracle expansion, evaluator
normalization, backward slice, and
represented join have the declared class-preserving combined overhead
\(\Psi_{\mathrm{src},5}\).

For the quantitative domination, normalize a record in a fixed refined cell
and retain its qualification indicator and complete source visibility.  The
resulting record belongs to the class defining
\cref{eq-affordable-coefficient-interaction-coordinate} with the same exact
state--coefficient signature.  Application of \(\mathfrak M_n\) and the
supremum over the fixed refined cell prove
\cref{eq-universal-source-cell-structural-interaction-bound}.  Taking the
finite maximum and using the common overhead majorant
\(\Psi_{\mathrm{all}}\) gives
\cref{eq-universal-master-structural-interaction-bound}.  The source-profile
comparison
\cref{eq-source-master-combined-equivalence} gives
\cref{eq-universal-combined-structural-interaction-bound}.

Finally,
\cref{prop-derived-channels-create-no-structure,lem-exact-lift-identity}
preserve source assignment under orientation, valid lifting, and boundary
readout.  The residual assertion is
\cref{thm-exhaustive-residual-normal-form,thm-no-uniform-residual-correlation}.

\end{proof}

\subsection{Boolean-complete mixed ancestry}
\label{subsec:boolean-complete-mixed-ancestry}

\begin{theorem}[Affine/product evaluator normalization and charge
conservation in mixed ancestral cells]
\label{thm-boolean-complete-mixed-ancestral-cells}

Fix a presented family with finite binary encodings and a saturated budget
\(B\).  Consider a normalized prospective source record whose active
state-dependent fields or instance-dependent coefficients are evaluated from
the declared binary interfaces by a finite represented Boolean circuit.
Evaluator normalization gives the same denotation at cost at most
\(\Psi_5(B)\) using affine gates and product gates over \(\mathbb F_2\).
The complete ancestral set of
\cref{def-state-coefficient-source-provenance} assigns every active affine
gate the \(\mathsf{Add}\) tag and every active product gate the
\(\mathsf{Mult}\) tag.

Consequently, the exact state--coefficient classification has the following
structural alternatives.

\begin{enumerate}[label=\textup{(\roman*)}]
\item A record whose complete active ancestry uses only affine gates and has
      no additional authenticated opcode lies in a pure-\(\mathsf{Add}\)
      signature cell.  If an authenticated non-Boolean opcode is also
      present, its tag remains and the full record lies in the corresponding
      mixed cell.
\item A record with an active product gate lies in a mixed signature cell
      containing \(\mathsf{Mult}\).  If its declared state interface has
      \(\mathsf{Add}\) provenance, then
      \[
      \mathsf{Add}\in S_{\mathrm{st}},
      \qquad
      \mathsf{Mult}\in S_{\mathrm{st}}\cup S_{\mathrm{cf}}.
      \]
\item Coordinate access, addition of the constant one, and multiplication
      implement XOR, negation, and conjunction.  They therefore normalize
      the gate syntax of each finite Boolean circuit already carried by the
      represented record.  The union is taken over the exact cells occupied
      by those budgeted records.  Expressivity of the affine/product basis
      does not establish membership of an unrepresented map in a budgeted
      cell.
      Authenticated comparison, symmetry, and stateful execution retain
      their \(\mathsf{Ord}\), \(\mathsf{Sym}\), and \(\mathsf{Filt}\) tags
      in addition to the Boolean implementation tags.
\end{enumerate}

For every exact signature pair, its numerical coordinate remains the
complete same-record quantity
\[
\operatorname{StructAct}_{\mathfrak M,
S_{\mathrm{st}},S_{\mathrm{cf}},n}^{\mathrm{sat}}(B)
=
\sup_{(\rho,\mathscr C,Q)}
\mathfrak M_n\!\left(I\longmapsto Q_I V_{\mathscr C,I}\right)
\]
from \cref{eq-affordable-coefficient-interaction-coordinate}.  Zero
incremental contextual charge of an internal deterministic constructor
preserves this complete coordinate by attributing the target projection to
its charged joint ancestry; the
qualification gate, source utility, and all geometric or quotient gates
remain in \(Q_I V_{\mathscr C,I}\).  Thus the mixed cells give the exhaustive
operational provenance of the represented Boolean evaluator.  Product gates
record implementation and cost; they are not independent structural
sources.  The numerical value is the complete qualified source product in
the existing master mode, and independent source attribution is performed by the
charge-bearing core of \cref{def-charge-bearing-core-signature}.
In particular, the gate identities neither construct a circuit nor place its
denotation in the budget slice.  Membership still requires the complete
uniform derivation and cost prescribed by
\cref{def-budgeted-saturated-observable-closure,def-budgeted-derivability-closure}.

\end{theorem}

\begin{proof}

By
\cref{thm-expanded-derivation-equivalence,thm-five-family-abs-provenance},
evaluator normalization preserves denotation and retains the charged
operational implementation.  Over \(\mathbb F_2\),
\[
u\mathbin{\mathrm{XOR}}v=u+v,
\qquad
\mathbin{\mathrm{NOT}}u=u+1,
\qquad
u\mathbin{\mathrm{AND}}v=uv.
\]
These identities show that the affine/product basis normalizes the already
represented finite Boolean circuit.  The first two gates have
\(\mathsf{Add}\) provenance and the third has \(\mathsf{Mult}\) provenance.
Because \(\mathcal A_{\mathrm{src},5}(\mathscr C)\) contains every tagged
non-interface vertex in the complete normalized ancestral DAG, an active
product gate and every additional authenticated opcode remain in the exact
signature.  By
\cref{def-budgeted-saturated-observable-closure,def-budgeted-derivability-closure},
the circuit occupies the budget slice only through the represented
derivation from which normalization started, with its complete cost.  This
proves the three structural alternatives.

The final display is the definition in
\cref{eq-affordable-coefficient-interaction-coordinate}.  The deterministic
no-creation identity of
\cref{lem-derived-constructor-zero-contextual-charge} attributes no
incremental conditional projection to a postprocessing once its complete
input algebra is exposed.  The exact contextual decomposition of
\cref{cor-five-family-abs-envelope} retains the earlier joint load, and the
source definition retains all applicable gates in
\(Q_I V_{\mathscr C,I}\).  Hence incremental charge and the complete
qualified source value obey the two distinct identities stated in the
theorem.  This proves the charge-conservation assertion.

\end{proof}

\begin{figure}[bp]
\centering
\begin{tikzpicture}[fw diagram,x=1cm,y=1cm]
  \node[fw box,minimum width=3.7cm] (record) at (0,3.25)
    {normalized saturated record};
  \node[fw source,minimum width=3.2cm] (source) at (0,2.15)
    {prospective source};
  \node[fw residual,minimum width=3.2cm] (res) at (3.7,2.15)
    {\(\Res\): no source cell};
  \node[fw provenance,minimum width=3.0cm] (state) at (-3.0,0.9)
    {state ancestry\\\(S_{\mathrm{st}}\)};
  \node[fw provenance,minimum width=3.0cm] (coeff) at (3.0,0.9)
    {coefficient ancestry\\\(S_{\mathrm{cf}}\)};
  \node[fw box,minimum width=4.2cm] (joint) at (0,-0.35)
    {exact five-family provenance pair};
  \node[fw geom,text width=9.3cm] (modes) at (0,-1.6)
    {\(\mathsf A\quad|\quad\mathsf X\quad|\quad
      \mathsf Q_{\mathrm{ex}}\quad|\quad
      \mathsf Q_{\mathrm{rev}}\quad|\quad
      \mathsf Q_{\mathrm{nr}}\)};
  \node[fw use,minimum width=4.2cm] (derived) at (0,-2.85)
    {derived uses: \(\Caus,\Lift,\Bdry\)};
  \draw[fw arrow] (record) -- (source);
  \draw[fw dashed arrow] (record) -- (res);
  \draw[fw arrow] (source) -- (state);
  \draw[fw arrow] (source) -- (coeff);
  \draw[fw arrow] (state) -- (joint);
  \draw[fw arrow] (coeff) -- (joint);
  \draw[fw arrow] (joint) -- (modes);
  \draw[fw arrow] (modes) -- (derived);
\end{tikzpicture}
\caption{Universal state--coefficient classification of affordable
prospective structural sources inside the saturated ledger.  The five master
modes are crossed with the two exact provenance signatures.  Causal
orientation, valid lifting, and boundary readout inherit their source cell;
the residual coordinate is the post-saturation non-source complement.}
\label{fig-universal-state-coefficient-source-decomposition}
\end{figure}

\subsection{Exhaustion and terminal-readout consequences}
\label{subsec:state-coefficient-exhaustion-consequences}

\begin{corollary}[Finite universal exhaustion criterion for affordable
structural sources]
\label{cor-universal-state-coefficient-source-exhaustion}

Fix a transfer class, a monotone sublinear positively homogeneous family
functional \(\mathfrak M_n\), and a budget sequence \(B_n\).  Suppose a bound
\(\varepsilon_n\) satisfies
\[
\operatorname{CoreAct}_{\mathfrak M,
K_{\mathrm{st}},K_{\mathrm{cf}},n}^{\mathrm{sat}}
\bigl(\Psi_{\mathrm{all}}(B_n,n)\bigr)
\le \varepsilon_n
\]
throughout the finite charge-bearing core index set, with empty cells assigned
value zero.
Then
\begin{align*}
\mathfrak U_{\mathfrak M,n}^{\mathrm{src,sat}}(B_n)
&\le\varepsilon_n,\\
I_{\mathfrak M,n}^{\mathrm{src,sat}}(B_n)
&\le2\varepsilon_n.
\end{align*}
In particular, a transfer-negligible common bound exhausts the affordable
target-aligned prospective source algebra.

\end{corollary}

\begin{proof}

Apply \cref{eq-global-core-reindexing} and then
\cref{eq-universal-master-structural-interaction-bound,%
eq-universal-combined-structural-interaction-bound}.  The finite maximum is
over charge-bearing core pairs; operational decorators remain in the complete
records defining those coordinates.

\end{proof}

\begin{corollary}[Separation of terminal verification from prospective
structural activity]
\label{cor-terminal-verifier-not-prospective-structural-activity}

Let \(v_I:X_I\to\{0,1\}\) be a represented terminal verifier, and suppose
that \(\sigma(v_I)\) has nonzero target projection.  The readout \(v_I\)
alone contributes zero to every structural-interaction coordinate in
\cref{eq-affordable-coefficient-interaction-coordinate}.  If a complete
temporally admissible source record \(\mathscr C\) uses \(v_I\), then its
contribution is
\[
Q_I V_{\mathscr C,I},
\]
with every applicable source gate and the complete ancestral derivation of
\(v_I\) retained in the unique state--coefficient cell of \(\mathscr C\).
In particular, target projection of a terminal verifier cannot replace
dynamic utility, compatibility, authorized orientation, target reach,
regularity, or temporal-source admissibility.

\end{corollary}

\begin{proof}

The defining supremum in
\cref{eq-affordable-coefficient-interaction-coordinate} ranges over complete
prospective source records and evaluates their qualified visibility, rather
than the projection mass of an isolated readout.  Hence an isolated verifier
is outside that class.  When it is an ancestor of a complete source record,
\cref{def-state-coefficient-source-provenance} retains every tagged vertex in
the normalized ancestral DAG of the qualified pair, including every
nonconstant qualification gate, and the same definition evaluates the full
product \(Q_I V_{\mathscr C,I}\).  These are exactly the two asserted cases.

\end{proof}

\begin{corollary}[An extremal readout requires represented intermediate
transport]
\label{cor-extremal-readout-requires-intermediate-transport}

Let \(D_I:X_I\to\mathbb R\) be a represented observable whose selected
minimum or maximum contains \(T_I\).  Extremum consistency supplies a
terminal-discrimination readout.  It supplies prospective ambient
\(\Geom\) activity precisely through a complete temporally admissible record
\(\mathscr C=(\mathfrak N,\Phi,D_I,J,\mathfrak K)\) that retains, before the
transition in which it is used, the represented transition rule, its
same-generator authority comparison, its authorized \(\Caus\) alignment,
its target reach, and its regularity comparison.  The contribution is the
complete qualified visibility \(Q_I V_{\mathscr C,I}\).

If no such record is present at the charged ceiling, \(D_I\) is a terminal
verifier and contributes zero to the prospective source profile.  If the
record is present, failure of any required gate sets its prospective value
to zero, while a quantitative regularity loss is retained in the regularity
factor of \(V_{\mathscr C,I}\).  Thus knowledge that the target is an
extremum cannot be substituted for represented control of the intermediate
transitions leading to that extremum.

\end{corollary}

\begin{proof}

Extremum consistency is one of the terminal-readout alternatives in
\cref{def-target-discrimination-certificate}.  By
\cref{cor-terminal-verifier-not-prospective-structural-activity}, an isolated
readout has zero prospective structural activity.  When \(D_I\) belongs to a
complete ambient record,
\cref{def-geom-extraction,def-authority-compatible-caus-alignment} retain the
transition, authority, alignment, reach, regularity, temporal, and cost gates
in the same qualified visibility.  The empty-record and failed-gate cases
follow from the defining supremum and product conventions.

\end{proof}

\section{Quantitative Normal Forms for Exact Geom Cells}
\label{subsec-quantitative-exact-geom-cells}
\label{sec:quantitative-exact-geom-cells}

The universal decomposition specializes to the following same-record
evaluations on the active-geometry branch.

\begin{corollary}[Cellwise quantitative form of the exact Geom decomposition]
\label{cor-exact-signature-geom-cellwise-quantitative-form}

At the normalized ceiling of
\cref{thm-canonical-exact-signature-geom-decomposition}, the complete
prospective Geom profile satisfies the exact finite identity
\begin{equation}
\label{eq-exact-signature-geom-cellwise-maximum}
G_n^{\mathrm{src,sat}}(\widehat B)
=
\max_{S\in\mathfrak P_{\Geom,5}}
\max\!\Bigl\{
G_{\mathsf X,S,n}^{\mathrm{src,sat},=}(\widehat B),
G_{Q,\mathrm{ex},S,n}^{\mathrm{src,sat},=}(\widehat B),
G_{Q,\mathrm{rev},S,n}^{\mathrm{src,sat},=}(\widehat B),
G_{Q,\mathrm{nr},S,n}^{\mathrm{src,sat},=}(\widehat B)
\Bigr\}.
\end{equation}
For every exact signature \(S\), the three quotient-supported coordinates
obey
\begin{align}
\label{eq-exact-signature-exact-quotient-cell-bound}
G_{Q,\mathrm{ex},S,n}^{\mathrm{src,sat},=}(\widehat B)
&\le m_n^{\mathrm{src,sat}}(\widehat B),\\
\label{eq-exact-signature-reversible-cell-bound}
G_{Q,\mathrm{rev},S,n}^{\mathrm{src,sat},=}(\widehat B)
&\le G_{Q,\mathrm{rev},n}^{\mathrm{src,sat}}(\widehat B),\\
\label{eq-exact-signature-resolvent-cell-bound}
G_{Q,\mathrm{nr},S,n}^{\mathrm{src,sat},=}(\widehat B)
&\le G_{Q,\mathrm{nr},n}^{\mathrm{src,sat}}(\widehat B).
\end{align}
The ambient cell has the exact same-record evaluation
\begin{equation}
\label{eq-exact-signature-ambient-cell-evaluation}
G_{\mathsf X,S,n}^{\mathrm{src,sat},=}(\widehat B)
=
\sup_{\operatorname{Prov}_{\Geom,5}(\mathfrak g)=S}
M_\Phi C_{\Phi,\mathfrak K}^{\mathrm{auth}}(D)
R_T(D)\frac1{1+\rho_\Phi(D)},
\end{equation}
where the supremum is over the temporally admissible normalized records at
that ceiling and is zero when the cell is empty.  Consequently it is bounded
by each of the three corresponding cell-restricted suprema of
\(C_{\Phi,\mathfrak K}^{\mathrm{auth}}(D)\), \(R_T(D)\), and
\((1+\rho_\Phi(D))^{-1}\).

Thus a family-specific Geom proof is complete precisely when it supplies a
bound for every coordinate appearing in
\cref{eq-exact-signature-geom-cellwise-maximum}.  A provenance label is not
an additional source coordinate: zero-charge or inherited constructors may
be contracted only through a theorem preserving the complete same-record
visibility in \cref{eq-exact-signature-ambient-cell-evaluation}.

\end{corollary}

\begin{proof}

The disjoint-union assertion of
\cref{thm-canonical-exact-signature-geom-decomposition} gives
\cref{eq-exact-signature-geom-cellwise-maximum}.  The exact-quotient
comparison is
\cref{thm-exact-quotient-geometry-visibility-dominated}; the other two
support bounds are inclusions of exact-signature subfamilies into their
master compensated coordinates.  Restricting
\cref{eq-prospective-ambient-geom-exact-profile} to one exact signature gives
\cref{eq-exact-signature-ambient-cell-evaluation}, and the factorwise bounds
follow because all displayed factors lie in \([0,1]\).

\end{proof}

\begin{theorem}[Reversible carriers and authenticated symmetry provenance]
\label{thm-reversible-carrier-authenticated-symmetry-separation}

Let \(\mathscr C\) be a dynamically qualified represented Geom source and
let \(\mathcal A_{\Geom,5}(\mathscr C)\) be its complete active ancestral set
from \cref{def-active-geom-source-frontier}.  The conductance identity
\[
c_\Phi(x,z)=c_\Phi(z,x)
\]
certifies the reversible Dirichlet carrier of \(\mathscr C\).  It contributes
no \(\mathsf{Sym}\) provenance tag by itself.

The source signature of \(\mathscr C\) contains \(\mathsf{Sym}\) precisely
when its complete active ancestral set contains a state-dependent symmetry constructor
with the authenticated action and equivariance record required by
\cref{def-five-family-abs-provenance}.  A data-dependent center, ordering,
relabeling, potential, conductance, or current lacking that authenticated
action retains the provenance of its represented construction and its
complete cost, and receives no \(\mathsf{Sym}\) tag.

Suppose every authenticated public action occurring in the complete active
ancestral set acts trivially on each state-dependent primitive interface used by the
record.  Then its orbit maps and averaging operations are
identity-equivalent on the active carrier, and every isotypic output is
either inherited from a differently tagged state-dependent argument or is
state-independent.  Hence the record contains no independent
\(\mathsf{Sym}\)-derived Geom source and produces no nonidentity
\(\mathsf{Sym}\)-derived orbit quotient.  Mixed records retain and are
charged through their remaining provenance tags.

\end{theorem}

\begin{proof}

Symmetry of \(c_\Phi\) is the defining condition of the static Dirichlet form
in \cref{def-static-dirichlet-geometry}; it is an edge-reversal identity and
contains no group-action datum.  The \(\mathsf{Sym}\) tag is assigned by
\cref{def-five-family-abs-provenance} only to a represented group action,
orbit map, equivariant operation, averaging operator, or isotypic projection
whose action on the active public object is authenticated.  The complete
active ancestral set in \cref{def-active-geom-source-frontier} exposes every
state-dependent ancestor of the active tuple.  These definitions prove the
first two assertions and retain the construction cost through
\cref{eq-geom-active-frontier-cost}.

When every authenticated action is the identity on the active primitive
interfaces, each orbit is a singleton and group averaging acts as the
identity on state-dependent inputs.  A nontrivial isotypic component then
has no state-dependent input from the authenticated action; any remaining
state dependence is inherited from another exposed argument and retains that
argument's provenance tag.  Pure interfaces preserve the union of those
tags by \cref{def-five-family-abs-provenance}.  Thus the symmetry
constructors add neither an independent active Geom datum nor a nonidentity
orbit partition.  The mixed-provenance conclusion follows from
\cref{thm-exhaustive-five-family-geom-source-classification}.

\end{proof}


\begin{figure}[tbp]
\centering
\begin{tikzpicture}[fw diagram,node distance=9mm]
  \node[fw geom,minimum width=3.1cm] (root)
    {valid \(\Geom\) source};
  \node[fw box,below left=13mm and 12mm of root,minimum width=4.8cm] (support)
    {support tag\\
     \(\mathsf X,\mathsf Q_{\mathrm{ex}},
       \mathsf Q_{\mathrm{rev}},\mathsf Q_{\mathrm{nr}}\)};
  \node[fw provenance,below right=13mm and 12mm of root,minimum width=5.5cm] (prov)
    {source signature\\
     \(\mathsf{Add},\mathsf{Mult},\mathsf{Ord},
       \mathsf{Sym},\mathsf{Filt}\)};
  \node[fw source,below=18mm of root,minimum width=4.0cm,yshift=-18mm] (cells)
    {twenty provenance restrictions\\
     \(\mathfrak S_{\Geom}^{4}\times\mathfrak J_{\Abs}^{5}\)};
  \draw[fw arrow] (root) -- (support);
  \draw[fw arrow] (root) -- (prov);
  \draw[fw arrow] (support) -- (cells);
  \draw[fw arrow] (prov) -- (cells);
\end{tikzpicture}
\caption{Every valid geometric source has one of four support configurations
and a nonempty five-family source signature.  The resulting twenty cells are
an exhaustive finite cover; mixed-provenance cells may overlap.  This is
\cref{thm-exhaustive-five-family-geom-source-classification}.}
\label{fig-exhaustive-five-family-geom-source-classification}
\end{figure}

\section{Operational Provenance and Charge-Bearing Structural Cores}
\label{subsec-operational-provenance-charge-cores}
\label{sec:operational-provenance-charge-cores}

The exact provenance signature records the authenticated opcodes in the
complete active ancestry.  Source attribution is a separate operation.
Constructors certified to create no new contextual charge retain their
operational provenance and complete cost, while their source credit remains
attached to the represented arguments from which they are derived.

\subsection{Inherited decorators and charge-bearing cores}
\label{subsec:inherited-decorators-charge-cores}

\begin{definition}[Certified inherited decorator]
\label{def-certified-inherited-decorator}

Fix a coefficient-complete normalized record \(\mathscr C\) to which the
ancestral construction of
\cref{def-state-coefficient-source-provenance} applies.  It may be a
prospective source record or a completed exact-quotient accounting record;
its temporal type is unchanged.  Write its complete ancestral set as
\(\mathcal A_{\mathrm{src},5}(\mathscr C)\), and use the fixed encoded
topological order of this finite DAG.

A tagged vertex
\(v\in\mathcal A_{\mathrm{src},5}(\mathscr C)\) is a
\emph{certified inherited decorator} when its retained record supplies one
of the following family-uniform certificates.

\begin{enumerate}[label=\textup{(\roman*)}]

\item A set \(P\) of earlier active arguments and a represented evaluator
      \(h_v\) satisfy
      \[
      f_v=h_v\circ J_P,
      \]
      and \(v\) has no primitive public read or independently authenticated
      structural source outside those arguments.  This includes the
      deterministic constructors covered by
      \cref{thm-affordable-composition-no-creation,%
lem-derived-constructor-zero-contextual-charge}.

\item The vertex is a filtration, transcript, or valid-Lift coordinate whose
      target contribution is assigned to an earlier represented base record
      by
      \cref{thm-filtration-provenance-charging,lem-exact-lift-identity}.
      Its charge-parent set is the complete retained set of those base-record
      arguments.

\item The vertex has authenticated \(\mathsf{Sym}\) provenance and its
      represented action is certified to act trivially on each active
      primitive interface used by the record, as in
      \cref{thm-reversible-carrier-authenticated-symmetry-separation}.  Its
      charge parents are the nontrivial active arguments from which its
      output is inherited.

\end{enumerate}

A certificate is either syntactic, as in the displayed represented
factorization, or a represented structural certificate whose description,
verifier, and verification cost belong to the complete record.  A semantic
assertion without a retained certificate does not make a vertex a decorator.
For each applicable clause, order its retained certified parent sets first by
inclusion and then lexicographically in the fixed encoded vertex order, and
select the lexicographically first inclusion-minimal set.  The collection is
finite because the ancestral DAG is finite.  Define
\(\operatorname{Par}_{\mathrm{src}}(v)\) as the union of the selected sets
over the applicable clauses.  This rule also applies when several
factorizations or structural certificates satisfy one clause.  Each selected
parent precedes \(v\), so the selection is canonical and charge inheritance
is acyclic.  The subscript \(\mathrm{src}\) is retained as notation and does
not retype a completed accounting record.  Write
\[
\operatorname{Dec}_{\mathrm{inh}}(\mathscr C)
\subseteq\mathcal A_{\mathrm{src},5}(\mathscr C)
\]
for the resulting canonical decorator set.

\end{definition}

\begin{definition}[Charge-bearing core signature]
\label{def-charge-bearing-core-signature}

For
\(v\in\mathcal A_{\mathrm{src},5}(\mathscr C)\), define its anchor set
recursively in the fixed topological order by
\[
\operatorname{Anch}_{\mathscr C}(v)
=
\begin{cases}
\{v\},
&v\notin\operatorname{Dec}_{\mathrm{inh}}(\mathscr C),
\\[1mm]
\displaystyle
\bigcup_{u\in\operatorname{Par}_{\mathrm{src}}(v)}
\operatorname{Anch}_{\mathscr C}(u),
&v\in\operatorname{Dec}_{\mathrm{inh}}(\mathscr C).
\end{cases}
\]
An empty charge-parent set gives an empty anchor set.  The
\emph{charge-bearing structural core} is
\[
\mathcal K_{\mathrm{src},5}(\mathscr C)
=
\bigcup_{v\in\mathcal A_{\mathrm{src},5}(\mathscr C)}
\operatorname{Anch}_{\mathscr C}(v).
\]
Split it according to the state--coefficient type of its retained vertices
and define
\begin{align}
\label{eq-core-state-signature}
K_{\mathrm{st}}(\mathscr C)
&=
\bigcup_{
 v\in\mathcal K_{\mathrm{src},5}(\mathscr C)
 \cap\mathcal A_{\mathrm{src},5}^{\mathrm{st}}(\mathscr C)}
\operatorname{Prov}_5(v),
\\
\label{eq-core-coefficient-signature}
K_{\mathrm{cf}}(\mathscr C)
&=
\bigcup_{
 v\in\mathcal K_{\mathrm{src},5}(\mathscr C)
 \cap\mathcal A_{\mathrm{src},5}^{\mathrm{cf}}(\mathscr C)}
\operatorname{Prov}_5(v).
\end{align}

The existing pair
\[
\operatorname{OpSig}(\mathscr C)
:=
\bigl(S_{\mathrm{st}}(\mathscr C),S_{\mathrm{cf}}(\mathscr C)\bigr)
\]
is the \emph{exact operational signature}; the pair
\[
\operatorname{CoreSig}(\mathscr C)
:=
\bigl(K_{\mathrm{st}}(\mathscr C),K_{\mathrm{cf}}(\mathscr C)\bigr)
\]
is the \emph{charge-bearing core signature}.

At a coefficient-complete normalized ceiling \(B\), define
\begin{align}
\label{eq-core-activity-refined-cell}
&\operatorname{CoreAct}_{
 \mathfrak M,K_{\mathrm{st}},K_{\mathrm{cf}}
 \mid S_{\mathrm{st}},S_{\mathrm{cf}},n}^{\mathrm{sat},=}(B)
\nonumber\\
&\quad:=
\sup
\left\{
\mathfrak M_n\!\left(I\longmapsto Q_I V_{\mathscr C,I}\right):
\begin{array}{l}
(\rho,\mathscr C,Q)\text{ is a complete normalized}\\
\qquad\text{prospective source record},\\
\operatorname{Cost}_I(\rho)\preceq B,\\
\operatorname{OpSig}(\mathscr C)
 =(S_{\mathrm{st}},S_{\mathrm{cf}}),\\
\operatorname{CoreSig}(\mathscr C)
 =(K_{\mathrm{st}},K_{\mathrm{cf}})
\end{array}
\right\},
\end{align}
with value zero for an empty class.  The record, qualification gate,
complete visibility, and complete cost in
\cref{eq-core-activity-refined-cell} are the original ones.  Core formation
is charge attribution in the structural ledger and does not contract the
represented derivation.

Finally, put
\begin{equation}
\label{eq-core-activity-aggregate-cell}
\operatorname{CoreAct}_{
 \mathfrak M,K_{\mathrm{st}},K_{\mathrm{cf}},n}^{\mathrm{sat}}(B)
=
\max_{\substack{
 S_{\mathrm{st}}\supseteq K_{\mathrm{st}}\\
 S_{\mathrm{cf}}\supseteq K_{\mathrm{cf}}}}
\operatorname{CoreAct}_{
 \mathfrak M,K_{\mathrm{st}},K_{\mathrm{cf}}
 \mid S_{\mathrm{st}},S_{\mathrm{cf}},n}^{\mathrm{sat},=}(B).
\end{equation}

\end{definition}

\begin{theorem}[Exact separation of operational provenance and
charge-bearing cores]
\label{thm-operational-core-provenance-separation}

Fix a coefficient-complete normalized record \(\mathscr C\) in the domain of
\cref{def-state-coefficient-source-provenance}: either a prospective source
or a completed terminal-compatible exact quotient retained for accounting.
Its operational and core signatures are canonical and
satisfy
\[
K_{\mathrm{st}}(\mathscr C)\subseteq S_{\mathrm{st}}(\mathscr C),
\qquad
K_{\mathrm{cf}}(\mathscr C)\subseteq S_{\mathrm{cf}}(\mathscr C).
\]
For the prospective-source specialization, each exact operational cell has
the finite refinement
\begin{equation}
\label{eq-operational-cell-core-refinement}
\operatorname{StructAct}_{
 \mathfrak M,S_{\mathrm{st}},S_{\mathrm{cf}},n}^{\mathrm{sat}}(B)
=
\max_{\substack{
 K_{\mathrm{st}}\subseteq S_{\mathrm{st}}\\
 K_{\mathrm{cf}}\subseteq S_{\mathrm{cf}}}}
\operatorname{CoreAct}_{
 \mathfrak M,K_{\mathrm{st}},K_{\mathrm{cf}}
 \mid S_{\mathrm{st}},S_{\mathrm{cf}},n}^{\mathrm{sat},=}(B).
\end{equation}
Consequently,
\begin{equation}
\label{eq-global-core-reindexing}
\max_{\substack{
 S_{\mathrm{st}}\ne\varnothing\\
 S_{\mathrm{cf}}\subseteq\mathfrak J_{\Abs}^{5}}}
\operatorname{StructAct}_{
 \mathfrak M,S_{\mathrm{st}},S_{\mathrm{cf}},n}^{\mathrm{sat}}(B)
=
\max_{K_{\mathrm{st}},K_{\mathrm{cf}}}
\operatorname{CoreAct}_{
 \mathfrak M,K_{\mathrm{st}},K_{\mathrm{cf}},n}^{\mathrm{sat}}(B).
\end{equation}

These are exact prospective-profile reindexing identities.  The same complete derivation cost
\(\operatorname{Cost}_I(\rho)\) and the same qualified visibility
\(Q_I V_{\mathscr C,I}\) occur on both sides.  A comparison, clock, stored
transcript, valid-Lift auxiliary, or certified trivial symmetry can enlarge
the operational signature and the complete cost, but it creates no
independent charge-bearing core index.

For a terminal-compatible exact quotient of either temporal type, view its complete normalized
record as \(\mathscr C_{\widehat\rho}\) and let
\((v_1,\ldots,v_k)\) be its canonical five-family contextual frontier.  The
exact structural-charge identity refines to
\begin{equation}
\label{eq-core-charge-identity}
V_I(q_{\widehat\rho})
=
\sum_{\substack{1\le\ell\le k\\
v_\ell\notin
\operatorname{Dec}_{\mathrm{inh}}(\mathscr C_{\widehat\rho})}}
\operatorname{ch}_I
 \bigl(E_{\widehat\rho,\Gamma_5,\ell}\bigr).
\end{equation}
Thus each nonzero contextual charge is carried by a core anchor; inherited
decorators retain their operational records and receive no independent
summand.

A reduced core signature is not a restricted operational grammar.  In
particular,
\[
\operatorname{CoreSig}(\mathscr C)
=(\{\mathsf{Add}\},\{\mathsf{Add}\})
\]
does not restrict the complete evaluator to its retained core anchors.  The
complete operational record and its cost remain present while qualified
visibility is charged to those anchors.  Gate expressivity supplies no
independent structural credit, and deterministic Boolean constructors cannot
make a failed qualification gate nonzero.

\end{theorem}

\begin{proof}

The complete normalized ancestral DAG is finite and has a fixed encoded
topological order.  Each charge parent in
\cref{def-certified-inherited-decorator} precedes its decorated vertex.
Therefore the anchor recursion terminates and is canonical.  Each anchor is
a vertex of the original complete ancestral set, which proves the two
signature inclusions.

Fix an exact operational signature pair in the prospective-source profile.
Core formation assigns each record in that exact cell one core pair.  The finitely many core-pair restrictions
are disjoint and cover the exact cell.  A supremum over this finite disjoint
union is the maximum of the restricted suprema, proving
\cref{eq-operational-cell-core-refinement}.  Taking the maximum over the
operational signatures and regrouping the same records by their core
signatures proves \cref{eq-global-core-reindexing}.

No represented vertex is removed, contracted, or re-evaluated.  The original
derivation, qualification gate, active source tuple, and operational
decorators remain in the record defining
\cref{eq-core-activity-refined-cell}.  Hence the complete cost and
\(Q_I V_{\mathscr C,I}\) are unchanged.

Consider the canonical contextual frontier of a terminal-compatible exact
quotient.  In case~\textup{(i)} of
\cref{def-certified-inherited-decorator}, the distinguished output is
measurable with respect to the joint algebra of its earlier charge parents.
Adjoining it leaves that algebra unchanged, so
\cref{thm-affordable-composition-no-creation,%
lem-derived-constructor-zero-contextual-charge} gives zero contextual
charge.  The filtration and valid-Lift cases have zero incremental charge,
with their target contribution assigned to the retained base record, by
\cref{thm-filtration-provenance-charging,lem-exact-lift-identity}.  For a
certified trivial symmetry, every orbit operation is identity-equivalent or
state-independent by
\cref{thm-reversible-carrier-authenticated-symmetry-separation}; adjoining
its output likewise leaves the earlier target-bearing algebra unchanged.

The contextual summands whose distinguished vertices are inherited
decorators are therefore zero.  Removing those zero summands from
\cref{eq-exact-qualified-structural-charge} proves
\cref{eq-core-charge-identity}.  Their costs remain in the complete carrier
ledger.

Finally, core formation records charge attribution rather than evaluator
contraction.  By
\cref{thm-boolean-complete-mixed-ancestral-cells}, deterministic
\(\mathsf{Mult}\)- and \(\mathsf{Ord}\)-tagged postprocessings in an already
represented finite Boolean evaluator carry zero incremental contextual
charge after their joint inputs are exposed.  Their complete implementation
cost remains in the record, and their complete qualified source value remains
\(Q_I V_{\mathscr C,I}\), proving the final assertion.

\end{proof}

\subsection{Budget-faithful Boolean normalization}
\label{subsec:boolean-normalization-budget-faithfulness}

\begin{theorem}[Budget-faithful Boolean normalization]
\label{thm-budget-faithful-boolean-implementation-separation}

Let \((\rho,\mathscr C,Q)\) be a complete family-uniform represented
prospective source record for a finite binary presentation, after expansion
of each presentation-admissible oracle interface, and suppose
\[
\operatorname{Cost}_I(\rho)\preceq B
\qquad(I\in\mathfrak I_n).
\]
Apply \(\operatorname{Exp}_5\circ\operatorname{Exp}\) to the complete
derivation of the qualified source pair, rather than only to its terminal
readout.  Denote the resulting record by
\((\widehat\rho,\widehat{\mathscr C},\widehat Q)\).  Then
\begin{equation}
\label{eq-budget-faithful-normalization-record}
\sem{\widehat\rho}_I=\sem{\rho}_I,
\qquad
\widehat Q_I=Q_I,
\qquad
V_{\widehat{\mathscr C},I}=V_{\mathscr C,I},
\qquad
\operatorname{Cost}_I(\widehat\rho)
\preceq\Psi_5(B).
\end{equation}
The cost on the right charges the complete normalized pair: its gates,
instance-dependent coefficients, tables, work state, storage, qualification
evaluator, and retained authentication records.

At the level of the single saturated budget slice, normalization has the
exact denotational image
\begin{equation}
\label{eq-budget-faithful-normalized-image}
\left\{
\bigl(\sem{\operatorname{Exp}_5(\operatorname{Exp}(\rho))}_I
\bigr)_{I\in\mathfrak I_n}:
\begin{array}{l}
\rho\in\mathsf{Der}_{\mathfrak I},\\[-1mm]
\operatorname{Cost}_I(\rho)\preceq B
\quad(I\in\mathfrak I_n)
\end{array}
\right\}
=
\mathcal R^{\mathrm{sat}}_{n,\le B}.
\end{equation}
The normalized representatives on the left have cost at most the
class-preserving ceiling \(\Psi_5(B)\).  This set is precisely the normalized
denotational image of the original budget-\(B\) slice and may be a proper
subset of \(\mathcal R^{\mathrm{sat}}_{n,\le\Psi_5(B)}\).

The implementation--source and exact-charge consequences of this same
normalization are stated separately in
\cref{lem:boolean-implementation-source-separation,%
cor:boolean-normalization-core-charge}.  Together the three statements
constitute the budget-faithful normalization result carried by the public
label above.

\end{theorem}

\begin{lemma}[Boolean implementation--source separation]
\label{lem:boolean-implementation-source-separation}

Each active affine or product gate that is an internal deterministic vertex
of the normalized implementation of the already represented Boolean map is a
certified inherited decorator.
Its retained direct-input set supplies an earlier certified parent set and
its gate evaluator supplies the factorization in
\cref{def-certified-inherited-decorator}; the canonical source-parent set is
then selected by the rule in that definition.  Hence
\begin{equation}
\label{eq-boolean-gate-anchor-inheritance}
\operatorname{Anch}_{\widehat{\mathscr C}}(v)
=
\bigcup_{u\in\operatorname{Par}_{\mathrm{src}}(v)}
\operatorname{Anch}_{\widehat{\mathscr C}}(u)
\end{equation}
for each such gate \(v\).  The gate remains in
\(\operatorname{OpSig}(\widehat{\mathscr C})\), with its complete cost, but
it supplies no independent core anchor.  An independently authenticated
primitive source retains its own core tag; the assertion concerns only the
deterministic implementation gates.

A product-free normalized Boolean evaluator is affine over \(\mathbb F_2\).
If an active product gate occurs, then
\[
\mathsf{Mult}\in
S_{\mathrm{st}}(\widehat{\mathscr C})
\cup S_{\mathrm{cf}}(\widehat{\mathscr C}),
\]
and the \(\mathsf{Mult}\) tag records implementation provenance.  Source
contribution remains the complete qualified value
\(Q_I V_{\mathscr C,I}\); \(\operatorname{CoreSig}\) then attributes that
value to its charge-bearing anchors.

Let
\[
(S_{\mathrm{st}},S_{\mathrm{cf}})
=\operatorname{OpSig}(\widehat{\mathscr C}),
\qquad
(K_{\mathrm{st}},K_{\mathrm{cf}})
=\operatorname{CoreSig}(\widehat{\mathscr C}).
\]
For a declared monotone sublinear positively homogeneous family functional
\(\mathfrak M_n\), the normalized record satisfies the same-record bound
\begin{align}
\label{eq-budget-faithful-normalized-core-membership}
\mathfrak M_n\!\left(I\longmapsto Q_I V_{\mathscr C,I}\right)
&=
\mathfrak M_n\!\left(
I\longmapsto\widehat Q_I V_{\widehat{\mathscr C},I}\right)
\notag\\
&\le
\operatorname{CoreAct}_{\mathfrak M,
K_{\mathrm{st}},K_{\mathrm{cf}}
\mid S_{\mathrm{st}},S_{\mathrm{cf}},n}^{\mathrm{sat},=}
\bigl(\Psi_5(B)\bigr).
\end{align}
Profile membership requires one fully charged representation.  A table
realization retains its complete encoded-table cost; a compact circuit
retains the complete description, access, evaluation, storage, retained-data,
and output cost of its represented DAG.  A charged conversion may supply a
second representation at the conversion's complete cost.  Normalization
preserves denotation and complete ancestry accounting; it does not infer an
expanded formula, truth table, or compact DAG from semantic existence.  Thus
a compact DAG is charged at its actual complete shared-circuit cost even when
a fully expanded formula is large, while semantic existence alone supplies
no record in the budget slice.

\end{lemma}

\begin{corollary}[Exact core charge after Boolean normalization]
\label{cor:boolean-normalization-core-charge}

If the terminal record is a terminal-compatible exact quotient, let
\((v_1,\ldots,v_k)\) be its canonical five-family contextual frontier.  The
internal implementation gates contribute no contextual-charge summand, and
\begin{equation}
\label{eq-budget-faithful-normalized-core-charge}
V_I(q_{\widehat\rho})
=
\sum_{\substack{1\le\ell\le k\\
v_\ell\notin
\operatorname{Dec}_{\mathrm{inh}}(\widehat{\mathscr C})}}
\operatorname{ch}_I
 \bigl(E_{\widehat\rho,\Gamma_5,\ell}\bigr).
\end{equation}
The total charge identity is invariant under affine/product evaluator
normalization.

\end{corollary}

\begin{proof}[Proof of
\cref{thm-budget-faithful-boolean-implementation-separation,%
lem:boolean-implementation-source-separation,%
cor:boolean-normalization-core-charge}]

Apply
\cref{thm-expanded-derivation-equivalence,thm-five-family-abs-provenance}
to the represented tuple containing the complete source record and its
qualification indicator.  Componentwise denotation preservation fixes the
source data, every qualification comparison, and every factor in its
visibility.  The monotone transfer map in
\cref{def-budget-structure} gives the cost comparison in
\cref{eq-budget-faithful-normalization-record}.

For the inclusion from left to right in
\cref{eq-budget-faithful-normalized-image}, denotation preservation replaces
the normalized record by its originating cost-\(B\) derivation.  Conversely,
the definitions in
\cref{def-budgeted-saturated-observable-closure,def-budgeted-derivability-closure}
give each member of \(\mathcal R^{\mathrm{sat}}_{n,\le B}\) a complete
uniform derivation \(\rho\) of cost at most \(B\); normalizing that witness
places the same denotation on the left.  The cost comparison from \(B\) to
\(\Psi_5(B)\) is the one-way enlargement in
\cref{thm-five-family-abs-provenance}.

An internal affine or product gate is a represented deterministic function
of its retained direct inputs and has no primitive read outside those
inputs.  It therefore satisfies clause~\textup{(i)} of
\cref{def-certified-inherited-decorator}, and the recursive definition in
\cref{def-charge-bearing-core-signature} gives
\cref{eq-boolean-gate-anchor-inheritance}.  The Boolean gate identities
in \cref{thm-boolean-complete-mixed-ancestral-cells} show that absence of a
product gate leaves only affine operations.  The complete ancestral-set
definition retains every active product gate in the operational signature.
Core formation removes no vertex or cost and assigns the same qualified
visibility to the unique exact operational/core cell.  Membership in that
cell proves
\cref{eq-budget-faithful-normalized-core-membership}.

Finally, internal deterministic gates have zero incremental contextual
charge by
\cref{thm-affordable-composition-no-creation,lem-derived-constructor-zero-contextual-charge}.
Substituting the certified decorator set in the exact core-charge identity
\cref{eq-core-charge-identity} proves
\cref{eq-budget-faithful-normalized-core-charge}.

\end{proof}

\Cref{fig-operational-core-provenance} separates implementation provenance
from the anchors that receive contextual source charge.

\begin{figure}[tbp]
\centering
\resizebox{.96\linewidth}{!}{%
\begin{tikzpicture}[fw diagram,node distance=8mm and 11mm]
  \node[fw source,text width=2.8cm] (repr)
    {charged table or\\compact shared DAG};
  \node[fw provenance,text width=2.9cm,right=of repr] (expand)
    {denotation-preserving\\Boolean normalization};
  \node[fw box,text width=2.8cm,right=of expand] (opsig)
    {operational signature\\all active gate tags};
  \node[fw use,text width=2.9cm,right=of opsig] (decor)
    {inherited decorators\\measurable from parents};
  \node[fw terminal,text width=2.8cm,right=of decor] (core)
    {core anchors\\contextual charge};
  \draw[fw arrow] (repr) -- (expand);
  \draw[fw arrow] (expand) -- (opsig);
  \draw[fw arrow] (opsig) -- (decor);
  \draw[fw arrow] (decor) -- node[above,font=\scriptsize]
    {zero increment} (core);
  \node[fw group,fit=(repr)(expand)(opsig)(decor)(core),
    label={[font=\scriptsize,fwInk]above:complete implementation cost retained}] {};
\end{tikzpicture}%
}
\caption{Boolean normalization records every implementation gate in the
operational signature.  A deterministic gate already measurable from its
charged parents is an inherited decorator: its representation cost remains,
but it creates no independent contextual source charge.}
\label{fig-operational-core-provenance}
\end{figure}

\subsection{Derivability filtration and cost frontiers}
\label{subsec:derivability-filtration-cost-frontiers}

\begin{lemma}[Well-posedness of the derivability filtration]\label{lem-derivability-well-posed}

The classes \(\mathcal D_{n,\le B}\) are nested in \(B\), and consequently
\(m_n^{\mathrm{sat}}(B)\), \(G_n^{\mathrm{sat}}(B)\), and
\(I_n^{\mathrm{sat}}(B)\) are nondecreasing in \(B\).  Below any chosen
finite-cost witness, every represented object has a Pareto-minimal
description cost. If the resource order is total and well ordered (in
particular, for a scalar discrete cost), this is a least cost.

\end{lemma}

\begin{proof}\phantomsection\label{proof-derivability-well-posed}

If \(B_1\preceq B_2\), every derivation of cost at most \(B_1\) is also a
derivation of cost at most \(B_2\), so
\(\mathcal D_{n,\le B_1}\subseteq\mathcal D_{n,\le B_2}\). At fixed
\(n\), each extraction coordinate is a supremum of nonnegative
visibility terms over its corresponding budget slice. The inclusion of
slices therefore proves monotonicity without any finiteness or
attainment argument. By the local-finiteness coding convention, the nonempty set of
descriptions of a represented object having cost at most any one of its
finite-cost descriptions is finite. Its finite set of costs has a
minimal element in the resource partial order, which proves the
Pareto-minimal assertion. Under a total well order that element can be
chosen least. All budget-slice arguments use an actual witness of cost
\(\preceq B\) and therefore require neither scalarization nor a global
least element.

\end{proof}

\begin{remark}[Cost-frontier convention]
\label{rem-cost-frontier-convention}

When the resource monoid is only partially ordered, subsequent phrases
such as ``least saturated cost'' refer to the upward-closed admissible
budget set, represented by its Pareto-minimal frontier. They do not
assert a least vector among incomparable resources. Statements using a
single number, logarithmic degree, infimum, or equality of threshold
costs are made after the declared scalar cost or monotone scalarization
has been fixed. The constructor-charge theorems themselves use only an
actual witness with cost \(\preceq B\).

\end{remark}

\begin{remark}[Derivability of committors and finite-horizon approximations]\label{rem-committor-derivability}

For a killed generator \(L_N\), the exact discounted committor
\(u_{\lambda,L}=(\lambda I-L_N)^{-1}k_T\) is legal under \(\mathcal B\)
exactly when it has a derivation in \(\mathcal D_{n,\le b(n)}\). Thus
the resolvent must be represented by an admissible derivation. A
finite-horizon truncation such as
\(\sum_{j\le M}a_jP_N^jk_T\) is admissible when the one-step rule,
coefficients, horizon, and terminal edge gate are derived data and the
finite computation lies within the resource envelope. Its saturated
computation cost is then charged explicitly. In the polynomial instance
this condition is \(b(n)=n^C\) for some fixed \(C\).

\end{remark}

\begin{convention}[Budgeted computation is presentation-relative]\label{conv-presentation-relative-budgeted-computation}

Throughout this part, a budget-admissible oracle-free computation is
relative to \(\mathcal R^{\mathcal B}_n\), the budgeted saturated
represented closure of the presented task. The algorithm may have any
uniform oracle-free code allowed by the resource model, but every
instance-dependent datum read by that code must be declared public or
generated by a prior vertex of the same represented derivation and its
legal lifted computation. This includes initialization data,
finite outcome-generation interfaces, lifted coordinates, auxiliary tables,
predicates, basis changes, quotient/fiber operations, and update maps.
Any generated datum becomes part of the full lifted observable
configuration and retains its derivation cost when used as structural
data. Measurability in \(\mathcal F^{\mathrm{env}}_{I,\le b(n)}\) alone
does not authorize a read or assign a budget.

If a transition rule uses a datum outside
\(\mathcal R^{\mathcal B}_n\), such as a withheld modulus, trapdoor,
target table, oracle answer, nonuniform advice string, or relation
recoverable only above the declared budget, then the resulting algorithm
is not budget-admissible and oracle-free for the presented task. Such a
rule uses oracle or advice data, performs over-budget recovery, or
enlarges the primitive signature. This convention is the computational
counterpart of \PartI's distinction between exposed and
presentation-inaccessible structure after imposing the saturated cost
filtration.

The complexity theorems below quantify over the presentation-relative
generators embedded by \PartII at the declared budget. Transfer-class
target arrival by such a generator forces a saturated source in the
observable accounting at the same budget, up to the transfer-class
distortion. Structurality follows from the \PartII channel decomposition
and the \PartIII cost filtration. For \(\mathcal B_{\mathrm{poly}}\), this convention
recovers ordinary presentation-relative polynomial computation and
inverse-polynomial target-arrival thresholds inside
\(\mathcal R_n^{\mathcal B_{\mathrm{poly}}}\).

\end{convention}

\section{Encoding Adequacy for Budgeted Presentations and the Polynomial
Instance}\label{encoding-adequacy-budgeted-presentations-and-the-polynomial-instance}
\label{sec:encoding-adequacy-budgeted-presentations}

The encoded presentation determines the input model and supplies all
public instance-dependent data. The budgeted saturated observable
closure is the closure of these data under the admissible budgeted
oracle-free constructions.

Budget completeness connects the abstract presentation to a uniform finite input encoding with comparable costs.
\begin{definition}[Budget-complete encoding]\label{def-presentation-complete-encoding}

An encoding \(\operatorname{enc}(I)\) of a task-family presentation is
complete for the declared budget structure \(\mathcal B\) when the
following simulations are uniformly charged by the resource monoid and
stay below \(b(n)\) whenever the represented derivation lies within the
declared budget. For \(\mathcal B_{\mathrm{poly}}\), this is the usual
polynomial-size encoding with polynomial-time simulations.

\textbf{Decoding.} The primitive public task data are decoded from
\(\operatorname{enc}(I)\) within the declared decoding budget:
state-size data, primitive observables, terminal predicates, declared
public parameters, static comparability data, and the declared
dynamic/lift constructor signatures.

\textbf{Constructor simulation.} Every admissible constructor used to
generate \(\mathcal R^{\mathcal B}_n\) has a uniform implementation
whose overhead is charged in \(\mathcal B\) from
\(\operatorname{enc}(I)\) and the current lifted configuration.

\textbf{Evaluator normalization.} Every state-dependent primitive,
lookup rule, and finite evaluator supplied by the encoding has either an
explicit charged table or a uniform charged implementation over a fixed
finite operational basis. Finite-precision arithmetic and Boolean
evaluators are compiled to this basis with class-preserving overhead.
An explicit table retains its encoded length and lookup cost. A
state-dependent access rule supplied without either representation is
an oracle interface in the
\hyperref[def-oracle-extended-closure]{oracle-extended presentation}.

\textbf{Transition representation.} Every uniform oracle-free
operational transition rule whose transcript charge is
\(\preceq b(n)\) is represented by its code, initializer, successor,
clock or budget counter, peak native-state capacity, finite
outcome/terminal transcript, and declared native-state effects. The
one-step kernel on the clocked transcript carrier produced by that rule
is therefore a legal forward observable transport datum. In the polynomial instance, this
includes exactly the uniform operational transition rules whose native
resources are polynomially bounded on \(\operatorname{enc}(I)\).

\textbf{Input closure.} A datum not decoded from
\(\operatorname{enc}(I)\), not generated by the admissible constructors,
and not present in the legal lifted configuration is outside the
budgeted observable algebra of this presentation. Supplying it changes the
primitive signature, gives advice, or gives oracle access.

\end{definition}

Encoding adequacy connects the operational representation to the declared computation model.
\begin{theorem}[Encoding adequacy]\label{thm-encoding-adequacy}

Fix a budget-complete encoding of the task family. A represented
operational transition rule is presentation-relative and lies within
\(\mathcal B\) exactly when its instance-dependent data have records in
\(\mathcal R^{\mathcal B}_n\), its lifted transcript is legally
represented, and its accumulated charge is \(\preceq b(n)\).
Consequently the budget-admissible algorithms for the encoded
presentation are exactly the budget-admissible generators embedded by
\PartII. Setting \(\mathcal B=\mathcal B_{\mathrm{poly}}\) recovers the
uniform polynomial-resource equivalence inside
\(\mathcal R_n^{\mathcal B_{\mathrm{poly}}}\).

\end{theorem}

\begin{proof}\phantomsection\label{proof-encoding-adequacy}

Let \(M\) be a uniform operational transition rule whose transcript
charge lies within \(\mathcal B\). At any represented step, its single
live native state contains the input-dependent working data, finite
control, work state, budget counter, and memory records. The finite
clocked carrier contains the outcome/terminal transcript; declared
effects expose only the required native-state readouts. The current
native state determines the represented conditional outcome law and successor. By transition
representation in the definition of budget-complete encoding, this
one-step rule is a legal forward
observable transport datum generated from
\(\mathcal R^{\mathcal B}_n\) and the lifted configuration. Hence
\(M\) is a presentation-relative budget-admissible transition rule, and
\PartII embeds each bounded slice as
\(L_M^{[H]}=\rho(P_M^{[H]}-I)\). For
\(\mathcal B_{\mathrm{poly}}\), this is the native polynomial-resource
simulation declared by the budget structure.

Conversely, let \(A\) be a presentation-relative budget-admissible
transition rule. Its transition is built from declared public data,
legal lifted coordinates, and admissible constructors in
\(\mathcal R^{\mathcal B}_n\), all within the declared resource
envelope. By decoding and constructor simulation, the encoded model can
decode the presentation, execute each constructor and lifted-coordinate
update, and realize the same one-step operational transition law with
charged overhead still below \(b(n)\). Thus \(A\) is a within-budget
uniform transition rule for the encoded presentation.

The equivalence concerns transition rules and their admissible
instance-dependent input data. A semantic function of the induced
future path has zero structural cost only when it has an explicit
representative. In particular, a
committor, first-hitting partition, occupation law, or resolvent value
of the same induced generator contributes to \(m^{\mathrm{sat}}\) or
\(G^{\mathrm{sat}}\) when the same denotation has a budget-admissible
saturated representative as a quotient, potential, geometry, lifted
interface, boundary readout, or admissible within-budget finite-horizon
computation. Derivability accounting requires a construction of the
semantic object; an admissible computation from
\(\mathcal R^{\mathcal B}_n\) is counted at its saturated cost.

\end{proof}

\begin{lemma}[Budgeted generation gives budgeted representation]\label{lem-budgeted-generation-gives-representation}

Fix a presentation-complete encoding and a budget-compatible resource
envelope. If a datum, relation, finite-horizon trajectory statistic,
first-hit quotient, target-contact statistic, boundary value, or
postprocessed observable is produced by a uniform oracle-free
computation from \(\mathcal R^{\mathcal B}_n\) and the legal clocked
transcript representation, with peak native-state and transcript charge
below \(b(n)\), then its
denotation has a budget-admissible saturated representative whose cost
lies within budget: the public code for the computation, the
horizon, the lifted coordinates it reads, and the finite postprocessing
rule. Conversely, if every budget-admissible representative of such a
denotation has cost above the declared ceiling, then the denotation has
not been produced as legal budgeted data for the original presentation.
Using it in a purported budget-admissible transition rule therefore
requires oracle or advice data, an enlarged primitive signature, or
over-budget recovery. Consequently every uniform public target alignment
admissible for the presented task belongs to the budgeted saturated
observable closure.

\end{lemma}

\begin{proof}\phantomsection\label{proof-budgeted-generation-gives-representation}

The forward implication is the definition of
\(\mathcal R^{\mathcal B}_n\) together with the
constructor-simulation clause of presentation completeness. An
admissible within-budget finite-horizon computation has a finite lifted
transcript whose charge is below \(b(n)\), reads only legal lifted
coordinates and declared public data, and applies only budget-admissible
oracle-free postprocessings. This transcript description is a
budget-admissible representative, with saturated cost bounded by the
declared ceiling. For the converse, suppose such a denotation were
produced by a presentation-relative within-budget computation from
the legal closure. The first part would give a within-budget
budget-admissible representative, contradicting the assumed lower bound
above \(b(n)\) on all budget-admissible representatives. Hence any use
of the denotation requires data outside the budgeted observable algebra,
an enlarged primitive signature, oracle or advice access, or computation
exceeding the declared budget.

\end{proof}

All \(L^2\) projections, normalized target correlations, reach ratios,
regularity ratios, and extraction quantities in this part use the
presentation reference measure on the base space and a Part
II-legitimate reference measure on any lifted configuration space. A
lifted reference law may be the canonical public full-support product
gauge on a padded finite configuration encoding, the public
initialization lift when it satisfies the support condition, or a
publicly represented transfer-class-bounded density change of one of
these gauges. The product gauge is sufficient for exact configuration
and first-hit accounting but supplies no structural Valid-Lift
projection or target-fraction identity.
Occupation measures, hitting-conditioned laws, invariant laws, Green
measures, and resolvent measures obtained from the same measured
operator qualify as reference measures for these quantities only when
they also have an independent admissible representation with controlled
density. Accordingly, a clock or program-counter coordinate aligned
under a valid public lifted law is exposed structure, whereas alignment
arising only after replacing that law by the algorithm's trajectory
distribution is not exposed by the presentation.

The function \(y_n\) is the centered target indicator. Low saturated
cost means that this indicator admits a representation in a
low-complexity admissible observable class. A large saturated hardness
cost means that every observable class within the declared budget has
negligible target contrast. Raw Walsh degree is one coordinate in this
saturated cost profile.

This part develops the discounted hitting functional, the
target-projected killed resolvent, the structural operator ledger with
full-process first-hit accounting, the Walsh/spectral bound, the combined extraction functional,
the saturated cost invariant \(B^*\) and its log-cost coordinate
\(d^*\), and the admissibility dichotomy. It proves hardness for
presentation-relative budget-admissible oracle-free generators in
regimes where the exposed channel representatives and comparison data
have been bounded: random targets, removed structure, obstructed
placement against the declared geometry, and proven saturated-cost lower-bound
constructions. These comparison data enter only the analysis of a
budget-admissible computation. A uniform
oracle-free procedure supplies only its transition rule; every
instance-dependent datum used by that rule must have a derivation in
\(\mathcal R^{\mathcal B}_n\), be a declared public parameter, or lie in
the legal lifted closure. \PartII embeds that rule as a generator, and
the present part estimates the exposed projections of that generator. The
problem-specific input supplied by an application is an explicitly proved
saturated structural negligibility bound for its complete
derivation-indexed quotient and geometric families. For a fixed
cryptographic instance whose gapping structure is admissible, \PartIII
reduces hardness to the statement that \(B^*\) lies above the declared
budget; in the polynomial instance this is the usual superpolynomial,
equivalently superlogarithmic, lower bound in log-cost notation. \PartIV
supplies the complementary quantitative learning analysis, and \PartVIII
records a separate application of the framework.

Admissibility is a computational provenance condition relative to the
budgeted saturated observable closure of the instance presentation. Part
II proves that every presentation-relative budget-admissible
oracle-free computation, represented on its full observable
configuration, induces a generator whose target-relevant transport is
decomposed into exposed structural components and a residual term. The
exposed structural components are \(\Geom\), \(\Caus\),
\(\Abs\), \(\Lift\), and \(\Bdry\). Among them,
the geometric/reversible component \(\Geom\) and the
quotient/lumping component \(\Abs\) provide the intrinsic
structure. The directed/drift component \(\Caus\), the
lifted-state component \(\Lift\), and the boundary/killing
component \(\Bdry\) respectively orient exposed edge
differences, reorganize valid lifted coordinates, and apply terminal or
support values to an already represented structure. The scalar
\(\Bdry\) data mark, kill, weight, or evaluate states and affect
transport only through a represented \(\Geom\), \(\Abs\),
or \(\Lift\) edge relation whose differences are oriented by
\(\Caus\). The residual term
\(\Res\) is residual accounting: it is the part of the
movement with no target-relevant representative in
\(\mathcal R^{\mathcal B}_n\) after the five exposed structural
projections have been applied. A proposed budget-admissible method using
an instance-dependent relation outside \(\mathcal R^{\mathcal B}_n\)
requires presentation-inaccessible information: oracle or advice
access, an over-budget recovery procedure, or an enlarged primitive
signature. If the relation is generated inside
\(\mathcal R^{\mathcal B}_n\), it appears as intrinsic
\(\Geom/\Abs\) structure or as a valid use of it after
lifting. The full target arrival of this decomposition is bounded by the
uniform saturated profile after applying the exact first-hit
construction; no separate residual-coupling bound is inferred here.

\begin{example}[SAT after the structural decomposition]\label{ex-paper3-sat-stopwatch}

For SAT, \PartI gave the assignment space, clause observables, defect
score, success sector, atoms, level sets, and possible exposed
quotients. \PartII classified every budget-admissible transport rule on
that space. The present analysis concerns the hitting time of the
satisfying sector. Its Laplace transform is expressed by the
target-projected killed resolvent of a budget-admissible operator and bounded by
the saturated public cost at which the observable algebra exposes either
a quotient or a regular aimed geometry.

\end{example}

\section{Arrival Time as the Complexity
Measure}\label{arrival-time-as-the-complexity-measure}
\label{sec:arrival-time-complexity-measure}

Fix the positive terminal sector defined in \PartI. Throughout this
section, write

\[
T=T_n=E_{I_n}^+,
\qquad
N=N_n=X_n\setminus T.
\]

The non-target sector \(N\) contains both nonterminal states and terminal
failures. Thus a terminal failure remains terminal even though it belongs
to \(N\).

Killed dynamics isolate the non-target evolution and the flux entering the target sector.
\begin{definition}[Killed dynamics at the target]\label{def-killed-dynamics}

Let \(X_n\) be finite. An admissible continuous-time operator \(L\) is
written in row-frontier convention, so a row probability vector evolves
by

\[
p_t=p_0e^{tL}.
\]

Assume \(L_{xy}\ge0\) for \(x\ne y\) and \(\sum_yL_{xy}=0\) after
adjoining an absorbing cemetery state if needed. The killed non-target
operator is the restriction

\[
L_N=(L_{xy})_{x,y\in N},
\qquad
K_N=-L_N.
\]

The target killing rate is

\[
k_T(x)=\sum_{z\in T}L_{xz},
\qquad x\in N.
\]

When the killed chain is stable, every eigenvalue of \(K_N\) has
positive real part. In the reversible case \(K_N\) is self-adjoint in
\(L^2(\mu_N)\), and the killed-target gap is

\[
\gamma_T(L)=\lambda_{\min}(K_N).
\]

\end{definition}

Killing removes from the non-target equation all mass that reaches
\(T\). The vector \(k_T\) gives the rate of transition from each
non-target state into the target sector. These are
the standard killed-chain and hitting-time conventions
\citep{Norris1997}.

The arrival functional measures the target flux generated by the killed process.
\begin{definition}[Arrival functional]\label{def-arrival-functional}

Let \(\mu_0\) be any initial probability measure, put
\(a_0=\mu_0(T)\), and let \(\mu_{0,N}\) be its restriction to \(N\).
Let \(\tau_T\) be the first hitting time of \(T\), with
\(\tau_T=0\) on an initially successful state, for the absorbed process
generated by \(L\). The arrival functional is

\[
\Phi_L(t)=\mathbb P_{\mu_0}(\tau_T\le t).
\]

The killed frontier on \(N\) is
\(q_t=\mu_{0,N}e^{tL_N}\). Its survival mass
and hit flux are

\[
\mathbb P_{\mu_0}(\tau_T>t)=q_t\mathbf 1_N,
\qquad
\Phi_L'(t)=q_tk_T=\mu_{0,N}e^{tL_N}k_T
\]

whenever \(\Phi_L\) is differentiable.

\end{definition}

The resulting complexity parameter is
\(\Phi_L(\mathsf t(b(n)))\), evaluated at the numerical time capacity
of the declared resource ceiling.

\begin{definition}[Discounted hitting functional]\label{def-stopwatch-score}

For discount rate \(\lambda>0\), define

\[
S_T(\lambda;L,\mu_0)
=
\int_0^\infty e^{-\lambda t}\,d\Phi_L(t)
=
\mathbb E_{\mu_0}\!\left[e^{-\lambda\tau_T}\mathbf 1_{\{\tau_T<\infty\}}\right].
\]

A budget ceiling \(b(n)\) of finite time capacity is read at discount
\(\lambda=\mathsf t(b(n))^{-1}>0\). The functional is
within budget when \(S_T(\mathsf t(b(n))^{-1};L,\mu_0)\) has the required
transfer-class margin. For \(\mathcal B_{\mathrm{poly}}\), this is the
usual inverse-polynomial reading at \(b(n)=n^{O(1)}\).
For an infinite time capacity, use the finite-slice readings or a
separately justified \(\lambda\downarrow0\) limit; the resolvent
definition itself always assumes \(\lambda>0\).

\end{definition}

The exponential discount weights target hits occurring before the
budget more heavily than those occurring substantially later. It thereby
reduces the hitting-time distribution to a scalar functional that is
comparable across operators.

\begin{proposition}[Within-budget arrival bounds the discounted hitting functional]\label{prop-budget-hit-stopwatch}

For every budget \(B>0\),

\[
S_T(1/B;L,\mu_0)
\ge
 e^{-1}\,\mathbb P_{\mu_0}(\tau_T\le B).
\]

\end{proposition}

\begin{proof}\phantomsection\label{proof-budget-hit-stopwatch}

On the event \(\{\tau_T\le B\}\), the factor \(e^{-\tau_T/B}\) is at
least \(e^{-1}\). Therefore

\[
\mathbb E[e^{-\tau_T/B}\mathbf 1_{\{\tau_T<\infty\}}]
\ge
\mathbb E[e^{-\tau_T/B}\mathbf 1_{\{\tau_T\le B\}}]
\ge e^{-1}\mathbb P(\tau_T\le B).
\]

\end{proof}

\begin{example}[Discounted hitting for SAT]\label{ex-arrival-sat}

For a SAT instance, \(T\) is the set of satisfying assignments and \(N\)
is the set of all other assignments together with any non-success
terminal states in the presentation. The quantity \(\Phi_L(t)\) is the
probability that the budget-admissible operator has hit a satisfying
assignment by time \(t\). The defect score assigns a scalar value to
each assignment, and the hitting functional measures arrival at zero
defect.

\end{example}

\section{Killed Resolvent Representation}\label{the-resolvent-stopwatch}
\label{sec:killed-resolvent-representation}

\PartII introduced budget-admissible generators as transport operators.
The Laplace transform of their target-hitting flux is a killed resolvent
\citep{Norris1997,Kato1995}.

The target-projected resolvent is the Laplace-domain representation of discounted target arrival.
\begin{definition}[Target-projected resolvent]\label{def-target-resolvent}

Let \(L_N\), \(K_N=-L_N\), and \(k_T\) be as in Definition
\hyperref[def-killed-dynamics]{Killed dynamics at the target}. For
\(\lambda>0\), the target-projected resolvent is

\[
A_T(\lambda;L,\mu_0)
=
\mu_0(T)
+\mu_{0,N}(\lambda I_N-L_N)^{-1}k_T
=
\mu_0(T)
+\mu_{0,N}(\lambda I_N+K_N)^{-1}k_T.
\]

\end{definition}

Here \((\lambda I_N-L_N)^{-1}\) integrates all discounted non-target
histories, while \(k_T\) restricts to histories terminating upon entry
into the target. The terminal edge gate identifies the final crossing,
whereas the killed resolvent accounts for the entire discounted history
preceding that crossing.

The full-history transport record exposes the complete represented path contribution preceding target contact.
\begin{definition}[Full-history target-contact transport]\label{def-full-history-target-transport}

Write the off-diagonal rates of \(L\) as \(q_L(x,z)\). The
terminal-contact edge gate is

\[
W_T^L(x,z)=q_L(x,z)\mathbf 1_T(z),
\qquad
k_T^L(x)=\sum_z W_T^L(x,z).
\]

For an initial frontier \(\mu_0\) and discount \(\lambda>0\), the
full-history target-contact transport is the nonnegative measure on
target-entering edges

\[
\mathsf H_{\lambda,L,\mu_0}^T(x,z)
=
\int_0^\infty e^{-\lambda t}\,(\mu_{0,N}e^{tL_N})(x)\,W_T^L(x,z)\,dt.
\]

Its total mass is the noninitial part of the target-projected resolvent:

\[
\sum_{x\in N}\sum_{z\in T}\mathsf H_{\lambda,L,\mu_0}^T(x,z)
=
\mu_{0,N}(\lambda I_N-L_N)^{-1}k_T^L
=
A_T(\lambda;L,\mu_0)-\mu_0(T).
\]

The measure \(\mathsf H_{\lambda,L,\mu_0}^T\) is an operator-level
hitting quantity rather than a representation constructor. It separates
the killed propagation \(e^{tL_N}\) from the terminal edge gate
\(W_T^L\). The \PartII decomposition applies to the generator driving
the propagation, and the displayed formula integrates the corresponding
component projections.

\end{definition}

\begin{lemma}[Killed-resolvent identity for the discounted hitting functional]\label{thm-resolvent-stopwatch-identity}

For every \(\lambda>0\),

\[
A_T(\lambda;L,\mu_0)=S_T(\lambda;L,\mu_0).
\]

\end{lemma}

This identity is the standard Laplace-transform representation of a
killed semigroup resolvent \citep{Norris1997,Kato1995}.

\begin{proof}\phantomsection\label{proof-resolvent-stopwatch-identity}

The killed frontier on \(N\) is \(\mu_{0,N}e^{tL_N}\), so the
instantaneous post-initial hit flux is
\(\mu_{0,N}e^{tL_N}k_T\). The atom at \(\tau_T=0\) has mass
\(a_0=\mu_0(T)\). Hence

\[
S_T(\lambda)
=
a_0+\int_0^\infty e^{-\lambda t}\mu_{0,N}e^{tL_N}k_T\,dt
=
a_0+\mu_{0,N}\left(\int_0^\infty e^{-t(\lambda I_N-L_N)}\,dt\right)k_T.
\]

Because \(L_N\) is a finite sub-Markov generator, its spectrum lies in
the closed left half-plane.  Hence \(\lambda I_N-L_N\) is invertible
for every \(\lambda>0\), the exponentially discounted matrix integral
converges, and it equals \((\lambda I_N-L_N)^{-1}\).  This includes an
absorbing non-target failure state, for which \(0\) may belong to the
spectrum of \(L_N\).

\end{proof}

\Cref{fig-arrival-resolvent-stopwatch} summarizes the exact identity in
\cref{thm-resolvent-stopwatch-identity}.

\begin{figure}[tbp]
\centering
\begin{tikzpicture}[fw diagram]
  \node[fw source,minimum width=2.2cm] (init) at (-6.0,0)
    {initial frontier\\\(\mu_{0,N}\)};
  \node[fw geom,minimum width=2.6cm] (kill) at (-3.1,0)
    {killed propagation\\\(e^{tL_N}\)};
  \node[fw terminal,minimum width=2.3cm] (gate) at (0,0)
    {target gate\\\(k_T\)};
  \node[fw use,minimum width=2.5cm] (disc) at (3.3,1.15)
    {discounted hits\\\(S_T(\lambda)\)};
  \node[fw provenance,minimum width=2.8cm] (res) at (3.3,-1.15)
    {killed resolvent\\\(A_T(\lambda)\)};
  \node[fw box,minimum width=2.0cm] (eq) at (6.3,0)
    {\(S_T=A_T\)};
  \draw[fw arrow] (init) -- (kill);
  \draw[fw arrow] (kill) -- (gate);
  \draw[fw arrow] (gate) -- (disc);
  \draw[fw arrow] (gate) -- (res);
  \draw[fw arrow] (disc) -- (eq);
  \draw[fw arrow] (res) -- (eq);
\end{tikzpicture}
\caption{Discounted target arrival and the target-projected killed resolvent
are the same scalar functional by
\cref{thm-resolvent-stopwatch-identity}.}
\label{fig-arrival-resolvent-stopwatch}
\end{figure}

This comparison record contains the charged data needed to transfer an operator estimate to target hitting.
\begin{definition}[Budgeted hitting-functional comparison data]\label{def-budgeted-stopwatch-comparison-data}

Budgeted hitting-functional comparison data for a budget-admissible operator
\(L\) under a declared budget structure \(\mathcal B\), with operative
ceiling \(b(n)\), are finite exposed data whose charge is
\(\preceq b(n)\) and that verify one of the following admissible arrival
readings for the already fixed operator.  The first reading is available
only when \(\mathsf t(b(n))<\infty\):

\[
A_T(\mathsf t(b(n))^{-1};L,\mu_0)\succeq
\beta_{\mathfrak C}(n)^{-1},
\] \[
\Phi_L(h(n))\succeq \beta_{\mathfrak C}(n)^{-1}
\quad\text{for some }1\le h(n)\le\mathsf t(b(n)),
\] \[
\gamma_T(L)\succeq \beta_{\mathfrak C}(n)^{-1}
\quad\text{together with exposed resolvent/survival comparison data.}
\]

When \(\mathsf t(b(n))=\infty\), one instead uses finite stopping
slices, or a separately proved \(\lambda\downarrow0\) limit; the
resolvent expression is never evaluated at \(\lambda=0\).

The comparison data may be a direct target-resolvent computation, a
killed-gap/Dirichlet comparison, or an exact quotient comparison. Its
ingredients must be generated by the observable algebra and the \PartII
decomposition into exposed structural channels plus residual accounting.
These data belong to the analysis of the fixed operator and are not
algorithmic inputs or conditions for membership in the declared
budgeted operational class. They assign a
budgeted hitting bound to an exposed component of \(L\).

\end{definition}

Hitting-functional comparison data are analytic objects. Budget-admissible membership is uniform execution under the
declared resource envelope using only \(\mathcal R^{\mathcal B}_n\)
and its legal lifted closure. After that execution rule is embedded as a
generator, the killed-resolvent identity identifies the operator quantity that
must be large. If the large quantity is carried by exposed structure in
the budgeted closure, the corresponding projection of the generator
supplies the comparison data. If it is not carried by exposed structure,
the remaining operator term belongs to \(\Res\). Any target
arrival of the complete process is then charged through the exact
first-hit quotient rather than through a standalone residual estimate.

\begin{convention}[Analytic role of comparison data]\label{conv-proof-side-comparison-data}

Throughout this part, hitting-functional comparison data establish a
quantitative bound for an already fixed operator. The structural
component, saturated cost, membership in the declared budgeted
operational class, and existence of
quotient or geometric structure are properties of the presented task
and its induced operator. An oracle-free procedure may output only a candidate solution; its
transition rule still determines \(L\). If no comparison data have been
established, the corresponding quantitative estimate remains
unavailable, while the underlying structure is unchanged.

For example, a public linear quotient makes XOR-type instances low-cost
abstraction instances as a structural fact. Their membership in
polynomial time is realized by Gaussian elimination, which implements
structure already present in the observable presentation.

\end{convention}

\begin{proposition}[Budget-admissible procedures are covered before comparison data enter]\label{prop-budgeted-procedures-covered-before-comparison-data}

Let \(A=(A_n)\) be a uniform oracle-free procedure for the presented
task family whose transcript and instance-dependent data lie within budget
under \(\mathcal B\), in the presentation-relative sense of Convention
\hyperref[conv-presentation-relative-budgeted-computation]{Budgeted
computation is presentation-relative}. For each \(n\), record the clocked
finite outcome/terminal transcript of \(A_n\) and charge the current
memory, tables, work state, and output state at peak native capacity. Its
one-step execution rule defines a Markov kernel \(P_{A,n}^{[H]}\) on this
transcript carrier and a generator
\(L_{A,n}^{[H]}=\rho(P_{A,n}^{[H]}-I)\).
This construction depends only on the transition rule of \(A\); the
procedure need not output comparison data, a quotient, an invariant, or
a structural decomposition.

\end{proposition}

\begin{proof}\phantomsection\label{proof-budgeted-procedures-covered-before-comparison-data}

The resource envelope makes the clocked transcript carrier finite. The
code and represented finite outcome interface of \(A_n\), together with
public data generated from \(\mathcal R^{\mathcal B}_n\), determine the
conditional transcript transition law, hence determine
\(P_{A,n}^{[H]}\). \PartII applies the exhaustive decomposition to
\(L_{A,n}^{[H]}\). If \(A\) has non-negligible target arrival at the declared
margin, the killed-resolvent identity gives a non-negligible
target-resolvent quantity for this induced generator. Comparison data, when
used, charge that operator quantity to an exposed structural projection.
If no exposed projection carries it, the contribution remains
\(\Res\) residual accounting, or it uses
presentation-inaccessible or oracle-dependent information outside the
budgeted saturated observable closure. A program
whose transition rule depends on such outside information is analyzed
instead as an oracle/advice computation, an over-budget recovery
procedure, or a computation after changing the primitive signature so
that the datum is public. Thus coverage of presentation-relative \(A\)
is independent of any comparison data.

\end{proof}

\begin{lemma}[Quantitative arrival--resolvent comparison and conditional gap route]\label{thm-resolvent-gap-arrival-proof-equivalence}

For every \(B,c>0\), the arrival distribution and discounted
hitting functional satisfy

\[
e^{-1}\Phi_L(B)
\le S_T(B^{-1};L,\mu_0)
\le \Phi_L(cB)+e^{-c}.
\]

Consequently \(\Phi_L(B)\ge\varepsilon\) implies
\(S_T(B^{-1})\ge e^{-1}\varepsilon\), while
\(S_T(B^{-1})\ge\varepsilon\) implies
\(\Phi_L(cB)\ge\varepsilon/2\) for
\(c=\log(2/\varepsilon)\). A killed-gap estimate is a third proof route
only when its declared comparison data also bound the relevant survival,
initial-overlap, and target-flux terms so as to imply one of these two
quantities. Under transfer-class bounds on those comparison factors and
on \(c\), the three proof routes yield the same transfer-class
non-negligibility conclusion.

\end{lemma}

\begin{proof}\phantomsection\label{proof-resolvent-gap-arrival-proof-equivalence}

The first inequality is Proposition
\hyperref[prop-budget-hit-stopwatch]{Within-budget arrival bounds the
discounted hitting functional}. For the second, split according to
\(\{\tau_T\le cB\}\):

\[
\begin{aligned}
S_T(B^{-1})
&=\mathbb E\!\left[e^{-\tau_T/B}\mathbf1_{\{\tau_T<\infty\}}\right]\\
&\le \mathbb P(\tau_T\le cB)
   +e^{-c}\mathbb P(cB<\tau_T<\infty)\\
&\le \Phi_L(cB)+e^{-c}.
\end{aligned}
\]

The two consequences follow by substitution. The resolvent reading
equals \(S_T\) by Theorem
\hyperref[thm-resolvent-stopwatch-identity]{Killed-resolvent identity for
the discounted hitting functional}. A gap alone contains no information
about initial overlap or flux into \(T\); hence the gap route is valid
only with the additional comparison data stated in the theorem. Those
data imply an arrival or resolvent bound, after which the displayed
inequalities apply.

\end{proof}

\begin{example}[SAT resolvent]\label{ex-resolvent-sat}

For SAT, \(k_T(x)\) records the rate at which the operator sends a
non-satisfying assignment into a satisfying assignment. The expression
\(\mu_{0,N}(\lambda I-L_N)^{-1}k_T\) sums all discounted non-satisfying
histories that end by entering a satisfying assignment. A polynomial SAT
channel supports this quantity, or an equivalent killed-gap or quotient
reading, through the exposed clause structure, declared geometry,
quotient, lift, or boundary data.

\end{example}

\section{Structural Operator Ledger and Full-Process Hitting}\label{each-component-through-the-same-stopwatch}
\label{sec:structural-operator-ledger}

\subsection{Channel coordinates and full-process accounting}
\label{subsec:channel-full-process-accounting}

\PartII classifies the full generator in an operator Hilbert space.
Probabilistic hitting functionals remain attached to the full positive
generator unless a component is separately verified to lie in the
Markov-generator cone.

The channel coordinates place every structural component on the common full-process operator ledger.
\begin{definition}[Channel operator coordinates]\label{def-channel-stopwatch-components}

Let a budget-admissible generator decompose as

\[
L=L_{\Geom}+L_{\Caus}+L_{\Abs}+L_{\Lift}+L_{\Bdry}+L_J.
\]

For a component \(\xi\), define its operator energy by

\[
\mathcal E_\xi(L)=\|L_\xi\|_{\mathscr H_{\mathrm{ch}}}^2.
\]

The Pythagorean identity
\eqref{eq-channel-pythagorean-accounting} gives
\(\|L\|_{\mathscr H_{\mathrm{ch}}}^2
=\sum_\xi\mathcal E_\xi(L)\);
no Pythagorean claim is made for the induced operator norm. A quantity of the
form \(A_T(\lambda;L-L_\xi,\mu_0)\) is a hitting functional only when
\(L-L_\xi\) separately satisfies the finite generator form.

\end{definition}

These coordinates measure the exact operator decomposition without
assigning standalone transition probabilities to its summands.

\begin{proposition}[Roles of the five structural components and the residual term]\label{prop-channel-ledger}

For finite components with verified hitting-functional data, their
analytic roles are inherited from the \PartII decomposition.
The term \(\Res\) is the residual component rather than an
exposed structural component.

\[
\begin{array}{c|c}
\text{component} & \text{role in }A_T(\lambda)\\
\hline
\Geom & \text{intrinsic local geometry: conductance and Poincar\'e/Cheeger scale}\\
\Caus & \text{orientation/use of an exposed potential built from intrinsic structure}\\
\Abs & \text{intrinsic represented quotient abstraction}\\
\Lift & \text{valid re-expression of already charged intrinsic structure}\\
\Bdry & \text{terminal/support coloring or boundary value; no movement by itself}\\
\Res & \text{orthogonal operator remainder; full-process first-hit accounting}
\end{array}
\]

The table records operator provenance. It does not identify contextual
target charges: every Lift/Bdry occurrence remains in the complete
carrier, and a smaller-source charge bound requires a represented
contextual transfer.

\end{proposition}

\begin{proof}\phantomsection\label{proof-channel-ledger}

The fixed-space and component decompositions of \PartII identify the
origin of each term. Local symmetric conductance belongs to the
geometric/reversible component \(\Geom\) and contributes through
the killed Poincar\'e or Cheeger gap. Represented nonlocal jumps constant on
quotient cells have algebraic quotient/lumping provenance. When the quotient
passes \cref{def-paper1-abs-usefulness-gate}, its qualified \(\Abs\)
coordinate reduces the dimension of the killed problem.
Directed exact one-form components constitute the directed/drift
component \(\Caus\), given by the monotone use of edge
differences of an exposed
potential built from the observable geometry or abstraction. Lift is
functorial composition or enlargement, so its target-directed part is
supported on the structure present before lifting at the level of
operator provenance. Its contextual charge is nevertheless retained
until a represented transfer bounds it. Boundary terms
color, kill, weight, or read terminal states, with terminal removal as
the zero-value case; they have no off-diagonal transition support, so
they affect movement only through differences across an already
represented Geom/Abs/Lift edge structure that is then oriented by
\(\Caus\). The residual term is orthogonal to the declared
structural potentials and quotient/lift subspaces. The killed-resolvent
identity of \Cref{the-resolvent-stopwatch} is applied to the full
positive generator. The exact first-hit quotient then supplies the
retrospective accounting-profile identity for target arrival; its source
status is governed separately by
\cref{def-pre-hit-temporally-admissible-source}.

\end{proof}

\begin{lemma}[Structural ledger and saturated accountability]\label{thm-two-discount}

After the \PartII decomposition, transfer-class target advantage is
assigned to the intrinsic problem structure: exposed local geometry
\((\Geom)\) and represented quotient structure
\((\Abs)\). Directed descent, valid lifted-state constructions,
and boundary readout respectively orient, reorganize, and evaluate these
components. Boundary readout consists of scalar state data and has no
independent transition mechanism. The residual component
\(\Res\) is the orthogonal remainder in the operator ledger;
no standalone Markov or baseline interpretation is assigned to it. For
every affordable full generator, all target arrival---including effects
arising from interactions among the six operator coordinates---is bounded
by the saturated accounting quotient profile at the cost of the corresponding
exact first-hit construction. Prospective source attribution is instead
governed by \cref{thm-prospective-selector-accountability}.

\end{lemma}

\begin{proof}\phantomsection\label{proof-two-discount}

The intrinsic local component is \(\Geom\): it supplies the
observable comparison relation, conductance, and killed gap. The
intrinsic quotient component is \(\Abs\): it supplies represented
compression compatible with the terminal structure. \(\Caus\)
uses these sources by orienting edge differences of an exposed
potential; \(\Lift\) composes or enlarges existing sources and
retains that structure in its complete carrier; \(\Bdry\) supplies
terminal/support colorings and boundary values on states. The
conditional charges of Lift and Bdry transfer to a smaller source only
through a represented comparison. The residual
term is only an orthogonal operator coordinate; it need not be a positive
generator, and no hitting probability is attributed to it in isolation.
Instead, apply the killed resolvent to the full positive generator and
construct its exact first-hit quotient. That quotient is a uniform
represented exact-quotient accounting record whose cost includes the entire
computation, including all interactions with \(\Res\). The
universal saturated-profile theorem therefore charges the resulting
target arrival as an audit identity without a componentwise positivity
assumption. By \cref{cor-first-hit-accounting-not-source}, this completed
record does not supply prospective source attribution; that attribution uses
the neutral baseline and selector advantage in
\cref{thm-prospective-selector-accountability}.
The roles of \(\Lift\) and \(\Bdry\) in the structural
ledger remain as stated: lifting is charged to its represented
construction, while boundary data supplies no movement by itself.

\end{proof}

\subsection{Transformed gradients and directed drift}
\label{subsec:transformed-gradients-directed-drift}

The transformed-gradient space fixes the represented directions eligible for geometric comparison.
\begin{definition}[Admissible transformed-gradient space]\label{def-transformed-gradient-space}

Fix a saturated budget \(B\) and a presentation-relative generator on
its full lifted configuration space. Let \(E_{\mathrm{tr},\le B}\) be
the represented transformed edge, quotient, or lifted comparability
structure generated at budget \(B\) from intrinsic
\(\Geom\)/\(\Abs\) data and valid \(\Lift\)
operations. Let \(\mathcal V_{\mathrm{tr},\le B}\) be the represented
scalar state-coloring space on that structure: pointwise Geom
potentials, quotient coordinates, lifted coordinates, and
\(\Bdry\) terminal/support colorings or boundary values, all
generated from the same budgeted observable algebra. The admissible information
flow is

\[
\begin{gathered}
\mathcal R^{\mathcal B}_n
\xrightarrow{\Geom/\Abs}
\text{represented edge or quotient structure}\\
\xrightarrow{\Lift}
\text{valid lifted comparability structure}\\[-1mm]
\text{and}\\[-1mm]
\Bdry\text{ supplies scalar terminal/support values on it.}
\end{gathered}
\]

The boundary/killing component \(\Bdry\) is scalar: it marks,
kills, weights, or evaluates states and terminal sectors. It contributes
no off-diagonal edge and no transport operator. Transport appears only after
\(E_{\mathrm{tr},\le B}\) supplies represented comparable pairs and a
scalar \(V\in\mathcal V_{\mathrm{tr},\le B}\) supplies differences
across those pairs. With represented nonnegative edge weights
\(a_B(x,z)\) supported on \(E_{\mathrm{tr},\le B}\), the associated
transformed-gradient edge space is

\[
\nabla_{E_{\mathrm{tr},\le B}}\mathcal V_{\mathrm{tr},\le B}
=
\{(x,z)\mapsto a_B(x,z)(V(z)-V(x)):V\in\mathcal V_{\mathrm{tr},\le B}\}.
\]

The edge inner product is the one supplied by the \PartII Hodge split on
the represented current support. A directed edge current is a
\(\Caus\) contribution at budget \(B\) exactly through its
orthogonal projection onto
\(\nabla_{E_{\mathrm{tr},\le B}}\mathcal V_{\mathrm{tr},\le B}\). This
definition states that the directed/drift component \(\Caus\)
follows gradients over represented geometric, quotient, or lifted-state
structure; \(\Bdry\) may assign the state values defining the
gradient but contributes no transition by itself.

\end{definition}

\Cref{fig-caus-hodge-projection} shows how the transformed-gradient space
separates authorized direction from the orthogonal operator remainder.

\begin{figure}[tbp]
\centering
\begin{tikzpicture}[fw diagram]
  \node[fw source,minimum width=3.1cm] (data) at (-4.7,1.1)
    {represented\\\(\Geom/\Abs/\Lift\) structure};
  \node[fw geom,minimum width=3.0cm] (grad) at (-1.0,1.1)
    {gradient space\\\(\nabla_E\mathcal V\)};
  \node[fw use,minimum width=3.0cm] (current) at (-1.0,-1.0)
    {directed current\\\(J_{\mathrm{dir}}\)};
  \node[fw terminal,minimum width=2.7cm] (caus) at (3.2,1.1)
    {authorized \(\Caus\)\\projection};
  \node[fw residual,minimum width=2.7cm] (res) at (3.2,-1.0)
    {cycle and\\orthogonal remainder};
  \draw[fw arrow] (data) -- (grad);
  \draw[fw arrow] (grad) -- (caus);
  \draw[fw arrow] (current) -- node[right,font=\scriptsize] {project} (caus);
  \draw[fw arrow] (current) -- (res);
\end{tikzpicture}
\caption{The Hodge projection of
\cref{def-transformed-gradient-space,thm-directed-drift-transformed-gradient-normal-form}.
Only the component of a represented current lying in the transformed-gradient
space receives \(\Caus\) status; the complementary components remain in their
cycle or residual coordinates.}
\label{fig-caus-hodge-projection}
\end{figure}

\begin{lemma}[Directed drift has transformed-gradient normal form]\label{thm-directed-drift-transformed-gradient-normal-form}

Let \(L_n\) be any presentation-relative uniform oracle-free
budget-admissible generator, represented on its finite \PartII carrier
(the clocked transcript carrier for an operational computation). Let
\(J_{\mathrm{dir}}\) be any directed or
antisymmetric target-relevant current appearing in the killed
full-history accounting after the symmetric Geom and exact Abs quotient
parts have been separated. For each saturated budget \(B\) there is an
orthogonal decomposition

\[
J_{\mathrm{dir}}
=
J_{\Caus,B}
+
J_{\mathrm{cyc},B}
+
J_{\mathrm{res},B}^{\mathrm{dir}},
\]

where \(J_{\Caus,B}\) lies in
\(\nabla_{E_{\mathrm{tr},\le B}}\mathcal V_{\mathrm{tr},\le B}\),
\(J_{\mathrm{cyc},B}\) is the divergence-free cycle component relative
to the represented transformed edge structure, and
\(J_{\mathrm{res},B}^{\mathrm{dir}}\) is orthogonal to every transformed
gradient generated at budget \(B\). The first component is a
\(\Caus\) channel. Its complete carrier charges the Geom/Abs
data, every Lift that supplies the represented edge structure, and every
Bdry scalar coloring or boundary value used in the potential. This is
cost accounting, not an automatic transfer of target charge between
extensions. The cycle
component contributes to target arrival only through represented
terminal contact across the same edge structure; a Bdry readout may
record or kill a terminal state, but it does not move mass. Its
contextual charge remains explicit unless a represented comparison
transfers that charge to the source it reads. Any transfer-class
target excess left in \(J_{\mathrm{res},B}^{\mathrm{dir}}\) is either
represented by a budget-admissible saturated observable and therefore
reclassified by \PartII into Geom/Caus/Abs/Lift/Bdry, or has no such
one-step representative at budget \(B\). In the latter case it remains
an orthogonal operator remainder; full-process arrival is charged at the
cost of the first-hit derivation. Presentation-inaccessible or
oracle-dependent data, nonuniform advice, and over-budget recovery are
inadmissible for the original presentation.

Consequently, a represented directed source at budget \(B\) consists of
a construction from the budgeted represented algebra of a
represented edge/quotient/lift geometry together with an admissible
scalar coloring whose edge differences point to the target. If no such
transformed gradient or represented quotient is available at that
budget, the directed operator contribution remains outside the
represented source coordinates. The universal profile theorem accounts
for any success of the complete computation at its enlarged
finite-horizon cost.

\end{lemma}

\begin{proof}\phantomsection\label{proof-directed-drift-transformed-gradient-normal-form}

Work on the finite edge Hilbert space of the directed current support,
with the positive edge weights used in the \PartII Hodge split. By
definition, \(E_{\mathrm{tr},\le B}\) is a represented
edge/quotient/lift support and \(\mathcal V_{\mathrm{tr},\le B}\) is a
finite-dimensional scalar state-coloring space on it. Therefore
\(\nabla_{E_{\mathrm{tr},\le B}}\mathcal V_{\mathrm{tr},\le B}\) is a
closed finite-dimensional subspace of the edge space. Let
\(J_{\Caus,B}\) be the orthogonal projection of
\(J_{\mathrm{dir}}\) onto this transformed-gradient subspace. This is
exactly the structural-gradient component in \PartII's Hodge split, with
the edge support enlarged only by admissible Abs and Lift
transformations and with Bdry entering only as scalar terminal/support
values already present in the saturated closure. Hence it is
\(\Caus\), and its carrier cost includes the represented edge
structure, scalar potential, and Geom/Abs data that produced them. This
provenance statement does not transfer contextual target charge.

Inside the orthogonal complement, take the finite Hodge decomposition
relative to the represented transformed graph with boundary: the
divergence-free part is \(J_{\mathrm{cyc},B}\), and the remaining part
is \(J_{\mathrm{res},B}^{\mathrm{dir}}\). A divergence-free cycle has no
net boundary flux in the represented quotient/geometry. A represented
Bdry coloring can record, stop, or weight terminal states, but since it
has no off-diagonal support it cannot turn a cycle into motion. If
terminal contact is accountable, it is because the same represented edge
structure supplies terminal-crossing differences or terminal-contact
edges already charged in Geom/Abs/Lift and then read by Bdry. Without
such represented terminal contact, the cyclic part has no accountable
target descent.

It remains to classify \(J_{\mathrm{res},B}^{\mathrm{dir}}\). By
construction it is orthogonal to every admissible transformed gradient
at budget \(B\) and has no represented quotient or boundary readout
already removed. If its target excess is represented by an admissible
observable relation, potential, quotient, lifted interface, scalar
boundary readout, represented propagation component, or legal
finite-horizon trajectory statistic, \PartII's
no-uniform-residual-correlation theorem reclassifies that representative
into one of the five structural components. The reclassified contribution
is then counted in \(m_n^{\mathrm{sat}}\) or \(G_n^{\mathrm{sat}}\) at
its least saturated carrier cost. Lift and Bdry remain represented
contextual extensions; moving their target charge to another source
requires a charged comparison satisfying the contextual-transfer
inequality. If no such
representative exists, the directed target excess is not legal
budget-admissible structure for the original presentation; it is
an orthogonal residual at this budget, or it uses
presentation-inaccessible or oracle-dependent data, nonuniform advice,
over-budget recovery, or data outside the budgeted represented algebra.
These alternatives exhaust the \PartII component theorem and the
budgeted derivability closure. Any target arrival of their full
composition is subject to \Cref{thm-universal-saturated-profile-accountability}.

\end{proof}

\begin{example}[Structural components for SAT]\label{ex-channel-stopwatch-sat}

A one-bit Hamming graph provides the geometric/reversible component for
SAT. Unit propagation or implication closure provides a directed/drift
component when the implications are exposed and target-aligned. Public
XOR clauses provide a quotient/lumping component
because the syndrome quotient collapses assignments into represented
affine cells. Adding auxiliary variables is Lift. Detecting \(D_\Phi=0\)
is boundary/killing data. Random restarts without structural alignment
belong to the residual term \(\Res\).

\end{example}

\section{The Walsh/Spectral Bound}\label{the-walshspectral-bound}
\label{sec:walsh-spectral-bound}

The spectral bound restricts a low-degree comparison vector to the
low-degree projection of the target. Consequently, low-degree target
mass bounds the discounted hitting estimate obtainable from
degree-\(d\) structure.

The degree projection measures how much target contrast lies in the declared low-degree observable band.
\begin{definition}[Degree projection and spectral mass]\label{def-degree-projection-spectral-mass}

On \(X=\{0,1\}^n\) with uniform measure, let

\[
\chi_S(x)=(-1)^{\sum_{i\in S}x_i},
\qquad S\subseteq\{1,\ldots,n\},
\]

be the Walsh characters \citep{Walsh1923,ODonnell2014}. For a centered
target contrast \(y\in L^2(X)\),
write

\[
y=\sum_{S\subseteq[n]}\widehat y(S)\chi_S,
\qquad
\widehat y(S)=\mathbb E_x[y(x)\chi_S(x)].
\]

The degree-\(\le d\) projection and spectral mass are

\[
P_dy=\sum_{|S|\le d}\widehat y(S)\chi_S,
\qquad
m(d)=\frac{\|P_dy\|_{L^2(\mu)}^2}{\|y\|_{L^2(\mu)}^2}
=
\frac{\sum_{|S|\le d}|\widehat y(S)|^2}{\sum_S|\widehat y(S)|^2}.
\]

For a general finite observable algebra, \(P_d\) denotes orthogonal
projection onto the chosen degree-\(\le d\) observable subspace
\(V_d\subseteq L^2(X_I,\mu_I)\). The invariant later optimizes over
admissible public basis changes.

\end{definition}

The quantity \(m(d)\) is the fraction of the target contrast contained
in the observable subspace of degree at most \(d\). If \(m(d)\) is
negligible, then degree-\(d\) observables have asymptotically equal means
on the target and its complement.

The comparison vector supplies the represented low-degree direction used in the Walsh bound.
\begin{definition}[Degree-d target comparison vector]\label{def-degree-d-target-comparison-vector}

A degree-\(d\) target comparison vector for a structural component
consists of an exposed unit vector \(b\in V_d\), a transfer-class
comparison factor \(C_{\mathrm{comp}}(n)\), and a verified discounted hitting
strength \(\Gamma_T(L)\) such that the comparison proves

\[
0\le \Gamma_T(L)
\le
C_{\mathrm{comp}}(n)
\left|\langle \bar y,b\rangle_{L^2(\mu)}\right|^2.
\]

Here \(\Gamma_T(L)\) may be the verified killed gap, a verified
target-resolvent lower-bound strength, or the equivalent verified
arrival strength from \Cref{the-resolvent-stopwatch}. If no such exposed \(b\) and
verification exist, then no degree-\(d\) budgeted hitting contribution
has been charged to that structural component.

\end{definition}

This condition is required when applying the theorem. A
presentation-relative budget-admissible algorithm supplies a uniform
transition rule whose
instance-dependent data are legal for the presented task. To prove a
budgeted hitting contribution at raw observable degree \(d\) in this
Walsh-bound subcase, the argument must exhibit an exposed low-degree
statistic in the induced generator and the corresponding inequality.
Without that proof, no raw degree-\(d\) contribution has been
established in the application; the objective saturated cost of the
component is unchanged.

The accounted degree condition ties a low-degree operator contribution to an exposed statistic and its charged derivation.
\begin{definition}[Degree-d accounted budget-admissible operator]\label{def-degree-d-admissible-operator}

A budget-admissible killed operator is accounted at effective degree at
most \(d\) when there is a proof that every target-directed structural
component of \(L\) is generated by observables in \(V_d\), after any
charged public basis change, and every claimed budgeted hitting
contribution has a degree-\(d\) target comparison vector. This is an
analytic classification of the operator rather than a membership
condition for the algorithm that induced it.

\end{definition}

\begin{lemma}[Walsh/spectral comparison-vector bound]\label{thm-gap-spectral-bound}

For every degree-\(d\) accounted budget-admissible operator \(L\),

\[
\Gamma_T(L)
\le
C_{\mathrm{comp}}(n)m(d),
\qquad
C_{\mathrm{comp}}\in\mathfrak C.
\]

In particular, a transfer-class verified hitting estimate at degree \(d\)
requires transfer-class target mass in the degree-\(\le d\) observable
band. In the polynomial instance this is the usual inverse-polynomial
statement.

\end{lemma}

\begin{proof}\phantomsection\label{proof-gap-spectral-bound}

Let \(b\in V_d\) be the exposed unit vector in the comparison. Since
\(P_d\) is orthogonal projection onto \(V_d\),

\[
\langle \bar y,b\rangle
=
\langle P_d\bar y,b\rangle.
\]

By Cauchy-Schwarz and \(\|b\|\le1\),

\[
\left|\langle \bar y,b\rangle\right|^2
\le
\|P_d\bar y\|^2
=
\frac{\|P_dy\|^2}{\|y\|^2}
=m(d).
\]

Substituting this inequality into the verified comparison inequality
gives the displayed bound.

\end{proof}

\begin{remark}[Basis robustness]\label{rem-basis-robustness}

The relevant degree is the degree in the public observable algebra after
admissible public basis changes and represented quotients are charged. A
raw-coordinate Walsh calculation is a diagnostic; the invariant uses the
basis that the presentation exposes.

\end{remark}

\begin{example}[XOR-SAT and basis dependence]\label{ex-walsh-bound-xor}

An XOR-SAT instance may have a target contrast that occupies degree
three in raw coordinate Walsh characters when one equation involves
three variables. The public syndrome observable represents that same
equation as one syndrome coordinate. In the public syndrome basis, the
target-facing proof vector is degree one, so the admissible invariant
reads the problem through the syndrome quotient.

\end{example}

\section{Quantitative Geom--Abs Interaction: Quotient Defect and Local
Compensation}\label{quotient-defect-and-geom-form-comparison}
\label{sec:quotient-defect-local-compensation}

\PartI defined static lumpability as compatibility between represented
quotient cells and terminal structure. \PartII expresses quotient
dynamics through intertwining. The present section measures how
accurately a represented quotient describes the killed dynamics, and
when local geometry controls the modes omitted by that quotient.

\subsection{Quotient defects and exact resolvent identities}
\label{subsec:quotient-defects-resolvent-identities}

We first compare the ambient and quotient generators through their represented
lift, projection, and lumpability defect.

\begin{definition}[Quotient lift, projection, and quotient generator]\label{def-quotient-lifts}

Let \(q:X_I\to Q\) be a represented quotient of saturated cost at most
\(B\). Let \(H=L^2(X_I,\mu_I)\) and \(H_Q=L^2(Q,\nu_Q)\), where
\(\nu_Q=q_*\mu_I\). The lift and projection are

\[
U_q:H_Q\to H,
\qquad
(U_qf)(x)=f(q(x)),
\] \[
P_q:H\to H_Q,
\qquad
(P_qg)(a)=\mathbb E_{\mu_I}[g(X)\mid q(X)=a].
\]

For a killed generator \(L\) on \(X_I\), the quotient generator induced
by \(q\) is

\[
L_Q=P_qLU_q.
\]

\end{definition}

The lift turns a quotient observable into a state observable by making
it constant on every quotient cell. The projection averages a state
observable over each cell. With the displayed conditional-expectation
projection, \(U_q\) is an isometry and \(P_q=U_q^*\). The quotient
generator is the compression of the original generator to the quotient
subspace.

The lumpability defect measures the failure of quotient lifting to intertwine the full and compressed generators.
\begin{definition}[Lumpability defect]\label{def-lumpability-defect}

The lumpability defect of \(q\) is the operator

\[
E_q=LU_q-U_qL_Q:H_Q\to H.
\]

Exact lumpability means \(E_q=0\), equivalently

\[
LU_q=U_qL_Q.
\]

\end{definition}

\begin{proposition}[Unit transfer for a canonical exact quotient]
\label{prop-canonical-exact-quotient-unit-transfer}

Let \(q:X\to Q\) be a finite terminal-compatible quotient on a finite
carrier equipped with a full-support reference law \(\mu\), equip
\(Q\) with the canonical law \(\nu_Q=q_*\mu\), and let
\(U_qf=f\circ q\).  Let \(L\) and \(L_Q\) be Markov or sub-Markov
generators satisfying

\[
LU_q=U_qL_Q.
\]

Suppose that \(T=q^{-1}(T_Q)\), that the target sectors are absorbing,
and that the quotient initialization is
\(\bar\mu_0=q_*\mu_0\).  Then \(U_q\) is isometric,
\(P_q=U_q^*\), and

\begin{equation}
\label{eq-canonical-exact-quotient-semigroup-transfer}
e^{tL}U_q=U_qe^{tL_Q}
\qquad(t\ge0).
\end{equation}

The full and quotient target-hitting laws, and hence their discounted
target functionals, agree exactly:

\begin{equation}
\label{eq-canonical-exact-quotient-hitting-transfer}
\Pr_{\mu_0}\{\tau_T\le t\}
=
\Pr_{\bar\mu_0}\{\tau_{T_Q}\le t\},
\qquad
A_T(\lambda;L,\mu_0)
=A_T^Q(\lambda;L_Q,\bar\mu_0).
\end{equation}

Consequently the exact-quotient comparison and normalization factor is
one for the canonical record:

\begin{equation}
\label{eq-canonical-exact-quotient-unit-factor}
C_q=1.
\end{equation}

Representation, simulation, and fiber-evaluation overhead contributes
to saturated cost and does not introduce a numerical transfer loss.

\end{proposition}

\begin{proof}

By the definition of the pushforward law,

\[
\|U_qf\|_{L^2(\mu)}^2
=
\int_X|f\circ q|^2\,d\mu
=
\int_Q|f|^2\,d\nu_Q
=
\|f\|_{L^2(\nu_Q)}^2.
\]

Thus \(U_q^*U_q=I\), and the conditional-expectation formula in
\Cref{def-quotient-lifts} gives \(P_q=U_q^*\).  Exact
intertwining passes to every power of the generators and therefore to
their exponential series, proving
\eqref{eq-canonical-exact-quotient-semigroup-transfer}.  Terminal
compatibility gives
\(\mathbf 1_T=U_q\mathbf 1_{T_Q}\).  Since both sectors are absorbing,
their occupation probabilities at time \(t\) equal their first-hit
probabilities by time \(t\).  Hence

\[
\begin{aligned}
\Pr_{\mu_0}\{\tau_T\le t\}
&=\mu_0e^{tL}\mathbf 1_T\\
&=\mu_0e^{tL}U_q\mathbf 1_{T_Q}\\
&=\bar\mu_0e^{tL_Q}\mathbf 1_{T_Q}\\
&=\Pr_{\bar\mu_0}\{\tau_{T_Q}\le t\}.
\end{aligned}
\]

Taking the Laplace--Stieltjes transform proves the functional identity.
All norm, initialization, dynamic, and terminal comparisons used by the
canonical exact-quotient record are equalities, which proves
\eqref{eq-canonical-exact-quotient-unit-factor}.

\end{proof}

Exact and approximate lumpability in this finite-state setting agree
with the usual Markov-chain aggregation notions
\citep{KemenySnell1976,Buchholz1994}. The defect is an operator
\(H_Q\to H\); geometric compensation must therefore control its action
on a space of quotient observables.

The component defect record resolves lumpability failure into symmetric and causal parts.
\begin{definition}[Symmetric and causal quotient defects]
\label{def-symmetric-causal-quotient-defects}

Let the killed generator selected from a Caus authority envelope be

\[
L=S+A,
\qquad S^*=S,
\qquad A^*=-A,
\]

where diagonal killing is included in \(S\).  Define the compressed
operators

\[
S_Q=P_qSU_q,
\qquad
A_Q=P_qAU_q,
\]

and the component defects

\[
E_q^{\mathsf s}=SU_q-U_qS_Q,
\qquad
E_q^{\mathsf a}=AU_q-U_qA_Q.
\]

Then

\begin{equation}
\label{eq-symmetric-causal-defect-sum}
L_Q=S_Q+A_Q,
\qquad
E_q=E_q^{\mathsf s}+E_q^{\mathsf a}.
\end{equation}

The first term is the symmetric or geometric compatibility defect.  The
second is the causal transport defect of the selected imposed/control
operator.  A reversible realization has \(A=0\) and hence
\(E_q^{\mathsf a}=0\).  An imposed nonreversible realization must retain
\(A_Q=P_qAU_q\); choosing a different quotient skew operator changes the
dynamic realization and is not a quotient comparison for the original
authority envelope.

\end{definition}

\begin{proposition}[Orthogonality of the compressed authority split]
\label{prop-compressed-authority-split}

With the conditional-expectation projection \(P_q=U_q^*\), one has

\[
S_Q^*=S_Q,
\qquad
A_Q^*=-A_Q,
\qquad
\operatorname{Re}\operatorname{tr}(S_Q^*A_Q)=0.
\]

Thus quotient compression preserves the reversible/nonreversible
operator split, and the two compressed components are orthogonal in the
real Hilbert--Schmidt inner product. Exact full lumpability is equivalent to

\[
E_q^{\mathsf s}=-E_q^{\mathsf a};
\]

componentwise causal compatibility is the stronger condition
\(E_q^{\mathsf s}=E_q^{\mathsf a}=0\).  Every authority-preserving
certificate records the component defects even when their sum cancels.

\end{proposition}

\begin{proof}

Since \(P_q=U_q^*\),

\[
S_Q^*=U_q^*S^*U_q=S_Q,
\qquad
A_Q^*=U_q^*A^*U_q=-A_Q.
\]

Moreover,

\[
\overline{\operatorname{tr}(S_Q^*A_Q)}
=\operatorname{tr}(A_Q^*S_Q)
=-\operatorname{tr}(A_QS_Q)
=-\operatorname{tr}(S_QA_Q),
\]

where cyclicity of the finite-dimensional trace is used in the last
equality. Hence \(\operatorname{tr}(S_Q^*A_Q)\) is purely imaginary and
its real part is zero.
Equation~\eqref{eq-symmetric-causal-defect-sum} gives the remaining
statements.  Recording the two terms prevents an exact full
intertwining caused by cancellation from being used as evidence that
the imposed causal and symmetric components descend separately.

\end{proof}

\begin{proposition}[Exact resolvent defect identity]\label{prop-defect-resolvent-perturbation}

For every \(\lambda>0\) in the common resolvent set of \(L\) and
\(L_Q\),

\[
(\lambda I-L)^{-1}U_q-U_q(\lambda I-L_Q)^{-1}
=
(\lambda I-L)^{-1}E_q(\lambda I-L_Q)^{-1}.
\]

Consequently,

\[
\left\|(\lambda I-L)^{-1}U_q-U_q(\lambda I-L_Q)^{-1}\right\|
\le
\|(\lambda I-L)^{-1}\|\,\|E_q\|\,\|(\lambda I-L_Q)^{-1}\|.
\]

\end{proposition}

The displayed equality is the resolvent identity in perturbation form
\citep{Kato1995}.

\begin{proof}\phantomsection\label{proof-defect-resolvent-perturbation}

From the definition of \(E_q\),

\[
(\lambda I-L)U_q-U_q(\lambda I-L_Q)
= -E_q.
\]

Multiplying on the left by \((\lambda I-L)^{-1}\) and on the right by
\((\lambda I-L_Q)^{-1}\) gives the identity. The norm bound is
submultiplicativity.

\end{proof}

The normalized defect measures the target-relevant loss in a nonreversible quotient resolvent comparison.
\begin{definition}[Normalized nonreversible resolvent defect]
\label{def-normalized-resolvent-defect}

Let \(W_Q\le H_Q\) and \(W\le H\) be represented comparison windows,
with orthogonal projections \(\Pi_Q\) and \(\Pi_W\). Set

\[
b_\lambda(q)
=
\|\Pi_WU_q(\lambda I-L_Q)^{-1}\Pi_Q\|,
\]

and

\[
d_\lambda(q)
=
\|\Pi_W(\lambda I-L)^{-1}E_q
  (\lambda I-L_Q)^{-1}\Pi_Q\|.
\]

The normalized resolvent defect is

\[
\Delta_\lambda(q;W_Q,W)
=
\begin{cases}
d_\lambda(q)/b_\lambda(q),&b_\lambda(q)>0,\\
0,&b_\lambda(q)=d_\lambda(q)=0,\\
\infty,&b_\lambda(q)=0<d_\lambda(q).
\end{cases}
\]

For a target-arrival comparison, the certificate must additionally
specify a represented functional \(\ell_T\in W^*\) and input
\(v_T\in W_Q\) corresponding to its initial-law and terminal-flux
reading. Define

\[
\begin{aligned}
b_\lambda^T(q)
&=
\left|
\ell_T\!\left(
\Pi_WU_q(\lambda I-L_Q)^{-1}\Pi_Qv_T
\right)
\right|,\\
d_\lambda^T(q)
&=
\left|
\ell_T\!\left(
\Pi_W(\lambda I-L)^{-1}E_q
(\lambda I-L_Q)^{-1}\Pi_Qv_T
\right)
\right|,
\end{aligned}
\]

and define \(\Delta_\lambda^T=d_\lambda^T/b_\lambda^T\) with the same
zero-denominator conventions as above. Its target-transfer stability
factor is

\[
S_{q,T}^{\mathrm{nr}}(\lambda)
=
\max\{0,1-\Delta_\lambda^T(q;\ell_T,v_T)\}.
\]

All windows, resolvent bounds, and retained operator data used to
evaluate these numbers belong to the charged comparison certificate. In
particular, \(\Delta_\lambda\) or \(\Delta_\lambda^T\) is admissible at
budget \(B\) only when these data have a derivation of cost at most
\(B\); the exact resolvents are not free constructors.

\end{definition}

\begin{corollary}[General quotient reconstruction bound]
\label{cor-defect-caps-quotient-gap}

On the represented windows of
\Cref{def-normalized-resolvent-defect},

\[
\begin{aligned}
&\|\Pi_W\bigl((\lambda I-L)^{-1}U_q
-U_q(\lambda I-L_Q)^{-1}\bigr)\Pi_Q\|\\
&\hspace{30mm}
\le
\Delta_\lambda(q;W_Q,W)b_\lambda(q).
\end{aligned}
\]

Consequently,

\[
\|\Pi_W(\lambda I-L)^{-1}U_q\Pi_Q\|
\le
(1+\Delta_\lambda(q;W_Q,W))b_\lambda(q),
\]

and, when \(\Delta_\lambda<1\), the same norm is at least
\((1-\Delta_\lambda)b_\lambda(q)\). Thus quotient resolvent estimates
transfer to the full process with the explicit additive loss
\(d_\lambda(q)\). A lower transfer certificate is nontrivial precisely
when \(\Delta_\lambda<1\), with loss
\((1-\Delta_\lambda)^{-1}\); the upper reconstruction bound has factor
\(1+\Delta_\lambda\). For a general
nonreversible generator this statement is the available conclusion; a
killed-gap comparison requires additional self-adjoint or sector
hypotheses.

For the represented target functional in
\cref{def-normalized-resolvent-defect}, the
same identity and the reverse triangle inequality also give

\[
\left|
\ell_T\!\left(\Pi_W(\lambda I-L)^{-1}U_q\Pi_Qv_T\right)
\right|
\ge
\bigl(1-\Delta_\lambda^T(q;\ell_T,v_T)\bigr)
b_\lambda^T(q)
\]

whenever \(\Delta_\lambda^T<1\). This scalar inequality, rather than
the operator-norm reconstruction bound alone, is the lower-transfer
record used in a target-visibility certificate.

\end{corollary}

\begin{proof}\phantomsection\label{proof-defect-caps-quotient-gap}

Project the identity in
\Cref{prop-defect-resolvent-perturbation} onto \(W_Q\) and
\(W\), then apply the definitions of \(d_\lambda\) and
\(\Delta_\lambda\). The zero-denominator conventions cover the two
degenerate cases. Applying the represented functional to the same
identity gives a scalar quotient term plus an error term of modulus
\(d_\lambda^T\); reverse triangle and
\(d_\lambda^T=\Delta_\lambda^Tb_\lambda^T\) prove the final inequality.

\end{proof}

\section{Reversible Defect Compensation}
\label{sec:reversible-defect-compensation}

In the reversible setting, Schur-complement control converts the defect bound
into quotient-to-ambient resolvent estimates.

\subsection{Compensated defects and Schur complements}
\label{subsec:compensated-defects-schur-complements}

\begin{definition}[Reversible compensated defect]
\label{def-reversible-compensated-defect}

Let \(N=X\setminus T\), let \(H_N=L^2(N,\mu|_N)\), and suppose the
killed operator \(K=-L_N\) is self-adjoint and positive semidefinite.
Let \(U:H_Q\to H_N\) be the isometric lift of a terminal-compatible
quotient, and set \(\Pi=UU^*\). Relative to
\(H_N=\operatorname{ran}U\oplus\operatorname{ran}U^\perp\), define

\[
A=U^*KU,
\qquad
B=(I-\Pi)KU,
\qquad
D=(I-\Pi)K(I-\Pi)|_{\operatorname{ran}U^\perp}.
\]

Then

\[
K=
\begin{pmatrix}
A&B^*\\
B&D
\end{pmatrix}.
\]

For \(\lambda>0\), the reversible compensated-defect parameter is

\begin{equation}
\label{eq-reversible-compensated-defect}
\theta_\lambda(q)
=
\left\|(D+\lambda I)^{-1/2}
B(A+\lambda I)^{-1/2}\right\|^2.
\end{equation}

The associated form-conditioning factor is

\[
S_q^{\mathrm{rev}}(\lambda)=\max\{0,1-\theta_\lambda(q)\}.
\]

This factor controls the Loewner comparison in the theorem below. It is
not, by itself, a lower bound for an arbitrary cross pairing
\(\langle u,(K+\lambda I)^{-1}v\rangle\).

The operator \(A\) measures quotient geometry, \(D\) measures
within-fiber geometry, and \(B\) is the coupling between them. Since
\(L_Q=-A\) and \(E_q=-B\), with the latter embedded in
\(\operatorname{ran}U^\perp\), exact lumpability is equivalent to
\(B=0\).

\end{definition}

\begin{lemma}[Schur-complement compensation theorem]
\label{thm-schur-complement-compensation}

Under \Cref{def-reversible-compensated-defect},

\begin{equation}
\label{eq-schur-complement-resolvent}
U^*(K+\lambda I)^{-1}U
=
\left(
A+\lambda I-B^*(D+\lambda I)^{-1}B
\right)^{-1}.
\end{equation}

If \(\theta_\lambda(q)<1\), then

\begin{equation}
\label{eq-schur-complement-order}
(A+\lambda I)^{-1}
\preceq
U^*(K+\lambda I)^{-1}U
\preceq
\frac{1}{1-\theta_\lambda(q)}(A+\lambda I)^{-1}.
\end{equation}

In particular, with
\(R_q=(A+\lambda I)^{-1}\) and
\(R_{\mathrm{full}}^q=U^*(K+\lambda I)^{-1}U\),

\[
0\preceq R_{\mathrm{full}}^q-R_q
\preceq
\frac{\theta_\lambda(q)}{1-\theta_\lambda(q)}R_q.
\]

Consequently, for \(u,v\in H_Q\),

\[
|\langle u,(R_{\mathrm{full}}^q-R_q)v\rangle|
\le
\frac{\theta_\lambda(q)}{1-\theta_\lambda(q)}
\langle u,R_qu\rangle^{1/2}
\langle v,R_qv\rangle^{1/2}.
\]

\end{lemma}

\begin{proof}

The upper-left block of the inverse of \(K+\lambda I\) is the inverse
of its Schur complement, which gives
\eqref{eq-schur-complement-resolvent}. Let \(X=A+\lambda I\) and
\(Y=D+\lambda I\). By
\eqref{eq-reversible-compensated-defect},

\[
0\preceq B^*Y^{-1}B\preceq\theta_\lambda(q)X.
\]

Hence

\[
(1-\theta_\lambda(q))X
\preceq X-B^*Y^{-1}B\preceq X.
\]

Both sides are positive definite when \(\theta_\lambda(q)<1\).
Inverting reverses Loewner order and proves
\eqref{eq-schur-complement-order}. The operator difference is positive
and bounded above by
\(\theta_\lambda(q)(1-\theta_\lambda(q))^{-1}R_q\). Cauchy--Schwarz for its
quadratic form gives the final estimate.

\end{proof}

\subsection{Target-functional and killed-gap transfer}
\label{subsec:target-functional-killed-gap-transfer}

\begin{corollary}[Target-functional reversible transfer]
\label{cor-target-functional-reversible-transfer}

With \(R_q\) and \(R_{\mathrm{full}}^q\) as above, fix represented
target-comparison vectors \(u_T,v_T\in H_Q\), and put

\[
b_{\lambda,\mathrm{rev}}^T
=|\langle u_T,R_qv_T\rangle|,
\qquad
d_{\lambda,\mathrm{rev}}^T
=|\langle u_T,(R_{\mathrm{full}}^q-R_q)v_T\rangle|.
\]

Define
\(\vartheta_\lambda^T=d_{\lambda,\mathrm{rev}}^T/
b_{\lambda,\mathrm{rev}}^T\) with the zero-denominator conventions of
\Cref{def-normalized-resolvent-defect}. If
\(\vartheta_\lambda^T<1\), then

\[
|\langle u_T,R_{\mathrm{full}}^qv_T\rangle|
\ge
(1-\vartheta_\lambda^T)b_{\lambda,\mathrm{rev}}^T.
\]

Moreover, the Schur-complement theorem gives the computable estimate

\[
d_{\lambda,\mathrm{rev}}^T
\le
\frac{\theta_\lambda(q)}{1-\theta_\lambda(q)}
\langle u_T,R_qu_T\rangle^{1/2}
\langle v_T,R_qv_T\rangle^{1/2}.
\]

\end{corollary}

\begin{proof}

The first assertion is the reverse triangle inequality. The second is
the final bilinear estimate in
\Cref{thm-schur-complement-compensation}.

\end{proof}

\begin{corollary}[Two-scale killed-gap bound]
\label{cor-two-scale-killed-gap}

Suppose, in addition, that

\[
A\succeq\gamma_Q I,
\qquad
D\succeq\gamma_\perp I,
\qquad
\|B\|\le\varepsilon,
\]

where \(\gamma_Q,\gamma_\perp>0\). Then

\[
\theta_\lambda(q)
\le
\frac{\varepsilon^2}
{(\gamma_Q+\lambda)(\gamma_\perp+\lambda)},
\]

and

\[
K\succeq\gamma_-I,
\qquad
\gamma_-
=
\frac{\gamma_Q+\gamma_\perp
-\sqrt{(\gamma_Q-\gamma_\perp)^2+4\varepsilon^2}}{2}.
\]

In particular, \(\gamma_->0\) precisely when
\(\varepsilon^2<\gamma_Q\gamma_\perp\). Thus a non-lumpable quotient
is controlled only when the coupling is small relative to both the
quotient gap and the within-fiber geometric gap.

\end{corollary}

\begin{proof}

For \(f=Uu+v\), with \(v\perp\operatorname{ran}U\),

\[
\langle f,Kf\rangle
\ge
\gamma_Q\|u\|^2+
\gamma_\perp\|v\|^2-
2\varepsilon\|u\|\|v\|.
\]

The least eigenvalue of the resulting two-dimensional quadratic form is
\(\gamma_-\). Its determinant is
\(\gamma_Q\gamma_\perp-\varepsilon^2\), which gives the positivity
criterion. The bound on \(\theta_\lambda\) follows from
\(\|(D+\lambda I)^{-1/2}\|\le
(\gamma_\perp+\lambda)^{-1/2}\),
\(\|(A+\lambda I)^{-1/2}\|\le
(\gamma_Q+\lambda)^{-1/2}\), and \(\|B\|\le\varepsilon\).

\end{proof}

\begin{corollary}[Represented local geometry compensates the defect]
\label{cor-represented-geometry-compensates-defect}

Let \(A_{\Phi_Q}\) and \(D_{\Phi_\perp}\) be the operators of
represented quotient and within-fiber Dirichlet forms. Suppose that,
on the certified windows,

\[
A+\lambda I
\succeq
(1-\delta_Q)(A_{\Phi_Q}+\lambda I),
\qquad
D+\lambda I
\succeq
(1-\delta_\perp)(D_{\Phi_\perp}+\lambda I),
\]

where \(0\le\delta_Q,\delta_\perp<1\), and that

\[
A_{\Phi_Q}\succeq g_QI,
\qquad
D_{\Phi_\perp}\succeq g_\perp I,
\qquad
\|B\|\le\varepsilon.
\]

Then

\begin{equation}
\label{eq-represented-geometry-theta-bound}
\theta_\lambda(q)
\le
\frac{\varepsilon^2}
{(1-\delta_Q)(1-\delta_\perp)
 (g_Q+\lambda)(g_\perp+\lambda)}.
\end{equation}

If the right-hand side is smaller than one, the represented geometries
yield a valid Schur-complement compensation certificate. Its transfer
loss is at most the reciprocal of one minus the right-hand side.

\end{corollary}

\begin{proof}

The form comparisons imply

\[
\|(A+\lambda I)^{-1/2}\|
\le[(1-\delta_Q)(g_Q+\lambda)]^{-1/2}
\]

and the analogous estimate for \(D+\lambda I\). Insert these two
bounds and \(\|B\|\le\varepsilon\) into
\eqref{eq-reversible-compensated-defect}. The Schur-complement loss then
follows from \Cref{thm-schur-complement-compensation}.

\end{proof}

\subsection{Coercive regularity and ambient consequences}
\label{subsec:coercive-regularity-ambient-consequences}

\begin{theorem}[Coercive regularity forces reversible quotient compensation]
\label{thm-coercive-regularity-forces-compensated-quotient}

Let \(L\) be a represented reversible killed generator on the nonterminal
space \(H_N\), put \(K=-L_N\), and suppose that the same represented
comparison record certifies

\begin{equation}
\label{eq-coercive-regularity-window}
0<\gamma I\preceq K\preceq\Lambda I
\qquad(0<\gamma\le\Lambda).
\end{equation}

Let \(q\ne\mathrm{id}\) be a represented terminal-compatible quotient for
this generator, with isometric lift \(U\), orthogonal projection
\(\Pi=UU^*\), and represented block data \(A,B,D\) from
\cref{def-reversible-compensated-defect}.  For every represented discount
\(\lambda>0\), set

\[
r_{\gamma,\Lambda,\lambda}
=
\frac{\gamma+\lambda}{\Lambda+\lambda}.
\]

Then its existing reversible compensated-defect parameter satisfies

\begin{align}
\label{eq-coercive-regularity-theta-bound}
\sqrt{\theta_\lambda(q)}
&\le 1-r_{\gamma,\Lambda,\lambda},\\
\label{eq-coercive-regularity-stability-bound}
S_q^{\mathrm{rev}}(\lambda)
=1-\theta_\lambda(q)
&\ge
2r_{\gamma,\Lambda,\lambda}
-r_{\gamma,\Lambda,\lambda}^{\,2}
\ge r_{\gamma,\Lambda,\lambda}.
\end{align}

In particular, \(E_q=0\) gives an exact quotient record.  If \(E_q\ne0\),
then \(q\) has a valid reversible compensated record and

\begin{equation}
\label{eq-coercive-regularity-transfer-loss}
\bigl(S_q^{\mathrm{rev}}(\lambda)\bigr)^{-1}
\le
\frac{\Lambda+\lambda}{\gamma+\lambda}.
\end{equation}

Consequently, whenever the quotient, its compressed and within-fiber block
data, the comparison in \cref{eq-coercive-regularity-window}, and the
quotient target-functional record have total saturated cost at most \(B\),
and the right side of
\cref{eq-coercive-regularity-transfer-loss} belongs to the declared transfer
class, every dynamically qualified target-visible use of this coarsening is
classified by the existing exact or reversible compensated nonidentity
quotient coordinate.  It cannot remain in the identity-supported ambient
coordinate.

\end{theorem}

\begin{proof}

Relative to
\(H_N=\operatorname{ran}U\oplus\operatorname{ran}U^\perp\), write

\[
K+\lambda I
=
\begin{pmatrix}
X&B^*\\
B&Y
\end{pmatrix},
\qquad
X=A+\lambda I,
\qquad
Y=D+\lambda I.
\]

Both \(X\) and \(Y\) are positive definite.  Define

\[
Z=Y^{-1/2}BX^{-1/2},
\qquad
\mathcal D=
\begin{pmatrix}
X&0\\0&Y
\end{pmatrix}.
\]

Block normalization gives the exact identity

\begin{equation}
\label{eq-coercive-normalized-block-identity}
\mathcal D^{-1/2}(K+\lambda I)\mathcal D^{-1/2}
=
\begin{pmatrix}
I&Z^*\\Z&I
\end{pmatrix}.
\end{equation}

The least eigenvalue of the self-adjoint operator on the right is
\(1-\|Z\|\).  Indeed, the singular-value decomposition of \(Z\) reduces
each nontrivial invariant block to
\(\left(\begin{smallmatrix}1&\sigma\\\sigma&1\end{smallmatrix}\right)\),
whose eigenvalues are \(1\pm\sigma\); any unmatched directions have
eigenvalue one.  By
\cref{def-reversible-compensated-defect},
\(\|Z\|=\sqrt{\theta_\lambda(q)}\).

The coercive comparison gives
\(K+\lambda I\succeq(\gamma+\lambda)I\).  The two diagonal blocks are
compressions of \(K+\lambda I\), so

\[
\mathcal D\preceq(\Lambda+\lambda)I.
\]

For every vector \(w\), therefore,

\begin{align*}
\left\langle
w,\mathcal D^{-1/2}(K+\lambda I)\mathcal D^{-1/2}w
\right\rangle
&\ge
(\gamma+\lambda)\|\mathcal D^{-1/2}w\|^2\\
&\ge
\frac{\gamma+\lambda}{\Lambda+\lambda}\|w\|^2.
\end{align*}

Taking the least eigenvalue in
\cref{eq-coercive-normalized-block-identity} yields

\[
1-\sqrt{\theta_\lambda(q)}
\ge r_{\gamma,\Lambda,\lambda},
\]

which proves \cref{eq-coercive-regularity-theta-bound}.  Since
\(0<r_{\gamma,\Lambda,\lambda}\le1\),

\[
1-\theta_\lambda(q)
\ge
1-(1-r_{\gamma,\Lambda,\lambda})^2
=2r_{\gamma,\Lambda,\lambda}
-r_{\gamma,\Lambda,\lambda}^{\,2}
\ge r_{\gamma,\Lambda,\lambda}.
\]

This is
\cref{eq-coercive-regularity-stability-bound,eq-coercive-regularity-transfer-loss}.
The identity \(E_q=-B\) from
\cref{def-reversible-compensated-defect} separates the exact case from the
nonzero-defect case.  In the latter case the strict inequality
\(\theta_\lambda(q)<1\), the transfer bound, and the represented
target-functional comparison give precisely the reversible dynamic-stability
record of
\cref{def-master-compensated-quotient-geometry-certificate}.  The master
normal form in
\cref{thm-affordable-quotient-geometry-master-normal-form} then assigns the
record to the exact or reversible compensated nonidentity quotient family at
its complete saturated cost.

\end{proof}

\begin{corollary}[Strict meaning of ambient geometry after quotient
exhaustion]
\label{cor-strict-ambient-geometry-after-quotient-exhaustion}

Fix a saturated budget \(B\), and suppose that the exact and reversible
compensated nonidentity quotient envelopes are bounded by \(\alpha_n\) as in
\cref{eq-prospective-quotient-branch-envelope}.  Let \(q\ne\mathrm{id}\)
be a represented terminal-compatible coarsening of a reversible geometric
record satisfying \cref{eq-coercive-regularity-window}, with
\((\Lambda+\lambda)/(\gamma+\lambda)\) in the transfer class.  Assume that
the quotient and every compressed, within-fiber, utility, target-functional,
and active-geometry datum in
\cref{def-master-compensated-quotient-geometry-certificate} have total
saturated cost at most \(B\).  Then

\begin{equation}
\label{eq-strict-ambient-coarsening-visibility-bound}
\overline A_q\,
M_{\mathfrak G_q}\,
C_{\mathfrak G_q}\,
\frac{R_{\mathfrak G_q}}{1+\rho_{\mathfrak G_q}}
\le
\alpha_n\,
\frac{\Lambda+\lambda}{\gamma+\lambda}.
\end{equation}

Put
\begin{equation}
\label{eq-coercive-quotient-obstruction-scale}
\beta_n(q)
=
\min\!\left\{
1,
\alpha_n\frac{\Lambda+\lambda}{\gamma+\lambda}
\right\}.
\end{equation}
For arbitrary thresholds \(a,c,r,g\in(0,1]\) satisfying
\begin{equation}
\label{eq-coercive-asymmetric-gate-threshold}
acrg>\beta_n(q),
\end{equation}
the same represented coarsening cannot satisfy all five conditions
\begin{equation}
\label{eq-coercive-asymmetric-gate-exclusion}
M_{\mathfrak G_q}=1,
\qquad
\overline A_q>a,
\qquad
C_{\mathfrak G_q}>c,
\qquad
R_{\mathfrak G_q}>r,
\qquad
\frac1{1+\rho_{\mathfrak G_q}}>g.
\end{equation}
In particular, every such coarsening satisfies the symmetric obstruction
\begin{equation}
\label{eq-coercive-symmetric-gate-obstruction}
M_{\mathfrak G_q}=0
\quad\text{or}\quad
\min\!\left\{
\overline A_q,
C_{\mathfrak G_q},
R_{\mathfrak G_q},
\frac1{1+\rho_{\mathfrak G_q}}
\right\}
\le \beta_n(q)^{1/4}.
\end{equation}

Thus, after quotient exhaustion, an ambient geometry can remain above scale
\(\alpha_n(\Lambda+\lambda)/(\gamma+\lambda)\) only when every affordable
nonidentity spectral or level-set coarsening loses that scale in a represented
quotient-utility, target-reach, descended-geometry, or target-functional
gate.  Coercive regularity itself cannot be the failed compensation gate.
This is the quantitative meaning of quotient-incompressible ambient geometry.

\end{corollary}

\begin{proof}

If \(E_q=0\), the master stability factor is one.  If \(E_q\ne0\),
\cref{thm-coercive-regularity-forces-compensated-quotient} gives
\[
S_q^{\mathrm{rev}}(\lambda)
\ge
\frac{\gamma+\lambda}{\Lambda+\lambda}.
\]
In either case the master normal form places the coarsening in the exact or
reversible compensated nonidentity quotient envelope.  By
\cref{eq-supported-geom-visibility}, its visibility is
\[
V_{\mathscr C}
=
\overline A_q S_q(\lambda)
M_{\mathfrak G_q}C_{\mathfrak G_q}
\frac{R_{\mathfrak G_q}}{1+\rho_{\mathfrak G_q}}
\le\alpha_n.
\]
Dividing by the displayed stability lower bound proves
\cref{eq-strict-ambient-coarsening-visibility-bound}.  Hence any coarsening
that retains a larger unconditioned geometric product would contradict the
nonidentity quotient-envelope bound.

Every factor in
\cref{eq-strict-ambient-coarsening-visibility-bound} lies in \([0,1]\), and
the mixing factor is binary.  If
\cref{eq-coercive-asymmetric-gate-exclusion} held under
\cref{eq-coercive-asymmetric-gate-threshold}, their product would exceed
\(\beta_n(q)\), contradicting
\cref{eq-strict-ambient-coarsening-visibility-bound,eq-coercive-quotient-obstruction-scale}.
Taking \(a=c=r=g=\beta_n(q)^{1/4}\) gives
\cref{eq-coercive-symmetric-gate-obstruction}.

\end{proof}

\Cref{fig-quotient-defect-compensation-trichotomy} summarizes the three
admissible routes from an ambient generator to a dynamically qualified
quotient certificate.

\begin{figure}[tbp]
\centering
\resizebox{.94\linewidth}{!}{%
\begin{tikzpicture}[fw diagram,node distance=8mm and 13mm]
  \node[fw source] (ambient) {ambient generator \(L\)};
  \node[fw provenance,right=of ambient,text width=2.8cm] (inter)
    {represented \(U,P\)\\and quotient \(L_Q\)};
  \node[fw box,right=of inter] (defect) {defect \(E_q=LU-UL_Q\)};
  \node[fw use,above right=9mm and 15mm of defect] (exact)
    {exact\\\(E_q=0\)};
  \node[fw use,right=15mm of defect] (rev)
    {reversible\\Schur complement};
  \node[fw use,below right=9mm and 15mm of defect] (gen)
    {general target\\resolvent defect};
  \node[fw terminal,right=20mm of rev,text width=2.8cm] (stab)
    {represented stability\\\(S_q(\lambda)\)};
  \draw[fw arrow] (ambient) -- (inter);
  \draw[fw arrow] (inter) -- (defect);
  \draw[fw arrow] (defect) -- (exact);
  \draw[fw arrow] (defect) -- (rev);
  \draw[fw arrow] (defect) -- (gen);
  \draw[fw arrow] (exact) -- (stab);
  \draw[fw arrow] (rev) -- (stab);
  \draw[fw arrow] (gen) -- (stab);
\end{tikzpicture}%
}
\caption{A quotient is dynamically qualified through a charged
route: exact intertwining, reversible Schur-complement compensation, or a
general target-functional resolvent estimate.  All routes terminate in the
same stability factor and use the same represented ambient generator.}
\label{fig-quotient-defect-compensation-trichotomy}
\end{figure}

\subsection{Causal alignment under the authority envelope}
\label{subsec:causal-authority-alignment}

The directed contribution is quantified by the alignment of the selected
authorized current with the represented potential.

\begin{definition}[Quantitative Caus alignment]
\label{def-quantitative-caus-alignment}

Let \(\mathcal H_E\) be the represented weighted edge space used by the
Hodge decomposition, let \(J\in\mathcal H_E\) be a represented directed
current, and let \(V\) be a represented scalar potential. Define

\[
\operatorname{Align}_{\Caus}(J,V)
=
\begin{cases}
\displaystyle
\frac{\bigl(\langle J,-\nabla V\rangle_E\bigr)_+^2}
{\|J\|_E^2\|\nabla V\|_E^2},
&\|J\|_E\|\nabla V\|_E>0,\\[3mm]
0,&\text{otherwise},
\end{cases}
\]

where \(a_+=\max\{a,0\}\). This number lies in \([0,1]\). It is one
for the canonical descending current \(J=-\nabla V\), and it vanishes
when the represented current has no descending component along \(V\).
The construction cost of \(V\), the edge structure, the weights, and
the current is included in the saturated certificate.

\end{definition}

This alignment factor retains only causal descent supported by the declared authority record.
\begin{definition}[Authority-compatible Caus alignment]
\label{def-authority-compatible-caus-alignment}

Let \(\mathfrak K\) be a causal provenance record for a
target-normalized potential \(V\).  If \(\mathfrak K\) verifies that its
current \(J_{\mathfrak K}\) is authorized by
\Cref{def-authorized-caus-use-current} and is transported
through every quotient, coordinate change, and lift in the certificate,
set

\[
\operatorname{Align}^{\mathrm{auth}}_{\Caus}
(\mathfrak K,V)
=
\operatorname{Align}_{\Caus}(J_{\mathfrak K},V).
\]

Set the value to zero when the current identity, target polarity, or
dynamic transport record is absent.  The authority-compatible local
factor is

\[
C_{\Phi,\mathfrak K}^{\mathrm{auth}}(V)
=
\chi_{\Phi}^{\mathrm{auth}}(V)
\operatorname{Align}^{\mathrm{auth}}_{\Caus}
(\mathfrak K,V),
\]

where \(\chi_{\Phi}^{\mathrm{auth}}(V)=1\) only when the represented
local-consistency, drift, or carrier comparison is evaluated for the
same full generator that authorizes \(\mathfrak K\). More precisely, the
record must contain and verify a displayed quantitative implication for
that generator---for example a drift margin
\(-LV\ge a>0\) on the nonterminal sector, a resolvent/form comparison
with declared constants, or the authorized carrier inequality
\eqref{eq-authorized-carrier-interface}. The margin and every
comparison loss must lie in the declared transfer class. Mere
construction of \(V\), or the formal identity
\(J=-\nabla_\Phi V\), does not set the indicator to one. This factor
lies in \([0,1]\).

\end{definition}

\begin{lemma}[Caus drift survives a controlled quotient defect]
\label{thm-caus-drift-survives-defect}

Let \(L\) be the generator of a finite Markov process with absorbing
target \(T\), and let \(q:X\to Q\) be a represented
terminal-compatible quotient.  Assume that \(L\) is selected from a
declared Caus authority envelope, that \(L_Q=P_qLU_q\), and that the
current paired with \(V\) has a causal provenance record
\(\mathfrak K\) transported from the same selected generator. Let
\(V:Q\to[0,\infty)\) vanish on the quotient target. Suppose that, on
every nonterminal quotient state,

\[
-L_QV\ge a>0,
\]

and that the pointwise defect action satisfies

\[
\|E_qV\|_{\infty,N}\le\delta<a.
\]

Then the lifted potential satisfies

\begin{equation}
\label{eq-caus-drift-defect-margin}
-L(UV)\ge a-\delta
\qquad\text{on }N=X\setminus T.
\end{equation}

If \(\tau_T\) is the first hitting time of \(T\), then

\[
\mathbb E_x\tau_T
\le
\frac{(UV)(x)}{a-\delta},
\qquad
\mathbb E_{\mu_0}\tau_T
\le
\frac{\mathbb E_{\mu_0}[UV]}{a-\delta}.
\]

The potential gives a \(\Caus\) certificate when
\(C_{\Phi,\mathfrak K}^{\mathrm{auth}}(V)>0\). Exact lumpability is the
case \(\delta=0\); geometric compensation is the case in which
represented local regularity makes the positive margin \(a-\delta\)
affordable and verifiable.  A current with no authority-compatible
transport record cannot be substituted into this conclusion.

\end{lemma}

\begin{proof}

The defect identity gives

\[
L(UV)=U(L_QV)+E_qV.
\]

The two hypotheses therefore imply
\eqref{eq-caus-drift-defect-margin}. Dynkin's formula for the stopped
process gives

\[
\mathbb E_x[(UV)(X_{t\wedge\tau_T})]
=(UV)(x)
+\mathbb E_x\int_0^{t\wedge\tau_T}L(UV)(X_s)\,ds.
\]

The left-hand side is nonnegative, so
\((a-\delta)\mathbb E_x[t\wedge\tau_T]\le(UV)(x)\).
Monotone convergence as \(t\to\infty\) proves the first hitting-time
bound; integration against \(\mu_0\) proves the second
\citep{Norris1997}.

\end{proof}

\section{Sectorial and Continuum Interfaces}
\label{sec:sectorial-continuum-interfaces}

The sectorial interface transports the same causal-authority record to the
conditional continuum realization.

\begin{definition}[Sectorial Caus-authority interface]
\label{def-sectorial-caus-authority-interface}

Let \(H=L^2(X,\mu)\), let \(\mathcal D\subseteq H\) be a dense common
form domain, and let

\[
\mathcal E
=
\mathcal E^{\mathsf s}
+\mathcal E^{\mathsf a,\mathrm{imp}}
+\mathcal E^{\mathsf a,\mathrm{ctrl}}
+\mathcal E^{\mathsf b}
\]

be a represented closed sectorial Markov form on \(\mathcal D\).
Here \(\mathcal E^{\mathsf s}\) is symmetric and nonnegative,
\(\mathcal E^{\mathsf a,\mathrm{imp}}\) is the imposed antisymmetric
form,
\(\mathcal E^{\mathsf a,\mathrm{ctrl}}
\in\mathfrak E^{\mathsf a,\mathrm{ctrl}}_{I,\le B}\) belongs to a
declared affordable control family, and \(\mathcal E^{\mathsf b}\) is
the symmetric nonnegative boundary or killing form.  The selected
antisymmetric form must satisfy a represented sector estimate

\[
\left|
\mathcal E^{\mathsf a,\mathrm{imp}}(u,v)
+\mathcal E^{\mathsf a,\mathrm{ctrl}}(u,v)
\right|
\le
K_{\mathrm{sec}}\,
\mathcal E^{\mathsf s}_1(u,u)^{1/2}
\mathcal E^{\mathsf s}_1(v,v)^{1/2},
\qquad
\mathcal E^{\mathsf s}_1(u,u)
=\mathcal E^{\mathsf s}(u,u)+\|u\|_2^2.
\]

The associated generator uses the sign convention
\(\mathcal E(u,v)=-\langle Lu,v\rangle\) on its operator domain.

A form-compatible quotient has an isometric pullback
\(U_q:H_Q\to H\) with
\(U_q\mathcal D_Q\subseteq\mathcal D\), adjoint projection
\(P_q=U_q^*\), and compressed forms

\[
\mathcal E_Q^\sharp(u,v)
=\mathcal E^\sharp(U_qu,U_qv),
\qquad
\sharp\in\{\mathsf s,\mathsf a,\mathsf b\}.
\]

For \(u\in\mathcal D_Q\) and \(v\in\mathcal D\) with
\(P_qv\in\mathcal D_Q\), define the form defects

\[
\langle\mathfrak E_q^\sharp u,v\rangle
=
\mathcal E^\sharp(U_qu,v)
-\mathcal E_Q^\sharp(u,P_qv),
\qquad
\sharp\in\{\mathsf s,\mathsf a\}.
\]

They are elements of the dual form space \(\mathcal D^*\).  The
antisymmetric defect is the continuum causal transport defect.

A represented carrier interface consists of functionals
\(\mathcal K,\mathfrak C,\Phi,R\) on a common admissible profile domain
such that, along every declared admissible profile \(t\mapsto V(t)\),

\begin{equation}
\label{eq-authorized-carrier-interface}
\frac{d}{dt}\mathcal K(V(t))
\le
-\mathfrak C(V(t))+\Phi(V(t))+R(V(t)),
\qquad
\mathfrak C(V(t))\ge0.
\end{equation}

A quotient, rescaling, compactification, or coordinate change is legal
for this interface when it transports the imposed and controlled forms
by the displayed compression rule, preserves the target and declared
capacity, and descends
\eqref{eq-authorized-carrier-interface}.  Every discrepancy is recorded
in \(\mathfrak E_q^{\mathsf s}\), \(\mathfrak E_q^{\mathsf a}\), or
the represented remainder \(R\), with its derivation and comparison
cost charged.

\end{definition}

\begin{lemma}[Conditional continuum realization of Caus authority]
\label{thm-continuum-caus-authority-bridge}

Assume \Cref{def-sectorial-caus-authority-interface}.  The
closed form determines a sectorial generator whose antisymmetric
component has the Caus authority envelope

\[
\mathcal E^{\mathsf a,\mathrm{imp}}
+\mathfrak E^{\mathsf a,\mathrm{ctrl}}_{I,\le B}.
\]

Every legal form-compatible quotient or coordinate transformation
therefore produces the authority-compatible operator record required by
the master quotient--geometry construction: its symmetric geometry is the compression of
\(\mathcal E^{\mathsf s}\), its imposed Caus component is the
compression of \(\mathcal E^{\mathsf a,\mathrm{imp}}\), and its
component defects are
\(\mathfrak E_q^{\mathsf s}\) and
\(\mathfrak E_q^{\mathsf a}\).  A different transformed antisymmetric
form is admissible only as a declared control or as a new dynamic
presentation.

In a finite, Galerkin, finite-element, or other represented finite-rank
realization, this form record supplies the dynamic and authority slots
of a master certificate.  It becomes a master certificate of
\Cref{def-master-compensated-quotient-geometry-certificate}
only when the target-normalized potential, reach, utility, and required
comparison estimates are also represented.  In the continuum it
supplies the corresponding closed-form slots under the same condition.
Effective-resistance, sector, commutator, hypocoercive, positive-
commutator, or flux/dissipation estimates may supply its comparison
record when their domains, closures, constants, and transport maps are
represented.  The theorem transfers provenance and compatibility; the
compactness, capacity, rigidity, and realization hypotheses of a
specific continuum problem remain hypotheses of its application
theorem.

\end{lemma}

\begin{proof}

The representation theorem for closed sectorial forms gives the
generator on the common form domain \citep{Kato1995}.  The symmetric and antisymmetric
parts are fixed by polarization, and the sector bound makes the selected
antisymmetric form a form-bounded perturbation of
\(\mathcal E^{\mathsf s}_1\).  Compression by \(U_q\) gives the three
quotient forms in the definition.  Subtracting their pullbacks from the
ambient forms gives the displayed dual-form defects, so the imposed
antisymmetric part cannot change without a declared control or a change
of presentation.

In a finite-rank realization, the forms are represented matrices and
their component defects agree with
\Cref{def-symmetric-causal-quotient-defects}.  When the
remaining current and target-potential data are represented, their
causal gate is precisely
\Cref{def-authority-compatible-caus-alignment}. In the
continuum, the corresponding operator identities hold in the
form/dual-form pairing. Any additional
routing or carrier theorem is usable precisely when its represented
comparison data control these defects and descend through the same
interface.  This proves the conditional bridge.

\end{proof}

\begin{corollary}[Authorized flux--dissipation polarity]
\label{cor-authorized-flux-dissipation-polarity}

Suppose \(t\mapsto\mathcal K(V(t))\) is absolutely continuous, the
terms below are integrable on finite time intervals, and a legal carrier
interface satisfies, for some \(\beta\ge0\),

\[
\Phi_+(V)\le\beta\,\mathfrak C(V)
\]

on its reduced represented profile domain, with no positive-flux mode
in \(\ker\mathfrak C\).  Then

\[
\mathcal K(V(T))
+(1-\beta)\int_0^T\mathfrak C(V(t))\,dt
\le
\mathcal K(V(0))+\int_0^TR(V(t))\,dt.
\]

For \(\beta<1\), the imposed dynamics are draining with a quantitative
margin.  For \(\beta=1\), the comparison leaves the represented
critical carrier sector to an equality or rigidity analysis.  For
\(\beta>1\), this carrier supplies no draining conclusion.  These
alternatives are invariant under every legal authority-preserving
transformation.

\end{corollary}

\begin{proof}

Substitute
\(\Phi\le\Phi_+\le\beta\mathfrak C\) into
\eqref{eq-authorized-carrier-interface} and integrate.  Legality of the
transformation makes the descended carrier identity identical up to the
recorded remainder, so the same estimate holds after quotienting or
changing coordinates.

\end{proof}

\begin{remark}[Imposed dissipation and transport]
\label{rem-imposed-dissipation-and-transport}

For a dissipative evolution, viscosity or friction belongs to the
declared symmetric form.  Conservative transport belongs to the imposed
antisymmetric form when it is skew under the declared pairing and
boundary conditions; any remaining symmetric or residual part is
retained in its corresponding ledger. Constraints enter through their
represented boundary, lift, or remainder interfaces.  A legal
transformation may reveal a target-aligned quotient or potential, but
it preserves the sign and provenance of the carrier identity.  In
particular, a dissipative term cannot be replaced by a favorable
directed current.

\end{remark}

\subsection{Local geometric form comparison}
\label{subsec:local-geometric-form-comparison}

We next bound the resolvent effect of a represented perturbation of the local
Dirichlet geometry.

\begin{lemma}[Local Dirichlet-form error for the geometric component]\label{lem-declared-geom-form-error}

The comparison uses only represented objects already defined: the
\PartII \(\Geom\) component of a budget-admissible generator and
a declared geometry \(\Phi\) already eligible for
\(G_n^{\mathrm{sat}}(B)\) at its saturated public cost. The scalar below
is charged auxiliary comparison data. It does not define a structural
component, invariant, constructor, or admissible object, but it may be
used at budget \(B\) only when the common form domain, the compression
to \(W\), and a derivation or proof of the displayed bound are
represented within the declared comparison budget.

Let \(N=X\setminus T\) be the nonterminal sector and let
\(H_N=L^2(N,\mu|_N)\). For a budget-admissible killed generator \(L\),
let \(L^{\Geom}\) be its \PartII geometric component restricted
to \(N\). This is the nonnegative paired local-rate component of the
fixed-space split; it is not the symmetric part of an arbitrary
nonreversible generator. Write its Dirichlet form as

\[
\mathcal E_L^{\Geom}(f,g)
=
\frac12\sum_{x,z\in N}c_L^{\Geom}(x,z)(f(x)-f(z))(g(x)-g(z))\mu(x)\mu(z).
\]

Let \(\Phi=(E_\Phi,c_\Phi)\) be a declared geometry with Dirichlet form
\(\mathcal E_\Phi\). Let \(W\le H_N\) be a finite target-oriented
observable subspace, typically \(W=\mathcal M(D)|_N\). For
\(\lambda>0\), set the relative form error

\[
\delta_{\lambda,W}(L,\Phi)
=
\sup_{0\ne f\in W}
\frac{|\mathcal E_L^{\Geom}(f,f)-\mathcal E_\Phi(f,f)|}
{\mathcal E_\Phi(f,f)+\lambda\|f\|_{H_N}^2}.
\]

The denominator is strictly positive on \(W\setminus\{0\}\). This number
records the relative form error on the verified target-oriented window.
It contributes no separate term to \(I_n^{\mathrm{sat}}(B)\); the
geometric/reversible data remain the declared geometries, potentials, defects,
and value readouts in \(G_n^{\mathrm{sat}}(B)\) at their least saturated
public cost.

\end{lemma}

\begin{proof}\phantomsection\label{proof-declared-geom-form-error}

The \PartII fixed-space split represents the geometric component by a
symmetric nonnegative conductance \(c_L^{\Geom}\) on the exposed
local relation, so the displayed Dirichlet form is a finite symmetric
positive semidefinite form on \(H_N\). The declared geometry \(\Phi\)
gives another such form on the same finite space. For \(\lambda>0\), the
denominator \(\mathcal E_\Phi(f,f)+\lambda\|f\|^2\) is positive for
every nonzero \(f\in W\); hence the supremum defining
\(\delta_{\lambda,W}\) is finite. The lemma only introduces this scalar
comparison between two already represented forms. Since no new
observable, edge relation, quotient, or readout is introduced, the
scalar is auxiliary comparison data and makes no separate contribution to
the saturated invariant. Its evaluation or proof cost nevertheless
belongs to the comparison-data charge specified in the statement.

\end{proof}

\begin{proposition}[Geometric Dirichlet-form comparison of resolvents]\label{prop-geom-form-resolvent-comparison}

Let \(A_L^W\) and \(A_\Phi^W\) be the self-adjoint positive operators on
\(W\) representing the compressed forms
\(\mathcal E_L^{\Geom}|_W\) and \(\mathcal E_\Phi|_W\). If
\(\delta=\delta_{\lambda,W}(L,\Phi)<1\), then, as quadratic forms on
\(W\),

\[
(1-\delta)(A_\Phi^W+\lambda I)
\preceq
A_L^W+\lambda I
\preceq
(1+\delta)(A_\Phi^W+\lambda I).
\]

Consequently,

\[
(A_L^W+\lambda I)^{-1}
\preceq
\frac{1}{1-\delta}(A_\Phi^W+\lambda I)^{-1}.
\]

Thus a comparison proof for the already-declared geometry \(\Phi\)
transfers to the actual geometric component of \(L\) with loss at most
\((1-\delta)^{-1}\) on the verified window. The transfer is usable at
the declared budget when this loss lies in the transfer class. For
\(\mathcal B_{\mathrm{poly}}\), this is the condition
\(1-\delta\ge1/\operatorname{poly}(n)\). When the declared transfer
condition fails, the altered conductance cannot borrow the comparison
proof for \(\Phi\); it must be represented as its own Geom datum and
charged to \(G_n^{\mathrm{sat}}(B)\) at its least saturated public cost,
or it remains outside the exposed within-budget structure.

\end{proposition}

\begin{proof}\phantomsection\label{proof-geom-form-resolvent-comparison}

For every \(f\in W\), the definition of \(\delta\) gives

\[
|\langle f,(A_L^W-A_\Phi^W)f\rangle|
\le
\delta\langle f,(A_\Phi^W+\lambda I)f\rangle.
\]

Adding \(\langle f,(A_\Phi^W+\lambda I)f\rangle\) to the middle term
gives the two form inequalities. Since \(W\) is finite dimensional and
\(\lambda>0\), both shifted operators are strictly positive. Inverting
positive operators reverses Loewner order, which gives the displayed
resolvent comparison.

\end{proof}

\Cref{fig-geom-form-comparison-window} summarizes the local form comparison
used in \cref{prop-geom-form-resolvent-comparison}.

\begin{figure}[tbp]
\centering
\begin{tikzpicture}[fw diagram]
  \node[fw source,minimum width=3.0cm] (decl) at (-4.4,1.0)
    {declared form\\\(\mathcal E_\Phi|_W\)};
  \node[fw geom,minimum width=3.0cm] (actual) at (-4.4,-1.0)
    {actual Geom form\\\(\mathcal E_L^{\Geom}|_W\)};
  \node[fw provenance,minimum width=3.0cm] (delta) at (0,0)
    {represented comparison\\\(\delta_{\lambda,W}<1\)};
  \node[fw use,minimum width=3.2cm] (order) at (4.2,0)
    {resolvent transfer\\loss \((1-\delta)^{-1}\)};
  \draw[fw arrow] (decl) -- (delta);
  \draw[fw arrow] (actual) -- (delta);
  \draw[fw arrow] (delta) -- (order);
\end{tikzpicture}
\caption{A declared geometry controls the actual geometric resolvent only on
the represented target-oriented window \(W\), with the explicit relative-form
loss from \cref{lem-declared-geom-form-error,prop-geom-form-resolvent-comparison}.
The comparison introduces no additional structural channel.}
\label{fig-geom-form-comparison-window}
\end{figure}

\begin{corollary}[Form comparison caps nonlinear geometric gain]\label{cor-geom-form-target-resolvent}

Let \(y_W=\Pi_Wy\). Under the hypotheses of
\cref{prop-geom-form-resolvent-comparison},

\[
\frac{\langle y_W,(A_L^W+\lambda I)^{-1}y_W\rangle}{\|y\|^2}
\le
\frac{1}{1-\delta_{\lambda,W}(L,\Phi)}
\frac{\langle y_W,(A_\Phi^W+\lambda I)^{-1}y_W\rangle}{\|y\|^2}.
\]

If, in addition, represented hitting-functional comparison data bound
the full killed target resolvent by these compressed quadratic forms,
then \Cref{the-resolvent-stopwatch} converts the resulting full
resolvent estimate into the corresponding arrival estimate, with the
same form-comparison loss and the already declared transfer-class
constants.

\end{corollary}

\begin{proof}\phantomsection\label{proof-geom-form-target-resolvent}

Apply the resolvent inequality from Proposition
\hyperref[prop-geom-form-resolvent-comparison]{Geometric Dirichlet-form
comparison of resolvents} to the vector \(y_W\) and divide by
\(\|y\|^2\). The compressed quadratic-form inequality alone does not
identify a hitting probability. Under the additional comparison
hypothesis in the statement, it bounds the relevant full killed
resolvent; \Cref{thm-resolvent-gap-arrival-proof-equivalence}
then gives the asserted arrival comparison.

\end{proof}

\subsection{Nonlinear closure and the quotient spectral split}
\label{subsec:nonlinear-closure-quotient-spectral-split}

The final phase classifies nonquotient nonlinear constructions and records the
spectral split induced by a represented quotient.

\begin{proposition}[Nonquotient nonlinear constructions are geometric or residual]\label{prop-nonquotient-nonlinear-geometry-closure}

Let a presentation-relative uniform oracle-free budget-admissible
construction introduce a nonlinear observable, feature map, score,
metric, or update rule inside the lifted configuration space. Any finite
observable postprocessing, coarsening, thresholding, or relabeling of
that object is itself another admissible feature map at the
corresponding saturated cost. If such a postprocessing satisfies the
represented-fiber, terminal-event, exact-compatibility, quotient-utility, and
nontrivial-public-structure clauses of
\cref{def-paper1-abs-usefulness-gate}, it is a qualified \(\Abs\) quotient and
is counted by \(m_n^{\mathrm{sat}}(B)\) at the least saturated budget
containing its complete uniform record. A prospective use is further
restricted by \cref{def-pre-hit-temporally-admissible-source}. After all
qualified quotient
postprocessings have been removed, every remaining target-relevant
exposed contribution of the construction is accounted for by one of the
following alternatives.

\begin{enumerate}

\item

It induces a pointwise geometry, defect observable, directed potential,
or scalar terminal/support coloring of least saturated representation
cost \(k\). Then its Geom/Caus contribution is one of the terms counted
by \(G(k)\). A Bdry coloring or readout remains in the complete carrier;
its conditional target charge is bounded by the represented source only
when a contextual transfer inequality is supplied.

\item

It claims to improve an existing declared geometry \(\Phi\) only through
its actual geometric conductance. Then the comparison with \(\Phi\) is
valid on a target-oriented window \(W\) only with the form-comparison
loss \((1-\delta_{\lambda,W})^{-1}\) from
\cref{cor-geom-form-target-resolvent}. If
this loss is not in the declared transfer class, no budgeted hitting
comparison transfers from \(\Phi\); the altered conductance must be
counted as its own geometry at every budget in its admissible budget set or treated by
the remaining alternatives.

\item

It has no exposed Geom/Caus/Abs/Lift/Bdry provenance in the observable
or lifted algebra. Then, after the \PartII projections, its contribution
is the residual term \(\Res\), presentation-inaccessible or
oracle-dependent information, nonuniform advice, or over-budget
representation work.

\end{enumerate}

Hence the decomposition contains no additional nonlinear component
between valid quotient structure and geometry. A nonlinear construction
either enforces exact lumpability and
is charged to \(m_n^{\mathrm{sat}}(B)\) at the budget representing its
quotient, or it supplies exposed geometric/causal data together with
possible scalar boundary colorings and is charged to
\(G_n^{\mathrm{sat}}(B)\) exactly when
\(B\in\mathfrak B_{\mathrm{sat},n}(R_{\Geom})\), or it is
residual, inadmissible, or outside the declared budget.

Here every boundary or lifted occurrence remains in the complete
carrier. Its separate contextual charge is dominated by the displayed
Geom/Caus source only when a represented transfer comparison is part of
the record.

\end{proposition}

\begin{proof}\phantomsection\label{proof-nonquotient-nonlinear-geometry-closure}

Embed the computation as a finite Markov generator on its bounded
clocked transcript carrier. \PartII exhaustiveness decomposes the
induced movement into Geom, Caus, Abs, valid Lift, Bdry, and
\(\Res\). The preliminary closure under finite observable
postprocessing removes every valid Abs quotient derivable from the
nonlinear object. Valid Lift creates no primitive target channel and is
retained with the structure it lifts; Bdry supplies scalar terminal/support
colorings or boundary values and is retained with the source that made
those values exposed. This is carrier-cost and provenance accounting;
their contextual target charges transfer only through a represented
comparison. Therefore any remaining exposed target-relevant
movement is Geom or Caus, represented by a conductance, pointwise
observable, directed potential, or edge difference generated from the
lifted observable algebra; Bdry can only color or read the states on
which that represented movement is evaluated. By definition of saturated
representation degree, that datum enters \(G_n^{\mathrm{sat}}(B)\) exactly
for budgets \(B\in\mathfrak B_{\mathrm{sat},n}(R_{\Geom})\) when its
stated admissibility conditions hold. If the construction is presented
only as a perturbation of an existing geometry, the preceding
form-comparison lemma is exactly the finite-dimensional condition under
which the existing geometric comparison proof transfers. If the
condition fails, the perturbation does not yield an uncharged geometric
improvement; it must be represented and counted as its own geometry or
else remains outside the exposed within-budget structure. The residual
and presentation-inaccessible cases are the remaining alternatives in \PartII's
decomposition.

\end{proof}

\begin{proposition}[Quotient projection gives the spectral split]\label{prop-quotient-projection-spectral-split}

A degree-\(d\) quotient represents the component \(P_dy\) of the
centered target contrast. In the Walsh basis, the orthogonal complement
has squared mass

\[
\|(I-P_d)y\|^2=
\sum_{|S|>d}|\widehat y(S)|^2.
\]

The represented quotient captures the normalized mass

\[
\frac{\|P_dy\|^2}{\|y\|^2}=m(d).
\]

\end{proposition}

\begin{proof}\phantomsection\label{proof-quotient-projection-spectral-split}

The represented degree-\(d\) quotient subspace is the range of the
orthogonal projection \(P_d\). Hence \(P_dy\) is the component of \(y\)
represented by quotient observables of degree at most \(d\), while
\((I-P_d)y\) is orthogonal to every such observable. Parseval's identity
in the Walsh basis gives the displayed tail formula. Dividing by
\(\|y\|^2\) gives the normalized represented mass.

\end{proof}

\begin{example}[Clause buckets and lumpability]\label{ex-lumpability-sat}

Grouping SAT assignments by equal violated-clause count gives
represented score buckets when the score is available. Exact lumpability
requires more: all assignments in one bucket must have the same
quotient-visible transition behavior and the same terminal status after
projection. Assignments with the same defect can have different clause
patterns and different target access under a budget-admissible operator.
The defect operator \(E_q\) measures this within-bucket mismatch.

\end{example}

\section{Combined Geometric and Quotient Extraction
Functional}\label{the-combined-geomabs-extraction-functional}
\label{sec:combined-geometric-quotient-extraction}

The spectral bound handles represented quotient extraction. Pointwise
geometric extraction has a separate regularity condition. The combined
functional records the target-relevant information available under a
saturated public cost budget. The primary variable is a budget \(B\);
the notation \(d\) is used only as a cost level, for example
\(d=\lceil\log_2 B\rceil\). Neither \(B\) nor \(d\) is raw Walsh degree
or mere measurability.

\subsection{Saturated representation cost and the three-layer ledger}
\label{subsec:saturated-representation-ledger}

We first place every represented structural object and its execution record on
one saturated resource ledger.

\begin{definition}[Saturated budget-admissible representation cost]\label{def-saturated-representation-degree}

Fix a budget structure
\(\mathcal B=(\mathcal R,\operatorname{cost},b,\mathfrak C)\) and a cost
level \(B\preceq b(n)\). The saturated budget-admissible representation
class \(\mathcal R_{n,\le B}^{\mathrm{sat}}\) consists of all
target-relevant observable representatives generated from the presented
observable algebra, declared public parameters, and legal lifted closure
by the declared oracle-free constructors with public cost at most
\(B\). At the declared ceiling these representatives constitute
\(\mathcal R^{\mathcal B}_n\); for the polynomial instance
\(\mathcal B_{\mathrm{poly}}\), the union over polynomial ceilings
recovers \(\mathcal R_n^{\mathcal B_{\mathrm{poly}}}\), not the entire
ambient measurable envelope. When the representative
includes a forward transition rule, its uniform code is allowed, but
every instance-dependent datum read by that rule must lie in the same
cost-filtered budgeted closure. These representatives include pointwise
observables, geometries, potentials, represented quotients, valid lifted
interfaces, boundary readouts, forward observable transport kernels, and
generator components after \PartII channel projection. The charged
public cost includes description length, effective dimension,
dense-basis scan cost, represented fiber operations, charged public
basis-change data, lifted memory/interface data, comparison data,
terminal readout data, and finite postprocessings used as observable
structure.

The constructors describe the resulting fibers, potentials,
conductances, forward transition rules, interfaces, terminal value
readouts, finite-horizon trajectory statistics, or component-projected
generator components as observable objects. Executing a named solver
provides no constructor for its exact future hitting behavior,
first-hitting partition, committor, occupation law, or resolvent unless
that semantic object is represented explicitly. A finite-horizon
computation generated
from budget-admissible data is instead charged at its saturated
computation cost. In the general resource order, the saturated cost of a
target-relevant component \(R\) is the admissible-budget set and its
Pareto frontier from \Cref{def-budgeted-derivability-closure}.
After a monotone scalar cost \(\kappa\) has been fixed, its scalar
saturated cost is

\[
\operatorname{Cost}_{\mathrm{sat},n}^{\kappa}(R)
=
\inf_{\rho:\,\sem{\rho}_I=R_I\ \forall I\in\mathfrak I_n}
\ \sup_{I\in\mathfrak I_n}
\kappa\!\left(\operatorname{Cost}_I(\rho)\right).
\]

We omit \(\kappa\) after fixing that scalarization. The optional log-cost degree is
\(\deg_{\mathrm{sat},n}(R)=\lceil\log_2(1+\operatorname{Cost}_{\mathrm{sat},n}(R))\rceil\).
Thus \(\mathcal R_{n,\le d}^{\mathrm{sat}}\) means
\(\mathcal R_{n,\le 2^d}^{\mathrm{sat}}\) when log-cost notation is
used. This is an objective infimum over the saturated class and is
attained under the locally finite discrete scalar-cost convention. It
does not depend on whether the representative has been discovered or
whether optimality has been proved.

The budget frontier, or the scalar cost after scalarization, is the
budget-admissible representation cost. A
component may have high raw Walsh degree and low saturated cost when a
valid quotient, basis change, lift, forward transport, or legal
finite-horizon computation generated from the budgeted saturated
observable closure represents it compactly. Broad measurability in the
finite state space is insufficient: a set-theoretic description or an
exact semantic future value may exist without any budgeted saturated
representative. Conversely, if a denoted object is produced by a
presentation-relative within-budget computation from
\(\mathcal R^{\mathcal B}_n\) and the legal lifted closure, that
computation is a budget-admissible representative at its saturated
computation cost by Lemma
\hyperref[lem-budgeted-generation-gives-representation]{Budgeted
generation gives budgeted representation}. Therefore every admissible
budgeted datum belongs to the budgeted saturated closure. Broad
measurability gives an ambient set-theoretic description, whereas
membership in this closure gives admissible budgeted structure.

\end{definition}

The three ledgers separate displayed, represented, and saturated resource accounting.
\begin{definition}[Three layers of cost accounting]\label{def-three-cost-ledgers}

Every object used in the hitting-time bounds is evaluated through three
accounting layers, all
evaluated relative to the same source of legal instance-dependent data:
\(\mathcal R^{\mathcal B}_n\).

\textbf{Provenance.} The provenance account records whether the object is
generated from \(\mathcal R^{\mathcal B}_n\), declared public
parameters, public lifted coordinates, and oracle-free constructor
rules. For transition rules, this account records whether all
instance-dependent data read by the uniform code have that provenance.
It determines whether the object belongs to the budgeted observable
algebra of the original presentation or instead requires
presentation-inaccessible information, nonuniform advice, oracle access,
over-budget recovery, or an enlarged primitive signature.

\textbf{Execution.} The execution account records the resources needed by
a named transition rule, simulation, evaluator, or update procedure
under the budget-compatible resource envelope, after the provenance
account has fixed the admissible input data. This account determines whether a
program or induced kernel is executable within the declared budget for
the presented task. Budget-compatible execution relative to an uncharged
modulus, trapdoor, target table, over-budget recovered relation, or
advice string is budget-compatible execution for a different input
object, not for the original presentation.

\textbf{Structural/spectral representation.} The structural account
records the least saturated public cost \(B\), or equivalently the least
log-cost level \(d\), at which the denoted object itself is represented
in \(\mathcal R_{n,\le B}^{\mathrm{sat}}\). This account controls
\(m^{\mathrm{sat}}(B)\), \(G^{\mathrm{sat}}(B)\),
\(I^{\mathrm{sat}}(B)\), and \(B^*\).

These accounts are distinct aspects of one budgeted closure rather than
independent admissibility criteria. The underlying data source is
\(\mathcal R^{\mathcal B}_n\) and its legal lifted closure.
Budget-compatible execution does not assign zero saturated
structural/spectral cost to every semantic object attached to the
execution. A finite-horizon object generated by an admissible within-budget
simulation is legal budgeted data and is charged at the saturated cost
of that computation; an uncomputed semantic future object is not. The
execution account is evaluated only after inadmissible information has been
removed or charged: a fast update rule using data outside the budgeted
saturated observable closure is not admissible for the original task,
even if it becomes efficient after enlarging the primitive signature. Finite
charged composition of budget-admissible updates preserves legality when
all leaves and constructors remain inside
\(\mathcal R^{\mathcal B}_n\). If the composition produces a
quotient, potential, geometry, valid lifted interface, boundary readout,
forward observable transport component, or finite-horizon trajectory
statistic by an admissible within-budget computation, Lemma
\hyperref[lem-budgeted-generation-gives-representation]{Budgeted
generation gives budgeted representation} places that denotation in the
budgeted saturated closure at its computation cost. If it produces only
an uncomputed stopped-future or exact resolvent semantic object, it is
not an admissible budgeted datum.

\end{definition}

\Cref{fig-three-cost-ledgers} displays the three accounts of
\cref{def-three-cost-ledgers} as views of one budgeted closure.

\begin{figure}[tbp]
\centering
\begin{tikzpicture}[fw diagram]
  \node[fw source,minimum width=4.0cm] (public) at (0,2.5)
    {legal public data\\\(\mathcal R_n^{\mathcal B}\)};
  \coordinate (split) at (0,1.7);
  \node[fw provenance,minimum width=3.1cm] (prov) at (-4.2,0.5)
    {provenance\\generated from presentation};
  \node[fw use,minimum width=3.1cm] (exec) at (0,0.5)
    {execution\\resource envelope};
  \node[fw geom,minimum width=3.4cm] (struct) at (4.2,0.5)
    {structural/spectral\\saturated representation};
  \node[fw terminal,minimum width=5.0cm] (profiles) at (0,-1.4)
    {\(m^{\mathrm{sat}}(B),\ G^{\mathrm{sat}}(B),\ I^{\mathrm{sat}}(B),\ B^*\)};
  \draw[fw arrow] (public) -- (split);
  \foreach \n in {prov,exec,struct}
    \draw[fw arrow] (split) -| (\n);
  \draw[fw arrow] (prov) -- (profiles);
  \draw[fw arrow] (exec) -- (profiles);
  \draw[fw arrow] (struct) -- (profiles);
\end{tikzpicture}
\caption{Provenance, execution, and structural/spectral representation are
three accounts of the same legal data source and resource ledger, as defined
in \cref{def-three-cost-ledgers}.}
\label{fig-three-cost-ledgers}
\end{figure}

\begin{remark}[Structural projection of forward transport]\label{rem-forward-transport-not-inert}

A forward observable transport kernel is included in
\(\mathcal R_{n,\le B}^{\mathrm{sat}}\) so that an algorithm-induced
transition rule is part of the saturated accounting. The extraction
functional has no additional transport summand; it evaluates such a kernel by
applying the \PartII channel projection to the represented generator
that drives the killed propagation and to the terminal edge gate
\(W_T^L\). The full-history hitting functional then integrates those projected
histories as in Definition
\hyperref[def-full-history-target-transport]{Full-history target-contact
transport}. A represented quotient propagation contributes to
\(m_n^{\mathrm{sat}}(B)\); a represented geometric or causal propagation
contributes to \(G_n^{\mathrm{sat}}(B)\); valid Lift and Bdry retain
their full carrier fields and contextual charge unless a represented
transfer bounds that charge by the source they reorganize or read; and the
post-projection remainder is \(\Res\). This projection is
component accounting on the represented transition rule together with
the exposed terminal edge gate. The projection itself neither computes
a resolvent or committor nor changes the reference measure, and it does
not treat the uncomputed stopped future as a representative. Every
target-contact-bearing represented history is
counted through the existing Geom/Caus or Abs coordinate, while
uncomputed stopped-future semantics are handled by derivability
accounting below.

\end{remark}

This definition charges a committor according to the derivation that makes it available.
\begin{definition}[Committor derivability accounting]\label{def-no-free-committor}

For an induced killed generator \(L\), the discounted committor or
target value function

\[
u_{\lambda,L}=(\lambda I-L_N)^{-1}k_T
\]

is a semantic function on the lifted configuration graph
\citep{MetznerSchuetteVandenEijnden2009}. It is counted
in \(\mathcal R_{n,\le B}^{\mathrm{sat}}\), and hence in
\(G_n^{\mathrm{sat}}(B)\) or boundary accounting, precisely when the
same function has a representation of cost at most \(B\) by the
observable structural grammar: an exposed potential or geometry, a
represented quotient, a valid lifted interface, a terminal value
readout, or a legal finite-horizon computation built from
budget-admissible data. The operation
\(L\mapsto(\lambda I-L_N)^{-1}k_T\) becomes a budget-admissible
representation only through an admissible constructor. A finite-horizon
first-hit quotient, hit indicator, or truncated value produced by
simulating the represented one-step rule for a within-budget horizon is
a budget-admissible representative at its saturated computation cost.
Thus the derivability closure contains computed future observables and
requires an explicit derivation for each semantic object.

\end{definition}

\begin{lemma}[Truncated-Neumann derivability accounting]\label{thm-no-truncated-neumann-bypass}

Let \(P_N\) be a killed one-step kernel or \(L_N\) a killed generator
induced by a presentation-relative oracle-free budget-admissible rule,
and let \(k_T\) be the exposed terminal edge gate. For any coefficients
\(a_0,\ldots,a_M\) and horizon \(M=M(n)\) generated within the declared
budget, set

\[
u_M=\sum_{j=0}^M a_jP_N^jk_T.
\]

This expression has legal budgeted provenance exactly when \(P_N\),
\(k_T\), and the coefficients are generated inside the budgeted
saturated observable closure. A truncated finite-horizon expression that
is actually generated by an admissible within-budget simulation is itself a
budget-admissible representative at its saturated computation cost. It
contributes at budget \(B\) precisely when that cost and its quotient or
Geom/Caus/Bdry admissibility conditions place the denoted function \(u_M\) in
\(\mathcal R_{n,\le B}^{\mathrm{sat}}\). If such a representative
exists, its target contribution is already included in
\(m_n^{\mathrm{sat}}(B)\) or \(G_n^{\mathrm{sat}}(B)\). If no such
representative exists at budget \(B\), then the expression defines no
within-budget structural component, although it may enter at a larger
budget.

Consequently, a truncated Neumann approximation to

\[
(\lambda I-L_N)^{-1}k_T
\quad\text{or}\quad
(\lambda+\rho)^{-1}\sum_{j=0}^M\left(\frac{\rho}{\lambda+\rho}\right)^jP_N^jk_T
\]

remains subject to the saturated filtration. Either the approximant is
generated as a legal finite-horizon observable representative and is
counted by \(I_n^{\mathrm{sat}}(B)\) at its saturated computation cost,
or it is only an uncomputed stopped-future denotation and has no
low-cost observable representation. The theorem concerns the
approximant as a quotient, potential, or boundary value; the represented
generator, its \PartII component projections, and the full-history
hitting accounting remain part of the saturated accounting.

\end{lemma}

\begin{proof}\phantomsection\label{proof-no-truncated-neumann-bypass}

The expression \(u_M\) is obtained by finite composition and finite
linear postprocessing of budget-admissible data. If \(M\), the one-step
rule, terminal edge gate, and coefficients lie in the budgeted saturated
closure and their accumulated simulation charge lies within budget, then this
computation is a budget-admissible representative of \(u_M\) at its
saturated computation cost. By Definition
\hyperref[def-saturated-representation-degree]{Saturated
budget-admissible representation cost}, the cost charged to a denoted
object is the infimum over all budget-admissible representatives of that
denotation. Hence \(u_M\) belongs to
\(\mathcal R_{n,\le B}^{\mathrm{sat}}\) exactly when its least
budget-admissible representative has cost at most \(B\); if the explicit
finite computation lies within budget, the declared budget contains it; in
the polynomial instance this becomes the usual fixed-exponent statement.
If the representative is a quotient, its projection kernel is one of the
kernels optimized in \(m_n^{\mathrm{sat}}(B)\). A represented potential
or geometry is optimized in \(G_n^{\mathrm{sat}}(B)\). If the object is
a terminal value readout or valid lifted interface over such data, its
complete carrier is retained in the structural profile; it is dominated
by the smaller Geom/Caus datum only through a represented contextual
transfer. For budgets below its actual saturated cost,
\(u_M\) is not yet counted. The full exact resolvent or infinite
stopped-future value is different: without a budget-admissible finite
representative or other observable representation, semantic notation for
that exact future is not a low-cost observable.

\end{proof}

\begin{lemma}[Budgeted full-history composition bound]\label{thm-budgeted-full-history-gaplessness}

Fix a task family and a budget-compatible resource envelope \(B(n)\).
Let \(L_n\) be the transcript generator of a presentation-relative
uniform oracle-free computation, including the clock and every retained
outcome, branch, public-history, and terminal record, with internal
state, memory, and tables charged at peak native capacity. Suppose
the computation has horizon at most \(H_B(n)\), and let \(\alpha_I\) be
its probability of reaching the public target verifier within that
horizon. Let \(\Psi(B,n)\) be a uniform saturated-cost bound for the
canonical first-hit quotient furnished by
\Cref{lem:exact-affordable-first-hit}. Then

\begin{equation}
\label{eq-budgeted-full-history-profile-bound}
\alpha_I
\le
eH_B(n)
m_I^{\mathrm{sat}}(\Psi(B,n)).
\end{equation}

The inequality applies to every finite composition of represented
updates, postprocessings, outcome branches, valid lifted coordinates,
boundary readouts, and all five channel components together with
\(\Res\). It does not require a componentwise resolvent
expansion, a positive partial sum of channel projections, or a residual
baseline assumption. Consequently, a non-negligible full-history target
arrival forces non-negligible represented quotient visibility at the
charged enlarged saturated budget.

\end{lemma}

\begin{proof}\phantomsection\label{proof-budgeted-full-history-gaplessness}

Execute the operational rule once while retaining its finite outcome and
terminal transcript, then replay only the terminal flags with the
deadline coordinate as in \Cref{lem:exact-affordable-first-hit}. The transcript law is exactly the
law of the original computation. The execution ledger charges one peak
native state, including all memory, branching, postprocessing, and
component interaction, together with the accumulated transcript.
The lemma constructs a uniform exactly lumpable quotient of saturated
cost at most \(\Psi(B,n)\) and gives

\[
\alpha_I\le eH_B(n)V_I(q_{\mathcal A}).
\]

Since \(q_{\mathcal A}\) is one of the quotient families in the
definition of \(m_I^{\mathrm{sat}}(\Psi(B,n))\), this is
\eqref{eq-budgeted-full-history-profile-bound}. The argument is applied
before any channel expansion, so it includes \(\Res\) and
all cancellations automatically.

\end{proof}

\begin{remark}[Committor cases]\label{rem-committor-cases}

If the presentation exposes a low-cost committor, target value
observable, or hitting-time partition, then the problem has a low-cost
progress observable and \(G_n^{\mathrm{sat}}(B)\) or
\(m_n^{\mathrm{sat}}(B)\) is nonzero at that budget. If the budgeted
saturated observable closure derives such a value by an admissible
within-budget finite-horizon computation or by other allowed
constructors, the same conclusion holds, with the value charged at its
saturated computation cost. If the only description is semantic, namely
``the infinite future of this induced kernel'' or ``the exact resolvent
for this induced kernel'' with no legal finite representative at the
budget, then the value exists extensionally but has no admissible
budgeted representation. This distinguishes semantic existence from a
charged observable representation.

\end{remark}

\subsection{Ambient geometric extraction}
\label{subsec:ambient-geometric-extraction}

The ambient coordinate combines represented Dirichlet geometry, authorized
directed alignment, and target reach.

\begin{definition}[Static Dirichlet geometry]\label{def-static-dirichlet-geometry}

A budget-\(B\) pointwise geometry is a pair \(\Phi=(E_\Phi,c_\Phi)\),
where \(E_\Phi\subseteq X\times X\) is a symmetric declared
comparability relation and \(c_\Phi(x,z)=c_\Phi(z,x)\ge0\) is an exposed
conductance represented with saturated cost at most \(B\). Its static
Dirichlet form is

\[
\mathcal E_\Phi(f,g)
=
\frac12\sum_{x,z\in X}c_\Phi(x,z)(f(x)-f(z))(g(x)-g(z))\mu(x)\mu(z).
\]

This is a quadratic form measuring variation across declared comparable
states. Here it serves as a regularity functional for hitting-time
estimates; \PartII associates such forms with generators.

\end{definition}

The ambient extraction coordinate measures target-visible geometry and authorized causal alignment.
\begin{definition}[Ambient Geom/Caus extraction coordinate]\label{def-geom-extraction}

Fix a Caus authority envelope \(\mathfrak A\). Let
\(\mathcal G^{X,\mathfrak A}_{n,\le B}\) be the family of represented
records \((\mathfrak N,\Phi,D,J,\mathfrak K)\) of total saturated cost
at most \(B\), where \(\mathfrak N\) is the charged normalized-kernel
record of \Cref{lem-lossless-charged-uniformization}, \(\Phi\) is a pointwise geometry,
\(D:X\to[0,\infty)\) is a progress observable vanishing on the positive
terminal sector, \(J\) is a directed current on the represented edge
space of \(\Phi\), and \(\mathfrak K\) proves that \(J\) is authorized
by a generator selected from \(\mathfrak A\). In particular, the
canonical current \(J=-\nabla_\Phi D\) is admissible only when it is
formed from the same represented conductance and potential and the same
full generator supplies its local or carrier comparison, with the
complete construction cost charged.  Each record family is produced by
one derivation scheme fixed across all sizes, and its cost bound is
uniform on \(\mathfrak I_n\), exactly as in
\Cref{def-abs-extraction}. For
\(\operatorname{Var}_\mu(D)>0\), define

\[
\rho_\Phi(D)=\frac{\mathcal E_\Phi(D,D)}{\operatorname{Var}_\mu(D)}.
\]

Let \(M_\Phi\in\{0,1\}\) be the mixing indicator, equal to one when the
declared geometry has a Poincare or conductance lower bound whose
reciprocal lies in the declared transfer class. In the
\(\mathcal B_{\mathrm{poly}}\) instance this is the usual
inverse-polynomial scale. Let
\(\chi_\Phi^{\mathrm{auth}}(D)\in\{0,1\}\) be the local consistency
gate, equal to one when the positive-defect landscape has no represented
local obstruction away from \(T\) on the declared comparison window
for the same selected generator. Define the quantitative Caus factor

\[
C_{\Phi,\mathfrak K}^{\mathrm{auth}}(D)
=
\chi_\Phi^{\mathrm{auth}}(D)
\operatorname{Align}^{\mathrm{auth}}_{\Caus}
(\mathfrak K,D).
\]

Let \(\mathcal M(D)\) be the closed span of monotone functions of \(D\)
oriented toward \(D=0\), and define the target reach of \(D\) by

\[
R_T(D)=\frac{\|\Pi_{\mathcal M(D)}y\|^2}{\|y\|^2}.
\]

The conditioned ambient Geom/Caus visibility contribution is

\[
V^X_{\Phi,D,J,\mathfrak K}:=
M_\Phi C_{\Phi,\mathfrak K}^{\mathrm{auth}}(D)
\frac{R_T(D)}{1+\rho_\Phi(D)}.
\]

The authority-relative ambient geometric extraction coordinate is

\[
G_{X,n,\mathfrak A}^{\mathrm{sat}}(B)=
\sup_{(\mathfrak N,\Phi,D,J,\mathfrak K)
\in\mathcal G^{X,\mathfrak A}_{n,\le B}}
V^X_{\Phi,D,J,\mathfrak K}.
\]

Because the free envelope ranges over every budget-admissible
represented selected skew operator, one has
\(\mathcal C_{QG}^{\mathrm{all}}(B)
=\mathcal C_{QG}^{\mathrm{free}}(B)\).  Write
\(G_{X,n}^{\mathrm{sat}}(B)\).

\end{definition}

The reach \(R_T(D)\) measures alignment of the defect with the target,
\(\rho_\Phi(D)\) measures its regularity across declared comparable
states, and \(C_{\Phi,\mathfrak K}^{\mathrm{auth}}(D)\) measures the
descending fraction of the authorized represented current. The mixing
and consistency gates impose the stated local comparison conditions for
the selected generator. The corresponding target-arrival
conclusion follows from an applicable arrival theorem, such as
\Cref{thm-caus-drift-survives-defect} when a pointwise drift
margin is available.

\section{Quotient Extraction and the Derivation Normal Form}
\label{sec:quotient-extraction-normal-form}

We now define exact quotient visibility and expose its complete represented
derivation.

\begin{definition}[\PartI quotient-utility condition]\label{def-paper1-abs-usefulness-gate}

Fix a represented-fiber quotient \(Q\), an induced budget-admissible
operator \(L\), a legal initialization \(\mu_0\), and a hitting-time
discount \(\lambda\). The quotient-utility factor
\(A_Q=A_Q(\lambda;L,\mu_0)\) is the \PartI fiber-represented-observable
utility condition, transported to the comparison window. It is a number in \([0,1]\),
and it is set to zero unless all of the following structural conditions
hold.

\begin{enumerate}

\item

\textbf{Represented fibers.} \(Q\) is a fiber-represented observable in
the \PartI sense: fiber membership, construction of the compact fiber
description, and finite Boolean operations on represented fibers have
saturated cost at most the current budget.  No point selection from a
fiber is required.

\item

\textbf{Terminal quotient event.} The target and terminal/failure
sectors used by the hitting functional are unions of quotient cells.

\item

\textbf{Exact dynamic compatibility.} On the target window, the induced
dynamics satisfy the exact intertwining \(L_XU=UL_Q\). A quotient with
nonzero defect contributes only through the compensated Geom/Caus
certificate of \Cref{def-supported-geom-certificate}; its
defect loss is not absorbed into the pure \(\Abs\) coordinate.

\item

\textbf{Quotient utility.} The represented quotient dynamics,
represented fiber operations, or quotient-level observables have
strictly positive discounted target functional from the legal
initialization:
\(A_T^Q(\lambda;L_Q,\bar\mu_0)>0\).

\item

\textbf{Nontrivial public structure.} The quotient is generated from the
public observable algebra and, after empty fibers are removed, its
partition strictly refines the verifier partition:
\[
\sigma(\mathbf 1_T)\subsetneq\sigma(q).
\]
Equivalently, at least one represented nonterminal quotient distinction
is not determined by \(\mathbf 1_T\). A
large target-containing cell may contribute visibility only through the
normalized projection below; it need not include a selector or sampling
procedure.

\end{enumerate}

When these conditions hold, let \(A_T^Q(\lambda;L_Q,\bar\mu_0)\) be the
target-projected quotient hitting functional from the pushed-forward
initialization \(\bar\mu_0\), and let \(C_Q\ge1\) be the product of the
charged exact-quotient comparison and normalization factors. We set

\[
A_Q(\lambda;L,\mu_0)
=
\min\{1,\lambda C_Q^{-1}A_T^Q(\lambda;L_Q,\bar\mu_0)\}.
\]

For the canonical pushforward measure and initialization of a
terminal-compatible exact quotient, \Cref{prop-canonical-exact-quotient-unit-transfer} gives \(C_Q=1\).
A nonunit factor is retained only when the quotient record explicitly
uses a different represented gauge or an additional comparison map.

Thus \(A_Q\) is the normalized structural utility of a represented
quotient on the active comparison window. It records exposed quotient
alignment independently of any target solver.

\end{definition}

The quotient extraction coordinate measures the target-visible value of an exact represented quotient.
\begin{definition}[Quotient extraction coordinate]\label{def-abs-extraction}

Let \(\mathfrak I_n\) be a size-\(n\) task slice. Define
\(\mathcal Q_{\mathfrak I_n,\le B}^{\mathrm{sat}}\) to be the
families \(\mathbf Q=(\mathscr R_I)_{I\in\mathfrak I_n}\) produced by
one uniform exactly lumpable represented-fiber quotient derivation of
saturated cost at most \(B\) uniformly on the slice. The family is the
restriction to size \(n\) of one derivation scheme
\(\rho\in\mathsf{Der}_{\mathfrak I}\) fixed across all sizes and
instances; the program is not re-chosen for each \(n\). Here a carrier
record

\[
\mathscr R_I=
(\Omega_I,\mu_{\Omega,I},T_{\Omega,I},\mu_{0,\Omega,I},
\mathfrak N_{\Omega,I},q_I,Q_I)
\]

contains a finite represented carrier, a legal public full-support
reference law, an exposed terminal sector, a legal initialization, the
charged normalization record
\(\mathfrak N_{\Omega,I}=(\widetilde L_{\Omega,I},\zeta_I,c_I,
P_{\Omega,I},L_{\Omega,I},\lambda_I)\), with
\(L_{\Omega,I}=P_{\Omega,I}-I\), and the represented exact quotient.
The utility factor is evaluated using this same \(L_{\Omega,I}\) and
\(\lambda_I\).  Thus normalization belongs to the exact-profile cost
from the outset; it is not added later as a cost-free extension. The
carrier may be the base task carrier or any budget-generated clocked
transcript representation produced by the computation-closure
clause. Its encoding, initialization, update, target verifier,
reference gauge, quotient fibers, and comparison data are all charged
in the same derivation. Such a computation carrier need not satisfy the
Valid-Lift projection or target-fraction identities. Those identities
are required only if the record separately claims a structural
\(\Lift\) contribution or transports a comparison back to a
base task. This uniformity excludes a separate carrier, quotient
program, or fiber table for each instance.
For \(\mathbf Q\in
\mathcal Q_{\mathfrak I_n,\le B}^{\mathrm{sat}}\), let \(P_{Q_I}\) be
orthogonal projection in \(L^2(\Omega_I,\mu_{\Omega,I})\) onto the
quotient-observable subspace on instance \(I\), let

\[
y_{\Omega,I}=\mathbf 1_{T_{\Omega,I}}
-\mu_{\Omega,I}(T_{\Omega,I}),
\]

and let \(A_{Q_I}\in[0,1]\) be the \PartI quotient-utility
factor from Definition
\hyperref[def-paper1-abs-usefulness-gate]{Part I quotient-utility
condition}, evaluated on the current comparison window. The conditioned quotient
visibility contribution at \(I\) is

\[
V_I(\mathbf Q):=
A_{Q_I}\frac{\|P_{Q_I}y_{\Omega,I}\|^2}
{\|y_{\Omega,I}\|^2},
\]

with value zero when the carrier target has reference mass zero or one.
The base-carrier case has \(\Omega_I=X_I\) and
\(y_{\Omega,I}=y_I\).

When the instance and uniform family are fixed, write this quantity as
\(V_Q\).

The instance-level abstraction extraction coordinate is

\[
m_I^{\mathrm{sat}}(B)=
\sup_{\mathbf Q\in\mathcal Q_{\mathfrak I_n,\le B}^{\mathrm{sat}}}
V_I(\mathbf Q),
\qquad \sup\varnothing:=0.
\]

For a distinguished instance sequence \((I_n)\), retain the shorthand
\(\mathcal Q_{n,\le B}^{\mathrm{sat}}\) for the corresponding family
of quotients and write
\(m_n^{\mathrm{sat}}(B)=m_{I_n}^{\mathrm{sat}}(B)\). In log-cost
notation, \(m_n(d)=m_n^{\mathrm{sat}}(2^d)\).

\end{definition}

The exact quotient normal form expands a represented derivation into the record used by the profile.
\begin{definition}[Derivation-induced exact quotient normal form]
\label{def-derivation-induced-quotient-normal-form}

An exact quotient normal-form record is an expanded derivation
\(\widehat\rho\in\mathsf{Der}^{\mathrm{nf}}_{\mathfrak I}\) whose
terminal vertex denotes a finite map

\[
q_{\widehat\rho,I}:\Omega_I\longrightarrow Q_{\widehat\rho,I}
\]

on the base space or on a budget-generated clocked transcript carrier.
The latter is legal for quotient accounting by computation closure and
does not have to be a structural Valid Lift. The record contains the
carrier encoding, legal reference gauge, initialization, exposed
terminal sector, and the same charged normalized-kernel record required
by \Cref{def-abs-extraction} in addition to the
data below. The record
may use a finite public label set containing unused labels; an unused
label has the represented false fiber predicate, while every nonempty
fiber has the representation required below. The record
contains the quotient evaluation rule, represented nonempty fibers and
their finite Boolean operations, the terminal cell sets, the induced
quotient operator, and the comparison and normalization data entering
\(A_{q_{\widehat\rho,I}}\). These data are typed attachments to the
terminal quotient interface and are copied into every backward-cut
record, whether they are stored as separate vertices or as fields of the
terminal vertex. The fiber map and its induced operator
satisfy exact intertwining. Its profile contribution is evaluated by
\Cref{def-paper1-abs-usefulness-gate}, which assigns
\(A_{q_{\widehat\rho,I}}=0\) unless the remaining terminal, utility, and
nontriviality conditions hold.

Let \(\mathsf{NFQ}_{\mathfrak I_n,\le B}\) be the family of such uniform
normal-form records whose actual recorded witness cost satisfies
\(\operatorname{Cost}^{\mathrm{nf}}_I(\widehat\rho)\preceq B\) for
every \(I\in\mathfrak I_n\).

Thus membership bounds the chosen record itself, not merely the cost of
another record with the same denotation. Local finiteness makes
a Pareto-minimal normal-form witness available below any chosen
finite-cost witness. The budgeted statements below use the actual
normal-form witness of cost \(\preceq B\), not a globally least element.
Let
\(\mathcal Q^{\mathrm{nf}}_{\mathfrak I_n,\le B}\) be their quotient
denotations, and write
\(\mathsf{NFQ}_{\mathfrak I_n}:=\bigcup_B
\mathsf{NFQ}_{\mathfrak I_n,\le B}\).

\end{definition}

\begin{lemma}[Affordable exact quotient derivation normal form]
\label{thm-affordable-exact-quotient-normal-form}

For every presented finite task family and every saturated budget
\(B\),

\begin{equation}
\label{eq-exact-quotient-normal-form-identity}
\mathcal Q^{\mathrm{sat}}_{\mathfrak I_n,\le B}
=
\mathcal Q^{\mathrm{nf}}_{\mathfrak I_n,\le B}.
\end{equation}

Consequently, the instance-level quotient profile has the exact
derivation-indexed representation

\[
y_{\widehat\rho,I}
:=\mathbf1_{T_{\Omega_{\widehat\rho,I}}}
-\mu_{\Omega_{\widehat\rho,I}}
  (T_{\Omega_{\widehat\rho,I}}),
\]

where the corresponding summand is defined to be zero when
\(\|y_{\widehat\rho,I}\|_2=0\), as in \Cref{def-abs-extraction}.

\begin{equation}
\label{eq-derivation-indexed-quotient-profile}
m_I^{\mathrm{sat}}(B)
=
\sup_{\widehat\rho\in
\mathsf{NFQ}_{\mathfrak I_n,\le B}}
A_{q_{\widehat\rho,I}}
\frac{\|P_{q_{\widehat\rho,I}}y_{\widehat\rho,I}\|_2^2}
{\|y_{\widehat\rho,I}\|_2^2},
\end{equation}

where \(\sup\varnothing=0\). Thus the exact \(\Abs\) profile ranges
over the complete resource-filtered derivation calculus, rather than
over a semantic classification of the resulting pointwise fiber
algebras.

\end{lemma}

\begin{proof}

Let
\(\mathbf Q\in\mathcal Q^{\mathrm{sat}}_{\mathfrak I_n,\le B}\).
By \Cref{def-abs-extraction}, one uniform derivation of cost at
most \(B\) produces its quotient map, fibers, dynamic intertwining, and
utility data. \Cref{thm-expanded-derivation-equivalence} expands
that record without changing its denotation or cost, so
\(\mathbf Q\in\mathcal Q^{\mathrm{nf}}_{\mathfrak I_n,\le B}\).
Conversely, contraction of a normal-form quotient record gives an
admissible derivation with identical quotient data and cost. This proves
\eqref{eq-exact-quotient-normal-form-identity}. Substitution of the
identical indexing family into \Cref{def-abs-extraction}
gives \eqref{eq-derivation-indexed-quotient-profile}.

\end{proof}

\subsection{The charged backward cut}
\label{subsec:charged-backward-cut}

The backward cut isolates the first joint frontier carrying all
state-dependent information used by the terminal quotient.

\begin{definition}[Charged backward cut and joint-source load]
\label{def-charged-backward-source-cut}

Let
\(\widehat\rho\in\mathsf{NFQ}_{\mathfrak I_n,\le B}\)
be an expanded exact-quotient record, and let
\(\mathsf G_{\widehat\rho}\) be the backward slice of its terminal
quotient-map vertex.  A \emph{charged backward cut} is an ordered finite
antichain

\[
\Gamma=(v_1,\ldots,v_k)
\]
in \(\mathsf G_{\widehat\rho}\) that separates every represented
state-dependent input of the terminal slice from the terminal vertex.
The singleton containing the terminal quotient-map vertex is permitted
as a cut and separates these paths at their common terminal endpoint.
Every state-dependent retained lifted coordinate or derived datum read
below the cut is included among the cut outputs.  State-independent public
parameters may remain as represented coefficients in the downstream
record.

The record carries a unique finite carrier
\((\Omega_I,\mu_{\Omega,I},T_{\Omega,I})\). Throughout this backward-cut
and structural-charge subsection, write

\[
y_I:=y_{\Omega,I}
=\mathbf1_{T_{\Omega,I}}-\mu_{\Omega,I}(T_{\Omega,I}),
\]

and take every projection in \(L^2(\Omega_I,\mu_{\Omega,I})\). This
local shorthand includes the base carrier as a special case and does
not replace a computation-carrier contrast by the base-task contrast.

If \(f_{v,I}\) denotes the output of a cut vertex, write

\[
J_{\Gamma,I}
=
(f_{v_1,I},\ldots,f_{v_k,I}).
\]

The cut record consists of the union of the ancestral records of the
vertices in \(\Gamma\), their represented joint-fiber interface, and the
complete downstream subrecord terminating in
\(q_{\widehat\rho,I}\).  In particular, the downstream record retains
the quotient operator, terminal cells, fiber predicates, comparison
data, and the utility factor \(A_{q_{\widehat\rho,I}}\).

For a represented observable \(F\), set

\[
\mathsf p_I(F)
:=
\frac{\|P_Fy_I\|_2^2}{\|y_I\|_2^2},
\]

and define the joint-source load carried across \(\Gamma\) by

\begin{equation}
\label{eq-joint-source-load}
\operatorname{Load}_I(\widehat\rho,\Gamma)
:=
A_{q_{\widehat\rho,I}}\,\mathsf p_I(J_{\Gamma,I}).
\end{equation}

The factor \(A_{q_{\widehat\rho,I}}\) belongs to the complete terminal
record.  Thus a public terminal discriminator contributes projection
mass only with the dynamic utility supplied by the represented
transport and quotient record.

\end{definition}

The backward factorization exposes the charged sources on which a terminal quotient depends.
\begin{theorem}[Quantitative backward factorization through a charged
joint frontier]
\label{thm-quantitative-backward-source-factorization}

Let \(\widehat\rho\) and \(\Gamma\) be as in
\Cref{def-charged-backward-source-cut}. Topological
contraction of the downstream subrecord gives one represented
family-uniform map \(h_{\widehat\rho,\Gamma,I}\) such that

\begin{equation}
\label{eq-backward-cut-factorization}
q_{\widehat\rho,I}
=h_{\widehat\rho,\Gamma,I}\circ J_{\Gamma,I}.
\end{equation}

If materializing the cut tuple requires a separate interface, the
complete cut record has cost

\begin{equation}
\label{eq-backward-cut-cost}
\operatorname{Cost}_{I'}(\widehat\rho_\Gamma)
\preceq
B\oplus\operatorname{Cost}_{\mathrm{join}}(\Gamma)
=:\Psi_{\mathrm{cut}}(B,\Gamma),
\qquad I'\in\mathfrak I_n.
\end{equation}

When the tuple interface is already retained in \(\widehat\rho\), no
additional charge is incurred. The map \(\Psi_{\mathrm{cut}}\) is
class preserving whenever the cut arity and represented join interface
lie in the declared transfer class; in particular, it is polynomial for
polynomial-cost expanded records.

For every instance,

\begin{align}
\mathsf p_I(J_{\Gamma,I})
&=
\mathsf p_I(q_{\widehat\rho,I})
+\frac{
\|(P_{J_{\Gamma,I}}-P_{q_{\widehat\rho,I}})y_I\|_2^2}
{\|y_I\|_2^2},
\label{eq-backward-cut-pythagoras}\\
V_I(q_{\widehat\rho})
&\le \operatorname{Load}_I(\widehat\rho,\Gamma).
\label{eq-backward-cut-load-bound}
\end{align}

Moreover, put \(J_{\Gamma,\ell}=(f_{v_1},\ldots,f_{v_\ell})\) for
\(\ell\ge1\), let \(J_{\Gamma,0}\) be the constant observable, and
write \(P_{J_{\Gamma,0,I}}=P_{\Gamma,0}\) for projection onto
constants. Define

\begin{equation}
\label{eq-backward-cut-conditional-increment}
\Delta_{I,\ell}(\widehat\rho,\Gamma)
:=
\frac{
\|(P_{J_{\Gamma,\ell,I}}-P_{J_{\Gamma,\ell-1,I}})y_I\|_2^2}
{\|y_I\|_2^2}.
\end{equation}

Then

\begin{equation}
\label{eq-backward-cut-increment-sum}
\mathsf p_I(J_{\Gamma,I})
=\sum_{\ell=1}^k\Delta_{I,\ell}(\widehat\rho,\Gamma),
\qquad
V_I(q_{\widehat\rho})
\le
A_{q_{\widehat\rho,I}}
\sum_{\ell=1}^k\Delta_{I,\ell}(\widehat\rho,\Gamma).
\end{equation}

Consequently, when \(k\ge1\), some conditional join extension
\((J_{\Gamma,\ell-1},f_{v_\ell})\) satisfies

\begin{equation}
\label{eq-load-bearing-conditional-join}
A_{q_{\widehat\rho,I}}
\Delta_{I,\ell}(\widehat\rho,\Gamma)
\ge
\frac{V_I(q_{\widehat\rho})}{k}.
\end{equation}

The load-bearing object in
\eqref{eq-load-bearing-conditional-join} is the charged conditional
join, including its preceding context and its complete downstream
terminal record. No conclusion about one cut coordinate in isolation is
required.

\end{theorem}

\begin{proof}

The terminal backward slice is a finite directed acyclic graph.  Treat
the outputs at \(\Gamma\) as inputs and evaluate the remaining vertices
in topological order.  The separating property ensures that every
state-dependent input used by this evaluation occurs in
\(J_{\Gamma,I}\); represented public coefficients remain in the same
downstream record.  This produces
\(h_{\widehat\rho,\Gamma,I}\) and proves
\eqref{eq-backward-cut-factorization}.  The ancestral records, the
downstream evaluator, and its terminal interface are subrecords of
\(\widehat\rho\).  Only a newly materialized tuple interface adds a
constructor charge, which proves \eqref{eq-backward-cut-cost}.  In
particular, a formula, table, recurrence, branching rule, or fiber
selector used below the cut remains in the downstream subrecord with
its construction, storage, evaluation, and fiber-operation costs.

Equation \eqref{eq-backward-cut-factorization} gives
\(\sigma(q_{\widehat\rho,I})\subseteq
\sigma(J_{\Gamma,I})\).  Apply \Cref{thm-affordable-composition-no-creation} to obtain
\eqref{eq-backward-cut-pythagoras}.  Since
\(0\le A_{q_{\widehat\rho,I}}\le1\), multiplication by this same
terminal utility factor gives

\[
V_I(q_{\widehat\rho})
=A_{q_{\widehat\rho,I}}\mathsf p_I(q_{\widehat\rho,I})
\le
A_{q_{\widehat\rho,I}}\mathsf p_I(J_{\Gamma,I}),
\]

which is \eqref{eq-backward-cut-load-bound}.  Finally, the spaces
generated by
\(J_{\Gamma,0},J_{\Gamma,1},\ldots,J_{\Gamma,k}\) are nested.
Successive conditional-expectation differences are orthogonal, and
\(P_{\Gamma,0}y_I=0\) because \(y_I\) is centered.  Pythagoras therefore
gives the first identity in \eqref{eq-backward-cut-increment-sum}; the
second follows from \eqref{eq-backward-cut-load-bound}.  At least one
of the \(k\) nonnegative summands is at least their average, which
proves \eqref{eq-load-bearing-conditional-join}.

\end{proof}

\begin{remark}[Projection ancestry and dynamic source transfer]
\label{rem-projection-ancestry-dynamic-source-transfer}

\Cref{thm-quantitative-backward-source-factorization} is an
unconditional factorization and cost-accounting result.  The joint cut
observable need not be lumpable, terminal compatible, or dynamically
qualified.  For an exact terminal-compatible quotient whose terminal
sector is a union of quotient cells, the centered terminal contrast is
already quotient measurable; in that case
\(\mathsf p_I(q_{\widehat\rho,I})=1\), and the nontrivial quantity is
the utility factor \(A_{q_{\widehat\rho,I}}\).  Transferring that factor
to a proper ancestral source requires the represented dynamic comparison
in \eqref{eq-quantitative-dynamic-source-transfer}.  This distinction
also permits joint effects: two cut coordinates may have zero individual
increments before conditioning on their joint context.

\end{remark}

The compressed-source cover packages the finitely many source envelopes needed for a budgeted audit.
\begin{definition}[Quantitative compressed-source cover]
\label{def-quantitative-dynamic-source-cover}

Let \(\mathfrak C_{\mathrm{src}}\) be a finite set of represented
source signatures and, for each
\(c\in\mathfrak C_{\mathrm{src}}\), let
\(\mathcal S_I^{(c)}(B)\) be a family of
dynamically qualified represented source certificates of saturated cost
at most \(B\).  Such a certificate is the evaluation of one
family-uniform derivation record containing a terminal target-aligned
object and all dynamic, terminal, comparison, and normalization data
that define its recorded visibility \(V_I(R)\).  In particular, an
exact-quotient record is a source certificate with

\[
V_I(R)
=
A_{q_{R,I}}\mathsf p_I(q_{R,I}).
\]

Write

\[
s_I^{(c)}(B)
:=
\sup_{R\in\mathcal S_I^{(c)}(B)}V_I(R),
\qquad \sup\varnothing:=0.
\]

A quantitative compressed-source cover at budget \(B\) consists of a
class-preserving cost enlargement \(\Psi_{\mathrm{src}}\), a factor
\(K_B\) in the transfer class, an error \(\delta_B\ge0\), and one
family-uniform syntactic rule with the following quantifier order.  For
every
\(\widehat\rho\in\mathsf{NFQ}_{\mathfrak I_n,\le B}\), the rule
selects, independently of the instance, a charged backward cut
\(\Gamma\), a represented witness of a signature
\(c\in\mathfrak C_{\mathrm{src}}\), a uniform source record \(R\), and
a represented comparison record.  Their single transfer carrier is

\[
\mathfrak T_{\widehat\rho}
=
(\widehat\rho^\Gamma,R,\operatorname{Cmp}_{\widehat\rho,\Gamma,R}),
\]

and satisfies, for every \(I\in\mathfrak I_n\),

\[
\operatorname{Cost}_I(\mathfrak T_{\widehat\rho})
\preceq\Psi_{\mathrm{src}}(B),
\qquad
R_I\in\mathcal S_I^{(c)}(\Psi_{\mathrm{src}}(B)).
\]

This carrier includes the materialized joint cut, the source, and the
comparison verifier.  Whenever \(A_{q_{\widehat\rho,I}}>0\), the same
uniformly selected records satisfy

\begin{equation}
\label{eq-quantitative-dynamic-source-transfer}
\operatorname{Load}_I(\widehat\rho,\Gamma)
\le
K_B V_I(R)+\delta_B.
\end{equation}

In particular,
\eqref{eq-quantitative-dynamic-source-transfer} is a dynamic
source-transfer condition: projection monotonicity alone supplies
\eqref{eq-backward-cut-load-bound}, while lumpability, target utility,
and comparison loss for \(R\) are verified by the transfer certificate.
The selection therefore has quantifier order
\(\forall\widehat\rho\,\exists
(\Gamma,c,R,\operatorname{Cmp})\,\forall I\), rather than permitting an
instance-dependent witness. Records with
\(A_{q_{\widehat\rho,I}}=0\) contribute zero to the exact-quotient
profile and require no source witness.

\end{definition}

The finite-envelope theorem bounds all charged backward sources by one family-uniform cover.
\begin{theorem}[Finite compressed-source envelope]
\label{thm-finite-dynamic-source-envelope}

Under \Cref{def-quantitative-dynamic-source-cover},

\begin{equation}
\label{eq-finite-dynamic-source-envelope}
m_I^{\mathrm{sat}}(B)
\le
K_B\max_{c\in\mathfrak C_{\mathrm{src}}}
s_I^{(c)}(\Psi_{\mathrm{src}}(B))
+\delta_B.
\end{equation}

If \(K_B\) lies in the transfer class, \(\delta_B\) is negligible,
and every source envelope on the right is negligible, then
\(m_I^{\mathrm{sat}}(B)\) is negligible.

\Cref{thm-exact-affordable-algebra-structural-charging} below
provides unconditional constructor-level accounting without replacing a
contextual extension by a smaller source. The present hypothesis is
needed precisely when such an extension is compressed to an
independently dynamically qualified certificate.

\end{theorem}

\begin{proof}

Fix \(\widehat\rho\).  If
\(A_{q_{\widehat\rho,I}}=0\), then
\(V_I(q_{\widehat\rho})=0\) and the asserted upper bound is immediate.
Assume \(A_{q_{\widehat\rho,I}}>0\).  \Cref{thm-quantitative-backward-source-factorization} and the source
transfer give

\[
V_I(q_{\widehat\rho})
\le
\operatorname{Load}_I(\widehat\rho,\Gamma)
\le
K_BV_I(R)+\delta_B
\le
K_B\max_{c\in\mathfrak C_{\mathrm{src}}}
s_I^{(c)}(\Psi_{\mathrm{src}}(B))+\delta_B.
\]

Taking the supremum over
\(\widehat\rho\in\mathsf{NFQ}_{\mathfrak I_n,\le B}\) and applying
\Cref{thm-affordable-exact-quotient-normal-form} proves
\eqref{eq-finite-dynamic-source-envelope}.  The negligibility assertion
follows from closure of the transfer class under finite maxima and
multiplication by \(K_B\).

\end{proof}

\section{Affordable Frontiers and Structural Charging}
\label{sec:affordable-frontiers-structural-charging}

Having exposed the backward cut, we decompose its visibility into contextual
increments along an affordable conditional-algebra frontier.

The construction proceeds through the carrier in
\cref{sssec-frontier-carrier}, the exact charge identity in
\cref{sssec-exact-structural-charge}, the constructor covers in
\cref{sssec-constructor-covers}, filtration provenance in
\cref{sssec-filtration-provenance}, and lossless uniformization in
\cref{sssec-lossless-uniformization}.
The corresponding load-bearing results are
\cref{def-affordable-conditional-algebra-frontier,thm-exact-affordable-algebra-structural-charging,cor-universal-syntax-affordable-charge-cover,thm-filtration-provenance-charging,lem-lossless-charged-uniformization}.

\subsection{Frontier records and complete carrier cost}
\label{sssec-frontier-carrier}
\label{subsec:frontier-carrier-cost}

\begin{definition}[Affordable conditional-algebra frontier]
\label{def-affordable-conditional-algebra-frontier}

An affordable frontier consists of the factorization below, the complete
carrier/cost record of \cref{def:affordable-frontier-carrier-cost}, and the
contextual charges of \cref{def:affordable-frontier-contextual-charge}.  The
present definition retains the public label for this combined object.

\medskip\noindent\textbf{(i) Frontier factorization.}
Let
\(\widehat\rho\in\mathsf{NFQ}_{\mathfrak I_n,\le B}\).  An
\emph{affordable conditional-algebra frontier} for
\(\widehat\rho\) is a family-uniform expanded record
\(\widehat\rho^{\Gamma}\) with an ordered finite represented interface

\[
\Gamma=(v_1,\ldots,v_k),
\qquad k\ge0,
\]

whose entries are actual state-dependent vertices in the terminal
backward slice.  The entries may be comparable in the derivation order.
The record contains represented join vertices

\[
J_{\Gamma,0}=*,
\qquad
J_{\Gamma,\ell}
=
(J_{\Gamma,\ell-1},f_{v_\ell}),
\qquad 1\le\ell\le k,
\]

and one represented downstream contraction satisfying

\begin{equation}
\label{eq-affordable-frontier-factorization}
q_{\widehat\rho,I}
=
h_{\widehat\rho,\Gamma,I}
   \circ J_{\Gamma,k,I}
\end{equation}

on every instance.  Pure tuple, relabeling, lift-interface, comparison,
boundary, and terminal-readout vertices remain in the downstream
record whenever their state-dependent arguments occur at earlier
represented vertices.  Those arguments pass through the interface,
while public codes and state-independent parameters remain represented
coefficients of \(h_{\widehat\rho,\Gamma}\).  A declared
state-dependent primitive whose argument is not separately exposed may
instead occur as a frontier entry; a boundary or terminal primitive is
then witnessed by an explicit boundary or terminal constructor class
and cannot be reassigned to a generative class without a represented
transfer.
Frontier entries are strictly ancestral to the terminal quotient for a
proper internal cut. The terminal quotient vertex itself is always
permitted as the sole catch-all frontier entry, as already required by
\Cref{def-charged-backward-source-cut}; its complete ancestral
derivation and downstream identity interface remain in the carrier.
The charged opcode is its actual primitive, boundary, quotient,
postprocessing, or five-family evaluator schema. Thus the catch-all
case preserves every affordability
field rather than treating a derived terminal map as a free primitive.
When \(k=0\), the factorization requires the quotient to be constant;
such a record has zero target charge under the nontriviality gate.

\end{definition}

\begin{definition}[Complete carrier and cost allocation for an affordable frontier]
\label{def:affordable-frontier-carrier-cost}

Let \((\widehat\rho,\Gamma)\) be an affordable conditional-algebra
frontier in the sense of
\cref{def-affordable-conditional-algebra-frontier}.

\medskip\noindent\textbf{(ii) Complete carrier.}
The carrier \(\widehat\rho^{\Gamma}\) retains the complete original
record, all mixed ancestors of the frontier, every join interface, all
later frontier entries, and the full downstream quotient operator,
fibers, dynamics, terminal sectors, comparison data, normalization, and
utility factor. Its uniform charged cost ceiling
\(D(\widehat\rho,\Gamma)\) is, by definition, a represented resource
vector that dominates the sum of the allocated carrier fields below.
In particular,

\begin{equation}
\label{eq-affordable-frontier-carrier-cost}
\operatorname{Cost}_{I'}(\widehat\rho^{\Gamma})
\preceq
D(\widehat\rho,\Gamma)
\qquad(I'\in\mathfrak I_n).
\end{equation}

\medskip\noindent\textbf{(iii) Cost allocation.}
The complete carrier-cost record consists of charged upper-bound fields for
description, evaluation, storage, finite outcome generation, precision, fibers,
interfaces, comparisons, and retained global execution resources.  Each
field is assigned once to the opcode schema of its recorded origin;
nonlocal execution or comparison fields are assigned once to a
distinguished carrier-overhead schema.  Shared fields are not duplicated.
For each schema \(\sigma\), let

\[
B_{\sigma,I'}(\widehat\rho,\Gamma)
\]

be the resource aggregation of the upper-bound fields assigned to
\(\sigma\) on instance \(I'\).  This is an allocation of the bounds
already recorded in the charged carrier, not a second charge on the
displayed trace.  Sub-accumulation of represented cost gives the
complete affordability ledger

\begin{equation}
\label{eq-affordable-frontier-cost-ledger}
\operatorname{Cost}_{I'}(\widehat\rho^{\Gamma})
\preceq
\Lambda_{I'}(\widehat\rho,\Gamma)
:=
\mathop{\bigoplus}_{\sigma}
B_{\sigma,I'}(\widehat\rho,\Gamma)
\preceq D(\widehat\rho,\Gamma).
\end{equation}

\end{definition}

\begin{definition}[Contextual extension and structural charge]
\label{def:affordable-frontier-contextual-charge}

Let \((\widehat\rho,\Gamma)\) carry the complete carrier and cost allocation
of \cref{def:affordable-frontier-carrier-cost}.

\medskip\noindent\textbf{(iv) Contextual charge.}
The \(\ell\)-th \emph{contextual source extension} is the complete
carrier with the inclusion

\[
\sigma(J_{\Gamma,\ell-1,I})
\subseteq
\sigma(J_{\Gamma,\ell,I})
\]

distinguished; denote it by
\(E_{\widehat\rho,\Gamma,\ell}\).  Its cost is the full carrier cost
\eqref{eq-affordable-frontier-carrier-cost}, rather than the marginal
cost of \(v_\ell\).  Its target-aligned structural charge is

\begin{equation}
\label{eq-contextual-source-extension-charge}
\operatorname{ch}_I(E_{\widehat\rho,\Gamma,\ell})
:=
A_{q_{\widehat\rho,I}}
\frac{
 \|(P_{J_{\Gamma,\ell,I}}
    -P_{J_{\Gamma,\ell-1,I}})y_I\|_2^2}
 {\|y_I\|_2^2}.
\end{equation}

The complete affordable structural-charge accounting datum is

\begin{equation}
\label{eq-complete-affordable-structural-charge-record}
\mathbf{SC}_I(E_{\widehat\rho,\Gamma,\ell})
:=
\left(
\operatorname{ch}_I(E_{\widehat\rho,\Gamma,\ell}),
D(\widehat\rho,\Gamma),
\bigl(B_{\sigma,I}(\widehat\rho,\Gamma)\bigr)_{\sigma}
\right).
\end{equation}

This tuple is metamathematical bookkeeping evaluated at \(I\); it is
not an additional represented observable and is not charged as a new
derivation vertex.

Thus the charged object is the conditional algebra extension together
with its complete preceding context and downstream dynamics.  The
coordinate \(f_{v_\ell}\) in isolation is not the charged source, and
no visibility estimate may replace the carrier cost by its marginal
gate cost.

\end{definition}

This cover assigns each contextual charge to a finite affordable constructor schema.
\begin{definition}[Finite affordable constructor-charge cover]
\label{def-finite-affordable-constructor-charge-cover}

A \emph{finite affordable constructor-charge cover} consists of:

\begin{enumerate}

\item a single family-uniform syntactic rule \(\mathfrak F\), fixed
independently of the size, budget, instance, and target values, that
assigns every exact-quotient normal-form record a frontier realization
\(\widehat\rho^{\Gamma}=\mathfrak F(\widehat\rho)\); the bounds below
are required on records with positive utility;

\item class-preserving majorants \(\Psi_{\mathrm{alg}}(B)\) and
\(K_{\mathrm{cut}}(B)\) such that

\begin{equation}
\label{eq-affordable-frontier-majorants}
D(\widehat\rho,\Gamma)
\preceq\Psi_{\mathrm{alg}}(B),
\qquad
|\Gamma|\le K_{\mathrm{cut}}(B);
\end{equation}

\item a finite set \(\mathfrak J_{\mathrm{con}}\), fixed independently
of \(n\) and \(B\), of represented constructor-source predicates
\(\mathsf W_c\), where
\(\mathsf W_c(E_{\widehat\rho,\Gamma,\ell})\) is certified by the
actual opcode of the distinguished vertex \(v_\ell\), or by a
represented structural certificate specifically comparing the
conditional extension
\(\sigma(J_{\Gamma,\ell-1})\subseteq
\sigma(J_{\Gamma,\ell})\) with a class-\(c\) source.  An unrelated
opcode elsewhere in the carrier is not such a witness.

\end{enumerate}

Every contextual extension produced by \(\mathfrak F\) satisfies at
least one predicate \(\mathsf W_c\).  Predicates may overlap, and a
mixed derivation may carry charges from several constructor families.
No priority rule assigns the complete derivation to one family. Under
the \hyperref[thm-five-family-abs-provenance]{five-family provenance
normal form}, evaluator instructions
carry one or more predicates in
\(\mathfrak J_{\Abs}^5\). Boundary and terminal interfaces
retain their state-dependent arguments in the frontier and their
readout records in the downstream carrier. All other opcodes
in the retained upstream and downstream carrier remain in the cost
ledger but do not acquire the extension's target charge merely by
co-occurrence.

A predicate \(\mathsf W_c\) is not auxiliary metadata.  It is a
family-uniform decidable relation on the encoded distinguished
extension.  Reading its authenticated opcode is syntactic parsing.  If
membership instead uses a semantic structural certificate, the
certificate description, its verifier, and the cost of verification
occur in the ledger \eqref{eq-affordable-frontier-cost-ledger}.  If the
claimed comparison cannot be verified from those charged fields within
the carrier budget, \(\mathsf W_c\) is false.

For \(c\in\mathfrak J_{\mathrm{con}}\), define the aggregate contextual charge envelope

\begin{equation}
\label{eq-aggregate-contextual-charge-envelope}
\mathcal C_{I,\mathfrak F}^{(c)}(D)
:=
\sup_{\substack{\widehat\rho\in\mathsf{NFQ}_{\mathfrak I_n}\\
A_{q_{\widehat\rho,I}}>0,\ D(\widehat\rho,\Gamma)\preceq D\\
\widehat\rho^\Gamma=\mathfrak F(\widehat\rho)}}
\sum_{\substack{1\le\ell\le|\Gamma|\\
\mathsf W_c(E_{\widehat\rho,\Gamma,\ell})}}
\operatorname{ch}_I(E_{\widehat\rho,\Gamma,\ell}),
\qquad \sup\varnothing:=0,
\end{equation}

and the individual-extension envelope

\begin{equation}
\label{eq-individual-contextual-charge-envelope}
\mathfrak c_{I,\mathfrak F}^{(c)}(D)
:=
\sup_{\substack{\widehat\rho\in\mathsf{NFQ}_{\mathfrak I_n},\
A_{q_{\widehat\rho,I}}>0\\
D(\widehat\rho,\Gamma)\preceq D,
\ \widehat\rho^\Gamma=\mathfrak F(\widehat\rho)\\
1\le\ell\le|\Gamma|,
\ \mathsf W_c(E_{\widehat\rho,\Gamma,\ell})}}
\operatorname{ch}_I(E_{\widehat\rho,\Gamma,\ell}),
\qquad \sup\varnothing:=0.
\end{equation}

\end{definition}

\begin{remark}[Structural witnesses do not assign charge]
\label{rem-structural-witnesses-do-not-assign-charge}

The quantity
\(\operatorname{ch}_I(E_{\widehat\rho,\Gamma,\ell})\) is fixed by the
nested affordable algebras and by the terminal utility of the complete
quotient before any witness relation \(\mathsf W_c\) is evaluated.  A
relation \(\mathsf W_c\) only certifies the represented constructor
provenance of an already charged conditional extension so that its
charge can be included in an envelope.  It neither creates, reallocates,
nor discounts that charge.  Removing all class names leaves the exact
charge-conservation identity unchanged; an overlapping witness cover
only repeats nonnegative summands in an upper bound.  In every case the
affordability argument uses the complete carrier ledger, not the cost of
the witness name or of the distinguished opcode alone.

\end{remark}

The charging theorem identifies terminal quotient visibility with its affordable contextual source charges.
\subsection{Exact structural-charge identity}
\label{sssec-exact-structural-charge}
\label{subsec:exact-structural-charge}

\begin{theorem}[Exact affordable-algebra structural charging]
\label{thm-exact-affordable-algebra-structural-charging}

Let \(\widehat\rho^{\Gamma}\) be an affordable
conditional-algebra frontier. Then

\begin{align}
\sum_{\ell=1}^{|\Gamma|}
\operatorname{ch}_I(E_{\widehat\rho,\Gamma,\ell})
&=
V_I(q_{\widehat\rho})
+A_{q_{\widehat\rho,I}}
\frac{
 \|(P_{J_{\Gamma,|\Gamma|,I}}-P_{q_{\widehat\rho,I}})y_I\|_2^2}
 {\|y_I\|_2^2},
\label{eq-exact-affordable-algebra-charge-conservation}\\
V_I(q_{\widehat\rho})
&\le
\sum_{\ell=1}^{|\Gamma|}
\operatorname{ch}_I(E_{\widehat\rho,\Gamma,\ell}).
\label{eq-affordable-algebra-charge-domination}
\end{align}

For every positive-utility exact quotient in the saturated profile,
terminal compatibility gives
\(P_{q_{\widehat\rho,I}}y_I=y_I\). In that case the remainder in
\eqref{eq-exact-affordable-algebra-charge-conservation} is zero and

\begin{equation}
\label{eq-exact-qualified-structural-charge}
V_I(q_{\widehat\rho})
=
\sum_{\ell=1}^{|\Gamma|}
\operatorname{ch}_I(E_{\widehat\rho,\Gamma,\ell}).
\end{equation}

Under \Cref{def-finite-affordable-constructor-charge-cover},
the complete exact-quotient profile satisfies

\begin{align}
m_I^{\mathrm{sat}}(B)
&\le
\sum_{c\in\mathfrak J_{\mathrm{con}}}
\mathcal C_{I,\mathfrak F}^{(c)}
   (\Psi_{\mathrm{alg}}(B))
\label{eq-finite-aggregate-constructor-charge-envelope}\\
&\le
|\mathfrak J_{\mathrm{con}}|
\max_{c\in\mathfrak J_{\mathrm{con}}}
\mathcal C_{I,\mathfrak F}^{(c)}
   (\Psi_{\mathrm{alg}}(B)),
\notag
\end{align}

and also

\begin{equation}
\label{eq-finite-individual-constructor-charge-envelope}
m_I^{\mathrm{sat}}(B)
\le
K_{\mathrm{cut}}(B)
\max_{c\in\mathfrak J_{\mathrm{con}}}
\mathfrak c_{I,\mathfrak F}^{(c)}
   (\Psi_{\mathrm{alg}}(B)).
\end{equation}

Consequently, negligible aggregate constructor-charge envelopes at
every budget in the transfer class imply negligible saturated
exact-quotient visibility. The same conclusion follows from negligible
individual-extension envelopes when \(K_{\mathrm{cut}}\) belongs to the
transfer class.

\end{theorem}

\begin{proof}

The represented factorization
\eqref{eq-affordable-frontier-factorization} gives
\(\sigma(q_{\widehat\rho,I})\subseteq
\sigma(J_{\Gamma,|\Gamma|,I})\).  Projection Pythagoras therefore gives

\[
\mathsf p_I(J_{\Gamma,|\Gamma|,I})
=
\mathsf p_I(q_{\widehat\rho,I})
+
\frac{
 \|(P_{J_{\Gamma,|\Gamma|,I}}-P_{q_{\widehat\rho,I}})y_I\|_2^2}
 {\|y_I\|_2^2}.
\]

The represented join algebras form a nested chain beginning with the
constants.  Successive conditional-expectation differences are
orthogonal, so

\[
\mathsf p_I(J_{\Gamma,|\Gamma|,I})
=
\sum_{\ell=1}^{|\Gamma|}
\frac{
 \|(P_{J_{\Gamma,\ell,I}}
    -P_{J_{\Gamma,\ell-1,I}})y_I\|_2^2}
 {\|y_I\|_2^2}.
\]

Multiplication by the retained terminal utility factor proves
\eqref{eq-exact-affordable-algebra-charge-conservation} and
\eqref{eq-affordable-algebra-charge-domination}.  If the quotient is
terminal compatible, then \(y_I\) is quotient measurable. Hence
\(P_qy_I=y_I\), and the factorization also gives
\(P_{J_{\Gamma,|\Gamma|}}y_I=y_I\).  The remainder vanishes, proving
\eqref{eq-exact-qualified-structural-charge}.

Fix a positive-utility record of cost at most \(B\) and use the
frontier selected by \(\mathfrak F\).  Every extension is covered by at
least one actual constructor-source predicate.  Nonnegativity gives

\[
V_I(q_{\widehat\rho})
\le
\sum_{c\in\mathfrak J_{\mathrm{con}}}
\sum_{\ell:\,\mathsf W_c(E_{\widehat\rho,\Gamma,\ell})}
\operatorname{ch}_I(E_{\widehat\rho,\Gamma,\ell}).
\]

The carrier has cost at most \(\Psi_{\mathrm{alg}}(B)\), so each inner
sum is bounded by its aggregate envelope.  Taking the supremum over
\(\widehat\rho\) proves
\eqref{eq-finite-aggregate-constructor-charge-envelope}.  Alternatively,
each of at most \(K_{\mathrm{cut}}(B)\) extensions is bounded by the
largest individual-extension envelope, which proves
\eqref{eq-finite-individual-constructor-charge-envelope}.  Closure of
the transfer class under finite sums, maxima, and multiplication by
\(K_{\mathrm{cut}}\) proves the final assertions.

\end{proof}

\begin{example}[An exact eight-state affordable algebra]
\label{ex-eight-state-affordable-algebra}

We give a complete finite calculation of the objects in
\cref{def:affordable-frontier-contextual-charge,thm-exact-affordable-algebra-structural-charging}.
The carrier is

\[
\Omega=\{0,1\}^3,
\qquad x=(c,u,v),
\qquad \mu_\Omega=\operatorname{Unif}(\Omega),
\]

with target and failure sectors

\[
T=\{111\},
\qquad
F=\{100,101,110\}.
\]

Thus the states with \(c=1\) are terminal; \(111\) is success and the
other three are failure.  The public state observables are the three bit
reads \(c,u,v\), and the terminal verifier is
\(t(c,u,v)=cuv=\mathbf 1_T(c,u,v)\).  The legal initialization is

\[
\mu_0=\operatorname{Unif}\{000,001,010\}.
\]

The single toy presentation declares exactly three admissible one-step
transition records, indexed by \(p\in\{0,\tfrac12,1\}\).  They are

\begin{equation}
\label{eq-eight-state-kernel}
\begin{array}{c|cccc}
x&000&001&010&011\\ \hline
K_p(x,\cdot)&\delta_{100}&\delta_{101}
&p\delta_{111}+(1-p)\delta_{110}&\delta_{110},
\end{array}
\qquad
K_p(x,\cdot)=\delta_x\quad(c=1).
\end{equation}

The finite code for \(p\) is part of this transition record.  Every
terminal state reachable from the support of \(\mu_0\) is reached after
exactly one transition.

\medskip\noindent\textbf{Grammar, budget, and complete slice.}
Fix the following finite constructor grammar.  It may read any of
\(c,u,v\), form ordered tuples of distinct coordinate reads by binary joins,
apply the downstream verifier \(t\), and attach one of the three records
\(K_p\) together with the displayed initialization, reference law, and
discount normalization.  The verifier is a downstream terminal operation:
the grammar does not feed \(t\) back into an upstream tuple constructor.
The declared scalar cost ledger for the complete three-coordinate quotient
record is

\begin{center}
\begin{tabular}{@{}lr@{}}
\toprule
charged field & cost\\
\midrule
primitive reads \(c,u,v\) & \(3\)\\
two nontrivial tuple joins & \(2\)\\
target and terminal readout & \(1\)\\
kernel/clock, initialization, reference law, and discount record & \(2\)\\
represented fibers and exact-intertwining verification & \(1\)\\
\midrule
complete carrier cost \(B_0\) & \(9\)\\
\bottomrule
\end{tabular}
\end{center}

These are declared representation costs, not operation counts inferred
after the calculation.  In particular, the finite \(p\)-code is included in
the transition row, as is the one-step clock needed for the displayed hitting
functional.  Every contextual extension below retains the complete carrier
and therefore has cost \(D=9\), rather than the marginal unit cost of the
newly distinguished read.  At budget \(B_0\), the grammar admits no iteration,
additional clock extension, feedback of the verifier into a tuple, or other
quotient constructor; each omitted opcode is assigned declared cost greater
than \(B_0\).

For each of the three transition records, the affordable quotient-output
slice at \(B_0=9\), modulo equality of the induced partitions, consists of
the constant map, the verifier partition, the coordinate reads and their
proper coordinate tuples, and the full tuple

\[
q=(c,u,v):\Omega\longrightarrow\{0,1\}^3.
\]

This list exhausts the affordable derivation slice.  Every derivation attaches
one of the three \(K_p\) records to a constant, a verifier output, a coordinate
read, or an ordered binary join of distinct coordinate reads.  Different
orders for the same coordinate subset induce the same partition and hence
the same projection and terminal-compatibility test.  The constant fails the
nontriviality condition.  The verifier partition has zero qualified utility:
although it separates \(T\), its zero-cell combines \(F\) with nonterminal
states and hence does not represent the failure sector.  Every proper
coordinate tuple also fails terminal compatibility.  Indeed, for the three
two-coordinate tuples, respectively,

\[
(c,u):110\sim111,
\qquad
(c,v):101\sim111,
\qquad
(u,v):011\sim111,
\]

and deleting another coordinate cannot repair the failure.  The full tuple
\(q\) is the identity partition, so its fibers are represented and both
terminal sectors are unions of fibers.  If \(\bar K_p\) denotes the induced
quotient kernel, then \(K_pU_q=U_q\bar K_p\) holds exactly.  The quotient
strictly refines the verifier partition.  Write \(q_p\) for the complete
quotient record obtained by attaching \(K_p\) to this map.  Consequently
\(q_p\) is the sole potentially positive-utility quotient record for each
\(p\), with positive utility at \(p=\tfrac12,1\); at \(p=0\) the same exact
quotient map has utility zero.

\medskip\noindent\textbf{Backward cut and conditional algebras.}
Choose the ordered backward cut

\[
\Gamma=(c,u,v),
\qquad
J_0=*,\quad J_1=c,\quad J_2=(c,u),\quad
J_3=(c,u,v)=q.
\]

This cut separates every state-dependent path from the public reads to the
terminal quotient.  The kernel, normalization, terminal cells, fiber records,
and exact-intertwining check remain in the downstream carrier.  In the state
order \(000,001,010,011,100,101,110,111\), the centered target contrast and
its successive conditional expectations are

\begin{align}
y&=\mathbf 1_T-\mu_\Omega(T)
=\frac18(-1,-1,-1,-1,-1,-1,-1,7),
\label{eq-eight-state-projections-a}\\
P_{J_1}y&=\frac18(-1,-1,-1,-1,1,1,1,1),\notag\\
P_{J_2}y&=\frac18(-1,-1,-1,-1,-1,-1,3,3),\notag\\
P_{J_3}y&=y.\notag
\end{align}

With the \(L^2(\mu_\Omega)\) norm,

\begin{equation}
\label{eq-eight-state-projection-norms}
\|y\|_2^2=\frac7{64},
\qquad
\|P_{J_1}y\|_2^2=\frac1{64},
\qquad
\|P_{J_2}y\|_2^2=\frac3{64},
\end{equation}

and orthogonality of the conditional increments gives

\begin{equation}
\label{eq-eight-state-increment-norms}
\bigl\|(P_{J_1}-P_{J_0})y\bigr\|_2^2=\frac1{64},
\quad
\bigl\|(P_{J_2}-P_{J_1})y\bigr\|_2^2=\frac2{64},
\quad
\bigl\|(P_{J_3}-P_{J_2})y\bigr\|_2^2=\frac4{64}.
\end{equation}

The normalized projection increments are therefore
\(1/7,2/7,4/7\), whose sum is one.

\medskip\noindent\textbf{Utility, charges, and the rule-uniform bound.}
For an arbitrary admissible \(K_p\), the one-step success probability is

\begin{equation}
\label{eq-eight-state-success}
\alpha_p=\Pr_{\mu_0}\{\tau_T\le1\}=\frac p3.
\end{equation}

Take \(\lambda=1\).  Since success can occur only on the first transition,
the uniformized target functional in
\eqref{eq:first-hit-uniformization} is
\(A_T^q=\alpha_p/(1+\lambda)=p/6\).  The quotient uses the canonical
pushforward comparison, so \(C_q=1\), and hence

\[
A_q=\min\{1,\lambda C_q^{-1}A_T^q\}=\frac p6,
\qquad
V(q_p)=A_q\frac{\|P_qy\|_2^2}{\|y\|_2^2}=\frac p6.
\]

Writing \(E_\ell\) for the complete contextual carrier with
\(J_{\ell-1}\subset J_\ell\) distinguished, the three charges are exactly

\begin{equation}
\label{eq-eight-state-contextual-charges}
\begin{aligned}
\operatorname{ch}(E_1)&=\frac p6\frac17=\frac p{42},\\
\operatorname{ch}(E_2)&=\frac p6\frac27=\frac p{21},\\
\operatorname{ch}(E_3)&=\frac p6\frac47=\frac{2p}{21}.
\end{aligned}
\end{equation}

Thus the theorem's remainder vanishes and its charge-conservation identity
is visible numerically:

\begin{equation}
\label{eq-eight-state-charge-conservation}
\sum_{\ell=1}^3\operatorname{ch}(E_\ell)
=\frac p{42}+\frac p{21}+\frac{2p}{21}
=\frac p6
=V(q_p).
\end{equation}

Let \(m^{\mathrm{sat}}\) denote the saturated profile of this single toy
presentation.  Its supremum ranges over the three admissible \(K_p\) records.
Because the exhaustive affordable slice has no other positive-utility
quotient,

\begin{equation}
\label{eq-eight-state-saturated-profile}
m^{\mathrm{sat}}(B_0)
=\max_{p\in\{0,\frac12,1\}}V(q_p)
=\frac16.
\end{equation}

Combining \eqref{eq-eight-state-success} and
\eqref{eq-eight-state-charge-conservation} yields, simultaneously for every
admissible transition rule \(K_p\), the exact class bound

\begin{equation}
\label{eq-eight-state-rule-bound}
\Pr_{\mu_0}\{\tau_T\le1\}
=2V(q_p)
=2\sum_{\ell=1}^3\operatorname{ch}(E_\ell)
\le 2m^{\mathrm{sat}}(B_0)
=\frac13,
\end{equation}

with equality for \(K_1\).  By comparison, applying the generic universal
accountability estimate \eqref{eq-pointwise-algorithm-profile-bound} with
\(H=1\) and a declared first-hit compilation ceiling \(B_{\mathcal A}=B_0\)
gives the valid but coarser bound \(\alpha_p\le e/6\).  The sharper
\(1/3\) above comes from evaluating the entire finite affordable algebra and
its exact conditional charges, rather than replacing the one-step discount
by the universal \(e^{-1}\) lower bound.

\end{example}

\Cref{fig-contextual-charge} records the charge identity in
\Cref{thm-exact-affordable-algebra-structural-charging}.

\begin{figure}[htbp]
\centering
\begin{tikzpicture}[fw diagram]
  \node[fw box,minimum width=1.3cm] (j0) at (-5.2,0.6) {\(J_0\)};
  \node[fw box,minimum width=1.3cm] (j1) at (-3.2,0.6) {\(J_1\)};
  \node[font=\large] (dots) at (-1.35,0.6) {\(\cdots\)};
  \node[fw box,minimum width=1.3cm] (jk) at (0.5,0.6) {\(J_k\)};
  \node[fw provenance,minimum width=1.8cm] (q) at (3.0,0.6) {\(q\)};
  \draw[fw arrow] (j0) -- node[above,font=\scriptsize] {\(\operatorname{ch}_1\)} (j1);
  \draw[fw arrow] (j1) -- (dots);
  \draw[fw arrow] (dots) -- node[above,font=\scriptsize] {\(\operatorname{ch}_k\)} (jk);
  \draw[fw arrow] (jk) -- node[above,font=\scriptsize] {terminal} (q);
  \node[fw terminal,minimum width=5.1cm] (sum) at (-1.35,-1.0)
    {\(\displaystyle \sum_{\ell=1}^{k}\operatorname{ch}_I(E_\ell)=V_I(q)\)};
  \draw[fw arrow] (dots.south) -- (sum.north);
\end{tikzpicture}
\caption{For a terminal-compatible exact quotient, successive
contextual charges telescope to its visibility by
\cref{thm-exact-affordable-algebra-structural-charging}.}
\label{fig-contextual-charge}
\end{figure}

\begin{remark}[Order dependence and invariant total charge]
\label{rem-contextual-charge-order-dependence}

The individual quantities
\(\operatorname{ch}_I(E_{\widehat\rho,\Gamma,\ell})\) depend on the
family-uniform frontier and on its encoded order, just as conditional
information increments do.  No individual increment is claimed to be an
intrinsic contribution of \(v_\ell\) in isolation.  What is invariant
for a terminal-compatible exact quotient is their sum
\eqref{eq-exact-qualified-structural-charge}.  All envelope statements
therefore fix \(\mathfrak F\) and its order before the instance is
evaluated.

\end{remark}

\begin{corollary}[Contextual structural-charge transfer]
\label{cor-contextual-structural-charge-transfer}

Assume a finite affordable constructor-charge cover at budget \(B\).
Suppose a family-uniform rule assigns every contextual extension
\(E_{\widehat\rho,\Gamma,\ell}\) a signature
\(c_\ell\in\mathfrak C_{\mathrm{src}}\), a dynamically qualified
source record
\(R_\ell\in
\mathcal S_I^{(c_\ell)}(\Psi_{\mathrm{tr}}(B))\), and a represented
comparison such that the complete transfer carrier

\[
(\widehat\rho^\Gamma,R_\ell,
  \operatorname{Cmp}_{E_\ell,R_\ell})
\]

has cost at most \(\Psi_{\mathrm{tr}}(B)\) uniformly over the task
slice and

\begin{equation}
\label{eq-contextual-structural-charge-transfer}
\operatorname{ch}_I(E_{\widehat\rho,\Gamma,\ell})
\le
K_B'V_I(R_\ell)+\delta_B'
\end{equation}

for every instance.  Then

\begin{equation}
\label{eq-contextual-structural-charge-source-envelope}
m_I^{\mathrm{sat}}(B)
\le
K_{\mathrm{cut}}(B)K_B'
\max_{c\in\mathfrak C_{\mathrm{src}}}
s_I^{(c)}(\Psi_{\mathrm{tr}}(B))
+K_{\mathrm{cut}}(B)\delta_B'.
\end{equation}

Thus an incremental transfer preserves negligibility when
\(K_{\mathrm{cut}}\), \(K_B'\), and \(\Psi_{\mathrm{tr}}\) lie in the
declared transfer class and \(\delta_B'\) is negligible.

\end{corollary}

\begin{proof}

Use \eqref{eq-affordable-algebra-charge-domination}, apply
\eqref{eq-contextual-structural-charge-transfer} to each of at most
\(K_{\mathrm{cut}}(B)\) extensions, and take the supremum over
actual-cost records in
\(\mathsf{NFQ}_{\mathfrak I_n,\le B}\).  This proves
\eqref{eq-contextual-structural-charge-source-envelope}.  The final
claim follows from transfer-class closure.

\end{proof}

\subsection{Canonical constructor covers}
\label{sssec-constructor-covers}
\label{subsec:canonical-constructor-covers}

\begin{corollary}[Canonical leaf charge cover under succinct interfaces]
\label{cor-universal-syntax-affordable-charge-cover}

Fix a presented finite task family with the derivation coding of
\Cref{def-budgeted-derivability-closure}.  Assume its declared
resource envelope supplies class-preserving majorants
\(L_{\mathrm{tr}}(B)\) and \(J_{\mathrm{tuple}}(B,k)\) with the
following two properties: an expanded record of actual witness cost at
most \(B\) has at most \(L_{\mathrm{tr}}(B)\) encoded vertices, and a
nested represented tuple of at most \(k\) vertices from that record has
complete description, materialization, storage, evaluation,
joint-fiber-membership, and finite Boolean fiber-operation cost at most
\(J_{\mathrm{tuple}}(B,k)\).  These are the succinct-interface
hypotheses.

Let \(\pi_{\mathrm{op}}\) be a fixed finite schema map on the
normal-form opcodes declared by the presentation; opcode parameters and
public coefficients are data, not new schemas.  For an exact-quotient
record \(\widehat\rho\), list in encoded order all state-dependent
leaves of its terminal backward slice and call the list
\(\Gamma_{\mathrm{syn}}(\widehat\rho)\).  Public constants and
state-independent parameters remain coefficients of the downstream
contraction.  Define

\[
\mathsf W_\sigma(E_{\widehat\rho,\Gamma,\ell})
\quad\Longleftrightarrow\quad
\pi_{\mathrm{op}}(\operatorname{op}(v_\ell))=\sigma.
\]

Thus the predicate is anchored at the distinguished algebra extension.
An internal or terminal opcode on the downstream path remains charged
in the carrier ledger but does not receive the leaf's target charge.
Then
\(\widehat\rho^{\Gamma_{\mathrm{syn}}}
=\mathfrak F_{\mathrm{syn}}(\widehat\rho)\) and the anchored predicates
satisfy \Cref{def-finite-affordable-constructor-charge-cover}, with

\[
|\Gamma_{\mathrm{syn}}(\widehat\rho)|
\le K_{\mathrm{syn}}(B):=L_{\mathrm{tr}}(B),
\qquad
D(\widehat\rho,\Gamma_{\mathrm{syn}})
\preceq\Psi_{\mathrm{syn}}(B)
:=B\oplus
J_{\mathrm{tuple}}(B,L_{\mathrm{tr}}(B))
\]

for every record of cost at most \(B\). The standard finite
operational-history configuration model satisfies these hypotheses when the charged trace
dominates its encoded length and tuples are represented by references
to their charged component records; in that model both majorants are
polynomial.  Consequently,

\begin{equation}
\label{eq-universal-syntax-affordable-charge-cover}
m_I^{\mathrm{sat}}(B)
\le
\sum_{\sigma\in\operatorname{im}(\pi_{\mathrm{op}})}
\mathcal C_{I,\mathfrak F_{\mathrm{syn}}}^{(\sigma)}
  (\Psi_{\mathrm{syn}}(B)).
\end{equation}

Any application-specific frontier may instead select outputs of actual
derived structural constructors, together with the remaining leaves
needed for factorization.  A coarser semantic list is valid only through
anchored represented certificates for those selected extensions.  No
assignment of a complete record to one terminal class is involved.

\end{corollary}

\begin{proof}

Every state-dependent input to a vertex of an expanded record is an
earlier vertex of the same finite DAG.  Evaluation of the retained DAG
in topological order therefore gives a represented contraction of the
terminal quotient through the joint observable of its state-dependent
leaves, proving \eqref{eq-affordable-frontier-factorization}.  Every
distinguished leaf has exactly one opcode schema under
\(\pi_{\mathrm{op}}\), so the anchored predicates cover every extension.
A primitive state-dependent terminal quotient is itself such a leaf and
is admitted by the
\hyperref[def-affordable-conditional-algebra-frontier]{affordable
conditional-algebra frontier}; its actual schema supplies its source
designation.

The first succinct-interface hypothesis bounds the number of leaves.
The second bounds the added nested tuple, including its representation
and evaluation; retaining the original record contributes at most
\(B\).  These are precisely the displayed majorants.  The complete DAG
remains in every carrier, so every upstream and downstream operation is
still present in \eqref{eq-affordable-frontier-cost-ledger}.  \Cref{thm-exact-affordable-algebra-structural-charging} now gives
\eqref{eq-universal-syntax-affordable-charge-cover}.

\end{proof}

\begin{corollary}[Exhaustive five-family envelope for the saturated
exact-quotient profile]
\label{cor-five-family-abs-envelope}

Apply the five-family expansion of the
\hyperref[thm-five-family-abs-provenance]{universal five-family Abs
provenance normal form} to every
\(\widehat\rho\in\mathsf{NFQ}_{\mathfrak I_n,\le B}\). In encoded
topological order, let \(\Gamma_5\) contain the state-dependent typed
vertices in the backward slice of the terminal quotient map after pure
interfaces have exposed their arguments. Include the terminal quotient
vertex. Let \(\mathfrak F_5\) be this family-uniform frontier rule and
put

\[
\mathsf W_c(E_{\widehat\rho,\Gamma_5,\ell})
\quad\Longleftrightarrow\quad
c\in\operatorname{Prov}_5(v_\ell),
\qquad
c\in\mathfrak J_{\Abs}^5.
\]

Then \(\mathfrak F_5\) is a finite affordable constructor-charge cover.
There are class-preserving majorants \(\Psi_5^{\mathrm{alg}}\) and
\(K_5\), fixed by the encoded operational model, such that

\begin{align}
m_I^{\mathrm{sat}}(B)
&\le
\sum_{c\in\mathfrak J_{\Abs}^5}
\mathcal C_{I,\mathfrak F_5}^{(c)}
  \bigl(\Psi_5^{\mathrm{alg}}(B)\bigr),
\label{eq-five-family-aggregate-abs-envelope}\\
m_I^{\mathrm{sat}}(B)
&\le
K_5(B)
\max_{c\in\mathfrak J_{\Abs}^5}
\mathfrak c_{I,\mathfrak F_5}^{(c)}
  \bigl(\Psi_5^{\mathrm{alg}}(B)\bigr).
\label{eq-five-family-individual-abs-envelope}
\end{align}

For every positive-utility terminal-compatible exact quotient, its
visibility is exactly the sum of the contextual charges on this
frontier, and every summand belongs to at least one of the five
families. Consequently, negligibility of the five aggregate envelopes,
or of the five individual envelopes with \(K_5\) in the transfer class,
implies negligibility of the complete saturated exact-quotient accounting
profile.
Restricting the same frontier to temporally admissible records gives the
source profile of \cref{def-source-accounting-profile-separation}.  Exact
first-hit quotients belong to the accounting frontier and enter the source
frontier only through an independent pre-hit representation.

\end{corollary}

\begin{proof}

The five-family normal form gives a finite typed backward slice with
class-preserving expansion cost. The encoded topological order is
family uniform. Since the terminal quotient is included,
\(q_{\widehat\rho}\) factors through the represented join of
\(\Gamma_5\). Every selected state-dependent vertex has at least one
authenticated five-family tag, and interface vertices have exposed
their tagged arguments. The expanded-record length and represented
tuple-interface bounds give \(K_5\) and
\(\Psi_5^{\mathrm{alg}}\).

Apply the
\hyperref[thm-exact-affordable-algebra-structural-charging]{exact
affordable-algebra structural-charging theorem}
with \(\mathfrak J_{\mathrm{con}}=
\mathfrak J_{\Abs}^5\). Equations
\eqref{eq-five-family-aggregate-abs-envelope} and
\eqref{eq-five-family-individual-abs-envelope} are its aggregate and
individual bounds. Terminal compatibility gives the exact charge
identity. The
\hyperref[lem:exact-affordable-first-hit]{exact first-hit lemma} supplies a
normal-form accounting quotient for every represented finite-horizon
computation.  Temporal source admissibility is applied separately as in
\cref{cor-first-hit-accounting-not-source}.

\end{proof}

\section{No-Creation and Filtration Provenance}
\label{sssec-filtration-provenance}
\label{sec:no-creation-filtration-provenance}

\begin{lemma}[Derived evaluator constructors carry zero contextual
charge]
\label{lem-derived-constructor-zero-contextual-charge}

Let \(E_{\widehat\rho,\Gamma_5,\ell}\) be a contextual extension in the
\hyperref[cor-five-family-abs-envelope]{five-family frontier}, and
suppose its distinguished vertex \(v_\ell\) is the output of a
represented \(\mathsf{Mult}\) or \(\mathsf{Ord}\) constructor, or a
\(\mathsf{Filt}\)-tagged deterministic constructor whose
state-dependent arguments occur earlier in the frontier. Then

\[
\operatorname{ch}_I(E_{\widehat\rho,\Gamma_5,\ell})=0.
\]

This applies to conjunction, products, composition gates, polynomial
and ring operations, equality tests, represented fiber formation,
comparisons, branching and selection, weights, metric and quadratic
comparisons, and ordered reductions. The same conclusion applies to a
\(\mathsf{Filt}\)-tagged clock, stored copy, memory update, recurrence
step, or history coordinate that is a represented deterministic
function of earlier state-dependent arguments.

\end{lemma}

\begin{proof}

Every state-dependent argument of \(v_\ell\) is an earlier vertex in
the encoded topological order defining \(\Gamma_5\); state-independent
public data remain coefficients. Hence the observable at \(v_\ell\)
is a represented postprocessing of \(J_{\Gamma_5,\ell-1,I}\). The
\hyperref[thm-affordable-composition-no-creation]{affordable composition
and target-projection theorem} gives

\[
\sigma(v_\ell)
\subseteq
\sigma(J_{\Gamma_5,\ell-1,I}).
\]

Therefore adjoining \(v_\ell\) does not change the represented joint
algebra:

\[
\sigma(J_{\Gamma_5,\ell,I})
=
\sigma(J_{\Gamma_5,\ell-1,I}).
\]

The two conditional-expectation projections in
\eqref{eq-contextual-source-extension-charge} are equal, so the
contextual charge is zero.

\end{proof}

\begin{lemma}[Filtration provenance charging]
\label{thm-filtration-provenance-charging}

Let \(E_{\widehat\rho,\Gamma_5,\ell}\) be a contextual extension in the
\hyperref[cor-five-family-abs-envelope]{five-family frontier} whose
distinguished vertex has \(\mathsf{Filt}\) provenance. Retain the
complete carrier, including the generating update, verifier, lifted
interfaces, terminal records, and all earlier state-dependent
arguments. Then filtration provenance supplies no independent
generative source. More precisely:

\begin{enumerate}

\item If the distinguished observable is a represented deterministic
function of earlier frontier coordinates, then

\[
\operatorname{ch}_I(E_{\widehat\rho,\Gamma_5,\ell})=0.
\]

\item If it is an auxiliary coordinate of a
\hyperref[def-valid-lift]{valid Lift}, the represented base projection
occurs earlier in the retained interface and the lifted target contrast
is its pullback. The auxiliary coordinate therefore adds zero
contextual target charge; the target contribution remains attached to
the represented base record.

\item For a budget-generated outcome or terminal transcript, apply the
\hyperref[thm-channel-exhaustiveness]{component-exhaustiveness
projection} to its retained generating update and terminal edge gate.
Normalize its represented evaluator into the five-family backward slice of
\cref{thm-five-family-abs-provenance}.  Independent outcome randomness is a
valid-Lift auxiliary coordinate and adds zero target charge by the preceding
case.  The update, verifier, stored flags, terminal map, and every subsequent
first-hit postprocessing are represented deterministic functions of earlier
frontier coordinates and add zero contextual charge by the first case.  Hence
the exact charge identity of \cref{cor-five-family-abs-envelope} assigns every
nonzero contextual charge to an earlier authenticated vertex of the completed
derivation.  If that vertex is invoked prospectively, the temporal-source and
represented-transfer conditions of
\cref{def-pre-hit-temporally-admissible-source,%
cor-contextual-structural-charge-transfer} classify and compare its source
record.  Without those additional conditions, the charge remains completed
accounting.  The post-saturation \(\Res\) term contains no represented
structural source by
\cref{prop-family-uniform-residual-reclassification,thm-no-uniform-residual-correlation}.

\end{enumerate}

Thus a clock, memory record, stored flag, outcome history, recurrence,
or first-hit record can retain \(\mathsf{Filt}\) carrier provenance but
cannot receive an independent source credit from that tag.

\end{lemma}

\begin{proof}

In the first case, all state-dependent arguments of the distinguished
vertex occur in \(J_{\Gamma_5,\ell-1,I}\). The
\hyperref[lem-derived-constructor-zero-contextual-charge]{derived
constructor zero-charge lemma} gives the asserted equality.

For a valid Lift \(\pi:\widetilde X\to X\), the target is
\(\widetilde T=\pi^{-1}(T)\) and target-fraction preservation gives

\[
\mathbf 1_{\widetilde T}
-\widetilde\mu(\widetilde T)
=
\bigl(\mathbf 1_T-\mu(T)\bigr)\circ\pi.
\]

The represented Lift interface exposes \(\pi\) and its base
state-dependent arguments before any auxiliary coordinate. Hence the
lifted target contrast is already measurable with respect to the
preceding joint algebra. Adjoining an auxiliary Lift coordinate leaves
its conditional expectation unchanged, so its contextual charge is
zero. This is the conditional-algebra form of the
\hyperref[lem-exact-lift-identity]{exact Lift identity}.

It remains to consider an operational outcome or terminal record. The
\hyperref[def-budgeted-derivability-closure]{budgeted derivability
closure} retains its initializer, generating update, verifier, outcome
interface, transcript, and terminal map at their accumulated cost.
\Cref{thm-five-family-abs-provenance} expands this complete record into a
charged topological derivation.  Independent random coordinates are valid-Lift
auxiliaries, so the second case gives zero incremental target charge.  After
those coordinates and the earlier state have been exposed, the update,
verifier flag, stored transcript coordinate, and terminal map are deterministic
represented postprocessings.  The first case therefore gives zero incremental
charge at each such vertex.  The same conclusion applies to the first-hit map,
which is a deterministic postprocessing of the stored flags and deadline.

The exact telescoping identity in
\cref{cor-five-family-abs-envelope} leaves every nonzero charge on an earlier
authenticated vertex of the retained derivation.  When the vertex is used
prospectively, the
\hyperref[thm-channel-exhaustiveness]{component-exhaustiveness theorem},
\cref{thm-complete-generic-source-theory}, and the temporal and transfer
conditions in
\cref{def-pre-hit-temporally-admissible-source,%
cor-contextual-structural-charge-transfer} classify its source and supply the
applicable comparison.  Otherwise the telescoping identity remains an
accounting identity.  In either case,
\cref{prop-family-uniform-residual-reclassification,thm-no-uniform-residual-correlation}
exclude a represented source from the post-saturation residual.  Thus no
operation in the filtration chain assigns a new generative source to the
\(\mathsf{Filt}\) tag.  This proves the three alternatives.

\end{proof}

\begin{remark}[A verifier and first-hit quotient do not supply a
selector]
\label{rem-first-hit-no-selector}

Let \(V_I\) be a represented public verifier and let
\(q_{\mathcal A}\) be the
\hyperref[lem:exact-affordable-first-hit]{exact first-hit quotient} of
a represented computation. The quotient map is an affordable function
of the generated terminal flags and deadline, so

\[
\sigma(q_{\mathcal A})
\subseteq
\sigma(s_0,\ldots,s_H,r).
\]

The
\hyperref[thm-affordable-composition-no-creation]{affordable composition
theorem} therefore assigns it no target information beyond that
generated transcript. Represented fibers provide uniform membership
predicates and finite Boolean fiber operations. They do not provide a
map selecting a successful transcript or candidate from a success
fiber. Such a selector is a separate represented derivation and is
charged at its own saturated cost; without one, the quotient cannot
generate a verifier-accepting candidate.

The quotient utility retains the actual success probability
\(\alpha_I\) through
\(V_{q_{\mathcal A}}\ge\alpha_I/(eH)\). Consequently the first-hit
construction records success already produced by the computation and
does not create it. A terminal test using a latent target rather than a
represented public verifier belongs to oracle or advice accounting.

\end{remark}

\begin{remark}[Dynamic transfer from contextual charge]
\label{rem-dynamic-transfer-from-contextual-charge}

The structural charge
\(\operatorname{ch}_I(E_{\widehat\rho,\Gamma,\ell})\) retains the
terminal utility of the complete computation; it does not assert that
the frontier coordinate is independently lumpable.  A comparison with
a smaller dynamically qualified source requires a represented
inequality

\[
\operatorname{ch}_I(E_{\widehat\rho,\Gamma,\ell})
\le
K_c(D)V_I(R_E)+\delta_c(D),
\]

with the full comparison carrier charged, followed by
\Cref{cor-contextual-structural-charge-transfer}. Under the
hypotheses of \Cref{lem:exact-affordable-first-hit}, every
finite-horizon uniform represented computation supplies an exact
quotient record to which the structural charging theorem applies.

\end{remark}

\subsection{Lossless uniformization}
\label{sssec-lossless-uniformization}
\label{subsec:lossless-uniformization}

\begin{lemma}[Lossless charged uniformization]
\label{lem-lossless-charged-uniformization}

Let \(\widetilde L\) be any represented finite Markov or sub-Markov
generator, let \(\zeta>0\) be represented and charged, and write its generator-form data as
\(q(x,y)\) and \(\kappa(x)\).  Choose and charge a represented scalar
\(c\) satisfying

\[
c\ge\zeta,
\qquad
c\ge
\max_x\left\{\sum_{y\ne x}q(x,y)+\kappa(x)\right\}.
\]

Then

\[
P=I+c^{-1}\widetilde L,
\qquad
L=P-I=c^{-1}\widetilde L,
\qquad
\lambda=\zeta/c
\]

defines a represented Markov or sub-Markov kernel \(P\), a normalized
generator \(L=P-I\), and a discount \(0<\lambda\le1\).  For every
absorbing target \(T\) and every initial law \(\mu_0\),

\begin{equation}
\label{eq-lossless-uniformization-arrival}
A_T(\zeta;\widetilde L,\mu_0)
=
A_T(\lambda;L,\mu_0).
\end{equation}

The normalization changes no state, target, quotient fiber, or
off-diagonal transition support.  It may add holding steps, as every
uniformization does, but preserves the target-hitting functional exactly.
Its scalar, rate-table arithmetic, normalized kernel,
discount conversion, and every transported quotient, defect, geometry,
and authority record are included in the saturated cost.
Such a scalar exists analytically on every finite carrier.  It is a
represented datum when the carrier has a represented finite enumeration,
when the generator presentation supplies an exit-rate majorant, or when
the majorant is included explicitly in the certificate.  Its construction
and verification are charged.  This is not a restriction on DTM
representation: for a discrete kernel generator \(\widetilde L=P-I\),
one may take \(c=\max\{1,\zeta\}\), and the canonical algorithmic reading
uses \(\zeta\le1\) and \(c=1\).

\end{lemma}

\begin{proof}

For \(x\ne y\), \(P_{xy}=q(x,y)/c\ge0\), while

\[
P_{xx}
=1-\frac{\sum_{y\ne x}q(x,y)+\kappa(x)}{c}
\ge0,
\qquad
\sum_yP_{xy}=1-\frac{\kappa(x)}c\in[0,1].
\]

Thus \(P\) is sub-Markov, and it is Markov when \(\kappa=0\).  Let
\(\widetilde L_N\) and \(\widetilde k_T\) be the killed operator and
target flux for \(\widetilde L\).  The normalized data are
\(L_N=c^{-1}\widetilde L_N\) and
\(k_T=c^{-1}\widetilde k_T\).  Therefore

\[
(\lambda I-L_N)^{-1}k_T
=
\left(c^{-1}(\zeta I-\widetilde L_N)\right)^{-1}
\,c^{-1}\widetilde k_T
=
(\zeta I-\widetilde L_N)^{-1}\widetilde k_T.
\]

Adding the unchanged initial target atom proves
\eqref{eq-lossless-uniformization-arrival}.  Positive scalar rescaling
also preserves the generator's off-diagonal zero pattern, absorbing sectors, and every exact
intertwining identity.  All other transported numerical data are
recomputed for \(L\) and retained in the charged record.

\end{proof}

\section{The Master Compensated Quotient--Geometry Certificate}
\label{sec:master-compensated-certificate}

The master certificate assembles the exact quotient, geometric, causal, and
defect-compensation records on the common saturated ledger.

The construction proceeds through the certificate data in
\cref{sssec-master-certificate-data}, the saturated profiles in
\cref{sssec-master-profiles}, and the authority-restricted normal form in
\cref{sssec-master-normal-form}.
The profile definition and its authority controls are
\cref{def-four-mode-geom-abs-profile,thm-authority-monotonicity-master-profile,thm-affordable-quotient-geometry-master-normal-form,cor-no-causal-substitution-master-normal-form}.

\subsection{Certificate data and visibility}
\label{sssec-master-certificate-data}
\label{subsec:master-certificate-data}

\begin{definition}[Master compensated quotient--geometry certificate]
\label{def-supported-geom-certificate}
\label{def-master-compensated-quotient-geometry-certificate}

The master certificate consists of the structural record and slot typing in
parts~\textup{(i)--(iv)}, the dynamic qualification and quotient utility of
\cref{def:master-certificate-dynamic-qualification}, and the visibility and
charge rule of \cref{def:master-certificate-visibility}.  The present
definition retains both established public labels for the combined object.

A master compensated quotient--geometry certificate of total saturated
cost at most \(B\) is a tuple

\[
\mathscr C=(\mathfrak N,q,\mathfrak G_q,\mathfrak S,\mathfrak K_q),
\]

\medskip\noindent\textbf{(i) Uniformization and budget.}
Here

\[
\mathfrak N=(\widetilde L,\zeta,c,P,L,\lambda)
\]

is the lossless charged uniformization record of
\Cref{lem-lossless-charged-uniformization}.  Thus every master
record contains a represented normalized kernel form \(L=P-I\), has
\(0<\lambda\le1\), and preserves the original selected process's
discounted target functional by
\eqref{eq-lossless-uniformization-arrival}.  All quotient, geometry,
stability, and causal slots below refer to this same normalized
generator \(L\).  If an authority envelope was declared for
\(\widetilde L\), the record must also transport its actual current by
the fixed factor \(c^{-1}\); normalization cannot select a different
current or control.

Membership at budget \(B\) includes one chosen family-uniform
derivation of the complete tuple with actual cost \(\preceq B\); it is
not defined through an infimum in the resource partial order.

\medskip\noindent\textbf{(ii) Quotient and geometry.}
The map \(q:X\to Q\) is a represented terminal-compatible quotient,
with the identity quotient allowed.  A nonidentity support quotient is
admitted only when the same record satisfies the represented-fiber,
terminal-event, positive quotient-utility, and nontrivial-public-structure
conditions of
\cref{def-paper1-abs-usefulness-gate}.  In the exact-support case its
dynamic condition is exact intertwining.  In a nonzero-defect case that
condition is replaced, and only replaced, by the charged stability and
target-functional comparison in part~\textup{(v)} below.  Thus the support
map is an affordable target-qualified quotient record; a represented fiber
map that fails these existing gates is not admitted merely because a
geometry was written on its image.  The geometry slot
\(\mathfrak G_q\) is either active or null.  An active slot is the
represented tuple

\[
\mathfrak G_q=(\Phi_Q,\Phi_\perp,D_Q,J_Q),
\]

where \(\Phi_Q\) is a Dirichlet geometry on \(Q\), \(D_Q:Q\to\mathbb
R\) is a progress observable, \(J_Q\) is a quotient-edge current, and
\(\Phi_\perp\) is the within-fiber geometry used by a nonzero-defect
stability record.  In an exact active certificate \(\Phi_\perp\) may be
empty.  The potential is target-normalized:

\[
D_Q\ge 0,
\qquad
D_Q=0\quad\text{on the positive quotient terminal sector}.
\]

\medskip\noindent\textbf{(iii) Causal provenance.}
The causal slot \(\mathfrak K_q\) is a provenance record in the sense
of \Cref{def-authorized-caus-use-current}.  For an active
geometry slot it verifies that \(J_Q=J_{\mathfrak K_q}\), records the
selected authority envelope and generator, declares whether \(J_Q\) is
the transported structural-gradient projection of the actual selected
skew current or the auxiliary geometric current of
\((\Phi_Q,D_Q)\), and transports that same current through every
quotient, coordinate change, and lift used by the certificate. In the
first case it is a \(\Caus\) operator component. In the second
it only orients the \(\Geom\) descent comparison and is not
identified with the operator's Caus component unless an explicit
represented equality is included. An active certificate without the
appropriate provenance record has zero alignment factor.

\medskip\noindent\textbf{(iv) Null geometry.}
The null slot, denoted \(\mathfrak G_q=\varnothing\), is
permitted precisely when \(E_q=0\) and \(q\) is an exact quotient
record of \Cref{def-abs-extraction}; the identity quotient is
allowed.  Thus a null slot records pure exact quotient accounting and
does not assert any geometric comparison; in this case
\(\mathfrak K_q=\varnothing\).  An identity quotient with an active
geometry slot instead belongs to the ambient Geom branch.  These two
records are disjoint because the slot is part of the certificate.

\end{definition}

\begin{definition}[Dynamic qualification and utility of a master certificate]
\label{def:master-certificate-dynamic-qualification}

Let \(\mathscr C=(\mathfrak N,q,\mathfrak G_q,
\mathfrak S,\mathfrak K_q)\) have the structural slots of
\cref{def-master-compensated-quotient-geometry-certificate}.

\medskip\noindent\textbf{(v) Dynamic stability.}
The record \(\mathfrak S\) is one of the following dynamic
certificates:

\[
S_q(\lambda)=
\begin{cases}
1,&E_q=0,\\
1-\vartheta_\lambda^T(q),
 &E_q\ne0\text{ and the reversible hypotheses hold with }
   \vartheta_\lambda^T(q)<1,\\
\max\{0,1-\Delta_\lambda(q)\},
 &E_q\ne0\text{ and a general resolvent certificate is used},\\
0,&\text{otherwise}.
\end{cases}
\]

In the second line \(\vartheta_\lambda^T(q)\) is the represented
target-functional relative loss of
\Cref{cor-target-functional-reversible-transfer}; its
Schur-complement bound, target vectors, and quotient-arrival
identification are part of \(\mathfrak S\). In the third line
\(\Delta_\lambda(q)\) abbreviates the represented
target-functional defect
\(\Delta_\lambda^T(q;\ell_T,v_T)\) from
\Cref{def-normalized-resolvent-defect}. The windows,
functional, and input are part of \(\mathfrak S\) and must realize the
same quotient-arrival scalar used in \(\overline A_q\).
A nonzero-defect certificate is dynamically qualified only when
\(S_q(\lambda)>0\) and \(S_q(\lambda)^{-1}\) belongs to the declared
transfer class.  In particular, the null geometry slot cannot qualify a
nonzero defect.  The stability record is evaluated for the full
selected generator.  Its component ledger contains
\(E_q^{\mathsf s}\) and \(E_q^{\mathsf a}\) from
\Cref{def-symmetric-causal-quotient-defects}; the reversible
line is available only when the selected skew operator vanishes.
The discount scale is part of the resource record: a certificate used
at resource bound \(B\) must have \(\lambda^{-1}\) in the corresponding
time-transfer class.  Finally, dynamic qualification requires the
represented target vectors and terminal-flux functional to identify the
actual full-process discounted target functional.  Thus the reversible
and general records both certify

\begin{equation}
\label{eq-master-full-target-transfer}
A_T(\lambda;L,\mu_0)
\ge
S_q(\lambda)A_T^Q(\lambda;L_Q,\bar\mu_0).
\end{equation}

This is the scalar target-functional conclusion of
\Cref{cor-target-functional-reversible-transfer} or
\Cref{cor-defect-caps-quotient-gap}, respectively.  An
operator-norm comparison without the displayed target-functional
identification is not a dynamically qualified nonzero-defect record.

\medskip\noindent\textbf{(vi) Quotient utility.}
For a nonidentity quotient, let

\[
\overline A_q
=
\min\{1,\lambda\overline C_q^{-1}
A_T^Q(\lambda;L_Q,\bar\mu_0)\},
\]

where \(\overline C_q\) contains the represented-fiber, terminal,
normalization, and quotient-level comparison factors, but not the
dynamic defect loss already recorded by \(S_q\). For an identity
quotient with an active geometry slot set
\(\overline A_{\mathrm{id}}=1\).  For an identity quotient with a null
slot set \(\overline A_{\mathrm{id}}=A_{\mathrm{id}}\).  Thus every
exact null-slot quotient, including the identity, satisfies
\(\overline A_q=A_q\).
For a nonidentity exact-support record, the comparison data are the same
data used by \cref{def-paper1-abs-usefulness-gate}; hence
\(\overline C_q=C_q\) and \(\overline A_q=A_q\).  In a nonzero-defect
record, \(\overline A_q\) retains those same nondynamic qualification
factors, while the defect loss occurs exactly once in \(S_q\).

\end{definition}

\begin{definition}[Master-certificate reach, visibility, and charge]
\label{def:master-certificate-visibility}

Let \(\mathscr C\) be a structurally typed and dynamically qualified master
certificate in the sense of
\cref{def-master-compensated-quotient-geometry-certificate,%
def:master-certificate-dynamic-qualification}.

\medskip\noindent\textbf{(vii) Target reach.}
For an active geometry slot, let \(\mathcal M_Q(D_Q)\) be the monotone
target-oriented span on \(Q\).  When
\(\|y\|_2=0\), set every reach factor and \(V_{\mathscr C}\) to zero.
Otherwise, when
\(\operatorname{Var}_{\nu_Q}(D_Q)>0\), define

\[
R_T^q(D_Q)
=
\frac{\|\Pi_{U\mathcal M_Q(D_Q)}y\|^2}{\|y\|^2},
\qquad
\rho_{\Phi_Q}(D_Q)
=
\frac{\mathcal E_{\Phi_Q}(D_Q,D_Q)}
{\operatorname{Var}_{\nu_Q}(D_Q)}.
\]

When \(\operatorname{Var}_{\nu_Q}(D_Q)=0\), set the active-slot reach
factor and \(\rho_{\Phi_Q}(D_Q)\) to zero.  The resulting certificate
visibility is zero.

Write the four geometric factors associated with an active slot as

\[
M_{\mathfrak G_q}=M_{\Phi_Q},\qquad
C_{\mathfrak G_q}=C_{\Phi_Q,\mathfrak K_q}^{\mathrm{auth}}(D_Q),\qquad
R_{\mathfrak G_q}=R_T^q(D_Q),\qquad
\rho_{\mathfrak G_q}=\rho_{\Phi_Q}(D_Q).
\]

For the null slot use the conventions

\[
M_{\varnothing}=C_{\varnothing}=1,
\qquad
\rho_{\varnothing}=0,
\qquad
R_{\varnothing}
=
\frac{\|P_q^{\mathrm{amb}}y\|^2}{\|y\|^2},
\]

where \(P_q^{\mathrm{amb}}=U_qP_q\) is the orthogonal projection onto
the quotient-observable subspace of \(L^2(X,\mu)\).  The certificate
visibility is the single expression

\begin{equation}
\label{eq-supported-geom-visibility}
V_{\mathscr C}
=
\overline A_q S_q(\lambda)M_{\mathfrak G_q}
C_{\mathfrak G_q}
\frac{R_{\mathfrak G_q}}{1+\rho_{\mathfrak G_q}}.
\end{equation}

\medskip\noindent\textbf{(viii) Complete charge.}
Every quotient, fiber predicate, conductance, potential, current,
authority envelope, selected control, transport map, component defect,
comparison window, defect estimate, and normalization factor in this
expression is charged in the single derivation of \(\mathscr C\).  A
current that is target-aligned but not authorized by that derivation
cannot be inserted into the certificate.  A null geometry slot
contributes no additional cost, so its value is exactly

\[
V_{\mathscr C}
=A_q\frac{\|P_q^{\mathrm{amb}}y\|^2}{\|y\|^2}=V_q.
\]

\end{definition}

\Cref{fig-master-certificate-dependency-dag} displays the same-record gates
whose product defines master visibility.

\begin{figure}[tbp]
\centering
\resizebox{\linewidth}{!}{%
\begin{tikzpicture}[fw diagram,node distance=7mm and 9mm]
  \node[fw source,text width=2.6cm] (gen)
    {normalized selected\\generator \(L\)};
  \node[fw provenance,text width=2.7cm,right=of gen] (q)
    {quotient, fibers,\\terminal gates};
  \node[fw provenance,text width=2.7cm,right=of q] (geom)
    {geometry and Caus\\authority record};
  \node[fw use,text width=2.5cm,right=of geom] (stab)
    {dynamic stability\\\(S_q(\lambda)\)};
  \node[fw use,text width=2.6cm,right=of stab] (reach)
    {utility, reach, energy\\\(\bar A_q,M,C,R,\rho\)};
  \node[fw terminal,text width=2.5cm,right=of reach] (vis)
    {visibility\\\(V_{\mathscr C}\)};
  \draw[fw arrow] (gen) -- (q);
  \draw[fw arrow] (q) -- (geom);
  \draw[fw arrow] (geom) -- (stab);
  \draw[fw arrow] (stab) -- (reach);
  \draw[fw arrow] (reach) -- (vis);
  \node[fw residual,text width=4.0cm,below=11mm of geom] (missing)
    {missing qualification, unauthorized current,\\or uncharged comparison};
  \draw[fw dashed arrow] (missing) -| node[pos=.2,below,font=\scriptsize]
    {visibility \(0\)} (vis.south);
\end{tikzpicture}%
}
\caption{The master certificate is a dependency DAG, not a menu of
independent estimates.  Quotient utility, represented geometry, authorized
orientation, dynamic transfer, and target reach must occur in one affordable
record.  A missing gate cannot be supplied from another certificate.}
\label{fig-master-certificate-dependency-dag}
\end{figure}

\subsection{Saturated master profiles}
\label{sssec-master-profiles}
\label{subsec:saturated-master-profiles}

The master profile takes the budgeted supremum over the compensated certificate families.
\begin{definition}[Master saturated quotient--geometry profile]
\label{def-four-mode-geom-abs-profile}
\label{def-master-saturated-quotient-geometry-profile}

The master profile comprises the certificate partition below, the numerical
coordinates of \cref{def:master-profile-coordinates}, and the authority
restriction of \cref{def:authority-relative-master-profile}.  The present
definition retains the established public labels for the combined profile.

Fix a task family and a size slice \(\mathfrak I_n\).  In this
definition a symbol \(\mathscr C\) denotes a family
\((\mathscr C_I)_{I\in\mathfrak I_n}\) produced by one derivation
scheme \(\rho\in\mathsf{Der}_{\mathfrak I}\) fixed across all sizes and
instances.  Every displayed predicate holds pointwise, and
\(\operatorname{Cost}_{\mathrm{sat}}(\mathscr C)\preceq B\) means that
the same derivation has that cost bound for every
\(I\in\mathfrak I_n\).  For a distinguished instance sequence
\((I_n)\), write \(V_{\mathscr C}=V_{\mathscr C_{I_n}}\); an
instance-level or family-functional profile replaces this evaluation by
the corresponding declared functional.  Per-instance re-selection of
the certificate program is not permitted.

\medskip\noindent\textbf{(i) Certificate families.}
Partition the master certificates into the following five disjoint
families, according to the displayed data in the certificate record:

\[
\begin{aligned}
\mathcal C_X(B)
&=\{\mathscr C:q=\mathrm{id},\ \mathfrak G_q\ne\varnothing,\
  \operatorname{Cost}_{\mathrm{sat}}(\mathscr C)\preceq B\},\\
\mathcal C_A(B)
&=\{\mathscr C:E_q=0,\
  \mathfrak G_q=\varnothing,\
  \operatorname{Cost}_{\mathrm{sat}}(\mathscr C)\preceq B\},\\
\mathcal C_{Q,\mathrm{ex}}(B)
&=\{\mathscr C:q\ne\mathrm{id},\ E_q=0,\
  \mathfrak G_q\ne\varnothing,\
  \operatorname{Cost}_{\mathrm{sat}}(\mathscr C)\preceq B\},\\
\mathcal C_{Q,\mathrm{rev}}(B)
&=\{\mathscr C:q\ne\mathrm{id},\ E_q\ne0,\
  A=0,\
  S_q(\lambda)=1-\vartheta_\lambda^T(q)>0,\
  S_q(\lambda)^{-1}\in\mathfrak C,\
  \operatorname{Cost}_{\mathrm{sat}}(\mathscr C)\preceq B\},\\
\mathcal C_{Q,\mathrm{nr}}(B)
&=\{\mathscr C:q\ne\mathrm{id},\ E_q\ne0,\
  S_q(\lambda)=1-\Delta_\lambda(q)>0,\
  S_q(\lambda)^{-1}\in\mathfrak C,\
  \operatorname{Cost}_{\mathrm{sat}}(\mathscr C)\preceq B\}.
\end{aligned}
\]

The fourth and fifth families are disjoint as certificate records: they
are indexed by the chosen reversible Schur-complement record or general
resolvent record, respectively.  A reversible operator may be analyzed
with either record, but each charged certificate has one declared
stability record.

\end{definition}

\begin{definition}[Master-profile coordinates]
\label{def:master-profile-coordinates}

Fix the task family, size slice, budget, and five certificate families of
\cref{def-master-saturated-quotient-geometry-profile}.

\medskip\noindent\textbf{(ii) Profile coordinates.}
Define

\[
\begin{aligned}
G_{A,n}^{\mathrm{sat}}(B)
&=\sup_{\mathscr C\in\mathcal C_A(B)}V_{\mathscr C},\\
G_{X,n}^{\mathrm{sat}}(B)
&=\sup_{\mathscr C\in\mathcal C_X(B)}V_{\mathscr C},\\
G_{Q,\mathrm{ex},n}^{\mathrm{sat}}(B)
&=\sup_{\mathscr C\in\mathcal C_{Q,\mathrm{ex}}(B)}V_{\mathscr C},\\
G_{Q,\mathrm{rev},n}^{\mathrm{sat}}(B)
&=\sup_{\mathscr C\in\mathcal C_{Q,\mathrm{rev}}(B)}V_{\mathscr C},\\
G_{Q,\mathrm{nr},n}^{\mathrm{sat}}(B)
&=\sup_{\mathscr C\in\mathcal C_{Q,\mathrm{nr}}(B)}V_{\mathscr C},
\end{aligned}
\]

with \(\sup\varnothing=0\).  Set

\[
G_{Q,\mathrm{ql},n}^{\mathrm{sat}}(B)
=
\max\left\{
G_{Q,\mathrm{rev},n}^{\mathrm{sat}}(B),
G_{Q,\mathrm{nr},n}^{\mathrm{sat}}(B)
\right\},
\]

and retain the active-geometry profile

\begin{equation}
\label{eq-geom-three-mode-profile}
G_n^{\mathrm{sat}}(B)
=
\max\left\{
G_{X,n}^{\mathrm{sat}}(B),
G_{Q,\mathrm{ex},n}^{\mathrm{sat}}(B),
G_{Q,\mathrm{ql},n}^{\mathrm{sat}}(B)
\right\}.
\end{equation}

For a declared monotone family functional \(\mathfrak M_n\), replace
\(V_{\mathscr C}\) in each of the five displayed suprema by
\(\mathfrak M_n(I\mapsto V_{\mathscr C,I})\) and write
\[
G_{A,\mathfrak M,n}^{\mathrm{sat}},\quad
G_{X,\mathfrak M,n}^{\mathrm{sat}},\quad
G_{Q,\mathrm{ex},\mathfrak M,n}^{\mathrm{sat}},\quad
G_{Q,\mathrm{rev},\mathfrak M,n}^{\mathrm{sat}},\quad
G_{Q,\mathrm{nr},\mathfrak M,n}^{\mathrm{sat}}.
\]
Restricting the same record families to temporally admissible prospective
sources gives the corresponding superscript \(\mathrm{src,sat}\).  This is
family evaluation of the existing five coordinates, not an additional
source profile.

\end{definition}

\begin{definition}[Authority-relative master profile]
\label{def:authority-relative-master-profile}

Fix the task family, size slice, budget, and master-profile coordinates of
\cref{def:master-profile-coordinates}.

\medskip\noindent\textbf{(iii) Authority restriction.}
Let \(\mathcal C_{QG}^{\mathrm{all}}(B)\) be the union of the five
certificate families.  For a declared Caus authority envelope
\(\mathfrak A\), let

\[
\mathcal C_{QG}^{\mathfrak A}(B)
=
\left\{
\mathscr C\in\mathcal C_{QG}^{\mathrm{all}}(B):
\begin{array}{l}
\text{the selected generator belongs to }\mathfrak A,\\
\text{and }\mathfrak K_q\text{ is its transported authority record}
\end{array}
\right\}.
\]

For a null geometry slot, the second condition is replaced by
\(\mathfrak K_q=\varnothing\); its quotient and stability data must
still be computed for a generator selected from \(\mathfrak A\).  The
authority-relative master profile is

\begin{equation}
\label{eq-authority-relative-master-profile}
\mathfrak U_{n,\mathfrak A}^{\mathrm{sat}}(B)
=
\sup_{\mathscr C\in\mathcal C_{QG}^{\mathfrak A}(B)}
V_{\mathscr C}.
\end{equation}

For the free authority envelope, write

\[
\mathcal C_{QG}(B)=\mathcal C_{QG}^{\mathrm{free}}(B)
\]

and

\begin{equation}
\label{eq-master-quotient-geometry-profile}
\mathfrak U_n^{\mathrm{sat}}(B)
=
\mathfrak U_{n,\mathrm{free}}^{\mathrm{sat}}(B)
=
\sup_{\mathscr C\in\mathcal C_{QG}(B)}V_{\mathscr C}.
\end{equation}

\medskip\noindent\textbf{(iv) Regime interpretation.}
The imposed and hybrid profiles are restrictions of the same universal
master-certificate family, not additional structural regimes.

The first active term agrees with the ambient definition above.  The
five parameter regimes are pure ambient Geom/Caus, pure exact Abs,
Geom/Caus on exact Abs, reversible geometrically compensated
quasi-Abs, and general-resolvent geometrically compensated quasi-Abs.
In log-cost notation,
\(G_n(d)=G_n^{\mathrm{sat}}(2^d)\) and
\(\mathfrak U_n(d)=\mathfrak U_n^{\mathrm{sat}}(2^d)\).

\end{definition}

\begin{example}[Running example: an LPN parity quotient]
\label{ex-lpn-running-part3}

Continue \cref{ex-lpn-running-part2}. A represented parity map
\(q_u(x)=\langle u,x\rangle\) is counted by the exact quotient profile only
when its public derivation, fibers, induced dynamics, terminal data, and utility
record have total saturated cost at most \(B\). Its contribution is exactly
\(V_I(q_u)\) from \cref{def-abs-extraction}; no separate semantic quotient is
added to the profile.

\end{example}

\subsection{Authority restriction and master normal form}
\label{sssec-master-normal-form}
\label{subsec:master-authority-normal-form}

\begin{lemma}[Authority monotonicity of the master profile]
\label{thm-authority-monotonicity-master-profile}

Let \(\mathfrak A_1\) and \(\mathfrak A_2\) have the same symmetric
generator, boundary conditions, and terminal sector.  Suppose every
selected total skew operator legal under \(\mathfrak A_1\), together
with its provenance, has an identical representation of no greater
cost under \(\mathfrak A_2\).  Equivalently, as resource-filtered
represented families,

\[
\mathfrak A^{\Caus}_{1,I,\le B}
\subseteq
\mathfrak A^{\Caus}_{2,I,\le B}.
\]

Then

\[
\mathfrak U_{n,\mathfrak A_1}^{\mathrm{sat}}(B)
\le
\mathfrak U_{n,\mathfrak A_2}^{\mathrm{sat}}(B).
\]

The imposed envelope is contained in every hybrid envelope whose
control family contains zero.  Such a hybrid envelope is contained in
the free envelope whenever each of its total skew operators is
budget-admissibly represented there.  Under these conditions, the
imposed profile is bounded by the hybrid profile, which is bounded by
the free profile.

\end{lemma}

\begin{proof}

Every selected generator and causal provenance record legal under
\(\mathfrak A_1\) is legal under \(\mathfrak A_2\), with the same
derivation and cost.  Hence
\(\mathcal C_{QG}^{\mathfrak A_1}(B)\subseteq
\mathcal C_{QG}^{\mathfrak A_2}(B)\).  The inequality follows by
taking suprema.  The final assertion is the free/imposed/hybrid
specialization of this inclusion.

\end{proof}

The master normal form places every affordable quotient--geometry record in the common certificate space.
\begin{theorem}[Affordable quotient--geometry master normal form]
\label{thm-affordable-geom-abs-normal-form}
\label{thm-affordable-quotient-geometry-master-normal-form}

Fix a declared Caus authority envelope. Every dynamically qualified,
authority-preserving Geom/Abs structural certificate counted by the
normalized source profiles above, and hence already carrying the
charged record \(\mathfrak N\), has, after composing its represented
quotients and retaining the full accumulated cost, a master certificate
of \Cref{def-master-compensated-quotient-geometry-certificate}
in \(\mathcal C_{QG}^{\mathfrak A}(B)\). The intersections of the five
families in
\Cref{def-master-saturated-quotient-geometry-profile} with
\(\mathcal C_{QG}^{\mathfrak A}(B)\) form an exhaustive disjoint
partition of these authority-relative certificate records. Under the
free authority envelope, moreover,

\begin{equation}
\label{eq-master-profile-five-case-identity}
\mathfrak U_n^{\mathrm{sat}}(B)
=
\max\left\{
m_n^{\mathrm{sat}}(B),
G_{X,n}^{\mathrm{sat}}(B),
G_{Q,\mathrm{ex},n}^{\mathrm{sat}}(B),
G_{Q,\mathrm{rev},n}^{\mathrm{sat}}(B),
G_{Q,\mathrm{nr},n}^{\mathrm{sat}}(B)
\right\}.
\end{equation}

A directed use of an active geometry contributes only through its
authorized \(\Caus\) projection and the factor
\(\operatorname{Align}^{\mathrm{auth}}_{\Caus}\). Its current
is the transported current of the same selected generator; a different
target-aligned current is a new control, a new dynamic presentation, or
a residual term. Valid Lift and Bdry add no additional intrinsic
regime.

A non-lumpable quotient without a transfer-class stability factor is
not a dynamically qualified quotient comparison. It can enter the
profile only after a represented quotient refinement, a separately
charged geometry, or a compensated certificate satisfying
\Cref{def-supported-geom-certificate}.

\end{theorem}

\begin{proof}

Choose any represented derivation of the certificate having cost at
most \(B\), and proceed through
its constructor DAG, carrying the selected generator and its authority
record at every node. Primitive conductances and potentials have
identity support and give \(q=\mathrm{id}\) with an active geometry
slot. Their directed use is admitted only when the full selected
generator supplies the local or carrier comparison for the canonical
geometric current. A represented quotient constructor gives
\(q\ne\mathrm{id}\). Nested quotient constructors collapse to one
represented quotient;
\Cref{prop-composed-quotient-defect} gives its exact defect
and its symmetric and causal components, and charges every intermediate
fiber interface. If the composite defect vanishes, the geometry slot is
null precisely when the record uses only the quotient extraction
coordinate; this is pure exact Abs. The component ledger is retained
even when \(E_q^{\mathsf s}+E_q^{\mathsf a}=0\), so exactness by
cancellation does not erase the selected causal operator. An active slot
gives Geom/Caus on exact Abs. If the defect is nonzero, dynamic
qualification requires an active slot and a transfer-class stability
record for the full generator. The reversible Schur-complement record
requires \(A=0\); otherwise the general resolvent record gives the last
family.  When \(A=0\), either comparison may be selected and charged.

Finite joins, postprocessing, and computation closure may produce a new
potential or quotient, but the output and the entire construction trace
then form a new charged derivation to which the same alternatives apply.
The transformed-gradient normal form assigns every represented directed
use to \(\Caus\), while
\Cref{thm-caus-authority-no-substitution} restricts that use to
the Hodge projection of the actual selected current. An auxiliary
canonical geometric current may orient the Geom certificate but is not
substituted for the operator current without an explicit represented
identity. Exact Lift preserves the source measure,
target fraction, and current identity, while Bdry supplies only scalar
terminal data. Induction over the derivation DAG therefore yields the
authority-relative five-case partition.

It remains to identify the null-slot profile.  Erasing a null slot
changes neither the quotient derivation nor its cost, and the null-slot
conventions in
\Cref{def-master-compensated-quotient-geometry-certificate}
give \(V_{\mathscr C}=V_q\). Conversely, every exact quotient record
counted by \(m_n^{\mathrm{sat}}(B)\) already contains the same charged
normalization record and becomes a master certificate by the canonical
null-slot annotation, without adding a constructor or output datum.
Hence
\(G_{A,n}^{\mathrm{sat}}(B)=m_n^{\mathrm{sat}}(B)\). Taking suprema over
the disjoint partition proves
\eqref{eq-master-profile-five-case-identity}.

\end{proof}

\begin{corollary}[No causal substitution in master normal form]
\label{cor-no-causal-substitution-master-normal-form}

Expansion into, and contraction from, the authority-relative master
normal form preserve the selected total skew operator and the identity
of every authorized current.  Consequently,
\(\mathfrak U_{n,\mathfrak A}^{\mathrm{sat}}(B)\) optimizes only legal
transforms of the imposed dynamics and declared controls.  A different
target-aligned current can affect this profile only after it is admitted
and charged as a control in \(\mathfrak A\); otherwise its
authority-compatible factor is zero and it remains in the residual
ledger.

\end{corollary}

\begin{proof}

The expansion in
\Cref{thm-affordable-quotient-geometry-master-normal-form}
copies the generator and current provenance through every constructor.
Contraction composes the recorded transports and therefore recovers the
same selected operator and current.  The final assertion is
\Cref{thm-caus-authority-no-substitution} applied to the
resulting record.

\end{proof}

\section{Visibility Domination and First-Hit Accounting}
\label{sec:visibility-first-hit-accounting}

The reduction proceeds through exact-geometry domination in
\cref{sssec-visibility-domination}, represented path lifts and compensated
success in \cref{sssec-path-lift-compensation}, the two-source reductions in
\cref{sssec-two-source-reductions}, the canonical exact first-hit quotient in
\cref{sssec-canonical-first-hit}, and pointwise computation accountability in
\cref{sssec-pointwise-accountability}.
The resulting profile reductions are
\cref{cor-two-source-master-profile-reduction,cor-quantitative-composite-alignment,cor-geom-abs-five-class-compatibility,cor-pointwise-algorithm-profile-bound,cor-authority-relative-algorithm-accountability}.

\subsection{Exact quotient-geometry domination}
\label{sssec-visibility-domination}
\label{subsec:exact-quotient-geometry-domination}

\Cref{thm-exact-quotient-geometry-visibility-dominated,thm-dynamic-compensation-first-hit-accountability,lem:exact-affordable-first-hit}
reduce compensated dynamic success to exact
quotient visibility through a represented first-hit construction.

\begin{theorem}[Exact quotient geometry is visibility-dominated]
\label{thm-exact-quotient-geometry-visibility-dominated}

For every budget \(B\),

\begin{equation}
\label{eq-exact-quotient-geometry-dominated}
G_{Q,\mathrm{ex},n}^{\mathrm{sat}}(B)
\le
m_n^{\mathrm{sat}}(B).
\end{equation}

The comparison preserves temporal typing and family evaluation.  In
particular, fix a monotone family functional \(\mathfrak M_n\).  In the
following display the two symbols are the \(\mathfrak M_n\)-evaluated
restrictions of the existing prospective exact-support and exact-quotient
profiles; they introduce no additional profile coordinate:
\begin{equation}
\label{eq-prospective-family-exact-quotient-geometry-dominated}
G_{Q,\mathrm{ex},\mathfrak M,n}^{\mathrm{src,sat}}(B)
\le
m_{\mathfrak M,n}^{\mathrm{src,sat}}(B).
\end{equation}

Consequently, the master profile has the reduced representation

\begin{align}
\mathfrak U_n^{\mathrm{sat}}(B)
&=
\max\left\{
m_n^{\mathrm{sat}}(B),
G_{X,n}^{\mathrm{sat}}(B),
G_{Q,\mathrm{rev},n}^{\mathrm{sat}}(B),
G_{Q,\mathrm{nr},n}^{\mathrm{sat}}(B)
\right\}
\notag\\
&=
\max\left\{
m_n^{\mathrm{sat}}(B),
G_{X,n}^{\mathrm{sat}}(B),
G_{Q,\mathrm{ql},n}^{\mathrm{sat}}(B)
\right\}.
\label{eq-master-profile-reduced}
\end{align}

Thus geometry on an exact quotient remains an operational search mode,
but it introduces no larger target-visibility term than the exact
quotient from which it is constructed.

\end{theorem}

\begin{proof}

Let \(\mathscr C\in\mathcal C_{Q,\mathrm{ex}}(B)\).  Its nonidentity
support quotient is itself a qualified exact-quotient record by
\cref{def-master-compensated-quotient-geometry-certificate}: it retains the
same represented-fiber, terminal-event, positive-utility,
nontrivial-public-structure, normalization, and cost data.  Since
\(E_q=0\), one has \(S_q(\lambda)=1\),
\(\overline C_q=C_q\), and \(\overline A_q=A_q\).  Moreover,

\[
U_q\mathcal M_Q(D_Q)
\subseteq\operatorname{ran}U_q
=\operatorname{ran}P_q^{\mathrm{amb}}.
\]

Orthogonal projection onto a subspace cannot have greater norm than
projection onto a containing subspace, so

\[
R_T^q(D_Q)
\le
\frac{\|P_q^{\mathrm{amb}}y\|^2}{\|y\|^2}.
\]

The mixing indicator and the authority-compatible local factor lie in
\([0,1]\) by
\Cref{def-authority-compatible-caus-alignment}, and
\(\rho_{\Phi_Q}(D_Q)\ge0\).  Substitution in
\eqref{eq-supported-geom-visibility} therefore gives

\[
V_{\mathscr C}
\le
A_q\frac{\|P_q^{\mathrm{amb}}y\|^2}{\|y\|^2}
=V_q.
\]

The qualified exact quotient record, including the gate data just listed, is
a subrecord of \(\mathscr C\), hence its saturated cost is at most \(B\).
Taking the supremum proves
\eqref{eq-exact-quotient-geometry-dominated}.  If \(\mathscr C\) is a
prospective source, its quotient subrecord, nondynamic Abs gates, and the
same qualification indicator are available before the transition for which
the geometry is invoked.
Thus the pointwise comparison preserves temporal source admissibility.
Applying a monotone \(\mathfrak M_n\) and taking the restricted supremum
proves
\eqref{eq-prospective-family-exact-quotient-geometry-dominated}.  Combining
the first inequality
with \eqref{eq-master-profile-five-case-identity} and the definition of
\(G_{Q,\mathrm{ql},n}^{\mathrm{sat}}\) proves
\eqref{eq-master-profile-reduced}.

\end{proof}

\subsection{Represented path lifts and compensated success}
\label{sssec-path-lift-compensation}
\label{subsec:path-lift-compensation}

\begin{lemma}[Represented finite-horizon path lift]
\label{lem-represented-finite-horizon-path-lift}

Let \(P\) be a represented finite Markov or sub-Markov kernel on
\(X\), let \(T\subseteq X\) be absorbing, and let \(\mu_0\) be a
represented initial law. Assume that the declared public full-support
gauge gives \(T\) mass strictly between zero and one. Adjoin a cemetery
state in the sub-Markov case, with its represented deficit weight. For every
integer \(H\ge1\), there is a finite represented outcome/terminal
transcript carrier and an
exact first-hit quotient \(q_{P,H}\) such that

\begin{equation}
\label{eq-path-lift-first-hit-visibility}
\Pr_{\mu_0}\{\tau_T\le H\}
\le eH\,V_{q_{P,H}}.
\end{equation}

The carrier record has cost bounded by a monotone ledger
\(\Psi_{\mathrm{path}}(\operatorname{Cost}(P),H,n)\) consisting of
\(H\) executions of the represented transition interface and retention
of their realized outcomes, \(H+1\) target tests, peak live-state
capacity, the transcript coordinates,
the initialization, the target test, the public
full-support product gauge, the exact quotient fibers, and the deadline
interface. No extensional transition table, rationality,
dyadic-probability, reversibility, or Valid-Lift hypothesis is required.

\end{lemma}

\begin{proof}

On \(\widehat X=X\sqcup\{\partial\}\), complete a sub-Markov kernel by
sending missing mass to the absorbing cemetery state \(\partial\).
Choose a represented finite-outcome realization of \(\widehat P\) and
\(\mu_0\). Such a realization is always available on finite \(X\): the
initial outcome may be \(x_0\) with law \(\mu_0\), and the next outcome
at \(x\) may be the successor \(y\) with law \(\widehat P(x,y)\).
Write \(a_0\) for the initialization outcome and \(a_{t+1}\) for the
outcome of transition \(t\). The recursion of \Cref{def-represented-operational-transition-rule} determines

\[
x_0=x(a_0),
\qquad
x_{t+1}=x_t\cdot a_{t+1}.
\]

For each \(t\), record the target flag
\(s_t=\mathbf 1_T(x_t)\). Let \(\Xi_H\) be the finite padded syntactic
carrier of records

\[
\xi=(a_0,a_1,\ldots,a_H;s_0,s_1,\ldots,s_H),
\]

and put \(\Omega_H=\Xi_H\times\{0,\ldots,H\}\). Its legal
initialization at pointer \(k=0\) is obtained by sequentially executing
the initializer and transition interface once, retaining only the
outcomes and target flags. Equivalently, its law is the analytic identity

\[
\mathbb P_{\mu_0}(\xi)
=p_{\mathrm{init}}(a_0)
 \prod_{t=0}^{H-1}p(a_{t+1}\mid x_t(\xi))
 \prod_{t=0}^{H}
 \mathbf 1\{s_t=\mathbf 1_T(x_t(\xi))\}.
\]

Equip the entire syntactic transcript carrier with a public full-support
product gauge fixed by the padded record encoding and independently of
the legal transcript law. On every syntactic record,
whether or not it has positive legal probability, use the deterministic
replay update

\[
F_H(\xi,k)=
\begin{cases}
(\xi,k),&s_k=1\text{ or }k=H,\\
(\xi,k+1),&s_k=0\text{ and }k<H.
\end{cases}
\]

Under the legal initialization, this update has exactly the first-hit
law of the first \(H\) transitions of \(P\). Preparation uses one live
state \(x_t\) at a time. Replay reads only the stored flags and never
reconstructs or stores \((x_0,\ldots,x_H)\). The displayed product is an
analytic identity for the generated transcript law; its individual
weights need not be evaluated or materialized by a separate constructor.
Let \(P_H\) be the point-mass kernel induced by \(F_H\) and set
\(L_H=P_H-I\).  The path quotient carries the canonical exact-profile
normalization record

\[
\mathfrak N_H=(L_H,1/H,1,P_H,L_H,1/H),
\]

whose entries are already charged by the path update and deadline
ledger.

Adjoin the remaining deadline \(r=H-k\) and define

\[
q_{P,H}(\xi,k)
=\left(r,
\min\{0\le j\le r:s_{k+j}=1\}\right),
\]

with the second coordinate equal to \(\infty\) when there is no hit.
Its fibers are finite Boolean combinations of the represented stored
terminal flags. The deterministic shift sends \((r,j)\) to
\((r-1,j-1)\) for \(j\ge1\), sends \((r,\infty)\) to
\((r-1,\infty)\), and makes the success and deadline cells absorbing.
Thus the quotient intertwines exactly. With the canonical pushforward
quotient gauge, \Cref{prop-canonical-exact-quotient-unit-transfer} gives unit transfer.
The calculation in
\eqref{eq:first-hit-uniformization}--\eqref{eq:first-hit-utility-bound},
which does not require \(\mu_0(T)=0\), now gives

\[
V_{q_{P,H}}
\ge \frac{\Pr_{\mu_0}\{\tau_T\le H\}}{eH}.
\]

Sequential operational generation of the transcript, its retained
outcome and terminal records, one peak native state, the cemetery
interface, and the flag scan give the stated repeated-use cost ledger.
Empty or full reference
targets are excluded from this lemma because their centered contrast is
zero. In any theorem whose left-hand structural visibility is then zero,
the desired upper bound is immediate; an algorithmic full-target task is
handled by the canonical fresh start phase in
\Cref{lem:exact-affordable-first-hit}.

\end{proof}

The first-hit theorem charges complete target arrival through the represented dynamic certificate.
\begin{theorem}[First-hit accounting for dynamic compensation]
\label{thm-dynamic-compensation-first-hit-accountability}

Let
\(\mathscr C=(\mathfrak N,q,\mathfrak G_q,\mathfrak S,\mathfrak K_q)\)
be a dynamically qualified nonzero-defect master certificate for a
finite presented process. Its charged record
\(\mathfrak N=(\widetilde L,\zeta,c,P,L,\lambda)\) supplies
\(L=P-I\) and \(0<\lambda\le1\) by
\Cref{lem-lossless-charged-uniformization}. Let \(\tau_T\) be the
first target-hitting time of the full selected process, measured in
uniformized transitions, and put

\[
\alpha_H=\Pr_{\mu_0}\{\tau_T\le H\}.
\]

For every integer \(H\ge1\), let \(q_{\mathscr C,H}\) be the exact
first-hit quotient of the represented finite-horizon transcript
carrier, including the quotient, geometry, within-fiber data, selected
current, comparison record, path-weight rule, and remaining deadline.
Then

\begin{equation}
\label{eq-compensation-to-first-hit}
V_{\mathscr C}
\le
\lambda\bigl(\alpha_H+(1+\lambda)^{-H}\bigr)
\le
e\lambda H V_{q_{\mathscr C,H}}
+\lambda(1+\lambda)^{-H}.
\end{equation}

In particular, suppose a budget family and an integer \(H_B\ge1\) have
uniform bounds

\[
\operatorname{Cost}_{\mathrm{sat}}(\mathscr C)\preceq B,
\qquad
1\ge\lambda\ge\lambda_B>0
\]

and every resulting first-hit quotient has saturated cost at most
\(\Psi(B,H_B,n)\). Then

\begin{equation}
\label{eq-quasi-envelope-exact-accounting}
G_{Q,\mathrm{ql},n}^{\mathrm{sat}}(B)
\le
eH_Bm_n^{\mathrm{sat}}\!\left(\Psi(B,H_B,n)\right)
+(1+\lambda_B)^{-H_B}.
\end{equation}

The displayed finite-budget inequality is available whenever the stated
uniform \(\lambda_B\) is part of the budget slice.  In the polynomial
specialization set \(B=b_{\mathrm{poly}}\).  For a fixed numerical transfer
margin \(n^{-k}\), let
\[
H_k=\lceil n^{k+3}\rceil,
\qquad
B_k^{\mathrm{hit}}
=\Psi(b_{\mathrm{poly}},H_k,n)
\]
be the class-preserving path-lift ceiling supplied by the declared transfer
system.  No profile-wide lower bound on \(\lambda\) is needed: if
\(m_n^{\mathrm{sat}}(B_k^{\mathrm{hit}})\) is negligible at these named
path-lift ceilings, then the reversible and general-resolvent compensated
quasi-lumpable envelopes at the single declared ceiling
\(b_{\mathrm{poly}}\) are negligible.

\end{theorem}

\begin{proof}

If the carrier target has reference mass zero or one, then
\(V_{\mathscr C}=0\) by \Cref{def-master-compensated-quotient-geometry-certificate}, and every
asserted upper bound is immediate. Assume henceforth that its public
full-support gauge gives the target mass strictly between zero and one.

All active-geometry factors in
\eqref{eq-supported-geom-visibility} lie in \([0,1]\), and
\(\overline C_q\ge1\).  Therefore

\[
V_{\mathscr C}
\le
\overline A_qS_q(\lambda)
\le
\lambda S_q(\lambda)
A_T^Q(\lambda;L_Q,\bar\mu_0).
\]

Dynamic qualification includes the target-functional comparison
\eqref{eq-master-full-target-transfer}; hence

\begin{equation}
\label{eq-master-visibility-bounded-by-full-arrival}
V_{\mathscr C}
\le
\lambda A_T(\lambda;L,\mu_0).
\end{equation}

Uniformization identifies the last target functional with the discounted
first-hit transform

\[
A_T(\lambda;L,\mu_0)
=
\sum_{j\ge0}\Pr_{\mu_0}\{\tau_T=j\}(1+\lambda)^{-j}.
\]

Splitting this sum at \(H\) and using
\((1+\lambda)^{-j}\le(1+\lambda)^{-H}\) for \(j>H\) gives

\[
A_T(\lambda;L,\mu_0)
\le
\alpha_H+(1+\lambda)^{-H}.
\]

The uniformized kernel is represented because every datum in
\(\mathscr C\), its selected dynamics, and the uniformization is part of
the charged constructor record. \Cref{lem-represented-finite-horizon-path-lift} therefore constructs
\(q_{\mathscr C,H}\) on the represented transcript carrier, with all of these
data and the deadline included in its cost, and yields

\[
\alpha_H
\le
eH V_{q_{\mathscr C,H}}.
\]

Combining the preceding three inequalities proves
\eqref{eq-compensation-to-first-hit}.  The quotient
\(q_{\mathscr C,H}\) is exact and has cost at most
\(\Psi(B,H,n)\), so its visibility is bounded by the exact saturated
profile at that enlarged budget.  Taking the supremum over both
nonzero-defect master families, using \(\lambda\le1\), and observing
that \((1+\lambda)^{-H_B}\le(1+\lambda_B)^{-H_B}\) proves
\eqref{eq-quasi-envelope-exact-accounting}.

For completeness, the polynomial conclusion does not require a uniform
\(\lambda_B\). Fix the numerical transfer margin \(n^{-k}\). Every
certificate satisfies
\(V_{\mathscr C}\le\lambda\) by
\eqref{eq-master-visibility-bounded-by-full-arrival}. Split the budget
slice into

\[
\lambda<n^{-(k+2)}
\qquad\text{and}\qquad
\lambda\ge n^{-(k+2)}.
\]

The first subfamily has visibility at most \(n^{-(k+2)}\). For the
second take \(H_k=\lceil n^{k+3}\rceil\). Since \(0<\lambda\le1\),

\[
(1+\lambda)^{-H_k}
\le \exp(-\lambda H_k/2)
\le e^{-n/2}
\]

for all sufficiently large \(n\). The represented path lift at horizon
\(H_k\) has cost at most the named class-preserving ceiling
\(B_k^{\mathrm{hit}}=\Psi(b_{\mathrm{poly}},H_k,n)\). Hence this
subfamily is bounded by

\[
eH_k\,m_n^{\mathrm{sat}}(B_k^{\mathrm{hit}})+e^{-n/2},
\]

which is negligible by the stated bound at that path-lift ceiling. Closure
of the numerical transfer class over its fixed margins gives negligibility
of the full quasi-lumpable envelope. This argument keeps the one declared
polynomial ceiling and its charged path-lift enlargements; it does not
replace them by a quantifier over polynomial budgets.

\end{proof}

\subsection{Two-source and provenance reductions}
\label{sssec-two-source-reductions}
\label{subsec:two-source-provenance-reductions}

\begin{corollary}[Two-source reduction of the affordable master profile]
\label{cor-two-source-master-profile-reduction}

Under the uniform envelope hypotheses of
\Cref{thm-dynamic-compensation-first-hit-accountability},

\begin{equation}
\label{eq-two-source-master-profile-reduction}
\mathfrak U_n^{\mathrm{sat}}(B)
\le
\max\left\{
G_{X,n}^{\mathrm{sat}}(B),
m_n^{\mathrm{sat}}(B),
eH_Bm_n^{\mathrm{sat}}\!\left(\Psi(B,H_B,n)\right)
+(1+\lambda_B)^{-H_B}
\right\}.
\end{equation}

Hence, at \(B=b_{\mathrm{poly}}\), negligible ambient Geom at the declared
ceiling, negligible exact-quotient visibility there, and the path-lift
bounds
\(m_n^{\mathrm{sat}}(B_k^{\mathrm{hit}})=\operatorname{negl}(n)\)
from \cref{thm-dynamic-compensation-first-hit-accountability} imply a
negligible full master profile.  Exact quotient geometry is absorbed at the
declared ceiling; nonzero-defect compensated geometry is absorbed through
its charged exact first-hit quotient at the named class-preserving
enlargement.

\end{corollary}

\begin{proof}

Use the five-case identity
\eqref{eq-master-profile-five-case-identity}.  \Cref{thm-exact-quotient-geometry-visibility-dominated} bounds geometry
on an exact quotient by \(m_n^{\mathrm{sat}}(B)\), while
\eqref{eq-quasi-envelope-exact-accounting} bounds both compensated
nonzero-defect families.  Taking the maximum gives
\eqref{eq-two-source-master-profile-reduction}.  The polynomial
conclusion follows from the final assertion of
\Cref{thm-dynamic-compensation-first-hit-accountability}.

\end{proof}

\begin{corollary}[Quantitative accounting of composite alignment]
\label{cor-quantitative-composite-alignment}

Let \(R\) be a dynamically qualified target-aligned object obtained by
an authority-preserving uniform derivation from public observables,
including a represented finite join, product-containing postprocessing, quotient,
geometry, potential, or lifted computation. If its declared authority
envelope is \(\mathfrak A\), let \(B_R\) be any admissible uniform
budget for its chosen represented derivation, equivalently

\[
B_R\in\mathfrak B_{\mathrm{sat},n}(R).
\]

then its conditioned structural visibility satisfies

\[
\operatorname{Vis}_I(R)
\le
\mathfrak U_{n,\mathfrak A}^{\mathrm{sat}}(B_R)
\le
\mathfrak U_n^{\mathrm{sat}}(B_R).
\]

More generally, for every family \(\mathscr R_B\) of such derivations
with total saturated cost at most \(B\),

\[
\sup_{R\in\mathscr R_B}\operatorname{Vis}_I(R)
\le
\mathfrak U_{n,\mathfrak A}^{\mathrm{sat}}(B)
\le
\mathfrak U_n^{\mathrm{sat}}(B),
\]

when the derivations in \(\mathscr R_B\) share \(\mathfrak A\).

Thus a join of observables may expose target alignment, but the joined
fiber map, the combining rule, and the dynamic comparison enter at their
complete derivation cost. An extensional transformation with no
derivation at budget \(B\) contributes no term to the budget-\(B\)
profile.

\end{corollary}

\begin{proof}

\Cref{thm-affordable-geom-abs-normal-form} assigns the declared
derivation of \(R\) to one of the five master-certificate families. Its
conditioned visibility is therefore a term optimized by
\(\mathfrak U_{n,\mathfrak A}^{\mathrm{sat}}(B_R)\). Authority
monotonicity bounds this by the free profile. Taking the supremum proves
the family statement.
The final assertion is the minimality clause of the budgeted
derivability closure.

\end{proof}

\begin{corollary}[Compatibility with quotient derivation normal form]
\label{cor-geom-abs-five-class-compatibility}
\label{cor-geom-abs-derivation-compatibility}

The quotient supporting each of the last four regimes is the terminal
output of an expanded quotient derivation in the sense of
\Cref{def-derivation-induced-quotient-normal-form}. Adding a
represented conductance, potential, current, or stability comparison
extends that derivation and retains the same quotient fibers. It changes
the dynamic certificate but does not supply a second source of quotient
provenance.  The extension must also retain the selected generator and
causal-current provenance; otherwise it changes the authority envelope
or has zero authority-compatible causal factor. A quotient refinement
is a new terminal fiber map and enters at the saturated cost of its
complete expanded derivation.

\end{corollary}

\begin{proof}

\Cref{thm-expanded-derivation-equivalence} puts the represented
quotient and every attached dynamic datum in one cost-preserving expanded
record. The fiber map is unchanged when the record is extended by
quotient geometry or Caus data.  The causal data are copied from the
selected generator by
\Cref{thm-caus-authority-no-substitution}. The exact and
nonzero-defect cases are then precisely the alternatives in
\Cref{thm-affordable-geom-abs-normal-form}. A refinement changes
the terminal map and therefore has its own quotient record and cost.

\end{proof}

\section{Canonical Exact First-Hit Quotient}
\label{sssec-canonical-first-hit}
\label{sec:canonical-exact-first-hit-quotient}

\begin{definition}[Canonical normalized replay carrier and first-hit quotient]
\label{def:canonical-normalized-first-hit-replay}

Let \(\mathcal A_n\) be a presentation-relative uniform represented
operational computation with running time at most
\(H_0=H_0(n)\ge0\). Assume that its one-step update, peak native and work
state, outcome/terminal transcript, and public terminal verifier have costs bounded
by the declared resource envelope.

\medskip\noindent\textbf{(i) Start normalization.}
Apply the canonical start-state normalization: adjoin a fresh public
phase bit whose start value is outside the success event. The first
transition realizes the original initialization, sets the
phase bit to active, and thereafter runs the original update; the
terminal verifier is enabled only in the active phase. Write

\[
H:=H_0+1
\]

for the normalized horizon. Original success by time \(H_0\), including
success of the original initialization, equals normalized first-hit
success by time \(H\). The augmentation costs one transition, one phase
bit, the represented initializer, and constant description overhead.
Thus support outside success is a normalization available to every
finite-horizon computation, not an assumption on the procedure.

\medskip\noindent\textbf{(ii) Charged transcript carrier.}
Pass to the full operational transcript representation and execute the
represented transition rule sequentially for the normalized horizon.
At each step retain its finite outcome record and the public terminal
record produced by the verifier. Denote the resulting padded transcript by

\[
\xi=(a_1,\ldots,a_H;s_0,\ldots,s_H),
\]

where \(s_t\in\{0,1\}\) is the success flag at time \(t\). The execution
keeps one current native state and does not store its intermediate
values. Adjoin a replay pointer, represented equivalently by a
remaining-deadline coordinate

\[
r\in\{0,\ldots,H\}.
\]

Write the resulting replay configurations as
\((\xi,r)\in\Omega_I\); the current transcript position is \(H-r\).
Pad every finite record coordinate to its declared resource bound and
extend the update inertly on unused syntactic encodings. Equip this
finite transcript product space with the canonical public full-support product gauge
\(\widetilde\mu_I\), fixed by the configuration encoding before the
update is analyzed. This is an admissible configuration reference
measure by \Cref{def-reference-measure-legality}; it requires
no structural-Lift or target-fraction hypothesis.

\medskip\noindent\textbf{(iii) Replay dynamics.}
The augmented replay
update reads \(s_{H-r}\), decrements \(r\), and advances the pointer at
each nonabsorbing step. It never invokes the native successor or the
verifier. Its one-step update is the deterministic map

\[
F_I:\Omega_I\longrightarrow\Omega_I.
\]

Explicitly,

\[
F_I(\xi,r)=
\begin{cases}
(\xi,r),&s_{H-r}=1\text{ or }r=0,\\
(\xi,r-1),&s_{H-r}=0\text{ and }r\ge1.
\end{cases}
\]

Success and deadline expiration are made absorbing. A non-successful
early halt has terminal flags equal to zero thereafter, so replay
continues inertly until it enters the failure cell at \(r=0\). Let

\[
\widetilde T_I
=\{(\xi,r):s_{H-r}=1\}
\]

be the represented replay success event. Under the legal transcript law
this event is exactly the public verifier event of the original
computation. The fresh start flag is zero, and whenever the syntactic
success sector is nonempty the full-support product gauge gives

\[
0<\widetilde\mu_I(\widetilde T_I)<1.
\]

If the syntactic success sector is empty, the success probability is
zero and the asserted bound is immediate. Thus no nondegeneracy
assumption is imposed on the computation.

\medskip\noindent\textbf{(iv) First-hit quotient.}
For a transcript with
remaining deadline \(r\),
define

\[
\tau_r(\xi)
:=
\min\{0\le j\le r:F_I^j(\xi,r)\in\widetilde T_I\},
\]

with \(\tau_r(\xi)=\infty\) if the displayed set is empty, and set

\[
q_{\mathcal A}(\xi,r):=(r,\tau_r(\xi)).
\]

Its quotient state space is

\[
Q_{\mathcal A}
=
\{(r,j):0\le j\le r\le H\}
\cup
\{(r,\infty):0\le r\le H\}.
\]

\end{definition}

\begin{lemma}[Exact budget-admissible first-hit quotient]
\label{lem:exact-affordable-first-hit}

Let \(\mathcal A_n\), \(H\), \(\Omega_I\), \(F_I\),
\(\widetilde T_I\), and \(q_{\mathcal A}\) be the represented objects
constructed under the hypotheses of
\cref{def:canonical-normalized-first-hit-replay}.

\medskip\noindent\textbf{(v) Exactness, cost, and visibility.}
Then \(q_{\mathcal A}:\Omega_I\to Q_{\mathcal A}\) is a
budget-admissible, terminal-compatible, exactly lumpable represented
quotient, with saturated cost bounded by
\eqref{eq:first-hit-saturated-cost} below. As \(I\) varies, these maps
are produced by one uniform quotient derivation from the represented
rule \(\mathcal A_n\), the public verifier, and the declared horizon.
Write \(\mu_{0,I}\) for the induced legal transcript law at replay
position \(r=H\). If

\[
\alpha_I
:=
\Pr_{\mu_{0,I}}\{\tau_H<\infty\},
\]

then its visibility in the sense of Definition
\hyperref[def-abs-extraction]{Quotient extraction coordinate} satisfies

\begin{equation}
\label{eq:first-hit-visibility}
V_{q_{\mathcal A}}
\ge
\frac{\alpha_I}{eH}.
\end{equation}

For this canonical quotient the comparison and normalization factor in
\Cref{def-paper1-abs-usefulness-gate} is exactly one: the
quotient law is the pushforward of the lifted reference law, the
quotient lift is isometric, the dynamics intertwine exactly, and the
terminal functional is preserved.  Consequently, inverse-polynomial
success of a polynomial-resource computation produces inverse-polynomial
visibility at polynomial saturated cost without an additional
comparison hypothesis.

\end{lemma}

\begin{proof}\phantomsection\label{proof-exact-affordable-first-hit}

\medskip\noindent\emph{Step 1: replay preserves the operational law.}

Let \(z_0\) be the normalized fresh native start state. For a transcript
prefix \(a_{1:t}\), define the single current native state recursively by

\[
z_t=z(a_{1:t}),
\qquad
z_{t+1}=z_t\cdot a_{t+1}.
\]

Writing \(\operatorname{ver}_I\) for the public verifier, the
sequentially generated transcript has law

\[
\Pr\{\xi\}
=
\prod_{t=0}^{H-1}p(a_{t+1}\mid z_t)
\prod_{t=0}^{H}
\mathbf 1\{s_t=\operatorname{ver}_I(z_t)\},
\]

with any subprobability deficit represented by the declared failure
outcome. This identity follows inductively from the defining conditional
outcome law at the current native state. Hence preparation followed by
deterministic flag replay has exactly the original terminal and
first-hit law. No transition table or intermediate native-state list is
materialized; preparation applies the represented update to one live
state and appends only its outcome and terminal records.

\medskip\noindent\emph{Step 2: represented fibers and their cost.}

For \(0\le j\le r\), the fiber over \((r,j)\) is

\begin{equation}
\label{eq:first-hit-finite-fiber}
q_{\mathcal A}^{-1}(r,j)
=
\left\{
(\xi,r):
F_I^i(\xi,r)\notin\widetilde T_I\ \ (0\le i<j),
\quad
F_I^j(\xi,r)\in\widetilde T_I
\right\}.
\end{equation}

The no-hit fiber is

\begin{equation}
\label{eq:first-hit-overflow-fiber}
q_{\mathcal A}^{-1}(r,\infty)
=
\left\{
(\xi,r):
F_I^i(\xi,r)\notin\widetilde T_I\ \ (0\le i\le r)
\right\}.
\end{equation}

Given a label \((r,j)\), its fiber predicate is constructed uniformly
from the represented transcript interface and the deadline. Preparing
the transcript uses at most \(H\) applications of the represented
one-step update and \(H+1\) applications of the public verifier.
Membership reads at most \(H+1\) stored terminal flags and invokes
neither operation.
Consequently,

\begin{equation}
\label{eq:first-hit-saturated-cost}
\begin{aligned}
\operatorname{Cost}_{\mathrm{sat},n}(q_{\mathcal A})
\preceq{}&
H\odot c_{\mathrm{step}}
\oplus (H+1)\odot c_T
\oplus c_{\mathrm{native}}^{\mathrm{peak}}
\oplus c_{\mathrm{transcript}}
\oplus c_{\mathrm{init}} \\
&\oplus c_{\mathrm{input}}
\oplus c_{\mathrm{gauge}}
\oplus c_{\mathrm{interface}}
\oplus c_{\mathrm{desc}}
\oplus c_{\mathrm{clock}}.
\end{aligned}
\end{equation}

where the terms are uniform size-slice majorants for the represented
update, verifier, peak native/work state, accumulated outcome/terminal
transcript, initializer, public-input access,
reference-gauge description, task interface, procedure description,
and deadline coordinate in the declared resource monoid. In the numeric
polynomial instance this
specializes to

\[
O\!\left(
H(c_{\mathrm{step}}+c_T)
+c_{\mathrm{native}}^{\mathrm{peak}}+c_{\mathrm{transcript}}+c_{\mathrm{init}}
+c_{\mathrm{input}}+c_{\mathrm{gauge}}+c_{\mathrm{interface}}
+|\operatorname{desc}(\mathcal A_n)|+\log H
\right).
\]

Thus the quotient and its fiber predicates
belong to the budgeted derivability closure. Here a fiber representative
is a compact public fiber predicate or description on which finite
Boolean operations can be performed; no representative element of a
nonempty fiber is required.

\medskip\noindent\emph{Step 3: exact quotient dynamics.}

The remaining-deadline coordinate makes the quotient transition exact.
For \(r\ge1\),

\begin{equation}
\label{eq:first-hit-finite-transition}
(r,j)\longmapsto(r-1,j-1)
\qquad(1\le j\le r),
\end{equation}

and

\begin{equation}
\label{eq:first-hit-overflow-transition}
(r,\infty)\longmapsto(r-1,\infty).
\end{equation}

The cells \((r,0)\) are absorbing success cells, and
\((0,\infty)\) is the absorbing timeout/failure cell. Let \(P\) be the
transition kernel induced by \(F_I\), normalize the algorithm embedding
as \(L_{\Omega_I}=P-I\), and let \(U\) denote the lift of quotient
observables. The target and failure sectors are unions of quotient cells,
and the lifted update intertwines exactly with the quotient update:

\begin{equation}
\label{eq:first-hit-intertwining}
L_{\Omega_I}U=UL_Q.
\end{equation}

Equip the quotient with
\(\nu_Q=(q_{\mathcal A})_*\widetilde\mu_I\) and the pushed-forward
initialization.  The target is the pullback of the success labels, and
both success and timeout sectors are absorbing.  \Cref{prop-canonical-exact-quotient-unit-transfer}, applied to
\eqref{eq:first-hit-intertwining}, therefore identifies the full and
quotient hitting laws exactly and gives

\begin{equation}
\label{eq:first-hit-unit-comparison}
C_{q_{\mathcal A}}=1.
\end{equation}

Its quotient lumpability defect is zero, and all construction overhead
is already included in \eqref{eq:first-hit-saturated-cost}.

\medskip\noindent\emph{Step 4: nontriviality of the quotient record.}

If \(\alpha_I=0\), inequality \eqref{eq:first-hit-visibility} is
immediate. Suppose henceforth that \(\alpha_I>0\). Since the normalized
initial law is concentrated on the fresh start phase outside
\(\widetilde T_I\), some finite first-hit fiber with
\(j\ge1\) has positive initial mass. Together with the deadline labels,
this fiber strictly refines the two-cell verifier partition. Thus the
quotient satisfies the nontrivial-public-structure condition, while its
positive first-hit mass supplies nonzero quotient utility.

\medskip\noindent\emph{Step 5: discounted arrival and visibility.}

Let

\[
y_I=\mathbf 1_{\widetilde T_I}
-\widetilde\mu_I(\widetilde T_I).
\]

Since \(\widetilde T_I\) is a union of the cells \((r,0)\),
\(P^{\mathrm{amb}}_{q_{\mathcal A}}y_I=y_I\), where
\(P^{\mathrm{amb}}_{q_{\mathcal A}}=U_{q_{\mathcal A}}
P_{q_{\mathcal A}}\), and hence

\begin{equation}
\label{eq:first-hit-projection}
\frac{\|P^{\mathrm{amb}}_{q_{\mathcal A}}y_I\|_2^2}
{\|y_I\|_2^2}=1.
\end{equation}

Conditional on requiring \(j\) discrete transitions to hit the target,
uniformization for \(L_{\Omega_I}=P-I\) gives the discount factor
\((1+\lambda)^{-j}\) \citep{Jensen1953}. Consequently,

\begin{equation}
\label{eq:first-hit-uniformization}
A_T^Q(\lambda)
=
\sum_{j=0}^{H}
\Pr_{\mu_{0,I}}\{\tau_H=j\}(1+\lambda)^{-j}.
\end{equation}

Taking \(\lambda=1/H\) gives

\begin{equation}
\label{eq:first-hit-discount-bound}
A_T^Q(1/H)
\ge
(1+1/H)^{-H}\alpha_I
\ge e^{-1}\alpha_I.
\end{equation}

The exact-profile normalization record is the canonical charged tuple

\[
\mathfrak N_{\mathcal A}
=
(L_{\Omega_I},1/H,1,P,L_{\Omega_I},1/H).
\]

Its kernel, generator, discount, and unit majorant are already present
in the transition and clock ledger of
\eqref{eq:first-hit-saturated-cost}; no rate-table enumeration or
additional normalization is used.

By Definition
\hyperref[def-paper1-abs-usefulness-gate]{Part I quotient-utility
condition} and exact lumpability,

\begin{equation}
\label{eq:first-hit-utility-bound}
A_{q_{\mathcal A}}
=
\min\left\{
1,\lambda A_T^Q(\lambda)
\right\}
\ge
\frac{\alpha_I}{eH}.
\end{equation}

Combining \eqref{eq:first-hit-projection} and
\eqref{eq:first-hit-utility-bound} yields

\[
V_{q_{\mathcal A}}
=
A_{q_{\mathcal A}}
\frac{\|P^{\mathrm{amb}}_{q_{\mathcal A}}y_I\|_2^2}{\|y_I\|_2^2}
\ge
\frac{\alpha_I}{eH},
\]

which proves \eqref{eq:first-hit-visibility}.

\end{proof}

\Cref{fig-canonical-first-hit-replay} shows how a completed operational run
becomes an exact accounting quotient without being retyped as a prospective
source.

\begin{figure}[tbp]
\centering
\resizebox{.96\linewidth}{!}{%
\begin{tikzpicture}[fw diagram,node distance=8mm and 11mm]
  \node[fw source,text width=2.6cm] (run)
    {normalized run\\horizon \(H\)};
  \node[fw provenance,text width=2.9cm,right=of run] (trans)
    {padded outcomes and\\flags \(\xi=(a_{1:H},s_{0:H})\)};
  \node[fw use,text width=2.7cm,right=of trans] (replay)
    {replay state\\\((\xi,r)\)};
  \node[fw box,text width=2.8cm,right=of replay] (quot)
    {first-hit quotient\\\(q_{\mathcal A}=(r,\tau_r)\)};
  \node[fw terminal,text width=2.8cm,right=of quot] (vis)
    {exact visibility\\\(V_q\ge\alpha/(eH)\)};
  \draw[fw arrow] (run) -- (trans);
  \draw[fw arrow] (trans) -- (replay);
  \draw[fw arrow] (replay) -- (quot);
  \draw[fw arrow] (quot) -- (vis);
  \node[fw residual,text width=4.0cm,below=11mm of quot] (type)
    {completed terminal record: accounting type};
  \draw[fw dashed arrow] (type) -- node[right,font=\scriptsize]
    {not pre-hit authority} (quot);
\end{tikzpicture}%
}
\caption{The canonical replay retains the charged transcript and advances
only a public deadline pointer.  The quotient \((r,\tau_r)\) is exactly
lumpable and converts success probability to retrospective visibility.  Its
temporal type remains accounting unless a separate pre-hit source certificate
is supplied.}
\label{fig-canonical-first-hit-replay}
\end{figure}

\section{Temporal Source Admissibility and Prospective Accounting}
\label{subsec-temporal-source-admissibility}
\label{sec:temporal-source-admissibility}

The structural ledger distinguishes information available in time to
influence a search from labels constructed from its completed terminal
record.  The distinction is defined on the represented transcript carrier and
is independent of the task family.

\subsection{Pre-hit source and selector records}
\label{subsec:pre-hit-source-selector-records}

\begin{definition}[Pre-hit filtration and temporally admissible source]
\label{def-pre-hit-temporally-admissible-source}

Let \(\mathcal A\) be a represented finite-horizon computation with
normalized horizon \(H\), full clocked transcript carrier
\(\Omega_{\mathcal A}^{[H]}\), and public verifier flags
\(s_0,\ldots,s_H\).  The fresh-start normalization of
\cref{lem:exact-affordable-first-hit} gives \(s_0=0\).  Put

\[
\tau=\inf\{t\in\{1,\ldots,H\}:s_t=1\},
\qquad
\inf\varnothing=\infty.
\]

For \(0\le t<H\), let
\[
\operatorname{pref}_t:
\Omega_{\mathcal A}^{[H]}
\longrightarrow
\Omega_{\mathcal A}^{[t]}
\]
be the represented map retaining the legal transcript through time \(t\).
The complete pre-hit record \(\mathcal H_t\) contains the public
presentation, the represented effects of the current native state, the
outcomes and internal randomness revealed by time \(t\), the candidates
generated by time \(t\), the flags \(s_0,\ldots,s_t\), and the answers of
presentation-admissible oracle calls completed by time \(t\).  It contains no
later outcome, verifier flag, terminal record, or oracle answer.

A represented record \(R\) is \emph{available at time \(t\)} when a
represented prefix evaluator \(R_t\) satisfies

\begin{equation}
\label{eq-prefix-factorization}
R(\omega)
=
R_t\!\left(\operatorname{pref}_t\omega\right)
\end{equation}

for every pair consisting of a legal prefix in \(\{\tau>t\}\) and a legal
completion \(\omega\) of that prefix.

A record is a \emph{temporally admissible structural source} for the
transition after time \(t\) when:

\begin{enumerate}[label=\textup{(\roman*)}]
\item it is available at time \(t\);
\item its complete represented derivation and evaluation cost are charged
before that transition;
\item its derivation has no later verifier flag
\(s_{t+1},\ldots,s_H\), later terminal record, or later oracle answer as an
ancestor; and
\item it is an ancestor in the represented derivation DAG of the transition,
selector, potential, quotient comparison, target-discrimination readout, or
cell-ordering rule whose target progress it is invoked to explain.
\end{enumerate}

Earlier negative verifier flags may occur in \(\mathcal H_t\).  Their effect is
included in the remaining-search-space baseline of
\cref{def-represented-pre-hit-selector}; the flags receive no independent
structural-source charge.  A represented record that fails the four
conditions remains available for terminal or accounting purposes but is not a
source for an earlier transition.

Temporal admissibility is representation-sensitive.  If one denotation has a
retrospective representation and an independent prefix representation, source
admissibility is established only by the latter, with its complete cost and
provenance.  An inadmissible oracle that supplies a later answer changes the
presented family by
\cref{thm-oracle-admissibility-presentation-separation}; it does not make the
answer pre-hit data for the original presentation.

\end{definition}

\begin{definition}[Represented pre-hit selector and neutral baseline]
\label{def-represented-pre-hit-selector}

Use the charged clock refinement in which each transition performs at most
one public verifier test.  A transition that performs \(k>1\) tests is
expanded into \(k\) represented microsteps, and a transition that performs no
test selects a null symbol \(\bot\).  The refined horizon counts these
microsteps and remains within the class-preserving execution enlargement.

On \(\{\tau>t\}\), let \(\mathcal C_t(\mathcal H_t)\) be the represented finite
carrier of eligible next-test cells, candidates, and, when required, the null
symbol. Extend the public verifier \(V_I\) to this carrier by
\(V_I(\bot)=0\).

A represented algorithmic selector is a represented next-test sampler whose
conditional law is the probability kernel

\[
K_t(\,\cdot\mid\mathcal H_t)
\]

on \(\mathcal C_t(\mathcal H_t)\). Its sampler, random-bit use, and transition
implementation are temporally admissible in the sense of
\cref{def-pre-hit-temporally-admissible-source}; the kernel table need not be
materialized. A deterministic selector is the point-mass case. Let

\[
\nu_t(\,\cdot\mid\mathcal H_t)
\]

be a declared search-neutral baseline kernel on the same eligible set.
Examples include uniform enumeration, weighted enumeration, and another
explicitly declared target-neutral testing rule.

For either kernel \(M_t\in\{K_t,\nu_t\}\), write

\[
M_t(V_I=1\mid\mathcal H_t)
:=
\sum_{C\in\mathcal C_t(\mathcal H_t)}
M_t(C\mid\mathcal H_t)V_I(C).
\]

This is analytic notation for the acceptance probability of the next test. It
does not enumerate the carrier or construct the set of accepting candidates.

Define the step-\(t\) target-alignment advantage by

\begin{equation}
\label{eq-pre-hit-selector-advantage}
\operatorname{Adv}_t(K_t;\nu_t)
=
\mathbb E\!\left[
\mathbf 1_{\{\tau>t\}}
\left(
 K_t(V_I=1\mid\mathcal H_t)
 -
 \nu_t(V_I=1\mid\mathcal H_t)
\right)
\right].
\end{equation}

The expectation evaluates the represented sampler analytically; it gives the
computation no accepting candidate or verifier-positive fiber.

For a saturated budget \(B\), define the prospective selector profile

\begin{equation}
\label{eq-prospective-selector-profile}
\operatorname{Sel}^{\mathrm{sat}}_n(B)
=
\sup\left\{
\bigl(\operatorname{Adv}_t(K_t;\nu_t)\bigr)_+:
\begin{array}{l}
0\le t<H;\\
K_t\text{ is family-uniform};\\
K_t\text{ is temporally admissible};\\
\text{complete prefix and transition cost is at most }B;\\
\nu_t\text{ is the baseline}
\end{array}
\right\},
\end{equation}

with value \(0\) when the sampler class is empty. The baseline is an analytic
comparison rule and need not be executed by the computation. Instance and
task-family subscripts are added when more than one presentation is under
discussion.

\end{definition}

\subsection{Prospective selector accountability}
\label{subsec:prospective-selector-accountability}

\begin{theorem}[Prospective selector accountability and no-free
target-aligned rearrangement]
\label{thm-prospective-selector-accountability}

Let \(\mathcal A\) be a presentation-relative uniform represented
finite-horizon computation, including computations expressed through
presentation-admissible oracle interfaces.  Apply the one-test clock
refinement of \cref{def-represented-pre-hit-selector}.  Let \(K_t\) be its
next-test selector after time \(t\), let \(\nu_t\) be the declared neutral
baseline, and put

\[
\operatorname{Enum}_{\nu}(\mathcal A,H)
=
\sum_{t=0}^{H-1}
\mathbb E\!\left[
\mathbf 1_{\{\tau>t\}}
\nu_t(V_I=1\mid\mathcal H_t)
\right].
\]

Then

\begin{equation}
\label{eq-selector-hazard-decomposition}
\Pr\{\tau\le H\}
=
\operatorname{Enum}_{\nu}(\mathcal A,H)
+
\sum_{t=0}^{H-1}
\operatorname{Adv}_t(K_t;\nu_t).
\end{equation}

If every temporally admissible represented sampler within the
class-preserving enlarged budget \(\Psi(B,n)\) has advantage at most
\(\delta_n\ge0\), then every budget-\(B\) computation satisfies

\begin{equation}
\label{eq-selector-accountability-bound}
\Pr\{\tau\le H\}
\le
\operatorname{Enum}_{\nu}(\mathcal A,H)
+
H\operatorname{Sel}^{\mathrm{sat}}_n(\Psi(B,n))
\le
\operatorname{Enum}_{\nu}(\mathcal A,H)
+
H\delta_n.
\end{equation}

Conversely, if

\[
\Pr\{\tau\le H\}
>
\operatorname{Enum}_{\nu}(\mathcal A,H)+\varepsilon,
\]

then some \(0\le t<H\) satisfies

\begin{equation}
\label{eq-selector-extraction}
\operatorname{Adv}_t(K_t;\nu_t)
>
\frac{\varepsilon}{H}.
\end{equation}

Thus every improvement over the declared enumeration baseline exposes a
temporally admissible represented target-aligned sampler. The responsible
record is its pre-hit sampler and transition implementation. Any ranking,
potential, or quotient readout invoked as its structural source must occur in
that sampler's pre-hit ancestry. A label formed after target contact is not a
source for the transitions preceding that contact.

\end{theorem}

\begin{proof}\phantomsection
\label{proof-prospective-selector-accountability}

The events \(\{\tau=t+1\}\), \(0\le t<H\), are pairwise disjoint.  Conditional
on \(\mathcal H_t\) and \(\{\tau>t\}\), the next test is sampled from
\(K_t(\cdot\mid\mathcal H_t)\). The one-test normalization and the analytic
acceptance-mass notation give the conditional identity

\[
\Pr\{\tau=t+1\mid\mathcal H_t,\tau>t\}
=
K_t(V_I=1\mid\mathcal H_t).
\]

Therefore

\[
\Pr\{\tau=t+1\}
=
\mathbb E\!\left[
\mathbf 1_{\{\tau>t\}}
K_t(V_I=1\mid\mathcal H_t)
\right].
\]

Summing over \(t\) and adding and subtracting
\(\nu_t(V_I=1\mid\mathcal H_t)\) in each summand gives
\eqref{eq-selector-hazard-decomposition}.

The represented sampler inducing \(K_t\) is an ancestor of the transition
that makes the corresponding test and is evaluated before that transition.
Encoding adequacy, computation closure, and the explicit microstep refinement
place the sampler, its random-bit use, and its transition implementation in
the saturated derivability closure at the class-preserving enlarged budget
\(\Psi(B,n)\). No probability table or accepting set is added to that record.
The profile definition bounds each advantage by
\(\operatorname{Sel}^{\mathrm{sat}}_n(\Psi(B,n))\); the assumed uniform bound
then bounds the profile by \(\delta_n\). Summing at most \(H\) terms proves
\eqref{eq-selector-accountability-bound}.

Finally, if a sum of \(H\) real numbers exceeds \(\varepsilon\), at least one
summand exceeds \(\varepsilon/H\). Apply this observation to the advantage
sum in \eqref{eq-selector-hazard-decomposition}. The extracted record is the
represented pre-hit sampler inducing \(K_t\), proving
\eqref{eq-selector-extraction}.

\end{proof}

\subsection{Structural charging of selector success}
\label{subsec:selector-structural-charging}

\begin{theorem}[Prospective selector accounting through five-family quotient
charge]
\label{thm-prospective-selector-structural-charge}

Fix a represented task family, the five-family frontier rule
\(\mathfrak F_5\) of \cref{cor-five-family-abs-envelope}, and a saturated
budget \(B\) with finite time capacity.  Let \(H_B(n)\) be the largest
refined horizon admitted by \(B\).  Let \(\Psi_{\mathrm{hit}}(B,n)\) be the
monotone ceiling obtained by substituting the resource majorants in \(B\)
into the explicit first-hit ledger
\eqref{eq:first-hit-saturated-cost}, and set

\begin{equation}
\label{eq-prospective-selector-charge-budget}
D_B^{\mathrm{sel}}(n)
:=
\Psi_5^{\mathrm{alg}}
\!\left(\Psi_{\mathrm{hit}}(B,n)\right).
\end{equation}

For every presented instance \(I\), write
\(\operatorname{Sel}^{\mathrm{sat}}_{I,n}(B)\) for the pointwise profile
obtained from \eqref{eq-prospective-selector-profile}.  Every represented
selector \(K_t\) counted by this profile satisfies

\begin{equation}
\label{eq-step-selector-structural-charge-bound}
\bigl(\operatorname{Adv}_t(K_t;\nu_t)\bigr)_+
\le
eH_B(n)
\sum_{c\in\mathfrak J_{\Abs}^5}
\mathcal C_{I,\mathfrak F_5}^{(c)}
   \bigl(D_B^{\mathrm{sel}}(n)\bigr).
\end{equation}

Consequently,

\begin{equation}
\label{eq-saturated-selector-structural-charge-bound}
\operatorname{Sel}^{\mathrm{sat}}_{I,n}(B)
\le
eH_B(n)
\sum_{c\in\mathfrak J_{\Abs}^5}
\mathcal C_{I,\mathfrak F_5}^{(c)}
   \bigl(D_B^{\mathrm{sel}}(n)\bigr),
\end{equation}

Applying any monotone positively homogeneous family functional to the
pointwise inequality gives the corresponding size-slice statement.  Every
budget-\(B\) computation also satisfies the whole-transcript bound

\begin{equation}
\label{eq-prospective-success-structural-charge-bound}
\Pr\{\tau\le H_B\}
\le
eH_B(n)
\sum_{c\in\mathfrak J_{\Abs}^5}
\mathcal C_{I,\mathfrak F_5}^{(c)}
   \bigl(D_B^{\mathrm{sel}}(n)\bigr).
\end{equation}

The right-hand sides are contextual-charge envelopes of completed exact
first-hit accounting quotients.  They retain the complete pre-hit sampler and
all later terminal records, but their numerical value is accounting
visibility.  Matching provenance tags or zero incremental charge of terminal
postprocessing does not transfer that value to a prospective source profile.

A prospective-source conclusion follows when the complete record also
contains, for every nonzero frontier term, the represented pre-hit comparison
required by
\cref{rem-dynamic-transfer-from-contextual-charge,cor-contextual-structural-charge-transfer}.
Under that additional same-record comparison, the corresponding source
envelopes replace the accounting envelopes in the three displayed bounds.
Without it, \cref{eq-step-selector-structural-charge-bound,eq-saturated-selector-structural-charge-bound,eq-prospective-success-structural-charge-bound}
are exact accounting estimates only.  In either interpretation, negligibility
of the displayed right-hand side, after multiplication by the admitted
horizon factor, implies the corresponding negligibility conclusion.

\end{theorem}

\begin{proof}\phantomsection
\label{proof-prospective-selector-structural-charge}

\medskip\noindent\emph{Step 1: reduce selector advantage to next-hit mass.}

Fix a selector occurrence at time \(t\), and let \(C_{t+1}\) be the candidate
sampled from \(K_t(\cdot\mid\mathcal H_t)\).  Its unconditional next-hit
mass is

\begin{equation}
\label{eq-selector-next-hit-mass}
h_t
:=
\mathbb E\!\left[
\mathbf 1_{\{\tau>t\}}
K_t(V_I=1\mid\mathcal H_t)
\right]
=
\Pr\{\tau=t+1\}.
\end{equation}

Put

\[
b_t
:=
\mathbb E\!\left[
\mathbf 1_{\{\tau>t\}}
\nu_t(V_I=1\mid\mathcal H_t)
\right].
\]

By \eqref{eq-pre-hit-selector-advantage},
\(\operatorname{Adv}_t(K_t;\nu_t)=h_t-b_t\), and \(b_t\ge0\).  Hence

\begin{equation}
\label{eq-positive-selector-advantage-next-hit}
\bigl(\operatorname{Adv}_t(K_t;\nu_t)\bigr)_+
\le h_t.
\end{equation}

When \(h_t=0\), \eqref{eq-step-selector-structural-charge-bound} follows
immediately.  Assume henceforth that \(h_t>0\).

\medskip\noindent\emph{Step 2: construct the truncated first-hit quotient.}

Form the represented truncated computation
\(\mathcal A^{\langle t\rangle}\) by retaining the original initialization
and transitions through the next test, declaring success only at that final
test, and then stopping.  Its terminal flag is the represented Boolean
observable

\[
z_{t+1}
=
\mathbf 1_{\{s_0=\cdots=s_t=0\}}V_I(C_{t+1}).
\]

The definition of the selector profile charges the complete derivation of
\(\mathcal H_t\), the sampler, and the next transition.  The survival gate,
terminal flag, and stopping instruction are deterministic filtration
postprocessings of retained fields.  Thus
\(\mathcal A^{\langle t\rangle}\) lies in the same execution class, has
horizon \(t+1\le H_B(n)\), and its exact first-hit quotient
\(q_{\mathcal A^{\langle t\rangle}}\) has saturated cost at most
\(\Psi_{\mathrm{hit}}(B,n)\) by
\cref{lem:exact-affordable-first-hit,def-uniform-saturated-success-profile}.
Its success probability is exactly \(h_t\).

Apply \cref{lem:exact-affordable-first-hit} to the truncated computation.
The quotient is terminal compatible and satisfies

\begin{equation}
\label{eq-selector-truncated-first-hit-bound}
h_t
\le
e(t+1)V_I\!\left(q_{\mathcal A^{\langle t\rangle}}\right)
\le
eH_B(n)V_I\!\left(q_{\mathcal A^{\langle t\rangle}}\right).
\end{equation}

\medskip\noindent\emph{Step 3: expand its exact structural charge.}

Expand this quotient by the family-uniform five-family frontier
\(\mathfrak F_5\).  The composed cost bound is precisely
\eqref{eq-prospective-selector-charge-budget}.  Terminal compatibility and
the exact structural-charge identity give

\begin{equation}
\label{eq-selector-truncated-exact-charge}
V_I\!\left(q_{\mathcal A^{\langle t\rangle}}\right)
=
\sum_{\ell=1}^{|\Gamma_5|}
\operatorname{ch}_I
\!\left(E_{\mathcal A^{\langle t\rangle},\Gamma_5,\ell}\right).
\end{equation}

Every extension in \(\Gamma_5\) has at least one authenticated provenance
tag in \(\mathfrak J_{\Abs}^5\).  Since the charges are nonnegative,

\begin{align}
V_I\!\left(q_{\mathcal A^{\langle t\rangle}}\right)
&\le
\sum_{c\in\mathfrak J_{\Abs}^5}
\sum_{\ell:\,\mathsf W_c
 (E_{\mathcal A^{\langle t\rangle},\Gamma_5,\ell})}
\operatorname{ch}_I
\!\left(E_{\mathcal A^{\langle t\rangle},\Gamma_5,\ell}\right)
\notag\\
&\le
\sum_{c\in\mathfrak J_{\Abs}^5}
\mathcal C_{I,\mathfrak F_5}^{(c)}
   \bigl(D_B^{\mathrm{sel}}(n)\bigr).
\label{eq-selector-five-family-charge-cover}
\end{align}

The last inequality is the defining supremum
\eqref{eq-aggregate-contextual-charge-envelope}, applied at the composed
carrier cost.  Combining
\eqref{eq-positive-selector-advantage-next-hit},
\eqref{eq-selector-truncated-first-hit-bound}, and
\eqref{eq-selector-five-family-charge-cover} proves
\eqref{eq-step-selector-structural-charge-bound}.  Taking the supremum over
the selector records in \eqref{eq-prospective-selector-profile} proves
\eqref{eq-saturated-selector-structural-charge-bound}.

\medskip\noindent\emph{Step 4: retain the provenance scope.}

Independent outcome randomness is a valid-\(\Lift\) auxiliary coordinate and
has zero incremental target charge.  Once its state-dependent arguments are
present, every update, stored flag, survival gate, terminal flag, stopping
instruction, and first-hit postprocessing is deterministic and has zero
incremental charge by
\cref{lem-derived-constructor-zero-contextual-charge,thm-filtration-provenance-charging}.
This locates the contextual charge inside the complete accounting DAG.  It
does not prove that a charge-bearing ancestor was available before the
selector occurrence with a numerically dominating prospective visibility.
That latter conclusion requires the represented comparison of
\cref{rem-dynamic-transfer-from-contextual-charge,cor-contextual-structural-charge-transfer}.

\medskip\noindent\emph{Step 5: obtain the whole-computation bound.}

For the whole computation, the assertion is immediate when
\(\Pr\{\tau\le H_B\}=0\).  Otherwise apply the same argument once to its
canonical first-hit quotient \(q_{\mathcal A}\), rather than separately to
its \(H_B\) hazards.  The quotient then has positive utility.  The first-hit
visibility bound, exact structural-charge identity, and five-family cover
give

\[
\Pr\{\tau\le H_B\}
\le
eH_B(n)V_I(q_{\mathcal A})
\le
eH_B(n)
\sum_{c\in\mathfrak J_{\Abs}^5}
\mathcal C_{I,\mathfrak F_5}^{(c)}
   \bigl(D_B^{\mathrm{sel}}(n)\bigr).
\]

The same exact quotient expansion proves
\eqref{eq-prospective-success-structural-charge-bound} without an additional
horizon summation.  By
\cref{cor-first-hit-accounting-not-source}, the completed quotient remains an
accounting record.  A prospective interpretation of a frontier term requires
its separate represented pre-hit comparison, as stated in the theorem.

\end{proof}

An accounting application instantiates
\cref{thm-prospective-selector-structural-charge} by estimating the fixed
five contextual-charge envelopes on the right-hand side.  A source-exhaustion
application additionally supplies the represented transfer comparison in the
theorem.  The supremum in each envelope is already taken over the declared
saturated budget slice.

\subsection{Transition-rule provenance}
\label{subsec:transition-rule-provenance}

\begin{theorem}[No-free succinct transition rule and active-source routing]
\label{thm-no-free-succinct-transition-rule}

Let \(\mathcal A\) be a family-uniform represented finite-horizon computation
of cost at most \(B\), and let \(K_t(\,\cdot\mid\mathcal H_t)\) be the kernel
used for its transition after time \(t\).  Evaluator normalization and
computation closure produce one represented derivation \(\rho_{K_t}\) of that
kernel.  Let \(\widehat B\) be a common class-preserving majorant of its
encoding, evaluator normalization, expanded-derivation, and active-source
join overheads.  At this ceiling the following properties hold.

\begin{enumerate}[label=\textup{(\roman*)}]
\item The inputs of \(\rho_{K_t}\) are the declared public presentation, the
represented pre-hit record \(\mathcal H_t\), and the random coordinates
revealed by time \(t\).  Every state-dependent intermediate value and every
instance-dependent coefficient used by the kernel occurs in its normalized
ancestral DAG.  A table occurring in the representation is charged by its
encoded length, storage, lookup, and evaluation costs.

\item For each dynamically qualified Geom/Abs source record \(\mathscr C\)
invoked to explain the transition, join \(\rho_{K_t}\) with its canonical
source derivation.  In the pure exact-\(\Abs\) mode this is the complete
backward slice of the terminal quotient.  In each active-Geom mode it is the
represented derivation of
\(\operatorname{Act}_{\Geom}(\mathscr C)\), containing the potential, local
comparisons, selected generator, current, authority comparison,
qualification record, and every applicable Geom gate.  These are exactly the
two cases of \cref{def-state-coefficient-source-provenance}.  A succinct
evaluator may generate different transitions at exponentially many states,
but those differences and the source record used to explain them remain in
the corresponding charged joint derivation and its exact five-family
signature.

\item If \(\rho_{K_t}\) is unavailable from the prefix at time \(t\), it is
absent from the prospective selector profile.  If it is available and has
positive selector advantage, that advantage is counted by
\cref{thm-prospective-selector-accountability}.  Structural credit requires
the same-record source comparison of
\cref{cor-contextual-structural-charge-transfer}; a terminal or retrospective
readout supplies accounting only.

\item Every dynamically qualified structural use of the complete active tuple
belongs to exactly one support mode:
\[
\mathsf A,\qquad
\mathsf Q_{\mathrm{ex}},\qquad
\mathsf Q_{\mathrm{rev}},\qquad
\mathsf Q_{\mathrm{nr}},\qquad
\mathsf X.
\]
The mode \(\mathsf A\) is pure exact \(\Abs\) with a null geometry slot.
The next three modes retain a represented nonidentity quotient together with
the applicable exact or compensated Geom data.  The ambient mode
\(\mathsf X\) has identity support and retains the complete full-carrier Geom
record.  Its directed contribution is the authorized \(\Caus\) projection of
the selected current; the gradient-orthogonal component is \(\Res\) and has
zero authorized target alignment.
\end{enumerate}

Consequently, a family-uniform rule with positive advantage over its declared
neutral selector has no unindexed policy-table source.  Its prospective
structural contribution is carried by a represented exact or compensated
quotient record, by a complete identity-supported ambient Geom record, or by
no qualified structural source; neutral baseline contact remains in the
enumeration term of \cref{thm-prospective-selector-accountability}.
Instance-dependent transition data outside the declared public derivability
closure constitute charged advice, an oracle extension, or a change of task
presentation.

\end{theorem}

\begin{proof}

Encoding adequacy and budgeted generation give a represented derivation of
every transition used by \(\mathcal A\); see
\cref{lem-budgeted-generation-gives-representation}.  Evaluator
normalization, expanded-derivation equivalence, and affordable composition
retain the complete implementation and its accumulated cost:
\cref{thm-expanded-derivation-equivalence,%
thm-affordable-composition-no-creation}.  In particular, a lookup table is an
encoded input to that implementation and is charged at its represented size.
When a source record \(\mathscr C\) is invoked, the represented join
constructor combines the kernel record with the terminal quotient in mode
\(\mathsf A\), or with \(\operatorname{Act}_{\Geom}(\mathscr C)\) in one of
the four active-Geom modes.  The composition-cost identity
\eqref{eq-affordable-composition-cost} charges both derivations and the join.
Budget compatibility and the common normalization overhead
\(\Psi_{\mathrm{all}}\) of
\cref{thm-universal-state-coefficient-source-decomposition} provide the
declared class-preserving majorant \(\widehat B\).
Expanding every interface and taking the applicable terminal-quotient or
active-tuple backward slice gives exactly the two-case canonical ancestry and
coefficient-complete signature of
\cref{def-state-coefficient-source-provenance}; its active-Geom specialization
is \cref{def-active-geom-source-frontier}.  This proves \textup{(i)} and
\textup{(ii)}.

Temporal availability is exactly
\cref{def-pre-hit-temporally-admissible-source}.  The selector hazard identity
and extraction statement are
\cref{thm-prospective-selector-accountability}; the required numerical source
transfer is \cref{cor-contextual-structural-charge-transfer}.  Completed
terminal flags and first-hit records retain their accounting status by
\cref{cor-first-hit-accounting-not-source}.  This proves \textup{(iii)}.

The five master modes and their disjoint exhaustiveness are
\cref{thm-affordable-geom-abs-normal-form}.  The four active-Geom support
tags are refined by
\cref{thm-no-unclassified-geom-source-after-quotient-assignment}.  The active
current split and the absence of an independent residual source are
\cref{prop-derived-channels-create-no-structure,%
thm-exhaustive-residual-normal-form}.  In an active-Geom record, the
represented active potential \(D\) satisfies
\(\nabla D\in\mathcal C_{\mathrm{str}}\), and therefore
\[
\left\langle
 (I-\Pi_{\mathcal C_{\mathrm{str}}})F_J,
 -\nabla D
\right\rangle_E=0.
\]
The residual current consequently has zero value in
\(\operatorname{Align}_{\Caus}^{\mathrm{auth}}\) by
\cref{def-quantitative-caus-alignment,%
def-authority-compatible-caus-alignment}.  This proves \textup{(iv)} and the
conclusion.

\end{proof}

\begin{theorem}[Adaptive potential accounting on the represented transcript]
\label{thm-adaptive-potential-transcript-accounting}

Let \(Z_0,\ldots,Z_H\) be a represented finite-horizon transcript and let
\(D_0,\ldots,D_{H-1}\) be real represented pre-hit potentials, where \(D_t\)
is available before the transition from \(Z_t\) to \(Z_{t+1}\).  Then,
pathwise,
\begin{align}
\label{eq-adaptive-changing-potential-accounting}
\sum_{t=0}^{H-1}
\bigl(D_t(Z_{t+1})-D_t(Z_t)\bigr)
&=D_{H-1}(Z_H)-D_0(Z_0)\notag\\
&\quad+
\sum_{t=1}^{H-1}
\bigl(D_{t-1}(Z_t)-D_t(Z_t)\bigr).
\end{align}
For a fixed potential \(D_t=D\), the update sum vanishes and
\begin{equation}
\label{eq-adaptive-fixed-potential-telescoping}
\sum_{t=0}^{H-1}\bigl(D(Z_{t+1})-D(Z_t)\bigr)
=D(Z_H)-D(Z_0).
\end{equation}
Equivalently, if time, history, learned state, and current parameters are
included in the represented clocked carrier and one fixed lifted potential
\(\widetilde D\) is evaluated on that carrier, its increments telescope as
in \cref{eq-adaptive-fixed-potential-telescoping}.  The construction and
update of every field of \(\widetilde D\) remain in its ancestral record.

The identities supply no target comparison by themselves.  A target-directed
interpretation of either endpoint term or any update term requires that term
to occur in a temporally admissible qualified source record.  Its complete
prospective value is then evaluated by the existing target projection and
the exact or compensated quotient, ambient \(\Geom/\Caus\), or terminal
gates.  For a singleton target on \(\mathbb F_2^n\), the target indicator has
the full expansion
\begin{equation}
\label{eq-singleton-target-full-fourier-expansion}
\mathbf1_{\{s\}}(x)
=2^{-n}\sum_{a\in\mathbb F_2^n}\chi_a(x+s),
\end{equation}
so target-directed endpoint progress is not, in general, one distinguished
Fourier coefficient of \(D\).

\end{theorem}

\begin{proof}
Expanding the left side of
\cref{eq-adaptive-changing-potential-accounting}, group the two occurrences
at each intermediate state \(Z_t\).  The uncancelled intermediate term is
\(D_{t-1}(Z_t)-D_t(Z_t)\), which proves the identity.  The fixed-potential
case follows by cancellation.  Computation closure and
\cref{thm-paper4-dependency-principle} retain the clock, history, learned
state, parameter updates, potential evaluations, and their complete costs.
Temporal and structural assignments are
\cref{def-pre-hit-temporally-admissible-source,thm-no-free-succinct-transition-rule,thm-affordable-geom-abs-normal-form}.
Finally, character orthogonality gives
\cref{eq-singleton-target-full-fourier-expansion}.
\end{proof}

\section{Full-Carrier Controls and Core-Refined Accounting}
\label{sec:full-carrier-core-refined-accounting}

\begin{corollary}[Prospective routing of a represented nonlinear readout]
\label{cor-prospective-nonlinear-readout-routing}

In the setting of \cref{thm-no-free-succinct-transition-rule}, let
\(\mathscr C\) be a represented candidate source record and let
\(D_t=F_t\circ J_t\) be a represented nonlinear readout generated from its
pre-hit active frontier, with complete derivation retained in the normalized
joint transition--source record.  Its structural use has the following
exhaustive routing.

\begin{enumerate}[label=\textup{(\roman*)}]
\item If \(D_t\) is not an ancestor of the canonical prospective terminal
field of \(\mathscr C\)---the terminal quotient in mode \(\mathsf A\), or a
field of \(\operatorname{Act}_{\Geom}(\mathscr C)\) in an active-Geom
mode---then it is inactive in the prospective contribution of
\(\mathscr C\).  In particular,
it contributes nothing through \(\mathscr C\).  Any prospective use by a
transition, selector, potential, quotient comparison, cell-ordering rule, or
same-generator comparison belongs to the separate complete source record in
whose canonical ancestry it occurs.  The present readout may remain a
terminal or accounting record.

\item Suppose its represented fiber map \(q_{D_t}\) is nonidentity and
target-compatible.  Exact fiber dynamics with a null geometry slot place the
record in \(\mathsf A\); exact fiber dynamics with an active geometry slot
place it in \(\mathsf Q_{\mathrm{ex}}\).  A represented nonzero defect with a
valid reversible or general-resolvent stability record places it in
\(\mathsf Q_{\mathrm{rev}}\) or \(\mathsf Q_{\mathrm{nr}}\).  Coercive
regularity supplies the reversible compensation bound whenever the quotient,
block data, target comparison, and coercive comparison occur in the same
affordable record.  A fiber map that lacks target compatibility or a required
dynamic comparison receives no quotient-supported source value.

\item Every structural use not admitted by the quotient-supported case is
evaluated through the complete active transition tuple.  In particular, if
the transition uses full-state fields not determined by \(q_{D_t}\), or uses
\(D_t\) as a potential without admitting its fibers as a dynamically
qualified quotient, those fields and their derivations remain in the active
frontier.  A qualified identity-support record belongs to \(\mathsf X\) and
is evaluated by the complete same-record ambient gate product.  The scalar
readout \(D_t\) receives no separate source value.

\item In every support mode, only the authorized gradient projection of the
selected current contributes directed structure.  Absence of a required
same-generator comparison, target alignment, target reach, or mixing gate
sets the corresponding ambient product to zero.  A represented quantitative
regularity loss remains in the factor \(1/(1+\rho_\Phi(D_t))\).  The
gradient-orthogonal component is \(\Res\).
\end{enumerate}

Thus polynomial evaluation of a nonlinear scalar readout supplies neither
quotient dynamics nor target-oriented transport by itself.  Dynamically
sufficient fiber structure is charged to the exact or compensated quotient
modes.  When the complete master record has identity support, full-state
transport is charged to the ambient mode with every active field and gate
retained.  When another nonidentity quotient is active, the complete tuple is
charged to the exact or compensated support mode selected by
\cref{thm-affordable-geom-abs-normal-form}.

\end{corollary}

\begin{proof}

Part~\textup{(i)} is clause~\textup{(iv)} of
\cref{def-pre-hit-temporally-admissible-source}; the extremal-readout
specialization is
\cref{cor-extremal-readout-requires-intermediate-transport}.  For
part~\textup{(ii)}, apply the support alternatives in
\cref{thm-affordable-geom-abs-normal-form}; the coercive specialization is
\cref{thm-coercive-regularity-forces-compensated-quotient}.  Failure of target
compatibility or dynamic qualification excludes the readout from the
quotient-supported certificate class by
\cref{thm-affordable-geom-abs-normal-form}.  Failure of fiber sufficiency, or
use of the scalar as a full-state potential, preserves every additional
active argument by \cref{def-active-geom-source-frontier}; identity support
is the ambient mode of \cref{thm-exhaustive-geom-normal-form}.  This proves
part~\textup{(iii)}.  Part~\textup{(iv)} follows from
\cref{def-authority-compatible-caus-alignment,def-geom-extraction,%
thm-exhaustive-residual-normal-form}.  These cases are exhaustive by
\cref{thm-no-free-succinct-transition-rule}.

\end{proof}

\begin{corollary}[Quotient-free full-carrier control normal form]
\label{cor-quotient-free-full-carrier-control-normal-form}

Work at one saturated ceiling \(B\) and suppose that the four
quotient-supported prospective master coordinates vanish there:
\begin{equation}
\label{eq-quotient-supported-coordinates-vanish}
G_{A,n}^{\mathrm{src,sat}}(B)
=G_{Q,\mathrm{ex},n}^{\mathrm{src,sat}}(B)
=G_{Q,\mathrm{rev},n}^{\mathrm{src,sat}}(B)
=G_{Q,\mathrm{nr},n}^{\mathrm{src,sat}}(B)
=0.
\end{equation}
Every positive dynamically qualified prospective contribution of a
represented transition rule then has the ambient support mode \(\mathsf X\).
For such a record \(\mathscr C\), adjoining the complete pre-hit record to the
state presents its active tuple as a represented full-carrier controlled
process with:

\begin{enumerate}[label=\textup{(\roman*)}]
\item state equal to the represented pre-hit record \(\mathcal H_t\), selected
generator and kernel \(K_t(\,\cdot\mid\mathcal H_t)\), and complete transition
ancestry retained at the charged ceiling;
\item geometry given by the represented relation, conductance, potential,
local comparisons, reach comparison, mixing comparison, and regularity
comparison in \(\operatorname{Act}_{\Geom}(\mathscr C)\); and
\item direction given by the authorized \(\Caus\) projection of the current
selected from that same active record.
\end{enumerate}

Its conditioned directed visibility is exactly the already-defined ambient
record contribution

\begin{equation}
\label{eq-quotient-free-control-visibility}
V^X_{\Phi,D,J,\mathfrak K}
=
M_\Phi
C_{\Phi,\mathfrak K}^{\mathrm{auth}}(D)
\frac{R_T(D)}{1+\rho_\Phi(D)}.
\end{equation}

Thus the usual control or reinforcement-learning description introduces no
additional structural channel: the represented kernel is the transition
rule, the selected generator is the control authority, the potential is the
progress observable, and its authorized gradient current is the \(\Caus\)
component.  Likewise, smooth, spectral, hierarchical, multiscale, and
self-similar descriptions introduce no additional channel.  Such a
description contributes precisely through the represented Geom fields and
comparisons that it supplies in the active tuple, with its complete
derivation cost retained.  Recursive reuse may shorten that derivation; it
does not remove its public ancestors or any active Geom gate.

Consequently, after quotient exhaustion, a prospective structural
contribution is positive precisely through a represented full-carrier
ambient record for which every factor in
\eqref{eq-quotient-free-control-visibility} is positive.  Failure of temporal
availability excludes the record from the prospective profile.  Failure of
authority, same-generator consistency, target alignment, target reach, or
mixing sets its qualified contribution to zero.  Quantitative irregularity
is retained in \(1/(1+\rho_\Phi(D))\).  The remaining gradient-orthogonal
current belongs to \(\Res\).

\end{corollary}

\begin{proof}

The support alternatives in
\cref{thm-no-free-succinct-transition-rule} are exhaustive.  Let
\(\mathscr C\) have positive prospective visibility.  If its support were
\(\mathsf A\), \(\mathsf Q_{\mathrm{ex}}\),
\(\mathsf Q_{\mathrm{rev}}\), or \(\mathsf Q_{\mathrm{nr}}\), its visibility
would be bounded by the corresponding nonnegative coordinate in
\eqref{eq-quotient-supported-coordinates-vanish}, a contradiction.  Hence
\(\mathscr C\) has support \(\mathsf X\).  The active fields in
\textup{(i)}--\textup{(iii)} are exactly the ambient fields of
\cref{def-active-geom-source-frontier}.  A history-dependent finite-horizon
kernel becomes a state-dependent kernel when the represented pre-hit record
is used as the state; this is the clocked transcript representation already
retained by \cref{def-pre-hit-temporally-admissible-source}.

The value of the qualified ambient record is
\eqref{eq-quotient-free-control-visibility} by
\cref{def-geom-extraction}.  The same-generator authority and directional
projection are those of
\cref{def-authority-compatible-caus-alignment}.  Evaluator normalization and
no-creation retain the ancestry and cost of every recursive or multiscale
implementation by
\cref{thm-expanded-derivation-equivalence,%
thm-affordable-composition-no-creation}.  The exhaustive geometric normal
form assigns every represented relation and conductance to the same Geom
coordinate, independently of the descriptive vocabulary used for it; see
\cref{thm-exhaustive-geom-normal-form}.  Finally,
\cref{thm-exhaustive-residual-normal-form} assigns the orthogonal current to
\(\Res\).  These statements prove the corollary.

\end{proof}

\Cref{fig-succinct-transition-structural-routing} displays the support and
source routing of the preceding theorem and corollary.

\begin{figure}[tbp]
\centering
\begin{tikzpicture}[fw diagram,node distance=6mm and 8mm]
  \node[fw source,text width=3.0cm] (rule)
    {family-uniform represented transition rule};
  \node[fw box,text width=2.7cm,below=of rule] (prefix)
    {available in the pre-hit prefix?};
  \node[fw terminal,text width=2.5cm,left=of prefix] (account)
    {terminal or retrospective accounting};
  \node[fw box,text width=2.7cm,below=of prefix] (support)
    {active support};
  \node[fw provenance,text width=2.6cm,below left=of support] (exact)
    {pure exact \(\Abs\) or Geom on an exact quotient};
  \node[fw provenance,text width=2.8cm,below=of support] (comp)
    {reversible or resolvent compensated quotient Geom};
  \node[fw geom,text width=2.4cm,below right=of support] (ambient)
    {identity-supported full-carrier Geom};
  \node[fw box,text width=3.4cm,below=of ambient] (geomdata)
    {represented Geom data:\\ relation, conductance, potential, and
     comparisons\\[-1mm]
     \scriptsize smooth, spectral, hierarchical, multiscale, or self-similar};
  \node[fw use,text width=2.8cm,below=of geomdata] (control)
    {full-carrier controlled process};
  \node[fw box,text width=2.4cm,below=of control] (current)
    {authorized selected current};
  \node[fw residual,text width=2.1cm,below left=of current] (res)
    {\(\Res\): zero authorized alignment};
  \node[fw source,text width=2.2cm,below right=of current] (caus)
    {\(\Caus\): complete Geom gates};
  \draw[fw arrow] (rule) -- (prefix);
  \draw[fw dashed arrow] (prefix) -- node[above,font=\scriptsize] {no} (account);
  \draw[fw arrow] (prefix) -- node[right,font=\scriptsize] {yes} (support);
  \draw[fw arrow] (support) -- (exact);
  \draw[fw arrow] (support) -- (comp);
  \draw[fw arrow] (support) -- (ambient);
  \draw[fw arrow] (ambient) -- (geomdata);
  \draw[fw arrow] (geomdata) -- (control);
  \draw[fw arrow] (control) -- (current);
  \draw[fw dashed arrow] (current) -- (res);
  \draw[fw arrow] (current) -- (caus);
\end{tikzpicture}
\caption{Exhaustive routing of a succinct transition rule.  Succinct
representation exposes one charged derivation rather than a state-indexed
policy table.  The exact branch contains pure exact \(\Abs\) and Geom on an
exact quotient; the compensated branch contains reversible and
general-resolvent quotient--Geom.  The remaining support is
identity-supported ambient Geom.  In the ambient branch,
control and reinforcement-learning terminology describes the represented
kernel and selected generator; smooth, spectral, hierarchical, multiscale,
or self-similar terminology describes the represented Geom data.  These are
not additional source channels.  Only the authorized \(\Caus\) component of
the selected ambient current contributes directed structure; the remaining
current is \(\Res\).}
\label{fig-succinct-transition-structural-routing}
\end{figure}

\begin{definition}[Core-refined first-hit accounting profile]
\label{def-core-refined-first-hit-accounting-profile}

Fix a saturated ceiling \(B\), a finite horizon bound \(H_B\), and a
monotone sublinear positively homogeneous family functional
\(\mathfrak M_n\).  Apply the coefficient-complete five-family
normalization to the exact first-hit record of a represented computation in
the budget slice.  For exact operational and charge-bearing core signatures,
define

\begin{align}
\label{eq-core-refined-first-hit-accounting-profile}
&\operatorname{HitCore}_{
 \mathfrak M,K_{\mathrm{st}},K_{\mathrm{cf}}
 \mid S_{\mathrm{st}},S_{\mathrm{cf}},n}^{\mathrm{sat},=}(B)
\nonumber\\
&\quad:=
\sup
\mathfrak M_n\!\left(
 I\longmapsto V_I(q_{\mathcal A})
\right),
\end{align}
where the supremum ranges over the normalized first-hit records whose
complete carrier cost is at most \(B\), whose exact operational signature is
\((S_{\mathrm{st}},S_{\mathrm{cf}})\), and whose charge-bearing core
signature is \((K_{\mathrm{st}},K_{\mathrm{cf}})\).  An empty restriction
has value zero.  Replacing \(\mathfrak M_n\) by evaluation at one presented
instance \(I\) defines the pointwise notation
\(\operatorname{HitCore}_{I,n}^{\mathrm{sat}}\).  Put

\begin{equation}
\label{eq-core-refined-first-hit-maximum}
\operatorname{HitCore}_{\mathfrak M,n}^{\mathrm{sat}}(B)
=
\max_{\substack{K_{\mathrm{st}},K_{\mathrm{cf}}\\
S_{\mathrm{st}}\supseteq K_{\mathrm{st}}\\
S_{\mathrm{cf}}\supseteq K_{\mathrm{cf}}}}
\operatorname{HitCore}_{
 \mathfrak M,K_{\mathrm{st}},K_{\mathrm{cf}}
 \mid S_{\mathrm{st}},S_{\mathrm{cf}},n}^{\mathrm{sat},=}(B).
\end{equation}

This is an accounting refinement of the existing exact first-hit records.
It is not a new structural source and does not assert that an indexed cell is
nonempty.  A record occurs in the supremum only through its represented
constructor, complete carrier, qualification data, and actual saturated
cost.  In particular, a spectral datum occurs in this profile only when its
constructor and every certificate used to authenticate it occur in that
record at cost at most \(B\).

\end{definition}

\begin{theorem}[Core-refined prospective accountability]
\label{thm-core-refined-prospective-accountability}

Let \(D_B^{\mathrm{hit}}\) be a common class-preserving ceiling for the
microstep refinement, exact first-hit construction, five-family
normalization, and core-signature construction of a budget-\(B\) record.
Then

\begin{align}
\label{eq-core-refined-selector-accountability}
\operatorname{Sel}_{\mathfrak M,n}^{\mathrm{sat}}(B)
&\le
eH_B\,
\operatorname{HitCore}_{\mathfrak M,n}^{\mathrm{sat}}
  (D_B^{\mathrm{hit}}),\\
\label{eq-core-refined-success-accountability}
\operatorname{Succ}_{\mathfrak M,n}(B)
&\le
eH_B\,
\operatorname{HitCore}_{\mathfrak M,n}^{\mathrm{sat}}
  (D_B^{\mathrm{hit}}).
\end{align}

For each cell in
\cref{eq-core-refined-first-hit-accounting-profile}, its visibility has the
exact charge representation

\begin{equation}
\label{eq-core-refined-first-hit-charge-identity}
V_I(q_{\mathcal A})
=
\sum_{\substack{1\le\ell\le |\Gamma_5|\\
v_\ell\notin\operatorname{Dec}_{\mathrm{inh}}
 (\mathscr C_{\mathcal A})}}
\operatorname{ch}_I
 (E_{\mathcal A,\Gamma_5,\ell}).
\end{equation}

Thus the exact operational cell, core cell, cost, and contextual charge of a
successful first-hit record are retained simultaneously.  The theorem makes
no transfer from this accounting visibility to a smaller prospective source
coordinate.  Such a transfer is available only under the represented
comparison required by
\cref{cor-contextual-structural-charge-transfer}.

\end{theorem}

\begin{proof}

The coefficient-complete normalized ancestral record has one exact
operational signature and, by
\cref{thm-operational-core-provenance-separation}, one charge-bearing core
signature.  The finitely many restrictions in
\cref{eq-core-refined-first-hit-maximum} are disjoint and exhaust the
normalized first-hit accounting family.  No record is introduced by taking
this maximum.

For the success bound, apply the exact first-hit inequality of
\cref{lem:exact-affordable-first-hit} to the represented budget-\(B\)
transcript.  Its normalized first-hit record occurs in exactly one cell at
the ceiling \(D_B^{\mathrm{hit}}\), whence

\[
\Pr\{\tau\le H_B\}
\le eH_B V_I(q_{\mathcal A})
\le eH_B
\operatorname{HitCore}_{I,n}^{\mathrm{sat}}(D_B^{\mathrm{hit}}).
\]

Apply the family functional and then take the saturated supremum to obtain
\cref{eq-core-refined-success-accountability}.  For a selector occurrence,
use the truncated record ending at its next represented verifier call, as in
the proof of
\cref{thm-prospective-selector-structural-charge}.  Its positive advantage
is at most its next-hit mass, and the same first-hit inequality places that
mass in one core-refined cell.  Taking the saturated supremum gives
\cref{eq-core-refined-selector-accountability}.

Finally, terminal compatibility and
\cref{eq-core-charge-identity} give
\cref{eq-core-refined-first-hit-charge-identity}.  The accounting record
remains retrospective unless it has a separate pre-hit representation, by
\cref{cor-first-hit-accounting-not-source}.  Therefore no source-profile
comparison follows merely from the core label; the last assertion is exactly
the transfer requirement of
\cref{rem-dynamic-transfer-from-contextual-charge,%
cor-contextual-structural-charge-transfer}.

\end{proof}

\begin{corollary}[Uniform enumeration in a target-permutation-neutral cell
family]
\label{cor-target-permutation-neutral-enumeration}

Assume that initially there are \(N\) represented cells, exactly one of which
is target-containing, and that after every all-negative pre-hit transcript the
target cell is conditionally uniform over the remaining untested cells.
Assume that one verifier call tests at most one cell, and let \(\nu_t\) be
uniform on the remaining untested cells.

Every represented sampler has zero expected advantage over \(\nu_t\).
Every computation making \(H\le N\) distinct tests satisfies

\begin{equation}
\label{eq-uniform-enumeration-bound}
\Pr\{\tau\le H\}
=
\frac{H}{N}.
\end{equation}

With repeated tests or early stopping, \(H/N\) is an upper bound.  More
generally, success exceeding \(H/N\) implies that the pre-hit record breaks
target-permutation neutrality.  The represented datum, score, quotient
readout, potential, or sampler implementation producing that distinction is
target-aligned structure and is charged at its complete saturated cost.

\end{corollary}

\begin{proof}

After \(t\) distinct failures, the target is conditionally uniform among the
\(N-t\) remaining cells.  Hence, for every probability kernel \(K_t\),

\[
\mathbb E\!\left[
K_t(J\mid\mathcal H_t)
\mid
\mathcal H_t,\tau>t
\right]
=
\frac{1}{N-t}
\sum_{C\in\mathcal C_t}K_t(C\mid\mathcal H_t)
=
\frac{1}{N-t}.
\]

The same identity holds for the uniform baseline, so the expected advantage
is zero.  The survival recursion is

\[
\Pr\{\tau>t+1\mid\tau>t\}
=
1-\frac{1}{N-t}
=
\frac{N-t-1}{N-t}.
\]

Multiplying for \(t=0,\ldots,H-1\) yields
\(\Pr\{\tau>H\}=(N-H)/N\), proving
\eqref{eq-uniform-enumeration-bound}. More generally, expose the random ordered
list of distinct cells tested before stopping. Conditional target-permutation
neutrality makes each set of \(d\) distinct tested cells contain the target
with probability \(d/N\). Repetition and early stopping give \(d\le H\), so
their success probability is at most \(H/N\).

\end{proof}

\section{Retrospective Quotients and Source Profiles}
\label{sec:retrospective-quotients-source-profiles}

\begin{corollary}[Retrospective first-hit quotients are accounting records]
\label{cor-first-hit-accounting-not-source}

Let \(q_{\mathcal A}\) be the exact first-hit quotient of
\cref{lem:exact-affordable-first-hit}.  Its coordinates are represented by
postprocessing the completed terminal flags and remaining deadline:

\[
q_{\mathcal A}
=
F(s_0,\ldots,s_H,r).
\]

For an earlier transition after time \(t<H\), this representation depends on
later flags and fails temporal source admissibility whenever two legal
continuations of the same surviving prefix have different later first-hit
times.  In that case \(q_{\mathcal A}\) does not factor through
\(\operatorname{pref}_t\) and cannot influence the selector used after time
\(t\).

The quotient remains a valid saturated accounting record: it preserves the
success law already generated by the computation.  It enters a prospective
source profile only through an independent temporally admissible prefix
representation, charged at that representation's complete cost.  Executing
later transitions and verifier calls to construct the label retains those
operations in the completed derivation and does not move the label before
them in the derivation order.

Consequently, a retrospective first-hit record cannot close a source
exhaustion cycle

\[
\text{residual success}
\longrightarrow
q_{\mathcal A}
\longrightarrow
\text{claimed pre-hit source}.
\]

The charged source ancestry of the completed record is the five-family
frontier of \cref{thm-five-family-abs-provenance,cor-five-family-abs-envelope};
the first-hit postprocessing itself contributes zero incremental charge by
\cref{thm-filtration-provenance-charging}.

\end{corollary}

\begin{proof}

If two legal completions agree through time \(t\) but have different future
first-hit times, their images under \(q_{\mathcal A}\) differ.  No function of
their common prefix can equal \(q_{\mathcal A}\) on both completions, so
\eqref{eq-prefix-factorization} fails.

An independent prefix evaluator is assessed through its own represented
derivation.  An evaluator that obtains its result by executing later updates
and verifier calls is downstream of those operations. Treating the resulting
completed label as their source reverses the derivation order.  By
\cref{thm-five-family-abs-provenance}, the completed quotient factors through
the charged joint frontier of its derivation, and
\cref{thm-affordable-composition-no-creation,thm-filtration-provenance-charging}
assign no new target projection or contextual charge to the terminal
postprocessing.  This proves the asserted accounting/source distinction.

\end{proof}

\begin{corollary}[No retrospective first-hit source loophole]
\label{cor-no-retrospective-first-hit-source-loophole}

An exact first-hit denotation cannot serve as an uncharged structural source
against a saturated-profile exhaustion theorem.  Every represented use of that
denotation falls into one of the following two cases.

\begin{enumerate}[label=\textup{(\roman*)}]
\item The representation fails temporal source admissibility at the claimed
use.  In particular, its derivation may have a later transition outcome,
verifier flag, or terminal record as an ancestor, as in the completed flag
construction of \cref{lem:exact-affordable-first-hit}.  It may remain
accounting-admissible but is source-inadmissible for that use by
\cref{def-pre-hit-temporally-admissible-source,cor-first-hit-accounting-not-source}.
\item The same denotation has an independent pre-hit representation.  Then
that representation, its full derivation cost, and its complete ancestry
belong to the saturated closure.  The source is classified by
\cref{thm-complete-generic-source-theory,thm-five-family-abs-provenance} and is
charged by \cref{cor-five-family-abs-envelope}.
\end{enumerate}

The same denotation may possess representations of both kinds; each is charged
according to its own derivation. Thus defining an observable extensionally by
the first-hit time supplies no third case. Without a represented evaluator it
is unavailable in the presented resource class; with a represented evaluator
it falls into case \textup{(i)} or \textup{(ii)} at each claimed use.

\end{corollary}

\begin{proof}

Every use of an observable in the budgeted closure has a represented
derivation by \cref{def-budgeted-derivability-closure,thm-encoding-adequacy}.
Fix the transition for which source credit is claimed. If the derivation has a
later transition outcome, verifier flag, or terminal record as an ancestor,
it fails the ancestor and prefix-factorization conditions of
\cref{def-pre-hit-temporally-admissible-source}, giving case \textup{(i)}.
If it has no such ancestor and is available before the transition, its prefix
evaluator is the independent representation in case \textup{(ii)}. A record
that is not available before the transition falls under case \textup{(i)} by
the same definition. These alternatives exhaust the represented derivation at
the claimed use. The five-family normal form and its exact contextual-charge
identity classify and charge every case-\textup{(ii)} record.

\end{proof}

\begin{definition}[Source and accounting profiles]
\label{def-source-accounting-profile-separation}

Let
\(\mathcal Q_{\mathfrak I_n,\le B}^{\mathrm{src}}\) be the subfamily of
\(\mathcal Q_{\mathfrak I_n,\le B}^{\mathrm{sat}}\) whose target-aligned
quotient/readout record is temporally admissible for the transition it is
invoked to explain. A record invoked at several transitions satisfies the
condition at each claimed use; a primitive public record available before
execution satisfies it at time \(0\). Set

\[
m_I^{\mathrm{src,sat}}(B)
=
\sup_{\mathbf Q\in
\mathcal Q_{\mathfrak I_n,\le B}^{\mathrm{src}}}
V_I(\mathbf Q).
\]

Define \(G_I^{\mathrm{src,sat}}(B)\) by imposing the same temporal
condition on target-aligned potentials, geometries, rankings, current
comparisons, and readouts. For a distinguished instance sequence \((I_n)\),
put

\[
I_n^{\mathrm{src,sat}}(B)
=
m_{I_n}^{\mathrm{src,sat}}(B)
+
G_{I_n}^{\mathrm{src,sat}}(B).
\]

The superscript \(\mathrm{src,sat}\) on
\(G_X,G_{Q,\mathrm{ex}},G_{Q,\mathrm{rev}},G_{Q,\mathrm{nr}}\), and
\(\mathfrak U\) denotes the same restriction of their defining certificate
families to temporally admissible records.  Their algebraic maximum and
two-coordinate identities are unchanged under this common restriction.

When the instance sequence is fixed, write
\(m_n^{\mathrm{src,sat}}\) and \(G_n^{\mathrm{src,sat}}\) for the first two
terms.

The \emph{source profiles} range only over temporally admissible pre-hit
quotients, potentials, geometries, rankings, and target-discrimination
readouts. Prospective sampler implementations are measured by
\(\operatorname{Sel}^{\mathrm{sat}}_n\) in
\cref{def-represented-pre-hit-selector}. The \emph{accounting profiles} may also
contain finite-horizon terminal records, completed transcripts, truncated
future values, and exact first-hit quotients.  Thus

\[
m_I^{\mathrm{src,sat}}(B)
\le
m_I^{\mathrm{sat}}(B).
\]

An accounting record may certify the hitting law but supplies no causal
explanation of an earlier transition without a separate temporally admissible
representation.  Universal hardness therefore uses
\cref{thm-prospective-selector-accountability}: a negligible prospective
selector profile together with a negligible or explicitly bounded
enumeration baseline implies negligible target arrival.  The exact first-hit
quotient remains an accounting identity.

\end{definition}

\begin{figure}[tbp]
\centering
\begin{tikzpicture}[fw diagram]
  \node[fw box,minimum width=2.5cm] (prefix) at (-5.2,0)
    {pre-hit prefix \(\mathcal H_t\)};
  \node[fw source,minimum width=2.7cm] (sampler) at (-1.7,0)
    {represented sampler};
  \node[fw use,minimum width=2.5cm] (test) at (1.7,0)
    {next verifier test};
  \node[fw terminal,minimum width=2.6cm] (flags) at (5.1,0)
    {later flags};
  \node[fw geom,minimum width=2.7cm] (source) at (-1.7,1.45)
    {pre-hit source \(R_t\)};
  \node[fw provenance,minimum width=2.8cm] (firsthit) at (5.1,-1.45)
    {completed \(q_{\mathcal A}\)};
  \node[fw residual,minimum width=2.8cm] (account) at (1.7,-1.45)
    {accounting profile};
  \draw[fw arrow] (prefix) -- (sampler);
  \draw[fw arrow] (prefix.north) |- (source.west);
  \draw[fw arrow] (source) -- node[right,font=\scriptsize] {available} (sampler);
  \draw[fw arrow] (sampler) -- (test);
  \draw[fw arrow] (test) -- (flags);
  \draw[fw arrow,dashed] (flags) -- node[right,font=\scriptsize] {postprocess} (firsthit);
  \draw[fw arrow,dashed] (firsthit) -- (account);
\end{tikzpicture}
\caption{Temporal typing of the structural ledger. A source used by the next
transition factors through the pre-hit prefix and is charged before the test,
as required by \cref{def-pre-hit-temporally-admissible-source}. Later verifier
flags generate the exact first-hit quotient only in the accounting branch, as
proved in \cref{cor-first-hit-accounting-not-source}.}
\label{fig-temporal-source-accounting-split}
\end{figure}

\subsection{Pointwise computation accountability}
\label{sssec-pointwise-accountability}
\label{subsec:pointwise-computation-accountability}

\begin{corollary}[Pointwise algorithm-to-accounting-profile bound]
\label{cor-pointwise-algorithm-profile-bound}

Let \(B_{\mathcal A}(n)\) denote the right-hand side of
\eqref{eq:first-hit-saturated-cost}, including its fixed constructor
overhead. For every instance \(I\),

\begin{equation}
\label{eq-pointwise-algorithm-profile-bound}
\alpha_I(\mathcal A_n)
\le
eH(n)
m_I^{\mathrm{sat}}\bigl(B_{\mathcal A}(n)\bigr).
\end{equation}

This is an exact accounting implication within the declared uniform
presentation-relative resource envelope, including presentation-admissible
oracle interfaces.  It is not a source-exhaustion implication: the
retrospective first-hit quotient belongs to the accounting profile by
\cref{cor-first-hit-accounting-not-source}.

\end{corollary}

\begin{proof}
The uniform family \(q_{\mathcal A}\) belongs to
\(\mathcal Q_{\mathfrak I_n,\le B_{\mathcal A}}^{\mathrm{sat}}\).
Hence \(V_I(q_{\mathcal A})\le
m_I^{\mathrm{sat}}(B_{\mathcal A})\). Rearranging
\eqref{eq:first-hit-visibility} proves the result.
\end{proof}

\begin{corollary}[Authority-relative computation accountability]
\label{cor-authority-relative-algorithm-accountability}

Assume, in addition, that the full lifted generator of
\(\mathcal A_n\) is selected from a declared Caus authority envelope
\(\mathfrak A\), and that every dynamic interface in its represented
derivation preserves that selection.  Then, pointwise in the presented
instance,

\begin{equation}
\label{eq-authority-relative-algorithm-accountability}
\alpha_I(\mathcal A_n)
\le
eH(n)
\mathfrak U_{n,\mathfrak A}^{\mathrm{sat}}
\bigl(B_{\mathcal A}(n)\bigr).
\end{equation}

Thus the generic clocked-transcript audit remains valid when a mathematical
object imposes its dynamics.  The exact first-hit quotient is an
authority-relative accounting record and is not a prospective source unless
it also has the prefix representation required by
\cref{def-pre-hit-temporally-admissible-source}.

\end{corollary}

\begin{proof}

\Cref{lem:exact-affordable-first-hit} constructs an exact quotient
for the actual lifted generator.  With a null geometry and causal slot,
this is a pure-Abs master certificate.  The additional hypothesis puts
that certificate in
\(\mathcal C_{QG}^{\mathfrak A}(B_{\mathcal A})\).  Hence
\(V_I(q_{\mathcal A})\le
\mathfrak U_{n,\mathfrak A}^{\mathrm{sat}}(B_{\mathcal A})\), and
rearranging the first-hit visibility inequality proves the claim.

\end{proof}

\section{The Uniform Saturated Profile and DTM Accountability}
\label{sec:uniform-saturated-profile-accountability}

Finally, we aggregate the pointwise bounds uniformly over an instance family
and over every represented computation at the declared budget.

The family-level conclusion is organized into uniform profiles in
\cref{sssec-uniform-profiles}, universal accountability in
\cref{sssec-universal-accountability}, finite-slice and residual accounting in
\cref{sssec-finite-residual-accountability}, and the combined profile in
\cref{sssec-combined-profile}.
The family-level formulas are
\cref{cor-uniform-derivation-indexed-quotient-profile,thm-universal-saturated-profile-accountability,cor-arbitrary-dtm-finite-slice-accountability,thm-combined-extraction-bound,prop-structural-extraction-realizes-channel}.

\subsection{Uniform visibility and success profiles}
\label{sssec-uniform-profiles}
\label{subsec:uniform-visibility-success-profiles}

\begin{definition}[Uniform saturated accounting visibility and success profiles]
\label{def-uniform-saturated-success-profile}

Let \(\mathfrak M_n\) be a monotone positively homogeneous functional
on nonnegative functions on \(\mathfrak I_n\). The principal examples
are expectation under a declared instance distribution, supremum over a
size slice, and infimum over a size slice. Define the uniform quotient
visibility profile by

\begin{equation}
\label{eq-uniform-saturated-visibility-profile}
m_{\mathfrak M,n}^{\mathrm{sat}}(B)
=
\sup_{\mathbf Q\in
\mathcal Q_{\mathfrak I_n,\le B}^{\mathrm{sat}}}
\mathfrak M_n\!\left(I\mapsto V_I(\mathbf Q)\right).
\end{equation}

In the temporal typing of
\cref{def-source-accounting-profile-separation}, this original saturated
profile is the family-level accounting profile.

Let \(\mathsf{Alg}_{n,\le B}\) be the uniform presentation-relative
represented computations, including computations using
presentation-admissible oracle interfaces, whose code, initialization,
peak native state and memory, updates, outcome/terminal transcript, and
total work lie
in the resource envelope \(B\). Their
task-facing representation includes the canonical fresh start phase of
\Cref{lem:exact-affordable-first-hit}; this changes neither their
original output law nor their original success probability. Their
optimal affordable success profile is

\begin{equation}
\label{eq-uniform-affordable-success-profile}
\operatorname{Succ}_{\mathfrak M,n}(B)
=
\sup_{\mathcal A\in\mathsf{Alg}_{n,\le B}}
\mathfrak M_n\!\left(I\mapsto\alpha_I(\mathcal A)\right).
\end{equation}

The envelope supplies these bounds rather than assuming them. When
\(\mathsf t(B)<\infty\), set

\[
H_B(n)=1+\left\lfloor\mathsf t(B)\right\rfloor
\]

for the normalized horizon, and let \(\Psi_{\mathrm{hit}}(B,n)\) be the
monotone resource ceiling obtained by substituting the componentwise
majorants declared by \(B\) into the explicit ledger
\eqref{eq:first-hit-saturated-cost}. Equivalently, one may use the
coarser repeated-monoid bound that charges every per-step field by
\(B\), every retained field once by \(B\), and adds the public verifier,
gauge, interface, and clock charges. Axioms B1--B4 make this a uniform
ceiling for every member of \(\mathsf{Alg}_{n,\le B}\). Write
\(\Psi(B,n)=\Psi_{\mathrm{hit}}(B,n)\). Thus ``envelope comparison'' is
notation for a bound derived from the declared resource ledger, not an
additional computational or structural hypothesis. In the polynomial
resource envelope, \(H_B\) and \(\Psi(B,n)\) are polynomial after a
fixed exponent enlargement.

\end{definition}

\begin{corollary}[Uniform derivation-indexed quotient profile]
\label{cor-uniform-derivation-indexed-quotient-profile}

For every family functional \(\mathfrak M_n\) and saturated budget
\(B\),

\begin{equation}
\label{eq-uniform-derivation-indexed-profile}
m_{\mathfrak M,n}^{\mathrm{sat}}(B)
=
\sup_{\widehat\rho\in
\mathsf{NFQ}_{\mathfrak I_n,\le B}}
\mathfrak M_n\!\left(
I\longmapsto
A_{q_{\widehat\rho,I}}
\frac{\|P_{q_{\widehat\rho,I}}y_{\widehat\rho,I}\|_2^2}
{\|y_{\widehat\rho,I}\|_2^2}
\right).
\end{equation}

The right-hand side ranges over one problem-independent constructor
grammar. A task family enters only through its declared primitives,
reference measures, terminal data, and the denotations of those
constructors.

\end{corollary}

\begin{proof}

The indexing families in
\(\mathcal Q^{\mathrm{sat}}_{\mathfrak I_n,\le B}\) and
\(\mathcal Q^{\mathrm{nf}}_{\mathfrak I_n,\le B}\) agree by
\Cref{thm-affordable-exact-quotient-normal-form}. Substitute
their common normal-form records into
\eqref{eq-uniform-saturated-visibility-profile}.

\end{proof}

Universal accountability bounds any admissible computation by the saturated structural profile at its charged budget.
\subsection{Universal saturated-profile accountability}
\label{sssec-universal-accountability}
\label{subsec:universal-saturated-profile-accountability}

\begin{theorem}[Universal saturated accounting-profile bound]
\label{thm-universal-saturated-profile-accountability}

For every task family, family functional \(\mathfrak M_n\), and budget
ceiling \(B\) with finite time capacity \(\mathsf t(B)<\infty\), with
\(H_B\) and \(\Psi\) derived from the declared resource ledger as
above,

\begin{equation}
\label{eq-universal-saturated-profile-accountability}
\operatorname{Succ}_{\mathfrak M,n}(B)
\le
eH_B(n)
m_{\mathfrak M,n}^{\mathrm{sat}}\bigl(\Psi(B,n)\bigr).
\end{equation}

Consequently, for every target success scale \(\varepsilon(n)>0\),

\begin{equation}
\label{eq-universal-profile-lower-bound}
m_{\mathfrak M,n}^{\mathrm{sat}}\bigl(\Psi(B,n)\bigr)
<
\frac{\varepsilon(n)}{eH_B(n)}
\quad\Longrightarrow\quad
\operatorname{Succ}_{\mathfrak M,n}(B)<\varepsilon(n).
\end{equation}

The implication is an exact audit statement.  Source-exhaustion arguments use
\cref{thm-prospective-selector-accountability} rather than substituting a
retrospective first-hit quotient into the source profile.

\end{theorem}

\begin{proof}
Fix \(\mathcal A\in\mathsf{Alg}_{n,\le B}\). The pointwise bound and
the uniform envelope estimates give

\[
\alpha_I(\mathcal A)
\le
eH_B(n)V_I(q_{\mathcal A})
\]

for every \(I\). Monotonicity and positive homogeneity of
\(\mathfrak M_n\) therefore imply

\[
\mathfrak M_n(\alpha_I(\mathcal A))
\le
eH_B(n)
\mathfrak M_n(V_I(q_{\mathcal A})).
\]

The first-hit maps form one uniform represented quotient family of cost
at most \(\Psi(B,n)\), so the last functional is bounded by
\(m_{\mathfrak M,n}^{\mathrm{sat}}(\Psi(B,n))\). Taking the supremum
over \(\mathcal A\) proves
\eqref{eq-universal-saturated-profile-accountability};
\eqref{eq-universal-profile-lower-bound} is its contrapositive.
\end{proof}

\begin{corollary}[Uniform prospective-selector hardness criterion]
\label{cor-universal-prospective-selector-hardness}

Fix a size slice on which every budget-\(B\) computation has refined horizon
at most \(H_B(n)\).  Suppose that, pointwise on that slice,

\[
\operatorname{Enum}_{\nu}(\mathcal A,H_B)\le\beta_n
\]

for every such computation, and every temporally admissible represented
sampler of saturated cost at most \(\Psi(B,n)\) satisfies

\[
\operatorname{Adv}_t(K_t;\nu_t)\le\delta_n.
\]

Then every computation satisfies

\begin{equation}
\label{eq-uniform-prospective-selector-hardness}
\Pr\{\tau\le H_B\}
\le
\beta_n+H_B(n)\delta_n.
\end{equation}

Consequently, for every monotone positively homogeneous family functional
\(\mathfrak M_n\),

\[
\operatorname{Succ}_{\mathfrak M,n}(B)
\le
\mathfrak M_n(\mathbf 1)\,
\bigl(\beta_n+H_B(n)\delta_n\bigr).
\]

If \(\beta_n\) and \(H_B\delta_n\) are negligible, every budget-\(B\)
computation has negligible target-arrival probability.

\end{corollary}

\begin{proof}

Apply \cref{eq-selector-accountability-bound} pointwise and use the two
uniform hypotheses.  Monotonicity and positive homogeneity of
\(\mathfrak M_n\) give the family-functional bound.

\end{proof}

The prospective hardness map and retrospective audit are summarized in
\Cref{fig-first-hit-hardness}.

\begin{figure}[tbp]
\centering
\begin{tikzpicture}[fw diagram]
  \node[fw box,minimum width=2.0cm] (alg) at (-5.0,0.8) {\(\mathcal A\)};
  \node[fw source,minimum width=3.2cm] (sel) at (-1.6,0.8)
    {pre-hit \(K_t\) and \(\nu_t\)};
  \node[fw terminal,minimum width=4.4cm] (bound) at (3.1,0.8)
    {\(\Pr\{\tau\le H\}\le\operatorname{Enum}_{\nu}+H\operatorname{Sel}^{\mathrm{sat}}\)};
  \node[fw provenance,minimum width=2.4cm] (q) at (-1.6,-1.0)
    {completed \(q_{\mathcal A}\)};
  \node[fw residual,minimum width=3.2cm] (audit) at (3.1,-1.0)
    {accounting profile};
  \draw[fw arrow] (alg) -- (sel);
  \draw[fw arrow] (sel) -- (bound);
  \draw[fw arrow,dashed] (alg) -- node[below,font=\scriptsize] {after execution} (q);
  \draw[fw arrow,dashed] (q) -- node[below,font=\scriptsize] {audit only} (audit);
\end{tikzpicture}
\caption{Prospective target advantage is charged through the pre-hit selector
in \cref{thm-prospective-selector-accountability}.  The completed first-hit
quotient remains the retrospective accounting record of
\cref{cor-first-hit-accounting-not-source}.}
\label{fig-first-hit-hardness}
\end{figure}

\subsection{Finite-slice and residual accountability}
\label{sssec-finite-residual-accountability}
\label{subsec:finite-slice-residual-accountability}

\begin{corollary}[Finite-slice accountability for an arbitrary DTM]
\label{cor-arbitrary-dtm-finite-slice-accountability}

Let \(M\) be any deterministic Turing machine and \(I\) any finite
input. Fix also one exposed target verifier whose program is uniform on
the input-size slices. For every finite horizon \(H_0\), the clocked representation in
\Cref{cor-all-turing-machines-represented}, followed by the
canonical start normalization, has a finite resource ceiling
\(B_{M,H_0}(n)\), uniform over every input of encoded length \(n\), and
an exact first-hit quotient satisfying

\[
\Pr\{M(I)\text{ reaches the exposed target by time }H_0\}
\le
e(H_0+1)\,
m_I^{\mathrm{sat}}
\!\left(\Psi_{\mathrm{hit}}(B_{M,H_0}(n(I)),n(I))\right),
\]

where \(n(I)\) is the encoded input size.

No running-time, halting, reversibility, geometry, lumpability, or
Valid-Lift assumption is made on \(M\). If \(M\) is unbounded or
nonhalting, these inequalities hold for every finite slice. Its eventual
target-arrival probability is the monotone limit of the finite-horizon
probabilities. Passing a spectral, resolvent, or budget estimate through
that limit is a separate analytic question and is not part of the DTM
representation theorem.

\end{corollary}

\begin{proof}

Use the finite clocked point-mass kernel from
\Cref{cor-all-turing-machines-represented} and apply
\Cref{lem:exact-affordable-first-hit}.  For fixed \(M\), an input
of length \(n\) can expose or alter at most \(O_M(H_0)\) additional tape
cells in \(H_0\) steps.  The padded configuration width, transition
cost, and fixed verifier cost therefore have one worst-case majorant
depending only on \(M,n,H_0\), uniformly over the size slice.  The
explicit first-hit ledger gives the displayed ceiling
\(B_{M,H_0}(n)\). Monotone convergence of the events
\(\{\tau_T\le H_0\}\) gives the final assertion.

\end{proof}

\begin{remark}[Saturated accountability of the residual]
\label{rem-saturated-residual-accountability}

Apply Theorem
\hyperref[thm-channel-exhaustiveness]{Component exhaustiveness} to the
full lifted generator of \(\mathcal A\). Its five structural components
and the two residual coordinates classified by
\cref{thm-exhaustive-residual-normal-form} may interact through cancellation, memory,
iteration, and adaptive branching. \Cref{thm-universal-saturated-profile-accountability} evaluates their
combined target arrival through the single exact first-hit accounting
quotient. Thus a successful use of dynamics initially assigned to
\(\Res\) produces a represented accounting record at the charged
finite-horizon cost. It produces a prospective \(\Abs\) source only when an
independent pre-hit representation satisfies
\cref{def-pre-hit-temporally-admissible-source}. No positivity assumption on
component partial sums is needed; the prospective lower-bound direction uses
the explicit baseline in \cref{thm-prospective-selector-accountability}.

\end{remark}

The projection \(P_Qy\) is the target mass visible to the represented
quotient, and \(A_Q\) is the \PartI quotient-utility factor. Their
product defines structural visibility. A target-correlated quotient need
not provide a budget-admissible selector; its correlation remains part
of the visibility invariant. Any budget-instantiated
extraction consequence is supplied by the applicable arrival theorem for
the declared class rather than by the definition of
\(m_n^{\mathrm{sat}}(B)\).

The supremum over
\(\mathcal Q_{\mathfrak I_n,\le B}^{\mathrm{sat}}\) includes every
nonlinear or runtime-generated projection \(U\) produced by one uniform
represented derivation and satisfying the stated conditions. If such a
projection is generated without oracle access from the presented instance, has
represented fibers with saturated cost at most \(B\), is exactly lumpable on the
target window, and has non-negligible \PartI quotient utility
factor \(A_Q\), then it is one of the null-geometry certificates counted
by \(\mathfrak U_n^{\mathrm{sat}}(B)\). If it is pointwise geometric or
causal rather than represented-fiber structure, its conditioned target
visibility is counted by the identity-quotient branch of the same master
profile at every budget in its admissible budget set. If it only
perturbs an already declared geometry, the Dirichlet form comparison of
\Cref{quotient-defect-and-geom-form-comparison} is the condition under which the \PartIII geometric comparison
proof transfers; otherwise the perturbation must be counted as its own
geometry at every budget in its admissible budget set. High raw Walsh degree does not
exclude a compact valid representation. Conversely,
\(\mathcal F_I\)-measurability does not imply within-budget cost or
target utility: a measurable target distinction is charged to
\(\mathfrak U_n^{\mathrm{sat}}(B)\) only when
\(B\in\mathfrak B_{\mathrm{sat},n}(R)\) and its \PartI quotient-utility
factor is nonzero. If it
is not counted by \(\mathfrak U_n^{\mathrm{sat}}(B)\) at any declared
budget level, then
it is not exposed within-budget structure for the original presentation:
it is presentation-inaccessible, oracle-dependent, over-budget to
construct, or residual from the public presentation. Thus a nonlinear
change of variables can reclassify residual transport as an Abs
component only by supplying an exposed
represented target-relevant quotient, and it can turn residual motion into a
Geom/Caus component only by supplying exposed geometric or potential data;
either contribution is charged in \(\mathfrak U_n^{\mathrm{sat}}(B)\)
exactly at the budgets in its admissible budget set.

\section{Combined Extraction and Profile Decomposition}
\label{sssec-combined-profile}
\label{sec:combined-extraction-profile-decomposition}

The combined extraction functional joins the exact-quotient and active-geometry profiles at one budget.
\begin{definition}[Combined extraction functional]\label{def-extraction-functional}

The untyped quantities in this definition are the accounting profiles of
\cref{def-source-accounting-profile-separation}.  The budget-\(B\) accounting
extraction functional is

\[
I_n^{\mathrm{sat}}(B)=m_n^{\mathrm{sat}}(B)+G_n^{\mathrm{sat}}(B).
\]

The master-certificate definition gives the exact identity

\begin{equation}
\label{eq-master-and-combined-profile-equivalence}
\mathfrak U_n^{\mathrm{sat}}(B)
=\max\{m_n^{\mathrm{sat}}(B),G_n^{\mathrm{sat}}(B)\},
\qquad
\mathfrak U_n^{\mathrm{sat}}(B)
\le I_n^{\mathrm{sat}}(B)
\le2\mathfrak U_n^{\mathrm{sat}}(B).
\end{equation}

Thus the two-coordinate sum is retained for channelwise accounting.  Its
prospective restriction satisfies

\begin{equation}
\label{eq-source-master-combined-equivalence}
\mathfrak U_n^{\mathrm{src,sat}}(B)
=
\max\{m_n^{\mathrm{src,sat}}(B),G_n^{\mathrm{src,sat}}(B)\},
\qquad
\mathfrak U_n^{\mathrm{src,sat}}(B)
\le I_n^{\mathrm{src,sat}}(B)
\le2\mathfrak U_n^{\mathrm{src,sat}}(B).
\end{equation}

In log-cost notation \(I_n(d)=I_n^{\mathrm{sat}}(2^d)\). When the
surrounding text writes \(m(d)\), \(G(d)\), or \(I(d)\), the variable
\(d\) is this saturated cost level, not raw Walsh degree.

For a family functional \(\mathfrak M_n\), define
\(G_{\mathfrak M,n}^{\mathrm{sat}}(B)\) by taking the analogous supremum
over family-uniform represented Geom/Caus data. Define

\[
\mathfrak U_{\mathfrak M,n}^{\mathrm{sat}}(B)
=\max\left\{
m_{\mathfrak M,n}^{\mathrm{sat}}(B),
G_{\mathfrak M,n}^{\mathrm{sat}}(B)
\right\},
\]

and set

\[
I_{\mathfrak M,n}^{\mathrm{sat}}(B)
=m_{\mathfrak M,n}^{\mathrm{sat}}(B)
+G_{\mathfrak M,n}^{\mathrm{sat}}(B).
\]

The retrospective audit bound therefore has the channel-compatible form

\begin{equation}
\label{eq-universal-combined-profile-bound}
\operatorname{Succ}_{\mathfrak M,n}(B)
\le eH_B\,m_{\mathfrak M,n}^{\mathrm{sat}}(\Psi(B,n))
\le eH_B\,\mathfrak U_{\mathfrak M,n}^{\mathrm{sat}}(\Psi(B,n))
\le eH_B\,I_{\mathfrak M,n}^{\mathrm{sat}}(\Psi(B,n)).
\end{equation}

The first inequality is an accounting identity; prospective hardness uses
\cref{cor-universal-prospective-selector-hardness}.

\end{definition}

\begin{theorem}[Full cost-tagged saturated-profile decomposition]
\label{thm-full-saturated-profile-decomposition}

Fix a represented family, its resource algebra, and the budgeted
derivability closure and saturated cost of
\cref{def-budgeted-derivability-closure,def-saturated-representation-degree}.
Apply the complete-carrier channel decomposition of
\cref{thm-channel-exhaustiveness} to every represented operator record and
assign each complete record its upward-closed admissible budget set and
Pareto-minimal frontier, which exist by
\cref{lem-derivability-well-posed}.  The resulting global cost-tagged ledger
has the following exhaustive decomposition.

\begin{enumerate}[label=\textup{(\roman*)}]
\item Its prospective \(\Abs\) coordinate is
      \(m_n^{\mathrm{src,sat}}(B)\).  Every contributing source quotient has exactly
      one of the five provenance tags in
      \cref{thm-five-family-abs-provenance} and satisfies
      \cref{def-pre-hit-temporally-admissible-source}.  The accounting
      coordinate \(m_n^{\mathrm{sat}}(B)\) may additionally contain
      completed records such as the exact first-hit quotient.
\item Its prospective \(\Geom/\Caus\) coordinate is
      \(G_n^{\mathrm{src,sat}}(B)\).  Every contributing geometry has the
      edge-conductance form and one of the four represented configurations
      in \cref{thm-exhaustive-geom-normal-form}.  Its complete active tuple
      factors through the canonical charged source frontier of
      \cref{thm-exhaustive-five-family-geom-source-classification}; hence its
      unique support tag and nonempty five-family provenance signature place
      it in the exhaustive finite cover
      \(\mathfrak S_{\Geom}^{4}\times\mathfrak J_{\Abs}^{5}\).
      \(\Caus\) is the authorized oriented projection of that same
      represented source and introduces no additional source class.
\item Every \(\Lift\) record has the fiber-kernel and operator
      factorization of \cref{thm-exhaustive-lift-normal-form}, and every
      \(\Bdry\) record has the killing/value factorization of
      \cref{thm-exhaustive-bdry-normal-form}.  These derived channels retain
      their complete carrier cost and transfer target visibility only
      through a represented contextual identity, as required by
      \cref{prop-derived-channels-create-no-structure}.
\item After all represented channel-qualified denotations have been
      assigned their admissible budget sets, the post-saturation residual is
      exactly
      \[
      L^{\Res,\mathrm{sat}}
      =
      L_{\mathrm{dir}}^{\Res}
      +
      L_{\mathrm{nl}}^{\Res}
      \]
      from \cref{thm-exhaustive-residual-normal-form}.  In particular,
      \emph{non-structural} means orthogonal to the represented structural
      channel subspaces, not merely absent from the current execution
      record.  A family-uniform represented structural witness cannot remain
      in this post-saturation residual by
      \cref{prop-family-uniform-residual-reclassification,thm-no-uniform-residual-correlation}.
\end{enumerate}

Refining the two prospective source coordinates in items~\textup{(i)} and
\textup{(ii)} by their canonical complete ancestral sets gives exactly the
source-mode--state-provenance--coefficient-provenance decomposition of
\cref{thm-universal-state-coefficient-source-decomposition}.  Thus the
finite 4960-cell source classification is a refinement of this full saturated
ledger, while \(\Caus\), \(\Lift\), and \(\Bdry\) inherit those cells and
\(\Res\) contributes none.

\end{theorem}

\begin{corollary}[Prospective saturated-profile identities]
\label{cor:prospective-saturated-profile-identities}

Consequently the intrinsic prospective source profile is precisely

\begin{equation}
\label{eq-full-saturated-profile-decomposition}
I_n^{\mathrm{src,sat}}(B)
=
m_n^{\mathrm{src,sat}}(B)
+
G_n^{\mathrm{src,sat}}(B),
\end{equation}

and its master certificate satisfies

\begin{equation}
\label{eq-full-saturated-master-decomposition}
\mathfrak U_n^{\mathrm{src,sat}}(B)
=
\max\!\left\{
m_n^{\mathrm{src,sat}}(B),
G_{X,n}^{\mathrm{src,sat}}(B),
G_{Q,\mathrm{ex},n}^{\mathrm{src,sat}}(B),
G_{Q,\mathrm{rev},n}^{\mathrm{src,sat}}(B),
G_{Q,\mathrm{nr},n}^{\mathrm{src,sat}}(B)
\right\}.
\end{equation}

A structural record \(R\) for which
\(B\notin\mathfrak B_{\mathrm{sat},n}(R)\) remains tagged to its structural
channel in the global ledger and is inactive at budget \(B\).  The
execution-level filtered residual \(L^{\Res,B}\) may retain that
same block below its exposure cost solely as the budget-bookkeeping case of
\cref{thm-exhaustive-residual-normal-form}; it is not part of
\(L^{\Res,\mathrm{sat}}\).  Thus saturation removes every represented
structural component from \(\Res\) without treating an over-budget record as
affordable.

\end{corollary}

\begin{corollary}[Budgeted computation success and source-charge bounds]
\label{cor:budgeted-computation-source-charge-bounds}

For every budget-\(B\) computation \(\mathcal A\), prospective selector
accountability gives

\begin{equation}
\label{eq-full-saturated-profile-success-bound}
\Pr\{\tau_{\mathcal A}\le H_B\}
\le
\operatorname{Enum}_{\nu}(\mathcal A,H_B)
+
H_B\operatorname{Sel}^{\mathrm{sat}}_n(\Psi(B,n)).
\end{equation}

The quantitative five-family accounting-charge form is

\begin{equation}
\label{eq-full-saturated-profile-source-charge-bound}
\Pr\{\tau_{\mathcal A}\le H_B\}
\le
eH_B(n)
\sum_{c\in\mathfrak J_{\Abs}^5}
\mathcal C_{I,\mathfrak F_5}^{(c)}
\!\left(
\Psi_5^{\mathrm{alg}}
\!\left(\Psi_{\mathrm{hit}}(B,n)\right)
\right).
\end{equation}

The displayed contextual envelopes account for the completed first-hit
record at its saturated cost.  A prospective-source interpretation of a
frontier term requires the represented comparison in
\cref{cor-contextual-structural-charge-transfer}.  \(\Lift\) and \(\Bdry\)
propagate represented sources, while the post-saturation \(\Res\) component
supplies no additional structural coordinate.

\end{corollary}

\begin{proof}[Proof of
\cref{thm-full-saturated-profile-decomposition,%
cor:prospective-saturated-profile-identities,%
cor:budgeted-computation-source-charge-bounds}]\phantomsection
\label{proof-full-saturated-profile-decomposition}

\Cref{thm-channel-exhaustiveness} gives the exhaustive ordered carrier
split.  The normal forms
\cref{thm-five-family-abs-provenance,thm-exhaustive-geom-normal-form,thm-exhaustive-five-family-geom-source-classification,thm-exhaustive-lift-normal-form,thm-exhaustive-bdry-normal-form,thm-exhaustive-residual-normal-form}
classify the five non-causal blocks, and
\cref{prop-derived-channels-create-no-structure} identifies \(\Caus\),
\(\Lift\), and \(\Bdry\) as derived uses of represented source data.
Budget-set assignment is well posed by
\cref{lem-derivability-well-posed}.  Therefore a represented structural
denotation is charged to its structural channel exactly at budgets belonging
to its upward-closed admissible budget set.
The reclassification and residual-correlation results
\cref{prop-family-uniform-residual-reclassification,thm-no-uniform-residual-correlation}
then leave in the post-saturation residual exactly the two orthogonal
remainders of \cref{thm-exhaustive-residual-normal-form}.  This proves the
ledger decomposition and \eqref{eq-full-saturated-profile-decomposition}.

\Cref{thm-affordable-geom-abs-normal-form} restricted by temporal source
admissibility gives the five branches in
\eqref{eq-full-saturated-master-decomposition}; taking their maximum is the
definition of the prospective master certificate.  The two-sided equivalence
\eqref{eq-source-master-combined-equivalence} then identifies that
maximum with the two-coordinate sum up to the displayed constant factor.
Finally, \cref{thm-prospective-selector-accountability} gives
\eqref{eq-full-saturated-profile-success-bound}, while
\cref{thm-prospective-selector-structural-charge} gives the direct
five-envelope accounting estimate
\eqref{eq-full-saturated-profile-source-charge-bound}.  No residual structural
source remains after least-cost saturation, and
\cref{cor-first-hit-accounting-not-source} places the completed first-hit
record in the separate accounting layer.

\end{proof}

\Cref{fig-full-saturated-profile-decomposition} combines the complete
classification of every channel in one cost-tagged tree.

\begin{figure}[p]
\centering
\resizebox{\textwidth}{!}{%
\begin{tikzpicture}[
  every node/.style={font=\scriptsize,align=center},
  root/.style={draw,rounded corners,fill=black!8,text width=38mm,
    minimum height=8mm},
  modality/.style={draw,rounded corners,text width=19mm,minimum height=8mm,
    font=\scriptsize\bfseries},
  leaf/.style={draw,rounded corners,text width=19mm,minimum height=7mm},
  small leaf/.style={draw,rounded corners,text width=15mm,minimum height=6mm},
  split edge/.style={draw,-{Latex[length=1.4mm]}},
  branch/.style={draw}
]
% Root and the six exhaustive carrier coordinates.
\node[root] (root) at (0,12.8)
  {full cost-tagged\\saturated operator ledger};
\coordinate (topSplit) at (0,12.15);
\node[modality,fill=green!11] (abs) at (-6.0,10.95) {\(\Abs\)};
\node[modality,fill=cyan!11] (geom) at (-3.6,10.95) {\(\Geom\)};
\node[modality,fill=orange!13] (caus) at (-1.2,10.95) {\(\Caus\)};
\node[modality,fill=yellow!15] (lift) at (1.2,10.95) {\(\Lift\)};
\node[modality,fill=red!8] (bdry) at (3.6,10.95) {\(\Bdry\)};
\node[modality,fill=gray!14] (res) at (6.0,10.95)
  {post-saturation\\\(J_{\mathrm{res}}\)};
\draw[split edge] (root) -- (topSplit);
\foreach \n in {abs,geom,caus,lift,bdry,res}
  \draw[split edge] (topSplit) -| (\n);

% Abs: the exhaustive provenance classification.
\coordinate (absSpineTop) at (-7.12,10.15);
\coordinate (absSpineBottom) at (-7.12,5.55);
\draw[branch] (abs) -- (absSpineTop) -- (absSpineBottom);
\node[small leaf,fill=green!7] (add) at (-6.0,9.85) {\(\mathsf{Add}\)};
\node[small leaf,fill=green!7] (mult) at (-6.0,8.75) {\(\mathsf{Mult}\)};
\node[small leaf,fill=green!7] (ord) at (-6.0,7.65) {\(\mathsf{Ord}\)};
\node[small leaf,fill=green!7] (sym) at (-6.0,6.55) {\(\mathsf{Sym}\)};
\node[small leaf,fill=green!7] (filt) at (-6.0,5.45) {\(\mathsf{Filt}\)};
\foreach \n in {add,mult,ord,sym,filt}
  \draw[split edge] (absSpineTop |- \n.west) -- (\n.west);

% Geom: the four active-geometry configurations.
\coordinate (geomSpineTop) at (-4.72,10.15);
\coordinate (geomSpineBottom) at (-4.72,6.55);
\draw[branch] (geom) -- (geomSpineTop) -- (geomSpineBottom);
\node[leaf,fill=cyan!7] (gambient) at (-3.6,9.75)
  {ambient\\\(q=\mathrm{id}\)};
\node[leaf,fill=cyan!7] (gexact) at (-3.6,8.55)
  {exact \(\Abs\)\\\(E_q=0\)};
\node[leaf,fill=cyan!7] (grev) at (-3.6,7.35)
  {quasi \(\Abs\)\\reversible};
\node[leaf,fill=cyan!7] (gnr) at (-3.6,6.15)
  {quasi \(\Abs\)\\general resolvent};
\foreach \n in {gambient,gexact,grev,gnr}
  \draw[split edge] (geomSpineTop |- \n.west) -- (\n.west);

% Caus: actual-current projection and the three authority regimes.
\node[leaf,fill=orange!8] (actual) at (-1.2,9.65)
  {actual-current\\Hodge projection};
\draw[split edge] (caus) -- (actual);
\coordinate (causSpineTop) at (-2.32,8.95);
\coordinate (causSpineBottom) at (-2.32,6.15);
\draw[branch] (actual) -- (causSpineTop) -- (causSpineBottom);
\node[small leaf,fill=orange!7] (free) at (-1.2,8.55) {free};
\node[small leaf,fill=orange!7] (imposed) at (-1.2,7.35) {imposed};
\node[small leaf,fill=orange!7] (hybrid) at (-1.2,6.15) {hybrid};
\foreach \n in {free,imposed,hybrid}
  \draw[split edge] (causSpineTop |- \n.west) -- (\n.west);

% Lift: fiber disintegration and the exact operator factorization.
\node[leaf,fill=yellow!10] (fiber) at (1.2,9.75)
  {fiber laws\\\(\kappa_x\subseteq\pi^{-1}(x)\)};
\node[leaf,fill=yellow!10] (lop) at (1.2,8.35)
  {operator\\factorization};
\coordinate (liftDataSpine) at (0.08,9.35);
\draw[branch] (lift) -- (liftDataSpine) -- (liftDataSpine |- lop.west);
\draw[split edge] (liftDataSpine |- fiber.west) -- (fiber.west);
\draw[split edge] (liftDataSpine |- lop.west) -- (lop.west);
\coordinate (liftSpineTop) at (0.08,7.65);
\coordinate (liftSpineBottom) at (0.08,4.75);
\draw[branch] (lop) -- (liftSpineTop) -- (liftSpineBottom);
\node[leaf,fill=blue!7] (ldesc) at (1.2,7.25)
  {descending\\\(\Abs\)};
\node[leaf,fill=yellow!10] (liface) at (1.2,6.05)
  {represented\\\(\Lift\) interface};
\node[leaf,fill=gray!10] (ldefect) at (1.2,4.85)
  {defect/\\complement\\reclassified};
\foreach \n in {ldesc,liface,ldefect}
  \draw[split edge] (liftSpineTop |- \n.west) -- (\n.west);

% Bdry: killing/value blocks and the two value provenances.
\coordinate (bdrySpineTop) at (2.48,10.15);
\coordinate (bdrySpineBottom) at (2.48,8.25);
\draw[branch] (bdry) -- (bdrySpineTop) -- (bdrySpineBottom);
\node[leaf,fill=red!6] (kill) at (3.6,9.65)
  {homogeneous loss\\\(-\kappa f\)};
\node[leaf,fill=red!6] (value) at (3.6,8.35)
  {value readout\\\(\sum_\zeta b_\zeta v(\zeta)\)};
\draw[split edge] (bdrySpineTop |- kill.west) -- (kill.west);
\draw[split edge] (bdrySpineTop |- value.west) -- (value.west);
\coordinate (valueSpineTop) at (4.72,7.65);
\coordinate (valueSpineBottom) at (4.72,6.05);
\draw[branch] (value) -- (valueSpineTop) -- (valueSpineBottom);
\node[leaf,fill=red!5] (declared) at (3.6,7.05)
  {declared\\terminal data};
\node[leaf,fill=red!5] (derived) at (3.6,5.75)
  {represented\\derived value};
\draw[split edge] (valueSpineTop |- declared.east) -- (declared.east);
\draw[split edge] (valueSpineTop |- derived.east) -- (derived.east);

% The saturated residual contains only the two non-structural remainders.
\coordinate (resSpineTop) at (4.88,10.15);
\coordinate (resSpineBottom) at (4.88,8.25);
\draw[branch] (res) -- (resSpineTop) -- (resSpineBottom);
\node[leaf,fill=gray!11] (rdir) at (6.0,9.65)
  {Hodge-orthogonal\\directed remainder};
\node[leaf,fill=gray!11] (rnl) at (6.0,8.25)
  {quotient/Lift-orthogonal\\nonlocal remainder};
\draw[split edge] (resSpineTop |- rdir.west) -- (rdir.west);
\draw[split edge] (resSpineTop |- rnl.west) -- (rnl.west);
\end{tikzpicture}%
}
\caption{Exhaustive classification of the full cost-tagged saturated
operator ledger.  The \(\Abs\) column applies
\cref{thm-five-family-abs-provenance}; the \(\Geom\) column applies
\cref{thm-exhaustive-geom-normal-form}, and each of its four support forms
carries the nonempty five-family source signature classified by
\cref{thm-exhaustive-five-family-geom-source-classification}; the \(\Caus\) column applies
\cref{def-caus-authority-envelope}; the \(\Lift\) column applies
\cref{thm-exhaustive-lift-normal-form}; and the \(\Bdry\) column applies
\cref{thm-exhaustive-bdry-normal-form}.
After every represented structural denotation has been charged at its least
saturated cost, \(J_{\mathrm{res}}\) consists exactly of the two
non-structural orthogonal remainders in
\cref{thm-exhaustive-residual-normal-form}.}
\label{fig-full-saturated-profile-decomposition}
\end{figure}
\clearpage

\section{Ambient-Geometry Accountability}
\label{sec:ambient-geometry-accountability}

\begin{theorem}[Saturated ambient-geometry accountability after quotient
exhaustion]
\label{thm-saturated-ambient-geom-accountability}

Fix a represented task family, the free Caus authority envelope, and a
saturated budget \(B\).  Restrict the ambient record family
\(\mathcal G^{X,\mathrm{free}}_{n,\le B}\) of
\cref{def-geom-extraction} to temporally admissible prospective source
records.  Then

\begin{equation}
\label{eq-prospective-ambient-geom-exact-profile}
G_{X,n}^{\mathrm{src,sat}}(B)
=
\sup_{(\mathfrak N,\Phi,D,J,\mathfrak K)}
M_\Phi
C_{\Phi,\mathfrak K}^{\mathrm{auth}}(D)
\frac{R_T(D)}{1+\rho_\Phi(D)},
\end{equation}

where the supremum is over that single restricted family and has value zero
when the family is empty.  The same zero convention applies to the three
factorwise suprema below.  Every factor belongs to \([0,1]\).  Consequently,

\begin{align}
\label{eq-ambient-geom-factorwise-upper-bound}
G_{X,n}^{\mathrm{src,sat}}(B)
\le
\min\Biggl\{&
\sup_{(\mathfrak N,\Phi,D,J,\mathfrak K)}
C_{\Phi,\mathfrak K}^{\mathrm{auth}}(D),
\quad
\sup_{(\mathfrak N,\Phi,D,J,\mathfrak K)}R_T(D),
\notag\\[-1mm]
&
\sup_{(\mathfrak N,\Phi,D,J,\mathfrak K)}
\frac{1}{1+\rho_\Phi(D)}
\Biggr\}.
\end{align}

More sharply, for every \(0<\varepsilon<1\),

\begin{equation}
\label{eq-ambient-geom-epsilon-witness}
G_{X,n}^{\mathrm{src,sat}}(B)>\varepsilon
\end{equation}

holds only if one family-uniform represented record of complete saturated
cost at most \(B\) satisfies simultaneously

\begin{equation}
\label{eq-ambient-geom-four-witness-conditions}
M_\Phi=1,
\qquad
C_{\Phi,\mathfrak K}^{\mathrm{auth}}(D)>\varepsilon,
\qquad
R_T(D)>\varepsilon,
\qquad
\rho_\Phi(D)<\varepsilon^{-1}-1.
\end{equation}

Conversely, if no record in the saturated ambient family satisfies all four
inequalities in \eqref{eq-ambient-geom-four-witness-conditions}, then

\begin{equation}
\label{eq-ambient-geom-no-witness-bound}
G_{X,n}^{\mathrm{src,sat}}(B)\le\varepsilon.
\end{equation}

The witness in \eqref{eq-ambient-geom-four-witness-conditions} has identity
quotient support.  Its geometry therefore acts on the full presented state
space, and every step used by its Dirichlet form lies in the exposed
comparability relation \(E_\Phi\).  The condition \(M_\Phi=1\) supplies the represented transfer-class
Poincare or conductance comparison; positivity of
\(C_{\Phi,\mathfrak K}^{\mathrm{auth}}(D)\) supplies the local-consistency
gate and an authorized descending current for the same selected generator;
\(R_T(D)>\varepsilon\) supplies target reach; and the last inequality is the
exact regularity bound appearing in the ambient visibility coordinate.

Now suppose that, at the same budget, the four prospective nonambient
coordinates satisfy

\begin{equation}
\label{eq-prospective-quotient-branch-envelope}
\max\!\left\{
m_n^{\mathrm{src,sat}}(B),
G_{Q,\mathrm{ex},n}^{\mathrm{src,sat}}(B),
G_{Q,\mathrm{rev},n}^{\mathrm{src,sat}}(B),
G_{Q,\mathrm{nr},n}^{\mathrm{src,sat}}(B)
\right\}
\le\alpha_n
\qquad (\alpha_n\ge0).
\end{equation}

Then the full prospective source profiles obey

\begin{align}
\label{eq-no-abs-ambient-master-reduction}
\mathfrak U_n^{\mathrm{src,sat}}(B)
&\le
\max\!\left\{
G_{X,n}^{\mathrm{src,sat}}(B),\alpha_n
\right\},\\
\label{eq-no-abs-ambient-combined-reduction}
I_n^{\mathrm{src,sat}}(B)
&\le
2\max\!\left\{
G_{X,n}^{\mathrm{src,sat}}(B),\alpha_n
\right\}.
\end{align}

Thus quotient exhaustion leaves exactly one possible affordable structural
source: an identity-supported ambient record satisfying the joint local-gap,
authorized-descent, target-reach, and regularity conditions above.  A
represented transformation that preserves identity support is already in the
supremum in \eqref{eq-prospective-ambient-geom-exact-profile}.  A
transformation that creates nonidentity fibers is charged to one of the four
coordinates in \eqref{eq-prospective-quotient-branch-envelope}.

\end{theorem}

\begin{proof}

Temporal source restriction changes only the admissible record family in the
supremum of \cref{def-geom-extraction}.  The visibility of each remaining
identity-supported record is, by definition,

\[
V^X_{\Phi,D,J,\mathfrak K}
=
M_\Phi
C_{\Phi,\mathfrak K}^{\mathrm{auth}}(D)
R_T(D)
\frac{1}{1+\rho_\Phi(D)}.
\]

Taking the supremum proves
\eqref{eq-prospective-ambient-geom-exact-profile}.  The mixing indicator is
binary, the authority-compatible alignment and target-reach factors lie in
\([0,1]\), and \(\rho_\Phi(D)\ge0\).  Hence the product is bounded by each
of its three nonbinary factors.  Taking suprema over the same record family
proves \eqref{eq-ambient-geom-factorwise-upper-bound}.

Assume \eqref{eq-ambient-geom-epsilon-witness}.  By the defining property of
a supremum, some record in the restricted family has visibility greater than
\(\varepsilon\).  Its mixing indicator cannot be zero.  Since the remaining
three factors lie in \([0,1]\), their product can exceed \(\varepsilon\)
only when each factor exceeds \(\varepsilon\).  The regularity inequality is
equivalent to

\[
\frac{1}{1+\rho_\Phi(D)}>\varepsilon
\quad\Longleftrightarrow\quad
\rho_\Phi(D)<\varepsilon^{-1}-1.
\]

This proves \eqref{eq-ambient-geom-four-witness-conditions}.  Conversely, if
no record satisfies all four conditions, then every record has mixing
indicator zero or at least one of the remaining factors at most
\(\varepsilon\).  Its visibility is therefore at most \(\varepsilon\).
Taking the supremum proves \eqref{eq-ambient-geom-no-witness-bound}.

The interpretation of the four conditions follows directly from
\cref{thm-valid-geom-source-characterization,def-authority-compatible-caus-alignment,def-geom-extraction}.
In particular, the authority-compatible factor is positive only when the
local comparison and current provenance belong to the same selected
generator.  By \cref{thm-caus-authority-no-substitution}, a separately chosen
target-oriented current cannot be inserted into this record.  The master
normal form of \cref{thm-affordable-geom-abs-normal-form} charges every
represented change of variables, join, potential, conductance, and current.
It assigns the result to the identity-supported ambient family precisely when
the terminal quotient is the identity; otherwise it assigns the record to an
exact or compensated nonidentity quotient family.

Finally, apply the prospective five-branch identity
\eqref{eq-full-saturated-master-decomposition}.  Under
\eqref{eq-prospective-quotient-branch-envelope}, four entries in that maximum
are at most \(\alpha_n\), leaving only
\(G_{X,n}^{\mathrm{src,sat}}(B)\).  This proves
\eqref{eq-no-abs-ambient-master-reduction}.  The prospective combined-profile
comparison \eqref{eq-source-master-combined-equivalence} gives
\eqref{eq-no-abs-ambient-combined-reduction}.

\end{proof}

\begin{theorem}[Represented-current admission and no semantic target support]
\label{thm-represented-current-admission-no-semantic-support}

Fix a presented finite task family, a saturated budget \(B\), and the
prospective ambient family of
\cref{thm-saturated-ambient-geom-accountability}.  A current contributes to
\(G_{X,n}^{\mathrm{src,sat}}(B)\) only as the current entry of one complete
record

\[
\mathfrak g=(\mathfrak N,\Phi,D,J,\mathfrak K)
\in\mathcal G^{X,\mathrm{free}}_{n,\le B}
\]

whose represented derivation satisfies all of the following conditions.  In
this record, \(J=J_{\mathfrak K}\): the current entry is the current identified
and authenticated by the authority record \(\mathfrak K\).

\begin{enumerate}[label=\textup{(\roman*)}]
\item the normalized kernel, edge relation, conductance, potential, selected
      current, and authority comparison are derived from the declared public
      presentation and presentation-admissible oracle evaluators;
\item their complete derivation and evaluation cost is at most \(B\);
\item the record is available before the transition whose progress it is
      invoked to explain; and
\item \(\mathfrak K\) authenticates the current of the selected generator, or
      authenticates an auxiliary geometric current together with the
      same-generator local, drift, resolvent, or carrier comparison required
      by \cref{def-authorized-caus-use-current,def-authority-compatible-caus-alignment}.
\end{enumerate}

If any condition fails, that denotation is absent from the prospective
ambient record family.  Equivalently, under the zero-extension convention for
incomplete records, a missing current identity, target polarity, dynamic
transport, or same-generator comparison gives

\begin{equation}
\label{eq-no-semantic-current-zero-authorized-alignment}
\operatorname{Align}_{\Caus}^{\mathrm{auth}}(\mathfrak K,D)=0,
\qquad
V^X_{\Phi,D,J,\mathfrak K}=0.
\end{equation}

Let \(S_I\subseteq E_\Phi\) be a represented edge sector such that
\(\nabla_\Phi D\) is supported on \(S_I\).  For every admitted record with
nonzero current,

\begin{equation}
\label{eq-represented-current-supported-alignment-bound}
V^X_{\Phi,D,J,\mathfrak K}
\le
\operatorname{Align}_{\Caus}^{\mathrm{auth}}(\mathfrak K,D)
\le
\frac{\|\mathbf1_{S_I}J_{\mathfrak K}\|_E^2}
     {\|J_{\mathfrak K}\|_E^2}
\le1.
\end{equation}

The ratio in
\cref{eq-represented-current-supported-alignment-bound} is an analytic
evaluation of an already admitted current.  It neither constructs the
current nor admits an extensionally specified current into the saturated
profile.  Evaluating the predicate \(\mathbf1_{S_I}(x,z)\) likewise does not
construct a normalized kernel or a sampler supported on \(S_I\); such a
kernel contributes only with its represented normalization, sampling
procedure, random-bit use, selected-generator identity, and complete cost.

For a randomized family-uniform derivation, let \(Q_{I,\omega}\) be the
indicator that its output record satisfies the four admission conditions and
all remaining Geom qualification gates on instance \(I\) and seed \(\omega\).
The qualified-record convention gives the pointwise bounds

\begin{equation}
\label{eq-current-qualification-indicator-bound}
0
\le
Q_{I,\omega}V^X_{\Phi,D,J,\mathfrak K}
\le
Q_{I,\omega}
\frac{\|\mathbf1_{S_I}J_{\mathfrak K}\|_E^2}
     {\|J_{\mathfrak K}\|_E^2}
\le
Q_{I,\omega},
\end{equation}

with the middle ratio set to zero for a zero current.  Consequently, for
every monotone positively homogeneous family functional \(\mathfrak M_n\),

\begin{equation}
\label{eq-current-qualification-family-bound}
\mathfrak M_n\!\left(
I\longmapsto
\mathbb E_\omega[Q_{I,\omega}V^X_{\Phi,D,J,\mathfrak K}]
\right)
\le
\mathfrak M_n\!\left(
I\longmapsto\Pr_\omega\{Q_{I,\omega}=1\}
\right).
\end{equation}

If qualification entails a represented success event \(\mathcal E_{I,\omega}\),
the right-hand side is at most the same family functional applied to
\(\Pr_\omega(\mathcal E_{I,\omega})\).  Thus a current constructed from a
candidate target, decoded center, learned ranking, or selector is charged
through the complete derivation of that record and through the probability
that the resulting record is target-qualified.

An oracle or advice string that supplies the current without a represented
evaluator in the declared presentation changes the presented family as in
\cref{thm-oracle-admissibility-presentation-separation}.  A current produced
from later verifier flags is an accounting record rather than a prospective
source by \cref{def-pre-hit-temporally-admissible-source}.

\end{theorem}

\begin{proof}

The ambient profile in
\cref{eq-prospective-ambient-geom-exact-profile} is a supremum over complete
represented records in
\(\mathcal G^{X,\mathrm{free}}_{n,\le B}\).  Membership in that family gives
conditions \textup{(i)} and \textup{(ii)} and identifies the current entry
with the current authenticated by \(\mathfrak K\), while
\cref{def-pre-hit-temporally-admissible-source} gives condition
\textup{(iii)}.  The current and comparison clauses in
\cref{def-authorized-caus-use-current,def-authority-compatible-caus-alignment}
give condition \textup{(iv)} and set the authorized alignment to zero when
its record is missing.  The complete visibility product then proves
\cref{eq-no-semantic-current-zero-authorized-alignment}.

Because \(\nabla_\Phi D\) is supported on \(S_I\),

\[
\langle J_{\mathfrak K},-\nabla_\Phi D\rangle_E
=
\langle \mathbf1_{S_I}J_{\mathfrak K},-\nabla_\Phi D\rangle_E.
\]

For nonzero \(J_{\mathfrak K}\) and nonzero \(\nabla_\Phi D\),
\cref{def-quantitative-caus-alignment,def-authority-compatible-caus-alignment}
identify the authorized alignment with the normalized positive pairing below.
Cauchy--Schwarz on the represented weighted edge space gives

\begin{align*}
\operatorname{Align}_{\Caus}^{\mathrm{auth}}(\mathfrak K,D)
&=
\frac{
\bigl(\langle\mathbf1_{S_I}J_{\mathfrak K},
             -\nabla_\Phi D\rangle_E\bigr)_+^2
}{
\|J_{\mathfrak K}\|_E^2\|\nabla_\Phi D\|_E^2
}\\
&\le
\frac{
\|\mathbf1_{S_I}J_{\mathfrak K}\|_E^2
\|\nabla_\Phi D\|_E^2
}{
\|J_{\mathfrak K}\|_E^2\|\nabla_\Phi D\|_E^2
}
=
\frac{\|\mathbf1_{S_I}J_{\mathfrak K}\|_E^2}
     {\|J_{\mathfrak K}\|_E^2}.
\end{align*}

If \(\nabla_\Phi D=0\), the alignment is zero by definition and the same
inequality holds.  This proves the second inequality in
\cref{eq-represented-current-supported-alignment-bound}; the first follows
because every other visibility factor lies in \([0,1]\).  Restriction of a
weighted norm to an edge subset gives the last inequality.

The budgeted derivability closure contains outputs of represented
constructors, not arbitrary extensional tables.  The no-substitution result
\cref{thm-caus-authority-no-substitution} prevents replacement of the
selected current by another target-oriented vector, and
\cref{thm-affordable-composition-no-creation} retains the complete ancestry
of every represented postprocessing.  A Boolean sector predicate supplies
an evaluator only; normalization and sampling require the additional
represented operations listed in the theorem.  This proves the admission
and predicate statements.

Finally, multiply the pointwise inequalities by the qualification indicator.
Every visibility and support fraction lies in \([0,1]\), which proves
\cref{eq-current-qualification-indicator-bound}.  In particular, pointwise in
the instance,

\[
\mathbb E_\omega
\bigl[Q_{I,\omega}V^X_{\Phi,D,J,\mathfrak K}\bigr]
\le
\mathbb E_\omega[Q_{I,\omega}]
=
\Pr_\omega\{Q_{I,\omega}=1\}.
\]

Monotonicity of \(\mathfrak M_n\) proves
\cref{eq-current-qualification-family-bound}; positive homogeneity is the
declared profile convention and is preserved.  The success-event conclusion
is monotonicity under
\(Q_{I,\omega}\le\mathbf1_{\mathcal E_{I,\omega}}\).  Oracle separation and
temporal classification are the two cited framework results.

\end{proof}

\section{Authorized Perturbations and Obstruction Criteria}
\label{sec:authorized-perturbation-obstructions}

\begin{definition}[Authorized ambient perturbation witness]
\label{def-authorized-ambient-perturbation-witness}

Let
\[
\mathfrak g=(\mathfrak N,\Phi,D,J,\mathfrak K)
\in\mathcal G^{X,\mathrm{free}}_{n,\le B}
\]
be a temporally admissible prospective ambient record of
\cref{def-geom-extraction}, with identity quotient support, and write
\[
\mathfrak N=(\widetilde L,\zeta,c,P,L,\lambda)
\]
for its charged uniformization record.  Its \emph{authorized ambient
perturbation kernel} is the edge-restricted sample-and-hold kernel
\begin{equation}
\label{eq-authorized-ambient-perturbation-kernel}
\begin{aligned}
P_\Phi(x,z)
&=
\mathbf 1_{\{(x,z)\in E_\Phi\}}P(x,z),
&&z\ne x,\\
P_\Phi(x,x)
&=
P(x,x)
+
\sum_{\substack{z\ne x\\(x,z)\notin E_\Phi}}P(x,z).
\end{aligned}
\end{equation}
The killing mass of \(P\) is unchanged, so \(P_\Phi\) is sub-Markov, and it
is Markov whenever \(P\) is Markov.  Operationally, one samples one
\(P\)-transition and replaces an outcome outside \(E_\Phi\) by a holding
step.  The construction uses one represented kernel step and one represented
edge-membership test and is therefore charged within the same
class-preserving derivation.

The \emph{authorized ambient perturbation data} are the existing represented
data together with this derived kernel:
\[
\operatorname{Pert}(\mathfrak g)
=
\bigl(P,P_\Phi,E_\Phi,c_\Phi,D,J_{\mathfrak K},\mathfrak K\bigr).
\]
The tuple retains the complete derivation and cost of \(\mathfrak g\).
The normalized kernel \(P\) is the selected public transition kernel;
\(P_\Phi\) is its canonical exposed local perturbation; \(E_\Phi\) and
\(c_\Phi\) are the geometric comparison data; and
\(J_{\mathfrak K}\) is either the structural-gradient projection of the
actual selected current or the auxiliary geometric current accompanied by
the represented comparison required in
\cref{def-authorized-caus-use-current,def-authority-compatible-caus-alignment}.

Define the represented descending-edge predicate
\[
E_\downarrow(\mathfrak g)
=
\left\{
(x,z)\in E_\Phi:
J_{\mathfrak K}(x,z)\bigl(D(x)-D(z)\bigr)>0
\right\}.
\]
For \(0<\varepsilon<1\), the perturbation witness is
\emph{\(\varepsilon\)-effective} when
\begin{equation}
\label{eq-effective-ambient-perturbation-gates}
M_\Phi=1,\qquad
\chi_\Phi^{\mathrm{auth}}(D)=1,\qquad
\operatorname{Align}_{\Caus}^{\mathrm{auth}}(\mathfrak K,D)>\varepsilon,
\qquad
R_T(D)>\varepsilon,\qquad
\frac1{1+\rho_\Phi(D)}>\varepsilon.
\end{equation}

This is a predicate on the existing ambient certificate record.  It defines
neither a new structural channel nor an additional profile coordinate.
Raw displacement in a chosen presentation coordinate is not part of the
predicate: locality means support on the represented relation \(E_\Phi\),
with the selected-kernel comparison and its cost retained.

\end{definition}

\begin{theorem}[Perturbation characterization and obstruction after quotient
exhaustion]
\label{thm-ambient-perturbation-characterization-obstruction}

Fix a represented task family, saturated budget \(B\), and
\(0<\varepsilon<1\).  Every dynamically qualified, temporally admissible
identity-supported Geom/Caus source is an authorized ambient perturbation
witness in the sense of
\cref{def-authorized-ambient-perturbation-witness}.  Conversely, an authorized
ambient perturbation record satisfying the five certificate conditions of
\cref{thm-valid-geom-source-characterization} is precisely an
identity-supported ambient source record in
\(\mathcal G^{X,\mathrm{free}}_{n,\le B}\).

If
\begin{equation}
\label{eq-ambient-visibility-extracts-effective-perturbation}
V^X_{\Phi,D,J,\mathfrak K}>\varepsilon,
\end{equation}
then \(\operatorname{Pert}(\mathfrak g)\) is \(\varepsilon\)-effective.  Its authorized
current satisfies
\begin{equation}
\label{eq-effective-perturbation-gradient-alignment}
\left\langle
J_{\mathfrak K},-\nabla_\Phi D
\right\rangle_E
>
\sqrt{\varepsilon}\,
\|J_{\mathfrak K}\|_E
\|\nabla_\Phi D\|_E,
\end{equation}
and its descending edges carry at least the same positive pairing:
\begin{equation}
\label{eq-effective-perturbation-descending-edge-mass}
\left\langle
\mathbf 1_{E_\downarrow(\mathfrak g)}J_{\mathfrak K},
-\nabla_\Phi D
\right\rangle_E
\ge
\left\langle
J_{\mathfrak K},-\nabla_\Phi D
\right\rangle_E.
\end{equation}
In particular, the perturbation is represented before use; its counted
geometric current and potential differences are supported on exposed
comparisons; and its descent is authorized by the same selected generator
that supplies the local-consistency, drift, resolvent, or carrier comparison.

Conversely, every \(\varepsilon\)-effective perturbation witness satisfies
\begin{equation}
\label{eq-effective-perturbation-visibility-lower-bound}
V^X_{\Phi,D,J,\mathfrak K}>\varepsilon^3.
\end{equation}

More generally, let \(0\le\delta<1\).  If every temporally admissible ambient
record of cost at most \(B\) satisfies at least one of
\begin{equation}
\label{eq-ambient-perturbation-obstruction-cover}
\begin{gathered}
M_\Phi=0,\qquad
\chi_\Phi^{\mathrm{auth}}(D)=0,\qquad
\operatorname{Align}_{\Caus}^{\mathrm{auth}}(\mathfrak K,D)\le\delta,\\
R_T(D)\le\delta,\qquad
\frac1{1+\rho_\Phi(D)}\le\delta,
\end{gathered}
\end{equation}
then
\begin{equation}
\label{eq-ambient-perturbation-obstruction-bound}
G_{X,n}^{\mathrm{src,sat}}(B)\le\delta.
\end{equation}

If, at the same budget, the nonambient quotient envelope in
\cref{eq-prospective-quotient-branch-envelope} is at most \(\alpha_n\), then
\begin{align}
\label{eq-perturbation-obstruction-master-bound}
\mathfrak U_n^{\mathrm{src,sat}}(B)
&\le \max\{\alpha_n,\delta\},\\
\label{eq-perturbation-obstruction-combined-bound}
I_n^{\mathrm{src,sat}}(B)
&\le 2\max\{\alpha_n,\delta\}.
\end{align}

Thus quotient exhaustion and the perturbation obstruction cover jointly
bound the complete prospective structural profile.  A represented
transformation with nonidentity composite fibers is assigned to the exact or
compensated quotient envelope.  A transformation with identity support,
including one with large displacement in raw coordinates, remains in the
ambient family and must satisfy
\cref{eq-effective-ambient-perturbation-gates} to contribute above the
corresponding scale.

\end{theorem}

\begin{proof}

The ambient branch in
\cref{thm-affordable-geom-abs-normal-form,thm-exhaustive-geom-normal-form}
has identity quotient support and an active geometry slot.  Its charged
normalized kernel is \(\mathfrak N\), its local relation and conductance are
\(\Phi\), and its current provenance is \(\mathfrak K\).  These are exactly
the entries of \(\operatorname{Pert}(\mathfrak g)\).  The converse follows from the
necessary-and-sufficient certificate conditions in
\cref{thm-valid-geom-source-characterization}, together with the authority
condition of \cref{def-authority-compatible-caus-alignment}.  This proves the
record characterization.

Assume \eqref{eq-ambient-visibility-extracts-effective-perturbation}.  By
\cref{def-geom-extraction},
\[
V^X_{\Phi,D,J,\mathfrak K}
=
M_\Phi\,
\chi_\Phi^{\mathrm{auth}}(D)\,
\operatorname{Align}_{\Caus}^{\mathrm{auth}}(\mathfrak K,D)\,
R_T(D)\,
\frac1{1+\rho_\Phi(D)}.
\]
The mixing and consistency gates are binary and the other three factors lie
in \([0,1]\).  A product greater than \(\varepsilon\) forces both binary gates
to equal one and each remaining factor to exceed \(\varepsilon\).  Hence
\(\operatorname{Pert}(\mathfrak g)\) is \(\varepsilon\)-effective.

With the authority gate equal to one,
\cref{def-quantitative-caus-alignment,def-authority-compatible-caus-alignment}
give
\[
\operatorname{Align}_{\Caus}^{\mathrm{auth}}(\mathfrak K,D)
=
\frac{
\bigl(
\langle J_{\mathfrak K},-\nabla_\Phi D\rangle_E
\bigr)_+^2
}{
\|J_{\mathfrak K}\|_E^2\|\nabla_\Phi D\|_E^2
}.
\]
The strict lower bound by \(\varepsilon\) proves
\eqref{eq-effective-perturbation-gradient-alignment}.  On every edge outside
\(E_\downarrow(\mathfrak g)\), the summand in the edge pairing is
nonpositive.  Removing those summands can only increase the pairing, proving
\eqref{eq-effective-perturbation-descending-edge-mass}.  Temporal
availability, represented support, complete cost, and same-generator
authority are the defining properties of
\cref{def-pre-hit-temporally-admissible-source,def-authorized-caus-use-current,def-authority-compatible-caus-alignment}.

Conversely, under \eqref{eq-effective-ambient-perturbation-gates}, the two
binary gates equal one and the alignment, reach, and regularity factors each
exceed \(\varepsilon\).  Their product is greater than
\(\varepsilon^3\), proving
\eqref{eq-effective-perturbation-visibility-lower-bound}.

Now assume \eqref{eq-ambient-perturbation-obstruction-cover}.  Every ambient
visibility has a zero binary factor or one of its three quantitative factors
is at most \(\delta\).  Since all remaining factors lie in \([0,1]\), every
record has visibility at most \(\delta\).  Taking the supremum proves
\eqref{eq-ambient-perturbation-obstruction-bound}.  Finally,
\cref{eq-no-abs-ambient-master-reduction,eq-no-abs-ambient-combined-reduction}
with
\(G_{X,n}^{\mathrm{src,sat}}(B)\le\delta\)
gives
\cref{eq-perturbation-obstruction-master-bound,eq-perturbation-obstruction-combined-bound}.

\end{proof}

\begin{corollary}[Polynomial perturbation obstruction criterion]
\label{cor-polynomial-ambient-perturbation-obstruction}

Let \(B=B(n)\) be a budget schedule in the polynomial resource envelope, and
suppose the nonambient quotient envelope is negligible at \(B(n)\).  The
complete prospective source profile is negligible whenever, for every fixed
\(k\ge1\) and all sufficiently large \(n\), every temporally admissible
identity-supported perturbation record of cost at most \(B(n)\) satisfies at
least one of
\begin{equation}
\label{eq-polynomial-ambient-perturbation-obstruction}
\begin{gathered}
M_\Phi=0,\qquad
\chi_\Phi^{\mathrm{auth}}(D)=0,\qquad
\operatorname{Align}_{\Caus}^{\mathrm{auth}}(\mathfrak K,D)\le n^{-k},\\
R_T(D)\le n^{-k},\qquad
\frac1{1+\rho_\Phi(D)}\le n^{-k}.
\end{gathered}
\end{equation}

Equivalently, a nonnegligible prospective source profile produces one
family-uniform affordable perturbation record for which all five displayed
obstructions fail at an inverse-polynomial scale.  The record includes the
selected kernel, exposed perturbation relation, potential, current, authority
proof, and complete represented derivation.

\end{corollary}

\begin{proof}

Apply
\cref{thm-ambient-perturbation-characterization-obstruction} with
\(\delta=n^{-k}\).  The obstruction cover gives
\(G_{X,n}^{\mathrm{src,sat}}(B(n))\le n^{-k}\).  Since \(k\) is arbitrary,
the ambient profile is negligible.  The quotient envelope is negligible by
hypothesis, so
\cref{eq-perturbation-obstruction-master-bound,eq-perturbation-obstruction-combined-bound}
proves negligibility of the master and combined prospective profiles.  The
equivalent formulation is the contrapositive, using
\cref{eq-ambient-visibility-extracts-effective-perturbation}.

\end{proof}

\begin{corollary}[Profile-level criterion for absence of an affordable
ambient geometry]
\label{cor-saturated-ambient-geom-negligibility}

Let \(B=B(n)\) be a declared budget schedule.  The ambient prospective
coordinate \(G_{X,n}^{\mathrm{src,sat}}(B(n))\) is negligible if and only if,
for every fixed \(k\ge1\) and all sufficiently large \(n\), the saturated
ambient family contains no temporally admissible record satisfying

\begin{equation}
\label{eq-ambient-geom-inverse-polynomial-witness}
M_\Phi=1,
\qquad
C_{\Phi,\mathfrak K}^{\mathrm{auth}}(D)>n^{-k},
\qquad
R_T(D)>n^{-k},
\qquad
\frac{1}{1+\rho_\Phi(D)}>n^{-k}.
\end{equation}

If the quotient-branch envelope in
\eqref{eq-prospective-quotient-branch-envelope} is also negligible, then the
complete prospective source profile is negligible.  This is one saturated
profile evaluation: its quantifier ranges over the single family-uniform
derivation class fixed by the presentation and budget, rather than over a
separately enumerated collection of algorithms.

\end{corollary}

\begin{proof}

Suppose first that the ambient coordinate is negligible.  Fix \(k\).  For all
sufficiently large \(n\), negligibility gives

\[
G_{X,n}^{\mathrm{src,sat}}(B(n))
\le n^{-(3k+1)}.
\]

A record satisfying \eqref{eq-ambient-geom-inverse-polynomial-witness} would
have visibility strictly greater than \(n^{-3k}\), a contradiction.

Conversely, assume the displayed witness is absent for every fixed \(k\).
For a chosen \(k\) and all sufficiently large \(n\), every ambient record has
mixing indicator zero or at least one of the remaining three factors at most
\(n^{-k}\).  Its visibility is at most \(n^{-k}\), and so is the supremum.
Since \(k\) is arbitrary, the ambient coordinate is negligible.  The final
assertion follows from
\cref{eq-no-abs-ambient-master-reduction,eq-no-abs-ambient-combined-reduction}.

\end{proof}

For a same-carrier statistical presentation,
\cref{lem-learned-geom-target-reach-projection} bounds this complete ambient
supremum by the saturated learned projection envelope at the
class-preserving enlarged cost.  The represented-readout and selector
refinements are
\cref{prop-noise-corrected-actionable-geom-extraction,thm-saturated-geometry-learning-selection-accountability}.

\begin{lemma}[Combined geometric and quotient structural comparison bound]\label{thm-combined-extraction-bound}

Let \(\mathcal S_B\) be a class of budget-admissible represented
structural certificates. Suppose that each \(R\in\mathcal S_B\) comes
with charged comparison data and two nonnegative visibility terms
\(V_R^{\Abs}\) and \(V_R^{\Geom}\) satisfying

\[
\Gamma_T(R)
\le C_{\mathcal B}(n)
\bigl(V_R^{\Abs}+V_R^{\Geom}\bigr),
\qquad
V_R^{\Abs}\le m_n^{\mathrm{sat}}(B),
\qquad
V_R^{\Geom}\le G_n^{\mathrm{sat}}(B).
\]

Define
\(\Gamma_T^{\mathrm{str}}(B)=\sup_{R\in\mathcal S_B}\Gamma_T(R)\),
with value zero when \(\mathcal S_B\) is empty. Then

\[
\Gamma_T^{\mathrm{str}}(B)\le C_{\mathcal B}(n)I_n^{\mathrm{sat}}(B),
\qquad
I_n^{\mathrm{sat}}(B)=m_n^{\mathrm{sat}}(B)+G_n^{\mathrm{sat}}(B).
\]

Equivalently,

\[
\Gamma_T^{\mathrm{str}}(B)
\le2C_{\mathcal B}(n)\mathfrak U_n^{\mathrm{sat}}(B).
\]

Here \(C_{\mathcal B}\) is the transfer-class comparison factor supplied
by the declared structural comparison data; for
\(\mathcal B_{\mathrm{poly}}\) it is polynomial. This structural
comparison bounds target contact; the construction of a target-reaching
procedure requires the corresponding arrival theorem.

\end{lemma}

\begin{proof}\phantomsection\label{proof-combined-extraction-bound}

For every \(R\in\mathcal S_B\), the three assumed inequalities give

\[
\Gamma_T(R)
\le C_{\mathcal B}(n)
\bigl(m_n^{\mathrm{sat}}(B)+G_n^{\mathrm{sat}}(B)\bigr).
\]

Taking the supremum over \(R\) proves the first displayed bound. Since
\(a+b\le2\max\{a,b\}\) for \(a,b\ge0\), the second follows. The theorem
therefore records exactly what the charged comparison data prove; the
\PartII decomposition by itself does not supply the first assumed
inequality. Full-process target arrival, including interaction with
\(\Res\), is controlled instead by
\eqref{eq-universal-combined-profile-bound}.

\end{proof}

\begin{proposition}[Structural extraction identifies a component]\label{prop-structural-extraction-realizes-channel}

If a term in the supremum defining
\(\mathfrak U_n^{\mathrm{sat}}(B)\) has transfer-class size, then the
corresponding admissible quotient--geometry certificate is an
exposed target-aligned component at saturated cost at most \(B\). This
proposition identifies the exposed component; target-arrival consequences are
supplied by the applicable arrival theorem rather than by the visibility
definition.

\end{proposition}

\begin{proof}\phantomsection\label{proof-structural-extraction-realizes-channel}

For a null geometry slot, a transfer-class-size term is
\(V_q=A_q\|P_q^{\mathrm{amb}}y\|^2/\|y\|^2\), so the exact quotient
subspace contains target mass at the declared threshold after imposing
the \PartI utility condition. For an active slot, a transfer-class-size
term in \eqref{eq-supported-geom-visibility} supplies represented target
reach, controlled regularity, and quantitative descending-current
alignment. A non-lumpable quotient additionally has an explicit
transfer-class stability factor. These are structural visibility
properties. A pointwise drift certificate
gives the hitting-time conclusion of
\Cref{thm-caus-drift-survives-defect}; other arrival conclusions
require the comparison theorem declared by the certificate.

\end{proof}

\begin{example}[SAT extraction]\label{ex-combined-sat}

For XOR-SAT, \(m_n^{\mathrm{sat}}(B)\) is inverse-polynomial at the
polynomial budget needed to represent the public syndrome quotient, so
the Abs term in \(I_n^{\mathrm{sat}}(B)\) is non-negligible. For a random
planted-looking SAT instance with no exposed quotient and a defect
landscape whose Hamming geometry has off-target obstructions, the
corresponding raw-basis and Hamming-geometry fragment profiles are
negligible.  The full \(m_n^{\mathrm{sat}}(B)\) and
\(G_n^{\mathrm{sat}}(B)\) have that conclusion only after a proved
uniform cover of the complete saturated kernel and geometry families.

\end{example}

\section{The Saturated Cost
Invariant}\label{the-saturated-cost-invariant}
\label{sec:saturated-cost-invariant}

The central object is the first saturated public cost budget at which
the combined extraction functional reaches the declared transfer-class
threshold. The log-cost degree \(d^*\) is a convenient coordinate for
this budget; it is not raw Walsh degree.

The saturated cost invariant is the first budget at which the combined profile reaches the declared threshold.
\begin{definition}[Saturated cost invariant]\label{def-degree-invariant}

Fix a structural-visibility threshold
\(\eta_{\mathfrak C}(n)\succeq 1/\beta_{\mathfrak C}(n)\) and one
admissible scalarization \(\kappa:C\to[0,\infty]\) from
\Cref{def-budgeted-derivability-closure}. For a scalar
\(s\ge0\), interpret \(I_{n,\kappa}^{\mathrm{sat}}(s)\) as the supremum
over complete represented records having
\(\kappa(\operatorname{Cost})\le s\). In the
polynomial instance this is \(n^{-c}\). Define the first extraction
budget

\[
B^*_\eta(n)=
\inf\{s\ge0:I_{n,\kappa}^{\mathrm{sat}}(s)
\ge\eta_{\mathfrak C}(n)\}.
\]

If the set is empty, set \(B^*_\eta(n)=\infty\). The corresponding
log-cost invariant is

\[
d^*_\eta(n)=\left\lceil\log_2(1+B^*_\eta(n))\right\rceil.
\]

The notation \(B^*(n)\) or \(d^*(n)\) denotes this threshold up to
class-equivalent changes of \(\eta_{\mathfrak C}\).
All infima, logarithms, and numerical comparisons involving \(B^*\) or
\(d^*\) in this section use this same scalarization. The underlying
resource-valued invariant remains the Pareto frontier of admissible
ceilings.

\end{definition}

\Cref{fig-saturated-threshold-budget} gives the threshold reading of
\cref{def-degree-invariant} in the hard-side regime.

\begin{figure}[tbp]
\centering
\begin{tikzpicture}[fw diagram,x=1cm,y=1cm,scale=.94,transform shape]
  \draw[-{Stealth[length=2.2mm,width=1.5mm]},draw=fwInk]
    (0,0) -- (11.5,0) node[below] {scalarized budget \(B\)};
  \draw[-{Stealth[length=2.2mm,width=1.5mm]},draw=fwInk]
    (0,0) -- (0,4.3) node[left,align=center] {\(I_n^{\mathrm{sat}}(B)\)};
  \draw[draw=fwAmber,dashed,line width=0.8pt]
    (0,2.6) -- (11.0,2.6)
    node[right,font=\scriptsize] {\(\eta_{\mathfrak C}(n)\)};
  \draw[draw=fwBlue,line width=1.4pt]
    (0,0.35) -- (2.0,0.35) -- (2.0,0.8) -- (4.2,0.8) --
    (4.2,1.25) -- (6.4,1.25) -- (6.4,2.0) -- (8.3,2.0) --
    (8.3,3.2) -- (10.8,3.2);
  \draw[draw=fwCoral,dashed] (5.0,0) -- (5.0,2.1);
  \node[font=\scriptsize,fwCoral,below] at (5.0,0) {\(b(n)\)};
  \draw[draw=fwViolet,dashed] (8.3,0) -- (8.3,3.25);
  \node[font=\scriptsize,fwViolet,below] at (8.3,0) {\(B^*\)};
  \node[fw residual,minimum width=2.5cm,font=\scriptsize] at (3.1,3.35)
    {declared budget\\below threshold};
  \draw[fw arrow] (4.0,3.05) -- (5.0,2.2);
\end{tikzpicture}
\caption{Schematic hard-side reading of the saturated cost invariant:
\(B^*\) is the infimum budget at which the monotone saturated profile reaches
\(\eta_{\mathfrak C}(n)\), while the displayed ceiling satisfies
\(b(n)<B^*\).  The diagram asserts neither continuity nor general attainment.}
\label{fig-saturated-threshold-budget}
\end{figure}

The invariant \(B^*\) is the first saturated public budget at which the
observable closure contains either a target-aligned quotient or a
regular target-oriented geometry. It is minimized over all admissible
representations generated from the presentation. Below this budget,
every structural representation beneath the declared ceiling has
negligible target contrast at the transfer-class threshold. Raw Walsh
degree alone does not determine \(B^*\): a valid low-cost quotient,
lift, or forward transport may represent the same target distinction,
whereas measurability alone provides no low-cost representation.

\begin{proposition}[The saturated cost invariant is intrinsic to the problem presentation]\label{prop-dstar-intrinsic}

The invariants \(B^*\) and \(d^*\) depend only on the family of task
presentations, the admissible dynamic slot activated by \PartII, and the
saturated cost filtration generated from the observable algebra. They
are independent of the choice of method, program, solver, or learning
algorithm.

\end{proposition}

\begin{proof}\phantomsection\label{proof-dstar-intrinsic}

The quantities \(m_n^{\mathrm{sat}}(B)\) and \(G_n^{\mathrm{sat}}(B)\)
are suprema over all observable structures in the saturated admissible
representation class at public cost at most \(B\). This class is
generated by the presentation and the oracle-free constructor rules:
public basis changes, represented quotients, valid lifts, exposed
potentials, pointwise geometries, scalar boundary readouts, and forward
observable transport/generator components after \PartII projection. The
class is fixed before choosing a particular solver. \PartII proves that
every algorithm-induced movement lies in Geom, Caus, Abs, Lift, Bdry, or
\(\Res\). The target-amplifying exposed contributions
charged in \(I_n^{\mathrm{sat}}(B)\) are the two intrinsic coordinates:
Geom/Caus, which measures exposed local geometry and its oriented use,
and Abs, which measures represented quotient abstraction. Lift is
retained with the intrinsic structure it lifts, and Bdry is retained
with the scalar readout and the structure it colors. Their contextual
target charges are not reassigned without a represented comparison.
The label \(\Res\) supplies no additional structural
coordinate: a verified channel projection is reclassified at its
charged cost, while the full-process target effect is accounted for by
the exact first-hit quotient. Therefore minimizing the budget over the charged exposed
coordinates is a property of the presented problem and its admissible
structural closure.

\end{proof}

\begin{lemma}[Saturated cost is invariant under admissible presentation equivalence]\label{thm-dstar-presentation-invariance}

Let two task-family presentations \(\mathcal T\) and \(\mathcal T'\) be
admissibly equivalent in the sense of \PartI, and assume the equivalence
transports the \PartII dynamic constructors: budget-admissible kernels,
represented quotients, valid lifts, boundary readouts, the
Geom/Caus/Abs/Lift/Bdry projections. Let \(a(d)\) and \(b(d)\) be
log-cost distortions for saturated public representation cost. Then

\[
I_{\mathcal T}(d)\le I_{\mathcal T'}(b(d)),
\qquad
I_{\mathcal T'}(d)\le I_{\mathcal T}(a(d)).
\]

Consequently the threshold degrees are comparable up to the same
distortion:

\[
d^*_{\mathcal T',\eta}\le b(d^*_{\mathcal T,\eta}),
\qquad
d^*_{\mathcal T,\eta}\le a(d^*_{\mathcal T',\eta}).
\]

If \(a\) and \(b\) preserve budget admissibility, then largeness of \(d^*\)
beyond every within-budget band is invariant under the presentation
equivalence.

\end{lemma}

\begin{proof}\phantomsection\label{proof-dstar-presentation-invariance}

Take any structural visibility term counted by \(I_{\mathcal T}(d)\),
with \(d\) in log-cost notation. If it is a Geom/Caus term, the
presentation equivalence transports the pointwise observable,
comparability relation, potential, current, mixing and consistency
gates, Caus authority envelope and current provenance, Caus alignment,
quotient stability record, scalar boundary coloring/readout, and
structural reach datum to \(\mathcal T'\) with log-cost at most
\(b(d)\). If it is an Abs term, the equivalence
transports the represented fibers, terminal compatibility, lumpability
condition, and charged quotient structural data to a represented quotient of
log-cost at most \(b(d)\). Valid Lift and Bdry terms are transported
because their target-fraction and terminal-readout conditions are part
of the transported dynamic constructor data. The unitary transport in
\Cref{def-admissible-presentation-equivalence} preserves the
visibility diagnostic exactly. Taking the supremum gives
\(I_{\mathcal T}(d)\le I_{\mathcal T'}(b(d))\). The reverse
inequality is identical with \(a(d)\). The displayed
threshold inequalities follow directly from the definition of
\(d^*_{\eta}\). Since budget admissibility is preserved by hypothesis,
escaping every within-budget band is also preserved.

\end{proof}

\begin{proposition}[Closure under generated representations]\label{prop-no-broad-narrow-dilemma}

The saturated class used in \(B^*\) is the cost-filtered
closure of the presented observable algebra under the admissible
oracle-free constructors fixed by the task object; its budget-filtered
part is \(\mathcal R^{\mathcal B}_n\). If a budget-admissible
uniform rule generates a target-aligning quotient, potential, geometry,
valid lift, boundary readout, or forward observable transport component,
then the generated object belongs to this saturated class at the least
saturated public cost of its representation. If that cost is within the
declared budget and the relevant conditions hold, then
\(I_n^{\mathrm{sat}}(B)\) reaches the declared transfer-class margin at
that budget. If \(B^*\) lies above the declared ceiling, up to the
allowed transfer-class overheads, no such within-budget admissible
representation exists in the budgeted saturated closure.

\end{proposition}

\begin{proof}\phantomsection\label{proof-no-broad-narrow-dilemma}

The constructor grammar is fixed by the task presentation, the
valid-lift conditions, the quotient conditions, and the resource envelope. A
runtime-generated object produced by an oracle-free uniform rule has a
finite lifted configuration trace, so \PartII records it as
lifted/generated observable data. If it supplies represented fibers, an
exposed potential, a valid lifted interface, a scalar boundary
value/readout, or a forward observable transport component after the
component projection, it is one of the objects over which
\(m_n^{\mathrm{sat}}(B)\) or \(G_n^{\mathrm{sat}}(B)\) takes its
supremum, with \(B\) equal to the least saturated public representation
cost. Thus the invariant ranges over the fixed admissible closure,
including representations outside a raw Walsh cover. A lower bound for
the cost of every target-aligned object in this closure is a structural
representation-cost theorem over the admissible constructor class.

\end{proof}

\begin{proposition}[Saturated structural cost can be tested by budget]\label{prop-saturated-degree-audit}

Fix \(n\), a saturated public budget \(B\), and the saturated admissible
constructor grammar of the presentation. For every admissible quotient
\(Q\in\mathcal Q^{\mathrm{sat}}_{n,\le B}\), set the visibility
diagnostic kernel \(K_Q^{\mathrm{vis}}=A_QP_Q\). For every admissible
active-slot master certificate
\(\mathscr C=(\mathfrak N,q,\mathfrak G_q,\mathfrak S,\mathfrak K_q)\),
where
\(\mathfrak G_q=(\Phi_Q,\Phi_\perp,D_Q,J_Q)\), set

\[
K_{\mathscr C}^{\mathrm{vis}}
=
\overline A_qS_q(\lambda)M_{\Phi_Q}
C_{\Phi_Q,\mathfrak K_q}^{\mathrm{auth}}(D_Q)
\frac{1}{1+\rho_{\Phi_Q}(D_Q)}
\Pi_{U\mathcal M_Q(D_Q)}.
\]

Then the structural coordinates are exactly the suprema of the
corresponding normalized visibility forms:

\[
m_n^{\mathrm{sat}}(B)=
\sup_{Q\in\mathcal Q^{\mathrm{sat}}_{n,\le B}}
\frac{\langle y,K_Q^{\mathrm{vis}}y\rangle}{\|y\|^2},
\] \[
G_n^{\mathrm{sat}}(B)=
\sup_{\mathscr C:\operatorname{Cost}_{\mathrm{sat}}(\mathscr C)\preceq B}
\frac{\langle y,K_{\mathscr C}^{\mathrm{vis}}y\rangle}{\|y\|^2}.
\]

Consequently \(B^*>B\) is proved at size \(n\) by bounding the explicit
visibility-kernel family on the target contrast for every
budget-admissible representative at budget \(B\). If the
budget-admissible visibility-kernel family has a finite enumeration, a
semialgebraic description, or another exact parametrization, this is an
ordinary finite-dimensional linear-algebra or optimization problem.

\end{proposition}

\begin{proof}\phantomsection\label{proof-saturated-degree-audit}

For quotients, \(\langle y,K_Q^{\mathrm{vis}}y\rangle/\|y\|^2\) is the
normalized visibility diagnostic of the represented quotient subspace
after applying the Abs utility condition. For Geom/Caus,
\(\langle y,K_{\mathscr C}^{\mathrm{vis}}y\rangle/\|y\|^2\) is the
normalized target-oriented certificate visibility after the quotient
support, stability, mixing, authority-compatible Caus alignment,
consistency, and regularity conditions. These are precisely the terms used in the
definitions of \(m_n^{\mathrm{sat}}(B)\) and \(G_n^{\mathrm{sat}}(B)\).
Taking suprema over the saturated constructor class gives the displayed
identities.

\end{proof}

Raw Walsh spectra and algebraic-immunity diagnostics can rule out target
visibility in a declared raw basis or in a covered family of public
bases. A bound for the full set of executable extraction procedures
requires the tested kernel family to cover the saturated structural
closure. A large-\(B^*\) theorem for the present visibility invariant bounds the visibility
diagnostics above, budget by budget, over all budget-admissible
saturated quotients, valid lifts, nonlinear represented kernels, and
Geom/Caus data. For declared fragments, source-to-witness implications
give additional consequences. The hardness direction uses the
accountability implication from target arrival to counted visibility.

\subsection{Universal Derivability
Accountability}\label{full-derivability-characterization}
\label{subsec:universal-derivability-accountability}

The next two statements give the unconditional lower-bound direction:
successful admissible computation forces visibility in the uniform
saturated profile. They follow from the budgeted derivability closure
and the clocked-transcript first-hit quotient.

\begin{proposition}[Forced derivable witness]\label{prop-forced-derivable-witness}

Let \(A_n\) be a presentation-relative uniform oracle-free algorithm
running inside a budget-compatible resource envelope. If its target
arrival probability is \(\alpha(n)\), let \(h(n)\) be its normalized
horizon, including the fresh start phase, and let
\(B_{\mathrm{hit}}(n)\) be the explicit first-hit ceiling from
\eqref{eq:first-hit-saturated-cost}. Then

\[
I_n^{\mathrm{sat}}(B_{\mathrm{hit}}(n))
\ge m_n^{\mathrm{sat}}(B_{\mathrm{hit}}(n))
\ge \frac{\alpha(n)}{e h(n)}.
\]

In particular, successful budget-admissible arrival produces a
charged exact first-hit quotient in the derivability closure with no
additional numerical transfer factor. If \(\alpha\) has a numerical
transfer-class lower margin and \(h(n)\le\gamma(n)\) for some
\(\gamma\in\mathfrak C_{\mathrm{num}}\), the displayed exact inequality
gives the corresponding numerical transfer-class margin.  A finite but
otherwise unrestricted declared time capacity does not imply this last
closure property. The polynomial wording is recovered by setting
\(\mathcal B=\mathcal B_{\mathrm{poly}}\), where the horizon is itself
polynomially bounded.

\end{proposition}

\begin{proof}\phantomsection\label{proof-forced-derivable-witness}

Embed the run of \(A_n\) on its clocked \PartII transcript carrier. The
finite control, work tapes, and memory remain in the current native state
and are charged at peak capacity; the carrier records the clock and
finite outcome/terminal transcript. The one-step transition is a legal represented kernel by
Encoding Adequacy. Let \(h(n)\) be the charged stopping horizon supplied
by the budget-compatible resource envelope after start normalization.
\Cref{lem:exact-affordable-first-hit} constructs
an exact represented first-hit quotient \(q_A\) whose saturated cost is
at most \(B_{\mathrm{hit}}\) and
whose visibility satisfies

\[
V_{q_A}
\ge
\frac{\alpha(n)}{e h(n)}.
\]

Since
\(V_{q_A}\le m_n^{\mathrm{sat}}(B_{\mathrm{hit}})\le
I_n^{\mathrm{sat}}(B_{\mathrm{hit}})\), the claimed exact lower bound
follows. For the optional transfer-margin conclusion, division by
\(h\le\gamma\) preserves the numerical transfer class by its declared
multiplicative closure.  B4 alone is not used for that conclusion.

\end{proof}

\begin{lemma}[Full derivability profile criterion]\label{thm-full-derivability-characterization}

Fix a task family and one budget envelope of finite time capacity.  At its
declared ceiling \(B\),

\[
\operatorname{Succ}_{\mathfrak M,n}(B)
\le
eH_B(n)
m_{\mathfrak M,n}^{\mathrm{sat}}(\Psi(B,n)).
\]

Therefore a saturated-profile bound at the named enlarged ceiling controls
the complete success profile of the declared operational class.  Under
\(\mathcal B_{\mathrm{poly}}\), set \(B=b_{\mathrm{poly}}\).  The concrete
implication is
\[
m_{\mathfrak M,n}^{\mathrm{sat}}
 \bigl(\Psi(b_{\mathrm{poly}},n)\bigr)=\operatorname{negl}(n)
\quad\Longrightarrow\quad
\operatorname{Succ}_{\mathfrak M,n}(b_{\mathrm{poly}})
=\operatorname{negl}(n),
\]
with the polynomial horizon factor already included in the declared
numerical transfer system.

The converse requires an additional executable source-to-arrival
theorem. A visible quotient or geometry need not by itself supply a
selector, sampler, descent rule, or decoder.

\end{lemma}

\begin{proof}\phantomsection\label{proof-full-derivability-characterization}

This is \Cref{thm-universal-saturated-profile-accountability}. Its proof uses the
uniform exact first-hit quotient and therefore applies to the complete
positive process without a residual-baseline hypothesis. The final
sentence records the logical direction of the result: the proof
constructs visibility from an algorithm, not an algorithm from arbitrary
visible structure.

\end{proof}

\begin{remark}[Scope of the characterization]\label{rem-characterization-not-lower-bound-method}

The theorem reduces the declared-class success profile to the full uniform
saturated profile at one charged enlarged ceiling.  A
large \(B^*\) supplies a sufficient lower-bound certificate. Without a
source-to-arrival theorem, a small \(B^*\) records visible structure but
does not assert the existence of a solver.

\end{remark}

\subsection{Alignment Complexity and the Saturated
Budget}\label{alignment-complexity-and-the-saturated-budget}
\label{subsec:alignment-complexity-saturated-budget}

Alignment complexity records the least budget supporting the required target-oriented comparison.
\begin{definition}[Budgeted alignment complexity]\label{def-alignment-complexity}

Fix the same admissible scalarization \(\kappa\) used for \(B^*\). Let
\(\operatorname{KT}(\mathscr R\mid\operatorname{enc}(I_n))\) be the
least scalarized charge of one family-uniform program that, from
\(\operatorname{enc}(I)\), constructs a \emph{complete} exact-quotient
or master quotient--geometry certificate \(\mathscr R_I\) for every
instance in the family. The record includes its carrier, reference law,
initialization, generator, quotient fibers, terminal sectors, utility,
geometry, authority/current provenance, defect and stability data,
normalization, and every comparison used by its visibility. The charge
includes the program, running resources, and every generated
instance-dependent datum; a cheap terminal evaluator with an expensive
carrier is therefore not cheap alignment. For a threshold \(\zeta>0\),
define

\[
\begin{aligned}
\kappa_{\mathrm{align},\zeta}(n)
=\min\bigl\{\operatorname{KT}(\mathscr R\mid\operatorname{enc}(I_n)):;&
\mathscr R\text{ is a complete exact-quotient or master certificate}\\
&\text{with visibility at least }\zeta(n)
\bigr\}.
\end{aligned}
\]

If no such record exists, set
\(\kappa_{\mathrm{align},\zeta}(n)=\infty\). Write
\(\kappa_{\mathrm{align}}\) only when the threshold is displayed or
fixed in context.

\end{definition}

\begin{lemma}[Saturated budget and alignment complexity are class-equivalent]\label{thm-alignment-sandwich}

There are monotone scalar overhead maps \(\Pi_1,\Pi_2\), determined by
the constructor simulation overheads of the presentation and the
transfer class, such that

\[
\kappa_{\mathrm{align},\eta/2}(n)\le \Pi_1(B^*_{\eta}(n)),
\qquad
B^*_{\eta}(n)\le
\Pi_2(\kappa_{\mathrm{align},\eta}(n)).
\]

Thus \(B^*(n)\) lies above the declared budget if and only if
the corresponding alignment complexity does, up to the transfer-class
overheads and the class-equivalent constant change from \(\eta\) to
\(\eta/2\).
For \(\mathcal B_{\mathrm{poly}}\), this is the usual
polynomial/superpolynomial equivalence.
The first inequality is read with an attaining witness when the coding
is locally finite as in \Cref{lem-derivability-well-posed};
without attainment it holds after replacing \(B^*(n)\) by any
class-equivalent budget \(B_+(n)>B^*(n)\) for which
\(I_n^{\mathrm{sat}}(B_+(n))\ge\eta_{\mathfrak C}(n)\).

\end{lemma}

\begin{proof}\phantomsection\label{proof-alignment-sandwich}

Under local finiteness, if \(B^*_{\eta}(n)=B\), choose witnesses for
\(I_n^{\mathrm{sat}}(B)=m_n^{\mathrm{sat}}(B)+G_n^{\mathrm{sat}}(B)
\ge\eta\). At least one of the two complete records has visibility at
least \(\eta/2\); choose it and put
\(\widetilde B=B\). Without local finiteness, fix \(B_+>B^*\)
satisfying the condition in the statement, choose the witness at
\(B_+\), and put \(\widetilde B=B_+\). Its derivation is a finite
sequence of primitive reads, admissible constructors, and legal
computations of total cost at most \(\widetilde B\). Budget completeness
simulates each step with class-preserving overhead, so the resulting
complete target-aligned record has conditional charged complexity at most
\(\Pi_1(\widetilde B)\). This is \(\Pi_1(B^*)\) in the attained case
and the stated \(B_+\)-replacement otherwise.

Conversely, if \(\kappa_{\mathrm{align},\eta}(n)=K\), fix a charged
program of scalarized total cost at most \(K\) constructing the complete
certificate from \(\operatorname{enc}(I_n)\).
By Encoding Adequacy and the computation-closure clause of Definition
\hyperref[def-budgeted-derivability-closure]{Budgeted derivability
closure}, this program is a budget-admissible represented record of
scalarized cost at most \(\Pi_2(K)\). Its visibility is at least
\(\eta\), so \(B^*_{\eta}(n)\le \Pi_2(K)\).

\end{proof}

\subsection{Reversible Spectrally Compressed
Fragment}\label{reversible-compressed-fragment}
\label{subsec:reversible-spectrally-compressed-fragment}

The reversible fragment isolates the represented spectral subfamily used by the compressed audit.
\begin{definition}[Reversible spectrally compressed fragment]\label{def-reversible-compressed-fragment}

Fix a declared monotone numerical spectral-capacity map
\(\mathsf s:C\to[0,\infty]\), subject to the B4 compatibility condition
of \Cref{def-budget-structure}.
Let \(\mathsf{Rev}_{\mathrm{gate}}\) denote the class of budget-admissible
killed generators \(L\) such that \(K_N=-L_N\) is self-adjoint and
positive semidefinite in \(L^2(\mu_N)\), has spectrum in \([0,\Lambda]\)
with \(\Lambda\le\mathsf s(b(n))\), is represented in the saturated closure,
starts from \(d\mu_0=f_0\,d\mu_N\) with \(\|f_0\|_2\) controlled by the
declared transfer class, and whose induced reversible geometry and
boundary readouts satisfy the Geom/Caus/Bdry conditions with
transfer-class constants. For a depth bound \(q(n)\), let
\(\mathcal D^{(q)}\) be the subclosure allowing at most \(q(n)\)
dependent applications of represented kernels or maps, with breadth and
postprocessing allowed exactly when their total charge remains
budget-admissible. In the polynomial instance this recovers the usual
polynomial spectral, norm, breadth, and comparison bounds.

\end{definition}

\begin{lemma}[Reversible spectral compression]\label{lem-reversible-spectral-compression}

Let \(K_N\succeq0\) be self-adjoint with spectrum in \([0,\Lambda]\).
For every \(\lambda>0\) and \(0<\varepsilon<1\), there is an explicit
real polynomial \(p_q\) of degree

\[
q=O\!\left((1+\sqrt{\Lambda/\lambda})\,\log(2/\varepsilon)\right)
\]

such that
\(\sup_{s\in[0,\Lambda]}|(\lambda+s)^{-1}-p_q(s)|\le \varepsilon/\lambda\).

\end{lemma}

\begin{proof}\phantomsection\label{proof-reversible-spectral-compression}

Map \([0,\Lambda]\) affinely to \([-1,1]\). The function
\(s\mapsto(\lambda+s)^{-1}\) is analytic on a Bernstein ellipse whose
parameter \(\rho>1\) satisfies
\(\log\rho\ge c(1+\sqrt{\Lambda/\lambda})^{-1}\) for an absolute
constant \(c>0\). Truncating its Chebyshev expansion at degree \(q\)
therefore gives uniform error
\[
O\!\left(
\lambda^{-1}
\exp\!\left\{-\frac{cq}{1+\sqrt{\Lambda/\lambda}}\right\}
\right).
\]
Choosing
\(q=C(1+\sqrt{\Lambda/\lambda})\log(2/\varepsilon)\) with \(C\)
large proves the claim \citep{Trefethen2019}.

\end{proof}

\begin{lemma}[Reversible compressed-fragment hardness]\label{thm-reversible-fragment-hardness}

Fix \(\lambda\ge1/\beta_{\mathfrak C}(n)\),
\(\Lambda\le\mathsf s(b(n))\), and
a tolerance \(\varepsilon_{\mathfrak C}(n)\) negligible relative to
\(\mathfrak C\). Let
\(q(n)=O((1+\sqrt{\Lambda/\lambda})
\log(2/\varepsilon_{\mathfrak C}(n)))\),
with total approximation charge lying within \(\mathcal B\).
Assume a charged spectral-to-visibility comparison: for every
represented \(h=p(K_N)k_T\) in this fragment there is a certificate
\(\mathscr C_h\in\mathcal D^{(q)}_{\le b(n)}\) such that

\[
|\langle f_0,h\rangle|
\le C_{\mathrm{bridge}}(n)V_{\mathscr C_h}
+\varepsilon_{\mathrm{bridge}}(n),
\]

where \(C_{\mathrm{bridge}}\) lies in the transfer class and
\(\varepsilon_{\mathrm{bridge}}\) is negligible. If
\(I_n^{(q)}(b(n))=\operatorname{negl}_{\mathfrak C}(n)\), then every
\(L\in\mathsf{Rev}_{\mathrm{gate}}\) has negligible discounted target
arrival at discount \(\lambda\). For \(\mathcal B_{\mathrm{poly}}\) and
inverse-polynomial thresholds, this recovers the usual
\(O((1+\sqrt{\Lambda/\lambda})\log n)\) degree choice.

\end{lemma}

\begin{proof}\phantomsection\label{proof-reversible-fragment-hardness}

Assume, contrapositively, that
\(A_T(\lambda;L,\mu_0)=\langle f_0,(\lambda+K_N)^{-1}k_T\rangle\) has
transfer-class size. Choose \(\varepsilon_{\mathfrak C}\) negligible
relative to \(\mathfrak C\) after absorbing the transfer-class losses in
\(\|f_0\|_2\), \(\|k_T\|_2\), and \(\lambda^{-1}\). By Lemma
\hyperref[lem-reversible-spectral-compression]{Reversible spectral
compression}, the observable

\[
h=p_q(K_N)k_T
\]

is within negligible pairing error of the exact resolvent value. Since
\(K_N\), \(k_T\), and the coefficients of \(p_q\) are represented,
\(h\in\mathcal D^{(q)}_{\le b(n)}\). The defining conditions of
\(\mathsf{Rev}_{\mathrm{gate}}\) make this a represented reversible
geometry and boundary readout. The approximation estimate makes
\(|\langle f_0,h\rangle|\) transfer-class non-negligible. The assumed
bridge then gives

\[
V_{\mathscr C_h}
\ge
\frac{|\langle f_0,h\rangle|
      -\varepsilon_{\mathrm{bridge}}(n)}
     {C_{\mathrm{bridge}}(n)},
\]

which is transfer-class non-negligible. Hence
\(I_n^{(q)}(b(n))\ge V_{\mathscr C_h}\) is non-negligible,
contradicting the hypothesis. Without the displayed bridge, polynomial
approximation of the resolvent supplies an observable pairing but does
not by itself prove visibility or target reach.

\end{proof}

\begin{example}[Saturated cost for SAT presentations]\label{ex-dstar-sat}

XOR-SAT with public syndrome observables has polynomial \(B^*\), hence
logarithmic \(d^*\) in log-cost notation, because the public syndrome
quotient is represented compactly even when raw Walsh degree is high. A
random 3-SAT presentation without an exposed quotient and without a
regular aimed defect geometry has large saturated cost in the
random-target model. A planted SAT instance has low saturated cost when
the presentation exposes a low-cost component that recovers the planted
structure.

\end{example}

\section{Nonlinear Quotients and the Kernel Normal
Form}\label{nonlinear-quotients-and-the-kernel-normal-form}
\label{sec:nonlinear-quotients-kernel-normal-form}

\subsection{Quotient kernels and represented nonlinear features}
\label{subsec:quotient-kernels-represented-features}

A nonlinear projection used by a computation is accounted for by the
same saturated invariant \(B^*\). In a finite Boolean task space, every
represented nonlinear feature map determines a finite feature subspace
and hence a weighted projection operator. For \(\Abs\), the
counted object is the quotient projection associated to represented
fibers. A feature Gram matrix is a coordinate description of the same
quotient after projection onto the represented feature span. If the
nonlinear construction instead supplies a pointwise geometry, potential,
metric, score, or scalar terminal/support coloring, it is counted in the
Geom/Caus coordinate, while a Bdry readout remains a full-cost
contextual extension over the represented structure. Its target charge
is transferred to that structure only by a represented comparison. If it only modifies a declared geometry, the
Dirichlet form comparison of \Cref{quotient-defect-and-geom-form-comparison} is the perturbation bound that
determines whether the geometric comparison transfers. If the feature
map induces an exactly lumpable quotient, its projection has algebraic
quotient provenance. It is counted in \(m_n^{\mathrm{sat}}(B)\) at its least
saturated public cost precisely when its complete record satisfies
\cref{def-paper1-abs-usefulness-gate,def-abs-extraction}. A prospective source
use also satisfies \cref{def-pre-hit-temporally-admissible-source}. A
nonzero-defect quotient enters only
through the supported Geom/Caus coordinate when its compensation record
is dynamically qualified. This is the kernel-trick form of the
invariant.

The nonlinear quotient kernel is the coordinate-free projection determined by an affordable represented feature partition.
\begin{definition}[Budget-admissible nonlinear quotient kernel]\label{def-admissible-nonlinear-quotient-kernel}

Let \((X,\mathcal F,\mu,T)\) be a finite presented task. A
budget-admissible nonlinear quotient representation at saturated public
cost at most \(B\) is a finite feature map or projection rule \(U\)
generated by the observable algebra, valid lifted coordinates, and
resource-envelope operations at cost at most \(B\). The cost includes
the feature description, quotient labels, comparison data, charged
public basis changes, lifted memory, retained preprocessing tables, and
finite outcome/terminal transcript coordinates whenever those objects are used as
observable structure. Nonuniform advice and oracle data are inadmissible
coordinates unless they are charged to the presentation; otherwise they
are excluded from the
uniform oracle-free class. Observable finite postprocessings
\(h\circ U\), including coarsenings and threshold labels, are treated as
feature maps at their own saturated cost. The equality fibers of the
chosen \(U\) define a quotient

\[
q_U:X\to Q_U.
\]

It is admissible as an \(\Abs\) quotient when its fibers are
represented in the \PartI fiber-represented-observable sense, the target and
terminal sectors are quotient events, the quotient is exactly lumpable
on the target window, and the quotient has nonzero quotient-utility
factor on that window. Approximate lumpability is handled separately by
\Cref{def-supported-geom-certificate}: it contributes to the
supported Geom/Caus profile only with a represented transfer-class
stability factor, and every remaining defect is decomposed again by Part
II. Let \(\mathcal H_U\le L^2(X,\mu)\) be the subspace of functions
constant on the fibers of \(q_U\), and let \(P_U\) be the weighted
orthogonal projection onto \(\mathcal H_U\). The associated
quotient-visibility kernel is the positive self-adjoint idempotent

\[
K_U=P_U,
\qquad
\operatorname{Vis}_T(U)=\frac{\langle y,K_Uy\rangle_{L^2(\mu)}}{\|y\|^2_{L^2(\mu)}}.
\]

This projection kernel supplies the quotient visibility. The extraction
contribution also includes the quotient-utility factor \(A_U\),
which records the represented level-set control and quotient target
utility. Different feature coordinates, Gram matrices, or nonlinear
formulas with the same represented fibers give the same \(P_U\). The
degree of \(U\) is the least saturated representation cost of the
induced quotient projection \(K_U\), independently of discovery time or
the proof of optimality, together with its charged utility data relative
to the observable algebra.

\end{definition}

\Cref{fig-nonlinear-feature-routing} records the structural routing of the
nonlinear feature map in \cref{def-admissible-nonlinear-quotient-kernel}.

\begin{figure}[tbp]
\centering
\begin{tikzpicture}[fw diagram]
  \node[fw source,minimum width=3.0cm] (feat) at (-4.8,0)
    {represented feature\\\(U:X\to\mathcal U\)};
  \node[fw provenance,minimum width=2.8cm] (fib) at (-1.2,0)
    {equality fibers\\\(q_U:X\to Q_U\)};
  \node[fw geom,minimum width=2.8cm] (exact) at (2.5,1.2)
    {qualified exact fibers\\\(K_U=P_U\)};
  \node[fw use,minimum width=2.8cm] (approx) at (2.5,-1.2)
    {nonzero defect\\compensated Geom};
  \node[fw terminal,minimum width=2.5cm] (abs) at (6.2,1.2)
    {\(\Abs\) visibility};
  \node[fw residual,minimum width=2.5cm] (geom) at (6.2,-1.2)
    {Geom/Caus ledger};
  \draw[fw arrow] (feat) -- (fib);
  \draw[fw arrow] (fib) -- (exact);
  \draw[fw arrow] (fib) -- (approx);
  \draw[fw arrow] (exact) -- (abs);
  \draw[fw arrow] (approx) -- (geom);
\end{tikzpicture}
\caption{A nonlinear formula is classified through its represented structural
effect. Exact target-compatible fibers give the quotient projection kernel of
\cref{prop-kernel-trick-normal-form}; a nonzero lumpability defect remains in
the compensated geometric ledger with its stability data.}
\label{fig-nonlinear-feature-routing}
\end{figure}

\begin{proposition}[Kernel-trick normal form]\label{prop-kernel-trick-normal-form}

Every admissible nonlinear quotient representation \(U\) at cost at most
\(d\) has quotient visibility contribution

\[
V_U=A_U\frac{\langle y,K_Uy\rangle}{\|y\|^2},
\]

where \(A_U\in[0,1]\) is the same \PartI quotient-utility factor
used in \cref{def-abs-extraction}. Thus \(m_n^{\mathrm{sat}}(B)\ge V_U\) when \(U\) is
budget-admissible. The visibility form depends only on the induced
quotient subspace and its utility data, independently of the coordinates,
feature basis, Gram representation, discovery procedure, or algorithmic
construction of \(U\).

\end{proposition}

\begin{proof}\phantomsection\label{proof-kernel-trick-normal-form}

Since \(X\) is finite, the feature map \(U\) has finite image and its
equality fibers form a finite partition of \(X\). The
quotient-observable space is the closed finite-dimensional subspace of
\(L^2(X,\mu)\) consisting of functions constant on those fibers.
Weighted orthogonal projection onto that subspace is unique; different
feature coordinates for the same fibers produce the same projection
\(P_U\). A non-idempotent feature kernel contributes to \(\Abs\)
only through this projection onto the represented quotient subspace,
multiplied by the \PartI quotient-utility factor; otherwise its
target-relevant part is classified by the Geom/Caus/Bdry/residual
alternatives. If the quotient-utility conditions hold, \(q_U\) is one of
the quotients in \(\mathcal Q^{\mathrm{sat}}_{\le B}\). The definition
of \(m_n^{\mathrm{sat}}(B)\) is the supremum of the conditioned quotient
visibility over these represented target-relevant quotient subspaces,
hence the displayed inequality.

\end{proof}

\begin{lemma}[Nonlinear kernel test is algorithm-independent]\label{thm-nonlinear-kernel-audit-algorithm-independent}

Fix the presented observable task family, the budget-compatible resource
envelope, the admissibility rules of \PartII, and the structural
visibility threshold \(\eta(n)\). The invariant \(B^*\) is a property of
this observable task object. If the generator induced by any
presentation-relative uniform oracle-free budget-admissible algorithm
contains a target-relevant nonlinear projection that induces a target-relevant
\(\Abs\) quotient at saturated public cost at most \(B\), then
\(m_n^{\mathrm{sat}}(B)\) counts the induced quotient visibility kernel
\(K_U\) through its conditioned projection mass. If the same nonlinear
construction instead induces a pointwise potential, defect, score,
metric, conductance, or geometry, then \(G_n^{\mathrm{sat}}(B)\) counts
its conditioned Geom/Caus visibility exactly at budgets in its admissible
budget set. If it is
only approximately lumpable, it contributes through
\(G_{Q,\mathrm{ql},n}^{\mathrm{sat}}(B)\) precisely when the represented
data supply the compensated certificate of
\Cref{quotient-defect-and-geom-form-comparison}. If it only approximately
matches an existing geometry, \Cref{quotient-defect-and-geom-form-comparison} supplies the Dirichlet form
comparison; without a transfer-class comparison loss, the construction
must be counted as its own geometry at budgets in its admissible budget set. If it
fails the valid quotient, lift, or boundary conditions, then its exact admissible
part is decomposed again by \PartII, while any target information
outside the observable/lifted algebra is presentation-inaccessible or
oracle-dependent, and any remaining nonlocal transport belongs to the
residual term \(\Res\).

Consequently, for every budget \(B\),

\[
B^*_\eta(n)>B
\quad\Longleftrightarrow\quad
I_n^{\mathrm{sat}}(B)=m_n^{\mathrm{sat}}(B)+G_n^{\mathrm{sat}}(B)<\eta(n)
\]

for the chosen structural visibility threshold. In kernel language this
says that the quotient visibility kernels \(K_Q^{\mathrm{vis}}\) and the
supported Geom/Caus visibility kernels \(K_{\mathscr C}^{\mathrm{vis}}\) from
\cref{prop-saturated-degree-audit} have total conditioned target visibility below \(\eta(n)\).
The kernel notation is the normal form of the already-defined accounting
\(I_n^{\mathrm{sat}}(B)=m_n^{\mathrm{sat}}(B)+G_n^{\mathrm{sat}}(B)\).
This identity concerns the represented structural profile. For a
budget-\(B\) algorithm, full-process target arrival is related to that
profile by the exact first-hit quotient at the enlarged budget
\(\Psi(B,n)\), as in \Cref{thm-universal-saturated-profile-accountability}.

\end{lemma}

\begin{proof}\phantomsection\label{proof-nonlinear-kernel-audit-algorithm-independent}

\PartII embeds every presentation-relative uniform oracle-free
budget-admissible algorithm as a finite Markov generator on each bounded
clocked transcript carrier. Component exhaustiveness then places
every target-relevant part of that generator in \(\Geom\),
\(\Caus\), \(\Abs\), valid \(\Lift\),
\(\Bdry\), or \(\Res\). Valid Lift and Bdry retain
their full represented carrier cost and contextual charge. A target-charge
transfer to an intrinsic Geom/Caus or Abs source requires a represented
comparison with its declared loss and residual term. By the kernel-trick normal form, every nonlinear quotient
source is represented by a unique quotient projection kernel and is
counted in \(m_n^{\mathrm{sat}}(B)\) exactly when
\(B\in\mathfrak B_{\mathrm{sat},n}(R_Q)\).
Pointwise geometric or causal sources are counted in
\(G_n^{\mathrm{sat}}(B)\). Approximate quotients enter only after the
\Cref{quotient-defect-and-geom-form-comparison} quotient defect estimate is applied; the controlled quotient
part is counted at budgets in its admissible budget set and the defect is decomposed
by the same \PartII channels. Nonquotient nonlinear geometric
modifications enter only through the \Cref{quotient-defect-and-geom-form-comparison} form comparison: a small
relative form error transfers an existing geometric comparison proof
with the displayed loss, while a large form error yields no uncharged
improvement and must be represented as its own Geom/Caus datum
or reclassified as residual or external information. A family-uniform
residual target correlation whose budget-admissible representative
satisfies one of the five structural-channel predicates cannot remain
residual by \PartII's no-uniform-residual-correlation theorem. Mere
measurability or correlation is not enough. The residual class is
defined after every represented structural-channel projection in the
current saturated closure has been removed. At budget \(B\), the
reclassified contribution is counted precisely when
\(B\in\mathfrak B_{\mathrm{sat},n}(R)\); otherwise it is absent from that
budget slice. If the representative is generated by an admissible
within-budget finite-horizon computation, Lemma
\hyperref[lem-budgeted-generation-gives-representation]{Budgeted
generation gives budgeted representation} puts that cost in the declared
budget account. If no budgeted representative exists, a transition rule
using the datum is not presentation-relative budget-admissible for the
original task; it uses data outside the budgeted observable algebra,
including presentation-inaccessible or oracle-dependent data,
nonuniform advice, over-budget recovery, or an enlarged primitive
signature. The induced
operator's exact first-hitting partition, committor, occupation law, or
stopped-future map is represented only by a budget-admissible saturated
representative; a finite-horizon version actually generated by an
admissible within-budget computation is counted. Thus
\(m_n^{\mathrm{sat}}(B)+G_n^{\mathrm{sat}}(B)\) exhausts the represented
quotient and Geom/Caus structural denotations at budget \(B\). This is
not a componentwise bound on the effect of \(\Res\); the
complete target-arrival probability is charged through the first-hit
quotient at \(\Psi(B,n)\). The saturated
classes are nested in \(B\), hence \(I_n^{\mathrm{sat}}(B)\) is
nondecreasing. Since \(B^*_\eta(n)\) is the least budget with
\(I_n^{\mathrm{sat}}(B)\ge\eta_{\mathfrak C}(n)\), the displayed
equivalence follows. Under an admissible presentation equivalence, the
\Cref{thm-dstar-presentation-invariance} transports the same statement
with the stated cost and normalization distortions.

\end{proof}

The saturated kernel test is therefore intrinsic to the fixed
presentation and resource envelope: algorithms
matter only through the budget-admissible projections, kernels, scalar
readouts, and value functions their induced generators instantiate over
represented structure. Comparison data identify the objective component
term in the existing accounting and need not be produced by an
algorithm or adversary. Raw Walsh
spectra, nonlinearity, and algebraic-immunity diagnostics can support
this test only when they bound the full saturated accounting
\(I_n^{\mathrm{sat}}(B)=m_n^{\mathrm{sat}}(B)+G_n^{\mathrm{sat}}(B)\),
including nonlinear represented quotient projections.

\subsection{Universal nonlinear accounting}
\label{subsec:universal-nonlinear-accounting}

The universal accounting record collects the quotient, geometric, and residual treatment of a represented nonlinear construction.
\begin{definition}[Universal observable structural accounting]\label{def-observable-structural-accounting}

Fix a task family, a budget-compatible resource envelope, a saturated
public budget family \(\mathcal B_n\), and the admissible
presentation-relative oracle-free operator class obtained from the Part
II representation theorem. For every induced generator \(L_n\), the task-level
structural accounting is the following object-level decomposition.

First apply the \PartII component decomposition on the induced finite
carrier as an operator identity. For a represented computation this is
its clocked transcript carrier. Discounted hitting
quantities are evaluated on the full positive killed propagation, with
the exposed terminal edge gate applied as in Definition
\hyperref[def-full-history-target-transport]{Full-history target-contact
transport}. Represented quotient components contribute to
\(m_n^{\mathrm{sat}}(B)\) when \(B\) belongs to their admissible budget
sets; pointwise geometric or causal components contribute to
\(G_n^{\mathrm{sat}}(B)\) under the same membership criterion. Valid Lift
and Bdry are retained as
full-cost contextual extensions. Bdry contributes no off-diagonal
movement, but neither its charge nor a Lift charge is reassigned to a
represented Geom/Caus or Abs source without a represented contextual
transfer inequality. A committor,
first-hitting partition, occupation law, or stopped-future value of
\(L_n\) contributes when the same denotation has a budget-admissible
representative in \(\mathcal R^{\mathrm{sat}}_{n,\le B}\), including a
representative generated by an admissible within-budget finite-horizon
computation. A semantic object contributes only through such a
representation. The orthogonal remainder \(\Res\) receives
no standalone hitting probability unless it is separately verified to
be a positive kernel. Target arrival of the complete process is retained
retrospectively in its exact represented first-hit accounting quotient.
Prospective structural credit is assigned only under
\cref{def-pre-hit-temporally-admissible-source}, and non-enumerative target
arrival is charged by \cref{thm-prospective-selector-accountability}.

\end{definition}

\begin{proposition}[Structural accounting after operational embedding]\label{prop-structural-accounting-is-universal}

The operational embedding maps the declared presentation-relative uniform
oracle-free budget class into the clocked-transcript generator class.  The
accounting above acts on this image and therefore requires no separately
supplied quotient, comparison datum, or invariant.  A represented program
enters through its induced generator \(L_A\); once \(L_A\) is formed, the same
admissible component projections, full-process first-hit accounting, and
saturated representation costs apply. A rule that becomes budget-admissible only
after adding unpresented data is not in this operator class for the
original task.

\end{proposition}

\begin{proof}\phantomsection\label{proof-structural-accounting-is-universal}

\PartII embeds every budget-bounded oracle-free algorithm whose
instance-dependent data are generated from
\(\mathcal R^{\mathcal B}_n\) and its legal lifted closure as a
finite Markov generator on each bounded clocked transcript carrier. The
generator form, fixed-space split, Hodge projection, quotient
projection, valid-lift projection, and boundary split are finite
algebraic operations on that generator and on the fixed observable
grammar. Derivability accounting decides which semantic future values
have saturated representatives. Every generator in the class and every
finite-horizon object generated by an admissible within-budget
computation remains included. Therefore two
algorithms inducing the same generator receive the same accounting, and
a semantic value function that lacks a representation in the budgeted
observable grammar has no low-cost representation associated with a
named algorithm alone.

\end{proof}

\begin{lemma}[Accounting-profile bound at the declared ceiling]\label{lem-affordable-source}

At the declared ceiling \(B\), the operational image satisfies
\[
\operatorname{Succ}_{\mathfrak M,n}(B)
\le eH_B(n)
m_{\mathfrak M,n}^{\mathrm{sat}}(\Psi(B,n)).
\]
Hence a negligible right-hand profile at that named enlargement gives a
negligible target-contact profile after the declared numerical transfer.
For \(\mathcal B_{\mathrm{poly}}\), this is the single substitution
\(B=b_{\mathrm{poly}}\), with enlarged ceiling
\(\Psi(b_{\mathrm{poly}},n)\).

\end{lemma}

\begin{proof}\phantomsection\label{proof-affordable-source}

Apply \eqref{eq-universal-saturated-profile-accountability}. The exact
first-hit quotient supplies the required exact-quotient accounting visibility term
for the complete computation. This implication assumes negligibility of the
accounting profile itself; source-profile hardness is instead the prospective
criterion of \cref{cor-universal-prospective-selector-hardness}.

\end{proof}

\begin{remark}[Role of the fragment criterion]\label{rem-characterization-role}

The visibility invariant measures represented target visibility. The
first-hit construction supplies target-contact accountability for every
uniform finite-horizon computation. Fragment theorems may give sharper
constants or separate Geom/Caus estimates, but the universal lower-bound
direction no longer depends on a fragment-specific residual argument.

\end{remark}

\section{Charged Noisy Combination and Semantic Cancellation}
\label{sec-charged-noisy-combination}

The saturated ledger distinguishes cancellation in the semantic label of a
derived observable from the noisy public primitives used to construct it.
The distinction is intrinsic to the represented derivation and therefore
applies independently of a particular elimination schedule.

\begin{definition}[Charged binary noisy-combination record]
\label{def-charged-binary-noisy-combination-record}

Let \(G\) be a finite binary vector space with uniform latent state
\(\vartheta\in G\).  A binary noisy-character presentation consists of public
records
\[
Y_i=\chi_{\gamma_i}(\vartheta)\xi_i,
\qquad 1\le i\le q,
\]
where \(\gamma_i\in\widehat G\), the signs \(\xi_i\) are independent of
\(\vartheta\), and
\[
\mathbb E\xi_i=\rho_i\in[-1,1].
\]
A \emph{charged binary noisy-combination record} contains a represented
matrix
\[
\Lambda=(\lambda_1^{\mathsf T},\ldots,
\lambda_K^{\mathsf T})\in\mathbb F_2^{K\times q}
\]
and the derived characters
\[
Y_{\lambda_j}:=\prod_{i:(\lambda_j)_i=1}Y_i.
\]
For a nonzero row \(\lambda\), define its semantic label, provenance support,
and attenuation by
\begin{align}
\label{eq-generic-combination-semantic-label}
\Gamma(\lambda)&:=\sum_{i=1}^q\lambda_i\gamma_i,
&
\operatorname{Prov}(\lambda)&:=\{i:\lambda_i=1\},\notag\\
w(\lambda)&:=|\operatorname{Prov}(\lambda)|,
&
a(\lambda)&:=\prod_{i\in\operatorname{Prov}(\lambda)}\rho_i.
\end{align}
For the whole materialized record put
\begin{equation}
\label{eq-generic-combination-dependence}
d_i(\Lambda):=|\{j:(\lambda_j)_i=1\}|,
\qquad
D_2(\Lambda):=\sum_{i=1}^q d_i(\Lambda)^2.
\end{equation}
For coefficients \(c=(c_1,\ldots,c_K)\), also put
\begin{equation}
\label{eq-generic-combination-weighted-dependence}
D_c(\Lambda)
:=
\sum_{i=1}^q
\left(
\sum_{j:(\lambda_j)_i=1}|c_j|
\right)^2.
\end{equation}

The complete record contains the primitive acquisition interface, the rule
constructing \(\Lambda\), every row locator, the shared combination DAG, the
materialized outputs, retained storage, the comparison or decoder, and the
prospective transition using the result.  Its cost is the resource sum of
those represented fields.  A shared vertex is charged once in the DAG; every
output, stored reference, and evaluation using that vertex remains charged.

Semantic cancellation means that \(\Gamma(\lambda)\) has smaller represented
support in a declared basis of \(\widehat G\).  It does not alter
\(\operatorname{Prov}(\lambda)\), \(w(\lambda)\), or the construction record.

\end{definition}

\begin{theorem}[Exact transport and complete charging for noisy combinations]
\label{thm-charged-noisy-combination-transport}

For the record of
\cref{def-charged-binary-noisy-combination-record} and every
\(v\in\widehat G\),
\begin{equation}
\label{eq-generic-combination-exact-character-transport}
\mathbb E_{\vartheta,\xi}
\bigl[Y_\lambda\chi_v(\vartheta)\bigr]
=
\mathbf1_{\{\Gamma(\lambda)=v\}}a(\lambda).
\end{equation}
For two rows, write
\[
a(\lambda_j+\lambda_l)
=
\prod_{i:(\lambda_j+\lambda_l)_i=1}\rho_i.
\]
Their exact conditional covariance over the primitive noise is
\begin{equation}
\label{eq-generic-combination-exact-covariance}
\operatorname{Cov}_{\xi}
\left(Y_{\lambda_j},Y_{\lambda_l}\mid\vartheta\right)
=
\chi_{\Gamma(\lambda_j+\lambda_l)}(\vartheta)
\left(
a(\lambda_j+\lambda_l)-a(\lambda_j)a(\lambda_l)
\right).
\end{equation}
Consequently, for a materialized located block
\[
H=\sum_{j=1}^K c_jY_{\lambda_j},
\qquad \sum_{j=1}^K|c_j|^2\le1,
\]
one has
\begin{align}
\label{eq-generic-combination-located-block}
\left|\mathbb E[H\chi_v(\vartheta)]\right|
&\le
\left(
\sum_{j:\Gamma(\lambda_j)=v}|a(\lambda_j)|^2
\right)^{1/2}\notag\\
&\le
\sqrt{L_v}\,
\max_{j:\Gamma(\lambda_j)=v}|a(\lambda_j)|,
\end{align}
where \(L_v=|\{j:\Gamma(\lambda_j)=v\}|\).  If
\(|\rho_i|\le\rho<1\) and every aligned row has provenance weight at least
\(q_0\), then
\begin{equation}
\label{eq-generic-combination-capacity-attenuation}
\left|\mathbb E[H\chi_v(\vartheta)]\right|
\le\sqrt{L_v}\,\rho^{q_0}.
\end{equation}
Moreover, conditional on \(\vartheta\), for every \(t>0\),
\begin{equation}
\label{eq-generic-combination-dependence-tail}
\Pr_{\xi}\!\left\{
\sum_{j=1}^Kc_j
\bigl(Y_{\lambda_j}-\mathbb E_\xi[Y_{\lambda_j}\mid\vartheta]\bigr)
\ge t
\ \middle|\ \vartheta
\right\}
\le
\exp\!\left[-\frac{t^2}{2D_c(\Lambda)}\right].
\end{equation}

Every fan-in-two combination DAG producing one row of provenance weight
\(w\) contains at least \(w-1\) binary combination vertices in the ancestral
sub-DAG of that output.  For several outputs, the complete construction cost
dominates the number of distinct vertices in the union of their ancestral
sub-DAGs.  Thus sharing is charged at its actual represented cost and does not
erase provenance weight or attenuation.

Finally, any represented postprocessing of the complete materialized tuple is
measurable in the tuple algebra and obeys the target-projection monotonicity of
\cref{thm-affordable-composition-no-creation}.  A prospective use of that
postprocessing retains the complete combination record.  Its structural
assignment is the exact or compensated quotient-supported mode when it has
qualified nonidentity fibers, ambient \(\Geom/\Caus\) when it has identity
support and qualified transport, or terminal accounting when it supplies no
authorized intermediate transport.

\end{theorem}

\begin{proof}
Expanding the product gives
\[
Y_\lambda
=\chi_{\Gamma(\lambda)}(\vartheta)
 \prod_{i:\lambda_i=1}\xi_i.
\]
Orthogonality of characters under the uniform law on \(G\), followed by
independence of the \(\xi_i\), proves
\cref{eq-generic-combination-exact-character-transport}.  Multiplying two
derived characters replaces their provenance rows by
\(\lambda_j+\lambda_l\); subtracting the product of their conditional means
proves \cref{eq-generic-combination-exact-covariance}.  Applying
Cauchy--Schwarz to the aligned summands proves
\cref{eq-generic-combination-located-block}; the uniform reliability and
weight hypotheses give
\cref{eq-generic-combination-capacity-attenuation}.

Changing the primitive sign \(\xi_i\) changes the centered weighted sum by
at most
\[
2\sum_{j:(\lambda_j)_i=1}|c_j|.
\]
The bounded-difference inequality with the squared sensitivities summed over
\(i\) gives \cref{eq-generic-combination-dependence-tail}.

A fan-in-two DAG with \(g\) combination vertices can join at most \(g+1\)
distinct primitive leaves into one output.  Hence an output depending on
\(w\) distinct leaves requires \(g\ge w-1\).  Taking the union of ancestral
sub-DAGs counts a shared vertex once and proves the multi-output statement.
The cost assertion is then
\cref{def-budgeted-derivability-closure,lem-description-length-charged-resource,thm-paper4-dependency-principle}.
Projection monotonicity is
\cref{thm-affordable-composition-no-creation}; the exhaustive support and
temporal assignments are
\cref{thm-affordable-geom-abs-normal-form,def-pre-hit-temporally-admissible-source,cor-terminal-verifier-not-prospective-structural-activity}.
\end{proof}

\begin{definition}[Saturated located-combination profile]
\label{def-saturated-located-combination-profile}

At a charged ceiling \(B\), define
\begin{equation}
\label{eq-saturated-located-combination-profile}
\operatorname{Comb}_{\mathfrak M,n}^{\mathrm{loc,sat}}(B)
:=
\sup_R
\mathfrak M_n\!\left(
I\longmapsto
\sup_{v\ne0}
\left|\mathbb E_I[H_R\chi_v(\vartheta)]\right|
\right),
\end{equation}
where \(R\) ranges over family-uniform, temporally admissible records of
\cref{def-charged-binary-noisy-combination-record} whose complete cost is at
most \(B\), and \(H_R\) is a normalized materialized located block in that
record.  The supremum is zero when the class is empty.

This is a source-resolved subprofile of the existing saturated profile.  It
does not add a source modality.  Quotient, Geom/Caus, Lift, learning, and
terminal gates remain those of the complete record.

\end{definition}

\begin{corollary}[Capacity--attenuation closure at one charged ceiling]
\label{cor-combination-capacity-attenuation-closure}

Suppose every record contributing to
\(\operatorname{Comb}_{\mathfrak M,n}^{\mathrm{loc,sat}}(B_n)\) materializes
at most \(L_n\) target-aligned rows, each of attenuation at most \(a_n\),
outside a presentation event \(\mathcal E_n\) satisfying
\(\mathfrak M_n(\mathbf1_{\mathcal E_n})\le\varepsilon_n\).  Then
\begin{equation}
\label{eq-combination-capacity-attenuation-closure}
\operatorname{Comb}_{\mathfrak M,n}^{\mathrm{loc,sat}}(B_n)
\le
\min\{1,\sqrt{L_n}a_n+\sqrt{L_n}\varepsilon_n\}.
\end{equation}
In particular, polynomial \(L_n\) and exponentially decreasing
\(a_n,\varepsilon_n\) give a negligible saturated located-combination
profile.

\end{corollary}

\begin{proof}
On the complement of the exceptional event apply
\cref{eq-generic-combination-located-block}; on the exceptional event use
the same inequality with unit attenuation.  Apply the monotone sublinear
family functional and the union bound.  Truncation at one gives the displayed
formula.
\end{proof}

\begin{corollary}[Prospective accountability of a noisy-combination use]
\label{cor-noisy-combination-prospective-accountability}

Let a represented pre-hit selector use a complete record from
\cref{def-charged-binary-noisy-combination-record} at ceiling \(B\).  Its
positive step advantage obeys the existing saturated structural-charge
inequality
\begin{equation}
\label{eq-noisy-combination-selector-accountability}
\bigl(\operatorname{Adv}_t(K_t;\nu_t)\bigr)_+
\le
eH_B(n)
\sum_{c\in\mathfrak J_{\Abs}^5}
\mathcal C_{I,\mathfrak F_5}^{(c)}
\bigl(D_B^{\mathrm{sel}}(n)\bigr).
\end{equation}
The complete combination record is an ancestor of every nonzero term used by
the selector.  Its materialized located-character contribution is bounded by
\cref{eq-generic-combination-located-block}; any further contribution is the
already-defined exact-\(\Abs\), compensated, ambient-\(\Geom/\Caus\), Lift,
or terminal coordinate selected by the complete record's qualification
gates.  Thus combination, repeated cancellation, and reordering introduce no
additional source coordinate.

\end{corollary}

\begin{proof}
Apply
\cref{thm-prospective-selector-structural-charge,thm-no-free-succinct-transition-rule}
to the selector's complete derivation.  Complete ancestry and the located
bound are
\cref{thm-paper4-dependency-principle,thm-charged-noisy-combination-transport}.
The structural assignment is
\cref{thm-affordable-geom-abs-normal-form}.
\end{proof}

\begin{figure}[tbp]
\centering
\resizebox{.98\linewidth}{!}{%
\begin{tikzpicture}[fw diagram,node distance=8mm and 10mm]
  \node[fw source,text width=2.5cm] (primitive)
    {noisy public\\characters};
  \node[fw provenance,text width=2.6cm,right=of primitive] (dag)
    {charged combination\\DAG and locators};
  \node[fw use,text width=2.5cm,above right=7mm and 11mm of dag] (semantic)
    {semantic support\\may cancel};
  \node[fw residual,text width=2.5cm,below right=7mm and 11mm of dag] (prov)
    {noise provenance\\and attenuation remain};
  \node[fw box,text width=2.7cm,right=22mm of dag] (aggregate)
    {materialized capacity\\and dependence};
  \node[fw geom,text width=2.7cm,right=of aggregate] (use)
    {qualified prospective\\source/use};
  \draw[fw arrow] (primitive)--(dag);
  \draw[fw arrow] (dag)--(semantic);
  \draw[fw arrow] (dag)--(prov);
  \draw[fw arrow] (semantic)--(aggregate);
  \draw[fw arrow] (prov)--(aggregate);
  \draw[fw arrow] (aggregate)--(use);
\end{tikzpicture}%
}
\caption{Charged noisy combination separates semantic cancellation from
noise provenance.  Cancellation may simplify the derived character, while
the primitive leaves, attenuation, shared DAG, located outputs, dependence,
and prospective use remain in one saturated record.}
\label{fig-charged-noisy-combination}
\end{figure}

\section{The Trace-Support Admissibility
Dichotomy}\label{the-trace-support-admissibility-dichotomy}
\label{sec:trace-support-admissibility-dichotomy}

\PartI separated extensional observability from affordable represented
derivability. The present section applies that distinction to operators and
records a residual baseline diagnostic for components that are
separately verified to be nonnegative.

Trace-support admissibility distinguishes represented structural support from support visible only in an induced trajectory.
\begin{definition}[Trace-support admissibility]\label{def-trace-support-admissible}

For a finite operator \(L\), the off-diagonal trace support is

\[
\operatorname{supp}_{\operatorname{tr}}(L)=\{(x,z):x\ne z,\ L_{xz}\ne0\}.
\]

The operator is trace-support admissible at budget \(B\) when its support
relation and every quotient cell, potential level set, conductance, or
boundary set used to define its nonzero rates have uniform represented
derivations of total saturated cost at most \(B\). Such relations are
necessarily measurable in the ambient envelope
\(\mathcal F_I\otimes\mathcal F_I\), but measurability alone is not an
admissibility certificate.

\end{definition}

\Cref{fig-trace-support-admissibility} distinguishes a represented support
relation from a relation visible only after an induced trajectory is known.

\begin{figure}[tbp]
\centering
\begin{tikzpicture}[fw diagram]
  \node[fw box,minimum width=3.0cm] (support) at (-4.7,0)
    {transition support\\\(\operatorname{supp}_{\rm tr}(L)\)};
  \node[fw source,minimum width=3.0cm] (deriv) at (-0.9,1.1)
    {uniform represented\\derivation within \(B\)};
  \node[fw residual,minimum width=3.0cm] (trace) at (-0.9,-1.1)
    {trajectory-only\\distinction};
  \node[fw geom,minimum width=3.2cm] (admit) at (3.4,1.1)
    {admissible structural\\support};
  \node[fw terminal,minimum width=3.2cm] (external) at (3.4,-1.1)
    {generated and charged,\\or unavailable at \(B\)};
  \draw[fw arrow] (support) -- (deriv);
  \draw[fw arrow] (support) -- (trace);
  \draw[fw arrow] (deriv) -- (admit);
  \draw[fw arrow] (trace) -- (external);
\end{tikzpicture}
\caption{Trace support is structural at budget \(B\) precisely when its
distinctions have the represented derivation required by
\cref{def-trace-support-admissible}. Measurability or occurrence in a
trajectory alone supplies no admissibility certificate.}
\label{fig-trace-support-admissibility}
\end{figure}

An admissible operator uses only distinctions generated affordably from
the presentation. A jump rule whose cells have no uniform represented
derivation within the declared budget defines no budget-admissible
structural component for the presented problem.

Non-admissibility means that the relation is unavailable as structural
data at the declared budget. An ambient transition table may still
contain such a relation. If a
presentation-relative budget-admissible oracle-free computation can
generate and use the relation from the presented instance, then the
generated relation appears in the lifted observable configuration and
\PartII classifies its movement into the exposed components. If it
cannot be so generated, then any transition table that still points
along it has no exposed provenance; after the exposed projections it is
\(\Res\) rather than a structural component. This operator
classification alone gives no probabilistic baseline bound: the
orthogonal remainder need not be nonnegative, and its interaction with
the other components may affect the full process. Quantitative target
arrival is evaluated on the full positive generator and charged through
\Cref{thm-universal-saturated-profile-accountability}.

The residual baseline records a target-coupling bound only when the residual operator is separately positive.
\begin{definition}[Residual baseline target coupling]\label{def-residual-baseline-coupling}

Let \(J\) be a residual component that has separately been verified to
be a nonnegative transition kernel or rate table from
\(N=X\setminus T\) to \(X\). With
reference law \(\mu_N=\mu(\cdot\mid N)\), its target hit rate is

\[
\operatorname{Hit}_T(J)
=
\sum_{x\in N}\mu_N(x)\sum_{z\in T}J(x,z).
\]

The baseline target coupling at the same total residual activity is

\[
\operatorname{Base}_T(J)
=
\mu(T)\sum_{x\in N}\mu_N(x)\sum_{z\in X}J(x,z).
\]

The residual excess is

\[
\operatorname{Exc}_T(J)
=
\bigl(\operatorname{Hit}_T(J)-\operatorname{Base}_T(J)\bigr)_+.
\]

These quantities are defined only under the stated positivity
hypothesis. They do not apply automatically to the orthogonal component
\(L^{\Res}\) of \Cref{thm-channel-exhaustiveness}, and their definition implies no bound
on \(\operatorname{Exc}_T(J)\).

\end{definition}

The baseline is a comparison statistic for a separately executable kernel.
For the full positive process, the pre-hit hazard identity of
\cref{thm-prospective-selector-accountability} separates neutral target
contact from represented prospective alignment.  The completed first-hit
quotient remains an exact audit of the same process.

\begin{proposition}[Residual interaction is charged by saturation]\label{prop-residual-excess-reclassification}

Let \(L_{\mathcal A}\) be the full positive generator induced by a
uniform presentation-relative algorithm, and decompose it into the five
structural channels and \(L^{\Res}_{\mathcal A}\). For a
declared horizon \(H\), choose a declared pre-hit neutral selector
\(\nu_t\).  Every target-arrival probability \(\alpha_I\) of the full
process satisfies

\begin{equation}
\label{eq-residual-interaction-saturated-charge}
\alpha_I
\le
\operatorname{Enum}_{\nu}(\mathcal A,H)
+
H\operatorname{Sel}^{\mathrm{sat}}_n(B_{\mathcal A}),
\end{equation}

where \(B_{\mathcal A}\) is the complete finite-horizon derivation cost
after the class-preserving clock refinement. The bound includes every effect
of \(L^{\Res}_{\mathcal A}\), including its interaction
with the other components.  Independently, the retrospective audit satisfies

\[
\alpha_I
\le
eH\,m_I^{\mathrm{sat}}
\!\left(\operatorname{Cost}_{\mathrm{sat}}(q_{\mathcal A})\right)
\]

by \cref{lem:exact-affordable-first-hit}.  If the transition rule uses
instance-dependent data with no uniform represented derivation, then it
is not a presentation-relative algorithm in the declared resource
class.

\end{proposition}

\begin{proof}\phantomsection\label{proof-residual-excess-reclassification}

Apply the channel decomposition only as the exact operator identity of
\cref{thm-channel-exhaustiveness}; the transition implementing the pre-hit
sampler belongs to the full positive update before any component expansion.
The exact identity
\eqref{eq-selector-hazard-decomposition}, followed by the definition of the
prospective selector profile, gives
\eqref{eq-residual-interaction-saturated-charge}.  Applying
\cref{lem:exact-affordable-first-hit} after execution gives the separate
accounting inequality.  The derivability definition supplies the final
provenance statement.

\end{proof}

\begin{corollary}[No residual bypass of prospective accounting]\label{cor-residual-family-baseline}

For a uniform task family, including presentation-admissible oracle
interfaces, suppose every budget-\(B\) computation has
\(\operatorname{Enum}_{\nu}(\mathcal A,H_B)\le\beta_n\), and every
temporally admissible selector at the enlarged cost has advantage at most
\(\delta_n\).  Then the full-process success probability, including every
interaction with \(\Res\), satisfies

\[
\Pr\{\tau\le H_B\}
\le
\beta_n+H_B(n)\delta_n.
\]

Hence residual interaction supplies no bypass around the pre-hit selector
ledger.  Negligibility of the prospective selector profile and of the
declared neutral baseline implies negligible full-process success.  The
orthogonal residual itself need not satisfy a componentwise probabilistic
baseline.

\end{corollary}

\begin{proof}\phantomsection\label{proof-residual-family-baseline}

Apply
\cref{cor-universal-prospective-selector-hardness,prop-residual-excess-reclassification}
to the full positive process.

\end{proof}

This dichotomy separates gaps within the fixed presentation from gaps created by inadmissible extensions.
\begin{definition}[Admissible-gap and inadmissible-gap cases]\label{def-admissibility-dichotomy}

An instance is in the inadmissible-gap case at budget \(B\) when every
operator with transfer-class verified discounted hitting strength uses
at least one quotient cell, potential level set, conductance relation,
or trace-support relation lacking a uniform represented derivation at
that budget. It is in the admissible-gap case when some transfer-class
verified hitting component is trace-support admissible at budget \(B\).

\end{definition}

\begin{lemma}[Inadmissible-gap saturated-profile criterion]\label{thm-inadmissible-gap-hardness}

For an instance family in the inadmissible-gap case, every uniform
budget-\(B\) algorithm still satisfies

\[
\operatorname{Succ}_{\mathfrak M,n}(B)
\le
eH_B(n)
m_{\mathfrak M,n}^{\mathrm{sat}}(\Psi(B,n)).
\]

Thus negligibility of the saturated profile at the enlarged budget
implies negligibility of every budget-admissible discounted target
arrival. Trace-support inadmissibility alone is a provenance statement;
it does not imply a probabilistic baseline estimate for
\(\Res\).

\end{lemma}

\begin{proof}\phantomsection\label{proof-inadmissible-gap-hardness}

Apply \Cref{thm-universal-saturated-profile-accountability}. The first-hit
quotient is generated from the complete computation, so the conclusion
does not depend on whether its one-step transition initially lies in an
exposed channel or in the orthogonal residual. The negligibility
statement follows from
\eqref{eq-universal-profile-lower-bound}.

\end{proof}

\begin{lemma}[Admissible gaps are governed by saturated cost]\label{thm-admissible-gap-degree-governed}

In the admissible-gap case, structured movement is governed by the
growth of \(B^*\). A within-budget admissible gapping quotient or
regular aimed geometry puts \(B^*\) below the declared ceiling; if every
budget-admissible gapping structure has saturated cost above \(b(n)\),
then \(B^*\) is large relative to the declared budget. In the polynomial
instance this is the usual polynomial versus superpolynomial
distinction.

\end{lemma}

\begin{proof}\phantomsection\label{proof-admissible-gap-degree-governed}

When a gapping operator is budget-admissible, all of the structure it
uses has a represented derivation from the public or lifted observable presentation. Part
II's exhaustiveness theorem places that structure in Geom/Caus, Abs, or
Lift. Ambient measurability does not determine budget admissibility. The
least saturated public budget at which the component
contributes extraction at the declared threshold is exactly the
definition of \(B^*\). The hardness criterion of \Cref{the-saturated-cost-invariant} gives the
conclusion.

\end{proof}

\begin{example}[Withheld modulus]\label{ex-hidden-modulus-dichotomy}

In subset sum, a public residue map \(x\mapsto\sum_i a_ix_i\pmod q\) can
define an admissible quotient. If the modulus \(q\) is withheld and no
uniform represented construction derives the residue cells within the
declared budget, then the quotient that would gap the target is
inadmissible at that budget, regardless of its extensional measurability.
A rule that uses the withheld modulus leaves the
budgeted observable algebra of the original presentation; it becomes
admissible only after enlarging the primitive signature so that \(q\)
is public, or through oracle or advice access relative to the original
presentation. It is therefore not a budget-admissible oracle-free
computation for the original task object.

\end{example}

\section{Structural Negligibility Certificates for Budget-Admissible
Generators}\label{structural-floors-for-budget-admissible-generators}
\label{sec:structural-negligibility-certificates}

\subsection{The negligibility certificate}
\label{subsec:negligibility-certificate}

The following regimes yield hardness results once the instance presentation
and \PartII budget-admissibility theorem are fixed.  For a
presentation-relative uniform computation, including one using
presentation-admissible oracle interfaces, the clocked-transcript embedding
produces a generator that is decomposed by \PartII and evaluated by the full
saturated accounting profile. Transition
tables depending on data outside the budgeted saturated observable
closure use presentation-inaccessible or oracle-dependent data, require
over-budget recovery, or require changing the primitive signature so
that the data are public; transition tables whose target contact has
exposed provenance are counted through the corresponding
budget-admissible structure. For a fixed cryptographic instance, the
analysis reduces hardness to \(B^*\) lying above the declared budget; in
the polynomial instance this is the usual superpolynomial cost statement.
The required application theorem must bound the single master
quotient--geometry profile over its complete resource-filtered derivation
family. \PartIV supplies the complementary quantitative learning
analysis, and \PartVIII is a separate specialization.

The structural negligibility certificate bounds the single full saturated
master profile. Universal accountability then applies directly.
\begin{definition}[Structural negligibility certificate for a
family]\label{def-structural-floor-certificate}

Let \(\{I_n\}\) be a presentation-complete task family. A structural
negligibility certificate under the budget structure \(\mathcal B\) is the
family-specific bound

\[
\mathfrak U_n^{\mathrm{sat}}(\Psi(b(n),n))
:=
\sup_{\mathscr C\in
\mathcal C_{QG}(\Psi(b(n),n))}
V_{\mathscr C}
=
\operatorname{negl}_{\mathfrak C}(n);
\]

The supremum ranges over every family-uniform accounting master certificate
generated by the resource-filtered derivation calculus. The bound may be
proved by separate estimates on the
five parameter regimes, but completeness is imposed on their union.  For the
exact \(\Abs\) coordinate, the
\hyperref[cor-five-family-abs-envelope]{exhaustive five-family envelope}
supplies the exhaustive
additive, multiplicative, order, symmetry, and filtration provenance
cover, including mixed contextual charges. The component estimates and
channel normal forms are family-specific applications of the general
framework. Completed first-hit records are included in the accounting
profile, while
\cref{cor-no-retrospective-first-hit-source-loophole} prevents their
retrospective representation from being used as an independent source.

\end{definition}




Structural negligibility yields presentation-relative hardness through universal accountability.
\begin{theorem}[Structural negligibility implies presentation-relative
hardness]\label{thm-structural-floor-implies-hardness}

If \(\{I_n\}\) has a structural negligibility certificate under
\(\mathcal B\), then

\[
\mathfrak U_n^{\mathrm{sat}}(\Psi(b(n),n))
=
\operatorname{negl}_{\mathfrak C}(n),
\qquad
I_n^{\mathrm{sat}}(\Psi(b(n),n))
\le2\mathfrak U_n^{\mathrm{sat}}(\Psi(b(n),n))
=\operatorname{negl}_{\mathfrak C}(n).
\]

Moreover, every presentation-relative uniform generator lying within
\(\mathcal B\), including computations using presentation-admissible oracle
interfaces, has negligible target arrival.

\end{theorem}

\begin{proof}\phantomsection\label{proof-structural-floor-implies-hardness}

The master normal-form theorem places every dynamically qualified accounting
Geom/Abs certificate at the enlarged budget in
\(\mathcal C_{QG}(\Psi(b(n),n))\). The assumed supremum therefore bounds all
five parameter regimes simultaneously. Equation
\eqref{eq-master-and-combined-profile-equivalence} gives
\(I_n^{\mathrm{sat}}\le2\mathfrak U_n^{\mathrm{sat}}\), proving the
displayed negligibility statement. \PartII's exact-Lift identity and
non-generativity hierarchy determine structural provenance but do not
erase contextual charge. The master certificate retains every
\(\Caus\), \(\Lift\), and \(\Bdry\) field at
full carrier cost; a smaller-source transfer is used only when its
represented comparison is part of that certificate.

Encoding adequacy and Lemma
\hyperref[lem-budgeted-generation-gives-representation]{Budgeted
generation gives budgeted representation} place any uniform
computation whose transcript lies within budget inside the same
saturated budget accounting. The universal saturated-profile accountability
theorem \cref{thm-universal-saturated-profile-accountability} bounds every
target-arrival probability by the full exact-quotient accounting profile times the
declared horizon factor. Transfer-class closure preserves negligibility. The
exact first-hit quotient is the audit witness in that implication; by
\cref{cor-no-retrospective-first-hit-source-loophole}, its retrospective
representation supplies no independent source, and any independent pre-hit
representation is already included in the assumed saturated supremum.

\end{proof}

\subsection{Certificate regimes}
\label{subsec:negligibility-certificate-regimes}

The first regime treats target placement that is random relative to the represented structural family.
\begin{lemma}[Random target regime]\label{thm-random-regime}

Put \(N=2^n\), and let \(y\in\{-1,1\}^N\) have independent Rademacher
coordinates. Let \(\mathcal K_{n,\le B}^{\mathrm{sat}}\) be a family of
positive-semidefinite contractions that is independent of \(y\) and
contains the quotient and Geom/Caus visibility kernels under
consideration. Suppose it has an operator-norm \(\varepsilon\)-net of
cardinality \(M_B\). Define

\[
r_1(B)=\sup_{K\in\mathcal K_{n,\le B}^{\mathrm{sat}}}
\frac{\operatorname{tr}K}{N},
\qquad
r_2(B)=\sup_{K\in\mathcal K_{n,\le B}^{\mathrm{sat}}}
\frac{\|K\|_{\mathrm F}}{N}.
\]

There is an absolute constant \(C\) such that, with probability at
least \(1-\delta\),

\begin{equation}
\label{eq-random-target-kernel-bound}
\sup_{K\in\mathcal K_{n,\le B}^{\mathrm{sat}}}
\frac{\langle y,Ky\rangle}{\|y\|_2^2}
\le
r_1(B)
+C\left[
r_2(B)\sqrt{\log\frac{2M_B}{\delta}}
+\frac{\log(2M_B/\delta)}{N}
\right]
+\varepsilon.
\end{equation}

If the right-hand side is
\(\operatorname{negl}_{\mathfrak C}(n)\), then each structural
visibility coordinate indexed by this kernel family is negligible and
\(I_n^{\mathrm{sat}}(B)\) is at most twice the right-hand side. The same
conclusion holds for a random fixed-density target contrast only after
supplying the analogous sampling-without-replacement quadratic-form
estimate. In particular, cover entropy alone is insufficient: the
normalized trace and Frobenius terms must also be small. Any
subexponential-budget lower bound is conditional on these quantities
being negligible uniformly throughout that budget range.

\end{lemma}

\begin{proof}\phantomsection\label{proof-random-regime}

For a fixed real symmetric \(K\), the Hanson--Wright inequality gives,
for every \(t>0\),

\[
\mathbb P\!\left\{
y^\mathsf TKy-\operatorname{tr}K
>C\bigl(\|K\|_{\mathrm F}\sqrt t+\|K\|_{\mathrm{op}}t\bigr)
\right\}
\le2e^{-t}
\]

for an absolute \(C\) \citep{RudelsonVershynin2013}. Apply this estimate
to every member \(K_0\) of the declared net with
\(t=\log(2M_B/\delta)\), use \(\|K_0\|_{\mathrm{op}}\le1\), and take a
union bound. For an arbitrary \(K\) choose a net point \(K_0\) with
\(\|K-K_0\|_{\mathrm{op}}\le\varepsilon\). Since
\(\|y\|_2^2=N\),

\[
\frac{|y^\mathsf T(K-K_0)y|}{N}
\le\|K-K_0\|_{\mathrm{op}}
\le\varepsilon.
\]

Dividing the simultaneous fixed-net bounds by \(N\) and taking the
supremum proves \eqref{eq-random-target-kernel-bound}. If both the Abs
and Geom/Caus coordinates use kernels in this family, each is bounded
by the same right-hand side; their sum is therefore bounded by twice
that quantity.

\end{proof}

\begin{lemma}[Presentation-inaccessible structure regime]\label{thm-removed-structure-regime}

If every transfer-class gapping structure requires data absent from the
declared represented closure \(\mathcal R_n^{\mathcal B}\), then
no budget-admissible operator in the class has a transfer-class gapping
component.

\end{lemma}

\begin{proof}\phantomsection\label{proof-removed-structure-regime}

This is the provenance half of Theorem
\hyperref[thm-inadmissible-gap-hardness]{Inadmissible-gap saturated-profile
criterion}. A datum without a uniform represented derivation cannot be
used by a presentation-relative operator at the declared budget. No
ambient measurability or degree estimate supplies the missing derivation.

\end{proof}

\begin{lemma}[Adversarial placement against the declared geometry]\label{thm-adversarial-placement-regime}

Suppose the only budget-admissible non-residual exposed structural
component below saturated budget \(B\) is a pointwise geometry. Assume
either that every such record fails its consistency condition, or that
there is one negligible function \(\nu_B(n)\) such that, uniformly over
all budget-admissible geometry records at that budget,

\[
\sup_{(\mathfrak N,\Phi,D,J,\mathfrak K)
\in\mathcal G^{X,\mathrm{free}}_{n,\le B}}R_T(D)
\le\nu_B(n).
\]

Then
\(G_n^{\mathrm{sat}}(B)=\operatorname{negl}_{\mathfrak C}(n)\). If also
\(m_{\mathfrak M,n}^{\mathrm{sat}}(\Psi(B,n))\) is negligible, every
budget-\(B\) algorithm has target arrival negligible after the explicit
horizon factor \(eH_B\).

\end{lemma}

\begin{proof}\phantomsection\label{proof-adversarial-placement-regime}

Failure of the consistency condition sets the corresponding term in
\(G_n^{\mathrm{sat}}(B)\) to zero. In the second case the single
uniform bound controls the numerator \(R_T(D)\) for every remaining
record. Since
\((1+\rho_\Phi(D))^{-1}\le1\) and the mixing indicator is at most one, the
supremum defining \(G_n^{\mathrm{sat}}(B)\) is negligible. If
the uniform saturated quotient profile is also negligible at the
enlarged first-hit budget, \Cref{thm-universal-saturated-profile-accountability} bounds the full
target-arrival probability. This last step is applied to the complete
positive process and therefore includes residual interactions.

\end{proof}

\begin{lemma}[Saturated-cost lower-bound
regime]\label{thm-proven-degree-regime}

Assume a family has a proven saturated-cost threshold \(B_0(n)\): there
is one negligible function \(\nu(n)\) such that for its centered target
contrast,

\[
\sup_{B<B_0(n)}\,
\sup_{K\in\mathcal K_{n,\le B}^{\mathrm{sat}}}
\frac{\langle y,Ky\rangle}{\|y\|^2}
\le\nu(n).
\]

Then

\[
I_n^{\mathrm{sat}}(B)=\operatorname{negl}_{\mathfrak C}(n)\quad(B<B_0(n)),
\qquad
B^*\ge B_0(n).
\]

In the raw Boolean direct-coordinate subcase, if a raw Walsh support
threshold is \(r_0(n)\), then exhaustive coordinate scanning below that
threshold has cost
\(\operatorname{Cost}_W(n,r_0)=\sum_{k<r_0(n)}\binom{n}{k}\). This is
one charged coordinate inside \(B_0(n)\), rather than the complete
budget-admissibility definition.

\end{lemma}

\begin{proof}\phantomsection\label{proof-proven-degree-regime}

By \cref{prop-saturated-degree-audit}, the quotient and Geom/Caus structural visibility
coordinates are bounded above by the displayed target quadratic forms
over the budget-admissible saturated visibility kernels. The hypothesis
therefore gives
\(I_n^{\mathrm{sat}}(B)=\operatorname{negl}_{\mathfrak C}(n)\) for every
\(B<B_0(n)\). The definition of \(B^*\) implies \(B^*\ge B_0\) up to the
chosen transfer-class threshold; for \(\mathcal B_{\mathrm{poly}}\) this
is the inverse-polynomial threshold. In the raw Boolean Walsh subcase,
the dense-band dimension below raw degree \(r_0\) is
\(\sum_{k<r_0}\binom{n}{k}\), which is the displayed direct-scan cost.

\end{proof}

\begin{remark}[Saturated-cost threshold versus raw
degree]\label{rem-degree-floor-not-max-degree}

The hypothesis is a lower support theorem for the saturated
budget-admissible kernel family. A raw Walsh, nonlinearity, or
algebraic-immunity statement supplies this hypothesis only when it
controls every budget-admissible represented quotient, public basis
change, valid lift, and geometric/causal kernel in the saturated class.
Otherwise it remains a diagnostic for one coordinate system and does
not establish that \(B^*\) lies above the declared budget; in
\(\mathcal B_{\mathrm{poly}}\), this is the usual superpolynomial
lower-bound statement. Any application claiming such a threshold must
bound the complete derivation-indexed quotient family and all supported
geometric modes. Learning fragments provide quantitative diagnostics
for the subfamilies they cover.

\end{remark}

\section{Complexity Barriers as Presentation-Extension
Tests}\label{the-barrier-theorems-and-the-operator-invariant}
\label{sec:presentation-extension-barrier-tests}

\subsection{Oracle extensions and presentation identity}
\label{subsec:oracle-presentation-identity}

The invariant measures spectral transport relative to a task
presentation. The primitive signature is part of the task object, so an
oracle extension must be interpreted relative to that signature. An
oracle already generated by the saturated closure is a conservative
interface. An oracle not generated by that closure enlarges the public
presentation and defines a different presented task. Classical
relativization, natural-proofs, and algebrization results test claims
that remain valid under specified presentation extensions; they do not
identify oracle access as an unclassified source of computation.

\Cref{fig-presentation-barrier-tests} gives the common presentation-relative
form of the three barrier tests discussed in this section.

\begin{figure}[tbp]
\centering
\begin{tikzpicture}[fw diagram]
  \node[fw source,minimum width=3.2cm] (task) at (-4.8,0)
    {presented task\\\(\mathcal T,\mathcal D_n\)};
  \node[fw provenance,minimum width=2.8cm] (oracle) at (-0.9,1.4)
    {oracle interface};
  \node[fw geom,minimum width=2.8cm] (natural) at (-0.9,0)
    {property/test class};
  \node[fw use,minimum width=2.8cm] (alg) at (-0.9,-1.4)
    {algebraic interface};
  \node[fw terminal,minimum width=3.7cm] (test) at (3.5,0)
    {compile inside \(\mathcal D_n\),\\or obtain a new presentation};
  \draw[fw arrow] (task) -- (oracle);
  \draw[fw arrow] (task) -- (natural);
  \draw[fw arrow] (task) -- (alg);
  \draw[fw arrow] (oracle) -- (test);
  \draw[fw arrow] (natural) -- (test);
  \draw[fw arrow] (alg) -- (test);
\end{tikzpicture}
\caption{Relativization, natural-proofs, and algebrization are evaluated as
tests of a specified extension. A represented class-preserving interface
compiles into the original closure; an inadmissible interface changes the
presented derivation system, as formalized in
\cref{thm-oracle-admissibility-presentation-separation}.}
\label{fig-presentation-barrier-tests}
\end{figure}

The oracle-extended closure records precisely the presentation obtained by adjoining a charged query interface.
\begin{definition}[Oracle-extended derivability closure]\label{def-oracle-extended-closure}

For a presented task \(\mathcal T\) and an oracle \(O\), let
\(\mathcal T^O\) be the presentation obtained by adjoining the query
interface \(O\), and let \(\mathcal D^O_n\) be Definition
\hyperref[def-budgeted-derivability-closure]{Budgeted derivability
closure} with one extra primitive: oracle-query operations
\(\omega\mapsto O(g(\omega))\) for represented queries \(g\), charged
once per call plus the cost of deriving \(g\). All constructors,
minimality, computation-closure, and structural-provenance clauses are
unchanged. This query charge prices access to the declared interface;
it does not charge construction of \(O\). Consequently
\(\mathcal T^O\) is the same effective task as \(\mathcal T\) only when
the oracle interface has a class-preserving represented derivation from
the original presentation.

\end{definition}

This theorem distinguishes admissible oracle access from an oracle extension that changes the presentation.
\begin{theorem}[Oracle admissibility and presentation separation]
\label{thm-oracle-admissibility-presentation-separation}

Suppose the oracle evaluator \(x\mapsto O(x)\) has one family-uniform
represented derivation from \(\mathcal T\), including its retained data
and evaluation rule. If substituting this derivation for every oracle
call has class-preserving accumulated cost \(\Psi_O(B,n)\), then every
record in \(\mathcal D^O_{n,\le B}\) compiles to a record in
\(\mathcal D_{n,\le\Psi_O(B,n)}\) with the same denotation. In
particular,

\[
I_{\mathcal T^O,n}^{\mathrm{sat}}(B)
\le
I_{\mathcal T,n}^{\mathrm{sat}}\bigl(\Psi_O(B,n)\bigr).
\]

The extension is then conservative for the declared resource class. If
no such derivation exists at the declared budget, then

\[
\mathcal T^O\ne\mathcal T,
\qquad
\mathcal D^O_n\ne\mathcal D_n
\quad\text{as presented derivation systems}.
\]

Indeed, the oracle evaluator is a primitive of \(\mathcal T^O\) and is
not an admissible primitive of \(\mathcal T\). Every target-aligned
oracle output is classified and charged as a source in the saturated
profile of \(\mathcal T^O\). An instance-dependent oracle table not
produced by one family-uniform rule is nonuniform advice and likewise
defines a different presented family.

\end{theorem}

\begin{proof}

Replace each oracle-query vertex by the represented evaluator for
\(O\). Closure under composition preserves the denotation, and the
complete substituted derivation has cost at most \(\Psi_O(B,n)\).
Applying the substitution to every saturated certificate proves the
profile inequality. In the absence of a represented evaluator, the
\hyperref[prop-oracle-accounting]{Oracle accounting proposition} places
the output outside the original derivability closure. Once declared
primitive, the
\hyperref[thm-channel-exhaustiveness]{Component exhaustiveness theorem}
and the
\hyperref[thm-five-family-abs-provenance]{Universal five-family evaluator
provenance normal-form theorem} classify its state-dependent
contribution in the enlarged presentation.

\end{proof}

\Cref{fig-oracle-separation} shows the two cases of
\Cref{thm-oracle-admissibility-presentation-separation}.

\begin{figure}[tbp]
\centering
\begin{tikzpicture}[fw diagram]
  \node[fw box,minimum width=2.1cm] (oracle) at (-5.0,0) {oracle \(O\)};
  \node[fw source,minimum width=2.7cm] (eval) at (-1.9,1.1)
    {represented evaluator};
  \node[fw geom,minimum width=3.5cm] (same) at (3.3,1.1)
    {\(\mathcal D^O_{n,\le B}\to\mathcal D_{n,\le\Psi_O}\)};
  \node[fw residual,minimum width=2.7cm] (noeval) at (-1.9,-1.1)
    {no represented evaluator};
  \node[fw terminal,minimum width=3.5cm] (new) at (3.3,-1.1)
    {\(\mathcal T^O\ne\mathcal T\)};
  \draw[fw arrow] (oracle) -- node[above left,font=\scriptsize] {admissible} (eval);
  \draw[fw arrow] (eval) -- node[above,font=\scriptsize] {compile} (same);
  \draw[fw arrow] (oracle) -- node[below left,font=\scriptsize] {inadmissible} (noeval);
  \draw[fw arrow] (noeval) -- node[below,font=\scriptsize] {enrich} (new);
\end{tikzpicture}
\caption{An admissible oracle compiles into the original closure; an
inadmissible oracle defines an enriched presentation, as formalized in
\cref{thm-oracle-admissibility-presentation-separation}.}
\label{fig-oracle-separation}
\end{figure}

\begin{corollary}[Presentation-coarse classifications cannot separate
\(P\) from \(NP\)]
\label{cor-presentation-coarse-no-p-np-classification}

Let \(\mathfrak C\) be a classification of presented problem families
that assigns the same class to \(\mathcal T\) and \(\mathcal T^O\)
whenever \(O\) is inadmissible for \(\mathcal T\). For every presented
\(NP\) search or decision family \(\mathcal T\), adjoin an oracle
\(O_{\mathcal T}\) that returns a valid witness or the verified decision
bit. Then \(\mathcal T^{O_{\mathcal T}}\) has a uniform polynomial
solver in its enriched presentation. Therefore \(\mathfrak C\) cannot
be both sound and complete for polynomial-time solvability on presented
\(NP\) families unless \(P=NP\). Equivalently, a classification that
identifies a family with all of its inadmissible oracle enrichments
cannot distinguish the original \(P\)-versus-\(NP\) question.

\end{corollary}

\begin{proof}

The witness oracle supplies a polynomial query, verification, and
readout procedure in \(\mathcal T^{O_{\mathcal T}}\). If it is
admissible for \(\mathcal T\), the
\hyperref[thm-oracle-admissibility-presentation-separation]{Oracle
admissibility and presentation separation theorem}
compiles that procedure into a polynomial solver for \(\mathcal T\).
If it is inadmissible, the same theorem gives
\(\mathcal T^{O_{\mathcal T}}\ne\mathcal T\), while the assumed
classification assigns them the same value. Soundness and completeness
would then classify every presented \(NP\) family as polynomial-time
solvable. Hence such a classifier implies \(P=NP\); absent that
collapse, it cannot distinguish the two complexity classes.

\end{proof}

\subsection{Relativization, natural proofs, and algebrization}
\label{subsec:presentation-barrier-tests}

Relativized accountability applies the same derivability ledger after a declared oracle extension.
\begin{lemma}[Universal derivability accountability relativizes]\label{thm-framework-relativizes}

For every fixed oracle presentation \(\mathcal T^O\), the
\hyperref[thm-encoding-adequacy]{Encoding adequacy theorem}, the
\hyperref[prop-forced-derivable-witness]{Forced derivable witness
proposition}, and the
\hyperref[thm-full-derivability-characterization]{Full derivability
lower-bound criterion} hold verbatim over \(\mathcal D^O_n\). Hence
negligibility of the saturated profile of \(\mathcal T^O\) at the
charged enlarged budget rules out inverse-polynomial target arrival by
uniform polynomial computations in that oracle presentation. The
profile is recomputed after adjoining \(O\); the theorem asserts no
identity between the profiles of \(\mathcal T\) and \(\mathcal T^O\).

\end{lemma}

\begin{proof}\phantomsection\label{proof-framework-relativizes}

An oracle call is a declared primitive constructor of
\(\mathcal T^O\). Encoding Adequacy simulates ordinary and oracle steps
one by one. The forced-witness proof applies the legal clocked-transcript
embedding and first-hit quotient to the induced transition. The
algorithm-to-profile inequality therefore holds inside the enlarged
presentation. The
\hyperref[thm-oracle-admissibility-presentation-separation]{Oracle
admissibility and presentation separation theorem}
determines whether this is a conservative interface for
\(\mathcal T\) or a distinct task presentation.

\end{proof}

\begin{corollary}[Presentation-relative relativization test]\label{cor-bgs-test}

If, for every oracle \(A\), both the
\hyperref[thm-universal-saturated-profile-accountability]{Universal
saturated-profile accountability theorem} and
a negligible-profile proof applied to the standard presentation of an
\(NP^A\)-complete family, they would imply
\(P^A\ne NP^A\) for every oracle \(A\). The Baker--Gill--Solovay theorem
excludes that extension-invariant conclusion
\citep{BakerGillSolovay1975}. In the presentation-relative framework,
the failure occurs in the profile premise: an arbitrary oracle may add
a target-aligned primitive, including a solver or answer table, and
thereby make the saturated profile non-negligible. Conservative oracle
extensions preserve the original conclusion through
\hyperref[thm-oracle-admissibility-presentation-separation]{the Oracle
admissibility and presentation separation theorem}; nonconservative
extensions define different presented
tasks and require their own structural analysis.

\end{corollary}

\begin{proof}

Assume the displayed accountability theorem and its negligible-profile
premise held for the standard presentation of an \(NP^A\)-complete family
for every oracle \(A\).  The profile-to-arrival implication would then rule
out every polynomial-time oracle decider for that family and yield
\(P^A\ne NP^A\) for every \(A\).  Baker--Gill--Solovay provide an oracle for
which \(P^A=NP^A\), a contradiction.  Therefore the extension-uniform
conclusion can fail only because at least one required premise is not
preserved.  In this framework the presentation-separation theorem locates
that failure in the recomputed profile when the oracle contributes a
target-aligned primitive; for a conservative extension its simulation
hypothesis instead preserves the original profile conclusion.

\end{proof}

\begin{lemma}[Presentation-relative natural-proofs test]
\label{thm-natural-proofs-position}

The predicate \(\Phi_n(\mathcal T)\) that a presented task has
superpolynomial full-closure \(B^*\) is a property of the complete
presentation, not of an unadorned Boolean truth table. Its value may
change when advice, a circuit family, or another target-aligned
primitive is adjoined. The
\hyperref[thm-random-regime]{Random target regime theorem} proves
largeness for
the declared random sparse presentation under its cover-entropy
hypothesis. The Razborov--Rudich barrier applies only to a derived
Boolean-function property that is simultaneously polynomial-time
constructive from the truth table, large, and useful against the stated
nonuniform circuit class \citep{RazborovRudich1997}. A
fixed-presentation proof of saturated-profile negligibility does not
acquire those three properties from the generic framework.

\end{lemma}

\begin{proof}\phantomsection\label{proof-natural-proofs-position}

Largeness is the conclusion of the
\hyperref[thm-random-regime]{Random target regime theorem} for the
specified presented kernel family. Usefulness in the present framework
is relative to the declared uniform saturated closure. Passing to
\(P/\mathrm{poly}\) adjoins nonuniform circuit data and changes the
computational authority being classified. If one additionally constructs
a polynomial-time truth-table test whose acceptance is large and
excludes that nonuniform circuit class, the resulting truth-table
property satisfies the hypotheses of Razborov--Rudich. In the absence of
that additional construction, the theorem remains a
presentation-relative structural statement, and the natural-proofs
hypotheses are not present.

\end{proof}

The algebraic-oracle closure records the presentation obtained by adjoining algebraic oracle primitives.
\begin{definition}[Algebraic-oracle closure]\label{def-algebraic-oracle-closure}

For an algebraic oracle \(\widetilde O\), such as a low-degree extension
or interpolation operation, define the enriched presentation
\(\mathcal T^{\widetilde O}\) by adjoining evaluations of
\(\widetilde O\) on represented queries as primitive operations and
leaving the remaining closure rules unchanged. The extension is
conservative precisely when these evaluations have a class-preserving
represented derivation in \(\mathcal T\).

\end{definition}

\begin{lemma}[Algebraic-oracle accountability and presentation
separation]\label{thm-framework-algebrizes}

The \hyperref[thm-encoding-adequacy]{Encoding adequacy theorem}, the
\hyperref[prop-forced-derivable-witness]{Forced derivable witness
proposition}, and the
\hyperref[thm-full-derivability-characterization]{Full derivability
lower-bound criterion} hold verbatim for the presentation
\(\mathcal T^{\widetilde O}\). Therefore the
accountability direction algebrizes in the Aaronson--Wigderson sense
\citep{AaronsonWigderson2009}. Its conclusion is evaluated using the
saturated profile of the enriched presentation. When
\(\widetilde O\) has a class-preserving represented evaluator in
\(\mathcal T\), the extension compiles conservatively into
\(\mathcal T\). Otherwise
\(\mathcal T^{\widetilde O}\ne\mathcal T\), and every structural source
introduced by the algebraic oracle is charged in the enriched profile.
Thus algebrization preserves accountability inside each presentation;
it does not preserve a family-specific profile-negligibility premise
across nonconservative algebraic extensions.

\end{lemma}

\begin{proof}\phantomsection\label{proof-framework-algebrizes}

The \hyperref[thm-encoding-adequacy]{Encoding adequacy theorem} and the
\hyperref[lem:exact-affordable-first-hit]{Exact budget-admissible
first-hit quotient lemma} operate on the finite update induced by the
declared primitive operations. An
algebraic-oracle evaluation is one such operation in
\(\mathcal T^{\widetilde O}\), so the algorithm-to-profile implication
holds in that presentation. The profile premise must then be proved for
the enriched saturated closure. Applying
\hyperref[thm-oracle-admissibility-presentation-separation]{Oracle
admissibility and presentation separation theorem} to the evaluator
\(\widetilde O\) gives the
two stated cases.

\end{proof}

\begin{proposition}[Barrier interpretation for fixed presentations]
\label{prop-barrier-compliant-routes}

A full-closure lower bound for a fixed presentation follows from a
family-specific proof that its saturated profile is negligible at the
declared budget. No additional barrier hypothesis is required for that
fixed-presentation implication. Classical relativization,
natural-proofs, and algebrization barriers constrain stronger statements
that discard the presentation or require the same structural premise to
survive nonconservative changes of computational authority. A specified
extension is handled by
\hyperref[thm-oracle-admissibility-presentation-separation]{Oracle
admissibility and presentation separation theorem}: conservative interfaces compile into the
original closure, while nonconservative interfaces define a different
family whose new sources require their own analysis.

\end{proposition}

\begin{proof}\phantomsection\label{proof-barrier-compliant-routes}

The generic accountability implication is evaluated in whichever
presentation supplies the computation. A family-specific lower bound
proves its profile premise in the declared closure. A conservative
extension preserves that premise through the compilation inequality. A
nonconservative extension adds primitive information and changes the
family, so the original premise is not a statement about the enriched
task. The classical barriers arise precisely when these distinct
presentations are identified or when one demands one premise uniformly
across all of them. The
\hyperref[cor-presentation-coarse-no-p-np-classification]{presentation-coarse
classification corollary} shows why such identification cannot classify
polynomial solvability without collapsing \(P\) and \(NP\).

\end{proof}

\begin{remark}[Scope]\label{rem-barrier-scope}

The barrier results distinguish a fixed-presentation theorem from a
claim invariant under changes of computational authority.  The LPN
specialization of this scope statement is indexed in
\cref{sec-lpn-complexity-barriers}.

\end{remark}

\section{Examples: Discounted Hitting for SAT, Subset Sum, and Degree
Scans}\label{worked-the-stopwatch-on-sat-subset-sum-and-degree-scans}
\label{sec:spectral-complexity-examples}

The examples below are static spectral tests of the hitting-time analysis.
They compute how much target contrast is visible at low degree in
selected observable bases and how much is exposed by public comparison
data generated from the SAT presentation, without access to a target
oracle.

\begin{example}[XOR-SAT]\label{ex-xor-sat-dstar}

For public XOR-SAT, the syndrome map is a represented observable with
affine fibers. Projecting onto the public syndrome quotient captures the
satisfying sector exactly. A raw-coordinate Walsh diagnostic may place
the same contrast in higher coordinate degrees when equations involve
many variables, but the public quotient is the budget-admissible basis
and controls the invariant.

\end{example}

\begin{example}[Random 3-SAT model]\label{ex-random-3sat-jres}

In the random-target approximation to random 3-SAT, raw low-degree Walsh
projections of the centered success contrast are negligible with high
probability for the declared basis family. If the clause-defect geometry
also fails the regular target-orientation condition, then the tested raw
low-degree and clause-geometry fragment profiles are negligible at
polynomial saturated budgets. Negligibility of the full
\(m_n^{\mathrm{sat}}(B)\) or \(G_n^{\mathrm{sat}}(B)\) additionally
requires a uniform exhaustive cover of their saturated certificate families. Any
operator portion outside the tested exposed fragments remains in the
residual ledger, while its full target effect is accounted for by the
exact first-hit quotient.

\end{example}

\begin{example}[Subset sum with public or withheld modulus]\label{ex-subset-exposed-hidden}

If the residue map \(x\mapsto\sum_i a_ix_i\pmod q\) is public, the
residue quotient is budget-admissible when its representation fits the
declared ceiling; in \(\mathcal B_{\mathrm{poly}}\) this can make
\(B^*\) polynomial. If the modulus and residue map are withheld, the
same quotient is not generated in the declared budgeted saturated
observable closure of the original presentation. Its quotient is then
budget-inadmissible rather than an admissible high-cost representation.

\end{example}

\begin{example}[XOR, Gaussian elimination, and public closures]\label{ex-cheap-high-degree-objections}

Budget-admissible computations whose raw Walsh support is high are
handled by admissible public basis changes and generated observables; in
the polynomial instance these are the ordinary polynomial computations.
Gaussian elimination on XOR-SAT computes the public syndrome quotient,
so its target signal is charged to \(m_n^{\mathrm{sat}}(B)\) in the
syndrome basis rather than left in \(\Res\). The implication
graph for 2-SAT and forward chaining for Horn-SAT are generated from
exposed clauses and are charged as Caus over public graph/closure
geometry \citep{Schaefer1978,AspvallPlassTarjan1979,DowlingGallier1984},
with Bdry only marking terminal satisfied states. A public
weighted-sum or residue map for subset-sum gives an Abs quotient; the
same map, when absent from the presentation and not generated within the
declared budget from exposed data, is presentation-inaccessible and
budget-inadmissible rather than residual structure.

\end{example}

The raw Walsh scan records the degree statistic used by the finite diagnostic.
\begin{definition}[Raw Walsh degree scan]\label{def-delta-I-estimator}

For the dense raw Walsh diagnostic, the empirical raw-degree increment
is

\[
\Delta I(d)=\widehat I(d)-\widehat I(d-1),
\qquad
\widehat I(d)=\widehat m(d)+\widehat G(d).
\]

For a sample of \(M\) states and a raw degree band of dense scan cost
\(N_d=\sum_{k\le d}\binom{n}{k}\), Walsh coefficient estimates
concentrate at scale

\[
O\!\left(\sqrt{\frac{N_d+\log(1/\delta)}{M}}\right).
\]

The scan certifies flatness only for this raw-coordinate cost model, up
to degrees whose dense band dimension lies within budget and whose sampling
floor is below the claimed signal. It is a diagnostic coordinate inside
the saturated profile rather than the complete budget-admissibility
definition.

\end{definition}

\section{Summary: Hardness as Saturated Extraction
Cost}\label{summary-hardness-as-saturated-extraction-cost}
\label{sec:spectral-complexity-summary}

The \PartII structural decomposition yields a quantitative theory of
target-arrival time. Its Laplace transform is the target-projected killed
resolvent. For self-adjoint killed dynamics satisfying the stated
comparison hypotheses, the resolvent, killed-target gap, and arrival
functional are comparable on the declared window. General
nonreversible dynamics retain the exact normalized resolvent-defect
bound. The resulting estimate depends on the extraction available
within the saturated public budget. These hypotheses classify the
induced operator and need not be supplied by the algorithm that induces
it.

\paragraph{Dependency trail.}
The global encoding, resolvent, and channel coordinates are
\cref{def-paper3-global-notation,def-presentation-complete-encoding,def-target-resolvent,def-budgeted-stopwatch-comparison-data,def-channel-stopwatch-components,def-transformed-gradient-space}.
The spectral and saturated ledgers are
\cref{def-degree-projection-spectral-mass,def-degree-d-target-comparison-vector,def-degree-d-admissible-operator,def-three-cost-ledgers,def-no-free-committor}.
The later cost, kernel, admissibility, and oracle coordinates are
\cref{def-alignment-complexity,def-reversible-compressed-fragment,def-admissible-nonlinear-quotient-kernel,def-observable-structural-accounting,def-trace-support-admissible,def-residual-baseline-coupling,def-admissibility-dichotomy,def-algebraic-oracle-closure,def-delta-I-estimator}.
Resolvent and channel timing are provided by
\cref{prop-budgeted-procedures-covered-before-comparison-data,thm-two-discount,thm-directed-drift-transformed-gradient-normal-form,thm-gap-spectral-bound}.
The geometric, causal, and defect comparisons are
\cref{prop-compressed-authority-split,cor-two-scale-killed-gap,cor-represented-geometry-compensates-defect,thm-continuum-caus-authority-bridge,cor-authorized-flux-dissipation-polarity,lem-declared-geom-form-error,prop-nonquotient-nonlinear-geometry-closure,prop-quotient-projection-spectral-split}.
The saturated-history and source-envelope results are
\cref{thm-no-truncated-neumann-bypass,thm-budgeted-full-history-gaplessness,thm-finite-dynamic-source-envelope,cor-geom-abs-derivation-compatibility,thm-saturated-ambient-geom-accountability,cor-saturated-ambient-geom-negligibility}.
The cost-invariant, nonlinear-kernel, and universal-accounting consequences are
\cref{prop-dstar-intrinsic,prop-no-broad-narrow-dilemma,thm-alignment-sandwich,thm-reversible-fragment-hardness,prop-kernel-trick-normal-form,thm-nonlinear-kernel-audit-algorithm-independent,prop-structural-accounting-is-universal,lem-affordable-source,thm-admissible-gap-degree-governed}.
The certificate regimes and presentation-extension tests are
\cref{thm-removed-structure-regime,thm-adversarial-placement-regime,thm-proven-degree-regime,thm-framework-relativizes,cor-bgs-test,thm-natural-proofs-position,thm-framework-algebrizes,prop-barrier-compliant-routes}.
Conventions and qualifications used in these arguments are indexed by
\cref{rem-cost-frontier-convention,rem-committor-derivability,rem-basis-robustness,rem-imposed-dissipation-and-transport,rem-forward-transport-not-inert,rem-committor-cases,rem-projection-ancestry-dynamic-source-transfer,rem-structural-witnesses-do-not-assign-charge,rem-contextual-charge-order-dependence,rem-dynamic-transfer-from-contextual-charge,rem-saturated-residual-accountability,rem-characterization-not-lower-bound-method,rem-characterization-role,rem-degree-floor-not-max-degree,rem-barrier-scope}.
The corresponding worked examples are
\cref{ex-paper3-sat-stopwatch,ex-arrival-sat,ex-resolvent-sat,ex-channel-stopwatch-sat,ex-walsh-bound-xor,ex-lumpability-sat,ex-combined-sat,ex-dstar-sat,ex-hidden-modulus-dichotomy,ex-xor-sat-dstar,ex-random-3sat-jres,ex-subset-exposed-hidden,ex-cheap-high-degree-objections}.

The summary theorem collects the accountability and structural-negligibility consequences used by the application parts.
\begin{theorem}[Saturated-profile hardness summary]\label{thm-paper3-summary}

Fix a task family, its declared budget ceiling \(B\), and a family functional
\(\mathfrak M_n\).  Then

\[
\operatorname{Succ}_{\mathfrak M,n}(B)
\le
eH_B(n)
m_{\mathfrak M,n}^{\mathrm{sat}}(\Psi(B,n))
\le
eH_B(n)
I_{\mathfrak M,n}^{\mathrm{sat}}(\Psi(B,n)).
\]

Thus negligibility of the saturated profile at the displayed
class-preserving enlargement implies negligible target arrival for the
declared operational class.  For \(\mathcal B_{\mathrm{poly}}\), the
statement is evaluated at \(B=b_{\mathrm{poly}}\) and
\(\Psi(b_{\mathrm{poly}},n)\); no additional polynomial-budget quantifier is
introduced.

\end{theorem}

\begin{proof}\phantomsection\label{proof-paper3-summary}

Apply \Cref{thm-universal-saturated-profile-accountability} and then
\(m_{\mathfrak M,n}^{\mathrm{sat}}
\le I_{\mathfrak M,n}^{\mathrm{sat}}\). The first-hit proof operates on
the full positive generator before channel decomposition, so the result
includes all five channels, their interactions, and
\(\Res\) without a componentwise positivity assumption.

\end{proof}

The principal results of this part are the budgeted derivability closure,
the arrival functional, the discounted hitting functional, the
target-projected killed resolvent, the budgeted comparison-data
equivalence, the exact structural operator ledger, the full-process
first-hit accountability theorem, the directed-drift
transformed-gradient normal form, the
budget-parametrized full-history composition bound, and the spectral comparison-vector bound
with declared transfer-class comparison factor, the normalized
nonreversible defect bound, the reversible Schur-complement compensation
theorem, the Caus drift-margin theorem, the five-regime master
quotient--geometry normal form and exact-quotient domination theorem, the extraction functional
\(I_n^{\mathrm{sat}}(B)=m_n^{\mathrm{sat}}(B)+G_n^{\mathrm{sat}}(B)\),
the saturated ambient-geometry accountability and inverse-polynomial witness
criterion after quotient exhaustion,
the universal five-family Abs provenance normal form and its exhaustive
contextual-charge envelope,
the intrinsic cost invariant \(B^*\) and log-cost coordinate \(d^*\),
the universal derivability accountability theorem, the alignment-complexity
sandwich, the reversible compressed-fragment theorem, the admissibility
dichotomy, and hardness in the exposed-structure regimes.

For a fixed transfer-class-gap cryptographic instance under the
polynomial budget, the results give the reduction

\[
\begin{gathered}
B^*\text{ is superpolynomial in the saturated represented algebra}
\\[-1mm]
\Longrightarrow\\[-1mm]
\text{hardness against presentation-relative movement}.
\end{gathered}
\]

The full closure gives an unconditional presentation-relative
lower-bound criterion at a chosen budget: if \(B^*\) lies above the declared ceiling, up to
the allowed transfer-class overheads, target arrival is absent in the
legal derivability closure at that budget. Fragment and family-specific
theorems supply sufficient conditions by proving that the
relevant saturated subclosure has no transfer-class target-aligned
representative. \PartIV uses the same \(I_n^{\mathrm{sat}}(B)\)
invariant as a learning-capacity bound. \PartVIII applies the structural
decomposition to the LPN charged-envelope criterion, while \PartIV gives the
learning-capacity reading of the same invariant.

\begin{example}[SAT final reading]\label{ex-final-sat-summary}

The discounted hitting functional for SAT depends on its presentation.
XOR-SAT with public parity rows has polynomial \(B^*\),
hence logarithmic \(d^*\) in log-cost notation. A random-target SAT
model has superpolynomial \(B^*\) with high probability. A
subset-sum-derived SAT instance with a withheld modulus is
budget-inadmissible. The invariant is attached to the presented problem
rather than to the name SAT alone.

\end{example}

\clearpage
\part{Learning and Computability Bounds}\label{part:learning}\setcounter{section}{-1}

\section{Introduction}\label{introduction-structure-exhausted-now-measure-what-can-be-learned}
\label{sec:learning-computability-introduction}

\paragraph{Part contract.}
\emph{Inputs:} the saturated cost and extraction profiles from
\cref{part:spectral-complexity}.
\emph{Output:} operator, kernel, neural-tangent-kernel, and certified
generalization bounds.
\emph{First pass:} the computational evaluation procedures may be skimmed.

\PartsItoIII define the task presentation, characterize
budget-admissible dynamics, and introduce the saturated cost invariant.
This part develops the corresponding statistical theory. Effective
dimension measures the number of statistically accessible modes,
alignment measures their orientation relative to the target, and the
source residual measures target mass outside the selected spectral
subspace.

\PartII shows that every budget-admissible oracle-free computation on
its finite clocked transcript carrier decomposes into the five exposed
structural channels

\[
\Geom,\qquad
\Caus,\qquad
\Abs,\qquad
\Lift,\qquad
\Bdry,
\]

and the residual term \(\Res\). For a task family whose
exposed structural components satisfy the stipulated cost conditions,
\PartIII shows that target arrival at the declared transfer-class margin
requires saturated extraction at the same margin,

\[
I_n^{\mathrm{sat}}(B)=m_n^{\mathrm{sat}}(B)+G_n^{\mathrm{sat}}(B)
\]

under the declared budget structure \(\mathcal B\), with ceiling
\(B=b(n)\). This accounting includes public basis changes, represented
quotients, runtime-generated observables, admissible finite-horizon
statistics, boundary readouts, and valid lifts. A learning algorithm is
itself a budget-admissible oracle-free computation on the presented
data; ordinary polynomial-time learning corresponds to
\(\mathcal B_{\mathrm{poly}}\). Consequently, generalization within this
formalism requires a fixed or learned observable algebra that captures
target mass at a statistical cost within the same saturated represented
closure. The choice \(\mathcal B=\mathcal B_{\mathrm{poly}}\) yields the
polynomial-time specialization.

We consider three settings.

\begin{itemize}
\tightlist
\item
  For fixed features, the observable algebra is an RKHS.
\item
  For learned features, a neural tangent kernel describes the evolving
  tangent-feature space.
\item
  After training, the spectrum of the learned kernel determines an
  empirical effective degree.
\end{itemize}

The analysis uses standard results on RKHSs and Mercer
eigendecompositions, the representer theorem, effective-dimension bounds
for kernel ridge regression, kernel--target alignment, neural tangent
kernels, and spectral-margin bounds for neural networks
\citep{Aronszajn1950,Mercer1909,KimeldorfWahba1971,CaponnettoDeVito2007,
CristianiniEtAl2002,JacotEtAl2018,BartlettFosterTelgarsky2017}.

\begin{definition}[Imported objects from \PartsItoIII]\label{def-paper4-dependencies}

\PartIV uses the following objects with the meanings fixed earlier.

\textbf{\PartI.} A task presentation has state space \(X_I\), observable
algebra \(\mathcal F_I\), positive terminal sector \(E_I^+\),
terminal-failure sector \(E_I^-\) when declared, static family-uniform
observable constructions, pointwise observables, represented-fiber
observables, and the exposed/presentation-inaccessible distinction. The
five static admissibility criteria are terminal alignment, extremum or
threshold consistency, cost within the declared budget structure,
family-uniformity, and monotone terminal filtration.

\textbf{\PartII.} A budget-admissible oracle-free computation becomes a
positive generator on each finite clocked transcript carrier.
Component exhaustiveness decomposes every budget-admissible generator
into \(\Geom\), \(\Caus\), \(\Abs\),
\(\Lift\), \(\Bdry\), plus \(\Res\) residual
accounting. The term \(\Res\) contains no additional exposed
structural source.

\textbf{\PartIII.} Computation is presentation-relative to a declared
budget structure \(\mathcal B\), inside the budgeted saturated
represented closure \(\mathcal R_n^{\mathcal B}\), equivalently the
derivability filtration \(\mathcal D_{n,\le b(n)}\). The extraction
functional is
\(I_n^{\mathrm{sat}}(B)=m_n^{\mathrm{sat}}(B)+G_n^{\mathrm{sat}}(B)\).
The term \(m_n^{\mathrm{sat}}(B)\) is the supremum over
represented-fiber quotient sources after admissible public basis
changes, valid lifted/generated coordinates, terminal compatibility, and
exact-lumpability conditions. The term \(G_n^{\mathrm{sat}}(B)\) is the
maximum of ambient Geom/Caus, Geom/Caus on an exact quotient, and
Geom/Caus on a compensated quasi-lumpable quotient. The saturated cost invariant
\(B^*\) is the first public budget at which \(I_n^{\mathrm{sat}}(B)\)
reaches the declared transfer-class margin; its log-cost coordinate is
\(d^*=\lceil\log_2 B^*\rceil\) when the resource is numeric. Raw Walsh
degree is one charged diagnostic coordinate; the budget criterion is
defined by the complete saturated cost.

\end{definition}

\begin{lemma}[Statistical interpretation of saturated extraction]\label{thm-paper4-dependency-principle}

Each learning object considered below consists of a \PartIII saturated
representative together with a statistical sampling procedure. A fixed
kernel, learned kernel, empirical NTK, feature quotient, margin bound, or
post-training spectral profile contributes to structural learning when
its instance-dependent construction belongs to
\(\mathcal D_{n,\le B}\) for the original presentation and its readout
satisfies the terminal conditions of \PartI and the channel
admissibility conditions of \PartII. The statistical quantities
\(\operatorname{Tr}(A^{-1})\), \(N_{\mathrm{eff}}(\lambda)\),
\(\log\mathcal N\), source residual, margin, and concentration errors
measure the sample cost of estimating or selecting these
representatives. They introduce no additional structural channel beyond
\(I_n^{\mathrm{sat}}(B)\).

\end{lemma}

\begin{proof}\phantomsection\label{proof-paper4-dependency-principle}

By the encoding-adequacy result of \PartIII, a budget-admissible learner
on a budget-complete encoding defines a presentation-relative generator
on its clocked training transcript carrier. The polynomial specialization recovers
uniform polynomial-time learners. \PartII decomposes the resulting
target-relevant dynamics into Geom, Caus, Abs, Lift, Bdry, and residual
accounting. If the learner constructs an operator, kernel, quotient,
spectral projector, feature map, or terminal readout within the declared
budget, the budgeted-generation lemma of \PartIII supplies a saturated
representative and assigns its derivation cost. Otherwise the object
requires oracle or advice data, an over-budget construction, or a
different primitive signature. Standard learning bounds then quantify
the estimation, capacity, and selection costs of the admissible
representative. Hence every included generalization term measures a
\PartIII source; all remaining terms are residual or inadmissible at the
declared budget.

\end{proof}

\paragraph{Source-resolved convention for Part~IV.}
Every hypothesis class, operator, kernel, prior, posterior, learned feature,
decoder, and deployed selector below is restricted by its exact source index
when a target-qualified structural use is made.  The complete quantitative
compilation is
\cref{def-complete-source-resolved-learning-ledger,%
thm-source-resolved-saturated-learning-compilation}.  Thus each statistical
bound retains \(B,m,\delta\), its source cell, and its complete
construction cost.  A noisy target-dependent reveal additionally retains
its source-assigned channel coefficient and noise parameter.  Generic
population--empirical deviations remain in risk units; they enter a
target-alignment, recovery, selector, or arrival bound only through a
represented same-law, same-unit converter.  This convention applies to
every theorem in this part, including fixed and learned kernels, NTK feature
learning, effective dimension, PAC--Bayes, dependent samples, channel
learning, active querying, and the complete learning decomposition.

\section{The Observable Algebra Is a Reproducing-Kernel Hilbert
Space}\label{the-observable-algebra-is-a-reproducing-kernel-hilbert-space}
\label{sec:observable-algebra-rkhs}

The observable algebra of \PartI acquires a statistical structure when
its functions are embedded in a Hilbert space. A feature map specifies
the observable representation, and its kernel records inner products
between represented states.

The feature observable embeds represented task functions in an RKHS and induces their kernel.
\begin{definition}[Feature observable and induced kernel]\label{def-feature-rkhs-observable}

Let \((X,\mathcal F,\mu)\) be a task presentation and let
\(\Phi:X\to\mathcal H\) be a measurable feature observable into a
Hilbert space. The induced kernel is

\[
k_\Phi(x,z)=\langle \Phi(x),\Phi(z)\rangle_{\mathcal H}.
\]

The RKHS \(\mathcal H_k\) is the closure of finite linear combinations
of kernel sections \(k_x(\cdot)=k(x,\cdot)\) with inner product
determined by \(\langle k_x,k_z\rangle_{\mathcal H_k}=k(x,z)\)
\citep{Aronszajn1950}.

\end{definition}

The Mercer decomposition expresses the represented kernel in its observable spectral coordinates.
\begin{definition}[Mercer eigendecomposition]\label{def-mercer-observable-spectrum}

Assume \(k\) is square-integrable and positive definite. Its integral
operator is

\[
(T_k f)(x)=\int_X k(x,z)f(z)\,d\mu(z).
\]

Let \((\eta_j,\varphi_j)\) denote its Mercer eigenpairs
\citep{Mercer1909},
\(\eta_1\ge\eta_2\ge\cdots\ge0\), a function
\(f=\sum_j \widehat f_j\varphi_j\) has RKHS norm

\[
\|f\|_{\mathcal H_k}^2
=
\sum_{j:\eta_j>0}\frac{\widehat f_j^2}{\eta_j}.
\]

\end{definition}

The eigenfunctions of \(T_k\) provide the observable basis for kernel
learning, in the same way that \PartIII uses Walsh characters on the
Boolean cube and Laplacian eigenfunctions on continuous spaces.

\begin{proposition}[RKHS norm as Dirichlet regularity]\label{prop-rkhs-dirichlet-energy}

Let \(L_W\ge0\) be the self-adjoint generator of an exposed geometric
Dirichlet form from \PartIII, with eigenpairs \((\lambda_j,\psi_j)\).
Define the resolvent kernel by spectral weights

\[
k_W(x,z)=\sum_j \frac{1}{1+\lambda_j}\psi_j(x)\psi_j(z).
\]

Then for \(f=\sum_j a_j\psi_j\),

\[
\|f\|_{\mathcal H_{k_W}}^2
=
\sum_j (1+\lambda_j)a_j^2
=
\|f\|_{L^2(\mu)}^2+\mathcal E_W(f,f).
\]

\end{proposition}

\begin{proof}\phantomsection\label{proof-rkhs-dirichlet-energy}

The RKHS norm divides each squared coefficient by the Mercer eigenvalue.
Here the eigenvalue is \(\eta_j=(1+\lambda_j)^{-1}\). Hence
\(\|f\|_{\mathcal H_{k_W}}^2=\sum_j a_j^2/\eta_j=\sum_j(1+\lambda_j)a_j^2\).
The spectral theorem for the Dirichlet form gives
\(\mathcal E_W(f,f)=\sum_j\lambda_j a_j^2\).

\end{proof}

\begin{corollary}[Kernel formulation of the observable criteria]\label{cor-choosing-kernel-choosing-observable}

The static criteria of \PartI have the following kernel formulation.
Terminal alignment is the label mass carried by Mercer modes. Extremum
or threshold consistency requires kernel scores to agree with the
terminal interface. Effective dimension supplies the statistical
complexity measure. Family uniformity requires a public kernel rule for
the entire family, and monotone terminal filtration requires stability
of the induced score bands under terminal readout. RKHS/Dirichlet
regularity is the \PartIII geometric regularity coordinate expressed in
the Mercer basis, rather than an additional static criterion.

\end{corollary}

\begin{proof}\phantomsection\label{proof-choosing-kernel-choosing-observable}

\PartI evaluates an observable through target projection, extremum or
threshold consistency, affine cost, family-uniform construction, and
terminal filtration. A kernel selects a feature map and therefore a
closed observable span. The Mercer eigendecomposition identifies target
projection with label mass in the kernel eigenmodes. Threshold
consistency requires the score formed from these modes to induce the
same terminal readout as the target interface. Effective dimension
counts the modes available within the statistical and computational
budget. Family uniformity requires the kernel rule to be generated from
the presentation uniformly over the family, while monotone terminal
filtration corresponds to the nested spectral or score bands induced by
the kernel readout. Finally, \cref{prop-rkhs-dirichlet-energy} identifies the
RKHS norm with the Dirichlet energy associated with \(I+L_W\), namely
the Geom regularity coordinate of \PartIII.

\end{proof}

\Cref{fig-observable-rkhs-bridge} summarizes the Part I--Part IV bridge in
\cref{prop-rkhs-dirichlet-energy,cor-choosing-kernel-choosing-observable}.

\begin{figure}[tbp]
\centering
\begin{tikzpicture}[fw diagram]
  \node[fw source,minimum width=2.7cm] (alg) at (-4.8,1.2)
    {observable algebra\\\(\mathcal F_I\)};
  \node[fw use,minimum width=2.5cm] (feat) at (-1.7,1.2)
    {feature map\\\(\Phi\)};
  \node[fw geom,minimum width=2.6cm] (kern) at (1.4,1.2)
    {kernel\\\(k_\Phi\)};
  \node[fw provenance,minimum width=2.8cm] (mercer) at (4.6,1.2)
    {Mercer modes\\\((\eta_j,\varphi_j)\)};
  \node[fw terminal,minimum width=2.8cm] (mass) at (1.8,-0.8)
    {target spectral mass};
  \node[fw geom,minimum width=3.2cm] (energy) at (5.8,-0.8)
    {\(\|f\|_{\mathcal H_k}^2=\|f\|_2^2+\mathcal E_W(f,f)\)};
  \draw[fw arrow] (alg) -- (feat);
  \draw[fw arrow] (feat) -- (kern);
  \draw[fw arrow] (kern) -- (mercer);
  \draw[fw arrow] (mercer) -- (mass);
  \draw[fw arrow] (mercer) -- (energy);
\end{tikzpicture}
\caption{A represented feature observable selects a kernel and hence a
Mercer observable basis.  Target alignment is spectral mass in that basis,
while the resolvent-kernel RKHS norm is the represented Dirichlet regularity
of \cref{prop-rkhs-dirichlet-energy}.}
\label{fig-observable-rkhs-bridge}
\end{figure}

Admissibility requires the learning kernel and its construction to remain inside the charged presentation.
\begin{definition}[Admissible learning kernel]\label{def-admissible-learning-kernel}

A fixed kernel \(k_I\) is admissible for an instance \(I\) when its
feature rule, kernel evaluations, and required Gram or integral-operator
approximations have derivations in \(\mathcal D_{n,\le B}\) from
\((X_I,\mathcal F_I,E_I^+,E_I^-)\) by a family-uniform oracle-free
rule. A learned kernel is admissible when a budget-admissible
oracle-free training computation produces it from the observed sample
and labels, and the terminal kernel is a represented terminal output on
the clocked training transcript carrier, with its finite
outcome-generation interface, precision, and postprocessing charged in the same saturated
budget accounting. The polynomial case is the instance
\(\mathcal B_{\mathrm{poly}}\).

Observed labels are part of the learning sample. Unobserved labels,
target membership outside the sample, planted certificates, and
relations not generated from the budgeted saturated observable closure
are not available to the learner. If a spectral projector or feature
quotient produced by a budget-admissible kernel has
total cost within \(\mathcal B\) and satisfies the represented-fiber,
terminal-compatibility, and lumpability conditions of
\cref{def-paper1-abs-usefulness-gate,def-abs-extraction}, its captured mass is a candidate contribution to
\(m_n^{\mathrm{sat}}(B)\) at the charged cost. If it is pointwise
Geom/Caus structure, it contributes through \(G_n^{\mathrm{sat}}(B)\).
Otherwise the same empirical correlation is presentation-inaccessible,
requires an oracle, exceeds the discovery or application budget,
changes the primitive signature, or remains residual accounting.

\end{definition}

\section{Saturated Cost and Effective
Dimension}\label{saturated-cost-and-effective-dimension}
\label{sec:saturated-cost-effective-dimension}

\PartIII accounts for public basis changes, represented quotients,
generated coordinates, admissible finite-horizon statistics, boundary
readouts, and valid lifts in the saturated public budget. In the RKHS
setting, ridge regularization assigns each Mercer mode a continuous
statistical weight.

Effective dimension measures the statistically active Mercer modes at the chosen regularization scale.
\begin{definition}[Effective dimension]\label{def-effective-dimension}

For a kernel integral operator \(T_k\) with eigenvalues \(\eta_j\ge0\)
and ridge scale \(\lambda>0\), define

\[
N_{\mathrm{eff}}(\lambda)
=
\operatorname{Tr}\bigl(T_k(T_k+\lambda I)^{-1}\bigr)
=
\sum_j\frac{\eta_j}{\eta_j+\lambda}.
\]

For an empirical Gram matrix \(K\in\mathbb R^{m\times m}\), the
corresponding quantity is

\[
\widehat N_{\mathrm{eff}}(\lambda)
=
\operatorname{Tr}\bigl(K(K+\lambda m I)^{-1}\bigr),
\]

equivalently \(\sum_j \widehat\eta_j/(\widehat\eta_j+\lambda)\) for the
eigenvalues of \(K/m\).

\end{definition}

\begin{proposition}[Effective dimension as statistical mode count]\label{prop-effective-dimension-cost}

The function \(N_{\mathrm{eff}}(\lambda)\) interpolates between the rank
of the kernel and zero. If
\(\eta_j\asymp j^{-2\beta}\) with \(\beta>1/2\), then
\(N_{\mathrm{eff}}(\lambda)\asymp\lambda^{-1/(2\beta)}\). On a compact
\(q\)-dimensional geometric domain, Weyl's law for a Laplacian kernel
gives \(N_{\mathrm{eff}}(\lambda)\) of the same order as a spectral
cutoff count \(C_q\Lambda^{q/2}\) under the resolvent relation
\(\eta=(1+\Lambda)^{-1}\) \citep{Weyl1912}. Thus
\(N_{\mathrm{eff}}\) is the RKHS realization of the statistical mode
count in the saturated budget. The provenance and representation
conditions of \PartIII remain separate requirements.

\end{proposition}

\begin{proof}\phantomsection\label{proof-effective-dimension-cost}

The summand \(\eta_j/(\eta_j+\lambda)\) is close to one for modes with
\(\eta_j\gg\lambda\) and close to zero for modes with
\(\eta_j\ll\lambda\). Thus the trace is a smoothed count of modes above
the learnability threshold. Polynomial eigenvalue decay gives the
displayed inverse relation by solving \(j^{-2\beta}\approx\lambda\).
Weyl's law gives the continuous-domain count for Laplacian eigenvalues
below \(\Lambda\). The definition of a budget-admissible kernel supplies
the independent provenance condition required for interpreting this
count structurally.

\end{proof}

Budget-feasible bands identify which spectral target mass is available to the represented learner.
\begin{definition}[Budget-feasible spectral bands and candidate capture]\label{def-learning-affordability}

A budget-admissible kernel band is budget-feasible at ridge scale
\(\lambda\) when its total saturated learning charge does not exceed
\(b(n)\). For a fixed kernel this budget contains the derivation and
representation cost of \(k\), the Gram or operator approximation cost,
the numerical precision, and \(N_{\mathrm{eff}}(\lambda)\). For a
learned kernel it also contains the training transcript and the
operator-selection cost. For a centered target field
\(y=\sum_j b_j\varphi_j\), its kernel spectral capture at a hard cutoff
is

\[
m_k(\tau)=\frac{\sum_{\eta_j\ge\tau}b_j^2}{\sum_jb_j^2},
\]

and its ridge-smoothed capture is

\[
m_k(\lambda)=
\frac{\sum_j\left(\frac{\eta_j}{\eta_j+\lambda}\right)^2b_j^2}{\sum_jb_j^2}.
\]

The quantity \(m_k\) is a kernel-specific candidate capture curve. It
contributes to \(m_n^{\mathrm{sat}}(B)\) only when the kernel, spectral
band, and induced quotient or lift satisfy the admissibility,
represented-fiber, terminal-compatibility, and lumpability conditions within
total budget \(B\). It contributes to \(G_n^{\mathrm{sat}}(B)\) only
when the same data define an admissible pointwise Geom/Caus source.
Otherwise \(m_k\) remains an empirical correlation statistic rather
than counted structural extraction.

\end{definition}

\Cref{fig-effective-dimension-capture} separates statistical mode count from
target-aligned spectral capture in
\cref{def-effective-dimension,def-learning-affordability}.

\begin{figure}[tbp]
\centering
\begin{tikzpicture}[fw diagram]
  \node[fw source,minimum width=3.0cm] (modes) at (-5.0,0)
    {Mercer modes\\\((\eta_j,\varphi_j)\)};
  \node[fw geom,minimum width=3.2cm] (ridge) at (-1.6,1.1)
    {ridge weights\\\(\eta_j/(\eta_j+\lambda)\)};
  \node[fw provenance,minimum width=3.2cm] (labels) at (-1.6,-1.1)
    {target coefficients\\\(b_j^2/\|y\|^2\)};
  \node[fw use,minimum width=3.2cm] (neff) at (2.0,1.1)
    {effective dimension\\\(N_{\mathrm{eff}}(\lambda)\)};
  \node[fw terminal,minimum width=3.2cm] (capture) at (2.0,-1.1)
    {captured target mass\\\(m_k(\lambda)\)};
  \node[fw box,minimum width=3.4cm] (candidate) at (5.8,0)
    {budget-feasible\\structural candidate};
  \draw[fw arrow] (modes) -- (ridge);
  \draw[fw arrow] (modes) -- (labels);
  \draw[fw arrow] (ridge) -- (neff);
  \draw[fw arrow] (labels) -- (capture);
  \draw[fw arrow] (neff) -- (candidate);
  \draw[fw arrow] (capture) -- (candidate);
\end{tikzpicture}
\caption{Effective dimension counts ridge-active modes, whereas captured
mass measures their alignment with the target.  A kernel band contributes to
the saturated profile only through the budget and admissibility conditions of
\cref{def-learning-affordability}.}
\label{fig-effective-dimension-capture}
\end{figure}

This construction is the continuous RKHS analogue of the
Walsh-to-Laplacian passage in \PartIII. A spectral cutoff or ridge scale
is a coordinate of the saturated public budget, and
\(N_{\mathrm{eff}}\) gives the corresponding smoothed mode count.
Generator--observable alignment is an additional requirement: the
dominant kernel modes must carry target mass. A low-dimensional kernel
whose eigenspaces are misaligned with the target remains statistically
uninformative even when it lies within budget.

\section{Two-Sided Learning Bounds}\label{the-tightened-two-sided-learning-bound}
\label{sec:two-sided-learning-bounds}

\subsection{Structural Operator Bounds Before the Kernel
Refinement}\label{structural-operator-bounds-before-the-kernel-refinement}
\label{subsec:structural-operator-bounds}

The RKHS effective-dimension bound refines the operator-level
statistical bound associated with the saturated budget of \PartIII. The
operator formulation separates three contributions that persist in the
kernel and neural-network results: complexity relative to a fixed
budget-admissible operator, target--operator alignment, and the cost of
operator selection.

The structural energy ball fixes the represented function class used by the operator-level generalization bound.
\begin{definition}[Structural energy ball]\label{def-structural-energy-ball}

On an empirical task space \(S=\{x_1,\ldots,x_m\}\), let
\(A\succeq \eta I\) be a budget-admissible symmetric structural
operator. For \(f\in\mathbb R^m\), set \(\mathcal E_A(f)=f^TAf\) and

\[
\mathcal F_A(R)=\{f\in\mathbb R^m:f^TAf\le R^2\}.
\]

\end{definition}

\begin{lemma}[Exact Rademacher complexity of an energy ball]\label{lem-ellipsoid-rademacher}

For any sign vector \(\sigma\in\{\pm1\}^m\),

\[
\sup_{f\in\mathcal F_A(R)}\sigma^Tf
=
R\sqrt{\sigma^TA^{-1}\sigma}.
\]

Consequently the empirical Rademacher complexity satisfies

\[
\operatorname{Rad}_S^A(R)
=
\frac{R}{m}\mathbb E_\sigma\sqrt{\sigma^TA^{-1}\sigma}
\le
\frac{R}{m}\sqrt{\operatorname{Tr}(A^{-1})}.
\]

\end{lemma}

\begin{proof}\phantomsection\label{proof-ellipsoid-rademacher}

Lagrange multipliers give the optimizer
\(f=RA^{-1}\sigma/\sqrt{\sigma^TA^{-1}\sigma}\). Averaging the support
function over signs gives the equality. Jensen's inequality and
\(\mathbb E_\sigma\sigma\sigma^T=I\) give the trace bound.

\end{proof}

The fixed-operator bound controls generalization after the represented kernel has been selected.
\begin{theorem}[Fixed-operator generalization bound]\label{thm-fixed-operator-rademacher-bound}

For a loss taking values in \([0,1]\) that is
\(L_\ell\)-Lipschitz in its prediction argument, and for a fixed
budget-admissible operator \(A\), with probability at least
\(1-\delta\), uniformly over \(h\in\mathcal F_A(R)\),

\[
L(h)
\le
L_S(h)
+\frac{2L_\ell R}{m}\sqrt{\operatorname{Tr}(A^{-1})}
+3\sqrt{\frac{\log(2/\delta)}{2m}}
+\varepsilon_{\mathrm{task}}.
\]

The bound is specific to the declared operator; optimization over the
budget-admissible operator class is a separate model-selection problem.
Here \(\varepsilon_{\mathrm{task}}\ge0\) is the fixed approximation or task
residual declared by the benchmark; it is zero when the population and
empirical risks are compared without such an additional residual.
For a target-qualified deployment, restrict \(A\) and its output class to
their exact source cell \(\alpha\).  The displayed radius is then
\(\Gamma_{R,\alpha}^{\rm learn}(B,m,\boldsymbol\eta,\delta)\) in
\cref{def-complete-source-resolved-learning-ledger}; conversion to target
gain is governed by
\cref{eq-source-resolved-learning-target-gain-minimum}.

\end{theorem}

\begin{proof}\phantomsection\label{proof-fixed-operator-rademacher-bound}

Apply the standard uniform generalization inequality for an
\(L_\ell\)-Lipschitz loss
\citep{BartlettMendelson2002,LedouxTalagrand1991}: with probability at
least \(1-\delta\),
uniformly over a class \(\mathcal F\),
\(L(h)\le L_S(h)+2L_\ell\widehat{\mathfrak R}_S(\mathcal F)+3\sqrt{\log(2/\delta)/(2m)}\).
For \(\mathcal F_A(R)=\{h:h^TAh\le R^2\}\), the ellipsoid Rademacher
lemma gives
\(\widehat{\mathfrak R}_S(\mathcal F_A(R))\le Rm^{-1}\sqrt{\operatorname{Tr}(A^{-1})}\).
Substitution gives the displayed bound. The term
\(\varepsilon_{\mathrm{task}}\) is the fixed approximation or task
residual already present in the declared benchmark, so it adds linearly.

\end{proof}

These capacities compare total task complexity with the target-aligned represented component.
\begin{definition}[Task capacity and structural capacity]\label{def-taskcap-structcap}

For a target field \(y\), define the fixed-operator capacity score

\[
\operatorname{TaskCap}_A(y)
=
\frac1m\sqrt{(y^TAy)\operatorname{Tr}(A^{-1})}.
\]

The identity fallback and modal gap score are

\[
\operatorname{GapCap}_A(h)
=
\frac1m\sqrt{(h^TAh)\operatorname{Tr}(A^{-1})},
\qquad
\operatorname{GapCap}_I(h)=\frac{\|h\|_2}{\sqrt m}.
\]

The resulting bound is the minimum of the modal and identity bounds. The
identity operator provides a reference value for detecting operator
misalignment.

For a parameterized budget-admissible class
\(\Theta\ni\theta\mapsto A_\theta\), define

\[
\operatorname{StructCap}(y;\Theta)
=
\min_{\theta\in\Theta}
\left[
\operatorname{TaskCap}_{A_\theta}(y)
+
c\sqrt{\frac{\log\mathcal N(\Theta,\rho)}{m}}
\right].
\]

The first term measures target complexity after fixing the operator,
whereas the second accounts for data-dependent operator selection.

\end{definition}

The random-label score supplies the null calibration for empirical target alignment.
\begin{definition}[Random-label null score]\label{def-random-label-null-score}

For a structural operator \(A\) and target field \(y\), define

\[
Z_A(y)
=
\frac{\operatorname{Tr}(A)-y^TAy}{2\sqrt{\sum_{i<j}A_{ij}^2}+\eta}.
\]

This statistic is the standardized reduction of \(y^TAy\) relative to
the random-sign null mean. A large positive value indicates
operator--target alignment, whereas a value near zero is consistent with
the null scale.

\end{definition}

\begin{lemma}[Learned-operator generalization bound]\label{thm-learned-operator-union-bound}

Let \(\mathcal H(R)=\bigcup_{\theta\in\Theta}\mathcal F_{A_\theta}(R)\).
Suppose \(A_\theta\succeq\eta I\) uniformly,
\(\theta\mapsto A_\theta^{-1}\) is
\(\mathrm{Lip}\)-Lipschitz in operator norm, and
\(\mathcal N(\Theta,\rho)\) is a finite \(\rho\)-covering number. Then

\[
\operatorname{Rad}_S^{\mathcal H}(R)
\le
\inf_{\rho>0}
\left[
\frac{R}{m}\sqrt{\max_\theta\operatorname{Tr}(A_\theta^{-1})}
+
R\sqrt{\frac{\mathrm{Lip}\,\rho}{m}}
+
\frac{R}{m\sqrt\eta}
\sqrt{2\log\mathcal N(\Theta,\rho)}
\right].
\]

Consequently, for a \([0,1]\)-valued
\(L_\ell\)-Lipschitz loss, the generalization gap over
\(\mathcal H(R)\) has the same confidence form as the fixed-operator
theorem, with twice \(L_\ell\) times the displayed quantity plus the
sampling tail and declared task residual.

\end{lemma}

\begin{proof}\phantomsection\label{proof-learned-operator-union-bound}

Fix \(\rho>0\) and choose a \(\rho\)-net \(\Theta_\rho\). For each net
point, the ellipsoid lemma bounds the empirical Rademacher support
function by \(Rm^{-1}\sqrt{\operatorname{Tr}(A_\theta^{-1})}\). The
Lipschitz condition and
\(|\sqrt a-\sqrt b|\le\sqrt{|a-b|}\) give, for every sign vector
\(\sigma\) and every \(\theta\) within distance \(\rho\) of a net
point \(\theta_0\),

\[
\frac{R}{m}
\left|
\sqrt{\sigma^\mathsf TA_\theta^{-1}\sigma}
-\sqrt{\sigma^\mathsf TA_{\theta_0}^{-1}\sigma}
\right|
\le
R\sqrt{\frac{\mathrm{Lip}\rho}{m}},
\]

because \(\|\sigma\|_2^2=m\). Moreover
\(A_\theta\succeq\eta I\) implies
\(\sup_{f\in\mathcal F_{A_\theta}(R)}\|f\|_2\le R/\sqrt\eta\).
The finite-union Rademacher lemma therefore contributes
\(R\sqrt{2\log|\Theta_\rho|}/(m\sqrt\eta)\)
\citep{Massart2000}. Taking the infimum over
\(\rho\) gives the displayed Rademacher bound, and the same
symmetrization/contraction inequality used in the fixed-operator theorem
gives the stated generalization form.

\end{proof}

\Cref{fig-learned-operator-selection-cost} separates the fixed-operator
capacity term from the cost of selecting an operator from data.

\begin{figure}[tbp]
\centering
\begin{tikzpicture}[fw diagram]
  \node[fw source,minimum width=3.0cm] (family) at (-4.6,0)
    {represented family\\\(\{A_\theta:\theta\in\Theta\}\)};
  \node[fw provenance,minimum width=2.9cm] (select) at (-0.8,1.1)
    {selection radius\\\(\log\mathcal N(\Theta,\rho)\)};
  \node[fw geom,minimum width=2.9cm] (fixed) at (-0.8,-1.1)
    {fixed-operator size\\\(\operatorname{Tr}(A_\theta^{-1})\)};
  \node[fw use,minimum width=3.4cm] (bound) at (3.6,0)
    {learned-operator\\Rademacher bound};
  \draw[fw arrow] (family) -- (select);
  \draw[fw arrow] (family) -- (fixed);
  \draw[fw arrow] (select) -- (bound);
  \draw[fw arrow] (fixed) -- (bound);
\end{tikzpicture}
\caption{The bound in \cref{thm-learned-operator-union-bound} charges two
distinct quantities: capacity after fixing \(A_\theta\), and the represented
choice among operators. The covering term prevents data-dependent geometry
selection from being treated as free.}
\label{fig-learned-operator-selection-cost}
\end{figure}

\begin{corollary}[Rademacher bound for the \PartI empirical audit]\label{cor-rademacher-audit-ceiling}

Let \(\mathcal H_B\) be a budget-admissible readout class at budget
\(B\) and let \(a_h(Z)\in[0,1]\) be a bounded audit statistic such that
\(\operatorname{Ext}_I(h;T_I)\le\mathbb E a_h(Z)\) for every
\(h\in\mathcal H_B\). For an audit sample \(Z_1,\ldots,Z_M\), write
\(\widehat a_M(h)=M^{-1}\sum_i a_h(Z_i)\). Then with probability at
least \(1-\delta\),

\[
\sup_{h\in\mathcal H_B}\operatorname{Ext}_I(h;T_I)
\le
\sup_{h\in\mathcal H_B}\widehat a_M(h)
+
2\operatorname{Rad}_M(\{a_h:h\in\mathcal H_B\})
+
3\sqrt{\frac{\log(2/\delta)}{2M}}.
\]

Thus the \PartI empirical audit penalty may be taken as

\[
\mathfrak P_{\mathcal H_B}(M,\delta)
=
2\operatorname{Rad}_M(\{a_h:h\in\mathcal H_B\})
+
3\sqrt{\frac{\log(2/\delta)}{2M}},
\]

or by the learned-operator cover bound above when
\(\mathcal H_B=\bigcup_{\theta\in\Theta}\mathcal F_{A_\theta}(R)\). This
inequality converts empirical optimization over observable combinations
into a uniform upper bound on actionable extraction.

\end{corollary}

\begin{proof}\phantomsection\label{proof-rademacher-audit-ceiling}

Apply the standard uniform Rademacher generalization inequality to the
bounded audit class \(\{a_h:h\in\mathcal H_B\}\). Since
\(\operatorname{Ext}_I(h;T_I)\le\mathbb E a_h(Z)\) by hypothesis for
each \(h\), the same upper bound holds for the extraction supremum.
Substituting the fixed-operator or learned-operator Rademacher estimates
gives the corresponding budgeted audit penalty.

\end{proof}

\begin{corollary}[Operator-selection conservation]\label{cor-operator-selection-conservation}

Enlarging the operator class may reduce
\(\operatorname{TaskCap}_{A_\theta}(y)\), but it also increases the
selection term \(\sqrt{\log\mathcal N(\Theta,\rho)/m}\) unless the
instance determines an aligned subfamily of small covering number.
Universal approximation establishes representability of an aligned
operator but does not control the cost of selecting it.

\end{corollary}

\begin{proof}\phantomsection\label{proof-operator-selection-conservation}

The learned-operator theorem decomposes the cost of selecting an
operator into a fixed-operator capacity term, a Lipschitz approximation
term, and the covering term
\[
\frac{R}{m\sqrt\eta}
\sqrt{2\log\mathcal N(\Theta,\rho)}.
\]
Enlarging the operator family can reduce
\(\operatorname{TaskCap}_{A_\theta}(y)\) through minimization over
\(\theta\). The same enlargement cannot decrease the covering number
at a fixed resolution. It yields a net certificate improvement only
when the task-capacity reduction exceeds the additional approximation
and covering terms. Universal approximation controls approximation
after a suitable \(\theta\) is known; the theorem also charges the cost
of selecting that \(\theta\).

\end{proof}

\begin{lemma}[Condition-number lower bound for operator capacity]\label{thm-random-label-operator-hardness}

Let \(A_\theta\succ0\) for every \(\theta\in\Theta\), and set
\(\kappa=\sup_\theta\operatorname{cond}(A_\theta)<\infty\). Then,
deterministically for every target vector \(y\),

\[
\min_{\theta\in\Theta}\operatorname{TaskCap}_{A_\theta}(y)
\ge
\frac{\|y\|_2}{\sqrt{m\kappa}}.
\]

For a sign field the right-hand side is \(1/\sqrt\kappa\). For the
centered indicator of a fixed-density target sector it is
\(\sqrt{\nu(T)(1-\nu(T))/\kappa}\). This estimate concerns the capacity
functional only; it does not by itself give a statistical minimax lower
bound or require a covering argument.

\end{lemma}

\begin{proof}\phantomsection\label{proof-random-label-operator-hardness}

Let \(\lambda_{\min}\) and \(\lambda_{\max}\) be the extremal
eigenvalues of \(A_\theta\). Then

\[
y^\mathsf TA_\theta y\ge\lambda_{\min}\|y\|_2^2,
\qquad
\operatorname{tr}(A_\theta^{-1})\ge\frac{m}{\lambda_{\max}}.
\]

Substitution into the definition of \(\operatorname{TaskCap}\) gives
\[
\operatorname{TaskCap}_{A_\theta}(y)
\ge
\frac{\|y\|_2}{\sqrt m}
\sqrt{\frac{\lambda_{\min}}{\lambda_{\max}}}
\ge
\frac{\|y\|_2}{\sqrt{m\kappa}}.
\]
Taking the infimum over \(\theta\) proves the first claim. The two
specializations follow from
\(\|y\|_2^2=m\) for signs and
\(\|y\|_2^2=m\nu(T)(1-\nu(T))\) for a centered fixed-density indicator.

\end{proof}

\subsection{Kernel Ridge Bounds at a Saturated
Representative}\label{kernel-ridge-bounds-at-a-saturated-representative}
\label{subsec:kernel-ridge-saturated-representative}

For a fixed budget-admissible kernel, the learning bound is the
statistical form of the saturated extraction bound in \PartIII. The
effective dimension refines the trace/Rademacher estimate and attains
the corresponding minimax rate. Derivability, admissibility, and
alignment determine whether the kernel's spectral mass contributes to
\(I_n^{\mathrm{sat}}(B)\).

The kernel ridge estimator is the regularized learner evaluated by the effective-dimension bound.
\begin{definition}[Kernel ridge estimator]\label{def-krr-estimator}

Given samples \((x_i,y_i)_{i=1}^m\) from \(Y=f_*(X)+\xi\), with
\(\mathbb E[\xi\mid X]=0\) and \(\mathbb E[\xi^2\mid X]\le\sigma^2\),
kernel ridge regression chooses

\[
\widehat f_\lambda
=
\arg\min_{f\in\mathcal H_k}
\frac1m\sum_{i=1}^m(f(x_i)-y_i)^2+\lambda\|f\|_{\mathcal H_k}^2.
\]

\end{definition}

Empirical alignment measures the observed target mass carried by the learned kernel.
\begin{definition}[Empirical kernel-target alignment]\label{def-kernel-target-alignment}

Let \(H=I-m^{-1}\mathbf 1\mathbf 1^T\), \(K_c=HKH\), and \(y_c=Hy\). The
centered kernel-target alignment is

\[
A_m(K,y)
=
\frac{\langle K_c,y_cy_c^T\rangle_F}{\|K_c\|_F\,\|y_cy_c^T\|_F}.
\]

This is the centered kernel--target alignment of
\citet{CristianiniEtAl2002}. It measures the empirical association
between the pairwise kernel geometry and label agreement. The statistic
contains no information about unobserved labels and is used only to
select among kernels satisfying the preceding provenance conditions.

\end{definition}

The kernel-ridge theorem transfers a fixed-kernel oracle certificate into
the structural ledger without changing its probability event or units.
\begin{theorem}[Transfer of a certified effective-dimension bound]
\label{thm-krr-effective-dimension-bound}

Fix a kernel \(k\), ridge scale \(\lambda\), regression function
\(f_*\in L^2(\mu)\), and confidence event \(\mathcal E_{k,\lambda}\) with
\(\Pr(\mathcal E_{k,\lambda})\ge1-\delta\).  A \emph{fixed-kernel oracle
certificate} consists of represented constants
\(C_{\rm var},C_{\rm bias}<\infty\) and a proof, on that same event, of
\begin{equation}
\label{eq-fixed-kernel-oracle-certificate}
\|\widehat f_\lambda-f_*\|_{L^2(\mu)}^2
\le
C_{\rm var}\frac{\sigma^2N_{\mathrm{eff}}(\lambda)}{m}
+C_{\rm bias}B_k(\lambda;f_*).
\end{equation}
The event may be supplied by any declared operator-concentration theorem,
but its hypotheses, constants, confidence allocation, and covariance/noise
normalization belong to the complete record.  Then the saturated
fixed-kernel risk coordinate obeys, on \(\mathcal E_{k,\lambda}\),

\[
\|\widehat f_\lambda-f_*\|_{L^2(\mu)}^2
\le
C_{\rm var}\frac{\sigma^2N_{\mathrm{eff}}(\lambda)}{m}
+C_{\rm bias}B_k(\lambda;f_*),
\]

where
\(B_k(\lambda;f_*)=\|(I-T_k(T_k+\lambda I)^{-1})f_*\|_{L^2(\mu)}^2\)
is the source residual.  If a named source/eigenvalue class also supplies a
matching minimax lower certificate with constants
\(c_{\rm low},C_{\rm high}>0\), optimizing the displayed upper certificate
is rate-matched only for that declared class and those constants.

\end{theorem}

\begin{proof}\phantomsection\label{proof-krr-effective-dimension-bound}

Equation \eqref{eq-fixed-kernel-oracle-certificate} is already a population
bound on the named event and in the named norm.  Entering the represented
kernel, event, and constants into the saturated ledger changes neither side,
so the displayed structural coordinate follows verbatim.  Optimization over
\(\lambda\) is permitted only after the record charges that selection or the
certificate is simultaneous in \(\lambda\).  A matching lower certificate,
when separately supplied for the same source/eigenvalue class, gives the
stated rate comparison.  No unnamed concentration or minimax theorem is used
by this transfer step.

\end{proof}

\begin{corollary}[Source-condition learning rate]\label{cor-source-condition-rate}

Assume \(f_*=T_k^r g\) with \(0<r\le1\), source radius
\(\|g\|_{L^2}\le R_s\), and Mercer eigenvalues
\(\eta_j\asymp j^{-2\beta}\) with \(\beta>1/2\). Assume also that
\eqref{eq-fixed-kernel-oracle-certificate} holds at the selected
\(\lambda_m\), or simultaneously over the represented candidate grid, with
\(C_{\rm var}\) and \(C_{\rm bias}\) bounded independently of \(m\).  The
selection of \(\lambda_m\) must be predeclared, independently certified, or
charged in the complete record. Then
the source residual satisfies
\(B_k(\lambda;f_*)\le C_rR_s^2\lambda^{2r}\), while
\(N_{\mathrm{eff}}(\lambda)\asymp\lambda^{-1/(2\beta)}\). On the
oracle-certificate event, optimizing the fixed-kernel risk gives

\[
\lambda_m
\asymp
m^{-1/(2r+1/(2\beta))},
\qquad
\|\widehat f_{\lambda_m}-f_*\|_{L^2}^2
\lesssim
m^{-2r/(2r+1/(2\beta))}
=
m^{-4\beta r/(4\beta r+1)},
\]

with an implicit constant determined by
the declared spectral, source, \(C_{\rm var}\), and \(C_{\rm bias}\)
constants. This establishes the quantitative correspondence between extraction
degree and the source condition: greater smoothness relative to a
budget-admissible kernel yields a smaller effective degree and a faster
rate.

\end{corollary}

\begin{proof}\phantomsection\label{proof-source-condition-rate}

Write the target in the Mercer basis. The source condition
\(f_*=T_k^rg\) gives coefficients \(a_j=\eta_j^r g_j\), so the ridge
bias is bounded by \(C_rR_s^2\lambda^{2r}\). The eigenvalue law
\(\eta_j\asymp j^{-2\beta}\) gives
\(N_{\mathrm{eff}}(\lambda)=\sum_j\eta_j/(\eta_j+\lambda)\asymp\lambda^{-1/(2\beta)}\).
The certified kernel-ridge inequality therefore gives, on its named event,
risk of order
\(\lambda^{2r}+m^{-1}\lambda^{-1/(2\beta)}\). Equating the two terms
yields \(\lambda_m\asymp m^{-1/(2r+1/(2\beta))}\), and substituting this
value gives \(m^{-2r/(2r+1/(2\beta))}=m^{-4\beta r/(4\beta r+1)}\).

\end{proof}

\begin{corollary}[Alignment-calibrated kernel-selection diagnostic]\label{cor-structural-learning-proxy}

For kernel selection, define the empirical diagnostic

\[
\mathcal B_\lambda(K,y)
=
\frac{\sigma^2N_{\mathrm{eff}}(\lambda)}{m\,\max\{A_m(K,y),a_0\}}
+
\bigl(1-m_k(\lambda)\bigr)\|y_c\|_m^2
+
\sigma^2,
\]

where \(a_0>0\) is a declared minimum alignment floor. When \(K\) is
derivable under the declared budget structure, \(A_m(K,y)\) is bounded
below, and \(m_k(\lambda)\) is close to one at the declared total
learning budget, the score records a within-budget candidate for
captured structure. It is not an excess-risk upper bound unless a
separate theorem derives the alignment denominator and source term in
this form. Small \(A_m\) does not certify \PartIII extraction,
and alignment obtained only above the declared total budget is
inadmissible for the original presentation.

\end{corollary}

\begin{proof}\phantomsection\label{proof-structural-learning-proxy}

The effective-dimension, captured-mass, and noise entries are measurable
diagnostics motivated by the kernel-ridge decomposition. The alignment
denominator is not present in the kernel-ridge oracle inequality; it
records the separate \PartIII requirement that captured spectral mass
be oriented toward the target rather than concentrated in irrelevant
low modes. The floor
\(a_0\) keeps the certificate finite when the empirical alignment is
null. If the kernel is derivable under the declared budget structure and
the alignment and captured mass conditions hold, each term in
\(\mathcal B_\lambda\) is a charged statistical or structural
diagnostic. A risk or minimax conclusion requires the explicit
hypotheses of \cref{thm-kernel-minimax-hard-side}.

\end{proof}

\begin{lemma}[Kernel minimax lower bound]\label{thm-kernel-minimax-hard-side}

Fix a budget-reachable kernel class \(\mathcal K\) whose data-dependent
selection cost lies in the declared transfer class. Assume a uniform
minimax lower-bound theorem for the declared regression class: for each
fixed \(k\in\mathcal K\), every admissible estimator has excess risk at
least
\[
c_0\min\left\{
\frac{\sigma^2N_{\mathrm{eff},k}(\lambda)}m
+B_k(\lambda;f_*),\,1
\right\}
\]
for the relevant choice of \(\lambda\), with the same constant \(c_0>0\).
Suppose every reachable kernel either has
\(N_{\mathrm{eff},k}(\lambda)/m\ge c_1>0\), or has
\(B_k(\lambda;f_*)\ge c_1>0\), uniformly over its budget-feasible ridge
scales. Suppose also that a declared finite-net or information-theoretic
selection lower bound makes this fixed-kernel lower bound uniform over
the data-dependent choice of \(k\). Then no estimator in the stated
class learns the target to vanishing excess risk.

\end{lemma}

\begin{proof}\phantomsection\label{proof-kernel-minimax-hard-side}

For each fixed reachable kernel and budget-feasible ridge scale, one of
the two alternatives makes the assumed minimax lower bound at least
\(c_0c_1\) (capped by \(c_0\)). The assumed selection lower bound makes
this conclusion simultaneous for the learner's data-dependent kernel
choice. Hence the infimum excess risk over the declared learner class is
bounded away from zero. Alignment or empirical capture alone would not
imply this conclusion; the uniform source-residual/capacity and
selection hypotheses are essential.

\end{proof}

\begin{corollary}[Learning lower bound from family-specific structural
negligibility]\label{cor-family-floor-learning-hard-side}

Fix a presented task family and a budget-reachable kernel or
learned-kernel class \(\mathcal K\). Suppose a later family-specific
argument proves that every representative within budget in
\(\mathcal K\) has negligible counted structural capture in
\(I_n^{\mathrm{sat}}(B)\), unless its effective dimension or selection
cost exceeds the declared budget. Assume additionally a proved
capture-to-source bridge for this kernel class: negligible counted
capture implies a source residual bounded below by a fixed positive
constant. Then no budget-admissible learner
whose fixed or learned kernel lies in \(\mathcal K\) and satisfies the
hypotheses of the kernel minimax and selection theorem learns the
target to vanishing excess risk.

\end{corollary}

\begin{proof}\phantomsection\label{proof-family-floor-learning-hard-side}

Condition on the budget-admissible terminal kernel or on a net element
of the learned-kernel class. If its effective dimension exceeds the
declared statistical budget, the capacity term in the kernel minimax
lower bound cannot vanish within the available sample resources. If its
effective dimension lies within budget, the family-specific structural bound
says that the kernel has negligible counted capture, and the assumed
capture-to-source bridge bounds the source residual below. The
learned-operator selection hypothesis
accounts for data-dependent selection. Thus budget-admissible learning
would imply a structural representative within budget that the
family-specific bound excludes.

\end{proof}

Under the source-class, concentration, and minimax hypotheses stated
above, effective dimension provides matching-order upper and lower
bounds. Alignment determines whether the selected spectral subspace is
target-relevant: a low-capacity kernel misaligned with the target does
not contribute to \(I_n^{\mathrm{sat}}(B)\), while a highly aligned
kernel whose derivation, evaluation, selection, or effective dimension
exceeds the declared budget is inadmissible for the original
presentation.

\Cref{fig-two-sided-learning-certificate-pipeline} aligns the upper and lower
learning certificates on the same represented kernel family and budget.

\begin{figure}[tbp]
\centering
\resizebox{.96\linewidth}{!}{%
\begin{tikzpicture}[fw diagram,node distance=8mm and 11mm]
  \node[fw source,text width=2.8cm] (family)
    {fixed or selected\\represented kernel};
  \node[fw provenance,text width=3.4cm,above right=9mm and 14mm of family] (upperdata)
    {effective dimension\\source residual\\selection};
  \node[fw terminal,text width=2.8cm,right=16mm of upperdata] (upper)
    {upper oracle\\risk certificate};
  \node[fw provenance,text width=3.4cm,below right=9mm and 14mm of family] (lowerdata)
    {uniform source class\\minimax alternative};
  \node[fw terminal,text width=2.8cm,right=16mm of lowerdata] (lower)
    {lower capacity or\\source-residual bound};
  \draw[fw arrow] (family) -- (upperdata);
  \draw[fw arrow] (upperdata) -- (upper);
  \draw[fw arrow] (family) -- (lowerdata);
  \draw[fw arrow] (lowerdata) -- (lower);
  \node[fw box,text width=3.2cm,right=17mm of family] (budget)
    {same sample law\\loss and saturated\\budget};
  \draw[fw dashed arrow] (budget) -- (upper);
  \draw[fw dashed arrow] (budget) -- (lower);
\end{tikzpicture}%
}
\caption{Two-sided learning control requires common scope.  Effective
dimension and source residual give an upper oracle inequality, whereas a
uniform source-class or minimax argument gives the hard side.  Both are
interpretable together only after their kernel family, population law, loss,
and saturated budget have been matched.}
\label{fig-two-sided-learning-certificate-pipeline}
\end{figure}

\section{The Kernel Trick as the Lift
Channel}\label{the-kernel-trick-as-the-lift-channel}
\label{sec:kernel-trick-lift-channel}

The Lift channel of \PartII permits an admissible enlargement of the
observable configuration while preserving the original source of target
information. Kernel methods give the corresponding statistical
construction.

\begin{lemma}[Kernel trick as composition-bounded Lift]\label{thm-kernel-trick-lift}

Let \(k(x,z)=\langle\Psi(x),\Psi(z)\rangle\) be an admissible feature
lift, possibly infinite-dimensional, whose kernel evaluations are
budget-admissible public computations from the presentation or from the
learned lifted configuration.  For any empirical loss depending only on the
sample evaluations and any regularizer that is a strictly increasing function
of the RKHS norm, every empirical minimizer has the form

\[
\widehat f(x)=\sum_{i=1}^m\alpha_i k(x_i,x).
\]

Thus the computational lift is confined to the sample span, while
\(N_{\mathrm{eff}}(\lambda)\), rather than the ambient feature
dimension, controls statistical complexity. The lift contributes to the
structural count only when
the feature map preserves the observable/terminal interface required by
\PartII's valid-lift definition.

\end{lemma}

\begin{proof}\phantomsection\label{proof-kernel-trick-lift}

This is the representer theorem \citep{KimeldorfWahba1971}. Decompose any candidate
\(f\in\mathcal H_k\) into \(f=f_S+f_\perp\), where \(f_S\) lies in the
span of \(\{k_{x_i}\}_{i=1}^m\) and \(f_\perp\) is orthogonal to that
span. Since \(f_\perp(x_i)=\langle f_\perp,k_{x_i}\rangle=0\), it does
not change the empirical loss, while it increases the RKHS norm unless
it is zero. Hence the minimizer lies in the sample span. The KRR bound
in \Cref{the-tightened-two-sided-learning-bound} then charges the lifted span by the effective dimension of
its spectrum.

\end{proof}

\begin{corollary}[Budget lower bound for feature lifts]\label{cor-lift-hard-side}

Assume the uniform source-class, fixed-kernel minimax, and
data-dependent selection hypotheses of
\Cref{thm-kernel-minimax-hard-side}.  If every admissible aligned
kernel has
\(N_{\mathrm{eff}}(\lambda)/m\ge c_1>0\) at every budget-feasible
scale, or has the alternative source-residual lower bound in that
theorem, then samples within the budget cannot learn the declared source
class to vanishing excess risk. The conclusion applies to
infinite-coordinate feature lifts because it concerns effective
dimension and source residual, not ambient feature dimension.

\end{corollary}

\begin{proof}\phantomsection\label{proof-lift-hard-side}

Apply \Cref{thm-kernel-minimax-hard-side}.  Its uniform
fixed-kernel bound supplies the capacity/source-residual alternative,
and its selection hypothesis makes that alternative simultaneous for
the data-dependent terminal kernel.  The representer theorem identifies
the empirical sample span but is not itself used as a minimax lower
bound.  Infinite ambient dimension does not alter the stated uniform
effective-dimension and source-class hypotheses.

\end{proof}

The kernel trick is therefore an instance of valid Lift whenever the
kernel is publicly generated and terminal-compatible. Composition may
access a high- or infinite-dimensional feature representation, but the
target information available within budget is precisely the counted
spectral mass at the corresponding effective dimension. The remaining
mass is a source residual or part of \(\Res\), rather than
an additional source.

\section{Neural Networks and Learned Kernels}\label{neural-networks-the-learned-kernel-realization}
\label{sec:neural-networks-learned-kernels}

\subsection{Kernel evolution during training}
\label{subsec:kernel-evolution-training}

A budget-admissible oracle-free neural-network training procedure
constructs represented terminal outputs and readouts on its clocked
training transcript carrier;
ordinary polynomial-time training is the
\(\mathcal B_{\mathrm{poly}}\) specialization. The empirical neural
tangent kernel (NTK) is the Gram kernel of the tangent features used by
gradient descent.

The empirical neural tangent kernel records inner products of the network's tangent features during training.
\begin{definition}[Empirical neural tangent kernel]\label{def-ntk}

For a network \(f(x;\theta)\), the empirical NTK at parameter value
\(\theta\) is

\[
\Theta_\theta(x,z)
=
\left\langle\nabla_\theta f(x;\theta),\nabla_\theta f(z;\theta)\right\rangle.
\]

The Gram matrix \(\Theta_t\) generally evolves during training and is
approximately constant in the lazy-training regime
\citep{JacotEtAl2018}.

\end{definition}

The movement trajectory tracks how training changes the learned kernel and its target alignment.
\begin{definition}[Kernel movement and alignment trajectory]\label{def-kernel-movement}

On a sample, define

\[
\delta_t=\frac{\|\Theta_t-\Theta_0\|_F}{\|\Theta_0\|_F},
\qquad
A_t=A_m(\Theta_t,y).
\]

Small \(\delta_t\) together with nearly constant \(A_t\) is an
empirical diagnostic consistent with lazy training, but is not by
itself a prediction-linearization theorem. Order-one kernel movement accompanied by
increasing alignment is evidence of feature learning.

\end{definition}

\begin{lemma}[Lazy-training regime]\label{thm-lazy-ntk-bound}

Assume the complete network record contains a verified population prediction
comparison

\[
\|\widehat f_{\mathrm{NN}}-\widehat f_{\Theta_0}\|_{L^2(\mu)}
\le\varepsilon_{\mathrm{lazy}},
\]

where \(\widehat f_{\Theta_0}\) is the fixed-kernel flow or ridge
predictor.  Assume on the same event that \(\Theta_0\) carries the certified
fixed-kernel bound \eqref{eq-fixed-kernel-oracle-certificate}, with constants
\(C_{\rm var}\) and \(C_{\rm bias}\). Then

\[
\|\widehat f_{\mathrm{NN}}-f_*\|_{L^2(\mu)}^2
\le
2C_{\rm var}\frac{\sigma^2N_{\mathrm{eff}}(\Theta_0,\lambda)}{m}
+2C_{\rm bias}B_{\Theta_0}(\lambda;f_*)
+2\varepsilon_{\mathrm{lazy}}^2.
\]

The empirical diagnostic
\(\|\Theta_t-\Theta_0\|_F=o(\|\Theta_0\|_F)\) alone does not imply these
hypotheses; in particular, it supplies neither spectral-scale operator
control nor a population generalization theorem
\citep{JacotEtAl2018,LeeEtAl2019,AroraEtAl2019,ChizatEtAl2019}.

\end{lemma}

\begin{proof}\phantomsection\label{proof-lazy-ntk-bound}

By hypothesis, the network predictor is within
\(\varepsilon_{\mathrm{lazy}}\) in \(L^2\) of the fixed-kernel
predictor.  Thus
\[
\|\widehat f_{\mathrm{NN}}-f_*\|_{L^2(\mu)}^2
\le
2\|\widehat f_{\Theta_0}-f_*\|_{L^2(\mu)}^2
+2\varepsilon_{\mathrm{lazy}}^2.
\]
Substitution of \eqref{eq-fixed-kernel-oracle-certificate} gives the stated
bound. No conclusion is drawn from
Frobenius kernel movement alone.

\end{proof}

\subsection{Feature learning and selection cost}
\label{subsec:feature-learning-selection-cost}

The feature-learning bound charges movement of the learned operator together with its selection complexity.
\begin{lemma}[Feature-learning regime]\label{thm-feature-learning-bound}

Suppose budget-admissible oracle-free training selects a represented
terminal kernel \(\Theta_T\) from an admissibly parameterized
feature-map class \(\mathcal K_\Theta\) with covering number
\(N_\rho=\mathcal N(\mathcal K_\Theta,\rho)<\infty\). The terminal kernel, training
transcript, precision, and selection cover are charged as a \PartIII
saturated representative.  Let \(\mathcal K_\rho\) be the represented
\(\rho\)-net.  Assume that one simultaneous event \(\mathcal E_\rho\), of
probability at least \(1-\delta\), carries for every
\(K\in\mathcal K_\rho\) the fixed-kernel certificate
\begin{equation}
\label{eq-feature-net-simultaneous-certificate}
\operatorname{ExcessRisk}(K)
\le
C_{\rm var}\frac{\sigma^2N_{\mathrm{eff}}(K,\lambda)}{m}
+C_{\rm bias}B_K(\lambda;f_*)+r_{\rm sel}.
\end{equation}
Here \(C_{\rm var},C_{\rm bias}\), the confidence split, and the represented
radius \(r_{\rm sel}\) are part of the record.  For an independent audit
sample of size \(m_a\) and a loss in an interval of length \(L_{\rm loss}\),
one admissible simultaneous empirical-to-population radius is
\[
r_{\rm sel}
=L_{\rm loss}\sqrt{\frac{\log(2N_\rho/\delta)}{2m_a}}.
\]
Assume further that the fixed-kernel right side and the corresponding KRR
prediction risk are jointly \(L_{\mathrm{sel}}\)-Lipschitz in the declared
kernel operator metric.  Let represented nonnegative
\(\varepsilon_{\rm est}\) and \(\varepsilon_{\rm opt}\) bound kernel
estimation and optimization error in the same excess-risk units.  Then, on
\(\mathcal E_\rho\),

\[
\operatorname{ExcessRisk}
\le
C_{\rm var}\frac{\sigma^2N_{\mathrm{eff}}(\Theta_T,\lambda)}{m}
+C_{\rm bias}B_{\Theta_T}(\lambda;f_*)
+r_{\rm sel}
+L_{\mathrm{sel}}\rho
+\varepsilon_{\mathrm{est}}+\varepsilon_{\mathrm{opt}}.
\]

The first two terms are the learned kernel's certified capacity and source
residual.  The radius and stability term charge selection from the finite
cover. Any
target-relevant movement not represented by the terminal kernel, valid
lifts, or terminal readouts remains residual accounting rather than an
independent source.

\end{lemma}

\begin{proof}\phantomsection\label{proof-feature-learning-bound}

Choose \(K_T^\rho\in\mathcal K_\rho\) within distance \(\rho\) of
\(\Theta_T\).  On \(\mathcal E_\rho\), apply
\eqref{eq-feature-net-simultaneous-certificate} to \(K_T^\rho\).
Joint Lipschitz stability transports its right side and prediction risk to
\(\Theta_T\) at total cost at most \(L_{\mathrm{sel}}\rho\); the two named
algorithmic errors add in the same units.  For the displayed independent
audit option, Hoeffding's inequality gives failure probability at most
\(2\exp(-2m_ar^2/L_{\rm loss}^2)\) for one net point.  Substituting the
displayed \(r_{\rm sel}\) and taking a union bound over \(N_\rho\) points
gives failure probability at most \(\delta\).

\PartII supplies the operator embedding on the
clocked training transcript carrier. \PartIII includes the terminal kernel,
transcript, cover resolution, precision, and readout in the saturated
derivability accounting before the learned kernel can contribute to
extraction.

\end{proof}

\Cref{fig-kernel-training-regimes} organizes the two represented kernel
training regimes of \cref{thm-lazy-ntk-bound,thm-feature-learning-bound}.

\begin{figure}[tbp]
\centering
\begin{tikzpicture}[fw diagram]
  \node[fw source,minimum width=3.4cm] (train) at (0,2.5)
    {represented training transcript};
  \node[fw geom,minimum width=3.2cm] (lazy) at (-3.8,0.9)
    {lazy regime\\kernel \(\Theta_0\)};
  \node[fw provenance,minimum width=3.2cm] (feature) at (3.8,0.9)
    {feature learning\\terminal \(\Theta_T\)};
  \node[fw use,minimum width=3.2cm] (fixed) at (-3.8,-0.8)
    {fixed-kernel capacity\\and source residual};
  \node[fw use,minimum width=3.5cm] (selected) at (3.8,-0.8)
    {capacity + residual\\+ selection charge};
  \node[fw terminal,minimum width=4.1cm] (risk) at (0,-2.3)
    {excess-risk control};
  \draw[fw arrow] (train) -- node[above left,font=\scriptsize] {\(\Theta_t\simeq\Theta_0\)} (lazy);
  \draw[fw arrow] (train) -- node[above right,font=\scriptsize] {selected movement} (feature);
  \draw[fw arrow] (lazy) -- (fixed);
  \draw[fw arrow] (feature) -- (selected);
  \draw[fw arrow] (fixed) -- (risk);
  \draw[fw arrow] (selected) -- (risk);
\end{tikzpicture}
\caption{Under the hypotheses of
\cref{thm-lazy-ntk-bound,thm-feature-learning-bound}, lazy training is
controlled by the initial kernel, whereas feature learning uses the terminal
kernel and pays its represented selection charge.}
\label{fig-kernel-training-regimes}
\end{figure}

\begin{lemma}[Observable-success accountability and residual audit]\label{thm-observable-success-certifiability}

Let \(\{I_n\}\) be a presented task family with exposed target
sectors \(T_I\subseteq X_I\) and declared budget structure
\(\mathcal B\). Let \(\mathcal A\) be a uniform oracle-free represented
computation whose operational transcript, lifted coordinates, precision,
and output map have cost within \(\mathcal B\), with finite normalized horizon
\(H_B(n)\). If

\[
\Pr[\mathcal A(I)\in T_I]\ge \Gamma(n),
\]

then its clocked transcript representation induces the exact first-hit
record

\[
q_{\mathcal A}\in
\mathcal Q_{n,\le B_{\mathrm{hit}}}^{\mathrm{sat}},
\qquad
B_{\mathrm{hit}}=\Psi_{\mathrm{hit}}(b(n),n),
\]

and, pointwise in the instance,

\begin{equation}
\label{eq-observable-success-exact-first-hit-scale}
V_I(q_{\mathcal A})
\ge \Gamma_{\mathrm{hit}}(n)
:=\frac{\Gamma(n)}{eH_B(n)}.
\end{equation}

This counted exact-quotient witness is unconditional within the declared
finite resource envelope. In the polynomial instance,
\(B_{\mathrm{hit}}=n^{O_C(1)}\) for a time-\(n^C\) algorithm. Separately,
for every
residual audit class \(\mathcal H_{B}\) generated from the same lifted
transcript and represented at budget \(B_{\mathrm{hit}}\), put

\[
\mathcal P_{\mathcal H_B}(M,\delta,\rho)
=
c\left(
\sqrt{\frac{N_{\mathcal H_B}+\log\mathcal N(\mathcal H_B,\rho)+\log(1/\delta)}{M}}
+\rho
\right),
\]

where \(N_{\mathcal H_B}\) is the effective mode count of the audited
residual span, \(M\) is the number of audit samples or independent
family instances, and \(\rho\) is the operator-net resolution.  For the
following audit assertion, assume the independent-sample,
bounded-statistic, and Lipschitz/cover hypotheses of
\Cref{cor-rademacher-audit-ceiling}, together with a proved
class bound

\[
\mathfrak R_M(\mathcal H_B)
\le
c\left(
\sqrt{\frac{N_{\mathcal H_B}+\log\mathcal N(\mathcal H_B,\rho)}{M}}
+\rho
\right).
\]

With
probability at least \(1-\delta\) over the audit sample, every statistic
in that declared class has population audited capture at most its
empirical audited capture plus
\(\mathcal P_{\mathcal H_B}(M,\delta,\rho)\). Thus an audit resolves
effects at the forced-witness scale when
\(\mathcal P_{\mathcal H_B}=o(\Gamma_{\mathrm{hit}})\). If the penalty
is not smaller than that scale, this audit is inconclusive; failure of a
sufficient Rademacher bound is not an indistinguishability or hardness
theorem.

A terminal observable outside \(\mathcal D_{n,\le B_{\mathrm{hit}}}\) is
inadmissible for the original presentation. It requires oracle or advice
data, a changed primitive signature, or discovery, representation,
precision, or application cost above the declared budget.

\end{lemma}

\begin{proof}\phantomsection\label{proof-observable-success-certifiability}

\Cref{lem:exact-affordable-first-hit} applies to the complete
configuration record and gives both the cost
\(B_{\mathrm{hit}}=\Psi_{\mathrm{hit}}(b(n),n)\) and
\eqref{eq-observable-success-exact-first-hit-scale}. This step already
includes every exposed channel, residual interaction, memory effect,
and terminal readout; no componentwise transfer premise is used.

For a separately declared residual audit class, the explicitly assumed
Rademacher estimate and
\Cref{cor-rademacher-audit-ceiling} give the stated uniform
empirical-to-population inequality. Comparing its penalty to
\(\Gamma_{\mathrm{hit}}\) gives the audit-resolution statement and no
converse lower bound.

The final assertion follows from the admissibility dichotomy of \PartIII.
A terminal object absent from the derivability filtration of the
original presentation is not a budget-admissible representative: it
belongs to an oracle or advice extension, changes the primitive
signature, or requires an over-budget construction.

\end{proof}

\begin{corollary}[Parameter-cover selection cost]\label{cor-parameter-cover-selection-cost}

If the parameter domain has dimension \(p\), radius \(R_\Theta\), and
the map \(\theta\mapsto\Theta_\theta\) is \(L_K\)-Lipschitz in empirical
operator norm, then

\[
\log\mathcal N(\mathcal K_\Theta,\rho)
\le
p\log\left(1+\frac{2R_\Theta L_K}{\rho}\right).
\]

For neural networks, \(L_K\) may be upper-bounded by layer Lipschitz or
spectral-norm controls.  For the independent bounded-loss audit option in
\cref{thm-feature-learning-bound}, substitution into its simultaneous
radius gives the explicit certificate
\[
r_{\rm sel}
\le
L_{\rm loss}
\sqrt{\frac{
\log(2/\delta)
+p\log(1+2R_\Theta L_K/\rho)}{2m_a}}.
\]

\end{corollary}

\begin{proof}\phantomsection\label{proof-parameter-cover-selection-cost}

A radius-\(R_\Theta\) subset of \(\mathbb R^p\) has a Euclidean
\(\epsilon\)-net of cardinality at most \((1+2R_\Theta/\epsilon)^p\). If
\(\theta\mapsto\Theta_\theta\) is \(L_K\)-Lipschitz in empirical
operator norm, an \(\epsilon=\rho/L_K\) parameter net maps to a
\(\rho\)-net of terminal kernels. Taking logarithms gives the displayed
covering bound. Substituting it into the exact radius in
\cref{thm-feature-learning-bound} gives the second display.

\end{proof}

\begin{corollary}[Feature-learning capacity tradeoff]\label{cor-feature-learning-conservation}

Under the hypotheses of \Cref{thm-feature-learning-bound}, write

\[
B_\Theta(\lambda)
=B_\Theta(\lambda;f_*).
\]

Assume the initial kernel carries the same fixed-kernel certificate constants
\(C_{\rm var},C_{\rm bias}\) and no corresponding selection charges.  The
terminal-kernel upper certificate is no larger than the initial-kernel upper
certificate whenever

\[
C_{\rm var}\frac{\sigma^2}{m}
\left[
N_{\mathrm{eff}}(\Theta_0,\lambda)
-N_{\mathrm{eff}}(\Theta_T,\lambda)
\right]
+C_{\rm bias}
\left[B_{\Theta_0}(\lambda)-B_{\Theta_T}(\lambda)\right]
\ge
r_{\rm sel}
+L_{\mathrm{sel}}\rho
+\varepsilon_{\mathrm{est}}+\varepsilon_{\mathrm{opt}}.
\]

When the source terms agree and the other errors are negligible, this
reduces to

\[
N_{\mathrm{eff}}(\Theta_0,\lambda)
-N_{\mathrm{eff}}(\Theta_T,\lambda)
\ge
\frac{m\,r_{\rm sel}}{C_{\rm var}\sigma^2}.
\]

The alignment-adjusted quantity

\[
\Delta\left(\frac{N_{\mathrm{eff}}}{A_m}\right)
=
\frac{N_{\mathrm{eff}}(\Theta_0,\lambda)}{A_m(\Theta_0,y)}
-
\frac{N_{\mathrm{eff}}(\Theta_T,\lambda)}{A_m(\Theta_T,y)}
\]

remains the structural diagnostic of
\Cref{cor-structural-learning-proxy}; it is not substituted
into the risk inequality.

\end{corollary}

\begin{proof}\phantomsection\label{proof-feature-learning-conservation}

Subtract the terminal-kernel certificate in
\cref{thm-feature-learning-bound} from the initial fixed-kernel certificate.
The resulting difference is exactly the left side of the displayed
condition, while the terminal-only charges form its right side.  The
capacity-only specialization sets the source difference and all charges
except \(r_{\rm sel}\) to zero. No alignment denominator appears
in either oracle inequality.

\end{proof}

\subsection{Composite capacity control}
\label{subsec:composite-capacity-control}

The composite bound reports the stronger of the learned-kernel and network-capacity controls.
\begin{lemma}[Composite neural-network bound]\label{thm-composite-nn-bound}

Assume the hypotheses of \Cref{thm-feature-learning-bound} for the trained
network's terminal kernel.  Let \(\mathcal E_{\rm ker}\) be its declared
simultaneous event, of probability at least \(1-\delta_{\rm ker}\), on which
the same declared population risk \(R(f)\) is bounded by

\[
\mathcal B_{\mathrm{ker}}
=
C_{\rm var}\frac{\sigma^2N_{\mathrm{eff}}(\Theta_T,\lambda)}{m}
+C_{\rm bias}B_{\Theta_T}(\lambda;f_*)
+r_{\rm sel}
+L_{\mathrm{sel}}\rho
+\varepsilon_{\mathrm{est}}+\varepsilon_{\mathrm{opt}}.
\]

Independently, suppose a named represented margin theorem supplies a finite
quantity \(\mathcal B_{\rm margin}\) and an event
\(\mathcal E_{\rm margin}\), of probability at least
\(1-\delta_{\rm margin}\), such that
\[
R(f)\le\mathcal B_{\rm margin}
\quad\text{on }\mathcal E_{\rm margin}.
\]
That certificate must state its margin convention, loss conversion, norm
normalization, architectural hypotheses, constants, and confidence term;
for example it may instantiate a fully stated spectral-margin theorem
\citep{BartlettFosterTelgarsky2017}.  It must cover the actual data-dependent
trained weights, rather than only a fixed network chosen before the sample.

Then, with probability at least
\(1-\delta_{\rm ker}-\delta_{\rm margin}\),

\[
R(f)
\le
\min\{\mathcal B_{\mathrm{ker}},\mathcal B_{\mathrm{margin}}\}
.
\]
The minimum is permitted because both quantities bound the same population
risk on one simultaneous event.  Terms already included in one certificate
are not added again to the other.

\end{lemma}

\begin{proof}\phantomsection\label{proof-composite-nn-bound}

On \(\mathcal E_{\rm ker}\cap\mathcal E_{\rm margin}\), both displayed
inequalities hold for the same random trained network and the same
population risk.  Hence their minimum is an upper bound.  The union bound
gives
\(\Pr(\mathcal E_{\rm ker}\cap\mathcal E_{\rm margin})
\ge1-\delta_{\rm ker}-\delta_{\rm margin}\).

\end{proof}

\section{Estimating the Learned Effective
Degree}\label{measuring-the-effective-degree-a-network-learned}
\label{sec:estimating-learned-effective-degree}

The spectrum of the learned kernel provides a post-training estimate of
the effective degree captured by the model.

\begin{definition}[Label spectral profile and effective degree]\label{def-label-spectral-profile}

Let \(K=\sum_{i=1}^m s_i u_iu_i^T\) be the centered learned Gram matrix,
with \(s_1\ge\cdots\ge s_m\ge0\). For centered labels \(y_c\), define

\[
\beta_i=\frac{\langle y_c,u_i\rangle^2}{\|y_c\|^2},
\qquad
\sum_i\beta_i=1.
\]

The cumulative label mass in the top \(k\) modes is
\(M(k)=\sum_{i\le k}\beta_i\). The effective degree at tolerance
\(\varepsilon\) is

\[
d_{\mathrm{eff}}(\varepsilon)
=
\min\{k:M(k)\ge1-\varepsilon\}.
\]

The smooth participation form is
\(d_{\mathrm{part}}=(\sum_i\beta_i)^2/\sum_i\beta_i^2\).

\end{definition}

\begin{lemma}[Spectral upper bound conditional on profile concentration]\label{thm-degree-spectrum-upper}

Let \(M_{\mathrm{pop}}(k)\) and \(M_{\mathrm{emp}}(k)\) be cumulative
population and empirical label-mass profiles in corresponding ordered
spectral bands. Assume that an independently established concentration
event of probability at least \(1-\delta\) gives
\[
\sup_k|M_{\mathrm{pop}}(k)-M_{\mathrm{emp}}(k)|
\le\Delta_m.
\]
Then, on that event,

\[
d_{\mathrm{eff}}^{\mathrm{pop}}(\varepsilon+\Delta_m)
\le
\widehat d_{\mathrm{eff}}(\varepsilon),
\]

For bounded kernels, a concrete value of \(\Delta_m\) may be supplied
by an operator-concentration and Davis--Kahan estimate for hard
eigengap-separated bands, or by a Lipschitz functional-calculus estimate
for declared smoothed bands.

\end{lemma}

\begin{proof}\phantomsection\label{proof-degree-spectrum-upper}

Put \(k=\widehat d_{\mathrm{eff}}(\varepsilon)\). By definition,
\(M_{\mathrm{emp}}(k)\ge1-\varepsilon\). The assumed event gives
\(M_{\mathrm{pop}}(k)\ge1-\varepsilon-\Delta_m\), so the smallest
population prefix containing mass \(1-(\varepsilon+\Delta_m)\) has
size at most \(k\). This is the displayed inequality. Matrix Bernstein
and Davis--Kahan are standard ways to verify the stated event when their
boundedness and eigengap hypotheses hold
\citep{Tropp2012,DavisKahan1970}.

\end{proof}

\begin{lemma}[Generalization certifies captured mass]\label{thm-generalization-captures-mass}

Assume the centered target is normalized so that
\(\|f_*\|_{L^2(\mu)}^2=1\), and measure squared-loss test error \(r\)
and noise level \(\sigma^2\) in these units. Assume an independent test
bound gives \(R(\widehat f)\le r\), and assume a separately proved
lower approximation inequality
\[
R(\widehat f)
\ge
1-m_{\Theta_T}(\lambda)
-C\frac{\sigma^2N_{\mathrm{eff}}(\lambda)}{m}
-\sigma^2-o(1).
\]
Then the ridge-feasible band captured at least

\[
m_{\Theta_T}(\lambda)
\ge
1-r-C\frac{\sigma^2N_{\mathrm{eff}}(\lambda)}{m}-\sigma^2-o(1)
\]

of the normalized target mass, whenever the kernel-bound hypotheses of
\Cref{neural-networks-the-learned-kernel-realization} apply.

\end{lemma}

\begin{proof}\phantomsection\label{proof-generalization-captures-mass}

Combine the assumed upper and lower bounds:
\[
1-m_{\Theta_T}(\lambda)
-C\frac{\sigma^2N_{\mathrm{eff}}(\lambda)}m
-\sigma^2-o(1)
\le r.
\]
Rearranging proves the claim. The lower approximation inequality is
essential; an upper kernel-risk bound alone cannot be inverted to infer
captured mass.

\end{proof}

\begin{lemma}[Population--empirical prefix-index sandwich]\label{thm-effective-degree-sandwich}

Let \(M_{\mathrm{pop}}(k)\) and \(M_{\mathrm{emp}}(k)\) be the
population and empirical cumulative target masses in the first \(k\)
kernel modes, and suppose
\[
\sup_k|M_{\mathrm{pop}}(k)-M_{\mathrm{emp}}(k)|\le\Delta_m.
\]
For \(a\in[0,1]\), let \(k_{\mathrm{pop}}(a)\) and
\(k_{\mathrm{emp}}(a)\) be the smallest prefix indices whose respective
cumulative masses are at least \(a\), using the endpoint conventions
\(k(a)=0\) for \(a\le0\) and \(k(a)=\infty\) for \(a>1\). Then

\[
k_{\mathrm{emp}}(a-\Delta_m)
\le
k_{\mathrm{pop}}(a)
\le
k_{\mathrm{emp}}(a+\Delta_m).
\]

When \cref{thm-generalization-captures-mass} supplies the mass level
\(q=1-r-C\sigma^2N_{\mathrm{eff}}(\lambda)/m-\sigma^2-o(1)\), set
\(a=q\) to obtain empirical lower and upper indices for the population
prefix needed to contain that mass.

\end{lemma}

\begin{proof}\phantomsection\label{proof-effective-degree-sandwich}

If \(M_{\mathrm{pop}}(k)\ge a\), then
\(M_{\mathrm{emp}}(k)\ge a-\Delta_m\); applying this at
\(k=k_{\mathrm{pop}}(a)\) gives the left inequality. Conversely,
\(M_{\mathrm{emp}}(k)\ge a+\Delta_m\) implies
\(M_{\mathrm{pop}}(k)\ge a\); applying this at
\(k=k_{\mathrm{emp}}(a+\Delta_m)\) gives the right inequality.

\end{proof}

\begin{remark}[Scope of finite-sample degree error bars]
\label{prop-effective-degree-error-bars}
\label{rem:effective-degree-error-bar-scope}

Bounded kernels and sub-Gaussian labels alone do not determine a universal
finite-sample radius for a data-dependent spectral cutoff.  A quantitative
degree certificate must name, on one simultaneous event,

\begin{enumerate}[label=\textup{(\roman*)}]
\item an operator estimate
      \(\|\widehat T-T\|_{\mathrm{op}}\le\varepsilon_T(m,\delta_T)\);
\item either a population eigengap \(g_\tau>2\varepsilon_T\) at every hard
      cutoff used by the procedure, or a displayed Lipschitz constant for a
      smoothed spectral filter;
\item a label-moment estimate in the same population law and norm;
\item simultaneous confidence allocation over every selected cutoff,
      regularization level, architecture, and subsampling rule.
\end{enumerate}

For a hard cutoff, Davis--Kahan converts the first item to projector error at
most \(2\varepsilon_T/g_\tau\).  For a smoothed filter \(\varphi\), an
applicable functional-calculus theorem instead contributes its displayed
operator-Lipschitz constant times \(\varepsilon_T\).  Either term must be
added to the separately certified label-moment radius before
\cref{thm-effective-degree-sandwich} is applied.  Effective-dimension and
participation-ratio estimates likewise require, respectively, a stated
trace/resolvent perturbation bound and a certified positive denominator.
DKW, bootstrap, Nyström, and delta-method conclusions are available only
through those named hypotheses; this remark asserts no distribution-free
rate for them
\citep{DvoretzkyKieferWolfowitz1956,Massart1990,DavisKahan1970}.

Independent held-out error is separate.  For a \([0,1]\)-valued loss and
\(M\) models fixed before the held-out sample, Hoeffding and a union bound
give the explicit simultaneous radius
\[
\sqrt{\frac{\log(2M/\delta)}{2m_{\mathrm{test}}}}.
\]

\end{remark}

\begin{proof}\phantomsection\label{proof-effective-degree-error-bars}

Bounded kernels give operator-norm concentration of the empirical
Gram/integral operator by matrix Bernstein \citep{Tropp2012}. For a fixed spectral band,
applying the spectral projector to sub-Gaussian labels gives the stated
square-root concentration; hard cutoffs add the Davis--Kahan projector
perturbation term \citep{DavisKahan1970}, while smoothed projectors are Lipschitz and avoid
eigengap denominators. Uniform cutoff control follows by applying a
DKW-style envelope to the cumulative smoothed spectral distribution, or
by a union bound over the finite empirical spectrum. The map
\(K\mapsto\operatorname{Tr}(K(K+\lambda mI)^{-1})\) is Lipschitz on the
positive cone away from the numerical floor with constant controlled by
\(\lambda^{-1}\). The participation ratio is a smooth function of the
first two spectral moments, so the delta method gives its rate.
Independent test error is an average of bounded or sub-exponential
losses, controlled by Hoeffding or Bernstein.

\end{proof}

Sweeping architectures gives the problem-level degree relative to the
model class:

\[
d_{\mathrm{eff}}^{\mathrm{problem}}(\varepsilon)
=
\min_{\text{admissible architectures}} d_{\mathrm{eff}}^{\mathrm{captured}}(\varepsilon).
\]

If every admissible architecture has low alignment or requires high
effective degree before attaining an error level within budget, the
architecture sweep provides an empirical lower-bound diagnostic for that
model class.

\section{Computational Evaluation Procedure}\label{computable-recipes}
\label{sec:computational-evaluation-procedure}

The preceding quantities can be evaluated from empirical operators and
Gram matrices as follows.

\begin{itemize}
\tightlist
\item
  \textbf{Initialization analysis.} Compute \(\Theta_0\),
  \(N_{\mathrm{eff}}(\lambda)\), \(A_m(\Theta_0,y)\), and the label
  spectral profile before full training.
\item
  \textbf{Training-regime classification.} Track
  \(\delta_t=\|\Theta_t-\Theta_0\|_F/\|\Theta_0\|_F\) and \(A_t\).
  Approximately constant values characterize lazy training, whereas
  increasing alignment indicates feature learning.
\item
  \textbf{Post-training feature-learning bound.} Compare
  \(N_{\mathrm{eff}}/A\) at initialization and convergence, then
  subtract the selection penalty \(\sqrt{\log\mathcal N/m}\).
\item
  \textbf{Composite bound.} Report the minimum of the learned-kernel
  and spectral-margin bounds.
\item
  \textbf{Model-class lower-bound diagnostic.} Compare admissible
  architectures. If every learned kernel within budget has low alignment
  or requires high effective degree before capture, the target has no
  counted low-degree structure in that model class.
\item
  \textbf{Fixed-operator analysis.} For any declared
  operator \(A\), compute \(\operatorname{Tr}(A^{-1})\), \(y^TAy\),
  \(\operatorname{TaskCap}_A(y)\), and the fallback identity score. This
  determines whether the operator yields a structural reduction relative
  to the identity reference or only enlarges the hypothesis class.
\item
  \textbf{Random-label null comparison.} For an operator family
  with a polynomial net, compare \(y^TAy\) against the sign-null mean
  \(\operatorname{Tr}(A)\) and variance \(2\sum_{i\ne j}A_{ij}^2\). A
  target whose score stays at the null scale across the net has no
  counted operator-level discount in that class.
\item
  \textbf{Finite-sample observable analysis.} From the exposed feature
  matrix, quotient fibers, labels, kernel Gram matrix, and audit
  parameters, compute \PartI's \(\widehat{\operatorname{ProjMass}}\),
  \(\widehat{\operatorname{ActMass}}\), inert mass, threshold margin,
  monotone-filtration score, \(N_{\mathrm{eff}}\), captured kernel mass,
  Rademacher penalty, residual-order ceiling, null score, and certified
  upper bound. These computations implement the budgeted extraction
  accounting on finite data.
\end{itemize}

\section{Structural Generalization
Bounds}\label{structural-deployment-certificates}
\label{sec:structural-generalization-bounds}

This section extends the kernel, NTK, effective-dimension, and
feature-learning analysis. It quantifies statistical capacity,
structural KL divergence, effective-resistance defect, coefficient
selection, defect-geometry profiles, and uncertainty from learning the
geometry. Each quantity is computed from a represented operator,
unlabeled geometry, labeled sample, or declared model family and is
included in the same saturated accounting. These statistical bounds do
not replace a family-specific structural lower bound such as \PartVIII.

Whenever the required objects are represented, we use the
defect-restricted, harmonic-reduced formulation. Effective resistance is
the Poincar\'e component of the more general defect-geometry profile.
The global spectral-gap and isotropic-prior expressions serve as
comparison bounds when the reduced profile cannot be represented.

\subsection{Defect-restricted PAC--Bayes control}
\label{subsec:defect-restricted-pac-bayes}

We first restrict the PAC--Bayes comparison to the represented defect profile
selected by the structural ledger.

\begin{definition}[Defect-restricted structural PAC--Bayes divergence]\label{def-structural-pac-bayes-kl}

Let \(\mathbf M_{\mathcal D}\succeq0\) be the domain Dirichlet operator
computed from unlabeled input geometry. Let \(H\) be the represented
coupled hypothesis subspace after the Abs coordinates and harmonic Geom
slack are fixed, and let \(\Pi_E\) project to the represented
lumpability-defect directions. Let
\(R_{E,H}\subseteq H\cap\operatorname{ran}\Pi_E\) be the represented
retained subspace, and let \(P_{E,H}\) be its orthogonal projection.
Define the reduced structural precision, as an operator on
\(R_{E,H}\), by

\[
\mathbf M_{E,H}
=
\left.P_{E,H}\mathbf M_{\mathcal D}P_{E,H}\right|_{R_{E,H}}.
\]

Put \(R_+=\operatorname{ran}\mathbf M_{E,H}\). The structural prior is
the proper Gaussian measure on \(R_+\)
\[
P^{\mathsf{str}}_{E,H}
=\mathcal N_{R_+}(0,\mathbf M_{E,H}^{+}).
\]
There is no improper flat factor on null or harmonic directions. If a
model retains such directions, it must assign them a separate proper,
label-independent prior and include their KL contribution. For a
Gaussian posterior
\(Q=\mathcal N_{R_+}(\widehat w,\Sigma_Q)\), with \(\Sigma_Q\) positive
definite on \(R_+\), define

\[
\mathrm{KL}^{\mathsf{str}}_{E,H}(Q)
=
\frac12\left[
\widehat w^T\mathbf M_{E,H}\widehat w
+
\operatorname{tr}(\mathbf M_{E,H}\Sigma_Q)
-
\log\det{}_{+}(\mathbf M_{E,H}\Sigma_Q)
-d_{E,H}
\right],
\]

where \(d_{E,H}=\operatorname{rank}(\mathbf M_{E,H})\) and \(\det_+\) is
the pseudo-determinant on the positive spectral subspace. The
mean-dependent part is

\[
\frac12\widehat w^T\mathbf M_{E,H}\widehat w.
\]

For graph or quotient defects this mean term is the
effective-resistance energy
\(\frac12\langle E_q,\mathcal L_W^+E_q\rangle\) after the represented
identification between coefficient directions and defect currents.

\end{definition}

\begin{lemma}[PAC--Bayes bound with a defect-restricted prior]\label{thm-structural-pac-bayes}

Assume the loss is bounded in \([0,1]\) and the structural prior
\(P^{\mathsf{str}}_{E,H}\) is fixed independently of the labeled
sample, for example from a fixed public geometry or an independent
unlabeled sample. Then for every posterior \(Q\) on \(R_+\), with
probability at least \(1-\delta\) over the labeled sample,

\[
R(Q)
\le
\widehat R_S(Q)
+
\sqrt{
\frac{
\mathrm{KL}^{\mathsf{str}}_{E,H}(Q)+\log(2\sqrt m/\delta)
}{2(m-1)}
}.
\]

If the defect and harmonic projections are unavailable in the
representation, the same bound may be applied on the positive spectral
range of \(\mathbf M_{\mathcal D}\) with a proper Gaussian prior there,
together with proper priors for any retained null directions. No
ordering between the reduced and unreduced KL terms is asserted without
the factorization hypotheses of \Cref{prop-structural-pac-bayes-reductions}.
When source cells are selected or mixed, the prior and posterior use the
cell-mixture decomposition in
\cref{eq-source-resolved-learning-pac-bayes-mixture}; the cell-selection
divergence and every source-dependent prior construction remain charged.

\end{lemma}

\begin{proof}\phantomsection\label{proof-structural-pac-bayes}

By hypothesis, the prior is fixed independently of the labeled sample
to which the inequality is applied. It is therefore an admissible
PAC--Bayes prior. McAllester's
PAC-Bayes inequality for bounded losses gives the displayed inequality
for any fixed legal prior \citep{McAllester1999,Maurer2004}.
Substituting \(P^{\mathsf{str}}_{E,H}\) gives
the KL in the definition above. The Gaussian KL formula on a positive
semidefinite precision is the standard finite-dimensional formula
restricted to \(R_+\); it is a KL between two proper probability
measures on the same finite-dimensional space. The mean-dependent term is
\(\frac12\widehat w^T\mathbf M_{E,H}\widehat w\). When the represented
coordinates encode a defect current \(E_q\) routed through a graph
Laplacian \(\mathcal L_W\), solving the compensation equation
\(\mathcal L_W\phi=E_q\) gives energy
\(\langle E_q,\mathcal L_W^+E_q\rangle\), so the KL divergence consists
of the effective-resistance term together with covariance and
normalization terms.
An unreduced proper prior is obtained by choosing the full represented
positive spectral subspace. It is a valid alternative, but is not
automatically weaker in KL.

\end{proof}

\Cref{fig-defect-pac-bayes-pipeline} shows how the represented defect geometry
enters the PAC--Bayes certificate without adding an uncharged prior direction.

\begin{figure}[tbp]
\centering
\begin{tikzpicture}[fw diagram]
  \node[fw source,minimum width=3.0cm] (domain) at (-5.3,0)
    {unlabeled geometry\\\(\mathbf M_{\mathcal D}\)};
  \node[fw provenance,minimum width=3.0cm] (reduce) at (-1.4,0)
    {represented defect range\\\(R_{E,H}\)};
  \node[fw geom,minimum width=3.0cm] (prior) at (2.5,1.0)
    {proper structural prior\\\(P^{\mathsf{str}}_{E,H}\)};
  \node[fw use,minimum width=3.0cm] (post) at (2.5,-1.0)
    {posterior\\\(Q\)};
  \node[fw terminal,minimum width=2.7cm] (risk) at (6.1,0)
    {risk certificate};
  \draw[fw arrow] (domain) -- (reduce);
  \draw[fw arrow] (reduce) -- (prior);
  \draw[fw arrow] (reduce) -- (post);
  \draw[fw arrow] (prior) -- (risk);
  \draw[fw arrow] (post) -- (risk);
\end{tikzpicture}
\caption{The defect-restricted construction of
\cref{def-structural-pac-bayes-kl,thm-structural-pac-bayes}. The prior is
formed on the represented positive defect range; retained null directions
require their own proper prior and KL charge.}
\label{fig-defect-pac-bayes-pipeline}
\end{figure}

\begin{proposition}[KL comparison under reducing factorization]\label{prop-structural-pac-bayes-reductions}

Suppose an orthogonal decomposition
\(R_{\mathrm{full}}=R_{E,H}\oplus R_{\mathrm{rem}}\) reduces
\(\mathbf M_{\mathcal D}\), so the full proper Gaussian prior factors
as \(P_{\mathrm{full}}=P_{E,H}\otimes P_{\mathrm{rem}}\). Suppose also
that the full posterior factors as
\(Q_{\mathrm{full}}=Q_{E,H}\otimes Q_{\mathrm{rem}}\). Then

\[
\mathrm{KL}(Q_{E,H}\Vert P_{E,H})
\le
\mathrm{KL}(Q_{\mathrm{full}}\Vert P_{\mathrm{full}}),
\]

with equality exactly when
\(\mathrm{KL}(Q_{\mathrm{rem}}\Vert P_{\mathrm{rem}})=0\).
For a graph defect, the exact penalty is

\[
\mathcal R_W(E_q)=\langle E_q,\mathcal L_W^+E_q\rangle
=
\sum_{k\ge2}\lambda_k^{-1}\langle E_q,v_k\rangle^2.
\]

The global-gap expression \(\|E_q\|^2/\lambda_2\) is the comparison
upper bound

\[
\mathcal R_W(E_q)\le \frac{\|E_q\|^2}{\lambda_2};
\]

and is exact only when the defect lies in the slowest active
eigendirection. Isotropic PAC--Bayes is recovered when the
domain precision is isotropic on the represented subspace and the
defect/harmonic restrictions are trivial.

\end{proposition}

\begin{proof}\phantomsection\label{proof-structural-pac-bayes-reductions}

The KL chain rule for the assumed product measures gives

\[
\mathrm{KL}(Q_{\mathrm{full}}\Vert P_{\mathrm{full}})
=
\mathrm{KL}(Q_{E,H}\Vert P_{E,H})
+
\mathrm{KL}(Q_{\mathrm{rem}}\Vert P_{\mathrm{rem}}).
\]

Nonnegativity of KL proves the comparison and its equality condition.
Without the reducing and factorization hypotheses, compression of a
precision matrix need not order Gaussian KL divergences. Diagonalizing
\(\mathcal L_W=\sum_k\lambda_kv_kv_k^T\) gives the displayed
effective-resistance identity. Since \(\lambda_k\ge\lambda_2\) for
\(k\ge2\), the global-gap inequality follows, with equality exactly when
all defect mass is on the \(\lambda_2\) eigenspace. If
\(\mathbf M_{\mathcal D}=\sigma^{-2}I\) and no directions are removed,
the structural KL is the ordinary isotropic Gaussian KL.

\end{proof}

\subsection{Selection of structural loss coefficients}
\label{subsec:structural-loss-coefficient-selection}

The next bounds choose the loss coefficients from the same certified structural
quantities.

\begin{lemma}[Selection of loss coefficients]\label{thm-principled-lambda-selection}

Let \(\Lambda\) be a finite, declared set of admissible coefficient
vectors. For each \(\lambda\in\Lambda\), let the training computation
produce a posterior \(Q_\lambda\) from a represented loss

\[
\widehat R_S(\theta)+\sum_{j=1}^r\lambda_j C_j(\theta),
\]

where each \(C_j\) is a computable certificate term, for example
\(\mathcal R_W(E_q)\), effective dimension, harmonic capacity, or
connectivity. With probability at least \(1-\delta\), simultaneously for
all \(\lambda\in\Lambda\),

\[
R(Q_\lambda)
\le
\widehat R_S(Q_\lambda)
+
\sqrt{
\frac{
\mathrm{KL}^{\mathsf{str}}_{E,H}(Q_\lambda)+
\log(2\sqrt m|\Lambda|/\delta)
}{2(m-1)}
}.
\]

Thus the coefficient vector minimizing the simultaneous bound is

\[
\lambda^*
\in
\operatorname*{argmin}_{\lambda\in\Lambda}
\left[
\widehat R_S(Q_\lambda)
+
\sqrt{
\frac{
\mathrm{KL}^{\mathsf{str}}_{E,H}(Q_\lambda)+
\log(2\sqrt m|\Lambda|/\delta)
}{2(m-1)}
}
\right].
\]

For continuous coefficient domains, the admissible implementation is a
finite \(\varepsilon\)-net or certified optimizer whose transcript gives
a finite cover and a stability error. The cover size and stability error
are charged in the saturated budget.

\end{lemma}

\begin{proof}\phantomsection\label{proof-principled-lambda-selection}

Apply the structural PAC--Bayes theorem to each fixed
\(\lambda\in\Lambda\) with confidence \(\delta/|\Lambda|\), then use the
union bound. The resulting event holds for all candidate posteriors at
once. On that event, minimizing the right-hand side yields the smallest
upper bound among the declared candidates. The same labeled sample may
be used to choose \(\lambda^*\) because the factor \(|\Lambda|\)
accounts for this finite model-selection step. For a continuous
domain, replacing it by a net adds
\(\log\mathcal N(\Lambda,\varepsilon)\) and the certified Lipschitz
stability error of the objective. Hence the coefficient-selection cost
is explicit in the bound.

\end{proof}

\begin{corollary}[Two computable coefficient-selection bounds]\label{cor-two-computable-coefficient-certificates}

Two complementary forms apply to different estimators.

\textbf{PAC-Bayes posterior form.} For bounded losses and a posterior
\(Q_\lambda\), use the theorem's square-root bound with
\(\mathrm{KL}^{\mathsf{str}}_{E,H}\). This gives the
stochastic-posterior bound.

\textbf{Deterministic effective-resistance form.} For a deterministic
coupled model, suppose an independently established excess-risk theorem
gives, simultaneously over the declared finite coefficient set, the
following right-hand side. Then use

\[
\operatorname{Cert}_{\mathrm{ER}}(\theta_\lambda)
=
\widehat R_S(\theta_\lambda)
+c_{\mathrm{loss}}\,\mathcal R_W(E_q(\theta_\lambda))
+C\sqrt{\frac{N^{\Abs}_{\mathrm{eff}}(\theta_\lambda)+N^{\mathsf{harm}}_{\mathrm{eff}}+\log(|\Lambda|/\delta)}{m}},
\]

where \(c_{\mathrm{loss}}\) is the declared Lipschitz/Poincare
conversion constant for the loss geometry. The selected coefficient
minimizes this displayed bound over \(\Lambda\). This form retains the
exact pseudoinverse penalty \(\mathcal R_W(E_q)\). Replacing it by
\(\|E_q\|^2/\lambda_2\) gives the weaker global-gap bound.

\end{corollary}

\begin{proof}\phantomsection\label{proof-two-computable-coefficient-certificates}

The PAC--Bayes form is \cref{thm-principled-lambda-selection}. The deterministic conclusion
is conditional on the excess-risk inequality stated in the corollary;
it is not a consequence of the PAC--Bayes theorem. Once that inequality
has been proved for each fixed coefficient vector, a union bound over
the finite set \(\Lambda\) supplies the displayed logarithmic term and
permits minimization on the simultaneous event. Since
\(\mathcal R_W(E_q)\) is
computed by solving \(\mathcal L_W z=E_q\) on the orthogonal complement
of the harmonic kernel, it is computable for finite represented graphs
and is the smallest graph-energy penalty compatible with the defect
current. The global-gap expression follows by replacing every active
eigenvalue by \(\lambda_2\), hence it can only increase the penalty
unless the defect is worst-mode aligned.

\end{proof}

\subsection{Active querying from resistance certificates}
\label{subsec:resistance-active-querying}

Submodularity turns the represented resistance certificate into a greedy
querying guarantee.

\begin{lemma}[Greedy querying under a submodular resistance certificate]\label{thm-structural-active-learning}

For defect/current \(E_q\) and domain Laplacian
\(\mathcal L_{\mathcal D}\), the value of querying a point \(x\) is the
certified reduction

\[
\Delta(x)
=
\mathcal C(S)-\mathcal C(S\cup\{x\})
=
\langle E_q,\mathcal L_{\mathcal D,S}^{+}E_q\rangle
-
\langle E_q,\mathcal L_{\mathcal D,S\cup\{x\}}^{+}E_q\rangle.
\]

Assume, for the declared grounding model and admissible query family,
that \(F(S)=\mathcal C(\varnothing)-\mathcal C(S)\) is nonnegative,
monotone, and submodular. Then greedy querying achieves the standard
\(1-1/e\) approximation to the optimal batch of the same size, with
stopping criterion \(\max_x\Delta(x)\le c_{\mathrm{query}}\).

\end{lemma}

\begin{proof}\phantomsection\label{proof-structural-active-learning}

Grounding a queried point adds a boundary condition to the domain
Laplacian and lowers the effective resistance seen by the defect
current. Rayleigh monotonicity gives monotonicity for the standard
grounding order. Submodularity is an explicit additional hypothesis:
it is not implied by Rayleigh monotonicity for an arbitrary query
operation or current functional. Under that hypothesis, the
Nemhauser--Wolsey greedy theorem gives the batch approximation
\citep{NemhauserWolseyFisher1978}. The stopping rule is the value-of-information inequality
comparing the best marginal certificate drop with query cost.

\end{proof}

\subsection{Intrinsic harmonic dimension and saturation}
\label{subsec:intrinsic-harmonic-dimension}

The intrinsic harmonic dimension identifies the point at which a fixed
geometry exhausts its represented target projection.

\begin{definition}[Intrinsic harmonic dimension, ceiling, and coverage gap]\label{def-intrinsic-harmonic-dimension}

Fix a represented unlabeled data geometry \(\mathcal L_{\mathcal D}\),
self-adjoint and nonnegative on the empirical or population task space,
with orthonormal eigenbasis \((v_k)\). Let \(y_T\) be the centered
target contrast or centered regression field, normalized only when
explicitly stated. For a declared structural threshold \(\tau\) above
the residual/noise floor, define

\[
\mathcal V^*
=
\operatorname{span}\{v_k: |\langle v_k,y_T\rangle|^2/\|y_T\|^2\ge\tau\},
\qquad
 d^*=\dim\mathcal V^*,
\qquad
\mathcal C=\frac{\|P_{\mathcal V^*}y_T\|^2}{\|y_T\|^2}.
\]

The number \(d^*\) is the intrinsic harmonic dimension of the target
relative to the declared geometry, and \(\mathcal C\) is the structural
information ceiling. For a represented model/readout subspace
\(\mathcal R\), define the certified structural extraction and coverage
gap by

\[
\mathcal E_{\mathrm{sig}}(\mathcal R)
=
\frac{\|P_{\mathcal R}P_{\mathcal V^*}y_T\|^2}{\|y_T\|^2},
\qquad
\operatorname{Gap}(\mathcal R)
=
\mathcal C-\mathcal E_{\mathrm{sig}}(\mathcal R)
=
\frac{\|(I-P_{\mathcal R})P_{\mathcal V^*}y_T\|^2}{\|y_T\|^2}.
\]

The equality \(\operatorname{Gap}(\mathcal R)=0\) is equivalent only to
target-vector coverage,
\[
P_{\mathcal R}P_{\mathcal V^*}y_T=P_{\mathcal V^*}y_T.
\]
Full structural-subspace coverage is the stronger identity
\(P_{\mathcal R}P_{\mathcal V^*}=P_{\mathcal V^*}\), and it implies
zero gap. These quantities are computable
from a finite represented operator by eigensolve, projection, and a
declared threshold or spectral gap.

\end{definition}

\begin{lemma}[Projection bound for a fixed geometry]\label{thm-structural-ceiling-no-creation}

For the declared geometry \(\mathcal L_{\mathcal D}\) and threshold
\(\tau\), every budget-admissible learner has certified structural
extraction at most \(\mathcal C\). Target projection outside
\(\mathcal V^*\) is classified as residual or noise mass by this
criterion.

\end{lemma}

\begin{proof}\phantomsection\label{proof-structural-ceiling-no-creation}

The target decomposes orthogonally as
\(y_T=P_{\mathcal V^*}y_T+(I-P_{\mathcal V^*})y_T\). By definition, the
first term is the whole target-bearing component above the declared
structural threshold, and the second term lies below that threshold. \PartII and
\PartIII classify target arrival outside the represented structural
component as residual accounting, oracle/advice information, a changed
presentation, or a new saturated derivation. Therefore the mass counted
by this geometry's structural criterion is
\(P_{\mathcal V^*}y_T\). For any represented readout subspace
\(\mathcal R\), projection cannot increase norm, so
\(\mathcal E_{\mathrm{sig}}(\mathcal R)\le\|P_{\mathcal V^*}y_T\|^2/\|y_T\|^2=\mathcal C\).
Raw projection onto \((I-P_{\mathcal V^*})y_T\) may improve empirical
fit, but by construction it is below the structural threshold and is
charged as residual/noise unless a different geometry or source
certificate reclassifies it.

\end{proof}

\begin{lemma}[Saturation at the intrinsic harmonic dimension]\label{thm-saturation-optimality}

Let \(\mathcal R_K\) be a represented model/readout subspace with
\(\dim\mathcal R_K\le K\). The certified gap is exactly

\[
\mathcal C-\mathcal E_{\mathrm{sig}}(\mathcal R_K)
=
\frac{\|(I-P_{\mathcal R_K})P_{\mathcal V^*}y_T\|^2}{\|y_T\|^2}.
\]

Consequently, if \(K\ge d^*\) and \(\mathcal R_K\) covers
\(\mathcal V^*\), then
\(\mathcal E_{\mathrm{sig}}(\mathcal R_K)=\mathcal C\). By
\cref{thm-structural-ceiling-no-creation}, no budget-admissible learner for the same
declared geometry has larger certified structural extraction. Thus a
covering model of capacity at least \(d^*\) is optimal within the
declared geometry. If \(K>d^*\), any additional raw fit beyond
\(\mathcal C\) lies in the residual component and is an
overfitting/noise term for this certificate.

\end{lemma}

\begin{proof}\phantomsection\label{proof-saturation-optimality}

The displayed identity is the Pythagorean projection identity applied to
the structural target component \(P_{\mathcal V^*}y_T\). If
\(P_{\mathcal R_K}\mathcal V^*=\mathcal V^*\), then
\(P_{\mathcal R_K}P_{\mathcal V^*}y_T=P_{\mathcal V^*}y_T\), so
\(\mathcal E_{\mathrm{sig}}=\mathcal C\).
\Cref{thm-structural-ceiling-no-creation} proves
that \(\mathcal C\) upper-bounds every budget-admissible learner's
certified structural extraction for the same geometry, so equality is
global optimality within that declared presentation. Capacity below
\(d^*\) cannot cover a \(d^*\)-dimensional structural subspace. Capacity
above \(d^*\) can only add directions orthogonal to \(\mathcal V^*\);
their target overlap is residual by definition and enters the
PAC-Bayes/effective-resistance certificate as capacity or residual cost,
not as new certified signal.

\end{proof}

\begin{corollary}[Coverage criterion for model-family optimality]\label{cor-optimal-up-to-model-choice}

Fix a metric family and a represented learning rule whose reachable
readout spans include a subspace \(\mathcal R\) of dimension \(d^*\)
with \(P_{\mathcal R}P_{\mathcal V^*}=P_{\mathcal V^*}\).
Then choosing capacity \(K^*=d^*\), training or solving the terminal
reach problem inside that family, and verifying this projection identity
proves optimality within the
declared geometry: no
budget-admissible learner for the same declared geometry and task
presentation has a better certified structural extraction. If the gap
is positive, its value is the exact structural shortfall and quantifies
the additional coverage required from the model family.

\end{corollary}

\begin{proof}\phantomsection\label{proof-optimal-up-to-model-choice}

The saturation theorem shows that a covering subspace at capacity
\(d^*\) reaches \(\mathcal C\). The ceiling theorem shows that
\(\mathcal C\) cannot be exceeded by any budget-admissible learner in
the declared geometry. The displayed projection identity verifies full
subspace coverage and hence zero gap; conversely, zero gap alone
verifies only coverage of the single target vector. Otherwise the
Pythagorean identity gives the exact missing target mass. Thus, once the model family is chosen rich enough to
contain a covering subspace, the remaining optimization problem has a
certificate-verifiable optimum; the irreducible choice is the
model/metric family that makes coverage possible.

\end{proof}

\section{Certification of Learned Geometries}
\label{sec:certification-learned-geometries}

We now separate label-independent geometry certification from the subsequent
propagation of its bound to a selected model family.

\begin{lemma}[Validity of a label-independent geometry certificate]\label{thm-valid-learned-geometry}

Let \(\{\mathcal L_{\mathcal D}(\theta):\theta\in\Theta\}\) be a
finite-coverable metric family computed from unlabeled data. Learning
\(\theta\) by maximizing the label criterion \(\mathcal E(\theta;y)\)
does not produce a label-independent geometry certificate or a
label-independent PAC--Bayes prior: sufficient metric/readout capacity
can place arbitrary random labels in the leading modes. Such a
procedure may still be analyzed as label-dependent learned-operator
selection, with its selection complexity and generalization error
charged. To obtain the label-independent certificate used here, select
the learned geometry through a label-independent
self-consistency criterion, for example

\[
\theta^*
\in
\operatorname*{argmin}_{\theta\in\Theta}
\left[
\mathcal R_{\theta}(E_{\mathrm{self}}(\theta))
+
C\sqrt{\frac{N_{\mathrm{eff}}^{\mathrm{metric}}(\theta)+\log\mathcal N(\Theta,\rho)+\log(1/\delta)}{m_u}}
\right],
\]

where \(m_u\) is the unlabeled sample size,

\[
E_{\mathrm{self}}(\theta)=L_{\mathrm{data}}(\theta)U_\theta-U_\theta L_{Q,\theta},
\qquad
\mathcal R_{\theta}(E)=\langle E,\mathcal L_{\mathcal D}(\theta)^+E\rangle.
\]

The labels are used only after \(\theta^*\) and the geometry capacity
are fixed by unlabeled data. The subsequent label-dependent reach,
intrinsic harmonic dimension, and PAC--Bayes prior under \(\theta^*\)
are then based on a label-independent geometry selection.

\end{lemma}

\begin{proof}\phantomsection\label{proof-valid-learned-geometry}

If \(\theta\) is optimized against \(\mathcal E(\theta;y)\), the metric
family can use the labels to reorder the spectrum. At capacity
approaching the sample size, the label vector lies in the readout span
for any labeling, including random labels, so the certificate becomes
circular if presented as label-independent; it must instead be charged
as learned-operator selection. In contrast, the displayed self-defect objective depends only
on unlabeled samples and the metric family. The finite cover term
charges data-dependent selection over \(\Theta\). Since \(\theta^*\) is
independent of the labels, PAC-Bayes priors and spectral subspaces
computed from \(\mathcal L_{\mathcal D}(\theta^*)\) satisfy the required
prior-independence condition. Capping model capacity by the intrinsic
dimension read from the unlabeled geometry prevents the label-dependent
readout from expanding until it certifies arbitrary labels. Thus the
geometry is selected independently of the labels, and the subsequent
label reach is evaluated within that fixed geometry.

\end{proof}

\begin{lemma}[Two-level geometry certification and propagation to the model
certificate]\label{thm-two-level-geometry-certification}

Split or resample the unlabeled data used to learn geometry. Let
\(\epsilon_{\mathrm{metric}}\) satisfy the explicit high-probability
operator estimate
\[
\|\mathcal L_{\mathcal D}^{\mathrm{emp}}(\theta^*)
-\mathcal L_{\mathcal D}^{\mathrm{pop}}(\theta^*)\|_{\mathrm{op}}
\le\epsilon_{\mathrm{metric}},
\]
obtained from held-out unlabeled data or another declared concentration
theorem. If the structural eigengap
separating \(\mathcal V^*\) from the residual spectrum is \(g_*>0\), assume
\(2\epsilon_{\mathrm{metric}}<g_*\) and define the empirical cluster through
the corresponding isolated spectral window. Then the population and
empirical structural projectors differ by

\[
\|P_{\mathcal V^*}^{\mathrm{pop}}-P_{\mathcal V^*}^{\mathrm{emp}}\|_{\mathrm{op}}
\le
\frac{2\,\epsilon_{\mathrm{metric}}}{g_*}.
\]
For every model-certificate coordinate \(Q\) used after geometry selection,
the complete record must also name a represented constant \(L_Q<\infty\)
such that
\[
|Q(P)-Q(P')|\le L_Q\|P-P'\|_{\mathrm{op}}
\]
on the admitted projector class. Such a coordinate satisfies
\[
|Q(P_{\mathcal V^*}^{\mathrm{pop}})
 -Q(P_{\mathcal V^*}^{\mathrm{emp}})|
\le \frac{2L_Q\epsilon_{\mathrm{metric}}}{g_*}.
\]
For target energy \(Q(P)=\|Py_T\|_2^2\), one may take
\(L_Q=\|y_T\|_2^2\).

\end{lemma}

\begin{proof}\phantomsection\label{proof-two-level-geometry-certification}

The spectral-window hypothesis permits Davis--Kahan projector perturbation
to be applied to the same isolated cluster, giving
\(\|P_{\mathrm{pop}}-P_{\mathrm{emp}}\|_{\mathrm{op}}\le 2\|\Delta\mathcal L\|_{\mathrm{op}}/g_*\).
The declared Lipschitz inequality for \(Q\), followed by substitution of
this projector bound, proves the coordinatewise transfer inequality. For
target energy, orthogonality gives
\(\|Py_T\|_2^2=\langle y_T,Py_T\rangle\); hence
\(|Q(P)-Q(P')|\le\|y_T\|_2^2\|P-P'\|_{\mathrm{op}}\).

\end{proof}

\begin{lemma}[Finite-net termination within a declared geometry family]\label{thm-final-residue-model-family}

Fix a represented metric family \(\Theta\), a precision \(\rho>0\), and
a budget \(B\). If the budget-admissible part of \(\Theta\) has a finite
\(\rho\)-net, then exhaustive certificate evaluation over that net
terminates after at most \(\mathcal N(\Theta,\rho)\) evaluations. If the
certificate objective is \(L_\Theta\)-Lipschitz in the covering metric,
the best net value differs from the infimum over the covered family by
at most \(L_\Theta\rho\).

This is a termination statement only for the fixed represented family
and precision. It neither proves that recursive enlargement of metric
families terminates nor that the chosen family contains an intrinsic
population geometry.

\end{lemma}

\begin{proof}\phantomsection\label{proof-final-residue-model-family}

By hypothesis the net is finite, so evaluating its members requires at
most its cardinality many evaluations. For any \(\theta\) in the
covered family choose a net point \(\theta_0\) at distance at most
\(\rho\). Lipschitz continuity gives
\(|F(\theta)-F(\theta_0)|\le L_\Theta\rho\). Taking infima proves the
approximation statement. No assertion about families outside the
declared cover follows from this argument.

\end{proof}

\begin{definition}[Certified geometry evaluation protocol]
\label{cor-complete-methodology-optimality}
\label{def:certified-geometry-evaluation-protocol}

Combining
\cref{thm-structural-pac-bayes,thm-principled-lambda-selection,thm-structural-active-learning,thm-valid-learned-geometry,thm-final-residue-model-family}
yields the following procedure.

\begin{enumerate}

\item

Declare a finite-coverable metric/model family \(\Theta\).

\item

Learn \(\theta^*\) from unlabeled data by minimizing
self-defect/effective-resistance plus metric capacity, independently of
the labels.

\item

Certify metric generalization on held-out unlabeled data and carry the
perturbation term forward.

\item

Compute \(\mathcal V^*\), \(d^*\), and \(\mathcal C\) under the
certified geometry.

\item

Choose a model/readout family of capacity \(K=d^*\) or the smallest
capacity whose represented span covers \(\mathcal V^*\).

\item

Verify the projection identity
\(P_{\mathcal R}P_{\mathcal V^*}=P_{\mathcal V^*}\) when full
subspace coverage is claimed, and compute the target-vector gap
\(\mathcal C-\mathcal E_{\mathrm{sig}}\). A zero target-vector gap,
up to metric and sampling errors, proves saturation for this target; a
positive gap quantifies its missing structural mass.

\item

Select regularization coefficients using the structural PAC--Bayes or
deterministic effective-resistance bound above.

\end{enumerate}

The protocol reports full coverage only when every displayed certificate is
finite on one simultaneous event and the projection identity is verified.
Otherwise it reports the measured target-vector gap and does not infer which
unverified condition caused it.

\end{definition}

\begin{theorem}[Conditional completeness of the geometry evaluation protocol]
\label{thm:conditional-geometry-protocol-completeness}

Assume the protocol of
\cref{def:certified-geometry-evaluation-protocol} returns, on one event
\(\mathcal E\), a valid represented geometry, its intrinsic target subspace
\(\mathcal V^*\), and a represented readout space \(\mathcal R\) satisfying
\[
P_{\mathcal R}P_{\mathcal V^*}=P_{\mathcal V^*}.
\]
Let \(y\) be the target vector and define the geometry-accessible target
energy by
\(\mathcal C=\|P_{\mathcal V^*}y\|_2^2\).  Then the readout captures all of
that energy:
\[
\|P_{\mathcal R}P_{\mathcal V^*}y\|_2^2=\mathcal C.
\]
If the certified metric, projector, and empirical target-energy errors sum to
\(\varepsilon_{\rm cert}\) in the comparison norm named by the record, assume
the record also supplies a represented constant \(L_{\rm cert}<\infty\) for
which the reported gap functional is \(L_{\rm cert}\)-Lipschitz in that norm
on the admitted certificate class. The reported target-vector gap is then at
most
\[
L_{\rm cert}\varepsilon_{\rm cert}.
\]
Generalization and regularization
claims remain exactly those supplied by the separately invoked PAC--Bayes,
resistance, and model-family certificates; no claim is made for geometries or
readouts outside the declared family.

\end{theorem}

\begin{proof}

Apply the verified projection identity to \(y\):
\(P_{\mathcal R}P_{\mathcal V^*}y=P_{\mathcal V^*}y\), and take squared
norms. The exact target-vector gap is zero. The triangle inequality bounds
the distance between the exact and reported certificate tuples by
\(\varepsilon_{\rm cert}\); the declared Lipschitz inequality gives the
displayed bound. The final sentence is a
scope statement: the cited learning theorems apply on \(\mathcal E\) with
their own hypotheses and additive errors.

\end{proof}

\Cref{fig-learned-geometry-certification-workflow} records the independence
and simultaneity requirements of the certified geometry protocol.

\begin{figure}[tbp]
\centering
\resizebox{.96\linewidth}{!}{%
\begin{tikzpicture}[fw diagram,node distance=8mm and 12mm]
  \node[fw source,text width=2.7cm] (unlab) {unlabeled geometry\\sample};
  \node[fw provenance,text width=2.8cm,right=of unlab] (geom)
    {represented geometry\\and projector certificate};
  \node[fw source,text width=2.7cm,below=13mm of unlab] (lab)
    {labeled training\\sample};
  \node[fw use,text width=2.8cm,right=of lab] (read)
    {readout or posterior\\selection};
  \node[fw source,text width=2.7cm,below=13mm of lab] (hold)
    {held-out\\certification\\sample};
  \node[fw provenance,text width=2.8cm,right=of hold] (shift)
    {metric, projector, and\\deployment checks};
  \node[fw box,text width=3.0cm,right=18mm of read] (event)
    {one simultaneous\\certificate event};
  \node[fw terminal,text width=2.8cm,right=of event] (bound)
    {coverage identity and\\risk bound};
  \draw[fw arrow] (unlab) -- (geom);
  \draw[fw arrow] (lab) -- (read);
  \draw[fw arrow] (hold) -- (shift);
  \draw[fw arrow] (geom) -- (event);
  \draw[fw arrow] (read) -- (event);
  \draw[fw arrow] (shift) -- (event);
  \draw[fw arrow] (event) -- (bound);
\end{tikzpicture}%
}
\caption{A learned geometry may enter a population-risk certificate only
through an independence split or a uniform theorem covering its selection.
Geometry construction, labeled readout fitting, and held-out propagation
checks join only on one simultaneous event before coverage or deployment is
claimed.}
\label{fig-learned-geometry-certification-workflow}
\end{figure}

\subsection{Defect-Geometry Inequality
Profiles}\label{defect-geometry-inequality-profile-certificates}
\label{subsec:defect-geometry-inequality-profiles}

The following constants belong to the represented defect geometry, not
to an ambient smoothness class. The learned routing form
\(\mathcal L_D\) acts on the lumpability defect \(E_q\). Effective
resistance determines whether the signed defect admits finite-energy
routing. The full profile additionally measures entropy dissipation,
gradient contraction, target capacity, and stability under distribution
shift.

A certified defect profile is

\[
\Pi_D=
\bigl(\mathcal L_D,E_q,\rho_D,
\lambda_D,\alpha_{\rm MLSI}^D,\kappa_D,
\operatorname{Cap}_D,\Phi_D,\alpha_T^D\bigr),
\]

with every entry audited on the same active defect support and the same
quotient defect used to form
\(\mathcal R_D(E_q)=\langle E_q,\mathcal L_D^+E_q\rangle\). The
constants are lower or upper certified bounds in the direction used by
the theorem: lower bounds for rates such as \(\alpha_{\rm MLSI}^D\),
\(\kappa_D\), transport constants, and target capacities; upper bounds
for residual energies, entropy defects, and boundary-risk terms.

\begin{lemma}[Selection among defect-profile bounds]\label{thm-defect-profile-selection}

Fix a represented quotient, a sample, and a label-independent or
independently evaluated admissible defect geometry.  For each component
\(s\in S\), let \(C_s\ge0\) denote the single true contribution in the risk
units of the surrounding theorem.  Suppose a finite audited family
\(\mathcal J_s\) and nonnegative random bounds
\(B_{s,j}=\widehat C_{s,j}+r_{s,j}\) satisfy, on one event
\(\mathcal E_s\) of probability at least \(1-\delta_s\),
\[
C_s\le B_{s,j}
\qquad\text{simultaneously for every }j\in\mathcal J_s.
\]
Then, on \(\mathcal E=\bigcap_{s\in S}\mathcal E_s\),

\[
\mathcal C_s^{\rm prof}(\Pi_D)
=
\min_{j\in \mathcal J_s}B_{s,j}
\quad\text{satisfies}\quad
C_s\le\mathcal C_s^{\rm prof}(\Pi_D),
\qquad
\Pr(\mathcal E)\ge1-\sum_{s\in S}\delta_s.
\]

The radii \(r_{s,j}\) include every finite-family, sample-split, or uniform
selection penalty required to establish the simultaneous event.  They may
be zero for deterministic analytic bounds.  Candidates for different
components are never minimized against one another.

\end{lemma}

\begin{proof}

On \(\mathcal E_s\), the same scalar \(C_s\) is no larger than every
\(B_{s,j}\), hence it is no larger than their minimum.  The probability
claim is the union bound over \(s\).

\end{proof}

\begin{remark}[Typed defect-profile candidates]
\label{rem:typed-defect-profile-candidates}

The routing component always contains the Poincar\'e/effective-resistance bound

\[
\mathcal C_{{\rm route},{\rm Poincare}}(\Pi_D)=c_R\mathcal R_D(E_q),
\]

with \(c_R\) the conversion constant of the surrounding risk theorem.
The remaining inequalities apply only to the components they bound:

\begin{itemize}

\item

\textbf{MLSI/entropy.} With the declared normalization, a defect MLSI
such as
\(\operatorname{Ent}^{\mathcal B}_{\mu_D}(\rho_D)\le c_{\rm MLSI}\alpha_D^{-1}\mathcal E_D(\rho_D,\log \rho_D)\)
supplies a concentration or entropy-production bound for the actual
residual density \citep{Gross1975,BobkovTetali2006}.

\item

\textbf{Bakry--Emery.} A curvature certificate
\(\Gamma_{2,D}(f)\ge \kappa_D\Gamma_D(f)\), with \(\kappa_D>0\),
supplies a propagation/contraction bound, for example
\(\Gamma_D(P_tf)\le e^{-2\kappa_Dt}P_t\Gamma_D(f)\) under the declared
reversible semigroup calculus, or a static curvature scale after
multiplication by the theorem's conversion constant
\citep{BakryEmery1985}.

\item

\textbf{Capacity/Cheeger.} Target access, boundary access, and death
avoidance are capacity components. Target access may use inverse target
capacity; death avoidance uses polar excision when
\(\operatorname{Cap}_D(\Sigma^-)=0\), or a monotone
boundary-risk/death-capacity upper charge otherwise.

\item

\textbf{Transport entropy.} Deployment, replay, or off-policy shift is a
distribution-shift component, for example
\(\operatorname{Lip}_D(\ell)\sqrt{\mathrm{KL}(\nu\Vert\mu)/\alpha_T^D}\)
\citep{BobkovGotze1999,OttoVillani2000}.

\end{itemize}

Logically implied inequalities are counted once. If a Bakry--\'Emery
condition implies the MLSI or transport constant in the declared
setting, the induced constant is substituted into the corresponding
component.

\end{remark}

\begin{corollary}[Substitution of simultaneous defect-profile certificates]
\label{cor:defect-profile-certificate-substitution}

Suppose a represented risk theorem has already established, on the event
\(\mathcal E\) of \cref{thm-defect-profile-selection},

\[
R(h)
\le
\widehat R_S(h)
+
\sum_{s\in S}C_s
+E_{\rm op}+E_{\rm src},
\]

where all terms are in the same risk units.  Then, on that event,
\[
R(h)
\le
\widehat R_S(h)
+\sum_{s\in S}\mathcal C_s^{\rm prof}(\Pi_D)
+E_{\rm op}+E_{\rm src}.
\]
The routing coordinate occurs once in the index set \(S\); any candidate
that already contains its resistance contribution is treated as one complete
bound \(B_{s,j}\), rather than added to it again.

\end{corollary}

\begin{proof}

Substitute \(C_s\le\mathcal C_s^{\rm prof}(\Pi_D)\) from
\cref{thm-defect-profile-selection} term by term.

\end{proof}

Training losses alone do not establish the stated profile bounds. If the
profile constants are selected using the labeled risk-evaluation data,
the theorem requires an independent sample split, a label-independent
selection rule, or a uniform concentration result over the declared
finite family. Analytically established constants require only the
finite-family selection term.

The Bakry--\'Emery condition provides a contraction estimate for the
defect-induced Caus gradient and therefore describes active residual
dissipation. An MLSI term remains distinct in optimization because a
minibatch curvature loss evaluates only a finite function class, whereas
the inequality is universal over the admissible domain. The training
objective may therefore contain both terms:

\[
\mathcal L_{\rm geom}
=
\lambda_R\mathcal R_D(E_q)
+\lambda_{\Gamma_2}\mathcal L_{\Gamma_2}
+\lambda_{\rm MLSI}\mathcal L_{\rm MLSI}
+\mathcal L_{\rm norm}.
\]

Here \(\mathcal L_{\Gamma_2}\) targets defect-gradient contraction and
\(\mathcal L_{\rm MLSI}\) is a direct entropy witness for the actual
residual density. The theorem uses the constants established after
training rather than the inclusion of these terms in the objective.

Persistent large \(E_q\) after admissible resistance, MLSI, and
curvature optimization indicates that the quotient has identified states
that must remain distinct. In this case the admissible modification is
to refine \(q\), rather than to absorb the defect into
\(\mathcal L_D\).

\section{Learning Under Symmetric Label Corruption}
\label{learning-under-symmetric-label-corruption}
\label{sec:learning-symmetric-label-corruption}

The preceding Rademacher, effective-dimension, and PAC--Bayes bounds admit a
common noise correction for binary labels.  Symmetric label corruption
attenuates the clean prediction advantage by the factor \(1-2\eta\).  The
statistical capacity and represented model-selection costs remain those
already defined in this part.

\Cref{fig-label-noise-learning-flow} summarizes this transport from corrupted
sample risk to clean prediction risk.

\begin{figure}[!ht]
\centering
\begin{tikzpicture}[fw diagram,node distance=7mm and 9mm]
  \node[fw source,minimum width=2.5cm] (clean) {clean labels\\\(Y=f_*(X)\)};
  \node[fw residual,right=of clean,minimum width=2.2cm] (flip)
    {flip channel\\\(\eta\)};
  \node[fw box,right=of flip,minimum width=2.7cm] (sample)
    {corrupted sample\\\((X_i,\widetilde Y_i)\)};
  \node[fw use,below=of sample,minimum width=2.7cm] (learn)
    {represented learner};
  \node[fw terminal,left=of learn,minimum width=2.8cm] (denoise)
    {clean-risk bound\\\(/(1-2\eta)\)};
  \draw[fw arrow] (clean) -- (flip);
  \draw[fw arrow] (flip) -- (sample);
  \draw[fw arrow] (sample) -- (learn);
  \draw[fw arrow] (learn) -- (denoise);
\end{tikzpicture}
\caption{Symmetric corruption attenuates clean-label advantage by
\(1-2\eta\).  The represented learner pays the existing Rademacher,
effective-dimension, PAC--Bayes, and operator-selection terms; dividing the
resulting corrupted-risk certificate by the attenuation factor gives the
clean-risk bound.}
\label{fig-label-noise-learning-flow}
\end{figure}

\subsection{Risk transport through a binary flip channel}
\label{subsec:binary-flip-risk-transport}

The following definition covers homogeneous symmetric noise and its
instance-dependent Massart extension.

\begin{definition}[Binary label-corruption channel]
\label{def-binary-label-corruption-channel}

Let \(X\sim\mu\), let the clean binary target be
\(Y=f_*(X)\in\{-1,1\}\), and let

\[
\widetilde Y=Y\Xi,
\qquad
\Pr\{\Xi=-1\mid X=x\}=\eta(x).
\]

Assume

\[
0\le\eta(x)\le\bar\eta<\frac12,
\qquad
\bar\rho:=1-2\bar\eta>0.
\]

For a binary predictor \(h:X\to\{-1,1\}\), write

\[
R_0(h)=\Pr\{h(X)\ne Y\},
\qquad
R_{\mathrm{corr}}(h)=\Pr\{h(X)\ne\widetilde Y\}.
\]

In the homogeneous case \(\eta(x)\equiv\eta\), put
\(\rho=1-2\eta\) and write \(R_\eta\) for the corrupted-label risk.

More generally, a bounded loss
\(\ell:[-1,1]\times\{-1,1\}\to[0,1]\) is
\emph{flip-symmetric} when

\[
\ell(z,1)+\ell(z,-1)=1
\qquad(-1\le z\le1).
\]

Its clean and corrupted risks are denoted by
\(R_0^\ell(f)\) and \(R_{\mathrm{corr}}^\ell(f)\).  The binary
zero--one loss and the clipped linear loss
\(\ell(z,y)=(1-yz)/2\) satisfy this identity.

\end{definition}

\begin{lemma}[Exact transport of clean risk through label corruption]
\label{lem-exact-label-noise-risk-transport}

Under \cref{def-binary-label-corruption-channel}, every binary predictor
satisfies

\begin{equation}
\label{eq-massart-risk-transport}
R_{\mathrm{corr}}(h)-R_{\mathrm{corr}}(f_*)
=
\mathbb E\!\left[
 (1-2\eta(X))\mathbf 1_{\{h(X)\ne f_*(X)\}}
\right]
\ge
\bar\rho R_0(h).
\end{equation}

For homogeneous noise, the identity becomes

\begin{equation}
\label{eq-homogeneous-label-noise-identity}
R_\eta(h)=\eta+\rho R_0(h).
\end{equation}

For every flip-symmetric loss,

\begin{equation}
\label{eq-symmetric-loss-noise-identity}
R_{\mathrm{corr}}^\ell(f)
=
\mathbb E\!\left[
 \eta(X)+(1-2\eta(X))
 \ell(f(X),Y)
\right].
\end{equation}

In particular, homogeneous corruption gives

\[
R_\eta^\ell(f)=\eta+\rho R_0^\ell(f).
\]

\end{lemma}

\begin{proof}

Condition on \(X=x\).  If \(h(x)=f_*(x)\), the corrupted label
disagrees with \(h(x)\) with probability \(\eta(x)\).  If
\(h(x)\ne f_*(x)\), it disagrees with \(h(x)\) with probability
\(1-\eta(x)\).  Subtracting the conditional risk of \(f_*\) gives

\[
(1-2\eta(x))\mathbf 1_{\{h(x)\ne f_*(x)\}}.
\]

Integration proves \eqref{eq-massart-risk-transport}; the lower bound follows
from \(1-2\eta(x)\ge\bar\rho\).  Constant \(\eta(x)=\eta\) gives
\eqref{eq-homogeneous-label-noise-identity}.  For a flip-symmetric loss,

\begin{align*}
\mathbb E[\ell(f(X),\widetilde Y)\mid X,Y]
&=(1-\eta(X))\ell(f(X),Y)
  +\eta(X)\ell(f(X),-Y)\\
&=\eta(X)+(1-2\eta(X))\ell(f(X),Y),
\end{align*}

which proves \eqref{eq-symmetric-loss-noise-identity}.

\end{proof}

\subsection{Rademacher capacity and the tolerable-noise region}
\label{subsec:rademacher-noise-tolerance}

The exact transport identity converts the existing generalization radius into
a clean-label denoising guarantee.

\begin{theorem}[Rademacher bound under label corruption]
\label{thm-rademacher-label-corruption-bound}

Let \(\mathcal H\) be a represented class of binary predictors and let
\(S=((X_i,\widetilde Y_i))_{i=1}^m\) be an i.i.d.\ corrupted sample from
\cref{def-binary-label-corruption-channel}, where \(m\ge1\) and
\(0<\delta<1\).  Let

\[
\mathcal L_{\mathcal H}
=
\left\{
(x,\widetilde y)\longmapsto
\mathbf 1_{\{h(x)\ne\widetilde y\}}:
h\in\mathcal H
\right\}
\]

and define

\begin{equation}
\label{eq-corrupted-rademacher-radius}
\Gamma_m^{\mathrm{Rad}}(\mathcal H,\delta)
:=
2\widehat{\mathfrak R}_S(\mathcal L_{\mathcal H})
+3\sqrt{\frac{\log(2/\delta)}{2m}}.
\end{equation}

With probability at least \(1-\delta\), simultaneously for every
\(h\in\mathcal H\),

\begin{equation}
\label{eq-corrupted-risk-uniform-deviation}
\left|
R_{\mathrm{corr}}(h)-\widehat R_{\mathrm{corr},S}(h)
\right|
\le
\Gamma_m^{\mathrm{Rad}}(\mathcal H,\delta).
\end{equation}

Assume \(f_*\in\mathcal H\), and let \(\widehat h\) be an
\(\varepsilon_{\mathrm{opt}}\)-empirical minimizer:

\[
\widehat R_{\mathrm{corr},S}(\widehat h)
\le
\inf_{h\in\mathcal H}
\widehat R_{\mathrm{corr},S}(h)
+\varepsilon_{\mathrm{opt}}.
\]

On the same event,

\begin{equation}
\label{eq-massart-rademacher-clean-risk}
R_0(\widehat h)
\le
\frac{
2\Gamma_m^{\mathrm{Rad}}(\mathcal H,\delta)
+\varepsilon_{\mathrm{opt}}
}{\bar\rho}.
\end{equation}

Under homogeneous noise, every fixed \(h\in\mathcal H\) also satisfies

\begin{equation}
\label{eq-pointwise-noise-corrected-rademacher}
R_0(h)
\le
\left[
\frac{
\widehat R_{\eta,S}(h)-\eta
+\Gamma_m^{\mathrm{Rad}}(\mathcal H,\delta)
}{\rho}
\right]_{[0,1]},
\end{equation}

where \([u]_{[0,1]}=\min\{1,\max\{0,u\}\}\).

For a target-qualified use, apply these inequalities to each exact source
cell of \(\mathcal H\).  Source selection is charged by
\cref{eq-source-resolved-learning-rademacher-union}, and every corrupted
reveal retains the coefficient
\(\vartheta(\eta)=(1-2\eta)^2\) in
\cref{eq-source-resolved-learning-information-ledger}.  The inverse factor
\(\rho^{-1}\) in clean-risk correction and the forward information
contraction \(\vartheta(\eta)\) remain separate coordinates.

\end{theorem}

\begin{proof}

The standard two-sided empirical Rademacher inequality for a
\([0,1]\)-valued loss class gives
\eqref{eq-corrupted-risk-uniform-deviation}; it is the two-sided form of the
bound used in
\cref{thm-fixed-operator-rademacher-bound,cor-rademacher-audit-ceiling}.
On this event,

\begin{align*}
R_{\mathrm{corr}}(\widehat h)
&\le
\widehat R_{\mathrm{corr},S}(\widehat h)
+\Gamma_m^{\mathrm{Rad}}\\
&\le
\widehat R_{\mathrm{corr},S}(f_*)
+\varepsilon_{\mathrm{opt}}
+\Gamma_m^{\mathrm{Rad}}\\
&\le
R_{\mathrm{corr}}(f_*)
+\varepsilon_{\mathrm{opt}}
+2\Gamma_m^{\mathrm{Rad}}.
\end{align*}

Apply \eqref{eq-massart-risk-transport} and divide by \(\bar\rho\) to
obtain \eqref{eq-massart-rademacher-clean-risk}.  In the homogeneous case,
combine the upper half of \eqref{eq-corrupted-risk-uniform-deviation} with
\eqref{eq-homogeneous-label-noise-identity}, then clip the resulting risk
bound to \([0,1]\).

\end{proof}

\begin{corollary}[Capacity-controlled noise tolerance]
\label{cor-capacity-controlled-noise-tolerance}

In the setting of \cref{thm-rademacher-label-corruption-bound}, a sufficient
condition for clean risk at most \(\varepsilon>0\) is

\begin{equation}
\label{eq-noise-tolerance-from-capacity}
2\Gamma_m^{\mathrm{Rad}}(\mathcal H,\delta)
+\varepsilon_{\mathrm{opt}}
\le
(1-2\bar\eta)\varepsilon.
\end{equation}

Equivalently,

\[
\bar\eta
\le
\frac12-
\frac{
\Gamma_m^{\mathrm{Rad}}(\mathcal H,\delta)
+\varepsilon_{\mathrm{opt}}/2
}{\varepsilon}.
\]

The same conclusion holds for a represented score class, a flip-symmetric
loss, and an \(\varepsilon_{\mathrm{opt}}\)-empirical minimizer whenever the
class contains a comparator \(f_*\) with
\(R_0^\ell(f_*)=0\).  For a fixed structural energy ball
\(\mathcal F_A(R)\), an \(L_\ell\)-Lipschitz loss, and clipped scores in
\([-1,1]\), the contraction principle and
\cref{lem-ellipsoid-rademacher} then permit the explicit radius

\begin{equation}
\label{eq-fixed-operator-label-noise-radius}
\Gamma_{A,m}(R,\delta)
=
\frac{2L_\ell R}{m}
\sqrt{\operatorname{Tr}(A^{-1})}
+3\sqrt{\frac{\log(2/\delta)}{2m}}.
\end{equation}

For a learned operator class, the expression in
\cref{thm-learned-operator-union-bound} replaces the fixed-operator
Rademacher term.  Thus denoising without sample memorization occurs whenever
the fixed or selected operator radius and
\(\varepsilon_{\mathrm{opt}}\) are
\(o(1-2\bar\eta)\); operator selection retains the covering cost already
present in that theorem.

\end{corollary}

\begin{proof}

The first assertion is \eqref{eq-massart-rademacher-clean-risk} with its
right-hand side bounded by \(\varepsilon\).  Solving the resulting inequality
for \(\bar\eta\) gives the second display.  If
\(R_0^\ell(f_*)=0\), \eqref{eq-symmetric-loss-noise-identity} gives
\(R_{\mathrm{corr}}^\ell(f_*)=\mathbb E\eta(X)\) and

\[
R_{\mathrm{corr}}^\ell(f)-R_{\mathrm{corr}}^\ell(f_*)
\ge
\bar\rho R_0^\ell(f).
\]

The same empirical-minimizer argument therefore proves the loss version.  For
an \(L_\ell\)-Lipschitz loss, Rademacher contraction bounds the loss-class
complexity by \(L_\ell\) times the score-class complexity.
\Cref{lem-ellipsoid-rademacher} supplies
\(Rm^{-1}\sqrt{\operatorname{Tr}(A^{-1})}\).  The learned-operator
substitution follows identically from
\cref{thm-learned-operator-union-bound}.

\end{proof}

\subsection{Structural PAC--Bayes correction}
\label{subsec:structural-pac-bayes-noise}

The same attenuation factor converts the structural PAC--Bayes certificate
into a clean-risk certificate.

\begin{theorem}[PAC--Bayes bound under symmetric label noise]
\label{thm-structural-pac-bayes-label-noise}

Assume homogeneous corruption at a known rate \(0\le\eta<1/2\), a bounded
flip-symmetric loss, a sample size \(m\ge2\), a confidence level
\(0<\delta<1\), and a represented prior \(P\) fixed independently of the
labeled sample.  With probability at least \(1-\delta\), simultaneously for
every posterior \(Q\ll P\),

\begin{equation}
\label{eq-pac-bayes-label-noise-clean-risk}
R_0^\ell(Q)
\le
\left[
\frac{
\widehat R_{\eta,S}^\ell(Q)-\eta
+
\sqrt{
\dfrac{
\mathrm{KL}(Q\Vert P)
+\log(2\sqrt m/\delta)
}{2(m-1)}
}
}{1-2\eta}
\right]_{[0,1]}.
\end{equation}

For the structural prior \(P=P^{\mathsf{str}}_{E,H}\) of
\cref{def-structural-pac-bayes-kl}, the divergence in
\eqref{eq-pac-bayes-label-noise-clean-risk} is exactly
\(\mathrm{KL}^{\mathsf{str}}_{E,H}(Q)\).  Thus the noise correction preserves
the defect-restricted structural complexity term of
\cref{thm-structural-pac-bayes}.

For a declared finite coefficient family \(\Lambda\), the bound holds
simultaneously after replacing the logarithmic term by
\(\log(2\sqrt m|\Lambda|/\delta)\), exactly as in
\cref{thm-principled-lambda-selection}.  Minimizing this simultaneous
right-hand side gives the certified noise-corrected coefficient choice.

\end{theorem}

\begin{proof}

The McAllester inequality used in the proof of
\cref{thm-structural-pac-bayes} applies to any represented prior fixed
independently of the labeled sample; choosing
\(P=P^{\mathsf{str}}_{E,H}\) recovers that theorem's structural
specialization.  The prior-independence requirement is unchanged because the
flip channel modifies the labels rather than the public or independently
learned geometry.  The homogeneous flip-symmetric identity in
\cref{lem-exact-label-noise-risk-transport} gives
\(R_\eta^\ell(Q)=\eta+(1-2\eta)R_0^\ell(Q)\).  Subtract \(\eta\), divide
by \(1-2\eta\), and clip to \([0,1]\).  The simultaneous coefficient form
is the union-bound argument of \cref{thm-principled-lambda-selection}.

\end{proof}

\subsection{Effective dimension after label debiasing}
\label{subsec:effective-dimension-label-debiasing}

For squared-loss prediction, symmetric flips scale the regression function
without changing its source subspace.

\begin{corollary}[Debiased kernel ridge bound for binary labels]
\label{cor-debiased-krr-label-noise}

Let \(f_*:X\to\{-1,1\}\), let
\(\widetilde Y=f_*(X)\Xi\) under homogeneous noise, and put
\(\rho=1-2\eta\).  Then

\[
\mathbb E[\widetilde Y\mid X]=\rho f_*(X).
\]

Let \(\widehat g_\lambda\) be the kernel-ridge estimator trained on
\((X_i,\widetilde Y_i)_{i=1}^m\), and define the debiased predictor
\(\widehat f_\lambda=\widehat g_\lambda/\rho\).  Suppose a fixed-kernel
oracle certificate \eqref{eq-fixed-kernel-oracle-certificate} holds on an
event \(\mathcal E_{\eta,k,\lambda}\) for the regression target
\(g_*=\rho f_*\), with constants \(C_{\rm var},C_{\rm bias}\) and with
\(v_\eta^2\) a valid conditional sub-Gaussian variance proxy for
\(\widetilde Y-\rho f_*(X)\).  Then, on the same event,

\begin{equation}
\label{eq-debiased-krr-label-noise}
\|\widehat f_\lambda-f_*\|_{L^2(\mu)}^2
\le
C_{\rm var}\frac{v_\eta^2N_{\mathrm{eff}}(\lambda)}
     {\rho^2m}
+C_{\rm bias}B_k(\lambda;f_*).
\end{equation}

Since the centered flip variable is bounded, one may always take a universal
bounded sub-Gaussian proxy; its exact conditional variance is
\(1-\rho^2=4\eta(1-\eta)\).  Moreover,

\[
\Pr\{\operatorname{sign}(\widehat f_\lambda(X))\ne f_*(X)\}
\le
\|\widehat f_\lambda-f_*\|_{L^2(\mu)}^2.
\]

Consequently a sufficient effective-dimension condition for vanishing
denoising error is

\[
C_{\rm var}\frac{v_\eta^2N_{\mathrm{eff}}(\lambda)}
     {(1-2\eta)^2m}
+C_{\rm bias}B_k(\lambda;f_*)
=o(1).
\]

\end{corollary}

\begin{proof}

Conditional expectation gives
\(\mathbb E[\widetilde Y\mid X]=f_*(X)\mathbb E\Xi=\rho f_*(X)\).
Apply the named fixed-kernel certificate to \(g_*=\rho f_*\):

\[
\|\widehat g_\lambda-\rho f_*\|_2^2
\le
C_{\rm var}\frac{v_\eta^2N_{\mathrm{eff}}(\lambda)}m
+C_{\rm bias}B_k(\lambda;\rho f_*).
\]

The ridge source residual is quadratic, so
\(B_k(\lambda;\rho f_*)=\rho^2B_k(\lambda;f_*)\).  Divide the inequality
by \(\rho^2\) to obtain \eqref{eq-debiased-krr-label-noise}.  If the sign
of \(\widehat f_\lambda\) is incorrect, its distance from the binary value
\(f_*(X)\) is at least one.  The classification inequality follows by
Markov's inequality applied pointwise to the squared error.

\end{proof}

\section{Saturated Statistical Attacks and Population Alignment}
\label{subsec-saturated-statistical-attacks}
\label{sec:saturated-statistical-attacks-alignment}

The statistical bounds above apply simultaneously to the full represented
budget slice.  The following definition fixes that class before the training
sample is drawn.

\begin{definition}[Budget-saturated statistical readout class]
\label{def-budget-saturated-statistical-readout-class}

Fix a presented binary prediction family, a sample length \(m\), and a budget
\(B\).  Let \(\mathcal H_{n,B,m}^{\mathrm{sat,stat}}\) be the
sample-independent union of all scores

\[
h_{\rho,S,\omega}:X_I\longrightarrow[-1,1]
\]

obtained as follows.  The record \(\rho\) is one family-uniform derivation in
\(\mathsf{Der}_{\mathfrak I}\); \(S\) ranges over every legal encoded
training sample of length \(m\); \(\omega\) ranges over every represented
internal random seed; and the complete cost of training, retained state,
output representation, and score evaluation is at most \(B\).  Every score
retains its complete derivation record.  The union over legal \(S\) and
\(\omega\) is taken before the realized sample is selected, so the class is a
fixed statistical class.  A data-dependent kernel, feature map, posterior,
regularization choice, or terminal postprocessing belongs to this class only
through its complete represented construction.

Let \(T_I=\{x:f_*(x)=1\}\).  Define the saturated learned projection envelope

\begin{equation}
\label{eq-saturated-learned-projection-envelope}
M_{I,n}^{\mathrm{sat,stat}}(B,m)
=
\sup_{h\in\mathcal H_{n,B,m}^{\mathrm{sat,stat}}}
\operatorname{ProjMass}_{\sigma(h)}(T_I),
\end{equation}

with supremum zero when the class is empty.  The algebra \(\sigma(h)\) is
generated by the represented score itself.  It is distinct from the ambient
measurable envelope of all public data.

For the clipped linear loss

\[
\ell_h(x,\widetilde y)=\frac{1-\widetilde y h(x)}2,
\]

let \(\Gamma_{n,B,m,\delta}^{\mathrm{sat,stat}}\) denote the radius in
\eqref{eq-corrupted-rademacher-radius} for the loss class generated by
\(\mathcal H_{n,B,m}^{\mathrm{sat,stat}}\).  A learned-operator cover may
instead be inserted through \cref{thm-learned-operator-union-bound}.

\end{definition}

\begin{theorem}[Uniform accountability of budget-saturated statistical
attacks]
\label{thm-saturated-statistical-attack-alignment}

Assume homogeneous symmetric label corruption at rate
\(0\le\eta<1/2\), and put \(\rho=1-2\eta\).  For
\(h:X_I\to[-1,1]\), write

\[
c_0(h)=\mathbb E[h(X)f_*(X)],
\qquad
\widehat c_{\eta,S}(h)
=\frac1m\sum_{i=1}^m h(X_i)\widetilde Y_i.
\]

With probability at least \(1-\delta\), simultaneously for every
family-uniform represented statistical computation of complete cost at most
\(B\), every represented random seed, and its output
\(h\in\mathcal H_{n,B,m}^{\mathrm{sat,stat}}\),

\begin{equation}
\label{eq-saturated-statistical-correlation-transport}
\left|
\widehat c_{\eta,S}(h)-\rho c_0(h)
\right|
\le
2\Gamma_{n,B,m,\delta}^{\mathrm{sat,stat}}.
\end{equation}

Consequently,

\begin{equation}
\label{eq-saturated-statistical-clean-alignment-bound}
|c_0(h)|
\le
\frac{
|\widehat c_{\eta,S}(h)|
+2\Gamma_{n,B,m,\delta}^{\mathrm{sat,stat}}
}{1-2\eta}.
\end{equation}

If the clean binary target is balanced, then every such output also satisfies

\begin{equation}
\label{eq-saturated-statistical-projection-bound}
|c_0(h)|^2
\le
\operatorname{Var}_{\mu_I}(h)\,
\operatorname{ProjMass}_{\sigma(h)}(T_I)
\le
M_{I,n}^{\mathrm{sat,stat}}(B,m).
\end{equation}

Hence, on the same event,

\begin{equation}
\label{eq-universal-saturated-statistical-attack-bound}
\sup_{\mathcal A,\omega}
|c_0(h_{\mathcal A,S,\omega})|
\le
\min\!\left\{
\sqrt{M_{I,n}^{\mathrm{sat,stat}}(B,m)},
\frac{
\displaystyle
\sup_{h\in\mathcal H_{n,B,m}^{\mathrm{sat,stat}}}
|\widehat c_{\eta,S}(h)|
+2\Gamma_{n,B,m,\delta}^{\mathrm{sat,stat}}
}{1-2\eta}
\right\}.
\end{equation}

For sign-valued outputs in the balanced setting, the same projection envelope
gives the universal approximation floor

\begin{equation}
\label{eq-universal-saturated-statistical-risk-floor}
\inf_{\mathcal A,\omega}
R_0(h_{\mathcal A,S,\omega})
\ge
\frac{1-\sqrt{M_{I,n}^{\mathrm{sat,stat}}(B,m)}}2.
\end{equation}

Every represented use of such an output as a target selector, ranking,
potential, quotient readout, or transition rule retains the same derivation
record and adds its represented use cost.  The quantities
\(c_0(h)\) and \(\operatorname{ProjMass}_{\sigma(h)}(T_I)\) are analytic
population audits of that affordable score.  A downstream use enters the
\(\Geom/\Abs\) prospective source ledger only when the same complete pre-hit record
also satisfies the represented channel predicate, target-discrimination
comparison, temporal admissibility, and cost gates of
\cref{thm-paper4-dependency-principle,thm-full-saturated-profile-decomposition}.
Otherwise the audit remains a statistical-accounting coordinate and target
arrival is controlled by the prospective-selector ledger.  Statistical
estimation introduces no additional structural channel.

\end{theorem}

\begin{proof}

The sample-independent hull in
\cref{def-budget-saturated-statistical-readout-class} contains the output of
every represented budget-\(B\) learner for every realized sample and every
represented seed.  Apply the simultaneous score-class form of
\cref{thm-rademacher-label-corruption-bound,cor-capacity-controlled-noise-tolerance}
to the clipped linear losses.  Since

\[
R_\eta^\ell(h)
=\frac{1-\rho c_0(h)}2,
\qquad
\widehat R_{\eta,S}^\ell(h)
=\frac{1-\widehat c_{\eta,S}(h)}2,
\]

the uniform risk deviation is exactly
\eqref{eq-saturated-statistical-correlation-transport}.  Division by
\(\rho>0\) gives
\eqref{eq-saturated-statistical-clean-alignment-bound}.  The event is uniform
over the whole output class, so data-dependent selection of the learner,
operator, score, or postprocessing preserves the inequality.

For the projection statement, balance gives
\(\mu_I(T_I)=1/2\) and
\(y_{T_I}=f_*/2\).  If \(h\) is nonconstant, the normalized correlation in
\cref{thm-generated-algebra-projection-ceiling}, applied to
\(\mathfrak A_I=\sigma(h)\), is

\[
\frac{
|\langle h-\mathbb Eh,y_{T_I}\rangle|^2
}{
\|h-\mathbb Eh\|_2^2\|y_{T_I}\|_2^2
}
=
\frac{|c_0(h)|^2}{\operatorname{Var}_{\mu_I}(h)}.
\]

Thus the first inequality in
\eqref{eq-saturated-statistical-projection-bound} follows from that ceiling.
It also holds for constant \(h\), because balance gives \(c_0(h)=0\).
The score range implies \(\operatorname{Var}(h)\le1\), and taking the
saturated supremum gives the second inequality.  Combining the projection
and statistical bounds and taking the supremum over computations and seeds
proves \eqref{eq-universal-saturated-statistical-attack-bound}.
For a sign-valued output,
\(R_0(h)=(1-c_0(h))/2\); applying
\(|c_0(h)|\le\sqrt{M_{I,n}^{\mathrm{sat,stat}}(B,m)}\) proves
\eqref{eq-universal-saturated-statistical-risk-floor}.

Finally, \cref{thm-paper4-dependency-principle} places the learner and every
represented downstream use in the saturated derivability closure.
\Cref{thm-full-saturated-profile-decomposition} assigns every target-bearing
structural use to the existing source ledger at its complete cost.

\end{proof}

\begin{corollary}[Target alignment forces saturated learned structure]
\label{cor-statistical-alignment-forces-saturated-structure}

In the balanced setting of
\cref{thm-saturated-statistical-attack-alignment}, if a budget-\(B\)
statistical computation produces an output with
\(|c_0(h)|>\varepsilon\), then

\begin{equation}
\label{eq-statistical-alignment-structural-witness}
M_{I,n}^{\mathrm{sat,stat}}(B,m)>\varepsilon^2.
\end{equation}

On the simultaneous event of that theorem it also satisfies

\begin{equation}
\label{eq-statistical-alignment-empirical-witness}
|\widehat c_{\eta,S}(h)|
>
(1-2\eta)\varepsilon
-2\Gamma_{n,B,m,\delta}^{\mathrm{sat,stat}}.
\end{equation}

Thus a clean alignment of size \(\varepsilon\) requires a represented learned
score whose generated algebra has projection mass greater than
\(\varepsilon^2\), while its statistical visibility above the uniform
selection radius is attenuated by \(1-2\eta\).

\end{corollary}

\begin{proof}

The first conclusion is the contrapositive of
\eqref{eq-saturated-statistical-projection-bound}.  Rearranging
\eqref{eq-saturated-statistical-correlation-transport} proves the second.

\end{proof}

\begin{theorem}[Two-level provenance and uniform noise transport for learned
readouts]
\label{thm-two-level-learned-readout-provenance}

Adopt the balanced setting and notation of
\cref{def-budget-saturated-statistical-readout-class,thm-saturated-statistical-attack-alignment},
and put \(\rho=1-2\eta\).  For every realized sample \(S\) and every
represented output \(h\), define the scalar transport residual

\begin{equation}
\label{eq-learned-readout-transport-residual}
r_{\eta,S}(h)
:=
\widehat c_{\eta,S}(h)-\rho c_0(h).
\end{equation}

With probability at least \(1-\delta\), simultaneously over the complete
budget-saturated statistical class,

\begin{align}
\label{eq-two-level-transport-residual-bound}
|r_{\eta,S}(h)|
&\le 2\Gamma_{n,B,m,\delta}^{\mathrm{sat,stat}},\\
\label{eq-two-level-population-alignment-bound}
|c_0(h)|
&\le \sqrt{M_{I,n}^{\mathrm{sat,stat}}(B,m)},\\
\label{eq-two-level-empirical-alignment-bound}
|\widehat c_{\eta,S}(h)|
&\le
\rho\sqrt{M_{I,n}^{\mathrm{sat,stat}}(B,m)}
+2\Gamma_{n,B,m,\delta}^{\mathrm{sat,stat}}.
\end{align}

Conversely, an output satisfying \(|c_0(h)|>\varepsilon\) witnesses

\begin{equation}
\label{eq-two-level-learned-structure-witness}
M_{I,n}^{\mathrm{sat,stat}}(B,m)>\varepsilon^2,
\qquad
|\widehat c_{\eta,S}(h)|
>
\rho\varepsilon-2\Gamma_{n,B,m,\delta}^{\mathrm{sat,stat}}.
\end{equation}

These estimates have the following provenance interpretation.

\begin{enumerate}[label=\textup{(\roman*)}]
\item A family-uniform learner is one represented derivation rule.  Its
output may depend on the presented instance, the realized sample, and a
represented random seed; family uniformity does not require those outputs to
be identical.

\item Every retained sample table, empirical prediction vector, posterior,
optimization transcript, and memorized lookup record is a represented finite
coordinate.  Its construction, storage, and evaluation costs remain in the
lifted computation record.  Its population projection is an analytic audit.
It is charged to the existing \(\Geom/\Abs\) prospective source ledger only
when the represented use satisfies the same-record channel predicate, target
comparison, temporal, and cost gates; otherwise it remains statistical
accounting.  The orthogonal remainder is assigned to \(\Res\) by the
exhaustive residual normal form.  Thus memorization creates neither an
uncharged population certificate nor a sixth structural source.

\item If one family-uniform affordable rule has nonnegligible population
alignment, then \eqref{eq-two-level-learned-structure-witness} records its
nonnegligible analytic target projection.  Any represented use of that
alignment as a selector, ranking, potential, quotient readout, or transition
rule retains the learner's full ancestry.  It is classified as a prospective
source only after its relevant represented channel predicate and pre-hit
comparison are verified.  If it induces positive selector advantage, that
advantage is bounded by the existing five contextual-charge envelopes through
\cref{thm-prospective-selector-structural-charge}; conversion of those
contextual charges to source values additionally uses
\cref{cor-contextual-structural-charge-transfer}.

\item A target-bearing table supplied separately for each instance without
one family-uniform represented derivation is nonuniform advice, not an output
of \(\mathcal H_{n,B,m}^{\mathrm{sat,stat}}\).  An over-budget construction is
likewise absent from the budget slice.
\end{enumerate}

Hence family-uniform learning may discover affordable structure already
derivable from the presentation, but sample-dependent fitting does not create
an additional provenance channel.

\end{theorem}

\begin{proof}

Equation \eqref{eq-two-level-transport-residual-bound} is
\eqref{eq-saturated-statistical-correlation-transport} after the abbreviation
\eqref{eq-learned-readout-transport-residual}.  The projection estimate
\eqref{eq-saturated-statistical-projection-bound} gives
\eqref{eq-two-level-population-alignment-bound}.  The triangle inequality
then gives

\[
|\widehat c_{\eta,S}(h)|
\le
\rho|c_0(h)|+|r_{\eta,S}(h)|,
\]

which proves \eqref{eq-two-level-empirical-alignment-bound}.  The first
conclusion in \eqref{eq-two-level-learned-structure-witness} is the
contrapositive of \eqref{eq-saturated-statistical-projection-bound}; the
second follows by applying the reverse triangle inequality to
\eqref{eq-saturated-statistical-correlation-transport}.

It remains to verify the provenance statements.  By
\cref{def-budget-saturated-statistical-readout-class}, the union over legal
samples and represented seeds is formed from one family-uniform derivation
record and includes the complete cost of training, retained state, output
representation, and evaluation.  Consequently every empirical table and
memorized value used by the learner is a charged coordinate on its finite
lifted carrier.  Apply
\cref{thm-paper4-dependency-principle,thm-full-saturated-profile-decomposition}
to every represented population use that satisfies their channel predicates
and temporal comparisons.  Its qualified structural projection is assigned
to the existing channels, and
\cref{thm-exhaustive-residual-normal-form} assigns the remaining orthogonal
component to \(\Res\).  If that remainder acquires a family-uniform affordable
target-aligned channel representative, it is reclassified by
\cref{prop-family-uniform-residual-reclassification,thm-no-uniform-residual-correlation};
therefore it cannot remain simultaneously residual and an affordable
structural source.

Finally, a downstream selector is a represented use of the learned score and
retains the derivation as an ancestor.  Its positive advantage is bounded by
\cref{thm-prospective-selector-structural-charge}.  The family-uniformity and
budget exclusions are precisely the membership requirements in
\cref{def-budget-saturated-statistical-readout-class}.  This proves all four
provenance assertions.

\end{proof}

\Cref{fig-learned-readout-provenance-split} depicts the population-structural
and empirical-residual decomposition of a learned readout.

\begin{figure}[tbp]
\centering
\begin{tikzpicture}[fw diagram]
  \node[fw source,minimum width=3.2cm] (learn) at (-4.7,0)
    {complete learned readout\\training, state, evaluation};
  \node[fw geom,minimum width=3.1cm] (pop) at (-0.6,1.1)
    {population-aligned\\projection};
  \node[fw residual,minimum width=3.1cm] (emp) at (-0.6,-1.1)
    {population-orthogonal\\empirical remainder};
  \node[fw provenance,minimum width=3.0cm] (ledger) at (3.6,1.1)
    {existing \(\Geom/\Abs\)\\source cell};
  \node[fw residual,minimum width=3.0cm] (res) at (3.6,-1.1)
    {\(\Res\) memorization\\coordinate};
  \draw[fw arrow] (learn) -- (pop);
  \draw[fw arrow] (learn) -- (emp);
  \draw[fw arrow] (pop) -- (ledger);
  \draw[fw arrow] (emp) -- (res);
  \draw[fw dashed arrow] (res) to[bend right=28]
    node[right,font=\scriptsize,align=left] {target-aligned reuse\\is reclassified} (ledger);
\end{tikzpicture}
\caption{Learned-readout provenance from
\cref{thm-two-level-learned-readout-provenance}. Population-aligned use keeps
its complete structural ancestry; the orthogonal sample-dependent component
is memorization in \(\Res\). An affordable target-aligned reuse returns to the
existing source ledger.}
\label{fig-learned-readout-provenance-split}
\end{figure}

\begin{corollary}[Residual-noise-order specialization]
\label{cor-learned-readout-residual-noise-order}

In the setting of
\cref{thm-two-level-learned-readout-provenance}, suppose the certified readout
class has residual-noise order at least \(q\) in the sense of
\cref{def-extraction-audit-ceilings}.  Then, on the same simultaneous event,
every \(h\) in that certified class satisfies

\begin{equation}
\label{eq-learned-readout-residual-order-bounds}
|c_0(h)|\le(1-2\eta)^q,
\qquad
|\widehat c_{\eta,S}(h)|
\le(1-2\eta)^{q+1}
+2\Gamma_{n,B,m,\delta}^{\mathrm{sat,stat}}.
\end{equation}

\end{corollary}

\begin{proof}

The residual-noise certificate and
\cref{thm-budgeted-extraction-bound} give
\(\operatorname{Ext}_I(h;T_I)\le\rho^{2q}\).  The normalization calculation
in the proof of \cref{thm-saturated-statistical-attack-alignment} gives

\[
|c_0(h)|^2
\le
\operatorname{Var}_{\mu_I}(h)\operatorname{Ext}_I(h;T_I).
\]

Since \(h\in[-1,1]\), its variance is at most one.  Taking square roots proves
the first inequality in \eqref{eq-learned-readout-residual-order-bounds}.
Substitution into
\(|\widehat c_{\eta,S}(h)|\le\rho|c_0(h)|+2\Gamma\) proves the second.

\end{proof}

\subsection{Noise-indexed population alignment and learning accountability}
\label{subsec-noise-indexed-population-alignment}
\label{subsec:noise-indexed-population-alignment}

The sample-independent hull in
\cref{def-budget-saturated-statistical-readout-class} is the correct class for
a simultaneous concentration event.  Computational recovery, however, is
evaluated under the realized sample and the algorithm's represented
randomness.  The following profiles retain that quantifier order.

\begin{definition}[Average-seed population alignment and generalizing
projection profiles]
\label{def-average-seed-population-alignment-profile}

Fix a balanced presented binary prediction family, sample length \(m\),
budget \(B\), and a monotone positively homogeneous family functional
\(\mathfrak M_n\).  Let the supremum below range over one family-uniform
represented learner \(\mathcal A\) whose complete training, retained-state,
output, fresh-input evaluation, and any invoked decoder cost is at most \(B\).
For instance \(I\),
let \(S\) have the declared training-sample law, let \(\omega\) have the
learner's represented seed law, and write

\[
h_{I,S,\omega}=h_{\mathcal A,I,S,\omega},
\qquad
c_{0,I}(h)=\mathbb E_{X\sim\mu_I}[h(X)f_{*,I}(X)].
\]

Define

\begin{align}
\label{eq-average-seed-population-alignment-profile}
\operatorname{Align}_{\mathfrak M,n}^{\mathrm{sat,stat}}(B,m)
&=
\sup_{\mathcal A}
\mathfrak M_n\!\left(
 I\longmapsto
 \mathbb E_{S,\omega}|c_{0,I}(h_{I,S,\omega})|^2
\right),\\
\label{eq-average-seed-generalizing-projection-profile}
\operatorname{GenProj}_{\mathfrak M,n}^{\mathrm{sat,stat}}(B,m)
&=
\sup_{\mathcal A}
\mathfrak M_n\!\left(
 I\longmapsto
 \mathbb E_{S,\omega}\!\left[
 \operatorname{Var}_{\mu_I}(h_{I,S,\omega})
 \operatorname{ProjMass}_{\sigma(h_{I,S,\omega})}(T_I)
 \right]
\right),
\end{align}

with supremum zero when the learner class is empty.  These are evaluations of
the existing saturated statistical class, not additional structural channels.
The expectation over \((S,\omega)\) occurs inside the family functional and
before the supremum over the one uniform learner rule.  In particular, a seed
or sample on which the learner happens to output the target readout contributes
according to its probability under the declared execution law.

For homogeneous flip noise, append a subscript \(\eta\) to either profile.
For a scalarized budget sublevel \(s\), append \(\kappa\) and restrict the same
supremum to complete scalarized cost at most \(s\).

\end{definition}

\begin{theorem}[Population alignment, generalization, and memorization
separation]
\label{thm-population-alignment-generalization-memorization}

In the setting of
\cref{def-average-seed-population-alignment-profile},

\begin{equation}
\label{eq-average-seed-alignment-projection-domination}
\operatorname{Align}_{\mathfrak M,n}^{\mathrm{sat,stat}}(B,m)
\le
\operatorname{GenProj}_{\mathfrak M,n}^{\mathrm{sat,stat}}(B,m)
\le
\mathfrak M_n\!\left(
 I\longmapsto M_{I,n}^{\mathrm{sat,stat}}(B,m)
\right).
\end{equation}

On the simultaneous event of
\cref{thm-saturated-statistical-attack-alignment}, every represented learner
and every realized sample satisfy

\begin{equation}
\label{eq-realized-average-seed-alignment-certificate}
\begin{aligned}
\mathbb E_{\omega}|c_{0,I}(h_{I,S,\omega})|^2
\le
\mathbb E_{\omega}\min\!\Biggl\{
&\operatorname{Var}_{\mu_I}(h_{I,S,\omega})
 \operatorname{ProjMass}_{\sigma(h_{I,S,\omega})}(T_I),
\\[-1mm]
&
\left(
\frac{|\widehat c_{\eta,S}(h_{I,S,\omega})|
+2\Gamma_{n,B,m,\delta}^{\mathrm{sat,stat}}}
{1-2\eta}
\right)^2
\Biggr\}.
\end{aligned}
\end{equation}

Writing \(E_I\) for this simultaneous event, its probability bound and
\(|c_{0,I}(h)|\le1\) also give the execution-averaged certificate

\begin{equation}
\label{eq-execution-average-alignment-certificate}
\begin{aligned}
\mathbb E_{S,\omega}|c_{0,I}(h_{I,S,\omega})|^2
\le \delta
+\mathbb E_S\!\Bigl[\mathbf 1_{E_I}\,
\mathbb E_{\omega}\min\!\Biggl\{&
\operatorname{Var}_{\mu_I}(h_{I,S,\omega})
\operatorname{ProjMass}_{\sigma(h_{I,S,\omega})}(T_I),\\[-1mm]
&
\left(
\frac{|\widehat c_{\eta,S}(h_{I,S,\omega})|
+2\Gamma_{n,B,m,\delta}^{\mathrm{sat,stat}}}
{1-2\eta}
\right)^2
\Biggr\}\Bigr].
\end{aligned}
\end{equation}

Suppose a represented decoder \(D\) applied to the retained output has the
following target-alignment margin: whenever \(D(h_{I,S,\omega})\) is the
correct target, \(|c_{0,I}(h_{I,S,\omega})|\ge\gamma\), for some represented
\(0<\gamma\le1\).  Then its recovery probability obeys

\begin{equation}
\label{eq-decoder-recovery-from-average-alignment}
\mathfrak M_n\!\left(
 I\longmapsto
 \Pr_{S,\omega}\{D(h_{I,S,\omega})\in T_I\}
\right)
\le
\gamma^{-2}
\operatorname{Align}_{\mathfrak M,n}^{\mathrm{sat,stat}}(B,m).
\end{equation}

Every sample table, optimization transcript, empirical prediction vector, and
seed-dependent lookup record retains its complete represented storage and
evaluation cost.  Its population projection is the term measured by
\cref{eq-average-seed-generalizing-projection-profile}.  The orthogonal
population component is assigned to \(\Res\) by
\cref{thm-exhaustive-residual-normal-form}.  If a family-uniform affordable
use of that component has positive population target alignment and satisfies
the relevant represented channel predicate and pre-hit comparison, its
projection is reclassified into the existing \(\Geom/\Abs\) source ledger at
its complete cost.  Without those gates it remains an analytic population
audit.  Thus instance-specific memorization contributes no
population alignment unless the represented learner generalizes it under the
declared family law.

\end{theorem}

\begin{proof}

For every fixed \((I,S,\omega)\), the projection calculation in
\cref{thm-saturated-statistical-attack-alignment} gives

\[
|c_{0,I}(h)|^2
\le
\operatorname{Var}_{\mu_I}(h)
\operatorname{ProjMass}_{\sigma(h)}(T_I).
\]

Take expectation over \((S,\omega)\), apply the monotonicity of
\(\mathfrak M_n\), and then take the supremum over the same uniform learner
rules.  This proves the first inequality in
\eqref{eq-average-seed-alignment-projection-domination}.  Each realized output
belongs to \(\mathcal H_{n,B,m}^{\mathrm{sat,stat}}\), its variance is at most
one, and its projection mass is at most
\(M_{I,n}^{\mathrm{sat,stat}}(B,m)\).  The second inequality follows.

On the simultaneous event,
\cref{eq-saturated-statistical-clean-alignment-bound} and the pointwise
projection bound are both valid for every represented seed.  Squaring the
first, taking the minimum of the two pointwise ceilings, and averaging over
\(\omega\) proves
\eqref{eq-realized-average-seed-alignment-certificate}.
On \(E_I^c\), the score range and the binary target give
\(|c_{0,I}(h)|^2\le1\).  Splitting the sample expectation over
\(E_I\) and \(E_I^c\), and using \(\Pr(E_I^c)\le\delta\), proves
\eqref{eq-execution-average-alignment-certificate}.

The decoder-margin condition gives the pointwise inequality

\[
\mathbf 1_{\{D(h_{I,S,\omega})\in T_I\}}
\le
\gamma^{-2}|c_{0,I}(h_{I,S,\omega})|^2.
\]

Expectation, application of \(\mathfrak M_n\), and the defining supremum in
\eqref{eq-average-seed-population-alignment-profile} prove
\eqref{eq-decoder-recovery-from-average-alignment}.

Finally, the complete learner record belongs to the lifted computation carrier
by \cref{def-budget-saturated-statistical-readout-class}.  Apply
\cref{thm-full-saturated-profile-decomposition,thm-exhaustive-residual-normal-form}
to its population use.  The represented structural projection is charged to
the existing source channels, and its orthogonal complement is \(\Res\).
Reclassification of an affordable target-aligned residual representative is
\cref{prop-family-uniform-residual-reclassification}.  This proves the final
assertion.

\end{proof}

\section{Geometry-Aware Learning and Affordable Target Selection}
\label{subsec-geometry-aware-learning-selection}
\label{sec:geometry-aware-learning-selection}

The learning record and every geometric use of its output belong to one
saturated derivation. The next lemma compares the target-reach factor of a
learned potential with the generated-algebra projection already used by the
statistical ledger.

\begin{lemma}[Generated-algebra ceiling for a saturated geometric potential]
\label{lem-learned-geom-target-reach-projection}

Fix an instance \(I\), target \(T_I\), and a represented ambient record

\[
\mathfrak g=(\mathfrak N,\Phi,D,J,\mathfrak K)
\in\mathcal G^{X,\mathfrak A}_{n,\le B}
\]

from \cref{def-geom-extraction}. Define the bounded represented score

\begin{equation}
\label{eq-bounded-learned-geom-score}
h_D=\frac{1-D}{1+D}.
\end{equation}

Then \(h_D:X_I\to(-1,1]\), the maps \(D\mapsto h_D\) and
\(h_D\mapsto D=(1-h_D)/(1+h_D)\) generate the same observable algebra, and

\begin{equation}
\label{eq-learned-geom-reach-projection-ceiling}
R_{T_I}(D)
\le
\operatorname{ProjMass}_{\sigma(D)}(T_I)
=
\operatorname{ProjMass}_{\sigma(h_D)}(T_I).
\end{equation}

Consequently,

\begin{equation}
\label{eq-learned-geom-visibility-projection-ceiling}
V^X_{\Phi,D,J,\mathfrak K}
\le
\frac{
M_\Phi C_{\Phi,\mathfrak K}^{\mathrm{auth}}(D)
}{1+\rho_\Phi(D)}
\operatorname{ProjMass}_{\sigma(h_D)}(T_I)
\le
\operatorname{ProjMass}_{\sigma(h_D)}(T_I).
\end{equation}

If the presented task is a same-carrier statistical task and the complete
record \(\mathfrak g\) is generated from a length-\(m\) sample, then \(h_D\)
belongs to
\(\mathcal H_{n,\Psi(B,n),m}^{\mathrm{sat,stat}}\).  Let
\(G_{X,n,\operatorname{learn}(m)}^{\mathrm{src,sat}}(B)\) denote the
restriction of the ambient source profile to precisely those complete
same-carrier geometric records generated from a length-\(m\) sample under
this presentation.  Then

\begin{equation}
\label{eq-full-saturated-geom-learned-projection-bound}
G_{X,n,\operatorname{learn}(m)}^{\mathrm{src,sat}}(B)
\le
\sup_{I\in\mathfrak I_n}
M_{I,n}^{\mathrm{sat,stat}}\bigl(\Psi(B,n),m\bigr).
\end{equation}

No conclusion for primitive, nonlearned, or differently sampled records in
the unrestricted profile \(G_{X,n}^{\mathrm{src,sat}}(B)\) follows from this
statistical-hull inclusion.

For a distinguished instance sequence or a declared monotone family
functional, the same conclusion holds with the corresponding evaluation on
both sides, as in
\cref{def-master-saturated-quotient-geometry-profile}.

\end{lemma}

\begin{proof}

Every function of \(D\), and therefore every member of the closed monotone
span \(\mathcal M(D)\), is \(\sigma(D)\)-measurable. Thus

\[
\mathcal M(D)\subseteq L^2(X_I,\sigma(D),\mu_I).
\]

Orthogonal projection onto a closed subspace cannot have larger norm than
orthogonal projection onto a containing closed subspace. Applying this to
the centered target contrast \(y_{T_I}\) gives

\[
\|\Pi_{\mathcal M(D)}y_{T_I}\|_2
\le
\|P_{\sigma(D)}y_{T_I}\|_2.
\]

Division by \(\|y_{T_I}\|_2\) proves the first inequality in
\eqref{eq-learned-geom-reach-projection-ceiling}. The transform in
\eqref{eq-bounded-learned-geom-score} is one-to-one on \([0,\infty)\), with
the displayed inverse on its image. Hence
\(\sigma(h_D)=\sigma(D)\).  The forward transform is a finite represented
postprocessing, and its complete evaluation cost is included in
\(\Psi(B,n)\). This proves the equality in
\eqref{eq-learned-geom-reach-projection-ceiling}.

Insert that ceiling into the visibility definition in
\cref{def-geom-extraction}. Since \(M_\Phi\),
\(C_{\Phi,\mathfrak K}^{\mathrm{auth}}(D)\), and
\((1+\rho_\Phi(D))^{-1}\) lie in \([0,1]\), both inequalities in
\eqref{eq-learned-geom-visibility-projection-ceiling} follow.

For a same-carrier statistical presentation, the derivation producing
\(\mathfrak g\) already contains its sample, retained state, represented
random seed, potential, and evaluation rule. Appending the bounded transform
therefore produces one member of
\(\mathcal H_{n,\Psi(B,n),m}^{\mathrm{sat,stat}}\) by
\cref{def-budget-saturated-statistical-readout-class,thm-paper4-dependency-principle}.
Its projection mass is bounded by the defining supremum
\(M_{I,n}^{\mathrm{sat,stat}}(\Psi(B,n),m)\). Taking the saturated supremum
over the same family-uniform geometric derivations proves
\eqref{eq-full-saturated-geom-learned-projection-bound}. The other evaluation
conventions follow by pointwise application followed by the declared
monotone functional.

\end{proof}

\begin{corollary}[Qualified learned-current accountability]
\label{cor-qualified-learned-current-accountability}

Adopt the same-carrier statistical setting of
\cref{lem-learned-geom-target-reach-projection}.  Let
\(\mathfrak g=(\mathfrak N,\Phi,D,J,\mathfrak K)\) be one complete
represented learned geometric record of saturated cost at most \(B\), let
\(h_D=(1-D)/(1+D)\), and let \(S_I\) be a represented edge sector supporting
\(\nabla_\Phi D\).  For its qualification indicator \(Q_{I,S,\omega}\),
\begin{equation}
\label{eq-qualified-learned-current-pointwise-ceiling}
Q_{I,S,\omega}V^X_{\Phi,D,J,\mathfrak K}
\le
Q_{I,S,\omega}
\min\!\left\{
\operatorname{ProjMass}_{\sigma(h_D)}(T_I),
\frac{\|\mathbf1_{S_I}J_{\mathfrak K}\|_E^2}
     {\|J_{\mathfrak K}\|_E^2}
\right\},
\end{equation}
where the current ratio is zero when \(J_{\mathfrak K}=0\).

Suppose qualification entails correctness of a represented decoder whose
derivation and evaluation are included in the same cost-\(B\) record, and
that decoder has the target-alignment margin \(\gamma\) of
\cref{thm-population-alignment-generalization-memorization}.  Then every
monotone sublinear positively homogeneous family functional satisfies
\begin{align}
\label{eq-qualified-learned-current-family-ceiling}
&\mathfrak M_n\!\left(
I\longmapsto
\mathbb E_{S,\omega}
\bigl[Q_{I,S,\omega}V^X_{\Phi,D,J,\mathfrak K}\bigr]
\right)\nonumber\\
&\qquad\le
\gamma^{-2}
\operatorname{Align}_{\mathfrak M,n}^{\mathrm{sat,stat}}(B,m)
\le
\gamma^{-2}
\operatorname{GenProj}_{\mathfrak M,n}^{\mathrm{sat,stat}}(B,m)\\
&\qquad\le
\gamma^{-2}
\mathfrak M_n\!\left(
I\longmapsto M_{I,n}^{\mathrm{sat,stat}}(B,m)
\right).\nonumber
\end{align}
On the simultaneous noise-transport event, the same learned record also
satisfies the realized empirical certificate
\cref{eq-realized-average-seed-alignment-certificate}, including the
factor \(1-2\eta\) and the complete saturated selection radius.

Thus the admitted current, the potential generating its gradient component,
the learned target projection, and the decoder qualification probability are
bounded in one cost-preserving record and then evaluated through the single
saturated family supremum.

\end{corollary}

\begin{proof}

\Cref{eq-learned-geom-visibility-projection-ceiling} bounds the visibility by
\(\operatorname{ProjMass}_{\sigma(h_D)}(T_I)\), while
\cref{eq-represented-current-supported-alignment-bound} bounds it by the
current-energy fraction on \(S_I\).  Multiplication by the qualification
indicator proves
\cref{eq-qualified-learned-current-pointwise-ceiling}.

Qualification and the decoder hypothesis give
\[
Q_{I,S,\omega}
\le
\mathbf1_{\{D_{\mathrm{dec}}(h_{I,S,\omega})\in T_I\}}.
\]
Since visibility lies in \([0,1]\),
\cref{eq-current-qualification-family-bound} bounds the left side of
\cref{eq-qualified-learned-current-family-ceiling} by the family-functional
decoder success probability.  Apply
\cref{eq-decoder-recovery-from-average-alignment} and then
\cref{eq-average-seed-alignment-projection-domination}.  The realized
noise-transport statement is
\cref{eq-realized-average-seed-alignment-certificate}.

\end{proof}

The measurable projection ceiling becomes statistically testable through a
represented readout. The following proposition applies the existing
actionable-mass and noise-transport bounds to every such readout at once.

\begin{proposition}[Noise-corrected actionable extraction for learned
geometries]
\label{prop-noise-corrected-actionable-geom-extraction}

Adopt the balanced homogeneous-noise setting of
\cref{thm-saturated-statistical-attack-alignment}, put
\(\rho=1-2\eta>0\), and fix a budget \(B\). Let \(g:X_I\to[-1,1]\) be any
represented target-discrimination, monotone-potential, ranking, or quotient
readout produced by a budget-\(B\) learned geometric record. Include its
complete construction, retained training state, target-discrimination
certificate, and evaluation cost, so that

\[
g\in
\mathcal R^{\mathrm{cert}}_{I,\Psi(B,n)}(T_I)
\cap
\mathcal H_{n,\Psi(B,n),m}^{\mathrm{sat,stat}}.
\]

With probability at least \(1-\delta\), simultaneously for every such
readout with positive variance,

\begin{equation}
\label{eq-noise-corrected-actionable-geom-extraction}
\operatorname{Ext}_I(g;T_I)
\le
\min\!\left\{
\begin{array}{l}
\operatorname{ActMass}_{I,\Psi(B,n)}(T_I),\\[1mm]
\operatorname{ProjMass}_{\sigma(g)}(T_I),\\[1mm]
\displaystyle
\frac{
\bigl(|\widehat c_{\eta,S}(g)|
+2\Gamma_{n,\Psi(B,n),m,\delta}^{\mathrm{sat,stat}}\bigr)^2
}{\rho^2\operatorname{Var}_{\mu_I}(g)}
\end{array}
\right\}.
\end{equation}

For a constant readout the extraction is zero. If an empirical audit,
effective-dimension ceiling, or certified residual-noise order also applies,
the corresponding ceilings of
\cref{def-extraction-audit-ceilings,thm-budgeted-extraction-bound} enter the
same minimum.

\end{proposition}

\begin{proof}

Membership in the certified class and
\cref{thm-budgeted-extraction-bound} give the actionable-mass ceiling.
The generated-algebra projection ceiling
\cref{thm-generated-algebra-projection-ceiling} gives the second ceiling.
Balance implies \(y_{T_I}=f_*/2\), and hence, for nonconstant \(g\),

\begin{equation}
\label{eq-actionable-extraction-correlation-identity}
\operatorname{Ext}_I(g;T_I)
=
\frac{|c_{0,I}(g)|^2}{\operatorname{Var}_{\mu_I}(g)}.
\end{equation}

The simultaneous transport inequality
\eqref{eq-saturated-statistical-clean-alignment-bound}, applied at the
enlarged complete cost, bounds

\[
|c_{0,I}(g)|
\le
\frac{
|\widehat c_{\eta,S}(g)|
+2\Gamma_{n,\Psi(B,n),m,\delta}^{\mathrm{sat,stat}}
}{\rho}.
\]

Squaring and substituting in
\eqref{eq-actionable-extraction-correlation-identity} proves the third
ceiling. All three hold on the same simultaneous event, so their minimum is
valid. The additional audit ceilings are the remaining conclusions of
\cref{thm-budgeted-extraction-bound}.

\end{proof}

\begin{theorem}[Saturated geometry--learning--selection accountability]
\label{thm-saturated-geometry-learning-selection-accountability}

Let a balanced binary-label task family have homogeneous noise
\(0\le\eta<1/2\), and let \(\mathfrak M_n\) be a monotone sublinear positively
homogeneous family functional. Equip the task with its represented public
verifier and the one-test clock refinement of
\cref{def-represented-pre-hit-selector}. Consider one family-uniform
represented computation whose complete learner, retained state, learned
geometry, target-discrimination readout, decoder, pre-hit sampler, verifier
calls, and terminal output have saturated cost at most \(B\). Let \(H_B(n)\)
be its largest admitted refined horizon. Suppose the represented decoder has
target-alignment margin \(0<\gamma\le1\) in the sense of
\cref{thm-population-alignment-generalization-memorization}. Then

\begin{equation}
\label{eq-joint-learning-selection-success-bound}
\begin{aligned}
&\mathfrak M_n\!\left(
 I\longmapsto
 \Pr_{S,\omega}\{\text{the represented decoder reaches }T_I\}
\right)\\
&\qquad\le
\min\!\left\{
\begin{aligned}
&\gamma^{-2}
\operatorname{Align}_{\mathfrak M,n,\eta}^{\mathrm{sat,stat}}(B,m),\\[-1mm]
&eH_B(n)
\sum_{c\in\mathfrak J_{\Abs}^5}
\mathfrak M_n\!\left(
 I\longmapsto
 \mathcal C_{I,\mathfrak F_5}^{(c)}
 \bigl(D_B^{\mathrm{sel}}(n)\bigr)
\right)
\end{aligned}
\right\}.
\end{aligned}
\end{equation}

The second entry is the whole-transcript first-hit accounting-charge bound of
\cref{thm-prospective-selector-structural-charge}.  It becomes a prospective
source envelope only when the same record also supplies the represented
pre-hit comparison required by
\cref{cor-contextual-structural-charge-transfer}.

The first entry admits both existing uniform ceilings

\begin{equation}
\label{eq-joint-learning-selection-alignment-ceilings}
\operatorname{Align}_{\mathfrak M,n,\eta}^{\mathrm{sat,stat}}(B,m)
\le
\operatorname{GenProj}_{\mathfrak M,n,\eta}^{\mathrm{sat,stat}}(B,m)
\le
\mathfrak M_n\!\left(
 I\longmapsto M_{I,n}^{\mathrm{sat,stat}}(B,m)
\right),
\end{equation}

and, on the simultaneous event of
\cref{thm-saturated-statistical-attack-alignment}, the empirical certificate
in \cref{eq-realized-average-seed-alignment-certificate}. A same-carrier
learned geometric source also satisfies
\cref{eq-full-saturated-geom-learned-projection-bound}; a represented
cross-carrier decoder is controlled by its displayed margin \(\gamma\) and
the first entry of \eqref{eq-joint-learning-selection-success-bound}.

Every score evaluation and every target-aligned selection operation are
charged in the same record. A score with an affordable evaluator enters the
statistical class through a represented training and output derivation. A
decoder or selector enters the budget slice when its own derivation, state,
queries, and execution cost fit that slice. Empirical tables and memorized
lookup state remain charged lifted coordinates. Their population-orthogonal
component is \(\Res\); an affordable population-aligned use is reclassified
into the existing \(\Geom/\Abs\) source ledger only after the relevant
channel predicate and pre-hit comparison are verified.

\end{theorem}

\begin{proof}

The decoder-margin conclusion
\eqref{eq-decoder-recovery-from-average-alignment}, applied to the complete
joint record, proves the first entry of
\eqref{eq-joint-learning-selection-success-bound}. Its two projection
ceilings are exactly
\eqref{eq-average-seed-alignment-projection-domination}, which proves
\eqref{eq-joint-learning-selection-alignment-ceilings}. The realized-sample
bound is \eqref{eq-realized-average-seed-alignment-certificate}.

For every fixed instance, the complete decoder and its verifier calls form a
represented finite-horizon computation. Apply
\cref{thm-prospective-selector-structural-charge} at budget \(B\):

\[
\Pr_{S,\omega}\{\text{the decoder reaches }T_I\}
\le
eH_B(n)
\sum_{c\in\mathfrak J_{\Abs}^5}
\mathcal C_{I,\mathfrak F_5}^{(c)}
\bigl(D_B^{\mathrm{sel}}(n)\bigr).
\]

Apply \(\mathfrak M_n\), use positive homogeneity and subadditivity, and
obtain the second entry of
\eqref{eq-joint-learning-selection-success-bound}. Both entries bound the
same success quantity, so their minimum also bounds it.
The source interpretation of the second entry, when used, is exactly the
additional transfer in \cref{cor-contextual-structural-charge-transfer}.

By \cref{def-budget-saturated-statistical-readout-class}, class membership
includes training, retained state, output representation, and fresh-input
evaluation. By
\cref{def-pre-hit-temporally-admissible-source,def-represented-pre-hit-selector},
the selector record also includes every pre-hit state, sampler, random-bit
use, and transition implementation. The budgeted derivability closure
\cref{def-budgeted-derivability-closure} charges their composition as one
derivation. Thus inexpensive evaluation of a fixed score supplies no
target-aligned selector without the represented selection derivation and its
cost.

Finally, apply
\cref{thm-two-level-learned-readout-provenance,thm-exhaustive-residual-normal-form}.
The population structural projection is charged to the existing source
ledger, and the orthogonal empirical component is \(\Res\). Its affordable
target-aligned reuse is reclassified by
\cref{prop-family-uniform-residual-reclassification}. This proves the
provenance assertions.

\end{proof}

\section{Complete Learned-Source Exhaustion and Noise Degradation}
\label{sec:learned-source-exhaustion-noise-degradation}

\begin{corollary}[Complete structural exhaustion of learned sources]
\label{cor-complete-structural-exhaustion-learned-sources}

Adopt the setting of
\cref{thm-saturated-geometry-learning-selection-accountability}.  Every
temporally admissible, dynamically qualified source record produced or used
by the complete learner--decoder--selector computation has, after the
class-preserving normalization of Part~III, exactly one index
\[
(\chi,S_{\mathrm{st}},S_{\mathrm{cf}}),
\qquad
\chi\in
\{\mathsf A,\mathsf X,\mathsf Q_{\mathrm{ex}},
  \mathsf Q_{\mathrm{rev}},\mathsf Q_{\mathrm{nr}}\},
\]
with
\[
\varnothing\ne S_{\mathrm{st}}\subseteq\mathfrak J_{\Abs}^{5},
\qquad
S_{\mathrm{cf}}\subseteq\mathfrak J_{\Abs}^{5}.
\]
Thus the complete learned source record lies in one of the 4,960 existing
source-mode--state-provenance--coefficient-provenance cells of
\cref{thm-universal-state-coefficient-source-decomposition}.  The index
retains the learner, selected kernel or geometry, represented coefficients,
decoder, sampler, qualification gates, and every downstream structural use
at their complete joint cost.

Restricting the Part~III supremum to learned prospective source records of
complete cost at most \(B\) gives
\begin{equation}
\label{eq-learned-source-complete-structural-exhaustion}
\sup_{\substack{
\mathscr C\ \text{is a learned prospective source record}\\
\operatorname{Cost}(\mathscr C)\preceq B}}
\mathfrak M_n\!\left(I\longmapsto Q_I V_{\mathscr C,I}\right)
\le
\max_{\substack{
\varnothing\ne S_{\mathrm{st}}\subseteq\mathfrak J_{\Abs}^{5}\\
S_{\mathrm{cf}}\subseteq\mathfrak J_{\Abs}^{5}}}
\operatorname{StructAct}_{\mathfrak M,
S_{\mathrm{st}},S_{\mathrm{cf}},n}^{\mathrm{sat}}
\bigl(\Psi_{\mathrm{all}}(B,n)\bigr).
\end{equation}
Every learned source therefore retains the complete numerical visibility of
its existing master cell.  In particular:

\begin{enumerate}[label=\textup{(\roman*)}]
\item the pure exact-\(\Abs\) mode \(\mathsf A\) is evaluated by the
      prospective quotient profile;
\item \(\mathsf Q_{\mathrm{ex}}\), \(\mathsf Q_{\mathrm{rev}}\), and
      \(\mathsf Q_{\mathrm{nr}}\) retain, respectively, the exact,
      reversible-compensated, and general-resolvent quotient--Geom
      visibility, including utility, stability, mixing, authorized Caus
      alignment, target reach, and regularity;
\item the ambient mode \(\mathsf X\) retains the exact same-record product
      \begin{equation}
      \label{eq-learned-ambient-complete-geom-product}
      Q_I V^X_{\Phi,D,J,\mathfrak K}
      =
      Q_I M_\Phi
      C_{\Phi,\mathfrak K}^{\mathrm{auth}}(D)
      \frac{R_{T_I}(D)}{1+\rho_\Phi(D)}.
      \end{equation}
      For a same-carrier learned geometry, the generated-algebra ceiling
      gives the simultaneous refinement
      \begin{equation}
      \label{eq-learned-ambient-geom-and-statistical-minimum}
      Q_I V^X_{\Phi,D,J,\mathfrak K}
      \le
      Q_I\min\!\left\{
      M_\Phi
      C_{\Phi,\mathfrak K}^{\mathrm{auth}}(D)
      \frac{R_{T_I}(D)}{1+\rho_\Phi(D)},
      \operatorname{ProjMass}_{\sigma(h_D)}(T_I)
      \right\},
      \qquad h_D=\frac{1-D}{1+D};
      \end{equation}
\item authorized \(\Caus\), valid \(\Lift\), and \(\Bdry\) uses inherit
      the same provenance cell, while their complete represented costs remain
      in the record.  They inherit a numerical source value only through the
      represented contextual transfer required by that cell; provenance
      alone supplies no numerical transfer.  The post-saturation \(\Res\)
      component contributes no prospective source cell.
\end{enumerate}

Consequently the generalized learning bounds use the full Part~III
structural exhaustion.  The statistical alignment and generalization
ceilings in
\cref{eq-joint-learning-selection-alignment-ceilings} and the structural
source ceiling in
\cref{eq-learned-source-complete-structural-exhaustion} bound the same
complete computation through their respective valid comparisons.  The
success estimate retains their minimum as in
\cref{eq-joint-learning-selection-success-bound}; neither the statistical
bound nor a provenance label replaces any Geom/Abs gate.
The source-resolved noise, data, confidence, and complexity coordinates are
compiled recordwise by
\cref{thm-source-resolved-saturated-learning-compilation}; hence every
learned source cell carries its own statistical candidate, ordered-reveal
contraction, contextual charge, and deployment converter.

\end{corollary}

\begin{proof}

The complete learning computation belongs to the budgeted derivability
closure by
\cref{def-budget-saturated-statistical-readout-class,thm-paper4-dependency-principle}.
Apply
\cref{thm-universal-state-coefficient-source-decomposition} to each
temporally admissible qualified source record in that derivation.  This gives
the unique master mode and exact state--coefficient provenance pair, with all
instance-dependent coefficient derivations and qualification gates retained.
Since the learned records form a subfamily of the records defining the
existing saturated source profile, restriction of the supremum followed by
\cref{eq-universal-master-structural-interaction-bound} proves
\cref{eq-learned-source-complete-structural-exhaustion}.

The mode assignments and their numerical visibility factors are
\cref{thm-affordable-geom-abs-normal-form,thm-full-saturated-profile-decomposition}.
The exact cellwise quotient and ambient evaluations are
\cref{cor-exact-signature-geom-cellwise-quantitative-form}.  In the ambient
case, \cref{def-geom-extraction} gives
\cref{eq-learned-ambient-complete-geom-product}, and
\cref{lem-learned-geom-target-reach-projection} gives the second ceiling in
\cref{eq-learned-ambient-geom-and-statistical-minimum}.  Provenance
inheritance and the residual assignment are
\cref{prop-derived-channels-create-no-structure,thm-exhaustive-residual-normal-form};
numerical transfer uses \cref{cor-contextual-structural-charge-transfer} when
its hypotheses hold.
Finally,
\cref{thm-saturated-geometry-learning-selection-accountability} already
takes the minimum of the valid statistical and structural estimates for the
complete joint record.  This proves the final assertion.

\end{proof}

\subsection{Source-resolved quantitative learning ledger}
\label{subsec-source-resolved-quantitative-learning-ledger}

\begin{definition}[Complete source-resolved learning ledger]
\label{def-complete-source-resolved-learning-ledger}

Let \(R\) be a complete represented learning record of sample size \(m\),
confidence level \(1-\delta\), horizon \(H\), and saturated cost at most
\(B\).  The record includes data acquisition, feature and kernel
construction, optimization, retained state, output evaluation, every
decoder or selector, and deployment.  Normalize every temporally admissible
target-qualified source in its backward slice by
\cref{thm-universal-state-coefficient-source-decomposition}.  Its exact
index is
\[
\alpha=(\chi,S_{\rm st},S_{\rm cf}),
\qquad
\chi\in\{\mathsf A,\mathsf X,\mathsf Q_{\rm ex},
          \mathsf Q_{\rm rev},\mathsf Q_{\rm nr}\},
\]
with
\(\varnothing\ne S_{\rm st}\subseteq\mathfrak J_{\Abs}^{5}\) and
\(S_{\rm cf}\subseteq\mathfrak J_{\Abs}^{5}\).  Write
\(\mathfrak A_{\rm learn}\) for this \(4960\)-element index set.

For \(a\in\mathfrak J_{\Abs}^{5}\), let
\begin{equation}
\label{eq-source-resolved-learning-contextual-ledger}
\mathcal L_{I,R}^{(a)}
=
\sum_{\ell:\,\mathsf W_a(E_{R,\ell})}
\operatorname{ch}_I(E_{R,\ell})
\end{equation}
be the five-family contextual charge in the complete
learner--deployment backward slice.  Every selected likelihood, feature
locator, coefficient, optimizer, belief coordinate, acquisition score,
kernel, current, Lift map, boundary converter, and output rule remains in
the carrier of this sum.

Enumerate the target-dependent reveals actually available before deployment
as \(Y_1,\ldots,Y_{L_R}\), use the canonical anchored assignment of
\cref{cor-universal-syntax-affordable-charge-cover}, and write \(a_j\) for
the source tag of \(Y_j\).  Define
\begin{equation}
\label{eq-source-resolved-learning-information-ledger}
J_{I,R}^{(a)}
=
\sum_{\substack{1\le j\le L_R\\a_j=a}}
\mathbb E\min\{c_j,\vartheta_j a_j^{\rm info}\},
\qquad
J_{I,R}^{\rm src}
=\sum_{a\in\mathfrak J_{\Abs}^{5}}J_{I,R}^{(a)}.
\end{equation}
For a binary symmetric reveal of rate \(\eta_j\), the admissible channel
coefficient is
\(\vartheta_j=(1-2\eta_j)^2\).  Other observation, reward, transition, and
filter channels use their typed coefficients from
\cref{def-controlled-channel-information-ledger}.  Deterministic
postprocessing contributes zero information and retains its construction
cost.

For a time \(0\le t<H\) and common ceiling \(D\), let
\(\mathscr R_{t}^{\rm learn}(D)\) denote the complete family-uniform
learning-record derivations deployed at time \(t\) with saturated cost at
most \(D\), and write \(R_I\) for the record produced on instance \(I\).
Define the saturated cellwise envelopes
\begin{align}
\label{eq-saturated-source-resolved-learning-charge-envelope}
\mathcal L_{\infty,t}^{(a),\rm sat}(D)
&=
\sup_{R\in\mathscr R_t^{\rm learn}(D)}\sup_I
\mathcal L_{I,R_I}^{(a)},\\
\label{eq-saturated-source-resolved-learning-information-envelope}
J_{\infty,t}^{(a),\rm sat}(D)
&=
\sup_{R\in\mathscr R_t^{\rm learn}(D)}\sup_I
J_{I,R_I}^{(a)}.
\end{align}
The supremum uses one family-uniform represented record on the whole size
slice; only its numerical evaluation is then maximized over presented
instances.  Thus these envelopes are pointwise ceilings and require no
interchange of a square root with a family functional.

Let \(\mathcal H_\alpha(B,m)\) be the output or loss class of complete
represented records of saturated cost at most \(B\) whose selected
prospective source has exact index \(\alpha\).  On the same probability law
and in the same loss units, let
\(\Gamma_{R,\alpha}^{\rm learn}(B,m,\boldsymbol\eta,\delta)\) be the minimum
of the applicable fixed-operator Rademacher, learned-operator,
effective-dimension, structural PAC--Bayes, sequential-Rademacher,
martingale PAC--Bayes, mixing-block, and finite-description candidates for
that restricted class.  Its value is \(+\infty\) when no candidate applies.

Finally, let \(U_R^{\rm stat\to tgt}\) be the smallest bound obtained by a
represented same-law, same-unit converter from one of these statistical
loss bounds to the deployed target-contact advantage.  Let
\(U_R^{\rm dyn\to tgt}\) be the corresponding minimum of applicable
capacity, Blackwell, Caus, Lift, observation, active-acquisition, and
controlled-dynamics bounds.  An unavailable target converter has value one.
All objects are evaluated at one class-preserving common ceiling
\begin{equation}
\label{eq-complete-source-resolved-learning-cost}
D_B^{\rm learn}
=
\Psi_{\rm learn}(B,H,n,m,\boldsymbol\eta,\delta),
\end{equation}
which majorizes the complete record, normalized frontier, ordered reveals,
confidence certificates, channel corrections, and every converter used in
a displayed minimum.  This ceiling changes no public constructor algebra.

\end{definition}

\begin{theorem}[Source-resolved compilation of saturated learning]
\label{thm-source-resolved-saturated-learning-compilation}

Use the ledger of
\cref{def-complete-source-resolved-learning-ledger}.  The following
statements hold.

\begin{enumerate}[label=\textup{(\roman*)}]
\item The target-qualified prospective source records form the disjoint
      union of their exact cells
      \(\alpha\in\mathfrak A_{\rm learn}\).  Authorized \(\Caus\), valid
      \(\Lift\), and \(\Bdry\) processing retain the exact cell of their
      earlier source.  Their evaluation and deployment costs remain in
      \(D_B^{\rm learn}\).  The post-saturation \(\Res\) component has zero
      target-qualified learning gain; an affordable aligned representative
      is reclassified before the residual is formed.

\item For losses in \([-1,1]\), the empirical Rademacher complexity of the
      complete exact-cell union satisfies
      \begin{equation}
      \label{eq-source-resolved-learning-rademacher-union}
      \widehat{\mathfrak R}_S
      \!\left(\bigcup_{\alpha\in\mathfrak A_{\rm learn}}
      \mathcal H_\alpha(B,m)\right)
      \le
      \max_{\alpha\in\mathfrak A_{\rm learn}}
      \widehat{\mathfrak R}_S(\mathcal H_\alpha(B,m))
      +\sqrt{\frac{2\log|\mathfrak A_{\rm learn}|}{m}}.
      \end{equation}
      Empty cells are omitted.  For a PAC--Bayes prior
      \(P=\sum_\alpha w_\alpha P_\alpha\) and posterior
      \(Q=\sum_\alpha q_\alpha Q_\alpha\) with matching conditional cells,
      \begin{equation}
      \label{eq-source-resolved-learning-pac-bayes-mixture}
      \operatorname{KL}(Q\Vert P)
      =
      \operatorname{KL}(q\Vert w)
      +\sum_\alpha q_\alpha
       \operatorname{KL}(Q_\alpha\Vert P_\alpha).
      \end{equation}
      Thus source selection is paid either by the explicit finite-cell
      Rademacher term or by the PAC--Bayes cell-selection divergence.
      Effective-dimension and source-condition candidates are evaluated
      inside the selected exact cell; they are aggregated across cells only
      through a valid union, mixture, or represented selection bound.

\item Let \(K_R\) be the represented deployed selector, let \(\nu_R\) be its
      declared target-neutral baseline, and put
      \(g_R=(\operatorname{Adv}(K_R;\nu_R))_+\).  With
      \(\beta_R^{+}\) the applicable decoupled neutral excess under the same
      pre-hit law,
      \begin{equation}
      \label{eq-source-resolved-learning-target-gain-minimum}
      g_R
      \le
      \min\left\{
      1,\,
      U_R^{\rm stat\to tgt},\,
      \beta_R^{+}
       +\sqrt{\frac{\log2}{2}J_{I,R}^{\rm src}},\,
      eH_B(n)\sum_{a\in\mathfrak J_{\Abs}^{5}}
       \mathcal L_{I,R}^{(a)},\,
      U_R^{\rm dyn\to tgt}
      \right\}.
      \end{equation}
      Each entry bounds the same survival-weighted target advantage under
      the same law.  The minimum is taken for one complete record before any
      supremum.

\item Let \(\mathscr R_t^{\rm learn}(D_B^{\rm learn})\) be the complete
      family-uniform learning-record derivations deployed at time \(t\), and
      denote the right side of
      \cref{eq-source-resolved-learning-target-gain-minimum} by
      \(U_{I,R,t}^{\rm learn}\).  Let
      \(\operatorname{Enum}_{I,\nu}^{\rm sat,learn}(B,H)\) be the supremum
      of the declared neutral-contact sum over the same complete learning
      computations.  Then
      \begin{align}
      \label{eq-source-resolved-learning-saturated-selector}
      \operatorname{Sel}_{I,n}^{\rm sat,learn}(B)
      &\le
      \max_{0\le t<H}
      \sup_{R\in\mathscr R_t^{\rm learn}(D_B^{\rm learn})}
      U_{I,R_I,t}^{\rm learn},\\
      \label{eq-source-resolved-learning-saturated-arrival}
      \operatorname{Succ}_{I,n}^{\rm sat,learn}(B)
      &\le
      \operatorname{Enum}_{I,\nu}^{\rm sat,learn}(B,H)
      +\sum_{t=0}^{H-1}
      \sup_{R\in\mathscr R_t^{\rm learn}(D_B^{\rm learn})}
      U_{I,R_I,t}^{\rm learn}.
      \end{align}
      Both inequalities persist under a monotone family functional.
\end{enumerate}

Accordingly every learning bound has four typed numerical layers:
complete construction cost, source-resolved statistical capacity,
source-resolved noise/information contraction, and a represented converter
to the deployed target quantity.  A finite value in one layer never
substitutes for an unavailable value in another.

\end{theorem}

\begin{proof}

The exact-cell partition and channel inheritance in \textup{(i)} are
\cref{thm-universal-state-coefficient-source-decomposition,%
prop-derived-channels-create-no-structure,%
thm-exhaustive-residual-normal-form}.  Every cell restriction keeps the
complete represented record, so the common cost is
\cref{eq-complete-source-resolved-learning-cost}.

For \textup{(ii)}, apply the standard exponential-moment proof of the
Rademacher union bound to the finite collection of nonempty cell-restricted
classes.  Each empirical process is \(1/\sqrt m\)-sub-Gaussian for
losses in \([-1,1]\); maximizing over at most
\(|\mathfrak A_{\rm learn}|\) classes gives the displayed selection term.
The PAC--Bayes identity is the relative-entropy chain rule for the joint
cell index and within-cell hypothesis.  Restricting an operator or kernel
to one cell leaves its existing effective-dimension and source-condition
proof unchanged.

For \textup{(iii)}, the information candidate is the ordered-reveal
decoupling inequality of
\cref{thm-adaptive-acquisition-information-target-contact}, with the chain
rule split according to
\cref{eq-source-resolved-learning-information-ledger}.  The structural
candidate is
\cref{thm-prospective-selector-structural-charge} applied to the complete
learner--deployment slice.  The statistical and dynamical entries are
admitted only through their represented target converters, so all entries
bound the same scalar.  Taking their minimum is therefore valid.

Finally, take the record supremum only after the recordwise minimum and use
\cref{thm-prospective-selector-accountability}.  This proves
\cref{eq-source-resolved-learning-saturated-selector,%
eq-source-resolved-learning-saturated-arrival}.  Monotonicity gives the
family-functional form.

\end{proof}

\begin{corollary}[Noise--data--complexity learning consequence]
\label{cor-source-resolved-noise-data-complexity-learning}

At a common ceiling \(D_B^{\rm learn}\), suppose
\[
\mathcal L_{\infty,t}^{(a),\rm sat}
\le\varepsilon_{a,t}(n,B,m,\boldsymbol\eta,\delta),
\qquad
J_{\infty,t}^{(a),\rm sat}
\le\jmath_{a,t}(n,B,m,\boldsymbol\eta,\delta).
\]
Then every complete saturated learning record satisfies
\begin{equation}
\label{eq-source-resolved-noise-data-complexity-bound}
\begin{aligned}
\operatorname{Succ}_{\mathfrak M,n}^{\rm learn}(B)
&\le
\operatorname{Enum}_{\mathfrak M,\nu}(B,H)
+\sum_{t=0}^{H-1}
\min\left\{
\begin{aligned}
&1,\ U_{\mathfrak M,t}^{\rm stat\to tgt},\
\beta_{\mathfrak M,t}^{+}
+\sqrt{\frac{\log2}{2}\sum_a\jmath_{a,t}},\\[-1mm]
&eH_B\sum_a\varepsilon_{a,t},\
U_{\mathfrak M,t}^{\rm dyn\to tgt}
\end{aligned}
\right\}.
\end{aligned}
\end{equation}
For a polynomial horizon, negligible neutral contact and negligible
source-resolved inputs give negligible target arrival.  The conclusion
applies simultaneously to passive, active, structural Bayesian, sequential,
and controlled learning because each is a restriction of the complete
record class.

\end{corollary}

\begin{proof}

Apply
\cref{eq-source-resolved-learning-saturated-arrival} and substitute the two
source-resolved ceilings before taking the time sum.  The stated learning
classes are restrictions of the same saturated record class, so no further
union over learning paradigms is required.

\end{proof}

\begin{figure}[tbp]
\centering
\resizebox{.98\linewidth}{!}{%
\begin{tikzpicture}[fw diagram,node distance=8mm and 10mm]
  \node[fw source,text width=2.8cm] (pub)
    {public presentation\\data, budget, noise};
  \node[fw provenance,text width=3.1cm,right=of pub] (cell)
    {exact source cell\\support \(+\) five-family ancestry};
  \node[fw use,text width=2.8cm,above right=7mm and 11mm of cell] (stat)
    {statistical capacity\\\(m,B,\delta\)};
  \node[fw use,text width=2.8cm,right=11mm of cell] (info)
    {source information\\\(\vartheta(\eta),J^{(a)}\)};
  \node[fw geom,text width=2.8cm,below right=7mm and 11mm of cell] (dyn)
    {Geom/Caus/Lift\\same-unit converter};
  \node[fw terminal,text width=3.0cm,right=18mm of info] (min)
    {recordwise target bound\\minimum};
  \node[fw box,text width=2.8cm,right=of min] (sat)
    {saturated supremum\\then time sum};
  \draw[fw arrow] (pub)--(cell);
  \draw[fw arrow] (cell)--(stat);
  \draw[fw arrow] (cell)--(info);
  \draw[fw arrow] (cell)--(dyn);
  \draw[fw arrow] (stat)--(min);
  \draw[fw arrow] (info)--(min);
  \draw[fw arrow] (dyn)--(min);
  \draw[fw arrow] (min)--(sat);
\end{tikzpicture}%
}
\caption{Source-resolved learning compilation.  Each complete record keeps
its exact structural ancestry, statistical capacity, source-assigned noise
contraction, and dynamical converter.  Comparable target bounds are minimized
within the record before saturation, as required by
\cref{thm-source-resolved-saturated-learning-compilation}.}
\label{fig-source-resolved-learning-compilation}
\end{figure}

\begin{table}[tbp]
\centering
\small
\begin{tabular}{@{}p{.22\linewidth}p{.23\linewidth}p{.22\linewidth}p{.25\linewidth}@{}}
\toprule
Learning coordinate & Statistical candidate & Source/noise ledger & Target deployment \\
\midrule
Fixed or learned kernel, KRR, NTK
& Rademacher, effective dimension, operator selection
& exact cell, \(\mathcal L^{(a)}\), \(J^{(a)}\)
& represented decoder or selector margin
\\
Structural PAC--Bayes
& within-cell KL and cell-mixture KL
& prior/posterior ancestry and channel SDPI
& same-law posterior-to-action converter
\\
Dependent or sequential learning
& blocking, sequential Rademacher, martingale PAC--Bayes
& ordered reveals with typed contraction
& represented policy or next-test converter
\\
Active and curiosity learning
& observation resolution and learned-policy radius
& acquisition information and five-family charge
& next-contact, capacity, or master comparison
\\
Controlled and Bayesian learning
& model, Bellman, filter, and belief estimation
& complete source cell with Caus/Lift inheritance
& value, return, or arrival converter in matching units
\\
\bottomrule
\end{tabular}
\caption{Application of the source-resolved ledger to the learning master
bounds.  Each row retains \(B,m,\boldsymbol\eta,\delta\), and \(H\) when
present.  Statistical and structural quantities meet only in the final
column through a represented converter.}
\label{tab-source-resolved-learning-audit}
\end{table}

\Cref{fig-learning-structural-exhaustion} summarizes the exact Part~III
classification inherited by every prospective learned source.

\begin{figure}[tbp]
\centering
\begin{tikzpicture}[fw diagram]
  \node[fw source,minimum width=3.4cm] (record) at (-5.7,0)
    {complete learner--decoder\\selector record at cost \(B\)};
  \node[fw provenance,minimum width=3.0cm] (index) at (-1.8,0)
    {unique exact index\\\((\chi,S_{\rm st},S_{\rm cf})\)};
  \coordinate (leftbus) at (0,0);
  \node[fw geom,minimum width=2.7cm] (abs) at (2.0,2.0)
    {exact \(\Abs\)};
  \node[fw geom,minimum width=2.7cm] (ambient) at (2.0,1.0)
    {ambient \(\Geom\)};
  \node[fw use,minimum width=2.7cm] (exact) at (2.0,0)
    {exact-quotient \(\Geom\)};
  \node[fw use,minimum width=2.7cm] (rev) at (2.0,-1.0)
    {reversible compensation};
  \node[fw use,minimum width=2.7cm] (resolvent) at (2.0,-2.0)
    {resolvent compensation};
  \coordinate (rightbus) at (4.0,0);
  \node[fw terminal,minimum width=3.0cm] (min) at (6.2,0)
    {statistical/structural\\minimum};
  \draw[fw arrow] (record) -- (index);
  \draw[fw arrow] (index.east) -- (leftbus);
  \foreach \x in {abs,ambient,exact,rev,resolvent}
    \draw[fw arrow] (leftbus) |- (\x.west);
  \foreach \x in {abs,ambient,exact,rev,resolvent}
    \draw[draw=fwInk,line width=.75pt] (\x.east) -| (rightbus);
  \draw[fw arrow] (rightbus) -- (min.west);
\end{tikzpicture}
\caption{Complete structural exhaustion of learned sources. The complete
learned record receives one of the 4,960 existing Part~III indices from
\cref{cor-complete-structural-exhaustion-learned-sources}; its exact source
visibility is combined recordwise with the applicable statistical,
noise-information, and dynamical ceilings by
\cref{thm-source-resolved-saturated-learning-compilation}.}
\label{fig-learning-structural-exhaustion}
\end{figure}

\begin{corollary}[Finite-description radius for the complete saturated
output class]
\label{cor-finite-description-saturated-learning-radius}

Fix \(n,B,m\), a scalarized locally finite derivation coding as in
\cref{def-budgeted-derivability-closure}, and let

\[
N_{n,B,m}^{\mathrm{out}}
=
\sup_{I\in\mathfrak I_n}
\left|\mathcal H_{n,B,m}^{\mathrm{sat,stat}}(I)\right|.
\]

Set this value to \(+\infty\) unless the fixed-budget output codebook is
finite uniformly over \(I\in\mathfrak I_n\).  When it is finite, the clipped
linear-loss class generated by the complete saturated outputs satisfies,

\begin{equation}
\label{eq-finite-description-saturated-radius}
\Gamma_{n,B,m,\delta}^{\mathrm{sat,stat}}
\le
2\sqrt{\frac{2\log N_{n,B,m}^{\mathrm{out}}}{m}}
+3\sqrt{\frac{\log(2/\delta)}{2m}}.
\end{equation}

The cardinality counts complete output records, including represented
training state, selected geometry, retained coefficients and precision,
decoder data, and every sample-dependent lookup value used at evaluation.
When a declared coding estimate gives
\(\log N_{n,B,m}^{\mathrm{out}}\le L(B,n,m)\), the first term in
\eqref{eq-finite-description-saturated-radius} is at most
\(2\sqrt{2L(B,n,m)/m}\).

\end{corollary}

\begin{proof}

For each fixed instance, local coding in
\cref{def-budgeted-derivability-closure} uses a finite
alphabet, charges every encoded coefficient, retained table, precision
parameter, state value, and evaluation rule, and has finite scalar-cost
sublevels.  A uniform size-slice bound additionally requires the displayed
supremum over instances to be finite; otherwise the corollary is applied
pointwise in \(I\), or its right side is \(+\infty\).

Massart's finite-class lemma for a \([0,1]\)-valued loss class gives

\[
\widehat{\mathfrak R}_S(\mathcal L_{\mathcal H})
\le
\sqrt{\frac{2\log N_{n,B,m}^{\mathrm{out}}}{m}}.
\]

Substitution in the radius definition
\eqref{eq-corrupted-rademacher-radius} proves
\eqref{eq-finite-description-saturated-radius}. The final statement follows
by monotonicity of the square root.

\end{proof}

\begin{theorem}[Noise degradation and learning accountability]
\label{thm-noise-indexed-learning-accountability}

Let \((\mathfrak I_{n,m,\eta})_{0\le\eta<1/2}\) be a binary-label family under
homogeneous flip noise, and let \(0\le\eta_1\le\eta_2<1/2\).  Put

\[
q(\eta_1,\eta_2)
=\frac{\eta_2-\eta_1}{1-2\eta_1}.
\]

Let \(\psi_{\mathrm{flip}}(B,n,m)\) charge the represented independent
Bernoulli-\(q(\eta_1,\eta_2)\) degradation of a length-\(m\) sample and every
downstream retained use. Assume that \(\mathfrak M_n\) is subadditive in
addition to being monotone and positively homogeneous. Then

\begin{align}
\label{eq-noise-degradation-average-alignment}
\operatorname{Align}_{\mathfrak M,n,\eta_1}^{\mathrm{sat,stat}}
 \bigl(\psi_{\mathrm{flip}}(B,n,m),m\bigr)
&\ge
\operatorname{Align}_{\mathfrak M,n,\eta_2}^{\mathrm{sat,stat}}(B,m),\\
\label{eq-noise-degradation-generalizing-projection}
\operatorname{GenProj}_{\mathfrak M,n,\eta_1}^{\mathrm{sat,stat}}
 \bigl(\psi_{\mathrm{flip}}(B,n,m),m\bigr)
&\ge
\operatorname{GenProj}_{\mathfrak M,n,\eta_2}^{\mathrm{sat,stat}}(B,m).
\end{align}

For a finite horizon \(H_B\), define the saturated neutral-contact profile

\begin{equation}
\label{eq-noise-indexed-neutral-contact-profile}
\operatorname{Enum}_{\mathfrak M,n,\eta}^{\mathrm{sat}}(B,H_B)
=
\sup_{\mathcal A:\operatorname{Cost}(\mathcal A)\preceq B}
\mathfrak M_n\!\left(
 I\longmapsto
 \operatorname{Enum}_{\nu}(\mathcal A,I,H_B)
\right),
\end{equation}

where the same declared target-neutral baseline rule is used throughout the
family.  Write
\(\operatorname{Succ}_{\mathfrak M,n,\eta}(B)\) for
\cref{eq-uniform-affordable-success-profile} evaluated on
\(\mathfrak I_{n,m,\eta}\).  Define

\begin{equation}
\label{eq-noise-indexed-family-selector-profile}
\operatorname{Sel}_{\mathfrak M,n,\eta}^{\mathrm{sat}}(B)
=
\sup_{K}
\mathfrak M_n\!\left(
 I\longmapsto
 \bigl(\operatorname{Adv}_{t}(K_{I,t};\nu_{I,t})\bigr)_+
\right),
\end{equation}

where the supremum ranges over one family-uniform temporally admissible
represented selector derivation, including its represented time rule, of
complete saturated cost at most \(B\). Thus one selector program is used on
the entire size slice; the supremum is not chosen separately for each
instance. Every represented budget-\(B\)
learning or search computation satisfies

\begin{equation}
\label{eq-noise-indexed-prospective-learning-accountability}
\operatorname{Succ}_{\mathfrak M,n,\eta}(B)
\le
\operatorname{Enum}_{\mathfrak M,n,\eta}^{\mathrm{sat}}(B,H_B)
+H_B
\operatorname{Sel}_{\mathfrak M,n,\eta}^{\mathrm{sat}}
 \bigl(\Psi(B,n)\bigr).
\end{equation}

For the complete learning subclass, the selector term in
\cref{eq-noise-indexed-prospective-learning-accountability} has the
source-resolved evaluation
\begin{equation}
\label{eq-noise-indexed-source-resolved-selector-evaluation}
\operatorname{Sel}_{\mathfrak M,n,\eta}^{\mathrm{sat,learn}}(B)
\le
\max_{0\le t<H_B}
\sup_{R\in\mathscr R_t^{\rm learn}(D_B^{\rm learn})}
\mathfrak M_n\!\left(
 I\longmapsto U_{I,R_I,t}^{\rm learn}
\right),
\end{equation}
where \(U_{I,R,t}^{\rm learn}\) is the same-record minimum in
\cref{eq-source-resolved-learning-target-gain-minimum}.  In particular,
the binary noise factor \((1-2\eta)^2\), the five source charges, the
sample size, confidence, complete construction cost, and every
statistical-to-target or dynamics-to-target converter remain in the
selector profile rather than being discarded after a risk estimate.

Consequently, a negligible neutral-contact profile together with a negligible
prospective selector profile gives negligible target arrival. If
\(\mathfrak M_n\) is expectation under the declared family law, then conversely, if
one budget-\(B\) computation has family-averaged success exceeding its
family-averaged neutral contact by \(\varepsilon\), then its represented
pre-hit selectors have total positive advantage exceeding \(\varepsilon\),
and at least one refined step has family-averaged positive advantage exceeding
\(\varepsilon/H_B\).

The profiles in
\cref{eq-noise-degradation-average-alignment,eq-noise-degradation-generalizing-projection,eq-noise-indexed-prospective-learning-accountability}
range over the complete saturated budget class.  Hence they quantify the noise
dependence of all family-uniform represented learning rules at once; no
algorithm-specific source family is introduced.

\end{theorem}

\begin{proof}

An independent flip of rate \(q(\eta_1,\eta_2)\) following a flip of rate
\(\eta_1\) has total flip rate \(\eta_2\).  Given any budget-\(B\) learner for
the \(\eta_2\) presentation, first apply this represented degradation to an
\(\eta_1\) sample and then run the same learner.  The joint distribution of
its retained output, represented seed, and every downstream evaluation is
identical to the original \(\eta_2\) execution.  Its complete cost is at most
\(\psi_{\mathrm{flip}}(B,n,m)\).  The two defining suprema therefore give
\eqref{eq-noise-degradation-average-alignment} and
\eqref{eq-noise-degradation-generalizing-projection}.

For every fixed computation, the pointwise prospective hazard identity in
\cref{thm-prospective-selector-accountability}, followed by replacement of
each signed advantage by its positive part, gives

\[
\Pr_I\{\tau\le H_B\}
\le
\operatorname{Enum}_{\nu}(\mathcal A,I,H_B)
+\sum_{t=0}^{H_B-1}
\bigl(\operatorname{Adv}_{t}(K_{I,t};\nu_{I,t})\bigr)_+.
\]

Apply the monotone sublinear positively homogeneous functional
\(\mathfrak M_n\), and then take the supremum over the common budget class.
For each fixed \(t\), the transition of a uniform computation is one of the
uniform selector derivations in
\eqref{eq-noise-indexed-family-selector-profile}; no pointwise choice of a
different selector for each instance is made.  Subadditivity contributes at
most \(H_B\) selector terms, and the supremum of a sum is at most the sum of
the suprema.  This proves
\eqref{eq-noise-indexed-prospective-learning-accountability}.  For the linear
expectation functional, the converse is the family-functional form of the exact identity
\eqref{eq-selector-hazard-decomposition}: the sum of the positive parts of the
step advantages dominates their signed sum, and one of at most \(H_B\) terms
is at least their average.

Finally,
\cref{eq-noise-indexed-source-resolved-selector-evaluation} is
\cref{eq-source-resolved-learning-saturated-selector} restricted to the
noise-indexed family at the common ceiling.  This restriction changes
neither the public algebra nor the complete cost of a record.

\end{proof}

\begin{corollary}[Noise-indexed learning threshold criterion]
\label{cor-noise-indexed-learning-threshold-criterion}

Use the expectation functional \(\mathfrak M_{n,m,\eta}\) and fixed
scalarization \(\kappa\) of
\cref{def-noise-indexed-saturated-recovery-cost}.  Assume that every
computation counted by \(B_{\mathrm{rec}}^*(n,m,\eta;\alpha)\) has, within
the same complete scalar cost, a represented learner--decoder realization
counted by
\cref{def-average-seed-population-alignment-profile}, and that every
successful decoded output has alignment margin at least \(\gamma\).
For \(0<\theta\le1\), define

\begin{equation}
\label{eq-noise-indexed-population-alignment-cost}
B_{\mathrm{align}}^*(n,m,\eta;\theta)
=
\inf\left\{
s\ge0:
\operatorname{Align}_{\mathfrak M,n,\eta,\kappa}^{\mathrm{sat,stat}}(s,m)
\ge\theta
\right\},
\end{equation}

with value \(\infty\) when the set is empty.  Then

\begin{equation}
\label{eq-recovery-cost-dominated-by-alignment-cost}
B_{\mathrm{align}}^*(n,m,\eta;\alpha\gamma^2)
\le
B_{\mathrm{rec}}^*(n,m,\eta;\alpha).
\end{equation}

Here \(B_{\mathrm{rec}}^*\) is the noise-indexed saturated recovery cost of
\cref{def-noise-indexed-saturated-recovery-cost}.

At a scalar budget \(s\), the strict inequality

\[
\operatorname{Align}_{\mathfrak M,n,\eta,\kappa}^{\mathrm{sat,stat}}(s,m)
<\alpha\gamma^2
\]

excludes recovery probability \(\alpha\) for the complete represented learner
class at that budget.  Equations
\eqref{eq-noise-degradation-average-alignment} and
\eqref{eq-noise-degradation-generalizing-projection} make the affordable
alignment region downward closed under class-preserving flip degradation.

\end{corollary}

\begin{proof}

If a scalar-cost-\(s\) computation recovers with probability at least
\(\alpha\), the stated represented realization and
\cref{eq-decoder-recovery-from-average-alignment} give

\[
\operatorname{Align}^{\mathrm{sat,stat}}(s,m)
\ge\alpha\gamma^2.
\]

Taking the two infima proves
\eqref{eq-recovery-cost-dominated-by-alignment-cost}.  Its contrapositive gives
the budgeted exclusion.  Downward closure is
\cref{thm-noise-indexed-learning-accountability}.

\end{proof}

\Cref{fig-saturated-statistical-attack-accountability} displays the single
uniform route from a represented learner to target alignment and structural
charging.

\begin{figure}[tbp]
\centering
\begin{tikzpicture}[fw diagram,node distance=8mm and 6mm]
  \node[fw source,minimum width=2.5cm] (sample)
    {corrupted sample};
  \node[fw provenance,right=of sample,minimum width=3.0cm] (slice)
    {saturated readout class\\\(\mathcal H^{\mathrm{sat,stat}}_{n,B,m}\)};
  \node[fw use,right=of slice,minimum width=2.5cm] (score)
    {selected score \(h\)};
  \node[fw terminal,right=of score,minimum width=2.7cm] (align)
    {clean target alignment};
  \node[fw box,below=of slice,minimum width=3.0cm] (rad)
    {uniform radius \(\Gamma\)};
  \node[fw geom,below=of align,minimum width=3.0cm] (profile)
    {Geom/Abs source ledger};
  \draw[fw arrow] (sample) -- (slice);
  \draw[fw arrow] (slice) -- (score);
  \draw[fw arrow] (score) -- (align);
  \draw[fw arrow] (rad) -- node[right,font=\scriptsize]
    {simultaneous control} (slice);
  \draw[fw arrow] (align) -- node[right,font=\scriptsize]
    {represented use} (profile);
\end{tikzpicture}
\caption{Every budget-admissible statistical attack selects one readout from
the same saturated class.  Uniform concentration controls sample-dependent
selection.  Empirical tables remain charged coordinates; on the population
carrier, every target-bearing structural projection retains its complete
Geom/Abs provenance and the orthogonal remainder is \(\Res\).}
\label{fig-saturated-statistical-attack-accountability}
\end{figure}

\section{Noise-Indexed Recovery Budgets and the Three Threshold Scales}
\label{subsec-noise-indexed-saturated-thresholds}
\label{sec:noise-indexed-recovery-budgets}

The statistical radius, the represented construction cost, and the information
carried by the sample impose distinct thresholds.  The following notation fixes
their common binary-noise scale.  For \(0\le u\le1\), let

\begin{equation}
\label{eq-binary-entropy-definition}
H_2(u)=-u\log_2u-(1-u)\log_2(1-u),
\end{equation}

with \(0\log_2 0=0\).  For \(0<u,v<1\), let

\begin{equation}
\label{eq-bernoulli-divergence-definition}
D_{\mathrm{Ber}}(u\Vert v)
=u\log\frac uv+(1-u)\log\frac{1-u}{1-v}.
\end{equation}

Thus \(D_{\mathrm{Ber}}\) uses natural logarithms.  The notation
\(H_2^{-1}\) denotes the
inverse of the increasing restriction \(H_2:[0,1/2]\to[0,1]\).

For integers \(m\ge1\) and \(0\le k\le m\), put

\begin{equation}
\label{eq-binary-hamming-ball-count}
N_m(k)=\sum_{j=0}^k\binom mj,
\qquad N_m(-1)=0.
\end{equation}

\begin{lemma}[Binary Hamming-ball counting bounds]
\label{lem-binary-hamming-ball-cost-bounds}

For \(0\le k\le m/2\), with \(q=k/m\),

\begin{equation}
\label{eq-binary-hamming-ball-entropy-bounds}
\frac{2^{mH_2(q)}}{m+1}
\le
\binom mk
\le
N_m(k)
\le
2^{mH_2(q)}.
\end{equation}

For every fixed integer \(k\ge1\), as \(m\to\infty\),

\begin{equation}
\label{eq-fixed-weight-hamming-ball-asymptotic}
N_m(k)=\frac{m^k}{k!}\left(1+O_k(m^{-1})\right).
\end{equation}

\end{lemma}

\begin{proof}

For \(k=0\), one has
\(\binom m0=N_m(0)=2^{mH_2(0)}=1\).  Assume
\(1\le k\le m/2\), so \(0<q\le1/2\).  In the binomial law with success
probability \(q=k/m\), the value \(k\) is a mode.  Since its \(m+1\) point
masses sum to one,
\[
\binom mk q^k(1-q)^{m-k}\ge\frac1{m+1}.
\]
Using
\(q^{-k}(1-q)^{-(m-k)}=2^{mH_2(q)}\) gives the first inequality in
\cref{eq-binary-hamming-ball-entropy-bounds}.  For \(j\le k\) and
\(q\le1/2\),
\[
q^j(1-q)^{m-j}
\ge q^k(1-q)^{m-k}.
\]
Multiplying the sum defining \(N_m(k)\) by the right-hand side and bounding
it by the total binomial mass proves the upper type bound.  Finally,
\(\binom mk=m^k/k!+O_k(m^{k-1})\), while the terms with \(j<k\) are
\(O_k(m^{k-1})\).  This proves
\cref{eq-fixed-weight-hamming-ball-asymptotic}.

\end{proof}

\begin{definition}[Noise-indexed saturated recovery cost]
\label{def-noise-indexed-saturated-recovery-cost}

Let \((\mathfrak I_{n,m,\eta})_{0\le\eta<1/2}\) be a family of presented
binary-label tasks whose sample law is obtained by applying the homogeneous
flip channel of \cref{def-binary-label-corruption-channel} at rate \(\eta\).
Let \(\mathfrak M_{n,m,\eta}\) be expectation under the declared instance and
sample law.  Fix the scalarization \(\kappa\) used for the saturated cost
invariant of \cref{def-degree-invariant}, and write

\begin{equation}
\label{eq-scalarized-affordable-success-profile}
\operatorname{Succ}^{\kappa}_{\mathfrak M_{n,m,\eta},n}(s)
=
\sup_{\substack{\mathcal A\ \mathrm{represented}\
\kappa(\operatorname{Cost}(\mathcal A))\le s}}
\mathfrak M_{n,m,\eta}
 \!\left(I\mapsto\alpha_I(\mathcal A)\right).
\end{equation}

This is \cref{eq-uniform-affordable-success-profile} restricted by the scalar
cost sublevel set of the already fixed scalarization.  For a target recovery
probability \(0<\alpha\le1\), define

\begin{equation}
\label{eq-noise-indexed-recovery-cost}
B_{\mathrm{rec}}^*(n,m,\eta;\alpha)
=
\inf\left\{
s\ge0:
\operatorname{Succ}^{\kappa}_{\mathfrak M_{n,m,\eta},n}(s)\ge\alpha
\right\},
\end{equation}

with value \(\infty\) when the set is empty.  For a declared scalar budget
\(\sigma(n)\), put

\begin{equation}
\label{eq-noise-indexed-computational-threshold}
\eta_{\mathrm{comp}}^{\mathrm{sat}}
 (n,m;\sigma,\alpha)
=
\sup\left\{
0\le\eta<\frac12:
\operatorname{Succ}^{\kappa}_{\mathfrak M_{n,m,\eta},n}
 \bigl(\sigma(n)\bigr)\ge\alpha
\right\},
\end{equation}

with supremum zero when the set is empty.  This is the success-profile
analogue of the threshold construction in \cref{def-degree-invariant}.  The
success profile in \cref{eq-noise-indexed-recovery-cost} already takes the
supremum over the complete budget-saturated class.

\end{definition}

\begin{proposition}[Degradation monotonicity of saturated success]
\label{thm-noise-degradation-monotonicity}
\label{prop:noise-degradation-monotonicity}

Let \(0\le\eta_1\le\eta_2<1/2\), and set

\begin{equation}
\label{eq-bsc-degradation-rate}
q(\eta_1,\eta_2)
=\frac{\eta_2-\eta_1}{1-2\eta_1}.
\end{equation}

Appending an independent represented Bernoulli-\(q(\eta_1,\eta_2)\) flip to
an \(\eta_1\)-corrupted label produces an \(\eta_2\)-corrupted label.  If
\(\psi_{\mathrm{flip}}(s,n,m)\) is a scalar ceiling for the complete
class-preserving cost of this postprocessing after any computation of
scalarized cost at most \(s\), then

\begin{equation}
\label{eq-noise-degradation-success-monotonicity}
\operatorname{Succ}^{\kappa}_{\mathfrak M_{n,m,\eta_1},n}
 \bigl(\psi_{\mathrm{flip}}(s,n,m)\bigr)
\ge
\operatorname{Succ}^{\kappa}_{\mathfrak M_{n,m,\eta_2},n}(s).
\end{equation}

Consequently every budget ceiling closed under this class-preserving
postprocessing has a downward-closed affordable-noise region.  In particular,
\cref{eq-noise-indexed-computational-threshold} is its right endpoint.

\end{proposition}

\begin{proof}

Let \(Z_1\sim\operatorname{Bernoulli}(\eta_1)\) and
\(Z_q\sim\operatorname{Bernoulli}(q)\) be independent.  Direct substitution of
\cref{eq-bsc-degradation-rate} gives

\[
0\le q(\eta_1,\eta_2)<\frac12,
\]

because \(\eta_1\le\eta_2<1/2\), and

\[
\Pr\{Z_1\oplus Z_q=1\}
=\eta_1+q-2\eta_1q
=\eta_2.
\]

Apply any represented scalar-cost-\(s\) computation for the \(\eta_2\) experiment
after these additional flips.  Its input and output laws are unchanged, while
its complete cost is
at most \(\psi_{\mathrm{flip}}(s,n,m)\).  Taking the supremum over the saturated
budget slice proves \cref{eq-noise-degradation-success-monotonicity}.  Closure
of the declared ceiling under the same postprocessing proves downward
closure.

\end{proof}

\begin{theorem}[Statistical, computational, and information thresholds]
\label{thm-three-noise-threshold-comparison}

Let \(\mathcal T\) be a finite parameter set of size \(N\ge2\), let
\(\Theta\) be uniform on \(\mathcal T\), and, conditional on
\(\Theta=\theta\), let the clean binary target be \(f_\theta\).  Let
\(X_1,\ldots,X_m\) be i.i.d.\ copies of \(X\), independent of \(\Theta\).  The
observed labels are

\[
\widetilde Y_i=f_\Theta(X_i)\Xi_i,
\qquad
\Pr\{\Xi_i=-1\}=\eta,
\]

where the flips are i.i.d.\ and independent of \((\Theta,X_1,\ldots,X_m)\).
Let \(\mathcal H=\{f_\theta:\theta\in\mathcal T\}\),
and assume the clean separation

\begin{equation}
\label{eq-clean-hypothesis-separation}
\inf_{\substack{\theta,\theta'\in\mathcal T\\\theta'\ne\theta}}
\Pr\{f_{\theta'}(X)\ne f_\theta(X)\}
\ge\Delta>0.
\end{equation}

Then the following three statements hold.

\begin{enumerate}[label=\textup{(\roman*)}]
\item Let \(\widehat f\in\mathcal H\) be an
\(\varepsilon_{\mathrm{opt}}\)-empirical minimizer over \(\mathcal H\), and
let \(\Gamma_m^{\mathrm{Rad}}(\mathcal H,\delta)\) be the
simultaneous radius in \cref{eq-corrupted-rademacher-radius}.  With
probability at least \(1-\delta\), exact identification is certified by

\begin{equation}
\label{eq-noise-statistical-identification-threshold}
2\Gamma_m^{\mathrm{Rad}}(\mathcal H,\delta)
+\varepsilon_{\mathrm{opt}}
<
(1-2\eta)\Delta.
\end{equation}

Equivalently, a sufficient statistical-noise threshold is

\begin{equation}
\label{eq-explicit-statistical-noise-threshold}
\eta
<
\frac12-
\frac{2\Gamma_m^{\mathrm{Rad}}(\mathcal H,\delta)
+\varepsilon_{\mathrm{opt}}}{2\Delta}.
\end{equation}

\item Every estimator \(\widehat\Theta=\widehat\Theta(X^m,\widetilde Y^m)\),
without a computational restriction, satisfies

\begin{equation}
\label{eq-noise-information-success-bound}
\Pr\{\widehat\Theta=\Theta\}
\le
\min\left\{
1,
\frac{m(1-H_2(\eta))+1}{\log_2N}
\right\}.
\end{equation}

Hence success at least \(\alpha\) requires

\begin{equation}
\label{eq-noise-information-necessary-condition}
m(1-H_2(\eta))
\ge
\alpha\log_2N-1.
\end{equation}

\item If the right-hand side of
\cref{eq-noise-information-success-bound} is smaller than \(\alpha\), then
\(B_{\mathrm{rec}}^*(n,m,\eta;\alpha)=\infty\).  If
\cref{eq-noise-statistical-identification-threshold} holds and, for some
\(s\ge0\), the represented empirical minimizer has scalarized complete cost
at most \(s\), then
\(B_{\mathrm{rec}}^*(n,m,\eta;1-\delta)\le s\).
Thus the saturated computational threshold lies between an explicitly
represented achievable procedure and the information bound; statistical
identification locates the sample-resolution threshold independently of the
cost of constructing the empirical minimizer.
\end{enumerate}

\end{theorem}

\begin{proof}

By \cref{eq-massart-rademacher-clean-risk}, on the simultaneous event,

\[
R_0(\widehat f)
\le
\frac{2\Gamma_m^{\mathrm{Rad}}(\mathcal H,\delta)
+\varepsilon_{\mathrm{opt}}}{1-2\eta}.
\]

Under \cref{eq-noise-statistical-identification-threshold}, this risk is
strictly smaller than \(\Delta\).  The separation
\cref{eq-clean-hypothesis-separation} then forces
\(\widehat f=f_\Theta\), which proves (i).

For (ii), measure entropy in bits and write
\(S=(X^m,\widetilde Y^m)\).  Since \(X^m\) is independent of \(\Theta\), the
chain rule and the binary output alphabet give

\begin{align*}
I(\Theta;S)
&=I(\Theta;\widetilde Y^m\mid X^m)\\
&\le
\sum_{i=1}^m
\left(
H(\widetilde Y_i\mid X^m,\widetilde Y^{i-1})
-H(\widetilde Y_i\mid\Theta,X^m,\widetilde Y^{i-1})
\right)\\
&\le m(1-H_2(\eta)).
\end{align*}

Let \(p=\Pr\{\widehat\Theta=\Theta\}\).  Fano's entropy argument gives

\[
H(\Theta\mid S)
\le 1+(1-p)\log_2N,
\]

and therefore

\[
p\log_2N-1
\le I(\Theta;S)
\le m(1-H_2(\eta)).
\]

This proves \cref{eq-noise-information-success-bound} and
\cref{eq-noise-information-necessary-condition}.  Part (iii) follows from
the definition of \(B_{\mathrm{rec}}^*\): the information inequality ranges
over all estimators, while a represented empirical minimizer is one member of
the saturated budget slice.

\end{proof}

For a class of \(2^n\) binary targets with separation \(\Delta=1/2\), the
information boundary and the finite-class Rademacher certificate can be drawn
on the same sample-ratio scale.  The latter is a sufficient certificate; the
former bounds every estimator.

\begin{figure}[tbp]
\centering
\begin{tikzpicture}
\pgfmathsetmacro{\StatConstant}{16*(2*sqrt(2*ln(2))+3*sqrt(ln(40)/(2*512)))^2}
\begin{axis}[fw axis,ymode=log,log basis y=10,
  xmin=0,xmax=.48,ymin=1,ymax=1e5,
  xlabel={noise rate $\eta$},ylabel={sample ratio $\gamma=m/n$},
  legend pos=north west]
  \addplot[name path=lower,draw=none,domain=.0001:.48] {1};
  \addplot[name path=itpath,draw=none,domain=.0001:.48,samples=180]
    {1/(1+(x*ln(x)+(1-x)*ln(1-x))/ln(2))};
  \addplot[name path=statpath,draw=none,domain=.0001:.48,samples=180]
    {\StatConstant/(1-2*x)^2};
  \addplot[name path=upper,draw=none,domain=.0001:.48] {1e5};
  \addplot[fill=fwCoral!10,draw=none] fill between[of=lower and itpath];
  \addplot[fill=fwGray!70,draw=none] fill between[of=itpath and statpath];
  \addplot[fill=fwTeal!10,draw=none] fill between[of=statpath and upper];
  \addplot[draw=fwBlue,line width=1.25pt,domain=.0001:.48,samples=180]
    {1/(1+(x*ln(x)+(1-x)*ln(1-x))/ln(2))};
  \addlegendentry{$\gamma_{\rm IT}(\eta)=1/(1-H_2(\eta))$}
  \addplot[draw=fwAmber,line width=1.25pt,domain=.0001:.48,samples=180]
    {\StatConstant/(1-2*x)^2};
  \addlegendentry{$\gamma_{\rm stat}(\eta)$, $n=512$, $\delta=.05$}
  \node[font=\scriptsize,align=center,fwCoral] at (axis cs:.27,1.9)
    {information bound excludes\\success tending to one};
  \node[font=\scriptsize,align=center,fwInk] at (axis cs:.22,24)
    {not ruled out by information;\\certificate not yet active};
  \node[font=\scriptsize,align=center,fwTeal] at (axis cs:.21,2500)
    {exact identification\\certified};
\end{axis}
\end{tikzpicture}
\caption{Information and statistical sample thresholds for \(N=2^n\)
binary targets with clean separation \(1/2\).  The lower curve is the
asymptotic binary-symmetric-channel boundary
\(\gamma(1-H_2(\eta))=1\) from
\cref{eq-noise-information-necessary-condition} at success tending to one.
The upper curve is the sufficient finite-class exact-identification condition
\cref{eq-noise-statistical-identification-threshold}, displayed at
\(n=512\), \(\delta=.05\), and zero optimization error.}
\label{fig-noise-sample-threshold-plane}
\end{figure}

\section{Dependent Dynamical Samples and Sequential Structural Learning}
\label{sec-dependent-dynamical-learning}
\label{sec:dependent-dynamical-learning}

The preceding statistical certificates use independent examples.  A controlled
system instead supplies independent episodes, dependent observations along one
trajectory, or data selected predictably from an interaction history.  The
following results retain the same represented operator, source-projection, and
complete-cost requirements while replacing independent symmetrization by the
appropriate dynamical sampling certificate.

\subsection{Dynamical-sample certificates and blocking}
\label{subsec:dynamical-sample-blocking}

\begin{definition}[Represented dynamical-sample certificate]
\label{def-dynamical-sample-certificate}

Fix a budget-admissible represented loss class
\(\mathcal L_B=\{\ell_h:h\in\mathcal H_B\}\), with
\(0\le \ell_h\le1\).  A \emph{represented dynamical-sample certificate}
is one of the following records, together with its complete construction,
storage, and evaluation cost.

\begin{enumerate}[label=\textup{(\roman*)}]
\item An \emph{independent-episode certificate} consists of independent
episode records \(E_1,\ldots,E_M\) and represented normalized episode losses
\(\ell_h(E_i)\in[0,1]\).

\item A \emph{mixing-trajectory certificate} consists of a stationary
represented process \((Z_t)_{t=1}^N\), its absolute-regularity coefficients
\[
\beta(k)
=
\sup_{t\ge1}
\mathbb E\left[
 \left\|
 \mathcal L(Z_{t+k},Z_{t+k+1},\ldots\mid\mathcal F_t)
 -\mathcal L(Z_{t+k},Z_{t+k+1},\ldots)
 \right\|_{\mathrm{TV}}
\right],
\]
where \(\mathcal F_t=\sigma(Z_1,\ldots,Z_t)\), and represented losses
\(\ell_h(Z_t)\in[0,1]\).  A claimed coefficient is usable only when its
upper bound is part of the certificate.

\item A \emph{predictable-path certificate} consists of a filtration
\((\mathcal F_t)_{t=0}^T\), observations \(Z_t\in\mathcal F_t\), represented
predictable decisions, and losses \(\ell_t(h)\in[0,1]\) such that
\(\ell_t(h)\) is \(\mathcal F_t\)-measurable.  Its conditional risk is
\[
r_t(h)=\mathbb E[\ell_t(h)\mid\mathcal F_{t-1}].
\]
The hypothesis, operator, posterior, or policy used at time \(t\) is
available from \(\mathcal F_{t-1}\), in the temporal sense of
\cref{def-pre-hit-temporally-admissible-source}.
\end{enumerate}

For block length \(a\) in case \textup{(ii)}, put
\[
\mu_a=\left\lfloor\frac{N}{2a}\right\rfloor,
\qquad
\delta_a=\delta-2(\mu_a-1)\beta(a).
\]
The block length is certified at confidence \(1-\delta\) when
\(\mu_a\ge1\) and \(\delta_a>0\).  All optimizations over \(a\), over an
operator family, or over a posterior are performed inside the same saturated
derivation and retain their selection cost.

\end{definition}

\begin{figure}[tbp]
\centering
\resizebox{\linewidth}{!}{%
\begin{tikzpicture}[fw diagram,node distance=7mm and 9mm]
  \node[fw source,text width=2.6cm] (record) {represented\\trajectory data};
  \node[fw box,text width=2.5cm,above right=5mm and 9mm of record] (episode)
    {independent\\episodes};
  \node[fw geom,text width=2.5cm,right=11mm of record] (mix)
    {stationary\\mixing path};
  \node[fw provenance,text width=2.5cm,below right=5mm and 9mm of record] (seq)
    {predictable\\adaptive path};
  \node[fw use,text width=2.7cm,right=15mm of episode] (iid)
    {structural\\Rademacher};
  \node[fw use,text width=2.7cm,right=15mm of mix] (block)
    {coupled blocks\\and \(\mu_a\)};
  \node[fw use,text width=2.7cm,right=15mm of seq] (mart)
    {sequential radius\\or PAC--Bayes};
  \node[fw terminal,text width=2.8cm,right=18mm of block] (risk)
    {population risk,\\regret, deployment};
  \draw[fw arrow] (record) -- (episode);
  \draw[fw arrow] (record) -- (mix);
  \draw[fw arrow] (record) -- (seq);
  \draw[fw arrow] (episode) -- (iid);
  \draw[fw arrow] (mix) -- (block);
  \draw[fw arrow] (seq) -- (mart);
  \draw[fw arrow] (iid) -- (risk);
  \draw[fw arrow] (block) -- (risk);
  \draw[fw arrow] (mart) -- (risk);
\end{tikzpicture}
}
\caption{Each dynamical sampling regime has its own represented dependence
certificate.  Independent symmetrization, mixing blocks, and predictable-path
complexity feed the same typed population-risk coordinates only through their
proved conversion theorem.}
\label{fig-dynamical-sampling-regimes}
\end{figure}

For a block \(B=(Z_{(j-1)a+1},\ldots,Z_{ja})\), write
\(\bar\ell_h(B)=a^{-1}\sum_{Z_t\in B}\ell_h(Z_t)\).  Let
\(\mathfrak R_{\mu_a}^{\mathrm{odd}}(a)\) and
\(\mathfrak R_{\mu_a}^{\mathrm{even}}(a)\) denote the expected Rademacher
complexities of the block-average class on independent copies of the odd and
even block laws, respectively.

\begin{theorem}[Structural blocking bound for a mixing trajectory]
\label{thm-beta-mixing-structural-blocking}

Let case \textup{(ii)} of
\cref{def-dynamical-sample-certificate} hold and define
\[
R(h)=\mathbb E\ell_h(Z_1),
\qquad
\widehat R_N(h)=\frac1N\sum_{t=1}^N\ell_h(Z_t).
\]
For every certified block length \(a\), with probability at least
\(1-\delta\), simultaneously for all \(h\in\mathcal H_B\),
\begin{equation}
\label{eq-beta-mixing-structural-blocking}
\left|R(h)-\widehat R_N(h)\right|
\le
2\max_{p\in\{\mathrm{odd},\mathrm{even}\}}
 \mathfrak R_{\mu_a}^{p}(a)
+3\sqrt{\frac{\log(4/\delta_a)}{2\mu_a}}
+\frac{2a}{N}.
\end{equation}
Consequently the right-hand side may be minimized over the represented set of
certified block lengths.  If the block-average evaluation vectors lie in a
fixed structural energy ball \(\mathcal F_A(R_A)\), the two Rademacher terms
are bounded by
\[
\mathfrak R_{\mu_a}^{p}(a)
\le
\frac{R_A}{\mu_a}\sqrt{\operatorname{Tr}(A^{-1})},
\]
with the block operator and its precision charged as in
\cref{def-structural-energy-ball}.

\end{theorem}

\begin{proof}
Retain the first \(2\mu_a a\) observations and divide them into alternating
blocks of length \(a\).  The omitted suffix has length less than \(2a\), so
boundedness of the loss gives
\[
\sup_h\left|\widehat R_N(h)-\widehat R_{2\mu_a a}(h)\right|
\le \frac{2a}{N}.
\]
Berbee coupling, applied separately to the odd and even block sequences,
couples each sequence to \(\mu_a\) independent blocks with total failure
probability at most \((\mu_a-1)\beta(a)\).  The union of the two coupling
failures therefore has probability at most
\(2(\mu_a-1)\beta(a)\) \citep{Paulin2015}.

Conditional on successful coupling, apply the two-sided independent-sample
Rademacher inequality to each parity class at failure probability
\(\delta_a/2\).  Stationarity makes the expectation of every block average
equal to \(R(h)\).  The truncated empirical mean is the average of the odd
and even block means, so its deviation is bounded by the larger of their two
deviations.  Adding the coupling and independent-sample failure probabilities
gives \(\delta\) and proves
\cref{eq-beta-mixing-structural-blocking}.  The final display follows from
\cref{lem-ellipsoid-rademacher} applied to the block evaluation vector.
\end{proof}

The predictable regime is measured on history trees.  For a sign sequence
\(\epsilon=(\epsilon_1,\ldots,\epsilon_T)\), a depth-\(T\) tree
\(\mathbf z\) has node
\(z_t(\epsilon_{1:t-1})\) at time \(t\).

\subsection{Sequential structural complexity}
\label{subsec:sequential-structural-complexity}

\begin{definition}[Sequential structural Rademacher complexity]
\label{def-sequential-structural-rademacher}

For a represented class \(\mathcal H_B\) and a represented history tree
\(\mathbf z\), define
\begin{equation}
\label{eq-sequential-structural-rademacher}
\mathfrak R_T^{\mathrm{seq}}(\mathcal H_B;\mathbf z)
=
\frac1T\mathbb E_\epsilon
\sup_{h\in\mathcal H_B}
\sum_{t=1}^T
\epsilon_t h_t\bigl(z_t(\epsilon_{1:t-1})\bigr).
\end{equation}
The saturated sequential complexity is the supremum over represented legal
history trees at the declared budget.  A path operator is a predictable tree
\(\mathbf A\) whose realization \(A_\epsilon\succ0\) acts on the length-\(T\)
evaluation vector
\[
h(\epsilon)
=
\bigl(h_t(z_t(\epsilon_{1:t-1}))\bigr)_{t=1}^T.
\]
Its pathwise inverse-precision radius is
\begin{equation}
\label{eq-sequential-inverse-precision-radius}
\mathfrak q_T(\mathbf A)
=
\frac1T\mathbb E_\epsilon
\sqrt{\epsilon^\mathsf T A_\epsilon^{-1}\epsilon}.
\end{equation}
The sequential structural energy ball of radius \(R_A\) consists of the
represented trees satisfying
\(h(\epsilon)^\mathsf T A_\epsilon h(\epsilon)\le R_A^2\) on every path.

\end{definition}

\begin{lemma}[Pathwise support and fixed-operator sequential radius]
\label{lem-sequential-energy-support}

For every path \(\epsilon\) and every positive definite path operator,
\begin{equation}
\label{eq-sequential-energy-support}
\sup_{f^\mathsf T A_\epsilon f\le R_A^2}
\epsilon^\mathsf T f
=
R_A\sqrt{\epsilon^\mathsf T A_\epsilon^{-1}\epsilon}.
\end{equation}
Hence a sequential structural energy ball satisfies
\begin{equation}
\label{eq-sequential-energy-radius-bound}
\mathfrak R_T^{\mathrm{seq}}
\le R_A\mathfrak q_T(\mathbf A).
\end{equation}
If \(A_\epsilon=A\) is fixed over paths, then
\begin{equation}
\label{eq-sequential-fixed-operator-trace}
\mathfrak R_T^{\mathrm{seq}}
\le
\frac{R_A}{T}\sqrt{\operatorname{Tr}(A^{-1})}.
\end{equation}

\end{lemma}

\begin{proof}
Cauchy--Schwarz in the inner product generated by \(A_\epsilon\) gives
\[
\epsilon^\mathsf T f
\le
\sqrt{f^\mathsf T A_\epsilon f}
\sqrt{\epsilon^\mathsf T A_\epsilon^{-1}\epsilon}.
\]
Equality holds at
\(f=R_AA_\epsilon^{-1}\epsilon/
\sqrt{\epsilon^\mathsf T A_\epsilon^{-1}\epsilon}\), proving
\cref{eq-sequential-energy-support}.  The represented class is contained in
the pathwise ellipsoid, so expectation and division by \(T\) prove
\cref{eq-sequential-energy-radius-bound}.  For fixed \(A\), Jensen's
inequality and
\(\mathbb E_\epsilon\epsilon\epsilon^\mathsf T=I\) give
\[
\mathbb E_\epsilon
\sqrt{\epsilon^\mathsf T A^{-1}\epsilon}
\le \sqrt{\operatorname{Tr}(A^{-1})},
\]
which proves \cref{eq-sequential-fixed-operator-trace}.
\end{proof}

\subsection{Learned-operator and PAC--Bayes control}
\label{subsec:sequential-learned-operator-pac-bayes}

\begin{theorem}[Sequential learned-operator cover bound]
\label{thm-sequential-learned-operator-cover}

Let \(\{\mathbf A_\theta:\theta\in\Theta\}\) be a represented predictable
operator family with \(A_{\theta,\epsilon}\succeq\eta I\) on every path and
\[
\sup_\epsilon
\|A_{\theta,\epsilon}^{-1}-A_{\theta',\epsilon}^{-1}\|_{\mathrm{op}}
\le L_A d(\theta,\theta').
\]
For the union of its radius-\(R_A\) sequential energy balls,
\begin{equation}
\label{eq-sequential-learned-operator-cover}
\mathfrak R_T^{\mathrm{seq}}
\le
\inf_{\rho>0}\left\{
R_A\sup_{\theta\in\Theta}\mathfrak q_T(\mathbf A_\theta)
+R_A\sqrt{\frac{L_A\rho}{T}}
+\frac{4R_A}{\sqrt\eta}
 \sqrt{\frac{\log\mathcal N(\Theta,d,\rho)}{T}}
\right\}.
\end{equation}
For fixed operators, the first term may be replaced by
\(R_AT^{-1}\sup_\theta\sqrt{\operatorname{Tr}(A_\theta^{-1})}\).
The covering number, the construction of its net, and the selected operator
remain part of the saturated cost.

\end{theorem}

\begin{proof}
Fix a \(\rho\)-net.  On every path, the support identity and
\(|\sqrt u-\sqrt v|\le\sqrt{|u-v|}\) imply that replacing \(\theta\) by
a net point \(\theta_0\) changes the normalized support by at most
\[
\frac{R_A}{T}
\sqrt{\|A_{\theta,\epsilon}^{-1}
-A_{\theta_0,\epsilon}^{-1}\|_{\mathrm{op}}\|\epsilon\|_2^2}
\le R_A\sqrt{\frac{L_A\rho}{T}}.
\]
For each net point, every coordinate of an evaluation vector has absolute
value at most \(R_A/\sqrt\eta\).  The finite-union lemma for sequential
Rademacher complexity therefore contributes at most
\(4R_A\eta^{-1/2}\sqrt{\log\mathcal N(\Theta,d,\rho)/T}\).
Adding the fixed-net-point radius, the approximation term, and the finite-net
term, then minimizing over \(\rho\), proves
\cref{eq-sequential-learned-operator-cover}.  The trace specialization follows
from \cref{lem-sequential-energy-support}.
\end{proof}

The existing structural prior also admits a predictable-path certificate.
The next theorem uses the same proper prior and the same charged structural
KL as \cref{def-structural-pac-bayes-kl}.

\begin{theorem}[Martingale structural PAC--Bayes bound]
\label{thm-martingale-structural-pac-bayes}

Let case \textup{(iii)} of
\cref{def-dynamical-sample-certificate} hold.  Let
\(P=P^{\mathsf{str}}_{E,H}\) be fixed before the evaluated innovations, or be
constructed from public geometry or an independent sample, and let \(Q\) be
any posterior with \(Q\ll P\).  In this theorem write
\[
\mathrm{KL}^{\mathsf{str}}_{E,H}(Q)
:=\mathrm{KL}(Q\Vert P);
\]
for Gaussian \(Q\), this is exactly the formula in
\cref{def-structural-pac-bayes-kl}.  For any \(\lambda>0\), with probability at
least \(1-\delta\), simultaneously for all such \(Q\),
\begin{equation}
\label{eq-martingale-structural-pac-bayes}
\frac1T\sum_{t=1}^T
\mathbb E_{h\sim Q}\bigl[r_t(h)-\ell_t(h)\bigr]
\le
\frac{\mathrm{KL}^{\mathsf{str}}_{E,H}(Q)+\log(1/\delta)}{\lambda T}
+\frac{\lambda}{8}.
\end{equation}
More generally, suppose
\(|r_t(h)-\ell_t(h)|\le c\) and put
\[
V_T(h)=\sum_{t=1}^T
\mathbb E\left[(r_t(h)-\ell_t(h))^2\mid\mathcal F_{t-1}\right].
\]
For \(0<\lambda<3/c\), the predictable-variance refinement is
\begin{equation}
\label{eq-martingale-structural-pac-bayes-bernstein}
\frac1T\sum_{t=1}^T
\mathbb E_Q[r_t(h)-\ell_t(h)]
\le
\frac{\mathrm{KL}^{\mathsf{str}}_{E,H}(Q)+\log(1/\delta)}{\lambda T}
+\frac{\lambda\,\mathbb E_QV_T(h)}{2T(1-\lambda c/3)}.
\end{equation}
Any data-dependent choice of \(\lambda\) is made over a represented countable
grid and includes its union or prior weight in the logarithmic term.

\end{theorem}

\begin{proof}
For fixed \(h\), the increments
\(X_t(h)=r_t(h)-\ell_t(h)\) have conditional mean zero.  Conditional
Hoeffding's lemma gives
\[
\mathbb E\left[
 \exp\left(\lambda X_t(h)-\lambda^2/8\right)
 \middle|\mathcal F_{t-1}
\right]\le1.
\]
Iteration shows that the corresponding product is a supermartingale of mean
at most one.  Integrate it with respect to \(P\), apply Markov's inequality,
and then use the Donsker--Varadhan change-of-measure inequality
\[
\mathbb E_QF
\le \mathrm{KL}(Q\Vert P)+\log\mathbb E_Pe^F.
\]
The resulting event is simultaneous in \(Q\) and yields
\cref{eq-martingale-structural-pac-bayes} after division by \(\lambda T\).
The KL is precisely \(\mathrm{KL}^{\mathsf{str}}_{E,H}\) because the prior
is the proper defect-restricted prior of
\cref{def-structural-pac-bayes-kl}.

For the second assertion, the conditional Bernstein exponential inequality
gives, for \(0<\lambda<3/c\),
\[
\mathbb E\exp\left\{
\lambda\sum_tX_t(h)
-\frac{\lambda^2V_T(h)}{2(1-\lambda c/3)}
\right\}\le1.
\]
The same integration, Markov, and change-of-measure steps prove
\cref{eq-martingale-structural-pac-bayes-bernstein}.  A weighted union bound
over a represented grid proves the final assertion
\citep{SeldinEtAl2012}.
\end{proof}

\subsection{Online regret and deployment}
\label{subsec:online-regret-deployment}

\begin{theorem}[Sequential structural complexity and online regret]
\label{thm-structural-online-regret}

Let \(\mathcal H_B\subseteq[-1,1]^{\mathcal Z}\) be represented and let
\(\ell(\widehat y,y)\in[0,1]\) be convex and \(L_\ell\)-Lipschitz in
\(\widehat y\).  The minimax expected regret against \(\mathcal H_B\) over
represented depth-\(T\) histories satisfies
\begin{equation}
\label{eq-structural-online-regret}
\mathcal V_T(\mathcal H_B)
\le
2L_\ell T
\sup_{\mathbf z}
\mathfrak R_T^{\mathrm{seq}}(\mathcal H_B;\mathbf z).
\end{equation}
If the minimax strategy is implemented through a represented finite tree,
cover, or oracle in the declared presentation, its complete implementation
cost is charged.  Thus
\cref{lem-sequential-energy-support,thm-sequential-learned-operator-cover}
give budget-admissible regret certificates exactly when that represented
implementation lies within the budget.

\end{theorem}

\begin{proof}
Apply the minimax theorem at each history node and symmetrize the difference
between the learner's cumulative loss and the best fixed
\(h\in\mathcal H_B\) by an independent tangent history.  Sequential
contraction for an \(L_\ell\)-Lipschitz loss bounds the symmetrized loss tree
by \(L_\ell\) times the prediction tree, and the two tangent copies produce
the factor two.  The normalized prediction-tree supremum is
\cref{eq-sequential-structural-rademacher}; multiplication by \(T\) proves
\cref{eq-structural-online-regret} \citep{RakhlinSridharanTewari2015}.
The statistical minimax inequality specifies a strategy on the history tree.
The framework admits it as a computation only through a represented
implementation, whose tree, cover, or declared oracle and evaluation cost are
then included by saturated no-creation and computation closure.
\end{proof}

\begin{corollary}[Structural online-to-batch deployment]
\label{cor-structural-online-to-batch}

Suppose each observation is conditionally drawn from a common population law
independently of \(\mathcal F_{t-1}\), and an online learner produces
predictable \(h_1,\ldots,h_T\) with expected regret \(\mathcal R_T\).  If the
loss is convex, the represented average
\(\bar h_T=T^{-1}\sum_{t=1}^Th_t\) satisfies
\begin{equation}
\label{eq-structural-online-to-batch}
\mathbb E R(\bar h_T)-\inf_{h\in\mathcal H_B}R(h)
\le \frac{\mathbb E\mathcal R_T}{T}.
\end{equation}
Alternatively, choosing an index uniformly from \(\{1,\ldots,T\}\) gives the
same bound for the expected risk of the selected iterate without convexity.
Every retained iterate, averaging operation, and deployment evaluator keeps
its complete represented cost.

\end{corollary}

\begin{proof}
Independence of the next observation from the predictable iterate gives
\(\mathbb E\ell(h_t,Z_t)=\mathbb ER(h_t)\).  Taking expectations in the
regret inequality and dividing by \(T\) bounds the average of
\(R(h_t)\).  Convexity gives
\(R(\bar h_T)\le T^{-1}\sum_tR(h_t)\).  Uniform selection has risk exactly
that average.  These two observations prove
\cref{eq-structural-online-to-batch}.
\end{proof}

The execution-law separation of population structure from retained sample
records extends from independent examples to complete trajectories.

\subsection{Population alignment and source conditions}
\label{subsec:trajectory-population-alignment}

\begin{theorem}[Trajectory population alignment and memorization separation]
\label{thm-trajectory-population-memorization-separation}

Let \(\mathcal A\) be one family-uniform represented learner trained from any
certificate in \cref{def-dynamical-sample-certificate}, and let
\(h_{I,\mathsf S,\omega}\) denote its retained output.  Evaluate population
alignment on a fresh episode or fresh legal continuation \(E'\), independent
of the learner's private seed conditional on the declared public instance:
\[
c_I^{\mathrm{dyn}}(h)
=
\mathbb E\bigl[h(E')f_{*,I}(E')\mid I\bigr].
\]
Then
\begin{equation}
\label{eq-trajectory-population-projection}
|c_I^{\mathrm{dyn}}(h)|^2
\le
\operatorname{Var}(h(E')\mid I)
\operatorname{ProjMass}_{\sigma(h)}(T_I),
\end{equation}
with the target event interpreted on the fresh execution carrier.  Averaging
over the training trajectory and seed, and applying the monotone family
functional, gives the same alignment--projection domination as
\cref{eq-average-seed-alignment-projection-domination}.

A table indexed only by training times, visited states, realized innovations,
or one completed trajectory remains a charged lifted/transcript record.  Its
component orthogonal to the fresh-execution population projection lies in
\(\Res\) by \cref{thm-exhaustive-residual-normal-form}.  Any represented
family-uniform use with nonzero fresh-execution target projection is charged
to the existing qualified \(\Geom/\Abs\) source ledger at its complete
training, storage, selection, and deployment cost.  Thus trajectory
memorization contributes to population alignment only through the represented
rule that transfers it to fresh executions.

\end{theorem}

\begin{proof}
On the fresh execution carrier, center \(h(E')\) and the target indicator.
Orthogonal projection of the centered target indicator onto
\(L^2(\sigma(h))\), followed by Cauchy--Schwarz, gives
\cref{eq-trajectory-population-projection}; it is the same Hilbert-space
calculation used in
\cref{thm-population-alignment-generalization-memorization}.  Expectation over
the training record and learner seed preserves the pointwise inequality, and
monotonicity of the family functional gives the averaged statement.

The complete training path is a represented lifted computation record by
\cref{thm-full-saturated-profile-decomposition}.  Applying
\cref{thm-exhaustive-residual-normal-form} to its fresh-execution use splits
the represented population projection from its orthogonal component.  The
former retains its qualified \(\Geom/\Abs\) provenance; the latter is
\(\Res\).  If an affordable family-uniform postprocessing of the latter has a
nonzero fresh-execution projection, then
\cref{prop-family-uniform-residual-reclassification} reclassifies that
represented postprocessing at its complete cost.  This proves the provenance
and memorization claims.
\end{proof}

Finally, the effective-dimension/source calculation depends on the sampling
regime only through the certified variance scale.

\begin{corollary}[Dynamical source-condition rate at certified effective
sample size]
\label{cor-dynamical-source-condition-rate}

Assume the source and eigenvalue hypotheses of
\cref{cor-source-condition-rate}.  Let \(m_{\mathrm{dyn}}>0\) be a represented
effective sample size for which the declared dynamical covariance and noise
certificate proves the fixed-kernel oracle inequality
\begin{equation}
\label{eq-dynamical-krr-effective-sample-oracle}
\mathbb E\|\widehat f_\lambda-f_*\|_{L^2(\mu)}^2
\le
C\left(
\frac{\sigma^2N_{\mathrm{eff}}(\lambda)}{m_{\mathrm{dyn}}}
+B_k(\lambda;f_*)
\right).
\end{equation}
For independent normalized episodes one may take
\(m_{\mathrm{dyn}}=M\).  For a certified mixing block construction one may
take \(m_{\mathrm{dyn}}=\mu_a\) when the represented block-average feature
and covariance operators satisfy the empirical-operator concentration event
used to establish \eqref{eq-dynamical-krr-effective-sample-oracle}; its coupling tail and
truncation error from \cref{thm-beta-mixing-structural-blocking} are retained
separately.  For a predictable martingale-noise certificate with variance
proxy \(0<V_T<\infty\), the variance normalization supplies
\(m_{\mathrm{dyn}}=\sigma^2T^2/V_T\) whenever the represented empirical
covariance concentration event used to establish
\eqref{eq-dynamical-krr-effective-sample-oracle} is certified.

Then
\begin{equation}
\label{eq-dynamical-source-condition-rate}
\lambda_{\mathrm{dyn}}
\asymp
m_{\mathrm{dyn}}^{-1/(2r+1/(2\beta))},
\qquad
\mathbb E\|\widehat f_{\lambda_{\mathrm{dyn}}}-f_*\|_{L^2}^2
\lesssim
m_{\mathrm{dyn}}^{-4\beta r/(4\beta r+1)},
\end{equation}
plus the displayed mixing coupling/truncation term or any represented
operator-selection term.  The effective sample size changes statistical
resolution; it leaves the saturated derivation, target projection, and
operator-selection charges unchanged.

\end{corollary}

\begin{proof}
The source condition gives
\(B_k(\lambda;f_*)\le C_rR_s^2\lambda^{2r}\), and the eigenvalue law gives
\(N_{\mathrm{eff}}(\lambda)\asymp\lambda^{-1/(2\beta)}\).  Substitution in
\cref{eq-dynamical-krr-effective-sample-oracle} yields
\[
C'\left(
\lambda^{2r}
+m_{\mathrm{dyn}}^{-1}\lambda^{-1/(2\beta)}
\right).
\]
Balancing the two terms proves
\cref{eq-dynamical-source-condition-rate}.  Independent episodes give the
usual variance factor \(M^{-1}\).  Coupled blocks give \(\mu_a\) independent
block observations on the successful coupling event.  For martingale noise,
the normalized average has conditional variance at most \(V_T/T^2\), which
equals \(\sigma^2/m_{\mathrm{dyn}}\) by definition.  The stated empirical
covariance event converts that scalar variance certificate into the kernel
oracle inequality.  The coupling, truncation, and selection terms arise
outside the bias--variance balance and are therefore retained additively.
\end{proof}

\section{Dynamical Label Corruption at Saturated Computational Capacity}
\label{subsec-dynamical-label-corruption}
\label{sec:dynamical-label-corruption}

The binary corruption channel of
\cref{def-binary-label-corruption-channel} extends to episode, mixing, and
predictable-path records by conditioning at the time at which each label is
revealed.  The resulting clean-risk estimates retain both the sample size and
the complete computational ceiling.

\subsection{Complete dynamical label-noise records}
\label{subsec:dynamical-complete-label-noise-records}

\begin{definition}[Complete dynamical label-noise record]
\label{def-complete-dynamical-label-noise-record}

Fix a problem size \(n\), a represented computational ceiling \(b(n)\), and
\(m\) binary observations.  A complete dynamical label-noise record consists
of the following data.

\begin{enumerate}[label=\textup{(\roman*)}]
\item A filtration \((\mathcal F_t)_{t=0}^{m}\), represented inputs \(X_t\),
      and clean labels
      \(Y_t=f_{*,t}(X_t)\in\{-1,1\}\).  The rule \(f_{*,t}\) is fixed by the
      presented task and the public pre-\(t\) record.
\item Predictable represented binary decision rules
      \(h_t:\mathcal X_t\to\{-1,1\}\).  The random rule \(h_t\), including
      every private random bit used to choose it, is
      \(\mathcal F_{t-1}\)-measurable; it is evaluated at \(X_t\) before
      \(Y_t\) or its corruption is revealed.  Real-valued scores are treated
      separately through the flip-symmetric linear loss below.
\item Corrupted labels
      \[
      \widetilde Y_t=Y_t\Xi_t,\qquad
      \Pr\{\Xi_t=-1\mid\mathcal F_{t-1},X_t\}
      =\eta_t(X_t),
      \]
      with a represented rate certificate
      \(0\le\eta_t(X_t)\le\eta_+<1/2\).  Put
      \(\rho_{\min}=1-2\eta_+>0\).
\item One dynamical-sample certificate from
      \cref{def-dynamical-sample-certificate}, including every mixing,
      covariance, predictable-variance, operator, or history-tree constant
      used below.
\item The clean-label semantics, corruption sampler, hypothesis or policy
      class, selected operator, empirical optimization, retained state,
      debiasing, deployment, comparison, confidence allocation, and numerical
      precision, all at complete saturated cost at most \(b(n)\).
\end{enumerate}

In the homogeneous case the record contains an exact represented
\(\eta_t\equiv\eta\), and \(\rho=1-2\eta\).  A certified interval
\(\eta\in[\eta_-,\eta_+]\) permits the Massart conclusions with
\(\rho_{\min}\); an information upper bound uses the lower endpoint
\(\eta_-\).  For a stationary mixing record, the admissible state-dependent
extension has \(\eta_t(X_t)=\eta(X_t)\).  There are iid uniforms \(U_t\),
independent of the entire stationary clean process, such that \(\Xi_t\) is a
measurable function of \((X_t,Y_t,U_t)\).  History-dependent rates use the
predictable-path branch.

The record applies to binary target, failure, classification,
committor-training, and binary policy-supervision labels.  Reward,
transition, and state-observation discrepancies retain the model-error
coordinates of \cref{def-complete-dynamical-learning-record}.
The clean rule \(f_{*,t}\) is an analytic task comparator.  It is not an
available pre-reveal source or policy input unless a separate represented
derivation constructs it and pays its complete cost.

\end{definition}

\begin{theorem}[Conditional transport through dynamical label noise]
\label{thm-dynamical-label-noise-conditional-transport}

Let \(h_t\) belong to a complete dynamical label-noise record and define
\begin{align*}
r_{0,t}(h_t)
&=\mathbb E\!\left[
 \mathbf1_{\{h_t(X_t)\ne Y_t\}}\mid\mathcal F_{t-1}\right],\\
r_{{\rm corr},t}(h_t)
&=\mathbb E\!\left[
 \mathbf1_{\{h_t(X_t)\ne\widetilde Y_t\}}
 \mid\mathcal F_{t-1}\right].
\end{align*}
Then
\begin{equation}
\label{eq-dynamical-label-noise-conditional-transport}
r_{{\rm corr},t}(h_t)-r_{{\rm corr},t}(f_{*,t})
=
\mathbb E\!\left[
(1-2\eta_t(X_t))
\mathbf1_{\{h_t(X_t)\ne f_{*,t}(X_t)\}}
\mid\mathcal F_{t-1}\right]
\ge \rho_{\min}r_{0,t}(h_t).
\end{equation}
Consequently,
\begin{equation}
\label{eq-dynamical-label-noise-average-transport}
\frac1m\sum_{t=1}^m r_{0,t}(h_t)
\le
\frac{1}{\rho_{\min}m}
\sum_{t=1}^m
\bigl[
r_{{\rm corr},t}(h_t)-r_{{\rm corr},t}(f_{*,t})
\bigr].
\end{equation}

For a bounded flip-symmetric loss,
\begin{equation}
\label{eq-dynamical-flip-symmetric-loss-transport}
\mathbb E[
\ell(h_t(X_t),\widetilde Y_t)\mid
\mathcal F_{t-1},X_t,Y_t]
=
\eta_t(X_t)
+(1-2\eta_t(X_t))\ell(h_t(X_t),Y_t).
\end{equation}
Under homogeneous noise,
\cref{eq-dynamical-label-noise-conditional-transport,eq-dynamical-label-noise-average-transport}
are equalities after replacing \(\rho_{\min}\) by \(\rho\).

If the clean augmented process is stationary and absolutely regular with
coefficients \(\beta_{\rm clean}(k)\), and the stationary state-dependent
corruption uses the independent iid innovations specified in
\cref{def-complete-dynamical-label-noise-record}, then the observed process
\((X_t,\widetilde Y_t)_{t\in\mathbb Z}\) satisfies
\begin{equation}
\label{eq-corrupted-process-beta-contraction}
\beta_{\rm corr}(k)\le\beta_{\rm clean}(k).
\end{equation}

\end{theorem}

\begin{proof}

Condition on \(\mathcal F_{t-1},X_t\).  Correct prediction incurs corrupted
loss \(\eta_t(X_t)\), while incorrect prediction incurs
\(1-\eta_t(X_t)\).  Subtraction of the clean comparator's corrupted risk
gives
\cref{eq-dynamical-label-noise-conditional-transport}; summation proves
\cref{eq-dynamical-label-noise-average-transport}.  The identity
\(\ell(z,-y)=1-\ell(z,y)\) gives
\cref{eq-dynamical-flip-symmetric-loss-transport}.

For the mixing assertion, adjoin an independent innovation sequence to the
clean stationary process.  Independent products do not increase absolute
regularity, and the corrupted observation is a measurable image of that
product.  Total-variation distance contracts under measurable maps, which
proves \cref{eq-corrupted-process-beta-contraction}.

\end{proof}

\subsection{Saturated noise and sample profiles}
\label{subsec:dynamical-saturated-noise-sample-profiles}

\begin{definition}[Saturated dynamical noise and sample profiles]
\label{def-saturated-dynamical-noise-profiles}

The complete profile consists of the affordable output and generalization
coordinates below together with the corrupted-excess coordinate of
\cref{def:dynamical-corrupted-excess-profile}.  A finite-codebook bound for
the first coordinate is given separately in
\cref{prop:dynamical-finite-description-radius}.

Let \(\mathscr A_{n,b,m}^{\rm dyn}(\eta_+)\) be the family of complete
dynamical label-noise records of cost at most \(b(n)\) and certified rates at
most \(\eta_+\).  Let
\(\mathcal H_{n,b,m}^{\rm dyn,sat}(\eta_+)\) be the sample-independent hull
of all their complete outputs.  Its \(0\)--\(1\) branch contains the binary
decisions in \cref{def-complete-dynamical-label-noise-record}; its
flip-symmetric linear-loss branch contains represented scores
\(h:\mathcal X\to[-1,1]\) and is never inserted into a mismatch indicator
without a represented decision rule.
On the fresh-execution carrier of
\cref{thm-trajectory-population-memorization-separation}, define
\begin{equation}
\label{eq-saturated-dynamical-projection-envelope}
M_{n,b,m}^{\rm dyn,sat}
(\eta_+)
=
\sup_{\substack{I\in\mathfrak I_n\\
 h\in\mathcal H_{n,b,m}^{\rm dyn,sat}(I;\eta_+)}}
\operatorname{ProjMass}_{\sigma(h)}(T_I),
\end{equation}
with value zero when the indexed class is empty.  The training sample and
private seed range inside the fixed hull before the supremum is evaluated;
the target projection is always measured on an independent fresh execution.
For each \(\mathcal A\in\mathscr A_{n,b,m}^{\rm dyn}(\eta_+)\), let
\(\eta_{\mathcal A,+}\) be its own certified rate ceiling, and let
\(\Gamma_{\mathcal A}(m,\delta)\) be the minimum of the represented
two-sided upper bounds applicable to that record for the absolute
population-minus-empirical corrupted-loss deviation.  A minimum is allowed
only after the candidates have been
converted to the same population law and loss and placed on one declared
simultaneous event, with confidence splitting over every data-dependent
block length, operator, posterior, or model choice.  Eligible candidates are
the independent-episode radius, the mixing-block radius of
\cref{thm-beta-mixing-structural-blocking}, and the represented sequential or
martingale comparisons of
\cref{thm-sequential-learned-operator-cover,thm-martingale-structural-pac-bayes}.
The recordwise value is \(+\infty\) when no candidate applies.  Define
\begin{equation}
\label{eq-saturated-dynamical-generalization-radius}
\Gamma_{\rm dyn}^{\rm sat}(n,b(n),m,\eta_+,\delta)
=
\sup_{\mathcal A\in\mathscr A_{n,b,m}^{\rm dyn}(\eta_+)}
\Gamma_{\mathcal A}(m,\delta).
\end{equation}

\end{definition}

\begin{definition}[Saturated dynamical corrupted-excess profile]
\label{def:dynamical-corrupted-excess-profile}

Fix \((n,b,m,\eta_+,\delta)\) and the affordable record class and
simultaneous generalization events of
\cref{def-saturated-dynamical-noise-profiles}.

For each record, define \(\operatorname{Exc}_{\mathcal A}\) as the minimum
of the following applicable upper bounds on its average corrupted excess
relative to the analytic clean task rule, on that same event:
\begin{enumerate}[label=\textup{(\alph*)}]
\item \(2\Gamma_{\mathcal A}+\varepsilon_{\rm opt}
      +\varepsilon_{\rm app,corr}\) for a represented empirical minimizer,
      where \(\varepsilon_{\rm app,corr}\) is the corrupted-risk gap between
      the best admissible represented comparator and the analytic task rule;
\item the normalized represented corrupted-loss regret plus its certified
      conditional-to-realized comparison for a predictable online learner,
      plus the chosen comparator's corrupted-risk gap to the task rule;
\item the represented posterior-minus-comparator empirical excess plus both
      applicable martingale PAC--Bayes radii and that comparator's
      corrupted-risk gap to the task rule.
\end{enumerate}
The comparator is used only for analysis unless it is independently
constructed by the learner; every optimization error, prior, confidence
split, and selection operation belongs to the complete record.  The
recordwise excess is \(+\infty\) when none of the three candidates applies.
Because each recordwise bound uses \(\eta_{\mathcal A,+}\), not the outer
ceiling, enlarging \(\eta_+\) only enlarges the nested affordable class.
Thus the saturated profiles below are nondecreasing in \(\eta_+\).  The
saturated corrupted-excess envelope is
\begin{equation}
\label{eq-saturated-dynamical-corrupted-excess-envelope}
\operatorname{Exc}_{\rm corr}^{\rm sat}
(n,b(n),m,\eta_+,\delta)
=
\sup_{\mathcal A\in\mathscr A_{n,b,m}^{\rm dyn}(\eta_+)}
\operatorname{Exc}_{\mathcal A}(m,\delta).
\end{equation}
The supremum, rather than a cross-regime minimum, makes the bound uniform
over every affordable sample-certificate type.  Define
\begin{equation}
\label{eq-saturated-dynamical-clean-noise-bound}
U_{\rm noise}^{\rm clean}(n,b(n),m,\eta_+,\delta)
=
\min\left\{
1,\,
\frac{
\operatorname{Exc}_{\rm corr}^{\rm sat}
 (n,b(n),m,\eta_+,\delta)}
{1-2\eta_+}
\right\}.
\end{equation}

\end{definition}

\subsection{Finite-description generalization and clean-risk transport}
\label{subsec:dynamical-generalization-clean-risk}

\begin{proposition}[Finite-description dynamical generalization radius]
\label{prop:dynamical-finite-description-radius}

Work in the setting of
\cref{def-saturated-dynamical-noise-profiles} at fixed
\((n,b,m,\eta_+,\delta)\).

Let
\[
N_{n,b,m}^{\rm dyn,out}(\eta_+)
=
\sup_{I\in\mathfrak I_n}
|\mathcal H_{n,b,m}^{\rm dyn,sat}(I;\eta_+)|,
\]
and set this value to \(+\infty\) unless the fixed-budget output codebook is
finite.  Assume the relevant loss takes values in \([0,1]\), the codebook is
fixed before the effective sample is revealed, and one represented sample
certificate supplies a simultaneous comparison to \(m_{\rm eff}\)
independent normalized observations with additive error
\(\varepsilon_{\rm dep}\) and confidence \(\delta_{\rm eff}\).  Put
\[
L_{\rm out}(n,b(n),m,\eta_+)
=\log N_{n,b,m}^{\rm dyn,out}(\eta_+).
\]
For any candidate with represented effective sample size \(m_{\rm eff}\),
effective confidence \(\delta_{\rm eff}\), and additive dependence error
\(\varepsilon_{\rm dep}\), the finite-description fallback is
\begin{equation}
\label{eq-dynamical-noise-finite-description-radius}
\Gamma_{\rm dyn}^{\rm sat}
\le
2\sqrt{\frac{2L_{\rm out}(n,b(n),m,\eta_+)}{m_{\rm eff}}}
+3\sqrt{\frac{\log(2/\delta_{\rm eff})}{2m_{\rm eff}}}
+\varepsilon_{\rm dep}.
\end{equation}
Here \(m_{\rm eff}=m\) for independent normalized observations,
\(m_{\rm eff}=\mu_a\),
\(\delta_{\rm eff}=\delta_a/2\), and
\(\varepsilon_{\rm dep}=2a/m\) for the certified mixing-block candidate,
and the predictable candidate uses its declared history-tree or
predictable-variance normalization.  The displayed fallback applies when
the effective-size and confidence quantities hold uniformly over the whole
affordable class.  Otherwise it is applied within each sample-certificate
subclass and the resulting subclass bounds are combined by a supremum.

\end{proposition}

\begin{proof}

For a fixed finite codebook, Massart's finite-class lemma bounds its empirical
Rademacher complexity by
\(\sqrt{2L_{\rm out}/m_{\rm eff}}\).  The two-sided empirical
Rademacher inequality for a \([0,1]\)-valued loss class therefore gives the
first two terms in
\cref{eq-dynamical-noise-finite-description-radius} simultaneously over the
codebook.  The represented sample certificate transfers that independent
comparison to the declared dynamical law and adds
\(\varepsilon_{\rm dep}\).  Taking the supremum over record subclasses is
valid only after their confidence events have been made simultaneous, which
is exactly the final hypothesis.

\end{proof}

\begin{theorem}[Uniform dynamical clean-risk and alignment bounds]
\label{thm-uniform-dynamical-label-noise-bound}

Every represented learner in
\(\mathcal H_{n,b,m}^{\rm dyn,sat}(\eta_+)\) whose corrupted excess is
bounded by
\(\operatorname{Exc}_{\rm corr}^{\rm sat}\) satisfies, on the simultaneous
event declared by its complete record,
\begin{equation}
\label{eq-uniform-dynamical-clean-risk-bound}
\frac1m\sum_{t=1}^m r_{0,t}(h_t)
\le
U_{\rm noise}^{\rm clean}(n,b(n),m,\eta_+,\delta).
\end{equation}
For a homogeneous balanced target, a represented score
\(h:X\to[-1,1]\), and
\[
c_0(h)=\mathbb E[h(X)f_*(X)],
\qquad
\widehat c_{\eta}(h)=\frac1m\sum_{t=1}^m h(X_t)\widetilde Y_t,
\]
the same event gives
\begin{equation}
\label{eq-uniform-dynamical-clean-alignment-bound}
|c_0(h)|
\le
\min\left\{
\sqrt{M_{n,b,m}^{\rm dyn,sat}(\eta)},
\frac{
|\widehat c_{\eta}(h)|
+2\Gamma_{\rm dyn}^{\rm sat}(n,b(n),m,\eta,\delta)}
{1-2\eta}
\right\},
\end{equation}
where \(M_{n,b,m}^{\rm dyn,sat}(\eta)\) is the fresh-execution
generated-algebra
projection envelope of
\cref{thm-trajectory-population-memorization-separation}.

For homogeneous noise and a dynamical kernel-ridge certificate with
effective sample size \(m_{\rm dyn}\), let \(\widehat g_\lambda\) be kernel
ridge trained on the corrupted labels and define the represented debiased
estimator \(\widehat f_\lambda=\widehat g_\lambda/(1-2\eta)\).  Then
\begin{equation}
\label{eq-dynamical-label-noise-krr-bound}
\mathbb E\|\widehat f_\lambda-f_*\|_{L^2(\mu)}^2
\le
C\left(
\frac{v_\eta^2N_{\rm eff}(\lambda)}
     {(1-2\eta)^2m_{\rm dyn}}
+B_k(\lambda;f_*)
\right)
+\frac{\varepsilon_{\rm dep}+\varepsilon_{\rm sel}}
       {(1-2\eta)^2}.
\end{equation}
The dependence and selection errors are the represented additive terms of
\cref{cor-dynamical-source-condition-rate}.

\end{theorem}

\begin{proof}

The definition of the corrupted excess envelope and
\cref{eq-dynamical-label-noise-average-transport} give
\cref{eq-uniform-dynamical-clean-risk-bound}.  In the homogeneous
flip-symmetric linear loss,
\[
R_\eta^\ell(h)=\frac{1-(1-2\eta)c_0(h)}2,
\qquad
\widehat R_\eta^\ell(h)=\frac{1-\widehat c_\eta(h)}2.
\]
The simultaneous deviation event therefore gives the second entry of
\cref{eq-uniform-dynamical-clean-alignment-bound}; the first is
\cref{eq-trajectory-population-projection}.  For kernel ridge, apply
\cref{eq-dynamical-krr-effective-sample-oracle} to the regression target
\((1-2\eta)f_*\), define the displayed debiased estimator, divide the full
oracle inequality by \((1-2\eta)^2\), and use quadratic homogeneity of
the source residual.  This proves
\cref{eq-dynamical-label-noise-krr-bound}.

\end{proof}

\subsection{Dynamical threshold scales}
\label{subsec:dynamical-threshold-scales}

\begin{definition}[Dynamical sample, statistical-noise, and computational-noise thresholds]
\label{def-dynamical-noise-thresholds}

For \(0<\varepsilon,\alpha\le1\), define
\begin{align}
m_{\rm noise}^{\rm sat}
(n,b,\eta_+,\varepsilon,\delta)
&=
\inf\left\{
m_0\ge1:
\sup_{m'\ge m_0}
\operatorname{Exc}_{\rm corr}^{\rm sat}
 (n,b(n),m',\eta_+,\delta)
\le(1-2\eta_+)\varepsilon
\right\},
\label{eq-dynamical-noise-sample-threshold}\\
\eta_{\rm stat,dyn}^{\rm sat}
(n,m;b,\varepsilon,\delta)
&=
\sup\left\{
0\le\eta<\frac12:
\operatorname{Exc}_{\rm corr}^{\rm sat}
 (n,b(n),m,\eta,\delta)
\le(1-2\eta)\varepsilon
\right\}.
\label{eq-dynamical-statistical-noise-threshold}
\end{align}
Infima of empty sets are \(+\infty\), and suprema of empty sets are zero.

Let
\(
\mathscr A_{n,b,m}^{\rm dyn,hom}(\eta)
\subseteq\mathscr A_{n,b,m}^{\rm dyn}(\eta)
\)
be the complete records evaluated in the exact homogeneous experiment
\(\eta_t\equiv\eta\).  Define
\begin{equation}
\label{eq-dynamical-noise-success-profile}
\operatorname{Succ}_{\rm dyn}^{\rm sat}(n,m,\eta;b)
=
\sup_{\mathcal A\in
\mathscr A_{n,b,m}^{\rm dyn,hom}(\eta)}
\mathfrak M_{n,m,\eta}\!\left(\alpha_I(\mathcal A)\right),
\end{equation}
and define
\begin{equation}
\label{eq-dynamical-computational-noise-threshold}
\eta_{\rm comp,dyn}^{\rm sat}(n,m;b,\alpha)
=
\sup\left\{
0\le\eta<\frac12:
\operatorname{Succ}_{\rm dyn}^{\rm sat}(n,m,\eta;b)\ge\alpha
\right\}.
\end{equation}
For polynomial capacity write \(b_C(n)=n^C\) and use the superscript
\((C)\) for the three profiles obtained by substituting \(b_C\).

\end{definition}

\begin{theorem}[Dynamical statistical, computational, and information bounds]
\label{thm-dynamical-noise-three-thresholds}

The sample threshold in
\cref{eq-dynamical-noise-sample-threshold} certifies clean average risk at
most \(\varepsilon\) for every sample size at least that threshold.  When
the saturated numerator has a represented noise-independent value
\[
\operatorname{Exc}_{\rm corr}^{\rm sat}
(n,b(n),m,\eta,\delta)
=\overline{\operatorname{Exc}}_{\rm corr}^{\rm sat}
(n,b(n),m,\delta)
\qquad(0\le\eta<1/2),
\]
\begin{equation}
\label{eq-explicit-dynamical-statistical-noise-threshold}
\eta_{\rm stat,dyn}^{\rm sat}
=
\left[
\frac12-
\frac{
\overline{\operatorname{Exc}}_{\rm corr}^{\rm sat}
 (n,b(n),m,\delta)}
{2\varepsilon}
\right]_{[0,1/2]}.
\end{equation}
The same statements hold after the polynomial substitution \(b(n)=n^C\).
Because the affordable record classes are nested in their certified rate
ceiling and the recordwise numerators use their own rates, the statistical
set in \cref{eq-dynamical-statistical-noise-threshold} is downward closed.

Suppose next that \(\Theta\) is uniform on a finite target family
\(\mathcal T_n\) of size \(N_n\).  Let \(\mathsf Z_0\) be the initial public
record, put \(I_0=I(\Theta;\mathsf Z_0)\), with mutual information measured
in bits throughout this paragraph, and assume that a represented
adaptive policy selects each query from \(\mathsf Z_0\) and the previous
noisy labels.  If the label returned at every query passes through an
independent homogeneous binary symmetric channel of rate \(\eta\), then the
complete transcript \(\mathsf Z_m\) satisfies
\begin{equation}
\label{eq-adaptive-bsc-information-bound}
I(\Theta;\mathsf Z_m)
\le I_0+m\bigl(1-H_2(\eta)\bigr).
\end{equation}
Every estimator, irrespective of computational cost, therefore satisfies
\begin{equation}
\label{eq-adaptive-noise-information-success-bound}
\Pr\{\widehat\Theta=\Theta\}
\le
\min\left\{
1,\,
\frac{I_0+m(1-H_2(\eta))+1}{\log_2N_n}
\right\}.
\end{equation}
Success at least \(\alpha\) requires
\begin{equation}
\label{eq-adaptive-noise-information-sample-lower-bound}
m(1-H_2(\eta))
\ge\alpha\log_2N_n-I_0-1.
\end{equation}
For \(N_n=2^n\) and a target-independent public record, the right side is
\(\alpha n-1\).

Finally, homogeneous noise degradation by the represented causal flip
postprocessing of \cref{thm-noise-degradation-monotonicity} gives
\[
b^\sharp(n)=\psi_{\mathrm{flip}}(b(n),n,m)
\]
and
\[
\operatorname{Succ}_{\rm dyn}^{\rm sat}
(n,m,\eta_1;b^\sharp)
\ge
\operatorname{Succ}_{\rm dyn}^{\rm sat}(n,m,\eta_2;b)
\quad(0\le\eta_1\le\eta_2<1/2).
\]
Thus every ceiling closed under this postprocessing has a downward-closed
computational-noise region with endpoint
\(\eta_{\rm comp,dyn}^{\rm sat}\).

\end{theorem}

\begin{proof}

The first statements are
\cref{eq-saturated-dynamical-clean-noise-bound} and the definitions of the
two thresholds.  Solving the defining inequality for \(\eta\) proves
\cref{eq-explicit-dynamical-statistical-noise-threshold}.
Nested record classes make the left side of the statistical condition
nondecreasing in \(\eta\), while \((1-2\eta)\varepsilon\) is decreasing;
this proves downward closure.

For the information bound, first retain the contribution \(I_0\) of the
initial public record.  Then condition successively on that record, the
queries, actions, and previous noisy labels.  The next query is generated
from this conditioned record and contributes no additional conditional
mutual information about \(\Theta\).  The next binary label is one use of a
memoryless binary symmetric channel, so its conditional mutual information
is at most \(1-H_2(\eta)\) bits.  The chain rule proves
\cref{eq-adaptive-bsc-information-bound}; Fano's entropy inequality proves
\cref{eq-adaptive-noise-information-success-bound,eq-adaptive-noise-information-sample-lower-bound}.
The final assertion applies the same independent postprocessing coupling as
\cref{thm-noise-degradation-monotonicity} at each represented label read.

\end{proof}

\Cref{fig-dynamical-three-threshold-ledger} separates the statistical,
computational, and information coordinates evaluated on the same noisy
experiment.

\begin{figure}[tbp]
\centering
\resizebox{.96\linewidth}{!}{%
\begin{tikzpicture}[fw diagram,node distance=8mm and 11mm]
  \node[fw source,text width=2.7cm] (input)
    {experiment\\\((n,b(n),m,\eta)\)};
  \node[fw use,text width=3.0cm,above right=12mm and 17mm of input] (stat)
    {statistical radius\\clean-risk transport};
  \node[fw use,text width=3.0cm,right=17mm of input] (comp)
    {affordable success\\supremum};
  \node[fw use,text width=3.0cm,below right=12mm and 17mm of input] (info)
    {channel\\information and\\Fano ceiling};
  \node[fw terminal,text width=3.0cm,right=20mm of stat] (tstat)
    {sample/statistical\\noise threshold};
  \node[fw terminal,text width=3.0cm,right=20mm of comp] (tcomp)
    {computational\\noise threshold};
  \node[fw terminal,text width=3.0cm,right=20mm of info] (tinfo)
    {information sample\\lower bound};
  \draw[fw arrow] (input) -- (stat);
  \draw[fw arrow] (input) -- (comp);
  \draw[fw arrow] (input) -- (info);
  \draw[fw arrow] (stat) -- (tstat);
  \draw[fw arrow] (comp) -- (tcomp);
  \draw[fw arrow] (info) -- (tinfo);
\end{tikzpicture}%
}
\caption{The three dynamical thresholds share the problem size, sample
size, noise law, and computational ceiling, but they optimize different
objects.  A statistical excess radius, an affordable-success supremum, and a
channel-information bound cannot replace one another; they may only be
reported side by side after their experiment is matched.}
\label{fig-dynamical-three-threshold-ledger}
\end{figure}

\section{Dynamical Model-Channel Corruption Beyond Binary Labels}
\label{subsec-dynamical-model-channel-noise}
\label{sec:dynamical-model-channel-noise}

Binary label flips are one instance of a represented observation channel.
Reward reports, executed actions, state observations, and filter updates require
different correction operators and different error norms.  This section
defines their common statistical record.  The dynamical consequences of those
typed errors are stated in Part~V.

\subsection{Complete Typed Channel Records}
\label{subsec:dynamical-complete-channel-records}

\begin{definition}[Complete dynamical model-channel noise record]
\label{def-complete-dynamical-model-channel-noise-record}

Fix a problem size \(n\), a complete computational ceiling \(b(n)\), and a
sample of \(m\) logged transitions.  Let
\(\mathcal F_t^{\rm obs}\subseteq\mathcal G_t^{\rm lat}\) be, respectively,
the public learner filtration and the analytic latent filtration.  A
\emph{complete dynamical model-channel noise record} consists of one or more
of the following typed interfaces.  When several interfaces occur in one
system, the record also names the joint or composite object whose error is
estimated; marginal certificates are not treated as a joint certificate.

\begin{enumerate}[label=\textup{(\roman*)}]
\item A reward record contains a clean conditional-mean comparator
      \(r_{*,t}=\mathbb E[R_t\mid\mathcal G_{t-1}^{\rm lat},X_t,A_t]\),
      a public report \(\widetilde R_t\), and represented predictable
      coefficients satisfying
      \begin{equation}
      \label{eq-dynamical-affine-reward-channel}
      \mathbb E[\widetilde R_t\mid\mathcal F_{t-1}^{\rm obs},Z_t]
      =c_t r_{*,t}^{\rm obs}(Z_t)+d_t+u_t,
      \qquad |c_t|\ge c_{\min}>0,
      \end{equation}
      where
      \(r_{*,t}^{\rm obs}(Z_t)=
      \mathbb E[r_{*,t}\mid\mathcal F_{t-1}^{\rm obs},Z_t]\).
      The record gives the same-law bound
      \(\|u\|_{L^2(d^b)}\le\zeta_r\) and a bounded,
      sub-Gaussian, sub-exponential, or predictable-variance certificate for
      the centered reward innovation.  The identity observation has
      \(c_t=1,d_t=u_t=0\).

      The record declares a probability law \(d^b\) on the disjoint union
      \[
      \mathsf S=\bigsqcup_t\{t\}\times\mathsf S_t,
      \qquad s=(t,h_{t-1}^{\rm obs},z_t),
      \]
      of pre-report predictor inputs.  Thus the time index is part of the
      input.  An \(L^2(d^b)\) norm uses the \(\mathsf S\)-marginal, whereas a
      reward loss uses the joint law
      \(d^b(ds)Q_t(d\widetilde r\mid s)\).  The functions
      \(c,d,u,r_*^{\rm obs}\) are measurable on \(\mathsf S\), the report
      and its conditional mean are square-integrable, and
      \[
      \sigma(\mathcal F_{t-1}^{\rm obs},Z_t)
      \subseteq
      \sigma(\mathcal G_{t-1}^{\rm lat},X_t,A_t).
      \]
      This inclusion and the tower property identify
      \(r_{*,t}^{\rm obs}=\mathbb E[R_t\mid
      \mathcal F_{t-1}^{\rm obs},Z_t]\).

\item An execution record contains intended actions \(A_t\), executed
      actions \(A_t^\circ\), and a represented conditional channel
      \begin{equation}
      \label{eq-dynamical-execution-channel-record}
      K_t^A(da^\circ\mid
      \mathcal F_{t-1}^{\rm obs},Z_t,X_t,A_t).
      \end{equation}
      The physical analytic kernel may depend on latent \(X_t\); a
      learner-evaluable version requires a represented public interface.
      If \(A_t^\circ\) is not logged, only the composite conditional law of
      the observed next report given the declared public predictor input is
      potentially identifiable.  A latent reward--transition composite
      additionally requires paired latent data or a represented
      reconstruction or stability theorem.  No separate estimate of
      \(K_t^A\) is inferred.

\item An observation record contains a latent clean observation or state
      \(X_t\), a public observation \(O_t\), and a represented emission
      channel \(K_t^O(do\mid X_t,A_{t-1}^\circ,
      \mathcal F_{t-1}^{\rm obs})\).  A clean-state comparison also contains
      paired latent data, a known channel with a represented correction, or
      the reconstruction certificate of
      \cref{def-dynamical-observation-reconstruction}.  No latent-state error
      is inferred from observation-law risk alone.

\item A partial-observation or filter record contains either separately
      identified transition and emission channels or only their identifiable
      history kernel.  It declares which object is estimated, the represented
      filter or memory update, its input and output norms, and a represented
      loss-to-filter-error calibration.  A training loss without that
      calibration remains an empirical loss coordinate.
\end{enumerate}

In every case, the record also contains a dynamical-sample certificate from
\cref{def-dynamical-sample-certificate}; the channel class, estimator,
calibration or inverse, coverage event, behavior-to-deployment comparison,
confidence allocation, numerical precision, retained state, and all data
collection, training, selection, storage, and deployment costs.  Their total
saturated cost is at most \(b(n)\).  A requested norm whose coverage or
calibration is absent receives the value \(+\infty\).
For square loss, the record either proves uniform bounds on both arguments
\(g(S)\) and the reported reward \(\widetilde R\), and records the resulting
loss normalization to \([0,1]\), or supplies a separate
represented unbounded-loss empirical-process or martingale theorem with its
tail, variance or truncation, and selection penalties.  The bounded
dynamical-sample certificate is otherwise ineligible and the squared-loss
excess certificate equals \(+\infty\).  The return-MGF calculation below is
not, by itself, a uniform generalization theorem for a data-selected squared
loss.

The latent variables and clean comparators in this definition are analytic
objects.  A policy may use only \(\mathcal F_t^{\rm obs}\)-measurable
quantities.  A channel is target-independent exogenous noise only when,
conditional on the declared public pre-hit state, its innovation is generated
by a target-free represented constructor and satisfies the declared
conditional independence from the latent target.  Information already
present in that public state remains charged.  State-, history-, learned-,
or target-dependent channels are admissible only with their complete
represented ancestry and the same-record source or contact route used by the
claim; otherwise the proposed prospective use is ineligible.  Mean-zero
terminology alone supplies no target-neutrality or source exemption.

\end{definition}

\Cref{fig-dynamical-channel-record-filtrations} separates the analytic and
policy-visible coordinates of this complete record.

\begin{figure}[tbp]
\centering
\resizebox{\linewidth}{!}{%
\begin{tikzpicture}[fw diagram,node distance=5mm and 12mm]
  \node[fw source,text width=4.1cm] (lat)
    {analytic latent filtration \(\mathcal G_t^{\rm lat}\)\\
     \(X_t,r_{*,t},A_t^\circ,R_t,T,F\)};
  \node[fw use,text width=2.7cm,right=of lat,yshift=18mm] (rew)
    {affine reward\\\((c,d,u)\)};
  \node[fw use,text width=2.7cm,below=3mm of rew] (act)
    {execution\\\(K_t^A\)};
  \node[fw use,text width=2.7cm,below=3mm of act] (obs)
    {observation\\\(K_t^O,D_t\)};
  \node[fw use,text width=2.7cm,below=3mm of obs] (fil)
    {filter/calibration\\\(\widehat b_t\)};
  \node[fw provenance,text width=4.1cm,right=13mm of act,yshift=-7mm] (pub)
    {policy-visible filtration \(\mathcal F_t^{\rm obs}\)\\
     \(Z_t,A_t,E_t,\widetilde R_t,O_t,\widehat b_t\)};
  \draw[fw dashed arrow] (lat.east) -- (rew.west);
  \draw[fw arrow] (lat.east) -- (act.west);
  \draw[fw arrow] (lat.east) -- (obs.west);
  \draw[fw dashed arrow] (lat.east) -- (fil.west);
  \draw[fw arrow] (rew.east) -- (pub.west);
  \draw[fw arrow] (act.east) -- (pub.west);
  \draw[fw arrow] (obs.east) -- (pub.west);
  \draw[fw arrow] (fil.east) -- (pub.west);
  \node[fw group,fit=(rew)(act)(obs)(fil)(pub),
        label={[font=\scriptsize]above:
        represented ancestry and \(\operatorname{cost}_{\rm sat}\le b(n)\)}] {};
\end{tikzpicture}%
}
\caption{A complete dynamical channel record separates the analytic latent
filtration from the policy-visible filtration.  Reward, execution,
observation, and filter interfaces retain their own norms and converters.
Only represented public-measurable outputs may enter a prospective
computation; latent comparators remain analytic.}
\label{fig-dynamical-channel-record-filtrations}
\end{figure}

\subsection{Reward-Channel Transport and Return Concentration}
\label{subsec:dynamical-reward-channel-transport}

\begin{theorem}[Affine reward-channel transport]
\label{thm-dynamical-affine-reward-channel-transport}

Let a reward record satisfy
\cref{eq-dynamical-affine-reward-channel} under its declared behavior-input
law.  For \(\nu\in\{b,\pi\}\), let \(\mathbb P^\nu\) be the declared joint
law of \((S,R,\widetilde R)\), let \(d^\nu\) be its \(\mathsf S\)-marginal,
and fix the conditional-mean versions
\[
r_{*,\nu}^{\rm obs}(s)=\mathbb E_{\mathbb P^\nu}[R\mid S=s].
\]
Assume \(d^\pi\ll d^b\) whenever the deployment conclusion is requested.
Under the behavior law put
\[
g_*(s)=c(s)r_{*,b}^{\rm obs}(s)+d(s)+u(s).
\]
Let \(\mathcal D\) be the training and output-selection sigma-field.  For
every represented square-integrable predictor \(g\) measurable with respect
to \(\sigma(\mathcal D,S)\), require a fresh evaluation pair with the
conditional law
\[
\mathcal L(S,\widetilde R\mid\mathcal D)
=d^b(ds)Q^b(d\widetilde r\mid s),
\qquad
\int \widetilde r\,Q^b(d\widetilde r\mid s)=g_*(s)
\]
almost surely.  Thus the evaluation input as well as its report innovation
is fresh relative to selection; an in-sample identity requires a separate
uniform empirical-process theorem.
Then squared-loss excess against that fresh report satisfies, conditionally
on \(\mathcal D\) and hence also after averaging, the exact identity
\begin{equation}
\label{eq-reward-square-loss-excess-identity}
\mathbb E\!\left[
(g(S)-\widetilde R)^2-(g_*(S)-\widetilde R)^2
\,\middle|\,\mathcal D\right]
=\|g-g_*\|_{L^2(d^b)}^2.
\end{equation}
Define the represented debiased predictor
\(\widehat r=(g-d)/c\), pointwise in the predictable coefficients.  The
displayed conclusion applies to exact coefficients satisfying
\cref{eq-dynamical-affine-reward-channel}; estimated coefficients require
their own represented perturbation terms.  If the left side of
\cref{eq-reward-square-loss-excess-identity} is at most \(E_r\ge0\), then
\begin{align}
\|\widehat r-r_{*,b}^{\rm obs}\|_{L^2(d^b)}
&\le \frac{\sqrt{E_r}+\zeta_r}{c_{\min}},
\label{eq-reward-clean-l2-transport}\\
\|\widehat r-r_{*,\pi}^{\rm obs}\|_{L^2(d^\pi)}
&\le
\frac{\sqrt{W_{\pi\leftarrow b}}
      (\sqrt{E_r}+\zeta_r)}{c_{\min}}
+\Delta_{r,\pi\leftarrow b}
\label{eq-reward-deployment-l2-transport}
\end{align}
whenever the represented deployment comparison
\(d^\pi\le W_{\pi\leftarrow b}d^b\) holds and the record certifies
\[
\Delta_{r,\pi\leftarrow b}
=\|r_{*,\pi}^{\rm obs}-r_{*,b}^{\rm obs}\|_{L^2(d^\pi)}.
\]
This shift is zero only under a represented conditional-reward invariance
statement, for example when the predictor input contains the full
reward-relevant state and action.

For \(M\) independent episodes, let
\((\mathcal H_{e,t})_{t=0}^{H}\) be a nested within-episode filtration such
that \(\mathcal H_{e,t}\) is pre-report, contains all earlier report errors,
and \(\widetilde R_{e,t}-R_{e,t}\) is
\(\mathcal H_{e,t+1}\)-measurable.  Define
\begin{align*}
B_{e,t}
&=\mathbb E[\widetilde R_{e,t}-R_{e,t}\mid\mathcal H_{e,t}],\\
\xi_{e,t}
&=\widetilde R_{e,t}-R_{e,t}-B_{e,t}.
\end{align*}
Suppose \(|B_{e,t}|\le\beta_{r,t}\) almost surely and
\[
\mathbb E[\exp(\lambda\xi_{e,t})\mid\mathcal H_{e,t}]
\le \exp(\lambda^2\sigma_{r,t}^2/2)
\qquad(\lambda\in\mathbb R).
\]
For fixed deterministic (hence \(\mathcal H_{e,t}\)-predictable) return
weights \(w_t\), with
probability at least \(1-\delta\),
\begin{equation}
\label{eq-reward-return-subgaussian-deviation}
\left|
\frac1M\sum_{e=1}^M\sum_{t=0}^{H-1}
w_t(\widetilde R_{e,t}-R_{e,t})
\right|
\le
\sum_{t=0}^{H-1}|w_t|\beta_{r,t}
+\sqrt{
\frac{2\log(2/\delta)}M
\sum_{t=0}^{H-1}w_t^2\sigma_{r,t}^2
}.
\end{equation}
The same iterated-MGF proof applies to any single globally ordered
predictable sequence satisfying the displayed conditional MGF; independence
is then unnecessary.  A \(\beta\)-mixing replacement requires a separately
proved bounded or truncated block certificate with its coupling and
truncation terms.  Neither mixing nor predictability alone authorizes
substitution of a radius.  A data-dependent predictor or weight choice uses
a simultaneous class theorem rather than this fixed-choice display.  At
fixed horizon the logged transition count is \(m_{\rm tr}=MH\); both \(M\)
and \(m_{\rm tr}\) remain in the complete cost record.

\end{theorem}

\begin{proof}

Put \(\mathscr X=\sigma(\mathcal D,S)\).  The conditional-law hypothesis
gives \(S\mid\mathcal D\sim d^b\) and
\(g_*(S)=\mathbb E[\widetilde R\mid\mathscr X]\).  Hence, conditionally on
\(\mathcal D\),
\[
\mathbb E[(g(S)-g_*(S))(g_*(S)-\widetilde R)\mid\mathcal D]
=\mathbb E\!\left[(g(S)-g_*(S))
  \mathbb E[g_*(S)-\widetilde R\mid\mathscr X]
  \middle|\mathcal D\right]=0.
\]
The pointwise expansion
\[
(g-\widetilde R)^2-(g_*-\widetilde R)^2
=(g-g_*)^2+2(g-g_*)(g_* -\widetilde R)
\]
and \(S\mid\mathcal D\sim d^b\) therefore prove
\cref{eq-reward-square-loss-excess-identity}.  Moreover,
\[
\widehat r-r_{*,b}^{\rm obs}
=\frac{g-g_*+u}{c}.
\]
Thus
\[
\|\widehat r-r_{*,b}^{\rm obs}\|_{L^2(d^b)}
\le c_{\min}^{-1}
\bigl(\|g-g_*\|_{L^2(d^b)}+\|u\|_{L^2(d^b)}\bigr),
\]
and the first bound follows from the excess identity.  If
\(\rho_{\pi\leftarrow b}=d d^\pi/d d^b\le W_{\pi\leftarrow b}\), then for
every measurable \(f\),
\[
\|f\|_{L^2(d^\pi)}^2
=\int f^2\rho_{\pi\leftarrow b}\,d d^b
\le W_{\pi\leftarrow b}\|f\|_{L^2(d^b)}^2.
\]
This bounds the behavior comparator under \(d^\pi\).  The triangle
inequality with the certified conditional-reward shift proves
\cref{eq-reward-deployment-l2-transport}.

For one episode, iterated conditional expectation and the conditional MGF
hypothesis give
\[
\mathbb E\exp\!\left(\lambda\sum_{t=0}^{H-1}w_t\xi_{e,t}\right)
\le
\exp\!\left(\frac{\lambda^2}{2}
             \sum_{t=0}^{H-1}w_t^2\sigma_{r,t}^2\right).
\]
Independence across episodes makes the MGF of the sum over \(e\) the product
of these bounds.  After division by \(M\), the proxy is
\(M^{-1}\sum_tw_t^2\sigma_{r,t}^2\).  Applying Chernoff's bound to both
signs gives the square-root term in
\cref{eq-reward-return-subgaussian-deviation}.  Finally,
\[
\left|\frac1M\sum_{e,t}w_tB_{e,t}\right|
\le\sum_t|w_t|\beta_{r,t},
\]
which supplies its deterministic bias term.

\end{proof}

\Cref{fig-affine-reward-channel-transport} records the two distinct reward
conversions: population debiasing and episode-return concentration.

\begin{figure}[tbp]
\centering
\resizebox{\linewidth}{!}{%
\begin{tikzpicture}[fw diagram,node distance=7mm and 9mm]
  \node[fw source,text width=2.6cm] (clean)
    {behavior comparator\\\(r_{*,b}^{\rm obs}\)};
  \node[fw use,text width=2.7cm,right=of clean] (aff)
    {affine channel\\\(g_*=cr_*+d+u\)};
  \node[fw box,text width=2.6cm,right=of aff] (exc)
    {fresh-report excess\\\(E_r=\|g-g_*\|_2^2\)};
  \node[fw provenance,text width=2.7cm,right=of exc] (deb)
    {debias\\\(\widehat r=(g-d)/c\)};
  \node[fw terminal,text width=3.1cm,right=of deb] (dep)
    {deployment bound\\density ratio + \(\Delta_{r,\pi\leftarrow b}\)};
  \draw[fw arrow] (clean) -- (aff);
  \draw[fw arrow] (aff) -- (exc);
  \draw[fw arrow] (exc) -- (deb);
  \draw[fw arrow] (deb) -- (dep);
  \node[fw source,text width=2.7cm,below=12mm of aff] (eps)
    {\(M\) independent episodes\\nested report filtration};
  \node[fw use,text width=2.9cm,right=of eps] (mgf)
    {conditional MGF\\\(\sum_tw_t^2\sigma_{r,t}^2\)};
  \node[fw terminal,text width=3.0cm,right=of mgf] (ret)
    {return radius\\bias + \(M^{-1/2}\) fluctuation};
  \draw[fw arrow] (eps) -- (mgf);
  \draw[fw arrow] (mgf) -- (ret);
\end{tikzpicture}%
}
\caption{Affine reward correction is a same-law transport.  Fresh-report
orthogonality converts squared-loss excess to behavior-law \(L^2\) error;
density comparison and an explicit conditional-reward shift give deployment
error.  Episode-level return noise is controlled separately by its predictable
bias and conditional-MGF radius.}
\label{fig-affine-reward-channel-transport}
\end{figure}

\subsection{Observation Reconstruction and Identifiability}
\label{subsec:dynamical-observation-reconstruction}

\begin{definition}[Observation reconstruction and identifiability]
\label{def-dynamical-observation-reconstruction}

For a finite latent carrier \(X\) and observation carrier \(O\), let
\(\mathcal C_{\rm adm}\) be the represented set of dynamic contexts on which
the uniform conclusion is claimed, and fix
\(c=(t,h_{t-1}^{\rm obs},a_{t-1}^\circ)\).  Let
\(\mathsf O_c(do\mid x)\) be the emission kernel and
\(D_c(dx'\mid o)\) a represented decoding kernel.  Their affine pushforward
operators are
\(\mathsf O_{c,*}:\mathcal P(X)\to\mathcal P(O)\) and
\(D_{c,*}:\mathcal P(O)\to\mathcal P(X)\).  Their composition induces the
Markov reconstruction kernel
\[
K_{{\rm obs},c}(x,B)
=(D_{c,*}\mathsf O_{c,*}\delta_x)(B)
=\sum_{o\in O}\mathsf O_c(o\mid x)D_c(B\mid o),
\qquad B\subseteq X.
\]
We use
\(\|P-Q\|_{\rm TV}=\sup_B|P(B)-Q(B)|
=\tfrac12\|P-Q\|_1\).  A uniform dynamic row-total-variation
reconstruction certificate is
\begin{equation}
\label{eq-observation-reconstruction-row-tv}
\varepsilon_{\rm obs}^{\rm rec}
=\sup_{c\in\mathcal C_{\rm adm}}\sup_{x\in X}
\|K_{{\rm obs},c}(x,\cdot)-\delta_x\|_{\rm TV}.
\end{equation}
It gives, for every bounded latent function \(f\),
\begin{equation}
\label{eq-observation-reconstruction-function-bound}
\sup_{c\in\mathcal C_{\rm adm}}\|K_{{\rm obs},c}f-f\|_\infty
\le\varepsilon_{\rm obs}^{\rm rec}\operatorname{Osc}(f).
\end{equation}
An occupancy-averaged certificate must name its context occupancy and cannot
be inserted into this sup-norm display.  All context dependence and decoder
selection remain in the same complete record.

Alternatively, for a declared joint law
\(\mu(dx)\mathsf O(do\mid x)\), write
\(\mu\mathsf O_*\) for its observation marginal and define, up to
\(\mu\mathsf O_*\)-almost-sure equality, the conditional-expectation operator
\[
(\mathsf O_\mu f)(o)=\mathbb E_\mu[f(X)\mid O=o].
\]
A law-dependent stability certificate is a represented tuple
\((\mathcal D_{\rm obs},\kappa_{\rm obs},\Delta_{\rm obs})\) such that,
simultaneously for every \(f\in\mathcal D_{\rm obs}\),
\[
\|f\|_{L^2(\mu)}
\le\kappa_{\rm obs}\|\mathsf O_\mu f\|_{L^2(\mu\mathsf O_*)}
+\delta_{\rm obs}(f),
\qquad 0\le\delta_{\rm obs}(f)\le\Delta_{\rm obs}.
\]
A learned clean-error conversion may use this inequality only after the same
record proves that the estimator-minus-target difference belongs to
\(\mathcal D_{\rm obs}\) and that its observed error is exactly the displayed
conditional expectation in the displayed law and norm.  A deployment change
in the latent posterior or in \(\mathsf O_\mu\) requires a separate
represented shift term.  The decoder, operator, stability domain, constants,
and construction costs belong to the same record.  If neither branch
applies, the clean latent-coordinate error is \(+\infty\), although
observed-law prediction risk may remain finite.

\end{definition}

\begin{proof}[Proof of \cref{eq-observation-reconstruction-function-bound}]

For each \(c,x\), the dual characterization of total variation gives
\[
|K_{{\rm obs},c}f(x)-f(x)|
\le
\|K_{{\rm obs},c}(x,\cdot)-\delta_x\|_{\rm TV}
\operatorname{Osc}(f).
\]
Taking the two suprema proves the claim.

\end{proof}

\begin{remark}[Observation noise is not invertible in general]
\label{rem-observation-noise-nonidentifiability}

Let \(X=\{0,1\}\) and suppose the two latent states have identical emission
rows.  A single observation has the same law under the two states; under
repeated conditionally memoryless observations of a fixed hidden state, with
no other public kernel that distinguishes the states, induction gives the
same law for every public transcript.  For any decoder the two reconstruction rows therefore
equal one law \(q\).  The triangle inequality gives
\[
1=\|\delta_0-\delta_1\|_{\rm TV}
\le\|\delta_0-q\|_{\rm TV}+\|q-\delta_1\|_{\rm TV},
\]
so \(\varepsilon_{\rm obs}^{\rm rec}\ge1/2\).  Thus
emission entropy, sample size, and an observed-law generalization bound do
not by themselves provide a clean-state or target-gate bound.

\end{remark}

\Cref{fig-observation-reconstruction-identifiability} displays both the
valid reconstruction route and the obstruction created by merged emission
rows.

\begin{figure}[tbp]
\centering
\resizebox{.94\linewidth}{!}{%
\begin{tikzpicture}[fw diagram,node distance=8mm and 13mm]
  \node[fw source,text width=2.8cm] (px) {latent laws\\\(\mathcal P(X)\)};
  \node[fw use,text width=2.8cm,right=of px] (po)
    {emission\\\(\mathsf O_{c,*}\)};
  \node[fw provenance,text width=2.8cm,right=of po] (dec)
    {represented decoder\\\(D_{c,*}\)};
  \node[fw terminal,text width=3.0cm,right=of dec] (tv)
    {reconstruction row\\\(\varepsilon_{\rm obs}^{\rm rec}\)};
  \draw[fw arrow] (px) -- (po);
  \draw[fw arrow] (po) -- (dec);
  \draw[fw arrow] (dec) -- (tv);
  \node[fw box,text width=4.2cm,below=10mm of tv] (fun)
    {\(\|K_{{\rm obs},c}f-f\|_\infty
      \le\varepsilon_{\rm obs}^{\rm rec}\operatorname{Osc}(f)\)};
  \draw[fw arrow] (tv) -- (fun);
  \node[fw source,below=27mm of px] (x0) {\(x=0\)};
  \node[fw source,below=7mm of x0] (x1) {\(x=1\)};
  \node[fw residual,text width=3.1cm,right=16mm of x0,yshift=-5mm] (merge)
    {identical emission row\\one observation law \(q\)};
  \node[fw terminal,text width=3.0cm,right=of merge] (floor)
    {no inversion\\\(\max_x\|q-\delta_x\|_{\rm TV}\ge\tfrac12\)};
  \draw[fw arrow] (x0) -- (merge);
  \draw[fw arrow] (x1) -- (merge);
  \draw[fw arrow] (merge) -- (floor);
\end{tikzpicture}%
}
\caption{A represented decoder supplies a clean-state comparison through the
composite reconstruction kernel.  Row TV controls bounded latent functions
by oscillation.  Identical emission rows erase the binary latent distinction,
forcing reconstruction error at least \(1/2\); observed-law risk alone gives
no clean target-gate bound.}
\label{fig-observation-reconstruction-identifiability}
\end{figure}

\subsection{Saturated Typed Errors and Entropy Continuity}
\label{subsec:dynamical-typed-error-entropy}

\begin{definition}[Saturated typed channel-error profiles]
\label{def-saturated-dynamical-channel-error-profiles}

The fixed-law profile is defined below.  Its uniform instance--law envelope
and reward-square specialization are the separate coordinates in
\cref{def:typed-channel-family-envelope,%
def:typed-channel-reward-square-profile}.

Let
\(q\in\{r,A,O,{\rm fil},{\rm comp}\}\) denote reward, execution,
observation, filter, or identifiable composite-kernel error, and fix a
declared norm \(\mathsf N\).  For each fixed presented instance \(I\) and
data-generating law \(\mathbb P_I\), a complete record \(\mathcal R\) names
the clean object, the returned object, their common population law, and the
true discrepancy
\[
e_{q,\mathsf N}(\mathcal R)\in[0,+\infty]
\]
between them in \(\mathsf N\).  For every candidate
\(j\in\mathcal J_{q,\mathsf N}(\mathcal R)\), put
\[
B_{\mathcal R,j}
=\widehat e_{q,\mathsf N,j}(\mathcal R)
 +\Gamma_{q,\mathsf N,j}(\mathcal R;m,\delta_j).
\]
Any nonlinear converter is applied before this notation, so every
\(B_{\mathcal R,j}\) has the same estimand, law, and norm as
\(e_{q,\mathsf N}(\mathcal R)\).

Let \(\mathscr R_{I,q,\mathsf N}(b(n))\) be the affordable complete record
outputs defined on this fixed sample space and evaluated under this fixed
law.  Eligibility requires one pre-specified, output-selection-independent
sample event
\(\mathcal E_{I,q,\mathsf N}^{\rm sat}(\delta)\), with
\(\mathbb P_I(\mathcal E_{I,q,\mathsf N}^{\rm sat})\ge1-\delta\), on which
\begin{equation}
\label{eq-saturated-channel-global-event}
e_{q,\mathsf N}(\mathcal R)\le B_{\mathcal R,j}
\quad\text{for every affordable data-dependent output \(\mathcal R\)
and every eligible \(j\)}.
\end{equation}
A finite represented codebook may obtain this event from a pre-data global
allocation \(\sum_{(\mathcal R,j)}\delta_{\mathcal R,j}\le\delta\).
Otherwise it requires one uniform cover, PAC--Bayes, or sequential theorem
over the affordable output hull.  A per-record confidence allocation is not
sufficient after the record is selected from data.  Events belonging to
different data-generating laws are never intersected.  Define
\begin{equation}
\label{eq-recordwise-dynamical-channel-error}
\operatorname{Err}_{q,\mathsf N}(\mathcal R;m,\delta)
=
\min_{j\in\mathcal J_{q,\mathsf N}(\mathcal R)}
 B_{\mathcal R,j},
\end{equation}
where the minimum of an empty family is \(+\infty\).  The saturated profile
is the recordwise minimum followed by the affordable supremum
\begin{equation}
\label{eq-saturated-dynamical-channel-error}
\operatorname{Err}_{q,\mathsf N;I,\mathbb P_I}^{\rm sat}
(n,b(n),m,\delta)
=
\sup_{\mathcal R\in\mathscr R_{I,q,\mathsf N}(b(n))}
\operatorname{Err}_{q,\mathsf N}(\mathcal R;m,\delta).
\end{equation}
On \(\mathcal E_{I,q,\mathsf N}^{\rm sat}(\delta)\),
\begin{equation}
\label{eq-saturated-channel-true-error-bound}
\sup_{\mathcal R\in\mathscr R_{I,q,\mathsf N}(b(n))}
e_{q,\mathsf N}(\mathcal R)
\le \operatorname{Err}_{q,\mathsf N;I,\mathbb P_I}^{\rm sat}
(n,b(n),m,\delta).
\end{equation}
Indeed, \cref{eq-saturated-channel-global-event} makes every
\(B_{\mathcal R,j}\) an upper bound for the same number simultaneously, so
that number is at most their minimum; taking the affordable supremum proves
the display.  The supremum of an empty record class is zero.  An eligible
record with an empty candidate family makes the unrestricted worst-record
profile \(+\infty\); finiteness is therefore a theorem only for a declared
subclass in which every record has a finite converter.  Profiles in distinct norms are not
minimized or added.  A represented norm, density, channel-inversion,
filter-stability, Bellman, path-law, or Lift theorem must first convert them
to one output coordinate.  The polynomial-capacity specialization is
\(b_C(n)=n^C\).

\end{definition}

\begin{definition}[Uniform instance--law channel-error envelope]
\label{def:typed-channel-family-envelope}

Fix \((q,\mathsf N,n,b,m,\delta)\) and the fixed-law typed profiles of
\cref{def-saturated-dynamical-channel-error-profiles}.

When the presented instance and law are fixed, their subscripts are
suppressed.  If a deterministic number uniform over a represented family
\(\mathfrak P_n\) of instance--law pairs is required, its probability-space
correct definition is the confidence envelope
\[
\overline{\operatorname{Err}}_{q,\mathsf N}^{\rm sat}
(n,b(n),m,\delta)
=\inf\left\{u:\inf_{(I,\mathbb P_I)\in\mathfrak P_n}
\mathbb P_I\!\left(
\sup_{\mathcal R\in\mathscr R_{I,q,\mathsf N}(b(n))}
e_{q,\mathsf N}(\mathcal R)\le u\right)\ge1-\delta\right\}.
\]
A proposed finite value for this outer envelope requires a represented
distribution-uniform theorem.  It is not the pointwise supremum of random
confidence bounds living on unrelated sample spaces.

\end{definition}

\begin{definition}[Reward-square specialization of the typed channel profile]
\label{def:typed-channel-reward-square-profile}

Fix the reward channel, behavior law, and complete record class appearing in
\cref{def-saturated-dynamical-channel-error-profiles}.

For reward square loss, the clean behavior-law squared-\(L^2\) risk profile
on the fixed \((I,\mathbb P_I)\) is
\begin{equation}
\label{eq-saturated-clean-reward-error}
U_{r,L^2(d^b)^2;I,\mathbb P_I}^{\rm sat}(n,b(n),m,\delta)
=
\sup_{\mathcal R\in\mathscr R_{I,r,L^2(d^b)}(b(n))}
\frac{
\bigl(\sqrt{\operatorname{Exc}_r(\mathcal R;m,\delta)}
+\zeta_r(\mathcal R)\bigr)^2
}{c_{\min}(\mathcal R)^2},
\end{equation}
with the same complete record class as
\cref{eq-saturated-dynamical-channel-error}.  A deployment or sup-norm
version is obtained only through the record's represented comparison.
The square in the subscript is literal: this coordinate bounds
\(\|\widehat r-r_{*,b}^{\rm obs}\|_{L^2(d^b)}^2\), not the unsquared norm.
Here \(\operatorname{Exc}_r(\mathcal R;m,\delta)\ge0\) is the recordwise
simultaneous upper confidence bound on the public squared-loss excess in
\cref{eq-reward-square-loss-excess-identity}: it is the minimum of the
represented empirical excess-plus-deviation certificates in that same
behavior law and has value \(+\infty\) when no such certificate is present.
It is not replaced by an RMSE, a clean risk, or an error in another norm
without a represented converter.  The empty record class has value zero; a
missing excess certificate has value \(+\infty\).

\end{definition}

\Cref{fig-saturated-channel-profile-order} makes explicit the order of the
two optimizations in a saturated channel profile.

\begin{figure}[tbp]
\centering
\resizebox{\linewidth}{!}{%
\begin{tikzpicture}[fw diagram,node distance=7mm and 11mm]
  \node[fw source,text width=3.0cm] (records)
    {complete records \(\mathcal R\)\\\(\operatorname{cost}_{\rm sat}\le b(n)\)};
  \node[fw provenance,text width=3.4cm,right=of records] (event)
    {one uniform event\\all selected \((\mathcal R,j)\)};
  \node[fw use,text width=3.2cm,right=of event] (min)
    {within-record minimum\\\(\min_j\{\widehat e_j+\Gamma_j\}\)};
  \node[fw terminal,text width=3.0cm,right=of min] (sup)
    {affordable worst case\\\(\sup_{\mathcal R}\operatorname{Err}(\mathcal R)\)};
  \draw[fw arrow] (records) -- (event);
  \draw[fw arrow] (event) -- (min);
  \draw[fw arrow] (min) -- (sup);
  \node[fw residual,text width=4.2cm,below=10mm of min] (norm)
    {same population law and norm only\\different norms require a converter};
  \draw[fw dashed arrow] (norm) -- (min);
\end{tikzpicture}
}
\caption{The typed saturated profile first selects the strongest simultaneous
certificate within one complete record and then takes the affordable worst
case over records.  The order is \(\sup_{\mathcal R}\min_j\), and only bounds
for the same estimand, population law, and norm may compete.}
\label{fig-saturated-channel-profile-order}
\end{figure}

\begin{theorem}[Finite-alphabet channel estimation and entropy continuity]
\label{thm-finite-alphabet-channel-estimation-entropy}

Let \(0<\delta<1\).  Let
\(K:\mathsf U\to\mathcal P(\mathsf V)\) be a finite channel with
\(d_U=|\mathsf U|\ge1\), \(d_V=|\mathsf V|\ge2\), and let
\((U_i,V_i)_{i=1}^m\) be iid with joint law
\(p(u)K(v\mid u)\), where \(p_*:=\min_up(u)>0\).  The affordable record
contains a represented certified lower bound
\(0<\underline p_*\le p_*\); the radius below is evaluated with
\(\underline p_*\).  Let \(\widehat K\) be the empirical row estimator,
with an arbitrary row assigned when its count is zero.  If
\[
m\underline p_*\ge8\log(2d_U/\delta),
\]
then, with probability at least \(1-\delta\),
\begin{equation}
\label{eq-finite-channel-row-tv-estimation}
\sup_{u\in\mathsf U}
\|\widehat K(u,\cdot)-K(u,\cdot)\|_{\rm TV}
\le
\tau_m
:=
\min\left\{1,
\sqrt{
\frac{d_V\log2+\log(2d_U/\delta)}{m\underline p_*}
}
\right\}.
\end{equation}
All logarithms in this concentration radius are natural.  For a non-iid
record, this theorem may be invoked only through a named sample certificate
that supplies, on one event, a lower bound on every effective row count and a
simultaneous row-TV inequality with a displayed radius.  No iid multinomial
radius is inferred from mixing or predictability alone.  An unvisited row has
mathematical TV distance at most one under the arbitrary convention, but its
framework certificate sentinel is \(+\infty\): without coverage it is not an
identified uniform-row converter.  A data-dependent channel or alphabet
choice additionally uses the global simultaneous event in
\cref{eq-saturated-channel-global-event}.  Without a represented
\(\underline p_*\), the affordable uniform-row certificate is \(+\infty\).

Measure the following Shannon quantities in bits.  For an input law \(\nu\),
put
\[
h_K(\nu)=H_\nu(V\mid U),
\qquad
C(K)=\sup_{\nu\in\mathcal P(\mathsf U)}I_\nu(U;V).
\]
Let
\[
\mathfrak f_{d_V}(\tau)
=
\begin{cases}
H_2(\tau)+\tau\log_2(d_V-1),&
0\le\tau\le1-d_V^{-1},\\
\log_2d_V,&1-d_V^{-1}<\tau\le1.
\end{cases}
\]
On the event in \cref{eq-finite-channel-row-tv-estimation},
\begin{align}
\sup_\nu|h_{\widehat K}(\nu)-h_K(\nu)|
&\le\mathfrak f_{d_V}(\tau_m),
\label{eq-finite-channel-conditional-entropy-continuity}\\
|C(\widehat K)-C(K)|
&\le2\mathfrak f_{d_V}(\tau_m).
\label{eq-finite-channel-capacity-continuity}
\end{align}
Every estimated entropy or capacity retains the channel-estimation cost,
coverage event, sample size \(m\), confidence, and complete ceiling \(b(n)\).

\end{theorem}

\begin{proof}

Let \(N_u\) be the number of observations with input \(u\).  Since
\(N_u\sim\operatorname{Bin}(m,p(u))\), the multiplicative Chernoff bound gives
\[
\Pr\{N_u<mp(u)/2\}
\le e^{-mp(u)/8}
\le e^{-m\underline p_*/8}.
\]
Consequently,
\[
\Pr\{\min_uN_u<m\underline p_*/2\}
\le d_Ue^{-m\underline p_*/8}\le\delta/2.
\]
Condition on the complementary count event.  For one row, the multinomial
Bretagnolle--Huber--Carol bound gives
\[
\Pr\{\|\widehat K(u,\cdot)-K(u,\cdot)\|_1>\epsilon
       \mid (N_{u'})_{u'\in\mathsf U}\}
\le2^{d_V}e^{-N_u\epsilon^2/2}.
\]
If the square root defining \(\tau_m\) is at least one, the TV claim is
deterministic.  Otherwise substitute \(\epsilon=2\tau_m\) and
\(N_u\ge m\underline p_*/2\) to obtain
\[
\Pr\{\|\widehat K(u,\cdot)-K(u,\cdot)\|_{\rm TV}>\tau_m
       \mid (N_{u'})_{u'\in\mathsf U}\}
\le2^{d_V}e^{-m\underline p_*\tau_m^2}.
\]
The union over \(u\) is at most
\[
d_U2^{d_V}
\exp\{-d_V\log2-\log(2d_U/\delta)\}=\delta/2.
\]
Combining the count and row-estimation failures proves
\cref{eq-finite-channel-row-tv-estimation}.

To prove the required entropy continuity without an additional assumption,
let \(P,Q\) be two laws on \(\mathsf V\) at TV distance
\(\tau\le1-d_V^{-1}\), and take a maximal coupling \((X,Y)\) with
\(\Pr\{X\ne Y\}=\tau\).  For \(E=\mathbf1_{\{X\ne Y\}}\),
\[
H(X)-H(Y)=H(X\mid Y)-H(Y\mid X)\le H(X\mid Y)
\le H_2(\tau)+\tau\log_2(d_V-1).
\]
The last step first reveals \(E\); conditional on \(E=0\), \(X=Y\), and
conditional on \(E=1\) there are at most \(d_V-1\) possible values of
\(X\) once \(Y\) is known.  Interchanging \(X,Y\) controls the other sign.
For \(\tau>1-d_V^{-1}\), the trivial bound
\(|H(P)-H(Q)|\le\log_2d_V\) applies.  Hence, on the row-TV event, every row
satisfies
\[
|H(\widehat K(u,\cdot))-H(K(u,\cdot))|
\le\mathfrak f_{d_V}(\tau_m).
\]
Indeed, on \([0,1-d_V^{-1}]\),
\[
\frac{d}{d\tau}
\bigl[H_2(\tau)+\tau\log_2(d_V-1)\bigr]
=\log_2\!\frac{(1-\tau)(d_V-1)}{\tau}\ge0,
\]
and \(\mathfrak f_{d_V}\) is constant thereafter.  Thus replacing a row's
actual TV distance by the common upper bound \(\tau_m\) preserves the
inequality.
The definition of conditional entropy and averaging with \(\nu(u)\) yield
\[
|h_{\widehat K}(\nu)-h_K(\nu)|
\le\sum_u\nu(u)\mathfrak f_{d_V}(\tau_m)
=\mathfrak f_{d_V}(\tau_m),
\]
uniformly in \(\nu\), proving
\cref{eq-finite-channel-conditional-entropy-continuity}.  Moreover,
\[
\|\nu\widehat K-\nu K\|_{\rm TV}
\le\sum_u\nu(u)
\|\widehat K(u,\cdot)-K(u,\cdot)\|_{\rm TV}
\le\tau_m.
\]
Thus the output-entropy term and the conditional-entropy term in
\(I_\nu(U;V)=H_\nu(V)-H_\nu(V\mid U)\) each change by at most
\(\mathfrak f_{d_V}(\tau_m)\), so
\[
\sup_\nu|I_\nu(\widehat K)-I_\nu(K)|
\le2\mathfrak f_{d_V}(\tau_m).
\]
Finally,
\(\lvert\sup_\nu I_\nu(\widehat K)-\sup_\nu I_\nu(K)\rvert
\le\sup_\nu|I_\nu(\widehat K)-I_\nu(K)|\), which proves
\cref{eq-finite-channel-capacity-continuity}.

\end{proof}

\Cref{fig-finite-channel-estimation-continuity} summarizes the coverage,
row-estimation, and information-continuity chain.

\begin{figure}[tbp]
\centering
\resizebox{\linewidth}{!}{%
\begin{tikzpicture}[fw diagram,node distance=7mm and 9mm]
  \node[fw source,text width=2.8cm] (sample)
    {iid pairs \((U_i,V_i)\)\\represented \(\underline p_*\)};
  \node[fw use,text width=2.8cm,right=of sample] (counts)
    {coverage gate\\\(N_u\ge m\underline p_*/2\)};
  \node[fw provenance,text width=2.8cm,right=of counts] (rows)
    {all empirical rows\\\(\sup_u\|\widehat K_u-K_u\|_{\rm TV}\le\tau_m\)};
  \node[fw geom,text width=2.7cm,right=of rows] (ent)
    {entropy continuity\\\(\mathfrak f_{d_V}(\tau_m)\)};
  \node[fw terminal,text width=2.9cm,right=of ent] (cap)
    {capacity continuity\\\(2\mathfrak f_{d_V}(\tau_m)\)};
  \draw[fw arrow] (sample) -- (counts);
  \draw[fw arrow] (counts) -- (rows);
  \draw[fw arrow] (rows) -- (ent);
  \draw[fw arrow] (ent) -- (cap);
  \node[fw residual,text width=3.0cm,below=10mm of counts] (missing)
    {missing row coverage\\framework sentinel \(+\infty\)};
  \draw[fw dashed arrow] (counts) -- (missing);
\end{tikzpicture}%
}
\caption{Uniform finite-channel estimation requires coverage of every input
row.  On the simultaneous coverage event, the common row-TV radius propagates
to conditional entropy with loss \(\mathfrak f_{d_V}(\tau_m)\) and to
capacity with loss \(2\mathfrak f_{d_V}(\tau_m)\).}
\label{fig-finite-channel-estimation-continuity}
\end{figure}

\section{Attainable Risk and the Complete Quantitative Learning Decomposition}
\label{sec-attainable-risk-complete-learning-decomposition}

The statistical envelopes above control estimation after a represented
learning class has been selected.  The saturated budget also determines
which population risks are attainable by that class.  Keeping these two
quantities separate gives an exact approximation--estimation decomposition.

\begin{definition}[Saturated attainable risk and approximation floor]
\label{def-saturated-attainable-risk-approximation-floor}

Fix an instance \(I\), sample size \(m\), channel parameter \(\zeta\), and
budget \(B\).  Let
\(\mathscr L_{I,m,\zeta}^{\rm sat}(B)\) be the complete represented learning
records whose data acquisition, selected hypothesis or posterior, retained
state, optimization, evaluation, and deployment costs are at most \(B\).
For \(R\in\mathscr L_{I,m,\zeta}^{\rm sat}(B)\), let
\(h_R\) be its deployed readout and let \(R_{0,I}(h_R)\in[0,1]\) be its
clean population risk conditional on that readout.  The expectation below
averages the sample and private randomness of the record.  Define
\begin{align}
\label{eq-saturated-attainable-clean-risk}
\mathcal R_{I,m,\zeta}^{\rm att,sat}(B)
&=
\inf_{R\in\mathscr L_{I,m,\zeta}^{\rm sat}(B)}
\mathbb E_R R_{0,I}(h_R),\\
\label{eq-saturated-learning-approximation-floor}
\operatorname{App}_{I,m,\zeta}^{\rm sat}(B)
&=
\mathcal R_{I,m,\zeta}^{\rm att,sat}(B)-R_{0,I}^*,
\qquad
R_{0,I}^*=\inf_h R_{0,I}(h).
\end{align}
The final infimum is an analytic population benchmark.  It supplies no
readout to a represented learner.  Empty affordable classes have attainable
risk one.  A family-uniform version restricts the infimum to records produced
by one represented derivation scheme on the declared family slice:
\begin{equation}
\label{eq-family-saturated-attainable-clean-risk}
\mathcal R_{\mathfrak M,n,m,\zeta}^{\rm att,sat}(B)
=
\inf_{R\ {\rm family\text{-}uniform}}
\mathfrak M_{n,m,\zeta}
\left(I\mapsto\mathbb E_RR_{0,I}(h_R)\right).
\end{equation}

For a represented comparison class \(\mathcal H_B\), put
\begin{equation}
\label{eq-represented-class-approximation-floor}
\operatorname{App}_{I}(\mathcal H_B)
=\inf_{h\in\mathcal H_B}R_{0,I}(h)-R_{0,I}^*.
\end{equation}
The comparison class contains the construction and deployment cost of each
member.  Selection from the class is charged separately unless the
statistical event is simultaneous over the complete class.

\end{definition}

\begin{theorem}[Complete quantitative learning decomposition]
\label{thm-complete-quantitative-learning-decomposition}

Let \(\widehat h\in\mathcal H_B\) be the deployed output of a complete
learning record.  For every clean population loss,
\begin{equation}
\label{eq-exact-learning-approximation-estimation-decomposition}
R_{0,I}(\widehat h)-R_{0,I}^*
=
\operatorname{App}_{I}(\mathcal H_B)
+\left(
R_{0,I}(\widehat h)-\inf_{h\in\mathcal H_B}R_{0,I}(h)
\right).
\end{equation}
Consequently the left side is at least
\(\operatorname{App}_{I}(\mathcal H_B)\).

Assume a binary flip-symmetric loss under homogeneous symmetric label noise
\(0\le\eta<1/2\), put \(\rho=1-2\eta\), and suppose one simultaneous
event gives
\begin{equation}
\label{eq-complete-learning-corrupted-uniform-radius}
\sup_{h\in\mathcal H_B}
\left|R_{\eta,I}(h)-\widehat R_{\eta,S}(h)\right|
\le\Gamma_{I}(B,m,\eta,\delta).
\end{equation}
If the represented optimizer satisfies
\begin{equation}
\label{eq-complete-learning-empirical-optimization-error}
\widehat R_{\eta,S}(\widehat h)
\le\inf_{h\in\mathcal H_B}\widehat R_{\eta,S}(h)
+\varepsilon_{\rm opt},
\end{equation}
then, on that event,
\begin{equation}
\label{eq-complete-learning-oracle-sandwich}
\operatorname{App}_{I}(\mathcal H_B)
\le R_{0,I}(\widehat h)-R_{0,I}^*
\le
\operatorname{App}_{I}(\mathcal H_B)
+\frac{2\Gamma_{I}(B,m,\eta,\delta)
       +\varepsilon_{\rm opt}}{1-2\eta}.
\end{equation}
The radius may be the minimum of the Rademacher, structural PAC--Bayes,
effective-dimension, finite-description, mixing-block, sequential, and
martingale candidates already proved in this part, provided that every
candidate controls the same corrupted population loss on the same
simultaneous event.  Behavior-to-deployment, observation reconstruction,
and model-channel terms enter additively only after their represented
same-loss converters have been applied.
For a source-selected class, the simultaneous radius includes the exact-cell
selection term in
\cref{eq-source-resolved-learning-rademacher-union} or the mixture divergence
in \cref{eq-source-resolved-learning-pac-bayes-mixture}.  Each
target-dependent reveal retains its source-assigned noise contraction in
\cref{eq-source-resolved-learning-information-ledger}.  If the deployed
output is used for target selection, its success is additionally bounded by
\cref{eq-source-resolved-learning-target-gain-minimum}; the risk sandwich
itself remains in clean-risk units.

Taking the infimum over all complete records of cost at most \(B\) gives the
exact attainable-risk lower floor
\begin{equation}
\label{eq-complete-learning-saturated-risk-floor}
\mathbb E_RR_{0,I}(h_R)-R_{0,I}^*
\ge\operatorname{App}_{I,m,\eta}^{\rm sat}(B)
\end{equation}
for every such record.  Any record satisfying
\cref{eq-complete-learning-corrupted-uniform-radius,eq-complete-learning-empirical-optimization-error}
supplies the corresponding upper certificate in
\cref{eq-complete-learning-oracle-sandwich} at its complete cost.

\end{theorem}

\begin{proof}

Add and subtract \(\inf_{h\in\mathcal H_B}R_{0,I}(h)\) to obtain
\cref{eq-exact-learning-approximation-estimation-decomposition}.  The second
term is nonnegative by membership of \(\widehat h\) in \(\mathcal H_B\),
which proves the lower bound.

Let \(h_B^*\) be an \(\epsilon\)-minimizer of clean risk in
\(\mathcal H_B\).  Homogeneous symmetric corruption gives the exact
identity
\[
R_{\eta,I}(h)-R_{\eta,I}(h_B^*)
=\rho\bigl(R_{0,I}(h)-R_{0,I}(h_B^*)\bigr).
\]
On the simultaneous event,
\begin{align*}
R_{\eta,I}(\widehat h)
&\le\widehat R_{\eta,S}(\widehat h)+\Gamma_I\\
&\le\widehat R_{\eta,S}(h_B^*)
   +\varepsilon_{\rm opt}+\Gamma_I\\
&\le R_{\eta,I}(h_B^*)
   +2\Gamma_I+\varepsilon_{\rm opt}.
\end{align*}
Divide by \(\rho\), let \(\epsilon\downarrow0\), and substitute the exact
decomposition.  This proves
\cref{eq-complete-learning-oracle-sandwich}.  The candidate-minimum rule is
valid because all candidates bound the same quantity on one intersection
event.  The saturated lower floor follows directly from the infimum in
\cref{eq-saturated-attainable-clean-risk}.

\end{proof}

\begin{corollary}[Quantitative parameter limits]
\label{cor-complete-learning-parameter-limits}

In the setting of
\cref{thm-complete-quantitative-learning-decomposition}, the clean case
\(\eta=0\) removes the corruption factor.  Increasing \(m\) can reduce the
statistical radius but leaves the approximation floor unchanged.  As
\(\eta\uparrow1/2\), a corruption-transport certificate becomes neutral
unless its numerator vanishes at the same rate.  Increasing the saturated
budget can only decrease
\(\operatorname{App}_{I,m,\eta}^{\rm sat}(B)\), because it enlarges the
complete record class.

\end{corollary}

\begin{proof}

The first three assertions follow from
\cref{eq-complete-learning-oracle-sandwich}.  Budget monotonicity follows by
inclusion of the saturated cost sublevel sets in
\cref{eq-saturated-attainable-clean-risk}.

\end{proof}

\begin{figure}[tbp]
\centering
\resizebox{.96\linewidth}{!}{%
\begin{tikzpicture}[fw diagram,node distance=7mm and 10mm]
  \node[fw source,text width=2.7cm] (presentation)
    {presentation, budget\\and comparison class};
  \node[fw provenance,text width=2.8cm,right=of presentation] (app)
    {attainable risk\\structural approximation};
  \node[fw use,text width=2.7cm,above right=7mm and 10mm of app] (stat)
    {sample complexity\\and noise transport};
  \node[fw use,text width=2.7cm,below right=7mm and 10mm of app] (opt)
    {represented optimizer\\and deployment shift};
  \node[fw terminal,text width=3.0cm,right=17mm of app] (risk)
    {two-sided clean\\population-risk bound};
  \draw[fw arrow] (presentation)--(app);
  \draw[fw arrow] (app)--(stat);
  \draw[fw arrow] (app)--(opt);
  \draw[fw arrow] (stat)--(risk);
  \draw[fw arrow] (opt)--(risk);
\end{tikzpicture}%
}
\caption{Complete quantitative learning decomposition.  The saturated
budget fixes the attainable population-risk floor.  Statistical estimation,
noise transport, represented optimization, and deployment comparison bound
the excess above that floor.  Exact source-cell selection and source-assigned
noise contraction are supplied by
\cref{thm-source-resolved-saturated-learning-compilation}.}
\label{fig-complete-quantitative-learning-decomposition}
\end{figure}

\section{Summary}\label{summary}
\label{sec:learning-computability-summary}

This part interprets the saturated invariant in statistical terms.

\paragraph{Dependency trail.}
The imported objects, kernel spectrum, and admissible capacity sets are
\cref{def-paper4-dependencies,def-mercer-observable-spectrum,def-admissible-learning-kernel,def-learning-affordability}.
The operator and neural coordinates are
\cref{def-structural-energy-ball,def-taskcap-structcap,def-random-label-null-score,def-krr-estimator,def-ntk,def-kernel-movement,def-intrinsic-harmonic-dimension}.
The RKHS and effective-dimension bridge is
\cref{thm-paper4-dependency-principle,cor-choosing-kernel-choosing-observable,prop-effective-dimension-cost}.
The operator and kernel bounds are
\cref{lem-ellipsoid-rademacher,thm-learned-operator-union-bound,cor-operator-selection-conservation,thm-random-label-operator-hardness,cor-source-condition-rate,cor-family-floor-learning-hard-side,thm-kernel-trick-lift,cor-lift-hard-side}.
The neural-kernel and observable-success results are
\cref{thm-lazy-ntk-bound,thm-observable-success-certifiability,cor-parameter-cover-selection-cost,cor-feature-learning-conservation,thm-composite-nn-bound}.
The symmetric-label-corruption results are
\cref{lem-exact-label-noise-risk-transport,thm-rademacher-label-corruption-bound,cor-capacity-controlled-noise-tolerance,thm-structural-pac-bayes-label-noise,cor-debiased-krr-label-noise,def-budget-saturated-statistical-readout-class,thm-saturated-statistical-attack-alignment,cor-statistical-alignment-forces-saturated-structure}.
The empirical/population provenance split and its residual-noise-order
specialization are
\cref{thm-two-level-learned-readout-provenance,cor-learned-readout-residual-noise-order}.
The execution-law population profiles, memorization separation, prospective
learning accountability, and alignment-cost criterion are
\cref{def-average-seed-population-alignment-profile,thm-population-alignment-generalization-memorization,thm-noise-indexed-learning-accountability,cor-noise-indexed-learning-threshold-criterion}.
The dependent-data extension, sequential capacity, martingale posterior
certificate, and trajectory provenance split are
\cref{def-dynamical-sample-certificate,thm-beta-mixing-structural-blocking,def-sequential-structural-rademacher,lem-sequential-energy-support,thm-sequential-learned-operator-cover,thm-martingale-structural-pac-bayes,thm-structural-online-regret,cor-structural-online-to-batch,thm-trajectory-population-memorization-separation,cor-dynamical-source-condition-rate}.
The uniform dynamical label-noise extension, including its computational,
sample, statistical-noise, and information profiles, is
\cref{def-complete-dynamical-label-noise-record,thm-dynamical-label-noise-conditional-transport,def-saturated-dynamical-noise-profiles,thm-uniform-dynamical-label-noise-bound,def-dynamical-noise-thresholds,thm-dynamical-noise-three-thresholds}.
The reward, execution, observation, filter, and composite-channel statistical
extension is
\cref{def-complete-dynamical-model-channel-noise-record,thm-dynamical-affine-reward-channel-transport,def-dynamical-observation-reconstruction,def-saturated-dynamical-channel-error-profiles,thm-finite-alphabet-channel-estimation-entropy}.
The geometry-aware learning and affordable-selection bounds are
\cref{lem-learned-geom-target-reach-projection,prop-noise-corrected-actionable-geom-extraction,thm-saturated-geometry-learning-selection-accountability,cor-complete-structural-exhaustion-learned-sources,cor-finite-description-saturated-learning-radius}.
The common source-resolved compilation that assigns every statistical,
Bayesian, sequential, active, curiosity, and controlled-learning coordinate
to its exact structural source, noise contraction, complete cost, and
represented target converter is
\cref{def-complete-source-resolved-learning-ledger,thm-source-resolved-saturated-learning-compilation,cor-source-resolved-noise-data-complexity-learning};
its dependency diagram and audit table are
\cref{fig-source-resolved-learning-compilation,tab-source-resolved-learning-audit}.
The attainable-risk profile and complete approximation--estimation
decomposition are
\cref{def-saturated-attainable-risk-approximation-floor,thm-complete-quantitative-learning-decomposition,cor-complete-learning-parameter-limits}.
The noise-indexed recovery cost, degradation law, and statistical--information
threshold comparison are
\cref{def-noise-indexed-saturated-recovery-cost,thm-noise-degradation-monotonicity,thm-three-noise-threshold-comparison}.
The effective-degree and certification results are
\cref{thm-degree-spectrum-upper,thm-effective-degree-sandwich,prop-effective-degree-error-bars,cor-two-computable-coefficient-certificates,thm-saturation-optimality,cor-optimal-up-to-model-choice,thm-two-level-geometry-certification,cor-complete-methodology-optimality,thm-defect-profile-selection,thm-paper4-summary}.
The final SAT reading and the defect-geometry certificate profile are
located at \cref{ex-final-sat-summary,defect-geometry-inequality-profile-certificates}.

\begin{lemma}[Learning capacity summary]\label{thm-paper4-summary}

For a budget-admissible oracle-free learner on a presented task family,
generalization requires target structure captured by an observable
algebra within the declared budget. At the operator level, the
fixed-operator complexity is \(\operatorname{Tr}(A^{-1})\), target
regularity is \(y^TAy\), and a learned operator incurs the selection term
\(\sqrt{\log\mathcal N(\Theta,\rho)/m}\). For a fixed kernel, the
observable algebra is a budget-admissible RKHS. After declaring a
monotone numerical statistical-capacity map
\(\mathsf n_{\mathrm{stat}}:C\to[0,\infty]\), the capacity condition is
\(N_{\mathrm{eff}}(\lambda)\le
\mathsf n_{\mathrm{stat}}(b(n))\), and captured structure is
target-aligned spectral mass satisfying the provenance, representation,
compatibility, and family-specific admissibility conditions of Parts
I--III. For a neural network, the represented empirical tangent kernel
defines the observable algebra. The lazy-training regime uses the
initial kernel, whereas feature learning incurs the cost of selecting
the terminal kernel. The observable-success theorem gives every
successful finite-horizon computation an exact first-hit witness at
scale \(\Gamma/(eH_B)\). A separately declared Rademacher audit resolves
residual statistics only when its sufficient error bound is smaller than
that scale; otherwise it is inconclusive. The learned-kernel spectrum determines the effective
degree, with the stated finite-sample errors. The structural PAC--Bayes
bound uses the domain-geometry KL divergence as its complexity term and
identifies effective resistance with the Poincar\'e component of the
defect-geometry profile. Under homogeneous symmetric label corruption, exact
risk transport rescales the Rademacher and PAC--Bayes clean-risk certificates
by \((1-2\eta)^{-1}\). Debiased kernel ridge rescales its variance term by
\((1-2\eta)^{-2}\) and preserves the source bias; the capacity and represented
model-selection costs are unchanged.  The saturated statistical-attack
theorem applies these bounds simultaneously to every represented learner and
its data-dependent output in one budget slice; target-bearing uses retain
their full ancestry, and enter the prospective Geom/Abs ledger only after the
same-record channel and temporal gates are verified. The average-seed alignment and generalizing-projection
profiles evaluate the realized execution law rather than a worst seed in the
sample-independent concentration hull.
For dynamical samples, independent episodes retain the independent-sample
certificate, stationary absolutely regular trajectories use the certified
blocking bound, and adaptive paths use sequential structural Rademacher
complexity or the martingale structural PAC--Bayes certificate.  The
pathwise inverse-precision radius and operator-cover term charge the geometry
used along the history.  Online-to-batch conversion preserves the represented
deployment cost, and trajectory-indexed memorization contributes to a fresh
execution only through its qualified population projection.  Source-condition
rates retain their exponent with the independent sample count replaced by the
certified dynamical effective sample size.
For binary labels on dynamical samples, conditional corruption transport
adds the exact margin \((1-2\eta_+)^{-1}\), while the uniform excess envelope
retains problem size \(n\), computational ceiling \(b(n)\), and sample size
\(m\).  The associated sample, statistical-noise, and computational-noise
thresholds specialize to polynomial capacity by \(b_C(n)=n^C\).  Adaptive
information is bounded by the initial public information plus
\(m(1-H_2(\eta))\) bits.  General reward, execution, observation, filter, and
composite-kernel corruption uses complete typed records.  Affine reward
transport gives a debiased behavior-law \(L^2\) bound and, with a represented
conditional-reward shift, a deployment bound.  For fixed finite channels,
the iid row radius depends on
\((m,\delta,d_U,d_V,\underline p_*)\); \(n\) and \(b(n)\) index the complete
record and its computation.  Observation-law risk alone converts to
\(+\infty\) clean-state error, whereas a represented noninvertible decoder
may yield a finite but nonvanishing reconstruction floor, at least \(1/2\)
in the binary identical-row example.  These quantities remain separate typed
model-error coordinates.  No model error, entropy, or channel-capacity
coordinate enters \(\Geom/\Caus/\Lift/\Bdry\) by itself: target-aligned use
requires the same complete record's model-to-Bellman, path, or contact
converter, temporal gate, and charged ancestry.
\Cref{thm-saturated-geometry-learning-selection-accountability} bounds learned
target reach by the generated-algebra projection envelope and charges the
represented decoder and selector in the same derivation. Prospective learning
accountability separates neutral contact
from target-aligned selector advantage, while
population-orthogonal memorization remains in \(\Res\). The noise-indexed inverse success profile defines
the exact budgeted computational endpoint, degradation makes its affordable
region downward closed, and Fano information bounds supply an
algorithm-independent upper endpoint. Ridge,
structural-loss, and profile bounds are
selected by minimizing a simultaneous bound over the declared finite
set. The intrinsic-harmonic-dimension theorem gives the optimal capacity
and coverage gap for the declared geometry. Once the model family covers
the structural subspace, no budget-admissible learner has greater
certified structural extraction relative to that geometry. Geometry
learning is handled by a label-independent self-defect criterion, with
the geometry-estimation error propagated to the model-level bound.
Every learned prospective source retains the complete five-mode,
state-provenance, and coefficient-provenance classification of
\cref{cor-complete-structural-exhaustion-learned-sources}.  Thus learning
introduces no coarser structural exception: its source belongs to one of the
4,960 existing Part~III cells, with the exact quotient or ambient Geom
visibility and every qualification gate evaluated on the same complete
record.
The attainable-risk profile adds the complementary optimization direction:
it infimizes clean population risk over complete affordable learning records.
The resulting oracle sandwich separates the structural approximation floor
from estimation, homogeneous-noise transport, represented optimization, and
deployment shift on one simultaneous event.

\end{lemma}

\begin{proof}\phantomsection\label{proof-paper4-summary}

\PartII decomposes budget-admissible oracle-free dynamics into the
exposed structural channels and residual accounting. \PartIII shows that
non-negligible structured arrival requires non-negligible
\(I_n^{\mathrm{sat}}(b(n))\) at the declared budget. A generalizing
predictor has non-negligible target correlation on unobserved states.
The fixed-operator Rademacher theorem bounds this correlation through
the trace term, and the learned-operator theorem includes the cost of
data-dependent operator selection. The random-label theorem supplies the
corresponding null lower bound for budget-coverable operator classes.

For a fixed budget-admissible kernel, the counted correlations are the
Mercer modes of the RKHS after the quotient and lift admissibility
conditions are imposed. The kernel ridge theorem in
\Cref{the-tightened-two-sided-learning-bound} gives matching
effective-dimension, source-residual, and noise terms on the standard
source classes, together with the optimized source-condition rates. The
PAC--Bayes results in
\Cref{structural-deployment-certificates} establish the validity
of a label-independent data-geometry prior and identify its KL
divergence with Dirichlet/effective-resistance energy. The same section
extends the Poincar\'e bound to defect-geometry profiles under verified
MLSI, Bakry--\'Emery, capacity/Cheeger, or transport inequalities;
includes coefficient and profile selection through finite-cover terms;
proves saturation at the intrinsic harmonic dimension; and specifies a
label-independent procedure for estimating the data geometry.

For neural networks, the fixed-kernel statement applies to the initial
NTK in the lazy-training regime and to the represented terminal kernel,
together with its covering cost, in the feature-learning regime. The
observable-success theorem applies the same accounting to a successful
budget-admissible computation: a represented target event is classified
as source capture, residual capture controlled by Rademacher complexity,
or a capacity barrier. Finally, the learned-kernel spectral profile gives
the effective degree, while the test error lower-bounds the captured mass
as in \Cref{measuring-the-effective-degree-a-network-learned}.
The sample-independent saturated readout hull in
\cref{def-budget-saturated-statistical-readout-class} contains every
data-dependent output before the realized sample is selected.
\Cref{thm-saturated-statistical-attack-alignment} therefore applies one
simultaneous concentration event to all represented statistical attacks and
uses \cref{thm-generated-algebra-projection-ceiling} to bound their balanced
clean-target alignment by the learned projection envelope.
\Cref{thm-two-level-learned-readout-provenance} then keeps empirical
memorization on its charged finite carrier, assigns only the orthogonal
population remainder to \(\Res\), and routes every represented target-bearing
use through the existing source and selector accounting.
\Cref{thm-population-alignment-generalization-memorization} replaces the
worst-seed projection envelope by the execution-law average needed for recovery,
and \cref{thm-noise-indexed-learning-accountability} separates its neutral and
prospective contributions before taking a noise threshold.
\Cref{thm-beta-mixing-structural-blocking} replaces independent
symmetrization by coupled block symmetrization on an absolutely regular
trajectory.  On predictable paths,
\cref{lem-sequential-energy-support,thm-sequential-learned-operator-cover}
express the sequential complexity through the same inverse structural
precision and selection cover as the independent operator bounds, while
\cref{thm-martingale-structural-pac-bayes} uses the existing proper
defect-restricted prior in a conditional exponential supermartingale.
\Cref{thm-trajectory-population-memorization-separation} applies the existing
population projection and residual normal form on a fresh execution carrier.
These results establish the stated dependence of generalization on
captured structure within the declared budget.
\Cref{thm-three-noise-threshold-comparison} then places the represented
statistical certificate, exact saturated computational crossing, and
information converse on the same noise scale.
The general model-channel extension follows by fresh-report conditional-mean
orthogonality, represented decoder or stability comparisons, and finite-row
multinomial concentration.  On the single event
\(\mathcal E_{I,q,\mathsf N}^{\rm sat}\), every candidate bounds the same
true discrepancy for every affordable data-selected record; hence the true
record error is at most the candidate minimum, and taking the affordable
supremum proves
\cref{eq-saturated-channel-true-error-bound}.  These are precisely
\cref{thm-dynamical-affine-reward-channel-transport,def-dynamical-observation-reconstruction,def-saturated-dynamical-channel-error-profiles,thm-finite-alphabet-channel-estimation-entropy}.

\end{proof}

Every operational certificate term is computed from represented finite data
or supplied as a represented assumption with a coverage event.  Analytic
comparators and true coverage probabilities are not learner inputs.
For an operator, they include \(\operatorname{Tr}(A^{-1})\), \(y^TAy\),
\(\operatorname{TaskCap}\), and the random-label null statistic. For a
Gram matrix, they include \(N_{\mathrm{eff}}\), kernel--target
alignment, candidate captured mass, source residual, and participation
degree. A trained network additionally provides kernel movement,
terminal-kernel improvement, the parameter-cover selection term, and a
spectral-margin bound. Subject to the admissibility conditions, these
quantities instantiate the \(I_n^{\mathrm{sat}}(B)\) and \(B^*\)
accounting for learnability, feature learning, and model-class lower
bounds.

Standard references for the statistical pieces include Aronszajn's RKHS
theory, Mercer's theorem, the representer theorem of Kimeldorf-Wahba and
later RKHS treatments, Rademacher contraction and Dudley/covering
bounds, Hanson-Wright concentration, Caponnetto-De Vito
effective-dimension minimax rates for regularized least squares,
Cristianini-Shawe-Taylor-Kandola kernel-target alignment, Davis-Kahan
and matrix Bernstein spectral concentration, Jacot-Gabriel-Hongler
neural tangent kernels, infinite-width NTK analyses such as Arora et
al., and Bartlett-Foster-Telgarsky spectral-margin bounds
\citep{Aronszajn1950,Mercer1909,KimeldorfWahba1971,
BartlettMendelson2002,Dudley1967,RudelsonVershynin2013,
CaponnettoDeVito2007,CristianiniEtAl2002,DavisKahan1970,Tropp2012,
JacotEtAl2018,AroraEtAl2019,BartlettFosterTelgarsky2017}.

\clearpage
\part{Dynamical-Systems Geometry and Controlled Structural Flow}
\label{part:dynamical-systems}\setcounter{section}{-1}

\section{Introduction: Controlled Dynamics in the Saturated Framework}
\label{sec-controlled-dynamics-introduction}

\paragraph{Part contract.}
\emph{Inputs:} the represented task presentation and channel decomposition of
\cref{part:task-spaces,part:dynamics}, the saturated source and accountability
theorems of \cref{part:spectral-complexity}, and the learning certificates of
\cref{part:learning}.
\emph{Output:} a single framework-level geometry for represented controlled
dynamical systems and coordinatewise bounds for arrival, reach--avoid,
mixing, smoothing, concentration, hitting time, Bellman error, policy loss,
learned control, and optimal structural active sampling.
\emph{First pass:} read
\cref{def-controlled-dynamical-presentation,thm-controlled-policy-channel-decomposition,def-certified-dynamical-geometry,thm-controlled-functional-inequality-consequences,thm-controlled-committor-bellman,thm-universal-dynamical-certificate}.
For affordable directed flow and lifted carriers, read
\cref{thm-controlled-caus-action-aiming,thm-controlled-lift-capacity-schur,thm-controlled-valid-lift-block-resolvent,thm-controlled-clocked-lift-transfer}.
For binary target interfaces learned from corrupted labels, read
\cref{thm-controlled-noise-corrected-caus-aiming,thm-controlled-noisy-gate-lift-transfer}.
For reward, execution, observation, and stochastic partial-observation noise,
read
\cref{thm-controlled-execution-channel-comparison,thm-controlled-multichannel-path-comparison,thm-controlled-pomdp-clocked-embedding,thm-controlled-pomdp-filter-transfer,thm-controlled-multichannel-information-bound,def-controlled-typed-entropy-outputs}.
For learned control, continue with
\cref{def-complete-dynamical-learning-record,thm-master-structural-dynamical-learning-certificate}.
For adaptive data acquisition and candidate transport, read
\cref{def-complete-master-active-sampling-record,thm-active-optimality-saturated-construction,thm-master-optimal-structural-active-sampling-envelope,thm-learned-active-rule-finite-horizon-loss}.
For the description-length axis, universal simulation, and its quantitative
relation to the structural learner, read
\cref{def-length-bounded-active-selector-profiles,thm-length-bounded-levin-realization,thm-structural-learner-length-bounded-levin,cor-length-bounded-levin-relation-lpn}.

The preceding parts already contain the structural ingredients of a
dynamical-system analysis: positive generators, reversible conductance,
authorized directed current, represented quotients and lifts, boundary
readouts, killed resolvents, temporal source admissibility, and learned
geometry.  This part assembles those ingredients for controlled dynamics.
An action rule selects a represented generator; it creates no additional
structural channel.  Its geometry is evaluated by the same saturated
\(\Geom/\Caus/\Abs\) ledger used for search and learning.

The direct setting is a finite represented carrier.  A deterministic system
is a point kernel, a randomized system is a Markov kernel, and a
history-dependent controller becomes Markovian on its clocked transcript
carrier.  Continuous and stochastic systems enter through a represented
finite or clocked lift together with an explicit generator or form-comparison
error.  Thus every claimed bound has a finite charged record on which its
constants can be evaluated.

Constructor exhaustion and numerical evaluation play distinct roles.  The
channel theorems classify every valid Caus and Lift record.  The saturated
action, aiming, capacity-amplification, and lift-access profiles introduced
below evaluate uniform suprema over the affordable instances of those
classes.

\section{Represented Controlled Dynamical Systems}
\label{sec-represented-controlled-dynamical-systems}

\begin{definition}[Controlled dynamical presentation]
\label{def-controlled-dynamical-presentation}

Let \(I\) be a finite task presentation in the sense of
\cref{def-task-space-signature}.  A controlled dynamical presentation over
\(I\) consists of the following represented data:

\begin{enumerate}[label=\textup{(\roman*)}]
\item a finite action alphabet \(A_I\) and a nonempty legal-action set
      \(\mathcal A_I(x)\subseteq A_I\) for every native state \(x\in X_I\);
\item for every \(a\in A_I\), a stochastic kernel
      \(P_I^a(x,z)\) whose evaluator, retained coefficients, precision, and
      support are generated by the declared constructor grammar;
\item a represented initial law \(\mu_{0,I}\), positive and negative terminal
      sectors \(T_I^+\) and \(T_I^-\), and, when a return criterion is used, a
      bounded represented reward \(r_I:X_I\times A_I\to[-R_I,R_I]\);
\item the complete resource record for action testing, kernel evaluation,
      random-bit use, state update, reward and terminal readout, and storage.
\end{enumerate}

A finite-horizon represented policy is a family of kernels
\[
\pi_t(\,\cdot\mid\mathcal H_t)
\quad (0\le t<H)
\]
supported on \(\mathcal A_I(x_t)\), where \(\mathcal H_t\) is the represented
pre-transition history.  The policy record contains one family-uniform
derivation of these kernels and charges its evaluation, retained state,
randomness, and selection operations.  A stationary Markov policy has the
special form \(\pi(a\mid x)\).

For a stationary policy define
\[
P_I^\pi(x,z)
=
\sum_{a\in\mathcal A_I(x)}\pi(a\mid x)P_I^a(x,z),
\qquad
L_I^\pi=\rho(P_I^\pi-I),\quad \rho>0.
\]
The scale \(\rho\) and its represented precision belong to the same resource
record.  The presentation is budget-admissible at \(B\) when all data used by
the selected finite computation have complete saturated cost at most \(B\).

\end{definition}

\begin{theorem}[Clocked policy embedding and exhaustive channel decomposition]
\label{thm-controlled-policy-channel-decomposition}

Fix a controlled dynamical presentation, a horizon \(H\), and a represented
policy of complete cost at most \(B\).  There is a represented clocked
transcript carrier \(\Omega_{I,\pi}^{[H]}\) and a time-homogeneous Markov kernel
\(P_{I,\pi}^{[H]}\) on that carrier whose projected native-state law is exactly
the law generated by the original policy through time \(H\).  Its complete
construction cost is at most a class-preserving enlarged ceiling
\(\Psi_{\rm ctl}(B,H,n)\).

The normalized transcript generator has the exhaustive decomposition
\begin{equation}
\label{eq-controlled-policy-channel-decomposition}
L_{I,\pi}^{[H]}
=L^{\Geom}+L^{\Caus}+L^{\Abs}+L^{\Lift}+L^{\Bdry}+L^{\Res}.
\end{equation}
Here \(L^{\Geom}\) is the represented reversible carrier,
\(L^{\Caus}\) is the authorized Hodge projection of the selected actual
current, \(L^{\Abs}\) contains qualified represented quotient motion,
\(L^{\Lift}\) is the valid history/interface lift, \(L^{\Bdry}\) contains
terminal and reward readouts, and \(L^{\Res}\) is the post-exhaustion
orthogonal remainder.  The policy selection and every action-dependent
coefficient remain ancestors of the corresponding component.  Derived
randomization, clocks, Bellman updates, stored values, and policy mixtures
carry no independent structural source charge.

The same statement holds for every finite slice of an unbounded execution.
Deterministic dynamics are included by taking each kernel to be a point mass.

\end{theorem}

\begin{proof}

The full record at time \(t\) contains the native state, clock, policy memory,
revealed random bits, selected action, reward and terminal flags, and every
retained auxiliary coordinate.  Conditional on that record, one represented
policy evaluation followed by one represented environment-kernel evaluation
determines the next-record kernel.  This is the finite transcript-kernel
construction of \cref{thm-algorithm-kernel-embedding}.  Composition closure
and the execution ledger give the ceiling \(\Psi_{\rm ctl}(B,H,n)\), while
projection to the native coordinate preserves the original finite-horizon
law.

Apply \cref{thm-channel-exhaustiveness} to the positive transcript generator,
then apply
\cref{thm-reversible-geom-exhaustive-directed-classification,thm-exhaustive-lift-normal-form,thm-exhaustive-bdry-normal-form,thm-exhaustive-residual-normal-form}.
The policy evaluator and selected action are coordinates of the valid
transcript lift, so their derivations remain in the backward ancestry of every
selected transition.  No-creation under represented composition is
\cref{prop-derived-channels-create-no-structure}.  A deterministic update is
the stochastic kernel with one nonzero entry in each row.  Applying the same
finite construction at each horizon proves the final assertion.

\end{proof}

\Cref{fig-controlled-dynamics-channel-map} records the dependency of the
controlled generator on the public presentation and the selected policy.

\begin{figure}[tbp]
\centering
\begin{tikzpicture}[fw diagram,node distance=8mm and 10mm]
  \node[fw source,text width=2.6cm] (pres) {task presentation\\actions and kernels};
  \node[fw provenance,text width=2.5cm,right=of pres] (pol)
    {represented policy\\complete cost};
  \node[fw box,text width=2.4cm,right=of pol] (gen)
    {clocked selected\\generator};
  \node[fw geom,below right=7mm and -2mm of gen] (geom) {\(\Geom\)};
  \node[fw use,right=4mm of geom] (caus) {\(\Caus\)};
  \node[fw provenance,right=4mm of caus] (abs) {\(\Abs\)};
  \node[fw source,below=5mm of geom] (lift) {\(\Lift\)};
  \node[fw terminal,right=4mm of lift] (bdry) {\(\Bdry\)};
  \node[fw residual,right=4mm of bdry] (res) {\(J_{\rm res}\)};
  \draw[fw arrow] (pres) -- (pol);
  \draw[fw arrow] (pol) -- (gen);
  \draw[fw arrow] (gen) |- (geom);
  \draw[fw arrow] (gen) |- (caus);
  \draw[fw arrow] (gen) |- (abs);
  \draw[fw arrow] (gen) |- (lift);
  \draw[fw arrow] (gen) |- (bdry);
  \draw[fw arrow] (gen) |- (res);
\end{tikzpicture}
\caption{A represented policy selects one cost-tagged transcript generator.
The existing exhaustive channel projection then classifies every component;
control introduces no additional structural channel.}
\label{fig-controlled-dynamics-channel-map}
\end{figure}

\subsection{Noisy execution and partially observed carriers}
\label{subsec-controlled-noisy-execution-pomdp}

The selected generator is determined by the action that is executed, not by
the action requested by the controller.  Likewise, a partially observed
controller acts on its public filtration, whereas target and failure events
are evaluated on an analytic latent carrier.  The following records make
these distinctions explicit.

\begin{definition}[Represented action-execution channel]
\label{def-controlled-action-execution-channel}

Let \(K^A(da^\circ\mid x,a,h)\) be a represented execution channel, where
\(a\) is the intended action, \(a^\circ\) is the executed action, and \(h\)
is the public pre-transition history.  For a stationary state policy, define
the actual selected reward and kernel by
\begin{align}
r^{\pi,K}(x)
&=
\sum_a\pi(a\mid x)
\sum_{a^\circ}K^A(a^\circ\mid x,a)
r^{a^\circ}(x),
\label{eq-controlled-executed-reward}\\
P^{\pi,K}(x,z)
&=
\sum_a\pi(a\mid x)
\sum_{a^\circ}K^A(a^\circ\mid x,a)
P^{a^\circ}(x,z).
\label{eq-controlled-executed-kernel}
\end{align}
The history-dependent version is the corresponding clocked kernel on the
transcript carrier.  It includes the channel evaluator, random bits,
execution feedback, and retained channel state.

The kernel in \cref{eq-controlled-executed-kernel}, rather than the intended
kernel, supplies the invariant law, symmetric form, authorized current,
Caus action, capacity, committor, and target-arrival law.  A publicly fixed,
target-independent channel has inherited randomization provenance.  Every
state-, history-, learned-, or target-dependent channel retains its complete
ancestors and cost in the selected generator.  If executed actions are
unobserved, only the represented composite kernel is available for separate
estimation.

\end{definition}

\Cref{fig-executed-action-kernel} shows where actuator noise enters the
physical controlled law.

\begin{figure}[tbp]
\centering
\resizebox{\linewidth}{!}{%
\begin{tikzpicture}[fw diagram,node distance=7mm and 9mm]
  \node[fw source] (state) {state/history \((x,h)\)};
  \node[fw provenance,right=of state] (pi) {policy \(\pi\)};
  \node[fw use,right=of pi] (a) {intended \(A=a\)};
  \node[fw use,right=of a] (ka) {execution \(K^A\)};
  \node[fw terminal,right=of ka] (ao) {executed \(A^\circ\)};
  \node[fw geom,text width=2.8cm,right=of ao] (phys)
    {physical branch\\\(r^{a^\circ},P^{a^\circ}\)};
  \draw[fw arrow] (state) -- (pi);
  \draw[fw arrow] (pi) -- (a);
  \draw[fw arrow] (a) -- (ka);
  \draw[fw arrow] (ka) -- (ao);
  \draw[fw arrow] (ao) -- (phys);
  \node[fw box,text width=4.5cm,below=11mm of phys] (out)
    {\(P^{\pi,K}\Rightarrow\) invariant law, form, Caus current,
     capacity, committor, arrival};
  \draw[fw arrow] (phys) -- (out);
  \node[fw residual,text width=2.7cm,below=11mm of a] (int)
    {intended comparator\\\(r^a,P^a\)};
  \draw[fw dashed arrow] (a) -- (int);
\end{tikzpicture}%
}
\caption{The physical controlled kernel is selected by the executed action.
Averaging policy, execution channel, and physical transition produces
\(P^{\pi,K}\); this actual kernel, not the intended comparator, determines
the invariant law, Dirichlet form, current, capacity, committor, and arrival.}
\label{fig-executed-action-kernel}
\end{figure}

\begin{theorem}[Execution-channel Bellman, path, and Caus comparison]
\label{thm-controlled-execution-channel-comparison}

Fix a discount \(0<\beta<1\), a represented policy \(\pi\), and a retained
value class \(\mathcal V_B\).  Define the exact execution defect
\begin{equation}
\label{eq-controlled-execution-bellman-defect}
\varepsilon_A^{\rm Bell}
=
\sup_{v\in\mathcal V_B}
\left\|
\sum_a\pi(a\mid\cdot)
\sum_{a^\circ}K^A(a^\circ\mid\cdot,a)
\bigl[r^{a^\circ}+\beta P^{a^\circ}v
      -r^a-\beta P^av\bigr]
\right\|_\infty.
\end{equation}
If the fixed points lie in \(\mathcal V_B\), then
\begin{equation}
\label{eq-controlled-execution-value-comparison}
\|V^{\pi,K}-V^\pi\|_\infty
\le\frac{\varepsilon_A^{\rm Bell}}{1-\beta}.
\end{equation}
For a represented action metric \(d_A\), suppose every retained
\(v\in\mathcal V_B\) satisfies
\[
|r^{a^\circ}(x)-r^a(x)|
+\beta|(P^{a^\circ}-P^a)v(x)|
\le L_{A,v}d_A(a^\circ,a).
\]
Then
\begin{equation}
\label{eq-controlled-execution-metric-defect}
\varepsilon_A^{\rm Bell}
\le
\sup_{v\in\mathcal V_B}L_{A,v}
\sup_{x,a}\mathbb E_{K^A}d_A(A^\circ,a).
\end{equation}

If the actual and intended reversible comparison forms share the same
carrier and boundaries and satisfy the represented form comparison
\[
0<a_A\le b_A<\infty,
\qquad
a_A\mathcal E_0(f,f)
\le\mathcal E_K(f,f)
\le b_A\mathcal E_0(f,f)
\qquad(f|_T=1, f|_F=0),
\]
then actuator noise changes condenser capacity only within
\begin{equation}
\label{eq-controlled-execution-capacity-comparison}
a_A\operatorname{Cap}_0(T,F)
\le\operatorname{Cap}_K(T,F)
\le b_A\operatorname{Cap}_0(T,F).
\end{equation}
For a fixed state policy whose transition choice is independent of reward
reports, reward-report noise does not alter this physical state-carrier
capacity.  For a history-dependent controller, reward reports can change
later actions and therefore the actual clocked generator and its capacity.
Observation noise likewise acts through the actual joint/filter carrier or a
represented form/Lift comparison.

For a finite horizon, suppose the reference and actuator records have the
same initial law and the same conditional transition, clean-reward,
reward-report, observation, and feedback kernels once the executed action is
fixed.  Let a represented co-adapted coupling use identical downstream
innovations at every common history for which intended and executed actions
agree.  Put
\(\tau_{\rm mis}=\inf\{t:A_t^\circ\ne A_t\}\), and suppose deterministic
\(q_t\in[0,1]\) satisfy, at every common history before mismatch,
\[
\Pr\{A_t^\circ\ne A_t\mid
\mathcal H_t^{\rm cpl},\tau_{\rm mis}\ge t\}\le q_t.
\]
Every common-path
functional \(F\in[0,1]\), including target-before-failure contact, satisfies
\begin{equation}
\label{eq-controlled-execution-path-coupling}
|\mathbb E_KF-\mathbb E_0F|
\le
1-\prod_{t=0}^{H-1}(1-q_t)
\le\sum_{t=0}^{H-1}q_t.
\end{equation}
If the initial laws instead have TV distance at most
\(\varepsilon_{\rm init}\), the first right side becomes
\(1-(1-\varepsilon_{\rm init})\prod_{t<H}(1-q_t)\).
For history-dependent hazards, the exact survival probability may instead be
retained when their joint conditional mismatch law is represented; the union
bound remains valid with the expected hazards.  A discounted comparison uses
\cref{thm-controlled-valid-lift-block-resolvent} or a represented resolvent
identity; the finite-horizon sum is not inserted into a discounted slot.

Finally, suppose actual and reference authorized currents use the same
mobility \(a_t\), and the represented current difference
\(\Delta J_t=J_t^{\pi,K}-J_t^\pi\) has finite Caus action.  Then
\begin{equation}
\label{eq-controlled-execution-caus-perturbation}
\left|
\int_0^H
\bigl[
\mathcal P_{\Caus}(J_t^{\pi,K},D_t)
-\mathcal P_{\Caus}(J_t^\pi,D_t)
\bigr]dt
\right|
\le
\left(\int_0^H\mathcal A_{\Caus}(\Delta J_t;a_t)dt\right)^{1/2}
\left(\int_0^H\mathcal E_{\rm aim}(D_t;a_t)dt\right)^{1/2}.
\end{equation}
This is an aiming-power comparison.  Target contact still requires the
same-record aiming-to-contact converter and its neutral baseline.

\end{theorem}

\begin{proof}

Write \(\mathcal T_Kv=r^{\pi,K}+\beta P^{\pi,K}v\) and
\(\mathcal T_0v=r^\pi+\beta P^\pi v\).  Since \(\mathcal T_K\) is a
\(\beta\)-contraction and \(V^\pi\in\mathcal V_B\),
\begin{align*}
\|V^{\pi,K}-V^\pi\|_\infty
&\le \|\mathcal T_KV^{\pi,K}-\mathcal T_KV^\pi\|_\infty
 +\|(\mathcal T_K-\mathcal T_0)V^\pi\|_\infty\\
&\le\beta\|V^{\pi,K}-V^\pi\|_\infty
 +\varepsilon_A^{\rm Bell}.
\end{align*}
Rearrangement proves \cref{eq-controlled-execution-value-comparison}; only
the reference fixed point is needed in the retained class.  Averaging the
pointwise metric inequality first under \(K^A(\cdot\mid x,a)\) and then
under \(\pi(\cdot\mid x)\), followed by the two suprema, proves
\cref{eq-controlled-execution-metric-defect}.

For the path comparison put \(s_t=\Pr\{\tau_{\rm mis}\ge t\}\).  The
conditional hazard hypothesis gives
\[
s_{t+1}
=\Pr\{\tau_{\rm mis}\ge t, A_t^\circ=A_t\}
\ge s_t(1-q_t).
\]
Induction yields \(s_H\ge\prod_{t<H}(1-q_t)\).  The coupled full paths agree
on \(\{\tau_{\rm mis}\ge H\}\), because every downstream innovation is then
shared.  Therefore
\[
|\mathbb E_KF-\mathbb E_0F|
\le\Pr\{\tau_{\rm mis}<H\}
\le1-\prod_{t<H}(1-q_t)\le\sum_{t<H}q_t.
\]
A maximal initial coupling gives the stated
\(1-(1-\varepsilon_{\rm init})\prod_t(1-q_t)\) variant.

Let \(\mathcal A_{T,F}=\{f:f|_T=1,f|_F=0\}\).  By the variational
definition,
\[
\operatorname{Cap}_K(T,F)=\inf_{f\in\mathcal A_{T,F}}
\mathcal E_K(f,f).
\]
Taking the infimum of the lower form inequality and, separately, of the upper
form inequality proves
\(a_A\operatorname{Cap}_0\le\operatorname{Cap}_K
\le b_A\operatorname{Cap}_0\).

Finally, linearity of the Caus edge pairing in its current gives
\[
\mathcal P_{\Caus}(J_t^{\pi,K},D_t)
-\mathcal P_{\Caus}(J_t^\pi,D_t)
=-\sum_{e}\Delta J_t(e)\nabla_eD_t.
\]
Indeed, \(\mathcal P_{\Caus}(J,D)=-\sum_eJ(e)\nabla_eD\); no mobility
weight occurs in this pairing.  Inserting
\(a_t(e)^{-1/2}a_t(e)^{1/2}\) edgewise and applying Cauchy--Schwarz bounds
its absolute value by
\(\mathcal A_{\Caus}(\Delta J_t;a_t)^{1/2}
\mathcal E_{\rm aim}(D_t;a_t)^{1/2}\); Cauchy--Schwarz in time proves
\cref{eq-controlled-execution-caus-perturbation}.  This calculation uses the
same mobility and the same aiming field in both records and does not itself
bound target contact.

\end{proof}

\Cref{fig-execution-channel-comparison-coordinates} separates the four
numerical consequences of an execution-channel comparison.

\begin{figure}[tbp]
\centering
\resizebox{\linewidth}{!}{%
\begin{tikzpicture}[fw diagram,node distance=7mm and 12mm]
  \node[fw source,text width=2.8cm] (cmp) {\(K^A\) versus identity\\same complete record};
  \node[fw use,text width=3.0cm,above right=10mm and 12mm of cmp] (bell)
    {Bellman defect\\\(\varepsilon_A^{\rm Bell}/(1-\beta)\)};
  \node[fw geom,text width=3.0cm,below right=10mm and 12mm of cmp] (form)
    {form sandwich\\capacity sandwich};
  \node[fw provenance,text width=3.0cm,above right=10mm and 58mm of cmp] (path)
    {mismatch hazards\\\(1-\prod_t(1-q_t)\)};
  \node[fw terminal,text width=3.2cm,below right=10mm and 58mm of cmp] (caus)
    {current difference\\\(\sqrt{\operatorname{Act}\operatorname{Aim}}\)};
  \node[fw residual,text width=3.1cm,right=14mm of caus] (contact)
    {same-record contact converter\\required for arrival};
  \draw[fw arrow] (cmp) -- (bell);
  \draw[fw arrow] (cmp) -- (form);
  \draw[fw arrow] (cmp) -- (path);
  \draw[fw arrow] (cmp) -- (caus);
  \draw[fw dashed arrow] (caus) -- (contact);
\end{tikzpicture}%
}
\caption{Execution noise has four typed consequences: Bellman defects control
value error, form comparison controls capacity, first-mismatch coupling
controls finite-horizon path events, and current perturbation controls aiming
power.  Only a same-record aiming-to-contact theorem converts the last branch
to target arrival.}
\label{fig-execution-channel-comparison-coordinates}
\end{figure}

\begin{theorem}[Closed-loop total-variation and relative-entropy comparison]
\label{thm-controlled-multichannel-path-comparison}

Consider two complete closed-loop records on the same measurable
finite-horizon latent--public path coordinates.  Use the same within-step reveal order in both
records, and let \(j\) range over initialization and all reward, action,
observation, transition, and feedback kernels in that order.  Suppose a
represented common-history comparison gives either
\[
\sup_h\|Q_j^1(\cdot\mid h)-Q_j^0(\cdot\mid h)\|_{\rm TV}
\le\delta_j,\qquad 0\le\delta_j\le1,
\]
or
\[
\sup_h
\operatorname{KL}\!\left(
Q_j^1(\cdot\mid h)\Vert Q_j^0(\cdot\mid h)
\right)
\le\kappa_j,\qquad 0\le\kappa_j<\infty,
\]
in nats.  Then the full path laws satisfy, respectively,
\begin{align}
\|\mathbb P^1-\mathbb P^0\|_{\rm TV}
&\le1-\prod_j(1-\delta_j)
\le\sum_j\delta_j,
\label{eq-controlled-path-tv-ledger}\\
\operatorname{KL}(\mathbb P^1\Vert\mathbb P^0)
&\le\sum_j\kappa_j,
\qquad
\|\mathbb P^1-\mathbb P^0\|_{\rm TV}
\le\sqrt{\frac12\sum_j\kappa_j}.
\label{eq-controlled-path-kl-ledger}
\end{align}
Consequently every \(F\in[0,1]\) has expectation shift bounded by the
applicable right side, clipped at one.  The KL branch has value \(+\infty\)
when absolute continuity or a conditional KL certificate fails.  A
transport--entropy replacement requires a represented common path metric and
reference law.

These are path-law comparison errors, not structural source values.  Write
\(\varepsilon_{\rm path}\) for the applicable clipped TV or Pinsker right
side.  If the same target-before-failure event has a represented reference
upper bound \(U_{\rm ref}\) under \(\mathbb P^0\), then its probability under
\(\mathbb P^1\) is at most
\[
\min\{1,U_{\rm ref}+\varepsilon_{\rm path}\}.
\]
The complete-record target-contact bound is the minimum of this quantity and
the eligible prospective selector, Caus-contact, capacity-contact, and Lift
bounds for that same event.  The raw comparison error alone is not a contact
bound.

\end{theorem}

\begin{proof}

Enumerate the finite reveal slots as \(j=1,\ldots,J\), including
initialization.  Conditional on agreement through slot \(j-1\), a maximal
coupling of the two conditional rows agrees at slot \(j\) with probability at
least \(1-\delta_j\).  Induction therefore gives
\[
\Pr\{\text{all \(J\) coordinates agree}\}
\ge\prod_{j=1}^J(1-\delta_j).
\]
The coupling characterization of TV and
\(1-\prod_j(1-\delta_j)\le\sum_j\delta_j\) prove
\cref{eq-controlled-path-tv-ledger}.

When the required absolute continuity holds, the exact chain rule is
\[
\operatorname{KL}(\mathbb P^1\Vert\mathbb P^0)
=\sum_{j=1}^J
\mathbb E_{\mathbb P^1}
\operatorname{KL}\!\left(
Q_j^1(\cdot\mid H_{j-1})\Vert
Q_j^0(\cdot\mid H_{j-1})\right).
\]
The uniform common-history bound makes the right side at most
\(\sum_j\kappa_j\).  Pinsker's inequality with natural-log KL gives
\(\|\mathbb P^1-\mathbb P^0\|_{\rm TV}
\le\sqrt{\operatorname{KL}(\mathbb P^1\Vert\mathbb P^0)/2}\), proving
\cref{eq-controlled-path-kl-ledger}.  Finally, for \(0\le F\le1\), the
layer-cake representation
\(F=\int_0^1\mathbf1_{\{F>s\}}ds\) and the event definition of TV give
\(|\mathbb E_1F-\mathbb E_0F|\le
\|\mathbb P^1-\mathbb P^0\|_{\rm TV}\).

\end{proof}

\Cref{fig-multichannel-path-comparison-ledger} compares TV and KL accumulation
in the common reveal order.

\begin{figure}[tbp]
\centering
\resizebox{\linewidth}{!}{%
\begin{tikzpicture}[fw diagram,node distance=6mm and 10mm]
  \node[fw source] (p0) {reference \(\mathbb P^0\)};
  \node[fw box,right=of p0] (i0) {init};
  \node[fw box,right=of i0] (a0) {action};
  \node[fw box,right=of a0] (x0) {transition};
  \node[fw box,right=of x0] (r0) {reward/report};
  \node[fw box,right=of r0] (o0) {observation/feedback};
  \draw[fw arrow] (p0)--(i0);
  \draw[fw arrow] (i0)--(a0);
  \draw[fw arrow] (a0)--(x0);
  \draw[fw arrow] (x0)--(r0);
  \draw[fw arrow] (r0)--(o0);
  \node[fw source,below=12mm of p0] (p1) {actual \(\mathbb P^1\)};
  \node[fw box,right=of p1] (i1) {init};
  \node[fw box,right=of i1] (a1) {action};
  \node[fw box,right=of a1] (x1) {transition};
  \node[fw box,right=of x1] (r1) {reward/report};
  \node[fw box,right=of r1] (o1) {observation/feedback};
  \draw[fw arrow] (p1)--(i1);
  \draw[fw arrow] (i1)--(a1);
  \draw[fw arrow] (a1)--(x1);
  \draw[fw arrow] (x1)--(r1);
  \draw[fw arrow] (r1)--(o1);
  \draw[fw dashed arrow] (i1)--node[right,font=\scriptsize]{\(\delta_j,\kappa_j\)}(i0);
  \draw[fw dashed arrow] (a1)--(a0);
  \draw[fw dashed arrow] (x1)--(x0);
  \draw[fw dashed arrow] (r1)--(r0);
  \draw[fw dashed arrow] (o1)--(o0);
  \node[fw terminal,text width=5.1cm,right=12mm of o0,yshift=-6mm] (out)
    {path shift\\\(1-\prod_j(1-\delta_j)\) or
     \(\sqrt{\frac12\sum_j\kappa_j}\)};
  \draw[fw arrow] (o0) -- (out);
  \draw[fw arrow] (o1) -- (out);
  \node[fw residual,text width=3.4cm,below=8mm of out] (base)
    {target bound also needs\\same-event \(U_{\rm ref}\)};
  \draw[fw dashed arrow] (base) -- (out);
\end{tikzpicture}%
}
\caption{Closed-loop discrepancies accumulate in a common reveal order.
Sequential coupling gives the TV product bound, while the relative-entropy
chain rule gives a sum followed by Pinsker.  Both are same-event shifts; a
target upper bound additionally needs a represented reference-event bound.}
\label{fig-multichannel-path-comparison-ledger}
\end{figure}

\section{Stochastic POMDP Carriers and Filter Transfer}
\label{sec-controlled-pomdp-filter-transfer}

\begin{definition}[Represented stochastic POMDP record]
\label{def-represented-stochastic-pomdp-record}

A finite represented stochastic POMDP record has finite latent, action,
reward-report, observation, feedback, and retained-memory alphabets.  It
contains a latent state carrier
\(X\), intended and executed action carriers \(A,A^\circ\), transition
kernels \(P_t(dx'\mid x,a^\circ,e,h)\), a clean reward kernel \(P^R\),
reward-report channel \(K^R\), observation channel \(K^O\), a joint
execution-and-feedback channel \(K^{A,E}\), initial law, and raw latent
trigger sets \(T^{\rm raw},F^{\rm raw}\).  These trigger sets may overlap;
they are not themselves the terminal sectors of the task interface.  The special case
\(P_t(dx'\mid x,a^\circ,e,h)=P_t^{a^\circ}(dx'\mid x)\) is the usual
controlled Markov model.  Its pre-action public history is
\[
H_t^{\rm obs}
=(O_0,A_0,E_0,\widetilde R_0,O_1,\ldots,
  A_{t-1},E_{t-1},\widetilde R_{t-1},O_t),
\]
with \(H_0^{\rm obs}=(O_0)\),
where \(E_t\) denotes whatever execution feedback is actually revealed.
The controller's random action kernel and memory update are measurable with
respect to
\(\mathcal F_t^{\rm obs}=\sigma(H_t^{\rm obs},\text{public seed})\).
The analytic latent filtration \(\mathcal G_t^{\rm lat}\) also contains
\(X_{0:t}\), clean rewards, executed actions, and unrevealed environment
innovations, and
\(\mathcal F_t^{\rm obs}\subseteq\mathcal G_t^{\rm lat}\).
The coordinate \(x\) includes every persistent latent environment or channel
state on which a later row depends.  A persistent state omitted from \(x\)
is permitted only when the record proves that its conditional law is already
a function of the retained clocked coordinates.  This is the Markov
adequacy requirement, not a notational convention.

For a horizon \(H\), fix a represented tie map
\(\chi_{T^{\rm raw},F^{\rm raw}}^{\rm tie}:X
  \to\{\varnothing,\mathsf T,\mathsf F\}\) that returns
\(\varnothing\) off \(T^{\rm raw}\cup F^{\rm raw}\), returns the unique
winner label on \(T^{\rm raw}\triangle F^{\rm raw}\), and returns the
declared winner label on \(T^{\rm raw}\cap F^{\rm raw}\).  Augment
the latent record by the absorbing first-arrival winner
\[
W_0=\chi_{T^{\rm raw},F^{\rm raw}}^{\rm tie}(X_0),
\qquad
W_{t+1}=\begin{cases}
W_t,&W_t\ne\varnothing,\\
\chi_{T^{\rm raw},F^{\rm raw}}^{\rm tie}(X_{t+1}),&W_t=\varnothing.
\end{cases}
\]
and form the clocked carrier
\begin{equation}
\label{eq-controlled-pomdp-clocked-carrier}
\Omega_{\rm POMDP}^{[H]}
=\{(t,x,h,\mathfrak m,w):0\le t\le H,
\ w\in\{\varnothing,\mathsf T,\mathsf F\}\},
\end{equation}
where \(\mathfrak m\) is the represented controller/filter memory; the symbol
\(m\) remains the training-sample size.  The coordinate \(\mathfrak m\)
also contains every persistent public channel and policy-randomness state
needed by a later row.  With \(\mu_0\) the declared latent initial law, the
initial clocked law is the pushforward of
\begin{equation}
\label{eq-controlled-pomdp-initial-clocked-law}
\lambda_0(dx,do)=\mu_0(dx)K_0^O(do\mid x)
\end{equation}
under \(h=(o)\), the represented initial-memory constructor, clock zero, and
\(w=\chi_{T^{\rm raw},F^{\rm raw}}^{\rm tie}(x)\).  Any initial channel state is included in
\(x\) or \(\mathfrak m\) before this pushforward.  One transition is
the represented composition
\begin{equation}
\label{eq-controlled-pomdp-kernel-factorization}
\begin{aligned}
&\pi_t(da\mid h,\mathfrak m)
 K_t^{A,E}(da^\circ,de\mid x,a,h)
 P_t(dx'\mid x,a^\circ,e,h)\\
&\qquad {}
 P_t^R(dr\mid x,a^\circ,e,x',h)
 K_t^R(d\widetilde r\mid r,x,a^\circ,e,x',h)\\
&\qquad {}
 K_{t+1}^O(do'\mid x',a^\circ,e,r,\widetilde r,h).
\end{aligned}
\end{equation}
Every later factor is conditioned on the earlier outputs in this declared
reveal order unless a represented conditional independence removes an
argument.  The product is followed by the deterministic history, memory,
clock, and winner update.
The \(A_t^\circ\)-marginal of \(K_t^{A,E}\) is the execution kernel
\(K_t^A\) used by marginal execution comparisons.  When \(A_t^\circ\) and
\(E_t\) are correlated, path-law and information comparisons use the joint
kernel and do not multiply its two marginals.
All kernels in the factorization may be replaced by one represented
composite kernel when only that composite is identifiable.
Continuous outcome spaces enter the finite framework only through a declared
finite-precision discretization or continuum Lift with its represented
approximation and target-interface error.

The terminal interface on the augmented carrier is the disjoint pair
\begin{equation}
\label{eq-controlled-pomdp-disjoint-terminal-interface}
T_\Omega=\{(t,x,h,\mathfrak m,w):w=\mathsf T\},
\qquad
F_\Omega=\{(t,x,h,\mathfrak m,w):w=\mathsf F\}.
\end{equation}
Thus overlap is resolved before the task-space terminal interface is formed,
and \(T_\Omega\cap F_\Omega=\varnothing\).

The clean latent target and exact posterior used below are analytic
comparators.  They are not public policy inputs or prospective sources.  A
target gate, filter, belief, committor, or decoder becomes available to the
controller only through a represented
\(\mathcal F_t^{\rm obs}\)-measurable derivation with complete pre-hit cost.
The complete record charges the channel and model descriptions, random bits,
history or sufficient-memory representation, filtering, policy evaluation,
target/failure gates, precision, sample size \(m\), and every training and
deployment operation under the ceiling \(b(n)\).

\end{definition}

\Cref{fig-pomdp-one-step-factorization} expands one row of the clocked POMDP
kernel in its declared reveal order.

\begin{figure}[tbp]
\centering
\resizebox{\linewidth}{!}{%
\begin{tikzpicture}[fw diagram,node distance=6mm and 6mm]
  \node[fw source,text width=2.3cm] (s) {clocked state\\\((t,x,h,\mathfrak m,w)\)};
  \node[fw provenance,right=of s] (pi) {\(\pi_t\)};
  \node[fw use,right=of pi] (ae) {\(K_t^{A,E}\)};
  \node[fw geom,right=of ae] (p) {\(P_t\)};
  \node[fw box,right=of p] (pr) {\(P_t^R\)};
  \node[fw use,right=of pr] (kr) {\(K_t^R\)};
  \node[fw use,right=of kr] (ko) {\(K_{t+1}^O\)};
  \node[fw terminal,text width=2.5cm,right=of ko] (upd)
    {history, memory, clock, winner update};
  \draw[fw arrow] (s)--(pi);
  \draw[fw arrow] (pi)--(ae);
  \draw[fw arrow] (ae)--(p);
  \draw[fw arrow] (p)--(pr);
  \draw[fw arrow] (pr)--(kr);
  \draw[fw arrow] (kr)--(ko);
  \draw[fw arrow] (ko)--(upd);
  \node[fw residual,text width=6.0cm,below=10mm of p] (vars)
    {\(A_t\to(A_t^\circ,E_t)\to X_{t+1}\to R_t
      \to\widetilde R_t\to O_{t+1}\)};
  \draw[fw dashed arrow] (ae) -- (vars);
  \draw[fw dashed arrow] (ko) -- (vars);
\end{tikzpicture}%
}
\caption{One POMDP step is a represented composition in a fixed reveal order:
public policy, joint execution/feedback channel, latent transition, clean
reward, reward report, observation, and deterministic record update.  The
product is Markov on the full clocked carrier even for a history-dependent
controller.}
\label{fig-pomdp-one-step-factorization}
\end{figure}

\begin{theorem}[POMDP Markov embedding, source accounting, and target geometry]
\label{thm-controlled-pomdp-clocked-embedding}

The kernel factorization in
\cref{eq-controlled-pomdp-kernel-factorization} defines a positive represented
Markov kernel on \(\Omega_{\rm POMDP}^{[H]}\).  Its law projects exactly to
the declared latent and public POMDP law through time \(H\).  At the
class-preserving ceiling
\(\Psi_{\rm ctl}(b(n),H,n)\), its generator has the exhaustive decomposition
of \cref{eq-controlled-policy-channel-decomposition}.  The actual executed
kernel and observation-measurable policy remain ancestors of every selected
current and transition.

On the joint clocked carrier, the clean first-arrival event is the analytic
boundary \((T_\Omega,F_\Omega)\) of
\cref{eq-controlled-pomdp-disjoint-terminal-interface}, with the chosen tie
convention represented explicitly.  Selector and
first-arrival identities apply on this actual carrier under their temporal
hypotheses.  Native capacity, Caus action, and aiming apply only when their
own represented reference-law, reversibility or comparison-form, authority,
and boundary hypotheses hold; otherwise capacity is unavailable and
probability entries retain their neutral value.  A
target-dependent capacity, aiming, or gate computation enters an affordable
prospective minimum only when its boundary evaluator and target comparison
are represented in the same temporally admissible record; otherwise it
remains an external audit of the actual law.  Projection to an unclocked latent state need not be a
valid lift when histories with the same state choose different actions.  A
base-carrier form, capacity, or resolvent bound transfers only when the same
record supplies the required lumpability, intertwining, form comparison, or
clocked Lift defect.  Valid-lift provenance by itself is not a numerical
capacity comparison.

Stochastic observations and rewards therefore add no structural channel.
They alter the actual represented carrier, its public-filtration restriction,
and the numerical constants evaluated by the existing channels.

\end{theorem}

\begin{proof}

Fix a clocked state.  Sum the factors in
\cref{eq-controlled-pomdp-kernel-factorization} in reverse reveal order.
Each factor is a nonnegative probability kernel, so each successive integral
equals one; the deterministic history, memory, clock, and winner update is a
point mass.  The composed row is therefore nonnegative and has total mass
one.

Equality of path laws follows on cylinder events by induction.  At time zero,
\cref{eq-controlled-pomdp-initial-clocked-law} followed by the deterministic
memory and winner constructors is exactly the declared initial public--latent
law, with
\(W_0=\chi_{T^{\rm raw},F^{\rm raw}}^{\rm tie}(X_0)\).  Suppose the projected cylinder law is
correct through time \(t\).  Conditional on its terminal cylinder, the
conditional law of
\((A_t,A_t^\circ,E_t,X_{t+1},R_t,\widetilde R_t,O_{t+1})\)
is exactly the product in
\cref{eq-controlled-pomdp-kernel-factorization}; the remaining coordinates
are its deterministic record update.  The tower property therefore gives
the declared cylinder probability through time \(t+1\).  Projection proves
the asserted equality through \(H\).

Composition closure charges every factor, random bit, retained channel
state, numerical precision, policy evaluation, and history or memory update.
This gives the ceiling \(\Psi_{\rm ctl}(b(n),H,n)\).  Applying
\cref{thm-controlled-policy-channel-decomposition} to this actual clocked
kernel gives the exhaustive decomposition.  Backward ancestry closure keeps
the policy, execution, transition, reward, report, observation, and feedback
factors attached to every selected transition and current.

The winner coordinate is absorbing.  Hence \(W=\mathsf T\) is exactly the event
that the represented tie convention declared target as the first winner,
and \(W=\mathsf F\) is the analogous failure event.  This proves the boundary
identity.  Selector, capacity, Caus, and Lift conclusions are then invoked
separately only when the hypotheses listed in the theorem are present; the
Markov embedding alone supplies none of their numerical comparisons.

Finally, two clocked states can share the same latent state \(x\) while their
public histories induce different action mixtures.  Their next-state rows
then differ, so the projection to \(x\) is not lumpable in general.  A
numerical base-carrier transfer therefore requires the stated represented
comparison.  No-creation under represented composition is
\cref{prop-derived-channels-create-no-structure}.

\end{proof}

\Cref{fig-pomdp-clocked-carrier-projection} distinguishes exact path-law
projection from an unjustified unclocked numerical transfer.

\begin{figure}[tbp]
\centering
\begin{tikzpicture}[fw diagram,node distance=8mm and 13mm]
  \node[fw box,text width=6.2cm] (omega)
    {clocked carrier \(\Omega_{\rm POMDP}^{[H]}\)\\
     public \((t,h,\mathfrak m)\) \quad analytic \((x,w)\)\\
     absorbing sectors \(T_\Omega=\{W=\mathsf T\}\),
     \(F_\Omega=\{W=\mathsf F\}\)};
  \node[fw terminal,text width=3.4cm,right=of omega,yshift=9mm] (path)
    {exact projection\\declared POMDP path law};
  \node[fw residual,text width=4.2cm,right=of omega,yshift=-13mm] (latent)
    {projection to latent \(x\)\\not numerical without lumpability,
     intertwining, or form comparison};
  \draw[fw arrow] (omega) -- (path);
  \draw[fw dashed arrow] (omega) -- (latent);
\end{tikzpicture}
\caption{Winner augmentation makes first arrival an absorbing boundary event
on the actual clocked carrier.  Projection reproduces the declared POMDP path
law, but projection to the unclocked latent state need not preserve dynamics:
equal latent states can carry histories with different action mixtures.}
\label{fig-pomdp-clocked-carrier-projection}
\end{figure}

\Cref{fig-pomdp-target-affordability-gate} records the extra gates required
before analytic target geometry enters a prospective minimum.

\begin{figure}[tbp]
\centering
\resizebox{\linewidth}{!}{%
\begin{tikzpicture}[fw diagram,node distance=7mm and 10mm]
  \node[fw source,text width=3.2cm] (analytic)
    {analytic objects\\\(T_\Omega,F_\Omega,b_t,\operatorname{Cap}\)};
  \node[fw residual,text width=2.8cm,above right=9mm and 10mm of analytic] (audit)
    {external law audit};
  \node[fw provenance,text width=3.0cm,right=of analytic] (eval)
    {represented pre-hit\\evaluator};
  \node[fw provenance,text width=3.0cm,right=of eval] (cost)
    {complete target ancestry\\and cost};
  \node[fw use,text width=3.0cm,right=of cost] (conv)
    {same-record event/contact\\converter};
  \node[fw terminal,text width=2.8cm,right=of conv] (min)
    {prospective\\\(U_{\rm hit}\) minimum};
  \draw[fw dashed arrow] (analytic) -- (audit);
  \draw[fw arrow] (analytic) -- (eval);
  \draw[fw arrow] (eval) -- (cost);
  \draw[fw arrow] (cost) -- (conv);
  \draw[fw arrow] (conv) -- (min);
  \node[fw residual,text width=3.4cm,below=10mm of cost] (sent)
    {missing raw converter \(+\infty\)\\missing probability converter \(1\)};
  \draw[fw dashed arrow] (sent) -- (conv);
\end{tikzpicture}%
}
\caption{Analytic target geometry evaluates the actual law but is not an
affordable prospective source.  Entry into a probability minimum requires a
represented pre-hit evaluator, complete target-dependent ancestry, and a
same-record conversion to the declared event.}
\label{fig-pomdp-target-affordability-gate}
\end{figure}

\subsection{Approximate Filters and Transfer to Value and Arrival}
\label{subsec-controlled-pomdp-filter-transfer}

\begin{definition}[Represented belief/filter comparison]
\label{def-represented-pomdp-filter-comparison}

Let
\[
b_t(h)=\mathcal L((X_t,W_t)\mid H_t^{\rm obs}=h)
\]
be a version of the analytic exact belief on the augmented latent state.
The record declares one common clocked carrier and measurable interfaces on
which \(b_t(h)\) and the represented \(\widehat b_t(h)\) are compared for
every admitted public history \(h\in\mathcal H_t^{\rm adm}\).  Put
\[
e_t(h)=\|\widehat b_t(h)-b_t(h)\|_{\mathsf B},
\qquad
\bar e_t=\sup_{h\in\mathcal H_t^{\rm adm}}e_t(h).
\]
Here \(\mathcal H_t^{\rm adm}\) contains the union of histories reachable
under either compared clocked law.  The record fixes common measurable
versions of the two belief interfaces on that union.  At a history having
zero exact evidence but positive approximate probability, a finite
comparison therefore requires a represented extension or a positive-evidence
stability certificate; otherwise \(e_t(h)=+\infty\).  This convention
prevents a regular conditional distribution on a null history from being
chosen after seeing the desired bound.
A represented
filter \(\widehat b_t\) is admissible only when its update uses public
history and its model arithmetic, evidence normalization, retained
precision, and deployment evaluation are included in the complete cost.
It may update the augmented winner coordinate only when a represented
pre-hit target/failure evaluator, or a learned public gate with its error
interface, occurs in the filter's public ancestry.  Otherwise the represented
filter lives on the non-hit state/memory carrier and the exact winner remains
analytic evaluation coordinates.
In that non-hit branch, the \(b_t\) and \(\widehat b_t\) used in the norm
below denote the corresponding common-space state/memory marginal; the
augmented belief notation and augmented norm are used only in the represented-
gate branch.
An approximate-filter certificate contains a declared norm and constants
such that, for every admitted transition \(h\mapsto h'\),
\begin{equation}
\label{eq-controlled-filter-stability-recurrence}
e_{t+1}(h')
\le\alpha_t e_t(h)+\varepsilon_{{\rm fil},t}(h'),
\qquad
\bar e_0\le\varepsilon_0,
\end{equation}
where \(\alpha_t\) is uniform over the admitted histories.  With
\(\bar\varepsilon_{{\rm fil},t}
=\sup_{h'}\varepsilon_{{\rm fil},t}(h')\), this implies
\(\bar e_{t+1}\le\alpha_t\bar e_t+
\bar\varepsilon_{{\rm fil},t}\).  Induction gives
\begin{equation}
\label{eq-controlled-filter-stability-unrolled}
\bar e_t
\le
\left(\prod_{j=0}^{t-1}\alpha_j\right)\varepsilon_0
+\sum_{s=0}^{t-1}
\left(\prod_{j=s+1}^{t-1}\alpha_j\right)
\bar\varepsilon_{{\rm fil},s}.
\end{equation}
Empty products equal one.  If \(\alpha_t\le\alpha<1\) and
\(\bar\varepsilon_{{\rm fil},t}\le\bar\varepsilon_{\rm fil}\), then
\begin{equation}
\label{eq-controlled-filter-fast-mixing-bound}
\bar e_t\le\alpha^t\varepsilon_0
+\frac{1-\alpha^t}{1-\alpha}\bar\varepsilon_{\rm fil}.
\end{equation}
A Bayes update receives a finite \(\alpha_t\) only from a represented
evidence floor or a direct stability proof; normalization near zero evidence
is not assumed Lipschitz.

Suppose the exact- and approximate-belief policies satisfy
\begin{equation}
\label{eq-controlled-belief-policy-lipschitz}
\|\pi_t(\cdot\mid b_t(h))-\pi_t(\cdot\mid\widehat b_t(h))\|_{\rm TV}
\le L_{\pi,t}e_t(h).
\end{equation}
Let \(K_t^{\rm ex}(h,a;\cdot)\) and
\(\widehat K_t(h,a;\cdot)\) be the exact and represented approximate
same-carrier clocked rows after action \(a\).  For \(h=h(\omega)\), define
the three rows
\begin{align*}
P_t(\omega,\cdot)
&=\sum_a\pi_t(a\mid b_t(h))K_t^{\rm ex}(h,a;\cdot),\\
P_t^{\rm mid}(\omega,\cdot)
&=\sum_a\pi_t(a\mid\widehat b_t(h))K_t^{\rm ex}(h,a;\cdot),\\
\widehat P_t(\omega,\cdot)
&=\sum_a\pi_t(a\mid\widehat b_t(h))\widehat K_t(h,a;\cdot),
\end{align*}
and the remaining model defect
\[
\delta_{{\rm model},t}(h)
=\|P_t^{\rm mid}(\omega,\cdot)-\widehat P_t(\omega,\cdot)\|_{\rm TV}.
\]
Markov-kernel contraction applied to the first two rows gives
\[
\|P_t(\omega,\cdot)-P_t^{\rm mid}(\omega,\cdot)\|_{\rm TV}
\le
\|\pi_t(\cdot\mid b_t(h))-\pi_t(\cdot\mid\widehat b_t(h))\|_{\rm TV}
\le L_{\pi,t}e_t(h).
\]
The triangle inequality and the universal TV ceiling one therefore give
\begin{equation}
\label{eq-controlled-pomdp-lift-defect}
\|\widehat P_t(\omega,\cdot)-P_t(\omega,\cdot)\|_{\rm TV}
\le \epsilon_t(\omega)
:=\min\{1,L_{\pi,t}e_t(h)
             +\delta_{{\rm model},t}(h)\},
\qquad
\bar\epsilon_t:=\sup_\omega\epsilon_t(\omega).
\end{equation}
Thus \(\delta_{{\rm model},t}\) is defined, not inferred from a differently
typed model loss, and is required in that same row-total-variation norm.  The
record also contains the common-space comparisons
\[
\|\widehat\mu_0-\mu_0\|_{\rm TV}\le\varepsilon_{\rm init}
\quad\text{and}\quad
\mathbb E_{\mathbb P^{\widehat b}}\!\left|
\widehat{\mathbf1}_{T,F,H}(\omega)-\mathbf1_{T,F,H}(\omega)
\right|\le\varepsilon_{T,F,H}^{\rm gate}
\]
for the initial law and target/failure event disagreement.  The expectation
is under the approximate-filter law on the declared common path carrier.
Thus the event comparison below is the literal triangle inequality between
\(\mathbb E_{\mathbb P^b}\mathbf1_{T,F,H}\),
\(\mathbb E_{\mathbb P^{\widehat b}}\mathbf1_{T,F,H}\), and
\(\mathbb E_{\mathbb P^{\widehat b}}\widehat{\mathbf1}_{T,F,H}\).
Missing
probability interfaces take the neutral value one, whereas a requested raw
model or filter error without its norm converter is \(+\infty\).  Analytic
joint-carrier evaluation sets the gate comparison to zero only for that
analytic probability; it does not expose the latent gate to the controller
or create a prospective source.
Without a filter-stability and policy/kernel converter, the filter-to-value
and filter-to-hit errors are \(+\infty\), while probability coordinates take
their neutral value one.

\end{definition}

\Cref{fig-pomdp-filter-error-recursion} displays the exact product--sum
structure of the uniform history-wise filter certificate.

\begin{figure}[tbp]
\centering
\resizebox{.96\linewidth}{!}{%
\begin{tikzpicture}[fw diagram,node distance=7mm and 13mm]
  \node[fw source] (e0) {\(\bar e_0\)};
  \node[fw box,right=of e0] (e1) {\(\bar e_1\)};
  \node[fw box,right=of e1] (e2) {\(\bar e_2\)};
  \node[fw terminal,right=29mm of e2] (et) {\(\bar e_t\)};
  \draw[fw arrow] (e0)--node[above,font=\scriptsize]{\(\alpha_0\)}(e1);
  \draw[fw arrow] (e1)--node[above,font=\scriptsize]{\(\alpha_1\)}(e2);
  \draw[fw arrow] (e2)--node[above,font=\scriptsize]
    {\(\prod_{j=2}^{t-1}\alpha_j\)}(et);
  \node[fw use,below=10mm of e0] (d0) {\(\bar\varepsilon_{{\rm fil},0}\)};
  \node[fw use,below=10mm of e2] (d1) {\(\bar\varepsilon_{{\rm fil},1}\)};
  \draw[fw arrow] (d0)--(e1);
  \draw[fw arrow] (d1)--(e2);
  \node[fw provenance,text width=5.3cm,below=24mm of et] (mix)
    {if \(\alpha<1\): \(\bar e_t\le\alpha^t\varepsilon_0
      +(1-\alpha^t)\bar\varepsilon_{\rm fil}/(1-\alpha)\)};
  \draw[fw dashed arrow] (et)--(mix);
  \node[fw residual,text width=3.5cm,above=10mm of e2] (floor)
    {finite \(\alpha_t\) needs evidence floor or direct stability proof};
  \draw[fw dashed arrow] (floor)--(e2);
\end{tikzpicture}
}
\caption{Filter error propagates by multiplicative stability and additive
update defects.  Iteration gives the exact product--sum bound; under a uniform
contraction, initial error decays and accumulated defect approaches
\(\bar\varepsilon_{\rm fil}/(1-\alpha)\).}
\label{fig-pomdp-filter-error-recursion}
\end{figure}

\begin{theorem}[Filter error to value and target-arrival error]
\label{thm-controlled-pomdp-filter-transfer}

Under \cref{def-represented-pomdp-filter-comparison}, the exact-belief and
represented-filter clocked laws satisfy the following bound whenever their
initial-law total variation is at most \(\varepsilon_{\rm init}\):
\[
p_{T,H}^{b}
=\mathbb E_{\mathbb P^b}\mathbf1_{T,F,H},
\qquad
p_{T,H}^{\widehat b}
=\mathbb E_{\mathbb P^{\widehat b}}
 \widehat{\mathbf1}_{T,F,H}.
\]
\begin{equation}
\label{eq-controlled-pomdp-filter-hit-transfer}
|p_{T,H}^{b}-p_{T,H}^{\widehat b}|
\le
\min\left\{1,
\varepsilon_{\rm init}
+\sum_{t=0}^{H-1}\bar\epsilon_t
+\varepsilon_{T,F,H}^{\rm gate}
\right\}.
\end{equation}
Here the last term is zero for analytic evaluation on the joint carrier and
is the represented target/failure interface error when the deployed record
uses learned public gates.  Discounted arrival uses the corresponding
resolvent or fast-fiber comparison.  A represented co-adapted coupling can
replace \(\sum_t\bar\epsilon_t\) by the sharper predictable hazard
\[
\sum_{t=0}^{H-1}
\mathbb E\!\left[
\mathbf1_{\{\tau_{\rm mis}\ge t\}}\epsilon_t(\Omega_t)
\right],
\]
where
\(\tau_{\rm mis}=\inf\{t:\Omega_{t+1}^{b}\ne
\Omega_{t+1}^{\widehat b}\}\) is the first divergent transition in that
coupling; on \(\{\tau_{\rm mis}\ge t\}\), the two pre-transition states at
time \(t\) agree.
An unconditioned average row error is not a valid substitute.

For the value branch, let \(Q_B^*\) be the exact action-value function of a
declared restricted benchmark for which the Bellman and
performance-difference identities are certified.  Suppose
\[
|Q_B^*(b,a)-Q_B^*(b',a)|\le L_Q\|b-b'\|_{\mathsf B}
\]
for every retained action.  Let the deployed approximate-filter action
\(\widehat A_t\) be \(\xi_t\)-greedy at \(\widehat b_t\):
\[
\max_aQ_B^*(\widehat b_t,a)
-Q_B^*(\widehat b_t,\widehat A_t)\le\xi_t.
\]
Then, with expectation under the deployed approximate-filter policy, the
discounted benchmark loss is bounded by
\begin{equation}
\label{eq-controlled-pomdp-filter-value-transfer}
V_B^*(b_0)-V^{\widehat\pi}(b_0)
\le
\sum_{t\ge0}\beta^t
\mathbb E_{\widehat\pi}
\bigl[2L_Qe_t(H_t^{\rm obs})+\xi_t\bigr],
\end{equation}
whenever the series and all comparisons are certified.  Under
exact greediness \(\xi_t=0\) and
\cref{eq-controlled-filter-fast-mixing-bound}, the right side is at most
\begin{equation}
\label{eq-controlled-pomdp-filter-value-geometric}
2L_Q\left[
\frac{\varepsilon_0}{1-\beta\alpha}
+\frac{\bar\varepsilon_{\rm fil}}{1-\alpha}
\left(\frac1{1-\beta}-\frac1{1-\beta\alpha}\right)
\right].
\end{equation}

\end{theorem}

\begin{proof}

Let \(d_t\) be the TV distance between the two exact-gate clocked laws at
time \(t\).  Markov-kernel contraction and
\cref{eq-controlled-pomdp-lift-defect} give
\[
d_{t+1}
\le\|\mu_t^bP_t-\mu_t^{\widehat b}P_t\|_{\rm TV}
 +\|\mu_t^{\widehat b}P_t-
       \mu_t^{\widehat b}\widehat P_t\|_{\rm TV}
\le d_t+\bar\epsilon_t.
\]
Since \(d_0\le\varepsilon_{\rm init}\), induction yields
\(d_H\le\varepsilon_{\rm init}+\sum_{t<H}\bar\epsilon_t\).
The same recursion on path-prefix kernels bounds the TV distance of the full
clocked path laws.  Applying it to the exact target event and then using
\[
\left|\mathbb E_{\mathbb P^{\widehat b}}
(\mathbf1_{T,F,H}-\widehat{\mathbf1}_{T,F,H})\right|
\le\varepsilon_{T,F,H}^{\rm gate}
\]
proves \cref{eq-controlled-pomdp-filter-hit-transfer}.  Sequential maximal
coupling makes this refinement explicit.  First maximally couple the initial
laws, whose mismatch probability is at most \(\varepsilon_{\rm init}\).
Conditional on initial agreement and agreement through the state at time
\(t\), maximally couple the two rows at that shared state.  Its conditional
mismatch probability is at most \(\epsilon_t(\Omega_t)\), so
\[
\Pr\{\tau_{\rm mis}=t\}
\le\mathbb E\!\left[
\mathbf1_{\{\tau_{\rm mis}\ge t\}}\epsilon_t(\Omega_t)
\right].
\]
Summing the disjoint first-mismatch events and adding the initial and gate
terms proves the stated survival-weighted refinement.

Fix a realized pair \(b=b_t(H_t^{\rm obs})\),
\(\widehat b=\widehat b_t(H_t^{\rm obs})\).  If
\(a^*\in\arg\max_aQ_B^*(b,a)\) and
\(\widehat a=\widehat A_t\), then
\begin{align*}
V_B^*(b)-Q_B^*(b,\widehat a)
&\le |Q_B^*(b,a^*)-Q_B^*(\widehat b,a^*)|\\
&\quad +Q_B^*(\widehat b,a^*)-Q_B^*(\widehat b,\widehat a)\\
&\quad +|Q_B^*(\widehat b,\widehat a)-Q_B^*(b,\widehat a)|\\
&\le2L_Qe_t(H_t^{\rm obs})+\xi_t.
\end{align*}
The certified performance-difference identity is
\[
V_B^*(b_0)-V^{\widehat\pi}(b_0)
=\mathbb E_{\widehat\pi}\sum_{t\ge0}\beta^t
 [V_B^*(b_t)-Q_B^*(b_t,\widehat A_t)].
\]
Substitution proves \cref{eq-controlled-pomdp-filter-value-transfer}.
For exact greediness, insert the uniform bound
\(e_t(H_t^{\rm obs})\le\bar e_t\) from
\cref{eq-controlled-filter-fast-mixing-bound}.  The two sums are
\[
\sum_{t\ge0}(\beta\alpha)^t=\frac1{1-\beta\alpha},
\qquad
\sum_{t\ge0}\beta^t(1-\alpha^t)
=\frac1{1-\beta}-\frac1{1-\beta\alpha},
\]
which gives \cref{eq-controlled-pomdp-filter-value-geometric}.

\end{proof}

\Cref{fig-pomdp-filter-transfer} shows the two distinct converters from
history-wise filter error.

\begin{figure}[tbp]
\centering
\resizebox{\linewidth}{!}{%
\begin{tikzpicture}[fw diagram,node distance=7mm and 11mm]
  \node[fw source,text width=2.5cm] (e) {filter error\\\(e_t(h)\)};
  \node[fw use,text width=2.8cm,above right=10mm and 12mm of e] (pol)
    {policy TV\\\(L_{\pi,t}e_t(h)\)};
  \node[fw provenance,text width=3.0cm,right=of pol] (row)
    {clocked row defect\\\(\epsilon_t(\omega)\)};
  \node[fw terminal,text width=3.2cm,right=of row] (hit)
    {path telescoping\\initial + rows + gate};
  \node[fw use,text width=2.9cm,below right=10mm and 12mm of e] (q)
    {\(Q_B^*\)-Lipschitz\\\(2L_Qe_t+\xi_t\)};
  \node[fw terminal,text width=3.2cm,right=17mm of q] (val)
    {performance difference\\discounted value loss};
  \draw[fw arrow] (e)--(pol);
  \draw[fw arrow] (pol)--(row);
  \draw[fw arrow] (row)--(hit);
  \draw[fw arrow] (e)--(q);
  \draw[fw arrow] (q)--(val);
  \node[fw residual,text width=3.3cm,below=9mm of hit] (base)
    {arrival still needs\\represented reference baseline};
  \draw[fw dashed arrow] (base)--(hit);
\end{tikzpicture}%
}
\caption{Filter error affects deployment only through explicit converters.
Policy Lipschitzness gives a clocked-kernel row defect and hence a path-event
shift.  Exact benchmark action-value Lipschitzness plus the
performance-difference identity gives value loss.  The two output units are
not interchanged.}
\label{fig-pomdp-filter-transfer}
\end{figure}

\section{Adaptive Channel Information and Belief Entropy}
\label{sec-controlled-channel-information}

\subsection{Conditional channel-information ledger}
\label{subsec-controlled-channel-information-ledger}

\begin{definition}[Conditional channel-information and entropy ledger]
\label{def-controlled-channel-information-ledger}

Measure Shannon information in bits.  For a represented channel
\(K:\mathsf U\to\mathcal P(\mathsf V)\) and a declared nonempty represented
class of input laws \(\mathcal Q\), define
\begin{align}
C_{\mathcal Q}(K)
&=\sup_{\nu\in\mathcal Q}I_\nu(U;V),
\label{eq-controlled-conditional-channel-capacity}\\
\vartheta_{\rm KL}(K;\mathcal Q)
&=
\sup_{\substack{\nu,\nu'\in\mathcal Q\\
0<\operatorname{KL}(\nu\Vert\nu')<\infty}}
\frac{\operatorname{KL}(\nu K\Vert\nu'K)}
     {\operatorname{KL}(\nu\Vert\nu')}
\in[0,1].
\label{eq-controlled-channel-sdpi-coefficient}
\end{align}
The second display uses the same logarithm base in numerator and denominator.
If the admissible-pair set in the second supremum is empty, its value is
defined to be zero; with this convention
\(\vartheta_{\rm KL}(K;\mathcal Q)\in[0,1]\) also for singleton input-law
classes.
For a restricted class \(\mathcal Q\), an SDPI use additionally requires
\(\mathcal Q\) to contain, at each admitted public history, the conditional
marginal law of \(U\), every conditional law of \(U\) given the task index,
and the mixtures used in the mutual-information decomposition.  Taking
\(\mathcal Q=\mathcal P(\mathsf U)\) is the unrestricted default.
For a conditional history-dependent channel, the record contains predictable
upper bounds \(c_{q,t}\) and \(\vartheta_{q,t}\) valid for every admitted
public history and input law.  Here \(q\) indexes the within-step reveal order
of execution feedback, observation, reward report, and any other public
output.  A finite output alphabet gives the fallback
\(c_{q,t}\le\log_2|\mathsf V_{q,t}|\).  A continuous channel without a
represented amplitude, moment, power, or input-law restriction has capacity
\(+\infty\).

For a finite-alphabet complete POMDP record, enumerate physical channel
outputs in their actual reveal order.  Bundle every initially sampled
persistent channel state, policy seed, and other retained auxiliary random
state into \(S_0^{\rm ret}\).  If \(\mathcal F_0^{\rm clean}\) is generated
by the presentation, the clean initial task/latent state, and deterministic
initialization data, set
\begin{equation}
\label{eq-controlled-initial-retained-entropy}
\mathsf H_{\rm init}^{[H]}
=H(S_0^{\rm ret}\mid\mathcal F_0^{\rm clean}).
\end{equation}
This term is charged once.  It is zero only when the retained auxiliary state
is conditionally deterministic; otherwise the finite-alphabet fallback is
\(\log_2|\mathsf S_0^{\rm ret}|\).  It deliberately excludes clean initial
task/latent uncertainty, which remains the prior-entropy or initial-information
coordinate.  A claimed total trajectory entropy must add that prior term;
the present ledger does not silently identify injected randomness with total
path entropy.

If \(Z_{q,t}^{\rm ch}\) is one physical-channel output,
\(\mathcal G_{q,t}^{\rm pre}\) contains every earlier public and latent
reveal and retained channel state, and \(U_{q,t}\) is its declared clean
channel input.  With the initial slot separated, define the
physical-channel noise entropy in bits by
\begin{equation}
\label{eq-controlled-physical-channel-entropy}
\mathsf H_{\rm ch}^{[H]}
=\sum_{t,q}
H(Z_{q,t}^{\rm ch}\mid\mathcal G_{q,t}^{\rm pre},U_{q,t}).
\end{equation}
For the factorization in
\cref{eq-controlled-pomdp-kernel-factorization}, this is the sum of the
joint execution/feedback entropy
\(H(A_t^\circ,E_t\mid X_t,A_t,H_t^{\rm obs},\text{channel state})\),
the reward-report entropy conditional on the clean reward and every earlier
reveal, and the observation entropy conditional on its clean latent input and
every earlier reveal.  The generic reveal-order definition prevents
correlated outputs from being omitted or counted twice.  Conditioning on a
retained state does not erase its randomness: its initial entropy has already
been charged once in \cref{eq-controlled-initial-retained-entropy}.
Let \(\mathcal G_t^{A,\rm pre}\) be the full retained sigma-field immediately
before \((A_t^\circ,E_t)\) is drawn,
\(\mathcal G_t^{\rm tr,pre}=\mathcal G_t^{A,\rm pre}\vee
\sigma(A_t^\circ,E_t)\), and
\(\mathcal G_t^{R,\rm pre}=\mathcal G_t^{\rm tr,pre}\vee
\sigma(X_{t+1})\).  Let \(\mathcal F_t^{\pi,\rm pre}\) contain the public
history and every retained policy-randomness state immediately before the
fresh action draw, but not that draw.  The stochastic-transition,
clean-reward, and policy-kernel entropies are stored separately:
\begin{align}
\mathsf H_{\rm env}^{[H]}
&=\sum_{t=0}^{H-1}
H(X_{t+1}\mid\mathcal G_t^{\rm tr,pre}),\\
\mathsf H_R^{[H]}
&=\sum_{t=0}^{H-1}
H(R_t\mid\mathcal G_t^{R,\rm pre}),\\
H_\pi^{[H]}
&=\sum_{t=0}^{H-1}H(A_t\mid\mathcal F_t^{\pi,\rm pre}).
\label{eq-controlled-environment-policy-entropy}
\end{align}
Conditioning transition entropy after the joint execution/feedback draw
prevents actuator randomness from being counted again as environment
randomness.  The clean-reward row is separate from report-channel entropy;
their chain-rule sum is used only when one combined environment coordinate is
explicitly requested.  Conditioning policy entropy on retained randomness
prevents a persistent or global seed from being charged afresh at every time;
its one-time charge is \(\mathsf H_{\rm init}^{[H]}\).
All Shannon entropies in these displays use base two.  An absent or
proved deterministic slot contributes zero.  If an exact finite-alphabet
upper bound is unavailable, the represented fallback is the sum of the
applicable log-alphabet bounds, not zero.  The
quantity \(\mathsf H_{\rm ch}^{[H]}\) measures randomization injected by the
physical channels.  It
is distinct from the information \(C_{\mathcal Q}(K)\) transmitted through
them.  For continuous channels, the record uses relative entropy to a
declared reference or mutual information; unconstrained differential entropy
is not used as a coordinate-invariant budget, and an unconstrained upper
entry is \(+\infty\).

\end{definition}

\Cref{fig-controlled-entropy-ledger-separation} keeps physical randomization
and transmitted task information in separate coordinates.

\begin{figure}[tbp]
\centering
\resizebox{.96\linewidth}{!}{%
\begin{tikzpicture}[fw diagram,node distance=9mm and 20mm]
  \node[fw terminal,text width=3.4cm] (ledger)
    {typed ledger\\no automatic addition};
  \node[fw provenance,text width=3.1cm,above left=11mm and 22mm of ledger] (pi)
    {initial retained state\\\(\mathsf H_{\rm init}\) and policy draws \(H_\pi\)};
  \node[fw use,text width=3.5cm,above right=11mm and 22mm of ledger] (ch)
    {execution/feedback, report, observation\\\(\mathsf H_{\rm ch}\)};
  \node[fw geom,text width=2.9cm,below left=13mm and 31mm of ledger] (env)
    {latent transition\\\(\mathsf H_{\rm env}\)};
  \node[fw box,text width=2.8cm,below=13mm of ledger] (rew)
    {clean reward\\\(\mathsf H_R\)};
  \node[fw source,text width=3.0cm,below right=13mm and 31mm of ledger] (pub)
    {transmitted task information\\\(C_{\rm pub}\)};
  \draw[fw arrow] (pi)--(ledger);
  \draw[fw arrow] (ch)--(ledger);
  \draw[fw arrow] (env)--(ledger);
  \draw[fw arrow] (rew)--(ledger);
  \draw[fw arrow] (pub)--(ledger);
\end{tikzpicture}%
}
\caption{The physical entropy ledger charges initial retained randomness
once and separates subsequent policy randomization,
stochastic-transition entropy, clean-reward entropy, and channel-injected
entropy.  Public channel capacity is a fifth, non-entropic coordinate that
bounds transmitted task information.  The arrows denote typed registration,
not the temporal order of a POMDP step.}
\label{fig-controlled-entropy-ledger-separation}
\end{figure}

\subsection{Adaptive information and rate--distortion}
\label{subsec-controlled-adaptive-information-rate-distortion}

\begin{theorem}[Adaptive multichannel information and rate--distortion bounds]
\label{thm-controlled-multichannel-information-bound}

Let \(\Theta\) be a latent task index and let \(\mathsf Z_0\) be the initial
public record.  When \(\mathsf Z_0\) includes an estimator trained on the
\(m\) logged samples, the complete record either keeps
\(I(\Theta;\mathsf Z_0)\) explicitly or factors
\(\mathsf Z_0=(\mathsf Z_{\mathrm{prior}},\mathsf Z_{\mathrm{train}})\) into the prior
public presentation and \(m\) declared training reveal slots with capacities
\(c^{\rm train}_{q,s}\).  At each deployment time and reveal slot, suppose the complete record
contains a conditional Markov chain
\begin{equation}
\label{eq-controlled-conditional-channel-markov-chain}
\Theta\longrightarrow U_{q,t}\longrightarrow Z_{q,t}
\quad\text{conditional on the public past},
\end{equation}
and a capacity certificate \(c_{q,t}\) from
\cref{def-controlled-channel-information-ledger}.  The controller's intended
action is generated from the public past and contributes zero conditional
mutual information.  Unrevealed target-independent actuator randomness also
contributes zero.  Execution feedback that depends on a latent state is a
public observation channel and receives its own capacity slot.
Every realized history-conditional input law must lie in the declared class
\(\mathcal Q\); the SDPI branch also uses the conditional and mixture closure
specified in \cref{def-controlled-channel-information-ledger}.

Let \(\mathsf Z_H\) be the full public transcript, let
\(\mathsf Z_{1:H}^{\rm dep}\) denote its ordered deployment-reveal suffix so
that \(\mathsf Z_H=(\mathsf Z_0,\mathsf Z_{1:H}^{\rm dep})\), and abbreviate the
represented deployment information budget by
\[
C_{\rm pub}^{[H]}:=\sum_{t,q}c_{q,t}.
\]
Then
\begin{align}
I(\Theta;\mathsf Z_H)
&\le I(\Theta;\mathsf Z_0)
+\sum_{t,q}c_{q,t},
\label{eq-controlled-adaptive-multichannel-information}\\
H(\Theta\mid\mathsf Z_H)
&\ge H(\Theta\mid\mathsf Z_0)-\sum_{t,q}c_{q,t}.
\label{eq-controlled-adaptive-posterior-entropy}
\end{align}
If \(\Theta\) is uniform on a family of size \(N_n\ge2\), let
\(I_0^{(m)}\) be a represented upper bound on
\(I(\Theta;\mathsf Z_0)\).  The exact value is admissible; under the training
factorization one may take
\begin{equation}
\label{eq-controlled-training-information-budget}
I(\Theta;\mathsf Z_0)\le I_0^{(m)}
:=I_{\mathrm{public}}+\sum_{s=1}^{m}\sum_q c^{\mathrm{train}}_{q,s},
\qquad
I_{\mathrm{public}}\ge I(\Theta;\mathsf Z_{\mathrm{prior}}).
\end{equation}
Every estimator satisfies
\begin{equation}
\label{eq-controlled-multichannel-fano-success}
\Pr\{\widehat\Theta=\Theta\}
\le
\min\left\{1,
\frac{I_0^{(m)}+\sum_{t,q}c_{q,t}+1}{\log_2N_n}
\right\}.
\end{equation}
Success at least \(\alpha\) requires
\begin{equation}
\label{eq-controlled-multichannel-information-necessary}
\sum_{t,q}c_{q,t}
\ge\alpha\log_2N_n-I_0^{(m)}-1.
\end{equation}

More generally, let \(d(\theta,\widehat\theta)\) be a represented distortion
and define the conditional rate--distortion function
\begin{equation}
\label{eq-controlled-conditional-rate-distortion}
R_{\Theta\mid\mathsf Z_0}(D)
=
\inf_{\substack{P_{\widehat\Theta\mid\Theta,\mathsf Z_0}\\
\mathbb E d(\Theta,\widehat\Theta)\le D}}
I(\Theta;\widehat\Theta\mid\mathsf Z_0).
\end{equation}
Any transcript estimator with expected distortion at most \(D\) must satisfy
\begin{equation}
\label{eq-controlled-rate-distortion-converse}
R_{\Theta\mid\mathsf Z_0}(D)
\le\sum_{t,q}c_{q,t}.
\end{equation}
Equivalently, with the usual extended-real convention, define
\begin{equation}
\label{eq-controlled-rate-distortion-inverse-floor}
D_{\min}(C)
=\inf\{D\ge0:R_{\Theta\mid\mathsf Z_0}(D)\le C\}.
\end{equation}
Every transcript estimator obeys
\(\mathbb E d(\Theta,\widehat\Theta)\ge
D_{\min}(C_{\rm pub}^{[H]})\).  Turning this distortion floor into an event-
probability bound requires a represented distortion-to-event inequality.
The rate--distortion branch is unavailable, rather than zero, when the source
law, distortion, or conditional channel factorization is not represented.
It is a converse and supplies no achievability claim.
The computational ceiling \(b(n)\) alone is not an information bound on
\(I_0^{(m)}\); the displayed \(m\)-dependence requires the represented
training-channel factorization.
If a target-dependent physical channel fails the conditional Markov
factorization in
\cref{eq-controlled-conditional-channel-markov-chain}, this capacity branch
is unavailable; its dependence is not discarded as noise.

If the conditional channel also has an SDPI certificate, then each summand
may be sharpened recordwise to
\begin{equation}
\label{eq-controlled-channel-sdpi-information}
I(\Theta;Z_{q,t}\mid\text{public past})
\le
\min\left\{c_{q,t},
\vartheta_{q,t}
I(\Theta;U_{q,t}\mid\text{public past})
\right\}.
\end{equation}

\end{theorem}

\begin{proof}

Write \(\mathsf Z_{<q,t}\) for the public record immediately before slot
\((q,t)\).  The chain rule gives
\[
I(\Theta;\mathsf Z_H)
=I(\Theta;\mathsf Z_0)
 +\sum_{t,q}I(\Theta;Z_{q,t}\mid\mathsf Z_{<q,t}).
\]
For each admitted past \(h\), conditional data processing and the capacity
definition give
\[
I(\Theta;Z_{q,t}\mid h)
\le I(U_{q,t};Z_{q,t}\mid h)
\le c_{q,t}.
\]
Integrating over \(h\) and summing proves
\cref{eq-controlled-adaptive-multichannel-information}.  An intended action
generated from the conditioned public past and fresh target-independent
policy randomness has zero conditional mutual information and contributes no
slot.  Applying the identical argument to the declared training reveals
proves \cref{eq-controlled-training-information-budget}.  Subtracting the
information bound in the identity
\(H(\Theta\mid Z)=H(\Theta)-I(\Theta;Z)\) proves
\cref{eq-controlled-adaptive-posterior-entropy}.

Let \(P_e=\Pr\{\widehat\Theta\ne\Theta\}\).  For uniform \(\Theta\), Fano's
entropy calculation is
\[
H(\Theta\mid\widehat\Theta)
\le H_2(P_e)+P_e\log_2(N_n-1)
\le1+P_e\log_2N_n.
\]
Because \(\widehat\Theta\) is a transcript postprocessing,
\(I(\Theta;\widehat\Theta)\le I(\Theta;\mathsf Z_H)\).  Using
\(H(\Theta)=\log_2N_n\) and rearranging gives
\(1-P_e\le(I(\Theta;\mathsf Z_H)+1)/\log_2N_n\), which proves
\cref{eq-controlled-multichannel-fano-success}.  Requiring
\(1-P_e\ge\alpha\) and rearranging proves
\cref{eq-controlled-multichannel-information-necessary}.

Every transcript estimator induces a conditional test channel
\(P_{\widehat\Theta\mid\Theta,\mathsf Z_0}\).  Conditional data processing
and the deployment chain rule give
\[
I(\Theta;\widehat\Theta\mid\mathsf Z_0)
\le I(\Theta;\mathsf Z_{1:H}^{\rm dep}\mid\mathsf Z_0)
\le C_{\rm pub}^{[H]}.
\]
If its expected distortion is at most \(D\), that test channel is feasible
in \cref{eq-controlled-conditional-rate-distortion}; hence
\(R_{\Theta\mid\mathsf Z_0}(D)\le C_{\rm pub}^{[H]}\).  Its achieved
distortion belongs to the set in
\cref{eq-controlled-rate-distortion-inverse-floor} and is therefore at least
the infimum of that set, proving the inverse floor.

For the SDPI branch, fix a public past \(h\).  The admitted-class closure
allows the defining KL contraction to be applied to every conditional input
law and its mixture:
\begin{align*}
&\sum_\theta p(\theta\mid h)
 \operatorname{KL}\!\left(
 P_{U\mid\theta,h}K\Vert P_{U\mid h}K\right)\\
&\qquad\le\vartheta_{q,t}
\sum_\theta p(\theta\mid h)
 \operatorname{KL}\!\left(
 P_{U\mid\theta,h}\Vert P_{U\mid h}\right).
\end{align*}
The two sides are respectively
\(I(\Theta;Z_{q,t}\mid h)\) and
\(\vartheta_{q,t}I(\Theta;U_{q,t}\mid h)\).  Averaging over \(h\) and
taking the minimum with the capacity bound proves
\cref{eq-controlled-channel-sdpi-information}.

\end{proof}

\subsection{Pre-Hit Sequential Information Contraction}
\label{subsec-pre-hit-sequential-information-contraction}

\begin{theorem}[Pre-hit sequential SDPI and information contraction]
\label{thm-pre-hit-sequential-sdpi-information-contraction}

Let \(\Theta\) be a finite latent task index, let \(\mathsf Z_0\) be the
initial public record, and fix a temporally admissible prefix ending before a
represented selector is sampled.  Enumerate every target-dependent public
reveal in that prefix in its actual order as
\(Y_1,\ldots,Y_L\).  The list includes observations, reward reports,
execution feedback, verifier flags, and answers returned by declared oracle
interfaces.  Put \(\mathscr F_0=\sigma(\mathsf Z_0)\).  Immediately before
\(Y_j\), let \(Q_j\) collect the represented query, action, disclosed policy
randomness, and deterministic state updates performed since the preceding
reveal, and set
\[
\mathscr P_j=\mathscr F_{j-1}\vee\sigma(Q_j),
\qquad
\mathscr F_j=\mathscr P_j\vee\sigma(Y_j).
\]
Write \(H_j^{\rm pre}\) for the represented pre-reveal history that
generates \(\mathscr P_j\).
Assume
\begin{equation}
\label{eq-pre-hit-query-conditional-independence}
\Theta\perp Q_j\mid\mathscr F_{j-1}.
\end{equation}
Every target-dependent coordinate in \(Q_j\) must instead occur as an
earlier reveal slot.

For each \(j\), suppose the complete record contains a clean channel input
\(U_j\), a history-conditional channel \(K_j\), and the conditional Markov
chain
\begin{equation}
\label{eq-pre-hit-sequential-channel-markov}
\Theta\longrightarrow U_j\longrightarrow Y_j
\quad\text{conditional on }\mathscr P_j.
\end{equation}
Let \(c_j(h)\) and \(\vartheta_j(h)\) be, respectively, capacity and KL-SDPI
certificates from
\cref{def-controlled-channel-information-ledger}, valid at every admitted
pre-reveal history \(h\) for the required conditional input laws and their
mixtures.  Define
\begin{align}
a_j(h)
&:=I(\Theta;U_j\mid\mathscr P_j=h),
\label{eq-pre-hit-sequential-upstream-information}\\
d_j
&:=\mathbb E\!\left[
\min\{c_j(H_j^{\rm pre}),
       \vartheta_j(H_j^{\rm pre})a_j(H_j^{\rm pre})\}
\right].
\label{eq-pre-hit-sequential-contracted-increment}
\end{align}
All information quantities in this theorem are measured in bits.  Then
\begin{align}
I(\Theta;\mathscr F_L\mid\mathsf Z_0)
&=
\sum_{j=1}^{L}I(\Theta;Y_j\mid\mathscr P_j)
\le \sum_{j=1}^{L}d_j,
\label{eq-pre-hit-sequential-information-chain-rule}\\
I(\Theta;\mathscr F_L)
&\le I(\Theta;\mathsf Z_0)+\sum_{j=1}^{L}d_j.
\label{eq-pre-hit-sequential-total-information}
\end{align}
In particular, adaptively selecting \(Q_j\), \(U_j\), or \(K_j\) from the
represented past does not alter the bound; the certificate used at slot
\(j\) is the one valid at the realized pre-reveal history.

Suppose, more specifically, that scalar bounds
\(c_j\in[0,+\infty]\) and \(\vartheta_j\in[0,1]\) dominate the two
history-dependent certificates at every admitted history.  Let
\(h_j=H(\Theta\mid\mathscr F_j)\).  Then
\begin{equation}
\label{eq-pre-hit-sdpi-entropy-one-step}
h_j
\ge
h_{j-1}-\min\{c_j,\vartheta_j h_{j-1}\}
=
\max\{h_{j-1}-c_j,(1-\vartheta_j)h_{j-1}\}.
\end{equation}
Consequently, if \(\underline h_0\le H(\Theta\mid\mathsf Z_0)\) and
\begin{equation}
\label{eq-pre-hit-sdpi-entropy-recursion}
\underline h_j
:=
\underline h_{j-1}
-\min\{c_j,\vartheta_j\underline h_{j-1}\},
\end{equation}
then
\begin{equation}
\label{eq-pre-hit-sdpi-entropy-cumulative}
H(\Theta\mid\mathscr F_L)
\ge \underline h_L
\ge
\max\left\{
\underline h_0-\sum_{j=1}^{L}c_j,
\ \underline h_0\prod_{j=1}^{L}(1-\vartheta_j)
\right\}.
\end{equation}
These formulas permit \(c_j=+\infty\) and \(\vartheta_j=1\); no
small-increment hypothesis is used.

The complete prefix record retains the construction and evaluation cost of
every \(Q_j\) and \(U_j\), the selected channel and its state, the proof or
certificate producing \(c_j\) and \(\vartheta_j\), every random bit, every
memory update, and every returned answer.  All of these are ancestors of the
later selector and are charged before it according to
\cref{def-pre-hit-temporally-admissible-source}.  A
presentation-admissible oracle answer occupies a channel slot with its
declared provenance and cost.  An oracle answer outside that presentation
changes the presented family by
\cref{thm-oracle-admissibility-presentation-separation}.  Later verifier
flags, terminal records, and future oracle answers are absent from
\(\mathscr F_L\).

\end{theorem}

\begin{proof}

Apply the conditional chain rule in the displayed reveal order:
\begin{align*}
I(\Theta;\mathscr F_L\mid\mathsf Z_0)
&=\sum_{j=1}^{L}
  I(\Theta;Q_j,Y_j\mid\mathscr F_{j-1})\\
&=\sum_{j=1}^{L}
  I(\Theta;Y_j\mid\mathscr P_j).
\end{align*}
The second equality is
\cref{eq-pre-hit-query-conditional-independence}.  Fix a pre-reveal history
\(h\).  The admitted-class closure in
\cref{def-controlled-channel-information-ledger} and the conditional Markov
chain in \cref{eq-pre-hit-sequential-channel-markov} give
\begin{align*}
I(\Theta;Y_j\mid\mathscr P_j=h)
&\le c_j(h),\\
I(\Theta;Y_j\mid\mathscr P_j=h)
&\le \vartheta_j(h)
       I(\Theta;U_j\mid\mathscr P_j=h).
\end{align*}
Taking their minimum, averaging over the pre-reveal history, and summing
proves \cref{eq-pre-hit-sequential-information-chain-rule}.  Adding
\(I(\Theta;\mathsf Z_0)\) proves
\cref{eq-pre-hit-sequential-total-information}.

Under uniform scalar certificates, the elementary mutual-information bound
also gives
\[
I(\Theta;U_j\mid\mathscr P_j)
\le H(\Theta\mid\mathscr P_j)
=H(\Theta\mid\mathscr F_{j-1})=h_{j-1};
\]
the middle equality again uses
\cref{eq-pre-hit-query-conditional-independence}.  Since
\[
h_j=h_{j-1}-I(\Theta;Y_j\mid\mathscr P_j),
\]
the capacity and SDPI bounds prove
\cref{eq-pre-hit-sdpi-entropy-one-step}.  The map
\(x\mapsto x-\min\{c_j,\vartheta_jx\}\) is nondecreasing on
\([0,+\infty)\).  Induction therefore proves
\(h_j\ge\underline h_j\).  The same recursion separately implies
\(\underline h_j\ge\underline h_{j-1}-c_j\) and
\(\underline h_j\ge(1-\vartheta_j)\underline h_{j-1}\).  Iterating the two
inequalities proves \cref{eq-pre-hit-sdpi-entropy-cumulative}.

The final assertions follow from the definition of the complete pre-hit
record and from the oracle-presentation separation theorem cited in the
statement.

\end{proof}

\subsection{Posterior Concentration, Target Contact, and Selector Advantage}
\label{subsec-pre-hit-target-contact-selector}

\begin{corollary}[Information contraction to posterior concentration,
target contact, and selector advantage]
\label{cor-pre-hit-information-target-contact-selector}

Retain the record of
\cref{thm-pre-hit-sequential-sdpi-information-contraction}, and identify
\(\mathscr F_L\) with a complete pre-hit record \(\mathcal H_t\).  Let
\(K_t(\cdot\mid\mathcal H_t)\) be a temporally admissible next-test sampler
as in \cref{def-represented-pre-hit-selector}; its fresh random bits are
conditionally independent of \(\Theta\) given \(\mathcal H_t\).  Let
\(V_\theta(c)\in\{0,1\}\) be the represented verification relation for
target index \(\theta\), and let \(C_{t+1}\) be the candidate sampled from
\(K_t\).  Put
\begin{equation}
\label{eq-pre-hit-information-next-contact}
p_t
:=
\Pr\{\tau>t,\ V_\Theta(C_{t+1})=1\}.
\end{equation}

The candidate carrier is common to all target indices, or is padded by a
represented common carrier, so that \(V_\theta(c)\) is measurable for every
cross-pair \((\theta,c)\) appearing under the decoupled law below.

Let \(\mathbb P_t\) be the joint law of
\((\Theta,\mathcal H_t,C_{t+1})\).  Define its target-decoupled reference by
\begin{equation}
\label{eq-pre-hit-information-decoupled-law}
\mathbb Q_t(d\theta,dh,dc)
:=
P_\Theta(d\theta)P_{\mathcal H_t}(dh)K_t(dc\mid h),
\end{equation}
and let
\begin{equation}
\label{eq-pre-hit-information-decoupled-baseline}
\beta_t
:=
\mathbb Q_t\{\tau>t,\ V_\Theta(C_{t+1})=1\}.
\end{equation}
If
\begin{equation}
\label{eq-pre-hit-information-prefix-upper}
J_t^\uparrow
:=
I(\Theta;\mathsf Z_0)+\sum_{j=1}^{L}d_j,
\end{equation}
then, with the lower-semicontinuous endpoint convention for
\(D_{\mathrm{Ber}}\) from
\cref{eq-bernoulli-divergence-definition},
\begin{align}
D_{\mathrm{Ber}}(p_t\Vert\beta_t)
&\le (\log 2)J_t^\uparrow,
\label{eq-pre-hit-information-binary-contraction}\\
(p_t-\beta_t)_+
&\le
\sqrt{\frac{\log 2}{2}J_t^\uparrow}.
\label{eq-pre-hit-information-pinsker-contact}
\end{align}
The same information budget controls the posterior:
\begin{equation}
\label{eq-pre-hit-information-posterior-tv}
\mathbb E\left[
\left\|P_{\Theta\mid\mathcal H_t}-P_\Theta\right\|_{\rm TV}
\right]
\le
\sqrt{\frac{\log 2}{2}J_t^\uparrow}.
\end{equation}

Suppose \(\Theta\) is uniform on \(N\ge2\) indices and the target sectors
have represented overlap at most \(\Lambda_t\), meaning
\begin{equation}
\label{eq-pre-hit-information-target-overlap}
\sup_{h,c}
\sum_{\theta=1}^{N}V_\theta(c)
\le\Lambda_t.
\end{equation}
Then \(\beta_t\le\Lambda_t/N\), and hence
\begin{equation}
\label{eq-pre-hit-information-overlap-contact}
p_t
\le
\min\left\{1,
\frac{\Lambda_t}{N}
+\sqrt{\frac{\log 2}{2}J_t^\uparrow}
\right\}.
\end{equation}
For disjoint target sectors, \(\Lambda_t=1\).  For any declared neutral
baseline \(\nu_t\),
\begin{equation}
\label{eq-pre-hit-information-selector-bound}
\bigl(\operatorname{Adv}_t(K_t;\nu_t)\bigr)_+
\le p_t,
\end{equation}
so \cref{eq-pre-hit-information-overlap-contact} is also a prospective
selector bound.

There is an optional chi-square refinement requiring a separate represented
certificate.  If the complete record contains
\begin{equation}
\label{eq-pre-hit-information-chi-square-input}
\chi^2(\mathbb P_t\Vert\mathbb Q_t)
:=
\int\left(\frac{d\mathbb P_t}{d\mathbb Q_t}-1\right)^2d\mathbb Q_t
\le\Xi_t,
\end{equation}
then
\begin{equation}
\label{eq-pre-hit-information-chi-square-contact}
p_t
\le
\min\left\{1,
\beta_t+\sqrt{\beta_t(1-\beta_t)\Xi_t}
\right\}.
\end{equation}
The derivation and provenance of \(\Xi_t\) are part of the same charged
pre-hit record; no KL-to-chi-square conversion is inferred.

\end{corollary}

\begin{proof}

The sampler is a postprocessing of \(\mathcal H_t\) using conditionally
target-independent fresh randomness.  Therefore
\begin{equation}
\label{eq-pre-hit-information-joint-kl-identity}
\operatorname{KL}(\mathbb P_t\Vert\mathbb Q_t)
=(\log 2)I(\Theta;\mathcal H_t)
\le(\log 2)J_t^\uparrow.
\end{equation}
Applying the measurable map
\((\theta,h,c)\mapsto
\mathbf1_{\{\tau(h)>t\}}V_\theta(c)\)
and using data processing for relative entropy proves
\cref{eq-pre-hit-information-binary-contraction}.  Pinsker's inequality in
nats proves \cref{eq-pre-hit-information-pinsker-contact}.  Applying
conditional Pinsker directly to
\(P_{\Theta\mid\mathcal H_t}\) and \(P_\Theta\), followed by Jensen's
inequality, proves \cref{eq-pre-hit-information-posterior-tv}.

Under the decoupled law, \(\Theta\) is independent of the history and the
sampled candidate.  Hence
\begin{align*}
\beta_t
&=
\mathbb E_{\mathbb Q_t}\left[
\mathbf1_{\{\tau>t\}}
\frac1N\sum_{\theta=1}^{N}V_\theta(C_{t+1})
\right]\\
&\le\frac{\Lambda_t}{N}.
\end{align*}
Combining this estimate with
\cref{eq-pre-hit-information-pinsker-contact} proves
\cref{eq-pre-hit-information-overlap-contact}.  The neutral selector has
nonnegative next-hit mass.  The definition
\cref{eq-pre-hit-selector-advantage} therefore gives
\((\operatorname{Adv}_t)_+\le p_t\), proving
\cref{eq-pre-hit-information-selector-bound}.

For the chi-square branch, write
\(L_t=d\mathbb P_t/d\mathbb Q_t\) and let \(E_t\) be the next-contact event.
Since \(\mathbb E_{\mathbb Q_t}(L_t-1)=0\), Cauchy--Schwarz gives
\begin{align*}
p_t-\beta_t
&=\mathbb E_{\mathbb Q_t}
[(L_t-1)(\mathbf1_{E_t}-\beta_t)]\\
&\le
\sqrt{\chi^2(\mathbb P_t\Vert\mathbb Q_t)}
\sqrt{\beta_t(1-\beta_t)},
\end{align*}
which proves \cref{eq-pre-hit-information-chi-square-contact}.

\end{proof}

\begin{corollary}[Cumulative first contact and full-transcript identification]
\label{cor-pre-hit-cumulative-contact-full-transcript}

Let \(\Theta\) be uniform on \(N\ge2\) indices.  For every
\(t=0,\ldots,H-1\), retain a complete prefix record satisfying
\cref{cor-pre-hit-information-target-contact-selector}, with the quantities
\(J_t^\uparrow\) and \(\Lambda_t\) defined in
\cref{eq-pre-hit-information-prefix-upper,eq-pre-hit-information-target-overlap}.
Then
\begin{equation}
\label{eq-pre-hit-information-cumulative-hit}
\Pr\{\tau\le H\}
\le
\min\left\{1,
\sum_{t=0}^{H-1}
\left(
\frac{\Lambda_t}{N}
+\sqrt{\frac{\log 2}{2}J_t^\uparrow}
\right)
\right\}.
\end{equation}

If, in addition, the complete terminal record contains a represented decoder
that identifies \(\Theta\) on every verified hit, let
\(\mathsf Z_H^{\rm full}\) denote that terminal transcript.  Suppose every
target-dependent coordinate after \(\mathsf Z_0\), including terminal
verification, decoder-selection, feedback, and oracle coordinates, is exposed
as an ordered channel slot satisfying
\cref{thm-controlled-multichannel-information-bound}.  If
\(\mathcal R_H^{\rm full}\) is the resulting slot index set, with represented
capacity certificates \(c_{q,r}\), put
\begin{equation}
\label{eq-pre-hit-information-full-transcript-upper}
J_H^{\rm full}
:=
I(\Theta;\mathsf Z_0)
+\sum_{(q,r)\in\mathcal R_H^{\rm full}}c_{q,r}.
\end{equation}
No terminal target-dependent coordinate is omitted from this sum;
target-independent postprocessing and fresh target-independent randomness
contribute zero conditional mutual information.  Fano's inequality then gives
the independent full-transcript bound
\begin{equation}
\label{eq-pre-hit-information-hit-fano}
\Pr\{\tau\le H\}
\le
\min\left\{1,
\frac{J_H^{\rm full}+1}{\log_2N}
\right\}.
\end{equation}
The recordwise minimum of
\cref{eq-pre-hit-information-cumulative-hit,eq-pre-hit-information-hit-fano}
may be taken before saturation.  Saturation then takes the supremum over
complete records at the charged enlarged ceiling; each record retains its
own channel, oracle, overlap, information, and decoder certificates.

\end{corollary}

\begin{proof}

The first-hit events are disjoint, so summing their probabilities and using
\cref{eq-pre-hit-information-overlap-contact} proves
\cref{eq-pre-hit-information-cumulative-hit}.  Under the identification
gate, map a verified hit to its target index and assign an arbitrary index on
failure.  The decoder succeeds whenever a hit occurs.  Its mutual
information is at most that of \(\mathsf Z_H^{\rm full}\).  Applying
\cref{thm-controlled-multichannel-information-bound} to every slot in
\(\mathcal R_H^{\rm full}\) gives
\[
I(\Theta;\mathsf Z_H^{\rm full})
\le J_H^{\rm full}.
\]
The Fano calculation in
\cref{thm-controlled-multichannel-information-bound} proves
\cref{eq-pre-hit-information-hit-fano}.  Every conversion was performed on
one complete record, so taking its minimum before the affordable supremum
has the order stated in the corollary.

\end{proof}

\Cref{fig-adaptive-multichannel-information} shows where sample size and
deployment horizon enter the information converse.

\begin{figure}[tbp]
\centering
\resizebox{\linewidth}{!}{%
\begin{tikzpicture}[fw diagram,node distance=7mm and 10mm]
  \node[fw source,text width=3.0cm] (train)
    {\(m\) training reveal slots\\\(I_0^{(m)}\)};
  \node[fw provenance,text width=3.4cm,right=of train] (adapt)
    {adaptive deployment slots \((t,q)\)\\
     \(\Theta\to U_{q,t}\to Z_{q,t}\mid\text{past}\)};
  \node[fw use,text width=3.0cm,right=of adapt] (sum)
    {information ledger\\\(I\le I_0^{(m)}+\sum_{t,q}c_{q,t}\)};
  \draw[fw arrow] (train)--(adapt);
  \draw[fw arrow] (adapt)--(sum);
  \node[fw terminal,text width=2.8cm,above right=9mm and 11mm of sum] (fano)
    {Fano\\identification success};
  \node[fw terminal,text width=2.8cm,below right=9mm and 11mm of sum] (rd)
    {rate--distortion\\distortion floor};
  \node[fw geom,text width=2.8cm,right=of sum] (post)
    {posterior entropy\\lower bound};
  \draw[fw arrow] (sum)--(post);
  \draw[fw arrow] (sum)--(fano);
  \draw[fw arrow] (sum)--(rd);
\end{tikzpicture}%
}
\caption{Adaptive information is charged reveal by reveal.  Training and
deployment budgets give
\(I(\Theta;\mathsf Z_H)\le I_0^{(m)}+\sum_{t,q}c_{q,t}\); posterior entropy,
Fano identification, and rate--distortion are separate consequences of that
same budget.}
\label{fig-adaptive-multichannel-information}
\end{figure}

\subsection{Explicit Channel-Information Budgets}
\label{subsec-explicit-channel-information-budgets}

\begin{corollary}[Explicit channel-information budgets]
\label{cor-controlled-explicit-channel-capacities}

The following values are in bits per represented channel use.

\begin{enumerate}[label=\textup{(\roman*)}]
\item If \(|\mathsf V|=d\ge2\) and every admitted input law satisfies
      \(0\le h_-\le\log_2d\) and
      \(H(V\mid U)\ge h_-\), then
      \begin{equation}
      \label{eq-controlled-finite-output-entropy-capacity}
      C_{\mathcal Q}(K)\le\log_2d-h_-.
      \end{equation}

\item A \(q\)-ary symmetric channel that, conditional on an error, chooses
      each of the other \(q-1\) symbols uniformly, with crossover probability
      \(0\le\epsilon\le(q-1)/q\), has
      \begin{equation}
      \label{eq-controlled-qary-symmetric-capacity}
      C(K)=\log_2q-H_2(\epsilon)-\epsilon\log_2(q-1).
      \end{equation}
      An additive channel on a finite group \(G\) with independent noise
      \(N\) has \(C(K)=\log_2|G|-H(N)\).  A \(q\)-ary erasure channel of
      erasure rate \(\epsilon\) has
      \(C(K)=(1-\epsilon)\log_2q\).

\item For an additive scalar Gaussian reporting channel with average input
      power at most \(P\) and independent noise variance \(\sigma^2>0\),
      \begin{equation}
      \label{eq-controlled-gaussian-reward-capacity}
      C_{P}(K)\le\frac12\log_2(1+P/\sigma^2).
      \end{equation}
      Without the power constraint the capacity entry is \(+\infty\).
\end{enumerate}

Marginal memoryless formulas are not applied to correlated channel noise.
Such a record uses conditional history-state capacities or a represented
joint block-capacity certificate.  The finite-alphabet estimates of
\cref{thm-finite-alphabet-channel-estimation-entropy} may be inserted after
their continuity error is added.

\end{corollary}

\begin{proof}

For the first statement,
\(I(U;V)=H(V)-H(V\mid U)\le\log_2d-h_-\).

For the \(q\)-ary symmetric channel, every row has entropy
\[
H(V\mid U=u)=H_2(\epsilon)+\epsilon\log_2(q-1).
\]
Consequently every input law has
\(I(U;V)\le\log_2q-H_2(\epsilon)
-\epsilon\log_2(q-1)\).  Uniform input makes the output uniform and attains
that upper bound.  For the finite-group channel \(V=U+N\), independence gives
\(H(V\mid U)=H(N)\), while \(H(V)\le\log_2|G|\); uniform \(U\) makes \(V\)
uniform and proves equality.  For the erasure channel, conditioning on the
erasure indicator gives \(I(U;V)=(1-\epsilon)H(U)\), maximized by uniform
input.

For the scalar Gaussian channel \(Y=X+N\), independence and the power bound
give \(\operatorname{Var}(Y)\le P+\sigma^2\).  The maximum-entropy inequality
therefore yields, in bits,
\[
I(X;Y)=h(Y)-h(N)
\le\frac12\log_2\frac{2\pi e(P+\sigma^2)}{2\pi e\sigma^2}
=\frac12\log_2(1+P/\sigma^2).
\]
For the ordinary unrestricted power-constrained input class, a centered
Gaussian input attains equality; for a restricted represented input class,
the displayed upper bound remains valid.

\end{proof}

\subsection{Belief entropy and target-gate uncertainty}
\label{subsec-controlled-belief-entropy-target-gates}

\begin{theorem}[Belief-entropy balance and target-gate uncertainty]
\label{thm-controlled-belief-entropy-balance}

Assume finite latent alphabets, and let \(Z_{t+1}^{\rm pub}\) collect the
public observation, reward report, and execution feedback revealed between
times \(t\) and \(t+1\), with no additional public reveal omitted.  Assume
exactly
\[
\mathcal F_{t+1}^{\rm obs}
=\mathcal F_t^{\rm obs}\vee\sigma(A_t,Z_{t+1}^{\rm pub}).
\]
Assume the policy draw uses only the public pre-action record and
target-independent policy randomness, so that
\(A_t\perp X_t\mid\mathcal F_t^{\rm obs}\).  This includes the deterministic
case in which the relevant policy seed is already in
\(\mathcal F_t^{\rm obs}\).  Then
\begin{equation}
\label{eq-controlled-belief-entropy-exact-balance}
H(X_{t+1}\mid\mathcal F_{t+1}^{\rm obs})
=H(X_{t+1}\mid\mathcal F_t^{\rm obs},A_t)
-I(X_{t+1};Z_{t+1}^{\rm pub}
  \mid\mathcal F_t^{\rm obs},A_t).
\end{equation}
The pre-observation uncertainty obeys
\begin{align}
H(X_{t+1}\mid\mathcal F_t^{\rm obs},A_t)
&\le H(X_t\mid\mathcal F_t^{\rm obs})
+H(A_t^\circ,E_t\mid X_t,A_t,\mathcal F_t^{\rm obs})
\notag\\
&\quad
+H(X_{t+1}\mid X_t,A_t^\circ,E_t,
  \mathcal F_t^{\rm obs},A_t).
\label{eq-controlled-belief-entropy-injection}
\end{align}
Write the last two terms as \(h_{A,t}\) and \(h_{P,t}\), respectively.  Thus
a represented information floor \(i_{t+1}^-\) for the mutual-information
term gives the quantitative bounds
\begin{align}
H(X_{t+1}\mid\mathcal F_{t+1}^{\rm obs})
&\le H(X_t\mid\mathcal F_t^{\rm obs})
 +h_{A,t}+h_{P,t}-i_{t+1}^-,
\label{eq-controlled-belief-entropy-upper}\\
H(X_H\mid\mathcal F_H^{\rm obs})
&\le\min\left\{\log_2|X|,
H(X_0\mid\mathcal F_0^{\rm obs})
+\sum_{t=0}^{H-1}(h_{A,t}+h_{P,t}-i_{t+1}^-)
\right\}.
\label{eq-controlled-belief-entropy-cumulative}
\end{align}
The scalar right side may be clipped below at zero.  In contrast, a
channel-capacity upper bound \(c_{t+1}\) gives only
\begin{equation}
\label{eq-controlled-belief-entropy-capacity-floor}
H(X_{t+1}\mid\mathcal F_{t+1}^{\rm obs})
\ge
H(X_{t+1}\mid\mathcal F_t^{\rm obs},A_t)-c_{t+1}.
\end{equation}
An entropy upper bound cannot be obtained from a capacity upper bound alone.
Here \(c_{t+1}\) must directly certify the displayed conditional mutual
information, for example by summing the admitted conditional Markov-channel
slots.  Physical output entropy alone is not such a certificate.

For either augmented winner indicator
\(J_t\in\{\mathbf1_{\{W_t=T\}},\mathbf1_{\{W_t=F\}}\}\), put
\(q_t=\Pr\{J_t=1\mid\mathcal F_t^{\rm obs}\}\) and let
\(p_t^{\rm MAP}=\mathbb E\min\{q_t,1-q_t\}\) be the Bayes gate error.  With
binary entropy measured in bits,
\begin{equation}
\label{eq-controlled-binary-gate-entropy-error}
H_2^{-1}\!\left(
H(J_t\mid\mathcal F_t^{\rm obs})
\right)
\le p_t^{\rm MAP}
\le\frac12H(J_t\mid\mathcal F_t^{\rm obs}),
\end{equation}
where \(H_2^{-1}\) is the inverse on \([0,1/2]\).  A represented learned gate
adds its approximation and statistical errors to this analytic Bayes error
through the deployment comparison of
\cref{def-controlled-noisy-gate-profiles}.

\end{theorem}

\begin{proof}

By the stated filtration identity,
\begin{align*}
I(X_{t+1};Z_{t+1}^{\rm pub}\mid\mathcal F_t^{\rm obs},A_t)
&=H(X_{t+1}\mid\mathcal F_t^{\rm obs},A_t)\\
&\quad-H(X_{t+1}\mid\mathcal F_{t+1}^{\rm obs}),
\end{align*}
which rearranges to
\cref{eq-controlled-belief-entropy-exact-balance}.  For the injection bound,
monotonicity under adjoining variables and the chain rule give
\begin{align*}
H(X_{t+1}\mid\mathcal F_t^{\rm obs},A_t)
&\le H(X_t,A_t^\circ,E_t,X_{t+1}
        \mid\mathcal F_t^{\rm obs},A_t)\\
&=H(X_t\mid\mathcal F_t^{\rm obs})
 +H(A_t^\circ,E_t\mid X_t,A_t,\mathcal F_t^{\rm obs})\\
&\quad+H(X_{t+1}\mid X_t,A_t^\circ,E_t,
       \mathcal F_t^{\rm obs},A_t).
\end{align*}
In the chain-rule expansion, the policy conditional-independence assumption
gives
\(H(X_t\mid\mathcal F_t^{\rm obs},A_t)
=H(X_t\mid\mathcal F_t^{\rm obs})\); the last term is the entropy of the
declared transition row by the transition Markov property.  This proves
\cref{eq-controlled-belief-entropy-injection}.
Substitution of \(i_{t+1}^-\) and induction prove
\cref{eq-controlled-belief-entropy-upper,eq-controlled-belief-entropy-cumulative}.
Substitution of the upper information certificate instead proves
\cref{eq-controlled-belief-entropy-capacity-floor}.

Put \(e(q)=\min\{q,1-q\}\).  Pointwise symmetry and the elementary chord
bound give
\(H_2(q)=H_2(e(q))\ge2e(q)\).  Taking expectations yields
\(p_t^{\rm MAP}\le\tfrac12H(J_t\mid\mathcal F_t^{\rm obs})\).
Concavity gives
\[
H(J_t\mid\mathcal F_t^{\rm obs})
=\mathbb EH_2(e(q_t))
\le H_2(\mathbb Ee(q_t))=H_2(p_t^{\rm MAP}).
\]
Since \(H_2\) is increasing on \([0,1/2]\), inversion proves the lower bound
in \cref{eq-controlled-binary-gate-entropy-error}.

\end{proof}

\Cref{fig-belief-entropy-balance} emphasizes the opposite directions of an
information floor and a capacity ceiling.

\begin{figure}[tbp]
\centering
\resizebox{.95\linewidth}{!}{%
\begin{tikzpicture}[fw diagram,node distance=7mm and 10mm]
  \node[fw source,text width=2.8cm] (old) {previous posterior\\\(H(X_t\mid\mathcal F_t)\)};
  \node[fw use,text width=2.8cm,right=of old] (inj)
    {execution + transition\\\(h_{A,t}+h_{P,t}\)};
  \node[fw geom,text width=3.0cm,right=of inj] (pre)
    {pre-observation\\uncertainty};
  \node[fw provenance,text width=3.2cm,below=11mm of pre] (info)
    {revealed information\\\(I(X_{t+1};Z_{t+1}^{\rm pub}\mid\cdot)\)};
  \node[fw terminal,text width=2.8cm,right=25mm of pre] (post)
    {next posterior\\uncertainty};
  \draw[fw arrow] (old)--(inj);
  \draw[fw arrow] (inj)--(pre);
  \draw[fw arrow] (pre)--(post);
  \draw[fw arrow] (info)--node[below right,font=\scriptsize]{subtract}(post);
  \node[fw use,text width=3.5cm,below=13mm of inj] (dir)
    {information floor \(\Rightarrow\) posterior upper bound\\
     capacity ceiling \(\Rightarrow\) posterior lower bound};
  \draw[fw dashed arrow] (pre)--(dir);
  \node[fw residual,text width=3.4cm,below=10mm of post] (gate)
    {winner-bit entropy\\\(H_2^{-1}(H)\le p_{\rm MAP}\le H/2\)};
  \draw[fw dashed arrow] (post)--(gate);
\end{tikzpicture}%
}
\caption{Belief entropy has an exact update: execution and transition
uncertainty form the pre-observation reservoir, and public information removes
mutual information from it.  An information lower bound controls posterior
uncertainty from above, whereas a capacity upper bound controls it only from
below.}
\label{fig-belief-entropy-balance}
\end{figure}

\section{Typed Entropy, Affordable Success, and Blackwell Order}
\label{sec-controlled-typed-entropy-success}

\subsection{Typed entropy and affordable success}
\label{subsec-controlled-typed-entropy-affordable-success}

\begin{definition}[Typed entropy outputs and no-substitution rule]
\label{def-controlled-typed-entropy-outputs}

One complete controlled record stores the entropy vector
\begin{equation}
\label{eq-controlled-typed-entropy-vector}
\mathbf H_{\rm ctl}
=\bigl(
\mathsf H_{\rm init},\mathsf H_{\rm ch},\mathsf H_{\rm env},\mathsf H_R,C_{\rm pub},
\operatorname{KL}_{\rm path},
\operatorname{Ent}_\mu,
\operatorname{Act}_{\Caus},
\operatorname{KL}_{\rm PB},
\log N_{\rm out},
H_\pi
\bigr).
\end{equation}
The initial-retained-state, physical-channel, stochastic-transition,
clean-reward, public-capacity,
policy, and \(\log_2N_{\rm out}\) entries are in bits.  Path-law KL,
PAC--Bayes KL, and
MLSI relative entropy use nats unless an explicit conversion is applied.
Caus action has edge-action units; it equals a thermodynamic entropy-
production functional only under the logarithmic-mean flux hypotheses of the
corresponding Caus theorem.  It is not called entropy production merely
because it is nonnegative.  Conversion from nats to bits multiplies by
\(\log_2e\).

Entries are compared or added only through a theorem with a common output.
In particular:
\begin{enumerate}[label=\textup{(\roman*)}]
\item MLSI density decay does not bound observation-channel capacity without
      a represented channel/semigroup comparison;
\item channel or transition entropy does not bound reward, action, observation,
      or target
      distortion without a represented inverse, coupling, metric, or
      calibration;
\item Caus action enters target contact only through
      \begin{equation}
      \label{eq-controlled-entropy-aiming-contact-rule}
      \operatorname{Hit}_{\rm contact}
      \le
      B^{\rm contact}
      +K\sqrt{\operatorname{Act}_{\Caus}
              \operatorname{Aim}_{\Caus}}
      +\delta,
      \end{equation}
      after the represented aiming-to-contact conversion; and
\item a Lift transports this numerical bound to another carrier only through
      a represented form, intertwining, clocked-defect, or resolvent
      comparison.
\end{enumerate}
Missing aiming or contact converters give the probability-neutral value one;
missing upper-error converters give \(+\infty\), and missing lower-rate
certificates give zero.

\end{definition}

\Cref{fig-typed-entropy-conversion-gates} summarizes the no-substitution
rule for the controlled entropy vector.

\begin{figure}[tbp]
\centering
\resizebox{\linewidth}{!}{%
\begin{tikzpicture}[fw diagram,node distance=6mm and 11mm]
  \node[fw source,text width=3.6cm] (typed)
    {typed inputs\\bits; path/PAC--Bayes/MLSI nats; Caus action;
     description length};
  \node[fw provenance,text width=3.3cm,right=of typed] (conv)
    {represented conversion theorem\\same law, units, and record};
  \node[fw terminal,text width=3.0cm,above right=11mm and 13mm of conv] (fano)
    {capacity \(\to\) Fano/RD};
  \node[fw terminal,text width=3.0cm,right=13mm of conv] (pinsker)
    {path KL \(\to\) event shift};
  \node[fw terminal,text width=3.0cm,below right=11mm and 13mm of conv] (contact)
    {Caus + aiming + baseline\\\(\to\) contact};
  \draw[fw arrow] (typed)--(conv);
  \draw[fw arrow] (conv)--(fano);
  \draw[fw arrow] (conv)--(pinsker);
  \draw[fw arrow] (conv)--(contact);
  \node[fw residual,text width=3.7cm,below=12mm of typed] (blocked)
    {no cross-type addition or minimization\\without the middle gate};
  \draw[fw dashed arrow] (typed)--(blocked);
\end{tikzpicture}%
}
\caption{The controlled entropy vector is typed, not additive.  Shannon
bits, path-law and statistical KL, MLSI entropy, Caus action, and description
length reach a common output only through a represented conversion theorem.
In particular, Caus action reaches target contact only with aiming and a
baseline.}
\label{fig-typed-entropy-conversion-gates}
\end{figure}

\begin{definition}[Affordable multichannel success profile]
\label{def-controlled-affordable-multichannel-success}

Let \(\mathscr P_{n,b,m,H}^{\rm ch}\) be a represented class of complete
POMDP or channelized-control records of cost at most \(b(n)\), sample size
\(m\), and horizon \(H\).  Fix a represented family on which
\(\Theta\) is uniform with size \(N_n\ge2\).  For a target-identification
task, define
\begin{equation}
\label{eq-controlled-affordable-multichannel-success}
\operatorname{Succ}_{\rm ch}^{\rm sat}(n,b(n),m,H)
=\sup_{\mathcal R\in\mathscr P_{n,b,m,H}^{\rm ch}}
\Pr_{\mathcal R}\{\widehat\Theta=\Theta\}.
\end{equation}
For each record, let \(U_{\rm sel}(\mathcal R)\) be its prospective-selector
probability bound or a qualified-source bound already converted to the same
success probability by a represented same-record contact theorem; its value
is one when that converter is absent.  Let \(U_{\rm path}(\mathcal R)\) be the minimum of its
same-event model, path-TV/KL, Caus-contact, capacity, and Lift bounds after
conversion to the same success probability.  In particular, every path-law
candidate has the form
\[
\min\{1,U_{\rm ref}(\mathcal R)+\varepsilon_{\rm path}(\mathcal R)\},
\]
where \(U_{\rm ref}\) bounds the reference record's same target-success event
and \(\varepsilon_{\rm path}\) is the represented TV or Pinsker shift.
Model, capacity, Caus, and Lift candidates are admitted only after their own
baseline or contact term is included.  A missing baseline or converter gives
the value one.  Finally, put
\[
U_{\rm info}(\mathcal R)
=\min\left\{1,
\frac{I_0^{(m)}(\mathcal R)+C_{\rm pub}^{[H]}(\mathcal R)+1}
     {\log_2N_n}
\right\}.
\]
Recordwise unavailable probability bounds equal one.  Then
\begin{equation}
\label{eq-controlled-affordable-multichannel-minimum}
\operatorname{Succ}_{\rm ch}^{\rm sat}
\le
\min\left\{
\sup_{\mathcal R}U_{\rm sel}(\mathcal R),
\sup_{\mathcal R}U_{\rm path}(\mathcal R),
\sup_{\mathcal R}U_{\rm info}(\mathcal R)
\right\}.
\end{equation}
The selector term uses the complete target-aligned pre-hit ledger at
\(\Psi_{\rm ctl}(b(n),H,n)\); the information term is a complexity-free
converse.  For target hitting rather than explicit identification, the
information term is eligible only when the represented hit/output event
identifies \(\Theta\), for example through disjoint verified target sectors.
The output-to-\(\Theta\) decoder, sector verification, pre-hit or terminal
admissibility, and their evaluation costs belong to the complete record.  If
uniformity, \(N_n\), the decoder, either information budget, or the
identification gate is unavailable, the information numerator is assigned
\(+\infty\) and the clipped probability candidate is one.
For polynomial capacity substitute \(b_C(n)=n^C\) everywhere.
For a fixed clean model and channel tuple \(\mathbf K\), write
\(\mathscr P_{n,b,m,H}^{\rm ch}(\mathbf K)\) for the corresponding restricted
complete-record class and
\(\operatorname{Succ}_{\rm ch}^{\rm sat}(n,b(n),m,H;\mathbf K)\) for
\cref{eq-controlled-affordable-multichannel-success} with that class.  This
semicolon notation suppresses no channel, simulation, or garbling cost.

\end{definition}

\begin{proof}[Justification of
\cref{eq-controlled-affordable-multichannel-minimum}]

Fix one complete record and write its success probability as \(p_{\mathcal
R}\).  Prospective-selector accountability, including the complete channel,
filter, policy, and deployment ancestry, gives
\(p_{\mathcal R}\le U_{\rm sel}(\mathcal R)\).  Every admitted path candidate
first applies the TV, KL, Caus, capacity, or Lift comparison and then adds its
represented upper bound for the same reference event, so
\(p_{\mathcal R}\le U_{\rm path}(\mathcal R)\).  Under the represented
identification gate, \cref{eq-controlled-multichannel-fano-success} gives
\(p_{\mathcal R}\le U_{\rm info}(\mathcal R)\); without it the value one
makes the inequality automatic.  Therefore
\[
p_{\mathcal R}\le
\min\{U_{\rm sel}(\mathcal R),U_{\rm path}(\mathcal R),
       U_{\rm info}(\mathcal R)\}.
\]
Taking the affordable supremum and using
\(\sup_{\mathcal R}\min_jU_j(\mathcal R)
\le\min_j\sup_{\mathcal R}U_j(\mathcal R)\) proves the claim.

\end{proof}

\Cref{fig-affordable-multichannel-success-order} shows why the bound is
formed recordwise before saturation.

\begin{figure}[tbp]
\centering
\resizebox{\linewidth}{!}{%
\begin{tikzpicture}[fw diagram,node distance=7mm and 11mm]
  \node[fw source,text width=2.7cm] (r) {one complete record\\\(\mathcal R\)};
  \node[fw use,text width=2.8cm,above right=11mm and 12mm of r] (sel)
    {selector/source\\target-alignment gate};
  \node[fw use,text width=2.8cm,right=of r] (path)
    {path/model\\same-event baseline};
  \node[fw use,text width=2.8cm,below right=11mm and 12mm of r] (info)
    {information\\identification gate};
  \node[fw provenance,text width=3.0cm,right=20mm of path] (min)
    {recordwise\\\(p_{\mathcal R}\le\min_jU_j(\mathcal R)\)};
  \node[fw terminal,text width=3.2cm,right=of min] (sup)
    {saturation\\\(\sup_{\mathcal R}\min_jU_j\le
      \min_j\sup_{\mathcal R}U_j\)};
  \draw[fw arrow] (r)--(sel);
  \draw[fw arrow] (r)--(path);
  \draw[fw arrow] (r)--(info);
  \draw[fw arrow] (sel)--(min);
  \draw[fw arrow] (path)--(min);
  \draw[fw arrow] (info)--(min);
  \draw[fw arrow] (min)--(sup);
\end{tikzpicture}
}
\caption{Affordable success is bounded recordwise before saturation.
Selector, path/model, and information candidates enter only after their
target-alignment, same-event-baseline, and identification gates are satisfied;
unavailable probability candidates equal one.}
\label{fig-affordable-multichannel-success-order}
\end{figure}

\subsection{Blackwell order for represented channels}
\label{subsec-controlled-blackwell-order}

\begin{theorem}[Multichannel degradation is a Blackwell partial order]
\label{thm-controlled-multichannel-blackwell-monotonicity}

Let two complete channel records have the same clean input, and include their
initial and training public records in the following comparison.  Let every
public channel of the second be a represented causal garbling of the
corresponding channel of the first:
\[
K_{q,t}^{(2)}=K_{q,t}^{(1)}G_{q,t}
\]
in the declared history order.  Suppose the garbling, its random bits, state,
and precision enlarge the ceiling by the class-preserving map
\(b^\sharp=\psi_G(b,n,m,H)\).  For an actuator channel, additionally assume
that the full closed-loop factorization permits this garbling before the
physical transition, so that the second executed kernel is a represented
causal simulation from the first record.  Equivalently, the entire second
conditional path kernel, not only its reported actuator symbol, must be a
causal postprocessing of the first.  Then
\begin{equation}
\label{eq-controlled-blackwell-success-monotonicity}
\operatorname{Succ}_{\rm ch}^{\rm sat}
(n,b^\sharp(n),m,H;\mathbf K^{(1)})
\ge
\operatorname{Succ}_{\rm ch}^{\rm sat}
(n,b(n),m,H;\mathbf K^{(2)}).
\end{equation}
For the same compatible input-law class \(\mathcal Q\) at every history,
conditional capacity satisfies
\(C_{{\rm pub},2}^{[H]}\le C_{{\rm pub},1}^{[H]}\).  More precisely, for
each policy \(\pi_2\) in the second record, let \(\pi_{1\to2}\) be its
matched online simulation in the first record, and denote the resulting
transcripts by \(\mathsf Z_H^{2,\pi_2}\) and
\(\mathsf Z_H^{1,\pi_{1\to2}}\).  Then
\[
I(\Theta;\mathsf Z_H^{2,\pi_2})
\le I(\Theta;\mathsf Z_H^{1,\pi_{1\to2}}).
\]
This inequality is not asserted for two independently chosen policies.  For
rate--distortion, retain the common degraded initial/training side information
\(\mathsf Z_0^{(2)}\) in both simulated experiments and use the same
conditional source and distortion.  The common inverse then satisfies
\[
D_{\min}(C_{{\rm pub},2}^{[H]})
\ge D_{\min}(C_{{\rm pub},1}^{[H]}).
\]
If one instead conditions each experiment on its different native
\(\mathsf Z_0^{(i)}\), no comparison of the two conditional
rate--distortion functions is made without a separate side-information
ordering theorem.
No scalar
``total-noise'' threshold is declared for incomparable channel tuples.
Downward-closed multichannel computational regions are defined only along a
represented garbling chain.

\end{theorem}

\begin{proof}

Fix an algorithm admitted for the second record.  From the first initial and
training records, apply the represented initial/training garblings.  At every
later reveal, condition on the already simulated second-record history,
apply \(G_{q,t}\) to the first-record output, and feed the result to the same
policy, memory update, and estimator.  Induction over reveal slots shows that
the simulated transcript has exactly the second-record law.  The actuator
hypothesis ensures the induction includes the physical transition rather
than merely an action report.  Complete-cost closure charges every garbling
and gives \(b^\sharp\); taking the supremum over second-record algorithms
proves \cref{eq-controlled-blackwell-success-monotonicity}.

Under the matched simulation, the second transcript is a causal
postprocessing of the first transcript, including its initial and training
coordinates.  Data processing therefore gives the displayed matched-policy
information inequality.  Applying the same argument to each admitted
history and compatible input law gives \(C_2\le C_1\) slotwise and after
summation.  After both experiments retain the common degraded side
information \(\mathsf Z_0^{(2)}\), the conditional source, distortion, and
rate--distortion function are literally the same.  Its inverse floor is
nonincreasing in the capacity argument; hence the smaller second deployment
budget has the larger distortion floor.  Distinct native side information
would change the function itself, which is why the theorem makes no such
unqualified comparison.

\end{proof}

\Cref{fig-multichannel-blackwell-order} depicts the represented degradation
order and its actuator guard.

\begin{figure}[tbp]
\centering
\resizebox{.94\linewidth}{!}{%
\begin{tikzpicture}[fw diagram,node distance=8mm and 12mm]
  \node[fw source,text width=2.9cm] (k1) {more informative tuple\\\(\mathbf K^{(1)}\)};
  \node[fw provenance,text width=3.2cm,right=of k1] (g)
    {represented causal garbling\\\(G_{q,t}\), cost \(b\mapsto b^\sharp\)};
  \node[fw terminal,text width=2.9cm,right=of g] (k2)
    {degraded tuple\\\(\mathbf K^{(2)}\)};
  \draw[fw arrow] (k1)--(g);
  \draw[fw arrow] (g)--(k2);
  \node[fw use,text width=4.0cm,below=11mm of g] (sim)
    {online simulation of every \(\mathbf K^{(2)}\)-algorithm from
     \(\mathbf K^{(1)}\)};
  \draw[fw arrow] (g)--(sim);
  \node[fw geom,text width=4.4cm,right=12mm of sim] (mon)
    {comparison of \(\mathbf K^{(1)}\) to \(\mathbf K^{(2)}\):\\
     success and capacity \(\ge\); distortion floor \(\le\)};
  \draw[fw arrow] (sim)--(mon);
  \node[fw residual,text width=4.0cm,left=12mm of sim] (act)
    {actuator case: simulate full conditional physical path kernel,
     not only a reported action};
  \draw[fw dashed arrow] (act)--(sim);
\end{tikzpicture}%
}
\caption{Represented causal garbling defines a partial, not scalar,
degradation order.  Online simulation proves success monotonicity and data
processing proves information monotonicity.  For actuator noise, the
garbling must simulate the complete conditional physical path kernel before
transition.}
\label{fig-multichannel-blackwell-order}
\end{figure}

\subsection{Pre-hit decision value and sufficiency}
\label{subsec-controlled-prehit-decision-sufficiency}

\begin{definition}[Represented pre-hit target-decision value]
\label{def-controlled-prehit-target-decision-value}

Fix a presented task, a time \(0\le t<H\), its survival event
\(\{\tau>t\}\), and a complete pre-hit record \(\mathcal H_t\).  Let
\(\Theta\) denote the target parameter and let
\(\mathcal A_t(\mathcal H_t)\) be the represented carrier of decisions
available after that record.  A bounded target-decision utility is a common
measurable function

\[
u_t:\{(\theta,a):a\in\mathcal A_t\}\longrightarrow[0,1].
\]

The utility is an analytic evaluator of a decision.  It is not supplied to a
selector as pre-hit data.  Any represented use of a utility value in forming
the next decision belongs to the selector's pre-hit ancestry and is charged
under \cref{def-pre-hit-temporally-admissible-source}.

For a represented pre-hit experiment \(\mathsf E\), let
\(\mathscr K_t(\mathsf E,B)\) contain the family-uniform selectors whose
complete prefix, sampler, random-bit, precision, storage, and transition cost
is at most \(B\).  Define
\begin{equation}
\label{eq-controlled-prehit-target-decision-value}
\operatorname{DVal}^{\mathrm{sat}}_{u,t}(\mathsf E;B)
=
\sup_{K_t\in\mathscr K_t(\mathsf E,B)}
\mathbb E_{\mathsf E}\!\left[
 \mathbf 1_{\{\tau>t\}}
 \int u_t(\Theta,a)K_t(da\mid\mathcal H_t)
\right].
\end{equation}
For a declared represented baseline \(\nu_t\) on the same carrier, put
\begin{equation}
\label{eq-controlled-prehit-target-decision-advantage}
\operatorname{DAdv}^{\mathrm{sat}}_{u,t}
 (\mathsf E,\nu_t;B)
=
\sup_{K_t\in\mathscr K_t(\mathsf E,B)}
\left(
\mathbb E_{\mathsf E}\!\left[
 \mathbf 1_{\{\tau>t\}}
 \int u_t(\Theta,a)
      \bigl(K_t-\nu_t\bigr)(da\mid\mathcal H_t)
\right]
\right)_+.
\end{equation}
For the next-test utility \(u_t(\Theta,a)=V_I(a)\),
\cref{eq-controlled-prehit-target-decision-advantage} is the selector
advantage of \cref{eq-pre-hit-selector-advantage}; taking the supremum over
times gives the corresponding prospective selector profile.  A terminal
identification utility uses
\(u_t(\Theta,a)=\mathbf 1_{\{a=\Theta\}}\).

\end{definition}

\begin{theorem}[Represented pre-hit Blackwell domination]
\label{thm-controlled-prehit-blackwell-value-monotonicity}

Let \(\mathsf E_1\) and \(\mathsf E_2\) be two complete pre-hit experiments
for the same presented task, target \(\Theta\), verifier, and target-decision
utility \(u_t\).  Include the initial record, training record, and the
complete prefix through time \(t\).  Suppose that \(\mathsf E_2\) is
generated from \(\mathsf E_1\) by represented causal kernels
\(G_0,\ldots,G_t\): after each \(\mathsf E_1\)-reveal and before the next
decision, \(G_j\) produces the corresponding \(\mathsf E_2\)-reveal from the
already generated \(\mathsf E_2\)-history.  The kernels and their random bits
have no later verifier flag, terminal record, target answer, or later oracle
answer as an ancestor.  The resulting online simulation reproduces the joint
law of

\[
(\Theta,\mathbf 1_{\{\tau>t\}},\mathcal H_t^{(2)})
\]

in the second experiment.  The decision carriers are common; more generally,
a represented action map
\(\iota_t(\mathcal H_t^{(1)},\mathcal H_t^{(2)},a_2)\) may be used when it
preserves utility almost surely:
\[
u_t\!\left(\Theta,
 \iota_t(\mathcal H_t^{(1)},\mathcal H_t^{(2)},a_2)\right)
=u_t(\Theta,a_2).
\]
For an actuator decision, the simulation must reproduce the complete
conditional physical transition, as in
\cref{thm-controlled-multichannel-blackwell-monotonicity}.

Assume that the complete costs of the causal garbling, action map, random
bits, state, storage, precision, and evaluation enlarge \(B\) by the
class-preserving map

\[
B^\sharp=\psi_{\rm BW}(B,n,m,H).
\]

Then
\begin{equation}
\label{eq-controlled-prehit-blackwell-value-monotonicity}
\operatorname{DVal}^{\mathrm{sat}}_{u,t}(\mathsf E_2;B)
\le
\operatorname{DVal}^{\mathrm{sat}}_{u,t}(\mathsf E_1;B^\sharp).
\end{equation}

More precisely, if \(K_t^{(2)}\) is a selector in the second experiment, its
simulated selector in the first has kernel
\begin{equation}
\label{eq-controlled-prehit-blackwell-kernel}
K_t^{(1\to2)}(da_1\mid h_1)
=
\int G_{\le t}(dh_2\mid h_1)
     \int K_t^{(2)}(da_2\mid h_2)
     \mathbf 1_{\{\iota_t(h_1,h_2,a_2)\in da_1\}},
\end{equation}
with \(\iota_t(a_2)=a_2\) on a common carrier, and it satisfies the exact
identity
\begin{equation}
\label{eq-controlled-prehit-blackwell-utility-identity}
\mathbb E_{\mathsf E_1}\!\left[
 \mathbf 1_{\{\tau>t\}}
 \int u_t(\Theta,a)K_t^{(1\to2)}(da\mid\mathcal H_t^{(1)})
\right]
=
\mathbb E_{\mathsf E_2}\!\left[
 \mathbf 1_{\{\tau>t\}}
 \int u_t(\Theta,a)K_t^{(2)}(da\mid\mathcal H_t^{(2)})
\right].
\end{equation}

If \(\nu_t^{(2)}\) is the declared baseline in \(\mathsf E_2\), define its
represented pullback \(\nu_t^{(1\to2)}\) by the right-hand side of
\cref{eq-controlled-prehit-blackwell-kernel} with \(K_t^{(2)}\) replaced by
\(\nu_t^{(2)}\).  When the baseline's declared target-neutrality certificate
is preserved by this simulation,
\begin{equation}
\label{eq-controlled-prehit-blackwell-advantage-monotonicity}
\operatorname{DAdv}^{\mathrm{sat}}_{u,t}
 (\mathsf E_2,\nu_t^{(2)};B)
\le
\operatorname{DAdv}^{\mathrm{sat}}_{u,t}
 (\mathsf E_1,\nu_t^{(1\to2)};B^\sharp).
\end{equation}
For \(u_t=V_I\), this is Blackwell monotonicity of the represented pre-hit
selector profile, with the complete garbling cost charged before the
decision.

The comparison concerns decision value, not provenance metadata.  A
represented statistic, constructor label, or backward-cut classification
establishes a garbling only in the direction from the fuller record to that
derived record.  It therefore bounds the decision value of selectors that
factor through the derived record.  An upper bound for the fuller record
follows only from a reverse represented garbling, as in
\cref{cor-controlled-represented-prehit-sufficiency}, or from a separate
numerical bound on its target-decision value.

\end{theorem}

\begin{proof}

Fix a complete second-experiment record of cost at most \(B\) and its
selector \(K_t^{(2)}\).  Generate its initial, training, and online prefix by
applying \(G_0,\ldots,G_t\) in causal order to the first-experiment record.
All random bits used by \(G_j\) are generated before the decision following
slot \(j\).  Hence the simulated prefix and the kernel in
\cref{eq-controlled-prehit-blackwell-kernel} factor through
\(\operatorname{pref}_t\) and satisfy the four temporal-source conditions of
\cref{def-pre-hit-temporally-admissible-source}.  The cost hypothesis charges
the complete simulation and places the resulting selector below
\(B^\sharp\).

By the joint-law hypothesis, the simulated triple
\((\Theta,\mathbf 1_{\{\tau>t\}},\mathcal H_t^{(2)})\) has its native
second-experiment law.  Conditional on the simulated history, draw
\(a_2\) from \(K_t^{(2)}\), and apply \(\iota_t\).  Utility preservation and
the tower property give
\cref{eq-controlled-prehit-blackwell-utility-identity}.  Taking the supremum
over all cost-\(B\) second-experiment records proves
\cref{eq-controlled-prehit-blackwell-value-monotonicity}.

Apply the same construction to \(\nu_t^{(2)}\).  The selector and baseline
utility expectations are each preserved exactly, so their difference is
preserved exactly.  Taking the positive part and then the saturated supremum
proves
\cref{eq-controlled-prehit-blackwell-advantage-monotonicity}.  The final
paragraph follows from the direction of
\cref{eq-controlled-prehit-blackwell-value-monotonicity}: a statistic
computed from a record is its degradation, whereas an upper bound for the
record requires a simulation of the record from the statistic or a direct
bound on the record itself.

\end{proof}

\begin{corollary}[Represented pre-hit sufficient statistic]
\label{cor-controlled-represented-prehit-sufficiency}

Let \(\mathcal H_t\) be a complete pre-hit record and let
\(S_t=s_t(\mathcal H_t)\) be a represented statistic computed before the
time-\(t+1\) decision.  Suppose that:

\begin{enumerate}[label=\textup{(\roman*)}]
\item the extraction \(s_t\), including its storage, precision, and evaluation,
      is temporally admissible and has complete cost map
      \(B\mapsto\psi_s(B,n,m,H)\);
\item there is a represented pre-hit reconstruction kernel
      \(R_t(dh\mid S_t)\), of complete cost map
      \(B\mapsto\psi_R(B,n,m,H)\), such that on \(\{\tau>t\}\)
      \begin{equation}
      \label{eq-controlled-represented-prehit-sufficiency}
      \mathcal L(\mathcal H_t\mid\Theta,S_t,\tau>t)
      =R_t(\,\cdot\mid S_t)
      =\mathcal L(\mathcal H_t\mid S_t,\tau>t);
      \end{equation}
\item the reconstruction uses no later flag, terminal record, target answer,
      or later oracle answer, and preserves the represented decision carrier
      and target utility, possibly through the action map of
      \cref{thm-controlled-prehit-blackwell-value-monotonicity}.
\end{enumerate}

Then \(S_t\) and \(\mathcal H_t\) are equivalent for the survival-weighted
target decisions of
\cref{def-controlled-prehit-target-decision-value}, with their complete
representation costs retained:
\begin{align}
\label{eq-controlled-sufficient-statistic-value-equivalence-a}
\operatorname{DVal}^{\mathrm{sat}}_{u,t}(S_t;B)
&\le
\operatorname{DVal}^{\mathrm{sat}}_{u,t}
 (\mathcal H_t;\psi_s(B,n,m,H)),\\
\label{eq-controlled-sufficient-statistic-value-equivalence-b}
\operatorname{DVal}^{\mathrm{sat}}_{u,t}(\mathcal H_t;B)
&\le
\operatorname{DVal}^{\mathrm{sat}}_{u,t}
 (S_t;\psi_R(B,n,m,H)).
\end{align}
The same two inequalities hold for
\(\operatorname{DAdv}^{\mathrm{sat}}_{u,t}\) when the declared baseline is
transported through the two kernels.  At a saturated ceiling closed under
both class-preserving cost maps, the two target-decision profiles coincide.

Thus a represented sufficient statistic can carry the complete optimal
target-decision value of the pre-hit record.  A constructor or provenance
condition alone is not sufficient in this sense;
\cref{eq-controlled-represented-prehit-sufficiency} and the affordable
reconstruction are the conditions that permit an upper bound to pass from
the statistic to the original record.

\end{corollary}

\begin{proof}

The statistic \(S_t=s_t(\mathcal H_t)\) is a deterministic represented
garbling of \(\mathcal H_t\).  Applying
\cref{thm-controlled-prehit-blackwell-value-monotonicity} with cost map
\(\psi_s\) proves
\cref{eq-controlled-sufficient-statistic-value-equivalence-a}.

Conversely, on \(\{\tau>t\}\), sample
\(\widetilde{\mathcal H}_t\) from \(R_t(\cdot\mid S_t)\); on the failure
sector retain the represented survival flag and use an arbitrary fixed
padded history, because the decision payoff there is multiplied by zero.
By \cref{eq-controlled-represented-prehit-sufficiency}, the tower property
gives equality of every survival-weighted utility expectation under
\((\Theta,S_t,\widetilde{\mathcal H}_t)\) and
\((\Theta,S_t,\mathcal H_t)\).  Composing any cost-\(B\) history selector
with this represented, pre-hit, fully charged reconstruction therefore
produces a statistic selector of cost at most
\(\psi_R(B,n,m,H)\) with the same objective value.  Taking the supremum
proves \cref{eq-controlled-sufficient-statistic-value-equivalence-b}
directly.  Transporting the baseline through the same two simulations proves
the advantage statements.  If a common saturated ceiling absorbs both cost
maps, the two opposite inequalities have the same feasible ceiling and give
equality.

\end{proof}

\section{Reversible Carrier Geometry and Authorized Flow}
\label{sec-controlled-reversible-geometry}

The selected generator supplies a native Dirichlet geometry only when its
reference measure and reversible conductance are represented.  Otherwise a
separately represented comparison geometry is used with the existing
form-comparison loss.

\begin{definition}[Certified dynamical geometry]
\label{def-certified-dynamical-geometry}

Fix a controlled policy record and one selected generator \(L\) on a finite
carrier \(X\).  A certified dynamical geometry at budget \(B\) is an ambient
Geom/Caus record of \cref{def-geom-extraction}, augmented by represented
initial, target, and failure data.  It has one of the following two carrier
forms.

\begin{enumerate}[label=\textup{(\roman*)}]
\item \emph{Native carrier.}  A full-support represented probability measure
      \(\mu\) is invariant for \(L\).  Writing
      \(L=L^{\mathsf s}+A+L^{\mathsf b}\) in \(L^2(X,\mu)\), the symmetric
      conductance is
      \[
      w_L(x,z)=\frac12\bigl(\mu(x)L(x,z)+\mu(z)L(z,x)\bigr)\ge0
      \quad (x\ne z),
      \]
      and
      \[
      \mathcal E_L(f,g)
      =\frac12\sum_{x,z}w_L(x,z)
        (f(x)-f(z))(g(x)-g(z)).
      \]
      Equivalently, the conductance in
      \cref{def-static-dirichlet-geometry} is
      \(c_L(x,z)=w_L(x,z)/(\mu(x)\mu(z))\).
\item \emph{Comparison carrier.}  A represented pointwise geometry
      \(\Phi=(E_\Phi,c_\Phi)\) is related to the selected generator by a
      represented local, resolvent, or form-comparison record with the error
      and transfer loss prescribed by
      \cref{lem-declared-geom-form-error,prop-geom-form-resolvent-comparison}.
\end{enumerate}

The certificate retains the target-normalized potential \(D\), the selected
actual current and authority record \((J,\mathfrak K)\), target reach
\(R_T(D)\), regularity \(\rho_\Phi(D)\), and the mixing and local-consistency
gates of \cref{def-geom-extraction}.  It may additionally retain certified
constants, all evaluated on this same carrier:
\begin{align*}
\gamma_\Phi
&=\inf_{\operatorname{Var}_\mu f>0}
  \frac{\mathcal E_\Phi(f,f)}{\operatorname{Var}_\mu(f)},\\
N_{\rm eff,\Phi}(\lambda)
&=\operatorname{Tr}\!\left(H_\Phi(H_\Phi+\lambda I)^{-1}\right),
\qquad H_\Phi\ge0,
\end{align*}
together with the convention that \(H_\Phi\) is the nonnegative
self-adjoint operator associated with \(\mathcal E_\Phi\) in the represented
finite realization, and with a conductance lower bound, an MLSI constant, a
Bakry--\'Emery curvature bound, or a condenser-capacity value when the
corresponding represented proof or audited finite computation is included.

An unavailable lower-rate certificate is assigned value zero.  An unavailable
upper error or cost certificate is assigned value \(+\infty\).  These
conventions supply no positive mixing, alignment, reach, or affordability
claim.  Every retained constant, eigenvalue computation, comparison map, and
precision parameter contributes to the complete saturated cost.

\end{definition}

\begin{proposition}[Native geometry and Hodge authority]
\label{prop-controlled-native-geometry-hodge}

In the native-carrier case of
\cref{def-certified-dynamical-geometry}, \(\mathcal E_L\) is nonnegative and
\[
-\langle f,L^{\mathsf s}f\rangle_{L^2(\mu)}
=\mathcal E_L(f,f).
\]
The skew part contributes no quadratic energy:
\(\langle f,Af\rangle_{L^2(\mu)}=0\).  Its unique structurally actionable
component is the authorized Hodge projection in
\cref{def-hodge-channel-split,def-caus-authority-envelope}; the orthogonal
component belongs to \(J_{\rm res}\).  Therefore the selected policy can use
target-oriented directed flow only through the same-record Caus factor
\(C_{\Phi,\mathfrak K}^{\rm auth}(D)\).

\end{proposition}

\begin{proof}

Invariance gives zero row and column sums in the weighted realization.
Pairing the symmetric off-diagonal terms yields the displayed Dirichlet form,
which is nonnegative because every \(w_L(x,z)\) is nonnegative.  Skew
adjointness gives \(\langle f,Af\rangle=0\).  The remaining statements are
the uniqueness and target-pairing identities of
\cref{thm-reversible-geom-exhaustive-directed-classification,thm-caus-authority-no-substitution}.

\end{proof}

\begin{theorem}[Certified mixing and safe spectral compression]
\label{thm-controlled-certified-mixing}

Let \(L=L^{\mathsf s}+A\) have a represented invariant law \(\mu\), with
\(A^*=-A\), and suppose the native symmetric form has certified gap
\(\gamma_\Phi>0\).  For the Markov semigroup \(P_t=e^{tL}\),
\begin{equation}
\label{eq-controlled-l2-mixing}
\operatorname{Var}_\mu(P_tf)
\le e^{-2\gamma_\Phi t}\operatorname{Var}_\mu(f).
\end{equation}
If a certified MLSI with constant \(\alpha_\Phi>0\) is part of the same
record, its declared entropy normalization gives the corresponding entropy
decay
\begin{equation}
\label{eq-controlled-entropy-mixing}
\operatorname{Ent}_\mu(P_t^*g)
\le e^{-2\alpha_\Phi t}\operatorname{Ent}_\mu(g)
\end{equation}
for represented probability densities in its verified domain.  A certified
Bakry--\'Emery inequality supplies its separately stated gradient contraction
for the reversible comparison semigroup; it is transferred to the selected
generator only through a represented comparison theorem.

Let \(P_S\) be an orthogonal projection onto a represented safe or admissible
subspace and put \(H_S=P_SH_\Phi P_S|_{\operatorname{ran}P_S}\).  Then for
every \(\lambda>0\),
\begin{equation}
\label{eq-safe-effective-dimension-compression}
\operatorname{Tr}\!\left(H_S(H_S+\lambda I)^{-1}\right)
\le N_{\rm eff,\Phi}(\lambda).
\end{equation}

\end{theorem}

\begin{proof}

For \(f_t=P_tf\), invariance and skew adjointness give
\[
\frac{d}{dt}\operatorname{Var}_\mu(f_t)
=2\langle f_t,Lf_t\rangle_\mu
=-2\mathcal E_\Phi(f_t,f_t)
\le-2\gamma_\Phi\operatorname{Var}_\mu(f_t).
\]
Gr\"onwall's inequality proves
\cref{eq-controlled-l2-mixing}.  The entropy conclusion is the defining
semigroup consequence of the certified MLSI in the declared normalization;
the restriction on transfer is exactly
\cref{thm-defect-profile-selection}.  For
\cref{eq-safe-effective-dimension-compression}, order the eigenvalues of
\(H_\Phi\) and its compression decreasingly.  Cauchy interlacing gives
\(\eta_j(H_S)\le\eta_j(H_\Phi)\).  The function
\(u\mapsto u/(u+\lambda)\) is increasing on \([0,\infty)\); summing the
termwise inequalities proves the trace bound.

\end{proof}

\subsection{Audited Functional-Inequality Profiles}
\label{subsec-controlled-functional-inequality-profile}

The Dirichlet form of a certified dynamical geometry supports several
inequalities with different outputs.  Their constants are recorded with a
fixed normalization so that rate comparisons are meaningful and so that one
logical certificate is not charged more than once under different names.

\begin{definition}[Normalized functional-inequality profile]
\label{def-controlled-functional-inequality-profile}
For a certified carrier \(\Phi\), its normalized functional-inequality
profile is the typed tuple
\[
\mathfrak F_\Phi
=\bigl(\mathfrak F_\Phi^{\rm diff},
       \mathfrak F_\Phi^{\rm ent},
       \mathfrak F_\Phi^{\rm tar}\bigr)
\]
defined in
\cref{def-controlled-diffusive-functional-profile,def-controlled-entropic-functional-profile,def-controlled-target-nonreversible-profile}.
The three components share one carrier, normalization, proof record, and
cost ledger; storing the same certificate in two components creates no
second charge.
\end{definition}

\begin{definition}[Diffusive and spectral profile]
\label{def-controlled-diffusive-functional-profile}

Let \(\Phi\) be a native or comparison carrier from
\cref{def-certified-dynamical-geometry}, with probability measure \(\mu\),
self-adjoint operator \(H_\Phi\), Dirichlet form \(\mathcal E_\Phi\), and
reversible semigroup \(S_t=e^{-tH_\Phi}\).  In the comparison-carrier case
write \(w_\Phi(x,z)=\mu(x)\mu(z)c_\Phi(x,z)\); in the native case use the
conductance \(w_L\) from \cref{def-certified-dynamical-geometry}.  Put
\[
\operatorname{Osc}(f)=\max_X f-\min_X f,
\qquad
\operatorname{Ent}_\mu(g)
=\mu(g\log g)-\mu(g)\log\mu(g)
\]
for \(g\ge0\), with \(0\log0=0\).  For \(A\subsetneq X\), let
\[
\lambda_\Phi^{\rm D}(A)
=\inf\left\{
\frac{\mathcal E_\Phi(f,f)}{\|f\|_{2,\mu}^{2}}:
0\ne f,\ f|_{A^c}=0
\right\}.
\]
Also set
\[
\lambda_\Phi^{\rm sp}(A)
=\inf\left\{
\frac{\mathcal E_\Phi(f,f)}{\operatorname{Var}_\mu(f)}:
0<\operatorname{Var}_\mu(f),\ f|_{A^c}=0
\right\}.
\]
The following entries constitute the normalized functional-inequality
profile \(\mathfrak F_\Phi\).  An entry is retained only with its represented
proof, audited finite computation, or represented comparison certificate.

\begin{enumerate}[label=\textup{(\roman*)}]
\item The Poincar\'e constant \(\gamma_\Phi\) is the largest certified
      \(\gamma\ge0\) such that
      \[
      \gamma\operatorname{Var}_\mu(f)
      \le \mathcal E_\Phi(f,f).
      \]
      With \(Q_\Phi(A,A^c)=\sum_{x\in A,z\notin A}w_\Phi(x,z)\), define
      \[
      \mathsf h_\Phi
      =\inf_{0<\mu(A)\le1/2}\frac{Q_\Phi(A,A^c)}{\mu(A)},
      \qquad
      q_\Phi=\max_x\frac1{\mu(x)}\sum_{z\ne x}w_\Phi(x,z).
      \]
      Thus \(\mathsf h_\Phi\) has generator-rate units.

\item A weak Poincar\'e rate is a decreasing function
      \(\beta_{\rm W}:(0,1]\to(0,\infty]\) satisfying
      \[
      \operatorname{Var}_\mu(f)
      \le \beta_{\rm W}(r)\mathcal E_\Phi(f,f)
          +r\operatorname{Osc}(f)^2,
      \qquad 0<r\le1.
      \]
      A super-Poincar\'e rate is a decreasing function
      \(\beta_{\rm S}:(0,\infty)\to(0,\infty]\) satisfying
      \[
      \|f\|_{2,\mu}^{2}
      \le r\mathcal E_\Phi(f,f)
          +\beta_{\rm S}(r)\|f\|_{1,\mu}^{2},
      \qquad r>0.
      \]

\item A Nash certificate of dimension \(d>0\) and constant \(C_{\rm N}\)
      satisfies
      \[
      \|f\|_{2,\mu}^{2+4/d}
      \le C_{\rm N}\mathcal E_\Phi(f,f)
                    \|f\|_{1,\mu}^{4/d}
      \]
      on its declared semigroup-invariant mean-zero, killed, or transient
      domain \(\mathscr D_{\rm N}\).  For \(d>2\),
      a Sobolev certificate with constant \(C_{\rm Sob}\) satisfies
      \[
      \|f\|_{2d/(d-2),\mu}^{2}
      \le C_{\rm Sob}\mathcal E_\Phi(f,f)
      \]
      on the same declared domain.

\item The spectral profile and a Faber--Krahn minorant are
      \[
      \Lambda_\Phi^{\rm mix}(v)
      =\inf_{0<\mu(A)\le v}\lambda_\Phi^{\rm sp}(A),
      \quad 0<v<1,
      \qquad \Lambda_\Phi^{\rm mix}(v)=\gamma_\Phi\quad(v\ge1),
      \]
      and
      \[
      \lambda_\Phi^{\rm D}(A)
      \ge c_{\rm FK}\mu(A)^{-2/d}
      \qquad(0<\mu(A)<1).
      \]

\end{enumerate}
\end{definition}

\begin{definition}[Entropic, curvature, and transport profile]
\label{def-controlled-entropic-functional-profile}
Let \(\Phi,\mu,H_\Phi,\mathcal E_\Phi\) be as in
\cref{def-controlled-diffusive-functional-profile}.  The entropic component
\(\mathfrak F_\Phi^{\rm ent}\) consists of the following represented
certificates.
\begin{enumerate}[label=\textup{(\roman*)}]

\item A logarithmic Sobolev constant \(\alpha_{\rm LS}\) and a modified
      logarithmic Sobolev constant \(\alpha_{\rm MLS}\) satisfy,
      respectively,
      \begin{align*}
      \operatorname{Ent}_\mu(f^2)
      &\le \frac{2}{\alpha_{\rm LS}}\mathcal E_\Phi(f,f),\\
      \operatorname{Ent}_\mu(g)
      &\le \frac{1}{2\alpha_{\rm MLS}}
        \mathcal E_\Phi(g,\log g),
      \end{align*}
      for the represented domains declared in the certificate.

\item With
      \[
      \Gamma(f,g)=\frac12\bigl(fH_\Phi g+gH_\Phi f-H_\Phi(fg)\bigr),
      \qquad
      \Gamma_2(f)=-\frac12H_\Phi\Gamma(f)+\Gamma(f,H_\Phi f),
      \]
      a Bakry--\'Emery \(CD(\kappa,N)\) certificate means
      \[
      \Gamma_2(f)\ge \kappa\Gamma(f)+\frac1N(H_\Phi f)^2
      \]
      pointwise, where \(N\in[1,\infty]\) and the last term is zero for
      \(N=\infty\).  The sign convention is the one for the nonnegative
      operator \(H_\Phi\).

\item A transport--entropy certificate for a represented metric \(d_\Phi\)
      is either
      \[
      T_1(C_{T_1}):\quad
      W_{1,d_\Phi}(\nu,\mu)^2
      \le 2C_{T_1}\operatorname{KL}(\nu\Vert\mu),
      \]
      or
      \[
      T_2(C_{T_2}):\quad
      W_{2,d_\Phi}(\nu,\mu)^2
      \le 2C_{T_2}\operatorname{KL}(\nu\Vert\mu),
      \]
      for represented \(\nu\ll\mu\).

\end{enumerate}
\end{definition}

\begin{definition}[Target and nonreversible profile]
\label{def-controlled-target-nonreversible-profile}
Let \(\Phi,\mu,H_\Phi,\mathcal E_\Phi\) be as in
\cref{def-controlled-diffusive-functional-profile}.  The target and
nonreversible component \(\mathfrak F_\Phi^{\rm tar}\) consists of the
following represented certificates.
\begin{enumerate}[label=\textup{(\roman*)}]

\item The target, failure, and isocapacitary profiles are restrictions of the
      represented condenser capacity from
      \cref{def-controlled-condenser-capacity}:
      \begin{align*}
      \mathsf C_{T^+}(v)
      &=\inf_{0<\mu(A)\le v}
        \frac{\operatorname{Cap}_\Phi(A,T^+)}{\mu(A)},\\
      \mathsf C_{T^-}(v)
      &=\inf_{0<\mu(A)\le v}
        \frac{\operatorname{Cap}_\Phi(A,T^-)}{\mu(A)},\\
      \mathsf C_{\rm iso}(v)
      &=\inf_{0<\mu(A)\le v}
        \frac{\operatorname{Cap}_\Phi(A,A^c)}{\mu(A)}.
      \end{align*}
      Empty or intersecting condenser classes are omitted from the infimum.

\item For a selected generator \(L=-H_\Phi+A\), a strong sector constant
      \(C_{\rm sec}\) satisfies
      \[
      |\langle f,Ag\rangle_\mu|
      \le C_{\rm sec}\,
         \mathcal E_\Phi(f,f)^{1/2}
         \mathcal E_\Phi(g,g)^{1/2}.
      \]
      A represented hypocoercive certificate is a quadratic functional
      \(\mathscr H\) and constants \(0<m_{\mathscr H}\le M_{\mathscr H}\),
      \(\kappa_{\mathscr H}>0\), such that on the centered declared domain
      \[
      m_{\mathscr H}\|f\|_2^2\le\mathscr H(f)
      \le M_{\mathscr H}\|f\|_2^2,
      \qquad
      \frac{d}{dt}\mathscr H(P_tf)
      \le-2\kappa_{\mathscr H}\mathscr H(P_tf).
      \]
\end{enumerate}

Rate functions are stored at the finite precision used by the downstream
theorem.  Their representation, optimization grid, domain verification, and
all comparison losses contribute to saturated cost.  A missing lower
constant equals zero, a missing upper constant equals \(+\infty\), and a
missing rate function supplies no bound.

\end{definition}

\begin{theorem}[Quantitative consequences of the audited profile]
\label{thm-controlled-functional-inequality-consequences}

Let \(\mathfrak F_\Phi\) be the profile of
\cref{def-controlled-functional-inequality-profile}.  Each certified entry
has the following consequence in its declared domain.

\begin{enumerate}[label=\textup{(\roman*)}]
\item Poincar\'e gives
      \[
      \operatorname{Var}_\mu(S_tf)
      \le e^{-2\gamma_\Phi t}\operatorname{Var}_\mu(f).
      \]
      If \(q_\Phi>0\), the continuous-time Cheeger bounds are
      \[
      \frac{\mathsf h_\Phi^2}{2q_\Phi}
      \le\gamma_\Phi\le2\mathsf h_\Phi.
      \]

\item For every \(0<r\le1\), weak Poincar\'e gives
      \begin{equation}
      \label{eq-controlled-weak-poincare-decay}
      \operatorname{Var}_\mu(S_tf)
      \le e^{-2t/\beta_{\rm W}(r)}\operatorname{Var}_\mu(f)
         +r\bigl(1-e^{-2t/\beta_{\rm W}(r)}\bigr)
          \operatorname{Osc}(f)^2.
      \end{equation}
      For every \(r>0\), super-Poincar\'e gives
      \begin{equation}
      \label{eq-controlled-super-poincare-smoothing}
      \|S_tf\|_2^2
      \le e^{-2t/r}\|f\|_2^2
         +\beta_{\rm S}(r)(1-e^{-2t/r})\|f\|_1^2.
      \end{equation}
      Both right sides may be minimized over the represented admissible
      values of \(r\).

\item Nash gives the ultracontractive half-step estimate
      \begin{equation}
      \label{eq-controlled-nash-smoothing}
      \|S_t|_{\mathscr D_{\rm N}}\|_{1\to2}^2
      \le\left(\frac{dC_{\rm N}}{4t}\right)^{d/2},
      \qquad t>0,
      \end{equation}
      and self-adjointness gives
      \(\|S_t|_{\mathscr D_{\rm N}}\|_{1\to\infty}
      \le(dC_{\rm N}/(2t))^{d/2}\).
      The stated Sobolev inequality implies the Nash inequality with the same
      constant by interpolation between \(L^1\) and \(L^{2d/(d-2)}\).

\item For the killed semigroup \(S_t^A\),
      \begin{equation}
      \label{eq-controlled-killed-spectral-decay}
      \|S_t^A\|_{2\to2}\le e^{-t\lambda_\Phi^{\rm D}(A)}.
      \end{equation}
      Hence a Faber--Krahn certificate gives
      \(\|S_t^A\|_{2\to2}
      \le e^{-tc_{\rm FK}\mu(A)^{-2/d}}\).
      If the uniformized kernel \(K=I-H_\Phi/q_\Phi\) is lazy, the represented
      generator-rate spectral profile supplies the relative-uniform mixing
      certificate
      \begin{equation}
      \label{eq-controlled-spectral-profile-mixing}
      t_{\infty}(\varepsilon)
      \le\int_{4\mu_*}^{4/\varepsilon}
      \frac{2\,dv}{v\Lambda_\Phi^{\rm mix}(v)},
      \qquad \mu_*=\min_x\mu(x),
      \end{equation}
      whenever the displayed lazy-kernel normalization is valid.

\item Logarithmic Sobolev gives hypercontractivity
      \[
      \|S_tf\|_{q(t)}\le\|f\|_p,
      \qquad q(t)-1=(p-1)e^{2\alpha_{\rm LS}t},
      \quad p>1.
      \]
      Modified logarithmic Sobolev gives
      \begin{equation}
      \label{eq-controlled-mlsi-decay}
      \operatorname{Ent}_\mu(S_t^*g)
      \le e^{-2\alpha_{\rm MLS}t}\operatorname{Ent}_\mu(g)
      \end{equation}
      for probability densities \(g\).

\item \(CD(\kappa,\infty)\) gives the pointwise gradient estimate
      \begin{equation}
      \label{eq-controlled-bakry-emery-contraction}
      \Gamma(S_tf)
      \le e^{-2\kappa t}S_t\Gamma(f).
      \end{equation}
      A finite \(N\) retains the additional nonnegative dimension term in
      the integrated \(\Gamma_2\) estimate.

\item If \(F\) is \(L_F\)-Lipschitz for \(d_\Phi\), then
      \begin{align}
      |\nu(F)-\mu(F)|
      &\le L_F\sqrt{2C_{T_1}\operatorname{KL}(\nu\Vert\mu)},
      \label{eq-controlled-t1-shift}\\
      |\nu(F)-\mu(F)|
      &\le L_F\sqrt{2C_{T_2}\operatorname{KL}(\nu\Vert\mu)}.
      \label{eq-controlled-t2-shift}
      \end{align}
      The second line uses \(W_1\le W_2\).

\item A represented condenser capacity gives the exact Dirichlet and
      reach--avoid identities in
      \cref{thm-controlled-reach-avoid-dirichlet-drift}.  The profiles
      \(\mathsf C_{T^+}\), \(\mathsf C_{T^-}\), and
      \(\mathsf C_{\rm iso}\) therefore provide, respectively, certified
      target-access, failure-access, and killed-set energy scales at every
      represented mass level \(v\).

\item A form comparison
      \(a\mathcal E_\Phi\le\mathcal E_\Psi\le b\mathcal E_\Phi\)
      transfers every Rayleigh-quotient lower bound by the factor \(a\) and
      every energy upper bound by the factor \(b\).  The sector constant
      controls the skew bilinear term but supplies no decay rate by itself.
      A hypocoercive certificate gives
      \begin{equation}
      \label{eq-controlled-hypocoercive-decay}
      \|P_tf\|_2^2
      \le\frac{M_{\mathscr H}}{m_{\mathscr H}}
      e^{-2\kappa_{\mathscr H}t}\|f\|_2^2.
      \end{equation}
\end{enumerate}

All transfers to the selected nonreversible generator use the represented
comparison, sector, or hypocoercive record carried by the same certified
dynamical geometry.

\end{theorem}

\begin{proof}

The Poincar\'e decay is the energy derivative calculation in
\cref{thm-controlled-certified-mixing}.  Uniformization by
\(K=I-H_\Phi/q_\Phi\) and the discrete Cheeger inequalities give the two
Cheeger bounds after multiplying by \(q_\Phi\).

For weak Poincar\'e, set \(y(t)=\operatorname{Var}_\mu(S_tf)\).  Markov
contraction preserves oscillation, and
\[
y'(t)=-2\mathcal E_\Phi(S_tf,S_tf)
\le-\frac2{\beta_{\rm W}(r)}
\bigl(y(t)-r\operatorname{Osc}(f)^2\bigr).
\]
Solving this scalar differential inequality proves
\cref{eq-controlled-weak-poincare-decay}.  The same calculation with
\(y(t)=\|S_tf\|_2^2\), using \(L^1\)-contraction, proves
\cref{eq-controlled-super-poincare-smoothing}.

Under Nash, \(y(t)=\|S_tf\|_2^2\) satisfies
\[
y'(t)\le-\frac2{C_{\rm N}}
y(t)^{1+2/d}\|f\|_1^{-4/d}.
\]
Integration proves \cref{eq-controlled-nash-smoothing}; the semigroup
identity and self-adjoint duality give the \(L^1\)-to-\(L^\infty\) estimate.
The Sobolev implication follows from H\"older interpolation.

The variational definition of \(\lambda_\Phi^{\rm D}(A)\) and the spectral theorem
give \cref{eq-controlled-killed-spectral-decay}; substitution gives the
Faber--Krahn estimate.  The evolving-level-set argument for the lazy
uniformized kernel gives
\cref{eq-controlled-spectral-profile-mixing}; the factor \(q_\Phi\) is
already contained in the generator-rate identity
\(\Lambda_\Phi^{\rm mix}=q_\Phi\Lambda_K^{\rm mix}\)
\citep{GoelMontenegroTetali2006}.

The logarithmic Sobolev and modified logarithmic Sobolev conclusions follow
by differentiating the corresponding norms and entropy along \(S_t\).
The \(\Gamma_2\) interpolation
\(s\mapsto S_s\Gamma(S_{t-s}f)\) proves
\cref{eq-controlled-bakry-emery-contraction}.  Kantorovich duality followed
by the declared \(T_1\) or \(T_2\) inequality proves
\cref{eq-controlled-t1-shift,eq-controlled-t2-shift}.  The capacity statement
is the Dirichlet principle.  Form comparison follows directly from the
variational definitions.  Finally, Gr\"onwall's inequality applied to
\(\mathscr H(P_tf)\), followed by the two norm comparisons, proves
\cref{eq-controlled-hypocoercive-decay}.
The Nash and weak-Poincar\'e implications use the normalizations of
\citet{CarlenKusuokaStroock1987} and \citet{RocknerWang2001}.
The transport and hypocoercive entries use, respectively,
\citet{BobkovGotze1999,OttoVillani2000} and \citet{Villani2009}.

\end{proof}

\begin{definition}[Typed functional-output slots]
\label{def-controlled-functional-output-slots}

The profile outputs are assigned to the following slots:
\[
\begin{aligned}
\mathscr S_{\rm FI}=\{\,&L^2\text{-decay},\
\operatorname{Ent}\text{-decay},\ 1\!\to\!2\text{-smoothing},\\
&1\!\to\!\infty\text{-smoothing},\ \text{killed survival},\
\text{gradient contraction},\\
&\text{distribution shift},\ \text{target/failure access},\
\text{hypocoercive decay}\}.
\end{aligned}
\]
A candidate bound is eligible for a slot only after a represented theorem
has converted its constants to that slot's normalization, time scale, carrier,
domain, and comparison loss.  A logical implication is represented by a
directed edge from its source certificate to the induced bound.  The induced
bound retains the source certificate's provenance and is not a second
independent contribution.

\end{definition}

\begin{theorem}[Slotwise selection and no double counting]
\label{thm-controlled-functional-profile-selection}

Fix one certified dynamical geometry and one output slot
\(s\in\mathscr S_{\rm FI}\).  Let \(\mathcal J_s\) be the finite represented
set of audited bounds converted to the units of \(s\).  If every
\(j\in\mathcal J_s\) proves \(U_s\le C_{s,j}\), then
\begin{equation}
\label{eq-controlled-functional-slot-minimum}
U_s\le C_s^{\rm FI}:=\min_{j\in\mathcal J_s}C_{s,j}.
\end{equation}
For an empirically selected finite family, the right side receives exactly
the simultaneous selection penalty supplied by the relevant uniform or
sample-split theorem.  Deterministically certified constants receive no
selection penalty.

If one certificate implies several inequalities, each downstream slot may
use the induced bound, but a propagation theorem may charge the originating
error only once.  Bounds in different slots are added only after a represented
propagation theorem converts them to one common risk, return, hitting, or
deployment quantity.  Thus Poincar\'e, LSI, MLSI, curvature, transport,
capacity, and hypocoercivity are neither summed as dimensionally incomparable
numbers nor counted repeatedly along an implication chain.

\end{theorem}

\begin{proof}

Every member of \(\mathcal J_s\) is an upper bound for the same scalar in the
same units, so their minimum is an upper bound.  Simultaneous empirical
validity is exactly the event supplied by the cited selection theorem.  The
remaining statements follow from provenance preservation under represented
composition and no-creation,
\cref{prop-derived-channels-create-no-structure,thm-defect-profile-selection}:
an induced inequality is a deterministic descendant of its source record,
and cross-slot addition is valid only after a common-output comparison has
been proved.

\end{proof}


\begin{figure}[tbp]
\centering
\begin{tikzpicture}[fw diagram,node distance=6mm and 8mm]
  \node[fw geom,text width=2.3cm] (form) {represented\\Dirichlet form};
  \node[fw box,text width=2.2cm,above right=4mm and 9mm of form] (local)
    {Cheeger\\Poincar\'e};
  \node[fw box,text width=2.2cm,right=of form] (weak)
    {weak/super\\Poincar\'e};
  \node[fw box,text width=2.2cm,below right=4mm and 9mm of form] (smooth)
    {Nash, Sobolev\\spectral profile};
  \node[fw use,text width=2.2cm,right=13mm of local] (l2) {\(L^2\) decay};
  \node[fw use,text width=2.2cm,right=13mm of weak] (reg) {smoothing\\killed survival};
  \node[fw use,text width=2.2cm,right=13mm of smooth] (ent) {entropy\\gradient};
  \node[fw provenance,text width=2.3cm,above=6mm of l2] (curv)
    {LSI/MLSI\\curvature};
  \node[fw provenance,text width=2.3cm,below=6mm of ent] (cap)
    {transport\\capacity};
  \node[fw terminal,text width=2.4cm,right=12mm of reg] (slots)
    {typed output slots\\take valid minima};
  \draw[fw arrow] (form) -- (local);
  \draw[fw arrow] (form) -- (weak);
  \draw[fw arrow] (form) -- (smooth);
  \draw[fw arrow] (local) -- (l2);
  \draw[fw arrow] (weak) -- (reg);
  \draw[fw arrow] (smooth) -- (reg);
  \draw[fw arrow] (curv) -- (l2);
  \draw[fw arrow] (curv) -- (ent);
  \draw[fw arrow] (cap) -- (ent);
  \draw[fw arrow] (cap) -- (reg);
  \draw[fw arrow] (l2) -- (slots);
  \draw[fw arrow] (reg) -- (slots);
  \draw[fw arrow] (ent) -- (slots);
\end{tikzpicture}
\caption{The audited functional-inequality profile converts certified
constants into typed dynamical outputs.  Bounds compete by a minimum only
after conversion to the same slot; implication edges preserve provenance and
therefore introduce no duplicate charge.}
\label{fig-controlled-functional-inequality-profile}
\end{figure}

\begin{figure}[tbp]
\centering
\begin{tikzpicture}[fw diagram,node distance=7mm and 8mm]
  \node[fw box,text width=2.8cm] (selected) {selected generator\\and invariant law};
  \node[fw geom,text width=2.6cm,right=of selected] (carrier)
    {reversible carrier\\Dirichlet form};
  \node[fw use,text width=2.5cm,above right=5mm and 8mm of carrier] (rates)
    {gap, conductance\\MLSI, curvature};
  \node[fw provenance,text width=2.5cm,right=of carrier] (current)
    {authorized Hodge\\current};
  \node[fw terminal,text width=2.5cm,below right=5mm and 8mm of carrier] (target)
    {target reach\\regularity};
  \node[fw source,text width=2.8cm,right=18mm of current] (cert)
    {qualified dynamical\\certificate};
  \draw[fw arrow] (selected) -- (carrier);
  \draw[fw arrow] (carrier) -- (rates);
  \draw[fw arrow] (carrier) -- (current);
  \draw[fw arrow] (carrier) -- (target);
  \draw[fw arrow] (rates) -- (cert);
  \draw[fw arrow] (current) -- (cert);
  \draw[fw arrow] (target) -- (cert);
\end{tikzpicture}
\caption{A controlled dynamical certificate couples a reversible carrier to
the actual authorized current and target-visible regularity.  Every constant
is evaluated on the same represented record and charged at its complete
saturated cost.}
\label{fig-certified-dynamical-geometry}
\end{figure}

\section{Committors, Bellman Recursions, and Reach--Avoid Geometry}
\label{sec-controlled-committor-bellman}

\subsection{Controlled committors and Bellman identities}
\label{subsec-controlled-committor-bellman-identities}

\begin{definition}[Controlled reach--avoid committors]
\label{def-controlled-reach-avoid-committor}

Fix disjoint represented terminal sectors \(T^+,T^-\subseteq X\), put
\(N=X\setminus(T^+\cup T^-)\), and let \(P^\pi\) be a selected policy
kernel.  The horizon-\(h\) reach--avoid committor is
\[
u_h^\pi(x)
=\Pr_x^\pi\{\tau_{T^+}\le h,\ \tau_{T^+}<\tau_{T^-}\}.
\]
For \(0<\beta<1\), the discounted committor is
\[
u_\beta^\pi(x)
=\mathbb E_x^\pi\!\left[
\beta^{\tau_{T^+}}
\mathbf1_{\{\tau_{T^+}<\tau_{T^-}\}}
\right],
\]
with the convention \(\beta^\infty=0\).  Both functions have boundary values
one on \(T^+\) and zero on \(T^-\).  Their semantic existence does not assign
them saturated representation cost; derivability is governed by
\cref{def-no-free-committor}.

For a bounded reward \(r\) and discount \(0<\beta<1\), define
\[
(\mathcal T^\pi v)(x)
=r^\pi(x)+\beta(P^\pi v)(x),
\qquad
(\mathcal T v)(x)
=\max_{a\in\mathcal A(x)}
\{r(x,a)+\beta(P^av)(x)\}.
\]
All maxima are finite and represented action evaluation and comparison are
charged to the selected-policy derivation.

\end{definition}

\begin{theorem}[Exact controlled Bellman and committor identities]
\label{thm-controlled-committor-bellman}

The finite-horizon committors satisfy
\begin{align}
u_0^\pi&=\mathbf1_{T^+},\label{eq-controlled-committor-initial}\\
u_{h+1}^\pi
&=\mathbf1_{T^+}+\mathbf1_NP^\pi u_h^\pi.\label{eq-controlled-committor-recursion}
\end{align}
The discounted committor is the unique bounded solution of
\begin{equation}
\label{eq-controlled-discounted-committor-bellman}
u_\beta^\pi
=\mathbf1_{T^+}+\beta\mathbf1_NP^\pi u_\beta^\pi.
\end{equation}
The optimal finite-horizon and discounted committors are obtained by replacing
\(P^\pi\) on \(N\) by \(\max_{a\in\mathcal A(x)}P^a\) pointwise.  For a
bounded represented reward, \(\mathcal T^\pi\) and \(\mathcal T\) are
\(\beta\)-contractions in \(\|\cdot\|_\infty\), and their unique fixed
points are \(V^\pi\) and \(V^*\), respectively.

For a finite legal-action set and \(\tau>0\), the soft action value obeys the
exact Gibbs identity
\begin{equation}
\label{eq-controlled-soft-bellman-gibbs}
\tau\log\sum_{a\in\mathcal A(x)}e^{q_a/\tau}
=\max_{p\in\Delta(\mathcal A(x))}
\left\{\sum_ap_aq_a+\tau\mathcal H(p)\right\},
\end{equation}
whose maximizer is the represented Boltzmann kernel.  Its construction and
precision are part of the policy cost.

\end{theorem}

\begin{proof}

Condition on the first transition.  A state already in \(T^+\) contributes
one, a state in \(T^-\) contributes zero, and a state in \(N\) contributes
the expectation of the remaining-horizon committor.  This proves
\cref{eq-controlled-committor-initial,eq-controlled-committor-recursion}.
The same conditioning with one factor \(\beta\) proves
\cref{eq-controlled-discounted-committor-bellman}.  The right-hand side is a
\(\beta\)-contraction on bounded functions with fixed boundary values, so
Banach's fixed-point theorem gives existence and uniqueness.  Maximization at
each state gives the dynamic-programming recursion because the finite-horizon
continuation value depends on the past only through the clocked state.

Stochastic kernels are contractions in the supremum norm; taking a finite
pointwise maximum preserves the same Lipschitz constant.  Hence the reward
operators are \(\beta\)-contractions and their fixed points equal the usual
discounted returns by iteration.  Finally, Lagrange multipliers for the
simplex constraint give \(p_a\propto e^{q_a/\tau}\); substitution proves
\cref{eq-controlled-soft-bellman-gibbs}
\citep{Bellman1957,Puterman1994,BertsekasTsitsiklis1996}.

\end{proof}

\begin{theorem}[Temporal admissibility of controlled values]
\label{thm-controlled-no-free-value}

Let \(u\) denote any committor, value function, advantage, policy table,
first-hit partition, or completed reach--avoid record of a selected controlled
system.  It belongs to a prospective source profile at budget \(B\) exactly
through a temporally admissible represented derivation of the same denotation
whose complete cost is at most \(B\).  A Bellman recursion evaluated for
\(h\) steps is charged for those steps and retained values.  An exact solve is
charged for its represented solver, arithmetic, precision, and output.  A
terminal record computed from future verifier flags remains an accounting
record under \cref{cor-first-hit-accounting-not-source}.

Consequently, the identities in
\cref{thm-controlled-committor-bellman} create no prospective target source.
When a represented value is available before a transition, its quotient,
Geom/Caus, Lift, and boundary provenance is classified by
\cref{thm-controlled-policy-channel-decomposition} at its complete cost.

\end{theorem}

\begin{proof}

Apply \cref{def-no-free-committor,thm-no-truncated-neumann-bypass} to the
policy-induced kernel on the clocked carrier.  Availability before the
transition and exclusion of later verifier flags are precisely
\cref{def-pre-hit-temporally-admissible-source}.  The channel assignment then
follows from \cref{thm-controlled-policy-channel-decomposition}.  These
alternatives are exhaustive because the completed transcript is a valid Lift
and its post-exhaustion remainder is \(J_{\rm res}\).

\end{proof}

\subsection{No uncharged prospective value amplification}
\label{subsec-controlled-no-uncharged-value-amplification}

\begin{theorem}[No uncharged prospective value amplification]
\label{thm-no-uncharged-prospective-value-amplification}

Let \(R\) be a family-uniform, temporally admissible pre-hit record of
complete saturated cost at most \(B\).  Let \(a_R(I)\in[0,1]\) be a
prospective selector, reach--avoid, or arrival value in a declared output
normalization.  Let \(\widehat B\) be a class-preserving common majorant,
within the already declared saturated budget structure, for evaluator
normalization, the clocked policy embedding, the complete active-source
backward slice, and every comparison used below.

A structural upper bound for \(a_R\) is admitted only through at least one of
the following two forms.

\begin{enumerate}[label=\textup{(\roman*)}]
\item There are an earlier dynamically qualified master certificate
      \(\mathscr C_R\), a represented constant \(K_R\ge0\), and a
      represented error \(\delta_R\ge0\), all in the complete record, such
      that pointwise
      \begin{equation}
      \label{eq-inherited-prospective-value-comparison}
      a_R(I)\le K_RV_{\mathscr C_R,I}+\delta_R(I).
      \end{equation}
      The value is then comparison-dominated by that earlier source.

\item The complete pre-hit record \(R\), including its selected kernel,
      potential, current, authority data, target interface, and output
      converter, contains a dynamically qualified master certificate
      \(\mathscr C_R^{\rm dir}\) and a represented same-normalization
      comparison
      \begin{equation}
      \label{eq-direct-prospective-value-comparison}
      a_R(I)\le K_R^{\rm dir}V_{\mathscr C_R^{\rm dir},I}
                 +\delta_R^{\rm dir}(I).
      \end{equation}
      The complete record is then itself the charged prospective source and
      is assigned, at cost at most \(\widehat B\), to exactly one existing
      ambient, exact-quotient, reversible-compensated, or
      general-resolvent master-profile cell.
\end{enumerate}

If neither represented comparison exists, the framework makes no numerical
transfer from structural visibility to \(a_R\).  The record remains fully
classified by the channel decomposition, but it cannot be counted as a
comparison-dominated prospective structural value.  A later terminal or
first-hit record is accounting only; a gradient-orthogonal contribution is
\(\Res\); an undeclared input is presentation enrichment or advice; and an
over-budget derivation is absent from the budget slice.

For any collection \(\mathscr R(B)\) for which every positive value has one
of the two displayed comparisons, assign a record to the inherited class
when \textup{(i)} is present and otherwise to the direct class
\textup{(ii)}.  These two classes are disjoint and cover
\(\mathscr R(B)\).  Define no new source channel and use the existing
master-profile cells containing the corresponding certificates.
Then, pointwise,
\begin{equation}
\label{eq-no-uncharged-value-pointwise-envelope}
\sup_{R\in\mathscr R(B)}a_R(I)
\le
\max\left\{
\sup_{R\in\mathscr R_{\rm inh}(B)}
  \bigl(K_RV_{\mathscr C_R,I}+\delta_R(I)\bigr),
\sup_{R\in\mathscr R_{\rm dir}(B)}
  \bigl(K_R^{\rm dir}V_{\mathscr C_R^{\rm dir},I}
       +\delta_R^{\rm dir}(I)\bigr)
\right\}.
\end{equation}
The same inequality holds after applying a monotone family functional and
then taking the saturated supremum.  In particular, ancestry equality,
zero incremental constructor charge, Bellman recursion, valid lifting, or
postprocessing alone never supplies either displayed numerical comparison.

\end{theorem}

\begin{proof}

Temporal admissibility places the complete pre-hit derivation of \(R\) in
the saturated closure.  Normalize its selected transition by
\cref{thm-controlled-policy-channel-decomposition} and retain the active
backward slice, including every instance-dependent coefficient and output
converter.  If \eqref{eq-inherited-prospective-value-comparison} occurs in
that record, it is precisely a represented same-normalization transfer, so
the first alternative holds.

Otherwise a structural explanation must use the complete record producing
the value.  If that record contains the qualification and comparison in
\eqref{eq-direct-prospective-value-comparison}, apply
\cref{thm-affordable-geom-abs-normal-form,thm-no-unclassified-geom-source-after-quotient-assignment}.
They assign \(\mathscr C_R^{\rm dir}\) to exactly one existing master cell
while preserving its full derivation and cost.  This proves the second
alternative and shows that it introduces no additional channel.

If neither comparison is represented, structural ancestry provides no
inequality in the output normalization.  This is exactly the scope
restriction in
\cref{cor-contextual-structural-charge-transfer,thm-prospective-selector-structural-charge}:
contextual or accounting charge transfers to a prospective source value
only through the represented comparison.  The remaining assignments follow
from
\cref{cor-first-hit-accounting-not-source,thm-exhaustive-residual-normal-form,thm-oracle-admissibility-presentation-separation,def-budgeted-derivability-closure}.

For a collection covered by the two alternatives, split it into the
inherited and direct subcollections, apply the corresponding pointwise
inequality, and take the supremum.  A supremum over their union is the
maximum of the two suprema, proving
\eqref{eq-no-uncharged-value-pointwise-envelope}.  Monotonicity gives the
family-functional statement, and saturation preserves the same common
ceiling.  Finally,
\cref{prop-derived-channels-create-no-structure,thm-controlled-no-free-value}
shows that the listed derived operations preserve source ancestry but do not
create a numerical comparison.

\end{proof}


\begin{figure}[tbp]
\centering
\resizebox{\linewidth}{!}{%
\begin{tikzpicture}[fw diagram,node distance=8mm and 10mm]
  \node[fw source,text width=2.7cm] (value)
    {temporally admissible\\prospective value};
  \node[fw provenance,text width=2.7cm,above right=8mm and 12mm of value]
    (inh) {represented comparison\\to earlier source};
  \node[fw geom,text width=2.7cm,below right=8mm and 12mm of value]
    (dir) {complete record is\\qualified source};
  \node[fw use,text width=2.8cm,right=14mm of inh] (bound)
    {same-normalization\\numerical bound};
  \node[fw provenance,text width=2.8cm,right=14mm of dir] (cell)
    {existing master-profile\\cell at full cost};
  \node[fw terminal,text width=2.8cm,right=16mm of value] (out)
    {charged prospective\\output bound};
  \node[fw residual,text width=3.2cm,below=18mm of value] (none)
    {no represented comparison:\\no structural-value transfer};
  \draw[fw arrow] (value) -- (inh);
  \draw[fw arrow] (value) -- (dir);
  \draw[fw arrow] (inh) -- (bound);
  \draw[fw arrow] (dir) -- (cell);
  \draw[fw arrow] (bound) -- (out);
  \draw[fw arrow] (cell) -- (out);
  \draw[fw dashed arrow] (value) -- (none);
\end{tikzpicture}%
}
\caption{Prospective value has no uncharged intermediate location.  It is
either comparison-dominated by an earlier qualified source or charged as a
complete directly qualified source in an existing master-profile cell.  In
the absence of either represented comparison, provenance alone supplies no
numerical structural-value bound.}
\label{fig-no-uncharged-prospective-value}
\end{figure}

\subsection{Condenser capacity and reach--avoid bounds}
\label{subsec-controlled-condenser-capacity-reach-avoid}

\begin{definition}[Represented condenser capacity]
\label{def-controlled-condenser-capacity}

Let \(\Phi\) be a represented reversible carrier and let \(S,T\subseteq X\)
be disjoint represented sets.  Its condenser capacity is
\begin{equation}
\label{eq-controlled-condenser-capacity}
\operatorname{Cap}_\Phi(S,T)
=\inf\{\mathcal E_\Phi(f,f):f|_S=1,\ f|_T=0\}.
\end{equation}
The infimum is over the finite represented function space.  A capacity value
is a certified coordinate only when the boundary sets, minimizer or solver,
energy evaluation, and precision are in the saturated record.  For
reach--avoid, the relevant condenser is \((T^+,T^-)\), optionally after
killing outside a represented safe carrier.

\end{definition}

\begin{theorem}[Dirichlet reach--avoid principle and drift-time bound]
\label{thm-controlled-reach-avoid-dirichlet-drift}

Let the selected policy kernel be reversible with respect to its represented
invariant law and let \(h\) be its undiscounted target-before-failure
committor.  Then \(h=1\) on \(T^+\), \(h=0\) on \(T^-\), it is harmonic on
\(N\), and
\begin{equation}
\label{eq-controlled-dirichlet-principle}
\mathcal E_\Phi(h,h)
=\operatorname{Cap}_\Phi(T^+,T^-).
\end{equation}
For every represented initial law,
\begin{equation}
\label{eq-controlled-reach-avoid-exact}
\Pr_{\mu_0}\{\tau_{T^+}<\tau_{T^-}\}
=\langle\mu_0,h\rangle.
\end{equation}

Independently of reversibility, suppose a represented nonnegative potential
\(D\) vanishes on \(T^+\), and the same selected generator satisfies
\(-LD\ge a>0\) before \(T^+\).  Then
\begin{equation}
\label{eq-controlled-drift-time-bound}
\mathbb E_{\mu_0}\tau_{T^+}
\le\frac{\mathbb E_{\mu_0}D}{a}.
\end{equation}
For a represented quotient potential with defect bounded by \(\delta<a\),
the denominator becomes \(a-\delta\) as in
\cref{thm-caus-drift-survives-defect}.

\end{theorem}

\begin{proof}

First-step conditioning gives the boundary and harmonic equations for \(h\).
For any competitor \(f=h+g\) with \(g=0\) on \(T^+\cup T^-\), discrete
integration by parts and harmonicity give
\(\mathcal E_\Phi(h,g)=0\).  Hence
\[
\mathcal E_\Phi(f,f)
=\mathcal E_\Phi(h,h)+\mathcal E_\Phi(g,g)
\ge\mathcal E_\Phi(h,h),
\]
which proves \cref{eq-controlled-dirichlet-principle}.  Equation
\cref{eq-controlled-reach-avoid-exact} is the definition of the committor
integrated against \(\mu_0\).  The hitting-time estimate is Dynkin's formula
applied to \(D(X_{t\wedge\tau_{T^+}})\), followed by monotone convergence;
the quotient-defect refinement is exactly
\cref{thm-caus-drift-survives-defect}
\citep{FukushimaOshimaTakeda2011}.

\end{proof}

\begin{theorem}[Invariant boundary capacity, occupation, and global-drift obstruction]
\label{thm-controlled-target-capacity-drift-envelope}

Let \(P\) be a represented Markov kernel on a finite carrier \(X\), let
\(\mu\) be a represented full-support invariant probability measure, and let
\(\Phi_P\) be the native symmetric geometry obtained from \(L_P=P-I\) as in
\cref{def-certified-dynamical-geometry}.  For \(A,B\subseteq X\), write
\[
Q_P(A,B)=\sum_{x\in A}\sum_{z\in B}\mu(x)P(x,z).
\]
Fix a nonempty proper represented target \(T\subsetneq X\), and let
\(\tau_T=\inf\{t\ge0:X_t\in T\}\).  The following statements hold for this
same selected kernel and target record.

\begin{enumerate}[label=\textup{(\roman*)}]
\item The target-boundary capacity is the invariant entrance flux:
      \begin{equation}
      \label{eq-controlled-target-boundary-capacity-flux}
      \operatorname{Cap}_{\Phi_P}(T,T^c)
      =\mathcal E_{\Phi_P}(\mathbf1_T,\mathbf1_T)
      =Q_P(T^c,T)=Q_P(T,T^c).
      \end{equation}
      More generally, for every represented \(F\subseteq T^c\),
      \begin{equation}
      \label{eq-controlled-target-condenser-flux-upper}
      \operatorname{Cap}_{\Phi_P}(T,F)
      \le Q_P(T^c,T).
      \end{equation}

\item If \(D\ge0\), \(D|_T=0\), and
      \((I-P)D\ge a>0\) on \(T^c\), then
      \begin{equation}
      \label{eq-controlled-invariant-global-drift-obstruction}
      \frac{a}{\operatorname{Osc}(D)}
      \le \frac{\mu(T)}{1-\mu(T)}.
      \end{equation}
      Thus an invariant full-carrier chain can have uniform normalized drift
      toward a small target only at the target-mass scale.

\item Let \(\mu_0\ll\mu\) and put
      \(R_0=\|d\mu_0/d\mu\|_\infty\).  For every integer \(H\ge0\),
      \begin{equation}
      \label{eq-controlled-capacity-finite-horizon-hit}
      \Pr_{\mu_0}\{\tau_T\le H\}
      \le R_0\bigl[\mu(T)+H
             \operatorname{Cap}_{\Phi_P}(T,T^c)\bigr].
      \end{equation}
      If \(0\le t_0\le H\) and the time-\(t_0\) law has density \(f_{t_0}\)
      with respect to \(\mu\), then
      \begin{equation}
      \label{eq-controlled-capacity-burn-in-hit}
      \Pr_{\mu_0}\{\tau_T\le H\}
      \le \Pr_{\mu_0}\{\tau_T\le t_0\}
       +\|f_{t_0}\|_\infty
        \bigl[\mu(T)+(H-t_0)
        \operatorname{Cap}_{\Phi_P}(T,T^c)\bigr].
      \end{equation}

\item Suppose in addition that \(P\) is lazy and self-adjoint in
      \(L^2(X,\mu)\), with spectral gap \(\gamma_P>0\).  If
      \(f_0=d\mu_0/d\mu\), then
      \begin{align}
      \Pr_{\mu_0}\{\tau_T\le H\}
      &\le \sum_{t=0}^{H}\mu_0P^t(T)\notag\\
      &\le (H+1)\mu(T)
      +\sqrt{\mu(T)(1-\mu(T))}\,
       \|f_0-1\|_{2,\mu}
       \frac{1-(1-\gamma_P)^{H+1}}{\gamma_P}.
      \label{eq-controlled-gap-target-occupation}
      \end{align}
      In particular, a relative-uniform mixing certificate
      \(\|f_{t_0}-1\|_\infty\le\varepsilon\), including one obtained from
      \cref{eq-controlled-spectral-profile-mixing}, permits
      \(\|f_{t_0}\|_\infty\le1+\varepsilon\) in
      \cref{eq-controlled-capacity-burn-in-hit}.

\item Let \(F\subseteq T^c\) be represented, and let a represented comparison
      geometry \(\Psi\) have the same boundary sets and satisfy
      \(\mathcal E_{\Phi_P}(f,f)\le
        b\mathcal E_{\Psi}(f,f)\) for every represented \(f\) in the
      condenser domain.  Then
      \begin{equation}
      \label{eq-controlled-capacity-form-comparison-upper}
      \operatorname{Cap}_{\Phi_P}(T,F)
      \le b\operatorname{Cap}_{\Psi}(T,F).
      \end{equation}
\end{enumerate}

Every quantity in these inequalities belongs to the selected native or
comparison geometry already classified by
\cref{thm-controlled-policy-channel-decomposition}; the theorem introduces
no additional structural channel.

\end{theorem}

\begin{proof}

Invariance gives equality of the two stationary cut flows:
\[
Q_P(T,T^c)
=\sum_{x\in T}\mu(x)-Q_P(T,T)
=\sum_{z\in T}\mu(z)-Q_P(T,T)
=Q_P(T^c,T).
\]
For the native symmetric conductance,
\(\mathcal E_{\Phi_P}(\mathbf1_T,\mathbf1_T)\) is their common value.
The boundary conditions for the condenser \((T,T^c)\) force its only
competitor to be \(\mathbf1_T\), proving
\cref{eq-controlled-target-boundary-capacity-flux}.  The same indicator is
an admissible competitor for \((T,F)\), which proves
\cref{eq-controlled-target-condenser-flux-upper}.

Because \(0\le D\le\operatorname{Osc}(D)\) after the normalization
\(D|_T=0\), one has
\((I-P)D=-PD\ge-\operatorname{Osc}(D)\) on \(T\).  Invariance and the drift
hypothesis therefore give
\[
0=\mu((I-P)D)
\ge a(1-\mu(T))-\mu(T)\operatorname{Osc}(D).
\]
Rearrangement proves
\cref{eq-controlled-invariant-global-drift-obstruction}.

Under the stationary path law, a hit by time \(H\) occurs either at time zero
or after one of the \(H\) transitions from \(T^c\) into \(T\).  The union
bound and stationarity give
\[
\Pr_\mu\{\tau_T\le H\}
\le\mu(T)+H Q_P(T^c,T).
\]
The Radon--Nikodym derivative of the full path law started from \(\mu_0\)
with respect to the stationary path law is \(f_0(X_0)\), and hence is at
most \(R_0\).  This proves
\cref{eq-controlled-capacity-finite-horizon-hit}.  Apply the same argument
to the Markov chain started from its time-\(t_0\) law and add the event of a
hit through time \(t_0\) to obtain
\cref{eq-controlled-capacity-burn-in-hit}.

For the lazy self-adjoint kernel, put \(g_T=\mathbf1_T-\mu(T)\).  The gap
and Cauchy--Schwarz imply
\[
\mu_0P^t(T)
=\mu(T)+\langle P^t(f_0-1),g_T\rangle_\mu
\le\mu(T)+(1-\gamma_P)^t
 \|f_0-1\|_{2,\mu}\sqrt{\mu(T)(1-\mu(T))}.
\]
The event \(\{\tau_T\le H\}\) is contained in the union of the occupation
events \(\{X_t\in T\}\).  Summing the last display and evaluating the
geometric series proves
\cref{eq-controlled-gap-target-occupation}.  The relative-uniform statement
is immediate from \(f_{t_0}\le1+\varepsilon\).

Finally, let \(h_\Psi\) minimize the comparison condenser energy.  It is an
admissible competitor for the native condenser, so
\[
\operatorname{Cap}_{\Phi_P}(T,F)
\le\mathcal E_{\Phi_P}(h_\Psi,h_\Psi)
\le b\mathcal E_\Psi(h_\Psi,h_\Psi)
=b\operatorname{Cap}_\Psi(T,F).
\]
This proves \cref{eq-controlled-capacity-form-comparison-upper} and the
theorem.

\end{proof}

\subsection{Invariant-mass and finite-horizon transfer}
\label{subsec-controlled-invariant-mass-horizon-transfer}

\begin{theorem}[Invariant mass conservation sharpens global target drift]
\label{thm-controlled-invariant-mass-capacity-drift}

Let \(P\), \(\mu\), \(\Phi_P\), and the nonempty proper represented target
\(T\subsetneq X\) satisfy the hypotheses of
\cref{thm-controlled-target-capacity-drift-envelope}.  Let \(D\ge0\) be a
nonconstant represented potential with \(D|_T=0\), and put

\begin{equation}
\label{eq-controlled-invariant-off-target-drift-margin}
a_T(P,D)
=
\inf_{x\in T^c}(I-P)D(x).
\end{equation}

Then the stationary mass balance is the exact identity

\begin{equation}
\label{eq-controlled-invariant-drift-flux-identity}
\int_{T^c}(I-P)D\,d\mu
=
\int_T(P-I)D\,d\mu
=
\sum_{x\in T}\sum_{z\in T^c}
 \mu(x)P(x,z)D(z).
\end{equation}

Consequently,

\begin{equation}
\label{eq-controlled-invariant-drift-capacity-ceiling}
\frac{\bigl(a_T(P,D)\bigr)_+}{\operatorname{Osc}(D)}
\le
\frac{Q_P(T,T^c)}{1-\mu(T)}
=
\frac{\operatorname{Cap}_{\Phi_P}(T,T^c)}{1-\mu(T)}
\le
\frac{\mu(T)}{1-\mu(T)}.
\end{equation}

Thus a positive drift margin on the whole complement of \(T\) is paid for
by the stationary boundary flux of the same selected kernel.  The statement
uses the represented invariant law of that kernel; an invariant law or form
from a different record enters through a represented comparison such as
\cref{thm-controlled-density-capacity-hit-transfer} below.

\end{theorem}

\begin{proof}

Invariance gives \(\int_X(I-P)D\,d\mu=0\).  Since \(D=0\) on \(T\),

\[
-\int_T(I-P)D\,d\mu
=\int_TPD\,d\mu
=\sum_{x\in T}\sum_{z\in T^c}\mu(x)P(x,z)D(z),
\]

which proves \cref{eq-controlled-invariant-drift-flux-identity}.  The
definition of \(a_T(P,D)\) gives

\[
\bigl(a_T(P,D)\bigr)_+\mu(T^c)
\le
\int_{T^c}(I-P)D\,d\mu.
\]

Every term in the last sum of
\cref{eq-controlled-invariant-drift-flux-identity} has
\(0\le D(z)\le\operatorname{Osc}(D)\).  Hence

\[
\int_{T^c}(I-P)D\,d\mu
\le
\operatorname{Osc}(D)Q_P(T,T^c).
\]

Divide by
\(\operatorname{Osc}(D)\mu(T^c)>0\), and use
\cref{eq-controlled-target-boundary-capacity-flux}.  Finally,
\(Q_P(T,T^c)\le\mu(T)\), which proves the last comparison in
\cref{eq-controlled-invariant-drift-capacity-ceiling}.

\end{proof}

\begin{theorem}[Density, conductance, capacity, and finite-horizon transfer]
\label{thm-controlled-density-capacity-hit-transfer}

Let \((P,\mu,\Phi_P)\) and
\((\widetilde P,\widetilde\mu,\Phi_{\widetilde P})\) be two complete
represented reversible Markov records on the same finite carrier \(X\).
For \(x\ne z\), write

\[
c_P(x,z)=\mu(x)P(x,z),
\qquad
c_{\widetilde P}(x,z)
=\widetilde\mu(x)\widetilde P(x,z).
\]

Assume the same record contains constants
\(0<r_-\le r_+<\infty\) and
\(0<\kappa_-\le\kappa_+<\infty\) such that, pointwise,

\begin{align}
\label{eq-controlled-density-comparison}
r_-\mu(x)
&\le\widetilde\mu(x)\le r_+\mu(x),\qquad x\in X,\\
\label{eq-controlled-conductance-comparison}
\kappa_-c_P(x,z)
&\le c_{\widetilde P}(x,z)
\le\kappa_+c_P(x,z),\qquad x\ne z.
\end{align}

Then, for every represented set \(A\subseteq X\) and every pair of disjoint
represented plates \(S,F\subseteq X\),

\begin{align}
\label{eq-controlled-density-target-mass-transfer}
r_-\mu(A)
&\le\widetilde\mu(A)\le r_+\mu(A),\\
\label{eq-controlled-density-form-transfer}
\kappa_-\mathcal E_{\Phi_P}(f,f)
&\le\mathcal E_{\Phi_{\widetilde P}}(f,f)
\le\kappa_+\mathcal E_{\Phi_P}(f,f),\\
\label{eq-controlled-density-capacity-transfer}
\kappa_-\operatorname{Cap}_{\Phi_P}(S,F)
&\le\operatorname{Cap}_{\Phi_{\widetilde P}}(S,F)
\le\kappa_+\operatorname{Cap}_{\Phi_P}(S,F).
\end{align}

Fix a nonempty proper represented target \(T\), an integer \(H\ge0\), and
an initial law \(\rho_0\ll\widetilde\mu\).  For
\(1<p\le\infty\), let \(q=p/(p-1)\), with \(q=1\) for \(p=\infty\), and put
\(f_0=d\rho_0/d\widetilde\mu\).  The finite-horizon target probability of
the selected kernel \(\widetilde P\) satisfies

\begin{equation}
\label{eq-controlled-density-lp-hit-transfer}
\Pr_{\rho_0}^{\widetilde P}\{\tau_T\le H\}
\le
\min\!\left\{
1,
\|f_0\|_{L^p(\widetilde\mu)}
\left[
\min\!\left\{1,
r_+\mu(T)+H\kappa_+
\operatorname{Cap}_{\Phi_P}(T,T^c)
\right\}
\right]^{1/q}
\right\}.
\end{equation}

In particular, if \(\rho_0\le R_{\widetilde\mu}\widetilde\mu\), then

\begin{equation}
\label{eq-controlled-density-linfty-hit-transfer}
\Pr_{\rho_0}^{\widetilde P}\{\tau_T\le H\}
\le
\min\!\left\{1,
R_{\widetilde\mu}
\left[
r_+\mu(T)+H\kappa_+
\operatorname{Cap}_{\Phi_P}(T,T^c)
\right]
\right\}.
\end{equation}

If instead the available initialization certificate is
\(\rho_0\le R_\mu\mu\), then

\begin{equation}
\label{eq-controlled-base-density-hit-transfer}
\Pr_{\rho_0}^{\widetilde P}\{\tau_T\le H\}
\le
\min\!\left\{1,
\frac{R_\mu}{r_-}
\left[
r_+\mu(T)+H\kappa_+
\operatorname{Cap}_{\Phi_P}(T,T^c)
\right]
\right\}.
\end{equation}

The same comparison applies after a represented burn-in: if the time-
\(t_0\) law obeys \(\rho_{t_0}\le R_\mu^{(t_0)}\mu\), then the right side of
\cref{eq-controlled-base-density-hit-transfer}, with \(H\) replaced by
\(H-t_0\) and \(R_\mu\) by \(R_\mu^{(t_0)}\), is added to
\(\Pr_{\rho_0}^{\widetilde P}\{\tau_T\le t_0\}\).

All four constants are comparisons between the two complete represented
records.  Their derivation, evaluation precision, and cost occur in the
same saturated record as the transferred capacity or hitting bound.

\end{theorem}

\begin{proof}

Summing \cref{eq-controlled-density-comparison} over \(A\) proves
\cref{eq-controlled-density-target-mass-transfer}.  Reversibility gives

\[
\mathcal E_{\Phi_P}(f,f)
=\frac12\sum_{x,z}c_P(x,z)(f(x)-f(z))^2,
\]

and the analogous identity for \(\widetilde P\).  Apply
\cref{eq-controlled-conductance-comparison} termwise to prove
\cref{eq-controlled-density-form-transfer}.  Taking the infimum over the
same condenser boundary conditions proves
\cref{eq-controlled-density-capacity-transfer}.

Let \(\widetilde{\mathbb P}_{\widetilde\mu}\) denote the stationary path law
of \(\widetilde P\), and put
\(E_H=\{\tau_T\le H\}\).  The path law started from \(\rho_0\) has
Radon--Nikodym derivative \(f_0(X_0)\) with respect to
\(\widetilde{\mathbb P}_{\widetilde\mu}\).  H\"older's inequality gives

\[
\Pr_{\rho_0}^{\widetilde P}(E_H)
\le
\|f_0\|_{L^p(\widetilde\mu)}
\left[
\widetilde{\mathbb P}_{\widetilde\mu}(E_H)
\right]^{1/q}.
\]

Apply \cref{eq-controlled-capacity-finite-horizon-hit} with stationary
initial law, followed by
\cref{eq-controlled-density-target-mass-transfer,eq-controlled-density-capacity-transfer}:

\[
\widetilde{\mathbb P}_{\widetilde\mu}(E_H)
\le
\widetilde\mu(T)
+H\operatorname{Cap}_{\Phi_{\widetilde P}}(T,T^c)
\le
r_+\mu(T)+H\kappa_+
\operatorname{Cap}_{\Phi_P}(T,T^c).
\]

Clipping probabilities at one proves
\cref{eq-controlled-density-lp-hit-transfer}.  The case \(p=\infty\)
gives \cref{eq-controlled-density-linfty-hit-transfer}.  If
\(\rho_0\le R_\mu\mu\), then
\cref{eq-controlled-density-comparison} gives

\[
\frac{d\rho_0}{d\widetilde\mu}
\le\frac{R_\mu}{r_-},
\]

which proves \cref{eq-controlled-base-density-hit-transfer}.  Apply the
same argument conditionally from time \(t_0\), and add the event of a hit
through \(t_0\), to obtain the burn-in statement.

\end{proof}

\subsection{Gibbs--Metropolis specialization}
\label{subsec-controlled-gibbs-metropolis-specialization}

\begin{corollary}[Represented Gibbs--Metropolis tilt]
\label{cor-controlled-gibbs-metropolis-capacity-hit}

Let \(P_0\) be a complete represented reversible Markov kernel with
full-support invariant law \(\mu\) on the finite carrier \(X\).  Let
\(D:X\to\mathbb R\) be a represented bounded potential, let
\(\theta\ge0\), and write

\[
D_0=D-\min_XD,
\qquad
\Omega_D=\operatorname{Osc}(D),
\qquad
Z_{\theta,D}=\sum_x\mu(x)e^{-\theta D_0(x)}.
\]

Define the Gibbs law and the Metropolis kernel by

\begin{align}
\label{eq-controlled-gibbs-tilted-law}
\mu_{\theta,D}(x)
&=Z_{\theta,D}^{-1}e^{-\theta D_0(x)}\mu(x),\\
\label{eq-controlled-gibbs-metropolis-kernel}
P_{\theta,D}(x,z)
&=P_0(x,z)
\min\{1,e^{-\theta(D_0(z)-D_0(x))}\},
\qquad x\ne z,
\end{align}

with the diagonal chosen to make every row sum to one.  Then
\(P_{\theta,D}\) is reversible with respect to \(\mu_{\theta,D}\), and

\begin{align}
\label{eq-controlled-gibbs-density-ratio}
e^{-\theta\Omega_D}
&\le
\frac{\mu_{\theta,D}(x)}{\mu(x)}
\le e^{\theta\Omega_D},\\
\label{eq-controlled-gibbs-conductance-ratio}
e^{-\theta\Omega_D}
&\le
\frac{
 \mu_{\theta,D}(x)P_{\theta,D}(x,z)}
 {\mu(x)P_0(x,z)}
\le e^{\theta\Omega_D}
\qquad(P_0(x,z)>0).
\end{align}

Consequently, for every pair of disjoint represented plates \(S,F\),

\begin{equation}
\label{eq-controlled-gibbs-capacity-comparison}
e^{-\theta\Omega_D}
\operatorname{Cap}_{\Phi_{P_0}}(S,F)
\le
\operatorname{Cap}_{\Phi_{P_{\theta,D}}}(S,F)
\le
e^{\theta\Omega_D}
\operatorname{Cap}_{\Phi_{P_0}}(S,F).
\end{equation}

If the represented target satisfies \(T\subseteq\operatorname*{argmin}_X D\),
then

\begin{equation}
\label{eq-controlled-gibbs-target-mass}
\mu_{\theta,D}(T)
=\frac{\mu(T)}{Z_{\theta,D}}
\le\min\{1,e^{\theta\Omega_D}\mu(T)\}.
\end{equation}

For every represented target \(T\), every \(1<p\le\infty\),
\(q=p/(p-1)\), and
\(f_0=d\rho_0/d\mu_{\theta,D}\),

\begin{equation}
\label{eq-controlled-gibbs-lp-hit-bound}
\Pr_{\rho_0}^{P_{\theta,D}}\{\tau_T\le H\}
\le
\min\!\left\{1,
\|f_0\|_{L^p(\mu_{\theta,D})}
\left[
\min\!\left\{1,
e^{\theta\Omega_D}
\bigl(\mu(T)+H\operatorname{Cap}_{\Phi_{P_0}}(T,T^c)\bigr)
\right\}
\right]^{1/q}
\right\}.
\end{equation}

An initialization bound \(\rho_0\le R_\mu\mu\) gives the explicit
specialization

\begin{equation}
\label{eq-controlled-gibbs-base-density-hit-bound}
\Pr_{\rho_0}^{P_{\theta,D}}\{\tau_T\le H\}
\le
\min\!\left\{1,
R_\mu e^{2\theta\Omega_D}
\bigl(\mu(T)+H\operatorname{Cap}_{\Phi_{P_0}}(T,T^c)\bigr)
\right\}.
\end{equation}

If, in addition, \(T\subseteq\operatorname*{argmin}_X D\), the shared
normalizer cancels and the sharper bound is

\begin{equation}
\label{eq-controlled-gibbs-target-minimum-hit-bound}
\Pr_{\rho_0}^{P_{\theta,D}}\{\tau_T\le H\}
\le
\min\!\left\{1,
R_\mu e^{\theta\Omega_D}
\bigl(\mu(T)+H\operatorname{Cap}_{\Phi_{P_0}}(T,T^c)\bigr)
\right\}.
\end{equation}

Finally, let \(F\ge0\) be a nonconstant represented potential with
\(F|_T=0\), and put
\(a_T(P_{\theta,D},F)=\inf_{x\notin T}(I-P_{\theta,D})F(x)\).  Then

\begin{equation}
\label{eq-controlled-gibbs-drift-capacity-ceiling}
\frac{\bigl(a_T(P_{\theta,D},F)\bigr)_+}
     {\operatorname{Osc}(F)}
\le
\frac{\operatorname{Cap}_{\Phi_{P_{\theta,D}}}(T,T^c)}
     {1-\mu_{\theta,D}(T)}.
\end{equation}

Whenever \(e^{\theta\Omega_D}\mu(T)<1\), the right side is at most

\begin{equation}
\label{eq-controlled-gibbs-base-drift-ceiling}
\frac{
e^{\theta\Omega_D}\operatorname{Cap}_{\Phi_{P_0}}(T,T^c)}
{1-e^{\theta\Omega_D}\mu(T)}.
\end{equation}

The represented tilt record contains the construction of \(D_0\), the
normalizer or its certified comparison, the Metropolis acceptance rule, and
the precision used in the displayed constants.  The factor
\(e^{\theta\Omega_D}\) is a quantitative density and conductance loss; the
description length of \(D\) supplies no replacement for this factor.

\end{corollary}

\begin{proof}

For \(x\ne z\), detailed balance follows from

\begin{align*}
\mu_{\theta,D}(x)P_{\theta,D}(x,z)
&=
\frac{\mu(x)P_0(x,z)}{Z_{\theta,D}}
e^{-\theta D_0(x)}
\min\{1,e^{-\theta(D_0(z)-D_0(x))}\}\\
&=
\frac{\mu(x)P_0(x,z)}{Z_{\theta,D}}
e^{-\theta\max\{D_0(x),D_0(z)\}},
\end{align*}

which is symmetric in \(x,z\).  Since
\(0\le D_0\le\Omega_D\),

\[
e^{-\theta\Omega_D}\le Z_{\theta,D}\le1.
\]

This proves
\cref{eq-controlled-gibbs-density-ratio,eq-controlled-gibbs-conductance-ratio}.
Apply
\cref{thm-controlled-density-capacity-hit-transfer} with

\[
r_-=\kappa_-=e^{-\theta\Omega_D},
\qquad
r_+=\kappa_+=e^{\theta\Omega_D}
\]

to obtain
\cref{eq-controlled-gibbs-capacity-comparison,eq-controlled-gibbs-lp-hit-bound}.
When \(D_0=0\) on \(T\), summing
\cref{eq-controlled-gibbs-tilted-law} over \(T\) gives
\cref{eq-controlled-gibbs-target-mass}.  If
\(\rho_0\le R_\mu\mu\), then
\(d\rho_0/d\mu_{\theta,D}\le R_\mu e^{\theta\Omega_D}\); combine this with
the bracket in \cref{eq-controlled-gibbs-lp-hit-bound} for \(p=\infty\)
to prove \cref{eq-controlled-gibbs-base-density-hit-bound}.

For \(x\in T\) and \(z\notin T\), the target-minimum hypothesis gives
\[
\mu_{\theta,D}(x)P_{\theta,D}(x,z)
=\frac{e^{-\theta D_0(z)}}{Z_{\theta,D}}\mu(x)P_0(x,z)
\le\frac1{Z_{\theta,D}}\mu(x)P_0(x,z).
\]
Summing the target cut and using
\cref{eq-controlled-target-boundary-capacity-flux} yields
\(
\operatorname{Cap}_{\Phi_{P_{\theta,D}}}(T,T^c)
\le Z_{\theta,D}^{-1}
\operatorname{Cap}_{\Phi_{P_0}}(T,T^c)
\).
Moreover,
\(
d\rho_0/d\mu_{\theta,D}
\le R_\mu Z_{\theta,D}e^{\theta\Omega_D}
\).
The stationary entrance-count bound
\cref{eq-controlled-capacity-finite-horizon-hit} cancels
\(Z_{\theta,D}\), proving
\cref{eq-controlled-gibbs-target-minimum-hit-bound}.

Apply
\cref{thm-controlled-invariant-mass-capacity-drift} to the selected kernel
\(P_{\theta,D}\) and potential \(F\) to prove
\cref{eq-controlled-gibbs-drift-capacity-ceiling}.  The upper capacity
comparison and
\(\mu_{\theta,D}(T)\le e^{\theta\Omega_D}\mu(T)\) then give
\cref{eq-controlled-gibbs-base-drift-ceiling}.

\end{proof}

The capacity, density, and Gibbs comparisons above enter the native
finite-horizon coordinate of
\cref{cor-nonlinear-nonquotient-dynamical-envelope} and the typed dynamical
certificate of \cref{def-coordinatewise-dynamical-certificate}.  Each use
retains its selected invariant law, initialization density, conductance
comparison, and oscillation loss in one complete record.


\section{Affordable Caus Action and Capacity Under Valid Lifts}
\label{sec-controlled-caus-lift-capacity}

The capacity and drift estimates above apply to the selected native carrier.
The present section gives the corresponding accounting and comparison tools
for authorized directed flow and for valid lifted carriers.  All suprema below
range over complete records.  In particular, the selected dynamics, target
interface, reference law, comparison maps, numerical precision, and every
evaluation used by the bound occur in the same saturated record.

\subsection{Caus action, aiming, and target contact}
\label{subsec-controlled-caus-action}

\begin{definition}[Complete Caus target-flow record]
\label{def-controlled-complete-caus-flow-record}

Fix a finite represented carrier, an orientation \(E^+\) of its active
unoriented edges, and a budget \(B\).  A complete Caus target-flow record
consists of

\begin{enumerate}[label=\textup{(\roman*)}]
\item a selected positive generator or transition kernel, its legal reference
      law, and the authority envelope of
      \cref{def-caus-authority-envelope,def-authorized-caus-use-current};
\item the actual authorized current \(J=(J_e)_{e\in E^+}\), rather than an
      independently chosen current with the same formal type;
\item a represented target-normalized potential \(D\), available before the
      transition on which it is used, together with its target and failure
      gates;
\item positive represented edge mobilities \(a=(a_e)_{e\in E^+}\), the
      horizon or resolvent scale, and every comparison used to convert flow
      into contact, selector, or hitting units; and
\item the complete derivation, storage, arithmetic, solver, comparison, and
      precision costs of the preceding data.
\end{enumerate}

The record is \emph{prospective} when its target-aligned entries satisfy
\cref{def-pre-hit-temporally-admissible-source}.  With
\(\nabla_eD=D(e^+)-D(e^-)\), define
\begin{align}
\mathcal A_{\Caus}(J;a)
&=\sum_{e\in E^+}\frac{J_e^2}{a_e},
\label{eq-controlled-caus-action}\\
\mathcal E_{\rm aim}(D;a)
&=\sum_{e\in E^+}a_e(\nabla_eD)^2,
\label{eq-controlled-caus-aiming-energy}\\
\mathcal P_{\Caus}(J,D)
&=\langle J,-\nabla D\rangle_{E^+}.
\label{eq-controlled-caus-target-power}
\end{align}
The first quantity is the \emph{Caus action}, the second is the
\emph{target-aiming energy}, and the third is the signed target power.
When the mobility-compatible coordinates in
\cref{def-quantitative-caus-alignment,def-authority-compatible-caus-alignment}
are chosen as \(w_e=a_e\) and \(F_e=J_e/a_e\), the power is their
unnormalized authorized alignment numerator.  Any other coordinate choice
requires a represented comparison before the two pairings are identified.

Suppose additionally that the record contains positive oriented stationary
or instantaneous fluxes \(\varphi_e,\varphi_{\bar e}\), with
the represented orientation convention
\(J_e=\varphi_e-\varphi_{\bar e}\), and uses the logarithmic-mean mobility
\[
a_e=\frac{\varphi_e-\varphi_{\bar e}}
          {\log\varphi_e-\log\varphi_{\bar e}},
\]
with the continuous diagonal value.  Then \(\mathcal A_{\Caus}\) is the
edge entropy production
\[
\sum_{e\in E^+}
(\varphi_e-\varphi_{\bar e})
\log\frac{\varphi_e}{\varphi_{\bar e}}.
\]
The same terminology is permitted for an authorized Hodge projection only
when the record proves contraction in this inverse-mobility norm.  In every
other case the quantity in \cref{eq-controlled-caus-action} retains the name
Caus action.

\end{definition}

\begin{theorem}[Caus action--aiming inequality]
\label{thm-controlled-caus-action-aiming}

Every complete Caus target-flow record satisfies
\begin{equation}
\label{eq-controlled-caus-action-aiming-pointwise}
\bigl(\mathcal P_{\Caus}(J,D)_+\bigr)^2
\le
\mathcal A_{\Caus}(J;a)\mathcal E_{\rm aim}(D;a).
\end{equation}
For a measurable time-dependent record on \([0,H]\),
\begin{equation}
\label{eq-controlled-caus-action-aiming-integrated}
\int_0^H\mathcal P_{\Caus}(J_t,D_t)_+\,dt
\le
\left(\int_0^H\mathcal A_{\Caus}(J_t;a_t)\,dt\right)^{1/2}
\left(\int_0^H\mathcal E_{\rm aim}(D_t;a_t)\,dt\right)^{1/2}.
\end{equation}
The discrete-time version replaces both integrals by sums.

\end{theorem}

\begin{proof}

Write \(u_e=J_e/\sqrt{a_e}\) and
\(v_e=-\sqrt{a_e}\nabla_eD\).  Cauchy--Schwarz gives
\(|\langle u,v\rangle|^2\le\|u\|_2^2\|v\|_2^2\), which is
\cref{eq-controlled-caus-action-aiming-pointwise}.  Apply
Cauchy--Schwarz once more in time to obtain
\cref{eq-controlled-caus-action-aiming-integrated}.

\end{proof}

\begin{definition}[Saturated Caus action profiles]
\label{def-controlled-saturated-caus-action-profiles}

Let \(\mathscr R_{\Caus}^{\rm pre}(B,H)\) be the represented family of
prospective complete Caus target-flow records of cost at most \(B\) and
horizon \(H\).  Define
\begin{align*}
\operatorname{Act}^{\rm sat}_{\Caus}(B,H)
&=\sup_{R\in\mathscr R_{\Caus}^{\rm pre}(B,H)}
  \int_0^H\mathcal A_{\Caus,R}(t)\,dt,\\
\operatorname{Aim}^{\rm sat}_{\Caus}(B,H)
&=\sup_{R\in\mathscr R_{\Caus}^{\rm pre}(B,H)}
  \int_0^H\mathcal E_{{\rm aim},R}(t)\,dt,\\
\operatorname{Flow}^{\rm sat}_{\Caus}(B,H)
&=\sup_{R\in\mathscr R_{\Caus}^{\rm pre}(B,H)}
  \int_0^H\mathcal P_{\Caus,R}(t)_+\,dt.
\end{align*}
Empty suprema are zero.  A record with an unavailable upper action or aiming
certificate has value \(+\infty\) in the corresponding upper profile.  These
profiles are accounting sub-suprema of the prospective source class in
\cref{def-source-accounting-profile-separation}; their family-uniform
descriptions and the optimization used to evaluate them are included in the
saturated cost.

A represented \emph{aiming-to-contact converter} is a same-record inequality
\begin{equation}
\label{eq-controlled-aiming-contact-converter}
\operatorname{Hit}_{\rm contact}(R)
\le B_R^{\rm contact}
 +K_R\int_0^H\mathcal P_{\Caus,R}(t)_+\,dt+\delta_R,
\end{equation}
where the left side is the total contact in the prospective-selector,
hitting-probability, or discounted-arrival normalization.  A converter that
directly bounds total contact uses \(B_R^{\rm contact}=0\); a converter that
only bounds advantage uses its represented neutral-contact baseline, such as
\(\operatorname{Enum}_\nu\).  The quantities
\(B_R^{\rm contact},K_R,\delta_R\) are represented before the relevant
transitions.  A committor may supply this converter only through the
represented derivation required by \cref{thm-controlled-no-free-value}.

\end{definition}

\begin{corollary}[Uniform affordable Caus-flow bound]
\label{cor-controlled-uniform-caus-flow-bound}

For every \(B,H\),
\begin{equation}
\label{eq-controlled-saturated-caus-flow-bound}
\operatorname{Flow}^{\rm sat}_{\Caus}(B,H)
\le
\sqrt{\operatorname{Act}^{\rm sat}_{\Caus}(B,H)
      \operatorname{Aim}^{\rm sat}_{\Caus}(B,H)}.
\end{equation}
If all records in a represented subclass have
\(B_R^{\rm contact}\le B_B^{\rm contact}\), \(K_R\le K_B\), and
\(\delta_R\le\delta_B\) in
\cref{eq-controlled-aiming-contact-converter}, then their total contact is at
most
\begin{equation}
\label{eq-controlled-saturated-caus-contact-bound}
B_B^{\rm contact}
+K_B\sqrt{\operatorname{Act}^{\rm sat}_{\Caus}(B,H)
          \operatorname{Aim}^{\rm sat}_{\Caus}(B,H)}+\delta_B.
\end{equation}
Without a represented converter, the action and aiming profiles remain flow
coordinates and supply no probability bound.

\end{corollary}

\begin{proof}

Apply \cref{eq-controlled-caus-action-aiming-integrated} to every record and
then enlarge each factor to its saturated supremum.  Substitution into
\cref{eq-controlled-aiming-contact-converter} proves the second assertion.

\end{proof}

\begin{example}[Action does not determine aiming]
\label{ex-controlled-caus-action-not-aiming}

Let the oriented edges be \(e_1,e_2\), take \(a_{e_1}=a_{e_2}=1\), and let
\(-\nabla D=(1,0)\).  The currents \(J^{\parallel}=(1,0)\) and
\(J^{\perp}=(0,1)\) have the same Caus action, equal to one, whereas
\(\mathcal P_{\Caus}(J^{\parallel},D)=1\) and
\(\mathcal P_{\Caus}(J^{\perp},D)=0\).  These vectors may be realized on a
finite graph by a target-directed edge and a disjoint cycle edge.  Hence an
action or entropy-production bound evaluates total directed activity; the
aiming factor evaluates its concentration on the represented target
potential.

\end{example}

\subsection{Noise-corrected Caus aiming}
\label{subsec-controlled-noise-corrected-caus}

The action of an actual represented current is evaluated on the selected
record and requires no label correction.  Label noise enters when a learned
potential is used as a proxy for the latent clean target potential.

\begin{definition}[Clean-aiming comparison record]
\label{def-controlled-clean-aiming-comparison}

Augment a complete Caus target-flow record by a complete dynamical
label-noise record of
\cref{def-complete-dynamical-label-noise-record}.  A clean-aiming comparison
contains a latent clean target-normalized analytic comparator \(D_0\), the
represented learned potential \(\widehat D\), the same edge mobilities
\(a_t\), and one of
the following represented certificates.

\begin{enumerate}[label=\textup{(\roman*)}]
\item \emph{Linear homogeneous debiasing.}  The population corrupted
      potential is \(D_\eta=\rho D_0\), after the declared centering, and
      \[
      \left(
      \int_0^H
      \mathcal E_{\rm aim}(\widehat D_t-D_{\eta,t};a_t)\,dt
      \right)^{1/2}
      \le\varepsilon_{\rm aim}.
      \]
\item \emph{Calibrated potential comparison.}  A represented calibration
      theorem gives directly
      \[
      \left(
      \int_0^H
      \mathcal E_{\rm aim}(\widehat D_t-D_{0,t};a_t)\,dt
      \right)^{1/2}
      \le\varepsilon_{\rm cal}.
      \]
\item \emph{Binary hard-potential fallback.}  Both potentials take values in
      \(\{0,1\}\).  Let \(d^b\) be the represented behavior law and assume
      the deployment comparisons \(d_t^\pi\le W_t d^b\).  Under the
      represented deployment laws \(d_t^\pi\), assume
      \[
      \Pr_{d_t^\pi}\{\widehat D_t\ne D_{0,t}\}
      \le U_{{\rm gate},{\rm dep}}^{\rm clean}
      \qquad(0\le t\le H),
      \]
      where the latter is a represented behavior-to-deployment consequence
      of the dynamical clean-risk bound,
      \[
      U_{{\rm gate},{\rm dep}}^{\rm clean}
      \le
      \min\left\{1,
      \left(\sup_{0\le t\le H}W_t\right)
      U_{\rm noise}^{\rm clean}(n,b(n),m,\eta_+,\delta)
      \right\}.
      \]
      Let the represented time-dependent gradient operators satisfy
      \[
      \kappa_{\rm aim,H}
      =
      \int_0^H
      \|\nabla_t\|_{L^2(d_t^\pi)\to\ell^2(a_t)}^2\,dt
      <\infty.
      \]
\end{enumerate}

The noise rate, centering, calibration, gradient norm, deployment
comparison, and their computation and precision are part of the same
saturated record.  A general score or ranking potential enters this
definition through case \textup{(i)} or \textup{(ii)}; scalar label
correlation alone is not a gradient-energy comparison.
The comparator \(D_0\) is used to state and audit the clean-risk comparison;
it is not a prospective source or a policy input unless an independent
pre-hit represented derivation constructs it.  Source ancestry remains that
of \(\widehat D\) and its complete training record.

\end{definition}

\begin{theorem}[Dynamical noise correction for clean Caus aiming]
\label{thm-controlled-noise-corrected-caus-aiming}

In case \textup{(i)} of
\cref{def-controlled-clean-aiming-comparison},
\begin{equation}
\label{eq-controlled-noise-corrected-aiming-energy}
\left(
\int_0^H\mathcal E_{\rm aim}(D_{0,t};a_t)\,dt
\right)^{1/2}
\le
\frac{
\left(
\int_0^H\mathcal E_{\rm aim}(\widehat D_t;a_t)\,dt
\right)^{1/2}
+\varepsilon_{\rm aim}}
{\rho}.
\end{equation}
In case \textup{(ii)}, the right side is
\[
\left(
\int_0^H\mathcal E_{\rm aim}(\widehat D_t;a_t)\,dt
\right)^{1/2}
+\varepsilon_{\rm cal}.
\]
In case \textup{(iii)},
\begin{equation}
\label{eq-controlled-binary-aiming-error}
\int_0^H
\mathcal E_{\rm aim}(\widehat D_t-D_{0,t};a_t)\,dt
\le
\kappa_{\rm aim,H}U_{{\rm gate},{\rm dep}}^{\rm clean}.
\end{equation}
At each time the squared \(L^2(d_t^\pi)\) distance appearing in this
estimate is exactly
\[
\|\widehat D_t-D_{0,t}\|_{L^2(d_t^\pi)}^2
=
\Pr_{d_t^\pi}\{\widehat D_t\ne D_{0,t}\}.
\]

For each record, minimize its valid same-unit aiming errors over these three
cases; let \(\varepsilon_{\rm aim}^{\rm sat}\) be the supremum of that
minimum over the complete budget slice.  Let
\(\operatorname{Aim}_{\Caus,{\rm obs}}^{\rm sat}\) be the observed-potential
aiming profile.  Every represented clean aiming-to-contact converter then
gives
\begin{equation}
\label{eq-controlled-noise-corrected-caus-contact}
U_{\Caus}^{\rm noise}
\le
\min\left\{
1,\,
B_B^{\rm contact}
+K_B
\sqrt{\operatorname{Act}_{\Caus}^{\rm sat}}\,
\frac{
\sqrt{\operatorname{Aim}_{\Caus,{\rm obs}}^{\rm sat}}
+\varepsilon_{\rm aim}^{\rm sat}}
{\rho}
+\delta_B
\right\}
\end{equation}
in the linear homogeneous case.  The calibrated case omits the division by
\(\rho\).  In the hard-potential branch, preserve the recordwise order:
\begin{equation}
\label{eq-controlled-hard-gate-aiming-envelope}
\varepsilon_{{\rm aim},{\rm hard}}^{\rm sat}
=
\sup_{\mathcal R\in\mathscr R_{B,H}^{\rm hard}}
\sqrt{
\kappa_{\rm aim,H}(\mathcal R)
U_{{\rm gate},{\rm dep}}^{\rm clean}(\mathcal R)}
\le
\sqrt{
\left(\sup_{\mathcal R}\kappa_{\rm aim,H}(\mathcal R)\right)
\left(\sup_{\mathcal R}
U_{{\rm gate},{\rm dep}}^{\rm clean}(\mathcal R)\right)}.
\end{equation}
Here \(\mathscr R_{B,H}^{\rm hard}\) is the complete affordable hard-gate
record class at the displayed ceiling and horizon, and every unlabeled
supremum on the right ranges over that same class.
Use the left side as the hard-branch candidate in
\(\varepsilon_{\rm aim}^{\rm sat}\).  This does not multiply minima attained
by different records.
All profiles in this display use the same horizon, mobilities, target
normalization, deployment law, and class-preserving enlarged ceiling.
For a hard gate learned at ceiling \(b(n)\) from \(m\) labels with Massart
rate at most \(\eta_+\), substitute
the displayed behavior-to-deployment bound on
\(U_{{\rm gate},{\rm dep}}^{\rm clean}\).
The polynomial-capacity conclusion is the same display with
\(b(n)=b_C(n)=n^C\).

\end{theorem}

\begin{proof}

The map \(D\mapsto(\int\mathcal E_{\rm aim}(D;a))^{1/2}\) is a seminorm.
The triangle inequality applied to
\(\rho D_0=\widehat D-(\widehat D-D_\eta)\) proves
\cref{eq-controlled-noise-corrected-aiming-energy}; the calibrated case is
identical without rescaling.  The definition of \(\kappa_{\rm aim,H}\) gives
the first inequality in
\cref{eq-controlled-binary-aiming-error} after integration in time, and the
squared \(L^2\) distance between binary gates is their misclassification
probability.  Combine these
comparisons with
\cref{thm-controlled-caus-action-aiming,cor-controlled-uniform-caus-flow-bound}
and the represented aiming-to-contact converter to obtain
\cref{eq-controlled-noise-corrected-caus-contact}.

\end{proof}

\subsection{Variational capacity comparison under valid lifts}
\label{subsec-controlled-lift-capacity}

\begin{definition}[Complete form-lift comparison record]
\label{def-controlled-complete-form-lift-record}

Let \((\Omega,\mathcal G_{\rm obs},\pi,\Lambda)\) be a valid lift with fiber
laws \(\kappa_x\), pullback \(U\), averaging map \(P_\kappa\), and
\(Q_\pi=I-UP_\kappa\) as in \cref{thm-exhaustive-lift-normal-form}.  A complete
form-lift comparison record contains a full-support represented base law
\(\mu\), its disintegrated lifted law
\[
\widetilde\mu(\omega)=\mu(\pi(\omega))\kappa_{\pi(\omega)}(\omega),
\]
a nonnegative self-adjoint represented operator \(H_\Omega\) on
\(L^2(\operatorname{supp}\widetilde\mu,\widetilde\mu)\), and represented
disjoint base boundary sets
\(S,T\).  It also contains the induced form, all comparison constants,
boundary and target gates, pseudoinverse or solver, precision, and complete
construction cost.  In these Hilbert spaces \(U\) is an isometry,
\(P_\kappa=U^*\), and \(Q_\pi\) is the orthogonal fiber projection.

Let \(\mathcal Q_0\) be the subspace of \(\operatorname{ran}Q_\pi\) whose
members vanish on \(\pi^{-1}(S\cup T)\), and let \(Q_0\) denote its
orthogonal projection.  Set
\[
H_X=U^*H_\Omega U,
\qquad H_{F,0}=Q_0H_\Omega Q_0|_{\mathcal Q_0},
\qquad B_0=Q_0H_\Omega U.
\]
For a base function \(g\) with positive energy define the leakage ratio
\begin{equation}
\label{eq-controlled-lift-leakage-ratio}
\rho_\pi
=\sup_{\langle g,H_Xg\rangle>0}
\frac{\langle B_0g,H_{F,0}^+B_0g\rangle}
     {\langle g,H_Xg\rangle}.
\end{equation}
The supremum is restricted to the represented condenser domain when the
comparison is used only for \((S,T)\).

\end{definition}

\begin{theorem}[Schur-complement capacity theorem for valid lifts]
\label{thm-controlled-lift-capacity-schur}

For every complete form-lift comparison record, positivity implies
\(B_0g\in\operatorname{ran}H_{F,0}\) and \(0\le\rho_\pi\le1\).  Minimization over
fiber fluctuations gives the effective base operator
\begin{equation}
\label{eq-controlled-lift-effective-form}
H_{\rm eff}=H_X-B_0^*H_{F,0}^+B_0.
\end{equation}
For the pulled-back condenser \((\pi^{-1}S,\pi^{-1}T)\),
\begin{equation}
\label{eq-controlled-lift-capacity-sandwich}
(1-\rho_\pi)\operatorname{Cap}_{X}(S,T)
\le
\operatorname{Cap}_{\Omega}(\pi^{-1}S,\pi^{-1}T)
\le
\operatorname{Cap}_{X}(S,T),
\end{equation}
where \(\operatorname{Cap}_{X}\) is computed with the pullback form
\(g\mapsto\langle g,H_Xg\rangle\).  If a represented comparison form
\(\mathcal E_\Phi\) satisfies
\[
a_\pi\mathcal E_\Phi(g,g)
\le\langle g,H_Xg\rangle
\le b_\pi\mathcal E_\Phi(g,g),
\]
then
\begin{equation}
\label{eq-controlled-lift-external-capacity-sandwich}
(1-\rho_\pi)a_\pi\operatorname{Cap}_{\Phi}(S,T)
\le
\operatorname{Cap}_{\Omega}(\pi^{-1}S,\pi^{-1}T)
\le
b_\pi\operatorname{Cap}_{\Phi}(S,T).
\end{equation}
In particular, \(B_0=0\) gives exact equality between the lifted and pullback
capacities.

\end{theorem}

\begin{proof}

Relative to \(\operatorname{ran}U\oplus\mathcal Q_0\), the quadratic
form of \(H_\Omega\) at \(Ug+q\), with \(q\in\mathcal Q_0\), is
\[
\langle g,H_Xg\rangle+2\langle q,B_0g\rangle
+\langle q,H_{F,0}q\rangle.
\]
For a finite positive block matrix, \(B_0g\) lies in
\(\operatorname{ran}H_{F,0}\); otherwise a vector in \(\ker H_{F,0}\) paired with
\(B_0g\) would make the displayed quadratic polynomial negative for one sign
of its scalar coefficient.  Completing the square shows that the minimum
over \(q\) is attained at \(-H_{F,0}^+B_0g\) modulo \(\ker H_{F,0}\), with value
\(\langle g,H_{\rm eff}g\rangle\).  Positivity gives
\(0\le B_0^*H_{F,0}^+B_0\le H_X\), hence \(0\le\rho_\pi\le1\).

The lifted boundary conditions constrain \(g\) on \(S,T\) and constrain
\(q\) to vanish on the corresponding boundary fibers.  The same completion
of squares on that condenser domain gives
\((1-\rho_\pi)\langle g,H_Xg\rangle
\le\langle g,H_{\rm eff}g\rangle
\le\langle g,H_Xg\rangle\).  Taking the boundary infimum proves
\cref{eq-controlled-lift-capacity-sandwich}.  Form comparison before the
same infimum proves
\cref{eq-controlled-lift-external-capacity-sandwich}.

\end{proof}

\begin{definition}[Saturated lift-capacity amplification]
\label{def-controlled-saturated-lift-capacity-amplification}

For a represented class \(\mathscr L(B)\) of complete form-lift records,
require each \(R\in\mathscr L(B)\) to carry a represented external comparison
form \(\mathcal E_{\Phi,R}\) and a finite constant \(b_{\pi,R}\) satisfying
the upper comparison in
\cref{thm-controlled-lift-capacity-schur}.  If an application starts from a
larger form-lift class, \(\mathscr L(B)\) below denotes its
comparison-bearing subclass.  The supremum of an empty subclass is zero;
a record without such a comparison has the neutral external-capacity bound
\(+\infty\) and is not inserted into this finite profile.  Put
\[
\operatorname{LiftAmp}^{\rm sat}_{\mathscr L}(B)
=\sup_{R\in\mathscr L(B)} b_{\pi,R},
\qquad
\operatorname{LiftLeak}^{\rm sat}_{\mathscr L}(B)
=\sup_{R\in\mathscr L(B)}\rho_{\pi,R}.
\]
Then, at the class-preserving enlarged ceiling containing the lift and
comparison computations,
\begin{equation}
\label{eq-controlled-saturated-lift-capacity-upper}
\sup_{R\in\mathscr L(B)}\operatorname{Cap}_{\Omega,R}
\le
\operatorname{LiftAmp}^{\rm sat}_{\mathscr L}(B)
\sup_{R\in\mathscr L(B)}\operatorname{Cap}_{\Phi,R}.
\end{equation}
Capacities in this display use the represented common rate normalization.
The valid-lift identity fixes the target interface and preserves target
fractions; evaluation of
\(\operatorname{LiftAmp}^{\rm sat}_{\mathscr L}\) is the additional uniform
computation over instances of the lift class.

\end{definition}

\subsection{Dynamic lift transfer}
\label{subsec-controlled-dynamic-lift-transfer}

\begin{definition}[Dynamic lift blocks and structural subclasses]
\label{def-controlled-dynamic-lift-classes}

For a complete valid-lift record and a lifted generator \(L_\Omega\), use the
orthogonal disintegration above and define
\[
L_X=P_\kappa L_\Omega U,\qquad
E_+=Q_\pi L_\Omega U,\qquad
E_-=P_\kappa L_\Omega Q_\pi,\qquad
L_F=Q_\pi L_\Omega Q_\pi.
\]
Thus \(E_+\) is the lift-normal-form defect \(E_\pi\) in the orthogonal
disintegration; it is not a second structural record.
The following represented subclasses will be used.

\begin{enumerate}[label=\textup{(\roman*)}]
\item \(\mathscr L_{\rm int}\): exact represented intertwining,
      \(E_+=0\);
\item \(\mathscr L_{\rm form}\): form-neutral lifts, \(Q_\pi H_\Omega U=0\);
\item \(\mathscr L_{\rm fm}(M_F,\kappa_F,\varepsilon_+,\varepsilon_-)\):
      fast-mixing weakly coupled lifts satisfying, in one declared operator
      norm,
      \[
      \|e^{tL_F}Q_\pi\|\le M_Fe^{-\kappa_Ft},\qquad
      \|E_+\|\le\varepsilon_+,
      \quad\|E_-\|\le\varepsilon_-;
      \]
\item \(\mathscr L_{\rm all}\): all complete valid-lift records, with no
      numerical intertwining or fiber-mixing assumption.
\end{enumerate}

Every constant and norm conversion belongs to the same record.  A
Poincar\'e, sector, or hypocoercive certificate supplies \(M_F,\kappa_F\)
only through a represented theorem in the displayed norm.  MLSI and
statistical \(\beta\)-mixing estimates are not substituted without an
explicit conversion to that norm.

\end{definition}

\begin{theorem}[Block-resolvent transfer for valid lifts]
\label{thm-controlled-valid-lift-block-resolvent}

Let \(\lambda>0\) lie in the resolvent sets required below.  Put
\(R_F(\lambda)=(\lambda-L_F)^{-1}\) on \(\operatorname{ran}Q_\pi\).  Then
\begin{equation}
\label{eq-controlled-valid-lift-feshbach}
P_\kappa(\lambda-L_\Omega)^{-1}U
=\bigl[\lambda-L_X-E_-R_F(\lambda)E_+\bigr]^{-1}.
\end{equation}
For \(R\in\mathscr L_{\rm int}\), the projected lifted resolvent equals
\((\lambda-L_X)^{-1}\).  When these are the killed or absorbing generators
for transported target and failure gates, exact initialization and pullback
gates give equality of the corresponding discounted first-arrival
functionals.

For \(R\in\mathscr L_{\rm fm}\),
\begin{equation}
\label{eq-controlled-fast-fiber-resolvent}
\|R_F(\lambda)Q_\pi\|\le\frac{M_F}{\lambda+\kappa_F},
\qquad
\|E_-R_F(\lambda)E_+\|
\le\frac{\varepsilon_-M_F\varepsilon_+}
         {\lambda+\kappa_F}.
\end{equation}
Writing \(R_X=(\lambda-L_X)^{-1}\) and
\(\sigma_\lambda=\varepsilon_-M_F\varepsilon_+/(\lambda+\kappa_F)\),
if \(\|R_X\|\sigma_\lambda<1\), then
\begin{equation}
\label{eq-controlled-fast-fiber-schur-error}
\left\|P_\kappa(\lambda-L_\Omega)^{-1}U-R_X\right\|
\le
\frac{\|R_X\|^2\sigma_\lambda}
     {1-\|R_X\|\sigma_\lambda}.
\end{equation}
Pairing this estimate with represented initial and target functionals gives
the same bound multiplied by their dual norms, plus their represented
interface errors.  For a valid exact pullback these interface errors vanish.

\end{theorem}

\begin{proof}

The block matrix of \(\lambda-L_\Omega\) on
\(\operatorname{ran}U\oplus\operatorname{ran}Q_\pi\) is
\[
\begin{pmatrix}
\lambda-L_X&-E_-\\
-E_+&\lambda-L_F
\end{pmatrix}.
\]
Block Gaussian elimination gives
\cref{eq-controlled-valid-lift-feshbach}.  Exact intertwining sets \(E_+=0\).
For the fast-fiber class, integrate the semigroup estimate to obtain the
first inequality in \cref{eq-controlled-fast-fiber-resolvent};
submultiplicativity gives the second.  The resolvent identity followed by the
Neumann-series bound proves
\cref{eq-controlled-fast-fiber-schur-error}.  Pairing with the two boundary
functionals proves the arrival statement.

\end{proof}

\begin{theorem}[Clocked finite-horizon lift transfer]
\label{thm-controlled-clocked-lift-transfer}

Let \(P_t^\Omega\) and \(P_t^X\), \(0\le t<H\), be represented one-step
kernels on compatible clock slices, acting backward on observables, with
represented pullbacks
\(U_t:L^\infty(X_t)\to L^\infty(\Omega_t)\).
For a first-arrival conclusion these are the kernels stopped, killed, or made
absorbing at the transported target and failure gates.
Define the pre-transition intertwining defect
\[
E_t=P_t^\Omega U_{t+1}-U_tP_t^X.
\]
Then
\begin{align}
&P_0^\Omega\cdots P_{H-1}^\Omega U_H
-U_0P_0^X\cdots P_{H-1}^X\notag\\
&\qquad=
\sum_{t=0}^{H-1}
P_0^\Omega\cdots P_{t-1}^\Omega
E_t
P_{t+1}^X\cdots P_{H-1}^X,
\label{eq-controlled-clocked-lift-telescoping}
\end{align}
where an empty product is the identity on its clock slice.
In a supremum or total-variation norm in which the kernels are contractions,
the norm of the left side is at most \(\sum_{t=0}^{H-1}\|E_t\|\).  Exact
intertwining gives equality of the projected finite-horizon laws.  With exact
valid-lift initialization and target gates for the stopped dynamics, the same
estimates hold for finite-horizon first arrival; approximate represented
interfaces add their two errors.  For unrestricted kernels the conclusion is
an occupation-law comparison, not a first-arrival comparison.

Each \(E_t\), comparator, and interface is derived before transition \(t\)
and charged at the complete clocked-record cost.  Statistical blocking may
certify an estimator of these defects, but it does not replace the displayed
dynamical comparison.

\end{theorem}

\begin{proof}

For \(0\le t\le H\), put
\[
Z_t
=P_0^\Omega\cdots P_{t-1}^\Omega
 U_t
 P_t^X\cdots P_{H-1}^X,
\]
with the empty products interpreted as identities.  Thus the left side of
\cref{eq-controlled-clocked-lift-telescoping} is \(Z_H-Z_0\), while
\[
Z_{t+1}-Z_t
=P_0^\Omega\cdots P_{t-1}^\Omega
 \bigl(P_t^\Omega U_{t+1}-U_tP_t^X\bigr)
 P_{t+1}^X\cdots P_{H-1}^X.
\]
Summing these identities proves the displayed telescoping formula.  Kernel
and pullback contraction, followed by the triangle inequality, gives the norm
bound.  Pairing the resulting observable comparison with the represented
initial law and the transported terminal indicator proves the first-arrival
statement.

\end{proof}

\begin{theorem}[Fast-fiber and clocked-defect Lift certificate split]
\label{thm-controlled-lift-transfer-certificate-split}

Fix one complete valid-Lift record with transported target and failure
interfaces.  The following two numerical transfer branches use the existing
dynamic blocks of \cref{def-controlled-dynamic-lift-classes}.

\begin{enumerate}[label=\textup{(\roman*)}]

\item Suppose the record contains a fast-fiber certificate

\[
\|e^{tL_F}Q_\pi\|\le M_Fe^{-\kappa_Ft},
\qquad
\|E_+\|\le\varepsilon_+,
\qquad
\|E_-\|\le\varepsilon_-,
\]

at a represented discount \(\lambda>0\).  Put

\begin{equation}
\label{eq-controlled-certified-fast-fiber-parameters}
\sigma_\lambda
=\frac{\varepsilon_-M_F\varepsilon_+}{\lambda+\kappa_F},
\qquad
\Delta_{\rm fm}(\lambda)
=\frac{\|R_X(\lambda)\|^2\sigma_\lambda}
       {1-\|R_X(\lambda)\|\sigma_\lambda},
\end{equation}

where \(R_X(\lambda)=(\lambda-L_X)^{-1}\), and assume
\(\|R_X(\lambda)\|\sigma_\lambda<1\).  For every represented base row
functional \(\ell\) and base column \(g\),

\begin{equation}
\label{eq-controlled-certified-fast-fiber-pairing}
\left|
\ell\!\left[
P_\kappa(\lambda-L_\Omega)^{-1}U-R_X(\lambda)
\right]g
\right|
\le
\|\ell\|_*\,\|g\|\,\Delta_{\rm fm}(\lambda).
\end{equation}

Let \(a_{0,T}^\Omega\) and \(a_{0,T}^X\) denote the represented lifted and
base initial target atoms, respectively.  If represented discounted-arrival
functionals
\(A_T^\Omega(\lambda)\) and \(A_T^X(\lambda)\) satisfy the endpoint
interface comparisons

\begin{align*}
\left|A_T^\Omega(\lambda)-a_{0,T}^\Omega
-\ell P_\kappa(\lambda-L_\Omega)^{-1}Ug\right|
&\le\varepsilon_{\Omega,\lambda}^{\rm int},\\
\left|A_T^X(\lambda)-a_{0,T}^X-\ell R_X(\lambda)g\right|
&\le\varepsilon_{X,\lambda}^{\rm int},
\end{align*}

and
\[
\left|a_{0,T}^\Omega-a_{0,T}^X\right|
\le\varepsilon_{0,T}^{\rm int},
\]
then

\begin{equation}
\label{eq-controlled-certified-fast-fiber-arrival}
\left|A_T^\Omega(\lambda)-A_T^X(\lambda)\right|
\le
\varepsilon_{0,T}^{\rm int}
+\varepsilon_{\Omega,\lambda}^{\rm int}
+\varepsilon_{X,\lambda}^{\rm int}
+\|\ell\|_*\,\|g\|\,\Delta_{\rm fm}(\lambda).
\end{equation}

Exact initialization and exact pullback target contact set all three
interface errors to zero with \(\ell=\mu_{0,N}\) and \(g=k_T\) in the
target-resolvent normalization of \cref{def-target-resolvent}.

\item Suppose the record supplies represented clock slices, stopped or
absorbing one-step kernels, and the defects

\[
E_t=P_t^\Omega U_{t+1}-U_tP_t^X,
\qquad 0\le t<H,
\]

in a contraction norm.  If its initialization and terminal-gate interface
errors are bounded by
\(\varepsilon_{\rm init}\) and \(\varepsilon_{\rm gate}\), then

\begin{equation}
\label{eq-controlled-certified-clocked-defect-arrival}
\left|p_{T,H}^\Omega-p_{T,H}^X\right|
\le
\min\!\left\{1,
\varepsilon_{\rm init}
+\sum_{t=0}^{H-1}\|E_t\|
+\varepsilon_{\rm gate}
\right\}.
\end{equation}

This branch applies to every represented clocked record carrying these
defects, including the complement of the certified fast-fiber branch.  Its
right side contains no fiber decay rate, off-diagonal coupling norm, or
resolvent smallness condition.

\end{enumerate}

When both branches are certified, both conclusions hold in their respective
units.  A represented conversion between discounted arrival and the chosen
finite horizon permits the corresponding converted upper bounds to be
minimized.  Without such a conversion,
\cref{eq-controlled-certified-fast-fiber-arrival} remains a discounted-
arrival estimate and
\cref{eq-controlled-certified-clocked-defect-arrival} remains a
finite-horizon first-arrival estimate.

When neither numerical branch is certified, valid-Lift target-fraction
preservation, source inheritance, every separately certified form-capacity
statement whose hypotheses hold, and the subprofile relation of
\cref{def-controlled-saturated-lift-access-envelope} remain available.  The
numerical access coordinate then takes the neutral default prescribed by
\cref{def-coordinatewise-dynamical-certificate}.  Thus the certificate split
changes no structural Lift classification and introduces no additional
source channel.

\end{theorem}

\begin{proof}

Under the hypotheses of part~\textup{(i)},
\cref{eq-controlled-fast-fiber-resolvent} gives

\[
\|E_-R_F(\lambda)E_+\|\le\sigma_\lambda.
\]

The Feshbach identity and the Neumann resolvent estimate in
\cref{eq-controlled-fast-fiber-schur-error} therefore give

\[
\left\|
P_\kappa(\lambda-L_\Omega)^{-1}U-R_X(\lambda)
\right\|
\le\Delta_{\rm fm}(\lambda).
\]

Dual pairing with \(\ell\) and \(g\) proves
\cref{eq-controlled-certified-fast-fiber-pairing}.  Insert the two projected
pairings and the two initial target atoms between
\(A_T^\Omega(\lambda)\) and \(A_T^X(\lambda)\), and apply the triangle
inequality, to obtain
\cref{eq-controlled-certified-fast-fiber-arrival}.

For part~\textup{(ii)},
\cref{eq-controlled-clocked-lift-telescoping} is an exact telescoping
identity.  Kernel and pullback contraction bound its norm by
\(\sum_t\|E_t\|\).  Pairing with the represented initial law and terminal
indicator, then adding the initialization and gate-interface errors, proves
\cref{eq-controlled-certified-clocked-defect-arrival}.  This derivation uses
only the clocked record, so it remains valid whenever the fast-fiber
hypotheses or their smallness condition are unavailable.  The statements
about simultaneous certificates and typed minima follow from applying both
valid inequalities to the same complete record.

\end{proof}


\begin{definition}[Saturated lift-access envelope]
\label{def-controlled-saturated-lift-access-envelope}

Let \(\operatorname{LiftSel}^{\rm sat}_{\mathscr L}(B,H)\) be the supremum
of the prospective selector advantage over complete records in a represented
lift class \(\mathscr L\), and let
\(\operatorname{LiftAcc}^{\rm sat}_{\mathscr L}(B,\lambda)\) be the analogous
discounted-arrival supremum.  Both are sub-suprema of the existing
\(\operatorname{Sel}^{\rm sat}\) and target-resolvent profiles, so
\begin{equation}
\label{eq-controlled-liftsel-subprofile}
\operatorname{LiftSel}^{\rm sat}_{\mathscr L}(B,H)
\le\operatorname{Sel}^{\rm sat}(B).
\end{equation}
For \(\mathscr L_{\rm int}\), discounted access transfers exactly.  For a
clocked class with \(\sum_t\|E_t\|\le\varepsilon_H\), finite-horizon access
is at most the base access plus \(\varepsilon_H\) and the interface errors.
For \(\mathscr L_{\rm fm}\), discounted access is at most the base access
plus the right side of
\cref{eq-controlled-fast-fiber-schur-error} paired with its boundary norms.

For \(\mathscr L_{\rm all}\), validity and
\cref{thm-exhaustive-lift-normal-form} provide complete classification and
source inheritance.  Its numerical access envelope remains the displayed
supremum until a uniform form, intertwining, aiming, or fiber-resolvent
profile evaluates it.

\end{definition}

\begin{example}[Product, weakly coupled, and slow fibers]
\label{ex-controlled-lift-structural-tests}

For \(\Omega=X\times F\), projection to \(X\), product initialization, and
\(L_\Omega=L_X\otimes I+I\otimes L_F\), one has \(E_+=0\) and
\(Q_\pi H_\Omega U=0\).  Pulled-back capacity, projected resolvents, and target
arrival therefore equal their base values.  If off-diagonal couplings have
norms \(\varepsilon_+\) and \(\varepsilon_-\), while the fiber decay is
\(M_Fe^{-\kappa_Ft}\), the correction has scale
\(\varepsilon_-M_F\varepsilon_+/(\lambda+\kappa_F)\).  A represented family
with fixed couplings and \(\kappa_F\downarrow0\) loses uniform smallness on
long resolvent scales \(\lambda\downarrow0\); at fixed positive \(\lambda\),
the displayed denominator remains at least \(\lambda\).
Thus target-fraction preservation alone does not supply a dynamic lift-access
bound; the relevant structural comparison is the charged fiber resolvent.

\end{example}

\section{Noisy Target Interfaces and Lifted Capacity}
\label{sec-controlled-noisy-lift-gates}
\label{subsec-controlled-noisy-lift-gates}

A valid lift transports the latent clean target.  A target or failure flag
learned from corrupted labels is an approximate boundary interface until a
represented recovery or stability theorem relates it to that clean target.

\begin{definition}[Dynamical clean-gate and capacity-stability profiles]
\label{def-controlled-noisy-gate-profiles}

Let \(T\) and \(F\) be the disjoint latent clean target and failure events
fixed by the task semantics.  For \(G\in\{T,F\}\), let
\(\widehat g_G:X\to\{0,1\}\) be a represented learned gate and let
\[
U_{{\rm gate},G}^{\rm clean}
=
\Pr_{d^b}\{\widehat g_G(X)\ne\mathbf1_G(X)\}
\le
U_{{\rm noise},G}^{\rm clean}
(n,b(n),m,\eta_{G,+},\delta_G)
\]
under a represented behavior law \(d^b\), by
\cref{eq-saturated-dynamical-clean-noise-bound}.  An exact gate has error
zero.  If the time-\(t\) deployment law satisfies
\(d_t^\pi\le W_t d^b\), define
\begin{equation}
\label{eq-controlled-deployment-gate-error}
\varepsilon_{T,F,H}^{\rm dep}
=
\min\left\{
1,\,
\sum_{t=0}^{H}W_t
\left(U_{{\rm gate},T}^{\rm clean}
     +U_{{\rm gate},F}^{\rm clean}\right)
\right\}.
\end{equation}
The events \(T,F\) are analytic comparators in this definition, not pre-hit
inputs to the learned controller.  They become prospective sources only if an
independent represented derivation constructs them within the temporal rules.
For direct independent homogeneous flip reads of both gates at deployment,
define the fallback
\begin{equation}
\label{eq-controlled-direct-flip-gate-error}
\varepsilon_{T,F,H}^{\rm flip}
=
1-(1-\eta_T)^{H+1}(1-\eta_F)^{H+1}
\le(H+1)(\eta_T+\eta_F).
\end{equation}

Capacity is measured under the reversible reference law \(\mu\), which need
not equal \(d^b\).  A capacity comparison therefore also contains a
represented density bound \(\mu\le W_{\mu\leftarrow b}d^b\), or direct
bounds on the same quantities.  Put
\begin{equation}
\label{eq-controlled-reference-gate-error}
\widehat F=\{x:\widehat g_F(x)=1\},
\qquad
\widehat T=\{x:\widehat g_T(x)=1\}\setminus\widehat F.
\end{equation}
Record
\begin{align}
\varepsilon_F^\mu
&=
\min\{1,W_{\mu\leftarrow b}U_{{\rm gate},F}^{\rm clean}\}
\ge \mu(F\mathbin{\triangle}\widehat F),
\notag\\
\varepsilon_T^\mu
&=
\min\{1,W_{\mu\leftarrow b}
(U_{{\rm gate},T}^{\rm clean}+U_{{\rm gate},F}^{\rm clean})\}
\ge \mu(T\mathbin{\triangle}\widehat T).
\label{eq-controlled-reference-boundary-errors}
\end{align}
Direct-bound branches replace the displayed products by their certified
upper bounds.  Removing \(\widehat F\) resolves learned target--failure
conflicts; because \(T\cap F=\varnothing\), the additional target error is
at most the failure-gate error.

For a represented reversible form \(\Phi\) and
\(\varepsilon_T,\varepsilon_F\in[0,1]\), define the charged joint
boundary-stability modulus
\begin{equation}
\label{eq-controlled-capacity-boundary-stability-modulus}
\omega_{\Phi,T,F}^{\rm Cap}(\varepsilon_T,\varepsilon_F)
=
\sup_{\substack{
T',F'\ {\rm represented}\\
T'\cap F'=\varnothing\\
\mu(T\mathbin{\triangle}T')\le\varepsilon_T\\
\mu(F\mathbin{\triangle}F')\le\varepsilon_F}}
\left|
\operatorname{Cap}_\Phi(T,F)
-\operatorname{Cap}_\Phi(T',F')
\right|.
\end{equation}
When the failure gate is exact, write
\(
\omega_{\Phi,T}^{\rm Cap}(\varepsilon)
=\omega_{\Phi,T,F}^{\rm Cap}(\varepsilon,0)
\).
The represented family of admissible boundary pairs, its optimization, form
evaluation, and precision belong to the same cost record.  A missing
deployment or reference-law comparison gives the neutral value one for its
probability error.  A missing boundary-stability computation gives
\(+\infty\) for capacity error.

\end{definition}

\begin{theorem}[Noise transfer through target gates and valid lifts]
\label{thm-controlled-noisy-gate-lift-transfer}

Let a clocked valid lift satisfy the hypotheses of
\cref{thm-controlled-clocked-lift-transfer}, with initialization error
\(\varepsilon_{\rm init}\), one-step defects \(E_t\), and a learned target
gate from \cref{def-controlled-noisy-gate-profiles}.  Then the clean
finite-horizon first-arrival probability and the represented lifted
first-arrival probability satisfy
\begin{equation}
\label{eq-controlled-noisy-gate-clocked-arrival}
\left|
p_{T,H}^{\rm clean}-\widehat p_{T,H}^{\Omega}
\right|
\le
\min\left\{
1,\,
\varepsilon_{\rm init}
+\sum_{t=0}^{H-1}\|E_t\|
+\varepsilon_{T,F,H}^{\rm dep}
\right\}.
\end{equation}
The direct-read version replaces the last term by
\(\varepsilon_{T,F,H}^{\rm flip}\).

For discounted killed dynamics, let \(\mu_{0,N}\) and
\(\widehat\mu_{0,N}\) be the clean and represented initial row
functionals on the compared nonterminal space, and let \(k_T\) and
\(\widehat k_T\) be the corresponding clean and learned terminal-contact
columns of \cref{def-target-resolvent}.  Assume the same-norm represented
comparisons
\[
\|\mu_{0,N}-\widehat\mu_{0,N}\|_*
\le\varepsilon_{0,\lambda},
\qquad
\|k_T-\widehat k_T\|
\le\varepsilon_{T,\lambda},
\]
and
\(
|\mu_0(T)-\widehat\mu_0(\widehat T)|
\le\varepsilon_{0,T}
\).
The contact-vector comparison must be derived from the learned gate and the
represented transition or rate operator; classification risk alone is not
substituted for this norm.  Writing \(R_{\rm eff}\) for the effective lifted
killed resolvent in \cref{eq-controlled-valid-lift-feshbach}, put
\[
A_T^{\rm clean}(\lambda)
=\mu_0(T)+\mu_{0,N}R_{\rm eff}k_T,
\qquad
\widehat A_T^\Omega(\lambda)
=\widehat\mu_0(\widehat T)
 +\widehat\mu_{0,N}R_X\widehat k_T.
\]
Then
\begin{align}
\left|A_T^{\rm clean}(\lambda)-
      \widehat A_T^\Omega(\lambda)\right|
&\le
\varepsilon_{0,T}
+\varepsilon_{0,\lambda}\|R_{\rm eff}\|\,\|k_T\|
\notag\\
&\quad
+\|\widehat\mu_{0,N}\|_*
\|R_{\rm eff}-R_X\|\,\|k_T\|
\notag\\
&\quad
+\|\widehat\mu_{0,N}\|_*\|R_X\|\,
\varepsilon_{T,\lambda}.
\label{eq-controlled-noisy-gate-discounted-arrival}
\end{align}
Exact intertwining or the fast-fiber estimate of
\cref{thm-controlled-valid-lift-block-resolvent} supplies the middle
resolvent term.  When the failure gate is learned, this same-norm term also
contains its represented killed-generator interface error; equivalently, a
resolvent identity may bound that contribution by
\(\|R_{\rm eff}\|\,\|L_N-\widehat L_{\widehat N}\|\,\|R_X\|\), after the
declared common-space identification of the two killed generators.

Finally, let a complete form-lift record satisfy
\(\mathcal E_X\le b_\pi\mathcal E_\Phi\), and use the conflict-resolved
\(\widehat T\) of \cref{eq-controlled-reference-gate-error}.  Then
\begin{equation}
\label{eq-controlled-noisy-lift-capacity-bound}
\operatorname{Cap}_\Omega(\pi^{-1}\widehat T,\pi^{-1}\widehat F)
\le
b_\pi\left[
\operatorname{Cap}_\Phi(T,F)
+\omega_{\Phi,T,F}^{\rm Cap}
 (\varepsilon_T^\mu,\varepsilon_F^\mu)
\right].
\end{equation}
If either
\begin{enumerate}[label=\textup{(\alph*)}]
\item represented target and failure scores both have uniform error smaller
      than their clean classification margins, or
\item the finite carrier has minimum atom mass
      \(\mu_*=\min_x\mu(x)>0\) and
      \(\max\{\varepsilon_T^\mu,\varepsilon_F^\mu\}<\mu_*\),
\end{enumerate}
then \(\widehat T=T\) and \(\widehat F=F\), so the original clean-boundary
lift-capacity theorem applies directly to the realized gates.  Under
condition \textup{(b)}, every represented boundary pair admitted by the
modulus with those error tolerances equals \((T,F)\); hence additionally
\(\omega_{\Phi,T,F}^{\rm Cap}
(\varepsilon_T^\mu,\varepsilon_F^\mu)=0\).
Condition \textup{(a)} alone does not force this global supremum to vanish,
because its admissible class also contains boundary pairs other than the
realized margin-certified gates.
Every noise term in the three comparisons is therefore a function of
\((n,b(n),m,\eta_{T,+},\eta_{F,+},\delta_T,\delta_F)\) through
\(U_{\rm noise}^{\rm clean}\), followed only by the displayed charged
deployment, contact-vector, and reference-law converters.  Setting
\(b(n)=n^C\) gives the polynomial-capacity specialization.

\end{theorem}

\begin{proof}

Under the behavior-to-deployment comparison, the probability of gate error
at time \(t\) is at most
\(W_t(U_{{\rm gate},T}^{\rm clean}+U_{{\rm gate},F}^{\rm clean})\).
A union bound over the stopped path gives
\cref{eq-controlled-deployment-gate-error}.  Add this event to the exact
telescoping comparison of
\cref{eq-controlled-clocked-lift-telescoping} to prove
\cref{eq-controlled-noisy-gate-clocked-arrival}.  Independent direct flips
give
\(\Pr\{\text{at least one erroneous read}\}
=1-(1-\eta_T)^{H+1}(1-\eta_F)^{H+1}\).

For the discounted statement, add the initial target-atom discrepancy and
write the noninitial terms in the target-resolvent orientation
\(\mu_{0,N}Rk_T\).  Insert and subtract
\(\widehat\mu_{0,N}R_{\rm eff}k_T\) and
\(\widehat\mu_{0,N}R_Xk_T\), then apply the dual and primal norms and the
triangle inequality.  For capacity, form comparison gives
\[
\operatorname{Cap}_\Omega(\pi^{-1}\widehat T,\pi^{-1}\widehat F)
\le b_\pi\operatorname{Cap}_\Phi(\widehat T,\widehat F),
\]
and \cref{eq-controlled-reference-boundary-errors} together with the
definition of \(\omega_{\Phi,T,F}^{\rm Cap}\) gives
\cref{eq-controlled-noisy-lift-capacity-bound}.

A uniform error below the classification margin preserves every gate sign.
For a binary gate on a finite carrier, any nonempty error set has
\(\mu\)-mass at least \(\mu_*\).  Hence
\(\max\{\varepsilon_T^\mu,\varepsilon_F^\mu\}<\mu_*\) forces both error
sets to be empty.  Both cases give
\(\widehat T=T\), \(\widehat F=F\), and the exact realized-gate conclusion.
In the atom-mass case, every admissible error set in the definition of the
modulus is also empty, so the modulus itself vanishes.  The margin case
bypasses the modulus by applying the exact theorem to its realized gates.

\end{proof}

\begin{theorem}[Fast-mode singleton target reach and effective dimension]
\label{thm-controlled-fast-mode-target-reach}

Let \(X\) have \(N\ge2\) elements, let \(\mu\) be uniform, and let
\(H_\Phi\) be a represented target-independent nonnegative self-adjoint
operator whose kernel consists exactly of the constants.  On
\(\mathring L^2(X,\mu)=\mathbf1^\perp\), define
\[
f_s
=\frac{\mathbf1_{\{s\}}-N^{-1}}
       {\sqrt{N^{-1}(1-N^{-1})}},
\qquad
\Pi_{\ge\lambda}=\mathbf1_{[\lambda,\infty)}(H_\Phi),
\qquad
d_\lambda=\operatorname{rank}\Pi_{\ge\lambda}.
\]
Then
\begin{equation}
\label{eq-controlled-fast-mode-singleton-tight-frame}
\frac1N\sum_{s\in X}
 \|\Pi_{\ge\lambda}f_s\|_{2,\mu}^{2}
=\frac{d_\lambda}{N-1}
\le\frac{2N_{\rm eff,\Phi}(\lambda)}{N-1}.
\end{equation}
Consequently, for every \(r>0\),
\begin{equation}
\label{eq-controlled-fast-mode-singleton-count}
\frac1N\left|\left\{s\in X:
\|\Pi_{\ge\lambda}f_s\|_{2,\mu}^{2}\ge r
\right\}\right|
\le\frac{d_\lambda}{(N-1)r}.
\end{equation}
If the nonzero spectrum of \(H_\Phi\) is bounded below by \(\lambda\), then
\begin{equation}
\label{eq-controlled-gap-forces-large-effective-dimension}
N_{\rm eff,\Phi}(\lambda)\ge\frac{N-1}{2}.
\end{equation}
Thus a target-independent geometry that mixes every full-carrier direction
at scale \(\lambda\) necessarily retains effective dimension of order the
carrier size.  The singleton vectors in this theorem are analytic target
probes; their evaluation does not supply a represented target source.

\end{theorem}

\begin{proof}

The centered singleton vectors form a unit-norm tight frame on
\(\mathring L^2(X,\mu)\):
\[
\frac1N\sum_{s\in X}f_s\otimes f_s
=\frac1{N-1}I_{\mathring L^2}.
\]
Taking the trace after multiplication by \(\Pi_{\ge\lambda}\) proves the
equality in \cref{eq-controlled-fast-mode-singleton-tight-frame}.  Every
eigenvalue \(u\ge\lambda\) contributes
\(u/(u+\lambda)\ge1/2\) to
\(N_{\rm eff,\Phi}(\lambda)\), whence
\(d_\lambda\le2N_{\rm eff,\Phi}(\lambda)\).  Markov's inequality applied to
the tight-frame average proves
\cref{eq-controlled-fast-mode-singleton-count}.  If every nonconstant mode
has eigenvalue at least \(\lambda\), each of the \(N-1\) such modes
contributes at least \(1/2\) to the effective dimension, proving
\cref{eq-controlled-gap-forces-large-effective-dimension}.

\end{proof}

\begin{corollary}[Dynamical envelope for nonlinear nonquotient geometry]
\label{cor-nonlinear-nonquotient-dynamical-envelope}

Let a budget-admissible nonlinear construction remain after every qualified
nonidentity quotient postprocessing has been assigned as in
\cref{prop-nonquotient-nonlinear-geometry-closure}, and suppose it is used
prospectively on the full carrier.  At the class-preserving ceiling of its
complete record, its structural source is exactly the ambient
\(\Geom/\Caus\) record selected by
\cref{cor-prospective-nonlinear-readout-routing}.  For its selected native
kernel \(P\), target \(T\), horizon \(H\), invariant law \(\mu\), and initial
density \(f_0=d\mu_0/d\mu\le R_0\), one has simultaneously
\begin{equation}
\label{eq-nonlinear-nonquotient-capacity-envelope}
\Pr_{\mu_0}\{\tau_T\le H\}
\le R_0\bigl[\mu(T)+H
\operatorname{Cap}_{\Phi_P}(T,T^c)\bigr],
\end{equation}
and every represented nonnegative target-normalized global potential \(D\)
with \(\operatorname{Osc}(D)>0\) used by its authorized current satisfies
\begin{equation}
\label{eq-nonlinear-nonquotient-drift-envelope}
\inf_{x\notin T}\frac{(I-P)D(x)}{\operatorname{Osc}(D)}
\le\frac{\mu(T)}{1-\mu(T)}.
\end{equation}
When the same native record is lazy and self-adjoint,
\cref{eq-controlled-gap-target-occupation} supplies the additional
finite-horizon occupation bound.  When it is compared to a represented
reversible geometry, \cref{eq-controlled-capacity-form-comparison-upper}
transfers its condenser bound with the declared form loss.  Its audited
functional-inequality profile supplies the mixing, smoothing, killed-set,
transport, and hypocoercive bounds of
\cref{thm-controlled-functional-inequality-consequences}; a
target-independent full-carrier operator also satisfies
\cref{eq-controlled-fast-mode-singleton-tight-frame,eq-controlled-gap-forces-large-effective-dimension}.
The right sides of
\cref{eq-nonlinear-nonquotient-capacity-envelope,eq-controlled-gap-target-occupation}
are in the same finite-horizon probability unit as \(U_{\rm hit}\) in
\cref{def-coordinatewise-dynamical-certificate}; the coordinate therefore
uses their minimum whenever they occur in the same complete record.

The committor and Bellman values of the same record remain governed by
\cref{thm-controlled-committor-bellman,thm-controlled-no-free-value}, and a
learned record retains the simultaneous Part~IV and controlled-learning
bounds in the \(U_{\rm learn}\) coordinate of
\cref{def-coordinatewise-dynamical-certificate}.  All these estimates attach
to the one already classified ambient source and its authorized current.
Neither nonlinear evaluation nor dynamical iteration creates another
source coordinate.

\end{corollary}

\begin{proof}

The nonlinear closure proposition and prospective-routing corollary assign
the complete construction either to a qualified quotient mode or to the
full-carrier ambient mode.  Under the hypothesis of the corollary, the
quotient modes have already been assigned, so its prospective structure is
the latter mode.  The selected policy and kernel are ancestors of that mode
by \cref{thm-controlled-policy-channel-decomposition}; Hodge authority is
\cref{prop-controlled-native-geometry-hodge}.

Apply \cref{eq-controlled-capacity-finite-horizon-hit} to the selected
native kernel to obtain
\cref{eq-nonlinear-nonquotient-capacity-envelope}.  Put
\[
a_D=\inf_{x\notin T}(I-P)D(x).
\]
If \(a_D\le0\), \cref{eq-nonlinear-nonquotient-drift-envelope} is immediate.
If \(a_D>0\), apply
\cref{eq-controlled-invariant-global-drift-obstruction} with \(a=a_D\).
This proves \cref{eq-nonlinear-nonquotient-drift-envelope}.
The comparison, functional-inequality, fast-mode, committor, Bellman, and
learning conclusions are the cited same-record theorems.  Temporal
admissibility and complete-cost inheritance are
\cref{def-pre-hit-temporally-admissible-source,prop-derived-channels-create-no-structure}.
Thus every displayed numerical certificate evaluates an existing
\(\Geom/\Caus\) coordinate and none changes the exhaustive structural
classification.

\end{proof}


\begin{figure}[tbp]
\centering
\resizebox{\linewidth}{!}{%
\begin{tikzpicture}[fw diagram,node distance=7mm and 8mm]
  \node[fw source,text width=2.7cm] (nonlinear)
    {nonlinear full-carrier record};
  \node[fw geom,text width=2.7cm,right=of nonlinear] (kernel)
    {selected native or comparison geometry};
  \node[fw box,text width=2.15cm,above right=6mm and 8mm of kernel] (cap)
    {capacity and boundary flux};
  \node[fw use,text width=2.15cm,right=8mm of kernel] (drift)
    {drift and occupation};
  \node[fw provenance,text width=2.15cm,below right=6mm and 8mm of kernel] (fast)
    {functional and fast modes};
  \node[fw terminal,text width=2.6cm,right=15mm of drift] (target)
    {finite-horizon target bound};
  \draw[fw arrow] (nonlinear) -- (kernel);
  \draw[fw arrow] (kernel) -- (cap);
  \draw[fw arrow] (kernel) -- (drift);
  \draw[fw arrow] (kernel) -- (fast);
  \draw[fw arrow] (cap) -- (target);
  \draw[fw arrow] (drift) -- (target);
  \draw[fw arrow] (fast) -- (target);
\end{tikzpicture}%
}
\caption{A prospective nonlinear construction that supplies no qualified
quotient remains one ambient \(\Geom/\Caus\) record.  Capacity, invariant
drift, occupation, and functional-mode estimates evaluate that same source
before its target-arrival bound is taken.}
\label{fig-nonlinear-nonquotient-dynamical-envelope}
\end{figure}

\begin{figure}[tbp]
\centering
\resizebox{\linewidth}{!}{%
\begin{tikzpicture}[fw diagram,node distance=8mm and 9mm]
  \node[fw source] (kernel) {selected kernel \(P^\pi\)};
  \node[fw use,right=of kernel] (backup) {Bellman recursion};
  \node[fw provenance,right=of backup] (value) {represented committor};
  \node[fw terminal,above right=6mm and 10mm of value] (target) {target \(T^+\)};
  \node[fw residual,below right=6mm and 10mm of value] (fail) {failure \(T^-\)};
  \node[fw geom,right=18mm of value] (cap) {Dirichlet capacity};
  \draw[fw arrow] (kernel) -- (backup);
  \draw[fw arrow] (backup) -- (value);
  \draw[fw arrow] (value) -- (target);
  \draw[fw arrow] (value) -- (fail);
  \draw[fw arrow] (target) -- (cap);
  \draw[fw arrow] (fail) -- (cap);
\end{tikzpicture}
}
\caption{Bellman iteration produces a prospective committor only through its
charged pre-transition derivation.  On a reversible carrier the same
reach--avoid function is the condenser-energy minimizer.}
\label{fig-controlled-committor-capacity}
\end{figure}

\subsection{Observation regularity and sparse-reward exploration}
\label{subsec-controlled-observation-regularity-sparse-reward}

\begin{definition}[Complete sparse-reward observation record]
\label{def-complete-sparse-reward-observation-record}

Fix a finite represented carrier \(X\), a finite target index
\(\Theta\), target sectors \(T_\theta\subseteq X\), and a refined horizon
\(H\).  A complete sparse-reward observation record is
\[
\mathcal R_{\rm cur}
=
(P_{0:H-1},\phi,d_{\mathcal O},\kappa_\varepsilon,
  M_{0:H},\widehat F_{0:H-1},b_{0:H-1},
  K_{0:H-1},V,\mathscr C_{\rm comp}),
\]
with the following data.

\begin{enumerate}[label=\textup{(\roman*)}]
\item \(P_t\) is the actual executed clocked transition kernel after the
      action and execution channels have been composed.  The public target
      reward is sparse:
      \[
      r_t^{\rm ext}=0\quad\text{on }\{\tau>t\},
      \qquad
      \tau=\inf\{t:X_t\in T_\Theta\}.
      \]
      Any terminal reward or verifier flag is revealed only at its actual
      clock time.

\item \(\phi:X\to(\mathcal O,d_{\mathcal O})\) is a represented observation
      map at declared evaluation precision.  A represented quantizer
      \(\kappa_\varepsilon:\mathcal O\to\mathcal C_\varepsilon\) gives
      \(q_\varepsilon=\kappa_\varepsilon\circ\phi\).  Put
      \[
      M_\varepsilon(\phi)
      =
      \left|q_\varepsilon(X)\right|.
      \]
      The quantizer, its cells, and \(M_\varepsilon(\phi)\) are used
      numerically only when their evaluation or certified bound is included
      in the record.

\item \(M_t\) is the represented exploration memory before transition
      \(t+1\).  It includes every retained count, visitation sketch,
      replay item, learned embedding, random-network target, predictor
      parameter, ensemble state, diversity archive, and random seed used by
      the intrinsic signal.  The represented predictor
      \(\widehat F_t\), intrinsic reward \(b_t\), and selector \(K_t\) are
      measurable functions of the public pre-hit observation history,
      \(M_t\), and fresh target-independent randomness.  Their training,
      update, selection, storage, precision, and deployment operations are
      represented before the transition they influence.

\item \(V\) is the public verifier and target interface.  A target score,
      terminal flag, later observation, future predictor error, or completed
      first-hit record is unavailable before its clock time in the sense of
      \cref{def-pre-hit-temporally-admissible-source}.

\item If \(q_\varepsilon\) is used only as an observation, no quotient
      dynamics is asserted.  If cell dynamics is used in a target-arrival
      comparison, \(\mathscr C_{\rm comp}\) is the exact or nonzero-defect
      master certificate of
      \cref{def-master-compensated-quotient-geometry-certificate}.  It
      contains the represented quotient kernel, within-fiber geometry,
      symmetric and skew defects, target vectors, stability or resolvent
      comparison, discount scale, current authority, output converter, and
      evaluation precision.
\end{enumerate}

In the resource monoid, the complete cost is
\begin{align}
\label{eq-complete-sparse-reward-record-cost}
B_{\rm cur}
={}&B_{\rm pres}\oplus B_\phi\oplus B_d\oplus B_{\kappa,\varepsilon}
\oplus B_{M_\varepsilon}\oplus B_{\rm reg}\oplus B_{\rm info}\notag\\
&\oplus B_{\rm mem}\oplus B_{\rm pred}\oplus B_{\rm train}
\oplus B_{r^{\rm ext}}\oplus B_{\rm bonus}\oplus B_{\rm pol}\notag\\
&\oplus B_{\rm exec}\oplus B_{\rm clock}\oplus B_{\rm verify}
\oplus B_{\rm out}\oplus B_{\rm cap}\oplus B_{\rm comp},\\
\label{eq-complete-curiosity-regularity-cost}
B_{\rm reg}
={}&B_\mu\oplus B_P\oplus B_{\mathcal E_\phi}
\oplus B_{\operatorname{Var}_\phi}\oplus B_{\widehat{\mathcal E}_\phi}
\oplus B_{\widehat{\operatorname{Var}}_\phi}
\oplus B_{\Gamma_E}\oplus B_{\Gamma_V},\\
\label{eq-complete-curiosity-information-cost}
B_{\rm info}
={}&B_{\mathsf Z_0}\oplus B_{C_\varepsilon}
\oplus B_{\rm reveal}\oplus B_{\rm channel}
\oplus B_{\rm SDPI}\oplus B_{J^{\rm fib}}
\oplus B_{\beta^{\rm fib}}\oplus B_{\rm conf},\\
\label{eq-complete-curiosity-capacity-cost}
B_{\rm cap}
={}&B_{P^{\rm ref}}\oplus B_{\mu^{\rm ref}}
\oplus B_{\widetilde\mu}\oplus B_{r_\pm}
\oplus B_{\kappa_\pm}\oplus B_{\operatorname{Cap}}
\oplus B_{\rho_0}\oplus B_{f_0}\oplus B_p,\\
\label{eq-complete-compensated-observation-cost}
B_{\rm comp}
={}&B_q\oplus B_{\rm fib}\oplus B_{L_q}
\oplus B_{E_q^{\mathsf s}}\oplus B_{E_q^{\mathsf a}}
\oplus B_{\Phi_Q}\oplus B_{\Phi_\perp}\notag\\
&\oplus B_{D_Q}\oplus B_{J_Q}\oplus B_{\mathfrak K_q}
\oplus B_{\rm stab}\oplus B_{\rm res}
\oplus B_{\overline A_q}\oplus B_{S_q}\notag\\
&\oplus B_{M_{\mathfrak G_q}}\oplus B_{C_{\mathfrak G_q}}
\oplus B_{R_{\mathfrak G_q}}\oplus B_{\rho_{\mathfrak G_q}}
\oplus B_{\rm targ}\oplus B_{\rm auth}\notag\\
&\oplus B_{\rm prec}\oplus B_{\rm select}\oplus B_{\rm compare}.
\end{align}
An absent slot has zero resource charge only when it is not used.  Reuse,
adaptive updates, and repeated evaluations accumulate according to
\cref{def-budget-constructor-accounting}.
The regularity line charges the invariant law and executed kernel, population
and empirical energy and variance, and both simultaneous deviation
certificates.  The information line charges the initial record, target-cell
comparator, ordered reveals and their channels, contraction certificates,
within-cell information and neutral contact, and confidence bookkeeping.
The capacity line charges the reference kernel and law, selected invariant
law, density and conductance comparisons, condenser-capacity evaluation,
initial law and density, and the selected Hölder exponents.  The slots in
the final display charge, respectively, the support and fiber
records, quotient generator and both defects, quotient and within-fiber
forms, potential, current, causal provenance, stability and resolvent
records, quotient utility and dynamic factor, all four geometric visibility
factors, target interface, authority, precision, certificate selection, and
the final same-normalization comparison.  Shared ancestors are represented
once in the derivation DAG; an additional evaluation is not free merely
because its input was already stored.

Suppose \(P\) is a represented kernel with invariant law \(\mu\).  The
metric observation variance, Dirichlet energy, and regularity ratio are
\begin{align}
\label{eq-observation-metric-variance}
\operatorname{Var}_{\mu}^{d}(\phi)
&=
\frac12\sum_{x,z\in X}\mu(x)\mu(z)
d_{\mathcal O}(\phi(x),\phi(z))^2,\\
\label{eq-observation-metric-dirichlet-energy}
\mathcal E_{P}^{d}(\phi)
&=
\frac12\sum_{x,z\in X}\mu(x)P(x,z)
d_{\mathcal O}(\phi(x),\phi(z))^2,\\
\label{eq-observation-regularity-ratio}
\rho_\phi(P)
&=
\begin{cases}
\mathcal E_P^d(\phi)/\operatorname{Var}_\mu^d(\phi),
&\operatorname{Var}_\mu^d(\phi)>0,\\
0,&\operatorname{Var}_\mu^d(\phi)=0.
\end{cases}
\end{align}
These are analytic coordinates of the represented record.  A numerical use
requires their represented evaluation or a certified upper bound at the
cost in \cref{eq-complete-sparse-reward-record-cost}.

\end{definition}

\begin{theorem}[Observation regularity, fiber information, and sparse-reward
exploration]
\label{thm-observation-regularity-sparse-reward-exploration}

Let \(\mathcal R_{\rm cur}\) be a complete record from
\cref{def-complete-sparse-reward-observation-record} of cost
\(B_{\rm cur}\preceq B\).  Assume that \(\Theta\) is uniform on
\(\{1,\ldots,N\}\) and that the indexed target sectors are the disjoint
singletons \(T_\theta=\{x_\theta\}\).  Define
\[
C_\varepsilon=q_\varepsilon(x_\Theta).
\]
At time \(t\), let \(\mathcal H_t^{\rm obs}\) be the complete pre-hit
observation, memory, predictor, intrinsic-reward, action, and execution
history available before \(K_t\) is sampled.  Give the analytic comparator
the additional cell label \(C_\varepsilon\), and put
\begin{equation}
\label{eq-observation-within-cell-information}
J_{R,t}^{\rm fib}
=
I\!\left(
\Theta;\mathcal H_t^{\rm obs}
\mid\mathsf Z_0,C_\varepsilon
\right).
\end{equation}
Equivalently, an ordered-reveal record may replace the right side by the
certified upper bound
\[
J_{R,t}^{\rm fib,\uparrow}
=
\sum_{j=1}^{L_{R,t}}
\mathbb E\min\{c_{R,t,j},
\vartheta_{R,t,j}a_{R,t,j}\}
\]
from \cref{thm-pre-hit-sequential-sdpi-information-contraction}, after
conditioning every channel slot on
\((\mathsf Z_0,C_\varepsilon)\).
Every target-dependent coordinate available at time zero is included in
\(\mathsf Z_0\); every target-dependent coordinate revealed later is
included in one of the ordered channel slots.  Thus the displayed sum cannot
omit information used by the selector.
Throughout the theorem, \(J_{R,t}^{\rm fib,\uparrow}\) denotes either this
certified upper bound or the exactly evaluated and charged value
\(J_{R,t}^{\rm fib}\).  Its neutral value is \(+\infty\) when neither is
available.

Let \(\mathbb Q_{R,t}^{\rm fib}\) be the law obtained by retaining
\((\mathsf Z_0,C_\varepsilon,\mathcal H_t^{\rm obs})\) and resampling
\(\Theta\) from
\(\mathcal L(\Theta\mid\mathsf Z_0,C_\varepsilon)\), independently of the
observation history.  Define the within-cell neutral contact
\begin{equation}
\label{eq-observation-within-cell-neutral-contact}
\beta_{R,t}^{\rm fib}
=
\mathbb Q_{R,t}^{\rm fib}
\{\tau>t,\ X_{t+1}\in T_\Theta\}.
\end{equation}
Put
\[
p_{R,t}
=
\mathbb P_{R,t}\{\tau>t,\ X_{t+1}\in T_\Theta\}.
\]
Then the survival-weighted next-contact probability \(p_{R,t}\) and the
selector advantage satisfy
\begin{align}
\label{eq-observation-fiber-information-contact-bound}
\bigl(\operatorname{Adv}_t(K_t;\nu_t)\bigr)_+
&\le p_{R,t}\notag\\
&\le
\min\!\left\{1,\,
\beta_{R,t}^{\rm fib}
+\sqrt{\frac{\log2}{2}J_{R,t}^{\rm fib,\uparrow}}
\right\}.
\end{align}
On the survival sector the external reward history is deterministic, hence
\begin{equation}
\label{eq-sparse-external-reward-zero-information}
I\!\left(\Theta;r_{0:t}^{\rm ext}
\mid\mathsf Z_0,C_\varepsilon,\tau>t\right)=0.
\end{equation}
The intrinsic rewards, predictor states, and exploration memories are
causal postprocessings already contained in
\(\mathcal H_t^{\rm obs}\); data processing gives them no additional
information coordinate beyond \(J_{R,t}^{\rm fib,\uparrow}\), while their
construction and deployment costs remain in \(B_{\rm cur}\).
Suppose, in addition, that the initial record is target-permutation-neutral,
meaning
\[
\mathbb P\{\Theta=\theta\mid\mathsf Z_0=z\}=\frac1N
\quad(1\le\theta\le N)
\]
for every realized \(z\), and that, conditional on
\((\mathsf Z_0,C_\varepsilon)\), the only target-index restriction is
membership in the displayed observation fiber.  If
\(M_\varepsilon(\phi)\) is bounded pointwise by the represented value, then
\begin{equation}
\label{eq-observation-cell-count-neutral-contact}
\beta_{R,t}^{\rm fib}
\le
\frac{M_\varepsilon(\phi)}{N},
\end{equation}
and therefore
\begin{equation}
\label{eq-observation-cell-count-contact-bound}
\bigl(\operatorname{Adv}_t(K_t;\nu_t)\bigr)_+
\le
\min\!\left\{1,\,
\frac{M_\varepsilon(\phi)}{N}
+\sqrt{\frac{\log2}{2}J_{R,t}^{\rm fib,\uparrow}}
\right\}.
\end{equation}
For a horizon \(H\), summing the disjoint first-contact events gives
\begin{equation}
\label{eq-observation-cell-count-cumulative-hit}
\Pr\{\tau\le H\}
\le
\min\!\left\{1,\,
\sum_{t=0}^{H-1}
\left(
\beta_{R,t}^{\rm fib}
+\sqrt{\frac{\log2}{2}J_{R,t}^{\rm fib,\uparrow}}
\right)
\right\}.
\end{equation}
If, conditional on each target index \(\theta\), the same complete record
contains the reversible density, conductance, capacity, and initialization
comparison of \cref{thm-controlled-density-capacity-hit-transfer}, then,
with \(q=p/(p-1)\), it also supplies
\begin{equation}
\label{eq-observation-regularity-capacity-hit-bound}
U_{R}^{\rm cur,cap}(H)
=
\max_{1\le\theta\le N}
\min\!\left\{1,
\|f_{0,\theta}\|_{L^p(\widetilde\mu_\theta)}
\left[
\min\!\left\{1,
r_{\theta,+}\mu_\theta(T_\theta)
+H\kappa_{\theta,+}
\operatorname{Cap}_{\Phi_{P_\theta}}(T_\theta,T_\theta^c)
\right\}
\right]^{1/q}
\right\}.
\end{equation}
For that record,
\begin{equation}
\label{eq-observation-information-capacity-minimum}
\Pr\{\tau\le H\}
\le
\min\!\left\{
U_R^{\rm cur,cap}(H),
\sum_{t=0}^{H-1}
\left(
\beta_{R,t}^{\rm fib}
+\sqrt{\frac{\log2}{2}J_{R,t}^{\rm fib,\uparrow}}
\right)
\right\},
\end{equation}
with the second candidate clipped at one.  Every comparison used in
\cref{eq-observation-regularity-capacity-hit-bound} carries the cost
\cref{eq-complete-curiosity-capacity-cost}, including evaluation or
certification of the displayed maximum.

The regularity coordinate controls the magnitude of a one-step intrinsic
signal.  If, under a stationary transition \((X_0,X_1)\sim\mu(dx)P(x,dz)\),
\[
0\le b_\phi(X_0,X_1)
\le a_\phi+L_\phi
d_{\mathcal O}(\phi(X_0),\phi(X_1)),
\]
then
\begin{equation}
\label{eq-observation-regularity-linear-novelty-bound}
\mathbb E_{\mu P}b_\phi
\le
a_\phi
+L_\phi\sqrt{2\mathcal E_P^d(\phi)}
=a_\phi
+L_\phi\sqrt{
2\rho_\phi(P)\operatorname{Var}_\mu^d(\phi)}.
\end{equation}
If instead
\(
b_\phi\le a_\phi+L_\phi d_{\mathcal O}^2
\),
then
\begin{equation}
\label{eq-observation-regularity-quadratic-novelty-bound}
\mathbb E_{\mu P}b_\phi
\le
a_\phi+2L_\phi\mathcal E_P^d(\phi)
=a_\phi
+2L_\phi\rho_\phi(P)\operatorname{Var}_\mu^d(\phi).
\end{equation}
If these quantities are estimated from a represented trajectory sample and
the same record supplies simultaneous deviations
\begin{equation}
\label{eq-observation-regularity-generalization-event}
\mathcal E_P^d(\phi)
\le\widehat{\mathcal E}_P^d(\phi)+\Gamma_E,
\qquad
\operatorname{Var}_\mu^d(\phi)
\ge\widehat{\operatorname{Var}}_\mu^d(\phi)-\Gamma_V,
\end{equation}
then, on that event,
\begin{align}
\label{eq-observation-regularity-learned-ratio}
\rho_\phi(P)
&\le
\frac{\widehat{\mathcal E}_P^d(\phi)+\Gamma_E}
     {(\widehat{\operatorname{Var}}_\mu^d(\phi)-\Gamma_V)_+}
\quad\text{when the denominator is positive},\\
\label{eq-observation-regularity-learned-novelty}
\mathbb E_{\mu P}b_\phi
&\le
a_\phi+L_\phi
\sqrt{2(\widehat{\mathcal E}_P^d(\phi)+\Gamma_E)}
\end{align}
for the linear novelty bound, with
\(2L_\phi(\widehat{\mathcal E}_P^d(\phi)+\Gamma_E)\) replacing the square
root term in the quadratic case.  The deviations
\((\Gamma_E,\Gamma_V)\) are the applicable simultaneous structural
Rademacher, PAC--Bayes, mixing-block, or martingale bounds from
\cref{part:learning,thm-beta-mixing-structural-blocking,%
thm-martingale-structural-pac-bayes}; their sample size, confidence,
effective dimension, selection penalty, trajectory dependence, and
evaluation cost remain in \(\mathcal R_{\rm cur}\).
Thus observation regularity bounds the available novelty magnitude, while
\cref{eq-observation-fiber-information-contact-bound} bounds target
selection.  Both coordinates remain in the same complete record.

If \(q_\varepsilon\) exactly intertwines the selected dynamics and satisfies
the target-qualification gates, its cell dynamics is assigned to the exact
\(\Abs/\Geom\) mode.  If it does not intertwine and a cell-level target
comparison is invoked, the comparison is admissible only through
\(\mathscr C_{\rm comp}\).  Its dynamic factor is
\begin{equation}
\label{eq-observation-compensated-dynamic-factor}
S_{q_\varepsilon}(\lambda)
=
\begin{cases}
1-\vartheta_\lambda^T(q_\varepsilon),
&\text{under the reversible comparison},\\
\max\{0,1-\Delta_\lambda(q_\varepsilon)\},
&\text{under the general resolvent comparison},\\
0,&\text{without either represented comparison},
\end{cases}
\end{equation}
and its source visibility is exactly the applicable master-certificate
value from
\cref{def:master-certificate-dynamic-qualification,def:master-certificate-visibility}.
Every term used to build or evaluate this compensated geometry is charged
by \cref{eq-complete-compensated-observation-cost}.  A nonintertwining
observation with no such certificate remains an ambient observation record;
its set-theoretic cells supply no quotient-dynamical bound.

Finally, at the common class-preserving ceiling
\[
B^{\rm cur}
=
\Psi_{\rm cur}(B,H,n,\varepsilon,\mathbf m),
\]
where \(\Psi_{\rm cur}\) is a declared class-preserving common majorant of
all slots in
\cref{eq-complete-sparse-reward-record-cost,eq-complete-compensated-observation-cost}
and of the ordered-reveal sample counts \(\mathbf m\).  It does not enlarge
the public constructor algebra.  Then
the saturated curiosity-selector subprofile obeys
\begin{equation}
\label{eq-saturated-curiosity-selector-envelope}
\sup_{\substack{0\le t<H\\
\mathcal R_{\rm cur}:\,B_{\rm cur}\preceq B}}
\bigl(\operatorname{Adv}_t(K_t;\nu_t)\bigr)_+
\le
\sup_{\substack{0\le t<H\\
\mathcal R_{\rm cur}:\,B_{\rm cur}\preceq B^{\rm cur}}}
\min\!\left\{1,\,
\beta_{R,t}^{\rm fib}
+\sqrt{\frac{\log2}{2}J_{R,t}^{\rm fib,\uparrow}}
\right\},
\end{equation}
where every displayed quantity belongs to the same complete pre-hit record.
Independent same-record capacity, Blackwell, source, Caus, and Lift bounds
may subsequently be minimized with this coordinate under
\cref{thm-saturated-same-record-prospective-certificate-envelope}; they are
not used to prove the present inequality.

\end{theorem}

\begin{proof}

Conditional on the complete observation history, \(K_t\) uses only
target-independent fresh randomness.  Hence it is a causal postprocessing
of the experiment
\((\mathsf Z_0,C_\varepsilon,\mathcal H_t^{\rm obs})\).  Under
\(\mathbb Q_{R,t}^{\rm fib}\), the target index and observation history are
conditionally independent given \((\mathsf Z_0,C_\varepsilon)\).  The
conditional relative-entropy identity gives
\[
\operatorname{KL}
(\mathbb P_{R,t}\Vert\mathbb Q_{R,t}^{\rm fib})
=(\log2)I(\Theta;\mathcal H_t^{\rm obs}
\mid\mathsf Z_0,C_\varepsilon).
\]
Map the joint record to the next-contact indicator and apply data
processing followed by Pinsker's inequality.  This yields
\[
(p_{R,t}-\beta_{R,t}^{\rm fib})_+
\le
\sqrt{\frac{\log2}{2}J_{R,t}^{\rm fib,\uparrow}}.
\]
The neutral selector has nonnegative target-contact mass, so
\((\operatorname{Adv}_t)_+\le p_{R,t}\), proving
\cref{eq-observation-fiber-information-contact-bound}.
Equation \cref{eq-sparse-external-reward-zero-information} follows because
\(r_0^{\rm ext},\ldots,r_t^{\rm ext}\) are the constant zero vector on
\(\{\tau>t\}\).  The intrinsic-record assertion is the data-processing
inequality applied to the causal constructors in
\cref{def-complete-sparse-reward-observation-record}.

Under target-permutation neutrality, fix an initial record \(z\).
Conditional on a target cell \(c\), \(\Theta\) is uniform on
\[
F_c(z)=\{\theta:q_{\varepsilon,z}(x_\theta)=c\}.
\]
Even after giving \(c\) to the decoupled selector, its conditional exact
target probability is at most \(1/|F_c(z)|\).  Averaging over the
conditionally uniform
\(\Theta\) gives
\[
\frac1N\sum_{c:F_c(z)\ne\varnothing}
|F_c(z)|\frac1{|F_c(z)|}
=\frac{|\{c:F_c(z)\ne\varnothing\}|}{N}
\le\frac{M_\varepsilon(\phi)}{N}.
\]
The bound is pointwise in \(z\), so averaging over \(\mathsf Z_0\) preserves
it.
This proves
\cref{eq-observation-cell-count-neutral-contact,eq-observation-cell-count-contact-bound}.
Summing the disjoint first-contact events proves
\cref{eq-observation-cell-count-cumulative-hit}.
The independent capacity estimate is
\cref{eq-controlled-density-lp-hit-transfer}, conditional on each
\(\Theta=\theta\), for the same selected kernel, target, horizon, and
initialization.  Averaging the conditional probabilities is bounded by the
displayed maximum.  Taking the smaller of two upper bounds for the same
unconditional event proves
\cref{eq-observation-information-capacity-minimum}; its listed construction
costs are precisely \cref{eq-complete-curiosity-capacity-cost}.

Cauchy--Schwarz and
\(\mathbb E_{\mu P}d_{\mathcal O}^2=2\mathcal E_P^d(\phi)\)
prove \cref{eq-observation-regularity-linear-novelty-bound}.  Direct
expectation proves
\cref{eq-observation-regularity-quadratic-novelty-bound}.

The exact and nonzero-defect cell-dynamics assignments are
\cref{def-master-compensated-quotient-geometry-certificate,%
def:master-certificate-dynamic-qualification}.  Their complete construction
cost is the resource accumulation in
\cref{eq-complete-sparse-reward-record-cost,eq-complete-compensated-observation-cost}.
The absence of an unrepresented quotient-dynamical comparison is
\cref{cor-prospective-nonlinear-readout-routing,%
thm-no-uncharged-prospective-value-amplification}.
Finally, every budget-\(B\) record on the left belongs to the class at the
common majorant \(B^{\rm cur}\).  Apply
\cref{eq-observation-fiber-information-contact-bound} record by record and
then take the displayed supremum.  This proves
\cref{eq-saturated-curiosity-selector-envelope} without invoking a
downstream universal certificate.

\end{proof}

\begin{corollary}[Resolution--plateau tradeoff]
\label{cor-resolution-plateau-tradeoff}

Under the target-permutation-neutrality hypothesis of
\cref{thm-observation-regularity-sparse-reward-exploration}, suppose
\[
M_\varepsilon(\phi)\le N^\beta,
\qquad 0\le\beta\le1,
\qquad
J_{R,t}^{\rm fib,\uparrow}\le J_{\rm fib}
\quad(0\le t<H).
\]
Then
\begin{equation}
\label{eq-resolution-plateau-tradeoff}
\Pr\{\tau\le H\}
\le
\min\!\left\{1,\,
H\left(
N^{-(1-\beta)}
+\sqrt{\frac{\log2}{2}J_{\rm fib}}
\right)
\right\}.
\end{equation}
For \(N=2^n\), the resolution term is
\(2^{-(1-\beta)n}\).  The value of \(\beta\) is the evaluated logarithmic
image size
\[
\beta
=
\frac{\log M_\varepsilon(\phi)}{\log N};
\]
it is not inferred from representation cost.  In particular, the
represented identity observation has cost linear in a binary state
description and has \(M_\varepsilon(\phi)=N\) at exact resolution.

\end{corollary}

\begin{proof}

Substitute
\(
M_\varepsilon(\phi)/N\le N^{-(1-\beta)}
\)
and the uniform information bound into
\cref{eq-observation-cell-count-cumulative-hit}.

\end{proof}

\begin{figure}[tbp]
\centering
\resizebox{\linewidth}{!}{%
\begin{tikzpicture}[fw diagram,node distance=7mm and 9mm]
  \node[fw source,text width=2.8cm] (obs)
    {represented observation\\metric and resolution};
  \node[fw provenance,text width=2.8cm,right=of obs] (nov)
    {memory, predictor,\\novelty and policy};
  \node[fw geom,text width=2.8cm,above right=7mm and 10mm of nov] (reg)
    {regularity\\\(\rho_\phi(P)\)};
  \node[fw box,text width=2.8cm,right=12mm of nov] (fiber)
    {cell resolution and\\fiber information};
  \node[fw provenance,text width=3.0cm,below right=7mm and 10mm of nov] (comp)
    {exact or compensated\\cell dynamics};
  \node[fw terminal,text width=2.8cm,right=16mm of fiber] (hit)
    {same-record\\target-contact bound};
  \draw[fw arrow] (obs)--(nov);
  \draw[fw arrow] (nov)--(reg);
  \draw[fw arrow] (nov)--(fiber);
  \draw[fw arrow] (nov)--(comp);
  \draw[fw arrow] (reg)--(hit);
  \draw[fw arrow] (fiber)--(hit);
  \draw[fw arrow] (comp)--(hit);
\end{tikzpicture}%
}
\caption{Sparse-reward observation accounting.  Observation evaluation,
novelty learning, memory, policy deployment, and the executed kernel are
charged before use.  Nonintertwining cells enter a dynamical comparison only
through the fully costed compensated certificate.  Resolution, conditional
fiber information, and regularity are evaluated on the same record before
the target-contact envelope is saturated.}
\label{fig-observation-regularity-sparse-reward}
\end{figure}

\section{Optimal Structural Active Sampling}
\label{sec-optimal-structural-active-sampling}

Active sampling combines the choice of a public experiment with the choice
of the next state transition.  The experiment may reveal new data only
through a declared presentation interface.  The transition remains a
selected controlled kernel and is therefore governed by the master
quotient--geometry certificate.  This section optimizes target arrival over
the resulting completely represented records.  Information gain is one
upper certificate for that objective, rather than its definition.

\subsection{Complete active-acquisition records and cost}
\label{subsec-complete-active-acquisition-records}

\begin{definition}[Complete master active-sampling record]
\label{def-complete-master-active-sampling-record}

Fix a finite target index \(\Theta\), initial public record \(\mathsf Z_0\),
finite refined horizon \(H\), and resource ceiling \(B\).  At a surviving
pre-hit history \(\mathcal H_t\), a complete master active-sampling record
contains the following represented data.

\begin{enumerate}[label=\textup{(\roman*)}]
\item A predictable action
      \(A_t=(A_t^{\rm qry},A_t^{\rm tr})\).  The query component selects a
      declared experiment, feature, or public-data readout.  The transport
      component selects the next native or lifted transition.  Conditional
      on \(\mathcal H_t\), the action uses only represented fresh randomness
      independent of \(\Theta\).

\item The selected history-conditional experiment
      \(K_{t}^{\rm obs}(dy\mid u,\mathcal H_t,A_t^{\rm qry})\), its clean
      input \(U_t\), returned observation \(Y_{t+1}\), and every physical
      channel state.  An interface not generated by the declared
      presentation is an extension in the sense of
      \cref{thm-oracle-admissibility-presentation-separation}.

\item The executed transition kernel
      \(P_t^{A_t^{\rm tr}}\), verifier interface, stopping rule, and output
      converter.  Every observation used by a later transition occurs in
      the pre-hit filtration before that transition.

\item A dynamically qualified master certificate
      \[
      \mathscr C_t
      =(
      \mathfrak N_t,q_t,\mathfrak G_{q_t},
      \mathfrak S_t,\mathfrak K_{q_t})
      \]
      for the selected transition, together with a represented
      same-normalization comparison to the target-arrival coordinate when
      structural visibility is used numerically.

\item The acquisition score, its statistical or analytic certificate, the
      represented rule that selects or approximately selects the action,
      belief or sufficient-record update, memory, clock, precision, and all
      capacity, information, Blackwell, learning, Lift, and output
      comparisons invoked by the record.
\end{enumerate}

In the resource monoid its complete cost is
\begin{align}
\label{eq-complete-master-active-sampling-cost}
B_{\rm act}
={}&B_{\mathsf Z_0}\oplus B_{\rm hist}
\oplus B_{\rm query}\oplus B_{\rm score}
\oplus B_{\rm optimize}\oplus B_{\rm obs}
\oplus B_{\rm channel}\notag\\
&\oplus B_{\rm belief}\oplus B_{\rm memory}
\oplus B_{\rm kernel}\oplus B_{\mathscr C}
\oplus B_{\rm exec}\oplus B_{\rm clock}\notag\\
&\oplus B_{\rm verify}\oplus B_{\rm stop}
\oplus B_{\rm out}\oplus B_{\rm info}
\oplus B_{\rm cap}\oplus B_{\rm BW}
\oplus B_{\rm learn}.
\end{align}
The master-certificate charge \(B_{\mathscr C}\) is the complete charge in
\cref{def:master-certificate-visibility}.  A displayed semantic maximizer
of an acquisition score supplies no action rule.  Deployment of that
maximizer requires the represented selection or approximation record charged
by \(B_{\rm optimize}\).

\end{definition}

The master certificate already includes the pure geometric regime.
Indeed, \(q=\mathrm{id}\) with an active geometry slot is
\(\mathcal C_X\), whereas nonidentity exact and nonzero-defect records are
\(\mathcal C_{Q,\mathrm{ex}}\), \(\mathcal C_{Q,\mathrm{rev}}\), and
\(\mathcal C_{Q,\mathrm{nr}}\).  Thus compensated quasi-lumpability is the
general case and full-carrier geometry is its identity-support realization.

\begin{definition}[Optimal structural active-sampling profiles]
\label{def-optimal-structural-active-sampling-profile}

Let \(\mathscr R_{I,H}^{\rm act}(B)\) be the complete records of
\cref{def-complete-master-active-sampling-record} with
\(B_{\rm act}\preceq B\).  For \(R\in\mathscr R_{I,H}^{\rm act}(B)\), set
\[
A_R^H=\Pr_R\{\tau\le H\},
\qquad
a_{R,t}=\bigl(\operatorname{Adv}_t(K_{R,t};\nu_{R,t})\bigr)_+.
\]
Define
\begin{align}
\label{eq-optimal-structural-active-arrival-profile}
\operatorname{ActArr}_{I,H}^{\rm sat}(B)
&=\sup_{R\in\mathscr R_{I,H}^{\rm act}(B)}A_R^H,\\
\label{eq-optimal-structural-active-selector-profile}
\operatorname{ActSel}_{I,H}^{\rm sat}(B)
&=\max_{0\le t<H}
  \sup_{R\in\mathscr R_{I,H}^{\rm act}(B)}a_{R,t},
\end{align}
with supremum zero for an empty record class.  These are analytic saturated
profiles.  A deployed optimal policy is present only when its action selector
is represented in the record.

\end{definition}

\begin{definition}[Family-uniform active saturated profile]
\label{def-family-uniform-active-saturated-profile}

Let \(\mathfrak I_{n,m,\zeta}\) be a presented size-and-sample slice with
declared channel parameter \(\zeta\), and let \(\mathfrak M_{n,m,\zeta}\)
be its fixed family functional.  In the homogeneous binary-noise case,
\(\zeta=\eta\in[0,1/2)\).  Write
\(\mathscr R_{\mathfrak I,n,m,\zeta,H}^{\rm act}(B)\) for the complete
records of \cref{def-complete-master-active-sampling-record} produced by one
family-uniform represented derivation scheme and having complete saturated
cost at most \(B\) on the whole slice.  Define
\begin{equation}
\label{eq-family-uniform-active-arrival-profile}
\operatorname{ActArr}_{\mathfrak M,n,m,\zeta,H}^{\rm sat}(B)
=
\sup_{R\in\mathscr R_{\mathfrak I,n,m,\zeta,H}^{\rm act}(B)}
\mathfrak M_{n,m,\zeta}\!\left(I\mapsto A_{R,I}^{H}\right).
\end{equation}
For a master mode
\[
c\in\mathfrak C_{\rm master}
:=
\{X,A,Q_{\rm ex},Q_{\rm rev},Q_{\rm nr}\},
\]
let
\(\operatorname{ActArr}_{\mathfrak M,n,m,\zeta,H;c}^{\rm sat}(B)\)
denote the same supremum restricted, respectively, to
\[
\mathcal C_X,\quad \mathcal C_A,\quad \mathcal C_{Q,\mathrm{ex}},
\quad \mathcal C_{Q,\mathrm{rev}},\quad \mathcal C_{Q,\mathrm{nr}}.
\]

The active restriction of the saturated success profile is the supremum in
\cref{eq-uniform-affordable-success-profile} over these same complete
active-interface derivations, with null queries admitted.  Denote it by
\(\operatorname{Succ}_{\mathfrak M,n,m,\zeta,H}^{\rm act,sat}(B)\).
This restriction changes no cost convention: the action rule, query,
observation, transport, certificate, stopping rule, and output are the
single derivation whose cost is displayed in
\cref{eq-complete-master-active-sampling-cost}.

\end{definition}

\begin{theorem}[Exact resource-valued active-acquisition recursion]
\label{thm-budgeted-active-acquisition-recursion}

For a represented surviving history \(h\) at time \(t\), let
\(\mathscr S_{t,H}^{\rm act}(h;B)\) be the complete represented suffix
records of horizon \(H-t\) and cost at most \(B\).  Decompose a nonterminal
suffix record at its first transition into its represented action \(a\) and
its
represented continuation rule \(\varrho\), which maps each returned
observation and next history to the remaining suffix record.  Let
\(c_t^{\rm full}(a,\varrho;h)\) be the complete resource-monoid cost of this
combined record, including the construction and storage of \(\varrho\).
Define
\[
V_t^{\rm act}(h;B)
=
\sup_{R\in\mathscr S_{t,H}^{\rm act}(h;B)}A_R^{H-t}.
\]
Then the exact finite-horizon value satisfies
\begin{equation}
\label{eq-optimal-structural-active-resource-recursion}
V_t^{\rm act}(h;B)
=
\sup_{\substack{(a,\varrho)\text{ represented}\\
                 c_t^{\rm full}(a,\varrho;h)\preceq B}}
\left\{
p_T(h,a)
+\mathbb E\!\left[
\mathbf1_{\{\tau>t+1\}}
A_{\varrho(H_{t+1})}^{H-t-1}
\mid h,a
\right]
\right\},
\end{equation}
with the declared terminal conversion and \(V_H^{\rm act}=0\) before that
conversion.  Here \(p_T(h,a)\) is the one-step first-contact probability.

If a declared scalar cost is stage-decomposable,
\[
\kappa(c_t^{\rm full}(a,\varrho;h))
=c_t^\kappa(a,h)+\kappa(\operatorname{Cost}(\varrho)),
\]
then \cref{eq-optimal-structural-active-resource-recursion} is equivalently
the same recursion with the scalar constraint
\(c_t^\kappa(a,h)+\kappa(\operatorname{Cost}(\varrho))\le b\).
When a represented continuation optimizer simultaneously realizes the
suffix suprema at every returned history, this becomes the familiar form
\begin{equation}
\label{eq-optimal-structural-active-bellman}
V_t^{\rm act}(h,b)
=
\sup_{\substack{a\in\mathcal A_t(h)\text{ represented}\\
                 c_t^\kappa(a,h)\le b}}
\left\{
p_T(h,a)
+\mathbb E\!\left[
\mathbf1_{\{\tau>t+1\}}
V_{t+1}^{\rm act}(H_{t+1},b-c_t^\kappa(a,h))
\mid h,a
\right]
\right\}.
\end{equation}
The continuation optimizer and every comparison used to construct it are
part of the complete record.

\end{theorem}

\begin{proof}

Every nonterminal suffix execution has one first represented action and a
represented rule for its continuations after the returned observation.
Conversely, concatenating such an action and continuation rule produces one
suffix record, and computation closure charges their combined record by
\(c_t^{\rm full}\).  First contact and survival continuation are disjoint,
which proves
\cref{eq-optimal-structural-active-resource-recursion} by taking the
supremum over the two equivalent descriptions of the same suffix-record
class.  The scalar constraint follows from stage decomposability.  A single
represented continuation optimizer permits the suffix-record supremum to be
taken inside the conditional expectation, proving
\cref{eq-optimal-structural-active-bellman}.  Without that represented
selector, the resource-valued record recursion remains the exact statement.

\end{proof}

\begin{theorem}[Active optimality equals affordable saturated construction]
\label{thm-active-optimality-saturated-construction}

At every declared \((n,m,\zeta,H,B)\),
\begin{equation}
\label{eq-active-optimality-saturated-success-identity}
\operatorname{ActArr}_{\mathfrak M,n,m,\zeta,H}^{\rm sat}(B)
=
\operatorname{Succ}_{\mathfrak M,n,m,\zeta,H}^{\rm act,sat}(B).
\end{equation}
The common value has the exact structural decomposition
\begin{equation}
\label{eq-active-optimality-five-mode-arrival-decomposition}
\operatorname{ActArr}_{\mathfrak M,n,m,\zeta,H}^{\rm sat}(B)
=
\max_{c\in\mathfrak C_{\rm master}}
\operatorname{ActArr}_{\mathfrak M,n,m,\zeta,H;c}^{\rm sat}(B).
\end{equation}
Thus optimizing active acquisition is exactly optimizing the target-arrival
coordinate of the complete saturated profile over affordable constructions;
the construction and the selected master coordinate are carried by the same
record.

For every \(\varepsilon>0\) and nonempty record class, there is one
family-uniform represented record \(R_\varepsilon\) of cost at most \(B\)
such that
\begin{equation}
\label{eq-active-epsilon-optimal-record}
\mathfrak M_{n,m,\zeta}(A_{R_\varepsilon}^{H})
\ge
\operatorname{ActArr}_{\mathfrak M,n,m,\zeta,H}^{\rm sat}(B)-\varepsilon.
\end{equation}
This is an existence statement inside the represented record class.  An
affordable procedure that selects \(R_\varepsilon\) from a collection of
candidates is present precisely when that selection procedure is itself a
represented derivation charged by \(B_{\rm optimize}\).

\end{theorem}

\begin{proof}

The two sides of
\cref{eq-active-optimality-saturated-success-identity} take the same family
functional and the same supremum over the complete active-interface
derivations of
\cref{def-family-uniform-active-saturated-profile}; hence they are equal,
including their cost and uniformity requirements.  The five master families
are disjoint and exhaustive by
\cref{def-master-saturated-quotient-geometry-profile}.  Partitioning the
record class by its carried master slot and taking the supremum proves
\cref{eq-active-optimality-five-mode-arrival-decomposition}.  The
\(\varepsilon\)-optimal record follows from the definition of a supremum.
Derivation admissibility in
\cref{def:derivation-admissibility-criterion} gives the final assertion:
semantic selection among represented records is not itself a derivation.

\end{proof}

\begin{corollary}[Constructive active optimum on a represented finite tree]
\label{cor-constructive-active-optimum-finite-tree}

Assume that the refined history tree and every represented action set are
finite, the one-step probabilities and continuation values are represented
to a declared absolute precision, and comparison plus tie-breaking are
represented.  Backward induction applied to
\cref{eq-optimal-structural-active-resource-recursion} constructs a
represented policy.  If the value and comparison error at level \(t\) is at
most \(\varepsilon_t\), its arrival value is within
\begin{equation}
\label{eq-constructive-active-optimum-error}
\sum_{t=0}^{H-1}\varepsilon_t
\end{equation}
of the best record in the enumerated affordable tree.  Exact represented
arithmetic gives exact attainment.  Enumeration, value evaluation,
comparison, storage, precision, and deployment are all charged in
\cref{eq-complete-master-active-sampling-cost}; the corollary applies at
budget \(B\) exactly when this complete construction cost is at most \(B\).

\end{corollary}

\begin{proof}

At the last level choose a represented action within
\(\varepsilon_{H-1}\) of the largest represented terminal value.  Induct
backward using the disjoint first-contact recursion.  Conditional
expectation is monotone, so the loss accumulated at level \(t\) is at most
the local comparison error plus the continuation loss.  Induction gives
\cref{eq-constructive-active-optimum-error}.  Every operation used by this
backward induction is among the charged slots listed in
\cref{def-complete-master-active-sampling-record}.

\end{proof}

\subsection{Information, master visibility, and optimal arrival}
\label{subsec-active-information-master-visibility}

\begin{theorem}[Adaptive acquisition information and target contact]
\label{thm-adaptive-acquisition-information-target-contact}

Fix \(R\in\mathscr R_{I,H}^{\rm act}(B)\).  Enumerate every
target-dependent returned coordinate in its actual reveal order as
\(Y_1,\ldots,Y_L\), and let \(Q_j\) contain the query, selected experiment,
disclosed policy randomness, and deterministic updates preceding \(Y_j\).
Then
\begin{equation}
\label{eq-active-action-conditional-independence}
\Theta\perp Q_j\mid\mathscr F_{j-1}.
\end{equation}
If a target-dependent datum is used to construct \(Q_j\), it occurs in an
earlier reveal slot and is consequently already in \(\mathscr F_{j-1}\).
With the capacity and SDPI certificates of
\cref{def-controlled-channel-information-ledger},
\begin{align}
\label{eq-active-acquisition-information-identity}
I(\Theta;\mathscr F_L\mid\mathsf Z_0)
&=
\sum_{j=1}^{L}I(\Theta;Y_j\mid\mathscr P_j),\\
\label{eq-active-acquisition-information-upper}
I(\Theta;\mathscr F_L\mid\mathsf Z_0)
&\le
J_R^{\rm act,\uparrow}
:=
\sum_{j=1}^{L}
\mathbb E\min\{c_j(H_j^{\rm pre}),
\vartheta_j(H_j^{\rm pre})a_j(H_j^{\rm pre})\}.
\end{align}

Let \(\beta_{R,t}^{\rm act}\) be the next-contact probability under the
law which retains \((\mathsf Z_0,\mathcal H_t)\) and the selected kernel but
resamples \(\Theta\) independently from its declared reference conditional
law.  Then
\begin{align}
\label{eq-active-acquisition-next-contact}
\Pr_R\{\tau>t,X_{t+1}\in T_\Theta\}
&\le
\min\!\left\{1,
\beta_{R,t}^{\rm act}
+\sqrt{\frac{\log2}{2}J_{R,t}^{\rm act,\uparrow}}
\right\},\\
\label{eq-active-acquisition-cumulative-contact}
A_R^H
&\le
U_R^{\rm act,info}
:=
\min\!\left\{1,
\sum_{t=0}^{H-1}\left(
\beta_{R,t}^{\rm act}
+\sqrt{\frac{\log2}{2}J_{R,t}^{\rm act,\uparrow}}
\right)\right\}.
\end{align}
Every selected query, experiment evaluator, returned answer, update, and
action-optimization step remains charged by
\cref{eq-complete-master-active-sampling-cost}.

\end{theorem}

\begin{proof}

The selected action is a measurable function of the earlier public history
and fresh target-independent randomness, proving
\cref{eq-active-action-conditional-independence}.  Apply
\cref{thm-pre-hit-sequential-sdpi-information-contraction} to obtain
\cref{eq-active-acquisition-information-identity,eq-active-acquisition-information-upper}.
The conditional decoupling and Pinsker argument in
\cref{cor-pre-hit-information-target-contact-selector} gives
\cref{eq-active-acquisition-next-contact}.  Summing the disjoint first-hit
events proves \cref{eq-active-acquisition-cumulative-contact}.

\end{proof}

\begin{theorem}[Master optimal structural active-sampling envelope]
\label{thm-master-optimal-structural-active-sampling-envelope}

For a complete record \(R\), let \(U_R^{\rm act,cap}\) be every represented
same-record capacity or occupation upper bound for \(A_R^H\), with neutral
value one when absent.  Let \(U_R^{\rm act,BW}\) be every represented
Blackwell-transferred target-decision bound, again with neutral value one.
Finally, the master comparison in \(R\) has one of the two forms
\begin{equation}
\label{eq-active-master-visibility-comparison}
A_R^H
\le
U_R^{\rm act,master}
:=
\min\{1,K_RV_{\mathscr C_R}+\delta_R\},
\end{equation}
where \((K_R,\delta_R,\mathscr C_R)\) is either the inherited or direct
same-normalization record of
\cref{thm-no-uncharged-prospective-value-amplification}.  Write
\(R_{\le t}\) for the complete active record truncated after its next test
at time \(t\).  Put
\begin{equation}
\label{eq-active-source-resolved-learning-candidate}
\begin{aligned}
U_R^{\rm act,src}
&=
\min\left\{
1,\,
\sum_{t=0}^{H-1}
\min\left\{
\begin{aligned}
&1,\ U_{R,t}^{\rm stat\to tgt},\
\beta_{R,t}^+
+\sqrt{\frac{\log2}{2}J_{I,R_{\le t}}^{\rm src}},\\[-1mm]
&eH_B(n)\sum_{a\in\mathfrak J_{\Abs}^{5}}
\mathcal L_{I,R_{\le t}}^{(a)},\
U_{R,t}^{\rm dyn\to tgt}
\end{aligned}
\right\}
\right\},
\end{aligned}
\end{equation}
using the same complete active record and the ledger of
\cref{def-complete-source-resolved-learning-ledger}.  Put
\begin{equation}
\label{eq-active-same-record-certificate}
U_R^{\rm act}
=
\min\{U_R^{\rm act,info},U_R^{\rm act,cap},
       U_R^{\rm act,BW},U_R^{\rm act,master},
       U_R^{\rm act,src}\}.
\end{equation}
At a class-preserving common majorant
\[
B^{\rm act}=\Psi_{\rm act}(B,H,n,\mathbf m)
\]
of every slot in \cref{eq-complete-master-active-sampling-cost},
\begin{equation}
\label{eq-master-optimal-active-saturated-envelope}
\operatorname{ActArr}_{I,H}^{\rm sat}(B)
\le
\sup_{R\in\mathscr R_{I,H}^{\rm act}(B^{\rm act})}U_R^{\rm act}.
\end{equation}

Partition the records by the master-certificate slot.  Then the right side
of \cref{eq-master-optimal-active-saturated-envelope} equals the maximum of
its restrictions to
\begin{equation}
\label{eq-master-optimal-active-five-mode-decomposition}
\mathcal C_X,\quad\mathcal C_A,\quad
\mathcal C_{Q,\mathrm{ex}},\quad
\mathcal C_{Q,\mathrm{rev}},\quad
\mathcal C_{Q,\mathrm{nr}}.
\end{equation}
Consequently, exact \(\Abs\), exact quotient geometry, compensated
quasi-lumpability, and pure ambient geometry are evaluated by one theorem.
The pure ambient restriction is precisely \(q=\mathrm{id}\) with
\(\mathfrak G_q\ne\varnothing\).

\end{theorem}

\begin{proof}

The information theorem bounds \(A_R^H\) by
\(U_R^{\rm act,info}\).  The capacity and Blackwell candidates are admitted
only after conversion to the same target-arrival event and probability law,
so each also bounds \(A_R^H\).  The source-resolved learning candidate is
\cref{eq-source-resolved-learning-target-gain-minimum} applied at each
disjoint first-contact time of the same active record.  The master candidate is the represented
comparison in \cref{eq-active-master-visibility-comparison}.  Their
recordwise minimum proves \(A_R^H\le U_R^{\rm act}\).  Encoding adequacy and
computation closure place the complete record at \(B^{\rm act}\); taking the
supremum proves \cref{eq-master-optimal-active-saturated-envelope}.

The five certificate families are disjoint and exhaustive by
\cref{def-master-saturated-quotient-geometry-profile}.  A supremum over their
union is the maximum of the five restricted suprema, proving
\cref{eq-master-optimal-active-five-mode-decomposition}.

\end{proof}

\subsection{Complexity, sample, and noise frontiers}
\label{subsec-active-complexity-sample-noise-frontiers}

\begin{definition}[Active construction thresholds]
\label{def-active-construction-thresholds}

Fix the admissible scalarization \(\kappa\) of
\cref{def-degree-invariant}.  In this definition the record class without a
displayed budget is the union of its finite-cost slices.  Write
\begin{equation}
\label{eq-scalarized-active-arrival-profile}
\operatorname{ActArr}_{\mathfrak M,n,m,\zeta,H}^{\rm sat,\kappa}(s)
=
\sup_{\substack{R\in
\mathscr R_{\mathfrak I,n,m,\zeta,H}^{\rm act}\\
\kappa(\operatorname{Cost}_{\rm sat}(R))\le s}}
\mathfrak M_{n,m,\zeta}(A_R^H).
\end{equation}
For \(0<\alpha\le1\), define
\begin{align}
\label{eq-active-optimal-computational-cost}
B_{\rm act}^*(n,m,\zeta,H;\alpha)
&=
\inf\left\{s\ge0:
\operatorname{ActArr}_{\mathfrak M,n,m,\zeta,H}^{\rm sat,\kappa}(s)
\ge\alpha\right\},\\
\label{eq-active-optimal-sample-threshold}
m_{\rm act}^*(n,\zeta,s,H;\alpha)
&=
\inf\left\{m\ge1:
\operatorname{ActArr}_{\mathfrak M,n,m,\zeta,H}^{\rm sat,\kappa}(s)
\ge\alpha\right\}.
\end{align}
For homogeneous binary noise, define
\begin{equation}
\label{eq-active-optimal-noise-threshold}
\eta_{\rm act}^{\rm sat}(n,m;s,H,\alpha)
=
\sup\left\{0\le\eta<\frac12:
\operatorname{ActArr}_{\mathfrak M,n,m,\eta,H}^{\rm sat,\kappa}(s)
\ge\alpha\right\}.
\end{equation}
Infima of empty sets are \(+\infty\), and suprema of empty sets are zero.
These are restrictions of the existing recovery and dynamical threshold
definitions to the complete active-interface record class, rather than new
success criteria.

\end{definition}

\begin{theorem}[Quantitative monotonicity of the active optimum]
\label{thm-quantitative-active-optimum-monotonicity}

The scalarized profile in
\cref{eq-scalarized-active-arrival-profile} is nondecreasing in \(s\).  It
is nondecreasing in \(H\) whenever an earlier stopping rule may be retained
in a longer record.  If the declared sample presentations admit a
class-preserving embedding that appends observations which the record may
ignore, then it is nondecreasing in \(m\) at the compiled embedding cost.

For homogeneous binary noise, let \(0\le\eta_1\le\eta_2<1/2\), and let
\(\psi_{\rm flip}(s,n,m)\) be the complete scalar cost of the represented
degradation in \cref{thm-noise-degradation-monotonicity}.  Then
\begin{equation}
\label{eq-active-optimum-noise-degradation}
\operatorname{ActArr}_{\mathfrak M,n,m,\eta_1,H}^{\rm sat,\kappa}
\bigl(\psi_{\rm flip}(s,n,m)\bigr)
\ge
\operatorname{ActArr}_{\mathfrak M,n,m,\eta_2,H}^{\rm sat,\kappa}(s).
\end{equation}
Consequently the affordable active region is downward closed in the noise
rate whenever the declared budget absorbs this degradation cost.  The three
thresholds in \cref{def-active-construction-thresholds} are therefore the
computational-cost, sample, and noise sections of one saturated active
construction frontier.

\end{theorem}

\begin{proof}

Budget monotonicity is inclusion of scalar-cost sublevel sets.  A record may
retain its earlier stopping rule in a longer horizon, proving horizon
monotonicity.  Under the stated presentation embedding, append the extra
observations and run the original record without reading them; computation
closure charges the embedding and preserves its arrival law.

For the noise assertion, append independent flips of rate
\(q(\eta_1,\eta_2)=(\eta_2-\eta_1)/(1-2\eta_1)\) to the \(\eta_1\) record,
then run any \(\eta_2\) active construction.  The complete history, action,
and arrival law agree with the \(\eta_2\) execution, while the cost is at
most \(\psi_{\rm flip}(s,n,m)\).  Taking the saturated supremum proves
\cref{eq-active-optimum-noise-degradation}.

\end{proof}

\begin{corollary}[Learned acquisition-score transfer]
\label{cor-learned-active-acquisition-score-transfer}

Let \(g_R(h,a)\in[0,1]\) be a represented population acquisition score and
\(\widehat g_R(h,a)\) its represented estimate.  Suppose one complete
trajectory-learning record certifies, simultaneously on its represented
action class,
\[
\sup_{h,a}|\widehat g_R(h,a)-g_R(h,a)|\le\Gamma_R
\]
by either
\cref{def-sequential-structural-rademacher,thm-martingale-structural-pac-bayes}.
If the represented optimizer selects \(\widehat a(h)\) satisfying
\[
\widehat g_R(h,\widehat a(h))
\ge\sup_{a}\widehat g_R(h,a)-\varepsilon_{\rm opt},
\]
then
\begin{equation}
\label{eq-learned-active-acquisition-score-loss}
\sup_a g_R(h,a)-g_R(h,\widehat a(h))
\le2\Gamma_R+\varepsilon_{\rm opt}.
\end{equation}
The statistical radius, optimizer, and deployment costs are included in
\(B_{\rm learn}\oplus B_{\rm optimize}\).  The conclusion compares the
learned selector with the best represented action in the same record; the
value of that action remains bounded by
\cref{thm-master-optimal-structural-active-sampling-envelope}.

\end{corollary}

\begin{proof}

Insert \(\widehat g_R\) on both sides of the optimizer inequality.  The two
uniform deviations contribute \(2\Gamma_R\), and the represented
optimization defect contributes \(\varepsilon_{\rm opt}\).

\end{proof}

\begin{theorem}[Finite-horizon statistical loss of a learned active rule]
\label{thm-learned-active-rule-finite-horizon-loss}

Let \(Q_t^{\rm act,*}(h,a,b)\) be the analytic action value obtained by
taking represented action \(a\) at history \(h\), charging its declared
cost, and then taking the saturated suffix-record supremum in
\cref{eq-optimal-structural-active-resource-recursion}.  This analytic
quantity is not supplied to a deployed selector.  Suppose a complete
learning record constructs \(\widehat Q_t\) and a represented action
\(\widehat a_t(h)\) such that, on one simultaneous event,
\begin{align}
\label{eq-active-learned-q-uniform-radius}
\sup_{h,a}
|\widehat Q_t(h,a)-Q_t^{\rm act,*}(h,a,b_t)|
&\le\Gamma_t,\\
\label{eq-active-learned-q-optimization}
\widehat Q_t(h,\widehat a_t(h))
&\ge\sup_a\widehat Q_t(h,a)-\varepsilon_{{\rm opt},t}.
\end{align}
The suprema range over the represented histories and actions admitted by the
remaining budget \(b_t\).  Then the deployed record \(\widehat R\) obeys
\begin{equation}
\label{eq-learned-active-finite-horizon-value-loss}
V_0^{\rm act}(h_0,b_0)-A_{\widehat R}^{H}
\le
\mathbb E_{\widehat R}\!\left[
\sum_{t=0}^{H-1}
\mathbf1_{\{\tau>t\}}
\bigl(2\Gamma_t+\varepsilon_{{\rm opt},t}\bigr)
\right].
\end{equation}
Every \(\Gamma_t\) is admitted only after conversion to the same clean
action-value law.  Its sample certificate, confidence split, operator or
prior selection, score construction, optimization, and deployment costs are
included in \(B_{\rm learn}\oplus B_{\rm optimize}\).

\end{theorem}

\begin{proof}

Insert \(\widehat Q_t\) on both sides of
\cref{eq-active-learned-q-optimization}.  On the simultaneous event, the
conditional loss relative to the largest admitted analytic action value is
at most \(2\Gamma_t+\varepsilon_{{\rm opt},t}\).  Apply the exact
first-contact recursion to the optimal suffix value and the deployed suffix
value.  On target contact there is no continuation loss; on survival the
same argument applies at the next history.  Backward induction and the tower
property give
\cref{eq-learned-active-finite-horizon-value-loss}.

\end{proof}

\begin{corollary}[Explicit sample--noise active optimality loss]
\label{cor-explicit-sample-noise-active-optimality-loss}

Assume homogeneous binary noise \(0\le\eta<1/2\), a common optimization
error \(\varepsilon_{\rm opt}\), a \([0,1]\)-valued flip-symmetric
action-score loss, and one represented sample certificate
with effective sample size \(m_{\rm eff}\), dependence error
\(\varepsilon_{\rm dep}\), and confidence allocation \(\delta/H\).  If the
sample-independent affordable action-value codebook has logarithmic size
\(L_{\rm act}(n,b,m,\eta)\), the finite-description candidate of
\cref{prop:dynamical-finite-description-radius}, transported to the clean
flip-symmetric action-value law, is
\begin{equation}
\label{eq-active-clean-finite-description-radius}
\Gamma_{\rm act}^{\rm clean}
\le
\frac1{1-2\eta}
\left[
2\sqrt{\frac{
2L_{\rm act}(n,b_{\mathfrak U}^{\rm LPN},m,\eta)}
{m_{\rm eff}}}
+3\sqrt{\frac{\log(2H/\delta)}{2m_{\rm eff}}}
+\varepsilon_{\rm dep}
\right].
\end{equation}
On the same event,
\begin{equation}
\label{eq-active-explicit-sample-noise-value-loss}
V_0^{\rm act}-A_{\widehat R}^{H}
\le
H\left(2\Gamma_{\rm act}^{\rm clean}
+\varepsilon_{\rm opt}\right).
\end{equation}
The record may replace
\cref{eq-active-clean-finite-description-radius} by any smaller simultaneous
same-law candidate obtained from
\cref{def-sequential-structural-rademacher,thm-sequential-learned-operator-cover,thm-martingale-structural-pac-bayes}; the operator, cover, posterior,
prior, confidence allocation, and selection rule remain charged.

\end{corollary}

\begin{proof}

Use \(\delta_{\rm eff}=\delta/H\) in
\cref{eq-dynamical-noise-finite-description-radius}.  The homogeneous
flip-symmetric transport identity in
\cref{thm-dynamical-label-noise-conditional-transport} divides the corrupted
score radius by \(1-2\eta\), proving
\cref{eq-active-clean-finite-description-radius}.  Substitute the common
radius and optimization error into
\cref{eq-learned-active-finite-horizon-value-loss} and bound the number of
surviving terms by \(H\).

\end{proof}

\begin{figure}[tbp]
\centering
\resizebox{\linewidth}{!}{%
\begin{tikzpicture}[fw diagram,node distance=8mm and 11mm]
  \node[fw source,text width=2.7cm] (params)
    {\(n,m,\zeta,H,B\)\\presentation and budget};
  \node[fw provenance,text width=3.0cm,right=of params] (records)
    {complete family-uniform\\active records};
  \node[fw box,text width=4.2cm,right=of records] (modes)
    {\(\mathcal C_X,\mathcal C_A,\mathcal C_{Q,\rm ex},
       \mathcal C_{Q,\rm rev},\mathcal C_{Q,\rm nr}\)};
  \node[fw terminal,text width=3.0cm,right=of modes] (optimum)
    {saturated optimal\\target arrival};
  \node[fw use,text width=3.2cm,below=10mm of records] (learn)
    {represented optimizer\\and statistical radius};
  \draw[fw arrow] (params)--(records);
  \draw[fw arrow] (records)--(modes);
  \draw[fw arrow] (modes)--(optimum);
  \draw[fw arrow] (records)--(learn);
  \draw[fw arrow] (learn)-|(optimum);
\end{tikzpicture}%
}
\caption{Optimal active sampling is a saturated construction problem.  The
parameterized presentation and complete budget determine one family-uniform
record class.  Its five master cells partition the exact arrival optimum.
A learned optimizer approaches that optimum only through its represented,
completely charged score and statistical certificate.}
\label{fig-active-optimality-saturated-construction}
\end{figure}

\begin{figure}[tbp]
\centering
\resizebox{\linewidth}{!}{%
\begin{tikzpicture}[fw diagram,node distance=7mm and 10mm]
  \node[fw source,text width=2.7cm] (pub)
    {public record and\\pre-hit history};
  \node[fw provenance,text width=2.8cm,right=of pub] (opt)
    {represented score\\and action optimizer};
  \node[fw use,text width=2.7cm,above right=7mm and 10mm of opt] (query)
    {declared query\\and returned reveal};
  \node[fw geom,text width=2.7cm,below right=7mm and 10mm of opt] (move)
    {selected transport\\and master certificate};
  \node[fw box,text width=3.0cm,right=15mm of opt] (cert)
    {information, capacity, Blackwell,\\master visibility};
  \node[fw terminal,text width=2.7cm,right=of cert] (hit)
    {recordwise optimal\\target-arrival bound};
  \draw[fw arrow] (pub)--(opt);
  \draw[fw arrow] (opt)--(query);
  \draw[fw arrow] (opt)--(move);
  \draw[fw arrow] (query)--(cert);
  \draw[fw arrow] (move)--(cert);
  \draw[fw arrow] (cert)--(hit);
  \draw[fw arrow] (query) to[bend left=28] (opt);
\end{tikzpicture}%
}
\caption{Optimal structural active sampling.  The action optimizer, query,
returned reveal, transport kernel, and master certificate belong to one
charged pre-hit record.  Quasi-lumpable compensation is the general master
case; setting \(q=\mathrm{id}\) gives pure ambient geometry.}
\label{fig-master-optimal-structural-active-sampling}
\end{figure}

\section{Affordable Occupation Flows and the Bellman Structural Converse}
\label{sec-affordable-occupation-bellman-converse}

The active optimum is a supremum over complete represented records.  Its
occupation-flow representation gives a linear target-flux primal, while a
Bellman barrier gives the dual upper bound.  Both descriptions retain the
same clocked history, remaining budget, and complete record cost.

\begin{definition}[Stopped occupation flow of a complete active record]
\label{def-stopped-occupation-flow-complete-active-record}

Fix an instance \(I\), horizon \(H\), and a complete active record
\(R\in\mathscr R_{I,H}^{\rm act}(B)\).  Let \(\mathsf H_t^R\) be its
represented measurable carrier of surviving clocked histories at time \(t\),
including the remaining resource state.  Let \(\mathsf A_t^R(h)\) be the
represented action carrier available at \(h\); an action includes its query, transport
kernel, returned-data interface, master certificate, and charged
continuation selector.  Put
\begin{align}
\label{eq-stopped-occupation-history-mass}
d_{R,t}(C)
&=\Pr_R\{\tau>t,\mathcal H_t\in C\},\\
\label{eq-stopped-occupation-action-flow}
f_{R,t}(dh,da)
&=d_{R,t}(dh)K_{R,t}(da\mid h).
\end{align}
Write \(p_T(h,a)\) for first contact at the next transition and
\(P_N(h'\mid h,a)\) for the surviving history kernel.  Then
\begin{align}
\label{eq-stopped-occupation-initial-balance}
d_{R,0}&=\delta_{h_0},\\
\label{eq-stopped-occupation-action-balance}
f_{R,t}(dh,\mathsf A_t^R(h))&=d_{R,t}(dh),\\
\label{eq-stopped-occupation-flow-balance}
d_{R,t+1}(C)
&=\int f_{R,t}(dh,da)P_N(C\mid h,a).
\end{align}
The stopped occupation record is
\(\omega_R=(d_{R,t},f_{R,t})_{t<H}\).  Let
\begin{equation}
\label{eq-affordable-stopped-occupation-set}
\mathfrak O_{I,H}^{\rm sat}(B)
=\{\omega_R:R\in\mathscr R_{I,H}^{\rm act}(B)\}.
\end{equation}
This set contains only flows induced by complete records.  Its convex hull is
an analytic device; it does not supply an affordable mixture or policy.

\end{definition}

\begin{theorem}[Exact affordable occupation-flow primal]
\label{thm-exact-affordable-occupation-flow-primal}

Every complete active record satisfies
\begin{equation}
\label{eq-occupation-flow-exact-target-arrival}
A_R^H
=\sum_{t=0}^{H-1}\int
f_{R,t}(dh,da)p_T(h,a).
\end{equation}
Consequently
\begin{equation}
\label{eq-affordable-occupation-primal-identity}
\operatorname{ActArr}_{I,H}^{\rm sat}(B)
=\sup_{\omega\in\mathfrak O_{I,H}^{\rm sat}(B)}
\sum_{t=0}^{H-1}\int f_t(dh,da)p_T(h,a)
=\sup_{\omega\in\operatorname{co}\mathfrak O_{I,H}^{\rm sat}(B)}
\sum_{t=0}^{H-1}\int f_t(dh,da)p_T(h,a).
\end{equation}
The same identity holds for a family-uniform slice after restricting to one
family-uniform derivation scheme and applying its fixed family functional.

\end{theorem}

\begin{proof}

The event of first contact at time \(t+1\) has probability
\(\int f_{R,t}(dh,da)p_T(h,a)\).  These events are disjoint, so summing
proves \cref{eq-occupation-flow-exact-target-arrival}.  The first equality
in \cref{eq-affordable-occupation-primal-identity} is the definition of the
active saturated profile applied to this identity.  A linear functional has
the same supremum over a set and its convex hull, which proves the second
equality.  No convex mixture is inserted into the represented record class.

\end{proof}

\begin{definition}[Affordable Bellman support]
\label{def-affordable-bellman-support}

Let \(v=(v_t)_{t=0}^{H}\) be bounded functions on the represented surviving
history-and-resource carrier, with \(v_H=0\).  For a complete first-step
record \(r=(h,a)\), define its positive Bellman increment
\begin{equation}
\label{eq-recordwise-positive-bellman-increment}
\mathfrak d_t(r;v)
=\left[
p_T(h,a)+\int P_N(dh'\mid h,a)v_{t+1}(h')-v_t(h)
\right]_+.
\end{equation}
Let \(\mathscr A_{I,t}^{\rm sat}(h;B)\) contain exactly the first-step
records occurring in a complete suffix record of total cost at most \(B\).
The saturated Bellman-support coordinate is
\begin{equation}
\label{eq-saturated-affordable-bellman-support}
\mathfrak D_{I,t}^{\rm sat}(B;v)
=\sup_h\sup_{r\in\mathscr A_{I,t}^{\rm sat}(h;B)}
\mathfrak d_t(r;v),
\end{equation}
with value zero for an empty class.  The supremum is analytic.  A deployed
barrier or maximizing action is available only through its represented
derivation and complete cost.

For a master mode \(c\in\mathfrak C_{\rm master}\), write
\(\mathfrak D_{I,t;c}^{\rm sat}(B;v)\) for the restriction to first-step
records whose selected transition carries mode \(c\), and write
\(\mathscr A_{I,t;c}^{\rm sat}(h;B)\) for that restricted first-step record
class.

\end{definition}

\begin{theorem}[Bellman barrier duality for affordable arrival]
\label{thm-bellman-barrier-duality-affordable-arrival}

For every bounded barrier \(v\) as in
\cref{def-affordable-bellman-support},
\begin{equation}
\label{eq-affordable-bellman-barrier-upper-bound}
\operatorname{ActArr}_{I,H}^{\rm sat}(B)
\le
v_0(h_0)+\sum_{t=0}^{H-1}
\mathfrak D_{I,t}^{\rm sat}(B;v).
\end{equation}
Moreover,
\begin{equation}
\label{eq-affordable-bellman-barrier-exact-dual}
\operatorname{ActArr}_{I,H}^{\rm sat}(B)
=
\inf_{v_H=0}
\left\{
v_0(h_0)+\sum_{t=0}^{H-1}
\mathfrak D_{I,t}^{\rm sat}(B;v)
\right\},
\end{equation}
where the infimum ranges over bounded analytic functions on the complete
history-and-resource carrier.

The support decomposes exactly by the existing master classification:
\begin{equation}
\label{eq-affordable-bellman-support-five-mode-decomposition}
\mathfrak D_{I,t}^{\rm sat}(B;v)
=\max_{c\in\mathfrak C_{\rm master}}
\mathfrak D_{I,t;c}^{\rm sat}(B;v).
\end{equation}
If the represented rectangularity gate of
\cref{def-dynamical-generalization-envelope} holds, the infimum is the
ordinary finite occupation-measure/Bellman linear-program dual on the
represented statewise action correspondence.  Without that gate,
\cref{eq-affordable-bellman-barrier-exact-dual} uses the complete suffix
records of \cref{thm-budgeted-active-acquisition-recursion}; it performs no
statewise stitching of separately affordable policies.

\end{theorem}

\begin{proof}

For one record \(R\), integrate
\cref{eq-recordwise-positive-bellman-increment} against
\(f_{R,t}(dh,da)\) and sum over time.  The signed increment is bounded by its positive
part, hence
\begin{align*}
&\sum_t\int f_{R,t}(dh,da)p_T(h,a)\\
&\quad\le
\sum_t\int f_{R,t}(dh,da)
\left[v_t(h)-\int P_N(dh'\mid h,a)v_{t+1}(h')
      +\mathfrak D_{I,t}^{\rm sat}(B;v)\right].
\end{align*}
The flow balances in
\cref{eq-stopped-occupation-initial-balance,eq-stopped-occupation-action-balance,eq-stopped-occupation-flow-balance}
make the two value sums telescope to \(v_0(h_0)\); \(v_H=0\), and surviving
mass is at most one at every time.  Together with
\cref{eq-occupation-flow-exact-target-arrival}, this proves the recordwise
bound and then \cref{eq-affordable-bellman-barrier-upper-bound}.

For the reverse inequality in the infimum, take
\(v_t(h)=V_t^{\rm act}(h;B_h)\), the exact resource-valued suffix optimum in
\cref{thm-budgeted-active-acquisition-recursion}.  Its recursion makes every
admitted Bellman increment nonpositive, so every support term is zero and
\(v_0(h_0)=\operatorname{ActArr}_{I,H}^{\rm sat}(B)\).  The upper bound
already proves that no barrier has a smaller objective.  This establishes
\cref{eq-affordable-bellman-barrier-exact-dual}.

The master families partition the complete first-step records by
\cref{thm-active-optimality-saturated-construction}; a supremum over their
union is the maximum of the restricted suprema.  This proves
\cref{eq-affordable-bellman-support-five-mode-decomposition}.  Under
rectangularity, the balance equations and Bellman inequalities are the
finite primal and dual linear programs, and finite-dimensional strong
duality applies.  The complete-suffix formulation supplies the stated
nonrectangular reading.

\end{proof}

\begin{theorem}[Structural evaluation of affordable Bellman support]
\label{thm-structural-evaluation-affordable-bellman-support}

Fix \(v\), \(t\), and a common class-preserving ceiling
\(B^\sharp=\Psi_{\rm Bell}(B,H,n)\) charging the barrier evaluator, selected
transition, complete master certificate, and every comparison used below.
For a first-step record \(r\), let
\begin{equation}
\label{eq-recordwise-bellman-support-certificate-minimum}
U_{r,t}^{\rm Bell}(v)
=\min\{1,U_{r,t}^{\rm cap},U_{r,t}^{\Caus},U_{r,t}^{\rm info},
U_{r,t}^{\rm BW},U_{r,t}^{\rm obs},U_{r,t}^{\rm master}\},
\end{equation}
where a candidate is finite precisely when the same record contains a
represented comparison from that candidate to the positive increment
\(\mathfrak d_t(r;v)\).  An unavailable comparison has value one.  Then
\begin{align}
\label{eq-recordwise-bellman-support-certificate-bound}
\mathfrak d_t(r;v)
&\le U_{r,t}^{\rm Bell}(v),\\
\label{eq-saturated-bellman-support-master-bound}
\mathfrak D_{I,t}^{\rm sat}(B;v)
&\le
\max_{c\in\mathfrak C_{\rm master}}
\sup_{r\in\mathscr A_{I,t;c}^{\rm sat}(B^\sharp)}
U_{r,t}^{\rm Bell}(v).
\end{align}
The candidates are, respectively, the existing controlled-capacity arrival
comparison, authorized Caus aiming-to-contact comparison, sequential
information/contact comparison, Blackwell target-decision comparison,
sparse-reward observation-resolution comparison, and direct or inherited
master visibility comparison.  They are numerical views of the existing
record and create no structural channel.

\end{theorem}

\begin{proof}

Every finite candidate in
\cref{eq-recordwise-bellman-support-certificate-minimum} is required to
bound the same positive Bellman increment under the same history law.  Their
minimum therefore proves
\cref{eq-recordwise-bellman-support-certificate-bound}.  Encoding adequacy
and computation closure place the complete comparison record at
\(B^\sharp\).  Take the supremum first within each master class and then use
\cref{eq-affordable-bellman-support-five-mode-decomposition}.  This proves
\cref{eq-saturated-bellman-support-master-bound}.  The candidate origins
are the coordinatewise comparisons in
\cref{thm-master-optimal-structural-active-sampling-envelope,thm-universal-dynamical-certificate};
their same-normalization requirement is
\cref{thm-no-uncharged-prospective-value-amplification}.

\end{proof}

\begin{theorem}[Quantitative prospective contact-witness coercivity]
\label{thm-quantitative-prospective-contact-witness-coercivity}

Fix a presented instance \(I\), a refined horizon \(H\ge1\), a saturated
budget \(B\), and a complete active record
\(R\in\mathscr R_{I,H}^{\rm act}(B)\).  Use the one-test normalization of
\cref{def-represented-pre-hit-selector}, let \(K_{R,t}\) be the represented
next-test kernel of \(R\), and let \(\nu_t\) be a declared neutral baseline
on the same eligible candidates.  Put
\begin{align}
\label{eq-recordwise-excess-over-neutral-search}
\varepsilon_R
&=
\left[
A_R^H-\operatorname{Enum}_{\nu}(R,H)
\right]_+,\\
\label{eq-recordwise-positive-selector-advantage}
a_{R,t}
&=
\bigl(\operatorname{Adv}_t(K_{R,t};\nu_t)\bigr)_+,\\
\label{eq-recordwise-next-contact-hazard}
h_{R,t}
&=
\mathbb E_R\!\left[
\mathbf1_{\{\tau>t\}}
K_{R,t}(V_I=1\mid\mathcal H_t)
\right].
\end{align}
Then the following pointwise chain is exact up to taking positive parts and
suprema:
\begin{equation}
\label{eq-prospective-contact-witness-coercive-chain}
\varepsilon_R
\le
\sum_{t=0}^{H-1}a_{R,t}
\le
\sum_{t=0}^{H-1}h_{R,t}
\le
\sum_{t=0}^{H-1}\mathfrak D_{I,t}^{\rm sat}(B;0)
=
\sum_{t=0}^{H-1}
\max_{c\in\mathfrak C_{\rm master}}
\mathfrak D_{I,t;c}^{\rm sat}(B;0).
\end{equation}
Consequently, if \(\varepsilon_R>0\), there are a time
\(t_*<H\) and a master mode \(c_*\in\mathfrak C_{\rm master}\) such that
\begin{equation}
\label{eq-prospective-contact-witness-extraction-scale}
\mathfrak D_{I,t_*;c_*}^{\rm sat}(B;0)
\ge \frac{\varepsilon_R}{H}.
\end{equation}
More concretely, for every
\(0\le\gamma<\varepsilon_R/H\), the complete first-step record class in
that mode contains a record \(r=(h,a)\) satisfying
\begin{equation}
\label{eq-prospective-contact-first-step-witness}
p_T(h,a)>\gamma.
\end{equation}
The record is prospective: its history, selected transition, randomness,
memory, and complete normalized derivation are available before the tested
transition and are charged within \(B\).  Its transition and source ancestry
have the existing master-mode and channel assignments of
\cref{thm-no-free-succinct-transition-rule,thm-complete-generic-source-theory};
no terminal flag or completed first-hit quotient is inserted into that
ancestry.

At the compiled comparison ceiling
\(B^\sharp=\Psi_{\rm Bell}(B,H,n)\), the same conclusion has the
same-record certificate form
\begin{equation}
\label{eq-prospective-contact-witness-certificate-scale}
\max_{c\in\mathfrak C_{\rm master}}
\sup_{r\in\mathscr A_{I,t_*;c}^{\rm sat}(B^\sharp)}
U_{r,t_*}^{\rm Bell}(0)
\ge \frac{\varepsilon_R}{H}.
\end{equation}
Here an unavailable comparison has the neutral value one, exactly as in
\cref{eq-recordwise-bellman-support-certificate-minimum}.  Thus
\cref{eq-prospective-contact-witness-certificate-scale} does not turn an
unavailable comparison into a structural estimate: it identifies either an
evaluated same-record certificate at the displayed scale or the precise
comparison slot that remains unevaluated.

Conversely, if numbers \(\delta_{t,c}\in[0,1]\) satisfy
\begin{equation}
\label{eq-modewise-contact-support-input-bound}
\mathfrak D_{I,t;c}^{\rm sat}(B;0)\le\delta_{t,c}
\qquad
(0\le t<H,\ c\in\mathfrak C_{\rm master}),
\end{equation}
then every complete active record of cost at most \(B\) obeys
\begin{equation}
\label{eq-modewise-contact-support-arrival-bound}
A_R^H
\le
\operatorname{Enum}_{\nu}(R,H)
+\sum_{t=0}^{H-1}
\max_{c\in\mathfrak C_{\rm master}}\delta_{t,c}.
\end{equation}
These implications use the existing contact, Bellman-support, and master
coordinates.  They introduce no additional structural channel.
The implication chain is summarized in
\cref{fig-quantitative-prospective-contact-witness}.

\end{theorem}

\begin{proof}

The hazard identity
\cref{eq-selector-hazard-decomposition} gives
\[
A_R^H-\operatorname{Enum}_{\nu}(R,H)
=\sum_{t=0}^{H-1}
\operatorname{Adv}_t(K_{R,t};\nu_t).
\]
Taking the positive part and using
\([\sum_t z_t]_+\le\sum_t(z_t)_+\) proves the first inequality in
\cref{eq-prospective-contact-witness-coercive-chain}.  The neutral hazard is
nonnegative, so
\[
a_{R,t}
\le
h_{R,t}
=
\mathbb E_R\!\left[
\mathbf1_{\{\tau>t\}}p_T(\mathcal H_t,A_t)
\right].
\]
Every first-step record used by \(R\) occurs in a complete suffix record of
total cost at most \(B\).  With \(v=0\), its positive Bellman increment is
exactly \(p_T(h,a)\).  Hence the integrand is at most
\(\mathfrak D_{I,t}^{\rm sat}(B;0)\), proving the third inequality.  The
last equality is the exact master decomposition
\cref{eq-affordable-bellman-support-five-mode-decomposition}.

If \(\varepsilon_R>0\), the sum of the \(H\) nonnegative support terms is at
least \(\varepsilon_R\); one time therefore has support at least
\(\varepsilon_R/H\).  The finite maximum over master modes gives \(c_*\).
The definition of a supremum then gives
\cref{eq-prospective-contact-first-step-witness} for every strict lower
level \(\gamma<\varepsilon_R/H\).  Temporal availability and complete cost
are part of
\cref{def-complete-master-active-sampling-record,def-affordable-bellman-support};
transition-rule provenance and the master assignment are
\cref{thm-no-free-succinct-transition-rule,thm-complete-generic-source-theory}.
The exclusion of later flags is
\cref{def-pre-hit-temporally-admissible-source,cor-first-hit-accounting-not-source}.

Apply
\cref{eq-saturated-bellman-support-master-bound} at \(t_*\).  Together with
\cref{eq-prospective-contact-witness-extraction-scale}, it proves
\cref{eq-prospective-contact-witness-certificate-scale}.  The convention for
an unavailable comparison is part of
\cref{thm-structural-evaluation-affordable-bellman-support}; retaining that
convention proves the stated scope without adding a comparison.

Finally, substitute
\cref{eq-modewise-contact-support-input-bound} into
\cref{eq-prospective-contact-witness-coercive-chain} and undo the positive
part.  This gives
\cref{eq-modewise-contact-support-arrival-bound}.

\end{proof}

\begin{figure}[tbp]
\centering
\resizebox{.98\linewidth}{!}{%
\begin{tikzpicture}[fw diagram,node distance=7mm and 9mm]
  \node[fw source,text width=2.8cm] (excess)
    {arrival above\\neutral search};
  \node[fw provenance,text width=2.7cm,right=of excess] (hazard)
    {one pre-hit selector\\advantage \(\ge\varepsilon/H\)};
  \node[fw geom,text width=2.8cm,right=of hazard] (contact)
    {next-contact Bellman\\support at the same scale};
  \node[fw box,text width=2.7cm,right=of contact] (mode)
    {one existing\\master mode};
  \node[fw terminal,text width=2.9cm,right=of mode] (cert)
    {evaluated comparison\\or identified open slot};
  \node[fw residual,text width=3.0cm,below=9mm of contact] (retro)
    {later first-hit labels:\\accounting only};
  \draw[fw arrow] (excess)--(hazard);
  \draw[fw arrow] (hazard)--(contact);
  \draw[fw arrow] (contact)--(mode);
  \draw[fw arrow] (mode)--(cert);
  \draw[fw dashed arrow] (contact)--(retro);
\end{tikzpicture}%
}
\caption{Quantitative prospective witness extraction.  Excess target
arrival forces a pre-hit next-contact record at scale
\(\varepsilon/H\).  The complete record has one existing master-mode
assignment.  A later first-hit label cannot be used as its source.}
\label{fig-quantitative-prospective-contact-witness}
\end{figure}

\begin{corollary}[Complete learned-control sandwich]
\label{cor-complete-learned-control-sandwich}

Let \(\widehat R\) be a complete learned active record satisfying the
simultaneous action-value and optimization bounds of
\cref{thm-learned-active-rule-finite-horizon-loss}.  Then
\begin{equation}
\label{eq-complete-learned-control-sandwich}
A_{\widehat R}^{H}
\le\operatorname{ActArr}_{I,H}^{\rm sat}(B)
\le
\inf_{v_H=0}\left\{
v_0(h_0)+\sum_{t<H}\mathfrak D_{I,t}^{\rm sat}(B;v)
\right\},
\end{equation}
and
\begin{equation}
\label{eq-complete-learned-control-optimality-loss}
0\le
\operatorname{ActArr}_{I,H}^{\rm sat}(B)-A_{\widehat R}^{H}
\le
\mathbb E_{\widehat R}\sum_{t<H}
\mathbf1_{\{\tau>t\}}
\bigl(2\Gamma_t+\varepsilon_{{\rm opt},t}\bigr).
\end{equation}
Each \(\Gamma_t\) is the smallest simultaneous same-action-value consequence
of the independent, mixing, sequential-Rademacher, martingale PAC--Bayes,
effective-dimension, label-noise, model-channel, filtering, and occupancy
shift certificates in the complete record.  Thus the first display bounds
the best affordable structural controller, while the second measures the
statistical and optimization distance of the learned controller from that
optimum.

\end{corollary}

\begin{proof}

The first inequality is membership of \(\widehat R\) in the affordable
record class.  The equality in
\cref{thm-bellman-barrier-duality-affordable-arrival} proves the second.
The final display is
\cref{thm-learned-active-rule-finite-horizon-loss}.  Its same-law candidate
aggregation is exactly the rule in
\cref{def-dynamical-generalization-envelope,def-master-dynamical-leaf-candidate-aggregation}.

\end{proof}

\begin{figure}[tbp]
\centering
\resizebox{.98\linewidth}{!}{%
\begin{tikzpicture}[fw diagram,node distance=7mm and 10mm]
  \node[fw source,text width=2.7cm] (records)
    {complete affordable\\active records};
  \node[fw provenance,text width=2.7cm,right=of records] (flow)
    {stopped occupation\\and target flux};
  \node[fw geom,text width=2.7cm,right=of flow] (dual)
    {Bellman barrier\\and support};
  \node[fw box,text width=3.0cm,right=of dual] (modes)
    {five master modes\\same-record certificates};
  \node[fw terminal,text width=2.8cm,right=of modes] (arrival)
    {best affordable\\target arrival};
  \node[fw use,text width=2.8cm,below=9mm of dual] (learn)
    {estimation, noise\\and optimization loss};
  \draw[fw arrow] (records)--(flow);
  \draw[fw arrow] (flow)--(dual);
  \draw[fw arrow] (dual)--(modes);
  \draw[fw arrow] (modes)--(arrival);
  \draw[fw arrow] (learn)--(dual);
\end{tikzpicture}%
}
\caption{Affordable occupation-flow/Bellman duality.  Target arrival is the
linear flux of complete represented records.  Bellman support is evaluated
through the existing master certificates, while learning bounds measure the
distance from the best affordable controller.}
\label{fig-affordable-occupation-bellman-duality}
\end{figure}

\section{Optimal Structural Bayesian Learning}
\label{sec-optimal-structural-bayesian-learning}

Structural Bayesian learning equips a complete active record with a
represented belief update.  The target posterior and the PAC--Bayes
posterior have different types: the former is a probability law on the
task index, whereas the latter is a probability law on predictors or
policies.  This section keeps the two laws separate and charges every
operation that turns public observations into a deployed target belief.

\subsection{Prior provenance and represented task beliefs}
\label{subsec-structural-bayes-prior-provenance}

\begin{definition}[Complete structural Bayesian record]
\label{def-complete-structural-bayesian-record}

Fix a finite target-index carrier \(\mathsf\Theta\), a declared prior
\(\Pi_0\) of full support, and a complete active record \(R\) from
\cref{def-complete-master-active-sampling-record}.  Let
\(\mathcal F_t\) be its complete public pre-hit filtration and put
\begin{equation}
\label{eq-analytic-task-posterior}
\Pi_t(\theta)
=\Pr\{\Theta=\theta\mid\mathcal F_t\}.
\end{equation}
This conditional law is an analytic comparator.  A deployed belief is a
represented \(\mathcal F_t\)-measurable probability vector
\(\widehat\Pi_t\) on \(\mathsf\Theta\).  Its comparison error is
\begin{equation}
\label{eq-structural-bayes-filter-error}
e_t
=\lVert\widehat\Pi_t-\Pi_t\rVert_{\rm TV}.
\end{equation}

Suppose the action \(A_t^{\rm qry}\) selects a declared observation kernel
\(L_t(dy\mid\theta,\mathcal F_t,A_t^{\rm qry})\).  A represented Bayesian
update consists of represented nonnegative likelihood scores
\(\widehat L_t(y\mid\theta)\), a represented evidence normalizer, and the
update
\begin{equation}
\label{eq-represented-structural-bayes-update}
\widehat\Pi_{t+1}(\theta)
=
\frac{\widehat\Pi_t(\theta)\widehat L_t(Y_{t+1}\mid\theta)}
{\sum_{u\in\mathsf\Theta}
 \widehat\Pi_t(u)\widehat L_t(Y_{t+1}\mid u)}
\end{equation}
whenever the denominator has a represented positive-evidence certificate.
An alternative represented normalizer or approximate update is admitted
with its error in \cref{eq-structural-bayes-filter-error}.  A zero or
uncertified denominator gives infinite raw filter error and the neutral
value one after conversion to a probability coordinate.

The complete cost is the active cost in
\cref{eq-complete-master-active-sampling-cost}.  The slots
\[
B_{\rm query},B_{\rm obs},B_{\rm channel},B_{\rm belief},
\qquad
B_{\rm memory},B_{\rm score},B_{\rm optimize},B_{\rm learn}
\]
collectively charge the
experiment selector, returned observation, likelihood interface, prior
code, likelihood evaluation, evidence normalization, belief storage,
precision, posterior statistic, acquisition optimizer, and deployment.
These operations retain the existing master structural channel of their
complete active record.
Write
\begin{equation}
\label{eq-structural-bayes-common-ceiling}
B^{\rm Bayes}
=\Psi_{\rm Bayes}(B,H,n,\mathbf m,\boldsymbol\zeta)
\end{equation}
for a class-preserving common majorant of these existing cost slots, where
\(\mathbf m\) and \(\boldsymbol\zeta\) collect the declared sample and
channel parameters.  This is a common accounting ceiling; the primitive
signature and constructor grammar remain unchanged.

The notation \(\Pi_t,\widehat\Pi_t\) is reserved for task-index beliefs.
A PAC--Bayes posterior \(Q\) remains a law on hypotheses, operators, or
policies and retains the KL charge of
\cref{def-structural-pac-bayes-kl,thm-martingale-structural-pac-bayes}.

\end{definition}

\begin{theorem}[Prior provenance and neutral initialization]
\label{thm-structural-bayes-prior-provenance-neutrality}

Let \(G\) be a represented initialization record independent of \(\Theta\).
Conditional on \(G\), let \(K_G\) be any represented probability kernel on
a common candidate carrier \(\mathsf C\).  For a verification relation
\(V_\theta(c)\in\{0,1\}\), suppose
\begin{equation}
\label{eq-structural-bayes-prior-overlap}
\sup_{c\in\mathsf C}
\sum_{\theta\in\mathsf\Theta}V_\theta(c)\le\Lambda.
\end{equation}
If \(\Theta\) is uniform on \(N=|\mathsf\Theta|\), then
\begin{equation}
\label{eq-target-independent-prior-neutral-contact}
\mathbb E\sum_cK_G(c)V_\Theta(c)\le\frac{\Lambda}{N}.
\end{equation}
When \(\mathsf C=\mathsf\Theta\) and
\(V_\theta(c)=\mathbf1_{\{c=\theta\}}\), the contact is exactly \(1/N\),
independent of every subjective nonuniform belief selected by \(G\).

If an initialization rule reads a target-dependent public variable \(Z\),
then \(Z\) belongs to the first reveal slot of the public filtration and the
resulting law is \(\Pi_0(\cdot\mid Z)\).  Its channel, likelihood, update,
normalization, and retained precision are charged as the first observation
and belief update.  A target-dependent initialization without this ancestry
is outside the complete structural Bayesian record class.

\end{theorem}

\begin{proof}

Independence gives
\[
\mathbb E\sum_cK_G(c)V_\Theta(c)
=\mathbb E_G\sum_cK_G(c)
  \frac1N\sum_\theta V_\theta(c)
\le\frac\Lambda N.
\]
For the displayed singleton relation, the inner target sum equals one for
every candidate, so equality holds.  When \(Z\) depends on \(\Theta\), the declared reveal
order places \(Z\) before every rule that uses it.  The chain rule for
conditional probability identifies conditioning on \(Z\) with the first
posterior update.  Temporal source admissibility and the complete active
cost then charge the stated record components.

\end{proof}

\subsection{Posterior concentration and target contact}
\label{subsec-structural-bayes-posterior-contact}

\begin{definition}[Target vulnerability and affordable posterior profile]
\label{def-affordable-structural-posterior-concentration}

For a probability law \(q\) on \(\mathsf\Theta\), define its target
vulnerability and excess over the declared prior by
\begin{align}
\label{eq-target-vulnerability-functional}
\mathcal V_T(q)
&=\sup_{c\in\mathsf C}\sum_{\theta}q(\theta)V_\theta(c),\\
\label{eq-target-vulnerability-excess}
\mathcal C_T(q;\Pi_0)
&=\bigl[\mathcal V_T(q)-\mathcal V_T(\Pi_0)\bigr]_+.
\end{align}
For disjoint singleton target cells,
\(\mathcal V_T(q)=\max_\theta q(\theta)\).  The functional is analytic
unless the record contains its represented evaluator and comparison rule.

Let \(\mathscr R_{I,H}^{\rm Bayes}(B)\) be the complete structural Bayesian
records of cost at most \(B\).  Define
\begin{align}
\label{eq-saturated-affordable-posterior-concentration}
\operatorname{PostConc}_{I,H}^{\rm sat}(B)
&=\max_{t<H}\sup_{R\in\mathscr R_{I,H}^{\rm Bayes}(B)}
\mathbb E_R\!\left[
\mathbf1_{\{\tau>t\}}
\mathcal C_T(\widehat\Pi_t;\Pi_0)
\right],\\
\label{eq-saturated-structural-bayes-filter-error}
\operatorname{FilErr}_{I,H}^{\rm sat}(B)
&=\max_{t<H}\sup_{R\in\mathscr R_{I,H}^{\rm Bayes}(B)}
\mathbb E_R\!\left[
\mathbf1_{\{\tau>t\}}e_t
\right].
\end{align}
The family-uniform versions restrict to one family-uniform represented
derivation scheme and apply the declared family functional.

Let \(\mathscr R_{I,H}^{\rm Bayes,src}(B)\) be the restriction to records
in which the target-vulnerability evaluator, its comparison with
\(\mathcal V_T(\Pi_0)\), and the action using the resulting score are
represented before that action.  Write
\begin{equation}
\label{eq-saturated-prospective-posterior-source-profile}
\operatorname{PostSrc}_{I,H}^{\rm sat}(B)
=\max_{t<H}
\sup_{R\in\mathscr R_{I,H}^{\rm Bayes,src}(B)}
\mathbb E_R\!\left[
\mathbf1_{\{\tau>t\}}
\mathcal C_T(\widehat\Pi_t;\Pi_0)
\right].
\end{equation}
Thus \(\operatorname{PostConc}^{\rm sat}\)
is an analytic evaluation of represented beliefs, whereas
\(\operatorname{PostSrc}^{\rm sat}\) is its temporally admissible
prospective-source subprofile.

\end{definition}

\begin{theorem}[Structural Bayesian target-contact accountability]
\label{thm-structural-bayesian-target-contact-accountability}

Let \(R\in\mathscr R_{I,H}^{\rm Bayes}(B)\), and let
\(K_t(dc\mid\mathcal F_t)\) be its represented next-test kernel.  Then
\begin{equation}
\label{eq-structural-bayes-exact-next-contact-identity}
\Pr_R\{\tau>t,V_\Theta(C_{t+1})=1\}
=\mathbb E_R\!\left[
\mathbf1_{\{\tau>t\}}
\sum_cK_t(c\mid\mathcal F_t)
\sum_\theta\Pi_t(\theta)V_\theta(c)
\right]
\end{equation}
and hence
\begin{equation}
\label{eq-structural-bayes-next-contact-bound}
\Pr_R\{\tau>t,V_\Theta(C_{t+1})=1\}
\le
\mathbb E_R\!\left[
\mathbf1_{\{\tau>t\}}
\bigl(
\mathcal V_T(\Pi_0)
+\mathcal C_T(\widehat\Pi_t;\Pi_0)+e_t
\bigr)
\right].
\end{equation}
Consequently
\begin{equation}
\label{eq-structural-bayes-saturated-arrival-bound}
\operatorname{BayesArr}_{I,H}^{\rm sat}(B)
\le
H\left(
\mathcal V_T(\Pi_0)
+\operatorname{PostConc}_{I,H}^{\rm sat}(B)
+\operatorname{FilErr}_{I,H}^{\rm sat}(B)
\right),
\end{equation}
where
\begin{equation}
\label{eq-structural-bayes-arrival-profile}
\operatorname{BayesArr}_{I,H}^{\rm sat}(B)
=\sup_{R\in\mathscr R_{I,H}^{\rm Bayes}(B)}\Pr_R\{\tau\le H\}.
\end{equation}
The right side may be clipped at one.  The same statements hold after a
family-uniform restriction and its fixed family functional.

\end{theorem}

\begin{proof}

Condition on \(\mathcal F_t\).  Fresh selector randomness is conditionally
independent of \(\Theta\), so the conditional contact probability is the
inner double sum in
\cref{eq-structural-bayes-exact-next-contact-identity}.  For every candidate
\(c\), the function \(\theta\mapsto V_\theta(c)\) takes values in
\([0,1]\); therefore its expectations under \(\Pi_t\) and
\(\widehat\Pi_t\) differ by at most \(e_t\).  Averaging over \(K_t\) and
then bounding by the supremum over \(c\) gives
\[
\sum_cK_t(c\mid\mathcal F_t)
\sum_\theta\Pi_t(\theta)V_\theta(c)
\le\mathcal V_T(\widehat\Pi_t)+e_t.
\]
Use \cref{eq-target-vulnerability-excess} to obtain
\cref{eq-structural-bayes-next-contact-bound}.  First-contact events are
disjoint.  Sum over \(t<H\), take the saturated supremum, and use
\cref{eq-saturated-affordable-posterior-concentration,eq-saturated-structural-bayes-filter-error}
to prove \cref{eq-structural-bayes-saturated-arrival-bound}.

\end{proof}

\begin{theorem}[Bayes-factor gain, information, and degradation]
\label{thm-structural-bayes-factor-information-degradation}

Condition on a surviving history and represented action.  Let
\(L_t(dy\mid\theta)\) be the declared observation channel and
\(\overline L_t(dy)=\sum_\theta\Pi_t(\theta)L_t(dy\mid\theta)\).
On positive-evidence outcomes, the exact analytic posterior obeys
\begin{equation}
\label{eq-exact-structural-bayes-factor-update}
\Pi_{t+1}(\theta)
=\Pi_t(\theta)
\frac{dL_t(\cdot\mid\theta)}{d\overline L_t}(Y_{t+1}).
\end{equation}
Writing information in bits,
\begin{align}
\label{eq-structural-bayes-kl-information-identity}
\mathbb E\!\left[
D_{\rm KL}(\Pi_{t+1}\Vert\Pi_t)\mid\mathcal F_t,A_t
\right]
&=(\log2)I(\Theta;Y_{t+1}\mid\mathcal F_t,A_t),\\
\label{eq-structural-bayes-vulnerability-gain}
\mathbb E\!\left[
\bigl(\mathcal V_T(\Pi_{t+1})-
      \mathcal V_T(\Pi_t)\bigr)_+
\mid\mathcal F_t,A_t
\right]
&\le
\sqrt{\frac{\log2}{2}
I(\Theta;Y_{t+1}\mid\mathcal F_t,A_t)}.
\end{align}
For represented beliefs,
\begin{equation}
\label{eq-approximate-structural-bayes-vulnerability-gain}
\bigl(\mathcal V_T(\widehat\Pi_{t+1})-
      \mathcal V_T(\widehat\Pi_t)\bigr)_+
\le
\lVert\Pi_{t+1}-\Pi_t\rVert_{\rm TV}+e_t+e_{t+1}.
\end{equation}

If the complete channel is a represented garbling of another declared
channel, its expected target vulnerability is no larger than that of the
more informative channel.  If its SDPI coefficient is \(\vartheta_t\), the
information term in \cref{eq-structural-bayes-vulnerability-gain} may be
replaced by the corresponding contracted information bound from
\cref{thm-pre-hit-sequential-sdpi-information-contraction}.  A represented
chi-square comparison supplies the additional same-law target-contact
candidate in \cref{eq-pre-hit-information-chi-square-contact}; conversion
to vulnerability units requires the represented comparison of
\cref{thm-affordable-posterior-concentration-source-routing}.

\end{theorem}

\begin{proof}

Bayes' rule gives \cref{eq-exact-structural-bayes-factor-update}.  Taking
the posterior expectation of its log density ratio and then averaging over
the observation gives the standard conditional mutual-information identity,
with the factor \(\log2\) converting bits to nats.  The functional
\(\mathcal V_T\) is one-Lipschitz in total variation: it is the supremum of
expectations of functions taking values in \([0,1]\).  Pinsker's inequality
and Jensen's inequality therefore prove
\cref{eq-structural-bayes-vulnerability-gain}.  Apply the same Lipschitz
property twice and the triangle inequality to obtain
\cref{eq-approximate-structural-bayes-vulnerability-gain}.

Blackwell garbling preserves every decision rule available after the
garbled channel as a randomized rule after the original channel.  Taking
the best target decision proves vulnerability monotonicity; this is the
specialization of
\cref{thm-controlled-prehit-blackwell-value-monotonicity}.  The SDPI and
chi-square conclusions follow from the cited same-law comparisons.

\end{proof}

\begin{corollary}[Cumulative structural Bayesian concentration]
\label{cor-cumulative-structural-bayesian-concentration}

Let \(Y_1,\ldots,Y_L\) be the actual ordered target-dependent reveals in a
complete structural Bayesian record, and write
\[
i_j=I(\Theta;Y_j\mid\mathcal F_{j-1},A_{j-1})
\]
in bits.  If \(J_t^\uparrow\) is any represented upper bound for the total
information in the pre-hit record through time \(t\), then
\begin{equation}
\label{eq-cumulative-structural-bayes-concentration}
\mathbb E\mathcal C_T(\widehat\Pi_t;\Pi_0)
\le
\min\left\{
\sqrt{\frac{\log2}{2}J_t^\uparrow},
\sum_{j\le L_t}\mathbb E
\sqrt{\frac{\log2}{2}i_j}
\right\}
+\mathbb E e_t.
\end{equation}
Each \(i_j\) may be replaced by its channel-capacity or SDPI candidate in
\cref{thm-pre-hit-sequential-sdpi-information-contraction}.  Target-
independent computation between reveals contributes zero information; its
belief-update, optimization, memory, and deployment costs remain in the
complete record.

\end{corollary}

\begin{proof}

The total-information candidate is
\cref{eq-pre-hit-information-posterior-tv}, followed by the
one-Lipschitz vulnerability comparison and the filter error.  For the
second candidate, telescope
\(\mathcal V_T(\Pi_t)-\mathcal V_T(\Pi_0)\), bound the positive part by the
sum of positive one-step increments, and apply
\cref{eq-structural-bayes-vulnerability-gain} to each reveal.  Replacing the
exact posterior by \(\widehat\Pi_t\) adds \(e_t\).  The conditional
information identity assigns zero to deterministic target-independent
updates.

\end{proof}

\subsection{Source-resolved actionable Bayesian gain}
\label{subsec-source-resolved-structural-bayesian-gain}

The posterior of \cref{eq-analytic-task-posterior} is an analytic law.  The
operational quantity is the improvement in the next represented test.  This
subsection resolves that improvement into the structural sources of the
complete pre-hit record and retains the noise contraction of every
target-dependent reveal.

\begin{definition}[Actionable Bayesian gain and source-resolved charge]
\label{def-source-resolved-actionable-bayesian-gain}

Let \(R\in\mathscr R_{I,H}^{\rm Bayes}(B)\), and let
\(K_t(\,\cdot\mid\mathcal F_t)\) be its represented next-test kernel.  Put
\begin{align}
\label{eq-actionable-structural-bayesian-gain}
g_{R,t}^{\rm Bayes}
={}&\Biggl[
\mathbb E_R\!\left[
\mathbf1_{\{\tau>t\}}
\sum_{c}K_t(c\mid\mathcal F_t)
\sum_{\theta}\Pi_t(\theta)V_\theta(c)
\right]\notag\\
&\qquad-
\mathbb E_R[\mathbf1_{\{\tau>t\}}]
\mathcal V_T(\Pi_0)
\Biggr]_+ .
\end{align}
This is the target-contact improvement actually deployed by the record.  It
is distinct from the analytic excess
\(\mathcal C_T(\Pi_t;\Pi_0)\) and from the represented-belief audit
\(\mathcal C_T(\widehat\Pi_t;\Pi_0)\).

Truncate the execution immediately after this next test and apply the
five-family frontier \(\mathfrak F_5\) of
\cref{cor-five-family-abs-envelope}.  For
\(a\in\mathfrak J_{\Abs}^{5}\), let
\begin{equation}
\label{eq-bayesian-source-resolved-contextual-charge}
\mathcal B_{I,R,t}^{(a)}
=
\sum_{\ell:\,\mathsf W_a(E_{R,t,\ell})}
\operatorname{ch}_I(E_{R,t,\ell})
\end{equation}
be the sum of the contextual charges in that truncated record carrying
authenticated provenance tag \(a\).  The survival gate, next-test flag,
stopping instruction, and terminal postprocessing remain in the complete
record.  Their deterministic filtration vertices have zero incremental
charge by \cref{lem-derived-constructor-zero-contextual-charge,thm-filtration-provenance-charging}.

At a common ceiling \(D\), define
\begin{equation}
\label{eq-saturated-bayesian-source-charge-envelope}
\mathcal B_{I,t}^{(a),\rm sat}(D)
=
\sup_{R\in\mathscr R_{I,H}^{\rm Bayes}(D)}
\mathcal B_{I,R,t}^{(a)},
\qquad \sup\varnothing=0,
\end{equation}
and use the declared family functional for the family-uniform version.  These
are restrictions of the existing contextual-charge envelopes, not new
structural channels.

When the complete pre-hit record contains a prospective master source, give
it the unique exact index
\begin{equation}
\label{eq-bayesian-exact-source-index}
\alpha(R,t)=(\chi,S_{\rm st},S_{\rm cf}),
\quad
\chi\in\mathfrak S_{\rm src}^{5},
\quad
\varnothing\ne S_{\rm st}\subseteq\mathfrak J_{\Abs}^{5},
\quad
S_{\rm cf}\subseteq\mathfrak J_{\Abs}^{5},
\end{equation}
from \cref{thm-universal-state-coefficient-source-decomposition}.  The
corresponding exact-cell envelope is the restriction of
\cref{eq-saturated-bayesian-source-charge-envelope} to records with that
index.  Exact cells are disjoint.  The five single-tag envelopes may overlap
and are used only as a nonnegative cover.

\end{definition}

\begin{theorem}[Exact source-resolved accounting of actionable Bayesian gain]
\label{thm-exact-source-resolved-structural-bayesian-gain}

Let \(H_B(n)\) be the largest refined horizon admitted by the complete
budget, and let
\begin{equation}
\label{eq-source-resolved-bayesian-charge-ceiling}
D_B^{\rm Bayes}(n)
=
\Psi_5^{\rm alg}\!\left(
\Psi_{\rm hit}(B^{\rm Bayes},n)
\right)
\end{equation}
be the class-preserving ceiling for the represented Bayesian record, its
next-test truncation, and its normalized five-family frontier.  Then every
record and every \(t<H_B(n)\) satisfy
\begin{align}
\label{eq-bayesian-gain-posterior-comparison}
g_{R,t}^{\rm Bayes}
&\le
\mathbb E_R\!\left[
\mathbf1_{\{\tau>t\}}
\bigl(\mathcal C_T(\widehat\Pi_t;\Pi_0)+e_t\bigr)
\right],\\
\label{eq-bayesian-gain-five-source-charge-bound}
g_{R,t}^{\rm Bayes}
&\le
eH_B(n)
\sum_{a\in\mathfrak J_{\Abs}^{5}}
\mathcal B_{I,R,t}^{(a)}.
\end{align}
Consequently the saturated actionable-gain profile
\begin{equation}
\label{eq-saturated-actionable-bayesian-gain-profile}
\operatorname{BayesGain}_{I,H}^{\rm sat}(B)
=
\max_{t<H}
\sup_{R\in\mathscr R_{I,H}^{\rm Bayes}(B)}
g_{R,t}^{\rm Bayes}
\end{equation}
obeys
\begin{equation}
\label{eq-saturated-bayesian-gain-five-source-envelope}
\operatorname{BayesGain}_{I,H}^{\rm sat}(B)
\le
eH_B(n)
\max_{t<H}
\sum_{a\in\mathfrak J_{\Abs}^{5}}
\mathcal B_{I,t}^{(a),\rm sat}
\bigl(D_B^{\rm Bayes}(n)\bigr).
\end{equation}

Write \(R_{\le t}\) for the complete record truncated immediately after the
deployed next test.  The Bayesian quantities are the corresponding
restrictions of the complete learning ledger:
\begin{equation}
\label{eq-bayesian-complete-learning-ledger-identification}
\mathcal B_{I,R,t}^{(a)}
=\mathcal L_{I,R_{\le t}}^{(a)},
\qquad
J_{R,t}^{\rm Bayes,(a)}
=J_{I,R_{\le t}}^{(a)}.
\end{equation}
Likelihood evaluation, evidence normalization, belief storage, posterior
approximation, acquisition optimization, and deployment therefore remain in
the complete source-resolved record of
\cref{def-complete-source-resolved-learning-ledger}.  The analytic posterior
is not inserted as a represented source.

For records carrying a prospective master source, the source-resolved class
is also the disjoint union of the exact indices in
\cref{eq-bayesian-exact-source-index}.  Hence its supremum is the maximum of
at most \(5(2^5-1)2^5=4960\) exact-cell suprema.  This reindexing retains
every likelihood evaluator, locator, normalizer, belief coordinate,
acquisition score, optimizer, transition, coefficient derivation, and
qualification gate in the same complete record.

\end{theorem}

\begin{proof}

Condition on \(\mathcal F_t\).  The first term in
\cref{eq-actionable-structural-bayesian-gain} is the exact conditional
next-contact identity of
\cref{eq-structural-bayes-exact-next-contact-identity}.  The
one-Lipschitz target-vulnerability comparison used in
\cref{thm-structural-bayesian-target-contact-accountability} gives
\cref{eq-bayesian-gain-posterior-comparison} after subtracting the prior
baseline and taking the positive part.

For the second bound, use the truncated computation that retains the complete
pre-hit record and stops after the next test.  Its success probability
\(h_t\) dominates \(g_{R,t}^{\rm Bayes}\), because the subtracted prior
contact is nonnegative.  Steps 2 and 3 of
\cref{thm-prospective-selector-structural-charge} give
\[
h_t
\le eH_B(n)
\sum_{a\in\mathfrak J_{\Abs}^{5}}
\mathcal B_{I,R,t}^{(a)}.
\]
This proves \cref{eq-bayesian-gain-five-source-charge-bound}.  Taking the
record supremum at the common normalized ceiling proves
\cref{eq-saturated-bayesian-gain-five-source-envelope}.

For a prospective master source, coefficient-complete normalization assigns
one support mode and one exact state--coefficient signature by
\cref{thm-universal-state-coefficient-source-decomposition}.  Those cells
are disjoint and cover the normalized prospective source class.  A supremum
over their union is therefore their maximum.  The complete-record assertion
is part of
\cref{def-complete-structural-bayesian-record,def-state-coefficient-source-provenance}.

\end{proof}

\begin{corollary}[Pre-hit status of Bayesian source charges]
\label{cor-prehit-status-bayesian-source-charges}

In the setting of
\cref{thm-exact-source-resolved-structural-bayesian-gain}, every nonzero
summand of \cref{eq-bayesian-source-resolved-contextual-charge} is assigned
to an authenticated vertex available before the next test.  Its complete
downstream record contains the represented likelihood use, belief update,
acquisition score, selected kernel, and target-contact normalization.
Consequently its source tag and exact master/signature index are determined
without using a later first-hit label.

More precisely, independent random coordinates are valid-Lift auxiliaries
and add zero charge.  Conditional on those coordinates and the earlier
state, normalization, storage, comparison, selection, survival, verification,
and stopping vertices are deterministic postprocessings and add zero charge.
Hence the charge-bearing frontier is the pre-hit ancestry of the deployed
update.  Caus, Lift, and Bdry retain the source index of that ancestry, while
its gradient-orthogonal complement is \(\Res\) and contributes zero
target-qualified source charge.

\end{corollary}

\begin{proof}

Apply \cref{thm-filtration-provenance-charging} to the complete truncated
record used in the proof of
\cref{thm-exact-source-resolved-structural-bayesian-gain}.  Its second and
third cases remove independent auxiliary randomness and every deterministic
descendant from the incremental charge ledger.  The remaining nonzero
vertices occur before the selected test and belong to the normalized
backward slice of the likelihood-to-action chain.  Their five-family tags
are fixed by \cref{thm-five-family-abs-provenance}, and their exact master and
state--coefficient indices are fixed by
\cref{thm-universal-state-coefficient-source-decomposition}.  Source
inheritance and the residual assignment are
\cref{prop-derived-channels-create-no-structure,thm-exhaustive-residual-normal-form}.

\end{proof}

\begin{theorem}[Noise-contracted source information for Bayesian updates]
\label{thm-noise-contracted-source-resolved-bayesian-gain}

Enumerate the target-dependent coordinates actually revealed before the
next test as \(Y_1,\ldots,Y_{L_t}\).  Fix the canonical anchored source
assignment of \cref{cor-universal-syntax-affordable-charge-cover}; write
\(a_j\in\mathfrak J_{\Abs}^{5}\) for the assigned source of \(Y_j\).  With
the channel-capacity, available-information, and SDPI quantities of
\cref{def-controlled-channel-information-ledger}, set
\begin{equation}
\label{eq-source-resolved-bayesian-information-ledger}
J_{R,t}^{\rm Bayes,(a)}
=
\sum_{\substack{1\le j\le L_t\\a_j=a}}
\mathbb E\min\{c_j,\vartheta_j a_j^{\rm info}\},
\qquad
J_{R,t}^{\rm Bayes,src}
=\sum_{a\in\mathfrak J_{\Abs}^{5}}
J_{R,t}^{\rm Bayes,(a)}.
\end{equation}
The superscript on \(a_j^{\rm info}\) distinguishes the available-information
scalar from the source tag \(a_j\).  Deterministic target-independent
postprocessing contributes zero to every summand.

Let \(\beta_{R,t}^{\rm Bayes}\) be the next-contact probability after
retaining the public pre-hit record and the selected kernel while resampling
the target independently from its declared reference conditional law, and
put
\begin{equation}
\label{eq-source-resolved-bayesian-neutral-excess}
\beta_{R,t}^{\rm Bayes,+}
=
\left[
\beta_{R,t}^{\rm Bayes}
-\mathbb E_R[\mathbf1_{\{\tau>t\}}]\mathcal V_T(\Pi_0)
\right]_+.
\end{equation}
Then
\begin{equation}
\label{eq-noise-contracted-source-resolved-bayesian-gain}
g_{R,t}^{\rm Bayes}
\le
\min\left\{
1,
\beta_{R,t}^{\rm Bayes,+}
+\sqrt{\frac{\log2}{2}J_{R,t}^{\rm Bayes,src}},
eH_B(n)\sum_{a\in\mathfrak J_{\Abs}^{5}}
\mathcal B_{I,R,t}^{(a)}
\right\}.
\end{equation}
Every structural Rademacher, sequential Rademacher, martingale PAC--Bayes,
Blackwell, capacity, or dynamical candidate converted to this same gain under
the same conditional law may be added to the recordwise minimum.  An
unavailable candidate has neutral value one.

For a homogeneous binary symmetric reveal of noise rate \(\eta\), the
channel candidate may use
\begin{equation}
\label{eq-binary-noise-source-sdpi-candidate}
\vartheta_j(\eta)=(1-2\eta)^2.
\end{equation}
Thus the square-root information term carries one factor \(1-2\eta\).
Noise correction of a learned clean score remains the inverse factor
\((1-2\eta)^{-1}\) in
\cref{thm-structural-pac-bayes-label-noise,thm-saturated-statistical-attack-alignment}; the contraction and correction factors have different directions and are retained separately.

\end{theorem}

\begin{proof}

The ordered-reveal chain rule and SDPI bound are
\cref{thm-adaptive-acquisition-information-target-contact}.  Its conditional
decoupling inequality bounds next contact by
\(\beta_{R,t}^{\rm Bayes}+\sqrt{(\log2)J_{R,t}^{\rm Bayes,src}/2}\).
Subtract the prior baseline, take the positive part, and combine with
\cref{eq-bayesian-gain-five-source-charge-bound}.  Every additional listed
candidate bounds the same record, event, normalization, and conditional law,
so their minimum is valid.  The binary symmetric SDPI candidate is the
standard maximal-correlation coefficient of the declared flip channel; its
square root gives the displayed vulnerability attenuation.  The inverse
noise factor belongs to the separate clean-risk conversion already proved in
the cited learning theorems.

\end{proof}

\begin{corollary}[Adaptive composition and no-source Bayesian consequence]
\label{cor-source-resolved-bayesian-adaptive-composition}

For every complete structural Bayesian record,
\begin{equation}
\label{eq-source-resolved-bayesian-arrival-composition}
\Pr_R\{\tau\le H\}
\le
H\mathcal V_T(\Pi_0)
+\sum_{t=0}^{H-1}g_{R,t}^{\rm Bayes}.
\end{equation}
Consequently, if at the common compiled ceiling
\begin{equation}
\label{eq-source-resolved-bayesian-negligible-inputs}
\mathcal B_{\mathfrak M,t}^{(a),\rm sat}
\le\varepsilon_{a,t}
\qquad
\bigl(a\in\mathfrak J_{\Abs}^{5},\ 0\le t<H\bigr),
\end{equation}
then
\begin{equation}
\label{eq-source-resolved-bayesian-no-source-bound}
\operatorname{BayesArr}_{\mathfrak M,H}^{\rm sat}(B)
\le
H\mathcal V_T(\Pi_0)
+eH_B(n)\sum_{t=0}^{H-1}
\sum_{a\in\mathfrak J_{\Abs}^{5}}\varepsilon_{a,t}.
\end{equation}
The same conclusion may use the smaller recordwise minimum in
\cref{eq-noise-contracted-source-resolved-bayesian-gain}.  Hence a
polynomial horizon preserves negligibility of the neutral contact and all
source-resolved gain inputs.

\end{corollary}

\begin{proof}

Subtract and add the prior-vulnerability baseline in each exact next-contact
identity and discard negative excess.  Summing the disjoint first-contact
events proves \cref{eq-source-resolved-bayesian-arrival-composition}.  Apply
\cref{eq-saturated-bayesian-gain-five-source-envelope} at each time and the
monotone family functional to obtain
\cref{eq-source-resolved-bayesian-no-source-bound}.

\end{proof}

\begin{figure}[tbp]
\centering
\resizebox{\linewidth}{!}{%
\begin{tikzpicture}[fw diagram,node distance=7mm and 9mm]
  \node[fw source,text width=2.3cm] (prior)
    {neutral prior\\contact};
  \node[fw provenance,text width=2.7cm,right=of prior] (reveals)
    {ordered noisy reveals\\source-resolved SDPI};
  \node[fw use,text width=2.6cm,right=of reveals] (update)
    {likelihood, normalizer\\and belief update};
  \node[fw geom,text width=2.7cm,right=of update] (action)
    {acquisition score\\and next selector};
  \node[fw terminal,text width=2.5cm,right=of action] (gain)
    {actionable target\\gain};
  \node[fw residual,text width=4.4cm,below=10mm of update] (ledger)
    {five source charges and exact master/signature cells};
  \draw[fw arrow] (prior)--(reveals);
  \draw[fw arrow] (reveals)--(update);
  \draw[fw arrow] (update)--(action);
  \draw[fw arrow] (action)--(gain);
  \draw[fw arrow] (reveals)--(ledger);
  \draw[fw arrow] (update)--(ledger);
  \draw[fw arrow] (ledger)-|(gain);
\end{tikzpicture}%
}
\caption{Source-resolved structural Bayesian gain.  The operational scalar is
the improvement in the represented next test.  Every noisy reveal retains its
source and contraction coefficient; deterministic updating adds no new source
charge; and the complete deployed rule is evaluated by same-record
information, learning, dynamical, and contextual-charge certificates.}
\label{fig-source-resolved-structural-bayesian-gain}
\end{figure}

\subsection{Optimal acquisition and structural routing}
\label{subsec-optimal-structural-bayes-acquisition}

\begin{theorem}[Optimal structural Bayesian acquisition]
\label{thm-optimal-structural-bayesian-acquisition}

Restrict the suffix-record class in
\cref{thm-budgeted-active-acquisition-recursion} to complete structural
Bayesian records and write
\(V_t^{\rm Bayes}(h,\widehat\pi;B)\) for its target-arrival supremum.
Then
\begin{equation}
\label{eq-optimal-structural-bayes-resource-recursion}
V_t^{\rm Bayes}(h,\widehat\pi;B)
=
\sup_{(a,\varrho,\widehat U)\,:\,
c_t^{\rm Bayes}(a,\varrho,\widehat U;h,\widehat\pi)\preceq B}
\left\{
p_T(h,a)+
\mathbb E\!\left[
\mathbf1_{\{\tau>t+1\}}
A_{\varrho(H_{t+1},\widehat U(\widehat\pi,Y_{t+1}))}^{H-t-1}
\mid h,a
\right]
\right\},
\end{equation}
where \(\widehat U\) is the represented belief update and
\(c_t^{\rm Bayes}\) is the complete active cost including its likelihood,
normalization, memory, comparison, continuation, and deployment slots.
The stopped occupation-flow primal and Bellman-barrier dual in
\cref{thm-exact-affordable-occupation-flow-primal,thm-bellman-barrier-duality-affordable-arrival}
hold after restricting their record and first-step sets to these Bayesian
records.

Always
\begin{equation}
\label{eq-bayesian-active-profile-inclusion}
\operatorname{BayesArr}_{I,H}^{\rm sat}(B)
\le\operatorname{ActArr}_{I,H}^{\rm sat}(B).
\end{equation}

\end{theorem}

\begin{proof}

Decompose a nonterminal structural Bayesian suffix record into its first
action, represented belief update, and represented continuation rule.  The
converse concatenation is a complete structural Bayesian suffix record and
has cost \(c_t^{\rm Bayes}\).  The disjoint first-contact decomposition
therefore proves \cref{eq-optimal-structural-bayes-resource-recursion}, by
the proof of \cref{thm-budgeted-active-acquisition-recursion}.  Restricting
the occupation and Bellman record sets preserves their balance identities,
telescoping proof, and exact suffix-value barrier.

Every structural Bayesian record is a complete active record, proving
\cref{eq-bayesian-active-profile-inclusion}.

\end{proof}

\begin{theorem}[Universal Bayesianization of represented active computation]
\label{thm-universal-bayesianization-active-computation}

Let \(R\in\mathscr R_{I,H}^{\rm act}(B)\) be a complete active record.  There
is a family-uniform compiler \(R\mapsto\mathfrak B(R)\) and a
class-preserving resource majorant
\begin{equation}
\label{eq-universal-bayesianization-ceiling}
B^{\rm uBayes}
=
\Psi_{\rm uBayes}(B,H,n)
\end{equation}
such that \(\mathfrak B(R)\in
\mathscr R_{I,H}^{\rm Bayes}(B^{\rm uBayes})\) and the following data are
identical in \(R\) and \(\mathfrak B(R)\):
\begin{enumerate}[label=\textup{(\roman*)}]
\item the public pre-hit filtration and represented sufficient history;
\item every query, observation, action, fresh-randomness, and transition
      kernel;
\item the verifier, stopping rule, output converter, and continuation rule;
\item the complete execution law, first-hit law, and target-arrival
      probability.
\end{enumerate}

The compiler attaches the declared full-support prior \(\Pi_0\), the
represented constant belief \(\widehat\Pi_t=\Pi_0\), the constant likelihood
score one, the evidence normalizer one, and the identity belief update.  The
original action and continuation rules retain the complete represented history
as their sufficient-memory input and need not read the attached belief.  The
analytic posterior
\[
\Pi_t=\Pr\{\Theta\in\cdot\mid\mathcal F_t\}
\]
is therefore available for conditional-probability identities without being
materialized by the compiled computation.

Conversely, forgetting the belief-specific slots of a complete structural
Bayesian record produces a complete active record with the same execution law
and no larger cost.  Consequently
\begin{align}
\label{eq-active-to-universal-bayesian-ceiling}
\operatorname{ActArr}_{I,H}^{\rm sat}(B)
&\le
\operatorname{BayesArr}_{I,H}^{\rm sat}
\bigl(\Psi_{\rm uBayes}(B,H,n)\bigr),\\
\label{eq-universal-bayesian-to-active-ceiling}
\operatorname{BayesArr}_{I,H}^{\rm sat}(B)
&\le
\operatorname{ActArr}_{I,H}^{\rm sat}(B).
\end{align}
The same inequalities hold for one family-uniform derivation and every
declared monotone positively homogeneous family functional.  Hence the two
profiles coincide on every transfer-equivalence class of ceilings; in
particular,
\begin{equation}
\label{eq-polynomial-bayesian-active-ceiling-equality}
\operatorname{BayesArr}^{\rm sat}_{\rm poly}
=
\operatorname{ActArr}^{\rm sat}_{\rm poly}.
\end{equation}

\end{theorem}

\begin{proof}

The definition of a complete structural Bayesian record in
\cref{def-complete-structural-bayesian-record} begins with a complete active
record.  Retain every field of \(R\).  Append a reference to the fixed prior,
the constant likelihood and normalizer, and the identity update.  Their codes
are independent of the target and of the history.  Evaluation and storage over
the \(H\) clocked steps have a uniform resource bound obtained by adding their
existing belief, likelihood, normalization, memory, and clock slots to
\cref{eq-complete-master-active-sampling-cost}.  Axioms B1--B4 and computation
closure give the class-preserving majorant
\cref{eq-universal-bayesianization-ceiling}.

The attached fields are not ancestors of the original action or continuation
vertices.  Evaluation of the retained record therefore produces exactly the
same conditional kernels at every history.  Induction on the clock gives
equality of the finite-dimensional path laws, hence equality of the first-hit
and target-arrival laws.  The analytic posterior is the conditional law of
this common path measure and requires no represented posterior evaluator.

Deleting the appended belief-specific vertices leaves the original complete
active record, which proves the converse map.  Taking the appropriate
suprema gives
\cref{eq-active-to-universal-bayesian-ceiling,eq-universal-bayesian-to-active-ceiling}.
Both compilers use one family-uniform syntax transformation, so the same proof
commutes with the family functional.  Mutual class-preserving domination gives
the transfer-class equality and its polynomial specialization.

\end{proof}

\begin{corollary}[Universal represented-computation Bayesian ceiling]
\label{cor-universal-represented-computation-bayesian-ceiling}

Let \(\mathsf{Alg}_{n,\le B}\) be the represented computation class of
\cref{def-uniform-saturated-success-profile}, with finite refined horizon
\(H_B(n)\).  There is a class-preserving compiled ceiling
\begin{equation}
\label{eq-universal-computation-bayesian-ceiling-budget}
B^{\rm compBayes}
=
\Psi_{\rm uBayes}
\bigl(\Psi_{\rm act}(B,H_B,n,\mathbf 0),H_B,n\bigr)
\end{equation}
for which
\begin{equation}
\label{eq-universal-computation-bayesian-ceiling}
\operatorname{Succ}_{\mathfrak M,n}(B)
\le
\operatorname{BayesArr}_{\mathfrak M,n,H_B}^{\rm sat}
\bigl(B^{\rm compBayes}\bigr).
\end{equation}
The reverse inclusion holds by forgetting the Bayesian and active-interface
annotations.  Thus the represented-computation, active, and structural
Bayesian optima coincide on the declared saturated transfer class.  In the
polynomial specialization,
\begin{equation}
\label{eq-universal-computation-active-bayesian-polynomial-ceiling}
\operatorname{Succ}_{\mathfrak M,n}^{\rm poly,sat}
=
\operatorname{ActArr}_{\mathfrak M,n}^{\rm poly,sat}
=
\operatorname{BayesArr}_{\mathfrak M,n}^{\rm poly,sat}.
\end{equation}

Moreover, at the common source-charge ceiling of
\cref{thm-exact-source-resolved-structural-bayesian-gain},
\begin{equation}
\label{eq-universal-bayesian-source-charge-ceiling}
\operatorname{Succ}_{\mathfrak M,n}(B)
\le
H_B\mathfrak M_n(\mathbf1)\mathcal V_T(\Pi_0)
+eH_B^2
\sum_{a\in\mathfrak J_{\Abs}^{5}}
\mathcal B_{\mathfrak M}^{(a),\rm sat}
\bigl(D_{B^{\rm compBayes}}^{\rm Bayes}(n)\bigr).
\end{equation}

\end{corollary}

\begin{proof}

Use the operational embedding of
\cref{thm-algorithm-kernel-embedding,cor-all-turing-machines-represented}.
Regard its public history as the sufficient-memory state, take a null query,
and retain its represented transition, verifier, stopping, and output rules.
Channel exhaustiveness and computation closure supply the complete active
record at the majorant \(\Psi_{\rm act}\); presentation-admissible oracle
interfaces remain fields of that same record.  Apply
\cref{thm-universal-bayesianization-active-computation} and take the
family-uniform supremum to prove
\cref{eq-universal-computation-bayesian-ceiling}.  Forgetting annotations
preserves the execution law and proves the reverse transfer-class inclusion.

Finally apply the prospective selector charge theorem
\cref{thm-prospective-selector-structural-charge} in its source-resolved
Bayesian form
\cref{eq-source-resolved-bayesian-arrival-composition,eq-saturated-bayesian-gain-five-source-envelope}
to the compiled record.  There are at most \(H_B\) pre-hit gains, each bounded
by the five-source envelope with factor \(eH_B\).  Monotonicity and positive
homogeneity of the family functional give
\cref{eq-universal-bayesian-source-charge-ceiling}.

\end{proof}

\begin{figure}[tbp]
\centering
\resizebox{0.96\linewidth}{!}{%
\begin{tikzpicture}[fw diagram,node distance=8mm and 10mm]
  \node[fw source,text width=2.5cm] (comp)
    {represented\\computation};
  \node[fw provenance,text width=2.6cm,right=of comp] (active)
    {null-query active\\history record};
  \node[fw use,text width=2.8cm,right=of active] (bayes)
    {constant-belief\\Bayesianization};
  \node[fw geom,text width=2.7cm,right=of bayes] (selector)
    {unchanged pre-hit\\selector};
  \node[fw terminal,text width=2.8cm,right=of selector] (bound)
    {five-source\\success bound};
  \draw[fw arrow] (comp)--(active);
  \draw[fw arrow] (active)--(bayes);
  \draw[fw arrow] (bayes)--(selector);
  \draw[fw arrow] (selector)--(bound);
  \draw[fw arrow,dashed,bend left=32]
    (bayes) to node[below,font=\scriptsize] {forget annotations} (active);
\end{tikzpicture}%
}
\caption{Universal Bayesianization preserves the execution law and the
represented pre-hit selector.  The exact posterior is an analytic conditional
law; the compiled computation carries only the constant represented prior and
identity update.  Prospective structural charging therefore bounds the common
computation, active, and Bayesian ceiling.}
\label{fig-universal-bayesian-computation-ceiling}
\end{figure}

\begin{theorem}[Structural provenance of affordable posterior concentration]
\label{thm-affordable-posterior-concentration-source-routing}

Suppose a complete structural Bayesian record represents and deploys a
positive value of
\(\mathcal C_T(\widehat\Pi_t;\Pi_0)\) through its represented likelihood,
normalizer, belief, acquisition score, optimizer, and next-test kernel.  The
deployed target-aligned quantity is the actionable gain
\(g_{R,t}^{\rm Bayes}\) of
\cref{eq-actionable-structural-bayesian-gain}.  Its complete pre-hit
derivation is temporally admissible, and its source-resolved contextual
charge is bounded by
\cref{thm-exact-source-resolved-structural-bayesian-gain}.  At the common ceiling
\(B^{\rm Bayes}=\Psi_{\rm Bayes}(B,H,n,\mathbf m,\boldsymbol\zeta)\), its complete backward cut belongs
to exactly one of
\begin{equation}
\label{eq-posterior-concentration-five-mode-routing}
\mathcal C_X,\quad\mathcal C_A,\quad
\mathcal C_{Q,\mathrm{ex}},\quad
\mathcal C_{Q,\mathrm{rev}},\quad
\mathcal C_{Q,\mathrm{nr}}.
\end{equation}
Consequently
\begin{equation}
\label{eq-posterior-concentration-five-mode-profile}
\operatorname{PostSrc}_{I,H}^{\rm sat}(B^{\rm Bayes})
=\max_{c\in\mathfrak C_{\rm master}}
\operatorname{PostSrc}_{I,H;c}^{\rm sat}(B^{\rm Bayes}).
\end{equation}

The corresponding actionable-gain class has the quantitative envelope
\begin{equation}
\label{eq-posterior-source-routed-actionable-gain}
\operatorname{BayesGain}_{I,H}^{\rm sat}(B)
\le
eH_B(n)\max_{t<H}
\sum_{a\in\mathfrak J_{\Abs}^{5}}
\mathcal B_{I,t}^{(a),\rm sat}
\bigl(D_B^{\rm Bayes}(n)\bigr),
\end{equation}
and each complete prospective source record has the unique exact
source--state--coefficient index of
\cref{eq-bayesian-exact-source-index}.  Thus the numerical ledger has no
unrestricted direct-source remainder: a positive deployed gain is evaluated
through its five contextual source charges, while the master index records
its exact support and mixed provenance.

For one record, every finite information, Blackwell, filter, capacity,
observation-resolution, or master-source candidate must be converted to the
same target-vulnerability excess under the same conditional law.  The
smallest such candidate is an upper bound for that record's concentration;
an unavailable converter has neutral value one.  Thus posterior
concentration is a numerical view of the existing master profile and not an
additional provenance family.

\end{theorem}

\begin{proof}

The represented concentration statistic, the score derived from it, and the
next-test kernel are available before the action.  Their complete derivation
is charged, so the deployed record satisfies
\cref{def-pre-hit-temporally-admissible-source}.  Apply channel
exhaustiveness, backward ancestry, and the master classification in
\cref{thm-complete-generic-source-theory,thm-exhaustive-five-family-geom-source-classification}.
The five master families are disjoint and exhaustive, so partitioning the
record class and taking the saturated supremum proves
\cref{eq-posterior-concentration-five-mode-profile}.  Every listed finite
candidate bounds the same record and coordinate by its defining comparison;
their minimum is therefore valid.  Missing comparisons contribute the
neutral value and supply no numerical conclusion.

The quantitative gain bound is
\cref{eq-saturated-bayesian-gain-five-source-envelope}.  The unique refined
index follows from
\cref{thm-universal-state-coefficient-source-decomposition}.  These two
decompositions have different roles: contextual charge bounds the deployed
gain, while the exact master/signature index records the complete source
carrier without assigning a new mode.

\end{proof}

\begin{corollary}[Learned structural Bayesian control]
\label{cor-learned-structural-bayesian-control}

Let \(\widehat R\) be a learned structural Bayesian record satisfying the
simultaneous action-value and optimization inequalities of
\cref{thm-learned-active-rule-finite-horizon-loss}, with their analytic
action-value comparator restricted to the structural Bayesian suffix
optimum in \cref{eq-optimal-structural-bayes-resource-recursion}.  Let
\(\Delta_{\rm fil}^{H}\) be the represented path/event comparison in
\cref{thm-controlled-pomdp-filter-transfer}.  Then
\begin{align}
\label{eq-learned-structural-bayes-upper}
A_{\widehat R}^{H}
&\le\operatorname{BayesArr}_{I,H}^{\rm sat}(B),\\
\label{eq-learned-structural-bayes-optimality-loss}
\operatorname{BayesArr}_{I,H}^{\rm sat}(B)-A_{\widehat R}^{H}
&\le
\mathbb E_{\widehat R}\sum_{t<H}
\mathbf1_{\{\tau>t\}}
\bigl(2\Gamma_t+\varepsilon_{{\rm opt},t}\bigr)
+\Delta_{\rm fil}^{H}.
\end{align}
At a compiled ceiling that contains the universal Bayesianization of the
complete active class, the same inequality is
\begin{equation}
\label{eq-learned-bayesian-universal-ceiling-loss}
0\le
\operatorname{ActArr}_{I,H}^{\rm sat}(B_0)-A_{\widehat R}^{H}
\le
\mathbb E_{\widehat R}\sum_{t<H}
\mathbf1_{\{\tau>t\}}
\bigl(2\Gamma_t+\varepsilon_{{\rm opt},t}\bigr)
+\Delta_{\rm fil}^{H},
\end{equation}
where \(B=\Psi_{\rm uBayes}(B_0,H,n)\).  After the operational compilation of
\cref{cor-universal-represented-computation-bayesian-ceiling}, the first term
is the saturated represented-computation optimum at its corresponding
transfer-equivalent ceiling.
The radii may be supplied by structural Rademacher, sequential, or
martingale PAC--Bayes bounds.  Their PAC--Bayes posterior \(Q\) remains
separate from the task belief \(\widehat\Pi_t\), and both costs occur in the
same complete record.

\end{corollary}

\begin{proof}

The first inequality is membership in the structural Bayesian record class.
Apply \cref{thm-learned-active-rule-finite-horizon-loss} to the exact-belief
reference execution and then
\cref{thm-controlled-pomdp-filter-transfer} to compare that execution with
the represented filter.  The triangle inequality gives
\cref{eq-learned-structural-bayes-optimality-loss}.  Substitute the profile
equality of \cref{thm-universal-bayesianization-active-computation} to obtain
\cref{eq-learned-bayesian-universal-ceiling-loss}.

\end{proof}

\begin{figure}[tbp]
\centering
\resizebox{\linewidth}{!}{%
\begin{tikzpicture}[fw diagram,node distance=7mm and 9mm]
  \node[fw source,text width=2.4cm] (prior)
    {declared task prior\\\(\Pi_0\)};
  \node[fw provenance,text width=2.7cm,right=of prior] (obs)
    {represented action\\and observation};
  \node[fw use,text width=2.7cm,right=of obs] (factor)
    {charged likelihood\\and Bayes factor};
  \node[fw geom,text width=2.7cm,right=of factor] (belief)
    {represented belief\\\(\widehat\Pi_t\)};
  \node[fw box,text width=2.8cm,right=of belief] (select)
    {acquisition and\\transition rule};
  \node[fw terminal,text width=2.6cm,right=of select] (contact)
    {target contact\\or arrival};
  \node[fw residual,text width=3.6cm,below=10mm of belief] (modes)
    {existing five master modes\\and same-record certificates};
  \draw[fw arrow] (prior)--(obs);
  \draw[fw arrow] (obs)--(factor);
  \draw[fw arrow] (factor)--(belief);
  \draw[fw arrow] (belief)--(select);
  \draw[fw arrow] (select)--(contact);
  \draw[fw arrow] (belief)--(modes);
  \draw[fw arrow] (modes)-|(contact);
\end{tikzpicture}%
}
\caption{Optimal structural Bayesian learning.  Target-dependent data enter
through represented observation channels, and every likelihood,
normalization, belief, acquisition, and transition operation is charged.
Posterior concentration is evaluated through the existing master
classification before it contributes to target contact.}
\label{fig-optimal-structural-bayesian-learning}
\end{figure}

\subsection{Optimal structural information extraction}
\label{subsec-optimal-structural-information-extraction}

The next result identifies the complete computational ceiling with optimal
structural Bayesian acquisition and resolves its target gain into the existing
structural modalities.  The optimization is over represented acquisition
records at one charged ceiling; the analytic posterior supplies conditional
identities and is materialized only when its evaluator is represented and
charged.

\begin{definition}[Optimal structural extraction certificate]
\label{def-optimal-structural-extraction-certificate}

Fix a finite target-index carrier \(\mathsf\Theta\), target vulnerability
\(\mathcal V_T\) from
\cref{def-affordable-structural-posterior-concentration}, a horizon \(H\),
and a complete structural Bayesian record \(R\).  At time \(t<H\), define
the same-record information candidate
\begin{equation}
\label{eq-optimal-extraction-information-candidate}
U_{R,t}^{\rm info}
=\min\left\{1,
\beta_{R,t}^{\rm Bayes,+}
+\sqrt{\frac{\log2}{2}J_{R,t}^{\rm Bayes,src}}
\right\},
\end{equation}
and the source candidate
\begin{equation}
\label{eq-optimal-extraction-source-candidate}
U_{R,t}^{\rm src}
=\min\left\{1,
eH_B(n)\sum_{a\in\mathfrak J_{\Abs}^{5}}
\mathcal B_{I,R,t}^{(a)}
\right\}.
\end{equation}
Let \(U_{R,t}^{\Caus}\) be the probability-normalized same-record Caus
contact coordinate of
\cref{def-coordinatewise-dynamical-certificate}.  Let
\(U_{R,t}^{\Lift}\) be the minimum of one,
\(U_{\Lift}^{H}\), \(U_{\Lift,\mathrm{noise}}^{H}\), and every represented
probability-normalized finite-horizon comparison derived from
\(U_{\Lift}^{\lambda}\) or
\(U_{\Lift,\mathrm{noise}}^{\lambda}\) in that same record.  Missing
comparisons have their neutral values.  The boundary candidate
\(U_{R,t}^{\Bdry}\) is the exact represented next-contact probability when
the boundary readout is available before the transition and is one otherwise.
The residual candidate is zero on the gradient-orthogonal component and is
not used as a bound for the qualified component.

The complete extraction certificate is
\begin{equation}
\label{eq-optimal-structural-extraction-certificate}
U_{R,t}^{\rm opt}
=\min\left\{
1,U_{R,t}^{\rm info},U_{R,t}^{\rm src},U_{R,t}^{\Caus},
U_{R,t}^{\Lift},U_{R,t}^{\Bdry},U_{R,t}^{\rm Bell}
\right\},
\end{equation}
where \(U_{R,t}^{\rm Bell}\) is the same-record Bellman-support minimum of
\cref{eq-recordwise-bellman-support-certificate-minimum} evaluated at the
zero barrier for the next-contact event.  Each finite entry
must bound the same next-contact gain under the same surviving-history law.

\end{definition}

\begin{theorem}[Bayes--Fisher optimality of a represented observation]
\label{thm-bayes-fisher-optimal-represented-observation}

Condition on a surviving pre-hit history and a represented acquisition
action.  Let \(\pi\) be the current full-support task belief, restrict to the
represented positive-likelihood support of the returned observation, let
\(\ell_y(\theta)=\log L(y\mid\theta)\) be the represented log-likelihood of
the returned observation, and put
\[
\pi^y(\theta)
=\frac{\pi(\theta)e^{\ell_y(\theta)}}
       {\sum_u\pi(u)e^{\ell_y(u)}}.
\]
For every probability law \(q\ll\pi\),
\begin{equation}
\label{eq-bayes-kl-variational-identity}
\mathbb E_q\ell_y-D_{\rm KL}(q\Vert\pi)
=\log\mathbb E_\pi e^{\ell_y}
-D_{\rm KL}(q\Vert\pi^y).
\end{equation}
Hence \(\pi^y\) is the unique maximizer of represented likelihood gain minus
KL displacement.  Along
\[
\pi_u^y(\theta)
=\frac{\pi(\theta)e^{u\ell_y(\theta)}}
       {\sum_v\pi(v)e^{u\ell_y(v)}},
\qquad 0\le u\le1,
\]
the exact Fisher--Rao flow and action are
\begin{align}
\label{eq-bayes-fisher-rao-flow}
\dot\pi_u^y(\theta)
&=\pi_u^y(\theta)
  \bigl(\ell_y(\theta)-\mathbb E_{\pi_u^y}\ell_y\bigr),\\
\label{eq-bayes-fisher-rao-action}
\lVert\dot\pi_u^y\rVert_{\rm FR}^2
&=\operatorname{Var}_{\pi_u^y}(\ell_y).
\end{align}
Thus the exact Bayesian update is the Fisher--Rao steepest likelihood flow
for the fixed represented observation.  Its likelihood construction,
normalization, precision, storage, and deployment costs are those of
\cref{def-complete-structural-bayesian-record}.

\end{theorem}

\begin{proof}
Substitute
\(
\log(\pi^y(\theta)/\pi(\theta))
=\ell_y(\theta)-\log\mathbb E_\pi e^{\ell_y}
\)
in \(D_{\rm KL}(q\Vert\pi^y)\) and rearrange.  Relative entropy is
nonnegative and strictly convex on the support of \(\pi\), proving the
variational assertion.  Differentiating the normalized exponential path
gives \cref{eq-bayes-fisher-rao-flow}; substitution in
\(\lVert v\rVert_{\rm FR}^2=\sum_\theta v(\theta)^2/\pi_u^y(\theta)\)
gives \cref{eq-bayes-fisher-rao-action}.  The cost assertion is the complete
Bayesian-record definition.
\end{proof}

\begin{theorem}[Optimal structural extraction equals the saturated ceiling]
\label{thm-optimal-structural-extraction-saturated-ceiling}

Let \(B\) be a declared represented-computation ceiling with refined horizon
\(H_B(n)\), and let \(B^{\rm compBayes}\) be
\cref{eq-universal-computation-bayesian-ceiling-budget}.  At finite budget the
two directions are the compiled-ceiling comparisons of
\cref{cor-universal-represented-computation-bayesian-ceiling}.  On the
polynomial saturated transfer class they give the exact identity
\begin{equation}
\label{eq-optimal-extraction-three-profile-equality}
\operatorname{Succ}_{\mathfrak M,n}^{\rm poly,sat}
=\operatorname{ActArr}_{\mathfrak M,n}^{\rm poly,sat}
=\operatorname{BayesArr}_{\mathfrak M,n}^{\rm poly,sat}.
\end{equation}
The Bayesian member is optimized by the resource-valued Bellman recursion in
\cref{eq-optimal-structural-bayes-resource-recursion}; each fixed observation
is processed optimally by
\cref{thm-bayes-fisher-optimal-represented-observation}.  For every complete
record and every \(t<H_B(n)\),
\begin{equation}
\label{eq-optimal-extraction-recordwise-gain-bound}
g_{R,t}^{\rm Bayes}
\le U_{R,t}^{\rm opt}.
\end{equation}
Consequently,
\begin{equation}
\label{eq-optimal-extraction-arrival-bound}
\operatorname{Succ}_{\mathfrak M,n}(B)
\le
H_B\mathfrak M_n(\mathbf1)\mathcal V_T(\Pi_0)
+\mathfrak M_n\!\left(
\left(
\sum_{t=0}^{H_B-1}
\sup_{R\in\mathscr R_{I,H_B}^{\rm Bayes}(B^{\rm compBayes})}
U_{R,t}^{\rm opt}
\right)_{I\in\mathfrak I_n}
\right).
\end{equation}

The structural roles in this bound are exhaustive.  \(\Geom\) supplies the
reversible information geometry; exact and compensated \(\Abs\) supply
qualified sufficient fibers; \(\Caus\) converts the qualified geometry into
target-directed current; \(\Lift\) transports belief, memory, and recurrent
state with its comparison defects; and \(\Bdry\) supplies the represented
target interface.  The gradient-orthogonal \(\Res\) component has zero
target-qualified gain.  Add, Mult, Ord, Sym, and Filt are the exact
constructor provenance of these same source records and create no additional
modality.

\end{theorem}

\begin{proof}
The profile equality is
\cref{cor-universal-represented-computation-bayesian-ceiling}.  The Bellman
recursion is exact because every complete suffix record decomposes into a
first acquisition, returned observation, represented update, and complete
continuation, and conversely those data concatenate to a legal suffix record.
The fixed-observation optimizer is
\cref{thm-bayes-fisher-optimal-represented-observation}.

The information and source entries of
\cref{eq-optimal-structural-extraction-certificate} bound the deployed gain
by
\cref{thm-noise-contracted-source-resolved-bayesian-gain,thm-exact-source-resolved-structural-bayesian-gain}.
The Caus, Lift, boundary, and Bellman entries enter only through their
same-record target-contact comparisons in
\cref{def-coordinatewise-dynamical-certificate,thm-structural-evaluation-affordable-bellman-support}.
Taking the minimum proves
\cref{eq-optimal-extraction-recordwise-gain-bound}.  Summing the disjoint
first-contact gains and using
\cref{eq-source-resolved-bayesian-arrival-composition} proves
\cref{eq-optimal-extraction-arrival-bound}.  Channel exhaustiveness, Hodge
orthogonality, source inheritance, and the five-family theorem give the final
modality statement.
\end{proof}

\begin{corollary}[Noise, sample, and target-size extraction bound]
\label{cor-noise-sample-target-size-optimal-extraction}

Let \(N=|\mathsf\Theta|\).  Suppose at most \(m\) target-dependent
observations are revealed before deployment.  For reveal \(j\), let \(c_j\)
be its represented conditional channel-capacity bound in bits,
\(\vartheta_j\) its represented SDPI coefficient, and \(a_j^{\rm info}\)
its available conditional information.  Then every complete record satisfies
\begin{equation}
\label{eq-general-channel-optimal-extraction-information}
J_{R,t}^{\rm Bayes,src}
\le
\min\left\{
\log_2N,
\sum_{j\le m}c_j,
\sum_{j\le m}\vartheta_j a_j^{\rm info}
\right\}.
\end{equation}
Every target-dependent advice, oracle answer, reward, verifier outcome, or
feedback coordinate appears as one of the ordered reveals with its own
capacity, available-information, derivation-cost, and source tag.  Selecting
or postprocessing a coordinate already present in the pre-hit record adds zero
conditional information and retains its complete computational cost.

For \(m\) binary symmetric reveals with common crossover probability
\(0\le\eta\le1/2\),
\begin{equation}
\label{eq-bsc-optimal-extraction-information}
J_{R,t}^{\rm Bayes,src}
\le
J_{m,\eta,N}^{\uparrow}
:=
\min\left\{
\log_2N,
m\bigl(1-H_2(\eta)\bigr),
(1-2\eta)^2\sum_{j\le m}a_j^{\rm info}
\right\}.
\end{equation}
For singleton targets under the uniform prior,
\begin{equation}
\label{eq-bsc-optimal-extraction-vulnerability}
g_{R,t}^{\rm Bayes}
\le
\min\left\{
U_{R,t}^{\rm src},U_{R,t}^{\Caus},U_{R,t}^{\Lift},
U_{R,t}^{\Bdry},U_{R,t}^{\rm Bell},
\beta_{R,t}^{\rm Bayes,+}
+\sqrt{\frac{\log2}{2}J_{m,\eta,N}^{\uparrow}}
\right\}.
\end{equation}

\end{corollary}

\begin{proof}
Mutual information with \(\Theta\) is at most \(H(\Theta)\le\log_2N\).
The ordered-reveal chain rule, channel-capacity bound, and sequential SDPI
bound give the other two candidates in
\cref{eq-general-channel-optimal-extraction-information}.  For the binary
symmetric channel, substitute
\(c_j=1-H_2(\eta)\) and
\(\vartheta_j=(1-2\eta)^2\).  Insert the resulting information bound in
\cref{eq-optimal-extraction-information-candidate}, then minimize with the
remaining same-record structural candidates to obtain
\cref{eq-bsc-optimal-extraction-vulnerability}.
\end{proof}

\begin{figure}[tbp]
\centering
\resizebox{.98\linewidth}{!}{%
\begin{tikzpicture}[fw diagram,node distance=7mm and 8mm]
\node[fw source,text width=2.5cm] (pub) {public presentation\\and channel};
\node[fw provenance,text width=2.5cm,right=of pub] (src) {source-resolved\\information};
\node[fw geom,text width=2.5cm,right=of src] (flow) {Bayes--Fisher\\posterior flow};
\node[fw use,text width=2.5cm,right=of flow] (bell) {Bellman-optimal\\acquisition};
\node[fw terminal,text width=2.5cm,right=of bell] (gain) {target vulnerability\\and arrival};
\node[fw residual,text width=8.2cm,below=9mm of flow] (modal)
  {\(\Geom/\Abs\) source \(\longrightarrow\) \(\Caus\) action
   \(\longrightarrow\) \(\Lift/\Bdry\) transfer; \(\Res\) contributes zero};
\draw[fw arrow] (pub)--(src);
\draw[fw arrow] (src)--(flow);
\draw[fw arrow] (flow)--(bell);
\draw[fw arrow] (bell)--(gain);
\draw[fw arrow] (src)--(modal);
\draw[fw arrow] (modal)-|(gain);
\end{tikzpicture}%
}
\caption{Optimal structural information extraction.  Bayesian natural-gradient
flow processes each represented observation, Bellman recursion chooses among
affordable acquisitions, and the complete target gain remains bounded by the
existing modality-resolved certificate.}
\label{fig-optimal-structural-information-extraction}
\end{figure}

\subsection{Compact latent beliefs and optimal represented compression}
\label{subsec-compact-latent-beliefs}

The represented belief in
\cref{def-complete-structural-bayesian-record} may be stored through a finite
latent state.  The latent state is part of the pre-hit computation: its
encoder, update, decoder, and deployment rule retain the complete source
ancestry and cost of the public record from which they are constructed.

\begin{definition}[Complete compact latent-belief record]
\label{def-complete-compact-latent-belief-record}

Let \(R\in\mathscr R_{I,H}^{\rm Bayes}(B)\) be a complete structural Bayesian
record, and let \(\mathcal H_t\) be its complete public pre-hit history.  A
complete compact latent-belief extension of \(R\) contains:
\begin{enumerate}[label=\textup{(\roman*)}]
\item a finite represented code carrier \(\mathcal Z_t\) and a temporally
      admissible represented encoder kernel
      \(E_t(dz\mid\mathcal H_t)\), whose auxiliary randomness is independent
      of the target conditional on \(\mathcal H_t\);
\item a represented recurrent update
      \(F_t(dz'\mid z,Y_{t+1},A_t^{\rm qry},\mathcal H_t)\) that produces
      \(Z_{t+1}\) before every transition that uses it;
\item a represented decoder \(D_t:Z_t\mapsto\widehat\Pi_t^Z\), together with
      every likelihood score, normalizer, posterior statistic, acquisition
      score, and deployed selector computed from \(Z_t\);
\item the complete encoder, update, storage, precision, decoder, optimization,
      verification, and deployment derivation.
\end{enumerate}

Condition on the survival sector \(\{\tau>t\}\).  When this event has positive
probability, define the analytic latent posterior, retained target information,
and sufficiency defect by
\begin{align}
\label{eq-compact-latent-analytic-posterior}
\Pi_t^Z(\theta)&=\Pr\{\Theta=\theta\mid Z_t,\tau>t\},\\
\label{eq-compact-latent-retained-information}
I_{R,t}^{\rm lat}&=I(\Theta;Z_t\mid\tau>t),\\
\label{eq-compact-latent-sufficiency-defect}
\Delta_{R,t}^{\rm suf}
&=I(\Theta;\mathcal H_t\mid Z_t,\tau>t),\\
\label{eq-compact-latent-decoder-error}
e_{R,t}^{\rm dec}
&=\mathbb E\!\left[
 \lVert\widehat\Pi_t^Z-\Pi_t^Z\rVert_{\rm TV}
 \mid\tau>t\right].
\end{align}
All four quantities are zero when the survival event is null.  Mutual
information is measured in bits.  Put
\begin{equation}
\label{eq-compact-latent-code-capacity}
b_{R,t}^{\rm lat}=\left\lceil\log_2|\mathcal Z_t|\right\rceil.
\end{equation}
If the represented code has \(d_t\) coordinates at \(p_t\) bits of precision
and a codebook or architecture description of length \(b_t^{\rm code}\), then
the record uses the certified bound
\begin{equation}
\label{eq-compact-latent-coordinate-capacity}
b_{R,t}^{\rm lat}\le d_tp_t+b_t^{\rm code}.
\end{equation}

The complete latent cost is
\begin{align}
\label{eq-complete-compact-latent-cost}
B_R^{\rm lat}
={}&B_R^{\rm Bayes}\oplus B_E\oplus B_F\oplus B_{\mathcal Z}
\oplus B_{\rm code}\oplus B_{\rm prec}\oplus B_{\rm mem}\notag\\
&\oplus B_D\oplus B_{\rm like}\oplus B_{\rm norm}
\oplus B_{\rm score}\oplus B_{\rm optimize}\oplus B_{\rm verify}
\oplus B_{\rm deploy}.
\end{align}
Shared ancestors occur once in the derivation DAG; every repeated evaluation
or update retains its clocked cost.  The source index is the exact source index
of the complete backward slice in
\cref{eq-bayesian-exact-source-index}.  Encoding and deterministic updating
add no target information and create no new source tag.

\end{definition}

\begin{theorem}[Latent information, sufficiency, and decision loss]
\label{thm-compact-latent-information-sufficiency}

For every complete compact latent-belief record and every nonnull survival
sector,
\begin{equation}
\label{eq-compact-latent-information-chain}
I(\Theta;\mathcal H_t\mid\tau>t)
=I_{R,t}^{\rm lat}+\Delta_{R,t}^{\rm suf}.
\end{equation}
If \(J_{R,t}^{\rm Bayes,src}\) is the source-resolved ordered-reveal upper
bound for the same conditioned history, then
\begin{equation}
\label{eq-compact-latent-information-capacity-bound}
I_{R,t}^{\rm lat}
\le
\min\left\{
b_{R,t}^{\rm lat},
J_{R,t}^{\rm Bayes,src},
I(\Theta;\mathcal H_t\mid\tau>t)
\right\}.
\end{equation}
Moreover,
\begin{equation}
\label{eq-compact-latent-posterior-sufficiency-tv}
\mathbb E\!\left[
\left\|
\Pr\{\Theta\in\cdot\mid\mathcal H_t,\tau>t\}-\Pi_t^Z
\right\|_{\rm TV}
\,\middle|\,\tau>t\right]
\le
\sqrt{\frac{\log2}{2}\Delta_{R,t}^{\rm suf}}.
\end{equation}

Let \(u(\theta,a)\in[0,1]\) be a represented target utility, and let
\(\operatorname{DVal}_{u,t}(\mathcal H_t)\) and
\(\operatorname{DVal}_{u,t}(Z_t)\) denote the optimal survival-conditioned
decision values with the same represented action carrier.  Then
\begin{equation}
\label{eq-compact-latent-decision-deficiency}
0\le
\operatorname{DVal}_{u,t}(\mathcal H_t)
-\operatorname{DVal}_{u,t}(Z_t)
\le
\sqrt{\frac{\log2}{2}\Delta_{R,t}^{\rm suf}}.
\end{equation}
Using the represented decoder instead of the analytic latent posterior adds at
most \(e_{R,t}^{\rm dec}\) to the right side.  In particular,
\(\Delta_{R,t}^{\rm suf}=0\) gives analytic decision-value equivalence.  When
the conditional reconstruction kernel is represented and charged, this is
exactly the represented value equivalence of
\cref{cor-controlled-represented-prehit-sufficiency}.

\end{theorem}

\begin{proof}
Conditional on survival, the encoder construction gives the Markov chain
\(\Theta\to\mathcal H_t\to Z_t\).  The mutual-information chain rule therefore
proves \cref{eq-compact-latent-information-chain}.  Data processing gives the
history-information bound, while
\(I(\Theta;Z_t\mid\tau>t)\le H(Z_t\mid\tau>t)\le\log_2|\mathcal Z_t|\)
gives the code-capacity bound.  The ordered-reveal inequality is
\cref{eq-source-resolved-bayesian-information-ledger} on the same conditioned
record.

The conditional-information identity also gives
\[
\Delta_{R,t}^{\rm suf}
=\mathbb E\!\left[
D_{\rm KL,2}\!\left(
\Pr\{\Theta\in\cdot\mid\mathcal H_t,\tau>t\}
\,\middle\|\,\Pi_t^Z
\right)\middle|\tau>t\right],
\]
where \(D_{\rm KL,2}\) is relative entropy in bits.  Pinsker's inequality and
Jensen's inequality prove
\cref{eq-compact-latent-posterior-sufficiency-tv}.  The expectation of any
\([0,1]\)-valued action utility changes by at most total variation.  Apply this
first to the history-optimal action and then optimize over latent actions to
obtain \cref{eq-compact-latent-decision-deficiency}.  The triangle inequality
adds the decoder error.  At zero conditional information, the conditional law
of the history given \((\Theta,Z_t)\) equals its law given \(Z_t\), which is the
sufficiency identity in
\cref{eq-controlled-represented-prehit-sufficiency}.  The represented version
uses that corollary's charged reconstruction hypothesis.
\end{proof}

\begin{definition}[Saturated compact latent-belief profile]
\label{def-saturated-compact-latent-belief-profile}

Let \(\mathscr R_{I,H}^{\rm lat}(B,b)\) be the complete compact latent-belief
records with \(B_R^{\rm lat}\preceq B\) and
\(\max_{t<H}b_{R,t}^{\rm lat}\le b\).  Define
\begin{equation}
\label{eq-saturated-compact-latent-arrival-profile}
\operatorname{LatBayesArr}_{I,H}^{\rm sat}(B,b)
=\sup_{R\in\mathscr R_{I,H}^{\rm lat}(B,b)}\Pr_R\{\tau\le H\},
\qquad\sup\varnothing=0.
\end{equation}
The maximum useful belief update is the restriction of
\cref{eq-saturated-actionable-bayesian-gain-profile} to the same records:
\begin{equation}
\label{eq-saturated-compact-latent-gain-profile}
\operatorname{LatBayesGain}_{I,H}^{\rm sat}(B,b)
=\max_{t<H}\sup_{R\in\mathscr R_{I,H}^{\rm lat}(B,b)}
g_{R,t}^{\rm Bayes}.
\end{equation}
Usefulness here means deployment through the represented next-test kernel in
\cref{eq-actionable-structural-bayesian-gain}.  Posterior mass stored in the
latent state without such a deployment contributes no actionable gain.
The family-uniform profile uses one represented encoder, update, decoder, and
deployment scheme on the whole size slice and then applies the declared family
functional.

\end{definition}

\begin{theorem}[Optimal compact latent acquisition and saturated completeness]
\label{thm-optimal-compact-latent-acquisition}

The profile in
\cref{eq-saturated-compact-latent-arrival-profile} is monotone in \(B\) and
\(b\), and
\begin{equation}
\label{eq-compact-latent-profile-inclusion}
\operatorname{LatBayesArr}_{I,H}^{\rm sat}(B,b)
\le\operatorname{BayesArr}_{I,H}^{\rm sat}(B).
\end{equation}
The same inclusion gives
\begin{equation}
\label{eq-compact-latent-gain-profile-inclusion}
\operatorname{LatBayesGain}_{I,H}^{\rm sat}(B,b)
\le\operatorname{BayesGain}_{I,H}^{\rm sat}(B).
\end{equation}
Restricting
\cref{eq-optimal-structural-bayes-resource-recursion} to records in
\(\mathscr R_{I,H}^{\rm lat}(B,b)\) gives the exact resource-valued Bellman
recursion that jointly optimizes acquisition, encoding, recurrent updating,
decoding, and continuation under the remaining cost and code-capacity budgets.

Conversely, encode the complete represented sufficient history and internal
state of a structural Bayesian record.  Its code length is the storage length
already charged in that record, its update is the original clocked state
update, and its decoder may be the original represented belief field or the
constant represented belief supplied by universal Bayesianization.  Hence a
class-preserving compiler \(\Psi_{\rm lat}\) satisfies
\begin{equation}
\label{eq-compact-latent-saturated-completeness}
\operatorname{BayesArr}_{I,H}^{\rm sat}(B)
\le
\sup_{b:\,B_{\rm code}(b)\preceq\Psi_{\rm lat}(B,H,n)}
\operatorname{LatBayesArr}_{I,H}^{\rm sat}
\bigl(\Psi_{\rm lat}(B,H,n),b\bigr).
\end{equation}
Here \(B_{\rm code}(b)\) denotes the represented storage, precision, update,
and evaluation charge of a code with capacity \(b\), as included in
\cref{eq-complete-compact-latent-cost}.
Together with
\cref{cor-universal-represented-computation-bayesian-ceiling}, this gives
\begin{equation}
\label{eq-compact-latent-polynomial-saturated-ceiling}
\operatorname{LatBayesArr}_{\mathfrak M,n}^{\rm poly,sat}
=\operatorname{BayesArr}_{\mathfrak M,n}^{\rm poly,sat}
=\operatorname{ActArr}_{\mathfrak M,n}^{\rm poly,sat}
=\operatorname{Succ}_{\mathfrak M,n}^{\rm poly,sat},
\end{equation}
where the latent profile takes the union over code capacities admitted by the
same polynomial saturated cost.
Because the compiler preserves the represented next-test kernel and its
pre-hit law, it also gives
\begin{equation}
\label{eq-compact-latent-polynomial-gain-ceiling}
\operatorname{LatBayesGain}_{\mathfrak M,n}^{\rm poly,sat}
=\operatorname{BayesGain}_{\mathfrak M,n}^{\rm poly,sat}.
\end{equation}

\end{theorem}

\begin{proof}
Class inclusion and monotonicity prove
\cref{eq-compact-latent-profile-inclusion}.  First-step decomposition and
concatenation preserve both remaining resources, so the proof of
\cref{thm-optimal-structural-bayesian-acquisition} gives the restricted exact
recursion.  For the converse, the complete clocked state is a finite represented
code.  Storing that code and applying the original transition and selector
reproduces the path law exactly; encoding adequacy and computation closure give
the class-preserving majorant \(\Psi_{\rm lat}\).  This proves
\cref{eq-compact-latent-saturated-completeness}.  Apply universal
Bayesianization and take the polynomial saturated transfer class to obtain the
arrival equality.  The compiler also leaves the next-test kernel, target
utility, survival event, and conditional law unchanged, so it preserves
\(g_{R,t}^{\rm Bayes}\) exactly.  This proves
\cref{eq-compact-latent-polynomial-gain-ceiling}.
\end{proof}

\begin{theorem}[Unified latent structural certificate]
\label{thm-unified-compact-latent-structural-certificate}

For one complete compact latent-belief record, put
\begin{equation}
\label{eq-compact-latent-information-target-candidate}
U_{R,t}^{\rm lat,info}
=\min\left\{1,
\beta_{R,t}^{\rm Bayes,+}
+\sqrt{\frac{\log2}{2}
\min\{b_{R,t}^{\rm lat},J_{R,t}^{\rm Bayes,src}\}}
\right\}.
\end{equation}
Let \(U_{R,t}^{\rm lat,stat}\) be the smallest same-record target-gain
consequence of the source-resolved Rademacher, effective-dimension,
PAC--Bayes, sequential, or martingale candidates for the complete
encoder--decoder--selector class.  Define
\begin{equation}
\label{eq-unified-compact-latent-certificate}
U_{R,t}^{\rm lat}
=\min\left\{
U_{R,t}^{\rm opt},
U_{R,t}^{\rm lat,info},
U_{R,t}^{\rm lat,stat},
U_{R,t}^{\Caus},U_{R,t}^{\Lift},U_{R,t}^{\Bdry},U_{R,t}^{\rm Bell}
\right\},
\end{equation}
where an unavailable candidate has neutral value one.  Then
\begin{equation}
\label{eq-unified-compact-latent-gain-bound}
g_{R,t}^{\rm Bayes}\le U_{R,t}^{\rm lat}.
\end{equation}

The assignment of the complete record is exhaustive.  The represented latent
metric, kernel, and belief simplex belong to \(\Geom\).  Encoder fibers belong
to exact or compensated \(\Abs\) precisely when their corresponding gates
hold.  A deployed target-directed belief or selector current is authenticated
\(\Caus\).  Recurrent state, memory, decoder comparison, and filter transfer
belong to \(\Lift\).  Support masks and the target decoder belong to \(\Bdry\).
The gradient-orthogonal component belongs to \(\Res\) and has zero
target-qualified gain.  Add, Mult, Ord, Sym, and Filt remain the constructor
provenance of these same records.

For a learned latent policy, let \(\Gamma_{R,t}^{\rm lat}\) be its applicable
same-record generalization radius and \(\varepsilon_{R,t}^{\rm opt}\) its
represented optimization error.  Relative to the complete-history optimum,
the finite-horizon loss is bounded by
\begin{equation}
\label{eq-learned-compact-latent-policy-loss}
\begin{aligned}
0\le
\operatorname{BayesArr}_{I,H}^{\rm sat}(B)-A_R^H
\le{}&
\sum_{t<H}\mathbb E\!\left[\mathbf1_{\{\tau>t\}}
\left(
\sqrt{\frac{\log2}{2}\Delta_{R,t}^{\rm suf}}
+e_{R,t}^{\rm dec}
+2\Gamma_{R,t}^{\rm lat}
+\varepsilon_{R,t}^{\rm opt}
\right)\right]
+\Delta_{\rm fil}^{H}.
\end{aligned}
\end{equation}

\end{theorem}

\begin{proof}
The information candidate is conditional data processing followed by the
target-contact Pinsker comparison in
\cref{thm-noise-contracted-source-resolved-bayesian-gain}.  The source,
dynamical, and Bellman candidates are the same-record entries of
\cref{thm-optimal-structural-extraction-saturated-ceiling}; the statistical
candidate is
\cref{thm-source-resolved-saturated-learning-compilation} applied to the
complete encoder--decoder--selector output class.  Taking their minimum proves
\cref{eq-unified-compact-latent-gain-bound}.

The channel assignment follows from
\cref{thm-complete-generic-source-theory,thm-exhaustive-five-family-geom-source-classification}
and from source inheritance under Caus, Lift, and Bdry.  Finally combine the
decision deficiency in
\cref{eq-compact-latent-decision-deficiency}, the decoder triangle inequality,
the learned active-rule loss in
\cref{thm-learned-active-rule-finite-horizon-loss}, and the represented filter
transfer in \cref{thm-controlled-pomdp-filter-transfer}.  Summing over surviving
times proves \cref{eq-learned-compact-latent-policy-loss}.
\end{proof}

\begin{corollary}[Noise, samples, and code capacity for latent beliefs]
\label{cor-noise-sample-code-capacity-latent-beliefs}

Let \(N=|\mathsf\Theta|\), and suppose that the complete record contains at
most \(m\) binary symmetric target-dependent reveals of common noise rate
\(0\le\eta\le1/2\).  Let \(D_{R,t}^{\rm ext}\) be the conditional information
in all other target-dependent advice, oracle, reward, verifier, or feedback
coordinates available before deployment.  If each clean binary input contains
at most one available bit, then
\begin{equation}
\label{eq-noise-sample-code-latent-information}
I_{R,t}^{\rm lat}
\le
\min\left\{
b_{R,t}^{\rm lat},\log_2N,
m\bigl(1-H_2(\eta)\bigr)+D_{R,t}^{\rm ext},
m(1-2\eta)^2+D_{R,t}^{\rm ext},
J_{R,t}^{\rm Bayes,src}
\right\}.
\end{equation}
For uniform targets, exact identification with probability at least \(\alpha\)
requires
\begin{equation}
\label{eq-latent-code-fano-necessary}
b_{R,t}^{\rm lat}\ge\alpha\log_2N-1.
\end{equation}
Consequently a \(d_t\)-coordinate, \(p_t\)-bit representation requires
\begin{equation}
\label{eq-latent-coordinate-fano-necessary}
d_tp_t+b_t^{\rm code}\ge\alpha\log_2N-1.
\end{equation}

\end{corollary}

\begin{proof}
Apply
\cref{cor-noise-sample-target-size-optimal-extraction} and data processing
through the encoder, then intersect its information bound with the code-capacity
bound in \cref{eq-compact-latent-information-capacity-bound}.  The external
coordinates enter their ordered reveal slots.  Fano's inequality gives
\(\alpha\le(I_{R,t}^{\rm lat}+1)/\log_2N\); combine this with the code and
coordinate capacities.
\end{proof}

\begin{figure}[tbp]
\centering
\resizebox{.98\linewidth}{!}{%
\begin{tikzpicture}[fw diagram,node distance=7mm and 8mm]
\node[fw source,text width=2.6cm] (hist) {public pre-hit\\history};
\node[fw use,text width=2.5cm,right=of hist] (acq) {optimal\\acquisition};
\node[fw geom,text width=2.5cm,right=of acq] (enc) {charged encoder\\and latent state};
\node[fw provenance,text width=2.5cm,right=of enc] (flow) {Bayes--Fisher\\belief flow};
\node[fw use,text width=2.5cm,right=of flow] (dec) {decoder and\\selector};
\node[fw terminal,text width=2.4cm,right=of dec] (hit) {target\\contact};
\node[fw residual,text width=7.2cm,below=10mm of flow] (cert)
 {code capacity, sufficiency, source, learning, Caus, Lift, Bdry, and Bellman certificates};
\draw[fw arrow] (hist)--(acq);
\draw[fw arrow] (acq)--(enc);
\draw[fw arrow] (enc)--(flow);
\draw[fw arrow] (flow)--(dec);
\draw[fw arrow] (dec)--(hit);
\draw[fw dashed arrow] (hist)--(cert);
\draw[fw dashed arrow] (enc)--(cert);
\draw[fw dashed arrow] (flow)--(cert);
\draw[fw dashed arrow] (dec)--(cert);
\end{tikzpicture}%
}
\caption{A compact latent belief is a completely charged component of optimal
active acquisition.  Its code capacity and sufficiency defect quantify
compression; its deployed gain remains bounded by the existing structural,
statistical, and dynamical certificates.}
\label{fig-compact-latent-belief-pipeline}
\end{figure}

\section{Bellman Error, Policy Loss, and the Universal Certificate}
\label{sec-universal-dynamical-certificate}

\begin{theorem}[Represented Bellman residual and policy loss]
\label{thm-controlled-bellman-residual-policy-loss}

Let \(0<\beta<1\), let rewards be bounded, and let \(v\) be a represented
value approximation.  For a fixed represented policy \(\pi\),
\begin{equation}
\label{eq-fixed-policy-bellman-residual-bound}
\|V^\pi-v\|_\infty
\le
\frac{\|\mathcal T^\pi v-v\|_\infty}{1-\beta}.
\end{equation}
For the optimal Bellman operator,
\begin{equation}
\label{eq-optimal-bellman-residual-bound}
\|V^*-v\|_\infty
\le
\frac{\|\mathcal T v-v\|_\infty}{1-\beta}.
\end{equation}
If \(\pi_v\) is a represented policy greedy with respect to \(v\), then
\begin{equation}
\label{eq-greedy-policy-loss-bound}
\|V^*-V^{\pi_v}\|_\infty
\le
\frac{2\beta}{(1-\beta)^2}
\|\mathcal T v-v\|_\infty.
\end{equation}
Every residual evaluation, action comparison, greedy selector, and retained
coefficient is included in the saturated cost of the certificate.  A sampled
residual enters these deterministic inequalities only through a simultaneous
Part~IV generalization certificate for the represented residual class.

\end{theorem}

\begin{proof}

For any \(\beta\)-contraction \(F\) with fixed point \(v_F\),
\[
\|v_F-v\|_\infty
\le\|Fv_F-Fv\|_\infty+\|Fv-v\|_\infty
\le\beta\|v_F-v\|_\infty+\|Fv-v\|_\infty.
\]
Rearrangement proves
\cref{eq-fixed-policy-bellman-residual-bound,eq-optimal-bellman-residual-bound}.
Greediness gives \(\mathcal T^{\pi_v}v=\mathcal T v\).  Insert \(v\) between
the two fixed points, use the contraction twice, and obtain the standard
estimate
\[
\|V^*-V^{\pi_v}\|_\infty
\le\frac{2\beta}{1-\beta}\|V^*-v\|_\infty.
\]
Combining it with \cref{eq-optimal-bellman-residual-bound} proves
\cref{eq-greedy-policy-loss-bound}.  The cost and sampled-residual statements
are applications of
\cref{def-budgeted-derivability-closure,def-budget-saturated-statistical-readout-class}.

\end{proof}

\begin{definition}[Coordinatewise dynamical certificate]
\label{def-coordinatewise-dynamical-certificate}

For one controlled record of budget \(B\), horizon \(H\), target/failure
pair, discount parameters, and selected policy, the coordinatewise dynamical
certificate is
\[
\begin{aligned}
\mathbf U_{\rm dyn}
={}&(U_{\rm hit},U_{\rm disc},U_{\rm time},U_{\rm ra},
  U_{\Caus},U_{\Lift}^{H},U_{\Lift}^{\lambda},U_{\rm noise},\\
&\quad U_{\Caus}^{\rm noise},U_{\Lift,{\rm noise}}^{H},
  U_{\Lift,{\rm noise}}^{\lambda},U_{\rm Cap}^{\rm noise},
  \mathbf U_{\rm ch},U_{\rm POMDP}^{H},U_{\rm info},U_{\rm cur}^{H},
  \mathbf H_{\rm ctl},U_{\rm mix},U_{\rm Bell},U_{\rm learn}).
\end{aligned}
\]
It packages existing estimates and is not an additional structural channel.
Its entries are defined as follows.

\begin{enumerate}[label=\textup{(\roman*)}]
\item \(U_{\rm hit}\) is the minimum of one, the prospective-selector bound
      \(\operatorname{Enum}_\nu+H
      \operatorname{Sel}^{\rm sat}(\Psi_{\rm ctl}(B,H,n))\), and every
      same-record native capacity, occupation, Caus-contact, or finite-horizon
      Lift bound, together with every same-record sparse-reward observation
      bound, already converted to first-arrival probability units.
\item \(U_{\rm disc}\) is the minimum of one and every represented upper
      comparison for the target-projected killed resolvent
      \(A_T(\lambda;L,\mu_0)=S_T(\lambda;L,\mu_0)\) admitted by
      \cref{def-target-resolvent,def-master-compensated-quotient-geometry-certificate},
      including \(U_{\Lift}^{\lambda}\) and
      \(U_{\Lift,{\rm noise}}^{\lambda}\) when their same-normalization
      comparisons are represented.
\item \(U_{\rm time}=\mathbb E_{\mu_0}D/a\) when the same selected
      generator has a represented drift certificate \(-LD\ge a>0\), and
      \(U_{\rm time}=+\infty\) otherwise.  A quotient defect replaces \(a\)
      by its certified positive residual margin.
\item \(U_{\rm ra}=\langle\mu_0,h\rangle\) when a represented reach--avoid
      committor is certified, and \(U_{\rm ra}=1\) otherwise.
\item \(U_{\Caus}\) is the minimum of one and every probability-normalized
      bound in \cref{eq-controlled-saturated-caus-contact-bound} supplied by
      a represented aiming-to-contact converter.  Its value is one when only
      action or aiming energy is available.
\item \(U_{\Lift}^{H}\) is the minimum of one and the base first-arrival
      bound plus the charged clocked defect and interface errors from
      \cref{thm-controlled-clocked-lift-transfer}.  Its default is one.
      The discounted coordinate \(U_{\Lift}^{\lambda}\) is the analogous
      minimum using exact intertwining or
      \cref{eq-controlled-fast-fiber-schur-error}; its default is one.
\item \(U_{\rm noise}\) is the maximum, over the represented binary
      supervised coordinates in the record, of the clean risk bounds
      \(U_{{\rm noise},q}^{\rm clean}
      (n,b(n),m,\eta_{q,+},\delta_q)\) from
      \cref{eq-saturated-dynamical-clean-noise-bound}.  It retains the
      complete computational ceiling \(b(n)\), sample size \(m\), and every
      noise margin \(1-2\eta_{q,+}\); its default is one.
\item \(U_{\Caus}^{\rm noise}\) is the minimum of one and the applicable
      clean-target Caus contact bounds in
      \cref{eq-controlled-noise-corrected-caus-contact}.
      The finite-horizon coordinate \(U_{\Lift,{\rm noise}}^H\) is
      \(\widehat p_{T,H}^{\Omega}\) plus the right side of
      \cref{eq-controlled-noisy-gate-clocked-arrival}, clipped at one.
      The discounted coordinate \(U_{\Lift,{\rm noise}}^\lambda\) is the
      represented learned-gate discounted arrival plus the right side of
      \cref{eq-controlled-noisy-gate-discounted-arrival}, again clipped at
      one after a declared probability normalization.  Their defaults are
      one.
\item \(U_{\rm Cap}^{\rm noise}\) is the right side of
      \cref{eq-controlled-noisy-lift-capacity-bound} in capacity units.  Its
      default is \(+\infty\).  It is not minimized with a probability until
      a represented capacity-to-arrival comparison is supplied.
\item The typed channel vector is
      \[
      \mathbf U_{\rm ch}
      =\bigl(
      U_{r,\mathsf N_r},U_{A,\mathsf N_A},
      U_{O,\mathsf N_O},U_{{\rm fil},\mathsf N_{\rm fil}},
      U_{{\rm comp},\mathsf N_{\rm comp}},U_{\rm path}
      \bigr),
      \]
      whose raw entries are the applicable profiles in
      \cref{eq-saturated-dynamical-channel-error}.  Their declared norms are
      part of the coordinate.  \(U_{\rm path}\) is a represented reference
      event probability plus the applicable probability-normalized shift in
      \cref{eq-controlled-path-tv-ledger,eq-controlled-path-kl-ledger},
      clipped at one; raw path distance alone is not an event bound.
      Raw entries default to \(+\infty\), and the path probability defaults
      to one.
\item \(U_{\rm POMDP}^{H}\) is the minimum of one and every same-record
      clocked POMDP first-arrival bound.  Its filter candidate is
      \[
      \min\{1,U_b^H+\Delta_{\rm fil}^H\},
      \]
      where \(U_b^H\) is a represented upper bound for the same exact-belief
      target event on the same reference record and
      \(\Delta_{\rm fil}^H\) is the comparison in
      \cref{eq-controlled-pomdp-filter-hit-transfer}.  The exact analytic
      probability is eligible as \(U_b^H\) only when its evaluation or proof
      and target/failure interface are represented and charged; otherwise
      \(U_b^H=1\).  Its default is therefore one.
      \(U_{\rm info}\) is the Fano converse from
      \cref{thm-controlled-multichannel-information-bound} when its success
      event identifies the target, and is one otherwise.  A rate--distortion
      converse enters a probability coordinate only after a represented
      inversion to a distortion floor and a distortion-to-event converter.
      \(\mathbf H_{\rm ctl}\) is the typed entropy vector in
      \cref{eq-controlled-typed-entropy-vector}; it is not a scalar error
      coordinate.
\item \(U_{\rm cur}^{H}\) is the minimum of one and every same-record
      sparse-reward observation bound in
      \cref{eq-observation-cell-count-cumulative-hit}.  In particular, a
      target-permutation-neutral record supplies
      \[
      U_{\rm cur}^{H}
      \le
      \min\!\left\{1,\sum_{t=0}^{H-1}
      \left(
      \frac{M_\varepsilon(\phi)}{N}
      +\sqrt{\frac{\log2}{2}J_{R,t}^{\rm fib,\uparrow}}
      \right)\right\}.
      \]
      Observation, quantization, memory, predictor training, intrinsic
      reward, policy deployment, executed transitions, verifier use, and
      every exact or compensated cell-dynamics comparison carry the costs
      in
      \cref{eq-complete-sparse-reward-record-cost,eq-complete-compensated-observation-cost}.
      Its default is one.
\item \(U_{\rm mix}(t)=e^{-2\gamma_\Phi t}\) for variance and
      \(e^{-2\alpha_\Phi t}\) for entropy when the corresponding same-record
      constant is certified; the default value is one.
\item \(U_{\rm Bell}\) is the appropriate right-hand side of
      \cref{eq-fixed-policy-bellman-residual-bound,eq-optimal-bellman-residual-bound,eq-greedy-policy-loss-bound};
      its value is \(+\infty\) without a represented residual certificate.
\item \(U_{\rm learn}\) is the minimum of the valid simultaneous structural,
      Rademacher, PAC--Bayes, effective-dimension, and geometry-selection
      bounds of \cref{part:learning}, after conversion to one declared risk or
      alignment unit and inclusion of every selection term.
\end{enumerate}

The entries are never added across incompatible units.  In particular,
\(U_{\Caus}\) enters \(U_{\rm hit}\) only after the contact converter,
\(U_{\Lift}^{H}\) competes only in finite-horizon probability units, and
\(U_{\Lift}^{\lambda}\) competes only in discounted-arrival units.  A minimum is taken
only between upper bounds for the same coordinate on the same represented
record.  The same rule applies to the noise-corrected Caus and Lift
coordinates; \(U_{\rm Cap}^{\rm noise}\) remains separate unless its
capacity units have been converted explicitly.  The entries of
\(\mathbf U_{\rm ch}\) and \(\mathbf H_{\rm ctl}\) likewise remain typed
until a represented reward, Bellman, path, filter, aiming, or Lift theorem
converts them to an existing coordinate.
The curiosity coordinate competes in \(U_{\rm hit}\) only after
\cref{thm-observation-regularity-sparse-reward-exploration} has converted
observation resolution and conditional fiber information to the same
finite-horizon target-arrival probability.

\end{definition}

\begin{theorem}[Universal coordinatewise bound for represented dynamical systems]
\label{thm-universal-dynamical-certificate}

Every budget-admissible deterministic, randomized, controlled, or
history-dependent finite-horizon dynamical system on a presented task has a
coordinatewise certificate
\(\mathbf U_{\rm dyn}\) at the class-preserving ceiling
\(\Psi_{\rm ctl}(B,H,n)\).  Its coordinates satisfy
\begin{align}
\Pr\{\tau_{T^+}\le H,\ \tau_{T^+}<\tau_{T^-}\}
&\le U_{\rm hit},
\label{eq-universal-dynamics-hit-bound}\\
S_T(\lambda;L,\mu_0)&\le U_{\rm disc},
\label{eq-universal-dynamics-discount-bound}\\
\mathbb E_{\mu_0}\tau_{T^+}&\le U_{\rm time},
\label{eq-universal-dynamics-time-bound}\\
\Pr_{\mu_0}\{\tau_{T^+}<\tau_{T^-}\}&\le U_{\rm ra}.
\label{eq-universal-dynamics-reach-avoid-bound}
\end{align}
The certified mixing, Bellman, policy-loss, and learned-control conclusions
are the corresponding entries of
\cref{thm-controlled-certified-mixing,thm-controlled-bellman-residual-policy-loss,thm-saturated-geometry-learning-selection-accountability}.
The same first-arrival probability is bounded by \(U_{\Caus}\) and
\(U_{\Lift}^{H}\) whenever their converters are present, and by
\(U_{\rm cur}^{H}\) whenever a complete sparse-reward observation record is
present.  The discounted
arrival is bounded by \(U_{\Lift}^{\lambda}\).  Accordingly,
\(U_{\rm hit}\le\min\{U_{\Caus},U_{\Lift}^{H},U_{\rm cur}^{H}\}\) and
\(U_{\rm disc}\le U_{\Lift}^{\lambda}\) for those same-record candidates.
For binary target or failure interfaces learned from corrupted labels, the same
statements hold with
\(U_{\Caus}^{\rm noise},U_{\Lift,{\rm noise}}^H\), and
\(U_{\Lift,{\rm noise}}^\lambda\); the clean supervised risk is bounded by
\(U_{\rm noise}\), and the learned-boundary lifted capacity is bounded by
\(U_{\rm Cap}^{\rm noise}\).
In particular,
\(U_{\rm hit}\le\min\{U_{\Caus}^{\rm noise},
U_{\Lift,{\rm noise}}^H\}\) and
\(U_{\rm disc}\le U_{\Lift,{\rm noise}}^\lambda\) when the corresponding
learned-interface records apply.
For a complete action-, reward-, observation-, or POMDP-channel record, the
selected generator is the actual executed clocked generator of
\cref{def-controlled-action-execution-channel,thm-controlled-pomdp-clocked-embedding}.
The same first-arrival probability is additionally bounded by
\(U_{\rm path}\) and \(U_{\rm POMDP}^{H}\) when their comparisons are
present.  Thus \(U_{\rm hit}\) may include those two same-unit candidates.
The target-identification coordinate is bounded by \(U_{\rm info}\) only
under the identification gate in
\cref{def-controlled-affordable-multichannel-success}.  Reward noise changes
target arrival only through the policy/history kernel or an explicit path
comparison; raw reward error is not a direct arrival candidate.

Moreover, put \(B^\sharp=\Psi_{\rm ctl}(B,H,n)\).  The unconditional
five-family accounting estimate is
\begin{equation}
\label{eq-universal-dynamics-selector-source-charge}
\operatorname{Sel}^{\rm sat}_{I,n}(B^\sharp)
\le
eH_{B^\sharp}(n)\sum_{c\in\mathfrak J_{\Abs}^{5}}
\mathcal C_{I,\mathfrak F_5}^{(c)}
\bigl(D^{\rm sel}_{B^\sharp}(n)\bigr).
\end{equation}
Here \(H_{B^\sharp}\), \(D^{\rm sel}_{B^\sharp}\), and the
contextual-charge envelopes are exactly those of
\cref{thm-prospective-selector-structural-charge}.  This estimate charges the
complete first-hit accounting record.  When every nonzero frontier term also
has the represented pre-hit same-record comparison required by
\cref{rem-dynamic-transfer-from-contextual-charge,cor-contextual-structural-charge-transfer},
the source-envelope version of that theorem assigns the corresponding
nonneutral improvement to its qualified prospective \(\Abs\) or
\(\Geom/\Caus\) record.  Qualified \(\Lift\) and \(\Bdry\) uses retain the
inherited source and their complete downstream cost; \(J_{\rm res}\) supplies
no prospective structural source.

The theorem is total: when a geometric or statistical certificate is absent,
the conventions of \cref{def-certified-dynamical-geometry,def-coordinatewise-dynamical-certificate}
return the neutral bound for that coordinate.  A numerical improvement is
claimed only from an evaluated, temporally admissible, target-aligned record
inside the declared budget.

\end{theorem}

\begin{proof}

Use \cref{thm-controlled-policy-channel-decomposition} to place the complete
execution on its represented clocked carrier.  The exact temporal hazard
identity and its saturated supremum in
\cref{thm-prospective-selector-accountability} prove
\cref{eq-universal-dynamics-hit-bound}.  Apply the pointwise accounting
inequality of \cref{thm-prospective-selector-structural-charge} at
\(B^\sharp=\Psi_{\rm ctl}(B,H,n)\).  Its definitions of
\(H_{B^\sharp}\) and \(D^{\rm sel}_{B^\sharp}\) give
\cref{eq-universal-dynamics-selector-source-charge}.  Under the stated
same-record comparison, the final paragraph of that theorem and
\cref{cor-contextual-structural-charge-transfer} give the prospective-source
form.

The Caus coordinate follows from
\cref{cor-controlled-uniform-caus-flow-bound} after the represented
aiming-to-contact conversion.  Exact or approximate Lift coordinates follow
from
\cref{thm-controlled-valid-lift-block-resolvent,thm-controlled-clocked-lift-transfer}.
All three are candidates for an existing arrival coordinate, so slotwise
selection in \cref{thm-controlled-functional-profile-selection} permits only
the stated same-unit minima.

The sparse-reward observation coordinate is
\cref{thm-observation-regularity-sparse-reward-exploration}.  Its
finite-horizon contact estimate is already in probability units, and its
complete ledger includes every observation, novelty-learning, memory,
policy, execution, verification, and compensated-geometry operation.
Consequently it is an additional same-record candidate for
\(U_{\rm hit}\), without treating a nonintertwining observation partition as
an exact quotient.

The binary clean-risk coordinate is
\cref{thm-uniform-dynamical-label-noise-bound}.  Its conversion to clean Caus
aiming is \cref{thm-controlled-noise-corrected-caus-aiming}; its finite-time,
discounted, and capacity transfer through a learned target gate is
\cref{thm-controlled-noisy-gate-lift-transfer}.  These results give the
noise-corrected coordinates without changing the native deterministic form,
resolvent, or capacity inequalities.

The general channel coordinates are
\cref{def-saturated-dynamical-channel-error-profiles,thm-controlled-execution-channel-comparison,thm-controlled-multichannel-path-comparison}.
Their POMDP transfer is
\cref{thm-controlled-pomdp-clocked-embedding,thm-controlled-pomdp-filter-transfer},
and their information converse is
\cref{thm-controlled-multichannel-information-bound}.  The no-substitution
rule in \cref{def-controlled-typed-entropy-outputs} prevents Shannon, MLSI,
Caus, PAC--Bayes, description, and policy entropies from entering unlike
coordinates.

The equality of discounted hitting and the target-projected killed resolvent
is \cref{def-target-resolvent}; taking the minimum over valid represented
upper comparisons proves \cref{eq-universal-dynamics-discount-bound}.
The drift and reach--avoid conclusions are
\cref{thm-controlled-reach-avoid-dirichlet-drift}; the default values
\(+\infty\) and one make the inequalities valid when those records are absent.
The remaining coordinate bounds are the cited mixing, Bellman, and learning
theorems applied to the same complete record.

Finally, source inheritance and post-saturation residual orthogonality are
\cref{prop-derived-channels-create-no-structure,thm-exhaustive-residual-normal-form}.
Temporal admissibility is
\cref{def-pre-hit-temporally-admissible-source,thm-controlled-no-free-value}.
Thus every positive structural contribution has one existing source
classification and every unavailable certificate reduces to its stated
neutral value.

\end{proof}

\begin{theorem}[Saturated same-record prospective certificate envelope]
\label{thm-saturated-same-record-prospective-certificate-envelope}

Fix a presented task, a budget \(B\), a refined horizon \(H\), and a time
\(0\le t<H\).  Let \(\mathscr R_{I,t}(B)\) be the collection of complete
temporally admissible pre-hit records that induce a next-test selector at
time \(t+1\), including the prefix law, selected kernel, target interface,
neutral baseline, and every comparison used below.  For
\(R\in\mathscr R_{I,t}(B)\), write
\[
a_R
=
\bigl(\operatorname{Adv}_t(K_R;\nu_R)\bigr)_+.
\]
Evaluate the following candidates under the same prefix law and in the same
survival-weighted next-contact units.  An unavailable candidate has the
neutral value one.

\begin{enumerate}[label=\textup{(\roman*)}]
\item \(U_R^{\rm src}\) is the minimum of one and every represented
      inherited or direct same-normalization comparison in
      \cref{eq-inherited-prospective-value-comparison,eq-direct-prospective-value-comparison}.

\item \(U_R^{\rm cap}\) is the minimum of one and every capacity or
      occupation bound for the next-contact event supplied by the selected
      kernel and prefix initialization.  In particular, a represented
      reversible comparison record from
      \cref{thm-controlled-density-capacity-hit-transfer}, used with
      one remaining transition, supplies
      \begin{equation}
      \label{eq-saturated-certificate-capacity-candidate}
      U_R^{\rm cap}
      \le
      \sigma_{R,t}\min\!\left\{1,
      \|f_{R,t}\|_{L^p(\widetilde\mu_R)}
      \left[
      \min\!\left\{1,
      r_{R,+}\mu_R(T_I)
      +\kappa_{R,+}
       \operatorname{Cap}_{\Phi_{P_R}}(T_I,T_I^c)
      \right\}
      \right]^{1/q}
      \right\},
      \end{equation}
      where
      \[
      \sigma_{R,t}=\Pr_R\{\tau>t\},\qquad
      \rho_{R,t}=\mathcal L_R(X_t\mid\tau>t),\qquad
      f_{R,t}=\frac{d\rho_{R,t}}{d\widetilde\mu_R},
      \]
      \(1<p\le\infty\), and \(q=p/(p-1)\).  More precisely, the
      conditional finite-horizon bound is multiplied by
      \(\sigma_{R,t}\) because \(a_R\) uses survival-weighted units.

\item If the target is indexed by a uniform latent variable on \(N\) values,
      \(U_R^{\rm info}\) is the minimum of one and every information or
      chi-square contact bound in
      \cref{eq-pre-hit-information-overlap-contact,eq-pre-hit-information-chi-square-contact}.
      The KL--SDPI candidate is explicitly
      \begin{equation}
      \label{eq-saturated-certificate-information-candidate}
      U_R^{\rm info}
      \le
      \min\!\left\{1,
      \frac{\Lambda_{R,t}}{N}
      +\sqrt{\frac{\log2}{2}
      \left(I(\Theta;\mathsf Z_0)+
      \sum_{j=1}^{L_{R,t}}d_{R,t,j}\right)}
      \right\}.
      \end{equation}

\item \(U_R^{\rm BW}\) is the minimum of one and every target-decision
      bound transferred from a represented pre-hit reference experiment by
      \cref{thm-controlled-prehit-blackwell-value-monotonicity}.  The
      reference experiment, causal garbling, action map, and complete
      simulation cost are part of \(R\).

\item \(U_R^{\rm dyn}\) is the minimum of one and every same-record Caus,
      finite-horizon Lift, path, POMDP, sparse-reward observation, or
      learned-interface bound already converted to the same next-contact
      units in
      \cref{def-coordinatewise-dynamical-certificate,thm-universal-dynamical-certificate}.
      For a complete observation record, its next-contact candidate is
      \begin{equation}
      \label{eq-saturated-certificate-curiosity-candidate}
      U_R^{\rm dyn}
      \le
      \min\!\left\{1,
      \beta_{R,t}^{\rm fib}
      +\sqrt{\frac{\log2}{2}J_{R,t}^{\rm fib,\uparrow}}
      \right\},
      \end{equation}
      with the complete ledger of
      \cref{def-complete-sparse-reward-observation-record}.

\item For a complete master active-sampling record,
      \(U_R^{\rm acq}\) is the minimum of one and the active next-contact
      candidate
      \begin{equation}
      \label{eq-saturated-certificate-active-acquisition-candidate}
      U_R^{\rm acq}
      \le
      \min\!\left\{1,
      \beta_{R,t}^{\rm act}
      +\sqrt{\frac{\log2}{2}J_{R,t}^{\rm act,\uparrow}}
      \right\},
      \end{equation}
      together with every same-record master, capacity, and Blackwell
      comparison already converted to this survival-weighted next-contact
      unit.  Its complete ledger is
      \cref{def-complete-master-active-sampling-record}; optimization over
      that complete record class is exactly the saturated active-success
      coordinate of
      \cref{thm-active-optimality-saturated-construction}.
\end{enumerate}

Put
\begin{equation}
\label{eq-saturated-recordwise-prospective-certificate}
U_R^{\rm cert}
=
\min\{U_R^{\rm src},U_R^{\rm cap},U_R^{\rm info},
       U_R^{\rm BW},U_R^{\rm dyn},U_R^{\rm acq}\}.
\end{equation}
Let \(B^{\rm cert}=\Psi_{\rm cert}(B,H,n,\mathbf m)\) be any declared
class-preserving common majorant of the existing cost maps for the selected
kernel, clock refinement, source comparison, capacity record, ordered
reveals, Blackwell simulation, sparse-reward observation and compensated
cell dynamics, master active sampling, Lift, and output conversion.  Here
\(\mathbf m\) denotes the sample or reveal counts already present in the
task presentation; it is omitted when no such counts occur.  Then
\begin{align}
\label{eq-saturated-selector-certificate-envelope}
\operatorname{Sel}^{\rm sat}_{I,n}(B)
&\le
\max_{0\le t<H}
\sup_{R\in\mathscr R_{I,t}(B^{\rm cert})}
U_R^{\rm cert},\\
\label{eq-saturated-success-certificate-envelope}
\Pr\{\tau\le H\}
&\le
\operatorname{Enum}_{\nu}(\mathcal A,H)
+\sum_{t=0}^{H-1}
\sup_{R\in\mathscr R_{I,t}(B^{\rm cert})}
U_R^{\rm cert}.
\end{align}
The second right-hand side may be clipped at one.  The same inequalities hold
after applying a monotone family functional.  The minimum in
\cref{eq-saturated-recordwise-prospective-certificate} is taken recordwise,
before either supremum.  Thus coordinates belonging to different selected
kernels, prefix laws, or presentation events are never combined.

\end{theorem}

\begin{proof}

For each \(R\), the neutral baseline has nonnegative next-contact mass, so
\cref{eq-pre-hit-information-selector-bound} gives \(a_R\le p_R\), where
\(p_R\) is the survival-weighted next-contact probability.  The two source
comparisons bound \(a_R\) directly by
\cref{thm-no-uncharged-prospective-value-amplification}.  Every available
capacity, information, Blackwell, or dynamical candidate bounds either
\(p_R\) or \(a_R\) under the same prefix law and in the same units.  Hence
\[
a_R\le U_R^{\rm cert}.
\]
The explicit capacity candidate is
\cref{eq-controlled-density-lp-hit-transfer} with one remaining transition;
the explicit information candidate is
\cref{eq-pre-hit-information-overlap-contact}.  Blackwell transfer is
\cref{eq-controlled-prehit-blackwell-value-monotonicity}, and the remaining
dynamical candidates are the same-unit entries admitted by
\cref{def-coordinatewise-dynamical-certificate}; the sparse-reward entry is
\cref{eq-observation-fiber-information-contact-bound}.  The active-sampling
entry is \cref{eq-active-acquisition-next-contact}; its master comparison is
admitted only after the same next-contact normalization has been represented.

Encoding adequacy and computation closure place every complete record and
comparison at the common majorant \(B^{\rm cert}\).  Taking the supremum over
records and then the maximum over times proves
\cref{eq-saturated-selector-certificate-envelope}.  Substitute the
recordwise inequality into the exact hazard decomposition
\cref{eq-selector-hazard-decomposition} and sum the \(H\) advantages to
obtain \cref{eq-saturated-success-certificate-envelope}.  Monotonicity gives
the family-functional statement.  Since every candidate is first evaluated
on the same record, the order of minimum and supremum is forced and no
cross-record certificate is introduced.

\end{proof}

\begin{theorem}[Exhaustive sparse-reward observation-to-transport accounting]
\label{thm-exhaustive-sparse-reward-observation-transport-accounting}

Fix a complete record \(\mathcal R_{\rm cur}\) from
\cref{def-complete-sparse-reward-observation-record}, a pre-hit time
\(t<H\), and the represented selector \(K_t\) used for the next transition.
At the common charged ceiling, the structural use of
\(q_\varepsilon=\kappa_\varepsilon\circ\phi\) belongs to exactly one of the
following master-profile modes.

\begin{enumerate}[label=\textup{(\roman*)}]
\item \emph{Exact cell dynamics.}  The cell map satisfies the target,
      public-structure, utility, and exact-intertwining gates.  Its record is
      assigned to the exact qualified \(\Abs/\Geom\) coordinate.

\item \emph{Compensated cell dynamics.}  The cell map satisfies the
      nondynamic qualification gates and has nonzero intertwining defect.
      A reversible or general-resolvent master certificate supplies the
      target-functional comparison.  Its record is assigned to the
      corresponding compensated coordinate, and its admissibility requires
      \[
      B_{\rm cur}\preceq B,
      \]
      with \(B_{\rm cur}\) and \(B_{\rm comp}\) evaluated term by term in
      \cref{eq-complete-sparse-reward-record-cost,eq-complete-compensated-observation-cost}.

\item \emph{Ambient full-carrier dynamics.}  The observation cells carry no
      admitted cell dynamics.  The executed kernel remains on the native
      carrier, and the observation history enters through the complete
      pre-hit selector record.  Its next-contact advantage obeys
      \begin{equation}
      \label{eq-exhaustive-sparse-reward-ambient-contact}
      \bigl(\operatorname{Adv}_t(K_t;\nu_t)\bigr)_+
      \le
      \min\!\left\{
      U_R^{\rm cert},
      \beta_{R,t}^{\rm fib}
      +\sqrt{\frac{\log2}{2}J_{R,t}^{\rm fib,\uparrow}}
      \right\},
      \end{equation}
      with both candidates clipped at one.  When a same-record capacity
      comparison is present, \(U_R^{\rm cert}\) includes
      \cref{eq-observation-regularity-capacity-hit-bound}; otherwise that
      capacity slot has its neutral value.
\end{enumerate}

The three modes are exhaustive and disjoint at the level of the selected
master certificate.  Predictor errors, counts, novelty bonuses, memories,
archives, and intrinsic rewards are causal postprocessings of the represented
pre-hit record and contribute no additional source kind.  Their complete
construction and deployment costs remain in
\cref{eq-complete-sparse-reward-record-cost}.  The terminal external reward
and completed first-hit record enter the accounting profile at their clock
times and supply no earlier source.

Restricting \cref{eq-saturated-selector-certificate-envelope} to complete
sparse-reward records gives
\begin{equation}
\label{eq-saturated-sparse-reward-complete-accounting}
\sup_{\substack{0\le t<H\\
R\in\mathscr R_{I,t}^{\rm cur}(B)}}
\bigl(\operatorname{Adv}_t(K_R;\nu_R)\bigr)_+
\le
\sup_{\substack{0\le t<H\\
R\in\mathscr R_{I,t}^{\rm cur}(B^{\rm cert})}}
U_R^{\rm cert},
\end{equation}
where \(\mathscr R_{I,t}^{\rm cur}\) is the subset of the existing complete
pre-hit records carrying
\cref{def-complete-sparse-reward-observation-record}.  The common ceiling
charges the observation, regularity, information, capacity, source,
cell-dynamics, execution, verification, and output comparisons before the
supremum is taken.

\end{theorem}

\begin{proof}

The four-mode master decomposition in
\cref{thm-complete-generic-source-theory} consists of ambient geometry,
exact qualified cell dynamics, reversible compensated dynamics, and
general-resolvent compensated dynamics.  Combining the latter two produces
the three displayed alternatives.  Their defining slots are part of the
certificate, so the alternatives are disjoint.  The exact and compensated
admission conditions and their complete charge are
\cref{def-master-compensated-quotient-geometry-certificate,def:master-certificate-dynamic-qualification,def:master-certificate-visibility}.

In the ambient mode,
\cref{eq-observation-fiber-information-contact-bound} bounds the actual
pre-hit selector.  Independently,
\cref{thm-saturated-same-record-prospective-certificate-envelope} bounds the
same survival-weighted advantage.  Their recordwise minimum proves
\cref{eq-exhaustive-sparse-reward-ambient-contact}.  Its capacity candidate
is \cref{eq-observation-regularity-capacity-hit-bound} after the stated
same-record conversion.

The sparse-reward identity is
\cref{eq-sparse-external-reward-zero-information}.  Constructor no-creation
and temporal admissibility are
\cref{prop-derived-channels-create-no-structure,def-pre-hit-temporally-admissible-source};
they assign derived novelty coordinates to their represented ancestry and
place terminal records in the accounting ledger.  Restricting the record set
in \cref{eq-saturated-selector-certificate-envelope} and retaining its
common cost majorant proves
\cref{eq-saturated-sparse-reward-complete-accounting}.

\end{proof}

\begin{figure}[tbp]
\centering
\begin{tikzpicture}[fw diagram,node distance=6mm and 9mm]
  \node[fw source,text width=3cm] (record) {complete controlled record\\at budget \(B\)};
  \node[fw provenance,text width=3cm,right=of record] (source)
    {saturated \(\Geom/\Abs\)\\source ledger};
  \node[fw box,text width=2.5cm,above right=6mm and 11mm of source] (hit)
    {arrival and\\discounted hit};
  \node[fw geom,text width=2.5cm,right=11mm of source] (flow)
    {mixing, time\\reach--avoid};
  \node[fw use,text width=2.5cm,below right=6mm and 11mm of source] (learn)
    {Bellman, policy\\learning};
  \draw[fw arrow] (record) -- (source);
  \draw[fw arrow] (source) -- (hit);
  \draw[fw arrow] (source) -- (flow);
  \draw[fw arrow] (source) -- (learn);
\end{tikzpicture}
\caption{The universal dynamical certificate is coordinatewise.  All entries
use one complete saturated record, while estimates in different units remain
separate.}
\label{fig-universal-dynamical-certificate}
\end{figure}

\section{Learned Policies and Continuous-System Transfer}
\label{sec-learned-controlled-continuous-transfer}

\subsection{Represented trajectory learning records}
\label{subsec-represented-trajectory-learning-records}

\begin{definition}[Complete dynamical learning record]
\label{def-complete-dynamical-learning-record}

Fix a controlled presentation at problem size \(n\) and a total budget
\(B=b(n)\).  The general ceiling \(b\) is retained symbolically; the
polynomial specialization is \(b_C(n)=n^C\).  A complete
dynamical learning record contains:

\begin{enumerate}[label=\textup{(\roman*)}]
\item a represented behavior policy, its independent-episode, stationary-path,
      or predictable-interaction sample, and the corresponding Part~IV
      dynamical-sample certificate;
\item for every binary target, failure, classification,
      committor-training, or policy-supervision coordinate, a complete
      dynamical label-noise record from
      \cref{def-complete-dynamical-label-noise-record}, including its clean
      analytic comparator, corruption rate, sample size, computational
      ceiling, debiasing or calibration, and behavior-to-deployment
      comparison.  A noiseless coordinate uses the zero-cost identity
      channel \(\eta=0\);
\item for every corrupted reward report, executed action, noisy observation,
      filter, or identifiable composite transition, a complete typed
      model-channel noise record from
      \cref{def-complete-dynamical-model-channel-noise-record}.  It includes
      its clean analytic comparator, public and latent filtrations, channel
      factorization or declared composite, sample size, coverage, tail and
      entropy certificates, inverse or calibration, behavior-to-deployment
      comparison, and saturated error profile.  An actual stochastic reward
      law is part of the true model; only a biased or imperfectly estimated
      report is a reward-corruption error;
\item for a partial-observation problem, the represented POMDP carrier,
      observation-measurable policy, exact analytic belief comparison,
      represented filter or memory, filter-stability recurrence, target and
      failure interfaces, and every clocked Lift, form, path-law, or
      resolvent comparison used at deployment;
\item the learned transition and reward records
      \((\widehat P,\widehat r)\), the learned geometry
      \(\widehat\Phi\), and every retained value, committor, action-value,
      policy, and selector record;
\item represented evaluation losses for the one-step model, geometry
      self-defect, Bellman or committor residual, policy comparison, target
      contact, and failure contact;
\item a behavior-to-deployment comparison: either represented importance
      weights with an audited bound, or a represented transport, density-ratio,
      or occupancy comparison;
\item the complete costs of data collection, training, geometry and policy
      selection, stored state, precision, deployment, and verification.
\end{enumerate}

All losses used in the bounded-loss theorems are normalized to \([0,1]\).
Another loss enters only with its represented sub-Gaussian, sub-exponential,
or predictable-variance certificate.  The budget split
\[
B=B_{\rm data}\oplus B_{\rm geom}\oplus B_{\rm model}
  \oplus B_{\rm value}\oplus B_{\rm pol}\oplus B_{\rm noise}
  \oplus B_{\rm ch}\oplus B_{\rm fil}\oplus B_{\rm info}
  \oplus B_{\rm dep}
\]
is part of the record.  The split records resource use; it creates no new
structural channel.

\end{definition}

\begin{definition}[Action-safe defect precision and effective dimension]
\label{def-action-safe-defect-precision}

Let \(H_\Phi\succeq0\) be the represented symmetric operator of a certified
dynamical geometry.  Let \(P_{\mathcal A}\) project onto the represented span
of increments generated by legal actions, and let \(P_{\rm safe}\) project
onto the represented subspace killed at the failure sector \(T^-\).  Assume
\(\operatorname{ran}P_{\rm safe}\subseteq
\operatorname{ran}P_{\mathcal A}\) and define the compressions
\begin{align*}
H_\Phi^{\mathcal A}
&=\left.P_{\mathcal A}H_\Phi P_{\mathcal A}
  \right|_{\operatorname{ran}P_{\mathcal A}},\\
H_\Phi^{\mathcal A,{\rm safe}}
&=\left.P_{\rm safe}H_\Phi^{\mathcal A}P_{\rm safe}
  \right|_{\operatorname{ran}P_{\rm safe}}.
\end{align*}
For \(\lambda>0\), put
\begin{align}
N_{\rm eff}^{\mathcal A}(\lambda)
&=\operatorname{Tr}\!\left[
H_\Phi^{\mathcal A}(H_\Phi^{\mathcal A}+\lambda I)^{-1}\right],
\label{eq-action-effective-dimension}\\
N_{\rm eff}^{\mathcal A,{\rm safe}}(\lambda)
&=\operatorname{Tr}\!\left[
H_\Phi^{\mathcal A,{\rm safe}}
(H_\Phi^{\mathcal A,{\rm safe}}+\lambda I)^{-1}\right].
\label{eq-action-safe-effective-dimension}
\end{align}

Let \(P_D\) and \(P_H\) be the already represented defect and
harmonic-reduction projections of
\cref{def-structural-pac-bayes-kl}.  On their retained positive range, the
control precision is
\begin{equation}
\label{eq-dynamical-defect-precision}
\mathbf M_{D,H}^{\rm dyn}(\lambda)
=\left.P_HP_D(H_\Phi^{\mathcal A,{\rm safe}}+\lambda I)
P_DP_H\right|_{R_{D,H}}.
\end{equation}
Every projection, retained basis, and spectral computation is charged in the
complete dynamical learning record.

\end{definition}

\begin{theorem}[Constrained-safe spectral capacity]
\label{thm-constrained-safe-spectral-capacity}

For every \(\lambda>0\), the nested compressions in
\cref{def-action-safe-defect-precision} satisfy
\begin{equation}
\label{eq-constrained-safe-effective-dimension-order}
N_{\rm eff}^{\mathcal A,{\rm safe}}(\lambda)
\le N_{\rm eff}^{\mathcal A}(\lambda)
\le N_{\rm eff,\Phi}(\lambda).
\end{equation}
If the represented safe target satisfies the compressed source condition
\(P_{\rm safe}f_*=(H_\Phi^{\mathcal A,{\rm safe}})^rg_{\rm safe}\),
\(0<r\le1\), then its source residual obeys
\begin{equation}
\label{eq-constrained-safe-source-residual}
\left\|\left[I-H_\Phi^{\mathcal A,{\rm safe}}
(H_\Phi^{\mathcal A,{\rm safe}}+\lambda I)^{-1}\right]
P_{\rm safe}f_*\right\|_2^2
+\|(I-P_{\rm safe})f_*\|_2^2
\le C_r\|g_{\rm safe}\|_2^2\lambda^{2r}
+\|(I-P_{\rm safe})f_*\|_2^2.
\end{equation}
The last term is the exact action/safety approximation residual; sample size
does not alter it.

\end{theorem}

\begin{proof}

The eigenvalues of a self-adjoint compression interlace those of its parent
operator.  The map \(u\mapsto u/(u+\lambda)\) is nonnegative and increasing,
so summing the interlacing inequalities first for \(P_{\mathcal A}\) and then
for \(P_{\rm safe}\) proves
\cref{eq-constrained-safe-effective-dimension-order}.  On
\(\operatorname{ran}P_{\rm safe}\), the stated compressed source condition
and the calculation of \cref{cor-source-condition-rate} bound the ridge
residual by \(C_r\|g_{\rm safe}\|_2^2\lambda^{2r}\).  Orthogonal decomposition adds the
unrepresented component \(\|(I-P_{\rm safe})f_*\|_2^2\) to both sides, proving
\cref{eq-constrained-safe-source-residual}.

\end{proof}

\begin{corollary}[Saturated accountability for learned policies]
\label{cor-saturated-learned-policy-accountability}

Consider a represented computation that learns a geometry, value or
committor, policy, and pre-hit action selector from finite represented data.
Include its training rule, sample, retained state, selected kernel, decoder,
policy evaluator, action comparisons, deployment state, and verifier calls in
one cost-\(B\) record.  Then:

\begin{enumerate}[label=\textup{(\roman*)}]
\item every prospective target-bearing object receives exactly one existing
      source-mode--state-provenance--coefficient-provenance index from
      \cref{cor-complete-structural-exhaustion-learned-sources};
\item its same-carrier target reach is bounded simultaneously by the exact
      Geom product and the generated-algebra projection ceiling in
      \cref{eq-learned-ambient-geom-and-statistical-minimum};
\item its target-arrival probability is bounded by the minimum of the
      prospective selector/source estimate and every applicable simultaneous
      Part~IV statistical estimate, as in
      \cref{thm-saturated-geometry-learning-selection-accountability,thm-master-structural-dynamical-learning-certificate};
\item trajectory memorization and sample-indexed tables remain charged lifted
      coordinates, while their population-orthogonal component belongs to
      \(J_{\rm res}\).  A population-aligned prospective use is reclassified
      into the existing qualified \(\Geom/\Abs\) ledger at its complete cost
      only when the same record supplies the required channel predicate and
      pre-hit comparison.
\end{enumerate}

Thus learning changes which represented geometry and policy are selected; it
does not enlarge the structural classification or omit the cost of selecting
and deploying them.

\end{corollary}

\begin{proof}

The complete learning and deployment computation satisfies the hypotheses of
\cref{thm-saturated-geometry-learning-selection-accountability}.  Apply
\cref{cor-complete-structural-exhaustion-learned-sources} to its prospective
records and
\cref{thm-two-level-learned-readout-provenance,thm-exhaustive-residual-normal-form}
to its empirical and population components.  The learned policy is a
represented selector in
\cref{thm-controlled-policy-channel-decomposition}; hence all policy
evaluation and action-selection costs remain in the same derivation.  These
cited conclusions give the four assertions.

\end{proof}

\subsection{Learned-Control Coordinates and Bellman Propagation}
\label{subsec-learned-control-bellman-propagation}

\begin{definition}[Learned-control certificate coordinates]
\label{def-learned-control-certificate-coordinates}

For one complete dynamical learning record, define
\[
\mathbf U_{\rm learn}^{\rm dyn}
=\Bigl(
\begin{aligned}
&U_{\rm model},U_{\rm geom},U_{\rm value},U_{\rm policy},
 U_{\rm shift},U_{\rm noise},\mathbf U_{\rm ch},U_{\rm POMDP},\\
&U_{\rm info},\mathbf H_{\rm ctl},U_{\rm safe},U_{\rm hit},
 U_{\rm Cap}^{\rm noise},U_{\rm return},U_{\rm regret}
\end{aligned}
\Bigr).
\]
Each entry is an upper bound in its named unit:

\begin{enumerate}[label=\textup{(\roman*)}]
\item \(U_{\rm model}\) bounds the represented joint actionwise one-step
      Bellman discrepancy \(\varepsilon_M\) of
      \cref{thm-dynamical-model-bellman-policy-propagation} on the retained
      value class \(\mathcal V_B\).  An optimal-operator defect alone is a
      value candidate but not a learned-greedy policy candidate.
\item \(U_{\rm geom}\) bounds the represented form or operator discrepancy
      \(\|\widehat H_\Phi-H_\Phi\|_{\rm op}\), together with the induced
      projector error at a certified eigengap.
\item \(U_{\rm value}\) bounds a value, action-value, or committor error in
      its declared norm after the model, geometry, and Bellman residual have
      been propagated by a represented contraction, resolvent, or functional-
      inequality certificate.
\item \(U_{\rm policy}\) bounds deployment-policy loss relative to the
      budget-reachable optimum.
\item \(U_{\rm shift}\) bounds the behavior-to-deployment risk change.
\item \(U_{\rm noise}\) bounds clean binary supervised risk after explicit
      corruption transport, maximized over the binary supervised coordinates
      in the record on their simultaneous confidence event.  It records
      \(n\), the complete computational ceiling, sample size \(m\), and each
      noise rate or Massart ceiling.
\item \(\mathbf U_{\rm ch}\) stores reward, execution, observation, filter,
      composite-kernel, and path-law errors with their declared norms.
      \(U_{\rm POMDP}\) stores filter-to-hit and filter-to-value consequences
      only after their represented converters.  \(U_{\rm info}\) stores the
      eligible target-identification converse, and
      \(\mathbf H_{\rm ctl}\) stores the noninterchangeable entropy outputs
      of \cref{def-controlled-typed-entropy-outputs}.
\item \(U_{\rm safe}\) bounds failure-before-target probability;
      \(U_{\rm hit}\) bounds target-before-failure arrival.
\item \(U_{\rm Cap}^{\rm noise}\) bounds learned-boundary lifted condenser
      capacity and remains in capacity units.
\item \(U_{\rm return}\) bounds discounted-return suboptimality, and
      \(U_{\rm regret}\) bounds normalized sequential regret.
\end{enumerate}

A coordinate takes the minimum of all valid certified upper bounds expressed
in the same unit.  Its neutral value is \(+\infty\) for an unbounded error or
loss coordinate and one for a probability coordinate.  Distinct coordinates
are joined only by an explicit propagation theorem.

\end{definition}

\begin{theorem}[Model, Bellman, and policy propagation]
\label{thm-dynamical-model-bellman-policy-propagation}

Let \(\mathcal T\) and \(\widehat{\mathcal T}\) be true and represented
learned discounted Bellman operators with common discount
\(0<\beta<1\), and let \(\mathcal T^a,\widehat{\mathcal T}^{\,a}\) be
their one-step action maps.  Let \(v\in\mathcal V_B\) be represented, put
\[
\varepsilon_M
=\sup_{w\in\mathcal V_B}
  \sup_{x}\sup_{a\in\mathcal A_B(x)}
  |\mathcal T^aw(x)-\widehat{\mathcal T}^{\,a}w(x)|,
\qquad
\varepsilon_B=\|\widehat{\mathcal T}v-v\|_\infty,
\]
and assume that the fixed points and iterates used below remain in
\(\mathcal V_B\).  For a history- or belief-action correspondence, the two
inner suprema are over its represented one-step selectors, equivalently over
the associated one-step policy operators.  This joint defect implies both
\(\|\mathcal Tw-\widehat{\mathcal T}w\|_\infty\le\varepsilon_M\) and
\(\|\mathcal T^\pi w-\widehat{\mathcal T}^{\,\pi}w\|_\infty
\le\varepsilon_M\) for every admitted selector \(\pi\).

Suppose, in addition, that a complete channel record has
converted its typed errors to the same deployment-law Bellman sup norm and
has a represented disjoint telescoping sequence of action maps
\(\mathcal T^{(0),a},\ldots,\mathcal T^{(5),a}\), from the true to the
learned map.  For \(1\le j\le5\), put
\[
\Delta_j
=\sup_{w\in\mathcal V_B}\sup_x\sup_{a\in\mathcal A_B(x)}
\left|\mathcal T^{(j),a}w(x)-\mathcal T^{(j-1),a}w(x)\right|,
\]
and require the same record to identify, in the declared order,
\[
(\Delta_1,\ldots,\Delta_5)
=\bigl(\varepsilon_r^{\rm Bell},
\beta\varepsilon_P^{\rm Bell},
\varepsilon_A^{\rm Bell},
\varepsilon_O^{\rm Bell},
\varepsilon_{\rm fil}^{\rm Bell}\bigr).
\]
Then the actionwise triangle inequality gives
\begin{equation}
\label{eq-channel-to-bellman-model-error}
\varepsilon_M
\le
\varepsilon_r^{\rm Bell}
+\beta\varepsilon_P^{\rm Bell}
+\varepsilon_A^{\rm Bell}
+\varepsilon_O^{\rm Bell}
+\varepsilon_{\rm fil}^{\rm Bell}
=:\varepsilon_M^{\rm ch}.
\end{equation}
An absent component is zero only when the corresponding interface is exact;
an unavailable converter makes the channel candidate \(+\infty\).  Components
with overlapping ancestry occur in one chosen telescoping step, rather than
being counted in two marginal steps.  In particular, if
\(\varepsilon_A^{\rm Bell}\) already includes the reward and transition
change caused by execution, those changes are excluded from the separate
reward and transition steps.  Then
\begin{equation}
\label{eq-learned-value-propagation}
\|V^*-v\|_\infty
\le\frac{\varepsilon_B+\varepsilon_M}{1-\beta}.
\end{equation}
If \(\widehat\pi_v\) is an admitted represented greedy selector for
\(\widehat{\mathcal T}\) at \(v\), then
\begin{equation}
\label{eq-learned-policy-propagation}
\|V^*-V^{\widehat\pi_v}\|_\infty
\le
\frac{2\beta(\varepsilon_B+\varepsilon_M)}{(1-\beta)^2}
+\frac{2\varepsilon_M}{1-\beta}.
\end{equation}

For a fixed policy, if
\(\|r^\pi-\widehat r^\pi\|_\infty\le\varepsilon_r\) and
\(\|(P^\pi-\widehat P^\pi)w\|_\infty\le\varepsilon_P\) on the retained
value class, then
\begin{equation}
\label{eq-fixed-policy-model-propagation}
\|V^\pi-\widehat V^\pi\|_\infty
\le\frac{\varepsilon_r+\beta\varepsilon_P}{1-\beta}.
\end{equation}
The same inequalities hold for a represented committor contraction, with its
declared contraction modulus in place of \(\beta\).  For an undiscounted
killed problem they hold with the represented killed-resolvent norm in place
of \((1-\beta)^{-1}\).
Raw \(L^2\) reward error, channel total variation, filter error, entropy, and
capacity are not substituted into
\cref{eq-channel-to-bellman-model-error} without a represented same-norm
converter.

\end{theorem}

\begin{proof}

The true residual satisfies
\[
\|\mathcal Tv-v\|_\infty
\le\|\mathcal Tv-\widehat{\mathcal T}v\|_\infty
  +\|\widehat{\mathcal T}v-v\|_\infty
\le\varepsilon_M+\varepsilon_B.
\]
Apply the contraction residual argument of
\cref{thm-controlled-bellman-residual-policy-loss} to obtain
\cref{eq-learned-value-propagation}.  Pointwise greediness and the two parts
of the joint defect give
\[
\mathcal Tv
\le\widehat{\mathcal T}v+\varepsilon_M
=\widehat{\mathcal T}^{\,\widehat\pi_v}v+\varepsilon_M
\le\mathcal T^{\widehat\pi_v}v+2\varepsilon_M.
\]
Since \(\mathcal Tv\ge\mathcal T^{\widehat\pi_v}v\), this proves
\(\|\mathcal Tv-\mathcal T^{\widehat\pi_v}v\|_\infty
\le2\varepsilon_M\).  The approximate-greedy one-step comparison gives
\[
\|V^*-V^{\widehat\pi_v}\|_\infty
\le\frac{2\beta}{1-\beta}\|V^*-v\|_\infty
   +\frac{2\varepsilon_M}{1-\beta};
\]
substitution proves \cref{eq-learned-policy-propagation}.  For a fixed
policy, subtract the two fixed-point equations, use the stochastic-kernel
contraction, and rearrange to obtain
\cref{eq-fixed-policy-model-propagation}.  The committor and killed-resolvent
statements repeat the same calculation with the stated inverse bound.
Finally, applying the actionwise triangle inequality to
\(\mathcal T^{(0),a}-\mathcal T^{(5),a}\), then taking the displayed
suprema, gives \cref{eq-channel-to-bellman-model-error}.  Substitution into
the preceding contraction inequalities gives every channelized
specialization.  Without an actionwise or one-step-policy converter, the
policy candidate is \(+\infty\), even if an optimal-operator value candidate
is finite.

\end{proof}

\subsection{Generalization Envelopes and Budget-Reachable Value}
\label{subsec-dynamical-generalization-budget-value}

\begin{definition}[Dynamical generalization envelope and budget-reachable value]
\label{def-dynamical-generalization-envelope}

For a represented bounded loss class \(\mathcal F\) in a complete dynamical
learning record, let
\[
\mathfrak G_{\rm dyn}(\mathcal F;B,m,\delta)
\]
be the minimum of the valid Part~IV independent-episode, mixing-block,
sequential-Rademacher, martingale PAC--Bayes, and effective-dimension upper
bounds for its population-minus-empirical loss.  Each candidate includes its
operator-selection, confidence, dependence, and behavior-to-deployment terms.
The dynamical candidates are supplied by
\cref{thm-beta-mixing-structural-blocking,thm-sequential-learned-operator-cover,thm-martingale-structural-pac-bayes,cor-dynamical-source-condition-rate};
the independent-episode candidate is
\cref{thm-fixed-operator-rademacher-bound,thm-learned-operator-union-bound} on
the represented episode-loss class.
Candidates are compared only after they bound the same population law and the
same loss.  The value is \(+\infty\) when no candidate applies.

For a represented flip-symmetric binary loss class, define the clean
supervised-risk envelope
\begin{equation}
\label{eq-dynamical-clean-generalization-envelope}
\mathfrak G_{\rm dyn}^{\rm clean}
(\mathcal F;B,m,\eta_+,\delta)
=
\min\left\{1,
U_{\rm noise}^{\rm clean}(n,B,m,\eta_+,\delta),
U_{\rm cal}^{\rm clean}(\mathcal F)
\right\}.
\end{equation}
Here \(U_{\rm cal}^{\rm clean}(\mathcal F)\) ranges over represented
same-law, same-unit calibration consequences of
\cref{thm-uniform-dynamical-label-noise-bound}; an unavailable candidate is
\(+\infty\).  Unlike \(\mathfrak G_{\rm dyn}\), this quantity bounds the
clean supervised risk or excess itself, not merely a population--empirical
deviation.  It is therefore inserted into an operator, Bellman, committor,
or geometry coordinate only through a represented loss-to-coordinate
propagation theorem.  Reward, transition, and state-observation losses that
are not binary label losses continue to use \(\mathfrak G_{\rm dyn}\), or the
typed channel profiles of
\cref{def-saturated-dynamical-channel-error-profiles} when a represented
correction applies.  The latter enter a model or Bellman coordinate only
through \cref{eq-channel-to-bellman-model-error}.

Let \(\Pi_B\) be the represented policies of complete saturated cost at most
\(B\), including data collection and deployment.  The budget-reachable value
and budget gap are
\begin{equation}
\label{eq-budget-reachable-control-value}
V_B^*(x)=\sup_{\pi\in\Pi_B}V^\pi(x),
\qquad
\Delta_B^{\rm ctrl}=\|V^*-V_B^*\|_\infty.
\end{equation}
This pointwise supremum is only a comparison benchmark unless a rectangularity
gate is present.  The Bellman branch is available only when the record
supplies a represented state/history action correspondence
\(\mathcal A_B\) such that (a) every measurable selector used by the
argument, including each greedy selector, belongs to \(\Pi_B\) with its
common evaluation and deployment cost charged, and (b) \(\Pi_B\) is closed
under state/history-wise one-step continuation.  Under this represented
rectangularity gate, the restricted operator \(\mathcal T_B\) is a
\(\beta\)-contraction and its fixed point is \(V_B^*\).  Without it,
\(V_B^*\) remains the displayed pointwise benchmark, but the restricted
Bellman fixed-point, greedy-policy, and return-propagation candidates are
\(+\infty\), unless a separately named rectangular closure is constructed
and charged at its enlarged ceiling.

For a partial-observation record, let
\(\Pi_{B,{\rm obs}}\subseteq\Pi_B\) be the represented policies measurable
with respect to the public observation filtration, including their filters,
memory, and deployment costs.  Define
\begin{equation}
\label{eq-budget-reachable-pomdp-value}
V_{B,{\rm obs}}^*(x)
=\sup_{\pi\in\Pi_{B,{\rm obs}}}V^\pi(x),
\qquad
\Delta_B^{\rm POMDP}
=\|V_B^*-V_{B,{\rm obs}}^*\|_\infty.
\end{equation}
The observation-measurable class requires the analogous represented
belief/history rectangularity gate before \(V_{B,{\rm obs}}^*\) is called a
Bellman fixed point.  In particular, statewise or beliefwise stitching of
individually affordable policies is not free.
For every \(\widehat\pi\in\Pi_{B,{\rm obs}}\), the exact three-term
decomposition is
\begin{equation}
\label{eq-pomdp-budget-statistical-decomposition}
V^*-V^{\widehat\pi}
=(V^*-V_B^*)+(V_B^*-V_{B,{\rm obs}}^*)
+(V_{B,{\rm obs}}^*-V^{\widehat\pi}).
\end{equation}
The middle term is an approximation/information restriction, not an
estimation or filter error.  In a fully observed record take
\(\Pi_{B,{\rm obs}}=\Pi_B\) and \(\Delta_B^{\rm POMDP}=0\).
For every learned \(\widehat\pi\in\Pi_B\), the exact decomposition is
\begin{equation}
\label{eq-control-budget-statistical-decomposition}
V^*-V^{\widehat\pi}
=(V^*-V_B^*)+(V_B^*-V^{\widehat\pi}).
\end{equation}
The exact gaps \(\Delta_B^{\rm ctrl}\) and \(\Delta_B^{\rm POMDP}\) are
analytic comparison quantities.  An operational certificate stores
represented upper bounds
\[
\Delta_B^{\rm ctrl}\le\overline\Delta_B^{\rm ctrl},
\qquad
\Delta_B^{\rm POMDP}\le\overline\Delta_B^{\rm POMDP},
\]
including the cost of evaluating or proving them.  A missing upper
certificate has value \(+\infty\).  The exact gaps may still be displayed in
an oracle decomposition, but they are not free certificate inputs.

\end{definition}

\section{Master Learning Certificate and Continuous-System Transfer}
\label{sec-master-dynamical-learning-certificate}

\begin{definition}[Master leaf event and same-norm candidate aggregation]
\label{def-master-dynamical-leaf-candidate-aggregation}

Fix a complete dynamical learning record \(\mathcal R\) of sample size
\(m\) and a represented learned policy \(\widehat\pi_v\) greedy for
\(\widehat{\mathcal T}\) at \(v\).  Let \(\mathcal I(\mathcal R)\) be the
finite represented set of every stochastic leaf certificate eligible for a
displayed minimum: generic generalization, geometry, binary corruption,
channel estimation and entropy continuity, filter stability, path
comparison, gate or interface estimation, and estimated information.  Assign
leaf probabilities \((\delta_i)_{i\in\mathcal I(\mathcal R)}\) with
\(\sum_i\delta_i\le\delta\).  Each data-dependent leaf uses a pre-specified
event uniform over its affordable output hull and independent of which output
is selected.  Work on
\[
\mathcal E_{\mathcal R}=\bigcap_{i\in\mathcal I(\mathcal R)}\mathcal E_i,
\qquad
\Pr(\mathcal E_{\mathcal R})\ge1-\delta.
\]
Deterministic identities have \(\delta_i=0\), and a candidate whose event is
not included is unavailable.  Internal allocations, including
\(\delta_{\rm ch}\), partition these leaves and are not counted a second
time.

For the Bellman branch, assume the applicable fully observed or
observation-measurable rectangularity gate of
\cref{def-dynamical-generalization-envelope}.  Then \(\mathcal T\) is the
corresponding restricted Bellman operator and its fixed point is
\(V_{B,{\rm obs}}^*\), with the fully observed convention
\(V_{B,{\rm obs}}^*=V_B^*\).  The learned-greedy branch also requires the
joint actionwise defect in
\cref{thm-dynamical-model-bellman-policy-propagation}.  If either gate is
absent, \(U_{\rm value},U_{\rm policy}\), and the Bellman-propagated
\(U_{\rm return}\) below equal \(+\infty\); the other coordinates remain
available.  On \(\mathcal E_{\mathcal R}\), define
\(U_M^{\rm noise},U_B^{\rm noise},U_\Phi^{\rm noise}\) to be the smallest
same-norm upper bounds obtained by propagating
\(\mathfrak G_{\rm dyn}^{\rm clean}\) through a represented
loss-to-coordinate calibration theorem for the model, Bellman, and geometry
coordinates, respectively.  A coordinate that is not a binary supervised
loss, or lacks such a propagation theorem, receives the value \(+\infty\).
Define \(U_M^{\rm ch},U_B^{\rm ch},U_\Phi^{\rm ch}\) analogously as the
smallest same-law, same-norm consequences of the typed reward, execution,
observation, composite-kernel, and filter profiles after the represented
converters in
\cref{thm-dynamical-affine-reward-channel-transport,thm-controlled-execution-channel-comparison,thm-controlled-pomdp-filter-transfer,eq-channel-to-bellman-model-error}.
An unavailable conversion has value \(+\infty\).
Define
\begin{align*}
\overline\varepsilon_M
&=\min\left\{
\widehat\varepsilon_M
  +\mathfrak G_{\rm dyn}(\mathcal F_{\rm model};B,m,\delta_M),
U_M^{\rm noise},U_M^{\rm ch}\right\},\\
\overline\varepsilon_B
&=\min\left\{
\widehat\varepsilon_B
  +\mathfrak G_{\rm dyn}(\mathcal F_{\rm Bell};B,m,\delta_B),
U_B^{\rm noise},U_B^{\rm ch}\right\},\\
\overline\varepsilon_\Phi
&=\min\left\{
\widehat\varepsilon_\Phi
  +\mathfrak G_{\rm dyn}(\mathcal F_{\rm geom};B,m,\delta_\Phi),
U_\Phi^{\rm noise},U_\Phi^{\rm ch}\right\}.
\end{align*}
Here each empirical term and class is the represented coordinate from
\cref{def-complete-dynamical-learning-record}.  Then the coordinates of
\(\mathbf U_{\rm learn}^{\rm dyn}\) use this one simultaneous event and
these three same-norm candidate minima.
\end{definition}

\begin{theorem}[Master structural learning certificate for controlled systems]
\label{thm-master-structural-dynamical-learning-certificate}
Under the record, rectangularity, confidence, and same-norm aggregation
hypotheses of
\cref{def-master-dynamical-leaf-candidate-aggregation}, the core coordinates
of \(\mathbf U_{\rm learn}^{\rm dyn}\) may be chosen as follows:
\begin{align}
U_{\rm model}
&=\overline\varepsilon_M,
\label{eq-master-dyn-model}\\
U_{\rm geom}
&=\left(\overline\varepsilon_\Phi,
\frac{2\overline\varepsilon_\Phi}{g_*}\right),
\label{eq-master-dyn-geometry}\\
U_{\rm value}
&=\frac{\overline\varepsilon_B+\overline\varepsilon_M}{1-\beta},
\label{eq-master-dyn-value}\\
U_{\rm policy}
&=\frac{2\beta(\overline\varepsilon_B+
   \overline\varepsilon_M)}{(1-\beta)^2}
  +\frac{2\overline\varepsilon_M}{1-\beta},
\label{eq-master-dyn-policy}\\
U_{\rm return}
&=\overline\Delta_B^{\rm ctrl}
 +\overline\Delta_B^{\rm POMDP}+U_{\rm policy}.
\label{eq-master-dyn-return}
\end{align}
The first entry in \cref{eq-master-dyn-geometry} is the operator coordinate.
The second is the projector coordinate at a certified gap \(g_*>0\), and is
assigned \(+\infty\) when no positive gap is certified.
Every \(U_Q^{\rm noise}\) is also \(+\infty\) unless its calibration bounds
the same norm as the generic candidate.  Thus binary label transport never
replaces a transition, reward, or operator norm by an unlike risk.
The same rule applies to \(U_Q^{\rm ch}\): a channel-entropy, row-TV,
\(L^2\), filter, or path-law number enters only after conversion to the norm
of that candidate.
Every target-qualified learned coordinate additionally carries the exact
source index, contextual-charge vector, and noise-information vector of
\cref{def-complete-source-resolved-learning-ledger}.  Model, Bellman, and
policy errors remain in their displayed norms; the source-resolved target
candidate enters only the arrival coordinate
\cref{eq-master-dyn-hit} or another coordinate with a represented
same-unit converter.

\end{theorem}

\begin{proof}
On the simultaneous leaf event, each candidate in a displayed minimum is an
upper bound for the same population object in the same norm.  Their minima
therefore remain upper bounds.  The model and geometry coordinates are the
definitions of the aggregated errors, while
\cref{thm-dynamical-model-bellman-policy-propagation} gives the value and
policy coordinates.  The exact return decomposition in
\cref{def-dynamical-generalization-envelope}, followed by the represented
upper certificates for its two analytic gaps, gives
\cref{eq-master-dyn-return}.
\end{proof}

\begin{proposition}[Auxiliary channel, information, safety, and access coordinates]
\label{prop-master-dynamical-auxiliary-coordinates}
Under the hypotheses of
\cref{def-master-dynamical-leaf-candidate-aggregation}, the remaining
coordinates are
\begin{align}
U_{\rm shift}
&=\min\{U_{\rm iw},U_{T_1},U_{T_2},U_{\rm occ}\},
\label{eq-master-dyn-shift}\\
U_{\rm noise}
&=\max_{q\in\mathcal Q_{\rm bin}}
U_{{\rm noise},q}^{\rm clean}
(n,B,m,\eta_{q,+},\delta_q),
\label{eq-master-dyn-noise}\\
\mathbf U_{\rm ch}
&=\Bigl(
\begin{aligned}
&\operatorname{Err}_{r,\mathsf N_r}^{\rm sat}(n,B,m,\delta_r),
 \operatorname{Err}_{A,\mathsf N_A}^{\rm sat}(n,B,m,\delta_A),\\
&\operatorname{Err}_{O,\mathsf N_O}^{\rm sat}(n,B,m,\delta_O),
 \operatorname{Err}_{{\rm fil},\mathsf N_{\rm fil}}^{\rm sat}(n,B,m,\delta_{\rm fil}),\\
&\operatorname{Err}_{{\rm comp},\mathsf N_{\rm comp}}^{\rm sat}(n,B,m,\delta_{\rm comp}),
 U_{\rm path}(n,B,m,H,\delta_{\rm path})
\end{aligned}
\Bigr),
\label{eq-master-dyn-channel-vector}\\
U_{\rm POMDP}
&=\bigl(U_{\rm POMDP}^{H},U_{\rm POMDP}^{\rm value}\bigr),
\qquad
U_{\rm info}
=\min\left\{1,
\frac{I_0^{(m)}+C_{\rm pub}^{[H]}+1}{\log_2N_n}
\right\},
\label{eq-master-dyn-pomdp-information}\\
U_{\rm safe}
&=\min\{1,U_{\rm fail\text{-}comm},U_{\rm death\text{-}cap},
          U_{\rm fail\text{-}gate}^{\rm noise},U_{\rm safe}^{\rm path},
          U_{\rm POMDP}^{F,H}\},
\label{eq-master-dyn-safe}\\
U_{\rm hit}
&=\min\{1,U_{\rm comm},U_{\rm selector},U_{\rm resolvent},
          U_{\Caus},U_{\Caus}^{\rm noise},U_{\Lift}^{H},
          U_{\Lift,{\rm noise}}^{H},U_{\rm hit}^{\rm path},
          U_{\rm POMDP}^{H},U_{\rm learn}^{\rm src}\},
\label{eq-master-dyn-hit}\\
U_{\rm Cap}^{\rm noise}
&=b_\pi\left[
\operatorname{Cap}_\Phi(T,F)
+\omega_{\Phi,T,F}^{\rm Cap}
(\varepsilon_T^\mu,\varepsilon_F^\mu)
\right],
\label{eq-master-dyn-noisy-capacity}\\
U_{\rm regret}
&=U_{\rm seq}.
\label{eq-master-dyn-regret}
\end{align}
The capacity entry is assigned \(+\infty\) when its reference-law
comparison, form comparison, or boundary-stability modulus is unavailable.
The source-resolved learning entry has neutral value one unless its complete
learning ledger and target converter are present in the same controlled
record.
Here \(\mathcal Q_{\rm bin}\) is the finite represented index set of binary
supervised coordinates in the complete record, with the declared confidence
allocations \(\delta_q\).
The maximum in \cref{eq-master-dyn-noise} is zero when the record has no
binary supervised coordinate; the identity channel \(\eta=0\) is used when
such a coordinate is present but uncorrupted.
Each \(\operatorname{Err}_{q,\mathsf N_q}^{\rm sat}\) in
\cref{eq-master-dyn-channel-vector} retains its displayed
\((n,B,m,\delta_q)\) arguments and its own norm.  The allocation satisfies
\[
\delta_r+\delta_A+\delta_O+\delta_{\rm fil}
+\delta_{\rm comp}+\delta_{\rm path}\le\delta_{\rm ch},
\]
and \(\delta_{\rm ch}\) participates in the master simultaneous confidence
budget.  The path entries
\(U_{\rm hit}^{\rm path}\) and \(U_{\rm safe}^{\rm path}\) are a represented
reference probability plus the applicable TV or KL path shift, clipped at
one; raw path distance is not itself an upper bound on either event.
Similarly,
\[
U_{\rm POMDP}^{H}=\min\{1,U_b^H+\Delta_{\rm fil}^H\},
\qquad
U_{\rm POMDP}^{F,H}=\min\{1,U_b^{F,H}+\Delta_{\rm fil}^{F,H}\},
\]
where each \(U_b\) is a represented and charged upper bound for the same
exact-belief reference event.  An uncomputed analytic probability has value
one, so it cannot improve an affordable prospective minimum.
\(U_{\rm POMDP}^{\rm value}\) is
\cref{eq-controlled-pomdp-filter-value-transfer} in value units only when its
restricted benchmark, performance-difference identity, and evaluation
interface are represented; otherwise it is \(+\infty\).
The information entry is
\cref{eq-controlled-multichannel-fano-success} when target identification is
the declared success event and the complete record contains the represented,
costed output-to-target-index decoder and any disjoint-sector verification;
it is one otherwise.  A missing information budget has value \(+\infty\)
before probability clipping.  The entropy vector is
\(\mathbf H_{\rm ctl}\) from
\cref{eq-controlled-typed-entropy-vector} and remains coordinatewise.
Here \(U_{\rm iw}\) is the bounded-importance-weight certificate;
\(U_{T_1},U_{T_2}\) are the applicable transport-entropy bounds;
\(U_{\rm occ}\) is another represented occupancy comparison;
\(U_{\rm fail\text{-}comm}\) and \(U_{\rm comm}\) are the propagated
failure and target committor bounds; \(U_{\rm death\text{-}cap}\) is zero
under certified polar excision and otherwise is the represented monotone
failure-capacity bound; \(U_{\rm selector}\) and \(U_{\rm resolvent}\) are
the existing prospective-selector and killed-resolvent bounds;
\(U_{\rm learn}^{\rm src}\) is the same-record target-gain minimum in
\cref{eq-source-resolved-learning-target-gain-minimum};
\(U_{\Caus}\) is admitted only through a represented contact converter;
\(U_{\Lift}^{H}\) and \(U_{\Lift}^{\lambda}\) are admitted only after
conversion to the arrival normalization used by this learned record; and
\(U_{\rm seq}\) is the Part~IV sequential structural-regret bound.
The noise-corrected entries are those of
\cref{thm-controlled-noise-corrected-caus-aiming,thm-controlled-noisy-gate-lift-transfer}.
In particular, \(U_{\rm fail\text{-}gate}^{\rm noise}\) is the deployment
failure-gate error added to a represented clean failure probability, and
\(U_{\Lift,{\rm noise}}^H\) is the learned lifted arrival plus the clocked
interface error, both clipped at one.  Raw discounted entries
\(U_{\Lift}^{\lambda}\) and \(U_{\Lift,{\rm noise}}^\lambda\) remain in the
discounted coordinate \(U_{\rm disc}\); they do not enter the finite-horizon
minimum without an explicit resolvent-to-horizon converter.
\end{proposition}

\begin{proof}
Each displayed entry is the clipped same-event consequence of the theorem
or definition cited in its accompanying paragraph.  The channel vector uses
the disjoint internal confidence allocation, the path coordinates add the
path distance to a charged reference event, the POMDP coordinates use the
filter-transfer theorem, and the information coordinate uses the represented
decoder form of Fano's inequality.  Unavailable conversions take their
declared neutral value, so none can improve the corresponding minimum.
\end{proof}

\begin{proposition}[Explicit deployment-shift and committor candidates]
\label{prop-master-dynamical-shift-committor-candidates}
Under the hypotheses of
\cref{def-master-dynamical-leaf-candidate-aggregation}, the following
represented formulas instantiate the shift and committor entries in
\cref{prop-master-dynamical-auxiliary-coordinates}.  For a represented
deployment loss (F) with
\(\operatorname{Lip}_{d_\Phi}(F)=L_F\), behavior occupancy \(d^b\),
deployment occupancy \(d^\pi\), and represented density ratio
\(w= d^\pi/d^b\le W\), the eligible shift entries are
\begin{align}
U_{\rm iw}
&=\widehat{\mathbb E}_{d^b}[wF]
  +\mathfrak G_{\rm dyn}(w\mathcal F;B,m,\delta_{\rm iw}),
\label{eq-master-dyn-importance-weight}\\
U_{T_p}
&=\widehat{\mathbb E}_{d^b}[F]
  +\mathfrak G_{\rm dyn}(\mathcal F;B,m,\delta_{\rm shift})
  +L_F\sqrt{2C_{T_p}\operatorname{KL}(d^\pi\Vert d^b)},
\qquad p\in\{1,2\},
\label{eq-master-dyn-transport-shift}\\
U_{\rm occ}
&=\widehat{\mathbb E}_{d^b}[F]
  +\mathfrak G_{\rm dyn}(\mathcal F;B,m,\delta_{\rm shift})
  +\|d^\pi-d^b\|_{\rm TV}.
\label{eq-master-dyn-occupancy-shift}
\end{align}
All three lines bound deployment risk.  The final term in the last two lines
is the certified behavior-to-deployment change for \([0,1]\)-valued loss.

For learned target and failure committors \(\widehat h_+,\widehat h_-\),
let \(\kappa_\pm\) be certified killed-resolvent norms and let
\(\overline\varepsilon_\pm\) be their empirical residuals plus the applicable
dynamical generalization envelopes.  Then
\begin{align}
U_{\rm comm}
&=\min\{1,\langle\mu_0,\widehat h_+\rangle
             +\kappa_+\overline\varepsilon_+\},
\label{eq-master-dyn-committor-hit}\\
U_{\rm fail\text{-}comm}
&=\min\{1,\langle\mu_0,\widehat h_-\rangle
             +\kappa_-\overline\varepsilon_-\},
\label{eq-master-dyn-committor-failure}\\
U_{\rm selector}
&=\min\{1,\operatorname{Enum}_\nu
             +H\operatorname{Sel}^{\rm sat}(\Psi_{\rm ctl}(B,H,n))\},
\label{eq-master-dyn-selector-hit}\\
U_{\rm resolvent}
&=\min\{1,eA_{T^+}(H^{-1};L,\mu_0)\}.
\label{eq-master-dyn-resolvent-hit}
\end{align}
Certified polarity sets \(U_{\rm death\text{-}cap}=0\).  At positive failure
capacity, this entry is the upper probability supplied by the represented
capacity-to-boundary-risk comparison, and equals one when that comparison is
absent.  Positive capacity alone is not substituted for a probability.

Each committor, propagation, shift, mixing, and access entry uses the smallest
same-unit consequence of the audited functional-inequality profile.  A
functional inequality with no represented conversion to that coordinate is
excluded from its minimum.  Thus the master certificate uses the full
profile without adding logically dependent inequalities or unlike units.

\end{proposition}

\begin{proof}

On \(\mathcal E_{\mathcal R}\), every stochastic leaf inequality holds for
the actually selected data-dependent output.  Therefore each generic
population loss is at most its empirical value plus its envelope.  For a
binary supervised coordinate,
\cref{thm-uniform-dynamical-label-noise-bound} and its represented same-norm
calibration supply the alternative noise-corrected candidate.  The typed
channel profiles and their same-norm reward, execution, observation, and
filter converters supply the \(U_Q^{\rm ch}\) candidates through
\cref{eq-channel-to-bellman-model-error}.  Since these candidates bound the
same population object in the same norm on the same event, their minimum is
again an upper bound.  This proves the three barred-error definitions and
\cref{eq-master-dyn-model}.  Geometry selection and projector transfer are
\cref{thm-two-level-geometry-certification}, which proves
\cref{eq-master-dyn-geometry}.

Under the represented rectangularity gate,
\(V_{B,{\rm obs}}^*\) is the restricted Bellman fixed point; without that
gate the theorem has already assigned the Bellman branch \(+\infty\).
Under the joint actionwise model defect, the learned-greedy selector satisfies
the two-sided one-step comparison proved in
\cref{thm-dynamical-model-bellman-policy-propagation}.  Substituting the
barred model and residual errors into that theorem proves
\cref{eq-master-dyn-model,eq-master-dyn-geometry,eq-master-dyn-value,eq-master-dyn-policy}.
The exact identity \cref{eq-pomdp-budget-statistical-decomposition}, the
triangle inequality, and the two represented gap certificates, with
\(\overline\Delta_B^{\rm POMDP}=0\) in the fully observed case, prove
\cref{eq-master-dyn-return}.

The shift entries are alternative comparisons between the same behavior and
deployment risks, using the transport consequences of the functional profile
and bounded importance weighting.  The committor residual identities follow
by applying the killed inverse to the represented residual.  The
clean-interface safe and hit entries follow from the reach--avoid
committor identity, condenser capacity, prospective-selector accountability,
and killed-resolvent identity in
\cref{thm-controlled-reach-avoid-dirichlet-drift,thm-prospective-selector-accountability,thm-resolvent-stopwatch-identity}.
The learned-interface Caus, Lift, failure-gate, and capacity entries are
\cref{thm-controlled-noise-corrected-caus-aiming,thm-controlled-noisy-gate-lift-transfer}.
For a POMDP target or failure entry, first bound the exact-belief reference
event by its represented same-record baseline \(U_b\), then add the path or
filter shift from
\cref{thm-controlled-multichannel-path-comparison,thm-controlled-pomdp-filter-transfer}.
If the baseline is merely analytic and uncomputed, its value is one and the
candidate is neutral.  Thus no latent posterior or target gate enters an
affordable minimum for free.  The information entry follows from
\cref{thm-controlled-multichannel-information-bound} only after the
represented output-to-\(\Theta\) decoder and target-sector verification make
the success event an identification event; otherwise its value is one.
The sequential coordinate is the Part~IV online bound applied to the
represented policy-loss class,
\cref{thm-structural-online-regret,cor-structural-online-to-batch}.  Slotwise functional-inequality selection is
the dynamical specialization of \cref{thm-defect-profile-selection}.  Taking
minima of upper bounds for the same quantity proves
\cref{eq-master-dyn-shift,eq-master-dyn-noise,eq-master-dyn-channel-vector,eq-master-dyn-pomdp-information,eq-master-dyn-safe,eq-master-dyn-hit,eq-master-dyn-noisy-capacity,eq-master-dyn-regret}.

All selected geometries, priors, losses, policies, comparisons, and deployment
operations belong to the complete record.  Saturated accounting therefore
charges their derivation and selection costs before any coordinate is used.
Every target-aligned prospective improvement still terminates in the
represented selector/source, Caus-contact, path-baseline, or Lift comparison
of that record.  Analytic latent beliefs, physical channel entropy, and Fano
converses are comparison quantities, not new structural sources.

\end{proof}

\Cref{fig-master-multichannel-certificate-routing} gives the detailed
multichannel routing behind the master certificate.

\begin{figure}[tbp]
\centering
\resizebox{\linewidth}{!}{%
\begin{tikzpicture}[fw diagram,node distance=7mm and 11mm]
  \node[fw source,text width=4.0cm] (in)
    {complete \((n,B,m,H,\delta)\) record\\
     typed channel profiles; path TV/KL; beliefs; information budgets;
     \(\mathbf H_{\rm ctl}\); selector/Caus/Lift records};
  \node[fw provenance,text width=4.0cm,right=of in] (gate)
    {conversion gates\\reward debiasing and shift; same-norm Bellman
     calibration; execution/path comparison; filter stability;
     identification; aiming/contact; Lift comparison};
  \node[fw terminal,text width=4.2cm,right=of gate] (out)
    {coordinatewise outputs\\\(\overline\varepsilon_M,
     \overline\varepsilon_B,\overline\varepsilon_\Phi\); value, policy,
     return; \(U_{\rm hit},U_{\rm safe},U_{\rm POMDP},U_{\rm info}\);
     typed \(\mathbf H_{\rm ctl}\)};
  \draw[fw arrow] (in)--(gate);
  \draw[fw arrow] (gate)--(out);
  \node[fw residual,text width=9.0cm,below=12mm of gate] (rule)
    {same law and unit \(\Rightarrow\) minimum; missing raw error
     \(=+\infty\); missing probability converter or baseline \(=1\);
     every target-aligned improvement retains complete ancestry};
  \draw[fw dashed arrow] (gate)--(rule);
\end{tikzpicture}%
}
\caption{The multichannel master certificate routes every raw coordinate
through a same-law, same-norm converter before combination.  Model, Bellman,
value, return, hit, safety, POMDP, and information bounds take minima only
within one output unit and one complete record; the entropy vector remains
typed.}
\label{fig-master-multichannel-certificate-routing}
\end{figure}

\begin{figure}[tbp]
\centering
\begin{tikzpicture}[fw diagram,node distance=6mm and 8mm]
  \node[fw source,text width=2.6cm] (data) {episode, mixing, or\\sequential record};
  \node[fw geom,text width=2.6cm,right=of data] (geom)
    {action-safe geometry\\and full profile};
  \node[fw provenance,text width=2.5cm,above right=5mm and 9mm of geom] (model)
    {model and\\Bellman bounds};
  \node[fw use,text width=2.5cm,right=12mm of geom] (policy)
    {value and\\policy loss};
  \node[fw terminal,text width=2.5cm,below right=5mm and 9mm of geom] (access)
    {shift, safety,\\and target access};
  \node[fw box,text width=2.5cm,right=15mm of policy] (out)
    {return, arrival,\\and regret};
  \draw[fw arrow] (data) -- (geom);
  \draw[fw arrow] (geom) -- (model);
  \draw[fw arrow] (geom) -- (policy);
  \draw[fw arrow] (geom) -- (access);
  \draw[fw arrow] (model) -- (out);
  \draw[fw arrow] (policy) -- (out);
  \draw[fw arrow] (access) -- (out);
\end{tikzpicture}
\caption{The learned-control certificate propagates one completely charged
trajectory record through action-safe geometry.  Generalization, functional
inequalities, policy loss, safety, and target access remain typed until a
theorem converts them to return, arrival, or regret units.}
\label{fig-master-dynamical-learning-certificate}
\end{figure}

\subsection{Continuous-System and Discretization Transfer}
\label{subsec-controlled-continuum-transfer}

\begin{corollary}[Represented continuum and discretization transfer]
\label{thm-controlled-continuum-lift-transfer}

Let \(L\) be a generator on a declared state space with bounded resolvent at
\(\lambda>0\), let \(\widehat L\) be a represented finite or clocked
generator, and let \(U\) be a represented lift of finite observables into the
declared space.  Suppose the common represented domain satisfies
\begin{equation}
\label{eq-controlled-generator-intertwining-defect}
E_U=LU-U\widehat L
\end{equation}
and the target and failure gates are transported by \(U\).  For first-arrival
claims, \(L\) and \(\widehat L\) are the corresponding killed or absorbing
generators.  Then the exact
resolvent identity is
\begin{equation}
\label{eq-controlled-lift-resolvent-identity}
(\lambda I-L)^{-1}U
-U(\lambda I-\widehat L)^{-1}
=(\lambda I-L)^{-1}E_U(\lambda I-\widehat L)^{-1}.
\end{equation}
Consequently,
\begin{equation}
\label{eq-controlled-lift-resolvent-error}
\left\|(\lambda I-L)^{-1}U
-U(\lambda I-\widehat L)^{-1}\right\|
\le
\|(\lambda I-L)^{-1}\|\,\|E_U\|\,
\|(\lambda I-\widehat L)^{-1}\|.
\end{equation}

If \(P_t\) and \(\widehat P_t\) are represented semigroups for which the
Duhamel integral is defined, then
\begin{equation}
\label{eq-controlled-lift-semigroup-identity}
P_tU-U\widehat P_t
=\int_0^tP_{t-s}E_U\widehat P_s\,ds.
\end{equation}
Every discounted-arrival, mixing, or finite-time certificate for
\(\widehat L\) transfers through these identities with the displayed norm or
form-comparison loss and the represented target-gate comparison.  The lift,
discretization, generator evaluation, error proof, and precision are charged
at their complete saturated cost.  Without this represented transfer record,
the finite certificate makes no claim about \(L\).

This corollary covers represented discretizations of deterministic flows,
stochastic kernels, controlled diffusions, and other continuous models
whenever their finite lift and defect estimate satisfy the displayed
hypotheses.

\end{corollary}

\begin{proof}

From \cref{eq-controlled-generator-intertwining-defect},
\[
(\lambda I-L)U
=U(\lambda I-\widehat L)-E_U.
\]
Multiply on the left and right by the two resolvents and rearrange to obtain
\cref{eq-controlled-lift-resolvent-identity}; submultiplicativity proves
\cref{eq-controlled-lift-resolvent-error}.  Differentiating
\(P_{t-s}U\widehat P_s\) in \(s\) and integrating from zero to \(t\) proves
\cref{eq-controlled-lift-semigroup-identity}.  Target-function transfer then
uses the represented gate comparison.  Cost and provenance are those of a
valid Lift in
\cref{thm-exhaustive-lift-normal-form,cor-lift-creates-nothing}; the
form-comparison alternative is
\cref{lem-declared-geom-form-error,prop-geom-form-resolvent-comparison}.
The block form and its exact-intertwining specialization are
\cref{thm-controlled-valid-lift-block-resolvent}; the present identity is the
two-carrier defect comparison used when the declared and represented
generators live on different approximation spaces.

\end{proof}

\begin{figure}[tbp]
\centering
\resizebox{\linewidth}{!}{%
\begin{tikzpicture}[fw diagram,node distance=8mm and 10mm]
  \node[fw source,text width=2.5cm] (data) {represented\\transition data};
  \node[fw provenance,text width=2.5cm,right=of data] (learn)
    {learned geometry\\and value};
  \node[fw use,text width=2.5cm,right=of learn] (policy)
    {represented policy\\and selector};
  \node[fw geom,text width=2.6cm,right=of policy] (struct)
    {existing saturated\\source cell};
  \node[fw terminal,text width=2.5cm,right=of struct] (bound)
    {structural and\\statistical minimum};
  \draw[fw arrow] (data) -- (learn);
  \draw[fw arrow] (learn) -- (policy);
  \draw[fw arrow] (policy) -- (struct);
  \draw[fw arrow] (struct) -- (bound);
\end{tikzpicture}
}
\caption{A learned policy remains one complete represented computation.
Geometry selection, value fitting, policy construction, and deployment are
charged before the valid structural and statistical bounds are combined.}
\label{fig-learned-policy-accountability}
\end{figure}

\begin{figure}[tbp]
\centering
\resizebox{\linewidth}{!}{%
\begin{tikzpicture}[fw diagram,node distance=8mm and 10mm]
  \node[fw box,text width=2.6cm] (cont) {declared continuous\\or stochastic system};
  \node[fw source,text width=2.5cm,right=of cont] (lift)
    {represented finite\\or clocked Lift};
  \node[fw provenance,text width=2.6cm,right=of lift] (defect)
    {generator defect\\\(E_U=LU-U\widehat L\)};
  \node[fw geom,text width=2.6cm,right=of defect] (finite)
    {finite dynamical\\certificate};
  \node[fw terminal,text width=2.5cm,right=of finite] (transfer)
    {transferred bound\\with explicit loss};
  \draw[fw arrow] (cont) -- (lift);
  \draw[fw arrow] (lift) -- (defect);
  \draw[fw arrow] (defect) -- (finite);
  \draw[fw arrow] (finite) -- (transfer);
\end{tikzpicture}
}
\caption{Continuous-system conclusions pass through a represented lift and an
explicit intertwining defect.  The finite certificate is transferred only
with its proved comparison loss.}
\label{fig-controlled-continuum-transfer}
\end{figure}

\section{Certified Neural Implementation of the Saturated Optimum}
\label{sec-certified-neural-implementation}

The saturated active and structural-Bayesian profiles are analytic optima over
complete represented records.  A trained network supplies one such record.
This section gives a finite-sample certificate for the distance between that
record and the saturated optimum.  The certificate combines a held-out lower
bound for the deployed record with a Bellman upper bound for the complete
saturated class.  Its optimization component is then specialized to a
represented weighted Adam update.

\subsection{Implementation records and error coordinates}
\label{subsec-neural-implementation-records}

\begin{definition}[Complete neural implementation record]
\label{def-complete-neural-implementation-record}

Fix an instance \(I\), horizon \(H\), and saturated active budget \(B\).  A
\emph{complete neural implementation record} consists of:
\begin{enumerate}[label=\textup{(\roman*)}]
\item a complete active or structural-Bayesian record in the sense of
      \cref{def-complete-master-active-sampling-record,def-complete-structural-bayesian-record};
\item represented neural evaluators for every learned score, value, barrier,
      belief, observation, and transition kernel used by that record;
\item the architecture, parameter precision, initialization, training sample,
      loss, optimizer transcript, stopping rule, and selected terminal weights;
\item a represented deployment compiler producing the pre-hit selector and
      every state update from those terminal weights; and
\item represented evaluation and upper-certificate records of the types
      defined below.
\end{enumerate}
The complete cost is
\begin{equation}
\label{eq-complete-neural-implementation-cost}
B_{\rm impl}
=B_{\rm active}\oplus B_{\rm arch}\oplus B_{\rm par}
 \oplus B_{\rm train}\oplus B_{\rm optimize}\oplus B_{\rm precision}
 \oplus B_{\rm deploy}\oplus B_{\rm eval}\oplus B_{\rm upper}.
\end{equation}
Every selected feature, audit cover, random seed, gradient, moment estimate,
preconditioner, likelihood, normalization, and comparison used by the record
is charged in the corresponding slot.  Write
\(\mathscr N_{I,H}(B;B_{\rm impl})\) for the complete neural implementation
records of total cost at most \(B_{\rm impl}\) whose deployed active component
has cost at most \(B\).  Let
\begin{align}
\label{eq-neural-class-arrival-optimum}
J_{\mathcal N}^{*}(B_{\rm impl})
&=\sup_{R_\theta\in\mathscr N_{I,H}(B;B_{\rm impl})}A_{R_\theta}^{H},\\
\label{eq-saturated-arrival-optimum-implementation}
J_{\rm sat}^{*}(B)
&=\operatorname{ActArr}_{I,H}^{\rm sat}(B).
\end{align}
For a deployed terminal parameter \(\widehat\theta\), define the analytic
coordinates
\begin{align}
\label{eq-neural-representation-error}
\varepsilon_{\rm rep}
&=J_{\rm sat}^{*}(B)-J_{\mathcal N}^{*}(B_{\rm impl}),\\
\label{eq-neural-training-error}
\varepsilon_{\rm train}
&=J_{\mathcal N}^{*}(B_{\rm impl})-A_{R_{\widehat\theta}}^{H},\\
\label{eq-neural-total-implementation-error}
\varepsilon_{\rm impl}
&=J_{\rm sat}^{*}(B)-A_{R_{\widehat\theta}}^{H}
=\varepsilon_{\rm rep}+\varepsilon_{\rm train}.
\end{align}
The budget comparison is made at a common class-preserving ceiling containing
the operational and neural records.  Thus both suprema refer to the same
primitive presentation and target event.

\end{definition}

\subsection{Parameter-indexed approximation and underfitting}
\label{subsec-parameter-indexed-underfitting}

\begin{definition}[Parameter-indexed saturated neural profile]
\label{def-parameter-indexed-saturated-neural-profile}

Let \((\mathcal N_p)_{p\ge0}\) be a represented architecture sequence in
which every record in \(\mathcal N_p\) has at most \(p\) trainable scalar
parameters at the declared precision.  Assume a represented padding map
\(\mathcal N_p\hookrightarrow\mathcal N_{p+1}\) that preserves the deployed
record.  At a common active ceiling \(B\) and implementation ceiling
\(B_{\rm impl}(p)\), define
\begin{align}
\label{eq-parameter-indexed-neural-optimum}
J_p^*(B)
&=\sup\left\{
A_{R_\theta}^{H}:
R_\theta\in
\mathscr N_{I,H}(B;B_{\rm impl}(p))\cap\mathcal N_p
\right\},\\
\label{eq-parameter-indexed-representation-error}
\varepsilon_{\rm rep}(p;B)
&=J_{\rm sat}^*(B)-J_p^*(B),\\
\label{eq-minimum-parameter-count}
p_{\min}(\epsilon;B)
&=\inf\left\{
p\in\mathbb N_0:
\varepsilon_{\rm rep}(p;B)\le\epsilon
\right\}.
\end{align}
The padding map makes \(J_p^*\) nondecreasing and
\(\varepsilon_{\rm rep}(p;B)\) nonincreasing.  The value
\(p_{\min}=+\infty\) is used when the set is empty.

For every \(p\), the complete parameter-explicit ledger is
\begin{equation}
\label{eq-parameter-explicit-implementation-ledger}
\mathfrak E_p
=\left(
\varepsilon_{\rm rep}(p),\
\varepsilon_{\rm train}(p),\
\{\Gamma_{t,p}^{(a)}\}_{t,a},\
\{\mathcal L_{p}^{(a)},J_{p}^{(a)}\}_{a},\
\Delta_{{\rm fil},p}^{H},\
\varepsilon_{{\rm num},p},\
\varepsilon_{{\rm opt},p},\
U_{{\rm Bell},p},\
L_{{\rm ev},p}
\right).
\end{equation}
Every coordinate is evaluated for the same architecture, precision, training
record, deployed selector, and audit law.  Dependence on the data size,
noise parameters, horizon, confidence allocation, and complete budget is
retained in the corresponding coordinate.

\end{definition}

\begin{theorem}[Certified underfitting and optimization intervals]
\label{thm-certified-underfitting-optimization-intervals}

Fix \(p\).  Suppose simultaneous represented certificates give
\begin{align}
\label{eq-saturated-optimum-confidence-interval}
L_{\rm sat}\le J_{\rm sat}^*(B)\le U_{\rm sat},\\
\label{eq-p-class-optimum-confidence-interval}
L_p\le J_p^*(B)\le U_p,\\
\label{eq-trained-run-confidence-interval}
L_{\rm run}\le A_{R_{\widehat\theta_p}}^H\le U_{\rm run}.
\end{align}
Then
\begin{align}
\label{eq-certified-underfitting-interval}
\max\{0,L_{\rm sat}-U_p\}
&\le\varepsilon_{\rm rep}(p;B)
\le U_{\rm sat}-L_p,\\
\label{eq-certified-training-error-interval}
\max\{0,L_p-U_{\rm run}\}
&\le\varepsilon_{\rm train}(p;B)
\le U_p-L_{\rm run}.
\end{align}
Consequently
\begin{align}
\label{eq-certified-sufficient-parameter-count}
p_{\rm suff}(\epsilon)
&=\inf\{p:U_{\rm sat}-L_p\le\epsilon\},\\
\label{eq-certified-insufficient-parameter-count}
p_{\rm fail}(\epsilon)
&=1+\sup\{p:L_{\rm sat}-U_p>\epsilon\}
\end{align}
satisfy
\begin{equation}
\label{eq-certified-minimum-parameter-bracket}
p_{\rm fail}(\epsilon)
\le p_{\min}(\epsilon;B)
\le p_{\rm suff}(\epsilon),
\end{equation}
with the empty-set conventions \(p_{\rm fail}=0\) and
\(p_{\rm suff}=+\infty\).

A represented solver compiled into \(\mathcal N_p\) supplies \(L_p\) and
therefore an upper certificate for \(p_{\min}\).  A Bellman upper certificate
restricted to \(\mathcal N_p\) supplies \(U_p\); together with a lower
certificate \(L_{\rm sat}\) for the unrestricted saturated optimum, it
supplies a lower certificate for \(p_{\min}\).

\end{theorem}

\begin{proof}

Subtract the upper and lower endpoints in
\cref{eq-saturated-optimum-confidence-interval,eq-p-class-optimum-confidence-interval}
to obtain \cref{eq-certified-underfitting-interval}.  Subtract the endpoints
in
\cref{eq-p-class-optimum-confidence-interval,eq-trained-run-confidence-interval}
to obtain \cref{eq-certified-training-error-interval}.  Nonnegativity gives
the clipped lower endpoints.

If \(U_{\rm sat}-L_p\le\epsilon\), then
\(\varepsilon_{\rm rep}(p;B)\le\epsilon\), proving the upper bracket.  If
\(L_{\rm sat}-U_p>\epsilon\), then
\(\varepsilon_{\rm rep}(p;B)>\epsilon\).  Monotonicity under the padding maps
excludes every smaller parameter count, proving the lower bracket.

\end{proof}

\subsection{Canonical finite-precision attention--belief architecture}
\label{subsec-canonical-attention-belief-architecture}

\begin{definition}[Canonical certified attention--belief architecture]
\label{def-canonical-certified-attention-belief-architecture}

Fix a finite horizon \(H\).  A canonical certified architecture
\(\mathcal C_{\Theta}^{\rm can}\) contains the following represented modules.
\begin{enumerate}[label=\textup{(\roman*)}]
\item A tokenizer maps the public presentation, current native state,
      surviving pre-hit history, resource record, and represented legal
      actions to \(N_0\) width-\(d_0\) tokens.  Illegal actions receive a
      represented mask.
\item One recurrent belief token stores a width-\(d_0\) compact latent state.
      Its update factors through the current pre-hit record.
\item Layer \(\ell\) is a pre-normalized transformer block with width
      \(d_\ell\), \(H_\ell\) attention heads, key width \(d_{k,\ell}\), value
      width \(d_{v,\ell}\), feed-forward width \(f_\ell\), temperature
      \(\tau_{\ell h}>0\), residual addition, and two affine
      normalizations.  The canonical residual stack has
      \(d_\ell=d_0\); a represented width-changing projection is counted in
      \(p_{\rm tok}\) when present.
\item The terminal representation feeds five represented heads:
      a masked policy score, a value function, a Bellman barrier, a centered
      geometry feature map \(\psi_{p_{\rm geom}}\), and a quotient label or
      quotient-logit map \(q\).
\item A finite-precision deployment compiler produces the pre-hit selector,
      belief update, transition record, and all audit evaluators from the
      stored parameters.
\end{enumerate}
The quotient head proposes a represented label.  Its exact or approximate
dynamical status is determined by the fiber-row audit of
\cref{def-fiber-transition-diameter}.

For dense affine maps with biases, the trainable parameter count is
\begin{align}
\label{eq-canonical-attention-parameter-count}
p_{\rm can}
={}&p_{\rm tok}+p_{\rm bel}
+\sum_{\ell=1}^{L}
\left(p_{{\rm att},\ell}+p_{{\rm ff},\ell}+4d_\ell\right)
+p_{\rm pol}+p_{\rm val}+p_{\rm bar}
+p_{\rm geom}+p_{\rm quot},\\
\label{eq-canonical-attention-block-parameter-count}
p_{{\rm att},\ell}
={}&H_\ell\!\left(
2d_\ell d_{k,\ell}+d_\ell d_{v,\ell}
+2d_{k,\ell}+d_{v,\ell}\right)
+H_\ell d_{v,\ell}d_\ell+d_\ell,\\
\label{eq-canonical-feedforward-parameter-count}
p_{{\rm ff},\ell}
={}&2d_\ell f_\ell+f_\ell+d_\ell .
\end{align}
The \(4d_\ell\) term stores the scale and shift parameters of the two
normalizations.  Tied parameters are counted once and their reuse is stored
in the architecture record.  Every scalar is stored at the declared
precision.

If \(E_\ell\) is the number of unmasked query--key pairs at layer \(\ell\),
one forward deployment has arithmetic cost bounded by
\begin{align}
\label{eq-canonical-attention-forward-cost}
\operatorname{Cost}_{\rm fwd}
\le C_{\rm op}\Biggl\{&
\operatorname{Cost}_{\rm tok}+\operatorname{Cost}_{\rm bel}
+\operatorname{Cost}_{\rm heads}
\notag\\
&+\sum_{\ell=1}^{L}
H_\ell\!\left[
N_\ell d_\ell(2d_{k,\ell}+d_{v,\ell})
+E_\ell(d_{k,\ell}+d_{v,\ell})
\right]
\notag\\
&+\sum_{\ell=1}^{L}
N_\ell d_\ell\left(H_\ell d_{v,\ell}+2f_\ell+1\right)
\Biggr\}.
\end{align}
where \(C_{\rm op}\) is the represented finite-precision arithmetic
conversion.  Parameter storage, activations, recurrent state, optimizer
moments, training, audit samples, and deployment are charged separately in
\cref{eq-complete-neural-implementation-cost}.

\end{definition}

\begin{theorem}[Canonical architecture implementation and compilation]
\label{thm-canonical-attention-architecture-implementation}

Every module in
\cref{def-canonical-certified-attention-belief-architecture} is a
finite composition of represented affine maps, comparisons, masks,
normalizations, pointwise maps, and finite-precision exponential
normalization.  Hence it defines a complete neural implementation record at
the cost in
\cref{eq-canonical-attention-parameter-count,eq-canonical-attention-forward-cost}.
Its deployed attention rows are complete attention records in
\cref{def-complete-attention-record}.

Conversely, let a complete bounded-branch clocked selector have \(M_t\)
represented legal actions at time \(t\).  At the class-preserving compilation
ceiling, a canonical architecture with \(N_0\ge\max_t M_t+1\) represents the
same pre-hit path law by tokenizing the legal actions and the recurrent
state, exposing the selector's already charged auxiliary seed, and assigning
the selected action the unique maximal finite-precision score.  The
architecture preserves the selector's source ancestry, latent code, mask,
tie-breaking rule, and complete cost.

\end{theorem}

\begin{proof}

The first assertion follows directly from the represented constructor list
and the complete cost definition.  The parameter formulas count the
\(Q,K,V,O\), feed-forward, normalization, tokenizer, recurrent, and output
maps once.  The operation bound counts each unmasked dot product, value
transport, and feed-forward multiplication.

For the converse, use the pathwise hard-attention compilation in
\cref{cor-attention-active-sampling-compilation}: for each represented seed,
the inverse-transform comparisons select one legal action.  Store the
surviving history and seed in the belief token, mask illegal actions, assign
the selected action a unique maximal score, and apply the original
tie-breaking convention.  Averaging over the same seed reproduces the
original kernel.  Constructor no-creation retains every source ancestor and
the compilation cost.

\end{proof}

\begin{figure}[tbp]
\centering
\resizebox{.98\linewidth}{!}{%
\begin{tikzpicture}[fw diagram,node distance=7mm and 9mm]
  \node[fw source,text width=2.4cm] (pub)
    {public presentation,\\state and history};
  \node[fw geom,text width=2.3cm,right=of pub] (tok)
    {tokens and\\belief state};
  \node[fw use,text width=2.4cm,right=of tok] (att)
    {masked attention\\and residual blocks};
  \node[fw terminal,text width=2.1cm,above right=6mm and 9mm of att] (pol)
    {policy and\\value heads};
  \node[fw provenance,text width=2.1cm,right=of att] (bar)
    {Bellman\\barrier};
  \node[fw geom,text width=2.1cm,below right=6mm and 9mm of att] (geo)
    {geometry and\\quotient heads};
  \node[fw box,text width=2.5cm,right=18mm of bar] (cert)
    {one charged empirical\\proof record};
  \draw[fw arrow] (pub)--(tok);
  \draw[fw arrow] (tok)--(att);
  \draw[fw arrow] (att)--(pol);
  \draw[fw arrow] (att)--(bar);
  \draw[fw arrow] (att)--(geo);
  \draw[fw arrow] (pol)--(cert);
  \draw[fw arrow] (bar)--(cert);
  \draw[fw arrow] (geo)--(cert);
\end{tikzpicture}%
}
\caption{Canonical certified implementation.  The deployed controller and
the Bellman, geometry, and quotient audits share the same public tokens,
belief state, parameters, transition law, precision, and cost ledger.}
\label{fig-canonical-certified-attention-architecture}
\end{figure}

\begin{proposition}[Held-out arrival lower certificate]
\label{prop-held-out-arrival-lower-certificate}

Condition on the complete training record and terminal weights
\(\widehat\theta\).  Run \(N_{\rm ev}\) independent represented evaluation
episodes under the deployed record and put
\[
\widehat A_{\rm ev}
=\frac1{N_{\rm ev}}\sum_{j=1}^{N_{\rm ev}}
\mathbf1_{\{\tau_j\le H\}}.
\]
For \(0<\delta_{\rm ev}<1\), define
\begin{equation}
\label{eq-held-out-arrival-lower-bound}
L_{\rm ev}
=\left[
\widehat A_{\rm ev}
-\sqrt{\frac{\log(1/\delta_{\rm ev})}{2N_{\rm ev}}}
\right]_+.
\end{equation}
Then, conditionally and hence unconditionally,
\begin{equation}
\label{eq-held-out-arrival-lower-confidence}
\Pr\{L_{\rm ev}\le A_{R_{\widehat\theta}}^{H}\}
\ge1-\delta_{\rm ev}.
\end{equation}
A sequential or dependent evaluation record may replace the square-root term
by a valid same-law martingale or mixing candidate from
\cref{thm-beta-mixing-structural-blocking,thm-martingale-structural-pac-bayes}.

\end{proposition}

\begin{proof}

Conditional on the training record, the evaluation indicators are independent
Bernoulli variables with mean \(A_{R_{\widehat\theta}}^{H}\).  The lower-tail
Hoeffding inequality gives
\[
\Pr\left\{
\widehat A_{\rm ev}-A_{R_{\widehat\theta}}^{H}
>\sqrt{\frac{\log(1/\delta_{\rm ev})}{2N_{\rm ev}}}
\,\middle|\,\widehat\theta
\right\}\le\delta_{\rm ev}.
\]
Clipping at zero preserves the lower bound.  Averaging over the training
record proves \cref{eq-held-out-arrival-lower-confidence}.

\end{proof}

\subsection{Empirical Bellman upper certificates}
\label{subsec-empirical-bellman-upper-certificates}

\begin{definition}[Saturated Bellman audit record]
\label{def-saturated-bellman-audit-record}

Let \(v=(v_t)_{t=0}^{H}\), \(v_H=0\), be a represented bounded Bellman
barrier on the complete surviving history-and-resource carrier.  For each
time \(t\), let \(\mu_t\) be a represented audit probability law on complete
first-step records \(r=(h,a)\).  A \emph{saturated coverage certificate} is
a represented number \(C_t<\infty\) such that every stopped first-step flow
of every \(R\in\mathscr R_{I,H}^{\rm act}(B)\) satisfies
\begin{equation}
\label{eq-saturated-audit-coverage}
f_{R,t}\ll\mu_t,
\qquad
\left\lVert\frac{df_{R,t}}{d\mu_t}\right\rVert_\infty\le C_t.
\end{equation}
Here \(f_{R,t}\) is the subprobability flow in
\cref{def-stopped-occupation-flow-complete-active-record}.  Let
\(r_{t,1},\ldots,r_{t,N_t}\) be an independent audit sample from \(\mu_t\),
independent of the construction of \(v\), and set
\begin{align}
\label{eq-audited-positive-bellman-residual}
d_t(r;v)
&=\left[p_T(h,a)+P_Nv_{t+1}(h,a)-v_t(h)\right]_+,\\
\label{eq-empirical-positive-bellman-residual}
\widehat d_t(v)
&=\frac1{N_t}\sum_{j=1}^{N_t}d_t(r_{t,j};v).
\end{align}
Let \(M_t\) be a represented bound \(0\le d_t(r;v)\le M_t\) on the support
of \(\mu_t\).  The coverage law, its sampler, barrier, transition and contact
evaluators, sample, precision, and comparisons are charged in
\(B_{\rm upper}\).

An exact verification record may instead supply represented numbers
\(\overline d_t(v)\) satisfying
\begin{equation}
\label{eq-exact-bellman-residual-verification}
\sup_{r\in\mathscr A_{I,t}^{\rm sat}(B)}d_t(r;v)
\le\overline d_t(v).
\end{equation}

\end{definition}

\begin{theorem}[Empirical saturated Bellman upper bound]
\label{thm-empirical-saturated-bellman-upper-bound}

Under \cref{def-saturated-bellman-audit-record}, allocate
\(\delta_{{\rm up},t}>0\) with
\(\sum_{t<H}\delta_{{\rm up},t}\le\delta_{\rm up}\), and put
\begin{align}
\label{eq-empirical-bellman-residual-upper-candidate}
u_t^{\rm emp}(v)
&=C_t\left[
\widehat d_t(v)
+M_t\sqrt{\frac{\log(1/\delta_{{\rm up},t})}{2N_t}}
\right],\\
\label{eq-complete-bellman-upper-certificate}
U_{\rm Bell}(v)
&=v_0(h_0)+\sum_{t=0}^{H-1}
\min\{\overline d_t(v),u_t^{\rm emp}(v)\},
\end{align}
where an unavailable candidate has value \(+\infty\).  Then, with probability
at least \(1-\delta_{\rm up}\),
\begin{equation}
\label{eq-empirical-saturated-bellman-upper-confidence}
J_{\rm sat}^{*}(B)
\le U_{\rm Bell}(v).
\end{equation}
The conclusion remains valid when a simultaneous structural Rademacher,
sequential, mixing, or martingale PAC--Bayes radius replaces the Hoeffding
radius, provided it bounds the same \(\mu_t\)-expectation for the selected
barrier and its construction and selection are included in the audit record.

\end{theorem}

\begin{proof}

Hoeffding's inequality and a union bound give, simultaneously for every
\(t<H\),
\begin{equation}
\label{eq-audit-mean-bellman-residual-bound}
\mathbb E_{\mu_t}d_t(r;v)
\le\widehat d_t(v)
+M_t\sqrt{\frac{\log(1/\delta_{{\rm up},t})}{2N_t}}.
\end{equation}
For any complete record \(R\), the coverage certificate gives
\[
\int d_t(r;v)f_{R,t}(dr)
\le C_t\mathbb E_{\mu_t}d_t(r;v)
\le u_t^{\rm emp}(v).
\]
The exact verification candidate gives the alternative bound
\(\int d_t\,df_{R,t}\le\overline d_t(v)\), because \(f_{R,t}\) has mass at
most one.  Integrating the signed Bellman inequality along the stopped flow
and telescoping as in
\cref{thm-bellman-barrier-duality-affordable-arrival} yields
\[
A_R^H
\le v_0(h_0)+\sum_{t<H}
\min\{\overline d_t(v),u_t^{\rm emp}(v)\}.
\]
Taking the supremum over \(R\in\mathscr R_{I,H}^{\rm act}(B)\) proves
\cref{eq-empirical-saturated-bellman-upper-confidence}.

\end{proof}

\begin{corollary}[Primal--dual empirical implementation certificate]
\label{cor-primal-dual-empirical-implementation-certificate}

Let \(L_{\rm ev}\) and \(U_{\rm Bell}(v)\) be constructed from independent
evaluation and upper-audit records with total failure probability
\(\delta=\delta_{\rm ev}+\delta_{\rm up}\).  Then, with probability at least
\(1-\delta\),
\begin{equation}
\label{eq-primal-dual-empirical-implementation-gap}
0\le\varepsilon_{\rm impl}
=J_{\rm sat}^{*}(B)-A_{R_{\widehat\theta}}^{H}
\le U_{\rm Bell}(v)-L_{\rm ev}.
\end{equation}
In particular, this certificate bounds the sum
\(\varepsilon_{\rm rep}+\varepsilon_{\rm train}\) without separately
estimating either coordinate.

If the hypotheses of
\cref{cor-learned-structural-bayesian-control} hold at the same compiled
ceiling, the implementation error also satisfies
\begin{equation}
\label{eq-combined-implementation-gap-candidates}
\varepsilon_{\rm impl}
\le\min\left\{
U_{\rm Bell}(v)-L_{\rm ev},
\mathbb E_{\widehat R}\sum_{t<H}\mathbf1_{\{\tau>t\}}
  (2\Gamma_t+\varepsilon_{{\rm opt},t})+\Delta_{\rm fil}^{H}
\right\}.
\end{equation}

\end{corollary}

\begin{proof}

On the simultaneous event,
\(L_{\rm ev}\le A_{R_{\widehat\theta}}^H\le J_{\rm sat}^{*}(B)\le
U_{\rm Bell}(v)\).  Subtraction proves
\cref{eq-primal-dual-empirical-implementation-gap}.  The second candidate is
\cref{eq-learned-structural-bayes-optimality-loss}; both candidates bound the
same nonnegative quantity on the same event, so their minimum is valid.

\end{proof}

\subsection{Scale-certified empirical closure of attenuating sources}
\label{subsec-scale-certified-empirical-attenuation}

The empirical upper audit becomes an asymptotic certificate when the same
record includes a verified scale law for structural capacity and target
attenuation.  The following definition records that additional datum at the
level of the saturated source cells.  The per-cell factorization, exponent
competition, and complete empirical-to-proof pipeline are displayed in
\cref{fig-scale-certified-source-cell,fig-capacity-attenuation-exponent-competition,fig-empirical-scale-certified-source-closure}.

\begin{definition}[Scale-certified attenuating source audit]
\label{def-scale-certified-attenuating-source-audit}

Fix a family functional \(\mathfrak M_n\), a budget schedule \(B_n\), and a
base size \(n_0\).  A \emph{scale-certified attenuating source audit}
consists of a finite set \(\mathfrak C_{\rm att}\) of existing qualified
source signatures such that, for every \(n\ge n_0\),
\begin{equation}
\label{eq-scale-audit-source-cover}
I_{\mathfrak M,n}^{\rm src,sat}(B_n)
\le
\sum_{c\in\mathfrak C_{\rm att}}
\sup_{R\in\mathcal S_{c,n}(B_n)}V_{c,n}(R).
\end{equation}
Here \(\mathcal S_{c,n}(B_n)\) denotes the restriction of the already-defined
complete record class to source cell \(c\), including all temporal, target,
dynamic, and cost qualification gates.  Thus the audit introduces no new
admissible record class.

For each \(c\) and \(n\), the audit supplies a represented Hilbert carrier
\(\mathcal H_{c,n}\), a coefficient record
\(z_{c,n}(R)\in\mathcal H_{c,n}\), a contraction
\(P_{c,n}:\mathcal H_{c,n}\to\mathcal H_{c,n}\), and a scalar target-transfer
functional \(T_{c,n}:\mathcal H_{c,n}\to\mathbb R\).  It also supplies
nonnegative numbers \(D_{c,n}\), \(a_{c,n}\), and \(e_{c,n}\) satisfying,
simultaneously for every \(R\in\mathcal S_{c,n}(B_n)\),
\begin{align}
\label{eq-scale-audit-factorization}
\lVert z_{c,n}(R)\rVert_{\mathcal H_{c,n}}&\le1,\\
V_{c,n}(R)
&\le
\left|T_{c,n}P_{c,n}z_{c,n}(R)\right|+e_{c,n},\\
\label{eq-scale-audit-capacity-attenuation}
\left\lVert T_{c,n}P_{c,n}\right\rVert_{\mathcal H_{c,n}\to\mathbb R}
&\le \sqrt{D_{c,n}}\,a_{c,n}.
\end{align}
The quantity \(D_{c,n}\) is the certified target-relevant capacity of the
budget-reachable carrier, \(a_{c,n}\) is its target-transfer attenuation,
and \(e_{c,n}\) is the certified factorization, tail, and numerical error.
For an orthogonal materialized band, one may take \(D_{c,n}\) to be its
rank and \(a_{c,n}\) the largest retained transfer multiplier.  For a ridge
band, the effective-dimension and tail certificates of
\cref{def-effective-dimension-extraction-audit,def-learning-affordability}
supply the corresponding continuous values.

At the base size, a simultaneous upper audit supplies
\begin{equation}
\label{eq-scale-audit-base-upper-bounds}
D_{c,n_0}\le\overline D_{c,0},
\qquad
a_{c,n_0}\le\overline a_{c,0},
\qquad
e_{c,n_0}\le\overline e_{c,0}.
\end{equation}
For all \(n\ge n_0\), a represented scale verifier supplies
\begin{align}
\label{eq-scale-audit-capacity-recurrence}
D_{c,n+1}&\le r_{c,n}D_{c,n},\\
\label{eq-scale-audit-attenuation-recurrence}
a_{c,n+1}&\le s_{c,n}a_{c,n},\\
\label{eq-scale-audit-error-envelope}
e_{c,n}&\le\overline e_{c,n}.
\end{align}
The complete audit cost contains the learned factorization, audit samples,
effective-dimension or operator calculation, numerical precision, interval
or symbolic scale verifier, source-cell cover, and every selection rule used
to choose them.

A \emph{proof-grade} audit has deterministic verification of
\cref{eq-scale-audit-source-cover,eq-scale-audit-factorization,eq-scale-audit-capacity-attenuation,eq-scale-audit-base-upper-bounds,eq-scale-audit-capacity-recurrence,eq-scale-audit-attenuation-recurrence,eq-scale-audit-error-envelope}.
An empirical simultaneous bound with failure probability \(\delta_{\rm aud}\)
gives the same conclusions on its declared audit event.  A randomly trained
record becomes proof-grade when its exported inequalities are checked by a
deterministic exact, interval, or symbolic verifier; randomness used to find
the certificate then has no role in its validity.

\end{definition}

\begin{figure}[tbp]
\centering
\resizebox{.98\linewidth}{!}{%
\begin{tikzpicture}[fw diagram,node distance=7mm and 10mm]
  \node[fw source,text width=2.5cm] (cell)
    {qualified saturated\\source cell \(c\)};
  \node[fw use,text width=2.2cm,right=of cell] (coeff)
    {coefficient record\\\(\lVert z\rVert\le1\)};
  \node[fw geom,text width=2.1cm,right=of coeff] (transfer)
    {represented\\transfer \(P\)};
  \node[fw provenance,text width=2.1cm,right=of transfer] (target)
    {target functional\\\(T\)};
  \node[fw terminal,text width=3.0cm,right=of target] (visible)
    {\(V_{c,n}\le\sqrt{D_{c,n}}\,a_{c,n}+e_{c,n}\)};
  \node[fw box,text width=2.25cm,below=8mm of transfer] (capacity)
    {capacity\\\(D_{c,n}\)};
  \node[fw box,text width=2.25cm,below=8mm of target] (attenuation)
    {attenuation\\\(a_{c,n}\)};
  \node[fw residual,text width=2.25cm,below=8mm of visible] (error)
    {certified error\\\(e_{c,n}\)};
  \draw[fw arrow] (cell)--(coeff);
  \draw[fw arrow] (coeff)--(transfer);
  \draw[fw arrow] (transfer)--(target);
  \draw[fw arrow] (target)--(visible);
  \draw[fw arrow] (capacity)--(visible);
  \draw[fw arrow] (attenuation)--(visible);
  \draw[fw arrow] (error)--(visible);
\end{tikzpicture}%
}
\caption{One scale-certified source cell.  Exhaustive source routing fixes the
cell; the represented factorization separates budget-reachable capacity,
target-transfer attenuation, and certified approximation error.}
\label{fig-scale-certified-source-cell}
\end{figure}

\begin{lemma}[Base-size audit with verified scale propagation]
\label{lem-base-audit-verified-scale-propagation}

Under \cref{def-scale-certified-attenuating-source-audit}, for every
\(n\ge n_0\),
\begin{align}
\label{eq-scale-propagated-capacity}
D_{c,n}
&\le
\overline D_{c,0}
\prod_{j=n_0}^{n-1}r_{c,j},\\
\label{eq-scale-propagated-attenuation}
a_{c,n}
&\le
\overline a_{c,0}
\prod_{j=n_0}^{n-1}s_{c,j},\\
\label{eq-scale-propagated-source-cell}
\sup_{R\in\mathcal S_{c,n}(B_n)}V_{c,n}(R)
&\le
\sqrt{\overline D_{c,0}}
\,\overline a_{c,0}
\prod_{j=n_0}^{n-1}\bigl(\sqrt{r_{c,j}}s_{c,j}\bigr)
+\overline e_{c,n}.
\end{align}
Empty products are one.

\end{lemma}

\begin{proof}
Iterating \cref{eq-scale-audit-capacity-recurrence} from the certified base
bound gives \cref{eq-scale-propagated-capacity}; iterating
\cref{eq-scale-audit-attenuation-recurrence} gives
\cref{eq-scale-propagated-attenuation}.  For a qualified record \(R\),
\cref{eq-scale-audit-factorization,eq-scale-audit-capacity-attenuation}
and \(\lVert z_{c,n}(R)\rVert\le1\) give
\[
V_{c,n}(R)
\le\sqrt{D_{c,n}}\,a_{c,n}+e_{c,n}.
\]
Substitute the two propagated bounds and
\cref{eq-scale-audit-error-envelope}, then take the supremum over the source
cell.
\end{proof}

\begin{theorem}[Empirical-to-asymptotic saturated-source closure]
\label{thm-empirical-asymptotic-saturated-source-closure}

Assume a scale-certified attenuating source audit.  Suppose that, for each
\(c\in\mathfrak C_{\rm att}\), there are constants
\(\beta_c\ge0\), \(\gamma_c>0\), and a positive subexponential function
\(p_c\), normalized by \(p_c(n_0)=1\), such that
\begin{equation}
\label{eq-scale-audit-exponent-rates}
r_{c,n}
\le
e^{\beta_c}\frac{p_c(n+1)}{p_c(n)},
\qquad
s_{c,n}\le e^{-\gamma_c}
\qquad(n\ge n_0).
\end{equation}
Then, on the simultaneous audit event,
\begin{equation}
\label{eq-empirical-asymptotic-source-profile-bound}
I_{\mathfrak M,n}^{\rm src,sat}(B_n)
\le
\sum_{c\in\mathfrak C_{\rm att}}
\left[
\sqrt{\overline D_{c,0}}\,\overline a_{c,0}
\sqrt{p_c(n)}
e^{-(\gamma_c-\beta_c/2)(n-n_0)}
+\overline e_{c,n}
\right].
\end{equation}
If \(\gamma_c>\beta_c/2\) for every cell and the finite sum of error
envelopes is negligible, the saturated prospective source profile is
negligible.  Prospective selector accountability then gives negligible
arrival whenever the declared neutral baseline is negligible.

For a proof-grade audit, the conclusion is a deterministic
computer-assisted theorem.  For a statistical audit, it holds with the
simultaneous confidence carried by the audit record.

\end{theorem}

\begin{proof}
The capacity recurrence in \cref{eq-scale-audit-exponent-rates} telescopes:
\[
\prod_{j=n_0}^{n-1}r_{c,j}
\le e^{\beta_c(n-n_0)}p_c(n),
\qquad
\prod_{j=n_0}^{n-1}s_{c,j}
\le e^{-\gamma_c(n-n_0)}.
\]
Insert these inequalities in
\cref{eq-scale-propagated-source-cell}, sum the source-cell bounds using
\cref{eq-scale-audit-source-cover}, and obtain
\cref{eq-empirical-asymptotic-source-profile-bound}.  Subexponential
\(p_c\) and \(\gamma_c-\beta_c/2>0\) make every leading term exponentially
small.  The negligibility and arrival conclusions follow from closure under
a fixed finite sum and
\cref{thm-prospective-selector-structural-charge,thm-prospective-selector-accountability}.
\end{proof}

\begin{figure}[tbp]
\centering
\begin{tikzpicture}[x=1.15cm,y=.78cm,>=Latex]
  \draw[->,semithick] (0,0)--(8.2,0) node[right] {size \(n\)};
  \draw[->,semithick] (0,-3.0)--(0,3.0)
    node[above,text width=2.5cm,align=center] {logarithmic\\contribution};
  \draw[very thick,orange!85!black] (.5,.25)--(7.5,2.35)
    node[right,black] {\(\tfrac12\log D\sim\tfrac{\beta_c}{2}n\)};
  \draw[very thick,teal!75!black] (.5,-.25)--(7.5,-2.65)
    node[right,black] {\(\log a\sim-\gamma_cn\)};
  \draw[very thick,red!75!black] (.5,0)--(7.5,-1.35)
    node[right,black] {net rate \(-(\gamma_c-\beta_c/2)n\)};
  \draw[densely dashed,gray] (4,-3)--(4,3);
  \node[fw box,text width=3.1cm,align=center] at (4.1,2.1)
    {\(\gamma_c>\beta_c/2\)\\attenuation dominates};
\end{tikzpicture}
\caption{Verified scale competition.  Capacity enters through
\(\sqrt{D_{c,n}}\), whereas target attenuation enters linearly through
\(a_{c,n}\).  The strict exponent gap
\(\gamma_c-\beta_c/2>0\) makes their product exponentially small.}
\label{fig-capacity-attenuation-exponent-competition}
\end{figure}

\begin{corollary}[Polynomial capacity loses to exponential attenuation]
\label{cor-polynomial-capacity-exponential-attenuation}

In \cref{thm-empirical-asymptotic-saturated-source-closure}, suppose
\(B_n\) is polynomial and the proof-grade audit certifies, for fixed
constants \(C_c\) and \(\gamma_c>0\),
\[
D_{c,n}\le(n+B_n)^{C_c},
\qquad
a_{c,n}\le e^{-\gamma_c n},
\qquad
\sum_c\overline e_{c,n}=\operatorname{negl}(n).
\]
Then
\begin{equation}
\label{eq-polynomial-capacity-exponential-attenuation}
I_{\mathfrak M,n}^{\rm src,sat}(B_n)
\le
\sum_c(n+B_n)^{C_c/2}e^{-\gamma_c n}
+\operatorname{negl}(n)
=\operatorname{negl}(n).
\end{equation}
The conclusion applies at a polynomial Levin ceiling because
\cref{lem-description-length-charged-resource,eq-code-reading-cost-domination}
charge the code, generated parameters, evaluation, storage, and deployment
inside \(B_n\).  The capacity inequality remains an independently verified
field of the audit; description length alone supplies no replacement for
that inequality.

\end{corollary}

\begin{proof}
Apply \cref{eq-scale-audit-factorization,eq-scale-audit-capacity-attenuation}
in every source cell and sum.  A polynomial factor times
\(e^{-\gamma_c n}\) is exponentially small.  The Levin statement follows
from the cited complete-cost domination.
\end{proof}

\begin{corollary}[Finite empirical discovery with deterministic universal
closure]
\label{cor-empirical-discovery-deterministic-universal-closure}

Fix the public presentation, budget schedule, and existing saturated source
cells.  A finite computation discharges the saturated-source estimate in
\cref{thm-empirical-asymptotic-saturated-source-closure} when its exported
record passes the following checks.
\begin{enumerate}[label=\textup{(\roman*)}]
\item The framework source-exhaustion theorem verifies the cover
      \cref{eq-scale-audit-source-cover}; sampled coverage is not used in
      place of this check.
\item Exact, rational-interval, or symbolic arithmetic verifies the
      factorization and norm inequalities
      \cref{eq-scale-audit-factorization,eq-scale-audit-capacity-attenuation}
      for the complete qualified cell, including the displayed error
      envelope.
\item A simultaneous empirical audit supplies candidate upper bounds at the
      base size, and an independent deterministic verifier proves the final
      base inequalities in \cref{eq-scale-audit-base-upper-bounds}.  The
      candidate may use the effective-dimension audit of
      \cref{def-effective-dimension-extraction-audit}, the Bellman upper audit
      of \cref{thm-empirical-saturated-bellman-upper-bound}, the simultaneous
      finite-network radius of
      \cref{cor-finite-network-simultaneous-audit-radius}, or the
      source-resolved learning radii of
      \cref{thm-parameter-source-resolved-empirical-target-gain}.
\item The same verifier proves the capacity, attenuation, and error scale
      laws
      \cref{eq-scale-audit-capacity-recurrence,eq-scale-audit-attenuation-recurrence,eq-scale-audit-error-envelope}
      and the rate comparison \cref{eq-scale-audit-exponent-rates}.
\item The code, generated parameters, precision, audit, verifier,
      evaluation, storage, and deployment all fit the common charged ceiling.
\item The strict exponent inequalities
      \(\gamma_c>\beta_c/2\) and the claimed negligible error envelopes are
      verified for every cell in the finite cover.
\item Any finite interval of sizes not covered by the propagated estimate is
      checked directly by the same outward-rounded upper audit.
\end{enumerate}
Then \cref{eq-empirical-asymptotic-source-profile-bound} is a deterministic
upper bound on the complete saturated prospective source profile.  The
training sample, optimizer, and random initialization may be used to find the
record, but do not occur in the validity of the resulting bound.  If item
\textup{(iii)} or item \textup{(iv)} is retained only as a confidence
statement, the conclusion has that confidence and is not an unconditional
computer-assisted theorem.

\end{corollary}

\begin{proof}
Items \textup{(i)}--\textup{(iv)} are exactly the hypotheses of
\cref{def-scale-certified-attenuating-source-audit}; item \textup{(v)}
places their construction in the declared derivability closure.  Item
\textup{(vi)} invokes
\cref{thm-empirical-asymptotic-saturated-source-closure}, and item
\textup{(vii)} joins the finite verified prefix to its propagated tail.
Deterministic verification makes the exported inequalities independent of
the random search that found them.
\end{proof}

\subsubsection{Proof-producing verification of structural scaling laws}
\label{subsubsec-proof-producing-structural-scaling-laws}

Finite-size experiments can discover a useful scaling law, but the
asymptotic assertion is supplied by a verified structural recurrence.  The
next definition records the recurrence, the complete cost of traversing it,
and the source-cell cover to which it applies.

\begin{definition}[Charged structural scaling-recurrence certificate]
\label{def-charged-structural-scaling-recurrence-certificate}

Fix a qualified saturated source cell \(c\), a size \(n\), and its declared
budget \(B_n\).  A \emph{charged structural scaling-recurrence certificate}
consists of the following represented data.
\begin{enumerate}[label=\textup{(\roman*)}]
\item A represented partially ordered state space \(\mathcal U_{c,n}\), an
      initial enclosure \(U_{c,n,0}\subseteq\mathcal U_{c,n}\), and a
      finite collection \(\mathfrak F_{c,n}\) of monotone interval
      transition maps.  The state stores only already-defined structural
      coordinates, including capacity, attenuation order, approximation
      defect, available rows or observations, provenance size, memory, and
      accumulated cost when those coordinates are present.
\item A represented recurrence cover: every record in
      \(\mathcal S_{c,n}(B_n)\) determines a finite path
      \begin{equation}
      \label{eq-structural-scaling-recurrence-path}
      u_0\in U_{c,n,0},\qquad
      u_{j+1}\in F_{c,n,j}(u_j),\quad
      F_{c,n,j}\in\mathfrak F_{c,n}.
      \end{equation}
      The cover is proved from the framework source-exhaustion and
      constructor-provenance theorems; empirical path sampling is not a
      substitute for this item.
\item A nonnegative progress functional \(\pi_{c,n}\), a nonnegative
      target-transfer envelope \(A_{c,n}\), and a nonnegative incremental
      complete-cost functional \(p_{c,n}\) satisfying along every covered
      path
      \begin{align}
      \label{eq-structural-scaling-progress-transfer}
      V_{c,n}(R)&\le A_{c,n}(u_J)+e_{c,n},\\
      \label{eq-structural-scaling-incremental-cost}
      \operatorname{Cost}(R)&\ge
      C_{c,n,0}+\sum_{j=0}^{J-1}
      p_{c,n}(u_j,u_{j+1}).
      \end{align}
      The incremental cost includes the recurrence evaluator, locator or
      collision rule, acquired data, generated records, arithmetic,
      precision, storage, history, verification, and deployment.
\item A finite verified prefix \(n_0\le n<N_0\), together with represented
      tail enclosures \(\overline U_{c,n,j}\), such that a symbolic or
      outward-rounded interval verifier proves, for every \(n\ge N_0\),
      \begin{equation}
      \label{eq-structural-scaling-inductive-enclosure}
      U_{c,n,0}\subseteq\overline U_{c,n,0},\qquad
      F_{c,n,j}(\overline U_{c,n,j})
      \subseteq\overline U_{c,n,j+1}.
      \end{equation}
\end{enumerate}
The training runs, learned recurrence proposal, finite prefix, verifier,
interval precision, and all proof objects are included in the common cost
record.  The certificate is \emph{proof-grade} when the recurrence cover,
all inclusions, and all numerical inequalities are checked
deterministically.

For a target progress level \(q\), define its charged crossing price by
\begin{equation}
\label{eq-structural-scaling-crossing-price}
\operatorname{Cross}_{c,n}(q)
=
\inf\left\{
C_{c,n,0}+\sum_{j=0}^{J-1}p_{c,n}(u_j,u_{j+1}):
\begin{array}{l}
(u_0,\ldots,u_J)\text{ satisfies }
\cref{eq-structural-scaling-recurrence-path},\\[-1mm]
\pi_{c,n}(u_J)\ge q
\end{array}
\right\},
\end{equation}
with infimum \(+\infty\) when no covered path reaches \(q\).

\end{definition}

\begin{theorem}[Verified recurrence propagation and crossing exclusion]
\label{thm-verified-recurrence-propagation-crossing-exclusion}

Let \cref{def-charged-structural-scaling-recurrence-certificate} be
proof-grade.  Then every qualified record in its covered source cell follows
one of the verified enclosures.  Consequently,
\begin{equation}
\label{eq-verified-recurrence-source-cell-bound}
\sup_{R\in\mathcal S_{c,n}(B_n)}V_{c,n}(R)
\le
\sup_{j,u\in\overline U_{c,n,j}}A_{c,n}(u)+e_{c,n}.
\end{equation}
Moreover, if \(B_n<\operatorname{Cross}_{c,n}(q)\), no record in that
budgeted cell reaches progress \(q\).  Hence
\begin{equation}
\label{eq-crossing-excluded-transfer-bound}
\sup_{R\in\mathcal S_{c,n}(B_n)}V_{c,n}(R)
\le
\sup_{\substack{j,u\in\overline U_{c,n,j}\\
                 \pi_{c,n}(u)<q}}
A_{c,n}(u)+e_{c,n}.
\end{equation}

If a learned formula \(\widehat u_{c,n,j}(\theta)\) supplies the interval
enclosures and the verifier establishes
\cref{eq-structural-scaling-inductive-enclosure} symbolically for all
\(n\ge N_0\), then the conclusion holds for all such sizes.  The empirical
prefix selects and calibrates the candidate formula; the verified recurrence
proves its asymptotic continuation.

\end{theorem}

\begin{proof}
The recurrence cover assigns a path to each qualified record.  Induction on
the path index in \cref{eq-structural-scaling-inductive-enclosure} places
every state of that path in its displayed enclosure.  Applying
\cref{eq-structural-scaling-progress-transfer} and taking the supremum gives
\cref{eq-verified-recurrence-source-cell-bound}.

If a budget-\(B_n\) path reached \(q\),
\cref{eq-structural-scaling-incremental-cost} would make its complete cost at
least the infimum in \cref{eq-structural-scaling-crossing-price}, contrary
to \(B_n<\operatorname{Cross}_{c,n}(q)\).  Restricting the first supremum to
states below \(q\) proves \cref{eq-crossing-excluded-transfer-bound}.  The
last assertion is the same induction with the learned formula used only as
the verifier's proposed enclosure.
\end{proof}

\begin{corollary}[A scaling breakpoint has a represented structural price]
\label{cor-scaling-breakpoint-structural-price}

Under \cref{thm-verified-recurrence-propagation-crossing-exclusion}, a change
from a verified subcritical envelope \(\pi_{c,n}<q_n\) to a regime with
\(\pi_{c,n}\ge q_n\) occurs only through a represented recurrence path whose
complete cost is at least \(\operatorname{Cross}_{c,n}(q_n)\).  A change of
recurrence grammar, acquisition interface, reuse pattern, precision, or
locator is a new represented transition map and carries its own code,
construction, evaluation, memory, and deployment charge.  Thus a
proof-grade lower bound
\begin{equation}
\label{eq-scaling-breakpoint-off-budget}
\operatorname{Cross}_{c,n}(q_n)>B_n
\end{equation}
excludes the breakpoint inside the declared saturated cell.

\end{corollary}

\begin{proof}
This is the crossing exclusion in
\cref{thm-verified-recurrence-propagation-crossing-exclusion}.  Constructor
provenance and complete-cost accounting place every modified recurrence
component in the path record before the modified transition is used.
\end{proof}

\begin{definition}[Disjoint collision--cancellation recurrence]
\label{def-disjoint-collision-cancellation-recurrence}

Consider a represented combining cell whose stage-\(j\) locator assigns
the current \(N_{j-1}\) records to a represented bucket set
\(\mathcal Q_j\) of cardinality \(Q_j\) and combines disjoint pairs inside
each bucket.  Write \(n_{j-1,a}\) for the occupancy of bucket
\(a\in\mathcal Q_j\).  Let \(d_j\) be the certified number of semantic
coordinates cancelled by equality of one bucket label, let \(w_j\) be
the maximum number of primitive noisy leaves in a resulting record, and let
\(D_j\) be the declared remaining semantic-carrier dimension.  The canonical
complete-pairing rule has the exact and guaranteed row counts
\begin{align}
\label{eq-generic-collision-exact-row-recurrence}
N_j&=\sum_{a\in\mathcal Q_j}
\left\lfloor\frac{n_{j-1,a}}2\right\rfloor,\\
\label{eq-generic-collision-guaranteed-row-recurrence}
N_j&\ge
\left\lfloor\frac{(N_{j-1}-Q_j)_+}{2}\right\rfloor,\\
\label{eq-generic-collision-semantic-provenance-recurrence}
D_j&=D_{j-1}-d_j,
&w_j&=2w_{j-1},\qquad w_0=1.
\end{align}
For independent primitive binary noise with signed mean \(\rho\), a disjoint
depth-\(j\) record has exact signed mean \(\rho^{w_j}=\rho^{2^j}\).
One conservative reverse sufficient-row recurrence for \(K\) final records is
\begin{equation}
\label{eq-generic-collision-reverse-sufficient-recurrence}
N_{j-1}^{\rm suff}=2N_j^{\rm suff}+Q_j+1,
\qquad N_r^{\rm suff}=K.
\end{equation}
The complete recurrence state also stores the locator-evaluation,
bucket-table, pairing, combination, provenance-DAG, acquisition, terminal
readout, memory, and description costs.  Alternative reuse or approximate
collision rules define different represented recurrence maps and retain
their dependence and compensation coordinates.

\end{definition}

\begin{theorem}[Exact transport and charged scaling certificate for
disjoint collision combining]
\label{thm-exact-scaling-certificate-disjoint-collision-combining}

For the cell in
\cref{def-disjoint-collision-cancellation-recurrence}, the occupancy
recurrence and noise transport are exact deterministic certificate fields.
The guaranteed forward recurrence and its reverse are constructive
sufficient-row fields.  After \(r\) disjoint rounds,
\begin{equation}
\label{eq-generic-collision-terminal-scaling}
\sum_{j=1}^{r}d_j=D_0-D_r,
\qquad
w_r=2^r,
\qquad
a_r=|\rho|^{2^r},
\qquad
N_0^{\rm suff}
=2^rK+\sum_{j=1}^{r}2^{j-1}(Q_j+1).
\end{equation}
The minimum complete cost of these sufficient schedules,
\begin{equation}
\label{eq-generic-collision-constructive-cost-optimization}
\operatorname{Build}^{\rm coll}_{n}(q,K)
=
\min_{\substack{r\ge0,\ Q_j,d_j\text{ represented}\\
                 \sum_{j\le r}d_j\ge q}}
\operatorname{Cost}_{\rm complete}
\left(r,(Q_j,d_j)_{j\le r},N_0^{\rm suff},K\right),
\end{equation}
is a constructive upper bound on the cost of reaching that cancellation
level inside the declared collision grammar.  A lower bound on the charged
crossing price is instead obtained by minimizing the incremental-cost
functional over the exact occupancy paths in
\cref{eq-generic-collision-exact-row-recurrence}.  In particular, any
deterministic certificate
\begin{equation}
\label{eq-generic-collision-crossing-lower-certificate}
\underline{\operatorname{Cross}}^{\rm coll}_{n}(q,K)
\le
\inf_{\substack{\text{covered exact occupancy paths}\\
                 \sum_jd_j\ge q,\ N_r\ge K}}
\left(C_{c,n,0}+\sum_jp_{c,n}(u_j,u_{j+1})\right)
\end{equation}
is a valid lower crossing certificate in
\cref{eq-structural-scaling-crossing-price}.
If a source-exhaustion theorem routes a complete source cell to this
recurrence grammar, the lower certificate bounds that complete cell;
otherwise it bounds precisely the collision cell named in the recurrence
cover.

\end{theorem}

\begin{proof}
Complete pairing gives
\cref{eq-generic-collision-exact-row-recurrence}.  At most one row per
bucket remains unpaired, which gives
\cref{eq-generic-collision-guaranteed-row-recurrence}.  Disjoint pairing doubles the
primitive leaf count and independence multiplies signed means, proving
\cref{eq-generic-collision-semantic-provenance-recurrence} and the first
three identities in \cref{eq-generic-collision-terminal-scaling}.  Repeated
substitution in
\cref{eq-generic-collision-reverse-sufficient-recurrence} gives the last
identity.  It constructs a feasible schedule and therefore an upper cost.
The exact occupancy recurrence supplies the path set used by
\cref{eq-generic-collision-crossing-lower-certificate}; applying
\cref{thm-verified-recurrence-propagation-crossing-exclusion} turns any
verified lower bound on that path infimum into an exclusion certificate.
\end{proof}

\begin{figure}[tbp]
\centering
\resizebox{.98\linewidth}{!}{%
\begin{tikzpicture}[fw diagram,node distance=7mm and 9mm]
  \node[fw source,text width=2.25cm] (runs)
    {finite certified\\training runs};
  \node[fw use,text width=2.35cm,right=of runs] (proposal)
    {learned scaling law\\and recurrence};
  \node[fw provenance,text width=2.45cm,right=of proposal] (cover)
    {source-cell cover and\\complete-cost recurrence};
  \node[fw geom,text width=2.35cm,right=of cover] (induct)
    {symbolic interval\\induction};
  \node[fw terminal,text width=2.45cm,right=of induct] (bound)
    {tail bound and\\crossing price};
  \draw[fw arrow] (runs)--(proposal);
  \draw[fw arrow] (proposal)--(cover);
  \draw[fw arrow] (cover)--(induct);
  \draw[fw arrow] (induct)--(bound);
  \node[fw residual,text width=3.0cm,below=9mm of induct] (charge)
    {locator, data, recurrence, memory,\\precision and verification charged};
  \draw[fw arrow] (charge)--(cover);
  \draw[fw arrow] (charge)--(induct);
\end{tikzpicture}%
}
\caption{Proof-producing scaling-law verification.  Empirical learning
proposes the recurrence and its envelope.  Structural source coverage,
complete-cost accounting, and symbolic induction establish the tail and the
price of entering a different scaling regime.}
\label{fig-proof-producing-scaling-law-verification}
\end{figure}

\begin{figure}[tbp]
\centering
\resizebox{.98\linewidth}{!}{%
\begin{tikzpicture}[fw diagram,node distance=7mm and 9mm]
  \node[fw source,text width=2.4cm] (learn)
    {trained empirical\\certificate search};
  \node[fw use,text width=2.4cm,right=of learn] (base)
    {simultaneous base-size\\upper candidates};
  \node[fw provenance,text width=2.5cm,right=of base] (verify)
    {exact source cover and\\inequality verification};
  \node[fw geom,text width=2.4cm,right=of verify] (scale)
    {capacity--attenuation\\scale propagation};
  \node[fw terminal,text width=2.4cm,right=of scale] (profile)
    {saturated source and\\arrival upper bound};
  \draw[fw arrow] (learn)--(base);
  \draw[fw arrow] (base)--(verify);
  \draw[fw arrow] (verify)--(scale);
  \draw[fw arrow] (scale)--(profile);
\end{tikzpicture}%
}
\caption{Empirical discovery and proof-grade closure.  Sampling and training
propose a certificate; deterministic source coverage, factorization, norm,
and scale checks establish the universal saturated-profile bound.}
\label{fig-empirical-scale-certified-source-closure}
\end{figure}

\subsection{Finite-network audit and selection}
\label{subsec-finite-network-implementation-audit}

\begin{corollary}[Finite-network simultaneous audit radius]
\label{cor-finite-network-simultaneous-audit-radius}

Let a represented network family have \(p\) real parameters in a Euclidean
ball of radius \(R_\Theta\), stored to the declared precision.  Suppose its
audited scalar \(z\mapsto\ell_\theta(z)\in[0,M]\) is
\(L_\Theta\)-Lipschitz in \(\theta\), uniformly in \(z\).  For an audit
sample of size \(N\), a parameter-net resolution \(\rho>0\), and confidence
\(\delta\), with probability at least \(1-\delta\), simultaneously for all
represented terminal weights in the family,
\begin{equation}
\label{eq-finite-network-simultaneous-audit-radius}
\mathbb E\ell_\theta
\le\widehat{\mathbb E}\ell_\theta
+M\sqrt{\frac{
p\log(1+2R_\Theta L_\Theta/\rho)+\log(1/\delta)}{2N}}
+2\rho.
\end{equation}
The network architecture, parameter norm, Lipschitz certificate, net
resolution, audit sample, and terminal-weight selection rule are part of the
complete implementation record.  The displayed radius may be replaced by
the smaller same-law candidate supplied by
\cref{thm-feature-learning-bound,thm-composite-nn-bound,thm-sequential-learned-operator-cover,thm-martingale-structural-pac-bayes}.

\end{corollary}

\begin{proof}

By \cref{cor-parameter-cover-selection-cost}, the parameter family has a
represented scalar-function \(\rho\)-net of cardinality at most
\((1+2R_\Theta L_\Theta/\rho)^p\).  Hoeffding's inequality and a union bound
give the square-root term simultaneously on the net.  Moving from a selected
parameter to its net point changes both the empirical and population means
by at most \(\rho\), producing \(2\rho\).

\end{proof}

\begin{theorem}[Certified deep-network implementation bridge]
\label{thm-certified-deep-network-implementation-bridge}

Let a complete neural implementation record use finite represented networks
for its action values, belief update, selector, and Bellman barrier.  Suppose:
\begin{enumerate}[label=\textup{(\roman*)}]
\item the deployed record and barrier lie within the common charged ceiling;
\item the evaluation lower bound in
      \cref{prop-held-out-arrival-lower-certificate} is available;
\item each Bellman residual has either the exact verification certificate in
      \cref{eq-exact-bellman-residual-verification} or the coverage certificate
      in \cref{eq-saturated-audit-coverage}; and
\item data-dependent network selection is covered by
      \cref{cor-finite-network-simultaneous-audit-radius} or by a smaller
      simultaneous same-law structural learning radius.
\end{enumerate}
Then the trained deep-network record satisfies the implementation bound
\cref{eq-primal-dual-empirical-implementation-gap}, with every empirical
radius in \cref{eq-empirical-bellman-residual-upper-candidate} replaced by
the selected simultaneous network radius.  If its learned action values also
satisfy
\cref{eq-active-learned-q-uniform-radius,eq-active-learned-q-optimization},
the sharper minimum in \cref{eq-combined-implementation-gap-candidates}
holds.  All terms are explicit functions of the horizon, evaluation and audit
sample sizes, confidence allocation, network parameter count and norm,
precision, coverage constants, structural learning radii, filtering error,
and represented optimization defects.

\end{theorem}

\begin{proof}

Apply \cref{cor-finite-network-simultaneous-audit-radius} to every selected
network scalar used in the Bellman audit, using a confidence allocation whose
sum is at most \(\delta_{\rm up}\).  This verifies the premise of
\cref{thm-empirical-saturated-bellman-upper-bound}.  Combine its upper bound
with \cref{prop-held-out-arrival-lower-certificate} and invoke
\cref{cor-primal-dual-empirical-implementation-certificate}.  The learned
action-value candidate follows from
\cref{thm-learned-active-rule-finite-horizon-loss,cor-learned-structural-bayesian-control}.

\end{proof}

\subsection{Empirical structural-channel attribution}
\label{subsec-empirical-structural-channel-attribution}

\begin{definition}[Audited parameter--source ledger]
\label{def-audited-parameter-source-ledger}

Let
\(\mathfrak J_{\Abs}^{5}
=\{\mathsf{Add},\mathsf{Mult},\mathsf{Ord},\mathsf{Sym},\mathsf{Filt}\}\).
For a complete \(p\)-parameter learner--deployment record \(R_p\), retain the
exact source index and contextual and information ledgers of
\cref{def-complete-source-resolved-learning-ledger}.  Write
\begin{align}
\label{eq-parameter-source-contextual-coordinate}
\mathcal L_{p}^{(a)}
&=\mathcal L_{I,R_p}^{(a)},\\
\label{eq-parameter-source-information-coordinate}
J_{p}^{(a)}
&=J_{I,R_p}^{(a)},\\
\label{eq-parameter-source-statistical-radius}
\Gamma_{t,p}^{(a)}
&=\Gamma_{R_p,\alpha(a)}^{\rm learn}
  (B,m,\boldsymbol\eta,\delta)
\end{align}
for each source \(a\), with the exact mixed-source cell \(\alpha(a)\)
retained whenever more than one source occurs in the backward slice.

An empirical source audit supplies estimates
\(\widehat{\mathcal L}_{p}^{(a)}\) and
\(\widehat J_{p}^{(a)}\), together with simultaneous represented radii
\(\xi_{\mathcal L,p}^{(a)}\) and \(\xi_{J,p}^{(a)}\) satisfying
\begin{align}
\label{eq-empirical-source-contextual-confidence}
\mathcal L_{p}^{(a)}
&\le\widehat{\mathcal L}_{p}^{(a)}
  +\xi_{\mathcal L,p}^{(a)},\\
\label{eq-empirical-source-information-confidence}
J_{p}^{(a)}
&\le\widehat J_{p}^{(a)}+\xi_{J,p}^{(a)}.
\end{align}
The audit sample, exact-cell selector, backward graph, reveal assignment,
noise conversion, and confidence allocation are fields of the complete
record.

\end{definition}

\begin{theorem}[Parameter- and source-resolved empirical target-gain bound]
\label{thm-parameter-source-resolved-empirical-target-gain}

On the simultaneous source-audit event of
\cref{def-audited-parameter-source-ledger}, the deployed target gain of
\(R_p\) satisfies
\begin{equation}
\label{eq-parameter-source-resolved-empirical-target-gain}
g_{R_p}
\le
\min\left\{
\begin{aligned}
&1,\quad U_{R_p}^{\rm stat\to tgt},\quad
U_{R_p}^{\rm dyn\to tgt},\\
&\beta_{R_p}^{+}
+\sqrt{\frac{\log2}{2}
\sum_{a\in\mathfrak J_{\Abs}^{5}}
\left(\widehat J_p^{(a)}+\xi_{J,p}^{(a)}\right)},\\
&eH_B(n)
\sum_{a\in\mathfrak J_{\Abs}^{5}}
\left(\widehat{\mathcal L}_p^{(a)}
      +\xi_{\mathcal L,p}^{(a)}\right)
\end{aligned}
\right\}.
\end{equation}
For an exact source cell \(\alpha\), restriction of every term to records
with that cell gives a cellwise bound \(U_{p,\alpha}^{\rm tgt}\).  The
Rademacher union and PAC--Bayes mixture penalties for selecting the observed
cell are exactly
\cref{eq-source-resolved-learning-rademacher-union,eq-source-resolved-learning-pac-bayes-mixture}.
Thus the audit reports both the source support learned by the network and a
simultaneous upper bound on the target gain carried by that support.

\end{theorem}

\begin{proof}

Insert
\cref{eq-empirical-source-contextual-confidence,eq-empirical-source-information-confidence}
into the recordwise minimum
\cref{eq-source-resolved-learning-target-gain-minimum}.  Monotonicity of the
sum and square root proves
\cref{eq-parameter-source-resolved-empirical-target-gain}.  Restricting the
same recordwise inequality to one exact cell preserves every comparison.
The selection penalties follow from the cited union and mixture identities.

\end{proof}

\begin{definition}[Charged source intervention family]
\label{def-charged-source-intervention-family}

Let \(\mathcal A=\mathfrak J_{\Abs}^{5}\).  A complete record carries a
\emph{charged source intervention family} when, for every
\(S\subseteq\mathcal A\), it represents a valid deployed record \(R_p^S\)
obtained by retaining precisely the source-gated vertices in \(S\) and
replacing every removed source input by its declared target-neutral
baseline.  Each intervention retains the common public presentation,
verifier, horizon, evaluation law, and nonremoved computation.  Its
recompilation, normalization, and deployment costs are charged.

Define the structural intervention value and Shapley allocation by
\begin{align}
\label{eq-source-intervention-value}
v_p(S)&=A_{R_p^S}^{H},\\
\label{eq-source-shapley-allocation}
\phi_{p,a}
&=\sum_{S\subseteq\mathcal A\setminus\{a\}}
\frac{|S|!\,(4-|S|)!}{5!}
\bigl(v_p(S\cup\{a\})-v_p(S)\bigr).
\end{align}
These coordinates are used only when all \(32\) interventions are valid
complete records.  The contextual and exact-cell ledger remains the
source-attribution certificate when that intervention gate is absent.

\end{definition}

\begin{theorem}[Simultaneous empirical source-intervention certificate]
\label{thm-simultaneous-empirical-source-intervention-certificate}

For every \(S\subseteq\mathcal A\), evaluate \(R_p^S\) on
\(N_{\rm int}\) independent episodes and let \(\widehat v_p(S)\) be its
empirical arrival frequency.  With
\begin{equation}
\label{eq-source-intervention-uniform-radius}
r_{\rm int}
=\sqrt{\frac{\log(64/\delta_{\rm int})}{2N_{\rm int}}},
\end{equation}
all \(32\) values satisfy
\begin{equation}
\label{eq-source-intervention-value-confidence}
|v_p(S)-\widehat v_p(S)|\le r_{\rm int}
\end{equation}
simultaneously with probability at least \(1-\delta_{\rm int}\).  If
\(\widehat\phi_{p,a}\) is obtained from
\cref{eq-source-shapley-allocation} by replacing \(v_p\) with
\(\widehat v_p\), then
\begin{align}
\label{eq-source-shapley-confidence}
|\phi_{p,a}-\widehat\phi_{p,a}|
&\le2r_{\rm int},\\
\label{eq-source-shapley-efficiency}
\sum_{a\in\mathcal A}\phi_{p,a}
&=v_p(\mathcal A)-v_p(\varnothing).
\end{align}
Hence the empirical run has two complementary rigorous descriptions:
the source ledger gives mechanism-wise upper certificates, while the
intervention family gives an interaction-aware allocation of the observed
arrival improvement.

\end{theorem}

\begin{proof}

Hoeffding's inequality gives failure probability at most
\(2e^{-2N_{\rm int}r_{\rm int}^2}=\delta_{\rm int}/32\) for one subset.
A union bound proves
\cref{eq-source-intervention-value-confidence}.  The Shapley weights for
one source are nonnegative and sum to one.  Each marginal difference uses
two intervention values, proving
\cref{eq-source-shapley-confidence}.  The efficiency identity follows by
the standard permutation form of the same weights: along every permutation,
the five marginal increments telescope from \(v_p(\varnothing)\) to
\(v_p(\mathcal A)\), and averaging preserves the identity.

\end{proof}

\subsection{Canonical primal--dual learning objective}
\label{subsec-canonical-primal-dual-learning-objective}

\begin{definition}[Canonical attention--belief training objective]
\label{def-canonical-attention-belief-training-objective}

Let \(R_\theta\) be the deployed record of
\cref{def-canonical-certified-attention-belief-architecture}.  Its population
primal loss is
\begin{equation}
\label{eq-canonical-primal-arrival-loss}
\mathcal L_{\rm arr}(\theta)
=-A_{R_\theta}^{H}
=-\Pr_{R_\theta}\{\tau_T\le H\}.
\end{equation}
For a represented critic \(v=(v_t)_{t\le H}\), \(v_H=0\), write
\begin{equation}
\label{eq-canonical-bellman-action-value}
Q_t^v(h,a)=p_T(h,a)+P_Nv_{t+1}(h,a).
\end{equation}
Let
\[
(\mathcal T_tv)(h)
=\max_{a\in\mathcal A_t(h)}Q_t^v(h,a)
\]
be the represented finite-action Bellman operator.
Given a represented sampling law \(\mu_t\), define
\begin{align}
\label{eq-canonical-squared-bellman-loss}
\mathcal L_{\rm Bell}(v)
&=\sum_{t<H}\mathbb E_{\mu_t}
\left[
(\mathcal T_tv)(h)-v_t(h)
\right]^2,\\
\label{eq-canonical-barrier-loss}
\mathcal L_{\rm bar}(v)
&=v_0(h_0)+
\sum_{t<H}C_t\,
\mathbb E_{\mu_t}
\left[
Q_t^v(h,a)-v_t(h)
\right]_+ .
\end{align}
The first loss trains a critic on its declared sampling law.  The second is
the population version of the saturated Bellman upper certificate and
retains the coverage constants \(C_t\).

Let \(\nu_t(\cdot\mid h)\) be the represented neutral reference row with
positive mass on every legal action, and put
\begin{equation}
\label{eq-canonical-bellman-gibbs-row}
\pi_{v,t}^{\tau}(a\mid h)
=
\frac{\nu_t(a\mid h)
\exp(Q_t^v(h,a)/\tau_t)}
{\sum_{a'}\nu_t(a'\mid h)
\exp(Q_t^v(h,a')/\tau_t)} .
\end{equation}
The policy and belief losses are
\begin{align}
\label{eq-canonical-attention-policy-loss}
\mathcal L_{\rm pol}(\theta;v)
&=
\sum_{t<H}\mathbb E_{f_{R_\theta,t}}
D_{\rm KL}\!\left(
\pi_{\theta,t}(\cdot\mid h)
\middle\Vert
\pi_{v,t}^{\tau}(\cdot\mid h)
\right),\\
\label{eq-canonical-channel-likelihood-loss}
\mathcal L_{\rm Bayes}(\theta)
&=
-\sum_{t<H}
\mathbb E\log
\left[
\sum_z q_{\theta,t}(z)
L_{\boldsymbol\eta,t}(Y_t\mid z,A_t)
\right].
\end{align}
Here \(L_{\boldsymbol\eta,t}\) is the likelihood of the declared observation
channel, including its noise parameters, and \(q_{\theta,t}\) is the
represented compact belief before the reveal.  Let
\(\mathcal L_{\rm num}\) be a represented upper loss for finite-precision
score, normalization, likelihood, and state-update discrepancies.  The
canonical trainable objective is
\begin{equation}
\label{eq-canonical-composite-training-objective}
\boxed{
\mathcal L_{\rm train}(\theta,v)
=
\mathcal L_{\rm pol}(\theta;v)
+\lambda_V\mathcal L_{\rm Bell}(v)
+\lambda_B\mathcal L_{\rm Bayes}(\theta)
+\lambda_{\rm num}\mathcal L_{\rm num}(\theta,v),
}
\end{equation}
with represented nonnegative weights.  The geometry and quotient heads are
trained on their declared empirical forms and are certified by the separate
same-law audits in
\cref{subsec-empirical-analytical-geometry,subsec-empirical-lumpability-certification}.

The empirical optimality objective is the primal--dual gap
\begin{equation}
\label{eq-canonical-primal-dual-certificate-loss}
\boxed{
\mathcal L_{\rm cert}(\widehat\theta,v)
=U_{\rm Bell}(v)-L_{\rm ev}(\widehat\theta).
}
\end{equation}
Its construction, data split, confidence allocation, and optimization are
fields of the complete implementation record.

\end{definition}

\begin{theorem}[Canonical loss calibration to saturated arrival]
\label{thm-canonical-loss-saturated-arrival-calibration}

Assume \(0<\nu_{\min,t}\le\nu_t(a\mid h)\) on every represented legal action
and surviving history.  Suppose
\begin{align}
\label{eq-canonical-q-calibration}
\|Q_t^v-Q_t^{\rm act,*}\|_\infty&\le\Gamma_{Q,t},\\
\label{eq-canonical-policy-kl-calibration}
D_{\rm KL}\!\left(
\pi_{\theta,t}(\cdot\mid h)
\middle\Vert
\pi_{v,t}^{\tau}(\cdot\mid h)
\right)&\le\kappa_t
\end{align}
uniformly on the surviving carrier.  Then the stepwise action defect obeys
\begin{equation}
\label{eq-canonical-stepwise-action-defect}
\max_a Q_t^{\rm act,*}(h,a)
-\mathbb E_{a\sim\pi_{\theta,t}(\cdot\mid h)}
Q_t^{\rm act,*}(h,a)
\le
2\Gamma_{Q,t}
+\tau_t\log\frac1{\nu_{\min,t}}
+\operatorname{Osc}(Q_t^v)\sqrt{\frac{\kappa_t}{2}}.
\end{equation}
Consequently,
\begin{align}
\label{eq-canonical-policy-arrival-loss}
J_{\rm sat}^*(B)-A_{R_\theta}^{H}
\le{}&
\varepsilon_{\rm rep}(p;B)
+\mathbb E_{R_\theta}
\sum_{t<H}\mathbf1_{\{\tau>t\}}
\left[
2\Gamma_{Q,t}
+\tau_t\log\frac1{\nu_{\min,t}}
+\operatorname{Osc}(Q_t^v)\sqrt{\frac{\kappa_t}{2}}
\right]\notag\\
&+\Delta_{\rm fil}^{H}+\varepsilon_{\rm num}.
\end{align}
On the simultaneous held-out and Bellman-audit event,
\begin{equation}
\label{eq-canonical-certificate-loss-optimality}
0\le
J_{\rm sat}^*(B)-A_{R_{\widehat\theta}}^H
\le
\mathcal L_{\rm cert}(\widehat\theta,v).
\end{equation}
Thus minimization of
\cref{eq-canonical-composite-training-objective} is connected to saturated
arrival through the displayed calibration coordinates, while
\cref{eq-canonical-primal-dual-certificate-loss} directly certifies the
remaining optimality gap.

\end{theorem}

\begin{proof}

The Gibbs variational identity
\cref{eq-attention-gibbs-variational-identity} gives
\[
\max_aQ_t^v(h,a)
-\mathbb E_{\pi_{v,t}^{\tau}}Q_t^v(h,a)
\le\tau_t\log\frac1{\nu_{\min,t}}.
\]
Pinsker's inequality and
\cref{eq-canonical-policy-kl-calibration} give
\[
\left|
\mathbb E_{\pi_{\theta,t}}Q_t^v
-\mathbb E_{\pi_{v,t}^{\tau}}Q_t^v
\right|
\le
\operatorname{Osc}(Q_t^v)\sqrt{\kappa_t/2}.
\]
Replacing \(Q_t^v\) by \(Q_t^{\rm act,*}\) in the maximum and expectation contributes
at most \(2\Gamma_{Q,t}\), proving
\cref{eq-canonical-stepwise-action-defect}.  Apply the finite-horizon
performance comparison in
\cref{thm-learned-active-rule-finite-horizon-loss}, add the represented
filter and numerical comparisons, and then add the parameter-class
representation gap.  This proves
\cref{eq-canonical-policy-arrival-loss}.

On the audit event,
\[
L_{\rm ev}(\widehat\theta)
\le A_{R_{\widehat\theta}}^H
\le J_{\rm sat}^*(B)
\le U_{\rm Bell}(v).
\]
Subtraction proves
\cref{eq-canonical-certificate-loss-optimality}.

\end{proof}

\subsection{Certified WAdam optimization}
\label{subsec-certified-wadam-optimization}

\begin{definition}[Certified WAdam training record]
\label{def-certified-wadam-training-record}

Let \(F:\mathbb R^p\to\mathbb R\) be the represented training objective.
A \emph{certified WAdam training record} stores the actual parameter iterates
\(\theta_k\), stochastic gradients, first and second moment states, bias
corrections, clipping operations, and positive preconditioners used by the
implementation.  Its terminal update is written in the common form
\begin{equation}
\label{eq-certified-wadam-update}
\theta_{k+1}=\theta_k-\alpha_k d_k,
\qquad
d_k=W_k^{-1}\widehat m_k,
\end{equation}
where \(W_k\) and \(\widehat m_k\) are the represented preconditioner and
bias-corrected first moment produced by the implemented Adam-type recursions.
Diagonal Adam, block-weighted Adam, and geometry-weighted variants are
instances of \cref{eq-certified-wadam-update}.  The optimizer is certified
through represented constants \(a_k,c_k>0\), \(b_k,s_k\ge0\) satisfying,
for the training filtration \(\mathcal G_k\),
\begin{align}
\label{eq-wadam-certified-gradient-alignment}
\mathbb E[\langle\nabla F(\theta_k),d_k\rangle\mid\mathcal G_k]
&\ge a_k\lVert\nabla F(\theta_k)\rVert_2^2-b_k,\\
\label{eq-wadam-certified-direction-second-moment}
\mathbb E[\lVert d_k\rVert_2^2\mid\mathcal G_k]
&\le c_k\lVert\nabla F(\theta_k)\rVert_2^2+s_k.
\end{align}
The complete record charges the computation and verification of these
constants.  Spectral bounds on \(W_k\), stochastic-gradient bias and
variance bounds, or direct interval verification of the two displayed
inequalities are admissible certificates when represented at their complete
cost.

\end{definition}

\begin{theorem}[WAdam descent and global optimization certificate]
\label{thm-wadam-descent-global-optimization-certificate}

Let \(\Theta_{\rm cert}\) be a represented region containing every iterate,
put \(F_*=\inf_{\theta\in\Theta_{\rm cert}}F(\theta)\), and assume \(F\) is
\(L_F\)-smooth on every segment \([\theta_k,\theta_{k+1}]\).  Suppose the certified
WAdam conditions
\cref{eq-wadam-certified-gradient-alignment,eq-wadam-certified-direction-second-moment}
hold.  Put
\begin{align}
\label{eq-wadam-effective-descent-coefficient}
q_k&=\alpha_ka_k-\frac{L_F\alpha_k^2c_k}{2},\\
\label{eq-wadam-additive-descent-error}
r_k&=\alpha_kb_k+\frac{L_F\alpha_k^2s_k}{2}.
\end{align}
If \(q_k>0\), then
\begin{equation}
\label{eq-wadam-stationarity-certificate}
\sum_{k=0}^{K-1}q_k
\mathbb E\lVert\nabla F(\theta_k)\rVert_2^2
\le F(\theta_0)-F_*+\sum_{k=0}^{K-1}r_k.
\end{equation}

Suppose additionally that the represented objective satisfies the
Polyak--\L{}ojasiewicz comparison
\begin{equation}
\label{eq-wadam-pl-certificate}
\lVert\nabla F(\theta)\rVert_2^2
\ge2\mu\bigl(F(\theta)-F_*\bigr)
\end{equation}
on the visited certified region, and \(0\le2\mu q_k\le1\).  Then
\begin{equation}
\label{eq-wadam-global-objective-gap}
\mathbb E[F(\theta_K)-F_*]
\le
\left[\prod_{k=0}^{K-1}(1-2\mu q_k)\right]
[F(\theta_0)-F_*]
+\sum_{j=0}^{K-1}r_j
  \prod_{k=j+1}^{K-1}(1-2\mu q_k).
\end{equation}
If the represented acquisition optimizer has a calibration
\begin{equation}
\label{eq-wadam-objective-to-action-calibration}
\varepsilon_{{\rm opt},t}
\le\lambda_t\bigl(F(\theta_K)-F_*\bigr)+\zeta_t,
\end{equation}
then substitution of \cref{eq-wadam-global-objective-gap} into
\cref{eq-combined-implementation-gap-candidates} gives a completely explicit
WAdam contribution to the saturated implementation error.
Moreover, if
\(\mathcal E_{\rm WAdam}(K)\) denotes the right side of
\cref{eq-wadam-global-objective-gap}, then for every
\(0<\delta_{\rm opt}<1\),
\begin{equation}
\label{eq-wadam-high-probability-objective-gap}
\Pr\left\{
F(\theta_K)-F_*
\le\frac{\mathcal E_{\rm WAdam}(K)}{\delta_{\rm opt}}
\right\}
\ge1-\delta_{\rm opt}.
\end{equation}
For a \(p\)-parameter architecture, write
\(\mathcal E_{\rm WAdam}(p,K)\); this notation retains the dependence of
\(F,L_F,\mu,(a_k,b_k,c_k,s_k)_{k<K}\), the initialization, precision, and
all optimizer-state costs on \(p\).

\end{theorem}

\begin{proof}

Smoothness and \cref{eq-certified-wadam-update} give
\[
F(\theta_{k+1})
\le F(\theta_k)
-\alpha_k\langle\nabla F(\theta_k),d_k\rangle
+\frac{L_F\alpha_k^2}{2}\lVert d_k\rVert_2^2.
\]
Take conditional expectations and use
\cref{eq-wadam-certified-gradient-alignment,eq-wadam-certified-direction-second-moment}
to obtain
\begin{equation}
\label{eq-wadam-one-step-descent}
\mathbb E[F(\theta_{k+1})\mid\mathcal G_k]
\le F(\theta_k)-q_k\lVert\nabla F(\theta_k)\rVert_2^2+r_k.
\end{equation}
Summing, taking expectations, and using
\(\mathbb EF(\theta_K)\ge F_*\) proves
\cref{eq-wadam-stationarity-certificate}.

Under \cref{eq-wadam-pl-certificate}, subtract \(F_*\) from
\cref{eq-wadam-one-step-descent} to obtain
\[
\mathbb E[F(\theta_{k+1})-F_*\mid\mathcal G_k]
\le(1-2\mu q_k)[F(\theta_k)-F_*]+r_k.
\]
Iteration proves \cref{eq-wadam-global-objective-gap}.  The calibrated
optimization error is then bounded by
\cref{eq-wadam-objective-to-action-calibration}, and
\cref{cor-primal-dual-empirical-implementation-certificate} completes the
stated substitution.  Since \(F(\theta_K)-F_*\ge0\), Markov's inequality
applied to \cref{eq-wadam-global-objective-gap} proves
\cref{eq-wadam-high-probability-objective-gap}.

\end{proof}

\subsection{Certified structural bias--variance ledger}
\label{subsec-certified-structural-bias-variance-ledger}

\begin{definition}[Deterministic bias and stochastic-radius coordinates]
\label{def-deterministic-bias-stochastic-radius-coordinates}

For one \(p\)-parameter canonical record, split every valid clean
action-value comparison as
\begin{equation}
\label{eq-action-value-bias-radius-split}
\Gamma_{Q,t}=\beta_{Q,t}+\sigma_{Q,t},
\qquad
\beta_{Q,t},\sigma_{Q,t}\ge0,
\end{equation}
where \(\beta_{Q,t}\) collects certified population approximation,
regularization, truncation, and deterministic model discrepancy, while
\(\sigma_{Q,t}\) is a simultaneous statistical radius obtained from a
represented iid, mixing, sequential, or martingale certificate.  Define
\begin{equation}
\label{eq-complete-wadam-action-coordinate}
\varepsilon_{\rm WAdam}(p,K)
=
\mathbb E_{R_{\widehat\theta_p}}
\sum_{t<H}\mathbf1_{\{\tau>t\}}
\left[
\lambda_{t,p}
\frac{\mathcal E_{\rm WAdam}(p,K)}{\delta_{\rm opt}}
+\zeta_{t,p}
\right]
\end{equation}
on the high-probability optimizer event, and define
\begin{align}
\label{eq-complete-deterministic-bias-coordinate}
B_{\rm det}(p,K,\boldsymbol\tau,b,H)
={}&
\varepsilon_{\rm rep}(p;B)
+\varepsilon_{\rm WAdam}(p,K)
+\Delta_{{\rm fil},p}^{H}
+\varepsilon_{{\rm num},p}
+\varepsilon_{{\rm trunc},H}
+\varepsilon_{{\rm suff},b}\notag\\
&+
\mathbb E_{R_{\widehat\theta_p}}
\sum_{t<H}\mathbf1_{\{\tau>t\}}
\left[
2\beta_{Q,t}
+\tau_t\log\frac1{\nu_{\min,t}}
+\operatorname{Osc}(Q_t^v)\sqrt{\frac{\kappa_t}{2}}
\right],\\
\label{eq-complete-stochastic-radius-coordinate}
V_{\rm stat}(p,N,\boldsymbol\eta,\delta)
={}&
2\mathbb E_{R_{\widehat\theta_p}}
\sum_{t<H}\mathbf1_{\{\tau>t\}}\sigma_{Q,t}.
\end{align}
Here \(\varepsilon_{\rm WAdam}(p,K)\) is the calibrated action-value
contribution obtained from
\cref{eq-wadam-global-objective-gap,eq-wadam-objective-to-action-calibration};
\(\varepsilon_{{\rm suff},b}\) is the compact-belief sufficiency defect at
code capacity \(b\); and \(\varepsilon_{{\rm trunc},H}\) is the certified
tail beyond the deployed horizon.  An absent coordinate receives its neutral
valid bound, so the ledger remains total.

For a predictable loss sequence, retain additionally the conditional
variance
\begin{equation}
\label{eq-canonical-predictable-variance-ledger}
\mathcal V_T(Q)
=
\mathbb E_{h\sim Q}
\sum_{t=1}^{T}
\mathbb E\!\left[
\bigl(r_t(h)-\ell_t(h)\bigr)^2
\middle|\mathcal F_{t-1}
\right].
\end{equation}
This is the variance coordinate used by the Bernstein PAC--Bayes radius in
\cref{eq-martingale-structural-pac-bayes-bernstein}.

\end{definition}

\begin{theorem}[Complete certified bias--variance decomposition]
\label{thm-complete-certified-bias-variance-decomposition}

On the simultaneous calibration, learning, filtering, numerical, and
optimizer event of a complete canonical record,
\begin{equation}
\label{eq-complete-certified-bias-variance-decomposition}
\boxed{
0\le
J_{\rm sat}^*(B)-A_{R_{\widehat\theta_p}}^H
\le
B_{\rm det}(p,K,\boldsymbol\tau,b,H)
+V_{\rm stat}(p,N,\boldsymbol\eta,\delta).
}
\end{equation}
The stochastic coordinate may use the smallest same-law candidate among the
structural Rademacher, PAC--Bayes, mixing-block, sequential, martingale, and
predictable-variance bounds already carried by the record.  In particular,
the martingale PAC--Bayes contribution at
\(0<\lambda<3/c\) is
\begin{equation}
\label{eq-canonical-pac-bayes-variance-radius}
\frac{
\mathrm{KL}^{\mathsf{str}}_{E,H}(Q)+\log(1/\delta)
}{\lambda T}
+
\frac{
\lambda\mathcal V_T(Q)
}{
2T(1-\lambda c/3)
}.
\end{equation}
If the entire right side of
\cref{eq-complete-certified-bias-variance-decomposition} is unavailable, the
primal--dual candidate
\(\mathcal L_{\rm cert}\wedge1\) remains a total certified gap.

\end{theorem}

\begin{proof}

Insert the split
\cref{eq-action-value-bias-radius-split} into
\cref{eq-canonical-policy-arrival-loss}.  Add the calibrated WAdam,
truncation, and compact-belief sufficiency defects, which compare the same
deployed record with the same saturated optimum.  Deterministic terms give
\cref{eq-complete-deterministic-bias-coordinate}; simultaneous statistical
terms give
\cref{eq-complete-stochastic-radius-coordinate}.  This proves
\cref{eq-complete-certified-bias-variance-decomposition}.  The displayed
predictable-variance candidate is
\cref{eq-martingale-structural-pac-bayes-bernstein}.  The final assertion is
\cref{eq-canonical-certificate-loss-optimality} clipped by the trivial
arrival range.

\end{proof}

\begin{proposition}[Exact clipping bias--variance identity]
\label{prop-exact-clipping-bias-variance-identity}

Let \(\mu\) be a represented audit law, let \(P\ll\mu\) be the deployment
law with represented density \(W=dP/d\mu\), and let \(0\le Y\le1\).  For
independent \((W_i,Y_i)_{i\le N}\) from \(\mu\), define
\[
J=\mathbb E_\mu[WY],
\qquad
\widehat J_C
=\frac1N\sum_{i=1}^{N}\min\{W_i,C\}Y_i .
\]
Then
\begin{align}
\label{eq-clipping-exact-bias}
b_C
:=J-\mathbb E\widehat J_C
&=\mathbb E_\mu[(W-C)_+Y],\\
\label{eq-clipping-exact-variance}
\operatorname{Var}(\widehat J_C)
&=\frac1N
\operatorname{Var}_\mu(\min\{W,C\}Y),\\
\label{eq-clipping-exact-mse}
\mathbb E(\widehat J_C-J)^2
&=b_C^2+
\frac1N\operatorname{Var}_\mu(\min\{W,C\}Y),\\
\label{eq-clipping-variance-upper}
\operatorname{Var}(\widehat J_C)
&\le\frac{C J}{N}.
\end{align}
Thus the clipping level has a represented exact bias term and a represented
second-moment variance term; its selection cost remains part of the audit
record.

\end{proposition}

\begin{proof}

Subtract \(\min\{W,C\}Y\) from \(WY\) to obtain
\cref{eq-clipping-exact-bias}.  Independence gives
\cref{eq-clipping-exact-variance}; the scalar bias--variance identity gives
\cref{eq-clipping-exact-mse}.  Finally,
\(\min\{W,C\}^2Y^2\le CWY\), proving
\cref{eq-clipping-variance-upper}.

\end{proof}

\begin{corollary}[Quantitative architecture and audit tradeoffs]
\label{cor-quantitative-architecture-audit-tradeoffs}

The complete canonical record has the following certified tuning curves.
\begin{enumerate}[label=\textup{(\roman*)}]
\item \emph{Parameter count.}
      The upper curve
      \[
      \varepsilon_{\rm rep}^{\rm up}(p)
      +\varepsilon_{\rm WAdam}^{\rm up}(p,K)
      +2\Gamma_p(N,\delta)
      \]
      combines the nonincreasing representation error of
      \cref{def-parameter-indexed-saturated-neural-profile} with the
      parameter-cover radius in
      \cref{eq-finite-network-simultaneous-audit-radius}.
\item \emph{Attention temperature.}
      If a represented score estimate satisfies
      \(\|\widehat s_t-s_t\|_\infty\le\epsilon_{s,t}\), then its Gibbs rows
      satisfy
      \begin{equation}
      \label{eq-temperature-score-kl-tv-tradeoff}
      D_{\rm KL}(P_{s_t}\Vert P_{\widehat s_t})
      \le\frac{2\epsilon_{s,t}}{\tau_t},
      \qquad
      \|P_{s_t}-P_{\widehat s_t}\|_{\rm TV}
      \le\sqrt{\frac{\epsilon_{s,t}}{\tau_t}}.
      \end{equation}
      The corresponding arrival candidate combines this sensitivity with
      the Gibbs approximation bias
      \[
      \sum_{t<H}\tau_t\log(1/\nu_{\min,t}).
      \]
\item \emph{Binary label noise.}
      Clean-risk radii acquire the factor
      \((1-2\eta_+)^{-1}\), and their squared or variance coordinates acquire
      \((1-2\eta_+)^{-2}\), by
      \cref{thm-rademacher-label-corruption-bound,cor-debiased-krr-label-noise}.
\item \emph{Horizon and dependence.}
      The represented tail \(\varepsilon_{{\rm trunc},H}\) is balanced
      against the accumulated predictable variance
      \(\mathcal V_H\) or the certified mixing effective sample size.
\item \emph{Latent capacity.}
      Increasing \(b\) decreases the certified sufficiency defect and
      increases the represented latent cover, PAC--Bayes, storage, and
      optimization coordinates of
      \cref{thm-unified-compact-latent-structural-certificate}.
\item \emph{Functional features.}
      Increasing \(p_{\rm geom}\) changes the statistical radii
      \(\epsilon_{C,p},\epsilon_{G,p}\) in
      \cref{eq-feature-covariance-explicit-radius,eq-feature-dirichlet-explicit-radius}
      and the analytical completeness remainder \(r_p\) in
      \cref{eq-full-poincare-from-feature-completeness}.
\item \emph{Quotient resolution.}
      A quotient with \(|Q|\) cells has statistical defect radius
      \[
      |Q|\sqrt{\frac{\log(2|X||Q|/\delta_{\rm lump})}{2N_{\min}}},
      \]
      together with its represented target-sufficiency and transfer defects.
\end{enumerate}

Let \(\Lambda\) be a represented finite or countable grid of complete
architecture, temperature, clipping, capacity, feature, and quotient
choices, with weights \(w_\lambda>0\),
\(\sum_{\lambda\in\Lambda}w_\lambda\le1\).  Evaluate every candidate at
confidence \(\delta w_\lambda\), let \(U_\lambda\) be its complete certified
upper bound, and select
\begin{equation}
\label{eq-certified-hyperparameter-selection}
\widehat\lambda\in\arg\min_{\lambda\in\Lambda}U_\lambda.
\end{equation}
Then, with probability at least \(1-\delta\), the selected record satisfies
its reported bound \(U_{\widehat\lambda}\).  This selection includes the full
architecture and audit choice rather than tuning one coordinate in
isolation.

\end{corollary}

\begin{proof}

The parameter, noise, horizon, latent, feature, and quotient statements are
the cited confidence or completeness comparisons.  For the temperature
statement, the uniform score error bounds the log-normalizer ratio by
\(\epsilon_{s,t}/\tau_t\), so the log density ratio is bounded by
\(2\epsilon_{s,t}/\tau_t\).  Taking expectation proves the KL bound and
Pinsker proves the total-variation bound.  A weighted union bound over
\(\Lambda\) proves the selected-candidate assertion.

\end{proof}

\begin{figure}[tbp]
\centering
\begin{tikzpicture}[fw diagram,node distance=7mm and 9mm]
  \node[fw source,text width=2.5cm] (knob)
    {architecture and\\audit choices};
  \node[fw provenance,text width=2.5cm,above right=7mm and 10mm of knob] (bias)
    {representation, optimizer,\\filter and temperature bias};
  \node[fw geom,text width=2.5cm,below right=7mm and 10mm of knob] (var)
    {sampling, noise and\\trajectory variance};
  \node[fw use,text width=2.6cm,right=22mm of knob] (sum)
    {complete certified\\upper curve};
  \node[fw terminal,text width=2.5cm,right=of sum] (sel)
    {selected deployed\\record};
  \draw[fw arrow] (knob)--(bias);
  \draw[fw arrow] (knob)--(var);
  \draw[fw arrow] (bias)--(sum);
  \draw[fw arrow] (var)--(sum);
  \draw[fw arrow] (sum)--(sel);
\end{tikzpicture}
\caption{Certified bias--variance selection.  Every architecture and audit
choice changes named deterministic and stochastic coordinates; simultaneous
selection minimizes their complete upper bound.}
\label{fig-certified-bias-variance-selection}
\end{figure}

\subsection{Empirical certification of analytical geometry}
\label{subsec-empirical-analytical-geometry}

\begin{definition}[Feature-restricted Dirichlet audit]
\label{def-feature-restricted-dirichlet-audit}

Let \(\Phi\) be a certified finite dynamical geometry with invariant law
\(\mu\), Dirichlet form \(\mathcal E_\Phi\), and analytical Poincar\'e
constant \(\gamma_\Phi\) from
\cref{def-controlled-diffusive-functional-profile}.  Let
\(\psi_p:X\to\mathbb R^p\) be a represented feature map, centered under
\(\mu\).  Define its covariance and Dirichlet matrices by
\begin{equation}
\label{eq-feature-covariance-dirichlet-matrices}
C_p=\mathbb E_\mu[\psi_p\psi_p^{\mathsf T}],
\qquad
G_p=\bigl(\mathcal E_\Phi(\psi_{p,i},\psi_{p,j})\bigr)_{i,j\le p}.
\end{equation}
Let \(\widehat C_p,\widehat G_p\) be represented unbiased audit estimates and
let \(\epsilon_{C,p},\epsilon_{G,p}\) satisfy, on one simultaneous event,
\begin{equation}
\label{eq-feature-matrix-confidence-event}
\|C_p-\widehat C_p\|_{\rm op}\le\epsilon_{C,p},
\qquad
\|G_p-\widehat G_p\|_{\rm op}\le\epsilon_{G,p}.
\end{equation}
Centering error is included in \(\epsilon_{C,p}\).

For independent audit samples, one explicit entrywise option is
\begin{align}
\label{eq-feature-covariance-explicit-radius}
\epsilon_{C,p}
&=pK_p^2
\sqrt{\frac{2\log(4p^2/\delta_{\rm FI})}{N_C}},\\
\label{eq-feature-dirichlet-explicit-radius}
\epsilon_{G,p}
&=pD_p^2
\sqrt{\frac{2\log(4p^2/\delta_{\rm FI})}{N_G}},
\end{align}
provided every centered coordinate obeys
\(|\psi_{p,i}|\le K_p\), every unbiased sampled Dirichlet increment has
coordinate magnitude at most \(D_p\), and the Dirichlet normalization is
included in \(D_p\).  A matrix concentration, mixing, or martingale radius
may replace these entrywise candidates.  The feature construction,
centering, edge sampler, dependence correction, precision, and selected
radius are charged in the implementation record.

\end{definition}

\begin{theorem}[Empirical Poincar\'e bracket and full-carrier transfer]
\label{thm-empirical-poincare-bracket-full-transfer}

On the event \cref{eq-feature-matrix-confidence-event}, suppose a represented
number \(\underline\gamma_p\ge0\) satisfies the semidefinite certificate
\begin{equation}
\label{eq-empirical-poincare-semidefinite-certificate}
\widehat G_p-\underline\gamma_p\widehat C_p
\succeq
\bigl(\epsilon_{G,p}
      +\underline\gamma_p\epsilon_{C,p}\bigr)I_p.
\end{equation}
Then every \(f_a=a^{\mathsf T}\psi_p\) satisfies
\begin{equation}
\label{eq-restricted-poincare-certified}
\mathcal E_\Phi(f_a,f_a)
\ge\underline\gamma_p\operatorname{Var}_\mu(f_a).
\end{equation}
For every represented \(a\) with
\(a^{\mathsf T}(\widehat C_p-\epsilon_{C,p}I)a>0\), the analytical constant
also satisfies
\begin{equation}
\label{eq-empirical-poincare-upper-test}
\gamma_\Phi
\le
\frac{a^{\mathsf T}(\widehat G_p+\epsilon_{G,p}I)a}
     {a^{\mathsf T}(\widehat C_p-\epsilon_{C,p}I)a}.
\end{equation}

Assume additionally that a represented linear projection
\(\Pi_p\) from the full Poincar\'e domain to
\(\operatorname{span}\psi_p\) has certified constants
\(A_p,D_p^{\rm pr},r_p\ge0\) such that, for every centered \(f\),
\begin{align}
\label{eq-feature-poincare-completeness}
\operatorname{Var}_\mu(f)
&\le
A_p\operatorname{Var}_\mu(\Pi_pf)
+r_p\mathcal E_\Phi(f,f),\\
\label{eq-feature-dirichlet-projection-stability}
\mathcal E_\Phi(\Pi_pf,\Pi_pf)
&\le D_p^{\rm pr}\mathcal E_\Phi(f,f).
\end{align}
If \(\underline\gamma_p>0\), then the true full-carrier constant has the
certified lower bound
\begin{equation}
\label{eq-full-poincare-from-feature-completeness}
\boxed{
\gamma_\Phi
\ge
\left(
\frac{A_pD_p^{\rm pr}}{\underline\gamma_p}+r_p
\right)^{-1}.
}
\end{equation}
Thus \(p\) controls both statistical estimation through
\cref{eq-feature-covariance-explicit-radius,eq-feature-dirichlet-explicit-radius}
and analytical underfitting through the completeness remainder \(r_p\).

\end{theorem}

\begin{proof}

The operator inequalities in
\cref{eq-feature-matrix-confidence-event} give
\[
G_p-\underline\gamma_pC_p
\succeq
\widehat G_p-\underline\gamma_p\widehat C_p
-(\epsilon_{G,p}
  +\underline\gamma_p\epsilon_{C,p})I_p
\succeq0.
\]
Taking the quadratic form in \(a\) proves
\cref{eq-restricted-poincare-certified}.  Conversely,
\(G_p\preceq\widehat G_p+\epsilon_{G,p}I\) and
\(C_p\succeq\widehat C_p-\epsilon_{C,p}I\).  The resulting Rayleigh
quotient is an admissible test in the full Poincar\'e infimum, proving
\cref{eq-empirical-poincare-upper-test}.

For a centered \(f\), apply
\cref{eq-restricted-poincare-certified} to \(\Pi_pf\), followed by
\cref{eq-feature-dirichlet-projection-stability}, in
\cref{eq-feature-poincare-completeness}:
\[
\operatorname{Var}_\mu(f)
\le
\left(
\frac{A_pD_p^{\rm pr}}{\underline\gamma_p}+r_p
\right)\mathcal E_\Phi(f,f).
\]
Taking the reciprocal proves
\cref{eq-full-poincare-from-feature-completeness}.

\end{proof}

\begin{theorem}[Generic empirical functional-inequality transfer]
\label{thm-generic-empirical-functional-inequality-transfer}

Let \(\mathscr D\) be the declared domain of one entry of
\(\mathfrak F_\Phi\).  Write its inequality as
\begin{equation}
\label{eq-generic-functional-inequality-form}
\mathcal D(f)\ge\kappa\,\mathcal V(f),
\qquad f\in\mathscr D,
\end{equation}
after using the normalization of
\cref{def-controlled-functional-inequality-profile}.  Let
\(\mathscr D_p\subseteq\mathscr D\) be a represented feature domain and
\(\mathcal N(f)\ge0\) a represented homogeneous control functional.
Suppose, simultaneously on the audit event,
\begin{align}
\label{eq-generic-functional-dissipation-deviation}
|\widehat{\mathcal D}(f)-\mathcal D(f)|
&\le\epsilon_{\mathcal D,p}\mathcal N(f),\\
\label{eq-generic-functional-variance-deviation}
|\widehat{\mathcal V}(f)-\mathcal V(f)|
&\le\epsilon_{\mathcal V,p}\mathcal N(f)
\end{align}
for all \(f\in\mathscr D_p\), and suppose
\begin{equation}
\label{eq-generic-empirical-functional-margin}
\widehat{\mathcal D}(f)
-\underline\kappa_p\widehat{\mathcal V}(f)
\ge
(\epsilon_{\mathcal D,p}
 +\underline\kappa_p\epsilon_{\mathcal V,p})\mathcal N(f)
\end{equation}
on that domain.  Then
\(\mathcal D(f)\ge\underline\kappa_p\mathcal V(f)\) on
\(\mathscr D_p\).

If a represented projection \(\Pi_p:\mathscr D\to\mathscr D_p\) satisfies
\begin{align}
\label{eq-generic-functional-completeness}
\mathcal V(f)
&\le A_p\mathcal V(\Pi_pf)+r_p\mathcal D(f),\\
\label{eq-generic-functional-projection-stability}
\mathcal D(\Pi_pf)
&\le D_p^{\rm pr}\mathcal D(f),
\end{align}
then the full analytical constant is at least
\begin{equation}
\label{eq-generic-full-functional-constant}
\kappa_\Phi
\ge
\left(
\frac{A_pD_p^{\rm pr}}{\underline\kappa_p}+r_p
\right)^{-1}.
\end{equation}
For any represented \(f\in\mathscr D_p\) satisfying
\(\widehat{\mathcal V}(f)>
\epsilon_{\mathcal V,p}\mathcal N(f)\), the same audit also gives the
analytical upper certificate
\begin{equation}
\label{eq-generic-functional-upper-test}
\kappa_\Phi
\le
\frac{\widehat{\mathcal D}(f)
      +\epsilon_{\mathcal D,p}\mathcal N(f)}
     {\widehat{\mathcal V}(f)
      -\epsilon_{\mathcal V,p}\mathcal N(f)}.
\end{equation}
The theorem applies to the Poincar\'e, weak/super-Poincar\'e at a fixed
represented parameter, Nash, Sobolev, spectral-profile, LSI, MLSI,
curvature, sector, hypocoercive, capacity, and transport slots whenever the
displayed audit and projection comparisons are supplied in that slot's
declared domain and normalization.

\end{theorem}

\begin{proof}

For \(f\in\mathscr D_p\), subtract the two deviation bounds from
\cref{eq-generic-empirical-functional-margin}; the error margin cancels and
gives \cref{eq-generic-functional-inequality-form} with
\(\kappa=\underline\kappa_p\).  Insert this restricted inequality and
\cref{eq-generic-functional-projection-stability} into
\cref{eq-generic-functional-completeness}.  The same reciprocal calculation
as in \cref{thm-empirical-poincare-bracket-full-transfer} proves
\cref{eq-generic-full-functional-constant}.  Finally,
\(\mathcal D(f)\le
\widehat{\mathcal D}(f)+\epsilon_{\mathcal D,p}\mathcal N(f)\) and
\(\mathcal V(f)\ge
\widehat{\mathcal V}(f)-\epsilon_{\mathcal V,p}\mathcal N(f)>0\).
The ratio \(\mathcal D(f)/\mathcal V(f)\) is an admissible test in the
infimum defining \(\kappa_\Phi\), which proves
\cref{eq-generic-functional-upper-test}.

\end{proof}

\subsection{Empirical certification of exact and quasi-lumpability}
\label{subsec-empirical-lumpability-certification}

\begin{definition}[Fiber transition diameter]
\label{def-fiber-transition-diameter}

Let \(P\) be a represented Markov kernel on a finite carrier \(X\), and let
\(q:X\to Q\) be a represented terminal-compatible quotient.  Define
\begin{align}
\label{eq-quotient-row-kernel}
P_q(x,c)&=P(x,q^{-1}(c)),\\
\label{eq-fiber-transition-diameter}
\delta_q(P)
&=\max_{\substack{x,x'\in X\\q(x)=q(x')}}
\|P_q(x,\cdot)-P_q(x',\cdot)\|_{\rm TV}.
\end{align}
For a continuous-time generator \(L\), use a represented uniformization
\(P=I+L/\omega\) and retain the rate \(\omega\) in every defect conversion.
The quotient is exactly lumpable precisely when
\(\delta_q(P)=0\).

\end{definition}

\begin{theorem}[Finite-sample lumpability certificate]
\label{thm-finite-sample-lumpability-certificate}

Suppose every \(x\in X\) has at least \(N_{\min}\) independent audited
outgoing transitions, and let \(\widehat P_q(x,\cdot)\) be the empirical
quotient-row distribution.  Put
\begin{align}
\label{eq-empirical-fiber-transition-diameter}
\widehat\delta_q
&=\max_{\substack{x,x'\in X\\q(x)=q(x')}}
\|\widehat P_q(x,\cdot)-\widehat P_q(x',\cdot)\|_{\rm TV},\\
\label{eq-lumpability-entrywise-radius}
r_q
&=\sqrt{\frac{\log(2|X||Q|/\delta_{\rm lump})}
                    {2N_{\min}}}.
\end{align}
Then, with probability at least \(1-\delta_{\rm lump}\),
\begin{equation}
\label{eq-true-lumpability-defect-bound}
\boxed{
\left|\delta_q(P)-\widehat\delta_q\right|\le|Q|r_q.
}
\end{equation}
For an implicit carrier, the same conclusion holds with
\(|Q|r_q\) replaced by \(2\Gamma_q^{\rm row}\), where
\(\Gamma_q^{\rm row}\) is any represented simultaneous total-variation
radius for the selected quotient-row class.  Its cover, feature locator,
adaptive sampling rule, and selection cost are charged.

If every represented quotient-row probability lies on a declared grid
\(\{0,h_q,2h_q,\ldots,1\}\), then
\begin{equation}
\label{eq-lumpability-grid-separation}
\widehat\delta_q+|Q|r_q<h_q
\quad\Longrightarrow\quad
\delta_q(P)=0.
\end{equation}
An exact symbolic or interval verification of all row equalities also gives
\(\delta_q(P)=0\) directly.

\end{theorem}

\begin{proof}

Hoeffding's inequality and a union bound over \(X\times Q\) give
\[
|P_q(x,c)-\widehat P_q(x,c)|\le r_q
\]
simultaneously.  Hence each empirical row is within
\(|Q|r_q/2\) in total variation of its true row.  The reverse and forward
triangle inequalities therefore bound the difference between every true
pairwise row distance and its empirical counterpart by \(|Q|r_q\).
Taking maxima preserves this uniform error bound and proves
\cref{eq-true-lumpability-defect-bound}.

Two distinct probability vectors on the common grid have total-variation
distance at least \(h_q\): one coordinate differs by at least \(h_q\), and
mass conservation supplies an equal total discrepancy of the opposite
sign.  Therefore \(\delta_q(P)\) is either zero or at least \(h_q\).
The strict upper bound in
\cref{eq-lumpability-grid-separation} forces the zero case.

\end{proof}

\begin{theorem}[Why the certified quotient is lumpable]
\label{thm-certified-fiber-diameter-lumpability-transfer}

Let \(\bar P\) be the quotient kernel obtained by choosing the common row
when \(\delta_q(P)=0\).  Then
\begin{equation}
\label{eq-exact-kernel-lumpability-intertwining}
PU_q=U_q\bar P.
\end{equation}
Consequently \(q(X_t)\) is Markov with kernel \(\bar P\), independently of
the representative inside each fiber.  For a terminal-compatible absorbing
target, the full and quotient first-hit laws agree by
\cref{prop-canonical-exact-quotient-unit-transfer}.

More generally, choose \(\bar P(c,\cdot)\) to be any convex average of the
rows \(P_q(x,\cdot)\) over \(q(x)=c\).  Then
\begin{align}
\label{eq-quasi-lumpability-operator-defect}
\|PU_q-U_q\bar P\|_{\infty\to\infty}
&\le2\delta_q(P),\\
\label{eq-quasi-lumpability-horizon-transfer}
\left|
\Pr\{\tau_T\le H\}
-\Pr\{\tau_{T_Q}^{\bar P}\le H\}
\right|
&\le\min\{1,H\delta_q(P)\}.
\end{align}
Here the quotient chain is started from the pushforward of the full chain's
initial law, and the target and quotient target are absorbing.  These are the
same initialization and terminal-compatibility conventions used by the
first-hit transfer theorem.
Under uniformization, the generator defect in
\cref{def-lumpability-defect} obeys
\begin{equation}
\label{eq-empirical-generator-lumpability-defect}
\|E_q\|_{\infty\to\infty}
\le2\omega\delta_q(P).
\end{equation}
Thus the empirical bound in
\cref{thm-finite-sample-lumpability-certificate} supplies either an exact
\(\Abs\) quotient when the defect is certified zero, or the explicit
quasi-lumpable defect input required by
\cref{quotient-defect-and-geom-form-comparison,prop-partial-lumpability-no-escape}.

\end{theorem}

\begin{proof}

If \(\delta_q(P)=0\), the value
\(P(x,q^{-1}(c'))\) depends only on \(q(x)\).  Define this value to be
\(\bar P(q(x),c')\).  For every \(f:Q\to\mathbb R\),
\[
PU_qf(x)
=\sum_{c'}P(x,q^{-1}(c'))f(c')
=U_q\bar Pf(x),
\]
which proves exact intertwining and the Markov property of the projected
process.

For the averaged row, every row in a fiber is within
\(\delta_q(P)\) in total variation of the average.  Expectations of
functions bounded by one therefore differ by at most \(2\delta_q(P)\),
proving \cref{eq-quasi-lumpability-operator-defect}.  Couple the projected
full transition and the quotient transition maximally at every step.  Their
one-step disagreement probability is at most \(\delta_q(P)\); a union bound
through \(H\) steps proves
\cref{eq-quasi-lumpability-horizon-transfer}.  Finally,
\(L=\omega(P-I)\) and \(\bar L=\omega(\bar P-I)\), so multiplying
\cref{eq-quasi-lumpability-operator-defect} by \(\omega\) proves
\cref{eq-empirical-generator-lumpability-defect}.

\end{proof}

\begin{corollary}[Canonical same-run geometry and quotient certificate]
\label{cor-canonical-same-run-geometry-quotient-certificate}

Let one canonical architecture expose the geometry feature head
\(\psi_{p_{\rm geom}}\) and quotient head \(q:X\to Q\).  Construct the
Dirichlet and quotient-row audits from independent represented samples of
the same deployed transition law.  Assume the hypotheses of
\cref{thm-empirical-poincare-bracket-full-transfer} hold with
\(\underline\gamma_p>0\).  On the simultaneous confidence event,
the true Poincar\'e constant satisfies
\begin{equation}
\label{eq-canonical-true-poincare-interval}
\left(
\frac{A_pD_p^{\rm pr}}{\underline\gamma_p}+r_p
\right)^{-1}
\le\gamma_\Phi
\le
\inf_{\substack{a\ {\rm represented}\\
a^{\mathsf T}(\widehat C_p-\epsilon_{C,p}I)a>0}}
\frac{
a^{\mathsf T}(\widehat G_p+\epsilon_{G,p}I)a
}{
a^{\mathsf T}(\widehat C_p-\epsilon_{C,p}I)a
},
\end{equation}
and the true fiber diameter satisfies
\begin{equation}
\label{eq-canonical-true-lumpability-interval}
\left[\widehat\delta_q-|Q|r_q\right]_+
\le\delta_q(P)
\le
\min\left\{1,\widehat\delta_q+|Q|r_q\right\}.
\end{equation}
The analogous interval for another functional slot is obtained from
\cref{eq-generic-full-functional-constant,eq-generic-functional-upper-test}
using the same feature head and that slot's declared forms.
The infimum in \cref{eq-canonical-true-poincare-interval} is \(+\infty\)
when no represented test vector has a positive certified denominator.

If the quotient-row probabilities lie on the represented grid of
\cref{eq-lumpability-grid-separation} and the upper endpoint in
\cref{eq-canonical-true-lumpability-interval} is smaller than \(h_q\), then
\(\delta_q(P)=0\).  Hence
\[
PU_q=U_q\bar P,
\]
the projected process is Markov, and terminal-compatible first-hit laws
transfer exactly.  Otherwise the upper endpoint is substituted for
\(\delta_q(P)\) in
\cref{eq-quasi-lumpability-operator-defect,eq-quasi-lumpability-horizon-transfer,eq-empirical-generator-lumpability-defect}.

\end{corollary}

\begin{proof}

The Poincar\'e lower and upper bounds are
\cref{eq-full-poincare-from-feature-completeness,eq-empirical-poincare-upper-test};
taking the infimum over available represented test vectors preserves the
upper bound.  The lumpability interval is
\cref{eq-true-lumpability-defect-bound} clipped to the range \([0,1]\).
Grid separation then gives zero fiber diameter, and
\cref{thm-certified-fiber-diameter-lumpability-transfer} gives exact
intertwining and first-hit transfer.  The same theorem gives the stated
quasi-lumpable substitutions.

\end{proof}

\begin{corollary}[End-to-end certified neural optimum]
\label{cor-end-to-end-certified-neural-optimum}

Consider a complete structural-Bayesian neural learner trained by a certified
WAdam record with \(p\) trainable parameters and terminal weights
\(\widehat\theta_p\).  At one common class-preserving ceiling, suppose the record
contains:
\begin{enumerate}[label=\textup{(\roman*)}]
\item the WAdam descent, PL, and objective-to-action calibration certificates
      of \cref{thm-wadam-descent-global-optimization-certificate};
\item the simultaneous structural learning radii \(\Gamma_{t,p}\), including
      their noise, dependence, selection, and confidence conversions;
\item the represented filter comparison \(\Delta_{{\rm fil},p}^{H}\);
\item the held-out arrival lower certificate \(L_{{\rm ev},p}\); and
\item the exact or coverage-based Bellman upper certificate
      \(U_{{\rm Bell},p}\).
\end{enumerate}
Then, with the declared simultaneous confidence,
\begin{align}
\label{eq-end-to-end-certified-neural-optimum}
0\le J_{\rm sat}^{*}(B)-A_{R_{\widehat\theta_p}}^H
\le\min\Biggl\{&U_{{\rm Bell},p}-L_{{\rm ev},p},\\
&\mathbb E_{\widehat R}\sum_{t<H}\mathbf1_{\{\tau>t\}}
\left(2\Gamma_{t,p}+
\lambda_{t,p}\frac{\mathcal E_{\rm WAdam}(p,K)}{\delta_{\rm opt}}
+\zeta_{t,p}\right)
+\Delta_{{\rm fil},p}^{H}\Biggr\},\notag
\end{align}
where the declared simultaneous failure probability includes
\(\delta_{\rm opt}\) from
\cref{eq-wadam-high-probability-objective-gap}.  The first candidate certifies the
complete representation-plus-training gap.  The second resolves the learned
record into statistical, optimizer, calibration, and filtering coordinates.
Every term in \cref{eq-end-to-end-certified-neural-optimum} is a coordinate
of the common parameter-explicit ledger \(\mathfrak E_p\) in
\cref{eq-parameter-explicit-implementation-ledger}; parameter count is
therefore not hidden inside an unspecified approximation term.

At the universal Bayesianized ceiling of
\cref{cor-universal-represented-computation-bayesian-ceiling},
\(J_{\rm sat}^{*}\) equals the saturated represented-computation optimum.
Thus a small certified primal--dual gap proves that the implemented learner
attains that universal ceiling to the displayed accuracy.

\end{corollary}

\begin{proof}

Use \cref{eq-wadam-high-probability-objective-gap} and
\cref{eq-wadam-objective-to-action-calibration} in the learned-control
candidate of
\cref{cor-primal-dual-empirical-implementation-certificate}.  The Bellman
candidate follows from the same corollary.  Finally apply the profile equality
in \cref{cor-universal-represented-computation-bayesian-ceiling} at the common
compiled ceiling.

\end{proof}

\begin{theorem}[Complete empirical structural proof package]
\label{thm-complete-empirical-structural-proof-package}

Fix a parameter count \(p\) and one complete implementation record.  Suppose
the record contains:
\begin{enumerate}[label=\textup{(\roman*)}]
\item the saturated, \(p\)-class, and trained-run intervals in
      \cref{eq-saturated-optimum-confidence-interval,eq-p-class-optimum-confidence-interval,eq-trained-run-confidence-interval};
\item the source-resolved audit of
      \cref{def-audited-parameter-source-ledger};
\item the primal--dual and WAdam certificates in
      \cref{cor-end-to-end-certified-neural-optimum};
\item the canonical loss calibration and bias--variance ledger in
      \cref{thm-canonical-loss-saturated-arrival-calibration,thm-complete-certified-bias-variance-decomposition};
\item a feature-restricted functional-inequality audit and its represented
      completeness comparison; and
\item a represented quotient \(q\), its quotient-row audit, and either its
      exact grid/separation record or its quasi-lumpability defect record.
\end{enumerate}
Allocate confidence so that the sum of all failure probabilities is at most
\(\delta_{\rm all}\).  Then, with probability at least
\(1-\delta_{\rm all}\), the same empirical run simultaneously certifies:
\begin{enumerate}[label=\textup{(\alph*)}]
\item the implementation gap in
      \cref{eq-end-to-end-certified-neural-optimum};
\item the calibrated primal--dual loss and separated deterministic-bias and
      stochastic-radius coordinates in
      \cref{eq-canonical-certificate-loss-optimality,eq-complete-certified-bias-variance-decomposition};
\item the underfitting, training, and minimum-parameter intervals in
      \cref{eq-certified-underfitting-interval,eq-certified-training-error-interval,eq-certified-minimum-parameter-bracket};
\item the exact learned source support and the source-resolved target-gain
      bound in
      \cref{eq-parameter-source-resolved-empirical-target-gain};
\item every audited analytical functional constant transferred by
      \cref{thm-empirical-poincare-bracket-full-transfer,thm-generic-empirical-functional-inequality-transfer};
      and
\item exact quotient dynamics and exact target-law transfer when
      \cref{eq-lumpability-grid-separation} holds, or the quantitative
      quasi-lumpability bounds
      \cref{eq-quasi-lumpability-operator-defect,eq-quasi-lumpability-horizon-transfer,eq-empirical-generator-lumpability-defect}
      otherwise.
\end{enumerate}
All conclusions use the same represented architecture, source graph,
transition law, quotient, feature map, precision, and common saturated
ceiling.  Therefore the empirical proof is a certificate about the true
analytical record rather than a collection of unrelated fitted diagnostics.

\end{theorem}

\begin{proof}

Each cited theorem holds on its represented audit event.  The union bound
gives the simultaneous probability.  All conversions are fields of the same
complete record, so their common-law and common-normalization hypotheses are
preserved.  Applying the cited conclusions on the intersection proves
\textup{(a)}--\textup{(f)}.

\end{proof}

\begin{figure}[tbp]
\centering
\resizebox{.98\linewidth}{!}{%
\begin{tikzpicture}[fw diagram,node distance=8mm and 10mm]
  \node[fw source,text width=2.7cm] (sat)
    {saturated represented-\\ computation optimum};
  \node[fw provenance,text width=2.8cm,right=of sat] (dual)
    {verified Bellman\\ upper certificate};
  \node[fw use,text width=2.8cm,below=10mm of dual] (train)
    {network, WAdam,\\ learning and filter errors};
  \node[fw geom,text width=2.8cm,right=of train] (policy)
    {deployed pre-hit\\ neural record};
  \node[fw terminal,text width=2.8cm,above=10mm of policy] (lower)
    {held-out arrival\\ lower certificate};
  \node[fw box,text width=2.9cm,right=of policy] (gap)
    {certified implementation\\ gap \(U_{\rm Bell}-L_{\rm ev}\)};
  \draw[fw arrow] (sat)--(dual);
  \draw[fw arrow] (dual)--(gap);
  \draw[fw arrow] (train)--(policy);
  \draw[fw arrow] (policy)--(lower);
  \draw[fw arrow] (lower)--(gap);
  \draw[fw arrow,dashed] (dual)--(train);
\end{tikzpicture}%
}
\caption{Certified neural realization of the saturated optimum.  The deployed
network supplies a primal lower certificate; a represented Bellman barrier
supplies an upper certificate for the complete saturated class.  Structural
generalization, WAdam optimization, filtering, precision, and deployment are
charged inside the common record.}
\label{fig-certified-neural-saturated-optimum}
\end{figure}

\begin{figure}[tbp]
\centering
\resizebox{.98\linewidth}{!}{%
\begin{tikzpicture}[fw diagram,node distance=8mm and 10mm]
  \node[fw source,text width=2.6cm] (run)
    {complete empirical\\training record};
  \node[fw provenance,text width=2.7cm,right=of run] (param)
    {parameter profile\\and underfitting};
  \node[fw use,text width=2.7cm,above right=8mm and 10mm of param] (source)
    {five-source ledger\\and interventions};
  \node[fw geom,text width=2.7cm,below right=8mm and 10mm of param] (fi)
    {functional constants\\and completeness};
  \node[fw box,text width=2.7cm,right=16mm of param] (lump)
    {exact or quasi-\\lumpability};
  \node[fw terminal,text width=2.8cm,right=of lump] (proof)
    {simultaneous analytical\\proof certificate};
  \draw[fw arrow] (run)--(param);
  \draw[fw arrow] (param)--(source);
  \draw[fw arrow] (param)--(fi);
  \draw[fw arrow] (source)--(lump);
  \draw[fw arrow] (fi)--(lump);
  \draw[fw arrow] (lump)--(proof);
\end{tikzpicture}%
}
\caption{Empirical structural proof package.  Parameter sufficiency,
source-resolved learning, analytical functional inequalities, and quotient
lumpability are certified from one completely charged record and one
simultaneous confidence allocation.}
\label{fig-complete-empirical-structural-proof}
\end{figure}

\section{Black-Box Structural Auditing of Trained Predictors}
\label{sec-black-box-structural-audit}

The preceding section certifies a represented training and deployment record.
This section treats the complementary setting in which the trained predictor
is available only through a query interface.  The audit compiles oracle calls,
a represented test law, label-channel information, and public interventions
into the same statistical, structural, geometric, and quotient coordinates
used throughout the framework.  Its conclusions concern the behavior of the
deployed predictor on the declared audit law.  A statement about the
unobserved training sample is produced precisely when the record contains a
represented transfer from the audit channel to the training channel.

\subsection{Oracle access and complete audit records}
\label{subsec-black-box-oracle-record}

\begin{definition}[Complete black-box predictor audit]
\label{def-complete-black-box-predictor-audit}

Let \(X\) be a represented input carrier and
\(Y=\{1,\ldots,K\}\) a finite label carrier.  A trained predictor is exposed
through one of the following represented interfaces:
\begin{enumerate}[label=\textup{(\roman*)}]
\item a label interface \(O_f(x)\in Y\);
\item a score interface \(O_f(x)\in[-M,M]^K\) at declared precision; or
\item a probability interface \(O_f(x)\in\Delta_K\), deterministic or
      sampled using represented oracle randomness.
\end{enumerate}
For label and score access, a represented readout maps the oracle output to a
probability vector \(p_f(x)\in\Delta_K\); a label is represented by its point
mass.  The readout and its numerical precision are part of the audit.

A \emph{complete black-box predictor audit} is a record
\begin{equation}
\label{eq-complete-black-box-record}
R_{\rm BB}
=
(O_f,\mathsf{Access},\mu,S_{\rm audit},\mathscr T,
  \mathscr D,Q,\boldsymbol\delta,\mathsf{Select}),
\end{equation}
where:
\begin{enumerate}[label=\textup{(\alph*)}]
\item \(\mu\) is the declared audit law and
      \(S_{\rm audit}=(Z_i)_{i=1}^N\) consists of independent audit units
      \(Z_i\sim\mu\); an audit unit contains its input and every clean,
      corrupted, repeated, or subgroup label declared by the protocol;
\item \(\mathscr T\) is a finite represented library of input
      interventions and perturbation kernels, with their coupling to the
      base input and their declared structural provenance;
\item \(\mathscr D\) is the represented family of diagnostics, losses,
      discretizers, feature maps, and quotient-row tests evaluated by the
      audit;
\item \(Q\) is the total number of oracle calls, including calls used for
      transformed inputs, repeated stochastic evaluation, diagnostic
      selection, and numerical calibration;
\item \(\boldsymbol\delta\) allocates the total confidence failure among all
      selected certificates; and
\item \(\mathsf{Select}\) stores every data-dependent diagnostic, feature,
      perturbation, binning, quotient, stopping decision, and comparison
      reported by the audit.
\end{enumerate}
All calls made from one \(Z_i\), including adaptive calls whose next input is
a represented function of earlier outputs in that block, form one bounded
audit block.  Blocks use independent audit examples and independent oracle
randomness.  The complete cost is
\begin{equation}
\label{eq-complete-black-box-audit-cost}
B_{\rm BB}
=B_{\rm access}\oplus B_{\rm data}\oplus B_{\rm query}
 \oplus B_{\rm intervention}\oplus B_{\rm diagnostic}
 \oplus B_{\rm selection}\oplus B_{\rm precision}
 \oplus B_{\rm certificate}.
\end{equation}
The black-box record reveals the represented input--output behavior of the
oracle.  Internal parameters, activations, source graph, optimizer state, and
training sample enter the record only when separately exposed.

\end{definition}

\begin{definition}[Test-law behavioral equivalence]
\label{def-black-box-behavioral-equivalence}

Two trained predictors \(f\) and \(g\) are \emph{behaviorally equivalent for
an audit record} when every transcript produced by the declared access
interface and every intervention in \(\mathscr T\) have the same law under
\(\mu\).  A black-box certificate is a function of this equivalence class.
It may identify a source in the internal training record only when an exposed
provenance map compiles the observed behavior into that source coordinate.

\end{definition}

\begin{theorem}[Finite-query simultaneous black-box audit]
\label{thm-finite-query-simultaneous-black-box-audit}

Let \(D_1,\ldots,D_M\) be diagnostics fixed before the audit sample is
revealed.  Suppose diagnostic \(j\) produces one block value
\(D_j(Z_i,O_f,U_i)\in[0,1]\), after any finite within-block adaptive query
sequence, and put
\begin{equation}
\label{eq-black-box-diagnostic-population-empirical}
d_j=\mathbb E D_j(Z_i,O_f,U_i),
\qquad
\widehat d_j=\frac1N\sum_{i=1}^N D_j(Z_i,O_f,U_i).
\end{equation}
For \(0<\delta_{\rm BB}<1\), define
\begin{equation}
\label{eq-black-box-simultaneous-radius}
r_{N,M,\delta_{\rm BB}}
=\sqrt{\frac{\log(2M/\delta_{\rm BB})}{2N}}.
\end{equation}
Then, with probability at least \(1-\delta_{\rm BB}\),
\begin{equation}
\label{eq-black-box-simultaneous-diagnostic-bound}
|d_j-\widehat d_j|
\le r_{N,M,\delta_{\rm BB}}
\qquad(1\le j\le M)
\end{equation}
simultaneously.  In particular, precision \(r_{N,M,\delta}\le\epsilon\)
is ensured by
\begin{equation}
\label{eq-black-box-audit-sample-complexity}
N\ge\frac{\log(2M/\delta_{\rm BB})}{2\epsilon^2}.
\end{equation}

If the diagnostic family is selected using the audit data, its fixed-family
cardinality in \cref{eq-black-box-simultaneous-radius} is replaced by the
applicable structural Rademacher, PAC--Bayes mixture, sequential Rademacher,
or martingale PAC--Bayes radius already carried by the complete record.  If a
diagnostic lies in \([a_j,b_j]\), its radius is multiplied by \(b_j-a_j\).

\end{theorem}

\begin{proof}
For each fixed \(j\), the block variables are independent and lie in
\([0,1]\), even though the calls inside one block may be adaptive.
Hoeffding's inequality gives
\[
\Pr\{|\widehat d_j-d_j|>r\}\le2e^{-2Nr^2}.
\]
At the radius in \cref{eq-black-box-simultaneous-radius}, the right side is
\(\delta_{\rm BB}/M\).  A union bound proves
\cref{eq-black-box-simultaneous-diagnostic-bound}; solving for \(N\) proves
\cref{eq-black-box-audit-sample-complexity}.  Rescaling proves the bounded
interval statement.  The selected-family replacements are exactly the
uniform and posterior-weighted selection theorems of
\cref{def-complete-source-resolved-learning-ledger,%
thm-sequential-learned-operator-cover,%
thm-martingale-structural-pac-bayes}.
\end{proof}

\subsection{Predictive performance, calibration, and law transfer}
\label{subsec-black-box-performance-calibration}

\begin{definition}[Canonical black-box predictive diagnostics]
\label{def-canonical-black-box-predictive-diagnostics}

For a clean audit label \(Y\), a reported label
\(\widehat Y_f=\arg\max_c p_f^c(X)\), and a fixed finite partition
\(\mathcal B\) of \([0,1]\), define
\begin{align}
\label{eq-black-box-clean-error}
R_0(f)&=\Pr\{\widehat Y_f\ne Y\},\\
\label{eq-black-box-confusion-entry}
J_{ac}(f)&=\Pr\{Y=a,\widehat Y_f=c\},\\
\label{eq-black-box-brier-risk}
R_{\rm Br}(f)&=\frac12\mathbb E
 \sum_{c=1}^K(p_f^c(X)-\mathbf1_{\{Y=c\}})^2,\\
\label{eq-black-box-binned-calibration}
\operatorname{Cal}_{\mathcal B}(f)
&=\sum_{B\in\mathcal B}\Pr\{C_f\in B\}
\left|
\mathbb E[\mathbf1_{\{\widehat Y_f=Y\}}\mid C_f\in B]
-\mathbb E[C_f\mid C_f\in B]
\right|,
\end{align}
where \(C_f=\max_c p_f^c(X)\).  Empty calibration bins contribute zero.
For \(\epsilon_{\rm pr}\in(0,1/K]\), the clipped normalized log loss is
\begin{equation}
\label{eq-black-box-normalized-log-risk}
R_{\log,\epsilon_{\rm pr}}(f)
=\mathbb E
\frac{-\log(p_f^Y(X)\vee\epsilon_{\rm pr})}
     {\log(1/\epsilon_{\rm pr})}.
\end{equation}
All displayed losses lie in \([0,1]\).  Classwise and subgroup risks use the
same definitions conditionally on a represented cell; their radius uses the
number of independent audit examples in that cell.

\end{definition}

\begin{corollary}[Simultaneous predictive audit]
\label{cor-simultaneous-black-box-predictive-audit}

Fix the label tie rule, probability readout, calibration partition, clipping
level, and a finite collection of represented subgroups before evaluation.
Empirical versions of
\cref{eq-black-box-clean-error,eq-black-box-confusion-entry,%
eq-black-box-brier-risk,eq-black-box-normalized-log-risk}
are within \(r_{N,M,\delta}\) of their population values simultaneously,
with one diagnostic allocated to every reported scalar.  Moreover, if
\(\widehat p_B=N_B/N\), \(\widehat a_B\) is the empirical correct mass in
bin \(B\), and \(\widehat c_B\) is its empirical confidence mass, then
\begin{equation}
\label{eq-black-box-calibration-confidence}
\operatorname{Cal}_{\mathcal B}(f)
\le
\sum_{B\in\mathcal B}|\widehat a_B-\widehat c_B|
+2|\mathcal B|r_{N,M,\delta},
\end{equation}
where the empirical masses are normalized by \(N\), so empty bins require no
conditional division.

\end{corollary}

\begin{proof}
Each risk and confusion entry is the expectation of a fixed \([0,1]\)-valued
diagnostic, so \cref{thm-finite-query-simultaneous-black-box-audit} applies.
For calibration, write the contribution of bin \(B\) as
\( |a_B-c_B|\), where
\[
a_B=\mathbb E[\mathbf1_{\{C_f\in B\}}
                    \mathbf1_{\{\widehat Y_f=Y\}}],
\qquad
c_B=\mathbb E[\mathbf1_{\{C_f\in B\}}C_f].
\]
Both summands are bounded diagnostics.  The reverse triangle inequality
gives \(|a_B-c_B|\le|\widehat a_B-\widehat c_B|+2r\); summing over bins proves
\cref{eq-black-box-calibration-confidence}.
\end{proof}

\begin{theorem}[Audit-law to deployment-law transfer]
\label{thm-black-box-audit-deployment-transfer}

Let \(\nu\) be a represented deployment law and let \(D\in[0,1]\) be one
black-box diagnostic.  Either of the following represented comparisons
transfers its audit-law certificate:
\begin{enumerate}[label=\textup{(\roman*)}]
\item if \(\|\nu-\mu\|_{\rm TV}\le\Delta_{\rm shift}\), then
\begin{equation}
\label{eq-black-box-tv-law-transfer}
\mathbb E_\nu D
\le\widehat{\mathbb E}_\mu D+r_{N,M,\delta}+\Delta_{\rm shift};
\end{equation}
\item if \(w=d\nu/d\mu\) is represented and
      \(0\le w\le W\), then the importance-weighted estimator satisfies
\begin{equation}
\label{eq-black-box-importance-law-transfer}
\mathbb E_\nu D
\le\frac1N\sum_{i=1}^Nw(Z_i)D_i
+W\sqrt{\frac{\log(2M/\delta)}{2N}}
\end{equation}
simultaneously over \(M\) fixed diagnostics.
\end{enumerate}

\end{theorem}

\begin{proof}
The total-variation dual inequality gives
\(|\mathbb E_\nu D-\mathbb E_\mu D|\le\Delta_{\rm shift}\), which combines
with \cref{eq-black-box-simultaneous-diagnostic-bound}.  In the second case,
\(\mathbb E_\mu[wD]=\mathbb E_\nu D\) and \(wD\in[0,W]\); the rescaled
Hoeffding and union-bound argument proves
\cref{eq-black-box-importance-law-transfer}.
\end{proof}

\subsection{Noise identification and clean-risk reconstruction}
\label{subsec-black-box-noise-identification}

\begin{definition}[Audit, effective, and training noise]
\label{def-black-box-three-noise-coordinates}

The black-box record keeps three typed noise coordinates.
\begin{enumerate}[label=\textup{(\roman*)}]
\item \emph{Audit-label noise} is the represented channel from a clean audit
      label \(Y\) to the observed audit label \(\widetilde Y\).
\item \emph{Effective task noise} is the disagreement channel between the
      observed label and a represented clean target on the audit law.
\item \emph{Training noise} is the corresponding channel in the sample used
      to construct the trained oracle.
\end{enumerate}
The first two are identified from paired or repeated audit labels, a known
channel, or another represented reconstruction in
\cref{def-complete-dynamical-model-channel-noise-record}.  The third is reported as a
compatible set unless the record supplies a represented training-to-audit
channel comparison.

\end{definition}

\begin{theorem}[Binary repeated-label noise interval]
\label{thm-black-box-binary-repeated-label-noise}

Let
\(\widetilde Y^{(1)}=Y\Xi_1\) and
\(\widetilde Y^{(2)}=Y\Xi_2\), where, conditionally on \((X,Y)\), the signs
\(\Xi_1,\Xi_2\) are independent and satisfy
\(\Pr\{\Xi_j=-1\}=\eta\le1/2\).  Put
\begin{equation}
\label{eq-black-box-repeated-label-disagreement}
d=\Pr\{\widetilde Y^{(1)}\ne\widetilde Y^{(2)}\}
=2\eta(1-\eta)
\end{equation}
and let \(\widehat d\) be the disagreement frequency in \(N\) independent
pairs.  With
\begin{align}
\label{eq-black-box-repeated-label-radius}
r_d&=\sqrt{\frac{\log(2/\delta_d)}{2N}},\\
\label{eq-black-box-repeated-label-clipped-endpoints}
d_-&=[\widehat d-r_d]_{[0,1/2]},
&d_+&=[\widehat d+r_d]_{[0,1/2]},\\
\label{eq-black-box-noise-inverse-map}
h(d)&=\frac{1-\sqrt{1-2d}}2,
\end{align}
one has, with probability at least \(1-\delta_d\),
\begin{equation}
\label{eq-black-box-repeated-label-noise-interval}
h(d_-)\le\eta\le h(d_+).
\end{equation}
Combining either endpoint with the exact risk transport in
\cref{lem-exact-label-noise-risk-transport} gives a simultaneous clean-risk
interval for every fixed audited predictor.

\end{theorem}

\begin{proof}
Conditional independence gives
\(d=\eta(1-\eta)+(1-\eta)\eta\).  Hoeffding's inequality gives
\(|d-\widehat d|\le r_d\).  The map \(\eta\mapsto2\eta(1-\eta)\) is increasing
on \([0,1/2]\), and \(h\) is its increasing inverse, proving the interval.
The risk statement follows by intersection with the predictive audit event
and substitution in \cref{eq-homogeneous-label-noise-identity}.
\end{proof}

\begin{corollary}[Joint binary noise and clean-risk confidence set]
\label{cor-black-box-joint-noise-clean-risk}

Suppose a noisy-label audit gives simultaneous intervals
\begin{equation}
\label{eq-black-box-noisy-risk-noise-rectangle}
R_\eta(f)\in[\widetilde r_-,\widetilde r_+],
\qquad
\eta\in[\eta_-,\eta_+]\subset[0,1/2).
\end{equation}
Define
\begin{equation}
\label{eq-black-box-clean-risk-compatible-set}
\mathcal I_0(f)
=
\left\{
\frac{u-v}{1-2v}:
u\in[\widetilde r_-,\widetilde r_+],\
v\in[\eta_-,\eta_+]
\right\}\cap[0,1].
\end{equation}
On the joint confidence event, \(R_0(f)\in\mathcal I_0(f)\).  Its reported
interval is the closed hull
\begin{equation}
\label{eq-black-box-clean-risk-compatible-hull}
\left[
\inf\mathcal I_0(f),\sup\mathcal I_0(f)
\right],
\end{equation}
whose endpoints are obtained by evaluating the rational map at the four
corners of the rectangle and clipping the resulting compatible values to
\([0,1]\).  If the noise confidence set
contains \(1/2\), the certified default is \([0,1]\).

\end{corollary}

\begin{proof}
The exact identity
\(R_\eta(f)=\eta+(1-2\eta)R_0(f)\) from
\cref{eq-homogeneous-label-noise-identity} gives
\(R_0(f)=(R_\eta(f)-\eta)/(1-2\eta)\) whenever \(\eta<1/2\).
Substitution of the simultaneous rectangle proves membership in
\(\mathcal I_0(f)\).  The rectangle is compact and the denominator is
positive, so the image is compact.  For fixed \(v\), the map is increasing in
\(u\); at either endpoint in \(u\), its derivative with respect to \(v\) has
constant sign \((2u-1)/(1-2v)^2\).  Its extrema therefore occur at corners.
At \(\eta=1/2\), the corrupted risk is independent of the clean
risk, yielding the full default interval.
\end{proof}

\begin{theorem}[Multiclass channel inversion with estimated noise]
\label{thm-black-box-multiclass-noise-inversion}

Let
\(C_{ab}=\Pr\{\widetilde Y=b\mid Y=a\}\) be a class-conditional corruption
matrix and let
\begin{align}
J_{ac}&=\Pr\{Y=a,\widehat Y_f=c\},&
\widetilde J_{bc}&=\Pr\{\widetilde Y=b,\widehat Y_f=c\}.
\end{align}
Then
\begin{equation}
\label{eq-black-box-multiclass-joint-transport}
\widetilde J=C^{\mathsf T}J.
\end{equation}
Let \(\widehat C\) be an invertible represented estimate and let
\(\widehat{\widetilde J}\) be the empirical noisy-label confusion matrix.
Define
\begin{equation}
\label{eq-black-box-multiclass-inverted-joint}
\widehat J=(\widehat C^{\mathsf T})^{-1}
             \widehat{\widetilde J}.
\end{equation}
On any simultaneous event satisfying
\begin{align}
\|\widehat{\widetilde J}-\widetilde J\|_1&\le\epsilon_J,
&
\|\widehat C^{\mathsf T}-C^{\mathsf T}\|_{1\to1}&\le\epsilon_C,
\end{align}
one has
\begin{equation}
\label{eq-black-box-multiclass-inversion-error}
\|\widehat J-J\|_1
\le
\|(\widehat C^{\mathsf T})^{-1}\|_{1\to1}
\,(\epsilon_J+\epsilon_C).
\end{equation}
Consequently the clean classification error obeys
\begin{equation}
\label{eq-black-box-multiclass-clean-risk-interval}
\left|
R_0(f)-\left(1-\sum_{a=1}^K\widehat J_{aa}\right)
\right|
\le
\|(\widehat C^{\mathsf T})^{-1}\|_{1\to1}
\,(\epsilon_J+\epsilon_C).
\end{equation}
The same linear transport applies to every bounded loss that is a represented
linear functional of the clean joint table.  The inverse norm records the
exact noise-amplification cost.

\end{theorem}

\begin{proof}
Condition on \((Y,\widehat Y_f)\) to obtain
\cref{eq-black-box-multiclass-joint-transport}.  Multiplication of
\cref{eq-black-box-multiclass-inverted-joint} by
\(\widehat C^{\mathsf T}\), followed by subtraction of
\(\widehat C^{\mathsf T}J\), gives
\[
\widehat C^{\mathsf T}(\widehat J-J)
=(\widehat{\widetilde J}-\widetilde J)
 +(C^{\mathsf T}-\widehat C^{\mathsf T})J.
\]
Since \(J\) is a probability table, \(\|J\|_1=1\).  Apply the induced norm
and the inverse to prove \cref{eq-black-box-multiclass-inversion-error}.
The diagonal functional has dual norm at most one, proving the risk interval.
\end{proof}

\begin{theorem}[Black-box nonidentifiability and certified training-noise transfer]
\label{thm-black-box-training-noise-identifiability}

Fix a black-box predictor oracle, audit law, and query protocol.  Without a
represented relation to the training record, its training-noise rate is not
identified: for every \(\eta_{\rm tr}\in[0,1]\), there is a training record
with that label-noise rate that returns the same oracle and hence the same
black-box audit transcript law.

Let \(\mathcal C_{\rm tr}(R_{\rm BB})\) denote the set of training channels
compatible with the exposed record.  If the record additionally supplies a
represented channel comparison
\begin{equation}
\label{eq-black-box-training-audit-channel-comparison}
\|C_{\rm tr}-C_{\rm audit}\|_{1\to1}\le\Delta_{\rm tr,aud},
\end{equation}
then every confidence set \(\mathfrak C_{\rm audit}\) for the audit channel
induces the training-channel confidence set
\begin{equation}
\label{eq-black-box-training-channel-confidence-set}
\mathfrak C_{\rm tr}
=\left\{
C:\inf_{C'\in\mathfrak C_{\rm audit}}
\|C-C'\|_{1\to1}\le\Delta_{\rm tr,aud}
\right\}.
\end{equation}
For binary homogeneous noise this reduces to
\begin{equation}
\label{eq-black-box-training-noise-transfer-interval}
\eta_{\rm tr}
\in
\bigl[
(\eta_{\rm audit}^- -\Delta_\eta)_+,
(\eta_{\rm audit}^+ +\Delta_\eta)\wedge1
\bigr]
\end{equation}
when the represented comparison is
\(|\eta_{\rm tr}-\eta_{\rm audit}|\le\Delta_\eta\).

\end{theorem}

\begin{proof}
For a fixed exposed oracle, consider a represented training compiler that
returns that oracle after reading and charging its input training record.
Feeding the compiler clean labels gives rate zero.  Passing the same labels
through any declared corruption channel before the compiler reads them gives
the prescribed rate and leaves the returned oracle unchanged.  Therefore all
black-box transcripts are identical, proving nonidentifiability from the
oracle and independent audit data alone.

On the event \(C_{\rm audit}\in\mathfrak C_{\rm audit}\),
\cref{eq-black-box-training-audit-channel-comparison} places
\(C_{\rm tr}\) in \cref{eq-black-box-training-channel-confidence-set} by
definition.  The binary interval is the triangle inequality followed by
intersection with \([0,1]\).
\end{proof}

\subsection{Behavioral structural channels and interventions}
\label{subsec-black-box-structural-channels}

\begin{definition}[Black-box structural intervention ledger]
\label{def-black-box-structural-intervention-ledger}

Let
\begin{equation}
\label{eq-black-box-structural-index}
\mathfrak A_{\rm BB}
=\{\mathsf{Add},\mathsf{Mult},\mathsf{Ord},\mathsf{Sym},
    \mathsf{Filt},\Geom,\Caus,\Lift,\Bdry\}.
\end{equation}
For every audited source \(a\in\mathfrak A_{\rm BB}\), a represented
intervention record contains:
\begin{enumerate}[label=\textup{(\roman*)}]
\item a coupling \((X,T_aX)\) on the audit law, generated without access to
      the current clean audit label unless the intervention is explicitly a
      label-conditioned analytic diagnostic;
\item an output action
      \(\rho_a(T_a):\Delta_K\to\Delta_K\), equal to the identity for an
      invariant intervention and to the represented label action for an
      equivariant intervention;
\item a target-preservation flag stating whether the clean label is unchanged
      or transformed by the same represented action; and
\item the query, construction, locator, selection, and numerical costs of the
      intervention.
\end{enumerate}
With total variation on \(\Delta_K\), define the behavioral defect and, when
the transformed clean label \(Y_a\) is represented, the target response
\begin{align}
\label{eq-black-box-source-behavioral-defect}
D_a(f)
&=\mathbb E\left[
\|p_f(T_aX)-\rho_a(T_a)p_f(X)\|_{\rm TV}
\right],\\
\label{eq-black-box-source-target-response}
A_a(f)
&=\mathbb E\left[
\ell(p_f(X),Y)-\ell(p_f(T_aX),Y_a)
\right].
\end{align}
The loss is represented and lies in \([0,1]\), so
\(A_a(f)\in[-1,1]\).  A filtration intervention uses two represented
prefixes that differ only in the audited reveal; a lift intervention compares
the base prediction with the represented pushforward of its lifted prediction;
a boundary intervention changes only the declared boundary record.  A Caus
intervention additionally contains a prospective transition rule compiled
from oracle outputs available before the label reveal.  Its target response
is therefore a pre-reveal behavioral alignment coordinate.

The ledger is behavioral: \(D_a\) and \(A_a\) describe dependence of the
deployed input--output map on a represented source intervention.  An exact
claim about internal source ancestry requires an exposed compiler from the
model record to the source-resolved ledger of
\cref{def-complete-source-resolved-learning-ledger}.

\end{definition}

\begin{theorem}[Simultaneous black-box structural signature]
\label{thm-simultaneous-black-box-structural-signature}

Fix \(L_a\) intervention draws per audit example for every
\(a\in\mathfrak A_{\rm BB}\), and average their bounded defects inside the
example block.  Let \(\widehat D_a\) and \(\widehat A_a\) be the resulting
empirical averages over \(N\) independent audit examples.  For
\begin{equation}
\label{eq-black-box-structural-signature-radius}
r_{\rm str}
=\sqrt{\frac{\log(4|\mathfrak A_{\rm BB}|/\delta_{\rm str})}{2N}},
\end{equation}
one has, with probability at least \(1-\delta_{\rm str}\),
\begin{align}
\label{eq-black-box-source-defect-interval}
D_a(f)&\in
\bigl[[\widehat D_a-r_{\rm str}]_+,
      (\widehat D_a+r_{\rm str})\wedge1\bigr],\\
\label{eq-black-box-source-alignment-interval}
A_a(f)&\in
\bigl[
\widehat A_a-2r_{\rm str},
\widehat A_a+2r_{\rm str}
\bigr]\cap[-1,1]
\end{align}
simultaneously for every source.  The complete query count of this audit is
\begin{equation}
\label{eq-black-box-structural-signature-query-count}
Q_{\rm str}
=N\sum_{a\in\mathfrak A_{\rm BB}}L_a(q_a^{\rm base}+q_a^{\rm trans}),
\end{equation}
where repeated base calls may be shared only when the oracle and coupling
record justify the shared transcript.

For a source family chosen from data, replace \(r_{\rm str}\) by its charged
source-resolved Rademacher or PAC--Bayes radius.  If a provenance compiler is
present, the corresponding exact source cell and target-gain upper bound are
then obtained from
\cref{thm-parameter-source-resolved-empirical-target-gain}.

\end{theorem}

\begin{proof}
The averaged intervention defect in one example block lies in \([0,1]\).
The two loss values defining the target response are separately in
\([0,1]\).  Apply
\cref{thm-finite-query-simultaneous-black-box-audit} to the defect and the two
loss expectations for every source.  The defect uses one radius; subtraction
of the two loss intervals uses two.  Summing the represented calls proves
\cref{eq-black-box-structural-signature-query-count}.  The selected-source
and compiled-provenance conclusions are the cited framework theorems.
\end{proof}

\begin{definition}[Black-box input-intervention value]
\label{def-black-box-input-intervention-value}

Let \(\mathcal A\subseteq\mathfrak A_{\rm BB}\) be a finite set of \(d\)
audited input sources.  A complete intervention family supplies, for every
\(S\subseteq\mathcal A\), a represented coupled input \(X^S\) in which the
sources in \(S\) are retained and the others are replaced by their declared
target-neutral baselines.  All \(2^d\) coupled constructions use the same
audit law, label action, oracle, loss, and numerical readout.  Define
\begin{align}
\label{eq-black-box-input-intervention-value}
v_f(S)&=1-\mathbb E\ell(p_f(X^S),Y^S),\\
\label{eq-black-box-input-shapley-value}
\phi_{f,a}
&=\sum_{S\subseteq\mathcal A\setminus\{a\}}
\frac{|S|!\,(d-|S|-1)!}{d!}
\bigl(v_f(S\cup\{a\})-v_f(S)\bigr).
\end{align}
These are input-behavior allocations.  They coincide with internal source
interventions only when the exposed model compiler proves that equivalence.

\end{definition}

\begin{theorem}[Simultaneous black-box intervention allocation]
\label{thm-simultaneous-black-box-intervention-allocation}

Evaluate every \(X^S\) on \(N_{\rm int}\) independent coupled audit units and
put
\begin{equation}
\label{eq-black-box-intervention-radius}
r_{\rm int}^{\rm BB}
=\sqrt{\frac{\log(2^{d+1}/\delta_{\rm int})}{2N_{\rm int}}}.
\end{equation}
Then, with probability at least \(1-\delta_{\rm int}\),
\begin{align}
\label{eq-black-box-intervention-value-interval}
|\widehat v_f(S)-v_f(S)|&\le r_{\rm int}^{\rm BB}
&&\text{for every }S\subseteq\mathcal A,\\
\label{eq-black-box-intervention-shapley-interval}
|\widehat\phi_{f,a}-\phi_{f,a}|&\le2r_{\rm int}^{\rm BB}
&&\text{for every }a\in\mathcal A,\\
\label{eq-black-box-intervention-shapley-efficiency}
\sum_{a\in\mathcal A}\phi_{f,a}
&=v_f(\mathcal A)-v_f(\varnothing).
\end{align}
The full deterministic-oracle query count is
\(Q_{\rm int}=N_{\rm int}2^d\); stochastic repetition multiplies this by the
represented number of calls per intervention.

\end{theorem}

\begin{proof}
Hoeffding's inequality with a union bound over \(2^d\) bounded losses proves
\cref{eq-black-box-intervention-value-interval}.  The Shapley weights for one
source are nonnegative and sum to one; every marginal contains two estimated
values, giving \cref{eq-black-box-intervention-shapley-interval}.  Along each
permutation of \(\mathcal A\), the marginal increments telescope from the
empty to the complete intervention, proving the efficiency identity.
\end{proof}

\subsection{Black-box geometry, alignment, and functional inequalities}
\label{subsec-black-box-geometry}

\begin{definition}[Black-box perturbation geometry]
\label{def-black-box-perturbation-geometry}

Let \(P\) be a represented Markov perturbation kernel with invariant audit
law \(\mu\), and suppose \((X,X')\sim\mu(dx)P(x,dx')\).  Let \(X''\sim\mu\)
be independent.  The output Dirichlet energy, variance, and labeled signed
progress are
\begin{align}
\label{eq-black-box-output-dirichlet-energy}
\mathcal E_{P,f}
&=\frac12\mathbb E\|p_f(X')-p_f(X)\|_2^2,\\
\label{eq-black-box-output-variance}
\mathcal V_f
&=\frac12\mathbb E\|p_f(X'')-p_f(X)\|_2^2,\\
\label{eq-black-box-perturbation-target-progress}
\mathcal C_{P,f}
&=\mathbb E\left[
\ell(p_f(X),Y)-\ell(p_f(X'),Y')
\right],
\end{align}
where \(Y'\) is the represented transformed clean label.  Each energy and
variance integrand lies in \([0,1]\); the two losses defining
\(\mathcal C_{P,f}\) lie in \([0,1]\).  When
\(\mathcal V_f>0\), define the observed Rayleigh ratio
\begin{equation}
\label{eq-black-box-output-rayleigh-ratio}
\rho_{P,f}=\mathcal E_{P,f}/\mathcal V_f.
\end{equation}

A prospective Caus record additionally specifies a kernel \(K_f\) whose row
is computed from oracle outputs and public information available before the
clean label reveal.  Replacing \(P\) by \(K_f\) in
\cref{eq-black-box-perturbation-target-progress} gives its audited Caus
alignment.  Label-conditioned perturbations remain analytic diagnostics and
are not prospective Caus records.

\end{definition}

\begin{theorem}[Finite-query geometric confidence region]
\label{thm-black-box-geometric-confidence-region}

Use \(N_E,N_V,N_C\) independent represented audit blocks to estimate
\(\mathcal E_{P,f},\mathcal V_f\), and the two loss expectations defining
\(\mathcal C_{P,f}\).  Allocate failure probabilities
\(\delta_E,\delta_V,\delta_C\), and put
\begin{align}
r_E&=\sqrt{\frac{\log(2/\delta_E)}{2N_E}},&
r_V&=\sqrt{\frac{\log(2/\delta_V)}{2N_V}},&
r_C&=\sqrt{\frac{\log(4/\delta_C)}{2N_C}}.
\end{align}
On an event of probability at least
\(1-\delta_E-\delta_V-\delta_C\),
\begin{align}
\label{eq-black-box-energy-confidence}
\mathcal E_{P,f}&\in
[E_-,E_+]
:=\bigl[[\widehat{\mathcal E}_{P,f}-r_E]_+,
        (\widehat{\mathcal E}_{P,f}+r_E)\wedge1\bigr],\\
\label{eq-black-box-variance-confidence}
\mathcal V_f&\in
[V_-,V_+]
:=\bigl[[\widehat{\mathcal V}_f-r_V]_+,
        (\widehat{\mathcal V}_f+r_V)\wedge1\bigr],\\
\label{eq-black-box-caus-progress-confidence}
\mathcal C_{P,f}&\in
[\widehat{\mathcal C}_{P,f}-2r_C,
 \widehat{\mathcal C}_{P,f}+2r_C]\cap[-1,1].
\end{align}
If \(V_->0\), then
\begin{equation}
\label{eq-black-box-rayleigh-confidence}
\frac{E_-}{V_+}
\le\rho_{P,f}\le
\frac{E_+}{V_-}.
\end{equation}
If \(P\) is reversible and the output coordinates lie in the declared
Dirichlet domain, the upper endpoint in
\cref{eq-black-box-rayleigh-confidence} is a represented test-function upper
certificate for the Poincar\'e constant of \(P\).  A lower certificate for the
full declared function domain follows when the audit also supplies the
uniform feature-domain deviations and completeness comparison required by
\cref{thm-generic-empirical-functional-inequality-transfer}.
For a deterministic oracle, the three independent audits use
\begin{equation}
\label{eq-black-box-geometric-query-count}
Q_{\rm geom}=2N_E+2N_V+2N_C
\end{equation}
oracle evaluations.  Represented stochastic repetition multiplies each
summand by its declared repeat count; shared calls are permitted only through
a represented common-block coupling.

\end{theorem}

\begin{proof}
Apply Hoeffding's inequality to the bounded energy and variance integrands and
to the two bounded loss expectations.  A union bound gives the simultaneous
event.  Division by positive interval endpoints proves
\cref{eq-black-box-rayleigh-confidence}.  For every nonconstant admissible
test function, the Poincar\'e constant is at most its Dirichlet-to-variance
ratio; summing over output coordinates preserves this Rayleigh quotient.
The full-domain lower transfer is exactly
\cref{eq-generic-full-functional-constant}.
\end{proof}

\subsection{Behavioral quotients and residual structure}
\label{subsec-black-box-quotient-residual}

\begin{definition}[Black-box output quotient]
\label{def-black-box-output-quotient}

Let \(b:\Delta_K\to Q_f\) be a represented finite discretizer, fixed or
selected with its charged complexity, and put
\begin{equation}
\label{eq-black-box-output-quotient-map}
q_f(x)=b(p_f(x)).
\end{equation}
Given a represented perturbation kernel \(P\), define its quotient-row map
and fiber diameter by
\begin{align}
P_{q_f}(x,c)&=P(x,q_f^{-1}(c)),\\
\delta_{q_f}(P)
&=\max_{q_f(x)=q_f(x')}
\|P_{q_f}(x,\cdot)-P_{q_f}(x',\cdot)\|_{\rm TV}.
\end{align}
The discretizer, bin locator, transition sampler, empty-cell convention, and
all row evaluations are fields of the complete black-box record.

\end{definition}

\begin{theorem}[Black-box exact and quasi-lumpability compilation]
\label{thm-black-box-lumpability-compilation}

Suppose first that the audited carrier \(X_{\rm aud}\) is finite and every
state has at least \(N_{\min}\) independent outgoing perturbation samples.
Then, with the radius of
\cref{eq-lumpability-entrywise-radius},
\begin{equation}
\label{eq-black-box-lumpability-confidence}
|\delta_{q_f}(P)-\widehat\delta_{q_f}|
\le |Q_f|r_{q_f}
\end{equation}
with probability at least \(1-\delta_{\rm lump}\).  On an implicit carrier,
the same conclusion holds with \(2\Gamma_{q_f}^{\rm row}\) whenever the
complete audit supplies the simultaneous quotient-row radius required by
\cref{thm-finite-sample-lumpability-certificate}.

If the grid-separation condition in
\cref{eq-lumpability-grid-separation} holds, then
\begin{equation}
\label{eq-black-box-exact-lumpability}
PU_{q_f}=U_{q_f}\bar P.
\end{equation}
Otherwise the certified upper endpoint
\begin{equation}
\label{eq-black-box-lumpability-upper-endpoint}
\overline\delta_{q_f}
=\min\{1,\widehat\delta_{q_f}+|Q_f|r_{q_f}\}
\end{equation}
gives
\begin{align}
\label{eq-black-box-quasi-lumpable-operator}
\|PU_{q_f}-U_{q_f}\bar P\|_{\infty\to\infty}
&\le2\overline\delta_{q_f},\\
\label{eq-black-box-quasi-lumpable-hit-transfer}
|\Pr\{\tau_T\le H\}-\Pr\{\tau_{T/Q_f}^{\bar P}\le H\}|
&\le\min\{1,H\overline\delta_{q_f}\}
\end{align}
for every represented terminal-compatible absorbing target.  Selection of
\(b\) on the same audit sample uses the row-class radius
\(\Gamma_{q_f}^{\rm row}\); it is not treated as a fixed-bin audit.
For a deterministic oracle on the finite audited carrier, caching every
base label and querying every sampled transition endpoint gives the explicit
upper count
\begin{equation}
\label{eq-black-box-quotient-query-count}
Q_{\rm quot}
\le |X_{\rm aud}|+\sum_{x\in X_{\rm aud}}N_x,
\end{equation}
where \(N_x\ge N_{\min}\) is the number of outgoing samples at \(x\).
Every additional candidate discretizer or stochastic repeat is charged in
the selection or repetition coordinate of \(B_{\rm BB}\).

\end{theorem}

\begin{proof}
The confidence statement is
\cref{thm-finite-sample-lumpability-certificate} applied to the represented
map in \cref{eq-black-box-output-quotient-map}.  Grid separation gives exact
lumpability.  The operator and first-hit comparisons follow from
\cref{thm-certified-fiber-diameter-lumpability-transfer} after substituting
the certified upper endpoint.  The selected-bin statement is the implicit
carrier clause of the same finite-sample theorem.
\end{proof}

\begin{definition}[Held-out structural surrogate and residual]
\label{def-black-box-held-out-structural-residual}

Split the audit sample into independent construction and evaluation parts.
Using only the construction part, build a represented surrogate
\(g_{\rm str}:X\to\Delta_K\) from the audited structural feature library,
intervention responses, quotient cells, or another source-resolved class in
the saturated grammar.  On the evaluation law define
\begin{align}
\label{eq-black-box-structural-residual-energy}
R_{\rm res}^{\rm BB}(f;g_{\rm str})
&=\frac12\mathbb E
\|p_f(X)-g_{\rm str}(X)\|_2^2,\\
\label{eq-black-box-residual-target-correlation}
C_{\rm res}^{\rm BB}(f;g_{\rm str})
&=\frac12\left|
\mathbb E\left\langle
p_f(X)-g_{\rm str}(X),
e_Y-\pi_Y
\right\rangle
\right|,
\end{align}
where \(\pi_Y=\mathbb E e_Y\) is estimated on the construction split or
specified by the audit law.  Both integrands have absolute value at most one.
The first coordinate measures output behavior not reconstructed by the
declared structural surrogate.  The second measures the part of that residual
that remains aligned with the clean audit target.  A provenance compiler may
route a target-orthogonal residual into \(J_{\rm res}\); the empirical norm
alone does not assert that internal classification.

\end{definition}

\begin{theorem}[Held-out black-box residual certificate]
\label{thm-held-out-black-box-residual-certificate}

Let the evaluation split contain \(N_{\rm res}\) independent examples and
let
\begin{equation}
\label{eq-black-box-residual-radius}
r_{\rm res}
=\sqrt{\frac{\log(4/\delta_{\rm res})}{2N_{\rm res}}}.
\end{equation}
Then, conditionally on the construction split and hence also unconditionally,
with probability at least \(1-\delta_{\rm res}\),
\begin{align}
\label{eq-black-box-residual-energy-confidence}
R_{\rm res}^{\rm BB}
&\in
\bigl[[\widehat R_{\rm res}^{\rm BB}-r_{\rm res}]_+,
      (\widehat R_{\rm res}^{\rm BB}+r_{\rm res})\wedge1\bigr],\\
\label{eq-black-box-residual-correlation-confidence}
C_{\rm res}^{\rm BB}
&\le
\widehat C_{\rm res}^{\rm BB}+r_{\rm res}.
\end{align}
If the surrogate is selected using the evaluation split, these radii are
replaced by the applicable represented structural Rademacher or PAC--Bayes
radius for the selected surrogate class.

\end{theorem}

\begin{proof}
Conditioned on the construction split, the surrogate is fixed.  The squared
simplex distance divided by two lies in \([0,1]\).  The inner product in
\cref{eq-black-box-residual-target-correlation}, divided by two, lies in
\([-1,1]\).  Hoeffding's inequality and a union bound prove the two claims.
The data-selected case is the uniform selected-class substitution already
used in \cref{thm-finite-query-simultaneous-black-box-audit}.
\end{proof}

\subsection{Master black-box certificate}
\label{subsec-master-black-box-certificate}

\begin{definition}[Complete quantitative black-box signature]
\label{def-complete-quantitative-black-box-signature}

The complete quantitative signature of one black-box audit is
\begin{equation}
\label{eq-complete-quantitative-black-box-signature}
\mathfrak C_{\rm BB}
=\left(
\begin{array}{c}
\mathcal I_{\rm risk},\mathcal I_{\rm cal},
\mathfrak C_{\rm audit},\mathfrak C_{\rm tr},
\{\mathcal I_{D_a},\mathcal I_{A_a}\}_{a\in\mathfrak A_{\rm BB}},\\
\mathcal I_{\mathcal E},\mathcal I_{\mathcal V},
\mathcal I_{\rho},\mathcal I_{\mathcal C},
\mathcal I_{\delta_q},
\mathcal I_{R_{\rm res}},\mathcal I_{C_{\rm res}},\\
N,Q,B_{\rm BB},\mu,\mathcal I_{\rm shift},
\boldsymbol\delta
\end{array}
\right).
\end{equation}
Every interval is tagged by its audit law, data split, query interface,
noise channel, selected diagnostic family, and confidence allocation.

\end{definition}

\begin{theorem}[End-to-end black-box structural certificate]
\label{thm-end-to-end-black-box-structural-certificate}

Consider one complete record from
\cref{def-complete-black-box-predictor-audit}.  Suppose it contains the
predictive diagnostics, one represented audit-noise identification record,
the structural interventions, a perturbation geometry, an output quotient,
and a held-out structural surrogate defined above.  Allocate confidence so
that
\begin{equation}
\label{eq-black-box-master-confidence-allocation}
\delta_{\rm pred}+\delta_{\rm noise}+\delta_{\rm str}
+\delta_{\rm geom}+\delta_{\rm lump}+\delta_{\rm res}
\le\delta_{\rm all}^{\rm BB}.
\end{equation}
Then, with probability at least \(1-\delta_{\rm all}^{\rm BB}\), the single
signature \(\mathfrak C_{\rm BB}\) simultaneously certifies:
\begin{enumerate}[label=\textup{(\roman*)}]
\item clean or noise-corrected predictive risk, calibration, confusion, and
      subgroup performance on the audit law;
\item the audit-label channel and its exact clean-risk amplification, and the
      compatible training-channel set whenever a training comparison is
      represented;
\item every behavioral structural defect and target response in
      \cref{eq-black-box-source-defect-interval,%
eq-black-box-source-alignment-interval};
\item the output energy, variance, Rayleigh ratio, and prospective Caus
      progress in
      \cref{eq-black-box-energy-confidence,%
eq-black-box-rayleigh-confidence,%
eq-black-box-caus-progress-confidence};
\item exact output-quotient dynamics when
      \cref{eq-lumpability-grid-separation} holds, and otherwise the explicit
      quasi-lumpable defect and first-hit loss in
      \cref{eq-black-box-quasi-lumpable-operator,%
eq-black-box-quasi-lumpable-hit-transfer}; and
\item the magnitude and clean-target correlation of output behavior not
      reconstructed by the declared structural surrogate.
\end{enumerate}
If a deployment-law comparison is represented, every bounded coordinate is
transported by \cref{thm-black-box-audit-deployment-transfer}.  If an exposed
model compiler maps the black-box behavior into the complete source-resolved
learning ledger, the source-qualified target gain additionally satisfies
\begin{equation}
\label{eq-black-box-compiled-target-gain}
g_f
\le
\min\left\{
1,
U_f^{\rm stat\to tgt},
U_f^{\rm dyn\to tgt},
\beta_f^+
+\sqrt{\frac{\log2}{2}
\sum_{a\in\mathfrak J_{\Abs}^{5}}J_f^{(a)}},
eH_B(n)\sum_{a\in\mathfrak J_{\Abs}^{5}}\mathcal L_f^{(a)}
\right\},
\end{equation}
with each unknown source coordinate replaced by its simultaneous empirical
upper endpoint.  The compilation requirement is satisfied automatically for
the white-box records of
\cref{def-complete-neural-implementation-record}; for a pure oracle record,
\cref{eq-black-box-compiled-target-gain} is included exactly when its
provenance compiler is exposed.

\end{theorem}

\begin{proof}
Apply
\cref{cor-simultaneous-black-box-predictive-audit,%
thm-black-box-binary-repeated-label-noise,%
thm-black-box-multiclass-noise-inversion,%
thm-simultaneous-black-box-structural-signature,%
thm-black-box-geometric-confidence-region,%
thm-black-box-lumpability-compilation,%
thm-held-out-black-box-residual-certificate}
on their respective audit splits and take the intersection of their events.
The union bound and
\cref{eq-black-box-master-confidence-allocation} give the stated confidence.
The law-transfer claim is
\cref{thm-black-box-audit-deployment-transfer}.  Under the provenance
compiler, substitute the empirical source upper endpoints into
\cref{eq-parameter-source-resolved-empirical-target-gain}; monotonicity of
the minimum, sum, and square root proves
\cref{eq-black-box-compiled-target-gain}.
\end{proof}

\begin{corollary}[Certified black-box comparison of two trained predictors]
\label{cor-certified-black-box-model-comparison}

Evaluate predictors \(f\) and \(g\) on the same paired audit units and let
\(\Delta_i=D_i(f)-D_i(g)\in[-1,1]\) be any fixed paired diagnostic
difference.  With
\begin{equation}
\label{eq-black-box-paired-model-radius}
r_{\rm pair}
=\sqrt{\frac{2\log(2M/\delta_{\rm pair})}{N}},
\end{equation}
one has simultaneously over \(M\) comparisons
\begin{equation}
\label{eq-black-box-paired-model-comparison}
\left|
\mathbb E\Delta_i-\frac1N\sum_{i=1}^N\Delta_i
\right|
\le r_{\rm pair}
\end{equation}
with probability at least \(1-\delta_{\rm pair}\).  The comparison therefore
inherits common-example cancellation without treating the two empirical
means as independent.

\end{corollary}

\begin{proof}
The paired differences are independent across audit units and lie in an
interval of length two.  Apply the rescaled Hoeffding inequality and a union
bound over the \(M\) fixed comparisons.
\end{proof}

\begin{figure}[tbp]
\centering
\resizebox{.98\linewidth}{!}{%
\begin{tikzpicture}[fw diagram,node distance=8mm and 10mm]
  \node[fw source,text width=2.5cm] (model) {trained model\\black-box oracle};
  \node[fw provenance,text width=2.5cm,right=of model] (access)
    {charged access\\and query record};
  \node[fw use,text width=2.7cm,right=of access] (data)
    {test law, labels,\\interventions};
  \node[fw geom,text width=2.7cm,right=of data] (audit)
    {risk, noise, structure,\\geometry, quotients};
  \node[fw terminal,text width=2.7cm,right=of audit] (cert)
    {simultaneous quantitative\\certificate};
  \draw[fw arrow] (model)--(access);
  \draw[fw arrow] (access)--(data);
  \draw[fw arrow] (data)--(audit);
  \draw[fw arrow] (audit)--(cert);
\end{tikzpicture}%
}
\caption{Black-box audit compilation.  Oracle access and every selected query
are charged before the test-law observations are converted into simultaneous
performance, noise, structural, geometric, and quotient coordinates.}
\label{fig-black-box-audit-compilation}
\end{figure}

\begin{figure}[tbp]
\centering
\begin{tikzpicture}[fw diagram,node distance=8mm and 10mm]
  \node[fw source,text width=2.5cm] (clean) {clean audit label\\$Y$};
  \node[fw residual,text width=2.5cm,right=of clean] (channel)
    {identified channel\\$C_{\rm audit}$};
  \node[fw use,text width=2.5cm,right=of channel] (observed)
    {observed label\\$\widetilde Y$};
  \node[fw geom,text width=2.7cm,below=of observed] (inverse)
    {risk reconstruction\\inverse-norm cost};
  \node[fw provenance,text width=2.7cm,left=of inverse] (transfer)
    {represented training--audit\\channel comparison};
  \node[fw terminal,text width=2.6cm,left=of transfer] (train)
    {training-noise\\compatible set};
  \draw[fw arrow] (clean)--(channel);
  \draw[fw arrow] (channel)--(observed);
  \draw[fw arrow] (observed)--(inverse);
  \draw[fw arrow] (inverse)--(transfer);
  \draw[fw arrow] (transfer)--(train);
\end{tikzpicture}
\caption{Typed noise reconstruction.  Audit noise is inferred from represented
anchors or repeated labels.  Training noise is transported only through a
represented comparison, yielding a confidence set rather than an inferred
training history.}
\label{fig-black-box-noise-reconstruction}
\end{figure}

\begin{figure}[tbp]
\centering
\resizebox{.98\linewidth}{!}{%
\begin{tikzpicture}[fw diagram,node distance=7mm and 7mm]
  \node[fw source,text width=2.5cm] (oracle) {oracle behavior\\on audit law};
  \node[fw provenance,text width=1.7cm,right=of oracle,yshift=17mm] (add) {Add};
  \node[fw provenance,text width=1.7cm,right=of oracle,yshift=8mm] (mult) {Mult};
  \node[fw provenance,text width=1.7cm,right=of oracle] (ord) {Ord};
  \node[fw provenance,text width=1.7cm,right=of oracle,yshift=-8mm] (sym) {Sym};
  \node[fw provenance,text width=1.7cm,right=of oracle,yshift=-17mm] (filt) {Filt};
  \node[fw geom,text width=2.1cm,right=19mm of ord,yshift=13mm] (geom) {Geom/Caus};
  \node[fw geom,text width=2.1cm,right=19mm of ord] (lift) {Lift/Bdry};
  \node[fw residual,text width=2.1cm,right=19mm of ord,yshift=-13mm] (res) {residual};
  \node[fw terminal,text width=2.8cm,right=18mm of lift] (signature)
    {behavioral structural\\signature};
  \foreach \n in {add,mult,ord,sym,filt}{\draw[fw arrow] (oracle)--(\n);}
  \foreach \n in {add,mult}{\draw[fw arrow] (\n)--(geom);}
  \draw[fw arrow] (ord)--(lift);
  \foreach \n in {sym,filt}{\draw[fw arrow] (\n)--(res);}
  \foreach \n in {geom,lift,res}{\draw[fw arrow] (\n)--(signature);}
\end{tikzpicture}%
}
\caption{Externally auditable structural channels.  Represented input
interventions produce behavioral defects and target-response intervals.
Internal provenance is added only through an exposed compiler.}
\label{fig-black-box-structural-intervention-tree}
\end{figure}

\begin{figure}[tbp]
\centering
\begin{tikzpicture}[fw diagram,node distance=8mm and 10mm]
  \node[fw source,text width=2.5cm] (x) {audit carrier\\$X$};
  \node[fw use,text width=2.5cm,right=of x] (pred)
    {oracle prediction\\$p_f(x)$};
  \node[fw provenance,text width=2.5cm,right=of pred] (bin)
    {represented bins\\$q_f(x)$};
  \node[fw geom,text width=2.6cm,below=of bin] (rows)
    {perturbation rows\\on output cells};
  \node[fw terminal,text width=2.7cm,left=of rows] (exact)
    {exact quotient\\or defect interval};
  \node[fw residual,text width=2.7cm,left=of exact] (transfer)
    {first-hit transfer\\loss $H\delta_{q_f}$};
  \draw[fw arrow] (x)--(pred);
  \draw[fw arrow] (pred)--(bin);
  \draw[fw arrow] (bin)--(rows);
  \draw[fw arrow] (rows)--(exact);
  \draw[fw arrow] (exact)--(transfer);
\end{tikzpicture}
\caption{Behavioral quotient audit.  Output cells become a dynamical
abstraction only after their perturbation rows satisfy the exact or
quasi-lumpability certificate.}
\label{fig-black-box-behavioral-quotient}
\end{figure}

\section{Length-Bounded Levin Search and the Structural Learner}
\label{sec-length-bounded-levin-structural-learner}

This section relates the saturated active profile to one represented universal
search record.  Description length remains a separate resource coordinate:
the universal record pays exponentially for an unknown record code, whereas a
structural learner may construct its deployed record from the presentation.
The comparison below records the code length, simulation overhead, verifier
error, learning error, and physical ceiling separately.

\subsection{Prefix codes and length-bounded saturated profiles}
\label{subsec-levin-prefix-codes-length-profiles}

\begin{definition}[Canonical code of a complete active record]
\label{def-canonical-complete-active-record-code}

Fix a prefix-free binary encoding
\(
\mathsf{code}:\mathscr R^{\rm act}\to\{0,1\}^{*}
\)
of the syntax of complete master active-sampling records from
\cref{def-complete-master-active-sampling-record}.  The code stores every
finite program, rational constant, precision declaration, continuation rule,
and output convention belonging to the record.  Public instance data are read
through the declared presentation interface and are not copied into the
program code.  Put
\begin{equation}
\label{eq-complete-active-record-description-length}
d(R):=|\mathsf{code}(R)|.
\end{equation}
The prefix condition gives
\begin{equation}
\label{eq-complete-active-record-kraft}
\sum_{R\in\mathscr R^{\rm act}}2^{-d(R)}\le1.
\end{equation}

The encoding is \emph{family-uniform}: one finite code describes the record
on every size slice, while the public presentation supplies \(n,m,\zeta\) and
the instance.  Let \(\mathsf{Int}\) be its represented universal interpreter.
There is a monotone class-preserving scalar majorant
\(\sigma_{\rm int}(s,d)\) such that interpreting the first \(s\) scalarized
steps of a length-\(d\) record costs at most
\(\sigma_{\rm int}(s,d)\).  The canonical finite-precision machine may use
\begin{equation}
\label{eq-universal-interpreter-majorant}
\sigma_{\rm int}(s,d)
\le c_{\rm int}(s+d+1)
\bigl(1+\lceil\log_2(s+d+2)\rceil\bigr)^{a_{\rm int}}
\end{equation}
for fixed machine constants \(c_{\rm int},a_{\rm int}\).  The scalarization
of complete cost includes reading the code, so for a fixed constant
\(c_{\rm read}\),
\begin{equation}
\label{eq-code-reading-cost-domination}
d(R)\le c_{\rm read}\bigl(1+\kappa(\operatorname{Cost}(R))\bigr).
\end{equation}

\end{definition}

\begin{definition}[Length-bounded active and selector profiles]
\label{def-length-bounded-active-selector-profiles}

For an instance \(I\), horizon \(H\), scalar ceiling \(b\), and
\(\ell\in\mathbb N\), define
\begin{align}
\label{eq-length-bounded-active-arrival-profile}
\operatorname{ActArr}_{I,H,\le\ell}^{\rm sat}(b)
&=
\sup_{\substack{R\in\mathscr R_{I,H}^{\rm act}\,;\
\kappa(\operatorname{Cost}(R))\le b,\ d(R)\le\ell}}
A_R^H,\\
\label{eq-length-bounded-active-selector-profile}
\operatorname{ActSel}_{I,H,\le\ell}^{\rm sat}(b)
&=
\max_{0\le t<H}
\sup_{\substack{R\in\mathscr R_{I,H}^{\rm act}\,;\
\kappa(\operatorname{Cost}(R))\le b,\ d(R)\le\ell}}
\bigl(\operatorname{Adv}_t(K_{R,t};\nu_{R,t})\bigr)_+.
\end{align}
The family-uniform versions apply the fixed functional
\(\mathfrak M_{n,m,\zeta}\) before taking the same supremum, as in
\cref{def-family-uniform-active-saturated-profile}.  Empty suprema are zero.

\end{definition}

\begin{lemma}[Monotonicity and exact exhaustion by code length]
\label{lem-length-bounded-profile-monotonicity-exhaustion}

For fixed \((I,H,b)\), both profiles in
\cref{def-length-bounded-active-selector-profiles} are nondecreasing in
\(\ell\), bounded above by their unrestricted counterparts, and satisfy
\begin{align}
\label{eq-length-profile-exhaustion}
\sup_{\ell\in\mathbb N}
\operatorname{ActArr}_{I,H,\le\ell}^{\rm sat}(b)
&=\operatorname{ActArr}_{I,H}^{\rm sat}(b),\\
\sup_{\ell\in\mathbb N}
\operatorname{ActSel}_{I,H,\le\ell}^{\rm sat}(b)
&=\operatorname{ActSel}_{I,H}^{\rm sat}(b).
\end{align}
Moreover, \cref{eq-code-reading-cost-domination} implies equality already
for every
\(
\ell\ge c_{\rm read}(1+b)
\).

\end{lemma}

\begin{proof}
The feasible record classes are nested in \(\ell\).  Every represented
record has one finite code, so the union of the length-bounded classes is the
unrestricted class.  A supremum over an increasing union equals the supremum
of the suprema.  Finally, every record of scalarized cost at most \(b\) has
code length at most \(c_{\rm read}(1+b)\) by
\cref{eq-code-reading-cost-domination}.
\end{proof}

\begin{definition}[Represented target gate and global gate error]
\label{def-levin-represented-target-gate-error}

Let \(V_I:X_I\to\{0,1\}\) be a represented public verifier of scalar cost
at most \(t_V(n,m)\), and put \(\widehat T_I=V_I^{-1}(1)\).  The clean
target gate is \(T_I\).  Define the pointwise gate defect
\begin{equation}
\label{eq-levin-pointwise-gate-defect}
e_V(I)=\mathbf1_{\{\widehat T_I\ne T_I\}}
\end{equation}
and its family value
\begin{equation}
\label{eq-levin-family-gate-error}
\varepsilon_V(n,m,\zeta)
=\mathfrak M_{n,m,\zeta}(I\mapsto e_V(I)).
\end{equation}
This global event controls every verifier call in one execution and therefore
is not multiplied by the horizon.  If only conditional per-call soundness and
completeness bounds are available, their predictable union bound is used in
place of \cref{eq-levin-family-gate-error} and is charged in the verifier
slot of the complete record.

\end{definition}

\subsection{The represented Levin record}
\label{subsec-represented-levin-record}

\begin{definition}[Length-bounded represented Levin search]
\label{def-length-bounded-represented-levin-search}

For \(\ell\in\mathbb N\cup\{\infty\}\), let
\(\mathfrak U_\ell\) be the complete active record that lazily decodes the
prefix tree, simulates each legal record code \(p\) with
\(|p|\le\ell\), and assigns branch \(p\) Kraft weight \(2^{-|p|}\).
Whenever a branch emits a candidate, \(\mathfrak U_\ell\) evaluates \(V_I\)
before declaring arrival.  Branch state, random bits, candidate emission,
verification, stopping, memory, and the scheduler clock are fields of the
complete record.

There are fixed constants \(c_{\rm sch},a_{\rm sch}\) such that a branch of
length \(d\) receives \(s\) interpreted branch steps within scalarized
universal cost
\begin{equation}
\label{eq-levin-scheduler-majorant}
\Lambda_{\mathfrak U}(d,s)
=c_{\rm sch}2^d(s+d+1)
\bigl(1+\lceil\log_2(s+d+2)\rceil\bigr)^{a_{\rm sch}}.
\end{equation}
This is obtained by the standard staged Kraft scheduler: at stage \(j\), it
executes every decoded prefix \(p\) with
\(|p|\le j\) for the stage share proportional to \(2^{-|p|}\); the Kraft
inequality bounds total simulated work, and the logarithmic factor pays for
prefix-tree traversal and resumable branch lookup.  Thus the description
penalty is exactly the factor \(2^d\); machine simulation and scheduling are
the displayed class-preserving factors.

The code of \(\mathfrak U_\ell\) contains the interpreter, scheduler,
verifier reference, and a self-delimiting numeral for \(\ell\).  Hence
\begin{equation}
\label{eq-levin-record-own-code-length}
d(\mathfrak U_\ell)
\le c_{\mathfrak U}
+2\left\lceil\log_2(\ell+2)\right\rceil
\quad(\ell<\infty),
\end{equation}
whereas \(d(\mathfrak U_\infty)\le c_{\mathfrak U}\).

\end{definition}

\begin{theorem}[Length-bounded Levin realization]
\label{thm-length-bounded-levin-realization}

Let \(q\in\mathscr R_{I,H}^{\rm act}\) satisfy
\(d(q)\le\ell\) and
\(\kappa(\operatorname{Cost}(q))\le b\).  Suppose a horizon-\(H\) execution
emits at most \(H\) candidates.  Define
\begin{align}
\label{eq-levin-branch-interpreted-cost}
s_q(b,n,m,H)
&=\sigma_{\rm int}(b,d(q))+H\,t_V(n,m),\\
\label{eq-levin-per-record-ceiling}
b_{\mathfrak U}(q;b,n,m,H)
&=\Lambda_{\mathfrak U}
\bigl(d(q),s_q(b,n,m,H)\bigr).
\end{align}
Let
\begin{equation}
\label{eq-levin-per-record-refined-horizon}
H_{\mathfrak U}(q;b,n,m,H)
=H_{b_{\mathfrak U}(q;b,n,m,H)}
\end{equation}
be the largest refined clock horizon admitted by this scalar ceiling.
Then, pointwise in the presentation,
\begin{equation}
\label{eq-levin-pointwise-realization}
A_{q,I}^{H}
\le
A_{\mathfrak U_\ell,I}^{H_{\mathfrak U}(q)}
\bigl(b_{\mathfrak U}(q;b,n,m,H)\bigr)
+e_V(I).
\end{equation}
Consequently, for every declared monotone subadditive family functional,
\begin{equation}
\label{eq-levin-family-realization}
\mathfrak M_{n,m,\zeta}(A_q^H)
\le
\mathfrak M_{n,m,\zeta}
\!\left(A_{\mathfrak U_\ell}^{H_{\mathfrak U}(q)}
(b_{\mathfrak U}(q;b,n,m,H))\right)
+\varepsilon_V(n,m,\zeta).
\end{equation}
For a linear expectation functional this is the usual average-instance
comparison.

\end{theorem}

\begin{proof}
Couple the branch indexed by \(\mathsf{code}(q)\) to the same represented
random seed as \(q\).  The interpreter reproduces its horizon-\(H\) path law
within \(\sigma_{\rm int}(b,d(q))\) branch cost.  At most \(H\) candidate
verifications add \(Ht_V\).  By
\cref{eq-levin-scheduler-majorant}, the universal record supplies this branch
cost by the ceiling in \cref{eq-levin-per-record-ceiling}.

On \(\{\widehat T_I=T_I\}\), the universal record declares a hit on the
coupled branch precisely when that branch first emits a target; an earlier
acceptance is also a true target hit.  Hence
\(A_{q,I}^{H}\le
A_{\mathfrak U_\ell,I}^{H_{\mathfrak U}(q)}\) on this event.  On its
complement both probabilities lie in \([0,1]\), which proves
\cref{eq-levin-pointwise-realization}.  Monotonicity and subadditivity of the
family functional give \cref{eq-levin-family-realization}.
\end{proof}

\begin{corollary}[Uniform length-class realization]
\label{cor-uniform-length-class-levin-realization}

Put
\begin{align}
\label{eq-levin-uniform-length-class-branch-cost}
s_\ell(b,n,m,H)
&=\sup_{0\le d\le\ell}\sigma_{\rm int}(b,d)+Ht_V(n,m),\\
\label{eq-levin-uniform-length-class-ceiling}
b_{\mathfrak U}(\ell;b,n,m,H)
&=\Lambda_{\mathfrak U}(\ell,s_\ell(b,n,m,H)).
\end{align}
Put
\begin{equation}
\label{eq-levin-uniform-length-class-horizon}
H_{\mathfrak U}(\ell;b,n,m,H)
=H_{b_{\mathfrak U}(\ell;b,n,m,H)}.
\end{equation}
Then
\begin{equation}
\label{eq-length-class-profile-levin-bound}
\operatorname{ActArr}_{\mathfrak M,n,m,\zeta,H,\le\ell}^{\rm sat}(b)
\le
\mathfrak M_{n,m,\zeta}\!\left(
A_{\mathfrak U_\ell}^{H_{\mathfrak U}(\ell)}
(b_{\mathfrak U}(\ell;b,n,m,H))
\right)
+\varepsilon_V(n,m,\zeta).
\end{equation}

\end{corollary}

\begin{proof}
For \(d(q)\le\ell\), monotonicity of
\(\sigma_{\rm int}\) and \(\Lambda_{\mathfrak U}\) places the simulated
branch below the uniform ceiling.  Apply
\cref{eq-levin-family-realization} and take the supremum over the defining
length class.
\end{proof}

\begin{proposition}[Sharp integration with saturated framework certificates]
\label{prop-levin-sharp-saturated-certificate-integration}

Let \(\{U_j(b)\}_{j\in\mathcal J}\) be any finite collection of valid
same-law, same-target upper certificates for
\(\operatorname{ActArr}_{\mathfrak M,n,m,\zeta,H}^{\rm sat}(b)\).
The collection may contain the prospective selector/source bound, master
active-sampling envelope, occupation-flow Bellman barrier, controlled
capacity or information candidate, and a simultaneous empirical Bellman
upper certificate furnished respectively by
\cref{thm-prospective-selector-structural-charge,thm-master-optimal-structural-active-sampling-envelope,thm-bellman-barrier-duality-affordable-arrival,thm-universal-dynamical-certificate,cor-primal-dual-empirical-implementation-certificate}.
Then
\begin{align}
\label{eq-levin-sharp-saturated-certificate-minimum}
\operatorname{ActArr}_{\mathfrak M,n,m,\zeta,H,\le\ell}^{\rm sat}(b)
\le
\min\Bigl\{&
\mathfrak M_{n,m,\zeta}\!\left(
A_{\mathfrak U_\ell}^{H_{\mathfrak U}(\ell)}
(b_{\mathfrak U}(\ell;b,n,m,H))\right)
+\varepsilon_V,\notag\\
&\inf_{j\in\mathcal J}U_j(b),\,1
\Bigr\}.
\end{align}
Only certificates expressed for the same family law, target event, horizon,
and scalarization enter the minimum.  Every certificate retains its own
complete construction and evaluation cost.

\end{proposition}

\begin{proof}
The first candidate is
\cref{eq-length-class-profile-levin-bound}.  The length-bounded record class
is a subset of the unrestricted active class, so every \(U_j(b)\) also bounds
its profile.  Arrival lies in \([0,1]\).  Taking the minimum of upper bounds
on the same scalar proves the result.
\end{proof}

\begin{corollary}[Single-record reduction for uniform computation]
\label{cor-levin-single-record-uniform-reduction}

Let \(q\) be one family-uniform record code, so \(d(q)\) is independent of
\(n\).  If \(b(n),H(n),t_V(n,m(n))\), the interpreter majorant, and the
scheduler majorant are polynomial on the declared slice, then
\(b_{\mathfrak U}(q;b,n,m,H)\) is polynomial.  Therefore
\begin{align}
\label{eq-levin-polynomial-single-record-implication}
&\mathfrak M_{n,m,\zeta}(A_q^H)
\text{ nonnegligible at polynomial cost}
\notag\\[-1mm]
&\hspace{25mm}\Longrightarrow\
\mathfrak M_{n,m,\zeta}(A_{\mathfrak U}^{H_{\mathfrak U}(q)})
\text{ nonnegligible at polynomial cost}.
\end{align}
provided \(\varepsilon_V\) is negligible.  Since \(\mathfrak U\) is itself
one family-uniform record, negligibility of its polynomial arrival is
equivalent to negligibility for the complete family-uniform active class.

\end{corollary}

\begin{proof}
The constant \(2^{d(q)}\) in
\cref{eq-levin-scheduler-majorant} is independent of \(n\); the remaining
factors are polynomial by hypothesis.  Apply
\cref{thm-length-bounded-levin-realization}.  The converse direction follows
from membership of \(\mathfrak U\) in the complete active class.
\end{proof}

\subsection{Certified relation to the structural learner}
\label{subsec-levin-structural-learner-relation}

\begin{definition}[Complete learned-control slack]
\label{def-complete-levin-learned-control-slack}

For a deployed learned record \(\widehat R\) satisfying the simultaneous
action-value hypotheses of
\cref{thm-learned-active-rule-finite-horizon-loss}, define
\begin{equation}
\label{eq-complete-levin-learning-slack}
\rho_{\rm learn}(\widehat R)
=
\mathbb E_{\widehat R}
\sum_{t=0}^{H-1}\mathbf1_{\{\tau>t\}}
\bigl(2\Gamma_t+\varepsilon_{{\rm opt},t}\bigr)
+\Delta_{\rm fil}^{H}.
\end{equation}
The radii \(\Gamma_t\) already include their represented approximation,
sample, dependence, noise-transport, selection, and same-law conversion
costs.  The parameter-explicit implementation may instead use the smaller
valid candidate in
\cref{eq-end-to-end-certified-neural-optimum}; its deterministic and
stochastic coordinates are separated by
\cref{thm-complete-certified-bias-variance-decomposition}.
For a canonical neural record, define the sharp same-record candidate
\begin{equation}
\label{eq-sharp-complete-levin-learning-slack}
\rho_{\rm learn}^{\sharp}(\widehat R)
=\min\left\{
\begin{array}{l}
\rho_{\rm learn}(\widehat R),\\
\bigl[U_{\rm Bell}-L_{\rm ev}\bigr]_+,\\
B_{\rm det}(p,K,\boldsymbol\tau,b,H)
 +V_{\rm stat}(p,N,\boldsymbol\eta,\delta),\\
\mathcal L_{\rm cert}\wedge1
\end{array}
\right\}.
\end{equation}
An unavailable candidate has value \(+\infty\).  Each finite entry bounds
the same saturated-arrival gap for the same record on the same simultaneous
event, so their minimum is valid.  For a non-neural active record set
\(\rho_{\rm learn}^{\sharp}=\rho_{\rm learn}\).

\end{definition}

\begin{theorem}[Structural learner versus length-bounded Levin search]
\label{thm-structural-learner-length-bounded-levin}

Fix \((n,m,\zeta,H,b,\ell)\), and train and deploy \(\widehat R\) at the
common ceiling
\(b_{\mathfrak U}(\ell;b,n,m,H)\) and refined horizon
\(H_{\mathfrak U}(\ell;b,n,m,H)\).  If its action-value certificate ranges
over the complete active suffix class at that ceiling, then on its declared
simultaneous event,
\begin{equation}
\label{eq-structural-learner-length-profile-lower}
\mathfrak M_{n,m,\zeta}
(A_{\widehat R}^{H_{\mathfrak U}(\ell)})
\ge
\operatorname{ActArr}_{\mathfrak M,n,m,\zeta,H,\le\ell}^{\rm sat}(b)
-\varepsilon_V(n,m,\zeta)
-\mathfrak M_{n,m,\zeta}(\rho_{\rm learn}^{\sharp}(\widehat R)).
\end{equation}
The upper comparison is membership at the common ceiling:
\begin{equation}
\label{eq-structural-learner-levin-sandwich}
\mathfrak M_{n,m,\zeta}
(A_{\widehat R}^{H_{\mathfrak U}(\ell)})
\le
\operatorname{ActArr}_{\mathfrak M,n,m,\zeta,H_{\mathfrak U}(\ell)}^{\rm sat}
\bigl(b_{\mathfrak U}(\ell;b,n,m,H)\bigr).
\end{equation}

\end{theorem}

\begin{proof}
The universal record \(\mathfrak U_\ell\) is one feasible active record at
the common ceiling.  Hence the saturated active optimum there dominates its
arrival.  The learned-control loss theorem, followed when necessary by the
represented filter transfer, gives a loss of at most
\(\rho_{\rm learn}^{\sharp}\) relative to that optimum.  The minimum is
valid by
\cref{eq-end-to-end-certified-neural-optimum,eq-complete-certified-bias-variance-decomposition,eq-canonical-certificate-loss-optimality}.
Combine this fact with
\cref{eq-length-class-profile-levin-bound}.  The upper inequality is the
definition of the saturated active profile.
\end{proof}

\begin{definition}[Length saturation at a declared ceiling]
\label{def-length-saturation-hypothesis}

For a displayed function \(\ell_*(n,m,\zeta,b)\) and error
\(\delta_{\rm len}\), write
\(\mathsf{LS}(\ell_*,\delta_{\rm len})\) for the quantitative statement
\begin{equation}
\label{eq-length-saturation-hypothesis}
0\le
\operatorname{ActArr}_{\mathfrak M,n,m,\zeta,H}^{\rm sat}(b)
-\operatorname{ActArr}_{\mathfrak M,n,m,\zeta,H,\le\ell_*}^{\rm sat}(b)
\le\delta_{\rm len}(n,m,\zeta,b).
\end{equation}

\end{definition}

\begin{corollary}[Learner comparison under length saturation]
\label{cor-structural-learner-unrestricted-levin-comparison}

Under \(\mathsf{LS}(\ell_*,\delta_{\rm len})\),
\begin{align}
\label{eq-structural-learner-unrestricted-profile-lower}
\mathfrak M_{n,m,\zeta}
(A_{\widehat R}^{H_{\mathfrak U}(\ell_*)})
\ge{}&
\operatorname{ActArr}_{\mathfrak M,n,m,\zeta,H}^{\rm sat}(b)
-\delta_{\rm len}
-\varepsilon_V
-\mathfrak M_{n,m,\zeta}(\rho_{\rm learn}^{\sharp}(\widehat R))
\end{align}
at the compiled ceiling
\(b_{\mathfrak U}(\ell_*;b,n,m,H)\).

\end{corollary}

\begin{proof}
Insert \cref{eq-length-saturation-hypothesis} into
\cref{eq-structural-learner-length-profile-lower}.
\end{proof}

\subsection{Certifiable length and structural advantage}
\label{subsec-levin-certifiable-length-structural-advantage}

\begin{definition}[Blind certifiable-length frontier]
\label{def-levin-certifiable-length-frontier}

For a physical scalar ceiling \(b_{\rm phys}\), define
\begin{equation}
\label{eq-levin-certifiable-length-frontier}
\ell_{\rm cert}(b_{\rm phys};b,n,m,H)
=\max\left\{\ell\in\mathbb N:
b_{\mathfrak U}(\ell;b,n,m,H)\le b_{\rm phys}\right\},
\end{equation}
with value \(-1\) when the set is empty.  If
\(s_\ell\le s\) throughout the relevant range, then
\begin{equation}
\label{eq-levin-certifiable-length-log-bounds}
\ell_{\rm cert}
\le\log_2\frac{b_{\rm phys}}{c_{\rm sch}(s+1)},
\qquad
\ell_{\rm cert}
=\log_2\frac{b_{\rm phys}}{s+1}
-O\bigl(\log\log(b_{\rm phys}+s+2)\bigr),
\end{equation}
whenever the right-hand range is nonempty and the machine constants are
fixed.  Thus blind universal simulation gains one description bit only when
its available compute is approximately doubled.

\end{definition}

\begin{definition}[True and witnessed effective length]
\label{def-true-witnessed-effective-length}

Let
\begin{equation}
\label{eq-levin-learner-parameter-vector}
\vartheta_{\rm learn}
=\bigl(
p,L,(H_\ell,d_\ell,d_{k,\ell},d_{v,\ell},f_\ell)_{\ell\le L},
N,K,\boldsymbol\tau,b,H,\varepsilon_{\rm opt}
\bigr)
\end{equation}
denote the completely charged architecture, sample, optimizer, temperature,
latent-capacity, horizon, and precision choices of
\cref{def-canonical-certified-attention-belief-architecture}.
For a deployed record \(\widehat R\), common ceiling \(b_{\rm phys}\), and
tolerance \(\gamma\ge0\), define
\begin{equation}
\label{eq-true-learner-effective-length}
\ell_{\rm eff}(\widehat R;b_{\rm phys},\gamma)
=\sup\left\{\ell:
\mathfrak M(A_{\widehat R}^{H})
\ge
\operatorname{ActArr}_{\mathfrak M,H,\le\ell}^{\rm sat}(b_{\rm phys})
-\gamma\right\}.
\end{equation}
For a finite represented benchmark set \(\mathcal A\) of records, define the
witnessed profile
\begin{equation}
\label{eq-witnessed-length-profile}
\operatorname{ActArr}_{\mathcal A,H,\le\ell}(b)
=\max_{\substack{a\in\mathcal A:\ d(a)\le\ell,\
\kappa(\operatorname{Cost}(a))\le b}}
\mathfrak M(A_a^H)
\end{equation}
and let \(\ell_{\rm eff}^{\mathcal A}\) be
\cref{eq-true-learner-effective-length} with this witnessed profile in place
of the saturated profile.  Since
\begin{equation}
\label{eq-witnessed-profile-below-saturated}
\operatorname{ActArr}_{\mathcal A,H,\le\ell}(b)
\le\operatorname{ActArr}_{\mathfrak M,H,\le\ell}^{\rm sat}(b),
\end{equation}
benchmark matching certifies \(\ell_{\rm eff}^{\mathcal A}\), not a lower
bound on the true \(\ell_{\rm eff}\).  Define the corresponding reported
advantages
\begin{equation}
\label{eq-levin-structural-length-advantages}
\Delta=\ell_{\rm eff}-\ell_{\rm cert},
\qquad
\Delta^{\mathcal A}
=\ell_{\rm eff}^{\mathcal A}-\ell_{\rm cert}.
\end{equation}
Both effective lengths depend on
\((n,m,\zeta,\vartheta_{\rm learn},b_{\rm phys},\gamma)\); this dependence is
suppressed only where the displayed record already fixes the coordinates.

\end{definition}

\begin{proposition}[Simultaneous empirical benchmark certificate]
\label{prop-simultaneous-witnessed-effective-length-certificate}

Let \(\mathcal A\) be finite.  Suppose simultaneous evaluation intervals
\([L_a,U_a]\) for \(a\in\mathcal A\) and
\([L_{\widehat R},U_{\widehat R}]\) for the learner hold with probability at
least \(1-\alpha\).  If
\begin{equation}
\label{eq-witnessed-effective-length-empirical-condition}
L_{\widehat R}
\ge
\max_{\substack{a\in\mathcal A:\ d(a)\le L,\
\kappa(\operatorname{Cost}(a))\le b_{\rm phys}}}U_a-\gamma,
\end{equation}
then on that event
\(
\ell_{\rm eff}^{\mathcal A}\ge L
\).
Independent binomial evaluations may use simultaneous Clopper--Pearson
intervals; dependent evaluations use the represented martingale or mixing
candidate already admitted by the framework.

\end{proposition}

\begin{proof}
On the simultaneous event, the learner's true arrival is at least
\(L_{\widehat R}\), while every eligible benchmark arrival is at most its
upper endpoint and hence at most the maximum in
\cref{eq-witnessed-effective-length-empirical-condition}.  Therefore the
learner is within \(\gamma\) of the witnessed maximum for every
\(\ell\le L\), which is the defining condition for
\(\ell_{\rm eff}^{\mathcal A}\ge L\).
\end{proof}

\begin{proposition}[Exact represented benchmark crossover]
\label{prop-exact-represented-levin-benchmark-crossover}

For a finite benchmark set \(\mathcal A\), each serialized record \(a\) has
the exact row
\begin{equation}
\label{eq-exact-represented-levin-benchmark-row}
\left(
d(a),\ \kappa(\operatorname{Cost}(a)),\ \mathfrak M(A_a^H),
\ \Lambda_{\mathfrak U}\!\left(
d(a),\sigma_{\rm int}(\kappa(\operatorname{Cost}(a)),d(a))+Ht_V
\right)
\right).
\end{equation}
The fourth coordinate is the physical ceiling at which blind universal
simulation is guaranteed to reproduce that benchmark.  Code length is
computed from the canonical serialized artifact; estimated source-code size
is not substituted for \(d(a)\).  Simultaneous evaluation intervals may
replace the third coordinate without changing the other three.

\end{proposition}

\begin{proof}
The first three entries are the definitions of description length, complete
cost, and arrival.  The final entry is
\cref{eq-levin-per-record-ceiling} with the benchmark's displayed cost.
\end{proof}

\subsection{Resource-matched benchmark frontiers}
\label{subsec-resource-matched-benchmark-frontiers}

The saturated profile and a represented benchmark family live in the same
resource order.  This permits comparisons at equal cost without collapsing
time, acquired data, memory, description length, and evaluation horizon into
an informal running-time exponent.

\begin{definition}[Represented benchmark frontier and matched-success cost]
\label{def-represented-benchmark-pareto-frontier}

Let \(\mathcal A\) be a represented family of complete active records for one
presented task family.  For a vector ceiling
\(\mathbf b=(b_{\rm time},b_{\rm data},b_{\rm mem},b_{\rm desc})\), put
\begin{equation}
\label{eq-represented-benchmark-arrival-frontier}
\operatorname{BenchArr}_{\mathcal A,H}(\mathbf b)
=
\sup_{a\in\mathcal A}
\left\{
\mathfrak M(A_a^H):
\begin{array}{l}
\operatorname{Cost}_{\rm time}(a)\le b_{\rm time},\quad
\operatorname{Cost}_{\rm data}(a)\le b_{\rm data},\\
\operatorname{Cost}_{\rm mem}(a)\le b_{\rm mem},\quad
d(a)\le b_{\rm desc}
\end{array}
\right\},
\end{equation}
with value zero when the feasible class is empty.  The matched-success Pareto
set at level \(\alpha\in(0,1)\) is
\begin{equation}
\label{eq-represented-benchmark-matched-success-pareto-set}
\mathcal P_{\mathcal A,H}(\alpha)
=
\operatorname{Min}_{\preceq}
\left\{
\operatorname{Cost}(a):a\in\mathcal A,
\ \mathfrak M(A_a^H)\ge\alpha
\right\},
\end{equation}
where \(\operatorname{Min}_{\preceq}\) retains the nondominated vectors in
the existing componentwise resource order.  These definitions use the same
presentation law, target, verifier, and horizon as the saturated profile.

\end{definition}

\begin{theorem}[Saturated headroom and simultaneous benchmark comparison]
\label{thm-saturated-pareto-benchmark-comparison}

Let \(\mathcal A\) be as in
\cref{def-represented-benchmark-pareto-frontier}, and let
\(\widehat R\) be a deployed represented learner.  At a common vector ceiling
\(\mathbf b\), define the saturated benchmark headroom
\begin{equation}
\label{eq-saturated-benchmark-headroom}
\operatorname{Head}_{\mathcal A}(\mathbf b)
=
\operatorname{ActArr}_{\mathfrak M,H}^{\rm sat}(\mathbf b)
-\operatorname{BenchArr}_{\mathcal A,H}(\mathbf b).
\end{equation}
Then \(\operatorname{Head}_{\mathcal A}(\mathbf b)\ge0\).

Suppose \([L_{\rm sat},U_{\rm sat}]\),
\([L_{\mathcal A},U_{\mathcal A}]\), and
\([L_{\widehat R},U_{\widehat R}]\) are simultaneous valid intervals for
the saturated profile, benchmark frontier, and learner arrival.  On their
common event,
\begin{align}
\label{eq-saturated-benchmark-headroom-interval}
\bigl[L_{\rm sat}-U_{\mathcal A}\bigr]_+
&\le \operatorname{Head}_{\mathcal A}(\mathbf b)
\le U_{\rm sat}-L_{\mathcal A},\\
\label{eq-learner-benchmark-additive-gain-interval}
L_{\widehat R}-U_{\mathcal A}
&\le
\mathfrak M(A_{\widehat R}^H)
-\operatorname{BenchArr}_{\mathcal A,H}(\mathbf b)
\le U_{\widehat R}-L_{\mathcal A}.
\end{align}
For a declared denominator floor \(\epsilon_0>0\), the multiplicative gain
obeys
\begin{equation}
\label{eq-learner-benchmark-multiplicative-gain-interval}
\frac{L_{\widehat R}}{\max\{U_{\mathcal A},\epsilon_0\}}
\le
\frac{\mathfrak M(A_{\widehat R}^H)}
     {\max\{\operatorname{BenchArr}_{\mathcal A,H}(\mathbf b),\epsilon_0\}}
\le
\frac{U_{\widehat R}}{\max\{L_{\mathcal A},\epsilon_0\}}.
\end{equation}

If the learner certificate gives
\begin{equation}
\label{eq-benchmark-learner-saturated-slack-assumption}
\mathfrak M(A_{\widehat R}^H)
\ge
\operatorname{ActArr}_{\mathfrak M,H}^{\rm sat}(\mathbf b)
-\rho_{\rm learn},
\end{equation}
then
\begin{equation}
\label{eq-benchmark-learner-headroom-realization}
\mathfrak M(A_{\widehat R}^H)
-\operatorname{BenchArr}_{\mathcal A,H}(\mathbf b)
\ge
\operatorname{Head}_{\mathcal A}(\mathbf b)-\rho_{\rm learn}.
\end{equation}
Thus the learner realizes the saturated improvement over the represented
benchmark family up to its already charged learning slack.

\end{theorem}

\begin{proof}
Every benchmark record is an admissible represented record in the saturated
class at the same ceiling.  Hence
\(\operatorname{BenchArr}_{\mathcal A,H}(\mathbf b)
\le\operatorname{ActArr}_{\mathfrak M,H}^{\rm sat}(\mathbf b)\), proving
nonnegativity.  Subtracting the endpoints of simultaneous intervals gives
\cref{eq-saturated-benchmark-headroom-interval,eq-learner-benchmark-additive-gain-interval}.
Monotonicity of division on \((0,\infty)\), after applying the declared floor,
gives \cref{eq-learner-benchmark-multiplicative-gain-interval}.  Finally,
subtract the benchmark frontier from
\cref{eq-benchmark-learner-saturated-slack-assumption} and use
\cref{eq-saturated-benchmark-headroom}.
\end{proof}

\begin{corollary}[Matched-success resource savings and Levin gain]
\label{cor-matched-success-resource-and-levin-gain}

Fix \(\alpha\in(0,1)\).  Let \(\mathbf c_{\widehat R}\) be the learner's
complete cost and let \(\mathbf c_a\in\mathcal P_{\mathcal A,H}(\alpha)\).
For each positive resource coordinate \(j\), define the matched-success
logarithmic saving
\begin{equation}
\label{eq-matched-success-log-resource-saving}
\operatorname{Save}_{j}(\widehat R;a,\alpha)
=\log_2\frac{c_{a,j}}{c_{\widehat R,j}}.
\end{equation}
It is positive precisely when the learner uses less of resource \(j\) than
that nondominated benchmark at the same certified arrival level.

If \(a\) has canonical code length \(d(a)\), the blind physical ceiling for
reproducing it is the fourth coordinate of
\cref{eq-exact-represented-levin-benchmark-row}.  When the simultaneous
condition in
\cref{eq-witnessed-effective-length-empirical-condition} holds through
\(L=d(a)\), the certified description-length reach beyond blind search is
\begin{equation}
\label{eq-benchmark-specific-levin-length-gain}
\Delta^{\mathcal A}
\ge d(a)-\ell_{\rm cert}.
\end{equation}

\end{corollary}

\begin{proof}
The first assertion is the definition of a logarithmic ratio.  The second is
\cref{prop-simultaneous-witnessed-effective-length-certificate} followed by
\cref{eq-levin-structural-length-advantages}.
\end{proof}

\begin{theorem}[Finite-grid nonidentifiability of the asymptotic length profile]
\label{thm-finite-grid-nonidentifiability-length-profile}

Fix any finite grid \(\mathcal G\subset\mathbb N^2\) of size--length pairs
and any finite collection of valid profile bounds evaluated only on
\(\mathcal G\).  Those bounds alone do not imply that, for every fixed
\(\ell\),
\begin{equation}
\label{eq-asymptotic-fixed-length-negligibility-property}
\operatorname{ActArr}_{\mathfrak M,n,H,\le\ell}^{\rm sat}
(\operatorname{poly}(n))
\text{ is negligible as }n\to\infty.
\end{equation}

\end{theorem}

\begin{proof}
Let \(n_{\max}\) and \(\ell_{\max}\) be the largest coordinates inspected
by \(\mathcal G\).  Extend any presentation outside the inspected slices by
adjoining one family-uniform record of length \(\ell_{\max}+1\) that returns
the neutral output for \(n\le n_{\max}\) and reaches the target at polynomial
cost for \(n>n_{\max}\).  This changes none of the finitely evaluated profile
values, but it violates
\cref{eq-asymptotic-fixed-length-negligibility-property} at the fixed length
\(\ell_{\max}+1\).  Hence the finite statements do not entail the asymptotic
quantifier.
\end{proof}

\subsection{LPN verifier and parameter-explicit specialization}
\label{subsec-levin-lpn-specialization}

\begin{corollary}[Length-bounded Levin relation for LPN]
\label{cor-length-bounded-levin-relation-lpn}

In the LPN presentation of \cref{def-lpn-task-object}, let
\(
V_I(x)=\mathbf1_{\{|Ax+b|_1\le\tau m\}}
\)
with
\(
\eta<\tau<H_2^{-1}(1-R)
\)
and \(R=n/m\).  A direct matrix--vector evaluation has
\begin{equation}
\label{eq-lpn-levin-verifier-cost}
t_V(n,m)\le c_Vmn
\end{equation}
bit operations for a fixed implementation constant \(c_V\), and
\begin{equation}
\label{eq-lpn-levin-global-verifier-error}
\varepsilon_V(n,m,\eta,\tau)
\le
\exp[-mD_{\rm Ber}(\tau\Vert\eta)]
+(2^n-1)2^{-m(1-H_2(\tau))}.
\end{equation}
Consequently, the uniform length-class simulation ceiling is
\begin{equation}
\label{eq-lpn-levin-length-class-ceiling}
b_{\mathfrak U}^{\rm LPN}(\ell;b,n,m,H)
=\Lambda_{\mathfrak U}\!\left(
\ell,
\sup_{d\le\ell}\sigma_{\rm int}(b,d)+c_VHmn
\right),
\end{equation}
Put
\begin{equation}
\label{eq-lpn-levin-compiled-horizon}
\widetilde H_{\mathfrak U}^{\rm LPN}
=H_{b_{\mathfrak U}^{\rm LPN}(\ell;b,n,m,H)}.
\end{equation}
At physical ceiling \(b_{\rm phys}\), the LPN blind frontier is therefore
\begin{align}
\label{eq-lpn-levin-certifiable-length-frontier}
\ell_{{\rm cert},{\rm LPN}}
&=\max\left\{\ell:
b_{\mathfrak U}^{\rm LPN}(\ell;b,n,m,H)\le b_{\rm phys}\right\},\\
\label{eq-lpn-levin-certifiable-length-scaling}
\ell_{{\rm cert},{\rm LPN}}
&=
\log_2\frac{b_{\rm phys}}
{\sup_{d\le\ell_{{\rm cert},{\rm LPN}}}
 \sigma_{\rm int}(b,d)+c_VHmn+1}
-O\!\left(\log\log(b_{\rm phys}+b+Hmn+2)\right).
\end{align}
In the direct-verification-dominated range this reduces to
\(\log_2 b_{\rm phys}-\log_2(Hmn)-O(\log\log b_{\rm phys})\).
Then \cref{eq-length-class-profile-levin-bound} becomes
\begin{align}
\label{eq-lpn-length-bounded-levin-profile-bound}
\operatorname{ActArr}_{{\rm LPN},n,m,\eta,H,\le\ell}^{\rm sat}(b)
\le{}&
\mathfrak M_{n,m,\eta}\!\left(
A_{\mathfrak U_\ell}^{\widetilde H_{\mathfrak U}^{\rm LPN}}
(b_{\mathfrak U}^{\rm LPN})\right)\notag\\
&+\exp[-mD_{\rm Ber}(\tau\Vert\eta)]
+(2^n-1)2^{-m(1-H_2(\tau))}.
\end{align}
For the learned LPN record, the right-hand side of
\cref{eq-structural-learner-length-profile-lower} additionally carries the
explicit noise-dependent action-value radius
\begin{equation}
\label{eq-lpn-levin-learner-radius}
\Gamma_{\rm act}^{\rm clean}
\le
\frac1{1-2\eta}
\left[
2\sqrt{\frac{2L_{\rm act}(n,b,m,\eta)}{m_{\rm eff}}}
+3\sqrt{\frac{\log(2\widetilde H_{\mathfrak U}^{\rm LPN}/\delta)}
                   {2m_{\rm eff}}}
+\varepsilon_{\rm dep}
\right]
\end{equation}
and finite-horizon learning loss at most
\begin{equation}
\label{eq-lpn-levin-learner-total-loss}
\widetilde H_{\mathfrak U}^{\rm LPN}
\bigl(2\Gamma_{\rm act}^{\rm clean}+\varepsilon_{\rm opt}\bigr)
+\Delta_{\rm fil}^{\widetilde H_{\mathfrak U}^{\rm LPN}}.
\end{equation}
For a canonical neural LPN implementation, the sharp complete learning
candidate is
\begin{equation}
\label{eq-lpn-levin-sharp-learning-slack}
\rho_{{\rm learn},{\rm LPN}}^\sharp
=
\min\left\{
\begin{array}{l}
\widetilde H_{\mathfrak U}^{\rm LPN}
(2\Gamma_{\rm act}^{\rm clean}+\varepsilon_{\rm opt})
 +\Delta_{\rm fil}^{\widetilde H_{\mathfrak U}^{\rm LPN}},\\
\bigl[U_{{\rm Bell},{\rm LPN}}-L_{{\rm ev},{\rm LPN}}\bigr]_+,\\
B_{{\rm det},{\rm LPN}}
(p,K,\boldsymbol\tau,b_{\mathfrak U}^{\rm LPN},
\widetilde H_{\mathfrak U}^{\rm LPN})
 +V_{{\rm stat},{\rm LPN}}(p,N,\eta,\delta),\\
\mathcal L_{{\rm cert},{\rm LPN}}\wedge1
\end{array}
\right\}.
\end{equation}
Every candidate retains the same \(n,m,\eta,H\), compiled horizon
\(\widetilde H_{\mathfrak U}^{\rm LPN}\), architecture, transition law,
target gate, and compiled ceiling.  The finite-description first row
may be replaced by any smaller simultaneous sequential-Rademacher,
martingale PAC--Bayes, mixing, or predictable-variance radius admitted by
\cref{thm-complete-certified-bias-variance-decomposition}.

\begin{samepage}
The length-bounded profile also inherits every unrestricted LPN upper
certificate.  In particular, define
\begin{equation}
\label{eq-lpn-levin-bellman-upper-candidate}
U_{{\rm Bell},{\rm LPN}}^{\rm sat}
=
\inf_{v_H=0}
\left\{
\mathfrak M_{n,m,\eta}(v_0(h_0))
+\sum_{t=0}^{H-1}
\max_{c\in\mathfrak C_{\rm master}}
\mathfrak U_{t,c}^{\rm LPN}(v)
\right\}
+\varepsilon_{\rm good}(n,m,\eta,\tau,\delta).
\end{equation}
Then the sharp integrated length-class upper bound is
\begin{align}
\label{eq-lpn-levin-integrated-length-profile-minimum}
\operatorname{ActArr}_{{\rm LPN},n,m,\eta,H,\le\ell}^{\rm sat}(b)
\le
\min\Bigl\{&
\mathfrak M_{n,m,\eta}\!\left(
A_{\mathfrak U_\ell}^{\widetilde H_{\mathfrak U}^{\rm LPN}}
(b_{\mathfrak U}^{\rm LPN})\right)
+\varepsilon_V(n,m,\eta,\tau),\notag\\
&U_{{\rm Bell},{\rm LPN}}^{\rm sat}
\Bigr\}.
\end{align}
\end{samepage}
The Bellman candidate expands mode by mode through
\cref{eq-lpn-nonambient-bellman-support-bound,eq-lpn-ambient-bellman-support-explicit-minimum};
hence its information, capacity, Blackwell, source, quotient-compensation,
direct-contact, and residual-Gibbs terms remain available to the minimum.
Thus \(n,m,\eta,H\), description length, action-class capacity, dependent
sample size, optimization, filtering, verifier error, and physical compute
all occur in displayed coordinates.
For a represented subset-Gaussian-elimination benchmark, the arrival entry
in \cref{eq-exact-represented-levin-benchmark-row} is exactly
\((1-\eta)^n\) on the full-rank event and equals
\(n^{-cR+o(1)}\) when \(\eta m=c\log n\), by
\cref{eq-lpn-subset-ge-success,eq-lpn-subset-ge-log-noise-window}.  A
represented located-spectral benchmark instead carries the exponent and
budget threshold in
\cref{eq-lpn-dynamics-located-spectral-exponent,eq-lpn-dynamics-located-spectral-threshold}.
Its complete fixed-data and streaming BKW realization, including sample reuse,
decoder error, time, data, memory, and canonical code length, is
\cref{thm-lpn-dependence-aware-bkw-recovery,cor-lpn-quantitative-improvement-over-bkw}.
These are benchmark-specific entries of the witnessed profile; the saturated
length-class comparison remains
\cref{eq-lpn-levin-integrated-length-profile-minimum}.

\end{corollary}

\begin{proof}
The verifier cost is the cost of computing \(Ax+b\) and its Hamming weight.
The global gate event and its binomial/random-code bound are
\cref{eq-lpn-verifier-good-event-probability}.  Substitute these values into
\cref{eq-levin-uniform-length-class-ceiling,eq-length-class-profile-levin-bound}.
The learning radius and finite-horizon propagation are
\cref{eq-active-clean-finite-description-radius,eq-active-explicit-sample-noise-value-loss},
with the represented filter comparison added as in
\cref{cor-learned-structural-bayesian-control}.  The minimum in
\cref{eq-lpn-levin-sharp-learning-slack} follows from
\cref{eq-sharp-complete-levin-learning-slack}.  Finally, a restricted
length class is contained in the unrestricted active class; combine
\cref{eq-lpn-length-bounded-levin-profile-bound} with
\cref{eq-lpn-bellman-support-family-bound} to obtain
\cref{eq-lpn-levin-integrated-length-profile-minimum}.
\end{proof}

\begin{figure}[tbp]
\centering
\resizebox{.98\linewidth}{!}{%
\begin{tikzpicture}[fw diagram,node distance=8mm and 10mm]
  \node[fw source,text width=2.8cm] (class)
    {active records\\cost \(b\), length \(\le\ell\)};
  \node[fw provenance,text width=2.8cm,right=of class] (levin)
    {Kraft scheduler\\and universal interpreter};
  \node[fw use,text width=2.8cm,above right=7mm and 10mm of levin] (gate)
    {represented target gate\\error \(\varepsilon_V\)};
  \node[fw geom,text width=2.8cm,below right=7mm and 10mm of levin] (learn)
    {structural learner\\slack \(\rho_{\rm learn}\)};
  \node[fw box,text width=3.0cm,right=16mm of levin] (ceil)
    {common ceiling\\\(\Lambda_{\mathfrak U}(\ell,s_\ell)\)};
  \node[fw terminal,text width=3.0cm,right=of ceil] (out)
    {certified length profile\\and \(\Delta^{\mathcal A}\)};
  \draw[fw arrow] (class)--(levin);
  \draw[fw arrow] (levin)--(gate);
  \draw[fw arrow] (levin)--(learn);
  \draw[fw arrow] (gate)--(ceil);
  \draw[fw arrow] (learn)--(ceil);
  \draw[fw arrow] (ceil)--(out);
\end{tikzpicture}%
}
\caption{Length-bounded Levin comparison.  Blind universal search pays the
Kraft factor \(2^\ell\) together with interpreter, scheduler, and verifier
costs.  The structural learner is compared at the same compiled ceiling and
carries its complete statistical, optimization, filtering, and deployment
slack.  Empirical comparisons certify the witnessed frontier
\(\ell_{\rm eff}^{\mathcal A}\); the true frontier retains the saturated
length-class supremum.}
\label{fig-length-bounded-levin-structural-learner}
\end{figure}

\section{Summary of the Dynamical-Systems Geometry}
\label{sec-controlled-dynamics-summary}

\begin{center}
\scriptsize
\begin{tabular}{>{\raggedright\arraybackslash}p{.23\linewidth}
                >{\raggedright\arraybackslash}p{.31\linewidth}
                >{\raggedright\arraybackslash}p{.36\linewidth}}
\toprule
Coordinate & Structural record & Framework result \\
\midrule
finite-horizon arrival & pre-hit selector and neutral baseline &
\cref{thm-prospective-selector-accountability,thm-prospective-selector-structural-charge} \\
discounted arrival & killed generator and target gate &
\hyperref[def-target-resolvent]{target resolvent};
\hyperref[def-master-compensated-quotient-geometry-certificate]{master certificate} \\
expected hitting time & authorized drift potential &
\cref{thm-controlled-reach-avoid-dirichlet-drift} \\
reach--avoid & committor and condenser geometry &
\cref{thm-controlled-committor-bellman,thm-controlled-reach-avoid-dirichlet-drift} \\
mixing and compression & invariant reversible carrier &
\cref{thm-controlled-certified-mixing} \\
functional rates & audited typed inequality profile &
\cref{thm-controlled-functional-inequality-consequences,thm-controlled-functional-profile-selection} \\
authorized target flow & Caus action, aiming energy, and contact converter &
\cref{thm-controlled-caus-action-aiming,cor-controlled-uniform-caus-flow-bound} \\
noise-corrected Caus flow & clean-target comparison and corrupted binary
supervision &
\cref{thm-uniform-dynamical-label-noise-bound,thm-controlled-noise-corrected-caus-aiming} \\
lifted condenser capacity & valid Lift and boundary-zero fiber form &
\cref{thm-controlled-lift-capacity-schur,def-controlled-saturated-lift-capacity-amplification} \\
lifted target access & intertwining defect and fiber resolvent &
\cref{thm-controlled-valid-lift-block-resolvent,thm-controlled-clocked-lift-transfer} \\
noisy Lift interface & target/failure deployment gates, killed contact
vector, and joint capacity stability &
\cref{def-controlled-noisy-gate-profiles,thm-controlled-noisy-gate-lift-transfer} \\
reward and execution noise & typed reward correction, actual executed kernel,
and Bellman/path converters &
\cref{thm-dynamical-affine-reward-channel-transport,def-controlled-action-execution-channel,thm-controlled-execution-channel-comparison} \\
observation and POMDP noise & public filtration, clocked latent carrier,
filter stability, and gate transfer &
\cref{thm-controlled-pomdp-clocked-embedding,thm-controlled-pomdp-filter-transfer} \\
multichannel entropy & conditional capacity, path KL/TV, posterior entropy,
Fano, and rate--distortion &
\cref{thm-controlled-multichannel-path-comparison,thm-controlled-multichannel-information-bound,thm-controlled-belief-entropy-balance,def-controlled-typed-entropy-outputs} \\
optimal active sampling & represented acquisition optimizer, declared
experiment, selected transport, and master quotient--geometry certificate &
\cref{thm-active-optimality-saturated-construction,thm-master-optimal-structural-active-sampling-envelope,thm-learned-active-rule-finite-horizon-loss} \\
affordable occupation duality & stopped history flow, target flux, and
resource-indexed Bellman barrier &
\cref{thm-exact-affordable-occupation-flow-primal,thm-bellman-barrier-duality-affordable-arrival,thm-structural-evaluation-affordable-bellman-support} \\
optimal structural Bayesian learning & declared task prior, represented
likelihood and belief, source-resolved actionable gain, and charged acquisition &
\cref{thm-structural-bayes-prior-provenance-neutrality,thm-structural-bayesian-target-contact-accountability,thm-exact-source-resolved-structural-bayesian-gain,thm-noise-contracted-source-resolved-bayesian-gain,thm-optimal-structural-bayesian-acquisition} \\
trajectory learning & mixing blocks, predictable history trees,
and martingale priors &
\cref{thm-beta-mixing-structural-blocking};
\cref{thm-sequential-learned-operator-cover};
\cref{thm-martingale-structural-pac-bayes} \\
action and safety capacity & action-generated and failure-killed operator
compressions &
\cref{thm-constrained-safe-spectral-capacity} \\
value and policy loss & represented Bellman residual &
\cref{thm-controlled-bellman-residual-policy-loss} \\
learned control & learner--geometry--policy record &
\hyperref[cor-saturated-learned-policy-accountability]{policy accountability};
\hyperref[thm-master-structural-dynamical-learning-certificate]{master certificate} \\
certified neural implementation & complete network and optimizer transcript,
held-out arrival audit, and saturated Bellman barrier &
\cref{cor-primal-dual-empirical-implementation-certificate,thm-certified-deep-network-implementation-bridge,thm-wadam-descent-global-optimization-certificate} \\
canonical neural controller & recurrent belief token, masked attention,
policy/value/barrier heads, and explicit parameter ledger &
\cref{def-canonical-certified-attention-belief-architecture,thm-canonical-attention-architecture-implementation,thm-canonical-loss-saturated-arrival-calibration} \\
bias--variance diagnosis & deterministic approximation ledger, predictable
variance, noise transport, and simultaneous tuning &
\cref{thm-complete-certified-bias-variance-decomposition,prop-exact-clipping-bias-variance-identity,cor-quantitative-architecture-audit-tradeoffs} \\
length-bounded universal search & prefix-free active-record code, Kraft
scheduler, verifier gate, and structural-learning slack &
\cref{thm-length-bounded-levin-realization,thm-structural-learner-length-bounded-levin,def-levin-certifiable-length-frontier} \\
empirical analytical proof & parameter profile, source audit, feature
Dirichlet form, and quotient-row diameter &
\cref{thm-certified-underfitting-optimization-intervals,thm-parameter-source-resolved-empirical-target-gain,thm-empirical-poincare-bracket-full-transfer,thm-finite-sample-lumpability-certificate,thm-complete-empirical-structural-proof-package} \\
continuous transfer & valid Lift and generator defect &
\cref{thm-controlled-continuum-lift-transfer} \\
\bottomrule
\end{tabular}
\end{center}

The universal conclusion is
\cref{thm-universal-dynamical-certificate}: every represented controlled
system is classified by the existing structural channels and receives a
coordinatewise certificate at its complete saturated cost.  Arrival, reach,
mixing, policy error, and
generalization retain their own exact identities and the strongest applicable
represented comparison.  Target-directed Caus flow enters an arrival bound
through an aiming-to-contact converter.  Lifted access enters through a form,
intertwining, clocked-defect, or fiber-resolvent comparison.  The unrestricted
valid-lift class remains the explicit saturated lift-access supremum of
\cref{def-controlled-saturated-lift-access-envelope}.
The stopped occupation-flow identity and Bellman barrier duality in
\cref{thm-exact-affordable-occupation-flow-primal,thm-bellman-barrier-duality-affordable-arrival}
also evaluate the best affordable active controller.  Their support terms
are derived coordinates over the five existing master modes.  The learned
controller is compared with this optimum by
\cref{cor-complete-learned-control-sandwich}; estimation and optimization
errors are not substituted for the structural optimum.
The implemented neural controller is connected to the same optimum by
\cref{cor-primal-dual-empirical-implementation-certificate}.  A held-out
execution audit provides the primal lower certificate, while an exact or
coverage-certified Bellman audit provides the saturated upper certificate.
The canonical implementation in
\cref{def-canonical-certified-attention-belief-architecture} fixes its
tokenizer, recurrent belief, masked attention, policy, value, barrier,
geometry, and quotient interfaces and counts their parameters and operations
explicitly.  Its trainable policy and likelihood losses are calibrated to
saturated arrival by
\cref{thm-canonical-loss-saturated-arrival-calibration}; its analytical
optimality objective is the primal--dual gap
\cref{eq-canonical-primal-dual-certificate-loss}.
The WAdam specialization in
\cref{thm-wadam-descent-global-optimization-certificate} converts its
represented descent, second-moment, PL, and objective-to-action comparisons
into the optimizer coordinate of the end-to-end gap.
The bias--variance ledger in
\cref{thm-complete-certified-bias-variance-decomposition} separates
representation, optimization, filtering, precision, truncation,
sufficiency, and temperature effects from the simultaneous statistical
radius.  Parameter count, temperature, label noise, horizon, clipping,
latent capacity, feature dimension, and quotient resolution are selected
through the complete simultaneous upper curve in
\cref{cor-quantitative-architecture-audit-tradeoffs}.
The parameter-indexed profile in
\cref{def-parameter-indexed-saturated-neural-profile} separates certified
underfitting from optimizer error and brackets the minimum sufficient
parameter count.  The same run can attribute learned target gain to the five
structural sources by
\cref{thm-parameter-source-resolved-empirical-target-gain}, transfer
feature-level Dirichlet estimates to true analytical functional-inequality
constants by
\cref{thm-empirical-poincare-bracket-full-transfer,thm-generic-empirical-functional-inequality-transfer},
and certify exact or quasi-lumpability by
\cref{thm-finite-sample-lumpability-certificate,thm-certified-fiber-diameter-lumpability-transfer}.
Exact lumpability is obtained when every pair of states carrying the same
label has the same probability of entering each next label: this is exactly
the row-diameter condition \(\delta_q(P)=0\), and it yields the intertwining
\(PU_q=U_q\bar P\).  A nonzero certified diameter remains an explicit
finite-horizon and generator defect rather than being treated as an exact
quotient.
The structural Bayesian restriction in
\cref{sec-optimal-structural-bayesian-learning} separates the declared task
prior, analytic task posterior, represented belief, and PAC--Bayes posterior.
Prior selection independent of the target retains neutral contact, while a
target-dependent initialization is the first charged observation and belief
update.  Posterior concentration bounds target contact by
\cref{thm-structural-bayesian-target-contact-accountability}, and every
prospectively used concentration readout retains the same five-mode
provenance by
\cref{thm-affordable-posterior-concentration-source-routing}.  The deployed
next-test improvement is the actionable gain of
\cref{def-source-resolved-actionable-bayesian-gain}; its exact five-source
charge, noise-contracted information ledger, and adaptive composition are
\cref{thm-exact-source-resolved-structural-bayesian-gain,thm-noise-contracted-source-resolved-bayesian-gain,cor-source-resolved-bayesian-adaptive-composition}.
When a binary target or failure interface is learned from noisy labels, the
same coordinates additionally carry the explicit clean-risk factor
\((1-2\eta_+)^{-1}\), the affordable ceiling \(b(n)\), and sample size \(m\).
Only represented same-norm calibration, deployment, contact-vector, and
reference-law comparisons propagate that statistical bound into Caus action,
Lift arrival, or condenser capacity.  The native deterministic capacity and
resolvent inequalities themselves are unchanged.
Reward, execution, observation, and filter corruption are now treated by the
same discipline.  Their raw \((n,b(n),m)\)-indexed errors remain typed;
executed actions define the actual generator, observation-measurable policies
define the affordable POMDP class, and filter errors enter arrival or return
only through the clocked or Bellman converters.  Physical-channel entropy,
stochastic-transition entropy, policy entropy, public channel capacity in
bits, path KL in nats, MLSI density entropy, Caus edge action, PAC--Bayes KL,
and description length remain distinct coordinates.  The
only target-contact use of Caus entropy production is the same-record product
with aiming followed by a contact converter.  Multichannel computational
monotonicity is asserted only along represented Blackwell garblings.

\clearpage
\part{Attention as Structural Gradient Flow}\label{part:attention}\setcounter{section}{-1}

\section{Represented Attention Records}
\label{sec-attention-records}

\paragraph{Part contract.}
\emph{Inputs:} the saturated constructor algebra of \PartsItoIII, the
source-resolved learning ledger of \PartIV, and the controlled-dynamics,
active-sampling, and structural-Bayesian certificates of \PartV.
\emph{Output:} an exact representation and source-charging theorem for
attention layers and their deployed transition kernels.
\emph{First pass:} read
\cref{def-complete-attention-record,thm-attention-gibbs-gradient-flow,thm-attention-structural-source-compilation,thm-attention-target-gain-envelope}.

Attention is used here as a represented mechanism for selecting and
transporting public values.  It is not adjoined as a new structural channel.
Every attention coordinate is compiled into the already exhaustive
\(\Geom,\Caus,\Abs,\Lift,\Bdry,\Res\) ledger and retains the five-family
constructor provenance of its score and value records.
The canonical finite-precision realization is
\cref{def-canonical-certified-attention-belief-architecture}.  Its exact
parameter and operation ledgers are
\cref{eq-canonical-attention-parameter-count,eq-canonical-attention-forward-cost},
and its policy, belief, critic, and certificate objectives are
\cref{def-canonical-attention-belief-training-objective}.

\begin{definition}[Complete represented attention record]
\label{def-complete-attention-record}

A complete represented attention record on a finite token carrier
\(\mathcal I\) is
\[
\mathcal R^{\rm att}
=\bigl(\mathcal R,\Theta_{\rm att},\operatorname{Tok},\nu,
Q,K,V,O,s,P_s,Y,\operatorname{Deploy}\bigr),
\]
where \(\mathcal R\) is the complete ancestral record and
\[
\Theta_{\rm att}
=(L,(H_\ell,N_\ell,d_\ell,d_{k,\ell})_{\ell\le L},
(\tau_{\ell h})_{\ell,h},p,R_Q,R_K,R_V,R_O)
\]
records depth, heads, token counts, widths, key widths, temperatures,
arithmetic precision, and represented parameter records.  For one head,
\(\nu(\cdot\mid i)\) is the represented masked reference kernel,
\[
s(i,j)=\langle Qz_i,Kz_j\rangle,
\qquad
P_s(j\mid i)
=\frac{\nu(j\mid i)e^{s(i,j)/\tau}}
       {\sum_k\nu(k\mid i)e^{s(i,k)/\tau}},
\qquad
Y_i=O\sum_jP_s(j\mid i)Vz_j.
\]
Multihead concatenation, residual addition, normalization, feed-forward
maps, recurrence, and a final selector are represented constructors in the
same record.  The complete attention cost
\(\operatorname{Cost}_{\rm att}(\mathcal R^{\rm att})\) charges tokenization,
masks, all parameter derivations, score evaluation, normalization at precision
\(p\), value transport, heads, residual and feed-forward paths, training,
selection, and deployment.  A target-dependent parameter, mask, or
initialization is charged as the first target-dependent ancestor.

The record is \emph{temporally admissible} when every row deployed after time
\(t\) factors through the pre-hit record of
\cref{def-pre-hit-temporally-admissible-source}.  It is
\emph{budget admissible} when its complete cost lies in the one declared
saturated ceiling.

\end{definition}

\begin{figure}[tbp]
\centering
\begin{tikzpicture}[fw diagram,node distance=6mm and 7mm,scale=.92,transform shape]
\node[fw source] (obs) {public tokens};
\node[fw geom,right=of obs] (score) {represented score};
\node[fw use,right=of score] (soft) {Gibbs row \(P_s\)};
\node[fw provenance,right=of soft] (val) {value transport};
\node[fw terminal,right=of val] (sel) {deployed selector};
\draw[fw arrow] (obs)--(score);
\draw[fw arrow] (score)--(soft);
\draw[fw arrow] (soft)--(val);
\draw[fw arrow] (val)--(sel);
\node[fw residual,below=of soft] (cost) {one complete cost and ancestry record};
\draw[fw dashed arrow] (obs)--(cost);
\draw[fw dashed arrow] (score)--(cost);
\draw[fw dashed arrow] (soft)--(cost);
\draw[fw dashed arrow] (val)--(cost);
\draw[fw dashed arrow] (sel)--(cost);
\end{tikzpicture}
\caption{A represented attention head.  Normalization and value transport do
not erase the ancestry or cost of the score used by the deployed selector.}
\label{fig-attention-complete-record}
\end{figure}

\section{Gibbs Geometry and Exact Structural Charging}
\label{sec-attention-gibbs-geometry}

\begin{theorem}[Gibbs variational identity and Fisher--Rao attention flow]
\label{thm-attention-gibbs-gradient-flow}

Fix a row \(i\), write \(\nu_j=\nu(j\mid i)\), and assume
\(\nu_j>0\) on its represented support.  For every probability vector
\(p\ll\nu\),
\begin{equation}
\label{eq-attention-gibbs-variational-identity}
\frac1\tau\mathbb E_p s-\operatorname{KL}(p\Vert\nu)
=\log\mathbb E_\nu e^{s/\tau}
-\operatorname{KL}(p\Vert P_s).
\end{equation}
Consequently \(P_s\) is the unique maximizer of the left-hand functional.
Along
\[
p_t(j)=\frac{\nu_je^{t s_j/\tau}}
              {\sum_k\nu_ke^{t s_k/\tau}},
\qquad 0\le t\le1,
\]
one has
\begin{align}
\label{eq-attention-fisher-rao-flow}
\dot p_t(j)&=\tau^{-1}p_t(j)
  \bigl(s_j-\mathbb E_{p_t}s\bigr),\\
\label{eq-attention-fisher-action}
\lVert\dot p_t\rVert_{\rm FR}^2
&=\tau^{-2}\operatorname{Var}_{p_t}(s),\\
\label{eq-attention-kl-cost}
\operatorname{KL}(P_s\Vert\nu)
&=\tau^{-1}\mathbb E_{P_s}s-\log\mathbb E_\nu e^{s/\tau}.
\end{align}
Thus softmax is a Fisher--Rao natural-gradient path on the represented row
simplex.  This statement does not identify an arbitrary parameter-training
trajectory with Euclidean gradient flow.

\end{theorem}

\begin{proof}
Substitute
\(\log(P_s(j)/\nu_j)=s_j/\tau-\log\mathbb E_\nu e^{s/\tau}\)
in \(\operatorname{KL}(p\Vert P_s)\) and rearrange.  Nonnegativity and strict
convexity of relative entropy give the maximizer.  Differentiating the
normalized exponential family proves
\cref{eq-attention-fisher-rao-flow}.  The Fisher--Rao metric is
\(\lVert v\rVert_{\rm FR}^2=\sum_jv_j^2/p_j\); substitution proves
\cref{eq-attention-fisher-action}.  Setting \(p=P_s\) in
\cref{eq-attention-gibbs-variational-identity} proves
\cref{eq-attention-kl-cost}.
\end{proof}

\begin{corollary}[Bayesian posterior as represented attention flow]
\label{cor-bayesian-posterior-attention-flow}

Let \(\mathcal C_t\) be a finite represented latent hypothesis or cell carrier
retained by a compact latent-belief record, let \(q_t\) be its current
full-support belief on \(\mathcal C_t\), and let \(L_t(y\mid z)>0\) be the
represented likelihood of the acquired observation.  The stored code for
\(q_t\), its carrier, and its decoder are included in the latent capacity and
complete cost of \cref{def-complete-compact-latent-belief-record}.
Choose one attention row with
\begin{equation}
\label{eq-bayes-attention-score-identification}
\nu_t(z)=q_t(z),
\qquad
s_t(z)=\tau\log L_t(y\mid z).
\end{equation}
Then its Gibbs row is the exact latent Bayesian update:
\begin{equation}
\label{eq-bayes-attention-posterior-identity}
P_{s_t}(z)
=
\frac{q_t(z)L_t(y\mid z)}
{\sum_uq_t(u)L_t(y\mid u)}.
\end{equation}
The Fisher--Rao flow and action in
\cref{eq-attention-fisher-rao-flow,eq-attention-fisher-action} are therefore
the Bayes--Fisher flow and action in
\cref{eq-bayes-fisher-rao-flow,eq-bayes-fisher-rao-action} on the same
represented latent simplex.

Suppose a represented score \(\widehat\ell_t\) approximates
\(\ell_t=\log L_t(y\mid\cdot)\) on that support with
\begin{equation}
\label{eq-bayes-attention-logit-error}
\lVert\widehat\ell_t-\ell_t\rVert_\infty
\le\varepsilon_{\ell,t}.
\end{equation}
Let \(\widehat q_{t+1}\) be the normalized exponential update using
\(\widehat\ell_t\), and let \(q_{t+1}\) be the exact update.  Then
\begin{align}
\label{eq-bayes-attention-kl-error}
D_{\rm KL}(q_{t+1}\Vert\widehat q_{t+1})
&\le2\varepsilon_{\ell,t},\\
\label{eq-bayes-attention-tv-error}
\lVert q_{t+1}-\widehat q_{t+1}\rVert_{\rm TV}
&\le\sqrt{\varepsilon_{\ell,t}}.
\end{align}
The likelihood evaluator, logit approximation, exponentiation, normalization,
precision, latent storage, and downstream use are charged respectively through
the complete latent-belief and attention records.

\end{corollary}

\begin{proof}
Substitute \cref{eq-bayes-attention-score-identification} in the Gibbs row of
\cref{def-complete-attention-record}; this proves
\cref{eq-bayes-attention-posterior-identity}.  The two exponential paths then
coincide, which gives the Fisher--Rao assertion.

Write \(Z\) and \(\widehat Z\) for the exact and approximate normalizers.
The uniform logit error gives
\(e^{-\varepsilon_{\ell,t}}Z\le\widehat Z
\le e^{\varepsilon_{\ell,t}}Z\).  Hence
\[
\left|
\log\frac{q_{t+1}(z)}{\widehat q_{t+1}(z)}
\right|
\le2\varepsilon_{\ell,t}
\]
on the common support.  Taking expectation under \(q_{t+1}\) proves the KL
bound.  Pinsker's inequality proves the TV bound.
\end{proof}

\begin{theorem}[Attention structural-source compilation]
\label{thm-attention-structural-source-compilation}

Every temporally and budget-admissible complete attention record has the
following exhaustive source assignment, at its complete charged ceiling:
\begin{center}
\begin{tabular}{p{.28\linewidth}p{.60\linewidth}}
\toprule
attention coordinate & framework coordinate \\
\midrule
token/key carrier and reversible row law & \(\Geom\) \\
qualified nonidentity score fibers & exact or compensated \(\Abs\) \\
target-oriented score gradient/current & authorized \(\Caus\) of that \(\Geom\) \\
values, heads, residual paths, depth, memory & \(\Lift\) \\
masks, support truncation, terminal interface & \(\Bdry\) \\
gradient-orthogonal transport & \(\Res\) \\
score and value constructors & \(\mathsf{Add},\mathsf{Mult},\mathsf{Ord},
\mathsf{Sym},\mathsf{Filt}\) provenance \\
\bottomrule
\end{tabular}
\end{center}
No row normalization, head concatenation, residual connection, or clocked
reuse creates an additional source coordinate.

\end{theorem}

\begin{proof}
Apply selected-generator exhaustiveness
\cref{thm-controlled-policy-channel-decomposition} to each deployed row.
The reversible and Hodge components give \(\Geom\), authorized \(\Caus\),
and \(\Res\).  A represented score fiber passes through the exact or
compensated quotient gates of
\cref{thm-no-unclassified-geom-source-after-quotient-assignment}; a fiber
that fails those gates remains full-carrier \(\Geom\), not \(\Abs\).
Value transport and stored head state are represented contextual extensions,
so \cref{thm-exhaustive-lift-normal-form} assigns them to \(\Lift\).
Masks and support changes are assigned by
\cref{thm-exhaustive-bdry-normal-form}.  Finally apply the five-family
provenance theorem and constructor no-creation to the finite ancestral DAG.
All derived vertices retain their ancestors and complete cost, proving both
exhaustiveness and no creation.
\end{proof}

\begin{figure}[tbp]
\centering
\begin{tikzpicture}[fw diagram,node distance=7mm and 8mm]
\node[fw source] (att) {attention record};
\node[fw geom,above right=of att] (g) {\(\Geom\)};
\node[fw use,right=of att] (c) {\(\Caus\)};
\node[fw provenance,below right=of att] (a) {\(\Abs\)};
\node[fw provenance,above=of g] (l) {\(\Lift\)};
\node[fw terminal,right=of c] (b) {\(\Bdry\)};
\node[fw residual,below=of a] (r) {\(\Res\)};
\draw[fw arrow] (att)--(g);
\draw[fw arrow] (att)--(c);
\draw[fw arrow] (att)--(a);
\draw[fw arrow] (att)--(l);
\draw[fw arrow] (att)--(b);
\draw[fw arrow] (att)--(r);
\end{tikzpicture}
\caption{Attention introduces no seventh structural modality.  Every useful
coordinate is assigned to the existing exhaustive ledger.}
\label{fig-attention-channel-compilation}
\end{figure}

\section{Attention Target-Gain Envelope}
\label{sec-attention-target-gain-envelope}

\begin{definition}[Source-resolved attention ledger]
\label{def-source-resolved-attention-ledger}

For a complete attention record \(R\) deployed at time \(t\), let
\(g_{R,t}^{\rm att}\) be its positive pre-hit selector advantage over the
declared neutral row.  Its source information is
\[
J_{R,t}^{\rm att,src}
=\sum_{a\in\mathfrak J_{\Abs}^{5}}J_{R,t}^{\rm att,(a)},
\]
with the same five-family coordinates as
\cref{def-complete-source-resolved-learning-ledger}.  Let
\(\mathcal L_{R,t}^{\rm att,(a)}\) be the corresponding qualified source
visibility, including the \(\Geom/\Abs\) support mode, and let
\(U_{R,t}^{\rm KL},U_{R,t}^{\rm stat\to tgt},U_{R,t}^{\rm dyn\to tgt}\)
be the already-defined relative-entropy, statistical-calibration, and
dynamical contact converters for this same record.  Define
\begin{equation}
\label{eq-source-resolved-attention-envelope}
U_{R,t}^{\rm att}
=\min\left\{
U_{R,t}^{\rm KL},
U_{R,t}^{\rm stat\to tgt},
\beta_{R,t}+\sqrt{\frac{\log2}{2}J_{R,t}^{\rm att,src}},
eH_B\sum_a\mathcal L_{R,t}^{\rm att,(a)},
U_{R,t}^{\rm dyn\to tgt}
\right\}.
\end{equation}
An unavailable converter contributes the neutral value one and therefore
cannot improve the minimum.

\end{definition}

\begin{theorem}[Source-resolved attention target-gain bound]
\label{thm-attention-target-gain-envelope}

Every temporally admissible complete attention record satisfies
\begin{equation}
\label{eq-attention-target-gain-bound}
0\le g_{R,t}^{\rm att}\le U_{R,t}^{\rm att}.
\end{equation}
At a saturated ceiling \(B\),
\begin{equation}
\label{eq-saturated-attention-source-domination}
\operatorname{AttSel}^{\rm sat}(B)
\le \operatorname{Sel}^{\rm sat}(B)
\le I^{\rm src,sat}(B).
\end{equation}
The same bounds hold for finite-precision hard attention by taking the
represented zero-temperature limit at its charged precision and tie-breaking
cost.

\end{theorem}

\begin{proof}
The five entries in \cref{eq-source-resolved-attention-envelope} are,
respectively, the represented change-of-measure bound, the source-resolved
learning calibration bound, Pinsker's information-to-contact conversion,
the contextual source-transfer bound, and the controlled-dynamics arrival
converter.  Each is applied to the same complete record and common law; hence
their minimum is valid.  The first inequality in
\cref{eq-saturated-attention-source-domination} is inclusion of represented
attention selectors in the temporally admissible selector profile.  The
second is prospective selector source routing from
\cref{thm-no-uncharged-prospective-value-amplification,thm-source-resolved-saturated-learning-compilation}.
Finite-precision argmax is a represented \(\mathsf{Ord}\) postprocessing;
constructor no-creation retains the score ancestry and charges precision and
tie breaking.
\end{proof}

\begin{corollary}[Attention cannot bypass optimal active sampling]
\label{cor-attention-active-sampling-compilation}

At the class-preserving compilation ceiling of
\cref{thm-master-optimal-structural-active-sampling-envelope}, every
attention-driven acquisition or transition rule is one feasible complete
active-sampling record.  Its arrival, posterior gain, Caus action, Lift
transfer, and sparse-reward value are bounded by the corresponding
coordinates of that master envelope.  Conversely, every bounded-branch
clocked selector can be represented by finite-precision hard attention over
its represented legal actions: expose the selector's already represented
auxiliary seed, apply its represented inverse-transform comparisons, and give
the selected action the unique maximal score.  Averaging over that same seed
recovers the original kernel.  The seed, comparisons, tie breaking, ancestry,
and compilation cost are retained.  Thus attention changes the representation of the selector,
not the saturated structural optimum.
That optimum is the structural Bayesian and represented-computation ceiling
of \cref{thm-optimal-structural-extraction-saturated-ceiling}; its explicit
noise, sample, and target-size bound is
\cref{cor-noise-sample-target-size-optimal-extraction}.

\end{corollary}

\begin{proof}
The forward statement is inclusion.  For the converse, tokenize the finitely
many represented legal actions and expose the selector's already charged
auxiliary seed.  Its represented inverse-transform comparisons select one
legal action; assign that action the unique maximal score, mask illegal
actions, and take represented argmax with the original tie-breaking rule.
This reproduces the selector pathwise for each seed and therefore reproduces
its transition law after averaging over that seed.  The construction is
a clocked \(\Lift\) of the original selector and retains all source ancestors
by \cref{thm-attention-structural-source-compilation}.
\end{proof}

\begin{theorem}[Joint compact-latent attention and optimal acquisition]
\label{thm-joint-compact-latent-attention-acquisition}

Let \(\operatorname{LatAttArr}_{I,H}^{\rm sat}(B,b)\) be the restriction of
\cref{eq-saturated-compact-latent-arrival-profile} to complete records whose
represented latent update or deployed selector is implemented by
\cref{def-complete-attention-record}.  Then
\begin{equation}
\label{eq-latent-attention-profile-inclusion}
\operatorname{LatAttArr}_{I,H}^{\rm sat}(B,b)
\le
\operatorname{LatBayesArr}_{I,H}^{\rm sat}(B,b).
\end{equation}
At the class-preserving finite-precision hard-attention compilation ceiling,
every bounded-branch compact latent selector has an attention representation
with the same pre-hit path law, source ancestry, and latent code.  Therefore
the two profiles coincide on their saturated transfer class.  Taking the union
over code capacities admitted by the polynomial saturated budget gives
\begin{equation}
\label{eq-latent-attention-polynomial-ceiling}
\operatorname{LatAttArr}_{\mathfrak M,n}^{\rm poly,sat}
=\operatorname{LatBayesArr}_{\mathfrak M,n}^{\rm poly,sat}
=\operatorname{Succ}_{\mathfrak M,n}^{\rm poly,sat}.
\end{equation}

For every one of these records, the deployed gain is bounded by the same-record
minimum.  Writing \(g_{R,t}^{\rm lat,att}\) for the actionable Bayesian gain
\cref{eq-actionable-structural-bayesian-gain} of its deployed attention
selector,
\begin{equation}
\label{eq-joint-latent-attention-gain-certificate}
g_{R,t}^{\rm lat,att}
\le
\min\left\{U_{R,t}^{\rm lat},U_{R,t}^{\rm att}\right\}.
\end{equation}
Thus optimal data acquisition selects the observation channel, the compact
latent record controls retained information and sufficiency, attention realizes
the represented posterior or selector flow, and the Bellman recursion selects
their optimal affordable continuation.

\end{theorem}

\begin{proof}
The first inequality is class inclusion.  For the converse, tokenize the
finitely many represented legal actions of each compact latent state, expose
the selector's already charged auxiliary seed, and use its represented
inverse-transform comparisons to assign the selected action the unique
maximal score.  Mask illegal actions and apply
\cref{cor-attention-active-sampling-compilation}.  The construction preserves
the selector kernel and code state and retains the original score ancestry and
compilation cost.  Combine this equality with
\cref{eq-compact-latent-polynomial-saturated-ceiling} to obtain the polynomial
identity.  Both entries in
\cref{eq-joint-latent-attention-gain-certificate} bound the same deployed gain
under the same conditional law by
\cref{thm-unified-compact-latent-structural-certificate,thm-attention-target-gain-envelope};
their minimum is therefore valid.
\end{proof}

\begin{figure}[tbp]
\centering
\begin{tikzpicture}[fw diagram,node distance=6mm and 6mm,scale=.92,transform shape]
\node[fw source] (pub) {public presentation};
\node[fw geom,right=of pub] (src) {qualified source ledger};
\node[fw use,right=of src] (att) {attention / active selector};
\node[fw terminal,right=of att] (hit) {pre-hit target gain};
\node[fw residual,below=of att] (audit) {retrospective terminal record};
\draw[fw arrow] (pub)--(src);
\draw[fw arrow] (src)--(att);
\draw[fw arrow] (att)--(hit);
\draw[fw dashed arrow] (hit)--(audit);
\draw[fw dashed arrow] (audit) to[bend left=24] node[below,font=\scriptsize]
  {no reverse source edge} (src);
\end{tikzpicture}
\caption{Prospective attention value is charged before deployment.  A later
first-hit record remains an accounting object and cannot supply an earlier
score.}
\label{fig-attention-temporal-accounting}
\end{figure}

\section{Summary of the Attention Compilation}
\label{sec-attention-summary}

Attention is a normalized, represented selection mechanism.  Its exact Gibbs
geometry supplies a quantitative KL and Fisher--Rao action, while the
structural compilation assigns every score, value, mask, residual path, and
deployed transition to the existing saturated ledger.  The resulting target
bound is the minimum of already certified information, learning, source, and
dynamical converters.  With a latent prior row and represented log-likelihood
scores, the Gibbs row is the exact Bayesian posterior update.  Joint
optimization with the compact latent record and the active Bellman recursion
therefore gives the optimal represented acquisition--encoding--attention
pipeline at the declared code and saturated cost ceilings.  Attention supplies
no additional source coordinate or normalization credit.
The canonical implementation and training interface are
\cref{def-canonical-certified-attention-belief-architecture,def-canonical-attention-belief-training-objective}.
Their saturated-arrival calibration and empirical bias--variance certificate
are
\cref{thm-canonical-loss-saturated-arrival-calibration,thm-complete-certified-bias-variance-decomposition}.

\clearpage
\part{Structural Meta-Complexity and Certificate Minimization}
\label{part:structural-metacomplexity}\setcounter{section}{-1}

\section{Introduction: Presentation-Relative Meta-Complexity}
\label{sec-structural-metacomplexity-introduction}

\paragraph{Part contract.}
\emph{Inputs:} the budgeted derivability closure and saturated cost of
\cref{def-budgeted-derivability-closure,def-saturated-representation-degree},
the prospective source and selector profiles of
\cref{def-source-accounting-profile-separation,def-represented-pre-hit-selector},
the exact source-cell decomposition of
\cref{thm-universal-state-coefficient-source-decomposition}, and the
length-bounded universal record of
\cref{def-canonical-complete-active-record-code,thm-length-bounded-levin-realization}.
\emph{Output:} an exact Pareto meta-frontier for the cost of crossing a
target-gain threshold, its generalized-inverse relation with saturated
profiles, typed decision/search/certification/deployment interfaces,
representation-minimization specializations, and a represented
hardness-magnification rule.
\emph{First pass:} read
\cref{def-strict-structural-meta-frontier,thm-strict-meta-profile-generalized-inverse,thm-structural-meta-accountability,thm-source-cell-meta-complexity-exhaustion,thm-represented-profile-magnification}.

Meta-complexity asks for the resources required to recognize, construct, or
certify a low-complexity representation.  Circuit minimization and
resource-bounded Kolmogorov complexity are classical instances of this
question \citep{AllenderGrochowVanMelkebeekMooreMorgan2017,LiuPass2020}.
The presentation-relative version developed here uses the complete structural
records already defined in \PartIII and \PartV.  Its threshold object is the least
saturated resource at which one prospective record attains a prescribed
target gain.  Description length, construction, data access, retained state,
evaluation, qualification, and deployment remain separate coordinates of the
same cost vector.

The resulting meta-problem is an inverse description of the saturated
profile.  The profile asks for the largest certified gain below a budget; the
meta-frontier asks for the budgets at which a certified gain threshold is
crossed.  Strict thresholds give an exact inverse without a compactness or
attainment convention.  Scalar threshold costs are obtained only after the
same admissible scalarization used by
\cref{def-degree-invariant,def-alignment-complexity} has been fixed.

\section{Structural Certificate Meta-Problems}
\label{sec-structural-certificate-meta-problems}

\subsection{Complete profile data and Pareto feasibility}
\label{subsec-complete-meta-profile-data}

\begin{definition}[Complete saturated meta-profile datum]
\label{def-structural-meta-profile-datum}

Fix a size slice, a presented family, and the resource order of
\cref{def-budgeted-derivability-closure}.  A \emph{complete saturated
meta-profile datum} is a pair
\(\mathfrak D_n=(\mathscr R_n,g_n)\), where:

\begin{enumerate}[label=\textup{(\roman*)}]
\item \(\mathscr R_n\) is one already-defined class of complete
      family-uniform represented records, each retaining its upward-closed
      admissible budget set \(\mathfrak B_{\mathrm{sat},n}(R)\);
\item \(g_n:\mathscr R_n\to[0,1]\) is the already-defined numerical
      coordinate evaluated on that complete record.
\end{enumerate}

Explicitly, for the size-\(n\) presentation slice \(\mathfrak I_n\),
\begin{equation}
\label{eq-complete-record-admissible-budget-set}
\mathfrak B_{\mathrm{sat},n}(R)
=
\left\{
B:\exists\rho\in\mathsf{Der}_{\mathfrak I}
\text{ such that }
\sem{\rho}_I=R_I
\text{ and }
\operatorname{Cost}_I(\rho)\preceq B
\text{ for every }I\in\mathfrak I_n
\right\}.
\end{equation}
It is upward closed by construction.  Its Pareto-minimal elements are the
complete cost frontier of \(R\) in the sense of
\cref{lem-derivability-well-posed,rem-cost-frontier-convention}.

Its saturated profile is
\begin{equation}
\label{eq-abstract-complete-saturated-meta-profile}
P_{\mathfrak D,n}^{\mathrm{sat}}(B)
=
\sup\left\{
g_n(R):R\in\mathscr R_n,\quad
B\in\mathfrak B_{\mathrm{sat},n}(R)
\right\},
\end{equation}
with empty supremum zero.  This definition is a wrapper for an existing
profile and adds no structural channel.  The principal instances are:

\begin{enumerate}[label=\textup{(\alph*)}]
\item prospective master certificates with
      \(g_n(R)=V_R\), giving
      \(\mathfrak U_n^{\mathrm{src,sat}}\);
\item represented pre-hit selector occurrences \((R,t)\) with
      \(g_n(R,t)=(\operatorname{Adv}_t(K_{R,t};\nu_{R,t}))_+\), giving
      \(\operatorname{Sel}_n^{\mathrm{sat}}\);
\item complete active records with \(g_n(R)=A_R^H\), giving
      \(\operatorname{ActArr}_{I,H}^{\mathrm{sat}}\);
\item untyped master accounting records with their accounting visibility,
      giving \(\mathfrak U_n^{\mathrm{sat}}\);
\item completed exact-quotient records with their accounting visibility,
      giving the exact-quotient retrospective accounting datum.
\end{enumerate}

The temporal type is part of \(\mathscr R_n\).  In particular, prospective
source, prospective selector, active-arrival, and retrospective accounting
data are distinct meta-profile data even when two records have the same
denotation.

\end{definition}

\begin{definition}[Strict structural meta-frontier and scalar threshold cost]
\label{def-strict-structural-meta-frontier}

For a complete datum \(\mathfrak D_n=(\mathscr R_n,g_n)\) and
\(0\le\theta<1\), define its strict Pareto feasibility region by
\begin{equation}
\label{eq-strict-structural-meta-frontier}
\mathfrak P_{\mathfrak D,n}^{>} (\theta)
=
\bigcup_{\substack{R\in\mathscr R_n\\g_n(R)>\theta}}
\mathfrak B_{\mathrm{sat},n}(R).
\end{equation}
It is upward closed in the resource order and decreasing in \(\theta\).

Fix an admissible monotone scalarization
\(\kappa:C\to[0,\infty]\).  The associated strict structural
meta-complexity is
\begin{equation}
\label{eq-strict-scalar-structural-meta-complexity}
\operatorname{SMC}_{\mathfrak D,n,\kappa}^{>}(\theta)
=
\inf\left\{
\operatorname{Cost}_{\mathrm{sat},n}^{\kappa}(R):
R\in\mathscr R_n,\quad g_n(R)>\theta
\right\},
\end{equation}
with value \(\infty\) when the class is empty.  The corresponding scalar
sublevel profile is
\begin{equation}
\label{eq-scalar-sublevel-meta-profile}
P_{\mathfrak D,n,\kappa}^{\mathrm{sat}}(b)
=
\sup_{\substack{R\in\mathscr R_n\\
\operatorname{Cost}_{\mathrm{sat},n}^{\kappa}(R)\le b}}
g_n(R).
\end{equation}

The non-strict feasible region \(\mathfrak P^{\ge}_{\mathfrak D,n}(\theta)\)
and its scalar threshold cost
\begin{equation}
\label{eq-nonstrict-scalar-structural-meta-complexity}
\operatorname{SMC}_{\mathfrak D,n,\kappa}^{\ge}(\theta)
=
\inf\left\{
\operatorname{Cost}_{\mathrm{sat},n}^{\kappa}(R):
R\in\mathscr R_n,\quad g_n(R)\ge\theta
\right\}
\end{equation}
are defined by replacing \(g_n(R)>\theta\) with \(g_n(R)\ge\theta\) in
\cref{eq-strict-structural-meta-frontier,eq-strict-scalar-structural-meta-complexity}
respectively.  Strict regions are used for exact profile inversion;
non-strict regions record attained certificate thresholds.

\end{definition}

\subsection{Exact generalized inverse}
\label{subsec-exact-meta-profile-inverse}

\begin{theorem}[Strict meta-frontier as the exact generalized inverse]
\label{thm-strict-meta-profile-generalized-inverse}

For every complete saturated meta-profile datum, budget \(B\), and
\(0\le\theta<1\),
\begin{equation}
\label{eq-exact-strict-profile-frontier-inverse}
P_{\mathfrak D,n}^{\mathrm{sat}}(B)>\theta
\quad\Longleftrightarrow\quad
B\in\mathfrak P_{\mathfrak D,n}^{>}(\theta).
\end{equation}
Consequently,
\begin{equation}
\label{eq-profile-recovered-from-meta-frontier}
P_{\mathfrak D,n}^{\mathrm{sat}}(B)
=
\sup\left(
\{\theta\in[0,1):B\in\mathfrak P_{\mathfrak D,n}^{>}(\theta)\}
\cup\{0\}
\right).
\end{equation}

For the scalarized quantities in
\cref{eq-strict-scalar-structural-meta-complexity,eq-scalar-sublevel-meta-profile},
\begin{align}
\label{eq-scalar-meta-inverse-implications}
\operatorname{SMC}_{\mathfrak D,n,\kappa}^{>}(\theta)>b
&\Longrightarrow
P_{\mathfrak D,n,\kappa}^{\mathrm{sat}}(b)\le\theta,\\
P_{\mathfrak D,n,\kappa}^{\mathrm{sat}}(b)>\theta
&\Longrightarrow
\operatorname{SMC}_{\mathfrak D,n,\kappa}^{>}(\theta)\le b,\\
\operatorname{SMC}_{\mathfrak D,n,\kappa}^{>}(\theta)<b
&\Longrightarrow
P_{\mathfrak D,n,\kappa}^{\mathrm{sat}}(b)>\theta.
\end{align}
Under the locally finite discrete scalar-cost convention of
\cref{def-budgeted-derivability-closure,def-saturated-representation-degree},
the relevant minimum is attained whenever it is finite, and the final two
lines combine to
\begin{equation}
\label{eq-attained-scalar-meta-profile-inverse}
P_{\mathfrak D,n,\kappa}^{\mathrm{sat}}(b)>\theta
\quad\Longleftrightarrow\quad
\operatorname{SMC}_{\mathfrak D,n,\kappa}^{>}(\theta)\le b.
\end{equation}

\end{theorem}

\begin{proof}
If the supremum in
\cref{eq-abstract-complete-saturated-meta-profile} exceeds \(\theta\), the
definition of a supremum supplies a feasible record \(R\) with
\(g_n(R)>\theta\).  Its admissible budget set contains \(B\), proving the
forward implication in
\cref{eq-exact-strict-profile-frontier-inverse}.  The reverse implication
follows directly from
\cref{eq-strict-structural-meta-frontier}.  Taking the supremum over strict
sublevels proves \cref{eq-profile-recovered-from-meta-frontier}.

If the scalar meta-complexity exceeds \(b\), no record of scalar cost at most
\(b\) has gain greater than \(\theta\), proving the first line of
\cref{eq-scalar-meta-inverse-implications}.  A record witnessing a scalar
profile value greater than \(\theta\) has cost at most \(b\), proving the
second line.  If the infimum in
\cref{eq-strict-scalar-structural-meta-complexity} is strictly below \(b\),
choose a feasible record with cost below \(b\); this proves the third line.
Local finiteness gives only finitely many derivation descriptions in every
finite scalar sublevel.  The well-posedness result
\cref{lem-derivability-well-posed} and the attainment convention in
\cref{def-saturated-representation-degree} therefore turn a finite infimum
into a minimum, proving
\cref{eq-attained-scalar-meta-profile-inverse}.
\end{proof}

\begin{proposition}[Family-uniform provenance with slice-dependent gates]
\label{prop-family-uniform-slice-dependent-meta-profile}

Let
\(\mathfrak D_n=(\mathscr R_n,g_n)\) be size slices of one complete
saturated meta-profile datum generated by the family-uniform grammar of
\cref{def-budgeted-derivability-closure,def:uniform-derivation-budget-filtration}.
Each \(R\in\mathscr R_n\) is the restriction of a record produced by one
fixed derivation program.  Its admissible budget set
\(\mathfrak B_{\mathrm{sat},n}(R)\), qualification gates, and value \(g_n(R)\)
are nevertheless evaluated on the size-\(n\) public presentation.  They may
therefore depend on \(n\), on represented rank or spectral data, and on the
declared size-dependent ceiling.

For any sequences of budgets \(B_n\) and strict thresholds
\(0\le\theta_n<1\), the exact crossing set is
\begin{equation}
\label{eq-family-uniform-exact-size-crossing-set}
\left\{n:P_{\mathfrak D,n}^{\mathrm{sat}}(B_n)>\theta_n\right\}
=
\left\{n:B_n\in
\mathfrak P_{\mathfrak D,n}^{>}(\theta_n)\right\}.
\end{equation}
At every member of this set there is a complete family-uniform record whose
size-\(n\) restriction fits \(B_n\) and has value greater than \(\theta_n\).
The witnessing record may vary with \(n\), because the saturated profile is a
slice-wise supremum over the fixed family-uniform record class.  Thus
family uniformity fixes the derivation rule and excludes instance-specific
advice; numerical exclusion of a size crossing is supplied by a uniform
upper bound on the corresponding saturated profile.

\end{proposition}

\begin{proof}
The fixed-program statement is the uniform-provenance clause of
\cref{def:uniform-derivation-budget-filtration}.  The definition
\cref{eq-complete-record-admissible-budget-set} evaluates the cost of that
program on every instance in the current size slice, and
\cref{def-structural-meta-profile-datum} evaluates the existing qualification
and gain coordinate on the same slice.  Hence these quantities remain
size-dependent even though the generating program is fixed.

Apply \cref{eq-exact-strict-profile-frontier-inverse} separately at each
\(n\) to obtain \cref{eq-family-uniform-exact-size-crossing-set}.  The
witness assertion is the strict-supremum argument in the proof of
\cref{thm-strict-meta-profile-generalized-inverse}.  Finally, the definition
of the profile takes a new supremum on every slice; it contains no
cross-size choice of one maximizing record.  The description/advice
distinction follows from
\cref{lem-description-length-charged-resource}.
\end{proof}

\begin{corollary}[Existing saturated threshold and alignment complexity]
\label{cor-existing-saturated-threshold-as-metacomplexity}

Apply \cref{thm-strict-meta-profile-generalized-inverse} to the prospective
master datum
\(\mathfrak U_n^{\mathrm{src,sat}}\).  Under the attained non-strict
convention, write
\(\operatorname{SMC}_{\mathrm{src},n,\kappa}^{\ge}(\theta)\) for its
least scalar cost at visibility at least \(\theta\), and define the typed
prospective threshold
\begin{equation}
\label{eq-prospective-source-threshold-cost}
B_{\eta}^{*,\mathrm{src}}(n)
=
\inf\left\{b:
I_{n,\kappa}^{\mathrm{src,sat}}(b)\ge\eta\right\}.
\end{equation}
Since
\begin{equation*}
\mathfrak U_n^{\mathrm{src,sat}}(B)
\le I_n^{\mathrm{src,sat}}(B)
\le2\mathfrak U_n^{\mathrm{src,sat}}(B),
\end{equation*}
the prospective threshold satisfies
\begin{equation}
\label{eq-existing-threshold-meta-sandwich}
\operatorname{SMC}_{\mathrm{src},n,\kappa}^{\ge}(\eta/2)
\le B_{\eta}^{*,\mathrm{src}}(n)
\le
\operatorname{SMC}_{\mathrm{src},n,\kappa}^{\ge}(\eta),
\end{equation}
with the same strict-threshold or
\(B_+>B_{\eta}^{*,\mathrm{src}}\) relaxation used in
\cref{thm-alignment-sandwich} when attainment is unavailable.  The
untyped invariant \(B^*_{\eta}\) of \cref{def-degree-invariant} is the
threshold of the separately defined combined accounting profile
\(I_n^{\mathrm{sat}}\), whose two summands need not be attained by one
record.  If
\(\operatorname{SMC}_{\mathrm{acct},n,\kappa}^{\ge}(\theta)\)
denotes the non-strict threshold cost for the single-record accounting
master profile \(\mathfrak U_n^{\mathrm{sat}}\), then
\cref{eq-master-and-combined-profile-equivalence} gives
\begin{equation}
\label{eq-accounting-threshold-meta-sandwich}
\operatorname{SMC}_{\mathrm{acct},n,\kappa}^{\ge}(\eta/2)
\le B^*_{\eta}(n)
\le
\operatorname{SMC}_{\mathrm{acct},n,\kappa}^{\ge}(\eta).
\end{equation}
Thus neither temporal use of a single-record meta-datum identifies the
two-coordinate sum with one certificate.  The
alignment complexity \(\kappa_{\mathrm{align},\eta}\) is the program-cost
scalarization of the complete-certificate frontier selected in
\cref{def-alignment-complexity}, up to the compiler overheads
\(\Pi_1,\Pi_2\) of \cref{thm-alignment-sandwich}.  Temporal restriction to
prospective records replaces that frontier by
\cref{eq-prospective-source-threshold-cost}.

\end{corollary}

\begin{proof}
The first inequality follows because
\(I_n^{\mathrm{src,sat}}(B)\ge\eta\) forces
\(\mathfrak U_n^{\mathrm{src,sat}}(B)\ge\eta/2\).  A master certificate of
visibility at least \(\eta\) makes
\(I_n^{\mathrm{src,sat}}(B)\ge\eta\), proving the second.  The same
argument with
\cref{eq-master-and-combined-profile-equivalence} proves
\cref{eq-accounting-threshold-meta-sandwich}.  The alignment-complexity
identification is precisely the two compiler directions in
\cref{thm-alignment-sandwich}.
\end{proof}

\begin{corollary}[Presentation-equivalence transfer]
\label{cor-structural-metacomplexity-presentation-transfer}

Suppose two presentations are admissibly equivalent with resource
distortions \(\Psi_{12}\) and \(\Psi_{21}\), and the equivalence transports
the complete record class and preserves its gain.  Then
\begin{equation}
\label{eq-meta-frontier-presentation-transfer}
B\in\mathfrak P_{\mathfrak D_1,n}^{>}(\theta)
\Longrightarrow
\Psi_{12}(B,n)\in\mathfrak P_{\mathfrak D_2,n}^{>}(\theta),
\end{equation}
and conversely with \(1,2\) interchanged.  Scalar meta-complexities are
class-equivalent under the induced scalar distortions.

\end{corollary}

\begin{proof}
Transport a witnessing complete record and its derivation through the
admissible presentation equivalence.  The transported cost is at most
\(\Psi_{12}(B,n)\), while denotation and gain are preserved.  Apply the
reverse compiler for the converse.  This is the threshold-resolved form of
\cref{thm-dstar-presentation-invariance}.
\end{proof}

\section{Temporal Meta-Complexity and Accountability}
\label{sec-temporal-metacomplexity-accountability}

\subsection{Source, selector, active, and accounting frontiers}
\label{subsec-temporally-typed-meta-frontiers}

\begin{definition}[Temporally typed structural meta-frontiers]
\label{def-temporally-typed-structural-meta-frontiers}

Use \cref{def-structural-meta-profile-datum} with the following existing
record classes and write:
\begin{align}
\label{eq-source-structural-meta-frontier}
\mathfrak P_{n}^{\mathrm{src},>}(\theta)
&:=\mathfrak P_{\mathfrak U^{\mathrm{src}},n}^{>}(\theta),\\
\label{eq-selector-structural-meta-frontier}
\mathfrak P_{I,H,n}^{\mathrm{sel},>}(\theta)
&:=\mathfrak P_{\operatorname{Sel},I,H,n}^{>}(\theta),\\
\label{eq-active-arrival-meta-frontier}
\mathfrak P_{I,H,n}^{\mathrm{act},>}(\theta)
&:=\mathfrak P_{\operatorname{ActArr},I,H,n}^{>}(\theta),\\
\label{eq-accounting-meta-frontier}
\mathfrak P_{I,H,n}^{\mathrm{acct},>}(\theta)
&:=\mathfrak P_{\mathrm{Acct},I,H,n}^{>}(\theta).
\end{align}
Here the source datum uses complete prospective master certificates, the
selector datum uses represented pre-hit selector occurrences, the active
datum uses complete finite-horizon control records, and the accounting datum
uses completed terminal-compatible records.  Every frontier inherits the
temporal type and complete cost of its underlying record class.

\end{definition}

\begin{theorem}[Structural meta-accountability]
\label{thm-structural-meta-accountability}

Let \(\mathcal A\) be a budget-\(B\), horizon-\(H\) represented computation
with neutral baseline \(\nu\), and let \(\Psi(B,n)\) be the charged
one-test-clock compilation ceiling of
\cref{thm-prospective-selector-accountability}.  For every \(\varepsilon>0\),
\begin{equation}
\label{eq-structural-meta-accountability-absence}
\Psi(B,n)\notin
\mathfrak P_{I,H,n}^{\mathrm{sel},>}(\varepsilon/H)
\quad\Longrightarrow\quad
\Pr\{\tau\le H\}
\le\operatorname{Enum}_{\nu}(\mathcal A,H)+\varepsilon.
\end{equation}
Conversely,
\begin{equation}
\label{eq-structural-meta-accountability-extraction}
\Pr\{\tau\le H\}
>\operatorname{Enum}_{\nu}(\mathcal A,H)+\varepsilon
\quad\Longrightarrow\quad
\Psi(B,n)\in
\mathfrak P_{I,H,n}^{\mathrm{sel},>}(\varepsilon/H).
\end{equation}

After scalarization, the sufficient lower-bound condition is
\begin{equation}
\label{eq-scalar-structural-meta-accountability}
\operatorname{SMC}_{I,H,n,\kappa}^{\mathrm{sel},>}(\varepsilon/H)
>
\kappa\!\left(\Psi(B,n)\right),
\end{equation}
with the boundary case governed by the exact Pareto statement
\cref{eq-structural-meta-accountability-absence}.

\end{theorem}

\begin{proof}
Absence from the strict selector frontier means, by
\cref{thm-strict-meta-profile-generalized-inverse}, that
\(\operatorname{Sel}_{I,n}^{\mathrm{sat}}(\Psi(B,n))\le\varepsilon/H\).
Substitute this inequality into
\cref{eq-selector-accountability-bound}.  This proves
\cref{eq-structural-meta-accountability-absence}.  The extraction statement
is \cref{eq-selector-extraction} together with the complete-cost conclusion
in the proof of \cref{thm-prospective-selector-accountability}.  The scalar
condition implies Pareto-frontier absence by
\cref{eq-scalar-meta-inverse-implications}.
\end{proof}

\begin{corollary}[Transfer-class meta-hardness criterion]
\label{cor-transfer-class-meta-hardness-criterion}

Let \(B(n)\) and \(H(n)\ge1\) belong to the declared transfer class.  Suppose
\(e_n\ge0\) and \(\varepsilon_n>0\) are negligible in that class and every
budget-\(B(n)\) execution satisfies
\begin{equation}
\label{eq-transfer-class-neutral-baseline-bound}
\operatorname{Enum}_{\nu}(\mathcal A,H(n))\le e_n,
\end{equation}
while
\begin{equation}
\label{eq-transfer-class-selector-meta-absence}
\Psi(B(n),n)\notin
\mathfrak P_{I,H(n),n}^{\mathrm{sel},>}
  (\varepsilon_n/H(n)).
\end{equation}
Then the complete success profile at budget \(B(n)\) is at most
\(e_n+\varepsilon_n\).

\end{corollary}

\begin{proof}
Apply \cref{thm-structural-meta-accountability} to every represented
execution in the saturated budget slice and take the defining supremum of
the success profile.
\end{proof}

\begin{remark}[Relation to prospective source and retrospective accounting]
\label{rem-meta-source-accounting-separation}

The selector frontier in \cref{thm-structural-meta-accountability} is
prospective and yields the unconditional success implication.  The source
frontier records complete \(\Geom/\Abs\) certificates.  A selector occurrence
first produces the contextual accounting charge of
\cref{thm-prospective-selector-structural-charge}; the same-record comparison
of \cref{cor-contextual-structural-charge-transfer} then compiles each
nonzero charge into a prospective source record.  This is a
selector-to-source map, so exclusion from the source frontier excludes the
corresponding selectors by contrapositive whenever that comparison record is
present.  The accounting frontier
records completed first-hit objects and yields the retrospective bounds of
\cref{thm-universal-saturated-profile-accountability}; temporal source status
is governed by
\cref{def-pre-hit-temporally-admissible-source,cor-no-retrospective-first-hit-source-loophole}.
Thus the three meta-frontiers preserve the source/selector/accounting typing
already present in the framework.

\end{remark}

\section{Source-Resolved Structural Meta-Complexity}
\label{sec-source-resolved-structural-metacomplexity}

\subsection{Exact cell frontiers}
\label{subsec-exact-source-cell-meta-frontiers}

\begin{definition}[Exact source-cell meta-complexity]
\label{def-exact-source-cell-meta-complexity}

Let
\begin{equation}
\label{eq-structural-meta-source-cell-index}
\mathfrak A_{\mathrm{src}}
=
\left\{
(\chi,S_{\mathrm{st}},S_{\mathrm{cf}}):
\begin{array}{l}
\chi\in\mathfrak S_{\mathrm{src}}^5,\\
\varnothing\ne S_{\mathrm{st}}\subseteq\mathfrak J_{\Abs}^5,\\
S_{\mathrm{cf}}\subseteq\mathfrak J_{\Abs}^5
\end{array}
\right\}.
\end{equation}
For \(\alpha\in\mathfrak A_{\mathrm{src}}\), let
\(\mathscr R_{\alpha,n}^{\mathrm{src,nf}}\) be the canonical normalized
complete prospective source records in that exact cell, with visibility
\(V_R\) and their complete normalized saturated costs.  Define
\begin{equation}
\label{eq-exact-source-cell-meta-frontier}
\mathfrak P_{\alpha,n}^{\mathrm{src,nf},>}(\theta)
=
\bigcup_{\substack{R\in\mathscr R_{\alpha,n}^{\mathrm{src,nf}}\\
V_R>\theta}}
\mathfrak B_{\mathrm{sat},n}(R),
\end{equation}
and define
\(\operatorname{SMC}_{\alpha,n,\kappa}^{\mathrm{src,nf},>}(\theta)\)
by \cref{eq-strict-scalar-structural-meta-complexity}.  Empty cells have
empty frontier and scalar cost \(\infty\).

\end{definition}

\begin{theorem}[Exact source-cell exhaustion of structural meta-complexity]
\label{thm-source-cell-meta-complexity-exhaustion}

The canonical normalized prospective source class is the disjoint union of
the \(4960\) cells in \(\mathfrak A_{\mathrm{src}}\).  Consequently,
\begin{equation}
\label{eq-exact-source-meta-frontier-union}
\mathfrak P_{n}^{\mathrm{src,nf},>}(\theta)
=
\bigcup_{\alpha\in\mathfrak A_{\mathrm{src}}}
\mathfrak P_{\alpha,n}^{\mathrm{src,nf},>}(\theta),
\end{equation}
and, for every scalarization,
\begin{equation}
\label{eq-exact-source-meta-cost-minimum}
\operatorname{SMC}_{n,\kappa}^{\mathrm{src,nf},>}(\theta)
=
\min_{\alpha\in\mathfrak A_{\mathrm{src}}}
\operatorname{SMC}_{\alpha,n,\kappa}^{\mathrm{src,nf},>}(\theta).
\end{equation}

If an unnormalized source record has cost at most \(B\) and visibility
greater than \(\theta\), its coefficient-complete normalization has the same
visibility, belongs to one exact cell, and has cost at most
\(\Psi_{\mathrm{src},5}(B)\).  Hence
\begin{equation}
\label{eq-source-meta-normalization-transfer}
B\in\mathfrak P_n^{\mathrm{src},>}(\theta)
\Longrightarrow
\Psi_{\mathrm{src},5}(B)
\in
\bigcup_{\alpha\in\mathfrak A_{\mathrm{src}}}
\mathfrak P_{\alpha,n}^{\mathrm{src,nf},>}(\theta).
\end{equation}
Every normalized cell record is itself a complete prospective source record
at its displayed cost, giving the reverse inclusion at that normalized
ceiling.

For the combined two-coordinate profile,
\cref{eq-source-master-combined-equivalence} gives the threshold comparison
\begin{equation}
\label{eq-combined-profile-source-meta-threshold-comparison}
I_n^{\mathrm{src,sat}}(B)>2\theta
\Longrightarrow
B\in\mathfrak P_n^{\mathrm{src},>}(\theta),
\qquad
B\in\mathfrak P_n^{\mathrm{src},>}(\theta)
\Longrightarrow
I_n^{\mathrm{src,sat}}(B)>\theta.
\end{equation}

\end{theorem}

\begin{proof}
The unique source-mode--state-provenance--coefficient-provenance index and
the cell count are
\cref{thm-universal-state-coefficient-source-decomposition}.  A strict
feasible record therefore belongs to exactly one cell, proving
\cref{eq-exact-source-meta-frontier-union}.  The infimum over a finite union
is the minimum of the cell infima, with empty cells assigned \(\infty\),
which proves \cref{eq-exact-source-meta-cost-minimum}.  The normalization
statement is the denotation-preserving compiler and cost majorant in that
same theorem.  Finally,
\(\mathfrak U^{\mathrm{src,sat}}\le I^{\mathrm{src,sat}}\le
2\mathfrak U^{\mathrm{src,sat}}\); apply
\cref{thm-strict-meta-profile-generalized-inverse} to obtain
\cref{eq-combined-profile-source-meta-threshold-comparison}.
\end{proof}

\begin{corollary}[Cellwise lower bounds close the complete source frontier]
\label{cor-cellwise-meta-lower-bounds-close-source-frontier}

Let \(B\) be a declared budget.  If
\begin{equation}
\label{eq-every-source-cell-meta-excludes-budget}
\Psi_{\mathrm{src},5}(B)
\notin
\mathfrak P_{\alpha,n}^{\mathrm{src,nf},>}(\theta)
\qquad(\alpha\in\mathfrak A_{\mathrm{src}}),
\end{equation}
then
\(\mathfrak U_n^{\mathrm{src,sat}}(B)\le\theta\) and
\(I_n^{\mathrm{src,sat}}(B)\le2\theta\).

\end{corollary}

\begin{proof}
The union in \cref{eq-exact-source-meta-frontier-union} is empty at the
normalized ceiling.  Apply
\cref{eq-source-meta-normalization-transfer,eq-exact-strict-profile-frontier-inverse}
and then \cref{eq-source-master-combined-equivalence}.
\end{proof}

\section{Representation, Extraction, and Operational Interfaces}
\label{sec-meta-representation-extraction-interfaces}

\subsection{Exact denotation and extraction}
\label{subsec-exact-denotation-extraction-meta}

\begin{definition}[Exact-denotation representation and extraction frontiers]
\label{def-exact-denotation-meta-frontiers}

Fix a represented evaluator grammar \(\mathscr E\), an output carrier
\(\mathcal Y\), and \(Y\in\mathcal Y\).  The exact-denotation
representation frontier is
\begin{equation}
\label{eq-exact-denotation-representation-frontier}
\mathfrak P_{\mathscr E,n}^{=}(Y)
=
\bigcup_{\substack{R\in\mathscr E\\\sem{R}=Y}}
\mathfrak B_{\mathrm{sat},n}(R).
\end{equation}
After scalarization, define
\begin{equation}
\label{eq-exact-denotation-representation-cost}
\operatorname{RepMC}_{\mathscr E,n,\kappa}(Y)
=
\inf_{\substack{R\in\mathscr E\\\sem{R}=Y}}
\operatorname{Cost}_{\mathrm{sat},n}^{\kappa}(R).
\end{equation}

This representation frontier is evaluated on the singleton semantic slice
\(\mathfrak I_Y=\{Y\}\).  A candidate evaluator may therefore hard-code a
description specialized to \(Y\), but its full description, evaluator, and
retained-state costs remain charged.  Family-uniformity in
\cref{eq-complete-record-admissible-budget-set} is uniformity over this
singleton; it does not require a uniform procedure that constructs the
evaluator from a varying public input.

For a public presentation \(I\), the extraction frontier
\(\mathfrak P_{\mathscr E,n}^{\mathrm{ext}}(Y\mid I)\) ranges over complete
represented computations that read only the declared interface of \(I\),
halt with an evaluator \(R\in\mathscr E\) satisfying \(\sem{R}=Y\), and
retain the complete construction and output costs.  Its scalar infimum is
\(\operatorname{ExtMC}_{\mathscr E,n,\kappa}(Y\mid I)\).
For a family of public presentations, the extraction derivation is uniform
over the declared family slice.  Thus extraction, unlike singleton
representation, charges the rule that finds the evaluator from the input.

The equality verifier may read \(Y\) when checking a completed output.  The
candidate evaluator receives only the inputs declared by \(\mathscr E\).
This separation prevents the verification interface from becoming a lookup
primitive of the candidate representation.

\end{definition}

\begin{proposition}[Output projection from extraction to representation]
\label{prop-meta-extraction-output-projection}

Suppose the charged output projection of a complete extraction record of
budget \(B\) produces its evaluator record at cost at most
\(\Psi_{\mathrm{out}}(B,n)\).  Then
\begin{equation}
\label{eq-extraction-implies-representation-frontier}
B\in\mathfrak P_{\mathscr E,n}^{\mathrm{ext}}(Y\mid I)
\Longrightarrow
\Psi_{\mathrm{out}}(B,n)\in
\mathfrak P_{\mathscr E,n}^{=}(Y).
\end{equation}
Thus extraction charges the search for a representation in addition to the
representation retained in its output.

\end{proposition}

\begin{proof}
Project the completed computation record to its represented output evaluator.
Computation closure in \cref{def-budgeted-derivability-closure} preserves the
denotation and charges the output through \(\Psi_{\mathrm{out}}\).
\end{proof}

\subsection{Decision, search, certification, and deployment}
\label{subsec-meta-operational-interfaces}

\begin{definition}[Four operational forms of a structural meta-problem]
\label{def-four-meta-operational-forms}

Fix an encoded size-\(n\) meta-input slice \(\Xi_n\).  Each
\(\xi\in\Xi_n\) names a complete datum \(\mathfrak D_\xi\), a budget
\(B_\xi\), and a threshold \(\theta_\xi\).  The four operational forms are:
\begin{enumerate}[label=\textup{(\roman*)}]
\item \emph{decision}: output the truth value of budget membership in the
      declared meta-frontier;
\item \emph{search}: output a witnessing complete record on a yes instance;
\item \emph{certification}: output the record together with its charged
      derivation, qualification data, and represented gain verifier;
\item \emph{deployment}: evaluate the record prospectively as an ancestor of
      the transition or selector whose gain it certifies.
\end{enumerate}

Each form is equipped with a family
\(\mathfrak D^{\mathsf A}=(\mathfrak D_n^{\mathsf A})_{n\ge1}\) of complete
meta-profile data,
\begin{equation}
\label{eq-operational-meta-profile-data}
\mathfrak D_n^{\mathsf A}
=
(\mathscr R_n^{\mathsf A},g_n^{\mathsf A}),
\qquad
\mathsf A\in
\{\mathsf{dec},\mathsf{search},\mathsf{cert},\mathsf{deploy}\}.
\end{equation}
For search, \(g_n^{\mathsf{search}}\) is the underlying gain of the valid
witness, and is zero for an invalid output.  For certification it is that
gain when the represented derivation, qualification, and gain verifier all
accept, and zero otherwise.  For deployment it is the prospective gain
actually realized by the represented deployed record.  A total decision
record \(D\) has the worst-slice correctness coordinate
\begin{equation}
\label{eq-total-meta-decision-gain}
g_n^{\mathsf{dec}}(D)
=
\inf_{\xi\in\Xi_n}
\mathbf1_{\left\{
D(\xi)
=
\mathbf1_{\{B_\xi\in
\mathfrak P_{\mathfrak D_\xi,n}^{>}(\theta_\xi)\}}
\right\}}.
\end{equation}
The decision record includes its terminating evaluation on the complete
declared input interface.  Each form therefore has its own represented
record class, numerical coordinate, and saturated cost.  A
compiler between two forms is admissible when it is represented, preserves
the required soundness/completeness relation, and carries an explicit cost
majorant.

\end{definition}

\begin{theorem}[Typed compiler calculus for meta-problems]
\label{thm-typed-meta-compiler-calculus}

Let \(\mathfrak D_n^{\mathsf A}\) and
\(\mathfrak D_n^{\mathsf B}\) be the typed data of two operational forms
from \cref{eq-operational-meta-profile-data}.  Suppose a represented compiler maps
every valid \(\mathsf A\)-record of budget \(B\) and gain greater than
\(\theta\) to a valid \(\mathsf B\)-record of budget at most
\(\Psi_{\mathsf A\to\mathsf B}(B,n)\) and gain greater than
\(\varphi(\theta)\).  Then
\begin{equation}
\label{eq-typed-meta-compiler-frontier-transfer}
B\in\mathfrak P_{\mathfrak D^{\mathsf A},n}^{>}(\theta)
\Longrightarrow
\Psi_{\mathsf A\to\mathsf B}(B,n)
\in\mathfrak P_{\mathfrak D^{\mathsf B},n}^{>}(\varphi(\theta)).
\end{equation}

Projection supplies certification-to-search and extraction-to-representation
compilers.  A represented verifier and a terminating no-case bound supply a
search-to-total-decision compiler.  A represented self-reduction supplies a
decision-to-search compiler.  Prospective deployment additionally requires
the prefix factorization and ancestry conditions of
\cref{def-pre-hit-temporally-admissible-source}.

\end{theorem}

\begin{proof}
Compose the witnessing derivation with the declared compiler.  Constructor
closure preserves uniformity, the compiler hypothesis preserves validity and
gain, and cost composition gives the displayed ceiling.  The listed special
cases use projection, verifier execution, self-reduction, and temporal prefix
evaluation, respectively.
\end{proof}

\section{Classical Representation-Minimization Specializations}
\label{sec-classical-metacomplexity-specializations}

\subsection{Circuit minimization}
\label{subsec-circuit-minimization-specialization}

\begin{definition}[Truth-table circuit presentation]
\label{def-truth-table-circuit-presentation}

Let \(f:\{0,1\}^k\to\{0,1\}\), put \(N=2^k\), and let
\(z_f\in\{0,1\}^{N}\) be its truth table.  The verification presentation
contains \(z_f\), the canonical order of \(\{0,1\}^k\), and equality of a
completed circuit truth table with \(z_f\).  The candidate grammar
\(\mathscr E_{\mathrm{circ}}\) consists of Boolean circuits over one fixed
finite gate basis.  A candidate circuit evaluator receives \(x\in\{0,1\}^k\)
and its own gate description; it has no truth-table lookup interface.

For a circuit record \(R\), write \(s(R)\) for its gate-count coordinate and
define
\begin{equation}
\label{eq-framework-circuit-representation-complexity}
\operatorname{CC}_{\mathrm{meta}}(z_f)
=
\min_{\substack{R\in\mathscr E_{\mathrm{circ}}\\\sem{R}=f}}s(R).
\end{equation}
Gate count is retained as one Pareto coordinate rather than used as a scalar
that discards the other costs.  For constants depending only on the finite
basis, every circuit of size at most \(s\) has a canonical prefix encoding
and evaluator satisfying
\begin{equation}
\label{eq-circuit-canonical-complete-cost}
L_{\mathrm{circ}}(s,k)
\le c_0+c_1s\left\lceil\log_2(k+s+2)\right\rceil,
\quad
T_{\mathrm{eval}}(s,k)\le c_2(k+s),
\quad
S_{\mathrm{eval}}(s,k)\le c_3(k+s).
\end{equation}
Let \(B_{\mathrm{circ}}(s,k)\) be the complete Pareto ceiling containing
these three bounds and the gate-count bound \(s(R)\le s\).  The truth-table
equality verifier has separate cost
\(O(2^kT_{\mathrm{eval}}(s,k))\), charged to certification rather than to
one candidate evaluation.

\end{definition}

\begin{theorem}[Circuit minimization specialization]
\label{thm-circuit-minimization-specialization}

For the presentation of \cref{def-truth-table-circuit-presentation},
\(\operatorname{CC}_{\mathrm{meta}}(z_f)\) is exactly the minimum circuit
size of \(f\) over the declared basis.  Hence the threshold language
\begin{equation}
\label{eq-framework-mcsp-threshold-language}
\left\{(z_f,s):
\operatorname{CC}_{\mathrm{meta}}(z_f)\le s\right\}
\end{equation}
is the corresponding Minimum Circuit Size Problem.  Its input length is
\(N=2^k\), while \(k\) is the semantic input dimension of the represented
function.

The exact-denotation Pareto frontier satisfies
\begin{equation}
\label{eq-mcsp-pareto-frontier-equivalence}
\operatorname{CC}_{\mathrm{meta}}(z_f)\le s
\quad\Longleftrightarrow\quad
B_{\mathrm{circ}}(s,k)
\in\mathfrak P_{\mathscr E_{\mathrm{circ}},k}^{=}(f).
\end{equation}
Thus gate count determines the classical threshold, while description,
evaluation, state, and verification costs remain visible in the complete
framework budget.

The extraction quantity
\(\operatorname{ExtMC}_{\mathscr E_{\mathrm{circ}}}(f\mid z_f)\)
additionally charges the computation that constructs a circuit from the truth
table.  The output projection in
\cref{prop-meta-extraction-output-projection} maps this extraction problem to
the representation-minimization problem.

\end{theorem}

\begin{proof}
The candidate grammar is precisely the declared circuit basis, exact
denotation is equality with \(f\) on all \(2^k\) inputs, and
\cref{eq-framework-circuit-representation-complexity} minimizes its retained
gate-count coordinate.  This is the standard minimum circuit size by
definition.  A size-at-most-\(s\) witness has the complete bounds in
\cref{eq-circuit-canonical-complete-cost}, proving the forward implication in
\cref{eq-mcsp-pareto-frontier-equivalence}.  Conversely, membership below
\(B_{\mathrm{circ}}(s,k)\) includes the coordinate constraint
\(s(R)\le s\), proving the reverse implication.  The final assertion is
\cref{prop-meta-extraction-output-projection}.
\end{proof}

\subsection{Universal programs and time-sensitive costs}
\label{subsec-program-minimization-specialization}

\begin{definition}[Declared universal-program meta-complexity]
\label{def-declared-universal-program-metacomplexity}

Fix a prefix universal interpreter \(\mathsf U\), a finite string
\(z=z_1\cdots z_N\), and one of the following represented output
interfaces:
\begin{enumerate}[label=\textup{(\roman*)}]
\item in \emph{stream mode}, \(\mathsf U(p)\) outputs all of \(z\) within
      \(t\) steps;
\item in \emph{bit mode}, \(\mathsf U(p,i)\) outputs \(z_i\) within
      \(t\) steps for every \(1\le i\le N\).
\end{enumerate}
Let \(\mathcal W_{\mathsf U}^{\mathrm{mode}}(z)\) be the set of pairs
\((p,t)\) satisfying the selected interface.  For any declared admissible
monotone scalarization \(c(|p|,t)\) with finite sublevels and description
length charged, put
\begin{equation}
\label{eq-declared-program-meta-complexity}
K_{\mathsf U}^{c,\mathrm{mode}}(z)
=
\inf_{(p,t)\in\mathcal W_{\mathsf U}^{\mathrm{mode}}(z)}c(|p|,t).
\end{equation}
The complete resource-vector form retains \((|p|,t,\text{space},\text{output
length})\) before scalarization.

The associated minimization language is
\begin{equation}
\label{eq-declared-program-minimization-language}
\operatorname{MinK}_{\mathsf U}^{c,\mathrm{mode}}
=
\left\{(z,s):K_{\mathsf U}^{c,\mathrm{mode}}(z)\le s\right\}.
\end{equation}

\end{definition}

\begin{theorem}[Program-minimization specialization]
\label{thm-program-minimization-specialization}

The exact-denotation representation complexity of
\cref{def-declared-universal-program-metacomplexity} equals
\(K_{\mathsf U}^{c,\mathrm{mode}}(z)\).  In bit mode, choosing
\(c(|p|,t)=|p|+t\) gives \(\operatorname{KT}_{\mathsf U}(z)\) for the
declared random-access interpreter, and
\cref{eq-declared-program-minimization-language} is its Minimum
\(\operatorname{KT}\) Problem.  In stream mode, choosing
\(c(|p|,t)=|p|+\lceil\log_2(1+t)\rceil\) gives the logarithmic-time
Levin-style output complexity.  These identifications are interface typed;
neither output convention is substituted for the other.  Machine changes
connected by represented universal simulation alter the quantities by their
explicit interpreter overheads.

The scheduler of \cref{def-length-bounded-represented-levin-search} acts on
the separate description-length/time resource vector.  Its branch cost
\cref{eq-levin-scheduler-majorant} supplies the represented simulation
comparison for every chosen scalarization; it does not identify the two
scalarizations above.

\end{theorem}

\begin{proof}
The records in the exact-denotation frontier are exactly the pairs \((p,t)\)
in \(\mathcal W_{\mathsf U}^{\mathrm{mode}}(z)\); minimizing the declared cost
gives the equality.  A represented universal interpreter composes with each
program and incurs its displayed simulation cost, proving the machine-transfer
statement.  The final assertion follows from
\cref{thm-length-bounded-levin-realization}.
\end{proof}

\section{Prefix Coding and Structural Incompressibility}
\label{sec-prefix-coding-structural-incompressibility}

\subsection{Unconditional counting bounds}
\label{subsec-unconditional-meta-counting}

\begin{theorem}[Prefix-code counting and random-table incompressibility]
\label{thm-prefix-code-random-table-incompressibility}

Let \(\ell\in\mathbb N\).  Fix independently of \(Z\) an evaluator class
\(\mathscr E\) with an injective prefix code \(c:\mathscr E\to\{0,1\}^*\),
and put
\(\mathscr E_{\le\ell}=\{R\in\mathscr E:|c(R)|\le\ell\}\).  Then
\begin{equation}
\label{eq-prefix-code-cardinality-bound}
\#\mathscr E_{\le\ell}\le2^\ell.
\end{equation}
Let \(Z\) be uniform on \(\{0,1\}^{N}\).  If every evaluator denotes one
Boolean table of length \(N\), then
\begin{equation}
\label{eq-random-table-exact-incompressibility}
\Pr\{Z\text{ has an exact code of length at most }\ell\}
\le2^{\ell-N}.
\end{equation}
For \(0<\varepsilon\le1/2\),
\begin{equation}
\label{eq-random-table-approximate-incompressibility}
\Pr\left\{
\exists R\in\mathscr E_{\le\ell}:
\frac1N\sum_{i=1}^{N}
\mathbf1_{\{R(i)=Z_i\}}
\ge\frac12+\varepsilon
\right\}
\le
\min\left\{1,2^\ell e^{-2\varepsilon^2N}\right\}.
\end{equation}

\end{theorem}

\begin{proof}
For a prefix-free code, Kraft's inequality gives
\(\sum_{R\in\mathscr E_{\le\ell}}2^{-|c(R)|}\le1\).  Since
\(|c(R)|\le\ell\), every summand is at least \(2^{-\ell}\), and hence
\(\#\mathscr E_{\le\ell}\le2^\ell\).  This is the same prefix-code
argument used for complete active records in
\cref{eq-complete-active-record-kraft}.  Each evaluator denotes at most one
exact table, while \(Z\) has \(2^N\) equiprobable values; a union bound gives
\cref{eq-random-table-exact-incompressibility}.  For a fixed evaluator, its
agreement count with \(Z\) is binomial \(\operatorname{Bin}(N,1/2)\).
Hoeffding's inequality bounds its upper tail by
\(e^{-2\varepsilon^2N}\).  Union over at most \(2^\ell\) evaluators and cap
the result by one.
\end{proof}

\begin{corollary}[Random truth-table circuit lower bound]
\label{cor-random-truth-table-circuit-lower-bound}

For the circuit presentation in
\cref{def-truth-table-circuit-presentation}, let \(L(s,k)\) be an upper bound
on the canonical prefix-code length of every circuit of size at most \(s\).
A uniform
random Boolean function on \(k\) bits satisfies
\begin{equation}
\label{eq-random-circuit-size-counting-bound}
\Pr\{\operatorname{CC}_{\mathrm{meta}}(Z)\le s\}
\le2^{L(s,k)-2^k}.
\end{equation}
The same substitution in
\cref{eq-random-table-approximate-incompressibility} gives the corresponding
approximation bound.

\end{corollary}

\begin{proof}
Apply \cref{thm-prefix-code-random-table-incompressibility} with
\(N=2^k\) and \(\ell=L(s,k)\).
\end{proof}

\subsection{Represented coding transfer}
\label{subsec-represented-meta-coding-transfer}

\begin{theorem}[Charged represented coding transfer]
\label{thm-charged-represented-coding-transfer}

Let \(\mathsf P\) be a represented distribution on a finite complete-record
class, and suppose the record contains a represented prefix encoder
\(c\), decoder \(D\), and cost majorant \(T_D\) satisfying
\begin{equation}
\label{eq-represented-meta-code-length-hypothesis}
D(c(R))=R,
\qquad
|c(R)|\le
\left\lceil-\log_2\mathsf P(R)\right\rceil+a
\end{equation}
for every \(R\) in the support.  Then
\begin{equation}
\label{eq-represented-coding-meta-cost-bound}
B_R^{\mathrm{code}}
:=
\operatorname{Cost}(D,c,\mathsf P)
\oplus |c(R)|
\oplus T_D(c(R)),
\qquad
B_R^{\mathrm{code}}
\in\mathfrak B_{\mathrm{sat},n}(R).
\end{equation}
Here \(\operatorname{Cost}(D,c,\mathsf P)\) is the complete charged cost of
the joint represented decoder, encoder, and distribution record; it includes
their descriptions, retained data, and evaluation interfaces.  The term
\(|c(R)|\) occupies the description-length coordinate, and
\(T_D(c(R))\) occupies the decoder execution coordinates.
Consequently any exact-denotation or structural meta-frontier containing the
decoded record also contains \(B_R^{\mathrm{code}}\).

\end{theorem}

\begin{proof}
The decoder, codeword, represented distribution data used by the decoder,
and its execution form one derivation of \(R\).  Constructor and computation
closure add precisely the displayed cost coordinates.  Frontier membership
then follows from the definition of the admissible budget set.
\end{proof}

\begin{remark}[Coding interface]
\label{rem-resource-bounded-coding-interface}

The coding conclusion of
\cref{thm-charged-represented-coding-transfer} is invoked through its
represented encoder and decoder.  General resource-bounded coding theorems
are separate meta-complexity results; efficient coding is itself tied to
computational lower bounds and one-way functions
\citep{HiraharaLuNanashima2024}.  The saturated ledger retains this issue as
the cost of the coding interface.

\end{remark}

\section{Learning, Compression, and Natural Properties}
\label{sec-meta-learning-compression-natural-properties}

\subsection{Learned records as feasible meta-certificates}
\label{subsec-learned-meta-certificates}

\begin{theorem}[Learned selector and source records enter their typed meta-frontiers]
\label{thm-source-resolved-learning-meta-transfer}

Let a complete learned record satisfy the hypotheses of
\cref{thm-source-resolved-saturated-learning-compilation}, have total budget
at most \(B\), and have certified prospective target gain greater than
\(\theta\) after all statistical, optimization, noise, filtering, and
deployment terms are included.  Let
\(\mathfrak P_{\mathrm{learn},n}^{\mathrm{sel},>}\) be the selector
frontier obtained by restricting the complete selector-occurrence datum to
these learned records.  Then
\begin{equation}
\label{eq-learned-record-meta-frontier-membership}
D_B^{\rm learn}
\in
\mathfrak P_{\mathrm{learn},n}^{\mathrm{sel},>}(\theta),
\end{equation}
where \(D_B^{\rm learn}=\Psi_{\rm learn}
(B,H,n,m,\boldsymbol\eta,\delta)\) is the complete common ceiling in
\cref{eq-complete-source-resolved-learning-cost}.  Therefore
\begin{equation}
\label{eq-meta-lower-bound-excludes-learned-gain}
D_B^{\rm learn}
\notin\mathfrak P_{\mathrm{learn},n}^{\mathrm{sel},>}(\theta)
\Longrightarrow
\text{every such certified learned record has gain at most }\theta.
\end{equation}

Suppose additionally that the complete learned record contains the
same-record comparisons required by
\cref{cor-contextual-structural-charge-transfer}.  If the resulting transfer
selects a prospective source record \(R_{\mathrm{src}}\), of exact source
visibility \(V(R_{\mathrm{src}})>\vartheta\), normalize it by
\cref{thm-universal-state-coefficient-source-decomposition} and let
\(\alpha\) be the resulting exact source index.  At the complete ceiling
\begin{equation*}
D_B^{\mathrm{learn,src}}
=
\Psi_{\mathrm{src},5}
\!\left(\Psi_{\mathrm{tr}}(D_B^{\rm learn})\right),
\end{equation*}
one has
\begin{equation}
\label{eq-learned-selected-source-meta-membership}
D_B^{\mathrm{learn,src}}
\in
\mathfrak P_{\alpha,n}^{\mathrm{src,nf},>}(\vartheta)
\subseteq
\mathfrak P_n^{\mathrm{src,nf},>}(\vartheta).
\end{equation}
This cellwise conclusion concerns the selected source record; one complete
learner may retain several source cells.

\end{theorem}

\begin{proof}
The learning compilation retains the deployed pre-hit selector and every
source, noise, statistical, optimization, and deployment coordinate at the
enlarged cost.  Its certified target gain is the selector-datum coordinate,
so \cref{def-strict-structural-meta-frontier} proves
\cref{eq-learned-record-meta-frontier-membership}; its contrapositive is
\cref{eq-meta-lower-bound-excludes-learned-gain}.  Under the additional
comparison, the selected source is a complete prospective record at the
transfer ceiling and its coefficient-complete normalization lies at
\(D_B^{\mathrm{learn,src}}\).  Apply
\cref{def-exact-source-cell-meta-complexity,thm-source-cell-meta-complexity-exhaustion}
to obtain \cref{eq-learned-selected-source-meta-membership}.
\end{proof}

Natural lower-bound properties can yield learning and compression procedures
under explicit access models \citep{CarmosinoEtAl2016,GoldbergKabanets2023}.
The following definition places constructivity, density, and usefulness in
the presentation-relative ledger.

\begin{definition}[Presentation-relative structural natural property]
\label{def-presentation-relative-structural-natural-property}

Fix a finite object carrier \(\mathcal Y_n\), a declared reference law
\(\mu_n\), a represented certificate class \(\mathscr R_{\alpha,n}\) with
declared class index \(\alpha\),
and a budget \(B\).  A property \(\mathcal P_n\subseteq\mathcal Y_n\) is
\((A,\lambda,B,\alpha)\)-\emph{structurally natural} when:

\begin{enumerate}[label=\textup{(\roman*)}]
\item its membership predicate and the declared sampling or access interface
      for \(\mu_n\) have a complete joint saturated cost at most \(A(n)\);
\item \(\mu_n(\mathcal P_n)\ge\lambda(n)\);
\item every denotation produced by a record in \(\mathscr R_{\alpha,n}\) of
      cost at most \(B(n)\) lies in \(\mathcal Y_n\setminus\mathcal P_n\).
\end{enumerate}

The class index \(\alpha\) may be an exact structural source cell, a circuit
class, or another already-declared certificate class.  A finite union is
handled by listing its indices and charging one common evaluator and access
record.

\end{definition}

\begin{proposition}[Natural property as a charged meta-distinguisher]
\label{prop-natural-property-meta-distinguisher}

An \((A,\lambda,B,\alpha)\)-structurally natural property gives a represented
test of cost \(A\) that accepts a \(\mu_n\)-sample with probability at least
\(\lambda\) and accepts no denotation in the cost-\(B\) cell
\(\mathscr R_{\alpha,n}\).  If a represented reduction converts this test
into a learner, compressor, or source certificate with overhead \(\Psi\),
\cref{thm-typed-meta-compiler-calculus} transfers its guarantee to the
corresponding meta-frontier at cost \(\Psi(A,B,n)\).

For the truth-table presentation, uniform \(\mu_n\), nonuniform circuit cell,
polynomial-time property evaluator in the truth-table length, and
inverse-polynomial \(\lambda\), the three clauses specialize to the classical
constructivity, largeness, and usefulness conditions of natural proofs
\citep{RazborovRudich1997}.

\end{proposition}

\begin{proof}
The first statement is the conjunction of the three defining clauses.  The
second composes the property evaluator with the declared reduction and
applies \cref{thm-typed-meta-compiler-calculus}.
\end{proof}

\section{Represented Hardness Magnification}
\label{sec-represented-hardness-magnification}

Hardness magnification converts a weak lower bound for a meta-problem into a
stronger lower bound for another computational class through an explicit
reduction.  Such phenomena are known for variants of circuit minimization
\citep{OliveiraPichSanthanam2018,OliveiraSanthanam2018,ChenEtAl2020}.  The
framework-level statement is a compiler theorem over complete profile data.

\begin{definition}[Represented profile reduction]
\label{def-represented-profile-reduction}

Let \(\mathfrak D^A_n=(\mathscr R^A_n,g^A_n)\) and
\(\mathfrak D^B_n=(\mathscr R^B_n,g^B_n)\) be complete meta-profile data.  A
represented profile reduction \((\mathcal C,\Psi,\varphi)\) consists of a
family-uniform compiler \(\mathcal C:\mathscr R^A_n\to\mathscr R^B_n\), a
monotone budget map \(\Psi\), and a strictly increasing gain map
\(\varphi:[0,1]\to[0,1]\) such that
\begin{equation}
\label{eq-represented-profile-reduction-conditions}
B\in\mathfrak B_{\mathrm{sat},n}(R)
\Longrightarrow
\Psi(B,n)\in\mathfrak B_{\mathrm{sat},n}(\mathcal C R),
\qquad
g_n^B(\mathcal C R)\ge\varphi(g_n^A(R)).
\end{equation}
The complete compiler record includes its instance map, oracle or sampling
interfaces, retained witness data, soundness relation, and resource
distortion.

\end{definition}

\begin{theorem}[Represented profile magnification]
\label{thm-represented-profile-magnification}

Every represented profile reduction satisfies
\begin{equation}
\label{eq-represented-magnification-frontier-transfer}
B\in\mathfrak P_{\mathfrak D^A,n}^{>}(\theta)
\Longrightarrow
\Psi(B,n)\in
\mathfrak P_{\mathfrak D^B,n}^{>}(\varphi(\theta)).
\end{equation}
Equivalently,
\begin{equation}
\label{eq-represented-magnification-contrapositive}
\Psi(B,n)\notin
\mathfrak P_{\mathfrak D^B,n}^{>}(\varphi(\theta))
\Longrightarrow
B\notin\mathfrak P_{\mathfrak D^A,n}^{>}(\theta).
\end{equation}
Let \(\kappa_A\) and \(\kappa_B\) be scalarizations of the two resource
orders, and suppose the monotone scalar majorant \(\psi\) satisfies
\begin{equation}
\label{eq-represented-magnification-scalar-majorant}
\kappa_B\!\left(\Psi(B,n)\right)
\le \psi\!\left(\kappa_A(B),n\right).
\end{equation}
Extend \(\psi\) by \(\psi(\infty,n)=\infty\).  Then
\begin{equation}
\label{eq-represented-magnification-scalar-cost}
\operatorname{SMC}_{\mathfrak D^B,n,\kappa_B}^{>}(\varphi(\theta))
\le
\inf_{\epsilon>0}
\psi\!\left(
\operatorname{SMC}_{\mathfrak D^A,n,\kappa_A}^{>}(\theta)+\epsilon,n
\right),
\end{equation}
with \(\epsilon=0\) when the source infimum is attained.

Thus a lower bound on the target meta-frontier magnifies to a lower bound on
the source profile precisely through the displayed compiler.  Source-cell
classification supplies the domain of the compiler; the numerical
magnification follows from
\cref{eq-represented-profile-reduction-conditions}.

\end{theorem}

\begin{proof}
Choose \(R\) witnessing
\(B\in\mathfrak P_{\mathfrak D^A,n}^{>}(\theta)\).  Then
\(g_n^A(R)>\theta\), and strict monotonicity gives
\(\varphi(g_n^A(R))>\varphi(\theta)\).  The compiler conditions place
\(\mathcal C R\) below \(\Psi(B,n)\) and give
\(g_n^B(\mathcal C R)>\varphi(\theta)\), proving
\cref{eq-represented-magnification-frontier-transfer}.  The contrapositive is
immediate.  For scalar costs, choose a source record within \(\epsilon\) of
the infimum when it is finite, apply the scalar cost majorant, and take the
infimum over \(\epsilon>0\).  If the source scalar meta-complexity is
infinite, the right-hand side is infinite by convention and the inequality is
automatic.
\end{proof}

\begin{definition}[Gap structural meta-problem]
\label{def-gap-structural-meta-problem}

Let \(\Xi_n\) be an encoded meta-input slice with a represented datum
assignment \(\xi\mapsto\mathfrak D_\xi\).  The input is the encoding of
\(\xi\), together with budget and rational-threshold encodings when those
parameters are not fixed for the problem family.  The verifier and access
interface are the ones retained by \(\mathfrak D_\xi\).

For budgets \(B_0\preceq B_1\) and thresholds
\(0\le\theta_0<\theta_1\le1\), define the promise problem
\begin{equation}
\label{eq-gap-structural-meta-problem}
\operatorname{GapSMC}
\bigl[(B_0,\theta_1),(B_1,\theta_0)\bigr]
\end{equation}
to have yes region
\(B_0\in\mathfrak P_{\mathfrak D_\xi,n}^{\ge}(\theta_1)\) and no region
\(B_1\notin\mathfrak P_{\mathfrak D_\xi,n}^{>}(\theta_0)\).  The promise gap records both
resource and gain separation.  Decision, search, certification, and
deployment versions use \cref{def-four-meta-operational-forms}.

\end{definition}

\section{Application Interface and Chapter Summary}
\label{sec-structural-metacomplexity-summary}

An application evaluates one of the already-defined strict frontiers.  The
source-resolved route evaluates all cells in
\cref{thm-source-cell-meta-complexity-exhaustion}.  The direct operational
route evaluates the selector frontier in
\cref{thm-structural-meta-accountability}.  The represented
selector-to-source comparison of
\cref{thm-prospective-selector-structural-charge,cor-contextual-structural-charge-transfer}
places a selector witness in a source cell at its charged transfer ceiling;
a source-cell lower bound then excludes that selector by contrapositive.
A completed first-hit record remains in the accounting frontier.

\begin{definition}[Family-functional selector meta-datum]
\label{def-family-functional-selector-meta-datum}

For a parameterized presented family with sample size \(m\), channel or
noise parameter \(\zeta\), and horizon \(H\ge1\), let
\(\mathscr R_{n,m,\zeta,H}^{\mathrm{sel}}\) be the class of pairs
\((R,t)\) such that \(0\le t<H\) and \(R\) is a complete family-uniform,
temporally admissible pre-hit selector record.  Write
\(K_{R,I,t}\) and \(\nu_{R,I,t}\) for its represented selector and declared
neutral baseline on instance \(I\).  For a declared monotone family
functional \(\mathfrak M_{n,m,\zeta}\) with values in \([0,1]\), define
\begin{equation}
\label{eq-family-functional-selector-meta-gain}
g_{\mathfrak M,n,m,\zeta,H}^{\mathrm{sel}}(R,t)
=
\mathfrak M_{n,m,\zeta}
\!\left(
I\longmapsto
\bigl(
\operatorname{Adv}_{I,t}(K_{R,I,t};\nu_{R,I,t})
\bigr)_+
\right).
\end{equation}
This gives the complete datum
\begin{equation}
\label{eq-family-functional-selector-meta-datum}
\mathfrak D_{\mathfrak M,n,m,\zeta,H}^{\mathrm{sel}}
=
\left(
\mathscr R_{n,m,\zeta,H}^{\mathrm{sel}},
g_{\mathfrak M,n,m,\zeta,H}^{\mathrm{sel}}
\right).
\end{equation}
Write
\(\mathfrak P_{\mathfrak M,n,m,\zeta,H}^{\mathrm{sel},>}\)
for its strict frontier and
\(\operatorname{SMC}_{\mathfrak M,n,m,\zeta,H,\kappa}^{\mathrm{sel},>}\)
for its scalar meta-complexity.

\end{definition}

By \cref{def-strict-structural-meta-frontier}, the parameter-indexed scalar
quantity is exactly
\begin{equation}
\label{eq-parameter-indexed-selector-metacomplexity}
\operatorname{SMC}_{\mathfrak M,n,m,\zeta,H,\kappa}^{\mathrm{sel},>}(\theta)
=
\inf\left\{
\operatorname{Cost}_{\mathrm{sat}}^{\kappa}(R,t):
\begin{array}{l}
(R,t)\in\mathscr R_{n,m,\zeta,H}^{\mathrm{sel}},\\
g_{\mathfrak M,n,m,\zeta,H}^{\mathrm{sel}}(R,t)>\theta
\end{array}
\right\}.
\end{equation}
Every statistical, noise, data-acquisition, representation, and deployment
term remains inside the complete record cost and its gain evaluation.  The
LPN specialization in \PartVIII uses
\cref{eq-parameter-indexed-selector-metacomplexity} with
\(\zeta=\eta\); its explicit meta-frontier consequence is
\cref{cor-lpn-structural-metacomplexity-lower-region}.

\begin{figure}[tbp]
\centering
\resizebox{.98\linewidth}{!}{%
\begin{tikzpicture}[fw diagram,node distance=8mm and 10mm]
  \node[fw source,text width=2.5cm] (pub) {presentation\\resource order};
  \node[fw provenance,text width=2.6cm,right=of pub] (records)
    {complete\\prospective records};
  \node[fw geom,text width=2.6cm,right=of records] (frontier)
    {Pareto frontier\\\(\mathfrak P^{>}(\theta)\)};
  \node[fw use,text width=2.6cm,right=of frontier] (profile)
    {profile\\inverse};
  \node[fw terminal,text width=2.7cm,right=of profile] (success)
    {selector\\accounting\\success bound};
  \draw[fw arrow] (pub)--(records);
  \draw[fw arrow] (records)--(frontier);
  \draw[fw arrow] (frontier)--(profile);
  \draw[fw arrow] (profile)--(success);

  \node[fw residual,text width=2.5cm,below=13mm of records] (repr)
    {exact\\denotation record};
  \node[fw box,text width=2.6cm,right=of repr] (classical)
    {MCSP / program\\minimization};
  \node[fw box,text width=2.6cm,right=of classical] (interfaces)
    {decision / search /\\certify / deploy};
  \node[fw box,text width=2.7cm,right=of interfaces] (magnify)
    {represented\\compiler\\magnification};
  \draw[fw arrow] (pub)--(repr);
  \draw[fw arrow] (repr)--(classical);
  \draw[fw arrow] (classical)--(interfaces);
  \draw[fw arrow] (interfaces)--(magnify);
  \draw[fw dashed arrow] (magnify)--(success);
\end{tikzpicture}%
}
\caption{Structural meta-complexity.  The upper row takes the exact Pareto
inverse of an existing saturated profile and feeds the prospective selector
accountability theorem.  The lower row separates representation minimization
from extraction and deployment; a magnification conclusion requires a
represented compiler with an explicit resource and gain map.}
\label{fig-structural-metacomplexity-dependency}
\end{figure}

\begin{center}
\small
\begin{tabular}{>{\raggedright\arraybackslash}p{.25\linewidth}
                >{\raggedright\arraybackslash}p{.32\linewidth}
                >{\raggedright\arraybackslash}p{.32\linewidth}}
\toprule
Meta-object & Underlying framework record & Exact result \\
\midrule
source frontier & prospective master certificate &
\cref{thm-source-cell-meta-complexity-exhaustion} \\
selector frontier & pre-hit selector occurrence &
\cref{thm-structural-meta-accountability} \\
active frontier & complete active record &
\cref{thm-strict-meta-profile-generalized-inverse} \\
accounting frontier & completed first-hit record &
\cref{thm-universal-saturated-profile-accountability,rem-meta-source-accounting-separation} \\
representation frontier & exact evaluator denotation &
\cref{thm-circuit-minimization-specialization,thm-program-minimization-specialization} \\
learning transfer & compiled learned selector and normalized source &
\cref{thm-source-resolved-learning-meta-transfer} \\
magnification & represented profile compiler &
\cref{thm-represented-profile-magnification} \\
\bottomrule
\end{tabular}
\end{center}

The meta-frontier is a reindexing of complete represented records by the
least resources needed to cross a numerical threshold.  Its universal
computational consequence is
\cref{thm-structural-meta-accountability}; its source decomposition is
\cref{thm-source-cell-meta-complexity-exhaustion}.  Classical meta-complexity
appears through exact-denotation grammars, while target-directed structural
meta-complexity retains temporal availability, qualification, and deployment.

\clearpage
\part{Resource-Bounded Structural Analysis of LPN}\label{part:lpn}\setcounter{section}{-1}

\section{Application to Learning Parity with
Noise}\label{sec-lpn-introduction}
\label{introduction-from-the-general-invariant-to-a-concrete-family}

\paragraph{Part contract.}
\emph{Inputs:} the presentation-relative saturated profile, exhaustive
Geom/Abs source theory, and accountability theorem from
\cref{part:spectral-complexity}, together with the uniform statistical-attack
bound from \cref{part:learning} and the universal controlled-system
certificate of \cref{part:dynamical-systems}.
The represented-attention compilation of \cref{part:attention} supplies the
source-preserving normal form for normalized active selectors.
The strict Pareto frontier and meta-accountability theorem of
\cref{part:structural-metacomplexity} convert the evaluated LPN selector
profile into an explicit certificate-extraction lower region.
\emph{Output:} an exact reduction of the prospective LPN source profile, a
separate core-refined first-hit accounting bound, an exact constant-noise
spectral perimeter, the hardness theorem under its isolated residual-ancestry
bound, a finite proof-grade empirical protocol that discharges precisely that
bound, the direct implication from fixed-noise LPN hardness to \(P\ne NP\),
the statistical denoising bound, and explicit polynomial, statistical, and
information-theoretic noise thresholds.
\emph{First pass:} the reader should retain the presentation, the saturated
statistical-alignment specialization, the quantitative phase diagram, the
constant-noise spectral ledger, the five source bounds, the core-reduction
theorem, and the final obstruction criterion.

\Cref{part:task-spaces,part:dynamics,part:spectral-complexity,part:learning,part:dynamical-systems,part:attention,part:structural-metacomplexity}
respectively supply task presentations, the budget-admissible channel
decomposition, the resource-saturated structural profile, structural learning
bounds, controlled-system certificates, the represented-attention
compilation, and the strict certificate-cost frontier.  We now apply these
constructions to
Learning Parity with Noise (LPN)
\citep{BlumFurstKearnsLipton1994,HopperBlum2001,BlumKalaiWasserman2003}.

All task-independent notions and implications used in this part are stated in
\cref{part:task-spaces,part:dynamics,part:spectral-complexity,part:learning,part:dynamical-systems,part:attention,part:structural-metacomplexity}.
Accordingly, the statements below either define the presented LPN family or
evaluate an already-defined framework coordinate for the public data
\((A,b,\eta)\).  General counting, source-routing, temporal-accountability,
learning, controlled-dynamics, geometric, quotient, and oracle-extension results are cited from the
preceding parts rather than restated here.
The latent specialization
\cref{thm-lpn-optimal-latent-useful-belief-update} applies optimal acquisition,
Bayes--Fisher updating, compact representation, source-resolved learning, and
attention to the same complete LPN record and bounds the maximum deployed
belief gain over their full polynomial saturated class.

We take the budget family
\(\mathcal B_{\mathrm{poly}}\) and analyze the standard binary LPN
presentation within the presentation-relative model, including every
oracle interface that is conservative in the sense of
\hyperref[def-lpn-admissible-oracle]{LPN oracle admissibility} below.
The principal ledger statements are the prospective-source reduction
\cref{eq-lpn-exact-saturated-source-obstruction} and the prospective success
bound
\[
\begin{aligned}
\operatorname{Succ}_{\mathfrak M,n,\eta}(\mathcal B_C)
\le{}&
H_C\left(2^{-n}
+\operatorname{Sel}_{\mathfrak M,n,\eta}^{\mathrm{sat}}
  (\widehat B_C)\right)\\
&+\varepsilon_{\mathrm{rig}}(n,m)
+\exp[-mD_{\mathrm{Ber}}(\tau\Vert\eta)]\\
&+(2^n-1)2^{-m(1-H_2(\tau))}.
\end{aligned}
\]

The last three terms explicitly account for exceptional public presentations.
The optimal-extraction specialization
\cref{cor-lpn-optimal-structural-extraction-ledger} additionally records the
complete \((n,m,\eta)\)-dependent information, SDPI, source, Caus, Lift,
boundary, and Bellman minimum for every represented acquisition rule.
The compact-belief specialization
\cref{cor-lpn-compact-latent-attention-ledger} refines the same record by latent
code capacity, precision, sufficiency defect, attention-flow error, and the
joint acquisition--encoding--attention optimum.
The exact-\(\Abs\),
exact-quotient Geom, and both compensated quotient--Geom coordinates vanish
by
\cref{cor-lpn-abs-source-exhaustion,cor-lpn-quotient-supported-geom-closure}.
The constant-noise ledger
\cref{thm-lpn-exact-constant-noise-spectral-perimeter} exponentially bounds
every explicitly materialized located spectrum and isolates compact
non-lumpable residual ancestry as the remaining ambient coordinate.  Under
\cref{thm-lpn-residual-ancestry-attention-bound}, the complete saturated LPN
source profile is bounded by \(2^{-n}+(1-\eta)^n\).  Complete controls
and lifts follow the
prospective-value routing of
\cref{thm-no-uncharged-prospective-value-amplification}; the displayed
selector profile retains every directly qualified complete ambient record.
The two ledger statements are
\cref{thm-lpn-saturated-floor,thm-lpn-one-way}.

Here \(b_{\mathrm{poly}}\) is the ceiling of the single declared
polynomial budget structure \(\mathcal B_{\mathrm{poly}}\), including
the class-preserving polynomial overheads already built into that
structure. The reduction proceeds componentwise on this saturated algebra.
The public instance supplies linear algebra and residual statistics.
The framework decomposition assigns every derived coordinate, readout,
and finite postprocessing to its represented prospective source,
post-saturation residual, or completed accounting record.  The complete
ambient coordinate is governed by the obstruction criterion in
\cref{thm-ambient-perturbation-characterization-obstruction}; its two
evaluated restrictions have the bounds just displayed.  First-hit visibility is retained
separately in \(\operatorname{HitCore}^{\mathrm{sat}}\) as a retrospective
audit.

The \hyperref[sec-lpn-structural-core-reduction]{LPN structural-reduction
section} applies the framework decomposition to the full saturated closure.
Primitive geometry, quotient support, derived channels, and residual
accounting are placed in the exhaustive master ledger.  Quotient-supported
source modes vanish, and the ambient coordinate is evaluated through the
framework gate product.
The transition-level dependency is fixed by
\cref{thm-no-free-succinct-transition-rule,%
cor-prospective-nonlinear-readout-routing,%
cor-quotient-free-full-carrier-control-normal-form}: every succinct LPN transition
retains one charged family-uniform derivation from the public presentation.
Its structural use has exact quotient, compensated quotient, or
identity-supported ambient support; terminal readouts and the residual
current receive no prospective target-aligned credit.  The LPN quotient seals
set the pure exact-\(\Abs\), exact-quotient Geom, and both compensated
quotient--Geom source coordinates to zero.  Hence
\cref{cor-quotient-free-full-carrier-control-normal-form} places every
remaining positive sequential-control or reinforcement-learning contribution
in the full-carrier ambient mode, where it is evaluated by the complete
Geom/Caus gate product.  A multiscale or self-similar representation
contributes only through the Geom data that it affordably derives from the
public LPN presentation.  The routing diagram is
\cref{fig-succinct-transition-structural-routing}.

The LPN-specific controlled-dynamics deployment is indexed in
\cref{subsec-lpn-controlled-nonlinear-ambient-geometry}.  Its capacity and
drift, Caus aiming, valid-Lift transfer, public functional profile, and
dynamical-learning coordinates are evaluated in
\cref{thm-lpn-controlled-target-capacity-drift,%
thm-lpn-caus-action-residual-aiming,%
thm-lpn-valid-lift-quantitative-transfer,%
thm-lpn-hypercube-functional-profile,%
thm-lpn-complete-controlled-learning-specialization}; their common-record
dependency is shown in
\cref{fig-lpn-controlled-nonlinear-ambient-dependency}.

Throughout this part, polynomial computation means computation in
\(\mathcal R_n^{\mathcal B_{\mathrm{poly}}}\), the polynomial
\hyperref[def-budgeted-saturated-observable-closure]{budgeted saturated
represented closure} of the presented instance, either directly or
through a presentation-admissible oracle interface. The
\hyperref[thm-oracle-admissibility-presentation-separation]{Oracle
admissibility and presentation separation theorem}
compiles every such interface into the same closure at class-preserving
cost. An oracle evaluator outside this closure belongs to an enriched
task presentation.

\section{The LPN Task Presentation}\label{sec-lpn-task-presentation}
\label{the-lpn-task-presentation}

Except in the noise-indexed comparison of
\cref{sec-lpn-noise-resource-phase}, fix a constant noise rate

\[
0<\eta<\frac12
\]

and a polynomial sample count

\[
m=m(n)=\operatorname{poly}(n).
\]

The LPN sampling experiment draws

\[
A\in \mathbb F_2^{m\times n},
\qquad
s\in\mathbb F_2^n,
\qquad
E\in\mathbb F_2^m,
\]

where the rows of \(A\) are independent uniform vectors, \(s\) is
uniform, and the coordinates of \(E\) are independent Bernoulli
\(\eta\). The public instance is

\[
I=(A,b,\eta),
\qquad
b=As+E.
\]

This is the standard random-sample search-LPN presentation
\citep{HopperBlum2001,Lyubashevsky2005}.

We write \(\mathcal C_A=\operatorname{Im}(A)\subseteq\mathbb F_2^m\)
for the public linear code.

The inversion task is to recover \(s\) from \((A,b,\eta)\).

\begin{definition}[LPN inversion task]\label{def-lpn-task-object}

The candidate space is

\[
X_I=\mathbb F_2^n
\]

with uniform reference measure \(\mu_I\). A point \(x\in X_I\)
represents a candidate secret. The success set is

\[
T_I=\{s\}.
\]

The public observable algebra is generated by \(A\), \(b\), \(\eta\),
candidate-coordinate projections on \(x\), public linear algebra over
\(\mathbb F_2\), the residual map

\[
r_I(x)=Ax+b=A(x+s)+E,
\]

finite postprocessings, admissible lifted coordinates, budget-admissible
boundary readouts, and forward transport data generated by the
\hyperref[def-budgeted-saturated-observable-closure]{budgeted saturated
represented closure} under \(\mathcal B_{\mathrm{poly}}\). The secret
\(s\) and the noise vector \(E\) are not primitive public observables.

\end{definition}

The LPN oracle interface is admissible exactly when it preserves the declared presentation and its charged costs.
\begin{definition}[Presentation-admissible oracle for LPN]
\label{def-lpn-admissible-oracle}

An oracle family \(O=(O_I)\) is presentation-admissible for the standard
LPN task under \(\mathcal B_{\mathrm{poly}}\) when its evaluator has one
family-uniform represented derivation from the
\hyperref[def-lpn-task-object]{LPN presentation}, including every
retained datum and its evaluation rule, and substitution of that
derivation for all oracle calls has class-preserving accumulated cost
within the declared polynomial ceiling. A finite collection of oracle
interfaces is admissible when their joint substitution has the same
property. Equivalently, adjoining these interfaces is a conservative
extension under
\hyperref[thm-oracle-admissibility-presentation-separation]{the Oracle
admissibility and presentation separation theorem}.

\end{definition}

The predicate \(x=s\) specifies the success set for the analysis and
defines the target contrast in the discounted hitting functional. The
public boundary readouts are generated by candidate statistics such as
the Hamming weight of \(r_I(x)\).

For this task presentation, the
\hyperref[def-generated-algebra-target-projection]{centered target
contrast} is

\[
y_s(x)=\mathbf 1_{\{x=s\}}-2^{-n}.
\]

The \hyperref[def-extraction-functional]{saturated extraction
functional} is

\begin{equation}
\label{eq-lpn-saturated-extraction-functional}
I_{\mathrm{LPN},n}^{\mathrm{sat}}(B)
=
 m_{\mathrm{LPN},n}^{\mathrm{sat}}(B)
 +G_{\mathrm{LPN},n}^{\mathrm{sat}}(B),
\end{equation}

where the suprema range over the
\hyperref[def-saturated-representation-degree]{Part~III saturated
representatives}
generated from the LPN presentation at public cost at most \(B\).
By \cref{def-extraction-functional}, this untyped quantity is the accounting
profile.  Its temporally admissible prospective restriction is
\(I_{\mathrm{LPN},n}^{\mathrm{src,sat}}(B)\), and completed first-hit
records remain in the accounting ledger rather than entering that source
restriction.

\begin{example}[Running example: the LPN profile]
\label{ex-lpn-running-part5}

The presentation in \cref{def-lpn-task-object} completes
\cref{ex-lpn-running-part1}. The transcript decomposition of
\cref{ex-lpn-running-part2} and every represented parity quotient from
\cref{ex-lpn-running-part3} now lie in the saturated LPN accounting profile
displayed above.  The remainder of this part evaluates its temporally
admissible Geom/Abs source restriction and retains completed first-hit
visibility in a separate core-refined accounting profile before applying the
accountability theorem.

\end{example}

\section{Statistical Attacks as Saturated LPN Learning}
\label{sec-lpn-saturated-statistical-attacks}

The LPN sample is a homogeneous binary label-corruption experiment.  The
framework theorem in
\cref{thm-saturated-statistical-attack-alignment} therefore controls the
complete saturated class of statistical readouts produced from the public
sample, including data-dependent operators, kernels, features, posteriors,
and terminal selectors.

\subsection{Noisy parity as a corrupted binary target}
\label{subsec-lpn-corrupted-parity-target}

\begin{example}[Learning parity from corrupted labels]
\label{ex-lpn-corrupted-label-learning}

Let \(a\) be uniform on \(\mathbb F_2^n\), let

\[
f_s(a)=(-1)^{\langle a,s\rangle},
\qquad
b=\langle a,s\rangle\oplus e,
\qquad
e\sim\operatorname{Bernoulli}(\eta),
\]

and set \(\widetilde Y=(-1)^b\).  Then

\[
\widetilde Y=f_s(a)(-1)^e,
\]

so \cref{def-binary-label-corruption-channel} applies with
\(\rho=1-2\eta\).  Consider the represented parity class

\[
\mathcal H_{\mathrm{par}}
=
\{f_t:t\in\mathbb F_2^n\},
\qquad
f_t(a)=(-1)^{\langle a,t\rangle}.
\]

Massart's finite-class lemma \citep{Massart2000} gives

\[
\widehat{\mathfrak R}_S
(\mathcal L_{\mathcal H_{\mathrm{par}}})
\le
\sqrt{\frac{2n\log2}{m}}.
\]

Hence the radius

\begin{equation}
\label{eq-lpn-parity-rademacher-radius}
\Gamma_{n,m,\delta}^{\mathrm{par}}
=
2\sqrt{\frac{2n\log2}{m}}
+3\sqrt{\frac{\log(2/\delta)}{2m}}
\end{equation}

may be substituted into \eqref{eq-massart-rademacher-clean-risk}.  Every
\(\varepsilon_{\mathrm{opt}}\)-empirical parity minimizer therefore satisfies

\begin{equation}
\label{eq-lpn-parity-clean-risk}
R_0(\widehat f)
\le
\frac{
2\Gamma_{n,m,\delta}^{\mathrm{par}}
+\varepsilon_{\mathrm{opt}}
}{1-2\eta}.
\end{equation}

Distinct parity functions disagree on exactly half of
\(\mathbb F_2^n\).  Thus

\begin{equation}
\label{eq-lpn-parity-exact-statistical-identification}
2\Gamma_{n,m,\delta}^{\mathrm{par}}
+\varepsilon_{\mathrm{opt}}
<
\frac{1-2\eta}{2}
\end{equation}

certifies \(\widehat f=f_s\).  In particular, if
\(\varepsilon_{\mathrm{opt}}\le(1-2\eta)/4\), there is a universal constant
\(C>0\) such that it suffices to take

\[
m
\ge
C\,
\frac{n+\log(1/\delta)}{(1-2\eta)^2}
\],

where \(C\) is fixed by \eqref{eq-lpn-parity-rademacher-radius}.

For the PAC--Bayes form, let \(P_{\mathrm{par}}\) be uniform on the
\(2^n\) parity vectors.  Then

\[
\mathrm{KL}(\delta_t\Vert P_{\mathrm{par}})=n\log2.
\]

For any posterior \(Q\) on \(\mathcal H_{\mathrm{par}}\), define the
right-hand side of \eqref{eq-pac-bayes-label-noise-clean-risk}, with
\(P=P_{\mathrm{par}}\), by
\(U_{\eta,S}(Q)\).  Since every \(t\ne s\) has clean risk \(1/2\),

\begin{equation}
\label{eq-lpn-pac-bayes-posterior-mass}
R_0(Q)=\frac{1-Q(s)}2,
\qquad
Q(s)\ge1-2U_{\eta,S}(Q).
\end{equation}

A point posterior with \(U_{\eta,S}(\delta_t)<1/2\) is therefore supported
on \(s\), while \(U_{\eta,S}(Q)<1/4\) gives \(Q(s)>1/2\).

Finally, let \(\widehat y(a)\in\mathbb F_2\) be the bit prediction
corresponding to a learned sign predictor and put

\[
\widehat e=b\oplus\widehat y(a).
\]

Then

\begin{equation}
\label{eq-lpn-noise-estimation-from-prediction}
\Pr\{\widehat e\ne e\}
=
\Pr\{\widehat y(a)\ne\langle a,s\rangle\}
=R_0(\widehat f).
\end{equation}

Replacing the corrupted label by the predicted clean parity therefore leaves
exactly the certified clean prediction error.  Under
\eqref{eq-lpn-parity-exact-statistical-identification}, the recovered parity
also reconstructs the noise bits on the observed rows exactly.

The finite-class and PAC--Bayes bounds are simultaneous sample-capacity
statements.  A learner producing \(\widehat f\) or \(Q\), its optimization
transcript, and every data-dependent feature or operator selection retain
their saturated derivation costs under
\cref{def-budget-saturated-statistical-readout-class,thm-paper4-dependency-principle,thm-learned-operator-union-bound}.

\end{example}

\subsection{Uniform target-alignment bound}
\label{subsec-lpn-uniform-statistical-alignment}

For a score \(h:\mathbb F_2^n\to[-1,1]\), put

\[
c_s(h)=\mathbb E_{a\sim\operatorname{Unif}(\mathbb F_2^n)}
\bigl[h(a)f_s(a)\bigr],
\qquad
\widehat c_{\eta,S}(h)
=\frac1m\sum_{i=1}^m h(a_i)(-1)^{b_i}.
\]

\begin{theorem}[Uniform saturated target-alignment bound for LPN
statistical attacks]
\label{thm-lpn-saturated-statistical-alignment}

Fix a budget \(B\), sample length \(m\), and confidence level
\(0<\delta<1\).  Let
\(\mathcal H_{\mathrm{LPN},n,B,m}^{\mathrm{sat,stat}}\) be the
specialization of
\cref{def-budget-saturated-statistical-readout-class} to the LPN
presentation.  With probability at least \(1-\delta\) over the public LPN
sample, simultaneously for every represented statistical attack of complete
cost at most \(B\), every represented seed, and its output
\(h\in\mathcal H_{\mathrm{LPN},n,B,m}^{\mathrm{sat,stat}}\),

\begin{equation}
\label{eq-lpn-uniform-statistical-alignment}
\left|
\widehat c_{\eta,S}(h)-(1-2\eta)c_s(h)
\right|
\le
2\Gamma_{\mathrm{LPN},n,B,m,\delta}^{\mathrm{sat,stat}}.
\end{equation}

On the event \(s\ne0\), the parity target is balanced and

\begin{equation}
\label{eq-lpn-uniform-structural-alignment}
|c_s(h)|^2
\le
M_{\mathrm{LPN},n}^{\mathrm{sat,stat}}(B,m).
\end{equation}

If the attack outputs a candidate \(\widehat s\), append the represented
parity readout \(h_{\widehat s}=f_{\widehat s}\).  This incurs only the
class-preserving cost of evaluating one binary inner product, and

\begin{equation}
\label{eq-lpn-recovery-is-parity-alignment}
c_s(f_{\widehat s})
=
\mathbf 1_{\{\widehat s=s\}}.
\end{equation}

For a randomized attack, averaging
\eqref{eq-lpn-recovery-is-parity-alignment} over its represented seed gives

\begin{equation}
\label{eq-lpn-randomized-recovery-alignment}
\Pr_{\omega}\{\widehat s=s\mid A,b\}
=
\mathbb E_{\omega}\,c_s(f_{\widehat s}).
\end{equation}

The appended parity evaluator and every downstream target use remain in the
same saturated LPN derivation at class-preserving overhead.  An affordable
statistical recovery attack is therefore included in LPN accounting.  The
appended post-output evaluator is a prospective Geom/Abs source only for a
later use at which it satisfies temporal source admissibility; statistical
denoising supplies no additional provenance channel.

\end{theorem}

\begin{proof}

The identity \(\widetilde Y=f_s(a)(-1)^e\) identifies the LPN sample with
the homogeneous corruption model of
\cref{def-binary-label-corruption-channel}.  Apply
\cref{thm-saturated-statistical-attack-alignment} to obtain
\eqref{eq-lpn-uniform-statistical-alignment}.  For \(s\ne0\), the character
\(f_s\) is balanced, so the projection conclusion of that theorem gives
\eqref{eq-lpn-uniform-structural-alignment}.  Since \(s\) is uniform,
\(\Pr\{s=0\}=2^{-n}\).

Walsh-character orthogonality gives

\[
\mathbb E_a[f_t(a)f_s(a)]
=
\mathbf 1_{\{t=s\}},
\]

which proves \eqref{eq-lpn-recovery-is-parity-alignment}; expectation over a
represented random seed proves
\eqref{eq-lpn-randomized-recovery-alignment}.  The parity evaluator is a
uniform represented postprocessing of \(\widehat s\) and \(a\).  Closure and
temporal typing follow from
\cref{thm-paper4-dependency-principle,thm-full-saturated-profile-decomposition,%
cor-first-hit-accounting-not-source}.

\end{proof}

\begin{theorem}[Exact LPN Walsh alignment and charged spectral location]
\label{thm-lpn-walsh-alignment-charged-location}

Fix \(A\in\mathbb F_2^{m\times n}\), let \(s\) be uniform on
\(\mathbb F_2^n\), let \(E\) have independent
\(\operatorname{Bernoulli}(\eta)\) coordinates, and put
\(b=As+E\) and \(\rho=1-2\eta\).  For
\(u\in\mathbb F_2^m\), \(v\in\mathbb F_2^n\), write
\[
\chi_u(z)=(-1)^{\langle u,z\rangle},
\qquad
\chi_v(t)=(-1)^{\langle v,t\rangle}.
\]
Then
\begin{equation}
\label{eq-lpn-exact-walsh-target-correlation}
\mathbb E_{s,E}\!\left[\chi_u(b)\chi_v(s)\mid A\right]
=
\mathbf1_{\{A^{\mathsf T}u=v\}}\rho^{|u|}.
\end{equation}

Let \(h_A:\mathbb F_2^m\to[-1,1]\) be a represented public readout and
normalize its Walsh coefficients by
\begin{equation}
\label{eq-lpn-public-readout-walsh-coefficients}
\widehat h_A(u)
:=
2^{-m}\sum_{z\in\mathbb F_2^m}h_A(z)\chi_u(z).
\end{equation}
Its exact target-character correlation is
\begin{equation}
\label{eq-lpn-exact-readout-target-correlation-spectrum}
\mathbb E_{s,E}\!\left[h_A(As+E)\chi_v(s)\mid A\right]
=
\sum_{u:A^{\mathsf T}u=v}
\widehat h_A(u)\rho^{|u|}.
\end{equation}

Suppose one represented derivation materializes a list
\(L_A\subseteq\mathbb F_2^m\), its indices, and the coefficients used in the
spectral readout
\[
h_{A,L}(z)=\sum_{u\in L_A}\widehat h_A(u)\chi_u(z).
\]
The derivation cost includes construction of \(L_A\), evaluation of its
membership or enumeration rule, construction or evaluation of the retained
coefficients, and evaluation of the displayed readout.  Put
\(\ell_A=|L_A|\).  If every target-aligned retained index has weight at least
\(q\), then
\begin{equation}
\label{eq-lpn-located-walsh-block-attenuation}
\left|
\mathbb E_{s,E}\!\left[h_{A,L}(As+E)\chi_v(s)\mid A\right]
\right|
\le
\sqrt{\ell_A}\,\rho^q.
\end{equation}
In particular, a charged ceiling that permits at most \(L_B\) explicitly
materialized coefficients gives the bound \(\sqrt{L_B}\rho^q\) for every
such located block at that ceiling.

For random \(A\) with independent uniform entries and fixed \(v\ne0\),
\begin{equation}
\label{eq-lpn-short-aligned-walsh-index-probability}
\Pr_A\!\left\{
\exists u:\ 1\le |u|<q,\ A^{\mathsf T}u=v
\right\}
\le
2^{-n}\sum_{k=1}^{q-1}\binom{m}{k}.
\end{equation}
Consequently, if \(n/m\to R\in(0,1)\),
\(q=\lfloor\delta m\rfloor\), \(0<\delta<1/2\), and
\(H_2(\delta)<R\), then
\begin{equation}
\label{eq-lpn-short-aligned-walsh-index-entropy-bound}
\Pr_A\!\left\{
\exists u:\ 1\le |u|<q,\ A^{\mathsf T}u=v
\right\}
\le
2^{-m(R-H_2(\delta)-o(1))}.
\end{equation}

The structural use of the located readout retains this complete derivation.
A locator whose selected indices depend on \(b\) is part of the represented
map \(h_A:b\mapsto h_A(b)\); its selection predicates and coefficient
evaluators are therefore included in the Walsh coefficients in
\cref{eq-lpn-public-readout-walsh-coefficients} and retain their complete
derivation cost.  The existence of an index satisfying
\(A^{\mathsf T}u=v\), or the analytic appearance of such an index in
\cref{eq-lpn-exact-readout-target-correlation-spectrum}, is not a represented
locator.  A target-qualified spectral record retains both the pre-hit rule
that selects its indices from the public presentation and the represented
comparison by which the selected readout is used as target-relevant.  In the
target-character specialization \(v=s\), the equality
\(A^{\mathsf T}u=s\) is an analytic test against the hidden target; it is not
available to the selection rule unless the public derivation independently
represents the corresponding comparison.  Such a locator receives the explicit block bound
\cref{eq-lpn-located-walsh-block-attenuation} precisely when its represented
output materializes the displayed list expansion.  Otherwise its complete
readout is evaluated by the exact spectrum
\cref{eq-lpn-exact-readout-target-correlation-spectrum} and by the existing
saturated statistical projection and Geom gates.
A target-qualified nonidentity fiber map generated by \(h_{A,L}\) belongs to
the existing \(\Abs\) profile.  An identity-supported use belongs to the
existing ambient \(\Geom\) profile and retains its mixing, authority,
\(\Caus\)-alignment, reach, regularity, and temporal gates.  A rule that
selects \(L_A\), a coefficient, a threshold, or a downstream transition is
part of the same represented derivation and supplies no independent source
coordinate.

\end{theorem}

\begin{proof}

Since \(b=As+E\),
\[
\chi_u(b)\chi_v(s)
=
(-1)^{\langle A^{\mathsf T}u+v,s\rangle}
(-1)^{\langle u,E\rangle}.
\]
The expectation of the first factor over uniform \(s\) is one when
\(A^{\mathsf T}u=v\) and zero otherwise.  Independence of the noise
coordinates gives
\[
\mathbb E_E(-1)^{\langle u,E\rangle}
=\prod_{i:u_i=1}\mathbb E(-1)^{E_i}
=\rho^{|u|},
\]
proving \cref{eq-lpn-exact-walsh-target-correlation}.  Expanding \(h_A\) in
the Walsh basis and applying that identity termwise proves
\cref{eq-lpn-exact-readout-target-correlation-spectrum}.

For the located block, only indices satisfying \(A^{\mathsf T}u=v\)
contribute.  Cauchy--Schwarz and Parseval give
\begin{align*}
\left|
\sum_{\substack{u\in L_A\\A^{\mathsf T}u=v}}
\widehat h_A(u)\rho^{|u|}
\right|
&\le
\rho^q\sqrt{\ell_A}
\left(\sum_{u\in L_A}|\widehat h_A(u)|^2\right)^{1/2}\\
&\le
\rho^q\sqrt{\ell_A}\,
\|h_A\|_{L^2(\operatorname{Unif}(\mathbb F_2^m))}
\le \rho^q\sqrt{\ell_A}.
\end{align*}
Every explicitly retained index and coefficient occupies a represented output
coordinate, so the output-size and evaluation ledgers bound their number by
the declared ceiling; this proves the \(L_B\) specialization.

For fixed nonzero \(u\), the vector \(A^{\mathsf T}u\) is uniform on
\(\mathbb F_2^n\).  Hence
\(\Pr\{A^{\mathsf T}u=v\}=2^{-n}\).  A union bound over the
\(\sum_{k=1}^{q-1}\binom{m}{k}\) vectors of the displayed weights proves
\cref{eq-lpn-short-aligned-walsh-index-probability}.  The standard entropy
bound
\(\sum_{k<\delta m}\binom{m}{k}\le2^{m(H_2(\delta)+o(1))}\)
proves \cref{eq-lpn-short-aligned-walsh-index-entropy-bound}.

Finally, the represented list-construction and coefficient-evaluation rules
are ancestors of every use of \(h_{A,L}\), so
\cref{thm-paper4-dependency-principle,thm-affordable-composition-no-creation}
retain their costs.  By \cref{def-lpn-task-object}, \(s\) is not a primitive
public observable.  Hence the analytic equality \(A^{\mathsf T}u=s\) cannot
be evaluated by a public locator unless its comparison is separately
derived from the declared presentation, in which case that derivation is an
ancestor and is charged by the same two results.  The support classification is
\cref{thm-affordable-geom-abs-normal-form}; temporal and downstream-use
classification is
\cref{def-pre-hit-temporally-admissible-source,thm-two-level-learned-readout-provenance,cor-contextual-structural-charge-transfer}.

\end{proof}

\Cref{fig-lpn-statistical-alignment-chain} records the distinct statistical,
spectral-location, and verification conversions used by the LPN alignment
argument.

\begin{figure}[tbp]
\centering
\resizebox{\linewidth}{!}{%
\begin{tikzpicture}[fw diagram,node distance=8mm and 8mm]
  \node[fw source,text width=2.4cm] (sample) {noisy rows\\\((A,b),\ m,\eta\)};
  \node[fw box,text width=2.3cm,right=of sample] (learner) {represented learner\\output \(h\)};
  \node[fw use,text width=2.5cm,right=of learner] (emp) {empirical noisy\\correlation};
  \node[fw geom,text width=2.5cm,right=of emp] (clean) {clean parity\\correlation\\divide by \(1-2\eta\)};
  \node[fw provenance,text width=2.4cm,right=of clean] (loc) {charged Walsh\\index locator};
  \node[fw terminal,text width=2.3cm,right=of loc] (verify) {candidate secret\\public verifier};
  \draw[fw arrow] (sample)--(learner);
  \draw[fw arrow] (learner)--(emp);
  \draw[fw arrow] (emp)--(clean);
  \draw[fw arrow] (clean)--(loc);
  \draw[fw arrow] (loc)--(verify);
  \node[fw residual,text width=4.8cm,below=8mm of clean] (ledger)
    {one budget ledger: fitting, evaluation, location, retained state, and verification};
  \draw[fw dashed arrow] (learner)--(ledger);
  \draw[fw dashed arrow] (loc)--(ledger);
\end{tikzpicture}%
}
\caption{LPN alignment is a chain of represented conversions.  Statistical
uniformity moves empirical noisy correlation to population correlation,
noise inversion supplies the factor \((1-2\eta)^{-1}\), and a separate
charged locator turns spectral mass into a candidate index before public
verification.  Formal Fourier mass alone skips neither the locator nor its
cost.}
\label{fig-lpn-statistical-alignment-chain}
\end{figure}

\begin{corollary}[Spectral mass requires charged location and target
qualification]
\label{cor-lpn-spectral-mass-requires-charged-location}

In the setting of
\cref{thm-lpn-walsh-alignment-charged-location}, a prospective spectral
source consists of one represented pre-hit derivation containing its index
locator, coefficient evaluator, target-relevance comparison, and downstream
use.  Its source value is evaluated from that complete record.  An
unmaterialized term of
\cref{eq-lpn-exact-readout-target-correlation-spectrum} contributes no
separate source value, and the condition \(A^{\mathsf T}u=s\) cannot be used
as a public location oracle.  Explicitly materialized located blocks obey
\cref{eq-lpn-located-walsh-block-attenuation}; a compact evaluator that does
not materialize its spectrum is evaluated as one complete readout through
the existing exact-\(\Abs\), compensated, terminal-readout, or ambient-
\(\Geom\) mode.

\end{corollary}

\begin{proof}

Temporal admissibility and complete ancestry are
\cref{def-pre-hit-temporally-admissible-source,thm-paper4-dependency-principle}.
The exact target correlation of the complete readout is
\cref{eq-lpn-exact-readout-target-correlation-spectrum}; its located-block
specialization is
\cref{eq-lpn-located-walsh-block-attenuation}.  The exhaustive support
assignment is \cref{thm-affordable-geom-abs-normal-form}, and isolated
terminal discrimination is separated from prospective structural activity
by \cref{cor-terminal-verifier-not-prospective-structural-activity}.

\end{proof}

\begin{theorem}[A formal diffuse spectrum creates no independent structural
source]
\label{thm-lpn-no-independent-diffuse-spectral-source}

Fix the LPN presentation of \cref{def-lpn-task-object}, a charged ceiling
\(B\), and a represented public residual readout
\(h_A:\mathbb F_2^m\to[-1,1]\).  Regard
\cref{eq-lpn-exact-readout-target-correlation-spectrum} as the exact analytic
evaluation of its target-character correlation.  The formal Walsh expansion
of \(h_A\) introduces no represented source vertex and no additional source
coordinate.  The explicit/implicit distinction and every prospective
consumer are charged by
\cref{thm-complete-spectral-aggregation-charging}.  Every prospective use of
\(h_A(b)\) has exactly one of the following existing assignments.

\begin{enumerate}[label=\textup{(\roman*)}]
\item If its represented derivation materializes indices and coefficients,
      their locator, target-relevance comparison, coefficient evaluation, and
      downstream use remain in the complete record.  A materialized block is
      governed by
      \cref{eq-lpn-located-walsh-block-attenuation}.
\item If the complete readout generates a target-qualified nonidentity fiber
      map, its dynamically qualified use belongs to the existing exact or
      compensated quotient-supported mode.
\item If the complete readout is used with identity support, its prospective
      value is the existing ambient-\(\Geom\) qualified visibility.  Its
      potential, represented local differences, transition rule,
      same-generator authority, authorized \(\Caus\) alignment, target reach,
      regularity, temporal provenance, and complete evaluation cost remain in
      that one record.  The Walsh expansion and the derived \(\Caus\) record
      add no independent orientation source.
\item If the readout supplies target discrimination or an extremum but no
      represented authorized intermediate transport, it is a terminal
      verifier and has zero prospective source activity.
\item If an evaluator explicitly processes \(N\) distinct spectral terms,
      its complete evaluation ledger contains at least \(N\) represented term
      evaluations.  An explicit expansion with superpolynomial \(N\) is
      outside the declared polynomial ceiling.
\end{enumerate}

In particular, the analytic restriction
\[
  \sum_{u:A^{\mathsf T}u=s}
  \widehat h_A(u)(1-2\eta)^{|u|}
\]
does not define a public aggregation rule: its indexing predicate contains
the nonpublic secret \(s\).  A public derivation that realizes the same
denotation by a compact circuit, factorization, recurrence, or another
represented rule is classified and charged by
items~\textup{(i)}--\textup{(iv)} according to that rule.  Consequently there
is no separate diffuse-aggregation source between the five-family
provenance ledger and the existing exact-\(\Abs\), compensated, ambient-
\(\Geom\), terminal-accounting, and cost coordinates.

\end{theorem}

\begin{proof}

By \cref{def-lpn-task-object}, the primitive public observables contain
\(A,b,\eta\) but not \(s\).  Hence \(A^{\mathsf T}u=s\) is available in the
analytic correlation calculation but is not a represented public selection
predicate.  By
\cref{def-pre-hit-temporally-admissible-source,thm-paper4-dependency-principle,%
thm-affordable-composition-no-creation}, every represented locator,
coefficient evaluator, comparison, and downstream transition remains in the
ancestral derivation with its complete cost.  This proves the charging
assertion in item~\textup{(i)}, and
\cref{thm-lpn-walsh-alignment-charged-location} gives its displayed bound.

The support of the complete readout is either nonidentity or identity.
\Cref{thm-affordable-geom-abs-normal-form} assigns a dynamically qualified
nonidentity record to the exact or compensated quotient-supported modes and
an identity-supported record to ambient \(\Geom\).  In the latter case,
\cref{def-geom-extraction,def-authority-compatible-caus-alignment} retain all
same-record Geom gates, while
\cref{prop-derived-channels-create-no-structure} assigns zero independent
source creation to the derived \(\Caus\) record.  This proves
items~\textup{(ii)} and~\textup{(iii)}.

\Cref{cor-extremal-readout-requires-intermediate-transport} proves
item~\textup{(iv)}.  For item~\textup{(v)}, each explicitly processed term is
an evaluated vertex of the represented derivation.  The evaluation ledger
therefore dominates the number of such vertices, so a superpolynomial
explicit expansion cannot lie below a polynomial ceiling.  These cases
exhaust whether spectral terms are materialized and, within a compact
readout, the existing support and temporal assignments.  No additional
source coordinate remains.

\end{proof}

\begin{corollary}[Fixed-ceiling suppression of located LPN spectral blocks]
\label{cor-lpn-fixed-ceiling-located-spectral-suppression}

Fix one declared ceiling whose explicitly materialized spectral output has
length at most \(L_B(n)\le n^C\), fix \(0<\eta<1/2\), and suppose
\(n/m\to R\in(0,1)\).  Choose \(0<\delta<1/2\) with
\(H_2(\delta)<R\), and put \(q=\lfloor\delta m\rfloor\).  For each fixed
nonzero target character \(v\), outside an event of probability at most
\begin{equation}
\label{eq-lpn-located-block-exceptional-probability}
2^{-m(R-H_2(\delta)-o(1))},
\end{equation}
every represented located block at that ceiling satisfies
\begin{equation}
\label{eq-lpn-fixed-ceiling-located-block-bound}
\left|
\mathbb E_{s,E}
\left[h_{A,L}(As+E)\chi_v(s)\mid A\right]
\right|
\le
n^{C/2}(1-2\eta)^{\lfloor\delta m\rfloor}.
\end{equation}
For fixed \(R,\eta,\delta,C\), the right-hand side is
\begin{equation}
\label{eq-lpn-fixed-ceiling-located-block-exponent}
\exp\!\left[
-\delta m\log\frac1{1-2\eta}
+\frac C2\log n+O(1)
\right]
=\exp[-\Theta(n)].
\end{equation}

A represented low-weight exception, when it occurs, consists of an explicitly
located \(u\) satisfying \(A^{\mathsf T}u=v\).  Its two-valued readout
\(\chi_u(b)\) has exact target-character correlation
\((1-2\eta)^{|u|}\).  When a represented downstream use turns this readout
into a nonidentity fiber map on the active candidate or transcript carrier and
passes the quotient-utility condition of
\cref{def-paper1-abs-usefulness-gate}, that complete use belongs to the
existing \(\Abs\) profile.  The quotient derivation includes the rule that
located \(u\), the character evaluator, the carrier map, and the quotient
dynamics.

\end{corollary}

\begin{proof}

On the complement of the event in
\cref{eq-lpn-short-aligned-walsh-index-entropy-bound}, every aligned index has
weight at least \(q\).  Apply
\cref{eq-lpn-located-walsh-block-attenuation} and
\(\ell_A\le L_B(n)\le n^C\) to obtain
\cref{eq-lpn-fixed-ceiling-located-block-bound}.  Taking logarithms gives
\cref{eq-lpn-fixed-ceiling-located-block-exponent}; since \(m=\Theta(n)\),
its negative linear term dominates its logarithmic term.  The last paragraph
uses \cref{eq-lpn-exact-walsh-target-correlation} and the exact-support branch
of \cref{thm-affordable-geom-abs-normal-form}.  The locator remains an ancestor
by \cref{thm-paper4-dependency-principle}.

\end{proof}

\section{Constant-Noise Spectral Closure Ledger}
\label{sec-lpn-constant-noise-spectral-closure}

This section consolidates the constant-noise part of the spectral argument.
Throughout it, \(0<\eta<1/2\) is fixed, \(n/m\to R\in(0,1)\), and the
ceiling is one declared polynomial saturated ceiling.  The vanishing-noise
regimes remain in
\cref{sec-lpn-noise-resource-phase,cor-lpn-quantitative-noise-phase-localization}.
The profile being evaluated here is prospective: completed first-hit records
remain in the accounting profile of
\cref{def-source-accounting-profile-separation}.
\Cref{fig-lpn-constant-noise-spectral-routing} displays the resulting
spectral routing.

\begin{lemma}[Residual-parity spectral fibers]
\label{lem-lpn-residual-parity-spectral-fibers}

Let one represented pre-hit derivation materialize
\(L_I=(u_1,\ldots,u_k)\in(\mathbb F_2^m)^k\).  Its natural residual-character
map on the LPN candidate carrier is
\begin{equation}
\label{eq-lpn-residual-parity-spectral-map}
q_{L,I}(x)
=
\bigl(\langle u_j,Ax+b\rangle\bigr)_{j=1}^{k}
=M_{L,A}x+c_{L,b},
\qquad
(M_{L,A})_{j,\cdot}=u_j^{\mathsf T}A.
\end{equation}
For every candidate \(x\),
\begin{equation}
\label{eq-lpn-residual-parity-target-fiber}
q_{L,I}^{-1}\bigl(q_{L,I}(s)\bigr)
=s+\ker M_{L,A}.
\end{equation}
On the verifier-good event of \cref{thm-lpn-public-verifier}, where
\(T_I=\{s\}\), the target is a union of \(q_{L,I}\)-fibers if and only if
\begin{equation}
\label{eq-lpn-spectral-terminal-rank-test}
\ker M_{L,A}=\{0\}
\quad\Longleftrightarrow\quad
\rank M_{L,A}=n.
\end{equation}
In this case \(q_{L,I}\) is injective and its represented partition is
identity-equivalent.  When \(\rank M_{L,A}<n\), the target fiber contains a
nonzero affine translate of \(s\), so the terminal-event clause of
\cref{def-paper1-abs-usefulness-gate} fails.

\end{lemma}

\begin{proof}
The affine form in \cref{eq-lpn-residual-parity-spectral-map} follows from
linearity of the binary inner product.  Therefore
\(q_{L,I}(x)=q_{L,I}(s)\) precisely when
\(M_{L,A}(x+s)=0\), proving
\cref{eq-lpn-residual-parity-target-fiber}.  A singleton target is a union of
fibers precisely when its own fiber is the singleton.  This is equivalent to
\(\ker M_{L,A}=\{0\}\), hence to full column rank.  Full column rank makes
the affine map injective.  This is the spectral-list specialization of the
additive fiber identity in \cref{thm-lpn-additive-quotient-identity}.
\end{proof}

\begin{proposition}[Spectral selection and the complete quotient gate]
\label{prop-spectral-selector-complete-abs-gate}

Let a family-uniform, temporally admissible spectral-selection record
construct a candidate- or transcript-carrier map \(q_I:\Omega_I\to Q_I\).
The record is an exact \(\Abs\) source precisely when the same complete
derivation supplies all five conditions of
\cref{def-paper1-abs-usefulness-gate}: represented fibers, terminal-compatible
cells, exact intertwining with the selected dynamics, positive normalized
quotient utility, and strict refinement of the verifier partition.  A
nonzero represented intertwining defect can contribute only through a
dynamically qualified compensated quotient--Geom certificate satisfying the
active-geometry and stability/transfer gates of
\cref{def-supported-geom-certificate,def-master-compensated-quotient-geometry-certificate}.
An identity-supported dynamically qualified use is evaluated in ambient
\(\Geom\).

The description length of the selector, its family uniformity, and its
coefficient-location cost establish budget admissibility of the objects they
actually construct.  They do not replace any of the five quotient gates.
Appending the public verifier to a spectral label supplies terminal
compatibility, but exact or compensated dynamics and positive prospective
utility remain properties of the resulting complete record.

\end{proposition}

\begin{proof}
A fiber-represented selected observable induces the static algebraic quotient
of \cref{def-algebraic-quotient}.  The exact \(\Abs\) coordinate is defined by
\cref{def-paper1-abs-usefulness-gate,def-abs-extraction}, which sets its
utility to zero unless all five displayed gates hold on one normalized
record.  The exact, nonzero-defect, and identity-support assignments are the
three corresponding branches of
\cref{thm-affordable-geom-abs-normal-form}.  Temporal availability and
complete ancestry are retained by
\cref{def-pre-hit-temporally-admissible-source,thm-paper4-dependency-principle}.
These definitions prove every assertion.
\end{proof}

\begin{theorem}[Coset-energy envelope for a complete bounded readout]
\label{thm-lpn-complete-readout-coset-energy-envelope}

In the setting of
\cref{thm-lpn-walsh-alignment-charged-location}, fix
\(v\ne0\), put \(\rho_\eta=1-2\eta\), and let
\(h_A:\mathbb F_2^m\to[-1,1]\) be any complete represented readout.  Then,
pointwise in \(A\),
\begin{equation}
\label{eq-lpn-complete-readout-coset-energy-pointwise}
\left|
\mathbb E_{s,E}
 \left[h_A(As+E)\chi_v(s)\mid A\right]
\right|^2
\le
\sum_{u:A^{\mathsf T}u=v}\rho_\eta^{2|u|}.
\end{equation}
For a random binary matrix \(A\),
\begin{equation}
\label{eq-lpn-complete-readout-coset-energy-average}
\mathbb E_A
\left|
\mathbb E_{s,E}
 \left[h_A(As+E)\chi_v(s)\mid A\right]
\right|
\le
\min\left\{
1,
2^{-n/2}\sqrt{(1+\rho_\eta^2)^m-1}
\right\}.
\end{equation}
Consequently, when \(n/m\to R\) and
\begin{equation}
\label{eq-lpn-coset-energy-decay-region}
R>\log_2(1+\rho_\eta^2),
\end{equation}
the right-hand side of
\cref{eq-lpn-complete-readout-coset-energy-average} is
\begin{equation}
\label{eq-lpn-coset-energy-decay-exponent}
\exp\!\left[
-\frac m2
\left(R-\log_2(1+\rho_\eta^2)-o(1)\right)\log 2
\right].
\end{equation}
The same bounds hold after averaging over represented randomness used to
choose \(h_A\), provided that randomness is independent of \((s,E)\)
conditional on \(A\).

\end{theorem}

\begin{proof}
Apply \cref{eq-lpn-exact-readout-target-correlation-spectrum} and
Cauchy--Schwarz:
\begin{align*}
\left|
\sum_{u:A^{\mathsf T}u=v}
\widehat h_A(u)\rho_\eta^{|u|}
\right|^2
&\le
\left(\sum_{u:A^{\mathsf T}u=v}|\widehat h_A(u)|^2\right)
\left(\sum_{u:A^{\mathsf T}u=v}\rho_\eta^{2|u|}\right)\\
&\le
\sum_{u:A^{\mathsf T}u=v}\rho_\eta^{2|u|},
\end{align*}
where Parseval and \(|h_A|\le1\) bound the first factor by one.  This proves
\cref{eq-lpn-complete-readout-coset-energy-pointwise}.

For every nonzero \(u\), \(A^{\mathsf T}u\) is uniform on
\(\mathbb F_2^n\), while \(u=0\) cannot satisfy
\(A^{\mathsf T}u=v\ne0\).  Hence
\begin{align*}
\mathbb E_A
\sum_{u:A^{\mathsf T}u=v}\rho_\eta^{2|u|}
&=2^{-n}\sum_{u\ne0}\rho_\eta^{2|u|}\\
&=2^{-n}\bigl((1+\rho_\eta^2)^m-1\bigr).
\end{align*}
Jensen's inequality for the square root and the unit bound on correlation give
\cref{eq-lpn-complete-readout-coset-energy-average}.  Substituting
\(n=Rm+o(m)\) gives
\cref{eq-lpn-coset-energy-decay-exponent}.  Conditioning on represented
randomness proves the final assertion.
\end{proof}

\begin{theorem}[Fixed-noise suppression of materialized located spectra]
\label{thm-lpn-constant-noise-materialized-spectral-closure}

Fix \(C<\infty\), choose \(0<\delta<1/2\) with
\(H_2(\delta)<R\), put \(q=\lfloor\delta m\rfloor\), and set
\(\rho_\eta=1-2\eta\).  Suppose the declared ceiling permits at most
\(L_B(n)\le n^C\) explicitly materialized spectral indices.  For every
fixed nonzero secret-space character \(v\in\mathbb F_2^n\), every
represented located block from
\cref{thm-lpn-walsh-alignment-charged-location} satisfies
\begin{align}
\label{eq-lpn-constant-noise-materialized-family-bound}
\mathbb E_A
\left|
\mathbb E_{s,E}\!\left[
h_{A,L}(As+E)\chi_v(s)\mid A
\right]
\right|
\le{}&
n^{C/2}\rho_\eta^{\lfloor\delta m\rfloor}
+n^{C/2}2^{-m(R-H_2(\delta)-o(1))}.
\end{align}
Each term on the right is \(\exp[-\Theta(n)]\) at fixed
\((R,\eta,\delta,C)\).

\end{theorem}

\begin{proof}
Fix \(v\ne0\).  Outside the short aligned-coset event in
\cref{eq-lpn-short-aligned-walsh-index-entropy-bound}, every
\(u\) satisfying \(A^{\mathsf T}u=v\) has weight at least
\(q\).  The located-block estimate
\cref{eq-lpn-located-walsh-block-attenuation} and
\(L_B(n)\le n^C\) give
\(n^{C/2}\rho_\eta^q\).  On the exceptional event, the \(q=0\) case of
\cref{eq-lpn-located-walsh-block-attenuation} bounds the correlation by
\(\sqrt{L_B(n)}\le n^{C/2}\); that event has probability at most
\(2^{-m(R-H_2(\delta)-o(1))}\).  This proves the fixed-nonzero-character
claim after expectation over \(A\), namely
\cref{eq-lpn-constant-noise-materialized-family-bound}.  Finally,
\(m=\Theta(n)\) and \(0<\rho_\eta<1\), so
\begin{equation*}
n^{C/2}\rho_\eta^{\lfloor\delta m\rfloor}
=
\exp\!\left[-\delta m\log\frac1{\rho_\eta}
+\frac C2\log n+O(1)\right]
=\exp[-\Theta(n)].
\end{equation*}
The exceptional exponent is positive because \(H_2(\delta)<R\).
Thus
\(n^{C/2}2^{-m(R-H_2(\delta)-o(1))}=\exp[-\Theta(n)]\) as well.
\end{proof}

\begin{theorem}[Rate-uniform semantic-cancellation bound for LPN]
\label{thm-lpn-rate-uniform-semantic-cancellation}

Let \(A\in\mathbb F_2^{m\times n}\) have independent uniform entries.  For
integers \(1\le q_0\le m\) and \(0\le k<n\), define the short-provenance,
small-semantic-support event
\begin{equation}
\label{eq-lpn-short-provenance-small-semantic-event}
\mathcal E_{q_0,k}
=
\left\{
\begin{array}{l}
\exists\,0\ne\lambda\in\mathbb F_2^m,
\ \exists S\subseteq[n]:\\[-1mm]
|\lambda|<q_0, |S|\le k,
\ 0\ne A^{\mathsf T}\lambda\in\mathbb F_2^S
\end{array}
\right\}.
\end{equation}
Then
\begin{equation}
\label{eq-lpn-short-provenance-small-semantic-probability}
\Pr_A(\mathcal E_{q_0,k})
\le
\left(\sum_{j=0}^{k}\binom nj2^{-(n-j)}\right)
\left(\sum_{r=1}^{q_0-1}\binom mr\right).
\end{equation}

Let a complete charged noisy-combination record materialize at most \(L_B\)
rows \(\lambda_j\), each with nonzero semantic label
\(A^{\mathsf T}\lambda_j\) supported on at most \(k\) secret coordinates.
For every normalized located character block \(H\) constructed from those
rows and every fixed positive noise \(0<\eta<1/2\),
\begin{align}
\label{eq-lpn-rate-uniform-combination-bound}
\mathbb E_A
\left|\mathbb E_{s,E}[H\chi_v(s)\mid A]\right|
\le{}&
\sqrt{L_B}(1-2\eta)^{q_0}
\notag\\
&+\sqrt{L_B}
\left(\sum_{j=0}^{k}\binom nj2^{-(n-j)}\right)
\left(\sum_{r=1}^{q_0-1}\binom mr\right),
\end{align}
uniformly for every nonzero semantic character \(v\) supported on the
declared terminal coordinate set.

Equivalently, define the finite-parameter margins
\begin{align}
\label{eq-lpn-combination-noise-margin}
\mathfrak E_{\mathrm{noise}}(n,m,\eta,q_0,L_B)
&:=
q_0\log\frac1{1-2\eta}-\frac12\log L_B,\\
\label{eq-lpn-combination-entropy-margin}
\mathfrak E_{\mathrm{comb}}(n,m,q_0,k,L_B)
&:=
(n-k)\log2
-\log\sum_{j=0}^{k}\binom nj
-\log\sum_{r=1}^{q_0-1}\binom mr
-\frac12\log L_B.
\end{align}
Then \cref{eq-lpn-rate-uniform-combination-bound} implies the explicit estimate
\begin{equation}
\label{eq-lpn-combination-two-margin-bound}
\mathbb E_A
\left|\mathbb E_{s,E}[H\chi_v(s)\mid A]\right|
\le
e^{-\mathfrak E_{\mathrm{noise}}}
+e^{-\mathfrak E_{\mathrm{comb}}}.
\end{equation}
It is negligible along any parameter schedule for which both margins grow
faster than \(\log n\).  This finite criterion retains \(n,m,\eta,q_0,k\),
and \(L_B\) without assuming a constant-rate limit.

Suppose \(n/m\to R\in(0,1)\), \(q_0=\lfloor\delta m\rfloor\), and
\(k/n\to\kappa\in[0,1/2)\).  If
\begin{equation}
\label{eq-lpn-semantic-cancellation-entropy-condition}
H_2(\delta)
<R\bigl(1-\kappa-H_2(\kappa)\bigr),
\end{equation}
then the exceptional probability in
\cref{eq-lpn-short-provenance-small-semantic-probability} is at most
\begin{equation}
\label{eq-lpn-semantic-cancellation-entropy-exponent}
2^{-m\left[
R(1-\kappa-H_2(\kappa))-H_2(\delta)-o(1)
\right]}.
\end{equation}
In particular, for \(k=o(n)\) and every fixed \(R>0\), one may choose
\(\delta>0\) with \(H_2(\delta)<R\).  If \(L_B\) is polynomial, both terms in
\cref{eq-lpn-rate-uniform-combination-bound} are then exponentially small.
This conclusion has no lower-rate restriction of the form
\cref{eq-lpn-coset-energy-decay-region}.

\end{theorem}

\begin{proof}
Fix nonzero \(\lambda\) and a coordinate set \(S\).  The random vector
\(A^{\mathsf T}\lambda\) is uniform on \(\mathbb F_2^n\), so
\[
\Pr_A\{0\ne A^{\mathsf T}\lambda\in\mathbb F_2^S\}
=\frac{2^{|S|}-1}{2^n}
\le2^{-(n-|S|)}.
\]
Union bounds over the displayed choices of \(S\) and \(\lambda\) prove
\cref{eq-lpn-short-provenance-small-semantic-probability}.

On \(\mathcal E_{q_0,k}^{\mathsf c}\), every materialized row with a
nonzero semantic label supported on at most \(k\) coordinates has provenance
weight at least \(q_0\).  Apply
\cref{eq-generic-combination-capacity-attenuation} with
\(\rho=1-2\eta\).  On \(\mathcal E_{q_0,k}\), use the same located-block
bound with unit attenuation.  Averaging over \(A\) and substituting
\cref{eq-lpn-short-provenance-small-semantic-probability} gives
\cref{eq-lpn-rate-uniform-combination-bound}.

For \(k/n\to\kappa<1/2\), the standard entropy estimates give
\begin{align*}
\sum_{j=0}^{k}\binom nj2^{-(n-j)}
&\le
2^{-n(1-\kappa-H_2(\kappa)-o(1))},\\
\sum_{r=1}^{q_0-1}\binom mr
&\le2^{m(H_2(\delta)+o(1))}.
\end{align*}
Using \(n=Rm+o(m)\) proves
\cref{eq-lpn-semantic-cancellation-entropy-exponent}.  When \(k=o(n)\),
\(\kappa=0\); since \(R>0\), continuity of \(H_2\) at zero permits a fixed
\(\delta>0\) with \(H_2(\delta)<R\).  Fixed positive noise makes
\((1-2\eta)^{\delta m}\) exponentially small, and a polynomial factor
does not change either exponential conclusion.
\end{proof}

\begin{corollary}[High-sample closure of the materialized additive-combination
sector]
\label{cor-lpn-high-sample-materialized-combination-closure}

Fix \(R\in(0,1)\), \(0<\eta<1/2\), and a polynomial saturated ceiling.
Consider the source-resolved LPN subprofile whose complete records

\begin{enumerate}[label=\textup{(\roman*)}]
\item construct their prospective target-bearing characters through the
      charged binary combination interface of
      \cref{def-charged-binary-noisy-combination-record};
\item materialize polynomially many located rows and coefficients; and
\item use a terminal coordinate support of size \(k=o(n)\), including every
      explicitly enumerated terminal table of polynomial cardinality.
\end{enumerate}

Then its saturated located-combination profile is exponentially small under
the random LPN presentation law.  This holds throughout the high-sample
region
\[
R\le\log_2\bigl(1+(1-2\eta)^2\bigr)
\]
as well as in the coset-energy decay region.

Every row locator, collision rule, shared combination DAG, sample reuse,
terminal table, decoder, verifier, and prospective transition remains in the
complete record.  A fixed program implementing the construction is admitted
at its actual execution, acquisition, memory, and output cost; fixed code
length alone supplies no polynomial-cost certificate.

\end{corollary}

\begin{proof}
Choose \(\delta>0\) with \(H_2(\delta)<R\) and apply
\cref{thm-lpn-rate-uniform-semantic-cancellation,cor-combination-capacity-attenuation-closure}.
Polynomial terminal-table cardinality implies \(k=O(\log n)=o(n)\).  Complete
cost and provenance follow from
\cref{thm-charged-noisy-combination-transport,lem-description-length-charged-resource,thm-paper4-dependency-principle}.
\end{proof}

\begin{figure}[tbp]
\centering
\resizebox{.98\linewidth}{!}{%
\begin{tikzpicture}[fw diagram,node distance=8mm and 9mm]
  \node[fw source,text width=2.3cm] (rows) {LPN rows\\
    \((A,b),m,\eta\)};
  \node[fw provenance,text width=2.55cm,right=of rows] (lambda)
    {charged combinations\\\(\lambda\)};
  \node[fw use,text width=2.45cm,above right=7mm and 10mm of lambda] (semantic)
    {remaining secret support\\\(\supp(A^{\mathsf T}\lambda)\)};
  \node[fw residual,text width=2.45cm,below right=7mm and 10mm of lambda] (noise)
    {noise support\\\(|\lambda|\)};
  \node[fw box,text width=2.55cm,right=20mm of lambda] (trade)
    {entropy--attenuation\\tradeoff};
  \node[fw terminal,text width=2.7cm,right=of trade] (bound)
    {rate-uniform bound\\at one ceiling};
  \draw[fw arrow] (rows)--(lambda);
  \draw[fw arrow] (lambda)--(semantic);
  \draw[fw arrow] (lambda)--(noise);
  \draw[fw arrow] (semantic)--(trade);
  \draw[fw arrow] (noise)--(trade);
  \draw[fw arrow] (trade)--(bound);
\end{tikzpicture}%
}
\caption{LPN semantic cancellation and noise provenance.  A combination may
cancel secret coordinates, but a short-noise-provenance combination reaching
a small terminal support is exponentially unlikely.  Outside that event,
the exact factor \((1-2\eta)^{|\lambda|}\) suppresses every polynomially
materialized located block, including in the high-sample region.}
\label{fig-lpn-semantic-cancellation-tradeoff}
\end{figure}

\begin{theorem}[Exact constant-noise spectral perimeter]
\label{thm-lpn-exact-constant-noise-spectral-perimeter}

At fixed \(0<\eta<1/2\), fixed rate, and a polynomial saturated ceiling,
let \(\mathfrak M_{n,m,\eta}\) denote the declared LPN expectation
functional.  The prospective spectral sector has the following exhaustive
routing.  Its correlation estimates are family-average statements for a
fixed nonzero target character under the unrestricted random-\(A\) law.

\begin{enumerate}[label=\textup{(\roman*)}]
\item The family-average correlation of every explicitly materialized
      located block is exponentially suppressed by
      \cref{thm-lpn-constant-noise-materialized-spectral-closure}.  When the
      block is produced by semantic cancellation into \(k=o(n)\) terminal
      coordinates, the rate-uniform estimate
      \cref{cor-lpn-high-sample-materialized-combination-closure} supplies the
      same conclusion throughout the high-sample region.
\item The natural residual-parity fiber map of a materialized selector either
      fails the terminal gate or is identity-equivalent, by
      \cref{lem-lpn-residual-parity-spectral-fibers,eq-lpn-spectral-terminal-rank-test}.
      A refinement using the
      verifier or another represented label enters exact \(\Abs\) only when
      the complete quotient gate of
      \cref{prop-spectral-selector-complete-abs-gate} holds.
\item An explicit expansion processing more than polynomially many terms is
      outside the ceiling by
      \cref{lem-description-length-charged-resource,thm-lpn-no-independent-diffuse-spectral-source}.
\item A compact evaluator may have exponentially many formal Walsh
      coefficients while retaining polynomial description and evaluation
      cost.  It is one complete represented readout, not an explicitly
      materialized coefficient list.  Its represented circuit,
      factorization, recurrence, or other compression rule is charged at its
      actual construction, evaluation, memory, precision, and retained-state
      cost.  Every coefficient locator, selected aggregation, comparison, and
      prospective transition is an ancestor of the same complete record.
      Its prospective use is assigned to an
      exact quotient, a compensated quotient, an identity-supported ambient
      \(\Geom/\Caus\) record, or a terminal readout by
      \cref{thm-complete-spectral-aggregation-charging,thm-lpn-no-independent-diffuse-spectral-source,rem-lpn-compact-factorization-no-spectral-locator}.
      Its complete bounded-readout character correlation obeys
      \cref{thm-lpn-complete-readout-coset-energy-envelope,eq-lpn-complete-readout-coset-energy-average}, which is
      exponentially small in the parameter region
      \cref{eq-lpn-coset-energy-decay-region}.
\end{enumerate}

For the final source-mode reduction, restrict the declared family functional
to the verifier, full-rank, and random-code-rigidity event of
\cref{def-lpn-named-events-restricted-functionals}, and denote the resulting
subprobability functional by \(\mathfrak M^*_{n,m,\eta}\).  On that event,
after the
exact and compensated quotient-supported LPN modes are removed by
\cref{cor-lpn-abs-source-exhaustion,cor-lpn-quotient-supported-geom-closure},
the only spectral form not covered by the exponential located-block estimate
is the compact, dynamically qualified, identity-supported residual-ancestry
ambient coordinate
\(G_{X,b,\mathfrak M^*,n}^{\mathrm{src,sat}}(B)\).  Its complete numerical
bound is exactly the premise isolated in
\cref{thm-lpn-residual-ancestry-attention-bound}; spectral location and
provenance classification alone do not change its value.  Outside
\cref{eq-lpn-coset-energy-decay-region}, the general coset-energy envelope
may equal the unit bound and therefore supplies no replacement estimate for
that sequential ambient coordinate.

\end{theorem}

\begin{proof}
Items~\textup{(i)} and~\textup{(ii)} are
\cref{thm-lpn-constant-noise-materialized-spectral-closure,cor-lpn-high-sample-materialized-combination-closure,lem-lpn-residual-parity-spectral-fibers,prop-spectral-selector-complete-abs-gate}.
For an explicit expansion, every processed term is an evaluated derivation
vertex, so the evaluation and output ledgers dominate the expansion length;
this is item~\textup{(v)} of
\cref{thm-lpn-no-independent-diffuse-spectral-source} and proves
item~\textup{(iii)}.  Evaluator normalization in
\cref{def-budgeted-derivability-closure} charges a compact factorization by
its represented derivation rather than by the cardinality of its formal
expansion.  The support and temporal routing in
\cref{thm-complete-spectral-aggregation-charging,thm-lpn-no-independent-diffuse-spectral-source}
therefore prove
item~\textup{(iv)}.  Removing the quotient-supported modes leaves the ambient
branch by
\cref{thm-lpn-complete-geom-structure-exhaustion}; its residual-independent
subbranch is already bounded by
\cref{lem-lpn-syndrome-free-ambient-geom-bound}.  The displayed
residual-ancestry coordinate is the remaining subbranch in
\cref{def-lpn-ambient-geom-ancestry-branches}.
\end{proof}

\begin{figure}[tbp]
\centering
\resizebox{\linewidth}{!}{%
\begin{tikzpicture}[fw diagram,node distance=7mm and 8mm]
  \node[fw source,text width=2.6cm] (pub) {public LPN record\\\((A,b,\eta)\)};
  \node[fw provenance,text width=2.7cm,right=of pub] (spectral)
    {represented spectral\\use};
  \node[fw use,text width=2.8cm,above right=of spectral] (list)
    {materialized list\\\(L_B\le n^C\)};
  \node[fw terminal,text width=3.25cm,font=\small,right=of list] (decay)
    {\(n^{C/2}(1-2\eta)^{\delta m}\)\\plus short-coset event};
  \node[fw provenance,text width=2.8cm,right=of spectral] (fibers)
    {represented fibers\\and dynamics};
  \node[fw source,text width=2.8cm,right=of fibers] (quot)
    {exact or compensated\\quotient gate};
  \node[fw geom,text width=2.8cm,below right=of spectral] (compact)
    {compact non-lumpable\\readout};
  \node[fw geom,text width=3.0cm,right=of compact] (ambient)
    {identity-supported\\ambient coordinate};
  \node[fw residual,text width=2.8cm,below=of compact] (outside)
    {explicit superpolynomial\\expansion: outside ceiling};
  \node[fw terminal,text width=3.0cm,below=of ambient] (terminal)
    {terminal readout:\\no prospective activity};
  \draw[fw arrow] (pub)--(spectral);
  \draw[fw arrow] (spectral)--(list);
  \draw[fw arrow] (list)--(decay);
  \draw[fw arrow] (spectral)--(fibers);
  \draw[fw arrow] (fibers)--(quot);
  \draw[fw arrow] (spectral)--(compact);
  \draw[fw arrow] (compact)--(ambient);
  \draw[fw arrow] (spectral)--(outside);
  \draw[fw arrow] (compact)--(terminal);
\end{tikzpicture}%
}
\caption{Exhaustive constant-noise spectral routing.  Explicitly located blocks obey
the exponential attenuation theorem.  Fiber compression reaches
\(\Abs\) only through the complete exact or compensated quotient gate.
Compact non-lumpable evaluators remain complete ambient records at their
actual cost or terminate without prospective activity; an explicitly
processed superpolynomial expansion is outside the ceiling.  The number of
terms in a formal Walsh expansion is not a separate source coordinate.}
\label{fig-lpn-constant-noise-spectral-routing}
\end{figure}

\section{Candidate Recovery and Prospective Selector Accountability}
\label{sec-lpn-candidate-recovery-selector-accountability}

\begin{corollary}[Noise-dependent accountability of LPN candidate recovery]
\label{cor-lpn-candidate-recovery-noise-accountability}

Use the simultaneous event of
\cref{thm-lpn-saturated-statistical-alignment}, and let a represented
randomized attack output \(\widehat s=\widehat s(A,b,\omega)\).  After the
represented parity evaluator has been included at its class-preserving cost,
put

\[
p_S
:=
\Pr_{\omega}\{\widehat s=s\mid A,b\},
\qquad
\rho:=1-2\eta,
\qquad
\Gamma_{\mathrm{LPN}}
:=
\Gamma_{\mathrm{LPN},n,B,m,\delta}^{\mathrm{sat,stat}}.
\]

On the event \(s\ne0\),

\begin{equation}
\label{eq-lpn-recovery-probability-projection-bound}
p_S
\le
\sqrt{M_{\mathrm{LPN},n}^{\mathrm{sat,stat}}(B,m)}.
\end{equation}

Character orthogonality also gives the exact execution-law identity

\begin{equation}
\label{eq-lpn-execution-law-recovery-alignment}
\mathbb E_S p_S
=
\mathbb E_{S,\omega}
\left|c_s(f_{\widehat s})\right|^2.
\end{equation}

Consequently, after the candidate parity evaluator and decoder have been
included in the complete cost,

\begin{align}
\label{eq-lpn-candidate-recovery-family-bound}
\mathfrak M_{n,m,\eta}\!\left(
 I\longmapsto\Pr_{S,\omega}\{\widehat s=s\}
\right)
&\le
\operatorname{Align}_{\mathfrak M,n,\eta}^{\mathrm{sat,stat}}(B,m),\\
\label{eq-lpn-learned-projection-envelope-evaluation}
\operatorname{Align}_{\mathfrak M,n,\eta}^{\mathrm{sat,stat}}(B,m)
&\le
\operatorname{GenProj}_{\mathfrak M,n,\eta}^{\mathrm{sat,stat}}(B,m)
\le
\mathfrak M_{n,m,\eta}\!\left(
 I\longmapsto M_{\mathrm{LPN},n}^{\mathrm{sat,stat}}(B,m)
\right).
\end{align}

The empirical noisy alignment of the candidate parity satisfies the
two-sided bound

\begin{equation}
\label{eq-lpn-candidate-empirical-alignment-sandwich}
\left(\rho p_S-2\Gamma_{\mathrm{LPN}}\right)_+
\le
\mathbb E_{\omega}
\left|\widehat c_{\eta,S}(f_{\widehat s})\right|
\le
\rho p_S+2\Gamma_{\mathrm{LPN}}.
\end{equation}

A sample lookup table used to evaluate
\(\widehat c_{\eta,S}(f_{\widehat s})\) remains a charged empirical record.
Its target-bearing projection is classified by the existing LPN ledger and
its orthogonal remainder is \(\Res\), as in
\cref{thm-two-level-learned-readout-provenance}.  If the learned output is
available before and used to prioritize later candidate tests, that use is a
prospective source and its positive nonbaseline advantage is bounded by
\cref{thm-prospective-selector-structural-charge}.  A post-output evaluation
remains accounting.  Neither case adds a memorization or learning channel to
the saturated profile.

\end{corollary}

\begin{proof}

For every represented seed, character orthogonality gives

\[
c_s(f_{\widehat s})
=
\mathbf 1_{\{\widehat s=s\}}
\]

by \eqref{eq-lpn-recovery-is-parity-alignment}.  On \(s\ne0\), apply
\eqref{eq-lpn-uniform-structural-alignment} to the appended parity readout and
average over \(\omega\):

\[
p_S
=
\mathbb E_\omega c_s(f_{\widehat s})
\le
\sqrt{M_{\mathrm{LPN},n}^{\mathrm{sat,stat}}(B,m)}.
\]

This proves \eqref{eq-lpn-recovery-probability-projection-bound}.  The uniform
identity \eqref{eq-lpn-execution-law-recovery-alignment} follows by averaging
the same pointwise equality over \((S,\omega)\).  The decoder margin is exactly
one, so
\cref{thm-population-alignment-generalization-memorization} gives
\eqref{eq-lpn-candidate-recovery-family-bound}; its projection domination
gives \eqref{eq-lpn-learned-projection-envelope-evaluation}.
The uniform
transport estimate \eqref{eq-lpn-uniform-statistical-alignment} gives, for
each represented seed,

\[
\rho c_s(f_{\widehat s})-2\Gamma_{\mathrm{LPN}}
\le
\left|\widehat c_{\eta,S}(f_{\widehat s})\right|
\le
\rho c_s(f_{\widehat s})+2\Gamma_{\mathrm{LPN}}.
\]

Expectation over \(\omega\), followed by the nonnegativity of the middle
term, proves
\eqref{eq-lpn-candidate-empirical-alignment-sandwich}.  The final provenance
statement is the LPN specialization of
\cref{thm-two-level-learned-readout-provenance}; downstream selector charging
is exactly \cref{thm-prospective-selector-structural-charge}.

\end{proof}

\begin{corollary}[Closed budget accounting for LPN learning rules]
\label{cor-lpn-closed-budget-learning-accounting}

Fix one declared saturated learning ceiling \(B\).  Let \(\mathcal D\) be one
family-uniform represented derivation from the LPN presentation.  Its complete
record consists of the feature or spectral-index selection rule, coefficient
and operator construction, optimization transcript, represented randomness,
retained state, output evaluator, and every downstream selector, quotient,
potential, or transition rule that uses the output.  By
\cref{def-budget-saturated-statistical-readout-class}, all these ancestors
belong to the already-defined complete cost
\(\operatorname{Cost}_{\mathrm{sat}}(\mathcal D)\).  The budget accounting
has the following two exhaustive cases.

\begin{enumerate}[label=\textup{(\roman*)}]
\item If
      \(\operatorname{Cost}_{\mathrm{sat}}(\mathcal D)\preceq B\), then the
      output belongs to
      \(\mathcal H_{\mathrm{LPN},n,B,m}^{\mathrm{sat,stat}}\).  A candidate
      output \(\widehat s\) satisfies
      \begin{align}
      \label{eq-lpn-closed-budget-success-projection-chain}
      \mathfrak M_{n,m,\eta}\!\left(
      I\longmapsto\Pr\{\widehat s=s\}
      \right)
      &\le
      \operatorname{Align}_{\mathfrak M,n,\eta}^{\mathrm{sat,stat}}(B,m)\\
      &\le
      \operatorname{GenProj}_{\mathfrak M,n,\eta}^{\mathrm{sat,stat}}(B,m)\\
      &\le
      \mathfrak M_{n,m,\eta}\!\left(
      I\longmapsto
      M_{\mathrm{LPN},n}^{\mathrm{sat,stat}}(B,m)
      \right).
      \end{align}
      Its empirical/population alignment obeys
      \cref{eq-lpn-candidate-empirical-alignment-sandwich}.  An explicitly
      located Walsh block additionally obeys
      \cref{eq-lpn-fixed-ceiling-located-block-bound}; a represented
      target-qualified nonidentity fiber use enters \(\Abs\), and an
      identity-supported use retains the complete ambient \(\Geom\) gate
      product.

\item If
      \(\operatorname{Cost}_{\mathrm{sat}}(\mathcal D)\not\preceq B\), then
      \(\mathcal D\) is absent from the budget-\(B\) statistical readout,
      selector, source, and success profiles.  Enlarging the ceiling to admit
      it charges the same complete record in the enlarged slice.
\end{enumerate}

Thus optimization, implicit aggregation, and spectral location are evaluated
as components of one represented derivation.  The saturated profiles contain
no additional decoder coordinate outside these two cases.

\end{corollary}

\begin{proof}

The membership assertion and complete-cost convention are the LPN
specialization of
\cref{def-budget-saturated-statistical-readout-class,thm-paper4-dependency-principle}.
The success and projection chain is
\cref{eq-lpn-candidate-recovery-family-bound,eq-lpn-learned-projection-envelope-evaluation};
the empirical bound is
\cref{eq-lpn-candidate-empirical-alignment-sandwich}.  Located spectral blocks
and their locator costs are
\cref{thm-lpn-walsh-alignment-charged-location,cor-lpn-fixed-ceiling-located-spectral-suppression}.
The support alternatives are
\cref{thm-affordable-geom-abs-normal-form}.  Budget exclusion in the second
case is the defining restriction of the saturated budget slice, and
\cref{thm-affordable-composition-no-creation} preserves the complete ancestry
when the ceiling is enlarged.

\end{proof}

\begin{theorem}[Prospective LPN learning and selector accountability]
\label{thm-lpn-prospective-learning-selector-accountability}

Fix a represented LPN computation class of budget \(B\), normalized horizon
at most \(H_B\), and homogeneous noise rate \(0\le\eta<1/2\).  Work on the
good-presentation event \(\mathcal G_{\rm ver}(\tau)\) of
\cref{thm-lpn-public-verifier}, on which the public
candidate verifier has accepting set \(\{s\}\), and write
\(\mathfrak M^{\rm ver}_{n,m,\eta}=\mathfrak M^{\mathcal G_{\rm ver}(\tau)}_{n,m,\eta}\)
for the corresponding restricted functional of
\cref{def-lpn-named-events-restricted-functionals}. Include every
public candidate-verifier call in the one-test clock refinement of
\cref{def-represented-pre-hit-selector}, and append one final verifier call to
the computation's output.  At each verifier step use the uniform kernel on
\(\mathbb F_2^n\) as the declared neutral baseline.  Then

\begin{equation}
\label{eq-lpn-prospective-neutral-selector-bound}
\operatorname{Succ}_{\mathfrak M^{\rm ver},n,\eta}(B)
\le
H_B2^{-n}
+H_B
\operatorname{Sel}_{\mathfrak M^{\rm ver},n,\eta}^{\mathrm{sat}}
 \bigl(\Psi(B,n)\bigr).
\end{equation}
This restricted statement, added to the verifier failure probability of
\cref{eq-lpn-verifier-good-event-probability}, bounds the unrestricted
\(\operatorname{Succ}_{\mathfrak M,n,\eta}(B)\) by
\(\mathfrak M_{n,m,\eta}\)-monotonicity, since success is bounded by one.

Put \(\widehat B=\Psi(B,n)\), let
\(\widehat H=H_{\widehat B}\), and let
\(D_{\widehat B}^{\mathrm{sel}}\) be the selector-charge ceiling in
\cref{eq-prospective-selector-charge-budget}.  The selector-to-charge
metatheorem gives

\begin{equation}
\label{eq-lpn-pointwise-selector-five-charge-bound}
\operatorname{Sel}_{I,n}^{\mathrm{sat}}(\widehat B)
\le
e\widehat H
\sum_{c\in\mathfrak J_{\Abs}^5}
\mathcal C_{I,\mathfrak F_5}^{(c)}
   \bigl(D_{\widehat B}^{\mathrm{sel}}\bigr),
\end{equation}

and therefore, after applying the family functional,

\begin{equation}
\label{eq-lpn-selector-five-charge-bound}
\operatorname{Sel}_{\mathfrak M^{\rm ver},n,\eta}^{\mathrm{sat}}(\widehat B)
\le
e\widehat H\,
\mathfrak M^{\rm ver}_{n,m,\eta}\!\left(
 I\longmapsto
 \sum_{c\in\mathfrak J_{\Abs}^5}
 \mathcal C_{I,\mathfrak F_5}^{(c)}
    \bigl(D_{\widehat B}^{\mathrm{sel}}\bigr)
\right).
\end{equation}

Consequently,

\begin{equation}
\label{eq-lpn-recovery-five-charge-bound}
\operatorname{Succ}_{\mathfrak M^{\rm ver},n,\eta}(B)
\le
H_B2^{-n}
+eH_B\widehat H\,
\mathfrak M^{\rm ver}_{n,m,\eta}\!\left(
 I\longmapsto
 \sum_{c\in\mathfrak J_{\Abs}^5}
 \mathcal C_{I,\mathfrak F_5}^{(c)}
    \bigl(D_{\widehat B}^{\mathrm{sel}}\bigr)
\right).
\end{equation}

For \(0\le\eta_1\le\eta_2<1/2\), the same LPN specialization of flip
degradation satisfies

\begin{align}
\label{eq-lpn-noise-degradation-execution-alignment}
\operatorname{Align}_{\mathfrak M,n,\eta_2}^{\mathrm{sat,stat}}(B,m)
&\le
\operatorname{Align}_{\mathfrak M,n,\eta_1}^{\mathrm{sat,stat}}
 \bigl(\psi_{\mathrm{flip}}(B,n,m),m\bigr),\\
\label{eq-lpn-noise-degradation-generalizing-projection}
\operatorname{GenProj}_{\mathfrak M,n,\eta_2}^{\mathrm{sat,stat}}(B,m)
&\le
\operatorname{GenProj}_{\mathfrak M,n,\eta_1}^{\mathrm{sat,stat}}
 \bigl(\psi_{\mathrm{flip}}(B,n,m),m\bigr).
\end{align}

After the parity evaluator is included in the scalarized computation cost,
the LPN decoder margin is one, and hence

\begin{equation}
\label{eq-lpn-alignment-recovery-cost-comparison}
B_{\mathrm{align}}^*(n,m,\eta;\alpha)
\le
B_{\mathrm{rec}}^*(n,m,\eta;\alpha).
\end{equation}

All quantities in these displays are specializations of the existing
saturated learning, selector, and five-family charge profiles.  They add no
LPN-specific source channel.

\end{theorem}

\begin{proof}

On \(\mathcal G_{\rm ver}(\tau)\), the appended verifier ensures that recovery
implies target contact within the refined horizon.  At every verifier step,
the uniform baseline assigns mass \(2^{-n}\) to the unique target \(s\); its
survival-weighted neutral contact is therefore at most \(2^{-n}\).  Summing at
most \(H_B\) such terms and applying
\cref{thm-noise-indexed-learning-accountability} to the restricted functional
\(\mathfrak M^{\rm ver}_{n,m,\eta}\) proves
\eqref{eq-lpn-prospective-neutral-selector-bound}.  The selector at every
step is represented before its corresponding verifier call and retains the
complete computation prefix, so it is temporally admissible by
\cref{def-pre-hit-temporally-admissible-source}.  Since
\(0\le\operatorname{Succ}_{\mathfrak M,n,\eta}(B)-
\operatorname{Succ}_{\mathfrak M^{\rm ver},n,\eta}(B)
\le\Pr(\mathcal G_{\rm ver}(\tau)^c)\) by
\cref{def-lpn-named-events-restricted-functionals} and success takes values
in \([0,1]\), adding
\eqref{eq-lpn-verifier-good-event-probability} to
\eqref{eq-lpn-prospective-neutral-selector-bound} bounds the unrestricted
success profile.

Apply \cref{thm-prospective-selector-structural-charge} pointwise at budget
\(\widehat B\) to obtain
\eqref{eq-lpn-pointwise-selector-five-charge-bound}, then apply the monotone
positively homogeneous expectation functional \(\mathfrak M^{\rm ver}_{n,m,\eta}\).
This proves \eqref{eq-lpn-selector-five-charge-bound}.  Substitution in
\eqref{eq-lpn-prospective-neutral-selector-bound} gives
\eqref{eq-lpn-recovery-five-charge-bound}.

Equations
\eqref{eq-lpn-noise-degradation-execution-alignment} and
\eqref{eq-lpn-noise-degradation-generalizing-projection} are
\cref{eq-noise-degradation-average-alignment,eq-noise-degradation-generalizing-projection}
specialized to the LPN flip channel.  Finally,
\eqref{eq-lpn-execution-law-recovery-alignment} verifies, with equality, the
margin-one learner--decoder realization required by
\cref{cor-noise-indexed-learning-threshold-criterion}; that corollary gives
\eqref{eq-lpn-alignment-recovery-cost-comparison}.

\end{proof}

The statistical scale in
\eqref{eq-lpn-uniform-statistical-alignment} is explicit: an alignment
\(\varepsilon\) can exceed the uniform selection radius only when

\[
(1-2\eta)\varepsilon
>
2\Gamma_{\mathrm{LPN},n,B,m,\delta}^{\mathrm{sat,stat}}.
\]

This comparison measures sample resolution.  Membership in the saturated
readout class separately enforces the complete computational cost of finding
and representing the aligned parity.

\section{Public Verification of the LPN Target}
\label{sec-lpn-public-verification}
\label{the-public-residual-landscape}
\label{public-verification-of-the-lpn-target}

The binary entropy \(H_2\), Bernoulli divergence
\(D_{\mathrm{Ber}}\), and the increasing inverse \(H_2^{-1}\) are defined in
\cref{eq-binary-entropy-definition,eq-bernoulli-divergence-definition}.

The public residual predicate at threshold \(\tau\) is

\[
V_I(x)=\mathbf 1\{|Ax+b|_1\le \tau m\}.
\]

The following task-level lemma identifies its accepting fiber with the
target of \hyperref[def-lpn-task-object]{the LPN inversion task}.

\begin{lemma}[Public verification of LPN recovery]
\label{lem-lpn-public-verifier}
\label{thm-lpn-public-verifier}

Fix \(0<\eta<1/2\), choose \(0<\delta<1/2-\eta\), and put
\(\tau=\eta+\delta\). If

\[
m\bigl(1-H_2(\tau)\bigr)
\ge n+\omega(\log n),
\]

then, with probability \(1-\operatorname{negl}(n)\) over the standard
random binary LPN presentation,

\[
V_I^{-1}(1)=\{s\}=T_I.
\]

More precisely,

\begin{equation}
\label{eq-lpn-verifier-good-event-probability}
\Pr\{V_I^{-1}(1)=\{s\}\}
\ge
1-\exp[-mD_{\mathrm{Ber}}(\tau\Vert\eta)]
-(2^n-1)2^{-m(1-H_2(\tau))}.
\end{equation}

\end{lemma}

\begin{proof}\phantomsection\label{proof-lpn-public-verifier}

For \(x\ne s\), write \(h=x+s\ne0\). Then

\[
Ax+b=A(x+s)+E=Ah+E.
\]

For the standard uniform row model, \(Ah\) is uniform on
\(\mathbb F_2^m\) and independent of \(E\). Hence \(Ah+E\) is uniform.
The Hamming-ball estimate gives
\citep{Gilbert1952,Varshamov1957}

\[
\Pr\{|Ah+E|_1\le\tau m\}
\le
2^{-m}\sum_{j=0}^{\lfloor\tau m\rfloor}\binom mj
\le 2^{-m(1-H_2(\tau))}.
\]

There are at most \(2^n-1\) nonzero choices of \(h\). A union bound
yields

\[
\Pr\{\exists x\ne s: |Ax+b|_1\le\tau m\}
\le
(2^n-1)2^{-m(1-H_2(\tau))}
\le 2^{n-m(1-H_2(\tau))}
=2^{-\omega(\log n)}.
\]

The displayed probability is negligible. Since \(\tau>\eta\), the
Chernoff bound gives
\(\Pr\{|E|_1>\tau m\}\le
\exp[-mD_{\mathrm{Ber}}(\tau\Vert\eta)]\). Therefore \(V_I(s)=1\)
and \(V_I(x)=0\) for every \(x\ne s\), outside a negligible event.
Adding the two failure probabilities proves
\eqref{eq-lpn-verifier-good-event-probability}.

\end{proof}

\begin{definition}[Named LPN presentation events and restricted functionals]
\label{def-lpn-named-events-restricted-functionals}
For a verifier threshold \(\tau\), write
\[
\begin{aligned}
\mathcal G_{\rm ver}(\tau)
  &:=\{V_I^{-1}(1)=\{s\}\},\\
\mathcal G_{\rm rig}
  &:=\{\operatorname{rank}(A)=n\}
    \cap\{\text{the random-code rigidity conclusion holds}\}.
\end{aligned}
\]
and \(\mathcal G_*(\tau)=\mathcal G_{\rm ver}(\tau)\cap
\mathcal G_{\rm rig}\).  The residual-flatness event is
\(\mathcal F_\delta\) of
\cref{eq-lpn-simultaneous-residual-flatness-event}; when all three are
needed, write \(\mathcal G_{*,\delta}(\tau)=
\mathcal G_*(\tau)\cap\mathcal F_\delta\).

For an event \(G\) and an integrable family quantity \(Z\), the
event-restricted functional and conditional functional are, respectively,
\begin{equation}
\label{eq-lpn-event-restricted-functionals}
\mathfrak M^{G}_{n,m,\eta}(Z)
=\mathbb E[Z\mathbf1_G],
\qquad
\overline{\mathfrak M}^{G}_{n,m,\eta}(Z)
=\mathbb E[Z\mid G]
\quad(\Pr(G)>0).
\end{equation}
Every unqualified phrase ``restricted to \(G\)'' below means the first,
subprobability functional.  Conditional expectations always carry the bar.
Consequently, for \(0\le Z\le1\),
\(\mathfrak M(Z)\le\mathfrak M^G(Z)+\Pr(G^c)\); no conditional
normalization is silently used.
\end{definition}

\Cref{fig-lpn-public-verifier} displays the two candidate fibers used in
\cref{thm-lpn-public-verifier}.

\begin{figure}[tbp]
\centering
\begin{tikzpicture}[fw diagram]
  \node[fw source,minimum width=3.1cm] (cand) at (0,2.5)
    {candidate \(x\in\mathbb F_2^n\)};
  \node[fw geom,minimum width=2.6cm] (secret) at (-3.6,1.0)
    {\(x=s\)};
  \node[fw residual,minimum width=2.6cm] (other) at (3.6,1.0)
    {\(x\ne s\)};
  \node[fw geom,minimum width=3.1cm] (enoise) at (-3.6,-0.5)
    {\(r_I(x)=E\)};
  \node[fw residual,minimum width=3.1cm] (uniform) at (3.6,-0.5)
    {\(r_I(x)=Ah+E\) uniform};
  \node[fw terminal,minimum width=3.2cm] (accept) at (-3.6,-2.0)
    {inside Hamming ball\\accept};
  \node[fw box,minimum width=3.2cm] (reject) at (3.6,-2.0)
    {outside Hamming ball\\reject w.h.p.};
  \node[fw use,minimum width=3.7cm] (fiber) at (0,-3.5)
    {\(V_I^{-1}(1)=\{s\}\)};
  \draw[fw arrow] (cand) -- (secret);
  \draw[fw arrow] (cand) -- (other);
  \draw[fw arrow] (secret) -- (enoise);
  \draw[fw arrow] (other) -- (uniform);
  \draw[fw arrow] (enoise) -- (accept);
  \draw[fw arrow] (uniform) -- (reject);
  \draw[fw arrow] (accept) -- (fiber);
  \draw[fw arrow] (reject) -- (fiber);
\end{tikzpicture}
\caption{For the public residual verifier, the true secret produces the
noise vector inside the threshold ball, whereas each false candidate produces
a uniform residual and is rejected with high probability in the sampling
regime of \cref{thm-lpn-public-verifier}.}
\label{fig-lpn-public-verifier}
\end{figure}

\section{Quantitative Noise--Resource Phase Diagram}
\label{sec-lpn-noise-resource-phase}

This section compares the same LPN presentation template across its noise coordinate.
Each value of \(\eta\) defines one presented family as in
\cref{def-lpn-task-object}; varying \(\eta\) does not enlarge its public
constructor grammar.  The comparison uses the single scalarization fixed in
\cref{def-degree-invariant} and, for the polynomial slice, the scalar ceiling
\(s_C(n)=n^C\) of the declared scalar-sublevel resource budget.  Thus its
admissible computation class is exactly
\(\{\mathcal A:\kappa(\operatorname{Cost}(\mathcal A))\le s_C(n)\}\).
The three thresholds are the represented-cost crossing of
\cref{def-noise-indexed-saturated-recovery-cost}, the statistical-resolution
condition of \cref{thm-three-noise-threshold-comparison}, and the information
carried by the binary flip channel.

\subsection{Exact finite-budget error enumeration}
\label{subsec-lpn-error-enumeration-threshold}

Use the task-independent Hamming-ball count \(N_m(k)\) from
\cref{eq-binary-hamming-ball-count,lem-binary-hamming-ball-cost-bounds}.  For
the LPN noise law, define its exact binomial mass by

\begin{equation}
\label{eq-lpn-error-ball-mass}
F_{m,\eta}(k)
=\sum_{j=0}^k\binom mj\eta^j(1-\eta)^{m-j}.
\end{equation}

Set \(F_{m,\eta}(-1)=0\).  This quantity is exactly
\(\Pr\{\operatorname{Bin}(m,\eta)\le k\}\).

Fix a verifier radius \(0<\tau<1/2\).  Let
\(c_{\mathrm{prep}}(n,m)\) be the scalarized represented cost of row-reducing
the public matrix and initializing the enumerator.  Let
\(c_{\mathrm{cand}}(n,m)\) be the scalarized represented cost of constructing
one error vector, solving the corresponding reduced linear system, and
applying the public verifier of \cref{thm-lpn-public-verifier}.  Both are
polynomial because they use only the public binary linear-algebra and Hamming
readouts of \cref{def-lpn-task-object}.  Define

\begin{equation}
\label{eq-lpn-error-enumerator-cost}
C_{\mathrm{EE}}(n,m,k)
=c_{\mathrm{prep}}(n,m)
+c_{\mathrm{cand}}(n,m)N_m(k)
\end{equation}

and the exact affordable radius

\begin{equation}
\label{eq-lpn-affordable-error-radius}
k_B(n,m;\tau)
=
\max\left(
\{-1\}\cup
\left\{
k\in\mathbb Z:\ 0\le k\le\lfloor\tau m\rfloor,\quad
C_{\mathrm{EE}}(n,m,k)\le B
\right\}
\right).
\end{equation}

\begin{proposition}[Exact-binomial success bound for affordable error
enumeration]
\label{prop-lpn-error-enumeration-success}

Assume \(m\ge n\).  For every \(B\), the scalarized success profile satisfies

\begin{equation}
\label{eq-lpn-error-enumeration-success-lower-bound}
\operatorname{Succ}^{\kappa}_{\mathfrak M_{n,m,\eta},n}(B)
\ge
F_{m,\eta}\bigl(k_B(n,m;\tau)\bigr)
-\varepsilon_{\mathrm{alg}}(n,m,\tau),
\end{equation}
where

\begin{equation}
\label{eq-lpn-error-enumeration-algebraic-failure}
\varepsilon_{\mathrm{alg}}(n,m,\tau)
=2^{n-m}
+(2^n-1)2^{-m(1-H_2(\tau))}.
\end{equation}

When \(k_B(n,m;\tau)\ge0\), a family-uniform represented LPN computation of
scalarized cost at most \(B\) witnesses
\cref{eq-lpn-error-enumeration-success-lower-bound}.

For \(0<\alpha<1-\varepsilon_{\mathrm{alg}}\), the exact certified
easy-side noise boundary is therefore

\begin{equation}
\label{eq-lpn-exact-enumeration-noise-boundary}
\eta_{\mathrm{EE}}(n,m;B,\tau,\alpha)
=
\sup\left\{
0\le\eta<\frac12:
F_{m,\eta}\bigl(k_B(n,m;\tau)\bigr)
\ge\alpha+\varepsilon_{\mathrm{alg}}(n,m,\tau)
\right\}.
\end{equation}

Set \(\eta_{\mathrm{EE}}=0\) when this set is empty.

For every \(\eta\) below this boundary,
\(B_{\mathrm{rec}}^*(n,m,\eta;\alpha)\le B\).

\end{proposition}

\begin{proof}

If \(k_B=-1\), the right-hand side of
\cref{eq-lpn-error-enumeration-success-lower-bound} is nonpositive, so the
profile inequality follows from its nonnegativity.  Suppose
\(k_B\ge0\).  Enumerate the binary vectors \(e'\) of weight at most
\(k_B(n,m;\tau)\).  For each vector, solve

\[
Ax=b+e'
\]

using the stored row reduction and apply \(V_I(x)\).  Return the first accepted
candidate.  The complete represented cost is bounded by
\cref{eq-lpn-error-enumerator-cost}.

The union bound over nonzero \(h\in\mathbb F_2^n\) gives

\[
\Pr\{\operatorname{rank}(A)<n\}
\le(2^n-1)2^{-m}
<2^{n-m}.
\]

On the full-column-rank event, if \(e'=E\), then \(Ax=b+E=As\) has the
unique solution \(x=s\).  If \(|E|\le k_B\), this error vector is enumerated
and \(|E|\le\tau m\), so the true candidate passes the verifier.  The
Hamming-ball argument in the proof of \cref{thm-lpn-public-verifier} gives

\[
\Pr\{\exists x\ne s:V_I(x)=1\}
\le(2^n-1)2^{-m(1-H_2(\tau))}.
\]

Outside these two algebraic failure events, recovery occurs whenever
\(|E|\le k_B\), an event of exact probability \(F_{m,\eta}(k_B)\).
Subtracting the two failure probabilities proves
\cref{eq-lpn-error-enumeration-success-lower-bound}.  The boundary statement
then follows from
\cref{def-noise-indexed-saturated-recovery-cost}.  Indeed, let
\(U_1,\ldots,U_m\) be independent uniform random variables on \([0,1]\) and
set \(X_i(\eta)=\mathbf 1_{\{U_i\le\eta\}}\).  Then
\(\sum_iX_i(\eta)\sim\operatorname{Bin}(m,\eta)\), and this sum is
nondecreasing in \(\eta\) for every realization.  Consequently
\(F_{m,\eta}(k)\) is nonincreasing in \(\eta\).  The set in
\cref{eq-lpn-exact-enumeration-noise-boundary} is therefore an initial
interval, and every \(\eta<\eta_{\mathrm{EE}}\) satisfies the displayed
success bound at level \(\alpha\).

\end{proof}

\Cref{fig-lpn-finite-budget-binomial-transition} displays one exact finite
instance of \cref{eq-lpn-error-enumeration-success-lower-bound} before the
separately bounded algebraic failure term is subtracted.

\begin{figure}[tbp]
\centering
\begin{tikzpicture}
\begin{axis}[fw axis,width=.86\linewidth,height=5.8cm,
  xmin=0,xmax=.05,ymin=0,ymax=1.02,
  xlabel={noise rate $\eta$},
  ylabel={$p_{\rm enum}=\Pr\{\operatorname{Bin}(256,\eta)\le4\}$},
  xtick={0,.01,.02,.03,.04,.05},ytick={0,.1,.25,.5,.75,.9,1}]
  \path[fill=fwGray!70] (axis cs:.009531622,0)
    rectangle (axis cs:.030981875,1);
  \addplot[draw=fwBlue,line width=1.4pt,domain=0:.05,samples=180]
    {(1-x)^256
      +256*x*(1-x)^255
      +32640*x^2*(1-x)^254
      +2763520*x^3*(1-x)^253
      +174792640*x^4*(1-x)^252};
  \draw[draw=fwAmber,dashed] (axis cs:0,.5)--(axis cs:.018221849,.5);
  \draw[draw=fwViolet,dashed]
    (axis cs:.018221849,0)--(axis cs:.018221849,.5);
  \fill[fwViolet] (axis cs:.018221849,.5) circle[radius=1.5pt];
  \node[font=\scriptsize,anchor=south west,fwViolet]
    at (axis cs:.018221849,.5) {$\eta_{1/2}=.018221849$};
  \node[font=\scriptsize,align=center,fwInk]
    at (axis cs:.02025675,.93) {$90\%$--$10\%$\\transition band};
\end{axis}
\end{tikzpicture}
\caption{Exact binomial success transition for enumeration through error
weight \(k_B=4\) with \(m=256\) samples.  The curve is
\(\sum_{j=0}^{4}\binom{256}{j}\eta^j(1-\eta)^{256-j}\).  Its \(90\%\),
\(50\%\), and \(10\%\) crossings are \(0.009531622\), \(0.018221849\),
and \(0.030981875\), respectively.  The general budget-dependent curve is
\cref{eq-lpn-error-enumeration-success-lower-bound} with the exact affordable
radius from \cref{eq-lpn-affordable-error-radius}.}
\label{fig-lpn-finite-budget-binomial-transition}
\end{figure}

\begin{corollary}[Polynomial enumeration window]
\label{cor-lpn-polynomial-enumeration-window}

Suppose \(m=\Theta(n^a)\) for some \(a>0\),
\(c_{\mathrm{cand}}(n,m)=\Theta(n^d)\), and
\(c_{\mathrm{prep}}(n,m)=O(n^{d_0})\) for some \(d<C\) and \(d_0<C\).
Under the fixed scalar ceiling \(s_C(n)=n^C\),
assume also

\[
m-n=\omega(\log n),
\qquad
m(1-H_2(\tau))-n=\omega(\log n),
\]

so that \(\varepsilon_{\mathrm{alg}}(n,m,\tau)\) is negligible.  Then,
away from the equality case \(d+ak=C\), the affordable radius is the largest
fixed integer satisfying

\begin{equation}
\label{eq-lpn-polynomial-affordable-weight}
\max\{d_0,d+ak\}<C.
\end{equation}

For a fixed affordable radius \(k_C\) and fixed \(0<\alpha<1\), let
\(\lambda_{k_C,\alpha}\) be the unique nonnegative solution of

\begin{equation}
\label{eq-lpn-poisson-threshold-equation}
e^{-\lambda}
\sum_{j=0}^{k_C}\frac{\lambda^j}{j!}
=\alpha.
\end{equation}

Then, with negligible algebraic failure,

\begin{equation}
\label{eq-lpn-constant-success-enumeration-scale}
\eta_{\mathrm{EE}}(n,m;s_C,\tau,\alpha)
=\frac{\lambda_{k_C,\alpha}}m+o(m^{-1}).
\end{equation}

For the nonnegligible-success criterion, if
\(\eta m=c\log n\) with fixed \(c>0\) and \(m=\Theta(n^a)\), even the
weight-zero member of the same computation succeeds with probability

\begin{equation}
\label{eq-lpn-zero-error-non-negligible-window}
(1-\eta)^m=n^{-c+o(1)}.
\end{equation}

If \(k_C\) is fixed and \(\eta m=\omega(\log n)\), then
\(F_{m,\eta}(k_C)=n^{-\omega(1)}\).  Thus bounded-weight enumeration has
constant-success scale \(m^{-1}\) and nonnegligible-success scale
\((\log n)/m\).

\end{corollary}

\begin{proof}

For fixed \(k\), combine
\cref{eq-fixed-weight-hamming-ball-asymptotic,eq-lpn-error-enumerator-cost}.
The candidate term has order \(n^{d+ak}\), giving
\cref{eq-lpn-polynomial-affordable-weight}; constants decide only the equality
case.  If \(\eta=\lambda/m\), the binomial law converges to
\(\operatorname{Poisson}(\lambda)\), which proves
\cref{eq-lpn-constant-success-enumeration-scale}.  More explicitly, for
fixed \(k_C\) the convergence of

\[
F_{m,\lambda/m}(k_C)
\longrightarrow
e^{-\lambda}\sum_{j=0}^{k_C}\frac{\lambda^j}{j!}
\]

is uniform for \(\lambda\) in compact intervals.  The limiting function is
continuous and strictly decreasing from \(1\) to \(0\), since its derivative
for \(\lambda>0\) is

\[
-e^{-\lambda}\frac{\lambda^{k_C}}{k_C!}.
\]

Its level-\(\alpha\) solution is therefore unique.  The negligible term
\(\varepsilon_{\mathrm{alg}}\) moves the corresponding \(\lambda\)-crossing
by \(o(1)\), and hence moves the noise crossing by \(o(m^{-1})\).

For \(\eta m=c\log n\), under this scaling \(\eta=o(1)\) and

\[
m\log(1-\eta)
=-\eta m+O(\eta^2m)
=-c\log n+o(1),
\]

which proves \cref{eq-lpn-zero-error-non-negligible-window}.  Finally,

\[
F_{m,\eta}(k_C)
\le
(k_C+1)m^{k_C}\exp[-\eta(m-k_C)].
\]

When \(\eta m=\omega(\log n)\) and \(m\) is polynomial, the right-hand
side is \(n^{-\omega(1)}\).

\end{proof}

\section{Asymptotic, Statistical, and Information Thresholds}
\label{sec-lpn-asymptotic-statistical-information-thresholds}
\label{subsec-lpn-three-threshold-formulas}

The preceding polynomial boundary is distinct from the point at which error
enumeration becomes more expensive than exhaustive secret search.

\begin{corollary}[Error-enumeration and secret-search crossover]
\label{cor-lpn-error-secret-crossover}

Let \(n/m\to R\in(0,1)\), fix \(0<\eta<1/2\), and let
\(k_m/m\to\eta\).  Then

\begin{equation}
\label{eq-lpn-typical-error-enumeration-exponent}
\log_2N_m(k_m)=mH_2(\eta)+o(m).
\end{equation}

Consequently typical-error enumeration and exhaustive search through the
\(2^n\) secrets have respective normalized exponents

\begin{equation}
\label{eq-lpn-two-exponential-attack-exponents}
E_{\mathrm{err}}(\eta;R)=\frac{H_2(\eta)}R,
\qquad
E_{\mathrm{secret}}=1.
\end{equation}

They cross at

\begin{equation}
\label{eq-lpn-exponential-attack-switch}
\eta_{\mathrm X}(R)=H_2^{-1}(R).
\end{equation}

This is an exponential-algorithm crossover.  For a continuous visualization
of the polynomial scale, suppose the scalarized per-candidate cost is
\(n^{d+o(1)}\) for some \(d<C\).  Under the scalar ceiling
\(s_C(n)=n^C\), define the leading-entropy count proxy

\begin{equation}
\label{eq-lpn-normalized-polynomial-budget}
\beta_{C,d}(n)=\frac{(C-d)\log_2n}{n}\longrightarrow0,
\end{equation}

and its formal intersection with the constant-rate entropy curve by

\begin{equation}
\label{eq-lpn-entropy-proxy-polynomial-crossing}
\eta_{\mathrm{poly}}^{\mathrm{ent}}(n;C,d,R)
=H_2^{-1}\!\left(R\beta_{C,d}(n)\right).
\end{equation}

The rigorous polynomial-budget intersection is the exact binomial-cost
condition in \cref{eq-lpn-affordable-error-radius}; for fixed error weight it
is \cref{eq-lpn-polynomial-affordable-weight}.  These formulas govern the
vanishing-noise scale plotted by the proxy.

\end{corollary}

\begin{proof}

Apply both sides of \cref{eq-binary-hamming-ball-entropy-bounds} and use
continuity of \(H_2\).  Division by \(n=Rm+o(m)\) proves
\cref{eq-lpn-two-exponential-attack-exponents}; strict monotonicity of \(H_2\)
on \([0,1/2]\) gives \cref{eq-lpn-exponential-attack-switch}.  The final two
displays define the continuous count proxy.  Exact polynomial affordability
was proved in
\cref{lem-binary-hamming-ball-cost-bounds,cor-lpn-polynomial-enumeration-window}.

\end{proof}

\begin{figure}[tbp]
\centering
\begin{tikzpicture}
\begin{axis}[fw axis,width=.9\linewidth,height=6.0cm,
  xmin=0,xmax=.22,ymin=0,ymax=2.35,
  xlabel={noise rate $\eta$},
  ylabel={base-$2$ exponent normalized by $n$},
  legend style={at={(.5,1.03)},anchor=south,legend columns=2,
    font=\scriptsize}]
  \addplot[draw=fwAmber,line width=1.25pt,domain=.00001:.22,samples=180]
    {(-(x*ln(x)+(1-x)*ln(1-x))/ln(2))/(1/3)};
  \addlegendentry{$E_{\rm err}(\eta)=H_2(\eta)/R$}
  \addplot[draw=fwBlue,line width=1.25pt,domain=0:.22] {1};
  \addlegendentry{$E_{\rm secret}=1$}
  \addplot[draw=fwCoral,line width=1.8pt,domain=.00001:.22,samples=180]
    {min(1,(-(x*ln(x)+(1-x)*ln(1-x))/ln(2))/(1/3))};
  \addlegendentry{$E_{\rm two}=\min\{E_{\rm err},E_{\rm secret}\}$}
  \addplot[draw=fwViolet,densely dashed,line width=1pt,domain=0:.22] {.03125};
  \addlegendentry{$\beta_{C,d}=(C-d)\log_2(n)/n=.03125$}
  \draw[draw=fwViolet,dashed]
    (axis cs:.000901267,0)--(axis cs:.000901267,.52);
  \node[font=\scriptsize,fwViolet,anchor=west]
    (polylabel) at (axis cs:.012,.22)
    {$\eta_{\rm poly}^{\rm ent}=.000901267$};
  \draw[draw=fwViolet,-{Latex[length=1.5mm]}]
    (polylabel.west)--(axis cs:.000901267,.03125);
  \draw[draw=fwAmber,dashed]
    (axis cs:.061490470,0)--(axis cs:.061490470,1.02);
  \node[font=\scriptsize,fwAmber,anchor=south west,rotate=90]
    at (axis cs:.061490470,.08) {$\eta_{\rm X}=.061490470$};
  \draw[draw=fwBlue,dashed]
    (axis cs:.173952331,0)--(axis cs:.173952331,2.02);
  \node[font=\scriptsize,fwBlue,anchor=south west,rotate=90]
    at (axis cs:.173952331,.08) {$\eta_{\rm IT}=.173952331$};
\end{axis}
\end{tikzpicture}
\caption{Cost exponents at the illustrative rate \(R=n/m=1/3\).
The attack exponents and proxy are
\cref{eq-lpn-two-exponential-attack-exponents,eq-lpn-exponential-attack-switch,eq-lpn-entropy-proxy-polynomial-crossing}.
Here \(\eta_{\rm X}=H_2^{-1}(R)=0.061490470\), while the information boundary
of \cref{thm-lpn-information-threshold} is
\(\eta_{\rm IT}=H_2^{-1}(1-R)=0.173952331\).  For the displayed finite
values \(n=256\), \(C=2\), and \(d=1\), \(\beta_{C,d}=.03125\) and its
entropy-proxy crossing is \(0.000901267\); the exact finite polynomial crossing uses the
binomial count in \cref{eq-lpn-affordable-error-radius}.}
\label{fig-lpn-exponent-thresholds}
\end{figure}

\begin{corollary}[Explicit LPN Rademacher threshold]
\label{cor-lpn-explicit-rademacher-noise-threshold}

Put

\begin{equation}
\label{eq-lpn-rademacher-numerator}
G_{n,\delta}
=2\sqrt{2n\log2}
+3\sqrt{\frac{\log(2/\delta)}2},
\end{equation}

so that
\(\Gamma_{n,m,\delta}^{\mathrm{par}}=G_{n,\delta}/\sqrt m\) in
\cref{eq-lpn-parity-rademacher-radius}.  An
\(\varepsilon_{\mathrm{opt}}\)-empirical parity minimizer is certified to
equal \(f_s\) whenever

\begin{equation}
\label{eq-lpn-explicit-rademacher-eta-threshold}
\eta
<
\eta_{\mathrm{Rad}}(n,m,\delta,\varepsilon_{\mathrm{opt}})
:=
\frac12-
\frac{2G_{n,\delta}}{\sqrt m}
-\varepsilon_{\mathrm{opt}}.
\end{equation}

Equivalently, when \(1-2\eta>2\varepsilon_{\mathrm{opt}}\), it suffices that

\begin{equation}
\label{eq-lpn-explicit-rademacher-sample-threshold}
m
>
\frac{16G_{n,\delta}^2}
{(1-2\eta-2\varepsilon_{\mathrm{opt}})^2}.
\end{equation}

For the complete budget-\(B\) saturated statistical class, a clean
correlation of magnitude \(\varepsilon>0\) rises above its simultaneous
selection radius only when

\begin{equation}
\label{eq-lpn-saturated-statistical-noise-threshold}
(1-2\eta)\varepsilon
>
2\Gamma_{\mathrm{LPN},n,B,m,\delta}^{\mathrm{sat,stat}},
\qquad\text{equivalently}\qquad
\eta
<
\frac12-
\frac{\Gamma_{\mathrm{LPN},n,B,m,\delta}^{\mathrm{sat,stat}}}
{\varepsilon}.
\end{equation}

This is the statistical-resolution boundary of the existing uniform
generalization theorem.  The represented cost of constructing the empirical
minimizer is still measured by
\cref{eq-noise-indexed-recovery-cost}.

\end{corollary}

\begin{proof}

Substitute
\(\Gamma_{n,m,\delta}^{\mathrm{par}}=G_{n,\delta}/\sqrt m\) into
\cref{eq-lpn-parity-exact-statistical-identification} and solve first for
\(\eta\), then for \(m\).  Rearranging the simultaneous deviation bound
\cref{eq-lpn-uniform-statistical-alignment} at clean alignment
\(\varepsilon\) gives
\cref{eq-lpn-saturated-statistical-noise-threshold}.

\end{proof}

\begin{theorem}[Sharp LPN information threshold away from equality]
\label{thm-lpn-information-threshold}

Under the standard random binary LPN law of
\cref{def-lpn-task-object}, in particular with \(s\) uniform on
\(\mathbb F_2^n\), let
\(P_{\mathrm e}=\Pr\{\widehat s\ne s\}\) for an arbitrary, possibly
randomized decoder of the LPN sample.  Then

\begin{equation}
\label{eq-lpn-fano-finite-converse}
H_2(P_{\mathrm e})
+P_{\mathrm e}\log_2(2^n-1)
\ge
n-m(1-H_2(\eta)).
\end{equation}

In particular, for \(0\le\delta<1-2^{-n}\), the condition
\(P_{\mathrm e}\le\delta\) requires

\begin{equation}
\label{eq-lpn-fano-noise-necessary-condition}
m(1-H_2(\eta))
\ge
n-H_2(\delta)-\delta\log_2(2^n-1).
\end{equation}

For an exponential strong converse, fix \(0<\epsilon<\eta<1/2\) and write

\[
L_\eta=\log_2\frac{1-\eta}{\eta}.
\]

Every decoder satisfies

\begin{equation}
\label{eq-lpn-information-strong-converse}
\Pr\{\widehat s=s\}
\le
2e^{-2m\epsilon^2}
+2^{m-n-m(H_2(\eta)-\epsilon L_\eta)}.
\end{equation}

Conversely, fix \(0<\delta<1-2^{-n}\) and suppose

\[
0<1-\frac{n+\log_2(2/\delta)}m<1.
\]

Define

\begin{equation}
\label{eq-lpn-explicit-verification-radius}
\tau_\delta(n,m)
=H_2^{-1}\!\left(
1-\frac{n+\log_2(2/\delta)}m
\right).
\end{equation}

If

\begin{equation}
\label{eq-lpn-explicit-information-achievability}
0\le\eta
\le
\tau_\delta(n,m)
-\sqrt{\frac{\log(2/\delta)}{2m}},
\end{equation}

then exhaustive candidate scanning with the public verifier recovers \(s\)
with probability at least \(1-\delta\).  The corresponding finite-sample
achievability endpoint is

\begin{equation}
\label{eq-lpn-finite-information-achievability-endpoint}
\eta_{\mathrm{IT}}^{\mathrm{ach}}(n,m,\delta)
=
H_2^{-1}\!\left(1-\frac{n+\log_2(2/\delta)}m\right)
-\sqrt{\frac{\log(2/\delta)}{2m}},
\end{equation}

whenever the inverse argument lies in \((0,1)\) and the displayed endpoint
is nonnegative.  Put

\[
K_{n,\delta}
=n-H_2(\delta)-\delta\log_2(2^n-1).
\]

Whenever \(1-K_{n,\delta}/m\in[0,1]\), the finite-sample converse endpoint is

\begin{equation}
\label{eq-lpn-finite-information-converse-endpoint}
\eta_{\mathrm{IT}}^{\mathrm{con}}(n,m,\delta)
=
H_2^{-1}\!\left(1-\frac{K_{n,\delta}}m\right).
\end{equation}

Every decoder with success probability at least \(1-\delta\) must have
\(\eta\le\eta_{\mathrm{IT}}^{\mathrm{con}}\), whereas exhaustive scanning
has success at least \(1-\delta\) throughout
\(0\le\eta\le\eta_{\mathrm{IT}}^{\mathrm{ach}}\).

Hence, when \(n/m\to R\in(0,1)\), the sharp information threshold away from
equality is

\begin{equation}
\label{eq-lpn-asymptotic-information-threshold}
\eta_{\mathrm{IT}}(R)=H_2^{-1}(1-R).
\end{equation}

Below this value exhaustive recovery succeeds with probability tending to
one; above it every decoder has exponentially small exact-recovery
probability.

\end{theorem}

\begin{proof}

Condition on the public matrix \(A\).  The secret is uniform on \(2^n\)
values and the channel from \(As\) to \(b\) consists of \(m\) independent
binary symmetric channels.  Therefore

\[
I(s;b\mid A)\le m(1-H_2(\eta)).
\]

Fano's entropy inequality gives

\[
H(s\mid A,b)
\le
H_2(P_{\mathrm e})
+P_{\mathrm e}\log_2(2^n-1).
\]

Since \(H(s)=n\), combining these inequalities proves
\cref{eq-lpn-fano-finite-converse}.  The function
\(p\mapsto H_2(p)+p\log_2(2^n-1)\) is increasing on
\([0,1-2^{-n})\), which gives
\cref{eq-lpn-fano-noise-necessary-condition}.
Solving this condition for \(\eta\) on the increasing branch of \(H_2\)
gives \cref{eq-lpn-finite-information-converse-endpoint}.

For the strong converse, let

\[
\mathcal T_\epsilon
=\left\{
e:\left|\frac{|e|}m-\eta\right|\le\epsilon
\right\}.
\]

Hoeffding's inequality gives
\(\Pr\{E\notin\mathcal T_\epsilon\}\le2e^{-2m\epsilon^2}\).  Every
\(e\in\mathcal T_\epsilon\) satisfies

\[
\Pr\{E=e\}
\le
2^{-m(H_2(\eta)-\epsilon L_\eta)}.
\]

For fixed \(A\) and received word \(y\), a randomized decoder assigns weights
to the possible secrets that sum to one.  Reindexing \(y=As+e\), summing over
the \(2^m\) received words, and using the uniform prior on \(s\) bounds the
success contribution of the typical set by

\[
2^{-n}2^m2^{-m(H_2(\eta)-\epsilon L_\eta)}.
\]

Adding the atypical probability proves
\cref{eq-lpn-information-strong-converse}.  If
\(R>1-H_2(\eta)\), choose

\[
0<\epsilon<
\min\left\{\eta,\frac{R-1+H_2(\eta)}{L_\eta}\right\};
\]

both terms then decay exponentially.

For achievability, the definition of \(\tau_\delta\) gives

\[
(2^n-1)2^{-m(1-H_2(\tau_\delta))}\le\frac\delta2.
\]

If \(\eta=0\), the true residual is identically zero and the displayed
false-candidate bound proves the claimed success probability directly.
Assume \(0<\eta<1/2\).

The Bernoulli Pinsker bound
\(D_{\mathrm{Ber}}(\tau_\delta\Vert\eta)
\ge2(\tau_\delta-\eta)^2\) and
\cref{eq-lpn-explicit-information-achievability} give

\[
e^{-mD_{\mathrm{Ber}}(\tau_\delta\Vert\eta)}\le\frac\delta2.
\]

Substitute both estimates into
\cref{eq-lpn-verifier-good-event-probability}.  On the resulting event the
verifier has the unique accepting candidate \(s\), so exhaustive scanning
recovers it.  This proves the endpoint statement in
\cref{eq-lpn-finite-information-achievability-endpoint}.  For any sequence
\(\delta_n\to0\) with \(\log(1/\delta_n)=o(m)\), both finite endpoints
converge to \(H_2^{-1}(1-R)\).  The achievability endpoint gives the lower
side of \cref{eq-lpn-asymptotic-information-threshold}; the strong converse
gives the upper side.

\end{proof}

\begin{corollary}[Rate asymptotics of the exponential and information
thresholds]
\label{cor-lpn-rate-asymptotic-thresholds}

The attack crossover \(\eta_{\mathrm X}(R)=H_2^{-1}(R)\) and information
threshold \(\eta_{\mathrm{IT}}(R)=H_2^{-1}(1-R)\) satisfy

\begin{align}
\label{eq-lpn-low-rate-threshold-asymptotics}
\eta_{\mathrm X}(R)
&\sim \frac{R}{\log_2(1/R)},
&
\eta_{\mathrm{IT}}(R)
&=\frac12-\sqrt{\frac{R\log 2}{2}}+O(R^{3/2})
&& (R\downarrow0),\\
\label{eq-lpn-high-rate-threshold-asymptotics}
\eta_{\mathrm X}(R)
&=\frac12-\sqrt{\frac{(1-R)\log 2}{2}}
  +O((1-R)^{3/2}),
&
\eta_{\mathrm{IT}}(R)
&\sim\frac{1-R}{\log_2(1/(1-R))}
&& (R\uparrow1).
\end{align}

\end{corollary}

\begin{proof}

As \(u\downarrow0\),

\[
H_2(u)=u\log_2(1/u)+\frac{u}{\log 2}+O(u^2),
\]

whose monotone inverse is
\(H_2^{-1}(y)\sim y/\log_2(1/y)\).  As \(x\to0\),

\[
H_2\!\left(\frac12-x\right)
=1-\frac{2x^2}{\log 2}+O(x^4),
\]

so
\(H_2^{-1}(1-y)=1/2-\sqrt{y\log 2/2}+O(y^{3/2})\).
Substitute \(y=R\) and \(y=1-R\).

\end{proof}

\subsection{The saturated computational boundary}
\label{subsec-lpn-saturated-computational-boundary}

For \(0<\vartheta\le1\), write

\begin{equation}
\label{eq-lpn-visibility-threshold-cost}
B^*_{\vartheta}(n;\mathrm{LPN}_\eta)
=
\inf\left\{
s\ge0:
I_{\mathfrak M_{n,m,\eta},n,\kappa}^{\mathrm{sat}}(s)
\ge\vartheta
\right\}.
\end{equation}

Here
\(I_{\mathfrak M_{n,m,\eta},n,\kappa}^{\mathrm{sat}}(s)\) is the family
profile of \cref{def-master-saturated-quotient-geometry-profile} restricted to
records with scalarized cost at most \(s\).  Thus
\cref{eq-lpn-visibility-threshold-cost} is precisely the saturated cost
invariant of \cref{def-degree-invariant} with the family and visibility level
shown explicitly.

\begin{theorem}[Saturated hard-side threshold for LPN]
\label{thm-lpn-three-threshold-phase-comparison}

Fix the declared polynomial scalar-sublevel resource ceiling
\(\mathcal B_C\), whose admissible computation class is

\[
\mathsf{Alg}_{n,\le\mathcal B_C}
=\left\{\mathcal A:
\kappa(\operatorname{Cost}(\mathcal A))\le s_C(n)=n^C
\right\},
\]

and let \(H_C(n)=H_{\mathcal B_C}(n)\) be its normalized polynomial horizon.
Let \(\Psi(\mathcal B_C,n)\) be the class-preserving
representation and first-hit enlargement in
\cref{thm-universal-saturated-profile-accountability}.  For every target
success level \(0<\alpha_n\le1\),

\begin{equation}
\label{eq-lpn-saturated-hard-side-threshold}
\kappa\!\left(\Psi(\mathcal B_C,n)\right)
<
B^*_{\alpha_n/(eH_C(n))}(n;\mathrm{LPN}_\eta)
\quad\Longrightarrow\quad
\operatorname{Succ}_{\mathfrak M_{n,m,\eta},n}(\mathcal B_C)<\alpha_n.
\end{equation}

The three quantitative comparison bounds are as follows.

\begin{enumerate}[label=\textup{(\roman*)}]
\item Under the hypotheses of
\cref{prop-lpn-error-enumeration-success}, for
\(0<\alpha_n<1-\varepsilon_{\mathrm{alg}}\) and
\(0\le\eta<\eta_{\mathrm{EE}}(n,m;s_C,\tau,\alpha_n)\), the represented
error enumerator gives
\(B_{\mathrm{rec}}^*(n,m,\eta;\alpha_n)\le s_C\).  Consequently,

\begin{equation}
\label{eq-lpn-enumeration-computational-threshold-lower-bound}
\eta_{\mathrm{EE}}(n,m;s_C,\tau,\alpha_n)
\le
\eta_{\mathrm{comp}}^{\mathrm{sat}}(n,m;s_C,\alpha_n).
\end{equation}

For fixed affordable radius this includes constant success at
\(\eta=\Theta(m^{-1})\) and nonnegligible success at
\(\eta=O((\log n)/m)\), as quantified in
\cref{cor-lpn-polynomial-enumeration-window}.

\item Exact empirical parity identification is statistically certified below
\(\eta_{\mathrm{Rad}}\) in
\cref{eq-lpn-explicit-rademacher-eta-threshold}.  This condition measures
sample resolution.  If the complete scalarized construction cost of the
empirical minimizer is at most \(s_C\) and
\(\eta_{\mathrm{Rad}}>0\), then

\begin{equation}
\label{eq-lpn-statistical-computational-threshold-lower-bound}
\eta_{\mathrm{Rad}}(n,m,\delta,\varepsilon_{\mathrm{opt}})
\le
\eta_{\mathrm{comp}}^{\mathrm{sat}}(n,m;s_C,1-\delta).
\end{equation}

\item Information permits high-probability recovery below
\(\eta_{\mathrm{IT}}(R)=H_2^{-1}(1-R)\) and forbids it above that value, away
from equality, by \cref{thm-lpn-information-threshold}.  At finite
\((n,m,\delta)\), whenever the converse endpoint is defined,

\begin{equation}
\label{eq-lpn-information-computational-threshold-upper-bound}
\eta_{\mathrm{comp}}^{\mathrm{sat}}(n,m;s_C,1-\delta)
\le
\eta_{\mathrm{IT}}^{\mathrm{con}}(n,m,\delta).
\end{equation}
\end{enumerate}

For a vanishing schedule \(\eta=\eta(n)\), the supremal affordable noise is
\cref{eq-noise-indexed-computational-threshold}; when the ceiling absorbs the
flip postprocessing of \cref{thm-noise-degradation-monotonicity}, it is the
right endpoint of the affordable interval.
\cref{eq-lpn-saturated-hard-side-threshold} is its framework-level hard-side
certificate.  Under the polynomial-enumeration hypotheses in item~\textup{(i)},
the present estimates locate the nonnegligible-success easy side through
\(\eta m=O(\log n)\).

\end{theorem}

\begin{proof}

If the inequality on the left of
\cref{eq-lpn-saturated-hard-side-threshold} holds, the definition in
\cref{eq-lpn-visibility-threshold-cost} gives

\[
m_{\mathfrak M_{n,m,\eta},n}^{\mathrm{sat}}
 \bigl(\Psi(\mathcal B_C,n)\bigr)
\le
I_{\mathfrak M_{n,m,\eta},n}^{\mathrm{sat}}
 \bigl(\Psi(\mathcal B_C,n)\bigr)
\le
I_{\mathfrak M_{n,m,\eta},n,\kappa}^{\mathrm{sat}}
 \bigl(\kappa(\Psi(\mathcal B_C,n))\bigr)
<\frac{\alpha_n}{eH_C(n)}.
\]

Applying \cref{thm-universal-saturated-profile-accountability}, equivalently
\cref{cor-pointwise-algorithm-profile-bound}, now proves the strict success
bound \(\alpha_n\).

For part~\textup{(i)}, the initial-interval conclusion proved in
\cref{prop-lpn-error-enumeration-success} gives, for every
\(0\le\eta<\eta_{\mathrm{EE}}\), one represented computation of scalarized
cost at most \(s_C\) and success at least \(\alpha_n\).  Hence

\[
[0,\eta_{\mathrm{EE}})
\subseteq
\left\{\eta:
\operatorname{Succ}^{\kappa}_{\mathfrak M_{n,m,\eta},n}(s_C)
\ge\alpha_n\right\}.
\]

Taking suprema and using
\cref{eq-noise-indexed-computational-threshold} proves
\cref{eq-lpn-enumeration-computational-threshold-lower-bound}; the two
asymptotic scales are those of
\cref{cor-lpn-polynomial-enumeration-window}.

For part~\textup{(ii)}, every
\(0\le\eta<\eta_{\mathrm{Rad}}\) satisfies the strict identification
inequality in
\cref{eq-lpn-explicit-rademacher-eta-threshold}.  The represented empirical
minimizer belongs to the scalar sublevel \(s_C\) by the stated complete-cost
bound, and it recovers the target with probability at least \(1-\delta\).
Thus the interval \([0,\eta_{\mathrm{Rad}})\) is contained in the defining
set for
\(\eta_{\mathrm{comp}}^{\mathrm{sat}}(n,m;s_C,1-\delta)\), which proves
\cref{eq-lpn-statistical-computational-threshold-lower-bound} after taking
suprema.

For part~\textup{(iii)}, if
\(\eta>\eta_{\mathrm{IT}}^{\mathrm{con}}(n,m,\delta)\), strict monotonicity
of \(H_2\) on \([0,1/2]\) makes
\cref{eq-lpn-fano-noise-necessary-condition} fail by a positive amount.
The monotonicity of

\[
p\longmapsto H_2(p)+p\log_2(2^n-1)
\]

on \([0,1-2^{-n})\) then gives a uniform error bound
strictly above \(\delta\): continuity and strict monotonicity give a number
\(\delta_\eta>\delta\), depending only on \((n,m,\eta,\delta)\), such that
every decoder satisfies \(P_{\mathrm e}\ge\delta_\eta\).  Hence every
decoder has success at most \(1-\delta_\eta<1-\delta\).  Therefore the
defining set for
\(\eta_{\mathrm{comp}}^{\mathrm{sat}}(n,m;s_C,1-\delta)\) is contained in
\([0,\eta_{\mathrm{IT}}^{\mathrm{con}}]\).  Taking its supremum proves
\cref{eq-lpn-information-computational-threshold-upper-bound}.

\end{proof}

\section{Structural Core Reduction for
LPN}\label{sec-lpn-structural-core-reduction}

This section applies the
\hyperref[thm-complete-generic-source-theory]{Part~III master-source
decomposition} and
\hyperref[thm-core-refined-prospective-accountability]{core-refined
first-hit accountability} to the standard binary LPN presentation. The
source and accounting clauses follow from the
\hyperref[thm-channel-exhaustiveness]{Part~II component decomposition}, the
\hyperref[lem-exact-lift-identity]{exact Lift identity}, and the LPN source
bounds proved below.

\subsection{Representation and Geometric Source Bounds}
\label{subsec-lpn-representation-geom-source-bounds}

The first phase identifies the declared polynomial LPN execution class with
its saturated closure and classifies the primitive and coupled geometric
cases.

\begin{lemma}[Representation of the polynomial LPN execution class with
admissible oracle interfaces]\label{lem-lpn-no-creation-final}

The declared family-uniform polynomial execution class on the LPN
presentation, including its presentation-admissible oracle interfaces,
embeds, by the
\hyperref[thm-encoding-adequacy]{encoding-adequacy theorem} and
\hyperref[lem-budgeted-generation-gives-representation]{budgeted
generation lemma}, as a
represented generator on the full lifted observable configuration at
polynomial saturated cost. Hence its target contribution is counted by the
generic saturated profile at the class-preserving enlarged cost.

\end{lemma}

\begin{proof}\phantomsection\label{proof-lpn-no-creation-final}

Apply the
\hyperref[thm-oracle-admissibility-presentation-separation]{Oracle
admissibility and presentation separation theorem}
to every presentation-admissible interface. Replacing each query vertex
by its family-uniform represented evaluator produces a uniform
polynomial computation in the original LPN saturated closure, with
class-preserving accumulated overhead.

Encoding adequacy simulates the input presentation, random transcript,
work state, clock, and output map as observable coordinates. The
transition rule of a uniform polynomial computation is a finite
constructor sequence using only public data and charged precision.
Therefore the induced generator and every terminal readout it uses lie
in the
\hyperref[def-budgeted-derivability-closure]{budgeted derivability
closure} at polynomial budget. Admission of a
future committor, latent noise vector, planted secret, or quotient to the
budgeted algebra is determined by such a represented derivation.

\end{proof}

\begin{lemma}[LPN specialization of the component
decomposition]\label{lem-lpn-component-decomposition}
\label{thm-lpn-five-term-partition-final}

Every represented target-relevant LPN term decomposes into
\(\Geom\), \(\Caus\), \(\Abs\),
\(\Lift\), \(\Bdry\), and residual accounting
\(\Res\). Only \(\Geom\) and \(\Abs\) are
generative sources. \(\Caus\) orients differences of a
represented potential, \(\Lift\) reorganizes an existing source,
and \(\Bdry\) scores or kills mass supplied by a source. Thus
\(\Caus\), \(\Lift\), and \(\Bdry\) charge back
to \(\Geom\) or \(\Abs\).

\end{lemma}

\begin{proof}\phantomsection\label{proof-lpn-five-term-partition-final}

This is the \hyperref[thm-channel-exhaustiveness]{Component
exhaustiveness theorem},
\hyperref[thm-exhaustive-residual-normal-form]{Exhaustive residual normal
form theorem}, and the
\hyperref[prop-channel-ledger]{Roles of the five structural components
and the residual term proposition}
applied to the LPN lifted
configuration. The \hyperref[def-pmp]{positive maximum principle} gives
the generator form. The exposed relation splits local
symmetric and local oriented rates. Quotient-compatible movement is
\(\Abs\). Auxiliary coordinates are valid lifts only when their
projection preserves target fraction. Terminal scores and
residual-weight predicates are boundary readouts. What remains after all
represented source projections is one of the two residual coordinates in
\cref{thm-exhaustive-residual-normal-form}. The decomposition is exhaustive
by the \hyperref[thm-channel-exhaustiveness]{component exhaustiveness
theorem}.  The following source bounds evaluate these branches in the
existing saturated \(\Geom/\Abs\) profile.

\end{proof}

\begin{lemma}[\PartII exact Lift identity specialized to
LPN]\label{lem-lpn-exact-lift-composition}

Let \(\pi:\widetilde X\to X\) be a valid lift with lifted measure
\(\widetilde\mu\) and target \(T\subseteq X\). Validity gives the exact
target-fraction identity

\[
\widetilde\mu(\pi^{-1}T)=\mu(T),
\qquad
\int_{\widetilde X} h\circ\pi\,d\widetilde\mu
=
\int_X h\,d\mu
\]

for every target-relevant base observable \(h\). These identities
compose exactly. Therefore any finite composition of valid lifts,
including recurrence, stacked auxiliary coordinates, and filtration
level-stacking, preserves target fraction exactly and charges to the
represented base source with zero target-fraction defect.

\end{lemma}

\begin{proof}\phantomsection\label{proof-lpn-exact-lift-composition}

This is the \hyperref[lem-exact-lift-identity]{exact Lift identity}
applied to the presented LPN task. Every admissible composition retains
the represented base projection and charges to that base source.

\end{proof}

\subsection{Capacity and arrival transfer under valid lifts}
\label{subsec-lpn-valid-lift-capacity-arrival-transfer}

\begin{proposition}[Native invariant capacity and access under an LPN lift]
\label{prop-lpn-valid-lift-native-capacity-access}

Let \(I=(A,b,\eta)\) be an LPN presentation on
\(X_I=\mathbb F_2^n\), let \(\mu_I\) be uniform, and let
\(T_I=\{s\}\).  Let
\((\Omega,\mathcal G_{\rm obs},\pi,\Lambda)\) be a complete valid-Lift
record, with
\(\widetilde\mu=\Lambda\mu_I\).  Then

\begin{equation}
\label{eq-lpn-valid-lift-target-mass}
\widetilde\mu(\pi^{-1}T_I)=\mu_I(T_I)=2^{-n}.
\end{equation}

Suppose first that \(P_\Omega\) is a represented normalized kernel whose
invariant law is \(\widetilde\mu\), and let \(\Phi_{P_\Omega}\) be its
native symmetric geometry.  Then

\begin{align}
\label{eq-lpn-valid-lift-native-capacity}
\operatorname{Cap}_{\Phi_{P_\Omega}}
 (\pi^{-1}T_I,\Omega\setminus\pi^{-1}T_I)
&=Q_{P_\Omega}(\pi^{-1}T_I,
                 \Omega\setminus\pi^{-1}T_I)
\le 2^{-n},\\
\label{eq-lpn-valid-lift-native-hit}
\Pr_{\widetilde\mu_0}\{\tau_{\pi^{-1}T_I}\le H\}
&\le R_0(H+1)2^{-n},
\qquad
R_0=\left\|\frac{d\widetilde\mu_0}
                         {d\widetilde\mu}\right\|_\infty .
\end{align}

\end{proposition}

\begin{proof}
The exact-Lift identity gives target mass \(2^{-n}\).  The
boundary-flux identity bounds the invariant target boundary by that mass,
and a union bound over the \(H+1\) invariant-time marginals, followed by the
density comparison \(\widetilde\mu_0\le R_0\widetilde\mu\), gives the hitting
bound.
\end{proof}

\begin{theorem}[Schur-complement capacity comparison for an LPN lift]
\label{thm-lpn-valid-lift-schur-capacity}

Under the LPN presentation and valid-Lift hypotheses of
\cref{prop-lpn-valid-lift-native-capacity-access}, let the same lift carry a
complete form-lift comparison record from
\cref{def-controlled-complete-form-lift-record}.  Let \(P\) be a normalized
uniform-invariant base kernel, let \(\Phi_P\) be its native symmetric
geometry, and suppose that on the condenser domain

\[
a_\pi\mathcal E_{\Phi_P}(g,g)
\le \langle g,H_Xg\rangle
\le b_\pi\mathcal E_{\Phi_P}(g,g).
\]

For every represented \(F\subseteq X_I\setminus T_I\),

\begin{align}
\label{eq-lpn-valid-lift-capacity-schur}
(1-\rho_\pi)a_\pi
 \operatorname{Cap}_{\Phi_P}(T_I,F)
&\le
 \operatorname{Cap}_{\Omega}(\pi^{-1}T_I,\pi^{-1}F)
\le
 b_\pi\operatorname{Cap}_{\Phi_P}(T_I,F)\notag\\
&\le b_\pi2^{-n}(1-P(s,s)).
\end{align}

For the opposite plate \(F=X_I\setminus T_I\), this becomes

\begin{equation}
\label{eq-lpn-valid-lift-capacity-complement}
(1-\rho_\pi)a_\pi2^{-n}(1-P(s,s))
\le
\operatorname{Cap}_{\Omega}
 (\pi^{-1}T_I,\pi^{-1}(X_I\setminus T_I))
\le
b_\pi2^{-n}(1-P(s,s)).
\end{equation}

If the lift is form-neutral, \(B_0=0\), and its pullback form is exactly
\(\mathcal E_{\Phi_P}\), then
\(\rho_\pi=0\), \(a_\pi=b_\pi=1\), and

\begin{equation}
\label{eq-lpn-form-neutral-capacity-equality}
\operatorname{Cap}_{\Omega}(\pi^{-1}T_I,\pi^{-1}F)
=\operatorname{Cap}_{\Phi_P}(T_I,F).
\end{equation}

Consequently, for every represented form-lift class \(\mathscr L(B)\)
whose comparison geometry is such a normalized LPN base geometry,

\begin{equation}
\label{eq-lpn-saturated-lift-capacity-bound}
\sup_{R\in\mathscr L(B)}
\operatorname{Cap}_{\Omega,R}
\le
2^{-n}\operatorname{LiftAmp}^{\rm sat}_{\mathscr L}(B),
\end{equation}

with the sharper bound \(2^{-n}\) on the native invariant subclass of
\cref{eq-lpn-valid-lift-native-capacity}.

\end{theorem}

\begin{proof}
Apply the block-form Schur complement theorem
\cref{thm-controlled-lift-capacity-schur} and then the two-sided base-form
comparison.  The singleton boundary-flux identity gives the last upper
bound.  The opposite-plate and form-neutral formulas are the corresponding
specializations, and taking the supremum over the represented class gives
\cref{eq-lpn-saturated-lift-capacity-bound}.
\end{proof}

\begin{theorem}[Clocked arrival transfer for a valid LPN lift]
\label{thm-lpn-clocked-lift-arrival-transfer}
Adopt the presentation, lift, and target notation of
\cref{prop-lpn-valid-lift-native-capacity-access}.

For the finite-horizon comparison, let \(P_t^\Omega,P_t^X\),
\(0\le t<H\), be the stopped or absorbing clocked kernels and let
\(U_t\) be their represented pullbacks.  Put

\[
E_t=P_t^\Omega U_{t+1}-U_tP_t^X.
\]

Write \(p_{T_I,H}^{\Omega}\) for first arrival at
\(\pi^{-1}T_I\) on the lifted carrier and \(p_{T_I,H}^{X}\) for first
arrival at \(T_I\) on the base carrier, under the compared represented
initializations.  Likewise, \(p_{\widehat T_I,H}^{\Omega}\) below denotes
first arrival at \(\pi^{-1}\widehat T_I\).

If the represented initialization and transported target interfaces have
errors \(\varepsilon_{\rm init}\) and \(\varepsilon_{\rm gate}\), then

\begin{equation}
\label{eq-lpn-clocked-lift-arrival-transfer}
\left|p_{T_I,H}^{\Omega}-p_{T_I,H}^{X}\right|
\le
\varepsilon_{\rm init}
+\sum_{t=0}^{H-1}\|E_t\|
+\varepsilon_{\rm gate}.
\end{equation}

Exact clocked intertwining and exact valid-Lift interfaces set the right side
to zero.  Let \(\widehat T_I=V_I^{-1}(1)\) be the public LPN verifier gate at
threshold \(\tau\), with \(\eta<\tau<1/2\).  Pointwise in the presentation,

\begin{equation}
\label{eq-lpn-clocked-lift-public-gate-transfer}
\left|p_{\widehat T_I,H}^{\Omega}-p_{T_I,H}^{X}\right|
\le
\varepsilon_{\rm init}
+\sum_{t=0}^{H-1}\|E_t\|
+\varepsilon_{\rm gate}
+\mathbf1_{\{\widehat T_I\ne T_I\}}.
\end{equation}

Therefore the LPN family functional satisfies

\begin{align}
\label{eq-lpn-clocked-lift-public-gate-family-transfer}
\mathfrak M_{n,m,\eta}\!\left(
 \left|p_{\widehat T_I,H}^{\Omega}-p_{T_I,H}^{X}\right|
\right)
&\le
\mathfrak M_{n,m,\eta}\!\left(
 \varepsilon_{\rm init}
 +\sum_{t=0}^{H-1}\|E_t\|
 +\varepsilon_{\rm gate}
\right)\notag\\
&\quad+
\exp[-mD_{\mathrm{Ber}}(\tau\Vert\eta)]
+(2^n-1)2^{-m(1-H_2(\tau))}.
\end{align}

\end{theorem}

\begin{proof}
The clocked telescoping identity
\cref{thm-controlled-clocked-lift-transfer} bounds the difference of the
stopped laws by the initialization defect, the sum of the intertwining
defects, and the target-interface defect.  Replacing the clean gate by the
public verifier adds its single symmetric-difference event.  Applying the
full LPN family functional and the verifier estimate gives the family
bound.
\end{proof}

\begin{theorem}[Fast-fiber resolvent transfer for a valid LPN lift]
\label{thm-lpn-fast-fiber-resolvent-transfer}
Adopt the presentation, lift, and target notation of
\cref{prop-lpn-valid-lift-native-capacity-access}.

Finally, suppose the complete dynamic lift belongs to
\(\mathscr L_{\rm fm}(M_F,\kappa_F,\varepsilon_+,\varepsilon_-)\), and put

\[
\sigma_\lambda
=\frac{\varepsilon_-M_F\varepsilon_+}
       {\lambda+\kappa_F}.
\]

For the clean killed target resolvents, assume that the represented initial
functional and the initial target atom, nonterminal row, and target-contact
pullback interfaces are exact in the sense of
\cref{thm-controlled-noisy-gate-lift-transfer}.  The operator comparison is

\begin{equation}
\label{eq-lpn-fast-fiber-lift-resolvent-transfer}
\left\|P_\kappa(\lambda-L_\Omega)^{-1}U
 -(\lambda-L_X)^{-1}\right\|
\le
\frac{\|R_X\|^2\sigma_\lambda}
     {1-\|R_X\|\sigma_\lambda},
\qquad
R_X=(\lambda-L_X)^{-1},
\end{equation}

whenever \(\|R_X\|\sigma_\lambda<1\).  In a normalized killed-kernel
supremum norm, \(\|R_X\|\le\lambda^{-1}\).  If the represented initial
functional and target-contact column have dual and primal norms at most one,
then \(\sigma_\lambda<\lambda\) gives

\[
A_{T_I}^{\Omega}(\lambda)
=A_{\pi^{-1}T_I}(\lambda;L_\Omega,\widetilde\mu_0),
\qquad
A_{T_I}^{X}(\lambda)
=A_{T_I}(\lambda;L_X,\mu_0),
\]

and

\begin{equation}
\label{eq-lpn-fast-fiber-lift-arrival-transfer}
\left|A_{T_I}^{\Omega}(\lambda)-A_{T_I}^{X}(\lambda)\right|
\le
\frac{\sigma_\lambda}{\lambda(\lambda-\sigma_\lambda)}.
\end{equation}

Exact dynamic intertwining gives \(\sigma_\lambda=0\).  Replacing the clean
gate by \(\widehat T_I\) adds
\(\mathbf1_{\{\widehat T_I\ne T_I\}}\) pointwise, and hence adds at most

\begin{equation}
\label{eq-lpn-valid-lift-verifier-error}
\exp[-mD_{\mathrm{Ber}}(\tau\Vert\eta)]
+(2^n-1)2^{-m(1-H_2(\tau))}
\end{equation}

after applying \(\mathfrak M_{n,m,\eta}\).  Thus the verifier error is one
presentation event for the complete clocked or resolvent comparison; it is
not multiplied by the horizon.

\end{theorem}

\begin{proof}
The Feshbach identity of
\cref{thm-controlled-valid-lift-block-resolvent} bounds the effective base
perturbation by \(\sigma_\lambda\).  The resolvent identity and a Neumann
series give the displayed operator norm whenever
\(\|R_X\|\sigma_\lambda<1\).  Pairing with the represented initial
functional and target-contact column yields the scalar arrival bound.
Replacing the target gate changes the comparison only on the verifier error
event, whose family probability is the displayed binomial/code bound.
\end{proof}

\begin{corollary}[Quantitative capacity and access under valid LPN lifts]
\label{cor-lpn-valid-lift-quantitative-transfer-synthesis}
\label{thm-lpn-valid-lift-quantitative-transfer}
The native invariant, Schur-capacity, clocked-arrival, and fast-fiber
resolvent conclusions are respectively
\cref{prop-lpn-valid-lift-native-capacity-access,thm-lpn-valid-lift-schur-capacity,thm-lpn-clocked-lift-arrival-transfer,thm-lpn-fast-fiber-resolvent-transfer}.
They concern different hypotheses and output units; none may be substituted
for another without its displayed converter.
\end{corollary}

\begin{proof}

The valid-Lift identity in
\cref{lem-lpn-exact-lift-composition} applied to the uniform LPN law gives
\cref{eq-lpn-valid-lift-target-mass}.  For the native invariant lifted
kernel, the boundary-flux identity
\cref{eq-controlled-target-boundary-capacity-flux} gives

\[
\operatorname{Cap}_{\Phi_{P_\Omega}}
 (\pi^{-1}T_I,\Omega\setminus\pi^{-1}T_I)
=Q_{P_\Omega}(\pi^{-1}T_I,
              \Omega\setminus\pi^{-1}T_I)
\le\widetilde\mu(\pi^{-1}T_I)=2^{-n}.
\]

Substitution in \cref{eq-controlled-capacity-finite-horizon-hit} proves
\cref{eq-lpn-valid-lift-native-hit}.

Apply \cref{thm-controlled-lift-capacity-schur} to the represented condenser
domain.  The upper base-capacity bound is
\cref{eq-controlled-target-condenser-flux-upper}; uniform invariance gives

\[
Q_P(T_I^c,T_I)=\mu_I(s)(1-P(s,s))
=2^{-n}(1-P(s,s)).
\]

This proves \cref{eq-lpn-valid-lift-capacity-schur}.  When the opposite plate
is all of \(T_I^c\), the indicator of \(T_I\) is the unique condenser
competitor, so the last display is the exact base capacity.  This proves
\cref{eq-lpn-valid-lift-capacity-complement}.  The form-neutral conclusion is
the \(B_0=0\) case of
\cref{eq-controlled-lift-capacity-sandwich}; exact pullback-form comparison
then proves \cref{eq-lpn-form-neutral-capacity-equality}.  Taking the
represented supremum and using \(1-P(s,s)\le1\) gives
\cref{eq-lpn-saturated-lift-capacity-bound} by
\cref{def-controlled-saturated-lift-capacity-amplification}.

The corrected observable-pullback form of
\cref{thm-controlled-clocked-lift-transfer} gives
\cref{eq-lpn-clocked-lift-arrival-transfer}.  On
\(\{\widehat T_I=T_I\}\), the stopped gates coincide, so the same estimate
applies to \(p_{\widehat T_I,H}^{\Omega}\).  On the complementary event both
arrival probabilities lie in \([0,1]\), which proves
\cref{eq-lpn-clocked-lift-public-gate-transfer}.  Apply the family functional
and use \cref{eq-lpn-verifier-good-event-probability} to obtain
\cref{eq-lpn-clocked-lift-public-gate-family-transfer}.

The fast-fiber bound
\cref{eq-controlled-fast-fiber-schur-error} is exactly
\cref{eq-lpn-fast-fiber-lift-resolvent-transfer}.  A normalized killed
semigroup is a supremum-norm contraction, hence its Laplace-transform
resolvent has norm at most \(\lambda^{-1}\).  Pair the operator estimate with
the represented initial and target functionals and substitute this norm bound.
The exact pullback interfaces eliminate the three interface-error terms in
\cref{eq-controlled-noisy-gate-discounted-arrival}, giving
\cref{eq-lpn-fast-fiber-lift-arrival-transfer}.  Exact
intertwining sets \(E_+=0\), hence \(\sigma_\lambda=0\).  The same
good-event split used for the clocked comparison proves
\cref{eq-lpn-valid-lift-verifier-error} for the public verifier gate.

\end{proof}

\Cref{fig-lpn-lift-schur-clocked-transfer} separates the two numerical
comparisons available for a valid lift: static elimination of fiber modes and
clocked transport of stopped laws.

\begin{figure}[tbp]
\centering
\resizebox{\linewidth}{!}{%
\begin{tikzpicture}[fw diagram,x=1cm,y=1cm]
  \node[fw source,text width=2.5cm] (base) at (0,2.0)
    {base condenser\\\((T,F,H_X)\)};
  \node[fw provenance,text width=2.8cm] (block) at (3.7,2.0)
    {lifted block form\\\(\bigl[\begin{smallmatrix}H_X&B\\B^*&H_F\end{smallmatrix}\bigr]\)};
  \node[fw geom,text width=2.9cm] (schur) at (7.5,2.0)
    {fiber elimination\\\(H_X-BH_F^{-1}B^*\)};
  \node[fw terminal,text width=3.1cm] (cap) at (11.5,2.0)
    {capacity sandwich\\\((1-\rho_\pi)a_\pi\operatorname{Cap}_X\le
      \operatorname{Cap}_\Omega\le b_\pi\operatorname{Cap}_X\)};
  \node[fw source,text width=2.5cm] (clockx) at (0,-1.0)
    {base stopped laws\\\(P_t^X\)};
  \node[fw provenance,text width=2.7cm] (pull) at (3.7,-1.0)
    {pullbacks \(U_t\)\\defects \(E_t\)};
  \node[fw use,text width=2.8cm] (clocko) at (7.5,-1.0)
    {lifted stopped laws\\\(P_t^\Omega\)};
  \node[fw terminal,text width=3.1cm] (arrival) at (11.5,-1.0)
    {arrival transfer\\\(\varepsilon_{\rm init}+\sum_t\|E_t\|+\varepsilon_{\rm gate}\)};
  \draw[fw arrow] (base)--(block);
  \draw[fw arrow] (block)--(schur);
  \draw[fw arrow] (schur)--(cap);
  \draw[fw arrow] (clockx)--(pull);
  \draw[fw arrow] (pull)--(clocko);
  \draw[fw arrow] (clocko)--(arrival);
  \draw[fw dashed arrow] (block)--node[right,font=\scriptsize]
    {fast-fiber scale \(\kappa_F\)} (pull);
\end{tikzpicture}%
}
\caption{Two distinct valid-Lift comparisons.  The upper row eliminates
fiber modes and compares condenser energies; the lower row telescopes
clocked intertwining defects and compares first-arrival probabilities.
Target-fraction preservation initializes both rows but does not set either
\(\rho_\pi\) or the clocked defects to zero.}
\label{fig-lpn-lift-schur-clocked-transfer}
\end{figure}

\subsection{Ambient ancestry and permutation bounds}
\label{subsec-lpn-ambient-ancestry-permutation-bounds}

\begin{definition}[Residual-independent and residual-ancestry ambient LPN records]
\label{def-lpn-ambient-geom-ancestry-branches}

Consider a temporally admissible LPN master certificate in the ambient family
\(\mathcal C_X(B)\) of
\cref{def-master-saturated-quotient-geometry-profile}; thus its terminal
quotient is the identity and its geometry slot is active.  Apply least-cost
source assignment as in \cref{thm-full-saturated-profile-decomposition}, so
that valid \(\Lift\), \(\Bdry\), and filtration postprocessings retain their
authenticated source ancestry.

The record is \emph{residual-independent} when the source-ancestor sub-DAG of
its complete qualified active record
\(\operatorname{Act}_{\Geom}(\mathscr C)\), including every represented
qualification gate, factors through \(A\), \(\eta\), the candidate
coordinates, public constants, and represented randomness, and has no read
of \(b\) or of the residual primitive \(r_I(x)=Ax+b\).  It is
\emph{residual-ancestry} when that complete source-ancestor sub-DAG contains
such a read.  This binary distinction records only whether the public
noisy-label primitive occurs in the complete qualified ancestry.  It is not
a structural-source class.
Write

\[
\mathcal C_{X,0}^{\mathrm{LPN}}(B),
\qquad
\mathcal C_{X,b}^{\mathrm{LPN}}(B)
\]

for the two record families.  They form a disjoint partition of the ambient
LPN family.  For the LPN expectation functional
\(\mathfrak M_{n,m,\eta}\), let

\[
G_{X,0,\mathfrak M,n}^{\mathrm{src,sat}}(B),
\qquad
G_{X,b,\mathfrak M,n}^{\mathrm{src,sat}}(B)
\]

denote the corresponding restrictions of the existing prospective ambient
profile.  These are provenance restrictions of \(G_{X,n}^{\mathrm{src,sat}}\),
not additional structural coordinates.

\end{definition}

\begin{corollary}[Five-family classification of residual-ancestry LPN
geometry]
\label{cor-lpn-residual-ancestry-five-family-geom-classification}

Every record in \(\mathcal C_{X,b}^{\mathrm{LPN}}(B)\) has the unique ambient
support tag \(\mathsf X\) and a nonempty source signature contained in
\[
\{\mathsf{Add},\mathsf{Mult},\mathsf{Ord},\mathsf{Sym},\mathsf{Filt}\}.
\]
Its complete active tuple, including every locality, conductance, potential,
gap, reach, regularity, current, authority, and comparison datum used by its
visibility, factors through the canonical charged joint frontier of
\cref{def-active-geom-source-frontier} at the class-preserving enlarged cost
\(\Psi_{\Geom,5}(B)\).  Mixed derivations retain every applicable tag.
Consequently, ``contains a read of \(b\) or \(r_I\)'' leaves no unclassified
geometric source: it is exhausted by the five existing provenance families
with ambient support, and introduces neither a sixth provenance family nor a
new structural channel.

\end{corollary}

\begin{proof}

The identity terminal quotient fixes the support tag \(\mathsf X\).  Apply
\cref{thm-exhaustive-five-family-geom-source-classification} to the complete
active derivation selected in
\cref{def-lpn-ambient-geom-ancestry-branches}.  Its factorization,
nonempty signature, and complete-cost bound are
\cref{eq-geom-active-frontier-factorization,eq-geom-active-frontier-cost}.

\end{proof}

\begin{lemma}[Target-permutation bound for residual-independent LPN geometry]
\label{lem-lpn-syndrome-free-ambient-geom-bound}

For every saturated budget \(B\),

\begin{equation}
\label{eq-lpn-syndrome-free-ambient-geom-bound}
G_{X,0,\mathfrak M,n}^{\mathrm{src,sat}}(B)
\le 2^{-n}.
\end{equation}

Equivalently, a residual-independent ambient geometry has no positive expected
advantage over target-neutral enumeration.  A selector whose structural
source ancestry remains residual-independent has zero expected advantage, under the
LPN presentation law, over the uniform kernel on the remaining untested
candidates.

\end{lemma}

\begin{proof}

Put \(N=2^n\).  Fix one family-uniform record in
\(\mathcal C_{X,0}^{\mathrm{LPN}}(B)\), and condition on \(A\) and on every
represented random coordinate in its source derivation.  By
\cref{def-lpn-task-object,def-lpn-ambient-geom-ancestry-branches}, its
potential \(D\) is then independent of the uniform secret \(s\).  Let

\[
Z=\{x\in\mathbb F_2^n:D(x)=0\},
\qquad k=|Z|.
\]

A qualified ambient record for the singleton target \(\{s\}\) requires
\(D(s)=0\), by \cref{def-geom-extraction}.  Equivalently for the
family-functional evaluation, extend the contribution of a derivation scheme
by zero on every instance on which its output is not a qualified record.  If
\(D\) is constant, its variance and ambient
visibility are zero.  Suppose that \(D\) is nonconstant and \(s\in Z\).
Since \(\mathcal M(D)\subseteq L^2(\sigma(D))\), orthogonal-projection
monotonicity gives

\[
R_{\{s\}}(D)
\le
\operatorname{ProjMass}_{\sigma(D)}(\{s\}).
\]

The fiber of \(D\) containing \(s\) is \(Z\).  For the centered singleton
contrast \(y_s=\mathbf 1_{\{s\}}-N^{-1}\), direct conditional expectation on
the fibers of \(D\) gives

\begin{equation}
\label{eq-lpn-zero-fiber-projection-mass}
\operatorname{ProjMass}_{\sigma(D)}(\{s\})
=
\frac{N-k}{k(N-1)}.
\end{equation}

Indeed, \(\mathbb E[y_s\mid\sigma(D)]\) equals \(k^{-1}-N^{-1}\)
on \(Z\) and \(-N^{-1}\) off \(Z\); its squared norm divided by
\(\|y_s\|_2^2=(N-1)/N^2\) is the right-hand side of
\eqref{eq-lpn-zero-fiber-projection-mass}.  Every other factor in the ambient
visibility of \cref{def-geom-extraction} lies in \([0,1]\).  Therefore,
conditional on the fixed source record,

\begin{align*}
\mathbb E_s\!\left[
 \mathbf 1_{\{s\in Z\}}
 V^X_{\Phi,D,J,\mathfrak K}
\right]
&\le
\frac{k}{N}\frac{N-k}{k(N-1)}\\
&=
\frac{N-k}{N(N-1)}
\le \frac1N.
\end{align*}

Removing the conditioning preserves the bound.  The expectation occurs
before the supremum over the one uniform derivation scheme in the definition
of the family-functional profile.  Taking that supremum proves
\eqref{eq-lpn-syndrome-free-ambient-geom-bound}.

For the operational statement, let \(\mathcal H_t^0\) be the represented
subrecord containing \(A\), the residual-independent source data, the selector's
randomness, the candidates already tested, and their negative outcomes.  The
ancestry condition makes the selector kernel measurable with respect to
\(\mathcal H_t^0\); in particular, it does not use \(b\).  Conditional on
\(\mathcal H_t^0\) and survival, the secret is uniform on the remaining
untested candidates, because \(s\) is initially uniform and each negative
outcome removes exactly the tested candidate.  If \(N_t\) candidates remain,
then

\[
\mathbb E\!\left[
 K_t(s\mid\mathcal H_t^0)
 \given \mathcal H_t^0,\tau>t
\right]
=
\frac1{N_t}
\sum_{x\in\mathcal C_t}K_t(x\mid\mathcal H_t^0)
=\frac1{N_t}.
\]

The same identity holds for the uniform baseline.  Averaging over
\(\mathcal H_t^0\) proves that the step advantage in
\cref{eq-pre-hit-selector-advantage} is zero.  This is the
target-permutation calculation of
\cref{cor-target-permutation-neutral-enumeration}, applied to the represented
source subfiltration through which the selector factors.

\end{proof}

\section{Residual Perturbations, Likelihood Quotients, and Direct-Target Obstructions}
\label{sec-lpn-residual-perturbation-quotient-obstructions}

\subsection{Residual Displacement and Simultaneous Separation}
\label{subsec-lpn-residual-displacement-separation}

\begin{lemma}[Exact residual displacement laws for LPN]
\label{lem-lpn-exact-residual-displacement-laws}

Let \(I=(A,b,\eta)\) be an LPN presentation and put

\[
W_I(x)=|r_I(x)|_1=|Ax+b|_1.
\]

Every ordered pair of candidate states has a unique displacement
\(u=z+x\in\mathbb F_2^n\), so \(z=x+u\), and

\begin{align}
\label{eq-lpn-exact-residual-perturbation}
r_I(x+u)&=r_I(x)+Au,\\
\label{eq-lpn-exact-weight-perturbation}
W_I(x+u)-W_I(x)
&=|Au|_1
-2\bigl|\operatorname{supp}r_I(x)
       \cap\operatorname{supp}(Au)\bigr|.
\end{align}

Writing \(v=x+s\), these identities become

\begin{equation}
\label{eq-lpn-hidden-displacement-coordinates}
r_I(x)=Av+E,
\qquad
r_I(x+u)=A(v+u)+E.
\end{equation}

The unique displacement sending \(x\) to the target is \(u=v=x+s\).
For fixed \(v,u\), chosen independently of the LPN draw, the following
distributional statements hold:

\begin{enumerate}[label=\textup{(\roman*)}]
\item If \(v=0\) and \(u\ne0\), then
      \(W_I(s)\sim\operatorname{Bin}(m,\eta)\),
      \(W_I(s+u)\sim\operatorname{Bin}(m,1/2)\), and
      \[
      \mathbb E\bigl[W_I(s+u)-W_I(s)\bigr]
      =m(1/2-\eta).
      \]
      Moreover,
      \begin{equation}
      \label{eq-lpn-target-fixed-perturbation-tail}
      \Pr\{W_I(s+u)\le W_I(s)\}
      \le
      \exp\!\left[-\frac m2(1/2-\eta)^2\right].
      \end{equation}
\item If \(v\ne0\) and \(u\notin\{0,v\}\), then
      \(W_I(x)\) and \(W_I(x+u)\) are independent
      \(\operatorname{Bin}(m,1/2)\) random variables.  Hence, for every
      \(a>0\),
      \begin{equation}
      \label{eq-lpn-off-target-fixed-perturbation-tail}
      \Pr\!\left\{
      \left|W_I(x+u)-W_I(x)\right|\ge am
      \right\}
      \le2\exp\!\left(-\frac{a^2m}{2}\right).
      \end{equation}
\item If \(v\ne0\) and \(u=v\), then \(x+u=s\) and the second residual
      in \eqref{eq-lpn-hidden-displacement-coordinates} is \(E\).
\end{enumerate}

\end{lemma}

\begin{proof}
The additive group gives the unique displacement \(z=x+u\), and linearity
gives \(r_I(x+u)=r_I(x)+Au\).  The weight identity follows coordinatewise.
For fixed nonzero, distinct vectors, the corresponding random linear forms
of a uniform row of \(A\) are independent fair bits; adding the independent
noise bit preserves the asserted binomial laws.  Hoeffding's inequality
gives the two fixed-pair tails.
\end{proof}

\begin{lemma}[Uniform residual flatness and kernel drift]
\label{lem-lpn-uniform-residual-flatness-kernel-drift}
Adopt the LPN presentation and residual notation of
\cref{lem-lpn-exact-residual-displacement-laws}.

For every \(\delta>0\), define the simultaneous residual-flatness event

\begin{equation}
\label{eq-lpn-simultaneous-residual-flatness-event}
\mathcal F_\delta
=
\left\{
\bigl||E|_1-\eta m\bigr|\le\delta m
\right\}
\cap
\bigcap_{0\ne w\in\mathbb F_2^n}
\left\{
\bigl||Aw+E|_1-m/2\bigr|\le\delta m
\right\}.
\end{equation}

It obeys the uniform estimate

\begin{equation}
\label{eq-lpn-simultaneous-residual-flatness-probability}
\Pr(\mathcal F_\delta^c)
\le
2e^{-2\delta^2m}
+2(2^n-1)e^{-2\delta^2m}
=2^{n+1}e^{-2\delta^2m}.
\end{equation}
Equivalently, for every \(\gamma>0\), the choice
\begin{equation}
\label{eq-lpn-simultaneous-residual-flatness-delta}
\delta_{n,m,\gamma}
=
\sqrt{\frac{(n+1)\log 2+\gamma}{2m}}
\end{equation}
gives \(\Pr(\mathcal F_{\delta_{n,m,\gamma}}^c)\le e^{-\gamma}\).
The target margin in
\eqref{eq-lpn-uniform-target-perturbation-margin} is positive whenever
\begin{equation}
\label{eq-lpn-positive-uniform-target-margin-sample-threshold}
m>
\frac{2((n+1)\log 2+\gamma)}{(1/2-\eta)^2}.
\end{equation}

Consequently, on \(\mathcal F_\delta\), simultaneously for every candidate
\(x\) and every displacement \(u\), including choices made from the complete
public presentation,

\begin{align}
\label{eq-lpn-uniform-off-target-perturbation-bound}
x\ne s,\ x+u\ne s
&\quad\Longrightarrow\quad
|W_I(x+u)-W_I(x)|\le2\delta m,\\
\label{eq-lpn-uniform-target-perturbation-margin}
x\ne s,\ x+u=s
&\quad\Longrightarrow\quad
W_I(x)-W_I(s)\ge
m(1/2-\eta-2\delta).
\end{align}
For any represented normalized kernel \(P\), including a kernel selected
from the complete public instance, put \(p_s=P(x,s)\).  If \(x\ne s\) and
\(\Delta_\delta=1/2-\eta-2\delta\), then every represented normalized
kernel, whether selected before or after observing \((A,b)\), satisfies

\begin{align}
\label{eq-lpn-uniform-all-kernel-drift-upper}
(P-I)W_I(x)
&\le-m\Delta_\delta p_s+2\delta m(1-p_s),\\
\label{eq-lpn-uniform-all-kernel-drift-lower}
(P-I)W_I(x)
&\ge-m(1/2-\eta+2\delta)p_s-2\delta m(1-p_s).
\end{align}
Equivalently, all family-parameter dependence in the normalized drift is
displayed by
\begin{equation}
\label{eq-lpn-uniform-normalized-kernel-drift}
\left|
\frac{(P-I)W_I(x)}m
+(1/2-\eta)P(x,s)
\right|
\le2\delta.
\end{equation}
Thus all off-target perturbations contribute at most \(2\delta m\) in
absolute residual-weight drift, while the only perturbation having the
target margin is the unique target-cancelling displacement.

\end{lemma}

\begin{proof}
Hoeffding's inequality applies to \(|E|_1\) and to every
\(|Aw+E|_1\), \(w\ne0\).  A union bound over the \(2^n\) residuals gives
the stated simultaneous event.  On that event the target and off-target
intervals give the two displacement bounds.  Averaging those bounds against
an arbitrary normalized kernel, while isolating its mass at \(s\), gives
the upper, lower, and normalized drift formulas.
\end{proof}

\begin{lemma}[Entropy-sharp simultaneous residual separation]
\label{lem-lpn-entropy-sharp-residual-separation}
Adopt the notation of
\cref{lem-lpn-exact-residual-displacement-laws,lem-lpn-uniform-residual-flatness-kernel-drift}.
There is also an entropy-sharp, rate-dependent form.  For every
\[
\eta<\tau_-<\tau_+<1/2,
\]
the simultaneous target-separation event satisfies
\begin{equation}
\label{eq-lpn-entropy-sharp-all-perturbation-separation}
\Pr\!\left\{
W_I(s)\le\tau_-m
\ \text{and}\
\min_{x\ne s}W_I(x)>\tau_+m
\right\}
\ge
1-e^{-mD_{\mathrm{Ber}}(\tau_-\Vert\eta)}
-(2^n-1)2^{-m(1-H_2(\tau_+))}.
\end{equation}
On this event, every target-cancelling perturbation has the quantitative
margin
\begin{equation}
\label{eq-lpn-entropy-sharp-target-perturbation-margin}
W_I(x)-W_I(s)>m(\tau_+-\tau_-)
\qquad(x\ne s).
\end{equation}
If the event is intersected with \(\mathcal F_\delta\), then every
represented normalized kernel satisfies
\begin{equation}
\label{eq-lpn-entropy-sharp-kernel-drift-upper}
\frac{(P-I)W_I(x)}m
<
-(\tau_+-\tau_-)P(x,s)
+2\delta(1-P(x,s)).
\end{equation}
The event supporting this joint perturbation estimate has probability at
least
\begin{equation}
\label{eq-lpn-joint-all-perturbation-event-probability}
1
-2^{n+1}e^{-2\delta^2m}
-e^{-mD_{\mathrm{Ber}}(\tau_-\Vert\eta)}
-(2^n-1)2^{-m(1-H_2(\tau_+))}.
\end{equation}
For a rate schedule \(n/m\to R\in(0,1)\), constants
\(\tau_-,\tau_+\) with exponentially small exceptional probability exist
precisely throughout the displayed construction whenever
\begin{equation}
\label{eq-lpn-entropy-sharp-perturbation-rate-window}
\eta<\tau_-<\tau_+<H_2^{-1}(1-R).
\end{equation}
More explicitly, if \(n/m\to R\in(0,1)\) and \(\gamma=o(n)\), then
\begin{equation}
\label{eq-lpn-simultaneous-perturbation-radius-rate-asymptotic}
\delta_{n,m,\gamma}
=
\sqrt{\frac{R\log2}{2}}+o(1).
\end{equation}
Thus the two-sided Hoeffding audit has a positive uniform target margin when
\begin{equation}
\label{eq-lpn-hoeffding-perturbation-positive-rate-window}
\eta
<
\frac12-\sqrt{2R\log2},
\end{equation}
whenever the right side is positive.  The entropy-sharp one-sided audit
replaces this sufficient window by
\(\eta<H_2^{-1}(1-R)\), through
\cref{eq-lpn-entropy-sharp-perturbation-rate-window}.

\end{lemma}

\begin{proof}
The lower-tail Chernoff bound controls the planted residual and the Hamming
ball estimate controls each nonplanted residual.  A union bound proves the
separation event.  Intersecting it with the uniform flatness event gives the
kernel inequality.  The interval
\(\eta<\tau_-<\tau_+<H_2^{-1}(1-R)\) is nonempty exactly in the stated
rate window; direct substitution gives the Hoeffding comparison window.
\end{proof}

\subsection{Base and Lifted Residual Drift}
\label{subsec-lpn-base-lifted-residual-drift}

\begin{proposition}[Exact base and lifted residual-drift decomposition]
\label{prop-lpn-base-lifted-residual-drift-decomposition}
Adopt the notation of
\cref{lem-lpn-exact-residual-displacement-laws}.  For such a kernel, the
corresponding one-step drift has the exact
instancewise expansion

\begin{equation}
\label{eq-lpn-all-perturbation-drift-expansion}
(P-I)W_I(x)
=
\sum_{u\in\mathbb F_2^n}
P(x,x+u)\bigl(W_I(x+u)-W_I(x)\bigr).
\end{equation}

If \(v=x+s\ne0\), the same sum separates the unique target-cancelling
perturbation from every other displacement:

\begin{align}
\label{eq-lpn-target-cancelling-drift-separation}
(P-I)W_I(x)
&=
P(x,s)\bigl(W_I(s)-W_I(x)\bigr)
\notag\\
&\quad+
\sum_{u\notin\{0,v\}}
P(x,x+u)\bigl(W_I(x+u)-W_I(x)\bigr),\\
\label{eq-lpn-target-cancelling-drift-remainder-bound}
\left|
(P-I)W_I(x)
-P(x,s)\bigl(W_I(s)-W_I(x)\bigr)
\right|
&\le
\max_{u\notin\{0,v\}}
\left|W_I(x+u)-W_I(x)\right|,
\end{align}

where the maximum is zero when its index set is empty.  These identities
also hold with \(P_\Phi\) in place of \(P\).  Thus the represented
terminal-contact term and off-target residual-weight variation are the two
exhaustive contributions to the drift of every legal identity-supported
perturbation kernel.  The terminal-contact term is accounted through
\cref{thm-lpn-no-free-direct-target-perturbation}; it is not an additional
structural source.

Thus \eqref{eq-lpn-exact-residual-perturbation} and
\eqref{eq-lpn-all-perturbation-drift-expansion} cover every legal
identity-supported candidate perturbation.  Data-dependent selection of
the weights \(P(x,x+u)\) remains part of the represented kernel record and
its saturated cost.

More generally, let \(\Omega\) be a represented clocked or lifted LPN
carrier with its existing candidate projection
\(\pi_X:\Omega\to\mathbb F_2^n\), and let \(P_\Omega\) be a represented
normalized kernel on \(\Omega\).  For \(\omega\in\Omega\), put
\(x=\pi_X(\omega)\) and
\[
u(\omega,\omega')
=\pi_X(\omega')+\pi_X(\omega).
\]
Then
\begin{equation}
\label{eq-lpn-lifted-all-perturbation-drift-expansion}
(P_\Omega-I)(W_I\circ\pi_X)(\omega)
=
\sum_{\omega'\in\Omega}
P_\Omega(\omega,\omega')
\bigl(W_I(x+u(\omega,\omega'))-W_I(x)\bigr).
\end{equation}
Transitions changing only memory, clock, or transcript coordinates have
\(u(\omega,\omega')=0\), so they contribute zero to this drift.
On \(\mathcal F_\delta\), if \(x\ne s\) and
\[
p_s^\Omega(\omega)
=
\sum_{\omega':\,\pi_X(\omega')=s}P_\Omega(\omega,\omega'),
\]
then \eqref{eq-lpn-uniform-all-kernel-drift-upper} and
\eqref{eq-lpn-uniform-all-kernel-drift-lower} hold with
\((P_\Omega-I)(W_I\circ\pi_X)(\omega)\) and
\(p_s^\Omega(\omega)\) in place of \((P-I)W_I(x)\) and \(p_s\).
The normalized and entropy-sharp bounds
\cref{eq-lpn-uniform-normalized-kernel-drift,eq-lpn-entropy-sharp-kernel-drift-upper}
hold under the same substitutions.

\end{proposition}

\begin{proof}
Expand \((P-I)W_I(x)\) over the displacement parametrization and isolate the
unique displacement \(u=x+s\).  The uniform remainder estimate follows by
bounding the convex combination by its largest absolute summand.  On a
lifted carrier, apply the same identity after the candidate projection;
transitions contained in one projection fiber have zero displacement and
therefore zero residual drift.  The base simultaneous bounds then apply to
the projected transition mass.
\end{proof}

\begin{corollary}[Exhaustive residual transport under LPN candidate perturbations]
\label{cor-lpn-exhaustive-identity-perturbation-transport}
\label{lem-lpn-exhaustive-identity-perturbation-transport}
The exact displacement laws, uniform and entropy-sharp simultaneous events,
and base/lifted drift decompositions are
\cref{lem-lpn-exact-residual-displacement-laws,lem-lpn-uniform-residual-flatness-kernel-drift,lem-lpn-entropy-sharp-residual-separation,prop-lpn-base-lifted-residual-drift-decomposition}.
Together they cover every represented identity-supported perturbation and
retain data-dependent kernel selection in the charged record.
\end{corollary}

\begin{proof}

The additive group of \(\mathbb F_2^n\) gives the unique representation
\(z=x+u\).  Since \(r_I(x)=Ax+b\), linearity gives

\[
r_I(x+u)=Ax+Au+b=r_I(x)+Au,
\]

which proves \eqref{eq-lpn-exact-residual-perturbation}.  For binary vectors
\(p,q\),

\[
|p+q|_1-|p|_1
=|q|_1-2|\operatorname{supp}p\cap\operatorname{supp}q|.
\]

Substitution of \(p=r_I(x)\) and \(q=Au\) proves
\eqref{eq-lpn-exact-weight-perturbation}.  Equation
\eqref{eq-lpn-hidden-displacement-coordinates} follows from
\(b=As+E\).  The equation \(x+u=s\) has the unique solution \(u=x+s=v\).

For a fixed nonzero vector \(w\), the coordinates of \(Aw\) are independent
uniform bits.  If \(v=0\) and \(u\ne0\), this observation gives the two
binomial marginals in part \textup{(i)}.  The increment is a sum of \(m\)
independent variables in \([-1,1]\), each with mean \(1/2-\eta\).
Hoeffding's inequality gives
\eqref{eq-lpn-target-fixed-perturbation-tail}.

If \(v\ne0\) and \(u\notin\{0,v\}\), the two vectors \(v\) and \(v+u\)
are distinct and nonzero, hence linearly independent over \(\mathbb F_2\).
For each row \(a_i\), the pair
\((\langle a_i,v\rangle,\langle a_i,v+u\rangle)\) is uniform on
\(\mathbb F_2^2\).  Adding the common independent bit \(E_i\) to both
coordinates preserves the uniform product law.  Independence across rows
gives the asserted independent binomial variables.  Their difference is a
sum of independent mean-zero variables in \([-1,1]\); the two-sided
Hoeffding bound proves
\eqref{eq-lpn-off-target-fixed-perturbation-tail}.  Part \textup{(iii)} is
the substitution \(u=v\) in
\eqref{eq-lpn-hidden-displacement-coordinates}.

Substituting \(z=x+u\) in the definition of \(PW_I-W_I\) gives
\eqref{eq-lpn-all-perturbation-drift-expansion}.  Isolating the term \(u=v\)
gives \eqref{eq-lpn-target-cancelling-drift-separation}; the remaining
kernel weights have sum at most one, so the triangle inequality gives
\eqref{eq-lpn-target-cancelling-drift-remainder-bound}.  The same algebra
applies to the sample-and-hold kernel \(P_\Phi\).

For each fixed nonzero \(w\), the vector \(Aw+E\) has independent uniform
coordinates.  Hoeffding's inequality and a union bound over the
\(2^n-1\) nonzero vectors give

\[
\Pr\!\left\{
\max_{0\ne w\in\mathbb F_2^n}
\bigl||Aw+E|_1-m/2\bigr|>\delta m
\right\}
\le2(2^n-1)e^{-2\delta^2m}.
\]

The same inequality applied to \(|E|_1\) gives
\(2e^{-2\delta^2m}\).  A second union bound proves
\eqref{eq-lpn-simultaneous-residual-flatness-probability}.  Substitution of
\eqref{eq-lpn-simultaneous-residual-flatness-delta} gives
\(2^{n+1}e^{-2\delta^2m}=e^{-\gamma}\).  The inequality
\(\delta_{n,m,\gamma}<(1/2-\eta)/2\) is equivalent to
\eqref{eq-lpn-positive-uniform-target-margin-sample-threshold}.

If \(x\ne s\),
then \(v=x+s\ne0\), so \eqref{eq-lpn-hidden-displacement-coordinates} and
\(\mathcal F_\delta\) place \(W_I(x)\) in
\([m/2-\delta m,m/2+\delta m]\).  This proves
\eqref{eq-lpn-uniform-off-target-perturbation-bound}; comparison with
\(W_I(s)=|E|_1\in[\eta m-\delta m,\eta m+\delta m]\) proves
\eqref{eq-lpn-uniform-target-perturbation-margin}.  In
\eqref{eq-lpn-target-cancelling-drift-separation}, the target term lies
between
\(-m(1/2-\eta+2\delta)p_s\) and
\(-m\Delta_\delta p_s\), while the total absolute contribution of all
other displacements is at most \(2\delta m(1-p_s)\).  This proves
\cref{eq-lpn-uniform-all-kernel-drift-upper,eq-lpn-uniform-all-kernel-drift-lower}.
Because \(\mathcal F_\delta\) is simultaneous over every nonzero \(w\), no
independence between the selected kernel and the public presentation is
used in these kernel bounds.

For the entropy-sharp form, Chernoff's inequality gives
\[
\Pr\{W_I(s)>\tau_-m\}
\le e^{-mD_{\mathrm{Ber}}(\tau_-\Vert\eta)}.
\]
For each \(x\ne s\), the residual \(A(x+s)+E\) is uniform on
\(\mathbb F_2^m\).  The binary Hamming-ball estimate and a union bound give
\[
\Pr\!\left\{\min_{x\ne s}W_I(x)\le\tau_+m\right\}
\le(2^n-1)2^{-m(1-H_2(\tau_+))}.
\]
This proves \cref{eq-lpn-entropy-sharp-all-perturbation-separation}.  The
difference of the two strict thresholds proves
\cref{eq-lpn-entropy-sharp-target-perturbation-margin}.  Isolating the
target term in \cref{eq-lpn-target-cancelling-drift-separation} and applying
\cref{eq-lpn-uniform-off-target-perturbation-bound} proves
\cref{eq-lpn-entropy-sharp-kernel-drift-upper}.  Combining the three
exceptional-probability bounds proves
\cref{eq-lpn-joint-all-perturbation-event-probability}.  If \(n/m\to R\), its
false-candidate exponent is positive exactly when
\(R<1-H_2(\tau_+)\).  Strict monotonicity of \(H_2\) on \([0,1/2]\) gives
\cref{eq-lpn-entropy-sharp-perturbation-rate-window}; both exceptional
terms then decay exponentially.  Substituting \(m=n/R+o(n)\) and
\(\gamma=o(n)\) into
\cref{eq-lpn-simultaneous-residual-flatness-delta} proves
\cref{eq-lpn-simultaneous-perturbation-radius-rate-asymptotic}.  The
condition \(1/2-\eta-2\delta_{n,m,\gamma}>0\) then gives
\cref{eq-lpn-hoeffding-perturbation-positive-rate-window}.

For the lifted statement, substitute
\(x=\pi_X(\omega)\) and
\(u(\omega,\omega')=\pi_X(\omega')+\pi_X(\omega)\) into
\(P_\Omega(W_I\circ\pi_X)-W_I\circ\pi_X\).  This proves
\eqref{eq-lpn-lifted-all-perturbation-drift-expansion}.  Grouping successors
by their projected displacement recovers
\eqref{eq-lpn-all-perturbation-drift-expansion}; the group with displacement
zero has zero increment.  On \(\mathcal F_\delta\), group instead the
successors whose candidate projection is \(s\).  Their total mass is
\(p_s^\Omega(\omega)\); every remaining nonzero projected displacement is
off-target.  The preceding target and off-target bounds therefore prove the
two asserted lifted-kernel inequalities.

\end{proof}

\Cref{fig-lpn-residual-perturbation-dichotomy} displays the exhaustive
candidate-displacement split behind the uniform drift bounds.

\begin{figure}[tbp]
\centering
\begin{tikzpicture}[fw diagram,node distance=8mm and 11mm]
  \node[fw source,text width=2.5cm] (state) {candidate \(x\)\\displacement \(u\)};
  \node[fw provenance,text width=2.7cm,right=of state] (identity)
    {exact identity\\\(r(x+u)=r(x)+Au\)};
  \node[fw terminal,text width=2.8cm,above right=8mm and 12mm of identity] (target)
    {\(u=x+s\)\\target-cancelling\\margin \(m(1/2-\eta-2\delta)\)};
  \node[fw residual,text width=2.8cm,below right=8mm and 12mm of identity] (flat)
    {\(u\ne x+s\)\\off-target\\variation at most \(2\delta m\)};
  \node[fw use,text width=2.7cm,right=18mm of identity] (kernel)
    {represented kernel\\weights \(P(x,x+u)\)};
  \node[fw geom,text width=2.8cm,right=of kernel] (drift)
    {drift = target mass\\+ flat remainder};
  \node[fw provenance,text width=3.1cm,below=27mm of kernel] (lift)
    {lifted transition projects to \(u\); within-fiber moves have \(u=0\)};
  \draw[fw arrow] (state)--(identity);
  \draw[fw arrow] (identity)--(target);
  \draw[fw arrow] (identity)--(flat);
  \draw[fw arrow] (target)--(kernel);
  \draw[fw arrow] (flat)--(kernel);
  \draw[fw arrow] (kernel)--(drift);
  \draw[fw dashed arrow] (lift)--(kernel);
\end{tikzpicture}
\caption{Every candidate transition has one displacement.  Exactly one
displacement reaches the secret and receives the target margin; all others
remain uniformly flat on \(\mathcal F_\delta\).  A data-dependent kernel may
choose the weights, but it cannot create a third drift term.  The same split
holds after projection of a clocked or lifted transition.}
\label{fig-lpn-residual-perturbation-dichotomy}
\end{figure}

\subsection{Residual Likelihood Quotients and Compact Factorization}
\label{subsec-lpn-residual-likelihood-quotients}

\begin{theorem}[Residual likelihood is an exact or compensated quotient record]
\label{thm-lpn-residual-likelihood-quotient-classification}

Fix an LPN presentation \(I=(A,b,\eta)\) with \(0<\eta<1/2\), and let
\(V_I\) be the public candidate verifier.  Define
\begin{equation}
\label{eq-lpn-residual-likelihood-factorization}
W_I(x)=|Ax+b|_1,
\qquad
\mathcal L_I(x)
=(1-\eta)^{m-W_I(x)}\eta^{W_I(x)}.
\end{equation}
Then \(\mathcal L_I=\ell_\eta\circ W_I\), where
\(\ell_\eta(k)=(1-\eta)^{m-k}\eta^k\) is strictly decreasing on
\(\{0,\ldots,m\}\).  Hence \(W_I\) and \(\mathcal L_I\) generate the same
represented fiber partition.  The terminal-compatible refinement
\begin{equation}
\label{eq-lpn-terminal-compatible-likelihood-quotient}
\widetilde q_{W,I}(x)=\bigl(V_I(x),W_I(x)\bigr)
\end{equation}
has at most \(m+2\) nonempty cells, is computable from the original LPN
presentation, and has complete saturated cost bounded by the residual
evaluation, Hamming-weight, verifier, and pair-encoding costs.

Let \(P_I\) be a represented normalized candidate kernel used with this
readout.  For \(x'\) satisfying \(W_I(x')=W_I(x)\), exact dynamic
compatibility of \(\widetilde q_{W,I}\) requires, for each quotient cell
\(C\),
\begin{equation}
\label{eq-lpn-likelihood-quotient-exact-lumpability-test}
\sum_{z:\widetilde q_{W,I}(z)=C}P_I(x,z)
=
\sum_{z:\widetilde q_{W,I}(z)=C}P_I(x',z).
\end{equation}
Writing \(z=x+u\), every term in this comparison is determined by
\begin{equation}
\label{eq-lpn-likelihood-quotient-intersection-dependence}
W_I(x+u)
=W_I(x)+|Au|_1
-2\bigl|\operatorname{supp}r_I(x)
\cap\operatorname{supp}(Au)\bigr|.
\end{equation}
Thus the scalar likelihood level alone determines the quotient transition law
precisely when the represented kernel makes the support-intersection
dependence in
\cref{eq-lpn-likelihood-quotient-intersection-dependence} constant after
aggregation over every destination cell.

Every temporally admissible target-qualified use of the factorized likelihood
has the following exhaustive classification.

\begin{enumerate}[label=\textup{(\roman*)}]
\item When \cref{eq-lpn-likelihood-quotient-exact-lumpability-test} holds and
      the quotient-utility condition is positive, the complete record belongs
      to the exact \(\Abs\) profile.
\item When the quotient has nonzero represented defect and its compensated
      comparison passes the existing defect, mixing, consistency, reach, and
      regularity gates, the complete record belongs to the reversible or
      general-resolvent compensated quotient mode.
\item When neither comparison is qualified, the likelihood is a represented
      terminal-discrimination readout with zero prospective quotient-source
      value.
\item A refinement retaining additional residual coordinates is evaluated as
      the refined quotient when its fibers remain nonidentity.  An injective
      refinement has identity support and belongs to ambient \(\Geom\); there
      the likelihood and weight computations are derived evaluators and the
      complete source value is exactly
      \cref{eq-lpn-identity-geom-complete-gate-product}.
\end{enumerate}

Consequently, polynomial evaluation of the factorized likelihood creates no
additional identity-supported source coordinate.  Its dynamically useful
compression is charged to the existing exact or compensated quotient modes;
an injective refinement is charged to the existing ambient mode with the same
represented kernel, local differences, authority, reach, regularity, and
complete cost.

\end{theorem}

\begin{proof}

The first factorization is immediate from the Bernoulli likelihood.  Since
\(\eta/(1-\eta)<1\), the ratio
\(\ell_\eta(k+1)/\ell_\eta(k)=\eta/(1-\eta)\) is strictly smaller than one;
therefore \(W_I\) and \(\mathcal L_I\) have identical fibers.  The map
\(W_I\) evaluates \(Ax+b\) and counts its nonzero coordinates.  Pairing it
with the represented verifier yields
\cref{eq-lpn-terminal-compatible-likelihood-quotient}; its first coordinate
isolates the target on the verifier good event, and its remaining cells are
indexed by at most \(m+1\) weight values.

Exact intertwining for a finite represented quotient is equivalent to equality
of the total transition mass from any two states in one fiber to every
destination fiber.  This is
\cref{eq-lpn-likelihood-quotient-exact-lumpability-test}.  The perturbation
identity
\cref{eq-lpn-exact-weight-perturbation} is exactly
\cref{eq-lpn-likelihood-quotient-intersection-dependence}, proving the stated
criterion.

The four alternatives are the exact-support, reversible compensated,
general-resolvent compensated, and identity-support modes of
\cref{thm-affordable-geom-abs-normal-form,def-master-compensated-quotient-geometry-certificate}.
The exact case additionally uses the quotient-utility gate
\cref{def-paper1-abs-usefulness-gate}.  Derived likelihood, weight,
comparison, and verifier evaluators retain their ancestry and add no
independent source by
\cref{thm-affordable-composition-no-creation,lem-derived-constructor-zero-contextual-charge}.
The identity-supported value is the complete product in
\cref{eq-lpn-identity-geom-complete-gate-product}.  These support modes are
exhaustive, which proves the final assertion.

\end{proof}

\begin{remark}[Compact evaluation does not locate an exponentially supported
spectrum]
\label{rem-lpn-compact-factorization-no-spectral-locator}

For \(0<\lambda<1\), the represented factorization
\begin{equation}
\label{eq-lpn-compact-exponential-spectrum-example}
F_{\lambda,I}(x)
=
\prod_{i=1}^{m}
\left(1+\lambda(-1)^{a_i\cdot x+b_i}\right)
\end{equation}
is evaluable using \(m\) residual tests and \(m-1\) multiplications.  Its
complete record retains the representation of \(\lambda\) and every
arithmetic and precision cost.  Its
formal expansion is
\begin{equation}
\label{eq-lpn-compact-exponential-spectrum-expansion}
F_{\lambda,I}(x)
=
\sum_{U\subseteq[m]}
\lambda^{|U|}
(-1)^{\sum_{i\in U}b_i}
(-1)^{(A^{\mathsf T}\mathbf1_U)\cdot x},
\end{equation}
which can contain exponentially many Walsh terms.  The compact evaluator
does not enumerate or locate those terms.  Moreover,
\begin{equation}
\label{eq-lpn-compact-factorization-is-weight-readout}
F_{\lambda,I}(x)
=
(1+\lambda)^{m-W_I(x)}(1-\lambda)^{W_I(x)},
\end{equation}
so it generates exactly the residual-weight fibers already classified in
\cref{thm-lpn-residual-likelihood-quotient-classification}.

Accordingly, \cref{eq-lpn-compact-exponential-spectrum-example} receives no
spectral source value from the number of terms in
\cref{eq-lpn-compact-exponential-spectrum-expansion}.  A represented
location rule and its target-relevance comparison are charged as required by
\cref{cor-lpn-spectral-mass-requires-charged-location,%
thm-lpn-no-independent-diffuse-spectral-source}.  A dynamically
qualified nonidentity use is charged to the existing exact or compensated
quotient mode.  An identity-supported use retains the complete ambient
\(\Geom\) gates.  When the factorization supplies only an extremal score,
with no represented authorized intermediate transport, it is a terminal
verifier by
\cref{cor-extremal-readout-requires-intermediate-transport}.

\end{remark}

\subsection{Direct-Target Perturbation Obstruction}
\label{subsec-lpn-direct-target-obstruction}

\begin{theorem}[No-free direct-target perturbation for LPN]
\label{thm-lpn-no-free-direct-target-perturbation}

Work on the verifier good event of
\cref{thm-lpn-public-verifier}, so that \(T_I=\{s\}\) and
\(V_I=\mathbf1_{\{s\}}\).  Let \(\mathcal A\) be a budget-\(B\)
computation and let \(K_t(\cdot\mid\mathcal H_t)\) be any of its temporally
admissible represented next-test samplers, after the one-test clock refinement of
\cref{def-represented-pre-hit-selector}.  Let \(\nu_t\) be its declared
target-neutral baseline.  Then the following identities and provenance
alternatives are exhaustive.

\begin{enumerate}[label=\textup{(\roman*)}]
\item The acceptance mass of the represented sampler is exactly its
      terminal-contact mass:
      \begin{equation}
      \label{eq-lpn-direct-target-terminal-gate}
      K_t(V_I=1\mid\mathcal H_t)
      =
      K_t(s\mid\mathcal H_t).
      \end{equation}
      For a time-homogeneous candidate kernel \(P\), this is the
      discrete-time specialization
      \[
      k_T^P(x)
      =
      \sum_{z\in T_I}P(x,z)
      =
      P(x,s)
      \]
      of the terminal-contact gate in
      \cref{def-killed-dynamics,def-full-history-target-transport}.

\item The one-step and finite-horizon target-contact probabilities are
      \begin{align}
      \label{eq-lpn-direct-target-one-step-hazard}
      \Pr\{\tau=t+1\}
      &=
      \mathbb E\!\left[
      \mathbf1_{\{\tau>t\}}
      K_t(s\mid\mathcal H_t)
      \right],\\
      \label{eq-lpn-direct-target-full-hazard}
      \Pr\{\tau\le H\}
      &=
      \sum_{t=0}^{H-1}
      \mathbb E\!\left[
      \mathbf1_{\{\tau>t\}}
      K_t(s\mid\mathcal H_t)
      \right].
      \end{align}
      Equivalently,
      \begin{equation}
      \label{eq-lpn-direct-target-baseline-decomposition}
      \Pr\{\tau\le H\}
      =
      \operatorname{Enum}_{\nu}(\mathcal A,H)
      +
      \sum_{t=0}^{H-1}
      \operatorname{Adv}_t(K_t;\nu_t).
      \end{equation}
      A uniform proposal on all \(2^n\) candidates has
      \(\nu_t(s\mid\mathcal H_t)=2^{-n}\).  Uniform sampling without
      replacement from \(N_t\) remaining candidates has conditional contact
      mass \(N_t^{-1}\).

\item For the fixed kernel \(P\), its \(H\)-step target value is the
      represented truncated committor
      \begin{equation}
      \label{eq-lpn-direct-target-truncated-committor}
      u_{H,P}
      =
      \sum_{j=0}^{H-1}P_N^j k_T^P.
      \end{equation}
      On a clocked transcript carrier the same formula applies to the
      homogeneous lifted kernel.  For the associated represented generator
      \(L=P-I\), the exact discounted committor is
      \[
      u_{\lambda,L}
      =
      (\lambda I-L_N)^{-1}k_T^L.
      \]
      Each expression contributes prospectively only through a represented
      derivation at its complete saturated cost, as prescribed by
      \cref{def-no-free-committor,thm-no-truncated-neumann-bypass}.

\item Any positive excess above the neutral baseline is already bounded by
      the completed five-family accounting charges.  If \(H_B(n)\) is the normalized
      budget horizon and \(D_B^{\mathrm{sel}}(n)\) is
      \cref{eq-prospective-selector-charge-budget}, then
      \begin{equation}
      \label{eq-lpn-direct-target-existing-charge-bound}
      \bigl(\operatorname{Adv}_t(K_t;\nu_t)\bigr)_+
      \le
      eH_B(n)
      \sum_{c\in\mathfrak J_{\Abs}^5}
      \mathcal C_{I,\mathfrak F_5}^{(c)}
      \bigl(D_B^{\mathrm{sel}}(n)\bigr).
      \end{equation}
      Thus \(K_t(s\mid\mathcal H_t)\) is an analytic target-contact
      coordinate of the represented transition rule.  It introduces no
      accounting coordinate outside the existing five-family ledger.  It
      contributes to a prospective \(\Geom\), \(\Abs\), or coupled source
      only through the represented pre-hit comparison required by
      \cref{cor-contextual-structural-charge-transfer}.

\item The standard LPN primitive signature in
      \cref{def-lpn-task-object} contains \(A,b,\eta\), candidate
      coordinates, the residual evaluator, and their declared constructors.
      It contains neither \(s\), \(E\), a target sampler, nor a committor.
      Consequently, a target-biased sampler used by the original
      presentation has a family-uniform represented derivation in the
      budgeted saturated closure and is charged by part~\textup{(iv)}.
      An evaluator or sampler without such a derivation is an added oracle
      primitive or nonuniform advice and defines a different presented
      family by
      \cref{conv-presentation-relative-budgeted-computation,prop-structural-accounting-is-universal,thm-oracle-admissibility-presentation-separation}.
\end{enumerate}

In particular, the verifier supplies an affordable test of a proposed
candidate but no target-generating sampler.  Under independent uniform
proposals, verifier rejection sampling satisfies
\begin{align}
\label{eq-lpn-uniform-verifier-rejection-success}
\Pr\{\tau\le H\}
&=
1-(1-2^{-n})^H
\le H2^{-n},\\
\label{eq-lpn-uniform-verifier-rejection-cost}
\mathbb E[\tau]
&=2^n.
\end{align}
Replacing the uniform proposal by a target-biased proposal invokes the
represented sampler and the charge in
\cref{eq-lpn-direct-target-existing-charge-bound}.

\end{theorem}

\begin{proof}

On the verifier good event,
\[
K_t(V_I=1\mid\mathcal H_t)
=
\sum_{z\in X_I}
K_t(z\mid\mathcal H_t)\mathbf1_{\{z=s\}}
=K_t(s\mid\mathcal H_t),
\]
which proves
\cref{eq-lpn-direct-target-terminal-gate}.  Conditional on
\(\mathcal H_t\) and survival through time \(t\), the next tested candidate
has law \(K_t\).  Hence its conditional hit probability is the right-hand
side of \cref{eq-lpn-direct-target-terminal-gate}.  Multiplication by the
survival indicator and expectation prove
\cref{eq-lpn-direct-target-one-step-hazard}.  The first-hit events are
pairwise disjoint, so summation proves
\cref{eq-lpn-direct-target-full-hazard}.  Adding and subtracting the declared
baseline in every summand gives
\cref{eq-lpn-direct-target-baseline-decomposition}, which is the LPN
specialization of
\cref{eq-selector-hazard-decomposition}.  The two displayed neutral masses
follow by counting the target in the full and remaining candidate carriers.

For a fixed \(P\), the probability of first entering \(T_I\) after \(j\)
surviving steps is \(P_N^jk_T^P\).  Summing \(j=0,\ldots,H-1\) proves
\cref{eq-lpn-direct-target-truncated-committor}.  The clock coordinate makes
a finite inhomogeneous computation homogeneous on its represented transcript
carrier.  The discounted formula and its derivability requirements are
\cref{def-target-resolvent,def-no-free-committor}; finite truncations are
charged by \cref{thm-no-truncated-neumann-bypass}.

The sampler, random-bit use, transition implementation, and complete prefix
are ancestors of the next test and lie in the saturated derivability closure
at the enlarged represented budget.  Applying
\cref{thm-prospective-selector-structural-charge} proves
\cref{eq-lpn-direct-target-existing-charge-bound}.  The primitive list and
the exclusion of \(s,E\) are explicit in
\cref{def-lpn-task-object}.  Computation closure therefore gives the two
presentation-relative alternatives in part~\textup{(v)}: a sampler in the
original task has a charged represented derivation, while adjoining an
underived evaluator or instance-dependent table changes the presentation by
\cref{thm-oracle-admissibility-presentation-separation}.

Finally, each independent uniform proposal succeeds with probability
\(2^{-n}\).  The failure probability for \(H\) proposals is
\((1-2^{-n})^H\), and the geometric waiting time has mean \(2^n\).
Bernoulli's inequality gives
\(1-(1-2^{-n})^H\le H2^{-n}\), proving
\cref{eq-lpn-uniform-verifier-rejection-success,eq-lpn-uniform-verifier-rejection-cost}.

\end{proof}

\section{Public-Ancestry Local Geometry and Gate Evaluation}
\label{sec-lpn-public-ancestry-gate-evaluation}

\subsection{Target-edge and ancestry normal forms}
\label{subsec-lpn-target-edge-ancestry-normal-forms}

\begin{corollary}[Target-conditioned edge predicates do not supply a
sampler]
\label{cor-lpn-target-edge-predicate-no-sampler}

On the verifier good event, the predicate
\[
(x,z)\longmapsto
\mathbf1_{E_\Phi}(x,z)
\mathbf1_{\{V_I(x)\ne V_I(z)\}}
\]
evaluates the target-incident part of the represented edge relation
\(E_\Phi\).  Its evaluation does not normalize or sample a transition
supported on that set.  A valid identity-supported geometric record using
these edges therefore contains a represented normalized kernel from the
original LPN saturated closure and the same-generator authority comparison of
\cref{def-authority-compatible-caus-alignment}.

With a uniform proposal, the sample-and-hold implementation reaches the
target with one-step probability at most \(2^{-n}\).  A larger target-contact
probability is the selector advantage in
\cref{eq-lpn-direct-target-existing-charge-bound}.  An auxiliary current
assigned to the target-incident edges without the represented kernel and
comparison fails the authority condition and has
\(\operatorname{Align}_{\Caus}^{\mathrm{auth}}=0\).

\end{corollary}

\begin{proof}

Since \(V_I=\mathbf1_{\{s\}}\), its values differ exactly when one endpoint is
\(s\).  Edge-membership evaluation is one represented Boolean computation.
The definition of a represented sampler in
\cref{def-represented-pre-hit-selector} additionally requires its sampling
implementation and random-bit use, while
\cref{def-authorized-ambient-perturbation-witness} forms the canonical
sample-and-hold kernel from the selected normalized kernel.  The uniform
contact probability and every excess above it are
\cref{eq-lpn-uniform-verifier-rejection-success,eq-lpn-direct-target-existing-charge-bound}.
Finally,
\cref{def-authority-compatible-caus-alignment} assigns zero alignment when
the current lacks the required same-generator authorization.

\end{proof}

\begin{theorem}[Public-ancestry normal form for qualified LPN local geometry]
\label{thm-lpn-public-ancestry-local-geometry-normal-form}

Fix an LPN presentation \(I=(A,b,\eta)\), a charged ceiling \(B\), and a
temporally admissible, dynamically qualified, identity-supported ambient
record in the original LPN saturated closure.  At a represented transition,
write \(x\in\mathbb F_2^n\) for the candidate coordinate and \(\xi\) for the
complete remaining represented state.  Put

\[
r=r_I(x)=Ax+b.
\]

The active evaluators that read \(b\) have exact representations using
only \(A,x,r,\xi\), obtained by the public substitution

\begin{equation}
\label{eq-lpn-public-ancestry-b-substitution}
b=r+Ax.
\end{equation}

In particular, the selected normalized transition kernel has the exact
candidate-displacement representation

\begin{equation}
\label{eq-lpn-public-ancestry-kernel-normal-form}
P_I\bigl((x,\xi),(x+u,\xi')\bigr)
=
\pi_I\bigl(u,\xi'\mid A,x,r,\xi\bigr),
\end{equation}

where \(\pi_I\) is the pushforward of the same represented sampler under
\((z,\xi')\mapsto(z+x,\xi')\).  Thus \(\pi_I\) is a reparameterization of
\(P_I\), not an additional kernel or structural source.  Sampling either
side and converting between \(z\) and \(u=z+x\) uses only the public Add
operation and retains the complete cost of \(P_I\).

The same substitution gives represented functions
\(\widehat D_I,\widehat E_I,\widehat c_I\) for the potential, exposed-edge
predicate, and conductance.  Their local values satisfy

\begin{align}
\label{eq-lpn-public-ancestry-potential-normal-form}
D_I(x,\xi)
&=\widehat D_I(A,x,r,\xi),\\
\label{eq-lpn-public-ancestry-local-difference-normal-form}
D_I(x+u,\xi')-D_I(x,\xi)
&=
\widehat D_I(A,x+u,r+Au,\xi')
-\widehat D_I(A,x,r,\xi),\\
\label{eq-lpn-public-ancestry-edge-normal-form}
\mathbf1_{E_\Phi}\bigl((x,\xi),(x+u,\xi')\bigr)
&=\widehat E_I(A,x,r,\xi,u,\xi'),\\
\label{eq-lpn-public-ancestry-conductance-normal-form}
c_\Phi\bigl((x,\xi),(x+u,\xi')\bigr)
&=\widehat c_I(A,x,r,\xi,u,\xi').
\end{align}

These identities give the following exhaustive structural assignment.

\begin{enumerate}[label=\textup{(\roman*)}]
\item A record whose active ancestral set contains no read of \(b\) or
      \(r_I\) belongs to the residual-independent ambient restriction.
\item A record whose active ancestral set contains such a read has
      \(\mathsf{Add}\) source ancestry.  Products, comparisons, mixtures,
      clocks, valid lifts, and derived \(\Caus\) records retain their exact
      operational signatures and costs but contribute no independent target
      alignment beyond this charged ancestry.  When the readout is expressed
      spectrally, its exact target correlation and the complete cost of the
      index-selection rule are those of
      \cref{thm-lpn-walsh-alignment-charged-location}.
\item A target-qualified nonidentity fiber map belongs to the existing
      \(\Abs\) profile.  With identity support, the complete target
      contribution is the existing ambient \(\Geom\) gate product evaluated
      on
      \cref{eq-lpn-public-ancestry-kernel-normal-form,eq-lpn-public-ancestry-local-difference-normal-form}.
\item Direct target contact is the single displacement \(u=x+s\).  Its
      transition mass is
      \(\pi_I(x+s,\xi'\mid A,x,r,\xi)\), summed over \(\xi'\).  This term is
      admitted only through the represented pre-hit sampler and its charged
      target-alignment comparison.  A verifier or target-edge predicate
      evaluates a proposed contact and supplies no sampler for this
      displacement.
\end{enumerate}

Consequently the LPN local-geometry ledger is exactly the qualified
visibility of the represented data in the displays above, with the complete
public derivation and cost retained.

\end{theorem}

\begin{proof}

The primitive public data are those of
\cref{def-lpn-task-object}.  At a transition with candidate coordinate \(x\),
the identity \(b=r_I(x)+Ax\) proves
\cref{eq-lpn-public-ancestry-b-substitution}.  Substitute this represented
identity at each active read of \(b\).  Affordable composition and
no-creation
\cref{thm-affordable-composition-no-creation,lem-lpn-no-creation-final}
retain the original active ancestry and charge the additional public Add
operations.

Sample \((z,\xi')\) from \(P_I((x,\xi),\cdot)\) and put \(u=z+x\).  Addition
by \(x\) is a bijection of \(\mathbb F_2^n\), so this is an exact pushforward
with the same normalization.  This proves
\cref{eq-lpn-public-ancestry-kernel-normal-form}.  The residual transport
identity

\[
r_I(x+u)=r_I(x)+Au
\]

then proves
\cref{eq-lpn-public-ancestry-potential-normal-form,eq-lpn-public-ancestry-local-difference-normal-form,eq-lpn-public-ancestry-edge-normal-form,eq-lpn-public-ancestry-conductance-normal-form}
by the same substitution in the represented evaluators.

The first two cases are the exact ancestry partition of
\cref{def-lpn-ambient-geom-ancestry-branches}, together with the inherited
charge identities
\cref{lem-derived-constructor-zero-contextual-charge,prop-derived-channels-create-no-structure}.
The support alternatives are the ambient and quotient-supported modes of
\cref{thm-affordable-geom-abs-normal-form}; target qualification of a
nonidentity fiber map is the \(\Abs\) admission condition.  Finally,
\(x+u=s\) holds exactly when \(u=x+s\).  The sampler and target-contact
statements are
\cref{cor-lpn-target-edge-predicate-no-sampler,thm-lpn-no-free-direct-target-perturbation}.

\end{proof}

\subsection{Identity-supported gate records}
\label{subsec-lpn-identity-supported-gate-records}

\begin{definition}[Identity-supported LPN gate record]
\label{def-lpn-identity-supported-gate-record}

Work on the good-presentation event of
\cref{thm-lpn-public-verifier}, and let

\[
\mathfrak g=(\mathfrak N,\Phi,D,J,\mathfrak K)
\in\mathcal C_{X,b}^{\mathrm{LPN}}(B)
\]

be any temporally admissible, dynamically qualified,
identity-supported, residual-ancestry ambient record.  Write
\(\Phi=(E_\Phi,c_\Phi)\) and
\(\mathfrak N=(\widetilde L,\zeta,c,P,L,\lambda)\).  Then every geometric
datum and every legal perturbation in the record is evaluated by the
following exhaustive list.

This theorem evaluates records already admitted to
\(\mathcal C_{X,b}^{\mathrm{LPN}}(B)\); it does not construct or admit such a
record.  When that family is empty, its saturated Geom coordinate is zero.

\begin{enumerate}[label=\textup{(\roman*)}]
\item \emph{Provenance and affordability.}  The relation \(E_\Phi\),
      conductance \(c_\Phi\), potential \(D\), selected normalized kernel
      \(P\), current \(J_{\mathfrak K}\), authority comparison, and their
      evaluation procedures belong to one represented derivation of total
      cost at most \(B\).  The record is available before the transition in
      which it is used.  Its \(b\)-dependent evaluators have the exact public
      ancestry normal forms of
      \cref{thm-lpn-public-ancestry-local-geometry-normal-form}.

\item \emph{All exposed perturbations.}  The candidate-coordinate projection
      of each charged edge has the unique form \((x,x+u)\).  Its sampler is
      the pushforward
      \cref{eq-lpn-public-ancestry-kernel-normal-form}, and its potential
      difference is
      \cref{eq-lpn-public-ancestry-local-difference-normal-form}.  On that
      projected edge the public residual changes by \(Au\), as in
      \eqref{eq-lpn-exact-residual-perturbation}.  The canonical legal local
      perturbation is the sample-and-hold kernel
      \(P_\Phi\) of
      \cref{def-authorized-ambient-perturbation-witness}; hence no edge,
      displacement, or rejected transition lies outside the charged
      perturbation record.  Clocked and lifted transitions obey
      \eqref{eq-lpn-lifted-all-perturbation-drift-expansion}.  On the
      simultaneous event \(\mathcal F_\delta\), every such perturbation,
      including a public-instance-dependent choice, obeys
      \cref{eq-lpn-uniform-off-target-perturbation-bound,eq-lpn-uniform-target-perturbation-margin};
      every base or lifted normalized kernel obeys
      \cref{eq-lpn-uniform-all-kernel-drift-upper,eq-lpn-uniform-all-kernel-drift-lower}.
      The exceptional probability and the positive-margin sample threshold
      are given explicitly by
      \cref{eq-lpn-simultaneous-residual-flatness-probability,eq-lpn-positive-uniform-target-margin-sample-threshold}.
      The separated terminal-contact term is exhausted by
      \cref{thm-lpn-no-free-direct-target-perturbation}: it is neutral
      contact, a completed accounting contact, a temporally admissible
      charged source use, or data outside the original LPN presentation.

\end{enumerate}
\end{definition}

\begin{theorem}[Local gates for identity-supported LPN geometry]
\label{thm-lpn-identity-geom-local-gates}
Let \(\mathfrak g=(\mathfrak N,\Phi,D,J,\mathfrak K)\) be an admitted record
in \(\mathcal C_{X,b}^{\mathrm{LPN}}(B)\) as in
\cref{def-lpn-identity-supported-gate-record}.  Its remaining local gates
are the following.
\begin{enumerate}[label=\textup{(\roman*)},start=3]

\item \emph{Transfer-scale mixing.}  The mixing gate is exactly \(M_\Phi\).
      It equals one precisely when the represented comparison data certify
      the required transfer-class Poincar\'e or conductance bound for the
      exposed geometry used by this same normalized record.

\item \emph{Authorized local consistency.}  The local gate is exactly
      \(\chi_\Phi^{\mathrm{auth}}(D)\).  It equals one only when the same
      selected generator supplies the represented local-consistency, drift,
      resolvent, or carrier comparison required by
      \cref{def-authority-compatible-caus-alignment}.  The potential or its
      formal gradient alone does not activate this gate.

\item \emph{Authorized descent.}  When the authority gate is active, the
      complete directed factor is
      \begin{equation}
      \label{eq-lpn-identity-geom-exact-alignment}
      \operatorname{Align}_{\Caus}^{\mathrm{auth}}(\mathfrak K,D)
      =
      \frac{
      \bigl(\langle J_{\mathfrak K},-\nabla_\Phi D\rangle_E\bigr)_+^2
      }{
      \|J_{\mathfrak K}\|_E^2\|\nabla_\Phi D\|_E^2
      },
      \end{equation}
      with value zero when either norm vanishes.  Its positive pairing is
      supported on the descending perturbations
      \(E_\downarrow(\mathfrak g)\) of
      \cref{def-authorized-ambient-perturbation-witness}.

\item \emph{Target reach.}  The exact reach satisfies
      \begin{equation}
      \label{eq-lpn-identity-geom-reach-ceiling}
      R_{\{s\}}(D)
      \le
      \operatorname{ProjMass}_{\sigma(D)}(\{s\}).
      \end{equation}
      Let \(V_I\) be the represented public verifier of
      \cref{thm-lpn-public-verifier} and put
      \(D_V=1-V_I\).  On the good-presentation event,
      \(D_V^{-1}(0)=\{s\}\) and
      \begin{equation}
      \label{eq-lpn-verifier-potential-unit-reach}
      R_{\{s\}}(D_V)=1.
      \end{equation}
      This unit reach records that the extremum identifies the target.  The
      isolated potential \(D_V\) is a terminal verifier.  It enters the
      prospective ambient profile only as part of the same pre-hit record
      that represents and certifies the intermediate transition,
      same-generator authority, authorized descent, and regularity required
      in parts~\textup{(i)}--\textup{(v)} and~\textup{(vii)}.

\item \emph{Dirichlet regularity.}  Every potential has the exact ratio
      \begin{equation}
      \label{eq-lpn-identity-geom-exact-regularity}
      \rho_\Phi(D)
      =
      \frac{
      \displaystyle
      \frac12\sum_{x\in\mathbb F_2^n}
      \sum_{u\in\mathbb F_2^n}
      c_\Phi(x,x+u)
      \bigl(D(x+u)-D(x)\bigr)^2\mu_I(x)\mu_I(x+u)
      }{\operatorname{Var}_{\mu_I}(D)}.
      \end{equation}
      For \(N=2^n\), the verifier potential has
      \begin{equation}
      \label{eq-lpn-verifier-potential-regularity}
      \operatorname{Var}_{\mu_I}(D_V)=\frac{N-1}{N^2},
      \qquad
      \rho_\Phi(D_V)
      =\frac1{N-1}\sum_{z\ne s}c_\Phi(s,z).
      \end{equation}
\end{enumerate}

\end{theorem}

\begin{proof}
The mixing and authority gates are their defining represented comparisons.
The directed factor is the normalized authorized current--gradient pairing.
Target reach is bounded by the target mass of the potential sigma-field, and
the verifier specialization is exact on the verifier event.  Expanding the
Dirichlet form of \(D\), and then of the singleton verifier potential, gives
the two regularity formulas.
\end{proof}

\begin{corollary}[Complete gate evaluation for identity-supported LPN geometry]
\label{cor-lpn-identity-geom-gate-evaluation}
\label{thm-lpn-identity-geom-gate-evaluation}
Under the hypotheses of
\cref{def-lpn-identity-supported-gate-record,thm-lpn-identity-geom-local-gates},
all provenance and local gates belong to the same represented record.

Consequently the complete visibility of the record is exactly

\begin{equation}
\label{eq-lpn-identity-geom-complete-gate-product}
V^X_{\Phi,D,J,\mathfrak K}
=
M_\Phi\,
\chi_\Phi^{\mathrm{auth}}(D)\,
\operatorname{Align}_{\Caus}^{\mathrm{auth}}(\mathfrak K,D)\,
R_{\{s\}}(D)\,
\frac1{1+\rho_\Phi(D)}.
\end{equation}

Parts \textup{(i)}--\textup{(vii)} are one conjunctive admissibility
record.  None is credited as target-aligned structure in isolation.  In
particular, an inverse-polynomial Poincar\'e gap activates only the carrier
gate \(M_\Phi\); it contributes zero to the LPN Geom profile whenever the
same record fails affordability, perturbation provenance, same-generator
consistency, authorized descent, target reach, or Dirichlet regularity.  The
target-aligned transfer quantity is exactly the complete product in
\cref{eq-lpn-identity-geom-complete-gate-product}.

The quantified additive class here is exactly the
\((\mathsf X,\mathsf{Add})\) provenance restriction of the saturated
prospective Geom profile.  Membership requires one family-uniform
represented pre-hit derivation of polynomial saturated cost, with the
complete active tuple and every gate evaluator retained as in
\cref{def-active-geom-source-frontier,thm-exhaustive-five-family-geom-source-classification}.
The simultaneous residual-flatness estimates in part~\textup{(ii)} control
the legal residual perturbations available to that tuple.  Every dynamically
qualified represented nonidentity level-set or spectral coarsening on the
same transfer window is charged to the exact or compensated
quotient-supported branches by
\cref{thm-coercive-regularity-forces-compensated-quotient,cor-lpn-quotient-supported-geom-closure}.
Consequently, an affordable target-aligned additive potential that survives
quotient exhaustion belongs to the identity-supported
\((\mathsf X,\mathsf{Add})\) cell and contributes only through the explicitly
evaluated product
\cref{eq-lpn-identity-geom-complete-gate-product}.  Its target alignment,
local transport, and regularity are evaluated together with the represented
cost, authority, and remaining gates of that same record.

For every \(0<\varepsilon<1\), a contribution larger than
\(\varepsilon\) therefore requires, in this same represented record,

\begin{equation}
\label{eq-lpn-identity-geom-simultaneous-gates}
M_\Phi=1,
\quad
\chi_\Phi^{\mathrm{auth}}(D)=1,
\quad
\operatorname{Align}_{\Caus}^{\mathrm{auth}}(\mathfrak K,D)>\varepsilon,
\quad
R_{\{s\}}(D)>\varepsilon,
\quad
\frac1{1+\rho_\Phi(D)}>\varepsilon.
\end{equation}

The verifier potential in
\eqref{eq-lpn-verifier-potential-unit-reach} evaluates the reach gate.
Activation of the mixing, local-consistency, and authorized-current gates
uses the complete selected geometric record in parts
\textup{(i)}--\textup{(v)} and \textup{(vii)}.
Absent that represented intermediate transport, the verifier potential is
assigned zero prospective source activity by
\cref{cor-extremal-readout-requires-intermediate-transport}; its knowledge of
the terminal extremum supplies no descent rule.

\end{corollary}

\begin{proof}

Part \textup{(i)} is the membership condition in
\cref{def-geom-extraction,def-pre-hit-temporally-admissible-source}.
Part \textup{(ii)} follows from
\cref{lem-lpn-exhaustive-identity-perturbation-transport,def-authorized-ambient-perturbation-witness}.
Parts \textup{(iii)} and \textup{(iv)} are the mixing and
same-generator consistency clauses of
\cref{def-geom-extraction,def-authority-compatible-caus-alignment}.
Part \textup{(v)} is
\cref{def-quantitative-caus-alignment,def-authority-compatible-caus-alignment};
the descending-edge support is
\cref{eq-effective-perturbation-descending-edge-mass}.

Every monotone function of \(D\) is \(\sigma(D)\)-measurable, so
\(\mathcal M(D)\subseteq L^2(\sigma(D))\).  Orthogonal-projection
monotonicity proves
\eqref{eq-lpn-identity-geom-reach-ceiling}.  On the verifier good event,
\(D_V=0\) at \(s\) and \(D_V=1\) elsewhere.  The target indicator
\(1-D_V=\mathbf1_{\{s\}}\) belongs to the closed span of monotone functions
oriented toward \(D_V=0\); hence the centered target contrast belongs to
\(\mathcal M(D_V)\), proving
\eqref{eq-lpn-verifier-potential-unit-reach}.

Substituting \(z=x+u\) into the Dirichlet form of
\cref{def-static-dirichlet-geometry} proves
\eqref{eq-lpn-identity-geom-exact-regularity}.  For \(D_V\), only edges with
exactly one endpoint \(s\) contribute.  Symmetry of \(c_\Phi\) and uniformity
of \(\mu_I\) give

\[
\mathcal E_\Phi(D_V,D_V)
=\frac1{N^2}\sum_{z\ne s}c_\Phi(s,z).
\]

Since \(D_V\) is one on \(N-1\) states and zero on one state, its variance
is \((N-1)/N^2\).  Division gives
\eqref{eq-lpn-verifier-potential-regularity}.

Finally, \eqref{eq-lpn-identity-geom-complete-gate-product} is the LPN
specialization of \cref{def-geom-extraction}.  All five factors lie in
\([0,1]\), with the first two binary.  A product larger than
\(\varepsilon\) forces both binary gates to equal one and each remaining
factor to exceed \(\varepsilon\), proving
\eqref{eq-lpn-identity-geom-simultaneous-gates}.  The last assertion follows
from the separate record requirements in
\cref{thm-valid-geom-source-characterization,def-authority-compatible-caus-alignment}.

\end{proof}

\Cref{fig-lpn-complete-ambient-gate-product} shows why no single favorable
geometric statistic activates an ambient target-aligned record.

\begin{figure}[tbp]
\centering
\resizebox{\linewidth}{!}{%
\begin{tikzpicture}[fw diagram,x=1cm,y=1cm]
  \node[fw source,text width=2.8cm] (record) at (0,0)
    {one complete ambient record\\\((\Phi,D,J,\mathfrak K,\rho)\)};
  \node[fw provenance,text width=2.0cm] (prov) at (3.4,3.2) {affordability\\and ancestry};
  \node[fw geom,text width=2.0cm] (mix) at (3.4,1.6) {transfer-scale\\mixing};
  \node[fw use,text width=2.0cm] (auth) at (3.4,0) {same-generator\\authority};
  \node[fw use,text width=2.0cm] (aim) at (3.4,-1.6) {authorized\\Caus alignment};
  \node[fw terminal,text width=2.0cm] (reach) at (3.4,-3.2) {target reach};
  \node[fw residual,text width=2.2cm] (reg) at (6.7,-3.2) {regularity\\\((1+\rho_\Phi)^{-1}\)};
  \node[fw terminal,text width=3.2cm] (product) at (9.8,0)
    {qualified visibility\\product of all gates};
  \draw[fw arrow] (record)--(prov);
  \draw[fw arrow] (record)--(mix);
  \draw[fw arrow] (record)--(auth);
  \draw[fw arrow] (record)--(aim);
  \draw[fw arrow] (record)--(reach);
  \draw[fw arrow] (record)--(reg);
  \draw[fw arrow] (prov)--(product);
  \draw[fw arrow] (mix)--(product);
  \draw[fw arrow] (auth)--(product);
  \draw[fw arrow] (aim)--(product);
  \draw[fw arrow] (reach)--(product);
  \draw[fw arrow] (reg)--(product);
\end{tikzpicture}%
}
\caption{The identity-supported ambient coordinate is conjunctive.  The
carrier, potential, current, authority comparison, reach certificate, and
regularity comparison belong to one affordable pre-hit record.  A spectral
gap, target extremum, or formal gradient in isolation contributes zero when
any other gate is absent.}
\label{fig-lpn-complete-ambient-gate-product}
\end{figure}

\subsection{Parameter-explicit verifier geometry}
\label{subsec-lpn-parameter-explicit-verifier-geometry}

\begin{corollary}[Parameter-explicit LPN verifier geometry]
\label{cor-lpn-parameter-explicit-verifier-geometry}

Fix \(0<\eta<\tau<1/2\), put \(N=2^n\), and let
\(\mathfrak g=(\mathfrak N,\Phi,D_V,J,\mathfrak K)\) be a record from
\cref{thm-lpn-identity-geom-gate-evaluation} using the public verifier
potential \(D_V=1-V_I\).  With probability at least
\begin{equation}
\label{eq-lpn-parameter-explicit-verifier-geometry-probability}
1-\exp[-mD_{\mathrm{Ber}}(\tau\Vert\eta)]
-(2^n-1)2^{-m(1-H_2(\tau))}
\end{equation}
over the LPN presentation, its complete visibility is
\begin{equation}
\label{eq-lpn-parameter-explicit-verifier-visibility}
V^X_{\Phi,D_V,J,\mathfrak K}
=
\frac{
M_\Phi\,
\chi_\Phi^{\mathrm{auth}}(D_V)\,
\operatorname{Align}_{\Caus}^{\mathrm{auth}}(\mathfrak K,D_V)
}{
1+\displaystyle\frac1{2^n-1}\sum_{z\ne s}c_\Phi(s,z)
}.
\end{equation}

The alignment factor in
\cref{eq-lpn-parameter-explicit-verifier-visibility} is evaluated only after
the complete target-qualified Geom record has been admitted.  It is the
derived \(\Caus\) factor of that same record and supplies no independent
source coordinate.  When the ambient record family is empty, the verifier
potential contributes zero by the empty-supremum convention in
\cref{thm-saturated-ambient-geom-accountability}.
Likewise, a represented evaluation of \(D_V(x)\), or the fact that its unique
minimum is \(s\), remains terminal verification unless the same charged
pre-hit record supplies the authorized intermediate transport and regularity
appearing in
\cref{eq-lpn-parameter-explicit-verifier-visibility}.

The dependence on the LPN family parameters is therefore explicit:
\(2^n\) is the candidate multiplicity, \(m\) is the number of public
equations, \(\eta\) enters through the Bernoulli divergence, and \(\tau\)
is the represented verifier radius.  The budget \(B\) restricts the record
through \(\mathcal C_{X,b}^{\mathrm{LPN}}(B)\); it creates no additional
probability penalty because
\cref{eq-lpn-simultaneous-residual-flatness-probability,eq-lpn-entropy-sharp-all-perturbation-separation}
hold simultaneously over all candidate displacements and therefore over
every kernel retained by the budget-\(B\) saturated closure.

\end{corollary}

\begin{proof}

The probability is
\cref{eq-lpn-verifier-good-event-probability}.  On that event,
\cref{eq-lpn-verifier-potential-unit-reach,eq-lpn-verifier-potential-regularity}
give
\[
R_{\{s\}}(D_V)=1,
\qquad
\rho_\Phi(D_V)
=\frac1{2^n-1}\sum_{z\ne s}c_\Phi(s,z).
\]
Substitution in
\cref{eq-lpn-identity-geom-complete-gate-product} proves
\cref{eq-lpn-parameter-explicit-verifier-visibility}.

\end{proof}

\begin{corollary}[Complete quantitative admissibility ledger for ambient LPN
geometry]
\label{cor-lpn-complete-quantitative-ambient-geom-ledger}

Fix \(0<\eta<\tau<1/2\), \(\delta>0\), and a saturated budget \(B\).
With probability at least

\begin{equation}
\label{eq-lpn-complete-geom-ledger-probability}
1
-\exp[-mD_{\mathrm{Ber}}(\tau\Vert\eta)]
-(2^n-1)2^{-m(1-H_2(\tau))}
-2^{n+1}e^{-2\delta^2m},
\end{equation}

the following statements hold simultaneously for every temporally
admissible residual-ancestry ambient record
\(\mathfrak g=(\mathfrak N,\Phi,D,J,\mathfrak K)\) of complete cost at most
\(B\).

\begin{enumerate}[label=\textup{(\roman*)}]
\item Every candidate edge is uniquely \((x,x+u)\), its public residual
      increment is \(Au\), and every normalized kernel in the record obeys
      \begin{equation}
      \label{eq-lpn-complete-ledger-residual-drift}
      \left|
      \frac{(P-I)W_I(x)}m
      +(1/2-\eta)P(x,s)
      \right|
      \le2\delta
      \qquad(x\ne s).
      \end{equation}
      The unique target-cancelling displacement has residual margin at least
      \(m(1/2-\eta-2\delta)\); every off-target displacement changes
      \(W_I\) by at most \(2\delta m\).

\item The complete admissible visibility, with all factors evaluated on
      this same charged record, satisfies
      \begin{equation}
      \label{eq-lpn-complete-ledger-factorwise-bound}
      \begin{aligned}
      V^X_{\Phi,D,J,\mathfrak K}
      &\le
      \min\Biggl\{
      M_\Phi,
      \chi_\Phi^{\mathrm{auth}}(D),
      \operatorname{Align}_{\Caus}^{\mathrm{auth}}(\mathfrak K,D),
      \operatorname{ProjMass}_{\sigma(D)}(\{s\}),\\
      &\hspace{38mm}
      \frac{1}{1+
      \mathcal E_\Phi(D,D)/\operatorname{Var}_{\mu_I}(D)}
      \Biggr\}.
      \end{aligned}
      \end{equation}
      Here \(M_\Phi=1\) only when the represented reciprocal gap lies in the
      transfer class; \(\chi_\Phi^{\mathrm{auth}}(D)=1\) only when the same
      selected generator supplies the required local comparison; and the
      alignment factor is the exact squared normalized pairing in
      \cref{eq-lpn-identity-geom-exact-alignment}.  These are joint factors,
      not independent sources.

\item For the public verifier potential \(D_V=1-V_I\), the reach and
      regularity factors are evaluated exactly.  For an already admitted
      target-qualified ambient Geom record, the full product is
      \begin{equation}
      \label{eq-lpn-complete-ledger-verifier-bound}
      V^X_{\Phi,D_V,J,\mathfrak K}
      =
      \frac{M_\Phi\chi_\Phi^{\mathrm{auth}}(D_V)
      \operatorname{Align}_{\Caus}^{\mathrm{auth}}(\mathfrak K,D_V)}{
      1+\displaystyle\frac1{2^n-1}
          \sum_{z\ne s}c_\Phi(s,z)
      }.
      \end{equation}
      The alignment factor is derived from this admitted geometry and has no
      standalone source value.  If the target-qualified ambient record is
      absent, this item contributes zero to the saturated profile.

\item For every represented next-test sampler, the target-incident mass is
      exactly \(K_t(s\mid\mathcal H_t)\).  Its positive excess over the
      declared neutral baseline is bounded by
      \begin{equation}
      \label{eq-lpn-complete-ledger-target-contact-charge}
      \bigl(\operatorname{Adv}_t(K_t;\nu_t)\bigr)_+
      \le
      eH_B(n)
      \sum_{c\in\mathfrak J_{\Abs}^5}
      \mathcal C_{I,\mathfrak F_5}^{(c)}
      \bigl(D_B^{\mathrm{sel}}(n)\bigr).
      \end{equation}
      A uniform proposal contributes exactly \(2^{-n}\) per test.

\end{enumerate}

These four conclusions jointly evaluate affordability and provenance, all
legal perturbations, transfer-scale mixing, same-generator consistency,
authorized descent, target reach, Dirichlet regularity, and direct target
contact.  The saturated ambient profile takes the supremum of the complete
same-record product; it does not take separate suprema of these properties
and combine them afterward.

\end{corollary}

\begin{proof}

Intersect the verifier good event of
\cref{eq-lpn-verifier-good-event-probability} with the simultaneous
residual-flatness event of
\cref{eq-lpn-simultaneous-residual-flatness-probability}.  A union bound gives
\cref{eq-lpn-complete-geom-ledger-probability}.  Part~\textup{(i)} is
\cref{eq-lpn-exact-residual-perturbation,eq-lpn-uniform-normalized-kernel-drift,eq-lpn-uniform-off-target-perturbation-bound,eq-lpn-uniform-target-perturbation-margin}.

For part~\textup{(ii)}, apply the exact product
\cref{eq-lpn-identity-geom-complete-gate-product}.  Every factor lies in
\([0,1]\), so the product is bounded by each factor.  Projection monotonicity
gives
\(R_{\{s\}}(D)\le
\operatorname{ProjMass}_{\sigma(D)}(\{s\})\), and the last factor is the
definition of \(\rho_\Phi(D)\).  The gap, consistency, and alignment
interpretations are
\cref{def-geom-extraction,def-authority-compatible-caus-alignment,eq-lpn-identity-geom-exact-alignment}.

For part~\textup{(iii)}, substitute the exact verifier reach and regularity
from
\cref{eq-lpn-verifier-potential-unit-reach,eq-lpn-verifier-potential-regularity}
into the same product.  The source assignment of the alignment factor is
\cref{prop-derived-channels-create-no-structure}.
Part~\textup{(iv)} is
\cref{eq-lpn-direct-target-terminal-gate,eq-lpn-direct-target-existing-charge-bound}.

\end{proof}

\section{Controlled Capacity and Caus Aiming for LPN}
\label{sec-lpn-controlled-capacity-caus-aiming}
\label{subsec-lpn-controlled-nonlinear-ambient-geometry}

The nonlinear ambient branch of
\cref{cor-lpn-complete-quantitative-ambient-geom-ledger} is evaluated on the
same represented kernel and the same native or comparison Dirichlet form.
This introduces no additional structural channel: capacity, drift,
functional inequalities, learned-operator error, and target occupation are
typed numerical outputs of the already qualified \(\Geom/\Caus\) record in
\cref{def-certified-dynamical-geometry,def-controlled-functional-output-slots}.

\subsection{Base-Carrier Capacity and Global Drift}
\label{subsec-lpn-base-capacity-global-drift}

Fix \(0<\eta<\tau<1/2\) in the verifier range of
\cref{thm-lpn-public-verifier} and \(\delta>0\).  Put \(N=2^n\) and define
the simultaneous presentation event

\begin{equation}
\label{eq-lpn-dynamical-good-event}
\mathcal G_{\tau,\delta}
=\{V_I^{-1}(1)=\{s\}\}\cap\mathcal F_\delta .
\end{equation}

The two estimates already proved in
\cref{thm-lpn-public-verifier,lem-lpn-exhaustive-identity-perturbation-transport}
give the explicit finite-parameter bound

\begin{equation}
\label{eq-lpn-dynamical-good-event-probability}
\Pr(\mathcal G_{\tau,\delta})
\ge
1-\exp[-mD_{\mathrm{Ber}}(\tau\Vert\eta)]
 -(2^n-1)2^{-m(1-H_2(\tau))}
 -2^{n+1}e^{-2\delta^2m}.
\end{equation}

\begin{theorem}[Controlled target capacity and drift for LPN]
\label{thm-lpn-controlled-target-capacity-drift}

Fix a presentation in \(\mathcal G_{\tau,\delta}\).  Let \(P\) be the
normalized selected kernel of a certified native dynamical geometry on
\(X_I=\mathbb F_2^n\), suppose that its invariant law is the uniform law
\(\mu_I\), and let \(\Phi_P\) be its native symmetric geometry.  Write

\begin{equation}
\label{eq-lpn-target-boundary-rate}
q_{P,s}=1-P(s,s).
\end{equation}

Then the target condenser, every represented opposite plate
\(F\subseteq X_I\setminus\{s\}\), and every initial law
\(d\mu_0/d\mu_I\le R_0\) satisfy

\begin{align}
\label{eq-lpn-singleton-capacity-exact}
\operatorname{Cap}_{\Phi_P}(\{s\},X_I\setminus\{s\})
&=2^{-n}q_{P,s},\\
\label{eq-lpn-singleton-capacity-opposite-plate}
\operatorname{Cap}_{\Phi_P}(\{s\},F)
&\le2^{-n}q_{P,s},\\
\label{eq-lpn-controlled-capacity-hit-bound}
\Pr_{\mu_0}\{\tau_s\le H\}
&\le R_0 2^{-n}(1+Hq_{P,s})
\le R_0(H+1)2^{-n}.
\end{align}

If a represented comparison geometry \(\Psi\) has the same plates and obeys
\(\mathcal E_{\Phi_P}(f,f)\le b\mathcal E_\Psi(f,f)\) on the condenser
domain, then
\begin{equation}
\label{eq-lpn-comparison-capacity-transfer}
\operatorname{Cap}_{\Phi_P}(\{s\},F)
\le b\operatorname{Cap}_\Psi(\{s\},F).
\end{equation}
The burn-in and relative-uniform versions are obtained by replacing \(R_0\)
with the certified post-burn-in density ratio in
\cref{thm-controlled-target-capacity-drift-envelope}.

More generally, let \(D\ge0\) be represented, let \(D(s)=0\), and suppose
the same selected kernel has a global normalized drift margin

\[
(I-P)D(x)\ge a>0
\qquad(x\ne s).
\]

Then

\begin{equation}
\label{eq-lpn-invariant-global-drift-ceiling}
\frac{a}{\operatorname{Osc}(D)}
\le\frac1{2^n-1}.
\end{equation}

For the public residual potential \(W_I(x)=|Ax+b|_1\), if

\begin{equation}
\label{eq-lpn-residual-global-descent-margin}
-\frac{(P-I)W_I(x)}m\ge a
\qquad(x\ne s),
\end{equation}

then, pointwise on the same event,

\begin{align}
\label{eq-lpn-residual-descent-forces-contact}
P(x,s)
&\ge
\frac{(a-2\delta)_+}{1/2-\eta},\\
\label{eq-lpn-residual-invariant-drift-ceiling}
a
&\le
\frac{(1/2-\eta+2\delta)q_{P,s}}{2^n-1}
\le\frac{1/2-\eta+2\delta}{2^n-1}.
\end{align}

Thus a globally useful nonlinear residual descent and its target-boundary
flux are evaluated in the same native record; neither is inferred from the
mere existence of the terminal residual score.

\end{theorem}

\begin{proof}

For the native symmetric conductance, invariance of \(\mu_I\) gives

\[
\mathcal E_{\Phi_P}(\mathbf1_{\{s\}},\mathbf1_{\{s\}})
=\mu_I(s)(1-P(s,s))=2^{-n}q_{P,s}.
\]

The boundary values of the condenser
\((\{s\},X_I\setminus\{s\})\) determine the unique competitor
\(\mathbf1_{\{s\}}\), proving
\cref{eq-lpn-singleton-capacity-exact}.  The same indicator is an admissible
competitor for \((\{s\},F)\), proving
\cref{eq-lpn-singleton-capacity-opposite-plate}.  The stationary entrance
count and initial-density comparison in
\cref{thm-controlled-target-capacity-drift-envelope} give
\cref{eq-lpn-controlled-capacity-hit-bound}.  The form-comparison conclusion
\cref{eq-controlled-capacity-form-comparison-upper} gives
\cref{eq-lpn-comparison-capacity-transfer}.

Apply the invariant-drift conclusion of
\cref{thm-controlled-target-capacity-drift-envelope} with
\(\mu_I(s)=2^{-n}\) to obtain
\cref{eq-lpn-invariant-global-drift-ceiling}.  On
\(\mathcal F_\delta\),
\cref{eq-lpn-uniform-normalized-kernel-drift} gives

\[
(1/2-\eta)P(x,s)
\ge -\frac{(P-I)W_I(x)}m-2\delta,
\]

which proves \cref{eq-lpn-residual-descent-forces-contact}.  Finally,
invariance makes the uniform mean of \((P-I)W_I\) zero.  Summing
\cref{eq-lpn-residual-global-descent-margin} over the \(N-1\) off-target
states and balancing the sum at \(s\) gives

\[
(N-1)a
\le \frac{(P-I)W_I(s)}m
\le (1/2-\eta+2\delta)(1-P(s,s)),
\]

where the last inequality follows directly from the two intervals in
\cref{eq-lpn-simultaneous-residual-flatness-event}.  Division proves
\cref{eq-lpn-residual-invariant-drift-ceiling}.

\end{proof}

\subsection{Residual Caus Energy and Contact Conversion}
\label{subsec-lpn-residual-caus-contact}

\begin{theorem}[Residual Caus energy and action for LPN]
\label{thm-lpn-residual-caus-energy-action}

Fix a presentation in \(\mathcal G_{\tau,\delta}\), assume
\(\Delta_\delta=1/2-\eta-2\delta>0\), and define the represented
target-normalized residual potential

\begin{equation}
\label{eq-lpn-clipped-residual-caus-potential}
D_{\delta,I}^{\rm res}(x)
=
\left(\frac{W_I(x)}m-\eta-\delta\right)_+.
\end{equation}

This potential is a represented scalar postprocessing of \(W_I\).  Its
source assignment is therefore the assignment of the residual-weight record
in \cref{thm-lpn-residual-likelihood-quotient-classification}; the Caus use
below is derived and creates no additional source by
\cref{prop-derived-channels-create-no-structure}.

Then

\begin{equation}
\label{eq-lpn-clipped-residual-potential-range}
D_{\delta,I}^{\rm res}(s)=0,
\qquad
\Delta_\delta
\le D_{\delta,I}^{\rm res}(x)
\le\frac12-\eta
\quad(x\ne s).
\end{equation}

Let \(R\) be a complete prospective Caus target-flow record on the LPN
candidate carrier, with selected normalized kernel \(P\), authorized current
\(J\), and represented positive edge mobilities \(a\).  Orient the
unoriented active edges arbitrarily and let \(\partial s\) denote those with
exactly one endpoint equal to \(s\).  The target-aiming energy of
\cref{eq-lpn-clipped-residual-caus-potential} satisfies the exact
parameter-dependent estimate

\begin{equation}
\label{eq-lpn-clipped-residual-aiming-split}
\mathcal E_{\rm aim}(D_{\delta,I}^{\rm res};a)
\le
4\delta^2
\sum_{\substack{e=\{x,z\}\ x,z\ne s}}a_e
+
(1/2-\eta)^2\sum_{e\in\partial s}a_e.
\end{equation}

Consequently,

\begin{equation}
\label{eq-lpn-clipped-residual-action-aiming}
\mathcal P_{\Caus}(J,D_{\delta,I}^{\rm res})_+
\le
\sqrt{\mathcal A_{\Caus}(J;a)}
\left(
4\delta^2
\sum_{\substack{e=\{x,z\}\ x,z\ne s}}a_e
+
(1/2-\eta)^2\sum_{e\in\partial s}a_e
\right)^{1/2}.
\end{equation}

\end{theorem}

\begin{proof}
On the simultaneous flatness event, off-target potential increments are at
most \(2\delta\), while target-boundary increments are at most
\(1/2-\eta\).  Splitting the aiming-energy sum over those two edge classes
gives the energy estimate.  The represented Caus action--aiming inequality
then gives the power bound.
\end{proof}

\begin{corollary}[Traffic and horizon specialization of residual Caus aiming]
\label{cor-lpn-traffic-horizon-caus-aiming}
Adopt the hypotheses and potential of
\cref{thm-lpn-residual-caus-energy-action}.

There is a canonical normalized specialization of this estimate.  Suppose
that \(P\) has uniform invariant law and that the complete record uses the
represented traffic mobility and current

\begin{equation}
\label{eq-lpn-traffic-caus-coordinates}
a_{\{x,z\}}
=\mu_I(x)P(x,z)+\mu_I(z)P(z,x),
\qquad
J^{\rm tr}_{\{x,z\}}
=\mu_I(x)P(x,z)-\mu_I(z)P(z,x).
\end{equation}

Take \(J=J^{\rm tr}\), or take its authorized Hodge projection when the
same record contains the inverse-mobility contraction required in
\cref{def-controlled-complete-caus-flow-record}.  Then

\begin{align}
\label{eq-lpn-traffic-caus-action}
\mathcal A_{\Caus}(J;a)
&\le1,\\
\label{eq-lpn-traffic-target-boundary-mobility}
\sum_{e\in\partial s}a_e
&=2^{1-n}q_{P,s}.
\end{align}

For the verifier potential \(D_V=1-V_I\),

\begin{align}
\label{eq-lpn-traffic-verifier-aim}
\mathcal E_{\rm aim}(D_V;a)
&=2^{1-n}q_{P,s},\\
\label{eq-lpn-traffic-verifier-power}
\mathcal P_{\Caus}(J,D_V)_+
&\le2^{(1-n)/2}\sqrt{q_{P,s}},
\end{align}

while the clipped residual potential satisfies

\begin{equation}
\label{eq-lpn-traffic-clipped-residual-aim}
\mathcal P_{\Caus}(J,D_{\delta,I}^{\rm res})_+
\le
\left[
4\delta^2
+(1/2-\eta)^2\,2^{1-n}q_{P,s}
\right]^{1/2}.
\end{equation}

For a clocked sequence of such complete records
\((P_t,J_t,a_t)_{t=0}^{H-1}\), the discrete-time action--aiming bound gives

\begin{align}
\label{eq-lpn-traffic-verifier-flow-horizon}
\sum_{t=0}^{H-1}
\mathcal P_{\Caus}(J_t,D_V)_+
&\le
H\left[
2^{1-n}\frac1H\sum_{t=0}^{H-1}q_{P_t,s}
\right]^{1/2}
\le H2^{(1-n)/2},\\
\label{eq-lpn-traffic-caus-flow-horizon}
\sum_{t=0}^{H-1}
\mathcal P_{\Caus}(J_t,D_{\delta,I}^{\rm res})_+
&\le
H\left[
4\delta^2
+(1/2-\eta)^2\,2^{1-n}
\frac1H\sum_{t=0}^{H-1}q_{P_t,s}
\right]^{1/2}.
\end{align}

In particular, choosing \(\delta=\delta_{n,m,\gamma}\) from
\cref{eq-lpn-simultaneous-residual-flatness-delta} yields

\begin{equation}
\label{eq-lpn-traffic-caus-flow-parameter}
\sum_{t=0}^{H-1}
\mathcal P_{\Caus}(J_t,D_{\delta,I}^{\rm res})_+
\le
H\left[
\frac{2((n+1)\log2+\gamma)}m
+(1/2-\eta)^2\,2^{1-n}
\frac1H\sum_{t=0}^{H-1}q_{P_t,s}
\right]^{1/2}.
\end{equation}

The event on which \cref{eq-lpn-traffic-caus-flow-parameter} holds has the
fully explicit probability

\begin{equation}
\label{eq-lpn-caus-parameter-event-probability}
\Pr\!\left(\mathcal G_{\tau,\delta_{n,m,\gamma}}\right)
\ge
1-\exp[-mD_{\mathrm{Ber}}(\tau\Vert\eta)]
 -(2^n-1)2^{-m(1-H_2(\tau))}
 -e^{-\gamma}.
\end{equation}

Its target margin is positive under the sample condition in
\cref{eq-lpn-positive-uniform-target-margin-sample-threshold}.  Thus the
sample ratio enters the off-target aiming term through
\(2((n+1)\log2+\gamma)/m\), the label-noise rate enters the target-boundary
term through \((1/2-\eta)^2\), and the two exceptional-presentation
exponents retain their exact \((n,m,\eta,\tau)\) dependence.

Every represented aiming-to-contact converter
\cref{eq-controlled-aiming-contact-converter} therefore gives its LPN
contact bound by adding its represented neutral-contact baseline
\(B_R^{\rm contact}\), multiplying the applicable right side of
\cref{eq-lpn-traffic-verifier-flow-horizon,eq-lpn-traffic-caus-flow-parameter}
by its represented \(K_R\) and adding its represented \(\delta_R\).

\end{corollary}

\begin{proof}
The inequality \((u-v)^2/(u+v)\le u+v\) bounds traffic action by one, and
uniform invariance gives the exact target-boundary mobility.  Substitution
in \cref{thm-lpn-residual-caus-energy-action}, followed by Cauchy--Schwarz
over time, proves the horizon formulas.  The stated event probability is the
union bound for verifier separation and residual flatness; the contact
formula is the declared aiming-to-contact converter.
\end{proof}

\begin{proposition}[Residual drift-to-contact conversion]
\label{prop-lpn-residual-drift-contact-conversion}
Adopt the hypotheses and potential of
\cref{thm-lpn-residual-caus-energy-action}.

The same clipped potential also supplies an exact drift-to-contact
conversion.  For \(x\ne s\),

\begin{equation}
\label{eq-lpn-clipped-residual-drift-contact}
\left|
(P-I)D_{\delta,I}^{\rm res}(x)
+D_{\delta,I}^{\rm res}(x)P(x,s)
\right|
\le2\delta(1-P(x,s)).
\end{equation}

Hence a pointwise descent margin
\((I-P)D_{\delta,I}^{\rm res}(x)\ge a\) implies

\begin{equation}
\label{eq-lpn-clipped-residual-contact-lower}
P(x,s)
\ge
\frac{(a-2\delta)_+}{\Delta_\delta}.
\end{equation}

If this margin holds for every \(x\ne s\) and \(P\) has uniform invariant
law, then

\begin{equation}
\label{eq-lpn-clipped-residual-global-drift}
a
\le
\frac{(1/2-\eta)q_{P,s}}{2^n-1}.
\end{equation}

\end{proposition}

\begin{proof}
Separate the target transition from all off-target transitions in the exact
kernel drift.  The flatness event bounds the latter contribution by
\(2\delta(1-P(x,s))\), yielding the pointwise identity and contact lower
bound.  Summing the descent inequality against the uniform invariant law
cancels total drift and leaves only the target boundary, which gives the
global ceiling.
\end{proof}

\begin{corollary}[Learned-parity transport to LPN Caus aiming]
\label{cor-lpn-learned-parity-caus-transport}
Adopt the verifier-good LPN presentation and normalized traffic records of
\cref{cor-lpn-traffic-horizon-caus-aiming}.

Finally, fix \(0<\delta_{\rm stat}<1\), and let a represented parity learner satisfying
\cref{eq-lpn-parity-clean-risk} output
\(\widehat s=\widehat s(A,b;\omega)\) before the controlled transitions.
Retain \(\omega\) as a static coordinate in its public product record, and
put

\[
D_0(x)=\mathbf1_{\{x\ne s\}},
\qquad
\widehat D_\omega(x)=\mathbf1_{\{x\ne\widehat s(\omega)\}}.
\]

On the simultaneous parity-learning event, define

\begin{equation}
\label{eq-lpn-parity-decoder-error-bound}
p_{\rm err}^{\rm par}
=
\min\left\{
1,
\frac{2\bigl(2\Gamma_{n,m,\delta_{\rm stat}}^{\rm par}
+\varepsilon_{\rm opt}\bigr)}{1-2\eta}
\right\}.
\end{equation}

Then

\begin{align}
\Pr_\omega\{\widehat s\ne s\}
&\le p_{\rm err}^{\rm par},
\label{eq-lpn-parity-decoder-error-probability}\\
\mathbb E_{\omega}
\|\widehat D_\omega-D_0\|_{L^2(\mu_I)}^2
&=2^{1-n}\Pr_\omega\{\widehat s\ne s\}
\le2^{1-n}p_{\rm err}^{\rm par}.
\label{eq-lpn-learned-singleton-potential-error}
\end{align}

For normalized uniform-invariant traffic records, the squared norm of the
gradient map from \(L^2(\mu_I)\) to \(\ell^2(a)\) is at most four.  Summing
the same represented comparison over a length-\(H\) record therefore gives
the calibrated-potential certificate

\begin{equation}
\label{eq-lpn-learned-caus-aiming-error}
\varepsilon_{\rm cal}^{2}
\le
4H\,2^{1-n}p_{\rm err}^{\rm par}.
\end{equation}

Combining this estimate with
\cref{thm-controlled-noise-corrected-caus-aiming} gives, for every
represented aiming-to-contact converter with
\(B_R^{\rm contact}\le B_B^{\rm contact}\), \(K_R\le K_B\), and
\(\delta_R\le\delta_B\),

\begin{equation}
\label{eq-lpn-noise-corrected-caus-contact}
U_{\Caus,\mathrm{LPN}}^{\rm noise}
\le
\min\left\{
1,
B_B^{\rm contact}
+K_BH2^{(1-n)/2}
\left(1+2\sqrt{p_{\rm err}^{\rm par}}\right)
+\delta_B
\right\}.
\end{equation}

If
\(2\Gamma_{n,m,\delta_{\rm stat}}^{\rm par}+\varepsilon_{\rm opt}
<(1-2\eta)/2\), then
\cref{eq-lpn-parity-exact-statistical-identification} gives
\(\widehat s=s\); the errors in
\cref{eq-lpn-learned-singleton-potential-error,eq-lpn-learned-caus-aiming-error}
then vanish.

\end{corollary}

\begin{proof}
The clean-risk transport bound controls decoder error.  Two distinct
singleton indicators differ on exactly two points, giving the exact
\(L^2(\mu_I)\) calibration.  The normalized traffic gradient map has squared
operator norm at most four, so summing over the horizon gives
\cref{eq-lpn-learned-caus-aiming-error}.  The framework noise-corrected Caus
theorem then yields the displayed contact certificate.
\end{proof}

\begin{corollary}[Caus action, residual aiming, and noisy target transport for LPN]
\label{cor-lpn-caus-action-residual-aiming}
\label{thm-lpn-caus-action-residual-aiming}
The residual-energy, traffic-horizon, drift-to-contact, and learned-noise
conclusions are respectively
\cref{thm-lpn-residual-caus-energy-action,cor-lpn-traffic-horizon-caus-aiming,prop-lpn-residual-drift-contact-conversion,cor-lpn-learned-parity-caus-transport}.
They share one complete represented Caus record and one declared confidence
allocation.
\end{corollary}

\begin{proof}

On \(\mathcal F_\delta\), the target residual satisfies
\(W_I(s)/m\le\eta+\delta\), whereas every off-target residual satisfies
\(1/2-\delta\le W_I(x)/m\le1/2+\delta\).  Substitution into
\cref{eq-lpn-clipped-residual-caus-potential} proves
\cref{eq-lpn-clipped-residual-potential-range}.  An edge joining two
off-target states therefore has potential increment at most \(2\delta\) in
absolute value, an edge in \(\partial s\) has increment between
\(\Delta_\delta\) and \(1/2-\eta\), and every self-loop has zero increment.
Splitting the aiming-energy sum over these edge classes proves
\cref{eq-lpn-clipped-residual-aiming-split}.  The Caus action--aiming
inequality
\cref{thm-controlled-caus-action-aiming} then proves
\cref{eq-lpn-clipped-residual-action-aiming}.

For the traffic coordinates, the scalar inequality
\((u-v)^2/(u+v)\le u+v\) for \(u,v\ge0\) gives

\[
\mathcal A_{\Caus}(J;a)
\le
\sum_{\{x,z\}}
\bigl[\mu_I(x)P(x,z)+\mu_I(z)P(z,x)\bigr]
\le1.
\]

The same estimate holds for the authorized Hodge projection by its recorded
inverse-mobility contraction.  Uniform invariance gives equality of the
stationary entrance and exit fluxes of \(\{s\}\), so

\[
\sum_{e\in\partial s}a_e
=2\mu_I(s)(1-P(s,s))
=2^{1-n}q_{P,s}.
\]

The verifier potential changes by one exactly on \(\partial s\), proving
\cref{eq-lpn-traffic-verifier-aim}.  Apply
\cref{thm-controlled-caus-action-aiming} and
\cref{eq-lpn-traffic-caus-action} to obtain
\cref{eq-lpn-traffic-verifier-power,eq-lpn-traffic-clipped-residual-aim}.
Summing the action and aiming bounds in time, then applying the discrete-time
form of
\cref{eq-controlled-caus-action-aiming-integrated}, proves
\cref{eq-lpn-traffic-verifier-flow-horizon,eq-lpn-traffic-caus-flow-horizon}.
The identity
\(4\delta_{n,m,\gamma}^2
=2((n+1)\log2+\gamma)/m\) proves
\cref{eq-lpn-traffic-caus-flow-parameter}.  The converter assertion is
\cref{cor-controlled-uniform-caus-flow-bound} applied to these same-record
estimates.

For \(x\ne s\), split the kernel average into the target state and the
off-target states.  On the latter set the positive-part operation in
\cref{eq-lpn-clipped-residual-caus-potential} is inactive, so

\[
(P-I)D_{\delta,I}^{\rm res}(x)
=-D_{\delta,I}^{\rm res}(x)P(x,s)
+\sum_{z\ne s}P(x,z)\frac{W_I(z)-W_I(x)}m.
\]

The simultaneous off-target bound in
\cref{eq-lpn-uniform-off-target-perturbation-bound} proves
\cref{eq-lpn-clipped-residual-drift-contact}.  If
\((I-P)D_{\delta,I}^{\rm res}(x)\ge a\), then

\[
a
\le
(1/2-\eta)P(x,s)+2\delta(1-P(x,s))
=2\delta+\Delta_\delta P(x,s),
\]

which proves \cref{eq-lpn-clipped-residual-contact-lower}.  Under uniform
invariance, the mean of \((P-I)D_{\delta,I}^{\rm res}\) is zero.  Summing
the global descent margin over the \(2^n-1\) off-target states and using

\[
(P-I)D_{\delta,I}^{\rm res}(s)
=\sum_{z\ne s}P(s,z)D_{\delta,I}^{\rm res}(z)
\le(1/2-\eta)q_{P,s}
\]

proves \cref{eq-lpn-clipped-residual-global-drift}.

For a parity output, character orthogonality gives
\(\mathbf1_{\{\widehat s\ne s\}}=2R_0(f_{\widehat s})\).
Apply \cref{eq-lpn-parity-clean-risk} and average over the represented seed
to obtain
\cref{eq-lpn-parity-decoder-error-probability}.  Two distinct singleton
potentials differ exactly at their two singleton states, proving
\cref{eq-lpn-learned-singleton-potential-error}.

For the traffic mobility, put
\(P^{\mathsf s}=(P+P^*)/2\).  Then

\[
\sum_{\{x,z\}}a_{\{x,z\}}(f(z)-f(x))^2
=2\langle f,(I-P^{\mathsf s})f\rangle_{\mu_I}
\le4\|f\|_{L^2(\mu_I)}^2,
\]

because the self-adjoint Markov operator \(P^{\mathsf s}\) has spectrum in
\([-1,1]\).  Summing over \(H\) clock steps and substituting
\cref{eq-lpn-learned-singleton-potential-error} gives the represented
calibrated-potential comparison in
\cref{def-controlled-clean-aiming-comparison}(ii).  The calibrated branch of
\cref{thm-controlled-noise-corrected-caus-aiming} proves
\cref{eq-lpn-learned-caus-aiming-error}.  The observed singleton potential
has aiming energy at most \(H2^{1-n}\), while
\cref{eq-lpn-traffic-caus-action} gives accumulated action at most \(H\).
Substitution into the same theorem proves
\cref{eq-lpn-noise-corrected-caus-contact}, including the represented
neutral-contact baseline \(B_R^{\rm contact}\le B_B^{\rm contact}\).  The
final assertion is
\cref{eq-lpn-parity-exact-statistical-identification}.

\end{proof}

\Cref{fig-lpn-caus-energy-contact-routing} separates the energy estimate from
the additional converter required to turn Caus power into target contact.

\begin{figure}[tbp]
\centering
\resizebox{\linewidth}{!}{%
\begin{tikzpicture}[fw diagram,node distance=8mm and 9mm]
  \node[fw source,text width=2.5cm] (record) {mobility \(a\), current \(J\),\\residual potential \(D^{\rm res}\)};
  \node[fw residual,text width=2.6cm,above right=7mm and 10mm of record] (off)
    {off-target energy\\\(4\delta^2\sum_{e\notin\partial s}a_e\)};
  \node[fw terminal,text width=2.6cm,below right=7mm and 10mm of record] (bdry)
    {target-boundary energy\\\((1/2-\eta)^2\sum_{e\in\partial s}a_e\)};
  \node[fw geom,text width=2.4cm,right=25mm of record] (energy) {aiming energy\\\(\mathcal E_{\rm aim}\)};
  \node[fw use,text width=2.5cm,right=of energy] (cs) {Cauchy--Schwarz\\\(\mathcal P_{\Caus}\le\sqrt{\mathcal A_{\Caus}\mathcal E_{\rm aim}}\)};
  \node[fw provenance,text width=2.4cm,right=of cs] (conv) {represented\\contact converter\\\((B_R,K_R,\delta_R)\)};
  \node[fw terminal,text width=2.2cm,right=of conv] (contact) {target-contact\\bound};
  \draw[fw arrow] (record)--(off);
  \draw[fw arrow] (record)--(bdry);
  \draw[fw arrow] (off)--(energy);
  \draw[fw arrow] (bdry)--(energy);
  \draw[fw arrow] (energy)--(cs);
  \draw[fw arrow] (cs)--(conv);
  \draw[fw arrow] (conv)--(contact);
\end{tikzpicture}%
}
\caption{Residual Caus routing for LPN.  Flat off-target edges and the
target boundary contribute different energy scales.  The action--aiming
inequality controls Caus power, but contact follows only after a represented
converter with its own baseline, gain, defect, ancestry, and cost.}
\label{fig-lpn-caus-energy-contact-routing}
\end{figure}

\subsection{Noise-Indexed Valid-Lift Interfaces}
\label{subsec-lpn-noise-indexed-lift-interface}

\begin{corollary}[Noise-indexed Caus and valid-Lift target interfaces for LPN]
\label{cor-lpn-noise-indexed-caus-lift-interfaces}

Fix a polynomial ceiling \(B=b_C(n)=n^C\), a horizon \(H\),
\(0<\delta_{\rm stat}<1\), and the represented class
\(\mathscr R_{B,H}^{\rm hard}\) of complete LPN dynamical label-noise
records at this ceiling and horizon.
Use the decoder error \(p_{\rm err}^{\rm par}\) from
\cref{eq-lpn-parity-decoder-error-bound} with this confidence allocation.
In the linear homogeneous
clean-aiming branch, \(\rho=1-2\eta\), so
\cref{eq-controlled-noise-corrected-caus-contact} specializes exactly to

\begin{equation}
\label{eq-lpn-linear-noise-corrected-caus-contact}
U_{\Caus}^{\rm noise}
\le
\min\left\{
1,
B_B^{\rm contact}
+K_B\sqrt{\operatorname{Act}_{\Caus}^{\rm sat}(B,H)}
\frac{
\sqrt{\operatorname{Aim}_{\Caus,{\rm obs}}^{\rm sat}(B,H)}
+\varepsilon_{\rm aim}^{\rm sat}}
{1-2\eta}
+\delta_B
\right\}.
\end{equation}

For the binary hard-potential branch, the same framework theorem gives

\begin{equation}
\label{eq-lpn-hard-gate-caus-noise-radius}
\begin{aligned}
(\varepsilon_{{\rm aim},{\rm hard}}^{\rm sat})^2
&=
\sup_{\mathcal R\in\mathscr R_{B,H}^{\rm hard}}
\kappa_{\rm aim,H}(\mathcal R)
U_{{\rm gate},{\rm dep}}^{\rm clean}(\mathcal R)\\
&\le
\left(\sup_{\mathcal R\in\mathscr R_{B,H}^{\rm hard}}
\kappa_{\rm aim,H}(\mathcal R)\right)
\left(\sup_{\mathcal R\in\mathscr R_{B,H}^{\rm hard}}
U_{{\rm gate},{\rm dep}}^{\rm clean}(\mathcal R)\right),
\end{aligned}
\end{equation}
where, record by record,
\[
U_{{\rm gate},{\rm dep}}^{\rm clean}(\mathcal R)
\le
\min\left\{
1,
\left(\sup_{0\le t\le H}W_t(\mathcal R)\right)
U_{\rm noise}^{\rm clean}(n,n^C,m,\eta,\delta_{\rm stat})
\right\}.
\]

Thus \(n\) and the polynomial ceiling enter through the saturated class,
\(m\) through the applicable dynamical generalization certificate, and
\(\eta\) through both that certificate and the exact attenuation factor
\((1-2\eta)^{-1}\).

Now fix one represented parity-output member of this class.  For the output in
\cref{thm-lpn-caus-action-residual-aiming}, use its existing public product
lift carrying the seed \(\omega\) and a uniform candidate \(x\).  On that
carrier, the clean and learned singleton gates are

\[
\mathsf T_0=\{(\omega,x):x=s\},
\qquad
\widehat{\mathsf T}
=\{(\omega,x):x=\widehat s(\omega)\}.
\]

Their exact reference-law error is

\begin{equation}
\label{eq-lpn-learned-target-gate-reference-error}
(\mathbb P_\omega\otimes\mu_I)
(\mathsf T_0\mathbin{\triangle}\widehat{\mathsf T})
=2^{1-n}\Pr_\omega\{\widehat s\ne s\}
\le2^{1-n}p_{\rm err}^{\rm par}.
\end{equation}

If the time-\(t\) deployment law has the represented comparison
\(d_t^\pi\le W_t(\mathbb P_\omega\otimes\mu_I)\), then
\cref{def-controlled-noisy-gate-profiles} gives

\begin{equation}
\label{eq-lpn-learned-target-gate-deployment-error}
\varepsilon_{T,H}^{\rm dep}
\le
\min\left\{
1,
2^{1-n}p_{\rm err}^{\rm par}
\sum_{t=0}^{H}W_t
\right\}.
\end{equation}

Every clocked valid lift in the same record therefore obeys

\begin{equation}
\label{eq-lpn-learned-gate-clocked-lift-transfer}
\left|p_{\mathsf T_0,H}^{\rm clean}
-\widehat p_{\widehat{\mathsf T},H}^{\Omega}\right|
\le
\min\left\{
1,
\varepsilon_{\rm init}
+\sum_{t=0}^{H-1}\|E_t\|
+\min\left\{1,
2^{1-n}p_{\rm err}^{\rm par}\sum_{t=0}^{H}W_t
\right\}
\right\}.
\end{equation}

Let \(\Phi\) be the native symmetric Dirichlet form of a normalized kernel
invariant under the public product law \(\mathbb P_\omega\otimes\mu_I\).
For an exact failure gate and a represented valid-Lift form comparison with
factor \(b_\pi\), the corresponding capacity statement is

\begin{equation}
\label{eq-lpn-learned-gate-lift-capacity-transfer}
\operatorname{Cap}_{\Omega}
(\pi^{-1}\widehat{\mathsf T},\pi^{-1}F)
\le
b_\pi\left[
2^{-n}
+\omega_{\Phi,\mathsf T_0}^{\rm Cap}
 (2^{1-n}p_{\rm err}^{\rm par})
\right].
\end{equation}

Under the exact-identification condition in
\cref{eq-lpn-parity-exact-statistical-identification}, one has
\(\widehat s=s\).  The realized reference and deployment gate errors are then
reset to zero, and the capacity comparison is applied with modulus argument
zero; consequently the gate and capacity-stability contributions in
\cref{eq-lpn-learned-target-gate-reference-error,eq-lpn-learned-target-gate-deployment-error,eq-lpn-learned-gate-clocked-lift-transfer,eq-lpn-learned-gate-lift-capacity-transfer}
then vanish exactly.

\end{corollary}

\begin{proof}

Substitute \(b_C(n)=n^C\) and \(\rho=1-2\eta\) into
\cref{thm-controlled-noise-corrected-caus-aiming}; its binary fallback and
the behavior-to-deployment comparison give
\cref{eq-lpn-hard-gate-caus-noise-radius}.  Two unequal singleton gates on
the uniform candidate carrier differ at exactly two states, while equal
gates have zero symmetric difference.  Averaging over the represented seed
proves \cref{eq-lpn-learned-target-gate-reference-error}; the decoder-error
bound is \cref{eq-lpn-parity-decoder-error-probability}.  The deployment
comparison in \cref{def-controlled-noisy-gate-profiles} proves
\cref{eq-lpn-learned-target-gate-deployment-error}.  Apply
\cref{thm-controlled-noisy-gate-lift-transfer} to obtain the clocked and
capacity inequalities, using
\cref{eq-controlled-target-condenser-flux-upper} on the product-law-invariant
base carrier for the clean singleton-capacity term
\((\mathbb P_\omega\otimes\mu_I)(\mathsf T_0)=2^{-n}\).
Exact parity identification makes the two gates identical; applying the
same theorem with realized gate error zero and modulus argument zero proves
the last assertion.

\end{proof}

\section{Hypercube Functional Geometry and Learned Operators}
\label{sec-lpn-hypercube-functional-learned-operators}

\subsection{Exact hypercube spectrum and functional outputs}
\label{subsec-lpn-hypercube-spectrum-functional-outputs}

\begin{theorem}[Exact spectrum and functional constants of the public hypercube]
\label{thm-lpn-hypercube-spectrum-functional-constants}

Let \(P_\square\) be the public coordinate-flip kernel

\begin{equation}
\label{eq-lpn-public-hypercube-kernel}
P_\square(x,x+e_i)=\frac1n,
\qquad 1\le i\le n,
\end{equation}

with uniform invariant law, and put \(H_\square=I-P_\square\).  Its Walsh
character \(\chi_S\) has eigenvalue \(2|S|/n\).  Consequently the audited
functional-inequality profile of
\cref{def-controlled-functional-inequality-profile} admits the following
simultaneous values.  All decay statements associated with these constants
refer to the continuous-time semigroup \(e^{-tH_\square}\); the profile does
not assert discrete-time convergence of the periodic kernel \(P_\square\).

\begin{align}
\label{eq-lpn-hypercube-gap-cheeger}
q_\square&=1,
&\mathsf h_\square&=\frac1n,
&\gamma_\square&=\frac2n,\\
\label{eq-lpn-hypercube-weak-super-poincare}
\beta_{\mathrm W}(r)&=\frac n2,
&
\beta_{\mathrm S}(r)&=
\begin{cases}
2^n,&0<r<n/2,\\
1,&r\ge n/2,
\end{cases}\\
\label{eq-lpn-hypercube-nash-sobolev}
d_{\mathrm N}&=n,
&C_{\mathrm N}&=2n,
&C_{\mathrm{Sob}}&=2n\quad(n>2),\\
\label{eq-lpn-hypercube-lsi-curvature}
\alpha_{\mathrm{LS}}&=\frac2n,
&\alpha_{\mathrm{MLS}}&=\frac2n,
&CD\!\left(\frac2n,\infty\right)&\ \text{holds},\\
\label{eq-lpn-hypercube-transport-sector}
C_{T_1}&=\frac n4,
&C_{\mathrm{sec}}&=0,
&(m_{\mathscr H},M_{\mathscr H},\kappa_{\mathscr H})
&=(1,1,2/n).
\end{align}

In particular, for \(S_t^\square=e^{-tH_\square}\),

\begin{equation}
\label{eq-lpn-hypercube-continuous-variance-decay}
\operatorname{Var}_{\mu_I}(S_t^\square f)
\le e^{-4t/n}\operatorname{Var}_{\mu_I}(f).
\end{equation}

\end{theorem}

\begin{proof}
Walsh diagonalization gives eigenvalues \(2|S|/n\), hence the gap and
variance decay.  Cube edge isoperimetry gives the Cheeger constant.  Tensor
product Poincar\'e, logarithmic-Sobolev, modified-logarithmic-Sobolev,
curvature, and Hamming concentration inequalities give the remaining exact
or certified constants with the displayed normalization.  Self-adjointness
sets the sector constant to zero and supplies the stated hypocoercive tuple.
\end{proof}

\begin{theorem}[Typed functional outputs of the public hypercube]
\label{thm-lpn-hypercube-typed-functional-outputs}
Adopt the carrier and constants of
\cref{thm-lpn-hypercube-spectrum-functional-constants}.

The remaining functional outputs are obtained from the same audited record.
For \(0<r\le1\), weak Poincar\'e gives

\begin{equation}
\label{eq-lpn-hypercube-weak-poincare-output}
\operatorname{Var}_{\mu_I}(S_t^\square f)
\le
e^{-4t/n}\operatorname{Var}_{\mu_I}(f)
+r\bigl(1-e^{-4t/n}\bigr)\operatorname{Osc}(f)^2.
\end{equation}

For every \(r>0\), the super-Poincar\'e output is

\begin{equation}
\label{eq-lpn-hypercube-super-poincare-output}
\|S_t^\square f\|_{2,\mu_I}^2
\le
e^{-2t/r}\|f\|_{2,\mu_I}^2
+\beta_{\mathrm S}(r)\bigl(1-e^{-2t/r}\bigr)
 \|f\|_{1,\mu_I}^2,
\end{equation}

where \(\beta_{\mathrm S}\) is the explicit function in
\cref{eq-lpn-hypercube-weak-super-poincare}.  On the declared mean-zero
Nash/Sobolev domain \(\mathscr D_{\mathrm N}\),

\begin{align}
\label{eq-lpn-hypercube-nash-smoothing-output}
\bigl\|S_t^\square|_{\mathscr D_{\mathrm N}}\bigr\|_{1\to2}^2
&\le\left(\frac{n^2}{2t}\right)^{n/2},\\
\bigl\|S_t^\square|_{\mathscr D_{\mathrm N}}\bigr\|_{1\to\infty}
&\le\left(\frac{n^2}{t}\right)^{n/2}.
\end{align}

For every represented proper set \(A\subsetneq X_I\), the killed semigroup
satisfies

\begin{equation}
\label{eq-lpn-hypercube-killed-output}
\|S_t^{\square,A}\|_{2\to2}
\le
\min\left\{
\exp\!\left[-\frac{2t}{n}\bigl(1-\mu_I(A)\bigr)\right],
\exp\!\left[-t c_{\mathrm{FK},\square}
                  \mu_I(A)^{-2/n}\right]
\right\}.
\end{equation}

The logarithmic-Sobolev, modified-logarithmic-Sobolev, and curvature
coordinates give, respectively,

\begin{align}
\label{eq-lpn-hypercube-hypercontractive-output}
\|S_t^\square f\|_{q(t),\mu_I}
&\le\|f\|_{p,\mu_I},
&q(t)-1&=(p-1)e^{4t/n},\qquad p>1,\\
\label{eq-lpn-hypercube-entropy-output}
\operatorname{Ent}_{\mu_I}(S_t^{\square,*}g)
&\le e^{-4t/n}\operatorname{Ent}_{\mu_I}(g),\\
\label{eq-lpn-hypercube-gradient-output}
\Gamma_\square(S_t^\square f)
&\le e^{-4t/n}S_t^\square\Gamma_\square(f).
\end{align}

If \(F\) is represented and \(L_F\)-Lipschitz for Hamming distance, then

\begin{equation}
\label{eq-lpn-hypercube-transport-output}
|\nu(F)-\mu_I(F)|
\le L_F\sqrt{\frac n2\operatorname{KL}(\nu\Vert\mu_I)}.
\end{equation}

Finally, on the centered hypocoercive domain,

\begin{equation}
\label{eq-lpn-hypercube-hypocoercive-output}
\|S_t^\square f\|_{2,\mu_I}^2
\le e^{-4t/n}\|f\|_{2,\mu_I}^2.
\end{equation}

For each typed output slot \(s_0\in\mathscr S_{\rm FI}\), let
\(\mathcal J_{s_0,\mathrm{LPN}}\) be the applicable members of these
audited hypercube bounds after conversion to that slot's units.  The
slotwise LPN certificate is therefore

\begin{equation}
\label{eq-lpn-hypercube-functional-slot-minimum}
C_{s_0,\mathrm{LPN}}^{\rm FI}
=\min_{j\in\mathcal J_{s_0,\mathrm{LPN}}}C_{s_0,j},
\end{equation}

with every induced entry retaining the single hypercube-profile provenance.

\end{theorem}

\begin{proof}
Apply the corresponding consequence in
\cref{thm-controlled-functional-inequality-consequences} to each constant in
\cref{thm-lpn-hypercube-spectrum-functional-constants}.  The killed bound
is the minimum of the spectral and Faber--Krahn routes.  All outputs are
converted to a common slot before their minimum is taken, so their single
profile provenance is retained.
\end{proof}

\subsection{Capacity, effective dimension, and carrier profile}
\label{subsec-lpn-hypercube-capacity-effective-dimension}

\begin{theorem}[Target capacity, effective dimension, and residual geometry of the hypercube]
\label{thm-lpn-hypercube-target-effective-dimension}
Adopt the public hypercube carrier of
\cref{thm-lpn-hypercube-spectrum-functional-constants}.

The kernel is self-adjoint, so its selected generator has zero antisymmetric
part and its native actual current has zero authorized Hodge component.
Consequently, for every target-normalized represented potential \(D\),

\begin{equation}
\label{eq-lpn-hypercube-zero-native-caus-alignment}
\operatorname{Align}_{\Caus}^{\mathrm{auth}}
(\mathfrak K_\square,D)=0
\end{equation}

when \(\mathfrak K_\square\) is the actual-current record of this selected
kernel.  This is the exact separation between its inverse-polynomial mixing
profile and its target-oriented current coordinate.

Here \(T_1\) uses ordinary Hamming distance.  For every represented
\(A\subsetneq X_I\),

\begin{align}
\label{eq-lpn-hypercube-spectral-dirichlet-profile}
\Lambda_\square^{\mathrm{mix}}(v)&\ge\frac2n,
&
\lambda_\square^{\mathrm D}(A)
&\ge\frac2n(1-\mu_I(A)),\\
\label{eq-lpn-hypercube-faber-krahn}
\lambda_\square^{\mathrm D}(A)
&\ge c_{\mathrm{FK},\square}\mu_I(A)^{-2/n},
&
c_{\mathrm{FK},\square}
&=\frac{2}{n2^n}(1-2^{-n})^{2/n}.
\end{align}

The target condenser and effective dimension are exact:

\begin{align}
\label{eq-lpn-hypercube-target-capacity}
\operatorname{Cap}_\square(\{s\},X_I\setminus\{s\})
&=2^{-n},\\
\label{eq-lpn-hypercube-effective-dimension}
N_{\mathrm{eff},\square}(\lambda)
&=\sum_{k=1}^n\binom nk
\frac{2k/n}{2k/n+\lambda}.
\end{align}

For \(0<\lambda\le1\),

\begin{equation}
\label{eq-lpn-hypercube-effective-dimension-lower}
N_{\mathrm{eff},\square}(\lambda)\ge2^{n-2}.
\end{equation}

Let

\[
f_s
=\frac{\mathbf1_{\{s\}}-2^{-n}}
{\sqrt{(2^n-1)/2^{2n}}}
\]

be the normalized centered singleton contrast and let
\(E_{\ge\lambda}=\mathbf1_{[\lambda,\infty)}(H_\square)\).  Then

\begin{equation}
\label{eq-lpn-hypercube-fast-target-mass}
\|E_{\ge\lambda}f_s\|_{2,\mu_I}^2
=\frac1{2^n-1}
\sum_{\substack{1\le k\le n\\2k/n\ge\lambda}}\binom nk.
\end{equation}

This is the equality case of the fast-mode target-reach accounting in
\cref{thm-controlled-fast-mode-target-reach}: a constant fraction of the
singleton target lies in fast hypercube modes only together with
exponentially many active modes.

For the public verifier potential \(D_V=1-V_I\) on the verifier-good event,
the same kernel has the exact local drift

\begin{equation}
\label{eq-lpn-hypercube-verifier-local-drift}
(I-P_\square)D_V(x)
=
\begin{cases}
-1,&x=s,\\
1/n,&d_H(x,s)=1,\\
0,&d_H(x,s)\ge2.
\end{cases}
\end{equation}

Hence its global off-target drift margin is zero.  More generally, if
\(W_I-\mu_I(W_I)=\sum_{v\ne0}\widehat W_I(v)\chi_v\), then

\begin{align}
\label{eq-lpn-hypercube-residual-walsh-energy}
\operatorname{Var}_{\mu_I}(W_I)
&=\sum_{v\ne0}\widehat W_I(v)^2,\\
\mathcal E_\square(W_I,W_I)
&=\sum_{v\ne0}\frac{2|v|_1}{n}\widehat W_I(v)^2,
\end{align}

so every nonconstant residual potential satisfies

\begin{equation}
\label{eq-lpn-hypercube-residual-regularity-range}
\frac2n
\le
\frac{\mathcal E_\square(W_I,W_I)}
{\operatorname{Var}_{\mu_I}(W_I)}
\le2.
\end{equation}

\end{theorem}

\begin{proof}
Self-adjointness removes native directed current.  The Walsh spectrum gives
the Dirichlet profile, effective-dimension sum, singleton spectral mass, and
residual regularity range.  The opposite-plate capacity equals the boundary
flux of the singleton and is \(2^{-n}\).  Pairing Walsh levels
\(k\ge n\lambda/2\) with their binomial multiplicities proves the effective
dimension and fast-mode assertions; direct neighbor counting gives the
verifier drift.
\end{proof}

\begin{corollary}[Complete functional profile of the public hypercube carrier]
\label{cor-lpn-hypercube-functional-profile}
\label{thm-lpn-hypercube-functional-profile}
The exact constants, typed functional outputs, and target/effective-dimension
geometry are
\cref{thm-lpn-hypercube-spectrum-functional-constants,thm-lpn-hypercube-typed-functional-outputs,thm-lpn-hypercube-target-effective-dimension}.
They are simultaneous outputs of one audited public-carrier record.
\end{corollary}

\begin{proof}

The Walsh diagonalization follows from
\(\chi_S(x+e_i)=(-1)^{\mathbf1_{\{i\in S\}}}\chi_S(x)\).
It gives \(q_\square=1\) and the gap \(2/n\).  The edge-isoperimetric
inequality on the cube gives \(\mathsf h_\square=1/n\).  Poincar\'e gives
the constant weak-Poincar\'e rate.  Also

\[
\|f\|_2^2\le2^n\|f\|_1^2,
\qquad
\|f\|_2^2\le\frac n2\mathcal E_\square(f,f)+\|f\|_1^2,
\]

which proves the displayed super-Poincar\'e rate.  On the mean-zero domain,
\(\|f\|_2\le2^{n/2}\|f\|_1\) and Poincar\'e give

\[
\|f\|_2^{2+4/n}
\le2n\,\mathcal E_\square(f,f)\|f\|_1^{4/n}.
\]

For \(n>2\), finite-space norm comparison gives
\(\|f\|_{2n/(n-2)}^2\le4\|f\|_2^2\le2n\mathcal E_\square(f,f)\)
on the same mean-zero domain.  This proves
\cref{eq-lpn-hypercube-nash-sobolev}.

Tensorization of the sharp two-point inequalities gives
\cref{eq-lpn-hypercube-lsi-curvature}; tensorization of the bounded-difference
subgaussian estimate gives \(T_1(n/4)\) for Hamming distance
\citep{Gross1975,BobkovGotze1999}.  The kernel is reversible, so the sector
constant is zero and the Poincar\'e energy itself supplies the displayed
hypocoercive triple.  The Poincar\'e consequence in
\cref{thm-controlled-functional-inequality-consequences} gives
\cref{eq-lpn-hypercube-continuous-variance-decay}.  Substituting the same
constants into the weak-Poincar\'e, super-Poincar\'e, Nash, killed-semigroup,
logarithmic-Sobolev, modified-logarithmic-Sobolev, curvature, transport, and
hypocoercive conclusions of that theorem proves
\cref{eq-lpn-hypercube-weak-poincare-output,eq-lpn-hypercube-super-poincare-output,eq-lpn-hypercube-nash-smoothing-output,eq-lpn-hypercube-killed-output,eq-lpn-hypercube-hypercontractive-output,eq-lpn-hypercube-entropy-output,eq-lpn-hypercube-gradient-output,eq-lpn-hypercube-transport-output,eq-lpn-hypercube-hypocoercive-output}.
The slotwise minimum is
\cref{thm-controlled-functional-profile-selection}.  Self-adjointness makes
the antisymmetric generator component zero; hence
\cref{prop-controlled-native-geometry-hodge,def-quantitative-caus-alignment}
give \cref{eq-lpn-hypercube-zero-native-caus-alignment}.

Poincar\'e gives \(\Lambda_\square^{\mathrm{mix}}(v)\ge2/n\).  If \(f\)
is supported in \(A\), Cauchy--Schwarz gives
\(\operatorname{Var}(f)\ge(1-\mu_I(A))\|f\|_2^2\); this proves the
Dirichlet estimate.  Since a proper subset has mass in
\([2^{-n},1-2^{-n}]\), minimizing
\((2/n)(1-v)v^{2/n}\) over those allowed masses gives the stated
Faber--Krahn minorant.

The capacity identity is
\cref{eq-lpn-singleton-capacity-exact} with \(q_{P_\square,s}=1\).
Summing the Walsh eigenvalues through the definition of effective dimension
proves \cref{eq-lpn-hypercube-effective-dimension}.  For
\(0<\lambda\le1\), all modes of weight at least \(n/2\) contribute at least
one half, proving \cref{eq-lpn-hypercube-effective-dimension-lower}.
Every nonconstant Walsh coefficient of the normalized singleton contrast has
squared magnitude \((2^n-1)^{-1}\), which proves
\cref{eq-lpn-hypercube-fast-target-mass}.

On the verifier-good event, \(D_V(s)=0\) and \(D_V=1\) elsewhere.
One-step substitution proves
\cref{eq-lpn-hypercube-verifier-local-drift}.  Finally, Parseval's identity
and the Walsh eigenvalues give
\cref{eq-lpn-hypercube-residual-walsh-energy}; their nonzero range
\([2/n,2]\) gives
\cref{eq-lpn-hypercube-residual-regularity-range}.

\end{proof}

\Cref{fig-lpn-hypercube-walsh-levels} summarizes the multiplicity calculation
that accompanies fast singleton modes on the public cube.

\begin{figure}[tbp]
\centering
\begin{tikzpicture}[fw diagram,x=1cm,y=1cm]
  \draw[->,draw=fwInk] (0,0)--(10.7,0) node[right,font=\scriptsize] {Walsh level \(k\)};
  \draw[->,draw=fwInk] (0,0)--(0,4.2) node[above,font=\scriptsize] {multiplicity};
  \foreach \x/\h in {1/.45,2/1.0,3/1.9,4/3.0,5/3.7,6/3.7,7/3.0,8/1.9,9/1.0,10/.45}
    \path[fill=fwBlue!28,draw=fwBlue] (\x-.34,0) rectangle (\x+.34,\h);
  \node[font=\scriptsize,align=center] at (5.5,4.15)
    {level multiplicity \(\binom nk\)};
  \draw[draw=fwCoral,line width=1pt,dashed] (3.5,0)--(3.5,3.8);
  \node[font=\scriptsize,fwCoral,anchor=west] at (3.6,3.45)
    {fast cutoff \(2k/n\ge\lambda\)};
  \node[fw source,text width=2.7cm] (singleton) at (1.4,-1.5)
    {singleton contrast\\mass at every\\Walsh level};
  \node[fw geom,text width=3.2cm] (neff) at (5.4,-1.5)
    {\(N_{\rm eff}(\lambda)=\sum_k\binom nk\frac{2k/n}{2k/n+\lambda}\)};
  \node[fw terminal,text width=2.6cm] (cap) at (9.4,-1.5)
    {capacity \(2^{-n}\)\\and \(N_{\rm eff}\ge2^{n-2}\)};
  \draw[fw arrow] (singleton)--(neff);
  \draw[fw arrow] (neff)--(cap);
\end{tikzpicture}
\caption{Walsh levels of the public hypercube.  Eigenvalue \(2k/n\) grows
with level, while multiplicity is binomial.  The centered singleton spreads
across all nonconstant levels, so placing a constant fraction in fast modes
requires exponentially many active modes even though the singleton
condenser capacity is only \(2^{-n}\).  Bar heights are schematic.}
\label{fig-lpn-hypercube-walsh-levels}
\end{figure}

\subsection{Nonlinear reweighting and learned-operator transfer}
\label{subsec-lpn-reweighting-learned-operator-transfer}

\begin{corollary}[Explicit nonlinear residual reweighting]
\label{cor-lpn-residual-metropolis-capacity-mixing}

Fix \(\theta\ge0\) and define the represented lazy local kernel

\begin{align}
\label{eq-lpn-residual-metropolis-kernel}
P_\theta(x,x+e_i)
&=\frac1{2n}
\min\{1,e^{-\theta(W_I(x+e_i)-W_I(x))}\},\\
P_\theta(x,x)
&=1-\sum_{i=1}^nP_\theta(x,x+e_i).
\end{align}

It is reversible for
\(\pi_\theta(x)=Z_\theta^{-1}e^{-\theta W_I(x)}\).  When this kernel is
retained as a certified dynamical geometry, the same record contains the
represented derivation of \(Z_\theta\), the invariant-law certificate, its
declared precision, and their complete cost, as required by
\cref{def-certified-dynamical-geometry}.  On
\(\mathcal F_\delta\), put
\(\Delta_-=1/2-\eta-2\delta>0\) and
\(\Delta_+=1/2-\eta+2\delta\).  Then

\begin{equation}
\label{eq-lpn-metropolis-target-mass-window}
\frac1{1+(2^n-1)e^{-\theta m\Delta_-}}
\le\pi_\theta(s)
\le
\frac1{1+(2^n-1)e^{-\theta m\Delta_+}},
\end{equation}

and

\begin{align}
\label{eq-lpn-metropolis-target-exit-capacity}
q_{\theta,s}
&\le\frac12e^{-\theta m\Delta_-},
&
\operatorname{Cap}_\theta(\{s\},X_I\setminus\{s\})
&\le\frac12\pi_\theta(s)e^{-\theta m\Delta_-}.
\end{align}

Its Poincar\'e gap obeys the explicit two-regime bound

\begin{equation}
\label{eq-lpn-metropolis-gap-upper}
\gamma_\theta
\le
\begin{cases}
e^{-\theta m\Delta_-},&\pi_\theta(s)\le1/2,\\
e^{4\theta\delta m}/(2^n-1),&\pi_\theta(s)>1/2.
\end{cases}
\end{equation}

In the normalization of
\cref{def-controlled-functional-inequality-profile}, every positive
logarithmic-Sobolev, modified-logarithmic-Sobolev, or
Bakry--\'Emery curvature rate certified for this same kernel is bounded
above by the right side of \cref{eq-lpn-metropolis-gap-upper}.  Each such
rate implies the Poincar\'e inequality at its own rate, whereas
\cref{eq-lpn-metropolis-gap-upper} is an upper bound for the optimal
Poincar\'e rate.

Thus this compact nonlinear reweighting is evaluated simultaneously by its
target mass, target-boundary capacity, and mixing rate.

\end{corollary}

\begin{proof}

Detailed balance follows from
\[
e^{-\theta W_I(x)}
\min\{1,e^{-\theta(W_I(z)-W_I(x))}\}
=e^{-\theta W_I(z)}
\min\{1,e^{-\theta(W_I(x)-W_I(z))}\}.
\]
On \(\mathcal F_\delta\), every off-target state satisfies
\(m\Delta_-\le W_I(x)-W_I(s)\le m\Delta_+\).  Summing the corresponding
Boltzmann ratios proves
\cref{eq-lpn-metropolis-target-mass-window}.  Every neighbor of \(s\) is
off-target and has residual increase at least \(m\Delta_-\), so summing the
\(n\) proposal rates proves
\cref{eq-lpn-metropolis-target-exit-capacity} by
\cref{eq-lpn-singleton-capacity-exact}.

If \(\pi_\theta(s)\le1/2\), the singleton conductance ratio is at most
\(q_{\theta,s}\); the upper Cheeger inequality gives
\(\gamma_\theta\le2q_{\theta,s}\).  If
\(\pi_\theta(s)>1/2\), use the complementary set.  Since

\[
1-\pi_\theta(s)
\ge\pi_\theta(s)(2^n-1)e^{-\theta m\Delta_+},
\]

its conductance ratio is at most
\(e^{\theta m(\Delta_+-\Delta_-)}/(2(2^n-1))\).  Applying the same upper
Cheeger inequality and using \(\Delta_+-\Delta_-=4\delta\) proves
\cref{eq-lpn-metropolis-gap-upper}.

\end{proof}

\begin{theorem}[Learned-operator transfer to LPN target capacity]
\label{thm-lpn-learned-operator-capacity-transfer}

Let \(H_\Phi\) be the native symmetric operator of a uniform-invariant LPN
kernel and let \(\widehat H_\Phi\ge0\) be a represented learned operator on
the same \(L^2(\mu_I)\) carrier.  Suppose the complete dynamical learning
record certifies

\begin{equation}
\label{eq-lpn-learned-operator-error}
\|H_\Phi-\widehat H_\Phi\|_{\mathrm{op}}
\le\varepsilon_\Phi.
\end{equation}

Define the learned target-boundary rate

\begin{equation}
\label{eq-lpn-learned-target-boundary-rate}
\widehat q_{\Phi,s}
=2^n\langle\mathbf1_{\{s\}},
\widehat H_\Phi\mathbf1_{\{s\}}\rangle_{\mu_I}.
\end{equation}

Then

\begin{align}
\label{eq-lpn-learned-capacity-transfer}
\operatorname{Cap}_\Phi(\{s\},X_I\setminus\{s\})
&\le2^{-n}(\widehat q_{\Phi,s}+\varepsilon_\Phi),\\
\label{eq-lpn-learned-capacity-hit-transfer}
\Pr_{\mu_0}\{\tau_s\le H\}
&\le R_0 2^{-n}
\bigl[1+H(\widehat q_{\Phi,s}+\varepsilon_\Phi)\bigr]
\end{align}

whenever \(d\mu_0/d\mu_I\le R_0\).  On the simultaneous event of
\cref{thm-master-structural-dynamical-learning-certificate}, one may take

\begin{equation}
\label{eq-lpn-learned-capacity-generalization-error}
\varepsilon_\Phi
=\widehat\varepsilon_\Phi
+\mathfrak G_{\mathrm{dyn}}
(\mathcal F_{\mathrm{geom}};B,m,\delta_\Phi).
\end{equation}

The minimum defining \(\mathfrak G_{\mathrm{dyn}}\) retains, in common
population-error units, the mixing-block estimate of
\cref{thm-beta-mixing-structural-blocking}, the predictable learned-operator
estimate of \cref{thm-sequential-learned-operator-cover}, and the martingale
PAC--Bayes estimate of \cref{thm-martingale-structural-pac-bayes}.  If a
certified eigengap \(g_*>0\) is also present, the learned fast-mode target
mass differs from the population fast-mode target mass by at most

\begin{equation}
\label{eq-lpn-learned-fast-mode-target-transfer}
\frac{2\varepsilon_\Phi}{g_*},
\end{equation}

in accordance with
\cref{thm-controlled-fast-mode-target-reach,eq-master-dyn-geometry}.

The appearances of \(s\) in
\cref{eq-lpn-learned-target-boundary-rate,eq-lpn-learned-fast-mode-target-transfer}
are analytic target audits and do not expose \(s\) to the learned operator.
When the learned dynamical record outputs a represented candidate
\(\widehat s\), append the parity readout
\(h_{\widehat s}(a)=(-1)^{\langle a,\widehat s\rangle}\).  Then
\cref{eq-lpn-recovery-is-parity-alignment,eq-lpn-randomized-recovery-alignment}
convert its recovery probability into clean parity alignment, and
\cref{eq-lpn-uniform-statistical-alignment} transports that alignment
through the LPN noise factor \(1-2\eta\).  The appended evaluator and its
complete cost remain in the same saturated learning record.

\end{theorem}

\begin{proof}

Put \(v_s=\mathbf1_{\{s\}}\), so
\(\|v_s\|_{2,\mu_I}^2=2^{-n}\).  Since the two condenser plates fill the
carrier, its unique Dirichlet competitor is \(v_s\), and therefore

\[
\operatorname{Cap}_\Phi(\{s\},X_I\setminus\{s\})
=\langle v_s,H_\Phi v_s\rangle
\le\langle v_s,\widehat H_\Phi v_s\rangle
 +\varepsilon_\Phi\|v_s\|_2^2.
\]

Substitution of
\cref{eq-lpn-learned-target-boundary-rate} proves
\cref{eq-lpn-learned-capacity-transfer}.  Apply
\cref{thm-controlled-target-capacity-drift-envelope} to obtain
\cref{eq-lpn-learned-capacity-hit-transfer}.  Equation
\cref{eq-lpn-learned-capacity-generalization-error} is the geometry
coordinate of
\cref{thm-master-structural-dynamical-learning-certificate}.  Finally, the
projector perturbation bound in \cref{eq-master-dyn-geometry}, evaluated on
the unit vector \(f_s\), proves
\cref{eq-lpn-learned-fast-mode-target-transfer}.

\end{proof}

\section{LPN Committors and Dynamical Learning}
\label{sec-lpn-committor-dynamical-learning}

\subsection{Committor and Bellman propagation}
\label{subsec-lpn-committor-bellman-propagation}

\begin{theorem}[LPN committor and same-norm Bellman propagation]
\label{thm-lpn-committor-same-norm-propagation}

Fix \(0<\eta<\tau<1/2\), \(0<\delta<1\), and confidence allocations
\(\delta_M,\delta_B,\delta_\Phi,\delta_{\rm gen},\delta_{\rm PB}>0\)
whose sum is at most \(\delta\).  Work on the verifier-good event of
\cref{thm-lpn-public-verifier}.  Let \(F\subseteq X_I\setminus\{s\}\) be a
represented failure sector, put
\(\mathsf N_I=X_I\setminus(\{s\}\cup F)\), and let \(P^\pi\) be a selected
LPN policy kernel.  Here and below \(u_\beta^*\) is the fixed point of the
optimal committor Bellman operator restricted to the same represented
budget-reachable policy class \(\Pi_{B,{\rm obs}}\), or \(\Pi_B\) in the
fully observed case, used in
\cref{eq-budget-reachable-control-value,thm-master-structural-dynamical-learning-certificate};
it is the reach--avoid specialization of \(V_{B,{\rm obs}}^*\).  The target
gate is then
\(V_I=\mathbf1_{\{s\}}\), and the finite-horizon and discounted committors
satisfy

\begin{align}
\label{eq-lpn-controlled-committor-recursion}
u_0^\pi
&=V_I,
&
u_{h+1}^\pi
&=V_I+\mathbf1_{\mathsf N_I}P^\pi u_h^\pi,\\
\label{eq-lpn-controlled-discounted-committor}
u_\beta^\pi
&=V_I+\beta\mathbf1_{\mathsf N_I}P^\pi u_\beta^\pi,
\qquad 0<\beta<1.
\end{align}

For every represented approximation \(v\),

\begin{equation}
\label{eq-lpn-controlled-committor-residual}
\|u_\beta^\pi-v\|_\infty
\le
\frac{
\|V_I+\beta\mathbf1_{\mathsf N_I}P^\pi v-v\|_\infty
}{1-\beta}.
\end{equation}

If a complete learned LPN control record supplies same-norm model and Bellman
errors \(\overline\varepsilon_M,\overline\varepsilon_B\) as in
\cref{thm-master-structural-dynamical-learning-certificate}, then

\begin{align}
\label{eq-lpn-learned-committor-value-propagation}
\|u_\beta^*-v\|_\infty
&\le
\frac{\overline\varepsilon_B+\overline\varepsilon_M}{1-\beta},\\
\label{eq-lpn-learned-committor-policy-propagation}
\|u_\beta^*-u_\beta^{\widehat\pi_v}\|_\infty
&\le
\frac{2\beta(\overline\varepsilon_B+\overline\varepsilon_M)}
     {(1-\beta)^2}
+\frac{2\overline\varepsilon_M}{1-\beta}.
\end{align}

The errors in these displays retain the exact three same-norm coordinates of
\cref{thm-master-structural-dynamical-learning-certificate}:

\begin{align}
\label{eq-lpn-master-same-norm-error-specialization}
\overline\varepsilon_M
&=
\min\left\{
\widehat\varepsilon_M+
\mathfrak G_{\rm dyn}(\mathcal F_{\rm model};B,m,\delta_M),
U_M^{\rm noise},U_M^{\rm ch}
\right\},\\
\overline\varepsilon_B
&=
\min\left\{
\widehat\varepsilon_B+
\mathfrak G_{\rm dyn}(\mathcal F_{\rm Bell};B,m,\delta_B),
U_B^{\rm noise},U_B^{\rm ch}
\right\},\\
\overline\varepsilon_\Phi
&=
\min\left\{
\widehat\varepsilon_\Phi+
\mathfrak G_{\rm dyn}(\mathcal F_{\rm geom};B,m,\delta_\Phi),
U_\Phi^{\rm noise},U_\Phi^{\rm ch}
\right\}.
\end{align}

Here each noise candidate is finite exactly when the complete record contains
the corresponding represented same-norm loss-to-coordinate calibration, and
each channel candidate is the smallest same-law, same-norm consequence of the
typed channel profiles and represented converters listed in
\cref{thm-master-structural-dynamical-learning-certificate}; an unavailable
candidate is \(+\infty\).  When a concrete generalization or PAC--Bayes
candidate is used in one coordinate, its \(\delta_{\rm gen}\) or
\(\delta_{\rm PB}\) allocation is charged within that coordinate's declared
\(\delta_M,\delta_B,\delta_\Phi\).  Data-dependent selection among candidates
uses the simultaneous confidence split required by
\cref{def-dynamical-generalization-envelope}.

\end{theorem}

\begin{proof}
The verifier event identifies the target gate with the singleton.  The
controlled Bellman identities give the committor recursions and fixed-policy
residual bound.  Applying the master same-norm propagation theorem gives the
learned value and policy inequalities; the simultaneous leaf event justifies
each displayed minimum.
\end{proof}

\subsection{Safe spectra and dependent-data candidates}
\label{subsec-lpn-safe-spectra-dependent-data}

\begin{theorem}[Action-safe LPN spectrum and kernel-ridge error]
\label{thm-lpn-action-safe-spectrum-krr}
Adopt the verifier-good presentation and public hypercube geometry of
\cref{thm-lpn-committor-same-norm-propagation,thm-lpn-hypercube-functional-profile}.

In the further target-only specialization \(F=\varnothing\), the legal
coordinate actions generate the full mean-zero Walsh space and there is no
failure killing.
Consequently

\begin{equation}
\label{eq-lpn-action-safe-projections}
P_{\mathcal A}=P_{\rm safe}=I
\quad\text{on }\mathbf1^\perp,
\end{equation}

and the action-safe effective dimension is exact:

\begin{equation}
\label{eq-lpn-action-safe-effective-dimension}
N_{\rm eff}^{\mathcal A,\mathrm{safe}}(\lambda)
=
\sum_{k=1}^{n}\binom nk
\frac{2k/n}{2k/n+\lambda}
\ge2^{n-2},
\qquad 0<\lambda\le1.
\end{equation}

For the normalized singleton contrast \(f_s\) in
\cref{eq-lpn-hypercube-fast-target-mass}, the exact action-safe ridge
residual is

\begin{equation}
\label{eq-lpn-action-safe-singleton-ridge-residual}
\left\|
\left[I-H_\square(H_\square+\lambda I)^{-1}\right]f_s
\right\|_{2,\mu_I}^2
=
\frac1{2^n-1}
\sum_{k=1}^{n}\binom nk
\left(\frac{\lambda}{2k/n+\lambda}\right)^2.
\end{equation}

On the row-input hypercube used for noisy parity learning, the clean parity
\(\chi_s\) is a Walsh eigenfunction.  Its corresponding residual is

\begin{equation}
\label{eq-lpn-action-safe-parity-ridge-residual}
\left\|
\left[I-H_\square(H_\square+\lambda I)^{-1}\right]\chi_s
\right\|_{2,\mu_I}^2
=
\left(\frac{\lambda}{2|s|_1/n+\lambda}\right)^2.
\end{equation}

For \(s\ne0\), this is the action-safe compressed residual.  For \(s=0\),
\(\chi_0=\mathbf1\) lies in the constant orthogonal complement of the
action-generated mean-zero range, and the same display retains the exact
approximation residual
\(\|(I-P_{\mathcal A})\chi_0\|_{2,\mu_I}^2=1\), as prescribed by
\cref{eq-constrained-safe-source-residual}.

For independent LPN rows, the debiased dynamical kernel-ridge estimator
\(\widehat f_\lambda\) of
\cref{eq-dynamical-label-noise-krr-bound}, applied to this row-input
hypercube and the clean parity \(\chi_s\), satisfies

\begin{align}
\label{eq-lpn-action-safe-parity-krr-bound}
\mathbb E\|\widehat f_\lambda-\chi_s\|_{2,\mu_I}^2
\le{}&
C\left[
\frac{v_\eta^2}{(1-2\eta)^2m}
\sum_{k=1}^{n}\binom nk
\frac{2k/n}{2k/n+\lambda}
+\left(\frac{\lambda}{2|s|_1/n+\lambda}\right)^2
\right]\notag\\
&+\frac{\varepsilon_{\rm sel}}{(1-2\eta)^2},
\end{align}

whenever the complete record contains the conditional sub-Gaussian variance,
covariance-concentration, effective-sample-size, and simultaneous selection
certificates required by that framework bound.  For a dependent record, the
same display replaces \(m\) by its certified \(m_{\rm dyn}\) and restores the
represented \(\varepsilon_{\rm dep}\) term.

\end{theorem}

\begin{proof}
Coordinate actions span every nonconstant Walsh character, so both safe
projections are the identity on \(\mathbf1^\perp\).  Walsh diagonalization
gives the effective dimension and singleton/parity ridge residuals.
Substitution into the framework noise-corrected kernel-ridge theorem gives
the independent-row bound; the dependent-data statement uses its certified
effective sample size and dependence defect.
\end{proof}

\begin{theorem}[LPN dependent-data generalization candidates]
\label{thm-lpn-dependent-generalization-candidates}
Adopt the complete represented LPN learning record and confidence allocation
of \cref{thm-lpn-committor-same-norm-propagation}.

The dynamical generalization candidates also specialize without changing
their sampling regimes.  For a length-\(m\) absolutely regular LPN record
and a certified block length \(a\), put

\begin{equation}
\label{eq-lpn-dynamical-block-parameters}
\mu_a=\left\lfloor\frac{m}{2a}\right\rfloor,
\qquad
\delta_a=\delta_{\rm gen}-2(\mu_a-1)\beta(a)>0.
\end{equation}

Then the mixing-block candidate is

\begin{equation}
\label{eq-lpn-dynamical-mixing-radius}
\Gamma_{\beta,a}^{\mathrm{LPN}}
=
2\max_{p\in\{\mathrm{odd},\mathrm{even}\}}
\mathfrak R_{\mu_a}^{p}(a)
+3\sqrt{\frac{\log(4/\delta_a)}{2\mu_a}}
+\frac{2a}{m}.
\end{equation}

For \(m\ge2\) independent LPN rows, \(\beta(a)=0\).  Taking \(a=1\) and
the \(2^n\)-element parity class gives the explicit candidate

\begin{equation}
\label{eq-lpn-dynamical-iid-block-radius}
\Gamma_{\beta,1}^{\mathrm{LPN}}
\le
2\sqrt{\frac{2n\log2}{\lfloor m/2\rfloor}}
+3\sqrt{\frac{\log(4/\delta_{\rm gen})}
              {2\lfloor m/2\rfloor}}
+\frac2m.
\end{equation}

For a represented predictable LPN operator family satisfying the hypotheses
of \cref{thm-sequential-learned-operator-cover}, write
\(\eta_{\rm op}>0\) for its pathwise operator lower bound, distinct from the
LPN label-noise rate \(\eta\).  Its sequential-radius candidate is

\begin{equation}
\label{eq-lpn-dynamical-sequential-radius}
\mathfrak R_{m,\mathrm{LPN}}^{\rm seq}
\le
\inf_{\varrho>0}
\left\{
R_A\sup_{\theta\in\Theta}\mathfrak q_m(\mathbf A_\theta)
+R_A\sqrt{\frac{L_A\varrho}{m}}
+\frac{4R_A}{\sqrt{\eta_{\rm op}}}
 \sqrt{\frac{\log\mathcal N(\Theta,d,\varrho)}{m}}
\right\}.
\end{equation}

For the represented LPN prediction class \(\mathcal H_B\) in the same
complete record, the online coordinate specializes to

\begin{equation}
\label{eq-lpn-dynamical-online-regret}
\frac{\mathcal V_m(\mathcal H_B)}m
\le
2L_\ell\sup_{\mathbf z}
\mathfrak R_m^{\rm seq}(\mathcal H_B;\mathbf z).
\end{equation}

Under independent LPN rows, the represented average predictor, or a
uniformly selected represented iterate, has population excess risk at most
the right side of \cref{eq-lpn-dynamical-online-regret}.  For the homogeneous
flip-symmetric binary loss covered by
\cref{eq-dynamical-label-noise-average-transport}, corrupted excess is
converted to clean parity risk by the same factor \((1-2\eta)^{-1}\) used
below.

For the uniform prior \(P_{\rm par}\) on the parity class and a point
posterior \(\delta_t\),
\(\operatorname{KL}(\delta_t\Vert P_{\rm par})=n\log2\).  Hence the
martingale PAC--Bayes candidate is

\begin{equation}
\label{eq-lpn-dynamical-pac-bayes-radius}
\Gamma_{\rm PB}^{\mathrm{LPN}}(\lambda)
=
\frac{n\log2+\log(1/\delta_{\rm PB})}{\lambda m}
+\frac{\lambda}{8},
\end{equation}

and optimization over \(\lambda>0\) gives

\begin{equation}
\label{eq-lpn-dynamical-pac-bayes-optimized}
\inf_{\lambda>0}\Gamma_{\rm PB}^{\mathrm{LPN}}(\lambda)
=
\sqrt{\frac{n\log2+\log(1/\delta_{\rm PB})}{2m}}.
\end{equation}

After conversion to one common bounded population loss, the definition of
\(\mathfrak G_{\rm dyn}\) takes the minimum of the applicable independent,
mixing-block, sequential, and martingale PAC--Bayes candidates above.  It also
admits the effective-dimension candidate exactly when the same complete record
satisfies the source, noise, and covariance hypotheses of
\cref{cor-dynamical-source-condition-rate}.

\end{theorem}

\begin{proof}
Apply the absolute-regularity blocking theorem with the displayed block
count.  Independent rows set the mixing coefficient to zero.  The sequential
cover, online-to-batch, and martingale PAC--Bayes theorems give the other
candidates; optimizing the elementary PAC--Bayes expression gives its closed
form.  Only candidates in the same bounded population-loss units enter the
minimum.
\end{proof}

\subsection{Complete learning transfer and stress tests}
\label{subsec-lpn-learning-transfer-stress-tests}

\begin{corollary}[Clean-risk transport for a complete LPN learning record]
\label{cor-lpn-complete-learning-clean-risk}
Under the hypotheses of
\cref{thm-lpn-committor-same-norm-propagation,thm-lpn-dependent-generalization-candidates},
the homogeneous LPN flip channel
preserves the dependence comparison and transports corrupted excess to clean
risk through

\begin{equation}
\label{eq-lpn-complete-dynamical-clean-risk}
U_{\rm noise}^{\rm clean}(n,B,m,\eta,\delta_{\rm gen})
=
\min\left\{
1,\,
\frac{
\operatorname{Exc}_{\rm corr}^{\rm sat}
(n,B,m,\eta,\delta_{\rm gen})}
{1-2\eta}
\right\}.
\end{equation}

Separately, the represented parity empirical minimizer evaluated in
\cref{ex-lpn-corrupted-label-learning} obeys the recordwise bound

\begin{equation}
\label{eq-lpn-complete-dynamical-parity-risk}
R_0(\widehat f)
\le
\min\left\{
1,\,
\frac{2\Gamma_{n,m,\delta_{\rm gen}}^{\rm par}
      +\varepsilon_{\rm opt}}
     {1-2\eta}
\right\},
\end{equation}

with \(\Gamma_{n,m,\delta_{\rm gen}}^{\rm par}\) given explicitly in
\cref{eq-lpn-parity-rademacher-radius}.

Every committor, operator, posterior, block length, selected policy, and
downstream evaluator in these displays belongs to the one complete saturated
LPN learning record.  The identities create no additional structural source.

\end{corollary}

\begin{proof}
The homogeneous flip channel preserves the dependence comparison and scales
binary excess risk by \(1-2\eta\).  Applying the saturated clean-risk
transport theorem gives the first formula.  The recordwise parity bound is
the uniform parity generalization estimate with the same confidence
allocation, clipped by the unit risk bound.
\end{proof}

\begin{corollary}[LPN committor, safe-spectrum, and dynamical-learning specialization]
\label{cor-lpn-complete-controlled-learning-specialization}
\label{thm-lpn-complete-controlled-learning-specialization}
The committor, action-safe spectral, dependent-generalization, and clean-risk
conclusions are
\cref{thm-lpn-committor-same-norm-propagation,thm-lpn-action-safe-spectrum-krr,thm-lpn-dependent-generalization-candidates,cor-lpn-complete-learning-clean-risk}.
All belong to one complete represented LPN learning record.
\end{corollary}

\begin{proof}

On the verifier-good event,
\cref{thm-lpn-public-verifier} identifies \(V_I\) with the singleton target
gate.  Substitution into
\cref{thm-controlled-committor-bellman} proves
\cref{eq-lpn-controlled-committor-recursion,eq-lpn-controlled-discounted-committor}.
The fixed-policy residual estimate is
\cref{thm-controlled-bellman-residual-policy-loss}; the learned value and
policy estimates are
\cref{thm-dynamical-model-bellman-policy-propagation}.  Substitution of the
master same-norm coordinates proves
\cref{eq-lpn-master-same-norm-error-specialization}.  Temporal admissibility
and complete cost are retained by
\cref{thm-controlled-no-free-value,cor-saturated-learned-policy-accountability}.

For every nonempty Walsh index \(S\), choose \(i\in S\).  The coordinate
difference satisfies
\(\chi_S(\,\cdot+e_i)-\chi_S=-2\chi_S\), so the legal coordinate-action
increments span every nonconstant Walsh mode.  This proves
\cref{eq-lpn-action-safe-projections}.  Apply
\cref{thm-constrained-safe-spectral-capacity} and the eigenvalues
\(2k/n\) from
\cref{thm-lpn-hypercube-functional-profile} to obtain
\cref{eq-lpn-action-safe-effective-dimension}.  The normalized singleton has
squared Walsh coefficient \((2^n-1)^{-1}\) on every nonconstant character;
the spectral multiplier of
\(I-H_\square(H_\square+\lambda I)^{-1}\) is
\(\lambda/(2k/n+\lambda)\).  Parseval proves
\cref{eq-lpn-action-safe-singleton-ridge-residual}.  Applying the same
multiplier to the single parity mode \(\chi_s\) proves
\cref{eq-lpn-action-safe-parity-ridge-residual}.
Substitution of
\cref{eq-lpn-action-safe-effective-dimension,eq-lpn-action-safe-parity-ridge-residual}
into \cref{eq-dynamical-label-noise-krr-bound}, with
\(m_{\rm dyn}=m\) and \(\varepsilon_{\rm dep}=0\) for independent rows,
proves \cref{eq-lpn-action-safe-parity-krr-bound}.

Set the trajectory length \(N=m\) in
\cref{thm-beta-mixing-structural-blocking} to obtain
\cref{eq-lpn-dynamical-block-parameters,eq-lpn-dynamical-mixing-radius}.
Independent rows have zero absolute-regularity coefficients, and the
finite-class Rademacher estimate
\(\mathfrak R_{\mu_1}\le
\sqrt{2n\log2/\mu_1}\) proves
\cref{eq-lpn-dynamical-iid-block-radius}.  Substitution of \(T=m\) and
\(\eta_{\rm op}\) for the operator lower bound in
\cref{thm-sequential-learned-operator-cover} proves
\cref{eq-lpn-dynamical-sequential-radius}.  The uniform parity prior gives
the displayed KL exactly; substitution in
\cref{thm-martingale-structural-pac-bayes} proves
\cref{eq-lpn-dynamical-pac-bayes-radius}.  Minimizing
\(c/\lambda+\lambda/8\), with
\(c=(n\log2+\log(1/\delta_{\rm PB}))/m\), proves
\cref{eq-lpn-dynamical-pac-bayes-optimized}.
The online-regret and online-to-batch conclusions are, respectively,
\cref{thm-structural-online-regret,cor-structural-online-to-batch} with
\(T=m\).  The clean-risk transport for their corrupted binary loss is
\cref{eq-dynamical-label-noise-average-transport}; this proves the statement
following \cref{eq-lpn-dynamical-online-regret}.

The common-loss minimum is exactly
\cref{def-dynamical-generalization-envelope}.
The corruption innovations do not increase absolute regularity by
\cref{thm-dynamical-label-noise-conditional-transport}.  The saturated
clean-risk conversion is
\cref{eq-saturated-dynamical-clean-noise-bound,thm-uniform-dynamical-label-noise-bound};
the recordwise statement
\cref{eq-lpn-complete-dynamical-parity-risk} is
\cref{eq-lpn-parity-clean-risk} with
\(\delta=\delta_{\rm gen}\), truncated by the trivial risk bound one.
Finally,
\cref{thm-controlled-policy-channel-decomposition,cor-saturated-learned-policy-accountability}
retain the one complete derivation and its existing structural source
classification.

\end{proof}

\begin{corollary}[Three exact LPN dynamical stress tests]
\label{cor-lpn-three-dynamical-stress-tests}

The controlled certificate gives the following parameter-explicit
evaluations.  Whenever \(R\) occurs below, assume
\(n/m\to R\in(0,1)\).

\begin{enumerate}[label=\textup{(\roman*)}]
\item \emph{Subset Gaussian elimination.}
      On \(\{\operatorname{rank}A=n\}\), let \(S(A)\) be the represented
      deterministic pivot-row set and put
      \(\widehat s=A_{S(A)}^{-1}b_{S(A)}\).  Then
      \begin{equation}
      \label{eq-lpn-subset-ge-success}
      \Pr_E\{\widehat s=s\mid A,\operatorname{rank}A=n\}
      =(1-\eta)^n,
      \end{equation}
      and hence
      \begin{equation}
      \label{eq-lpn-subset-ge-unconditional-success}
      \Pr\{\widehat s=s\}
      \ge(1-2^{n-m})(1-\eta)^n.
      \end{equation}
      If \(n/m\to R\in(0,1)\) and
      \(\eta m=c\log n\), then
      \begin{equation}
      \label{eq-lpn-subset-ge-log-noise-window}
      (1-\eta)^n=n^{-cR+o(1)}.
      \end{equation}

\item \emph{Located spectral or BKW-type materialization.}
      Put \(\rho=1-2\eta\).  If a charged located block has length
      \(L_B(n)\) and every aligned index has weight at least \(\delta m\),
      then
      \begin{equation}
      \label{eq-lpn-dynamics-located-spectral-budget}
      \left|
      \mathbb E[h_{A,L}(As+E)\chi_v(s)\mid A]
      \right|
      \le\sqrt{L_B(n)}\,\rho^{\lfloor\delta m\rfloor}.
      \end{equation}
      If \(n/m\to R\) and
      \(L_B(n)\le2^{\beta n+o(n)}\), the logarithm of the right side is at
      most
      \begin{equation}
      \label{eq-lpn-dynamics-located-spectral-exponent}
      n\left[
      \frac{\beta\log2}{2}
      -\frac\delta R\log\frac1{1-2\eta}
      \right]+o(n).
      \end{equation}
      It is exponentially negative whenever
      \begin{equation}
      \label{eq-lpn-dynamics-located-spectral-threshold}
      \beta
      <\frac{2\delta}{R\log2}\log\frac1{1-2\eta}.
      \end{equation}
      The exceptional-presentation probability remains the explicit term in
      \cref{eq-lpn-located-block-exceptional-probability}.

\item \emph{Threshold-guided residual dynamics.}
      Fix
      \(\eta<\tau_-<\tau<\tau_+<H_2^{-1}(1-R)\).  On the event in
      \cref{eq-lpn-entropy-sharp-all-perturbation-separation},
      \begin{equation}
      \label{eq-lpn-threshold-equals-target-indicator}
      \mathbf1_{\{W_I\le\tau m\}}=\mathbf1_{\{s\}}.
      \end{equation}
      Consequently every represented normalized kernel satisfies, for
      \(x\ne s\),
      \begin{align}
      \label{eq-lpn-threshold-guided-contact-identity}
      P\mathbf1_{\{W_I\le\tau m\}}(x)&=P(x,s),\\
      \label{eq-lpn-threshold-guided-drift-identity}
      (I-P)(1-\mathbf1_{\{W_I\le\tau m\}})(x)&=P(x,s),\\
      \label{eq-lpn-threshold-guided-capacity-identity}
      \operatorname{Cap}_{\Phi_P}(\{W_I\le\tau m\},
      \{W_I>\tau m\})&=2^{-n}q_{P,s}
      \end{align}
      for a uniform-invariant native kernel.  Thus the threshold readout's
      prospective transition value is exactly the already charged
      target-contact, drift, and capacity coordinate of the selected kernel.
\end{enumerate}

\end{corollary}

\begin{proof}

The pivot set depends only on \(A\).  Conditional on full rank,
\(\widehat s=s+A_{S(A)}^{-1}E_{S(A)}\), so equality holds precisely when
the \(n\) selected independent noise bits vanish.  This proves
\cref{eq-lpn-subset-ge-success}; the rank bound in
\cref{prop-lpn-error-enumeration-success} proves
\cref{eq-lpn-subset-ge-unconditional-success}.  Expanding
\(n\log(1-\eta)=-n\eta+O(n\eta^2)\) gives
\cref{eq-lpn-subset-ge-log-noise-window}.

The Cauchy--Schwarz attenuation in
\cref{eq-lpn-located-walsh-block-attenuation} proves
\cref{eq-lpn-dynamics-located-spectral-budget}.  Substitute
\(m=n/R+o(n)\) and the bound on \(L_B\), then take logarithms, to obtain
\cref{eq-lpn-dynamics-located-spectral-exponent}; its negativity is exactly
\cref{eq-lpn-dynamics-located-spectral-threshold}.

The strict separation
\(W_I(s)\le\tau_-m<\tau m<\tau_+m<W_I(x)\) for \(x\ne s\) proves
\cref{eq-lpn-threshold-equals-target-indicator}.  Applying \(P\) gives the
contact identity, subtraction from one gives the drift identity, and
\cref{eq-lpn-singleton-capacity-exact} gives the capacity identity.

\end{proof}

\subsection{Exact BKW traces and resource-matched comparison}
\label{subsec-lpn-exact-bkw-pareto-comparison}

The located-spectral estimate in
\cref{eq-lpn-dynamics-located-spectral-budget} records the attenuation of a
materialized block.  The generic semantic-cancellation theorem
\cref{thm-charged-noisy-combination-transport} distinguishes the terminal
support of \(\lambda^{\mathsf T}A\) from the noise-provenance support
\(|\lambda|\).  A complete BKW comparison also records how the block is
constructed, how often an original noisy equation is reused, how the terminal
score is decoded, and which data interface supplies the equations.  The
following trace representation makes these quantities explicit.  It uses the
collision-and-terminal-transform structure of BKW and its concrete variants
\citep{BlumKalaiWasserman2003,BogosTramerVaudenay2015,WiggersSamardjiska2021}
while retaining the manuscript's presentation-relative resource ledger.

\begin{definition}[Prefix-adapted LPN combination trace]
\label{def-lpn-prefix-adapted-bkw-trace}

Let \(q\) LPN equations be written as
\(b=As+E\), where \(A\in\mathbb F_2^{q\times n}\) and
\(E_i\stackrel{\rm iid}{\sim}\operatorname{Ber}(\eta)\).  Fix a nonempty
terminal coordinate block \(S\subseteq[n]\), with \(|S|=k\).  A
\emph{prefix-adapted combination trace} for \(S\) is a represented matrix
\begin{equation}
\label{eq-lpn-bkw-combination-trace-matrix}
\Lambda=(\lambda_1^{\mathsf T},\ldots,\lambda_K^{\mathsf T})
\in\mathbb F_2^{K\times q}
\end{equation}
such that:
\begin{enumerate}[label=\textup{(\roman*)}]
\item every \(\lambda_j\) is nonzero;
\item \((\Lambda A)_{[n]\setminus S}=0\);
\item the represented rule constructing \(\Lambda\) is measurable with
respect to \(A_{[n]\setminus S}\), its own charged randomness, and its
clocked construction record, and does not inspect \(A_S\), \(b\), or the
noise;
\item the construction cost, row locators, collision tables, combination
DAG, terminal score, comparison rule, and output map are retained in the
complete record.
\end{enumerate}
Put
\begin{equation}
\label{eq-lpn-bkw-trace-weight-degree-data}
w_j=|\lambda_j|,\qquad
d_i=|\{j:(\lambda_j)_i=1\}|,\qquad
D_2(\Lambda)=\sum_{i=1}^{q}d_i^2,\qquad
\mu_\eta(\Lambda)=\sum_{j=1}^{K}(1-2\eta)^{w_j}.
\end{equation}
The terminal score for \(z\in\mathbb F_2^S\) is
\begin{equation}
\label{eq-lpn-bkw-terminal-score}
Z_{\Lambda,S}(z)
=\sum_{j=1}^{K}
(-1)^{(\Lambda b)_j+\langle(\Lambda A)_{j,S},z\rangle}.
\end{equation}

For full-secret recovery, a BKW trace record is a finite collection
\(\mathfrak T=(S_g,\Lambda_g)_{g=1}^{G}\) whose terminal blocks partition
\([n]\), together with the represented concatenation of the block estimates.
The same acquired equations may be reused across different blocks; that reuse
is charged in each \(D_2(\Lambda_g)\) and in the complete execution ledger.

\end{definition}

\begin{lemma}[Exact BKW trace moments and dependence]
\label{lem-lpn-bkw-exact-trace-moments}

Write \(\rho=1-2\eta\).  Conditional on a realized prefix-adapted trace
\(\Lambda\) and on the inspected coordinates of \(A\),
\begin{align}
\label{eq-lpn-bkw-correct-score-mean}
\mathbb E_E[Z_{\Lambda,S}(s_S)]
&=\mu_\eta(\Lambda)=\sum_{j=1}^{K}\rho^{w_j},\\
\label{eq-lpn-bkw-exact-noise-covariance}
\operatorname{Cov}_E
\left((-1)^{\lambda_j^{\mathsf T}E},
      (-1)^{\lambda_l^{\mathsf T}E}\right)
&=\rho^{|\lambda_j+\lambda_l|}-\rho^{w_j+w_l}.
\end{align}
For each \(z\ne s_S\), conditional on the same trace data and averaging over
the uninspected block \(A_S\) and the noise,
\begin{equation}
\label{eq-lpn-bkw-wrong-score-mean}
\mathbb E_{A_S,E}[Z_{\Lambda,S}(z)]=0.
\end{equation}
Moreover, for every \(t>0\),
\begin{align}
\label{eq-lpn-bkw-correct-score-tail}
\Pr\{Z_{\Lambda,S}(s_S)\le\mu_\eta(\Lambda)-t\mid\Lambda\}
&\le \exp\!\left[-\frac{t^2}{2D_2(\Lambda)}\right],\\
\label{eq-lpn-bkw-wrong-score-tail}
\Pr\{Z_{\Lambda,S}(z)\ge t\mid\Lambda\}
&\le \exp\!\left[-\frac{t^2}{2D_2(\Lambda)}\right].
\end{align}

\end{lemma}

\begin{proof}
At the correct block value, clause~\textup{(ii)} of
\cref{def-lpn-prefix-adapted-bkw-trace} gives
\begin{equation*}
(\Lambda b)_j+\langle(\Lambda A)_{j,S},s_S\rangle
=\lambda_j^{\mathsf T}E.
\end{equation*}
Independence of the original noise bits yields
\(\mathbb E(-1)^{\lambda_j^{\mathsf T}E}=\rho^{w_j}\).  Multiplying the
two signs replaces their exponents by
\((\lambda_j+\lambda_l)^{\mathsf T}E\), which proves
\cref{eq-lpn-bkw-exact-noise-covariance}.

For \(z\ne s_S\), put \(v=z+s_S\ne0\).  Prefix adaptation leaves
\(A_S\) independent of the trace.  Hence
\((\langle A_{i,S},v\rangle+E_i)_{i\le q}\) are independent uniform bits,
and every nonzero \(\lambda_j\) has zero character mean.  This proves
\cref{eq-lpn-bkw-wrong-score-mean}.

Changing original bit \(E_i\) changes the correct score by at most \(2d_i\).
For a wrong score the same bound holds after replacing the independent input
bit by \(\langle A_{i,S},v\rangle+E_i\).  The bounded-difference inequality
therefore has squared sensitivity sum
\(4\sum_i d_i^2=4D_2(\Lambda)\).  Its two one-sided forms give
\cref{eq-lpn-bkw-correct-score-tail,eq-lpn-bkw-wrong-score-tail}.
\end{proof}

\begin{theorem}[Exact dependence-aware BKW recovery certificate]
\label{thm-lpn-dependence-aware-bkw-recovery}

Let \(\widehat s_{S_g}\) maximize
\(Z_{\Lambda_g,S_g}\), using the represented deterministic tie rule, and
concatenate the block estimates.  Conditional on successful construction of
the trace collection \(\mathfrak T\),
\begin{equation}
\label{eq-lpn-bkw-full-recovery-lower-bound}
\Pr\{\widehat s=s\mid\mathfrak T\}
\ge
1-\sum_{g=1}^{G}
\min\left\{1,
2^{k_g}
\exp\!\left[-
\frac{\mu_\eta(\Lambda_g)^2}{8D_2(\Lambda_g)}
\right]
\right\}.
\end{equation}
If trace construction fails with probability at most
\(\varepsilon_{\rm tr}\), the unconditional right side is reduced by
\(\varepsilon_{\rm tr}\).  Public verification further contributes the
global error
\begin{equation}
\label{eq-lpn-bkw-complete-arrival-certificate}
\varepsilon_V(n,q,\eta,\tau)
=
\exp[-qD_{\rm Ber}(\tau\Vert\eta)]
+(2^n-1)2^{-q(1-H_2(\tau))}
\end{equation}
when the trace is built from \(q\) independent presented equations.

\end{theorem}

\begin{proof}
For block \(g\), apply
\cref{eq-lpn-bkw-correct-score-tail} with
\(t=\mu_\eta(\Lambda_g)/2\).  Apply
\cref{eq-lpn-bkw-wrong-score-tail} with the same threshold to each of the
\(2^{k_g}-1\) wrong block values.  Outside the union of these events the
correct score is strictly larger than every wrong score.  A union bound over
the candidates and then over the terminal blocks proves
\cref{eq-lpn-bkw-full-recovery-lower-bound}.  The trace-construction and
verifier errors are separate events.  The verifier term is
\cref{eq-lpn-levin-global-verifier-error} with matrix height \(q\).
\end{proof}

\begin{definition}[Canonical materialized BKW schedule and exact cost]
\label{def-lpn-canonical-materialized-bkw-schedule}

For one terminal block of size \(k\), a canonical disjoint BKW schedule is an
integer tuple
\begin{equation}
\label{eq-lpn-bkw-schedule-parameters}
\theta_{\rm BKW}=(q,r,k,b_1,\ldots,b_r),\qquad
\sum_{j=1}^{r}b_j=n-k.
\end{equation}
At stage \(j\), the schedule buckets the current rows by the next \(b_j\)
coordinates and forms disjoint pairs within each bucket.  Put
\begin{align}
\label{eq-lpn-bkw-stage-dimensions}
D_0&=n,&D_j&=n-\sum_{h=1}^{j}b_h,\\
\label{eq-lpn-bkw-stage-row-recurrence}
N_0&=q,&
N_j&=\left\lfloor
\frac{(N_{j-1}-2^{b_j})_+}{2}
\right\rfloor.
\end{align}
The recurrence is a deterministic lower bound: among at most \(2^{b_j}\)
buckets, pairing leaves at most one unpaired row per bucket.  The final trace
has \(K=N_r\), mutually disjoint leaf supports of weight \(2^r\), and hence
\begin{equation}
\label{eq-lpn-bkw-disjoint-trace-parameters}
\mu_\eta=K\rho^{2^r},\qquad
D_2=K2^r.
\end{equation}

Fix implementation constants
\(c_{\rm key},c_{\rm pair},c_{\rm xor},c_{\rm dag},
c_{\rm hist},c_{\rm F},c_{\rm max},c_{\rm acq}>0\).
The canonical bit-operation and word-operation ledger is
\begin{align}
\label{eq-lpn-bkw-exact-reduction-time}
T_{\rm red}(\theta_{\rm BKW})
=\sum_{j=1}^{r}\Bigl[&
c_{\rm key}N_{j-1}b_j+c_{\rm pair}N_j
+c_{\rm xor}N_j(D_{j-1}+1)\notag\\
&+2c_{\rm dag}N_j
\left\lceil\log_2(N_{j-1}\vee2)\right\rceil
\Bigr],\\
\label{eq-lpn-bkw-exact-terminal-time}
T_{\rm term}(\theta_{\rm BKW})
&=c_{\rm hist}K(k+1)+c_{\rm F}k2^k+c_{\rm max}2^k,\\
\label{eq-lpn-bkw-exact-data-cost}
Q_{\rm BKW}(\theta_{\rm BKW})&=q,\\
\label{eq-lpn-bkw-exact-memory-cost}
M_{\rm BKW}(\theta_{\rm BKW})
&=\max_{1\le j\le r}
\Bigl\{
N_{j-1}(D_{j-1}+1)+N_j(D_j+1)\notag\\
&\hspace{34mm}
+2N_j\lceil\log_2(N_{j-1}\vee2)\rceil,
\ 2^k\lceil\log_2(K+1)\rceil
\Bigr\},\\
\label{eq-lpn-bkw-exact-total-time}
T_{\rm BKW}^{\rm fix}(\theta_{\rm BKW})
&=T_{\rm red}+T_{\rm term},\\
\label{eq-lpn-bkw-exact-stream-time}
T_{\rm BKW}^{\rm stream}(\theta_{\rm BKW})
&=c_{\rm acq}q+T_{\rm red}+T_{\rm term}.
\end{align}
The code length \(d_{\rm BKW}(\theta_{\rm BKW})\) is the canonical
self-delimiting code of the stage tag, the integers in
\cref{eq-lpn-bkw-schedule-parameters}, the collision and tie rules, the
terminal transform, precision, and the fixed interpreter identifier, as
specified by \cref{def-canonical-complete-active-record-code}.  A full-secret
schedule sums time and data actually acquired across its block traces, takes
peak memory, and includes every reused-data reference in the DAG ledger.

\end{definition}

\begin{corollary}[Closed-form disjoint-BKW sample threshold]
\label{cor-lpn-disjoint-bkw-sample-threshold}

Let \(G\) terminal blocks have size at most \(k\), and allocate total trace
failure \(\epsilon\in(0,1)\) uniformly across them.  A sufficient final row
count for each block is
\begin{equation}
\label{eq-lpn-bkw-required-final-rows}
K_{\rm req}(r,k,G,\eta,\epsilon)
=
\left\lceil
\frac{8\,2^r
\bigl(k\log 2+\log(G/\epsilon)\bigr)}
{(1-2\eta)^{2^{r+1}}}
\right\rceil.
\end{equation}
Starting from \(N_r=K_{\rm req}\), the reverse recursion
\begin{equation}
\label{eq-lpn-bkw-reverse-sample-recurrence}
N_{j-1}^{\rm req}=2N_j^{\rm req}+2^{b_j}+1,
\qquad j=r,r-1,\ldots,1,
\end{equation}
produces an explicit sufficient initial sample count
\(q_{\rm req}=N_0^{\rm req}\).  With these values,
\begin{equation}
\label{eq-lpn-bkw-disjoint-success-certificate}
\Pr\{\widehat s=s\}
\ge1-\epsilon-\varepsilon_V(n,q_{\rm req},\eta,\tau).
\end{equation}
All constants in the data, reduction, terminal, memory, and Levin costs are
then obtained by substituting the reverse recurrence into
\cref{eq-lpn-bkw-exact-reduction-time,eq-lpn-bkw-exact-terminal-time,eq-lpn-bkw-exact-memory-cost}.

\end{corollary}

\begin{proof}
Substitute \cref{eq-lpn-bkw-disjoint-trace-parameters} into
\cref{eq-lpn-bkw-full-recovery-lower-bound}.  Each block error is at most
\begin{equation*}
2^k\exp\!\left[-
\frac{K(1-2\eta)^{2^{r+1}}}{8\,2^r}
\right].
\end{equation*}
The choice in \cref{eq-lpn-bkw-required-final-rows} makes this at most
\(\epsilon/G\).  The reverse recurrence implies
\(N_j\ge N_j^{\rm req}\) in
\cref{eq-lpn-bkw-stage-row-recurrence}, completing the proof.
\end{proof}

\begin{corollary}[Proof-producing scaling audit for materialized BKW]
\label{cor-lpn-proof-producing-bkw-scaling-audit}

Fix \(n,m,\eta,\epsilon\), terminal block size \(k\), and number \(G\) of
terminal blocks.  For a required cancellation level \(q_{\rm can}\), let
\(\mathcal P_{\rm BKW}^{\rm fix}(q_{\rm can})\) be the exact occupancy paths
of the canonical materialized BKW cell that use at most \(m\) presented
rows, cancel at least \(q_{\rm can}\) semantic coordinates, and retain
\[
N_r\ge K_{\rm req}(r,k,G,\eta,\epsilon).
\]
Define the fixed-data charged crossing value
\begin{equation}
\label{eq-lpn-bkw-charged-crossing-value}
\mathfrak C_{\rm BKW}^{\rm fix}
(n,m,\eta,\epsilon;q_{\rm can})
=
\inf_{p\in\mathcal P_{\rm BKW}^{\rm fix}(q_{\rm can})}
\operatorname{Cost}_{\rm complete}(p),
\end{equation}
with infimum \(+\infty\) when the path set is empty.  The complete cost is
the acquisition, time, memory, description, provenance, terminal, and
verification ledger in
\cref{def-lpn-canonical-materialized-bkw-schedule}.  The streaming crossing
value removes the \(N_0\le m\) restriction and charges
\(c_{\rm acq}N_0\).

The canonical guaranteed-row schedules give the constructive candidate
\begin{equation}
\label{eq-lpn-bkw-constructive-schedule-cost}
\mathfrak U_{\rm BKW}^{\rm fix}
=
\min_{\theta_{\rm BKW}}
\left\{
\begin{array}{l|l}
\operatorname{Cost}_{\rm complete}(\theta_{\rm BKW})
&
\begin{array}{l}
\sum_{j=1}^{r}b_j\ge q_{\rm can},\quad
N_r\ge K_{\rm req}(r,k,G,\eta,\epsilon),\\
N_0\le m,\quad b_j\in\mathbb N,\quad
\theta_{\rm BKW}\text{ satisfies }
\cref{eq-lpn-bkw-stage-row-recurrence}
\end{array}
\end{array}
\right\},
\end{equation}
and
\begin{equation}
\label{eq-lpn-bkw-crossing-build-comparison}
\mathfrak C_{\rm BKW}^{\rm fix}
\le \mathfrak U_{\rm BKW}^{\rm fix}.
\end{equation}
Under the LPN substitutions
\begin{equation}
\label{eq-lpn-bkw-generic-scaling-substitution}
Q_j=2^{b_j},\qquad d_j=b_j,\qquad
\rho=1-2\eta,\qquad K=N_r,
\end{equation}
the exact occupancy paths and noise transport are those of
\cref{thm-exact-scaling-certificate-disjoint-collision-combining}.  Hence a
deterministic certificate of the lower inequality
\begin{equation}
\label{eq-lpn-bkw-crossing-exceeds-budget}
\mathfrak C_{\rm BKW}^{\rm fix}
(n,m,\eta,\epsilon;q_{\rm can})>B_n
\end{equation}
proves that no record in the canonical materialized BKW cell reaches that
cancellation level together with the displayed terminal-accuracy
certificate within budget \(B_n\).

A learning pipeline may propose a piecewise schedule and lower-bound
template from finitely many evaluated sizes.  It gives a proof for every
\(n\ge N_0\) when the exported certificate verifies:
\begin{enumerate}[label=\textup{(\roman*)}]
\item the exact bucket-occupancy, semantic-support, provenance, and complete
      cost recurrences;
\item the noise-dependent terminal requirement
      \cref{eq-lpn-bkw-required-final-rows};
\item a complete recurrence cover of the canonical materialized BKW cell and
      a branch-and-bound, dynamic-programming, or symbolic lower bound on
      every exact path in that cover; and
\item the uniform tail inequality
      \cref{eq-lpn-bkw-crossing-exceeds-budget}, including every breakpoint
      of the proposed piecewise law.
\end{enumerate}
The finite experiments discover and test the scaling template.  The exact
recurrence cover, lower path optimization, and symbolic induction establish
its continuation and the price of every scaling breakpoint.

\end{corollary}

\begin{proof}
The substitutions in
\cref{eq-lpn-bkw-generic-scaling-substitution} identify the bucket alphabet,
semantic cancellation, provenance size, and noise attenuation with
\cref{def-disjoint-collision-cancellation-recurrence}.  The
dependence-aware score certificate gives the terminal row requirement in
\cref{eq-lpn-bkw-required-final-rows}.  Guaranteed complete-pairing
schedules are members of the exact path set, which proves
\cref{eq-lpn-bkw-crossing-build-comparison}.  A verified lower bound on the
exact path infimum is the crossing certificate in
\cref{thm-verified-recurrence-propagation-crossing-exclusion}; that theorem
proves the exclusion and the symbolic tail assertion.
\end{proof}

\begin{definition}[Fixed-data and streaming optimized BKW envelopes]
\label{def-lpn-fixed-streaming-bkw-envelopes}

Let \(\mathcal A_{\rm BKW}^{\rm fix}(n,m,\eta)\) contain all represented
prefix-adapted BKW trace records that use only the \(m\) rows of the public
presentation.  Let \(\mathcal A_{\rm BKW}^{\rm stream}(n,\eta)\) contain the
same represented stage grammar with a conservative LPN sample interface,
where every acquired row is counted by
\cref{eq-lpn-bkw-exact-data-cost,eq-lpn-bkw-exact-stream-time}.  Define
\begin{align}
\label{eq-lpn-fixed-bkw-optimized-frontier}
A_{\rm BKW}^{*,\rm fix}(\mathbf b)
&=\operatorname{BenchArr}_{\mathcal A_{\rm BKW}^{\rm fix},H}(\mathbf b),\\
\label{eq-lpn-stream-bkw-optimized-frontier}
A_{\rm BKW}^{*,\rm stream}(\mathbf b)
&=\operatorname{BenchArr}_{\mathcal A_{\rm BKW}^{\rm stream},H}(\mathbf b).
\end{align}
The corresponding matched-success frontiers are
\(\mathcal P_{\mathcal A_{\rm BKW}^{\rm fix},H}(\alpha)\) and
\(\mathcal P_{\mathcal A_{\rm BKW}^{\rm stream},H}(\alpha)\).
An optimizer over the integer schedule parameters evaluates these definitions;
it may discard dominated schedules but may not omit acquisition, table,
combination, decoder, verification, or description costs.
The family also contains represented approximate-collision or coded-BKW
stages when their code, quantizer, compensation defect, and terminal decoder
are included in the complete record.  Such a stage uses the existing
compensated-quotient certificate of
\cref{def-master-compensated-quotient-geometry-certificate}; the exact trace certificate
\cref{thm-lpn-dependence-aware-bkw-recovery} applies to the zero-defect
subfamily.  Thus optimization over the full envelope retains modern BKW
variants, while the closed-form grid below reports the canonical exact
disjoint subfamily whose constants are evaluated without an approximation.

\end{definition}

\begin{theorem}[Complete structural charging of a BKW schedule]
\label{thm-lpn-complete-bkw-structural-charging}

Every record in \cref{def-lpn-fixed-streaming-bkw-envelopes} is one complete
record in the existing saturated profile.  Its row combinations have
\(\mathsf{Add}\) ancestry; a dynamically qualified block or fiber map is
charged to its exact or compensated \(\Abs\) coordinate; collision equality,
table ordering, and stage clocks retain their \(\mathsf{Sym}\),
\(\mathsf{Ord}\), and \(\mathsf{Filt}\) provenance; iteration is charged to
\(\Lift\); the terminal score and prospective choice retain their authorized
\(\Geom/\Caus\) gates; public target testing is charged to the verifier; and
nonstructural transcript storage remains in \(J_{\rm res}\).  Derived
constructors create no additional source charge.

For every row \(\lambda_{g,j}\), the semantic support is
\(\operatorname{supp}(A^{\mathsf T}\lambda_{g,j})\subseteq S_g\), while its
noise-provenance support is \(|\lambda_{g,j}|\).  The collision stages reduce
the former; the latter determines the exact attenuation
\((1-2\eta)^{|\lambda_{g,j}|}\) and retains the complete combination-DAG
cost.  Hence every polynomially materialized trace with
\(\max_g|S_g|=o(n)\) lies in the rate-uniform sector of
\cref{cor-lpn-high-sample-materialized-combination-closure}.

For a materialized trace collection put
\begin{equation}
\label{eq-lpn-bkw-located-spectrum-identification}
L_B=\sum_{g=1}^{G}K_g,\qquad
q_{\min}=\min_{g,j}|\lambda_{g,j}|,\qquad
\beta_{\rm BKW}=\frac{\log_2(L_B\vee1)}{n},\qquad
\delta_{\rm BKW}=\frac{q_{\min}}{q}.
\end{equation}
Then its located target-character coordinate satisfies
\begin{equation}
\label{eq-lpn-bkw-located-spectrum-complete-bound}
\left|\operatorname{Corr}_{\rm loc}(\mathfrak T;s)\right|
\le\min\!\left\{1,\sqrt{L_B}\,(1-2\eta)^{q_{\min}}\right\}.
\end{equation}
For fixed data \(q=m\), this is exactly
\cref{eq-lpn-dynamics-located-spectral-budget} with
\(\delta=\delta_{\rm BKW}\).  For streaming data, the identical formula
holds with the charged acquired height \(q\).  Every other upper certificate
available to the same complete record remains in the minimum prescribed by
the saturated framework.

\end{theorem}

\begin{proof}
The provenance assignments are instances of
\cref{thm-complete-generic-source-theory,cor-five-family-abs-envelope,thm-affordable-geom-abs-normal-form,prop-derived-channels-create-no-structure}.
The trace constructor, its locators, and its costs remain ancestors by
\cref{thm-paper4-dependency-principle,thm-affordable-composition-no-creation,thm-charged-noisy-combination-transport}.
Applying \cref{thm-lpn-walsh-alignment-charged-location} to the
\(L_B\) explicitly located characters, followed by Cauchy--Schwarz, gives
\cref{eq-lpn-bkw-located-spectrum-complete-bound}.  The fixed and streaming
statements differ only in the declared source of the \(q\) independent rows;
the latter source is included in the cost by definition.  The rate-uniform
statement follows from
\cref{thm-lpn-rate-uniform-semantic-cancellation}.
\end{proof}

\begin{corollary}[Quantitative saturated improvement relative to optimized BKW]
\label{cor-lpn-quantitative-improvement-over-bkw}

Let \(\diamond\in\{\mathrm{fix},\mathrm{stream}\}\), and let
\([L_{\rm BKW}^{\diamond},U_{\rm BKW}^{\diamond}]\) be a simultaneous valid
interval for the optimized BKW frontier at resource ceiling \(\mathbf b\).
Take \([L_{\rm sat},U_{\rm sat}]\) from the sharp integrated LPN minimum,
including
\cref{eq-lpn-levin-integrated-length-profile-minimum} and the complete LPN
Bellman/source candidates, and let
\([L_{\widehat R},U_{\widehat R}]\) be the learner interval.  Then
\begin{align}
\label{eq-lpn-bkw-saturated-headroom-interval}
\bigl[L_{\rm sat}-U_{\rm BKW}^{\diamond}\bigr]_+
&\le
\operatorname{ActArr}_{\rm LPN}^{\rm sat}(\mathbf b)
-A_{\rm BKW}^{*,\diamond}(\mathbf b)
\le U_{\rm sat}-L_{\rm BKW}^{\diamond},\\
\label{eq-lpn-bkw-learner-gain-interval}
L_{\widehat R}-U_{\rm BKW}^{\diamond}
&\le
\mathfrak M(A_{\widehat R}^H)-A_{\rm BKW}^{*,\diamond}(\mathbf b)
\le U_{\widehat R}-L_{\rm BKW}^{\diamond}.
\end{align}
At matched arrival \(\alpha\), the time, data, and memory improvements are
the three coordinates of
\cref{eq-matched-success-log-resource-saving}; the BKW-specific witnessed
Levin gain is
\begin{equation}
\label{eq-lpn-bkw-specific-levin-gain}
\Delta^{\mathcal A_{\rm BKW}}
\ge d_{\rm BKW}(\theta_{\rm BKW})
-\ell_{{\rm cert},{\rm LPN}}
\end{equation}
whenever the simultaneous benchmark-matching condition holds through that
canonical code length.

For the disjoint materialized schedule, every displayed quantity is an
explicit function of
\begin{equation}
\label{eq-lpn-bkw-complete-parameter-list}
(n,m,q,\eta,H,\tau,\epsilon,G,r,k,b_1,\ldots,b_r,
c_{\rm key},c_{\rm pair},c_{\rm xor},c_{\rm dag},c_{\rm hist},c_{\rm F},
c_{\rm max},c_{\rm acq})
\end{equation}
through
\cref{eq-lpn-bkw-required-final-rows,eq-lpn-bkw-reverse-sample-recurrence,eq-lpn-bkw-exact-reduction-time,eq-lpn-bkw-exact-terminal-time,eq-lpn-bkw-exact-memory-cost,eq-lpn-bkw-complete-arrival-certificate}.

\end{corollary}

\begin{proof}
Apply \cref{thm-saturated-pareto-benchmark-comparison} with
\(\mathcal A=\mathcal A_{\rm BKW}^{\diamond}\).  The matched-success and
description-length statements are
\cref{cor-matched-success-resource-and-levin-gain}.  The parameter list is a
collection of the arguments appearing in the cited exact formulas.
\end{proof}

\paragraph{Reproducible parameter grid.}
The canonical evaluator
\texttt{scripts/lpn\_bkw\_frontier.py} enumerates balanced integer block
schedules, retains the nondominated resource vectors, and writes the complete
grid to
\texttt{docs/source/01\_task\_space/outputs/lpn\_bkw\_frontier.csv}.
The following rows use unit implementation constants, target trace failure
\(10^{-2}\), \(m=4n\), terminal sizes
\(k\in\{1,2,4,8,16\}\), and at most twelve stages.  Each displayed row is
the minimum-time streaming point among the retained schedules at its
\((n,m,\eta)\); logarithms are base two.

\begin{table}[tbp]
\centering
\small
\begin{tabular}{rrrrrrrr}
\toprule
\(n\)&\(\eta\)&\(k\)&\(r\)&\(\log T\)&\(\log M\)&\(\log Q\)&success lower\\
\midrule
128&0.010&16&6&35.282&32.445&27.748&0.990001\\
128&0.125&16&4&42.567&39.727&35.031&0.990000\\
256&0.010&16&8&51.569&47.744&43.064&0.990000\\
256&0.125&16&5&65.483&61.634&56.957&0.990000\\
512&0.010&16&9&78.431&73.657&69.003&0.990000\\
512&0.125&16&6&102.754&97.946&93.285&0.990000\\
\bottomrule
\end{tabular}
\caption{Certified canonical disjoint-BKW streaming points.  The public
fixed-data presentations in these rows contain only \(m=4n\) equations and
do not meet the corresponding sufficient sample thresholds; the fixed-data
frontier is therefore empty for this certified schedule grid rather than being
silently supplemented with streaming samples.}
\label{tab-lpn-certified-bkw-grid}
\end{table}

\begin{figure}[tbp]
\centering
\begin{minipage}{.48\linewidth}
\centering
\begin{tikzpicture}
\begin{axis}[fw axis,width=.97\linewidth,height=5.6cm,
xlabel={noise \(\eta\)},ylabel={\(\log_2 K_{\rm req}\)},
xmin=0,xmax=.2,ymin=0,ymax=300,
legend style={font=\scriptsize,at={(.03,.97)},anchor=north west}]
\addplot[draw=fwBlue,line width=1.2pt,domain=0:.2,samples=160]
 {ln(8*2^4*(16*ln(2)+ln(16/.01)))/ln(2)
  -2^5*ln(1-2*x)/ln(2)};
\addlegendentry{\(r=4\)}
\addplot[draw=fwAmber,line width=1.2pt,domain=0:.2,samples=160]
 {ln(8*2^6*(16*ln(2)+ln(16/.01)))/ln(2)
  -2^7*ln(1-2*x)/ln(2)};
\addlegendentry{\(r=6\)}
\addplot[draw=fwCoral,line width=1.2pt,domain=0:.2,samples=160]
 {ln(8*2^8*(16*ln(2)+ln(16/.01)))/ln(2)
  -2^9*ln(1-2*x)/ln(2)};
\addlegendentry{\(r=8\)}
\end{axis}
\end{tikzpicture}
\end{minipage}\hfill
\begin{minipage}{.48\linewidth}
\centering
\begin{tikzpicture}
\begin{axis}[fw axis,width=.97\linewidth,height=5.6cm,
xlabel={noise \(\eta\)},ylabel={located threshold \(\beta_{\rm crit}\)},
xmin=0,xmax=.2,ymin=0,ymax=.5,
legend style={font=\scriptsize,at={(.03,.97)},anchor=north west}]
\addplot[draw=fwBlue,line width=1.2pt,domain=0:.2,samples=160]
 {2*.01/(.25*ln(2))*ln(1/(1-2*x))};
\addlegendentry{\(\delta=.01\)}
\addplot[draw=fwAmber,line width=1.2pt,domain=0:.2,samples=160]
 {2*.025/(.25*ln(2))*ln(1/(1-2*x))};
\addlegendentry{\(\delta=.025\)}
\addplot[draw=fwCoral,line width=1.2pt,domain=0:.2,samples=160]
 {2*.05/(.25*ln(2))*ln(1/(1-2*x))};
\addlegendentry{\(\delta=.05\)}
\end{axis}
\end{tikzpicture}
\end{minipage}
\caption{Noise dependence of the certified BKW resources.  Left: the exact
terminal-row threshold from
\cref{eq-lpn-bkw-required-final-rows} for \(k=16\), \(G=16\), and
\(\epsilon=.01\).  Right: the located-spectrum boundary from
\cref{eq-lpn-dynamics-located-spectral-threshold} at \(R=.25\).  Increasing
the number of reductions raises the combination order and therefore amplifies
the noise penalty explicitly.}
\label{fig-lpn-bkw-noise-resource-frontiers}
\end{figure}

\begin{figure}[tbp]
\centering
\resizebox{.98\linewidth}{!}{%
\begin{tikzpicture}[fw diagram,node distance=8mm and 10mm]
  \node[fw source,text width=2.5cm] (data)
    {LPN equations\\fixed \(m\) or charged \(q\)};
  \node[fw provenance,text width=2.5cm,right=of data] (trace)
    {prefix-adapted trace\\\(\Lambda,w,d,D_2\)};
  \node[fw geom,text width=2.5cm,above right=7mm and 10mm of trace] (noise)
    {exact noise ledger\\\(\mu_\eta,\operatorname{Cov}\)};
  \node[fw box,text width=2.5cm,below right=7mm and 10mm of trace] (cost)
    {stage ledger\\time, data, memory};
  \node[fw use,text width=2.7cm,right=17mm of trace] (decode)
    {terminal score\\and recovery interval};
  \node[fw box,text width=2.7cm,right=of decode] (opt)
    {nondominated optimized\\BKW frontier};
  \node[fw terminal,text width=2.8cm,right=of opt] (gain)
    {saturated headroom\\learner gain};
  \draw[fw arrow] (data)--(trace);
  \draw[fw arrow] (trace)--(noise);
  \draw[fw arrow] (trace)--(cost);
  \draw[fw arrow] (noise)--(decode);
  \draw[fw arrow] (cost)--(decode);
  \draw[fw arrow] (decode)--(opt);
  \draw[fw arrow] (opt)--(gain);
\end{tikzpicture}%
}
\caption{Exact resource-matched BKW comparison.  The combination trace keeps
sample reuse and its covariance visible; the stage record charges construction
and storage; optimization retains the nondominated fixed-data and streaming
frontiers; and the saturated comparison reports headroom and achieved learner
gain separately.}
\label{fig-lpn-exact-bkw-pareto-ledger}
\end{figure}

\Cref{cor-lpn-three-dynamical-stress-tests} supplies direct checks on the
subset-elimination, located-spectral, and threshold-guided regimes.
\Cref{fig-lpn-controlled-nonlinear-ambient-dependency} summarizes the common
qualified-record dependency of their controlled structural coordinates.

\begin{figure}[tbp]
\centering
\resizebox{\linewidth}{!}{%
\begin{tikzpicture}[fw diagram,x=1cm,y=1cm]
  \node[fw source,text width=2.4cm] (lpn) at (0,0)
    {public LPN record\\\((A,b,\eta)\)};
  \node[fw provenance,text width=2.6cm] (kernel) at (3.5,0)
    {one qualified\\dynamical record};
  \node[fw geom,text width=2.7cm] (cap) at (7.2,4)
    {capacity and drift\\\(2^{-n},\delta_{n,m,\gamma}\)};
  \node[fw geom,text width=2.7cm] (fi) at (7.2,2)
    {functional slots\\safe spectrum};
  \node[fw use,text width=2.7cm] (caus) at (7.2,0)
    {Caus action--aiming\\\((1-2\eta)^{-1}\)};
  \node[fw provenance,text width=2.7cm] (lift) at (7.2,-2)
    {valid Lift transfer\\\(\rho_\pi,E_t,\sigma_\lambda\)};
  \node[fw use,text width=2.7cm] (learn) at (7.2,-4)
    {Bellman and learning\\\(m,\Gamma_{\rm dyn}\)};
  \node[fw terminal,text width=2.8cm] (ledger) at (11.4,0)
    {existing prospective\\\(\Geom/\Caus\) ledger};
  \draw[fw arrow] (lpn) -- (kernel);
  \draw[fw arrow] (kernel) -- (cap);
  \draw[fw arrow] (kernel) -- (fi);
  \draw[fw arrow] (kernel) -- (caus);
  \draw[fw arrow] (kernel) -- (lift);
  \draw[fw arrow] (kernel) -- (learn);
  \draw[fw arrow] (cap) -- (ledger);
  \draw[fw arrow] (fi) -- (ledger);
  \draw[fw arrow] (caus) -- (ledger);
  \draw[fw arrow] (lift) -- (ledger);
  \draw[fw arrow] (learn) -- (ledger);
\end{tikzpicture}%
}
\caption{The LPN nonlinear ambient branch uses one qualified record.
Capacity and drift, the typed functional profile, Caus action, valid-Lift
transfer, and dynamical learning are numerical outputs of that same record.
Their explicit \(n,m,\eta\) dependence returns to the existing prospective
\(\Geom/\Caus\) ledger.}
\label{fig-lpn-controlled-nonlinear-ambient-dependency}
\end{figure}

\section{Readout-Agnostic Capacity, Information, and Blackwell Envelopes}
\label{sec-lpn-readout-agnostic-envelopes}

The next three bounds remove the syntactic form of a represented readout from
the numerical comparison.  They apply, respectively, to the invariant law of
one selected kernel, the ordered reveal channels of one pre-hit record, and
the represented comparison of two experiments.  Each conclusion retains its
own hypotheses and complete cost; none is inferred from provenance alone.

\subsection{Invariant-law capacity and represented tilts}
\label{subsec-lpn-invariant-law-capacity-tilts}

\begin{theorem}[LPN invariant-density capacity and drift envelope]
\label{thm-lpn-invariant-density-capacity-drift-envelope}

Fix the uniform law \(\mu_I\) on \(X_I=\mathbb F_2^n\).  Let \(P_0\) be a
complete represented reversible reference kernel with invariant law
\(\mu_I\), and put

\[
q_{0,s}=1-P_0(s,s).
\]

Let \((\widetilde P,\widetilde\mu,\Phi_{\widetilde P})\) be a complete
represented reversible selected record on the same carrier.  Suppose that
the record certifies the pointwise comparisons

\begin{align}
\label{eq-lpn-invariant-density-comparison}
r_-2^{-n}
&\le\widetilde\mu(x)\le r_+2^{-n},\qquad x\in X_I,\\
\label{eq-lpn-invariant-conductance-comparison}
\kappa_-2^{-n}P_0(x,z)
&\le\widetilde\mu(x)\widetilde P(x,z)
\le\kappa_+2^{-n}P_0(x,z),\qquad x\ne z.
\end{align}

Then

\begin{align}
\label{eq-lpn-invariant-target-mass-capacity}
\widetilde\mu(s)
&\le r_+2^{-n},
&
\operatorname{Cap}_{\Phi_{\widetilde P}}(\{s\},X_I\setminus\{s\})
&\le\kappa_+2^{-n}q_{0,s}.
\end{align}

For every horizon \(H\) and initial law
\(\rho_0\le R_{\widetilde\mu}\widetilde\mu\),

\begin{equation}
\label{eq-lpn-invariant-density-hit-envelope}
\Pr_{\rho_0}^{\widetilde P}\{\tau_s\le H\}
\le
\min\!\left\{
1,\,
R_{\widetilde\mu}
\bigl(r_++H\kappa_+q_{0,s}\bigr)2^{-n}
\right\}.
\end{equation}

If instead \(\rho_0\le R_\mu\mu_I\), the same conclusion holds with
\(R_{\widetilde\mu}\) replaced by \(R_\mu/r_-\).  Finally, if a
nonconstant represented potential \(F\ge0\), with \(F(s)=0\), has global
selected-kernel drift margin

\[
a_s(\widetilde P,F)
=\inf_{x\ne s}(I-\widetilde P)F(x)>0,
\]

then, whenever \(r_+2^{-n}<1\),

\begin{equation}
\label{eq-lpn-invariant-density-drift-envelope}
\frac{a_s(\widetilde P,F)}{\operatorname{Osc}(F)}
\le
\frac{\kappa_+2^{-n}q_{0,s}}
     {1-r_+2^{-n}}.
\end{equation}

Thus a nonlinear represented readout is absorbed into the selected
invariant law and conductance comparison: its syntactic form does not occur
in the bound.  The density, conductance, initialization, and potential
certificates, including their evaluation precision and complete cost, remain
in the same prospective dynamical record.

\end{theorem}

\begin{proof}

The reference singleton capacity is
\(
2^{-n}q_{0,s}
\)
by \cref{eq-lpn-singleton-capacity-exact}.  Apply
\cref{eq-controlled-density-target-mass-transfer,eq-controlled-density-capacity-transfer}
to obtain \cref{eq-lpn-invariant-target-mass-capacity}.  The
\(L^\infty\) and base-density cases of
\cref{thm-controlled-density-capacity-hit-transfer} give
\cref{eq-lpn-invariant-density-hit-envelope} and its stated variant.  Apply
\cref{eq-controlled-invariant-drift-capacity-ceiling} to
\((\widetilde P,\widetilde\mu)\), then use
\cref{eq-lpn-invariant-target-mass-capacity}; this proves
\cref{eq-lpn-invariant-density-drift-envelope}.

\end{proof}

\begin{corollary}[Noise-indexed LPN tilt and mixing tradeoff]
\label{cor-lpn-noise-indexed-tilt-capacity-tradeoff}

Retain the reference kernel \(P_0\) of
\cref{thm-lpn-invariant-density-capacity-drift-envelope}.  For any
represented bounded potential \(D\) and inverse temperature \(\theta\ge0\),
let \(P_{\theta,D}\) be its represented Gibbs--Metropolis tilt from
\cref{cor-controlled-gibbs-metropolis-capacity-hit}.  If
\(\rho_0\le R_\mu\mu_I\), then

\begin{equation}
\label{eq-lpn-generic-gibbs-hit-envelope}
\Pr_{\rho_0}^{P_{\theta,D}}\{\tau_s\le H\}
\le
\min\!\left\{
1,\,
R_\mu e^{2\theta\operatorname{Osc}(D)}
\bigl(1+Hq_{0,s}\bigr)2^{-n}
\right\}.
\end{equation}

If \(D(s)=\min_{x\in X_I}D(x)\), then

\begin{equation}
\label{eq-lpn-target-minimum-gibbs-hit-envelope}
\Pr_{\rho_0}^{P_{\theta,D}}\{\tau_s\le H\}
\le
\min\!\left\{
1,\,
R_\mu e^{\theta\operatorname{Osc}(D)}
\bigl(1+Hq_{0,s}\bigr)2^{-n}
\right\}.
\end{equation}

In particular, if
\(
\log R_\mu+\log(H+1)+2\theta\operatorname{Osc}(D)=o(n)
\), the right side of
\cref{eq-lpn-generic-gibbs-hit-envelope} is \(2^{-n+o(n)}\).  Under the
target-minimum hypothesis, the factor two may be omitted from this
condition by
\cref{eq-lpn-target-minimum-gibbs-hit-envelope}.  These conclusions use the evaluated
oscillation of the represented potential; affordability alone is not
substituted for that number.

For the residual Metropolis record
\(P_\theta=P_{\theta,W_I}\) of
\cref{cor-lpn-residual-metropolis-capacity-mixing}, work on
\(\mathcal F_\delta\), put

\[
\Delta_-=\frac12-\eta-2\delta>0,
\qquad
\Delta_+=\frac12-\eta+2\delta,
\]

and fix \(0<\varepsilon<1\).  In the subcritical-tilt regime

\begin{equation}
\label{eq-lpn-residual-tilt-subcritical-condition}
\theta m\Delta_+
\le(1-\varepsilon)n\log2,
\end{equation}

one has

\begin{align}
\label{eq-lpn-residual-tilt-subcritical-mass}
\pi_\theta(s)
&\le
\frac{2^{-\varepsilon n}}{1-2^{-n}},\\
\label{eq-lpn-residual-tilt-subcritical-hit}
\Pr_{\rho_0}^{P_\theta}\{\tau_s\le H\}
&\le
\min\!\left\{
1,\,
\frac{R_\theta2^{-\varepsilon n}}{1-2^{-n}}
\left(1+\frac H2e^{-\theta m\Delta_-}\right)
\right\}
\quad
\text{if }\rho_0\le R_\theta\pi_\theta.
\end{align}

In the supercritical-concentration regime

\begin{equation}
\label{eq-lpn-residual-tilt-supercritical-condition}
\theta m\Delta_-
\ge(1+\varepsilon)n\log2,
\end{equation}

the same record satisfies

\begin{align}
\label{eq-lpn-residual-tilt-supercritical-mass}
\pi_\theta(s)
&\ge\frac1{1+2^{-\varepsilon n}},\\
\label{eq-lpn-residual-tilt-supercritical-exit-capacity}
q_{\theta,s}
&\le\frac12\,2^{-(1+\varepsilon)n},
&
\operatorname{Cap}_\theta(\{s\},X_I\setminus\{s\})
&\le\frac12\,2^{-(1+\varepsilon)n}.
\end{align}

Moreover, if \(4\theta\delta m\le(1-\zeta)n\log2\) for some
\(0<\zeta<1\), then

\begin{equation}
\label{eq-lpn-residual-tilt-supercritical-gap}
\gamma_\theta
\le\frac{2^{-\zeta n}}{1-2^{-n}}.
\end{equation}

Equations
\cref{eq-lpn-residual-tilt-subcritical-condition,eq-lpn-residual-tilt-supercritical-condition}
locate the transition at the explicit scale
\(
\theta m(1/2-\eta)\asymp n\log2
\), with its finite-sample uncertainty carried by \(2\delta m\).

\end{corollary}

\begin{proof}

Apply \cref{eq-controlled-gibbs-base-density-hit-bound} with
\(
\mu_I(s)=2^{-n}
\)
and
\(
\operatorname{Cap}_{\Phi_{P_0}}(\{s\},X_I\setminus\{s\})
=2^{-n}q_{0,s}
\)
to obtain \cref{eq-lpn-generic-gibbs-hit-envelope}.
Under \(D(s)=\min D\), apply instead
\cref{eq-controlled-gibbs-target-minimum-hit-bound} to obtain
\cref{eq-lpn-target-minimum-gibbs-hit-envelope}.

Under \cref{eq-lpn-residual-tilt-subcritical-condition}, the upper target
mass bound in \cref{eq-lpn-metropolis-target-mass-window} gives
\[
\pi_\theta(s)
\le
\frac1{(2^n-1)e^{-\theta m\Delta_+}}
\le\frac{2^{-\varepsilon n}}{1-2^{-n}},
\]
which proves \cref{eq-lpn-residual-tilt-subcritical-mass}.  The stationary
entrance-count estimate
\cref{eq-controlled-capacity-finite-horizon-hit}, the initial-density
comparison, and
\cref{eq-lpn-metropolis-target-exit-capacity} give
\[
\Pr_{\rho_0}^{P_\theta}\{\tau_s\le H\}
\le R_\theta\pi_\theta(s)
\left(1+\frac H2e^{-\theta m\Delta_-}\right),
\]
which proves \cref{eq-lpn-residual-tilt-subcritical-hit}.

Under \cref{eq-lpn-residual-tilt-supercritical-condition}, the lower target
mass bound in \cref{eq-lpn-metropolis-target-mass-window} gives
\cref{eq-lpn-residual-tilt-supercritical-mass};
\cref{eq-lpn-metropolis-target-exit-capacity} gives
\cref{eq-lpn-residual-tilt-supercritical-exit-capacity}.  The mass bound is
strictly greater than one half.  Therefore the second branch of
\cref{eq-lpn-metropolis-gap-upper}, followed by
\(
e^{4\theta\delta m}\le2^{(1-\zeta)n}
\), proves \cref{eq-lpn-residual-tilt-supercritical-gap}.

\end{proof}

\begin{corollary}[Fast-fiber and clocked-defect access for LPN]
\label{cor-lpn-fast-fiber-clocked-access-split}

Fix one complete valid-Lift record over the LPN target \(T_I=\{s\}\), and
let its base selected kernel satisfy
\cref{thm-lpn-invariant-density-capacity-drift-envelope}.  Define the base
finite-horizon capacity envelope

\begin{equation}
\label{eq-lpn-lift-base-capacity-envelope}
U_{X,H}^{\rm Cap}
=
\min\!\left\{
1,\,
R_{\widetilde\mu}
\bigl(r_++H\kappa_+q_{0,s}\bigr)2^{-n}
\right\}.
\end{equation}

The two certificate branches of
\cref{thm-controlled-lift-transfer-certificate-split} specialize as follows.

\begin{enumerate}[label=\textup{(\roman*)}]
\item If the LPN lift has the represented fast-fiber data
      \((M_F,\kappa_F,\varepsilon_+,\varepsilon_-,\lambda)\) and satisfies
      the smallness condition in
      \cref{eq-controlled-certified-fast-fiber-parameters}, then
      \begin{equation}
      \label{eq-lpn-fast-fiber-discounted-access}
      A_{\{s\}}^\Omega(\lambda)
      \le
      A_{\{s\}}^X(\lambda)
      +\varepsilon_{0,s}^{\rm int}
      +\varepsilon_{\Omega,\lambda}^{\rm int}
      +\varepsilon_{X,\lambda}^{\rm int}
      +\|\ell\|_*\|g\|\Delta_{\rm fm}(\lambda),
      \end{equation}
      where
      \(\Delta_{\rm fm}(\lambda)\) is the explicit rational expression in
      \cref{eq-controlled-certified-fast-fiber-parameters}.  Any evaluated
      LPN base-access bound for \(A_{\{s\}}^X(\lambda)\) may be substituted
      in the same discounted-arrival units.

\item For represented stopped or absorbing clock slices, no fiber-mixing
      hypothesis is required.  With clocked defects
      \(E_t=P_t^\Omega U_{t+1}-U_tP_t^X\),
      \begin{equation}
      \label{eq-lpn-clocked-defect-capacity-access}
      p_{\{s\},H}^\Omega
      \le
      \min\!\left\{
      1,\,
      U_{X,H}^{\rm Cap}
      +\varepsilon_{\rm init}
      +\sum_{t=0}^{H-1}\|E_t\|
      +\varepsilon_{\rm gate}
      \right\}.
      \end{equation}
      For the residual Metropolis base record, the sharper subcritical
      expression in
      \cref{eq-lpn-residual-tilt-subcritical-hit} replaces
      \(U_{X,H}^{\rm Cap}\).  It depends on \((n,m,\eta)\) through
      \(2^{-\varepsilon n}\),
      \(\Delta_-=1/2-\eta-2\delta\), and the certified initialization
      density.
\end{enumerate}

When both branches are certified, they are retained separately.  They may
be minimized only after the same record supplies the represented conversion
between discounted arrival and the selected finite horizon.  Exact
valid-Lift initialization, target pullback, and intertwining set the corresponding
interface and defect terms to zero.

\end{corollary}

\begin{proof}

Part~\textup{(i)} is
\cref{eq-controlled-certified-fast-fiber-arrival} with the LPN singleton
target.  For part~\textup{(ii)}, combine
\cref{eq-controlled-certified-clocked-defect-arrival} with the base access
bound \cref{eq-lpn-invariant-density-hit-envelope}.  The residual
Metropolis substitution is
\cref{eq-lpn-residual-tilt-subcritical-hit}.  The unit-conversion and exact-
interface statements are the final clauses of
\cref{thm-controlled-lift-transfer-certificate-split}.

\end{proof}

\subsection{Sequential information contraction for the LPN presentation}
\label{subsec-lpn-sequential-information-contraction}

\begin{theorem}[LPN pre-hit information and target-contact envelope]
\label{thm-lpn-prehit-information-contact-envelope}

Let \(S\) be the uniform secret in the random binary LPN presentation of
\cref{def-lpn-task-object}, and take the initial public record to be
\(
\mathsf Z_0=(A,b)
\).
Consider one complete temporally admissible prefix before a represented
next-candidate selector is sampled.  Enumerate its target-dependent reveals
as in \cref{thm-pre-hit-sequential-sdpi-information-contraction}, and retain
the corresponding contracted increments \(d_1,\ldots,d_L\).  Define

\begin{equation}
\label{eq-lpn-prehit-information-budget}
J_{n,m,\eta,t}^{\mathrm{LPN}}
=
\min\!\left\{n,\,m\bigl(1-H_2(\eta)\bigr)\right\}
+\sum_{j=1}^{L}d_j.
\end{equation}

Then

\begin{equation}
\label{eq-lpn-prehit-prefix-information}
I(S;\mathcal H_t)
\le J_{n,m,\eta,t}^{\mathrm{LPN}}.
\end{equation}

Let \(K_t(\cdot\mid\mathcal H_t)\) be the represented next-candidate
kernel, let \(C_{t+1}\) be its sampled candidate, and put

\[
p_t^{\mathrm{contact}}
=\Pr\{\tau_s>t,\ C_{t+1}=S\},
\qquad
\tau_s=\inf\{t\ge1:C_t=S\}.
\]

The exact-target sectors are disjoint and have overlap one.  Hence

\begin{equation}
\label{eq-lpn-prehit-information-contact}
p_t^{\mathrm{contact}}
\le
\min\!\left\{
1,\,
2^{-n}
+\sqrt{\frac{\log2}{2}
J_{n,m,\eta,t}^{\mathrm{LPN}}}
\right\}.
\end{equation}

For a horizon \(H\), with the ordered reveal ledger recomputed at each
prefix,

\begin{equation}
\label{eq-lpn-prehit-information-cumulative-contact}
\Pr\{\tau_s\le H\}
\le
\min\!\left\{
1,\,
\sum_{t=0}^{H-1}
\left(
2^{-n}
+\sqrt{\frac{\log2}{2}
J_{n,m,\eta,t}^{\mathrm{LPN}}}
\right)
\right\}.
\end{equation}

For the terminal branch, let \(\mathcal R_H^{\rm full}\) index every
target-dependent reveal after \((A,b)\) through the complete terminal
transcript, including any verifier, hit-selection, feedback, reward, or
oracle coordinate, and let \(c_{q,r}\) be its represented conditional
capacity certificate.  Define
\begin{equation}
\label{eq-lpn-full-transcript-information-budget}
J_{n,m,\eta,H}^{\mathrm{LPN,full}}
=
\min\!\left\{n,\,m\bigl(1-H_2(\eta)\bigr)\right\}
+\sum_{(q,r)\in\mathcal R_H^{\rm full}}c_{q,r}.
\end{equation}
If the terminal record contains the candidate reached at an exact target
contact and all coordinates used to identify that contact occur in this
full ledger, the represented-decoder branch gives

\begin{equation}
\label{eq-lpn-prehit-information-fano-contact}
\Pr\{\tau_s\le H\}
\le
\min\!\left\{1,
\frac{J_{n,m,\eta,H}^{\mathrm{LPN,full}}+1}{n}\right\}.
\end{equation}

When no target-dependent reveal is adjoined after \((A,b)\), every retained
coordinate is a represented postprocessing of the public LPN experiment and
the sum in \cref{eq-lpn-prehit-information-budget} is zero.  In that case an
exact-secret output obeys, for every \(0<\epsilon<\eta<1/2\), the sharper
strong converse

\begin{equation}
\label{eq-lpn-prehit-information-strong-converse}
\Pr\{\widehat S=S\}
\le
\min\!\left\{
1,\,
2e^{-2m\epsilon^2}
+2^{m-n-m(H_2(\eta)-\epsilon L_\eta)}
\right\},
\qquad
L_\eta=\log_2\frac{1-\eta}{\eta}.
\end{equation}

Every post-public verifier flag, feedback value, reward, or admissible oracle
answer used before the selector occupies an ordered reveal slot.  Its query,
channel, contraction certificate, returned value, memory, precision, and
complete evaluation cost are charged in the same prefix.  Public
postprocessing and fresh target-independent randomness contribute no new
information slot.

\end{theorem}

\begin{proof}

Because \(A\) is independent of \(S\), the binary symmetric channel bound in
the proof of \cref{thm-lpn-information-threshold} gives
\[
I(S;A,b)=I(S;b\mid A)
\le\min\!\left\{n,\,m(1-H_2(\eta))\right\}.
\]
Apply
\cref{eq-pre-hit-sequential-total-information} to the ordered reveals.  This
proves \cref{eq-lpn-prehit-prefix-information}.

Use the analytic target relation
\(V_s(x)=\mathbf1_{\{x=s\}}\).  For every candidate \(x\), exactly one of
the \(2^n\) target indices verifies it, so the overlap coordinate in
\cref{eq-pre-hit-information-target-overlap} is
\(\Lambda_t=1\).  Substitution in
\cref{eq-pre-hit-information-overlap-contact} proves
\cref{eq-lpn-prehit-information-contact}.  The exact first-contact events
are disjoint; \cref{eq-pre-hit-information-cumulative-hit} therefore gives
\cref{eq-lpn-prehit-information-cumulative-contact}.  Retaining the
target-contact candidate gives the identification map required in
\cref{eq-pre-hit-information-hit-fano}.  The binary-symmetric-channel bound
for \(I(S;A,b)\), followed by
\cref{thm-controlled-multichannel-information-bound} on every slot in
\(\mathcal R_H^{\rm full}\), gives
\[
I(S;\mathsf Z_H^{\rm full})
\le J_{n,m,\eta,H}^{\mathrm{LPN,full}}.
\]
Substitution in \cref{eq-pre-hit-information-hit-fano} proves
\cref{eq-lpn-prehit-information-fano-contact}.

If there is no later target-dependent reveal, the complete output is a
postprocessing of \((A,b)\) and target-independent randomness.  Data
processing followed by
\cref{eq-lpn-information-strong-converse} proves
\cref{eq-lpn-prehit-information-strong-converse}.  The final charging
statement is the LPN specialization of the temporal and oracle clauses in
\cref{thm-pre-hit-sequential-sdpi-information-contraction}.

\end{proof}

\subsection{Blackwell envelopes for LPN pre-hit readouts}
\label{subsec-lpn-blackwell-prehit-readout-envelopes}

The represented Blackwell order compares LPN readouts by causal simulation.
It transfers a numerical target-decision bound only in the direction from a
reference experiment to its represented garblings.  The following
specialization records the three reference experiments used below.

\begin{theorem}[LPN target-decision bounds under represented garbling]
\label{thm-lpn-prehit-blackwell-target-decision-bounds}

Let
\(
u_{\rm id}(s,x)=\mathbf 1_{\{x=s\}}
\)
be the exact-secret decision utility of
\cref{def-controlled-prehit-target-decision-value}, and let
\(B^\sharp=\psi_{\rm BW}(B,n,m,H)\) be the complete cost enlargement of
\cref{thm-controlled-prehit-blackwell-value-monotonicity}.  The following
statements hold for the random binary LPN presentation of
\cref{def-lpn-task-object}.

\begin{enumerate}[label=\textup{(\roman*)}]
\item Let \(\mathsf E_0\) be a target-permutation-neutral,
      residual-independent pre-hit reference experiment.  Its time-\(t\)
      decision is made on the survival sector and can select one remaining
      candidate.  If \(\mathsf E\) is a represented causal garbling of
      \(\mathsf E_0\), then
      \begin{equation}
      \label{eq-lpn-blackwell-residual-independent-value}
      \operatorname{DVal}^{\mathrm{sat}}_{u_{\rm id},t}
       (\mathsf E;B)
      \le
      \operatorname{DVal}^{\mathrm{sat}}_{u_{\rm id},t}
       (\mathsf E_0;B^\sharp)
      \le 2^{-n}.
      \end{equation}

\item Work conditionally on a matrix \(A\) of full column rank.  Let
      \(S(A)\) be the represented deterministic pivot-row set from
      \cref{cor-lpn-three-dynamical-stress-tests}, and define the pivot
      experiment
      \begin{equation}
      \label{eq-lpn-blackwell-pivot-experiment}
      \mathsf E_{\rm piv}
      =\bigl(A,S(A),b_{S(A)}\bigr).
      \end{equation}
      Before any target-dependent verifier outcome is adjoined, its Bayes
      exact-secret decision value is
      \begin{equation}
      \label{eq-lpn-blackwell-pivot-bayes-value}
      \sup_{K}
      \mathbb E\!\left[
        \int u_{\rm id}(s,x)K(dx\mid\mathsf E_{\rm piv})
        \given A,\operatorname{rank}A=n
      \right]
      =(1-\eta)^n.
      \end{equation}
      Here the supremum ranges over all probability kernels, so it is also an
      upper bound for every saturated budget.  The represented rule
      \(x=A_{S(A)}^{-1}b_{S(A)}\) attains equality once the charged ceiling
      contains pivot extraction and binary elimination, of complete cost
      \begin{equation}
      \label{eq-lpn-blackwell-pivot-complete-cost}
      B_{\rm piv}(n,m)\preceq mn^2+n^3.
      \end{equation}
      Consequently, every represented causal garbling \(\mathsf E\) of
      \(\mathsf E_{\rm piv}\) satisfies
      \begin{equation}
      \label{eq-lpn-blackwell-pivot-garbling-bound}
      \operatorname{DVal}^{\mathrm{sat}}_{u_{\rm id},0}
       (\mathsf E;B)
      \le (1-\eta)^n.
      \end{equation}
      If the pivot experiment is completed by a represented failure symbol
      when \(\operatorname{rank}A<n\), then every decision rule on the
      unconditional experiment satisfies
      \begin{equation}
      \label{eq-lpn-blackwell-pivot-unconditional-bound}
      \operatorname{DVal}^{\mathrm{sat}}_{u_{\rm id},0}
       (\mathsf E;B)
      \le
      \min\!\left\{1,\,2^{n-m}+(1-\eta)^n\right\}.
      \end{equation}

\item Let \(\mathsf E_{\rm pub}=(A,b)\) be the complete public LPN
      experiment.  Every represented causal garbling \(\mathsf E\) of
      \(\mathsf E_{\rm pub}\), including its target-independent random bits
      and public postprocessing, obeys the finite Fano inequality
      \begin{equation}
      \label{eq-lpn-blackwell-public-fano}
      H_2(1-p_B)
      +(1-p_B)\log_2(2^n-1)
      \ge n-m\bigl(1-H_2(\eta)\bigr),
      \end{equation}
      where
      \(
      p_B=\operatorname{DVal}^{\mathrm{sat}}_{u_{\rm id},0}
      (\mathsf E;B)
      \).
      For every \(0<\epsilon<\eta<1/2\), the corresponding explicit bound is
      \begin{equation}
      \label{eq-lpn-blackwell-public-strong-converse}
      p_B
      \le
      \min\!\left\{
      1,
      2e^{-2m\epsilon^2}
      +2^{m-n-m\left(H_2(\eta)-\epsilon L_\eta\right)}
      \right\},
      \qquad
      L_\eta=\log_2\frac{1-\eta}{\eta}.
      \end{equation}
\end{enumerate}

All three comparisons retain the complete derivation and garbling costs.
Thus they apply at the one charged enlarged ceiling rather than adjoining a
free readout.

\end{theorem}

\begin{proof}

For (i), condition on the complete residual-independent pre-hit record and
on survival.  By
\cref{lem-lpn-syndrome-free-ambient-geom-bound}, the secret is uniform on the
remaining untested candidates.  If \(N_t\) candidates remain, the conditional
best exact decision value is \(N_t^{-1}\), while the survival probability is
at most \(N_t/2^n\).  Their product is at most \(2^{-n}\).  Apply
\cref{thm-controlled-prehit-blackwell-value-monotonicity} to every represented
causal garbling and charge its simulation at \(B^\sharp\).  This proves
\cref{eq-lpn-blackwell-residual-independent-value}.

For (ii), put \(A_S=A_{S(A)}\) and \(Y=b_{S(A)}\).  Conditional on \(A\) and
full column rank,
\[
Y=A_Ss+E_S,
\qquad E_S\sim\operatorname{Bernoulli}(\eta)^{\otimes n},
\]
and \(A_S\) is invertible.  Since \(s\) is uniform, Bayes' formula gives,
for every \(x,y\in\mathbb F_2^n\),
\begin{equation}
\label{eq-lpn-blackwell-pivot-posterior}
\Pr\{s=x\mid A,Y=y,\operatorname{rank}A=n\}
=
(1-\eta)^{n-|y+A_Sx|_1}\eta^{|y+A_Sx|_1}.
\end{equation}
For \(0\le\eta<1/2\), its maximum over \(x\) is \((1-\eta)^n\), attained
uniquely at \(x=A_S^{-1}y\); at \(\eta=1/2\), every posterior atom equals
\(2^{-n}=(1-\eta)^n\).  Integrating the largest posterior atom proves
\cref{eq-lpn-blackwell-pivot-bayes-value}.  Pivot extraction scans the
\(m\times n\) matrix and binary elimination uses cubic arithmetic cost,
which gives \cref{eq-lpn-blackwell-pivot-complete-cost}.  Equality of the
represented rule is also the exact calculation in
\cref{eq-lpn-subset-ge-success}.  Blackwell monotonicity now proves
\cref{eq-lpn-blackwell-pivot-garbling-bound}.  On rank failure, bound the
decision utility by one; the rank-failure estimate used in
\cref{eq-lpn-subset-ge-unconditional-success} is at most \(2^{n-m}\).
Adding the two events proves
\cref{eq-lpn-blackwell-pivot-unconditional-bound}.

For (iii), compose a decision rule for \(\mathsf E\) with the represented
causal garbling from \((A,b)\).  The result is a randomized decoder of the
complete LPN sample, with exactly the same exact-recovery probability by
\cref{eq-controlled-prehit-blackwell-utility-identity}.  Apply
\cref{eq-lpn-fano-finite-converse,eq-lpn-information-strong-converse} from
\cref{thm-lpn-information-threshold}.  The first application gives
\cref{eq-lpn-blackwell-public-fano}; the second gives
\cref{eq-lpn-blackwell-public-strong-converse}.  Data processing is already
quantified over every measurable postprocessing, so the conclusion is
independent of the represented garbling's syntactic form.

\end{proof}

\begin{corollary}[Direction of the LPN Blackwell comparisons]
\label{cor-lpn-blackwell-comparison-direction}

The pivot experiment is a represented statistic of the public experiment:
\begin{equation}
\label{eq-lpn-blackwell-public-to-pivot}
\mathsf E_{\rm pub}
\longrightarrow
\mathsf E_{\rm piv},
\qquad
(A,b)\longmapsto\bigl(A,S(A),b_{S(A)}\bigr).
\end{equation}
Hence
\begin{equation}
\label{eq-lpn-blackwell-public-pivot-value-order}
\operatorname{DVal}^{\mathrm{sat}}_{u_{\rm id},0}
 (\mathsf E_{\rm piv};B)
\le
\operatorname{DVal}^{\mathrm{sat}}_{u_{\rm id},0}
 (\mathsf E_{\rm pub};\psi_{\rm BW}(B,n,m,0)).
\end{equation}
An upper bound on the full public experiment passes back from
\(\mathsf E_{\rm piv}\) whenever an affordable represented pre-hit
reconstruction kernel satisfies
\cref{eq-controlled-represented-prehit-sufficiency}; in that case
\cref{cor-controlled-represented-prehit-sufficiency} supplies the reverse
comparison.  A separate direct public-record bound gives the other valid
route.  In the absence of either route, the pivot value \((1-\eta)^n\)
bounds the pivot experiment and its garblings, whereas the full public
experiment and its garblings are bounded by
\cref{eq-lpn-blackwell-public-fano,eq-lpn-blackwell-public-strong-converse}.
The complete-record bounds therefore apply directly to decision rules using
\((A,b)\).

\end{corollary}

\begin{proof}

The displayed map is represented before the decision and has polynomial
complete cost.  Apply
\cref{thm-controlled-prehit-blackwell-value-monotonicity} in the displayed
direction to obtain
\cref{eq-lpn-blackwell-public-pivot-value-order}.  The affordable
reconstruction criterion of
\cref{cor-controlled-represented-prehit-sufficiency} supplies the stated
sufficient reverse comparison.  The remaining bounds
are \cref{thm-lpn-prehit-blackwell-target-decision-bounds}.

\end{proof}

\begin{corollary}[Recordwise readout-agnostic LPN envelope]
\label{cor-lpn-recordwise-readout-agnostic-envelope}

Fix one complete prospective LPN record, one target-arrival event
\(E_H=\{\tau_s\le H\}\), and one common probability law
\(\mathbb P_{\mathcal R}\).  The law may condition on a fixed presented
instance or may include the random LPN presentation, but every candidate
below is first expressed under this same law.

A pointwise bound
\(\Pr(E_H\mid I)\le U(I)\) contributes
\(\mathbb E_{\mathbb P_{\mathcal R}}U(I)\).  If it is certified only on a
presentation event \(\mathcal G\), it contributes

\begin{equation}
\label{eq-lpn-common-law-good-event-conversion}
\mathbb P_{\mathcal R}(\mathcal G^c)
+\mathbb E_{\mathbb P_{\mathcal R}}
 \left[\mathbf1_{\mathcal G}U(I)\right].
\end{equation}

Thus the residual-Metropolis bound retains the explicit failure probability
of \(\mathcal F_\delta\), and the conditional pivot bound retains either the
full-rank conditioning or the rank-failure term in
\cref{eq-lpn-blackwell-pivot-unconditional-bound}.  Denote the resulting
common-law candidates as follows:

\begin{itemize}
\item \(\overline U_{\rm cap}\) is the converted applicable right side of
      \cref{eq-lpn-invariant-density-hit-envelope} or
      \cref{eq-lpn-residual-tilt-subcritical-hit};
\item \(\overline U_{\rm lift}\) is the converted right side of
      \cref{eq-lpn-clocked-defect-capacity-access} when the arrival occurs on
      a valid lifted carrier;
\item \(\overline U_{\rm info}\) is the smaller applicable right side of
      \cref{eq-lpn-prehit-information-cumulative-contact,eq-lpn-prehit-information-fano-contact};
\item \(\overline U_{\rm cur}\) is the applicable same-record minimum in
      \cref{eq-observation-information-capacity-minimum} for a complete
      sparse-reward observation record, with
      \(\mathsf Z_0=(A,b)\), \(\Theta=S\), and target
      \(T_S=\{S\}\);
\item \(\overline U_{\rm BW}\) is available when either the complete
      terminal experiment, including every target-dependent reveal through
      time \(H\), is a represented causal garbling of a bounded reference
      experiment and a represented terminal map converts target arrival into
      its decision utility, or the survival-weighted Blackwell decision
      bounds are summed over the disjoint first-contact times;
\item \(\overline U_{\Caus}\) is a converted contact bound obtained from the same
      selected record through
      \cref{eq-lpn-linear-noise-corrected-caus-contact} or a represented
      aiming-to-contact converter as in
      \cref{cor-lpn-traffic-horizon-caus-aiming}.
\end{itemize}

Set an unavailable candidate equal to the neutral value one.  A candidate is
available only when its hypotheses, interface, units, temporal ancestry, and
complete-cost certificate are present in this record and it bounds
\(\mathbb P_{\mathcal R}(E_H)\).  Then

\begin{equation}
\label{eq-lpn-recordwise-readout-agnostic-envelope}
\mathbb P_{\mathcal R}\{\tau_s\le H\}
\le
\min\{\overline U_{\rm cap},\overline U_{\rm lift},
\overline U_{\rm info},\overline U_{\rm cur},
\overline U_{\rm BW},\overline U_{\Caus}\}.
\end{equation}

The minimum is taken before saturation.  Saturation then takes the supremum
over complete records at the common charged ceiling.  In particular, no
capacity, information, sparse-reward observation, Blackwell, Lift, or Caus
certificate from one record is combined with a coordinate from another.

\end{corollary}

\begin{proof}

The pointwise and good-event conversions follow from the tower property and
the decomposition of \(E_H\) over \(\mathcal G\) and \(\mathcal G^c\).
For the second Blackwell route, sum the survival-weighted decision values
over the disjoint first-contact times.  After these conversions, each
available candidate is an upper bound for the same scalar probability under
\(\mathbb P_{\mathcal R}\), so their minimum is an upper bound.  The
neutral-default and order-of-operations statements are exactly the
coordinatewise convention of
\cref{def-coordinatewise-dynamical-certificate}.

\end{proof}

\begin{figure}[tbp]
\centering
\resizebox{\linewidth}{!}{%
\begin{tikzpicture}[fw diagram,node distance=8mm and 10mm]
  \node[fw source,text width=2.7cm] (pub)
    {public LPN record\\\((A,b;n,m,\eta)\)};
  \node[fw geom,text width=3.0cm] (cap) at (4.2,3.3)
    {invariant law and conductance\\capacity / drift};
  \node[fw use,text width=3.0cm] (cur) at (4.2,1.1)
    {observation regularity\\fiber information};
  \node[fw provenance,text width=3.0cm] (info) at (4.2,-1.1)
    {ordered pre-hit reveals\\SDPI / information};
  \node[fw provenance,text width=3.0cm] (bw) at (4.2,-3.3)
    {represented experiments\\Blackwell order};
  \node[fw terminal,text width=3.25cm] (nums) at (8.6,0)
    {recordwise bounds\\
     \(\overline U_{\rm cap},\overline U_{\rm lift},\overline U_{\rm info}\)\\
     \(\overline U_{\rm cur},\overline U_{\rm BW},\overline U_{\Caus}\)};
  \node[fw terminal,text width=2.9cm] (hit) at (12.3,0)
    {same-record minimum\\target-arrival envelope};
  \draw[fw arrow] (pub)--(cap);
  \draw[fw arrow] (pub)--(cur);
  \draw[fw arrow] (pub)--(info);
  \draw[fw arrow] (pub)--(bw);
  \draw[fw arrow] (cap)--(nums);
  \draw[fw arrow] (cur)--(nums);
  \draw[fw arrow] (info)--(nums);
  \draw[fw arrow] (bw)--(nums);
  \draw[fw arrow] (nums)--(hit);
\end{tikzpicture}%
}
\caption{Four recordwise routes for LPN.  Capacity sees the selected
invariant law and conductance, sparse-reward accounting sees observation
regularity and within-cell information, SDPI sees the ordered reveal
channels, and Blackwell order sees represented experiment degradation.  Each route
produces a number for one complete record and one common probability law;
only then are applicable bounds minimized and passed to the existing
target-arrival ledger.}
\label{fig-lpn-readout-agnostic-four-route-envelope}
\end{figure}

\begin{corollary}[Saturated quantitative prospective certificate envelope for LPN]
\label{cor-lpn-saturated-quantitative-prospective-certificate-envelope}

Fix \(n,m\ge1\), \(0<\eta<1/2\), budget \(B\), and refined horizon \(H\).
Put \(B_{\rm LPN}^{\rm cert}=\Psi_{\rm cert}(B,H,n,m)\).  For a complete
pre-hit selector record
\(R\in\mathscr R_{I,t}(B_{\rm LPN}^{\rm cert})\), set
\(\sigma_{R,t}=\Pr_R\{\tau_s>t\}\).  Whenever the corresponding same-record
certificates are present, define
\begin{align}
\label{eq-lpn-saturated-certificate-capacity-candidate}
C_{R,t}^{\rm cap}
&=
\sigma_{R,t}\min\!\left\{1,\,
\|f_{R,t}\|_{L^p(\widetilde\mu_R)}
\left[\min\!\left\{1,\,
2^{-n}(r_{R,+}+\kappa_{R,+}q_{0,R,s})
\right\}\right]^{1/q}\right\},\\
\label{eq-lpn-saturated-certificate-information-candidate}
C_{R,t}^{\rm info}
&=
\min\!\left\{1,\,2^{-n}
+\sqrt{\frac{\log2}{2}J_{n,m,\eta,t}^{\rm LPN}(R)}\right\},\\
J_{n,m,\eta,t}^{\rm LPN}(R)
&=
\min\{n,m(1-H_2(\eta))\}
+\sum_{j=1}^{L_{R,t}}d_{R,t,j},
\end{align}
where \(q_{0,R,s}=1-P_{0,R}(s,s)\),
\(f_{R,t}=d\mathcal L_R(X_t\mid\tau_s>t)/d\widetilde\mu_R\),
\(1<p\le\infty\), and \(q=p/(p-1)\).

For a complete sparse-reward observation record, let
\(C_{\varepsilon,R}=q_{\varepsilon,R}(s)\), let
\(\beta_{R,t}^{\rm fib}\) be the decoupled posterior contact in
\cref{eq-observation-within-cell-neutral-contact} with
\(\mathsf Z_0=(A,b)\), and let
\(J_{R,t}^{\rm fib,\uparrow}\) be the charged ordered-reveal upper bound on
\[
I\!\left(S;\mathcal H_{R,t}^{\rm obs}
\mid A,b,C_{\varepsilon,R}\right).
\]
Set
\begin{equation}
\label{eq-lpn-saturated-certificate-curiosity-candidate}
C_{R,t}^{\rm cur}
=
\min\!\left\{1,
\beta_{R,t}^{\rm fib}
+\sqrt{\frac{\log2}{2}J_{R,t}^{\rm fib,\uparrow}}
\right\}.
\end{equation}
This candidate has value one when the complete observation record or its
information comparison is unavailable.

For a complete master active-sampling record, let
\(\beta_{R,t}^{\rm act}\) and
\(J_{R,t}^{\rm act,\uparrow}\) be the decoupled next-contact and ordered
acquisition-information coordinates of
\cref{thm-adaptive-acquisition-information-target-contact}.  Set
\begin{equation}
\label{eq-lpn-saturated-certificate-active-candidate}
C_{R,t}^{\rm act}
=
\min\!\left\{1,
\beta_{R,t}^{\rm act}
+\sqrt{\frac{\log2}{2}J_{R,t}^{\rm act,\uparrow}}
\right\},
\end{equation}
with neutral value one when that complete record is unavailable.

Let \(C_{R,t}^{\rm learn}\) be the complete LPN specialization of
\cref{eq-source-resolved-learning-target-gain-minimum}:
\begin{equation}
\label{eq-lpn-saturated-certificate-source-resolved-learning-candidate}
C_{R,t}^{\rm learn}
=
\min\left\{
1,\,
U_{R,t}^{\rm stat\to tgt},\,
\beta_{R,t}^{+}
+\sqrt{\frac{\log2}{2}
\sum_{a\in\mathfrak J_{\Abs}^{5}}
J_{I,R_{\le t}}^{(a)}},\,
eH_B(n)\sum_{a\in\mathfrak J_{\Abs}^{5}}
\mathcal L_{I,R_{\le t}}^{(a)},\,
U_{R,t}^{\rm dyn\to tgt}
\right\}.
\end{equation}
For an LPN binary reveal, every summand uses
\(\vartheta_j(\eta)=(1-2\eta)^2\).  The common ceiling charges the public
record, learner, locator, optimizer, posterior or memory state, selected
kernel, Caus/Lift processing, verifier, and deployment converter.

Let \(C_{R,t}^{\rm BW}\) be the smallest applicable bound among
\begin{equation}
\label{eq-lpn-saturated-certificate-blackwell-candidates}
2^{-n},\qquad
\min\{1,2^{n-m}+(1-\eta)^n\},\qquad
\min\!\left\{1,\,
2e^{-2m\epsilon^2}
+2^{m-n-m(H_2(\eta)-\epsilon L_\eta)}\right\},
\end{equation}
where \(0<\epsilon<\eta\) and
\(L_\eta=\log_2((1-\eta)/\eta)\), with applicability determined by
\cref{thm-lpn-prehit-blackwell-target-decision-bounds}.  Let
\(C_{R,t}^{\rm src}\) be the smallest applicable same-record comparison
\begin{equation}
\label{eq-lpn-saturated-certificate-source-candidates}
\min\{1,K_R(2^{-n}+(1-\eta)^n)+\delta_R\},
\qquad
\min\{1,K_R^{\rm dir}V_{\mathscr C_R^{\rm dir},I}
             +\delta_R^{\rm dir}\}.
\end{equation}
The first requires
\cref{eq-inherited-prospective-value-comparison} and
\cref{thm-lpn-residual-ancestry-attention-bound}; the second requires
\cref{eq-direct-prospective-value-comparison}.  Finally, let
\(C_{R,t}^{\rm dyn}\) be the smallest applicable same-record next-contact
bound in
\cref{cor-lpn-recordwise-readout-agnostic-envelope,%
cor-lpn-fast-fiber-clocked-access-split,%
cor-lpn-noise-indexed-caus-lift-interfaces}.
Set every unavailable candidate to one and put
\[
C_{R,t}^{\rm LPN}
=\min\{C_{R,t}^{\rm cap},C_{R,t}^{\rm info},C_{R,t}^{\rm BW},
       C_{R,t}^{\rm src},C_{R,t}^{\rm dyn},C_{R,t}^{\rm cur},
       C_{R,t}^{\rm act},C_{R,t}^{\rm learn}\}.
\]
Then
\begin{align}
\label{eq-lpn-saturated-selector-certificate-envelope}
\operatorname{Sel}_{I,n,\eta}^{\rm sat}(B)
&\le
\max_{0\le t<H}\sup_{R\in\mathscr R_{I,t}(B_{\rm LPN}^{\rm cert})}
C_{R,t}^{\rm LPN},\\
\label{eq-lpn-saturated-success-certificate-envelope}
\operatorname{Succ}_{I,n,\eta}(B)
&\le
\min\!\left\{1,\,
H2^{-n}
+\sum_{t=0}^{H-1}
\sup_{R\in\mathscr R_{I,t}(B_{\rm LPN}^{\rm cert})}
C_{R,t}^{\rm LPN}\right\}.
\end{align}
Every term retains its full
\((n,m,\eta,H,B)\)-dependence and the applicable density, conductance,
SDPI, observation-resolution, within-cell information, Blackwell, Caus,
Lift, learning, and good-presentation parameters.
If the recordwise envelope is at most
\(\varepsilon_{n,m,\eta,B,H}\), then
\[
\operatorname{Sel}_{I,n,\eta}^{\rm sat}(B)
\le\varepsilon_{n,m,\eta,B,H},
\qquad
\operatorname{Succ}_{I,n,\eta}(B)
\le\min\{1,H(2^{-n}+\varepsilon_{n,m,\eta,B,H})\}.
\]

\end{corollary}

\begin{proof}

For the uniform LPN carrier,
\(\mu_I(\{s\})=2^{-n}\) and the reference singleton capacity is
\(2^{-n}q_{0,R,s}\).  Thus the capacity candidate is
\cref{eq-saturated-certificate-capacity-candidate}.  The information,
Blackwell, and source candidates follow from
\cref{eq-lpn-prehit-information-contact,%
thm-lpn-prehit-blackwell-target-decision-bounds,%
eq-lpn-complete-ambient-geom-source-bound,%
thm-no-uncharged-prospective-value-amplification}.
The sparse-reward observation candidate is
\cref{eq-observation-fiber-information-contact-bound} with the LPN target
index and initial public record substituted as displayed.
The active-acquisition candidate is
\cref{eq-active-acquisition-next-contact} on the same complete prefix.
The source-resolved learning candidate is
\cref{thm-source-resolved-saturated-learning-compilation} with the LPN
binary-noise coefficient and complete LPN record substituted.
Apply
\cref{thm-saturated-same-record-prospective-certificate-envelope}.
The uniform neutral next-test baseline contributes at most \(2^{-n}\) per
refined step.  Summation gives the success bound.

\end{proof}

\subsection{Sparse-reward observation and transport accounting for LPN}
\label{subsec-lpn-sparse-reward-observation-transport}

\begin{theorem}[Complete sparse-reward observation accounting for LPN]
\label{thm-lpn-complete-sparse-reward-observation-accounting}

Fix \(0<\eta<\tau<1/2\), \(\delta>0\), and a budget in the declared
polynomial saturated ceiling.  Work on the intersection of
\(\mathcal G_{\tau,\delta}\) with the full-rank and random-code-rigidity
event of \cref{thm-lpn-random-code-rigidity-final}.  Let \(S\) be the
uniform secret, set
\(X_I=\mathbb F_2^n\), \(T_S=\{S\}\), and take
\(\mathsf Z_0=(A,b)\).  Let \(\mathcal R_{\rm cur}^{\rm LPN}\) be a
complete record from
\cref{def-complete-sparse-reward-observation-record} whose observation is a
represented map
\[
\phi_I:\mathbb F_2^n\longrightarrow(\mathcal O,d_{\mathcal O})
\]
derived from the public LPN presentation, and put
\(q_{\varepsilon,I}=\kappa_\varepsilon\circ\phi_I\) and
\(C_{\varepsilon,I}=q_{\varepsilon,I}(S)\).  The public external reward is
the verifier value revealed after a candidate is tested.  On
\(\mathcal G_{\tau,\delta}\), it is zero throughout every surviving
pre-hit prefix.

For the complete observation history before the next selector, define
\begin{equation}
\label{eq-lpn-curiosity-within-cell-information}
J_{I,t}^{\rm cur}
=
I\!\left(
S;\mathcal H_{I,t}^{\rm obs}
\mid A,b,C_{\varepsilon,I}
\right).
\end{equation}
If every later target-dependent coordinate is placed in the conditional
ordered-reveal ledger, then
\begin{equation}
\label{eq-lpn-curiosity-ordered-reveal-bound}
J_{I,t}^{\rm cur}
\le
J_{I,t}^{\rm cur,\uparrow}
=
\sum_{j=1}^{L_{I,t}}
\mathbb E\min\{c_{I,t,j},
\vartheta_{I,t,j}a_{I,t,j}\}.
\end{equation}
Public postprocessing of \((A,b)\), intrinsic rewards, predictor updates,
exploration memories, and fresh target-independent randomness add no term to
this sum.  Their construction and deployment remain in the complete cost
\begin{align}
\label{eq-lpn-complete-curiosity-cost}
B_{\rm cur}^{\rm LPN}
={}&B_{(A,b,\eta)}\oplus B_{\phi_I}\oplus B_{d_{\mathcal O}}
\oplus B_{\kappa,\varepsilon}\oplus B_{M_\varepsilon}
\oplus B_{\rm reg}\oplus B_{\rm info}\notag\\
&\oplus B_{\rm mem}\oplus B_{\rm pred}\oplus B_{\rm train}
\oplus B_{r^{\rm ext}}\oplus B_{\rm bonus}\oplus B_{\rm pol}\notag\\
&\oplus B_{\rm exec}\oplus B_{\rm clock}\oplus B_{\rm verify}
\oplus B_{\rm out}\oplus B_{\rm cap}\oplus B_{\rm comp}.
\end{align}
Here the last four analytic ledgers are expanded in
\cref{eq-complete-curiosity-regularity-cost,eq-complete-curiosity-information-cost,eq-complete-curiosity-capacity-cost,eq-complete-compensated-observation-cost}.

Let \(\mathbb Q_{I,t}^{\rm fib}\) retain
\((A,b,C_{\varepsilon,I},\mathcal H_{I,t}^{\rm obs})\) and resample
\(S\) from
\(\mathcal L(S\mid A,b,C_{\varepsilon,I})\), independently of the
observation history.  Put
\begin{equation}
\label{eq-lpn-curiosity-decoupled-contact}
\beta_{I,t}^{\rm cur}
=
\mathbb Q_{I,t}^{\rm fib}
\{\tau_s>t,\ X_{t+1}=S\}.
\end{equation}
Write \(R_{\le t}\) for this complete curiosity record truncated immediately
after its next test.  Then the actual survival-weighted contact and selector advantage satisfy
\begin{equation}
\label{eq-lpn-curiosity-contact-bound}
\bigl(\operatorname{Adv}_t(K_t;\nu_t)\bigr)_+
\le
\Pr\{\tau_s>t,\ X_{t+1}=S\}
\le
\min\!\left\{1,
\beta_{I,t}^{\rm cur}
+\sqrt{\frac{\log2}{2}J_{I,t}^{\rm cur,\uparrow}},
eH_B(n)\sum_{a\in\mathfrak J_{\Abs}^{5}}
\mathcal L_{I,R_{\le t}}^{(a)}
\right\}.
\end{equation}
Consequently,
\begin{equation}
\label{eq-lpn-curiosity-cumulative-contact-bound}
\Pr\{\tau_s\le H\}
\le
\min\!\left\{1,
\sum_{t=0}^{H-1}
\min\left\{1,
\beta_{I,t}^{\rm cur}
+\sqrt{\frac{\log2}{2}J_{I,t}^{\rm cur,\uparrow}},
eH_B(n)\sum_{a\in\mathfrak J_{\Abs}^{5}}
\mathcal L_{I,R_{\le t}}^{(a)}
\right\}\right\}.
\end{equation}

The master-profile assignment is exact.  An exact qualified use of
\(q_{\varepsilon,I}\) lies in the LPN exact quotient coordinates, and a
nonzero-defect qualified use lies in one of the two compensated coordinates.
These three coordinates vanish by
\cref{eq-lpn-exact-abs-source-vanishes,eq-lpn-exact-quotient-geom-closed,eq-lpn-compensated-source-vanishes}.
Every remaining observation-driven transition is therefore an ambient
full-carrier record, with no cell-dynamical credit.  Its target-arrival
coordinate is additionally bounded, whenever the corresponding same-record
certificate is present, by
\begin{align}
\label{eq-lpn-curiosity-uniform-capacity-bound}
\Pr_{\mu_0}\{\tau_s\le H\}
&\le R_0(H+1)2^{-n},\\
\label{eq-lpn-curiosity-comparison-capacity-bound}
\Pr_{\rho_0}^{\widetilde P}\{\tau_s\le H\}
&\le
\min\!\left\{1,
R_{\widetilde\mu}
\bigl(r_++H\kappa_+q_{0,s}\bigr)2^{-n}
\right\}.
\end{align}
The first line applies to a uniform-invariant selected kernel with
\(d\mu_0/d\mu_I\le R_0\); the second applies under the represented invariant
density and conductance comparisons of
\cref{thm-lpn-invariant-density-capacity-drift-envelope}.  For the residual
Metropolis record, the parameter-explicit replacement is
\cref{eq-lpn-residual-tilt-subcritical-hit}, with
\(\Delta_-=1/2-\eta-2\delta\) and
\(\Delta_+=1/2-\eta+2\delta\).

Finally, the observation-derived intrinsic signal obeys, on the same
stationary record,
\begin{equation}
\label{eq-lpn-curiosity-regularity-bound}
\mathbb E_{\mu P}b_{\phi_I}
\le
a_{\phi_I}+L_{\phi_I}
\sqrt{2\rho_{\phi_I}(P)
\operatorname{Var}_{\mu}^{d}(\phi_I)},
\end{equation}
with the empirical replacement in
\cref{eq-observation-regularity-learned-ratio,eq-observation-regularity-learned-novelty}.
Thus novelty magnitude, within-cell target information, controlled capacity,
and compensated dynamics are evaluated on one record and charged before
saturation.  The candidate in
\cref{eq-lpn-curiosity-contact-bound} is exactly
\(C_{R,t}^{\rm cur}\) in
\cref{eq-lpn-saturated-certificate-curiosity-candidate}, and hence enters
the complete LPN envelope
\cref{eq-lpn-saturated-selector-certificate-envelope,eq-lpn-saturated-success-certificate-envelope}.

Under the random LPN presentation, the same conclusions hold after adding
the explicit complement probability
\begin{equation}
\label{eq-lpn-curiosity-good-event-complement}
\varepsilon_{\mathrm{rig}}(n,m)
+\exp[-mD_{\mathrm{Ber}}(\tau\Vert\eta)]
+(2^n-1)2^{-m(1-H_2(\tau))}
+2^{n+1}e^{-2\delta^2m}.
\end{equation}

\end{theorem}

\begin{proof}

On \(\mathcal G_{\tau,\delta}\), the public verifier accepts precisely
\(s\), so the external reward is the constant zero vector on every surviving
prefix.  Apply
\cref{eq-sparse-external-reward-zero-information}.  Conditional on
\((A,b,C_{\varepsilon,I})\), public postprocessings and intrinsic-record
updates are causal postprocessings.  The conditional sequential SDPI theorem
\cref{thm-pre-hit-sequential-sdpi-information-contraction} therefore gives
\cref{eq-lpn-curiosity-ordered-reveal-bound}.  The conditional
decoupling-and-Pinsker argument of
\cref{thm-observation-regularity-sparse-reward-exploration} gives
\cref{eq-lpn-curiosity-contact-bound}'s information candidate.  Its
source-charge candidate is
\cref{eq-source-resolved-learning-target-gain-minimum} for the same
truncated curiosity record.  Taking the recordwise minimum gives the full
display; summing the disjoint first-contact
events gives \cref{eq-lpn-curiosity-cumulative-contact-bound}.

The mode assignment is
\cref{thm-exhaustive-sparse-reward-observation-transport-accounting}.
Substitution of the established LPN vanishing coordinates leaves its ambient
mode.  The native uniform-capacity estimate is
\cref{eq-lpn-controlled-capacity-hit-bound}, and the comparison estimate is
\cref{eq-lpn-invariant-density-hit-envelope}.  Observation regularity is
\cref{eq-observation-regularity-linear-novelty-bound}.  Every operation used
in these substitutions occurs in \cref{eq-lpn-complete-curiosity-cost};
constructor reuse follows
\cref{def-budget-constructor-accounting}.

The good-event complement is the union of
\cref{eq-lpn-rigidity-failure-probability,eq-lpn-dynamical-good-event-probability}.
Apply the unit bound there and the tower property on the intersection to
obtain the final family statement.

\end{proof}

\section{Optimal Structural Active Sampling for LPN}
\label{sec-lpn-optimal-structural-active-sampling}

The standard LPN presentation supplies the complete public record
\((A,b,\eta)\) once.  Active selection of its rows, residual coordinates,
features, candidate tests, or native transitions is therefore computation on
that record, not a new sampling channel.  Fresh random samples and chosen
parity queries are separate declared interfaces and are evaluated below as
presentation extensions.

The framework optimum is specialized in
\cref{thm-lpn-optimal-active-construction-equivalence}: the null-query active
profile is the standard finite-horizon saturated computation profile up to
the completely charged clock-and-certificate compilation.  Its
computational, sample, and noise dependence is evaluated in
\cref{cor-lpn-parameter-explicit-active-acquisition-regimes,cor-lpn-explicit-learned-active-optimum-loss}.

\begin{definition}[LPN master active-sampling ceiling]
\label{def-lpn-master-active-sampling-ceiling}

For a budget \(B\) and refined horizon \(H\), let
\[
B_{\rm LPN}^{\rm act}
=\Psi_{\rm act}(B,H,n,m,\eta)
\]
be a class-preserving common majorant of
\cref{eq-complete-master-active-sampling-cost} and the existing LPN source,
capacity, Blackwell, Lift, Caus, learning, verifier, and good-event
conversions.  Let
\(\mathscr R_{{\rm LPN},H}^{\rm act}(B_{\rm LPN}^{\rm act})\) be the
corresponding complete records with initial public record
\(\mathsf Z_0=(A,b)\), target index \(S\), carrier
\(\mathbb F_2^n\), and target \(T_S=\{S\}\).
Write \(\mathcal G_{\rm rig}\) for the full-rank and random-code-rigidity
event in \cref{thm-lpn-random-code-rigidity-final}, and set
\begin{equation}
\label{eq-lpn-master-active-good-event}
\mathcal G_{\rm act}
=\mathcal G_{\tau,\delta}\cap\mathcal G_{\rm rig}.
\end{equation}

The five restricted record sets
\(
\mathscr R_{X}^{\rm act},
\mathscr R_{A}^{\rm act},
\mathscr R_{Q,\mathrm{ex}}^{\rm act},
\mathscr R_{Q,\mathrm{rev}}^{\rm act},
\mathscr R_{Q,\mathrm{nr}}^{\rm act}
\)
are obtained by requiring the master slot to belong to the corresponding
certificate family of
\cref{eq-master-optimal-active-five-mode-decomposition}.

\end{definition}

\begin{theorem}[Complete LPN master active-sampling envelope]
\label{thm-lpn-complete-master-active-sampling-envelope}

Fix \(0<\eta<\tau<1/2\), \(\delta>0\), and the ceiling of
\cref{def-lpn-master-active-sampling-ceiling}.  On the LPN good-presentation
event used in
\cref{thm-lpn-complete-sparse-reward-observation-accounting}, every complete
active record satisfies
\begin{equation}
\label{eq-lpn-master-active-recordwise-bound}
A_R^H
\le
U_R^{\rm LPN,act}
:=
\min\{U_R^{\rm act,info},U_R^{\rm act,cap},
       U_R^{\rm act,BW},U_R^{\rm act,master},
       U_R^{\rm act,src}\},
\end{equation}
where every candidate is evaluated for the same record and probability law.
For the information coordinate one may take
\begin{align}
\label{eq-lpn-master-active-information-budget}
J_{R,t}^{\rm LPN,act}
&=
\min\{n,m(1-H_2(\eta))\}
+\sum_{j=1}^{L_{R,t}}d_{R,t,j},\\
\label{eq-lpn-master-active-contact-candidate}
U_R^{\rm act,info}
&=
\min\!\left\{1,
\sum_{t=0}^{H-1}\left(
2^{-n}+\sqrt{\frac{\log2}{2}J_{R,t}^{\rm LPN,act}}
\right)\right\}.
\end{align}
The conditional within-cell candidate of
\cref{eq-lpn-curiosity-contact-bound} may replace
\cref{eq-lpn-master-active-contact-candidate} when it is smaller.
Here \(U_R^{\rm act,src}\) is
\cref{eq-active-source-resolved-learning-candidate} with
\(\vartheta_j(\eta)=(1-2\eta)^2\), the LPN five-family contextual charges,
and the complete cost in
\cref{def-lpn-master-active-sampling-ceiling}.

The four nonambient active-geometry modes have zero LPN prospective source
coordinate at the polynomial ceiling:
\begin{equation}
\label{eq-lpn-master-active-nonambient-vanishing}
\sup_{R\in
\mathscr R_{A}^{\rm act}\cup
\mathscr R_{Q,\mathrm{ex}}^{\rm act}\cup
\mathscr R_{Q,\mathrm{rev}}^{\rm act}\cup
\mathscr R_{Q,\mathrm{nr}}^{\rm act}}
V_{\mathscr C_R}
=0.
\end{equation}
Retain every additive comparison error rather than absorbing it into the
vanishing source coordinate:
\begin{equation}
\label{eq-lpn-master-active-nonambient-error}
\Delta_{I,H}^{\rm act,Q}
=
\sup_{R\in
\mathscr R_{A}^{\rm act}\cup
\mathscr R_{Q,\mathrm{ex}}^{\rm act}\cup
\mathscr R_{Q,\mathrm{rev}}^{\rm act}\cup
\mathscr R_{Q,\mathrm{nr}}^{\rm act}}
\min\{U_R^{\rm act,info},U_R^{\rm act,cap},
       U_R^{\rm act,BW},\min(1,\delta_R)\}.
\end{equation}
For the exact same-normalization source comparisons, \(\delta_R=0\) and
this coordinate is zero.
Consequently the optimal qualified active source profile reduces, with every
comparison error retained explicitly, to the identity-support master mode:
\begin{align}
\label{eq-lpn-master-active-reduction-to-identity-geom}
\operatorname{ActArr}_{I,H}^{\rm sat}(B)
&\le
\max\!\left\{
\sup_{R\in\mathscr R_X^{\rm act}(B_{\rm LPN}^{\rm act})}
U_R^{\rm LPN,act},
\Delta_{I,H}^{\rm act,Q}
\right\},
\qquad I\in\mathcal G_{\rm act},\\
\label{eq-lpn-master-active-family-average}
\mathbb E_I\operatorname{ActArr}_{I,H}^{\rm sat}(B)
&\le
\mathbb E_I\!\left[
\mathbf1_{\mathcal G_{\rm act}}
\max\!\left\{
\sup_{R\in\mathscr R_X^{\rm act}(B_{\rm LPN}^{\rm act})}
U_R^{\rm LPN,act},
\Delta_{I,H}^{\rm act,Q}
\right\}
\right]
+\varepsilon_{\rm good}(n,m,\eta,\tau,\delta),
\end{align}
where
\begin{equation}
\label{eq-lpn-master-active-good-event-error}
\varepsilon_{\rm good}
=
\varepsilon_{\rm rig}(n,m)
+e^{-mD_{\rm Ber}(\tau\Vert\eta)}
+(2^n-1)2^{-m(1-H_2(\tau))}
+2^{n+1}e^{-2\delta^2m}.
\end{equation}
For \(R\in\mathscr R_X^{\rm act}\), the master certificate has
\(q=\mathrm{id}\), \(S_q(\lambda)=1\), and
\begin{equation}
\label{eq-lpn-master-active-identity-visibility}
V_{\mathscr C_R}
=
M_{\mathfrak G_R}C_{\mathfrak G_R}
\frac{R_{\mathfrak G_R}}{1+\rho_{\mathfrak G_R}}.
\end{equation}
Accordingly, an inherited comparison retains its declared
family-functional normalization:
\begin{equation}
\label{eq-lpn-master-active-identity-comparison}
\mathfrak M_{n,m,\eta}(U_R^{\rm act,master})
\le
\min\{1,K_R(2^{-n}+(1-\eta)^n)
        +\mathfrak M_{n,m,\eta}(\delta_R)\}.
\end{equation}
A direct comparison instead retains its pointwise normalization:
\begin{equation}
\label{eq-lpn-master-active-direct-identity-comparison}
U_{R,I}^{\rm act,master}
\le
\min\!\left\{1,
K_R^{\rm dir}
M_{\mathfrak G_R}C_{\mathfrak G_R}
\dfrac{R_{\mathfrak G_R}}{1+\rho_{\mathfrak G_R}}
+\delta_{R,I}^{\rm dir}\right\}.
\end{equation}
Thus the remaining active coordinate is evaluated only through the
completely charged target alignment, reach, mixing, and regularity of the
selected full-carrier geometry; no quotient or observation-cell credit is
retained in this branch.

\end{theorem}

\begin{proof}

Apply \cref{thm-master-optimal-structural-active-sampling-envelope} with
\(\Theta=S\) and \(\mathsf Z_0=(A,b)\).  The initial LPN information bound
and ordered-reveal increment are
\cref{thm-lpn-prehit-information-contact-envelope}; summing its disjoint
next-contact bounds gives
\cref{eq-lpn-master-active-information-budget,eq-lpn-master-active-contact-candidate}.
The capacity, Blackwell, and master candidates are the same-record
substitutions already used in
\cref{cor-lpn-saturated-quantitative-prospective-certificate-envelope}.
The source-resolved candidate is
\cref{thm-source-resolved-saturated-learning-compilation} specialized to
the complete LPN active record.  It retains the separate Add, Mult, Ord,
Sym, and Filt charges and the binary-noise contraction.

The pure exact \(\Abs\), exact quotient-geometry, and both compensated
coordinates vanish by
\cref{eq-lpn-exact-abs-source-vanishes,eq-lpn-exact-quotient-geom-closed,eq-lpn-compensated-source-vanishes}.
This proves \cref{eq-lpn-master-active-nonambient-vanishing}.  The exhaustive
five-mode partition bounds every nonambient record by its retained
comparison error and leaves \(\mathcal C_X\) as the only nonzero source
mode.  This proves the pointwise bound
\cref{eq-lpn-master-active-reduction-to-identity-geom}.  The value of an
identity-support active master certificate is
\cref{eq-supported-geom-visibility} with
\(\overline A_{\rm id}=S_{\rm id}=1\), which gives
\cref{eq-lpn-master-active-identity-visibility}.  The inherited ancestor
bound is \cref{eq-lpn-complete-ambient-geom-source-bound}; substituting it
and the direct visibility in the corresponding comparison forms proves
\cref{eq-lpn-master-active-identity-comparison,eq-lpn-master-active-direct-identity-comparison}.
Finally apply the unit
bound on the complement of the good event; its probability is
\cref{eq-lpn-curiosity-good-event-complement}.  The tower property proves
\cref{eq-lpn-master-active-family-average}.

\end{proof}

\begin{theorem}[Fixed, passive-stream, and chosen-query LPN acquisition]
\label{thm-lpn-three-active-acquisition-interfaces}

The three active interfaces have the following quantitative accounting.

\begin{enumerate}[label=\textup{(\roman*)}]
\item In the fixed standard presentation, every selected row, coordinate,
      residual feature, or deterministic postprocessing \(F(A,b)\) satisfies
      \begin{equation}
      \label{eq-lpn-fixed-record-active-zero-information}
      I(S;F(A,b)\mid A,b,\mathcal H_t)=0.
      \end{equation}
      Its construction, score optimization, memory update, and use in the
      selected transition remain charged by
      \cref{eq-complete-master-active-sampling-cost}.  A later verifier
      answer occupies an ordered reveal slot.

\item Suppose a declared passive-stream interface returns
      \(Z_j=(a_j,y_j)\), where \(a_j\) is fresh uniform and independent of
      the past and
      \(y_j=\langle a_j,S\rangle+E_j\),
      \(E_j\sim\operatorname{Bernoulli}(\eta)\).  After \(q\) requested
      samples,
      \begin{equation}
      \label{eq-lpn-passive-stream-information-bound}
      I(S;Z_{1:q}\mid A,b)
      \le
      \min\{n,q(1-H_2(\eta))\}.
      \end{equation}

\item Suppose instead that a declared chosen-query interface accepts
      \(a_j\in\mathbb F_2^n\) selected from the pre-query history and returns
      \(\langle a_j,S\rangle+E_j\).  This is a presentation extension.
      Querying \(e_i\) independently \(q_i\) times and taking the majority
      value gives
      \begin{equation}
      \label{eq-lpn-chosen-query-majority-recovery}
      \Pr\{\widehat S\ne S\}
      \le
      \sum_{i=1}^{n}
      \exp\!\left[-\frac{q_i(1-2\eta)^2}{2}\right].
      \end{equation}
      Thus \(q_i\ge
      2(1-2\eta)^{-2}\log(n/\epsilon)\) makes the total error at most
      \(\epsilon\).
\end{enumerate}

\end{theorem}

\begin{proof}

Part~\textup{(i)} is conditional data processing: \(F(A,b)\) is measurable
with respect to the conditioning record.  For one passive sample,
\(a_j\) is independent of \(S\), while
\[
I(S;y_j\mid a_j,\mathcal H_j)
\le H(y_j\mid a_j,\mathcal H_j)-H(E_j)
\le1-H_2(\eta).
\]
The conditional chain rule proves
\cref{eq-lpn-passive-stream-information-bound}.  The chosen-query interface
is separated from the fixed presentation by
\cref{thm-oracle-admissibility-presentation-separation}.  For each basis
query, Hoeffding's inequality bounds the majority error by
\(\exp[-2q_i(1/2-\eta)^2]\).  A union bound gives
\cref{eq-lpn-chosen-query-majority-recovery} and the displayed sufficient
query count.

\end{proof}

\begin{corollary}[Parameter-explicit LPN active-acquisition regimes]
\label{cor-lpn-parameter-explicit-active-acquisition-regimes}

Fix a monotone scalarization \(\kappa\), write \(b=\kappa(B)\), and let
\(c_{\rm samp}>0\) be the charged scalar cost of one declared fresh-sample
request, including its returned row, label, storage, and channel update.
The number of fresh-sample requests in any such record is at most
\begin{equation}
\label{eq-lpn-budgeted-active-sample-count}
q_B
=
\min\!\left\{H,\left\lfloor\frac{b}{c_{\rm samp}}\right\rfloor\right\}.
\end{equation}
For the passive-stream extension, let
\[
D_{R,t}^{\rm post}
=\sum_{j=1}^{L_{R,t}}d_{R,t,j}
\]
contain the contracted information of every other target-dependent reveal.
Then
\begin{align}
\label{eq-lpn-active-total-information-n-m-eta}
I(S;\mathcal H_t)
&\le
J_{n,m,\eta,q_B,t}^{\rm act}
:=
\min\!\left\{
n,\,(m+q_B)(1-H_2(\eta))+D_{R,t}^{\rm post}
\right\},\\
\label{eq-lpn-active-contact-n-m-eta}
\Pr\{\tau_s\le H\}
&\le
\min\!\left\{1,\,
H2^{-n}
+\sum_{t=0}^{H-1}
\sqrt{\frac{\log2}{2}
J_{n,m,\eta,q_B,t}^{\rm act}}
\right\},\\
\label{eq-lpn-active-fano-n-m-eta}
\Pr\{\widehat S=S\}
&\le
\min\!\left\{1,\,
\frac{J_{n,m,\eta,q_B,H}^{\rm act}+1}{n}
\right\}.
\end{align}
The fixed standard presentation is the specialization \(q_B=0\).

For a chosen-query extension with equal allocation, put
\(
\underline q_B=\lfloor q_B/n\rfloor
\).
The basis-majority construction satisfies
\begin{equation}
\label{eq-lpn-chosen-query-budget-noise-bound}
\Pr\{\widehat S\ne S\}
\le
n\exp\!\left[
-\frac{\underline q_B(1-2\eta)^2}{2}
\right].
\end{equation}
In particular, error at most \(\epsilon\) is obtained whenever
\begin{equation}
\label{eq-lpn-chosen-query-budget-threshold}
q_B
\ge
\frac{2n}{(1-2\eta)^2}\log\frac n\epsilon+n.
\end{equation}

For the fixed-presentation identity-geometry mode, every same-record
uniform-invariant certificate supplies
\begin{equation}
\label{eq-lpn-active-uniform-capacity-n}
U_R^{\rm act,cap}
\le
\min\{1,R_0(H+1)2^{-n}\}.
\end{equation}
For a residual Gibbs--Metropolis certificate, put
\(
\Delta_\pm=1/2-\eta\pm2\delta
\).
If
\(
\theta m\Delta_+\le(1-\epsilon_0)n\log2
\)
and \(\rho_0\le R_\theta\pi_\theta\), then
\begin{equation}
\label{eq-lpn-active-residual-capacity-n-m-eta}
U_R^{\rm act,cap}
\le
\min\!\left\{
1,\,
\frac{R_\theta2^{-\epsilon_0 n}}{1-2^{-n}}
\left(1+\frac H2e^{-\theta m\Delta_-}\right)
\right\}.
\end{equation}
Thus every displayed active certificate depends explicitly on
\((n,m,\eta,H,B)\), with \((\tau,\delta)\) entering through the
good-presentation and residual-flatness conversions.

\end{corollary}

\begin{proof}

Equation \cref{eq-lpn-budgeted-active-sample-count} is the joint cost and
horizon restriction.  Combine the public-record information bound
\(
I(S;A,b)\le\min\{n,m(1-H_2(\eta))\}
\)
with \cref{eq-lpn-passive-stream-information-bound} and the conditional
chain rule.  This proves
\cref{eq-lpn-active-total-information-n-m-eta}.  The overlap-one contact
and Fano conversions are
\cref{thm-lpn-prehit-information-contact-envelope}, proving
\cref{eq-lpn-active-contact-n-m-eta,eq-lpn-active-fano-n-m-eta}.

Apply \cref{eq-lpn-chosen-query-majority-recovery} with
\(q_i=\underline q_B\) to obtain
\cref{eq-lpn-chosen-query-budget-noise-bound}; the displayed threshold
implies \(\underline q_B\ge
2(1-2\eta)^{-2}\log(n/\epsilon)\).
The two capacity bounds are respectively
\cref{eq-lpn-controlled-capacity-hit-bound,eq-lpn-residual-tilt-subcritical-hit}.

\end{proof}

\begin{theorem}[Optimal active construction for the standard LPN presentation]
\label{thm-lpn-optimal-active-construction-equivalence}

Fix \((n,m,\eta,H,B)\) in the standard LPN presentation and use the family
functional \(\mathfrak M_{n,m,\eta}\).  Let
\begin{equation}
\label{eq-lpn-optimal-active-arrival-profile}
\operatorname{ActArr}_{{\rm LPN},n,m,\eta,H}^{\rm std,sat}(B)
=
\sup_{R\in\mathscr R_{{\rm LPN},H}^{\rm act}(B)}
\mathfrak M_{n,m,\eta}(A_R^H),
\end{equation}
where the query component is null or a deterministic readout of the fixed
public record \((A,b)\).  Then
\begin{equation}
\label{eq-lpn-active-optimum-saturated-identity}
\operatorname{ActArr}_{{\rm LPN},n,m,\eta,H}^{\rm std,sat}(B)
=
\operatorname{Succ}_{{\rm LPN},n,m,\eta,H}^{\rm act,sat}(B),
\end{equation}
and its exact arrival decomposition is the five-mode maximum in
\cref{eq-active-optimality-five-mode-arrival-decomposition}.

Let \(\operatorname{Succ}_{{\rm LPN},n,m,\eta,H}^{\rm sat}(B)\) denote the
saturated success profile over all represented standard-presentation LPN
computations stopped by time \(H\).  There is a class-preserving complete
majorant \(\Psi_{\rm LPN}^{\rm act}\), obtained from
\cref{def-lpn-master-active-sampling-ceiling}, such that
\begin{equation}
\label{eq-lpn-active-computation-profile-sandwich}
\operatorname{ActArr}_{{\rm LPN},n,m,\eta,H}^{\rm std,sat}(B)
\le
\operatorname{Succ}_{{\rm LPN},n,m,\eta,H}^{\rm sat}(B)
\le
\operatorname{ActArr}_{{\rm LPN},n,m,\eta,H}^{\rm std,sat}
\bigl(\Psi_{\rm LPN}^{\rm act}(B,H,n,m,\eta)\bigr).
\end{equation}
Thus the standard-presentation active optimum and the complete finite-horizon
saturated success profile are the same affordable class up to the explicit
cost of compiling the clocked action, master certificate, and comparison
record.  The compilation introduces no fresh LPN sample or chosen parity
query.

After scalarization, let \(\psi_{\rm LPN}^{\rm act}\) majorize this compiled
cost, and let
\(B_{{\rm act},{\rm LPN}}^*(n,m,\eta,H;\alpha)\) denote
\cref{eq-active-optimal-computational-cost} restricted to the standard LPN
active record class.  For every \(\varepsilon>0\), the corresponding
threshold costs obey
\begin{equation}
\label{eq-lpn-active-recovery-cost-equivalence}
B_{\rm rec}^*(n,m,\eta;\alpha)
\le
B_{{\rm act},{\rm LPN}}^*(n,m,\eta,H;\alpha)
\le
\psi_{\rm LPN}^{\rm act}
\bigl(B_{\rm rec}^*(n,m,\eta;\alpha)+\varepsilon,
      H,n,m,\eta\bigr),
\end{equation}
whenever the recovery infimum is witnessed within horizon \(H\) at the
displayed \(\varepsilon\)-enlarged cost.  Hence polynomial affordability is
preserved in both directions by a polynomial class-preserving compilation.

\end{theorem}

\begin{proof}

The identity
\cref{eq-lpn-active-optimum-saturated-identity} is
\cref{thm-active-optimality-saturated-construction} on the standard LPN
slice.  Equation
\cref{eq-lpn-fixed-record-active-zero-information} shows that deterministic
selection inside \((A,b)\) creates no additional observation channel; its
computation and use remain in the complete cost.

Dropping certificate annotations from a complete active record leaves its
represented query, transition, stopping, and output computation, proving
the first inequality in
\cref{eq-lpn-active-computation-profile-sandwich}.  Conversely,
\cref{cor-all-turing-machines-represented,thm-algorithm-kernel-embedding}
turn every finite-horizon represented computation into a null-query clocked
active record.  Encoding adequacy, the universal dynamical certificate, and
\cref{def-lpn-master-active-sampling-ceiling} append its completely charged
master and comparison slots at
\(\Psi_{\rm LPN}^{\rm act}\), without changing the executed path law.  This
proves the second inequality.  Apply the two inequalities to the defining
infima in
\cref{eq-noise-indexed-recovery-cost,eq-active-optimal-computational-cost}
to obtain \cref{eq-lpn-active-recovery-cost-equivalence}.

\end{proof}

\begin{corollary}[Explicit learned-optimum loss for LPN active sampling]
\label{cor-lpn-explicit-learned-active-optimum-loss}

In the setting of
\cref{cor-lpn-parameter-explicit-active-acquisition-regimes}, suppose a
complete standard-presentation LPN record learns its represented
action-value rule through a bounded flip-symmetric score comparison from a
sample-independent represented codebook and a sample certificate with effective size
\(m_{\rm eff}\), dependence error \(\varepsilon_{\rm dep}\), confidence
\(1-\delta\), common optimization error \(\varepsilon_{\rm opt}\), and
affordable action-value codebook logarithm
\(L_{{\rm LPN},{\rm act}}(n,m,\eta,B)\).  Then, on its simultaneous event,
\begin{equation}
\label{eq-lpn-learned-active-optimum-loss-n-m-eta}
\begin{aligned}
V_{{\rm LPN},0}^{\rm act}-A_{\widehat R}^{H}
\le H\Bigg\{&
\frac{2}{1-2\eta}
\left[
2\sqrt{\frac{2L_{{\rm LPN},{\rm act}}(n,m,\eta,B)}{m_{\rm eff}}}
+3\sqrt{\frac{\log(2H/\delta)}{2m_{\rm eff}}}
+\varepsilon_{\rm dep}
\right]
+\varepsilon_{\rm opt}\Bigg\}.
\end{aligned}
\end{equation}
The same record simultaneously retains the information, target-contact,
chosen-interface, and capacity bounds in
\cref{eq-lpn-active-total-information-n-m-eta,eq-lpn-active-contact-n-m-eta,eq-lpn-chosen-query-budget-noise-bound,eq-lpn-active-residual-capacity-n-m-eta}.
For the fixed standard presentation, \(q_B=0\); passive fresh samples and
chosen parity queries retain their separate presentation-extension status.

\end{corollary}

\begin{proof}

Apply \cref{cor-explicit-sample-noise-active-optimality-loss} with the LPN
family functional and the displayed codebook.  The factor
\((1-2\eta)^{-1}\) is the homogeneous label-noise transport of
\cref{thm-uniform-dynamical-label-noise-bound}; the remaining parameter
bounds are the cited LPN active-acquisition estimates.

\end{proof}

\begin{theorem}[LPN occupation-flow and Bellman-support certificate]
\label{thm-lpn-occupation-flow-bellman-support-certificate}

Fix the standard-presentation LPN slice
\((n,m,\eta,H,B)\), let \(\rho=1-2\eta\), and use the completely compiled
ceiling
\[
B_{\rm LPN}^{\rm Bell}
=\Psi_{\rm Bell}
\bigl(\Psi_{\rm LPN}^{\rm act}(B,H,n,m,\eta),H,n\bigr).
\]
For a bounded family-uniform barrier \(v=(v_t)_{t\le H}\), \(v_H=0\),
define
\begin{equation}
\label{eq-lpn-modewise-bellman-support-certificate}
\mathfrak U_{t,c}^{\rm LPN}(v)
=
\sup_{R\in\mathscr R_{{\rm LPN},t;c}^{\rm act}
                 (B_{\rm LPN}^{\rm Bell})}
\mathfrak M_{n,m,\eta}
\left(U_{R,t}^{\rm Bell}(v)\right),
\qquad c\in\mathfrak C_{\rm master},
\end{equation}
where \(U_{R,t}^{\rm Bell}\) is the same-record minimum in
\cref{eq-recordwise-bellman-support-certificate-minimum}.  On the LPN good
event of
\cref{thm-lpn-complete-sparse-reward-observation-accounting}, the standard
active optimum satisfies
\begin{equation}
\label{eq-lpn-bellman-support-family-bound}
\operatorname{ActArr}_{{\rm LPN},n,m,\eta,H}^{\rm std,sat}(B)
\le
\inf_{v_H=0}
\left\{
\mathfrak M_{n,m,\eta}(v_0(h_0))
+\sum_{t=0}^{H-1}
\max_{c\in\mathfrak C_{\rm master}}
\mathfrak U_{t,c}^{\rm LPN}(v)
\right\}
+\varepsilon_{\rm good}(n,m,\eta,\tau,\delta).
\end{equation}
In particular, the zero barrier gives the completely explicit next-contact
specialization
\begin{equation}
\label{eq-lpn-zero-barrier-support-family-bound}
\operatorname{ActArr}_{{\rm LPN},n,m,\eta,H}^{\rm std,sat}(B)
\le
\sum_{t=0}^{H-1}
\max_{c\in\mathfrak C_{\rm master}}
\mathfrak U_{t,c}^{\rm LPN}(0)
+\varepsilon_{\rm good}(n,m,\eta,\tau,\delta).
\end{equation}

The four nonambient modes obey the explicit comparison-error bound
\begin{equation}
\label{eq-lpn-nonambient-bellman-support-bound}
\max_{c\in\{A,Q_{\rm ex},Q_{\rm rev},Q_{\rm nr}\}}
\mathfrak U_{t,c}^{\rm LPN}(0)
\le
\sup_{R\in\mathscr R_A^{\rm act}\cup
                 \mathscr R_{Q,{\rm ex}}^{\rm act}\cup
                 \mathscr R_{Q,{\rm rev}}^{\rm act}\cup
                 \mathscr R_{Q,{\rm nr}}^{\rm act}}
\mathfrak M_{n,m,\eta}
\min\{1,U_{R,t}^{\rm info},U_{R,t}^{\rm cap},
U_{R,t}^{\rm BW},\delta_{R,t}\}.
\end{equation}
For exact same-normalization comparisons \(\delta_{R,t}=0\), so this
quantity vanishes.

For the identity-support mode, every record retains the following explicit
candidate minimum:
\begin{equation}
\label{eq-lpn-ambient-bellman-support-explicit-minimum}
\mathfrak U_{t,X}^{\rm LPN}(0)
\le
\sup_{R\in\mathscr R_X^{\rm act}(B_{\rm LPN}^{\rm Bell})}
\mathfrak M_{n,m,\eta}
\min\left\{
\begin{array}{l}
1,\\
\displaystyle
2^{-n}+\sqrt{\frac{\log2}{2}
J_{R,t}^{\rm LPN,act}},\\[2mm]
\displaystyle R_0(H+1)2^{-n},\\[1mm]
\displaystyle
\frac{R_\theta2^{-\epsilon_0n}}{1-2^{-n}}
\left(1+\frac H2e^{-\theta m(1/2-\eta-2\delta)}\right),\\[2mm]
U_{R,t}^{\rm BW},\\[1mm]
\displaystyle
\min\!\left\{1,
K_R\bigl(2^{-n}+(1-\eta)^n\bigr)+\delta_{R,t}\right\},\\[2mm]
\displaystyle
\min\!\left\{1,
K_R^{\rm dir}M_{\mathfrak G_R}C_{\mathfrak G_R}
\frac{R_{\mathfrak G_R}}{1+\rho_{\mathfrak G_R}}
+\delta_{R,t}^{\rm dir}\right\}
\end{array}
\right\}.
\end{equation}
Here every entry is evaluated conditionally from the same surviving history.
For \(v=0\), the positive Bellman increment is exactly the next-contact
probability \(p_T(h,a)\).  A whole-horizon candidate enters this minimum only
through its represented re-rooted comparison at that history.  For a general
barrier \(v\), the corresponding entry in
\cref{eq-lpn-modewise-bellman-support-certificate} is finite only when the
record supplies the Bellman-increment converter required by
\cref{thm-structural-evaluation-affordable-bellman-support}.
The residual-Gibbs candidate is included when
\(\theta m(1/2-\eta+2\delta)\le(1-\epsilon_0)n\log2\) and the represented
initial-density comparison of
\cref{eq-lpn-active-residual-capacity-n-m-eta} holds.  Otherwise that entry
has the neutral value one.  The information budget is
\begin{equation}
\label{eq-lpn-bellman-support-information-budget}
J_{R,t}^{\rm LPN,act}
\le
\min\{n,m(1-H_2(\eta))\}
+\sum_{j=1}^{L_{R,t}}d_{R,t,j},
\end{equation}
and every ordered reveal, contraction coefficient, barrier evaluation,
transition, and comparison is included in
\(B_{\rm LPN}^{\rm Bell}\).

If a complete LPN learner produces \(\widehat R\) with the action-value
radius in
\cref{eq-lpn-learned-active-optimum-loss-n-m-eta}, then, on the same
simultaneous event,
\begin{align}
\label{eq-lpn-complete-learning-control-upper}
A_{\widehat R}^{H}
&\le
\operatorname{ActArr}_{{\rm LPN},n,m,\eta,H}^{\rm std,sat}(B),\\
\label{eq-lpn-complete-learning-control-lower}
A_{\widehat R}^{H}
&\ge
\operatorname{ActArr}_{{\rm LPN},n,m,\eta,H}^{\rm std,sat}(B)
-H\left(2\Gamma_{{\rm LPN},{\rm act}}^{\rm clean}
          +\varepsilon_{\rm opt}\right),
\end{align}
where
\begin{equation}
\label{eq-lpn-complete-learning-control-radius}
\Gamma_{{\rm LPN},{\rm act}}^{\rm clean}
\le
\frac1\rho\left[
2\sqrt{\frac{2L_{{\rm LPN},{\rm act}}(n,m,\eta,B)}{m_{\rm eff}}}
+3\sqrt{\frac{\log(2H/\delta)}{2m_{\rm eff}}}
+\varepsilon_{\rm dep}
\right].
\end{equation}
Thus \cref{eq-lpn-bellman-support-family-bound} controls the best affordable
LPN controller, while
\cref{eq-lpn-complete-learning-control-lower,eq-lpn-complete-learning-control-radius}
control the statistical and optimization distance of a learned controller
from that optimum.

\end{theorem}

\begin{proof}

Compile every standard-presentation computation into the complete active
record class by
\cref{thm-lpn-optimal-active-construction-equivalence}.  Apply
\cref{thm-bellman-barrier-duality-affordable-arrival} and then the structural
support evaluation
\cref{thm-structural-evaluation-affordable-bellman-support}.  The five-mode
decomposition gives the maximum in
\cref{eq-lpn-bellman-support-family-bound}.  On the complement of the good
event use the unit bound; its probability is
\cref{eq-lpn-master-active-good-event-error}.

The exact and compensated nonambient LPN sources vanish by
\cref{eq-lpn-master-active-nonambient-vanishing}.  Their numerical Bellman
increments therefore retain only the same-event candidates and comparison
errors, proving
\cref{eq-lpn-nonambient-bellman-support-bound}.  In the identity mode,
specialize to \(v=0\), so the positive Bellman increment is the conditional
next-contact probability.  Then
insert successively the information/contact bound
\cref{eq-lpn-master-active-contact-candidate}, the uniform-capacity and
residual-Gibbs bounds
\cref{eq-lpn-active-uniform-capacity-n,eq-lpn-active-residual-capacity-n-m-eta},
the inherited comparison
\cref{eq-lpn-master-active-identity-comparison}, and the direct master
visibility
\cref{eq-lpn-master-active-direct-identity-comparison}.  These candidates
bound that same conditional next-contact probability only through their
stored re-rooted same-normalization converters; their minimum proves
\cref{eq-lpn-ambient-bellman-support-explicit-minimum}.  The information
display is
\cref{eq-lpn-master-active-information-budget}.  The zero-barrier case of
\cref{eq-lpn-bellman-support-family-bound} is
\cref{eq-lpn-zero-barrier-support-family-bound}.

Finally, apply
\cref{cor-complete-learned-control-sandwich,cor-explicit-sample-noise-active-optimality-loss}
and substitute the LPN finite-description radius.  This proves
\cref{eq-lpn-complete-learning-control-upper,eq-lpn-complete-learning-control-lower,eq-lpn-complete-learning-control-radius}.

\end{proof}

\begin{corollary}[LPN specialization of prospective contact-witness
coercivity]
\label{cor-lpn-prospective-contact-witness-coercivity}

Fix the standard LPN presentation, parameters
\((n,m,\eta,H,B)\), and the compiled ceiling
\(B_{\rm LPN}^{\rm Bell}\) of
\cref{thm-lpn-occupation-flow-bellman-support-certificate}.  Let
\(R\) be one complete family-uniform active record and let \(\nu\) be the
declared target-neutral next-test baseline.  If
\begin{equation}
\label{eq-lpn-family-excess-arrival-hypothesis}
\mathfrak M_{n,m,\eta}(A_R^H)
>
\mathfrak M_{n,m,\eta}
  \bigl(\operatorname{Enum}_{\nu}(R,H)\bigr)
+\varepsilon,
\end{equation}
then some \(t_*<H\) satisfies
\begin{equation}
\label{eq-lpn-family-selector-witness-scale}
\mathfrak M_{n,m,\eta}
\bigl(\operatorname{Adv}_{t_*}(K_{R,t_*};\nu_{t_*})\bigr)
>
\frac{\varepsilon}{H}.
\end{equation}
The complete normalized pre-hit derivation of this selector reads \(b\) or
the residual primitive \(r_I(x)=Ax+b\).  Thus its transition record has
residual ancestry.  After the quotient-supported modes are assigned, an
ambient witness belongs to the residual-ancestry branch of
\cref{def-lpn-ambient-geom-ancestry-branches}; a residual-independent
derivation cannot satisfy
\cref{eq-lpn-family-selector-witness-scale}.

At the same time and compiled ceiling, the zero-barrier Bellman certificate
satisfies
\begin{equation}
\label{eq-lpn-contact-witness-certificate-scale}
\max_{c\in\mathfrak C_{\rm master}}
\mathfrak U_{t_*,c}^{\rm LPN}(0)
\ge\frac{\varepsilon}
{H\lvert\mathfrak C_{\rm master}\rvert}
=\frac{\varepsilon}{5H}.
\end{equation}
The exact and compensated quotient-supported modes are evaluated by
\cref{eq-lpn-nonambient-bellman-support-bound}; when their represented
same-normalization comparisons are exact, their contribution is zero.
Consequently a positive witness at the displayed scale is either

\begin{enumerate}[label=\textup{(\roman*)}]
\item a residual-ancestry ambient record evaluated by one of the finite
      candidates in
      \cref{eq-lpn-ambient-bellman-support-explicit-minimum}; or
\item a residual-ancestry ambient record for which that same-record
      next-contact comparison is unavailable and therefore has the neutral
      certificate value one.
\end{enumerate}

In case \textup{(i)}, every available comparison is already explicit in
\((n,m,\eta,H,B)\): its information term is bounded by
\cref{eq-lpn-bellman-support-information-budget}, its uniform-capacity term
is \(R_0(H+1)2^{-n}\), its represented residual-Gibbs term is the third
nontrivial entry of
\cref{eq-lpn-ambient-bellman-support-explicit-minimum}, and its inherited or
direct master comparison retains its complete source and converter cost.
Case \textup{(ii)} is not a new structural mechanism.  It is exactly the
absence of a represented comparison from the already charged ambient record
to its next-contact value.  Provenance classification, posterior
concentration, or a retrospective first-hit record does not fill that slot.

Conversely, if every residual-ancestry ambient first-step record at the
compiled ceiling has a represented same-normalization comparison bounded by
\(\delta_{t}^{b}(n,m,\eta,H,B)\), then
\begin{equation}
\label{eq-lpn-residual-ancestry-contact-arrival-converse}
\mathfrak M_{n,m,\eta}(A_R^H)
\le
\mathfrak M_{n,m,\eta}
  \bigl(\operatorname{Enum}_{\nu}(R,H)\bigr)
+\sum_{t=0}^{H-1}\delta_t^b(n,m,\eta,H,B)
+\varepsilon_{\rm good}(n,m,\eta,\tau,\delta).
\end{equation}
In particular, a bound
\(\delta_t^b\le(1-\eta)^n\) gives the explicit excess-arrival estimate
\begin{equation}
\label{eq-lpn-explicit-contact-witness-arrival-bound}
\mathfrak M_{n,m,\eta}(A_R^H)
\le
H2^{-n}+H(1-\eta)^n
+\varepsilon_{\rm good}(n,m,\eta,\tau,\delta)
\end{equation}
for the uniform full-carrier baseline.

\end{corollary}

\begin{proof}

Apply the family functional to the exact hazard identity
\cref{eq-selector-hazard-decomposition}.  If
\cref{eq-lpn-family-excess-arrival-hypothesis} holds, the sum of the \(H\)
family-averaged selector advantages exceeds \(\varepsilon\), so one term
satisfies
\cref{eq-lpn-family-selector-witness-scale}.  If the complete normalized
source derivation of that selector did not read \(b\) or \(r_I\), it would be
residual-independent.  Its expected advantage over the target-neutral
baseline would then be zero by
\cref{lem-lpn-syndrome-free-ambient-geom-bound}, contradicting
\cref{eq-lpn-family-selector-witness-scale}.  This proves the ancestry
claim without treating information in the public residual as an affordable
readout.

Apply
\cref{thm-quantitative-prospective-contact-witness-coercivity} pointwise and
then the monotone LPN family functional.  The pointwise master maximum is at
most the sum of its five mode coordinates.  Hence one fixed master mode
carries at least one fifth of the family-averaged witness, which gives the
factor \(5=\lvert\mathfrak C_{\rm master}\rvert\) in
\cref{eq-lpn-contact-witness-certificate-scale}.  Structural Bellman
evaluation at the compiled ceiling is
\cref{thm-lpn-occupation-flow-bellman-support-certificate}; it yields
\cref{eq-lpn-contact-witness-certificate-scale}.  The nonambient evaluation
is
\cref{eq-lpn-nonambient-bellman-support-bound}.  The remaining master mode is
ambient, and the residual-independent alternative has just been excluded at
positive family advantage.  The same-record minimum
\cref{eq-lpn-ambient-bellman-support-explicit-minimum} then gives exactly the
two alternatives in the statement: a finite evaluated converter or the
neutral value assigned to an unavailable converter.

For the converse, use the residual-independent zero-advantage identity and
the exact nonambient comparisons, substitute the assumed ambient
same-normalization bounds into
\cref{eq-modewise-contact-support-arrival-bound}, and account for the
complement of the LPN good event by the unit bound.  This proves
\cref{eq-lpn-residual-ancestry-contact-arrival-converse}.  The uniform
full-carrier neutral baseline is at most \(H2^{-n}\) by
\cref{eq-lpn-valid-lift-neutral-baseline}; substituting
\(\delta_t^b\le(1-\eta)^n\) proves
\cref{eq-lpn-explicit-contact-witness-arrival-bound}.

\end{proof}

\begin{figure}[tbp]
\centering
\resizebox{.98\linewidth}{!}{%
\begin{tikzpicture}[fw diagram,node distance=7mm and 10mm]
  \node[fw source,text width=2.6cm] (params)
    {LPN parameters\\\(n,m,\eta,H,B\)};
  \node[fw provenance,text width=2.8cm,right=of params] (flow)
    {complete standard\\occupation flows};
  \node[fw geom,text width=2.8cm,right=of flow] (support)
    {five-mode Bellman\\support maximum};
  \node[fw box,text width=3.2cm,right=of support] (numbers)
    {information, capacity,\\source and noise bounds};
  \node[fw terminal,text width=2.8cm,right=of numbers] (opt)
    {best affordable\\LPN arrival};
  \node[fw use,text width=2.8cm,below=9mm of support] (learn)
    {learned-rule radius\\and optimizer loss};
  \draw[fw arrow] (params)--(flow);
  \draw[fw arrow] (flow)--(support);
  \draw[fw arrow] (support)--(numbers);
  \draw[fw arrow] (numbers)--(opt);
  \draw[fw arrow] (learn)--(support);
\end{tikzpicture}%
}
\caption{Quantitative LPN learning/control certificate.  The Bellman-support
maximum evaluates the best affordable standard-presentation controller.
The learned-rule radius, including the factor \((1-2\eta)^{-1}\), measures
the additional statistical and optimization loss.}
\label{fig-lpn-learning-control-bellman-support}
\end{figure}

\subsection{Structural Bayesian learning for LPN}
\label{subsec-lpn-structural-bayesian-learning}

The declared LPN target prior is uniform.  The public matrix is independent
of the secret, whereas the noisy right-hand side is the first
target-dependent observation.  This reveal order fixes the prior provenance
before any represented belief or acquisition rule is constructed.

\begin{definition}[LPN structural belief and posterior-concentration profiles]
\label{def-lpn-structural-bayesian-profile}

Let \(S\sim\operatorname{Unif}(\mathbb F_2^n)\), so
\begin{equation}
\label{eq-lpn-declared-uniform-task-prior}
\Pi_0^{\rm LPN}(u)=2^{-n},
\qquad u\in\mathbb F_2^n.
\end{equation}
For \(0<\eta<1/2\), conditional on a presented \((A,b)\), the analytic
task posterior is
\begin{equation}
\label{eq-lpn-analytic-task-posterior}
\Pi_{A,b}^{\rm LPN}(u)
=
\frac{
\eta^{|Au+b|_1}(1-\eta)^{m-|Au+b|_1}}
{\displaystyle
 \sum_{v\in\mathbb F_2^n}
 \eta^{|Av+b|_1}(1-\eta)^{m-|Av+b|_1}}.
\end{equation}
This displayed law is an analytic comparator.  A prospective LPN belief
\(\widehat\Pi_{R,t}\) is available only through a complete structural
Bayesian record at the compiled ceiling
\begin{equation}
\label{eq-lpn-structural-bayes-compiled-ceiling}
B_{\rm LPN}^{\rm Bayes}
=\Psi_{\rm Bayes}(B_{\rm LPN}^{\rm Bell},H,n,m,\eta).
\end{equation}
Its target concentration and filter error are
\begin{align}
\label{eq-lpn-recordwise-posterior-concentration}
C_{R,t}^{\rm LPN}
&=\left[\max_{u\in\mathbb F_2^n}
\widehat\Pi_{R,t}(u)-2^{-n}\right]_+,\\
\label{eq-lpn-recordwise-posterior-filter-error}
e_{R,t}^{\rm LPN}
&=\left\lVert
\widehat\Pi_{R,t}-
\mathcal L(S\mid\mathcal F_{R,t})
\right\rVert_{\rm TV}.
\end{align}
Let \(\mathfrak M^*=\mathfrak M^{\mathcal G_{\rm act}}_{n,m,\eta}\)
be the declared family functional restricted to the LPN good event in
\cref{eq-lpn-master-active-good-event}.  Put
\begin{align}
\label{eq-lpn-saturated-posterior-concentration-profile}
\operatorname{PostConc}_{\mathfrak M^*,n,m,\eta,H}^{\rm LPN,sat}
(B_{\rm LPN}^{\rm Bayes})
&=\max_{t<H}\sup_R
\mathfrak M^*_{n,m,\eta}
\mathbb E_R\!\left[
\mathbf1_{\{\tau>t\}}C_{R,t}^{\rm LPN}
\right],\\
\label{eq-lpn-saturated-posterior-filter-profile}
\operatorname{FilErr}_{\mathfrak M^*,n,m,\eta,H}^{\rm LPN,sat}
(B_{\rm LPN}^{\rm Bayes})
&=\max_{t<H}\sup_R
\mathfrak M^*_{n,m,\eta}
\mathbb E_R\!\left[
\mathbf1_{\{\tau>t\}}e_{R,t}^{\rm LPN}
\right],
\end{align}
where \(R\) ranges over one family-uniform complete structural Bayesian
record scheme at the displayed ceiling.  Every likelihood score, feature
locator, evidence normalization, retained coefficient, precision choice,
belief update, acquisition score, optimizer, and deployed transition is
included in that record.
Write
\begin{equation}
\label{eq-lpn-saturated-prospective-posterior-source-profile}
\operatorname{PostSrc}_{\mathfrak M^*,n,m,\eta,H}^{\rm LPN,sat}
(B_{\rm LPN}^{\rm Bayes})
=\max_{t<H}\sup_{R\in
\mathscr R_{{\rm LPN},H}^{\rm Bayes,src}(B_{\rm LPN}^{\rm Bayes})}
\mathfrak M^*\mathbb E_R\!\left[
\mathbf1_{\{\tau>t\}}C_{R,t}^{\rm LPN}
\right],
\end{equation}
where the record restriction is the LPN specialization of
\cref{eq-saturated-prospective-posterior-source-profile}.

\end{definition}

\begin{theorem}[Quantitative structural Bayesian envelope for LPN]
\label{thm-lpn-quantitative-structural-bayesian-envelope}

Use the parameter range and good-presentation event of
\cref{thm-lpn-complete-master-active-sampling-envelope}.  Every complete
LPN structural Bayesian record satisfies
\begin{equation}
\label{eq-lpn-bayesian-recordwise-next-contact}
\Pr_R\{\tau>t,C_{t+1}=S\}
\le
\mathbb E_R\!\left[
\mathbf1_{\{\tau>t\}}
\bigl(2^{-n}+C_{R,t}^{\rm LPN}+e_{R,t}^{\rm LPN}\bigr)
\right].
\end{equation}
Consequently its family-uniform active Bayesian optimum obeys
\begin{align}
\label{eq-lpn-saturated-structural-bayesian-arrival}
\operatorname{BayesArr}_{\mathfrak M,n,m,\eta,H}^{\rm LPN,sat}(B)
\le{}&
H\left(
2^{-n}
+\operatorname{PostConc}_{\mathfrak M^*,n,m,\eta,H}^{\rm LPN,sat}
(B_{\rm LPN}^{\rm Bayes})
+\operatorname{FilErr}_{\mathfrak M^*,n,m,\eta,H}^{\rm LPN,sat}
(B_{\rm LPN}^{\rm Bayes})
\right)\\
&+\varepsilon_{\rm good}(n,m,\eta,\tau,\delta).
\notag
\end{align}

The actionable-gain form of the same bound is
\begin{align}
\label{eq-lpn-actionable-structural-bayesian-arrival}
\operatorname{BayesArr}_{\mathfrak M,n,m,\eta,H}^{\rm LPN,sat}(B)
\le{}&
H\left(
2^{-n}
+\operatorname{BayesGain}_{\mathfrak M^*,n,m,\eta,H}^{\rm LPN,sat}
(B_{\rm LPN}^{\rm Bayes})
\right)\notag\\
&+\varepsilon_{\rm good}(n,m,\eta,\tau,\delta).
\end{align}
The actionable inequality has no additive filter-error term.  Every
target-aligned effect of the represented filter belongs to the deployed
next-test gain and retains its complete source ancestry.

For one record, the existing ordered-reveal information budget gives
\begin{equation}
\label{eq-lpn-posterior-concentration-information-candidate}
\mathfrak M^*_{n,m,\eta}
\mathbb E_R\!\left[
\mathbf1_{\{\tau>t\}}C_{R,t}^{\rm LPN}
\right]
\le
\min\left\{1-2^{-n},
\sqrt{\frac{\log2}{2}J_{R,t}^{\rm LPN,act}}
+\mathfrak M^*_{n,m,\eta}
 \mathbb E_R[\mathbf1_{\{\tau>t\}}e_{R,t}^{\rm LPN}]
\right\},
\end{equation}
where
\begin{equation}
\label{eq-lpn-posterior-concentration-information-budget}
J_{R,t}^{\rm LPN,act}
\le\min\{n,m(1-H_2(\eta))\}
+\sum_{j=1}^{L_{R,t}}d_{R,t,j}.
\end{equation}
Selection of a row, coordinate, or feature from the fixed public record adds
zero conditional information after \((A,b)\), as in
\cref{eq-lpn-fixed-record-active-zero-information}; its complete
computational cost and structural provenance remain charged.

\end{theorem}

\begin{proof}

For singleton targets, the vulnerability in
\cref{eq-target-vulnerability-functional} is maximum posterior mass and its
uniform-prior value is \(2^{-n}\).  Apply
\cref{thm-structural-bayesian-target-contact-accountability} conditionally
on each surviving history to obtain
\cref{eq-lpn-bayesian-recordwise-next-contact}.  Sum the disjoint first-hit
events, take the family-uniform saturated supremum, and use the unit bound on
the complement of the LPN good event.  Its probability is
\cref{eq-lpn-master-active-good-event-error}, proving
\cref{eq-lpn-saturated-structural-bayesian-arrival}.  Apply
\cref{eq-source-resolved-bayesian-arrival-composition} to the same record,
then take the family-uniform saturated supremum and restore the complement of
the good event.  This gives
\cref{eq-lpn-actionable-structural-bayesian-arrival}.

The exact conditional posterior satisfies
\[
\max_u\Pi_{R,t}(u)-2^{-n}
\le\lVert\Pi_{R,t}-\Pi_0^{\rm LPN}\rVert_{\rm TV}.
\]
The target-vulnerability functional is one-Lipschitz, so replacing the
exact posterior by \(\widehat\Pi_{R,t}\) adds at most
\(e_{R,t}^{\rm LPN}\).  Apply
\cref{eq-pre-hit-information-posterior-tv} with the LPN information budget
\cref{eq-lpn-master-active-information-budget} to prove
\cref{eq-lpn-posterior-concentration-information-candidate,eq-lpn-posterior-concentration-information-budget}.
The fixed-record assertion is
\cref{thm-lpn-three-active-acquisition-interfaces}.

\end{proof}

\begin{theorem}[Source-resolved structural Bayesian gain for LPN]
\label{thm-lpn-posterior-concentration-five-mode-exhaustion}

At the ceiling \(B_{\rm LPN}^{\rm Bayes}\), every represented LPN posterior
concentration used by an acquisition or transition rule has the exact
decomposition below.  The fixed indices \((n,m,\eta,H)\), restricted
functional \(\mathfrak M^*\), and common ceiling are suppressed in this
display.  The decomposition and recordwise comparisons are unconditional.
The complete ambient numerical row and the saturated numerical conclusions
below use
\cref{thm-lpn-residual-ancestry-attention-bound} at this ceiling.
\begin{equation}
\label{eq-lpn-posterior-concentration-five-mode-decomposition}
\operatorname{PostSrc}_{\mathfrak M^*}^{\rm LPN,sat}
=\max_{c\in\{X,A,Q_{\rm ex},Q_{\rm rev},Q_{\rm nr}\}}
\operatorname{PostSrc}_{\mathfrak M^*;c}^{\rm LPN,sat}.
\end{equation}

The actionable gain of the deployed next-test kernel has the source-resolved
five-family charges of
\cref{def-source-resolved-actionable-bayesian-gain}.  These are the LPN
restriction of
\cref{def-complete-source-resolved-learning-ledger,%
eq-bayesian-complete-learning-ledger-identification}.  Refine them by the
master support mode and exact state--coefficient signature and suppress the
fixed indices in the notation.  On the LPN good event their complete table is
\begin{equation}
\label{eq-lpn-source-resolved-bayesian-gain-table}
\begin{array}{c|c}
\text{source cell} & \text{framework conclusion}\\ \hline
\mathsf A,\mathsf Q_{\rm ex},\mathsf Q_{\rm rev},\mathsf Q_{\rm nr}
  &0\\
\mathsf X:\ \mathsf{Add},\ \text{residual-independent}
  &2^{-n}\\
\mathsf X:\ \text{complete ambient record}
  &\le 2^{-n}+(1-\eta)^n
    \quad\text{by \cref{thm-lpn-residual-ancestry-attention-bound}}\\
\mathsf X:\ \mathsf{Mult}
  &0\ \text{new source; retain the complete source record}\\
\mathsf X:\ \mathsf{Ord}
  &0\ \text{new source; retain the complete source record}\\
\mathsf X:\ \mathsf{Sym}
  &\text{retain the authenticated complete source record}\\
\mathsf X:\ \mathsf{Filt}
  &\text{retain the represented complete source record}.
\end{array}
\end{equation}
Authorized \(\Caus\), valid \(\Lift\), and \(\Bdry\) retain the same row;
their represented current, interface, and boundary costs remain in the
complete record.  Their gradient-orthogonal component is \(\Res\) and has
zero target-qualified gain.  Therefore every record satisfies the exact
framework comparison
\begin{equation}
\label{eq-lpn-total-source-resolved-bayesian-charge}
g_{R,t}^{\rm Bayes}
\le
\min\left\{
U_{R,t}^{\rm cert},
eH_B(n)\sum_{a\in\mathfrak J_{\Abs}^{5}}
\mathcal B_{\mathfrak M^*,R,t}^{(a)}
\right\}.
\end{equation}
The complete ambient record, including every derived operation and its cost,
is bounded at the same ceiling by
\cref{thm-lpn-complete-geom-structure-exhaustion,cor-lpn-complete-ambient-geom-source-seal}.

For the ordered target-dependent reveals in the same record, the
source-resolved information ledger obeys
\begin{align}
\label{eq-lpn-source-resolved-bayesian-information-add}
J_{R,t}^{\rm Bayes,(Add)}
&\le
\min\{n,m(1-H_2(\eta))\}
+\sum_{j=1}^{L_{R,t}}d_{R,t,j},\\
\label{eq-lpn-source-resolved-bayesian-information-derived}
J_{R,t}^{\rm Bayes,(Mult)}
=J_{R,t}^{\rm Bayes,(Ord)}
=J_{R,t}^{\rm Bayes,(Filt)}
&=0
\end{align}
after conditioning on their earlier state-dependent arguments.  Authenticated
symmetry and derived channels inherit the information of their earlier
source.  Each homogeneous noisy parity reveal may use the SDPI coefficient
\((1-2\eta)^2\).

Let \(U_{R,t}^{\rm Bayes,info}\) be the resulting neutral-excess information
candidate from
\cref{eq-noise-contracted-source-resolved-bayesian-gain}, let
\(U_{R,t}^{\rm Bayes,learn}\) be the smaller applicable structural
Rademacher or martingale PAC--Bayes action-value candidate from
\cref{eq-lpn-complete-learning-control-radius}, and let
\(U_{R,t}^{\rm Bayes,dyn}\) be the same-record minimum of the applicable
capacity, Blackwell, observation-resolution, Caus, and Lift candidates in
\cref{cor-lpn-recordwise-readout-agnostic-envelope}.  Unavailable candidates
have value one.  Then
\begin{equation}
\label{eq-lpn-recordwise-posterior-certificate-minimum}
g_{R,t}^{\rm Bayes}
\le
\min\left\{
1,
U_{R,t}^{\rm Bayes,info},
U_{R,t}^{\rm Bayes,learn},
U_{R,t}^{\rm Bayes,dyn},
U_{R,t}^{\rm cert}
\right\}.
\end{equation}
All candidates use the same pre-hit record, conditional law, target-contact
normalization, and common charged ceiling.  An unavailable candidate has the
neutral value one.  The directly qualified ambient branch remains part of
the complete prospective Geom profile, as prescribed by
\cref{thm-no-uncharged-prospective-value-amplification}.
Every displayed candidate retains its explicit
\((n,m,\eta,H,B,\delta)\) arguments; each binary reveal uses
\((1-2\eta)^2\), while clean-risk correction uses the distinct inverse
factor \((1-2\eta)^{-1}\).

Taking the family-uniform saturated supremum and applying the complete
ambient seal at the same class-preserving ceiling gives
\begin{equation}
\label{eq-lpn-saturated-source-resolved-bayesian-gain}
\operatorname{BayesGain}_{\mathfrak M^*,n,m,\eta,H}^{\rm LPN,sat}(B)
\le
2^{-n}+(1-\eta)^n.
\end{equation}
The unrestricted family functional retains the good-presentation error
\(\varepsilon_{\rm good}(n,m,\eta,\tau,\delta)\).
Combining this estimate with
\cref{eq-lpn-actionable-structural-bayesian-arrival} gives
\begin{equation}
\label{eq-lpn-source-resolved-structural-bayesian-arrival}
\begin{aligned}
\operatorname{BayesArr}_{\mathfrak M,n,m,\eta,H}^{\rm LPN,sat}(B)
{}&\le
H\bigl(2^{1-n}+(1-\eta)^n\bigr)\\
&\qquad
+\varepsilon_{\rm good}(n,m,\eta,\tau,\delta).
\end{aligned}
\end{equation}
The saturated Bayesian class is a subprofile of the same complete saturated
source profile.  No additional profile coordinate is introduced; the
residual-ancestry row is evaluated by
\cref{thm-lpn-residual-ancestry-attention-bound}.

\end{theorem}

\begin{proof}

Apply
\cref{thm-affordable-posterior-concentration-source-routing} to the LPN
record class.  The exact record-class partition proves
\cref{eq-lpn-posterior-concentration-five-mode-decomposition}.  The exact and
compensated nonambient modes vanish by
\cref{eq-lpn-master-active-nonambient-vanishing,eq-lpn-exact-abs-source-vanishes,eq-lpn-exact-quotient-geom-closed,eq-lpn-compensated-source-vanishes}.
The charge-bearing vertices are prospective pre-hit ancestors by
\cref{cor-prehit-status-bayesian-source-charges}; hence the LPN source seals
apply to the charge ledger itself rather than to a retrospective first-hit
label.

For the ambient mode, the source tags are exhausted by
\cref{thm-lpn-no-multiplicative-source,thm-lpn-order-provenance,cor-lpn-symmetry-provenance-geom-closure,thm-lpn-filtration-provenance}.
Source inheritance under Caus, Lift, and Bdry and the residual assignment are
\cref{prop-derived-channels-create-no-structure,thm-lpn-valid-lift-quantitative-transfer,cor-lpn-no-nonreversible-geometric-escape}.
These results prove the provenance rows of the table.  The complete ambient
evaluation in
\cref{thm-lpn-complete-geom-structure-exhaustion,cor-lpn-complete-ambient-geom-source-seal}
bounds the resulting complete source record at the same saturated ceiling.
Apply
\cref{thm-exact-source-resolved-structural-bayesian-gain,thm-saturated-same-record-prospective-certificate-envelope}
to obtain \cref{eq-lpn-total-source-resolved-bayesian-charge}.  The complete
ambient branch is bounded directly by the complete LPN source seal.

The Add information budget is
\cref{eq-lpn-posterior-concentration-information-budget}.  Conditional on
the earlier arguments, deterministic Mult, Ord, and Filt vertices add zero
information and zero contextual charge by
\cref{lem-derived-constructor-zero-contextual-charge,thm-filtration-provenance-charging}.
The binary-noise coefficient is
\cref{eq-binary-noise-source-sdpi-candidate}.  This proves
\cref{eq-lpn-source-resolved-bayesian-information-add,eq-lpn-source-resolved-bayesian-information-derived}.

Apply the recordwise minimum in
\cref{thm-noise-contracted-source-resolved-bayesian-gain}, together with the
LPN learning and dynamical candidates cited in the statement, to obtain
\cref{eq-lpn-recordwise-posterior-certificate-minimum}.  Taking the
family-uniform record supremum proves
the certificate part of the claim.  The Bayesian record class is a
restriction of the complete saturated source class, so
\cref{eq-lpn-complete-ambient-geom-source-bound} proves
\cref{eq-lpn-saturated-source-resolved-bayesian-gain}.  Restriction to the
good event is removed with the explicit complement probability in
\cref{eq-lpn-master-active-good-event-error}.  Substitution into
\cref{eq-lpn-actionable-structural-bayesian-arrival} proves
\cref{eq-lpn-source-resolved-structural-bayesian-arrival}.

\end{proof}

\begin{corollary}[Optimal structural extraction ledger for LPN]
\label{cor-lpn-optimal-structural-extraction-ledger}

Let \(S\) be uniform on \(\mathbb F_2^n\), let the fixed public presentation
contain \(m\) independent binary-symmetric parity reveals of noise rate
\(0\le\eta\le1/2\), and let \(D_{R,t}^{\rm ext}\) be the total conditional
information in target-dependent advice, oracle answers, verifier outcomes, or
feedback revealed after \((A,b)\) and before the deployed test.  Every such
coordinate is included in the ordered reveal ledger with its complete source
and cost.  Then
\begin{equation}
\label{eq-lpn-optimal-extraction-information-bound}
J_{R,t}^{\rm LPN,opt}
\le
\min\left\{
n,
m\bigl(1-H_2(\eta)\bigr)+D_{R,t}^{\rm ext},
m(1-2\eta)^2+D_{R,t}^{\rm ext}
\right\}.
\end{equation}
Selection and postprocessing of the fixed public record contribute zero to
\(D_{R,t}^{\rm ext}\); their computation and structural provenance remain in
the complete record.

At the compiled ceiling of
\cref{thm-optimal-structural-extraction-saturated-ceiling}, the deployed
singleton-target gain satisfies
\begin{equation}
\label{eq-lpn-optimal-extraction-modality-minimum}
g_{R,t}^{\rm LPN}
\le
\min\left\{
U_{R,t}^{\rm src},U_{R,t}^{\Caus},U_{R,t}^{\Lift},
U_{R,t}^{\Bdry},U_{R,t}^{\rm Bell},
\beta_{R,t}^{\rm LPN,+}
+\sqrt{\frac{\log2}{2}J_{R,t}^{\rm LPN,opt}}
\right\}.
\end{equation}
Here the source coordinate is the exact LPN restriction of the five-family
and master-support ledger in
\cref{eq-lpn-source-resolved-bayesian-gain-table}; the Caus, Lift, boundary,
and Bellman coordinates are the LPN evaluations already collected in
\cref{eq-lpn-recordwise-posterior-certificate-minimum}.  The residual
coordinate is zero after target qualification.

The parameter dependence has the following exact limiting behavior:
\begin{enumerate}[label=\textup{(\roman*)}]
\item at \(\eta=0\), the information candidate permits \(n\) bits and is
      consistent with represented Gaussian elimination;
\item when \(\eta m=c\log n\), the subset-elimination source retains the
      value \((1-\eta)^n=n^{-cR+o(1)}\) from
      \cref{eq-lpn-subset-ge-log-noise-window};
\item at fixed \(0<\eta<1/2\), the channel contribution is explicitly
      bounded by the two noisy terms in
      \cref{eq-lpn-optimal-extraction-information-bound}, while target arrival
      uses the smaller complete structural coordinate in
      \cref{eq-lpn-optimal-extraction-modality-minimum};
\item at \(m(1-H_2(\eta))<n\), the information candidate recovers the finite
      LPN converse of \cref{thm-lpn-information-threshold}.
\end{enumerate}

\end{corollary}

\begin{proof}
Apply \cref{cor-noise-sample-target-size-optimal-extraction} with
\(N=2^n\), BSC capacity \(1-H_2(\eta)\), and SDPI coefficient
\((1-2\eta)^2\).  Each parity input contains at most one available bit, so
the SDPI sum is at most \(m(1-2\eta)^2\).  The conditional chain rule appends
\(D_{R,t}^{\rm ext}\), proving
\cref{eq-lpn-optimal-extraction-information-bound}.  Substitute this bound
in \cref{eq-bsc-optimal-extraction-vulnerability} and use the LPN
same-record candidates cited in the statement to obtain
\cref{eq-lpn-optimal-extraction-modality-minimum}.  The four regime statements
are direct substitutions and the cited LPN converse.
\end{proof}

\begin{corollary}[Compact latent-attention ledger for LPN]
\label{cor-lpn-compact-latent-attention-ledger}

Adopt the LPN presentation and ordered-reveal convention of
\cref{cor-lpn-optimal-structural-extraction-ledger}.  Let \(R\) be a complete
compact latent-attention LPN record, let \(b_{R,t}^{\rm lat}\) be its code
capacity, and let \(I_{R,t}^{\rm LPN,lat}=I(S;Z_t\mid\tau>t)\).  Then
\begin{equation}
\label{eq-lpn-compact-latent-information-bound}
I_{R,t}^{\rm LPN,lat}
\le
\min\left\{
b_{R,t}^{\rm lat},n,
m\bigl(1-H_2(\eta)\bigr)+D_{R,t}^{\rm ext},
m(1-2\eta)^2+D_{R,t}^{\rm ext},
J_{R,t}^{\rm LPN,opt}
\right\}.
\end{equation}
If the code has \(d_t\) coordinates at \(p_t\)-bit precision and exact secret
recovery has probability at least \(\alpha\), then
\begin{equation}
\label{eq-lpn-compact-latent-size-necessary}
d_tp_t+b_t^{\rm code}\ge\alpha n-1.
\end{equation}

Let \(U_{R,t}^{\rm LPN,lat}\) be the LPN restriction of
\cref{eq-unified-compact-latent-certificate}, and let
\(U_{R,t}^{\rm LPN,att}\) be the attention envelope in
\cref{eq-source-resolved-attention-envelope} for the same record.  The deployed
gain satisfies
\begin{equation}
\label{eq-lpn-compact-latent-attention-gain}
g_{R,t}^{\rm LPN,lat,att}
\le
\min\left\{
U_{R,t}^{\rm LPN,lat},
U_{R,t}^{\rm LPN,att},
U_{R,t}^{\rm src},U_{R,t}^{\Caus},U_{R,t}^{\Lift},
U_{R,t}^{\Bdry},U_{R,t}^{\rm Bell}
\right\}.
\end{equation}
All entries refer to the same pre-hit law and compiled ceiling.  Attention
scores derived from \((A,b)\), their normalizers, value paths, latent memory,
posterior decoder, and selected transitions retain the complete LPN source
ancestry.  Optimal acquisition is the Bellman optimization over these same
records, as in
\cref{thm-optimal-compact-latent-acquisition,thm-joint-compact-latent-attention-acquisition}.

The parameter regimes are consistent with the existing LPN phase ledger:
\begin{enumerate}[label=\textup{(\roman*)}]
\item at \(\eta=0\), the information and code candidates permit an \(n\)-bit
      exact solution representation;
\item at \(\eta=1/2\), the binary reveals contribute zero channel information;
\item when \(\eta m=c\log n\), the subset-elimination candidate retains
      \((1-\eta)^n=n^{-cR+o(1)}\);
\item at fixed \(0<\eta<1/2\), target arrival uses the minimum in
      \cref{eq-lpn-compact-latent-attention-gain}, rather than the information
      availability bound alone.
\end{enumerate}

\end{corollary}

\begin{proof}
Apply \cref{cor-noise-sample-code-capacity-latent-beliefs} with
\(N=2^n\) and the LPN binary-symmetric reveals.  Intersect its bound with
\cref{eq-lpn-optimal-extraction-information-bound}; this proves
\cref{eq-lpn-compact-latent-information-bound}.  The latent-size condition is
\cref{eq-latent-coordinate-fano-necessary}.  Apply
\cref{thm-unified-compact-latent-structural-certificate,thm-joint-compact-latent-attention-acquisition}
and the LPN same-record modality minimum
\cref{eq-lpn-optimal-extraction-modality-minimum} to obtain
\cref{eq-lpn-compact-latent-attention-gain}.  The regime statements are direct
substitutions and
\cref{eq-lpn-subset-ge-log-noise-window}.
\end{proof}

\begin{theorem}[Optimal latent representation and useful belief update for LPN]
\label{thm-lpn-optimal-latent-useful-belief-update}

Restrict to complete LPN latent records constructed from the public
presentation \((A,b)\), target-independent auxiliary randomness, and the
represented pre-hit transcript.  Thus \(D_{R,t}^{\rm ext}=0\): selecting,
reordering, or postprocessing public rows and deterministic verifier outcomes
adds no new target-dependent reveal, while every such operation retains its
complete cost and source ancestry.  Write
\begin{equation}
\label{eq-lpn-public-latent-information-ceiling}
J_{n,m,\eta}^{\rm lat}(b)
=\min\left\{
b,n,m\bigl(1-H_2(\eta)\bigr),m(1-2\eta)^2
\right\}.
\end{equation}
Every code-capacity-\(b\) record satisfies
\begin{equation}
\label{eq-lpn-every-public-latent-information-bound}
I(S;Z_t\mid\tau>t)
\le J_{n,m,\eta}^{\rm lat}(b).
\end{equation}
If its represented decoder identifies \(S\) with probability at least
\(\alpha\), then
\begin{equation}
\label{eq-lpn-public-latent-size-bound}
b\ge\alpha n-1,
\qquad
d_tp_t+b_t^{\rm code}\ge\alpha n-1.
\end{equation}

The useful part of a belief update is its deployed actionable gain
\(g_{R,t}^{\rm Bayes}\).  For every complete record it obeys the same-record
bound
\begin{equation}
\label{eq-lpn-public-latent-recordwise-useful-gain}
g_{R,t}^{\rm Bayes}
\le
\min\left\{
U_{R,t}^{\rm opt},
U_{R,t}^{\rm lat,stat},
U_{R,t}^{\rm att},
U_{R,t}^{\Caus},U_{R,t}^{\Lift},U_{R,t}^{\Bdry},U_{R,t}^{\rm Bell},
\beta_{R,t}^{\rm Bayes,+}
+\sqrt{\frac{\log2}{2}J_{n,m,\eta}^{\rm lat}(b)}
\right\}.
\end{equation}
Every entry refers to the same acquired data, latent encoder, recurrent
update, decoder, attention selector, pre-hit law, and charged ceiling.
Accordingly, the bound simultaneously optimizes the sampling schedule through
the resource-valued Bellman recursion, the represented posterior update
through the Bayes--Fisher flow, the latent representation through the
code-capacity and sufficiency certificates, and the learned deployment rule
through the source-resolved statistical certificate.

Let
\begin{equation}
\label{eq-lpn-public-latent-useful-gain-profile}
\operatorname{LatBayesGain}_{\mathfrak M^*,n,m,\eta,H}^{\rm LPN,sat}(B,b)
=\max_{t<H}\sup_{R\in\mathscr R_{{\rm LPN},H}^{\rm lat}(B,b)}
\mathfrak M^*_{n,m,\eta}\!\left(g_{R,t}^{\rm Bayes}\right)
\end{equation}
be the LPN specialization of
\cref{eq-saturated-compact-latent-gain-profile}.  On the good-presentation
event, the supremum ranges over one family-uniform encoder, update, decoder,
acquisition, and deployment scheme on the whole parameter slice, and
\begin{equation}
\label{eq-lpn-optimal-latent-useful-gain-bound}
\operatorname{LatBayesGain}_{\mathfrak M^*,n,m,\eta,H}^{\rm LPN,sat}(B,b)
\le
\operatorname{BayesGain}_{\mathfrak M^*,n,m,\eta,H}^{\rm LPN,sat}(B)
\le 2^{-n}+(1-\eta)^n.
\end{equation}
For the unrestricted presentation functional, the right side becomes
\begin{equation}
\label{eq-lpn-optimal-latent-useful-gain-bound-unrestricted}
2^{-n}+(1-\eta)^n
+\varepsilon_{\rm good}(n,m,\eta,\tau,\delta).
\end{equation}
Taking the union over every latent capacity admitted by the polynomial
saturated budget gives
\begin{equation}
\label{eq-lpn-optimal-latent-useful-gain-polynomial-ceiling}
\operatorname{LatBayesGain}_{\mathfrak M^*,n,m,\eta,H}^{\rm LPN,poly,sat}
=
\operatorname{BayesGain}_{\mathfrak M^*,n,m,\eta,H}^{\rm LPN,poly,sat}
\le 2^{-n}+(1-\eta)^n.
\end{equation}
Thus the displayed ceiling ranges over every family-uniform latent
representation constructible from the public presentation at saturated
polynomial cost; it is not restricted to a chosen embedding architecture or
learning rule.

Finally, the loss caused by compressing the complete pre-hit history into a
particular latent representation remains explicit.  For each presented LPN
instance \(I\), let \(A_{I,R}^H\) be its arrival value.  Then
\begin{equation}
\label{eq-lpn-latent-representation-deployment-loss}
\begin{aligned}
0\le{}&
\operatorname{BayesArr}_{I,H}^{\rm LPN,sat}(B)-A_{I,R}^H
\\
\le{}&
\sum_{t<H}\mathbb E_R\!\left[\mathbf1_{\{\tau>t\}}
\left(
\sqrt{\frac{\log2}{2}\Delta_{R,t}^{\rm suf}}
+e_{R,t}^{\rm dec}
+2\Gamma_{R,t}^{\rm lat}
+\varepsilon_{R,t}^{\rm opt}
\right)\right]
+\Delta_{\rm fil}^{H}.
\end{aligned}
\end{equation}
Hence code capacity controls retained information, the sufficiency and decoder
terms control representation loss, and
\cref{eq-lpn-optimal-latent-useful-gain-bound} controls the maximum portion of
the resulting belief update that can improve the next target test.

\end{theorem}

\begin{proof}
With \(D_{R,t}^{\rm ext}=0\),
\cref{eq-lpn-public-latent-information-ceiling,eq-lpn-every-public-latent-information-bound}
follow from \cref{eq-lpn-compact-latent-information-bound}; the coordinate and code
requirements are
\cref{eq-latent-code-fano-necessary,eq-latent-coordinate-fano-necessary}.
The recordwise minimum combines
\cref{eq-unified-compact-latent-certificate,eq-source-resolved-attention-envelope,eq-lpn-optimal-extraction-modality-minimum}
after substituting the public-only information ceiling.  These comparisons
refer to one record and one conditional law, so their minimum proves
\cref{eq-lpn-public-latent-recordwise-useful-gain}.

Class inclusion
\cref{eq-compact-latent-gain-profile-inclusion} gives the first comparison in
\cref{eq-lpn-optimal-latent-useful-gain-bound}; the complete LPN source
evaluation \cref{eq-lpn-saturated-source-resolved-bayesian-gain} gives the
second.  Restoring the complement of the good event gives
\cref{eq-lpn-optimal-latent-useful-gain-bound-unrestricted}.  The exact
compiler identity
\cref{eq-compact-latent-polynomial-gain-ceiling} proves
\cref{eq-lpn-optimal-latent-useful-gain-polynomial-ceiling}.  It preserves the
next-test law and therefore introduces no comparison loss.  The final display
is the LPN specialization of
\cref{eq-learned-compact-latent-policy-loss}.
\end{proof}

\begin{corollary}[Learned LPN structural Bayesian policy]
\label{cor-learned-lpn-structural-bayesian-policy}

Let a complete LPN learner construct a represented belief and action-value
rule with filter path error \(\Delta_{\rm fil}^{H}\), optimization error
\(\varepsilon_{\rm opt}\), and the clean action-value radius in
\cref{eq-lpn-complete-learning-control-radius}.  Then
\begin{equation}
\label{eq-learned-lpn-structural-bayesian-loss}
0\le
\operatorname{BayesArr}_{\mathfrak M,n,m,\eta,H}^{\rm LPN,sat}(B)
-\mathfrak M_{n,m,\eta}(A_{\widehat R}^{H})
\le
H\left(2\Gamma_{{\rm LPN},{\rm act}}^{\rm clean}
+\varepsilon_{\rm opt}\right)
+\mathfrak M_{n,m,\eta}(\Delta_{{\rm fil},\widehat R}^{H}).
\end{equation}
The factor \((1-2\eta)^{-1}\), effective sample size, dependence error,
confidence allocation, codebook size, and optimization cost are therefore
exactly those already displayed in
\cref{eq-lpn-complete-learning-control-radius}.  The PAC--Bayes posterior
over policies and the LPN task belief over \(S\) remain separately typed and
separately charged.

\end{corollary}

\begin{proof}

Apply \cref{cor-learned-structural-bayesian-control} and substitute
\cref{eq-lpn-complete-learning-control-radius}.  Bound the number of
surviving loss terms by \(H\).

\end{proof}

\begin{figure}[tbp]
\centering
\resizebox{\linewidth}{!}{%
\begin{tikzpicture}[fw diagram,node distance=7mm and 9mm]
  \node[fw source,text width=2.5cm] (secret)
    {uniform secret\\\(S\in\mathbb F_2^n\)};
  \node[fw provenance,text width=2.5cm,below=10mm of secret] (matrix)
    {target-independent\\matrix \(A\)};
  \node[fw use,text width=2.6cm,right=13mm of secret,yshift=-5mm] (rhs)
    {first target-dependent\\reveal \(b=AS+e\)};
  \node[fw geom,text width=2.8cm,right=of rhs] (belief)
    {charged likelihood,\\normalizer and belief};
  \node[fw box,text width=2.8cm,right=of belief] (action)
    {optimal structural\\acquisition rule};
  \node[fw terminal,text width=2.5cm,right=of action] (hit)
    {candidate contact\\with \(S\)};
  \node[fw residual,text width=4.0cm,below=10mm of belief] (profile)
    {five source charges\\information, learning and dynamics};
  \draw[fw arrow] (secret)--(rhs);
  \draw[fw arrow] (matrix)--(rhs);
  \draw[fw arrow] (rhs)--(belief);
  \draw[fw arrow] (belief)--(action);
  \draw[fw arrow] (action)--(hit);
  \draw[fw arrow] (belief)--(profile);
  \draw[fw arrow] (profile)-|(hit);
\end{tikzpicture}%
}
\caption{Structural Bayesian learning for LPN.  The uniform prior supplies
the neutral \(2^{-n}\) contact.  Every noisy reveal retains its source and
SDPI coefficient.  Likelihood, normalization, posterior approximation,
acquisition, and deployment are deterministic charged descendants; the
actionable next-test gain is bounded by the five source charges and the
same-record learning and dynamical certificates.}
\label{fig-lpn-structural-bayesian-learning}
\end{figure}

\begin{figure}[tbp]
\centering
\resizebox{\linewidth}{!}{%
\begin{tikzpicture}[fw diagram,node distance=8mm and 10mm]
  \node[fw source,text width=2.7cm] (lpn)
    {LPN target\\and public record};
  \node[fw geom,text width=3.0cm,above right=9mm and 12mm of lpn] (fixed)
    {fixed presentation\\select public data and transport};
  \node[fw use,text width=3.0cm,right=13mm of lpn] (passive)
    {passive sample stream\\fresh uniform rows};
  \node[fw provenance,text width=3.0cm,below right=9mm and 12mm of lpn] (chosen)
    {chosen-query extension\\selected parity rows};
  \node[fw terminal,text width=3.1cm,right=18mm of passive] (ledger)
    {master active profile\\information and target arrival};
  \draw[fw arrow] (lpn)--(fixed);
  \draw[fw arrow] (lpn)--(passive);
  \draw[fw arrow] (lpn)--(chosen);
  \draw[fw arrow] (fixed)--(ledger);
  \draw[fw arrow] (passive)--(ledger);
  \draw[fw arrow] (chosen)--(ledger);
\end{tikzpicture}%
}
\caption{The three LPN acquisition interfaces.  Selection inside the fixed
public record creates no new information channel.  Passive fresh samples and
chosen parity queries are declared extensions with different quantitative
bounds.  Each interface enters the same master active-sampling ledger.}
\label{fig-lpn-three-active-acquisition-interfaces}
\end{figure}

\section{Charged Decoder Geometry and Current Classification}
\label{sec-lpn-charged-decoder-current-classification}

\subsection{Post-Output Decoder Geometry and Level Quotients}
\label{subsec-lpn-postoutput-decoder-geometry}

\begin{theorem}[Charged post-output decoder geometry for LPN]
\label{thm-lpn-charged-postoutput-decoder-geometry}

Fix \(n\ge2\), \(0\le\eta<\tau<1/2\), an LPN instance
\(I=(A,b,\eta)\), and a family-uniform
represented decoder derivation \(\rho_{\mathcal D}\) of complete saturated
cost at most \(B_{\mathcal D}\).  Let

\[
\widehat s=\mathcal D_n(A,b;\omega)
\]

be its output, where the random seed \(\omega\), when present, is carried by
the public product lift.  Immediately after this output is represented, and
before the correction transitions below, define

\begin{align}
\label{eq-lpn-decoder-translation-potential}
\Theta_{\widehat s}(x)&=x+\widehat s,
&D_{\widehat s}(x)&=|x+\widehat s|_1,\\
\label{eq-lpn-decoder-correction-kernel}
P_{\widehat s}(x,x+e_i)
&=\frac1n\mathbf1_{\{x_i\ne\widehat s_i\}},
&P_{\widehat s}(x,x)&=1-\frac{D_{\widehat s}(x)}n.
\end{align}

All other transition probabilities are zero.  Let \(H\ge1\) be the charged
correction horizon, with public correction-phase initialization \(x_0=0\).
Evaluator normalization and computation closure produce
one derivation \(\rho_{\mathrm{corr}}\), containing
\(\rho_{\mathcal D}\), the stored output, the public verifier evaluation, the
potential and conductance evaluators, the fixed-space current, the correction
kernel, and its clocked transcript, with

\begin{equation}
\label{eq-lpn-decoder-geometry-complete-cost}
\operatorname{Cost}_I(\rho_{\mathrm{corr}})
\preceq
B_{\mathcal D}+mn+n+H(n+\log n).
\end{equation}

The implicit hypercube relation and its conductance are generated by the
coordinate index \(i\); no \(2^n\)-entry table occurs in this bound.  In
particular, for \(m=\Theta(n)\), \(B_{\mathcal D}\le n^C\), and

\[
H=\left\lceil
n\left(\log n+\log\frac1\delta\right)
\right\rceil,
\qquad 0<\delta<1,
\]

the complete cost in
\eqref{eq-lpn-decoder-geometry-complete-cost} is bounded by one fixed
polynomial in \(n\), uniformly over the family.

This bound constructs the represented hypercube carrier centered at
\(\widehat s\).  Target-oriented admissibility is a separate conjunct in the
same record.  By the progress-observable condition of
\cref{def-geom-extraction}, the potential must vanish on
\(T_I=\{s\}\), whereas

\begin{equation}
\label{eq-lpn-decoder-center-admissibility}
D_{\widehat s}(s)=d_H(s,\widehat s),
\qquad
D_{\widehat s}(s)=0
\Longleftrightarrow
\widehat s=s.
\end{equation}

Thus the carrier is an admissible target-oriented Geom record precisely on
decoder success.  When \(\widehat s\ne s\), it is centered on the wrong
state, fails target normalization, and contributes zero under the qualified
record convention.  Its spectral gap receives no separate target-aligned
credit.

On the event \(\{\widehat s=s\}\), the fixed-space split of
\(P_{\widehat s}-I\) supplies a dynamically qualified
identity-supported ambient record with the following exact gate values:

\begin{align}
\label{eq-lpn-decoder-geometry-exact-gates}
q_{\mathrm{amb}}&=\mathrm{id},
&M_{\Phi_{\widehat s}}&=1,
&\chi_{\Phi_{\widehat s}}^{\mathrm{auth}}
  (D_{\widehat s})&=1,\\
\operatorname{Align}_{\Caus}^{\mathrm{auth}}
  (\mathfrak K_{\widehat s},D_{\widehat s})&=1,
&R_{\{s\}}(D_{\widehat s})&=1,
&\rho_{\Phi_{\widehat s}}(D_{\widehat s})&=\frac1n.
\end{align}

Consequently,

\begin{equation}
\label{eq-lpn-decoder-geometry-exact-visibility}
V^X_{\Phi_{\widehat s},D_{\widehat s},
J_{\widehat s},\mathfrak K_{\widehat s}}
=\frac{n}{n+1}.
\end{equation}

Combining
\cref{eq-lpn-decoder-center-admissibility,eq-lpn-decoder-geometry-exact-visibility},
the qualified target-aligned contribution of this represented construction
is exactly

\begin{equation}
\label{eq-lpn-decoder-qualified-geom-contribution}
\mathbf1_{\{\widehat s=s\}}
V^X_{\Phi_{\widehat s},D_{\widehat s},
J_{\widehat s},\mathfrak K_{\widehat s}}
=
\frac{n}{n+1}\mathbf1_{\{\widehat s=s\}}.
\end{equation}

Thus the inverse-polynomial carrier gap, same-generator consistency,
authorized alignment, conditional target reach, and regularity factor are
all multiplied by the represented target-admissibility indicator.  No
factor from the incorrectly centered carrier enters the target-aligned
profile.

\end{theorem}

\begin{proof}
Retaining the decoder derivation, output, verifier, potential, implicit cube
relation, and clocked correction rule gives the displayed polynomial cost.
The potential vanishes at the true target exactly on decoder success.  On
that event direct evaluation of the selected correction kernel gives all
five gates and regularity \(1/n\); off that event target admissibility is
zero.  Multiplying the gates yields the exact qualified contribution.
\end{proof}

\begin{theorem}[Exact decoder level quotient and correction arrival]
\label{thm-lpn-decoder-level-quotient-arrival}
Adopt the charged decoder construction of
\cref{thm-lpn-charged-postoutput-decoder-geometry}.

The same record exposes the represented level-set quotient

\begin{equation}
\label{eq-lpn-decoder-level-quotient}
q_{\widehat s}:X_I\longrightarrow\{0,1,\ldots,n\},
\qquad
q_{\widehat s}(x)=D_{\widehat s}(x).
\end{equation}

On \(\{\widehat s=s\}\), this is a nonidentity terminal-compatible exact
quotient.  Its quotient chain is

\begin{equation}
\label{eq-lpn-decoder-level-quotient-chain}
w\longmapsto
\begin{cases}
w-1,&\text{with probability }w/n,\\
w,&\text{with probability }1-w/n,
\end{cases}
\qquad 1\le w\le n,
\end{equation}

and \(0\) is absorbing.  Thus its lumpability defect is zero, its canonical
comparison factor is one, and its target-projection ratio is one.  With the
displayed horizon and \(\lambda=1/H\),

\begin{align}
\label{eq-lpn-decoder-correction-arrival}
\Pr\{\tau_{\{s\}}>H\mid\widehat s=s\}
&\le ne^{-H/n}\le\delta,\\
\label{eq-lpn-decoder-level-quotient-visibility}
V_I(q_{\widehat s})
&\ge\frac{1-\delta}{eH}.
\end{align}

\end{theorem}

\begin{proof}
When \(\widehat s=s\), Hamming weight is a strong lumping of the correction
kernel and gives the displayed pure-death chain.  Each incorrect coordinate
survives \(H\) updates with probability at most \(e^{-H/n}\); a union bound
gives the arrival estimate.  Substitution into the exact quotient visibility
formula with \(\lambda=H^{-1}\) gives its lower bound.
\end{proof}

\begin{corollary}[Family accounting for charged decoder geometry]
\label{cor-lpn-charged-decoder-family-accounting}
Adopt the constructions of
\cref{thm-lpn-charged-postoutput-decoder-geometry,thm-lpn-decoder-level-quotient-arrival}.

Let

\[
p_{n,\eta}(\mathcal D)
=
\Pr_{A,s,E,\omega}\{\mathcal D_n(A,As+E;\omega)=s\},
\]

and put

\begin{equation}
\label{eq-lpn-decoder-verifier-exception}
\zeta_{n,m,\eta,\tau}
=
\exp[-mD_{\mathrm{Ber}}(\tau\Vert\eta)]
+(2^n-1)2^{-m(1-H_2(\tau))}.
\end{equation}

At the complete charged budget

\[
B_{\mathrm{corr}}
\succeq
B_{\mathcal D}+mn+n+H(n+\log n),
\]

the completed post-output audit therefore satisfies

\begin{align}
\label{eq-lpn-decoder-success-charged-to-geom}
&\mathbb E\!\left[
\mathbf1_{\{V_I^{-1}(1)=\{s\},\ \widehat s=s\}}
V^X_{\Phi_{\widehat s},D_{\widehat s},
J_{\widehat s},\mathfrak K_{\widehat s}}
\right]
=\frac{n}{n+1}
\Pr\!\left\{V_I^{-1}(1)=\{s\},\ \widehat s=s\right\},\\
\label{eq-lpn-decoder-success-charged-to-abs}
&\mathbb E\!\left[
\mathbf1_{\{V_I^{-1}(1)=\{s\},\ \widehat s=s\}}
V_I(q_{\widehat s})
\right]
\ge\frac{1-\delta}{eH}
\Pr\!\left\{V_I^{-1}(1)=\{s\},\ \widehat s=s\right\}.
\end{align}

Thus the decoder cost is an authenticated ancestor of both post-output
records, and the complete dependence on the noise law is retained in the
represented success mass \(p_{n,\eta}(\mathcal D)\) and the explicit verifier
exception \eqref{eq-lpn-decoder-verifier-exception}.  These displays audit
structure constructed after the decoder output.  They do not make that
record a prospective source for the decoding transitions that produced
\(\widehat s\).
For this particular represented construction, before taking the saturated
supremum, the family-functional target-aligned Geom contribution on the
verifier good event is exactly

\begin{equation}
\label{eq-lpn-decoder-qualified-family-geom-contribution}
\frac{n}{n+1}
\Pr\!\left\{
V_I^{-1}(1)=\{s\},\ \widehat s=s
\right\}.
\end{equation}

It lies in the explicit interval

\begin{equation}
\label{eq-lpn-decoder-qualified-family-geom-interval}
\frac{n}{n+1}
\bigl(p_{n,\eta}(\mathcal D)-\zeta_{n,m,\eta,\tau}\bigr)_+
\le
\frac{n}{n+1}
\Pr\!\left\{
V_I^{-1}(1)=\{s\},\ \widehat s=s
\right\}
\le
\frac{n}{n+1}p_{n,\eta}(\mathcal D).
\end{equation}

\end{corollary}

\begin{proof}
Intersect decoder success with the verifier event and substitute the exact
Geom visibility and quotient lower bound.  The verifier complement is
bounded by its explicit binomial/code probability.  Since the verifier event
can remove at most that probability from decoder success, the two-sided
family interval follows.  Every term retains the decoder derivation as an
ancestor, so the conclusion is retrospective.
\end{proof}

\begin{corollary}[Fully charged post-output LPN geometry and exact level-quotient audit]
\label{cor-lpn-charged-decoder-geometry-level-quotient}
\label{thm-lpn-charged-decoder-geometry-level-quotient}
The post-output geometry, exact level quotient, and family accounting are
\cref{thm-lpn-charged-postoutput-decoder-geometry,thm-lpn-decoder-level-quotient-arrival,cor-lpn-charged-decoder-family-accounting}.
They are retrospective consequences of one charged decoder record and are
not prospective sources for the computation that produced its output.
\end{corollary}

\begin{proof}

The derivation \(\rho_{\mathcal D}\) is retained as an ancestor when its
\(n\)-bit output is stored.  Matrix--vector multiplication and the public
verifier require at most \(O(mn)\) Boolean operations.  The translation and
Hamming potential require \(O(n)\) operations.  A correction step stores the
current \(n\)-bit candidate, generates an index using \(O(\log n)\) random
bits, and applies one comparison and update.  Charging the complete clocked
state record gives the conservative \(H(n+\log n)\) term.  The constructor
cost rule of
\cref{def-budgeted-derivability-closure} proves
\eqref{eq-lpn-decoder-geometry-complete-cost}.  The hypercube relation is the
uniform evaluator \((x,i)\mapsto x+e_i\), and the conductance coefficient has
an \(O(n)\)-bit encoding, so this construction contains no exponential
table.

Equation \eqref{eq-lpn-decoder-center-admissibility} is the defining identity
for Hamming distance.  If \(\widehat s\ne s\), then
\(D_{\widehat s}(s)>0\), so the potential fails the target-normalization
condition in \cref{def-geom-extraction}.  The centered carrier remains a
charged represented object, but it is outside the admissible target-oriented
record family and has zero target-aligned profile contribution.  This proves
the zero branch in
\cref{eq-lpn-decoder-qualified-geom-contribution}.

Assume \(\widehat s=s\).  Put \(D=D_{\widehat s}\) and \(w=D(x)\).  Direct
substitution in \eqref{eq-lpn-decoder-correction-kernel} gives

\begin{equation}
\label{eq-lpn-decoder-exact-drift}
(P_{\widehat s}-I)D(x)=-\frac{w}{n}\le-\frac1n
\qquad(x\ne s).
\end{equation}

The same-generator drift margin \(1/n\) activates the authority-compatible
local gate.  The symmetric part of the fixed-space split has rate \(1/(2n)\)
on every hypercube edge.  Its Walsh eigenvalue on a character indexed by
\(S\subseteq\{1,\ldots,n\}\) is \(-|S|/n\); hence its gap is \(1/n\) and
the carrier gate \(M_{\Phi_{\widehat s}}\) equals one within this admissible
record.  The carrier gate has no separate target-aligned contribution.  Its
directed current is a positive scalar
multiple of \(-\nabla_{\Phi_{\widehat s}}D\), so the authorized alignment is
one and its gradient-orthogonal residual projection is zero.  Since
\(D^{-1}(0)=\{s\}\), the centered target indicator belongs to the
target-oriented monotone span of \(D\), giving target reach one.

Under the uniform candidate law, \(D\) is binomial
\(\mathrm{Bin}(n,1/2)\), and therefore

\[
\operatorname{Var}_{\mu_I}(D)=\frac n4.
\]

For the symmetric rate \(1/(2n)\), every hypercube edge changes \(D\) by one,
so

\[
\mathcal E_{\Phi_{\widehat s}}(D,D)=\frac14,
\qquad
\rho_{\Phi_{\widehat s}}(D)=\frac1n.
\]

These identities prove
\eqref{eq-lpn-decoder-geometry-exact-gates}; substitution into
\cref{eq-lpn-identity-geom-complete-gate-product} proves
\eqref{eq-lpn-decoder-geometry-exact-visibility}.

The fibers of \(q_{\widehat s}\) are the represented Hamming spheres about
\(\widehat s\).  Equation
\eqref{eq-lpn-decoder-correction-kernel} shows that the transition law of
the level depends only on its current value \(w\), proving
\eqref{eq-lpn-decoder-level-quotient-chain} and exact lumpability.  The
target is the singleton level \(0\), while the positive levels strictly
refine the verifier partition.  Canonical pushforward gives comparison
factor one by
\cref{prop-canonical-exact-quotient-unit-transfer}, and the target indicator
is quotient-measurable, so the projection ratio is one.

An initially incorrect coordinate remains unselected after \(H\) independent
coordinate choices with probability at most \(e^{-H/n}\).  A union bound over
at most \(n\) such coordinates proves
\eqref{eq-lpn-decoder-correction-arrival}.  Uniformization with
\(\lambda=1/H\) then gives

\[
A_T^{q_{\widehat s}}(1/H)
\ge(1+1/H)^{-H}(1-\delta)
\ge e^{-1}(1-\delta).
\]

The quotient-utility normalization of
\cref{def-paper1-abs-usefulness-gate} and the unit projection ratio yield
\eqref{eq-lpn-decoder-level-quotient-visibility}.

Finally, the good-presentation probability in
\cref{eq-lpn-parameter-explicit-verifier-geometry-probability} is
\(1-\zeta_{n,m,\eta,\tau}\).  On that event, decoder success makes both
records qualified and gives the two pointwise contributions just computed.
The independent seed is a valid product lift, and unsuccessful or
unqualified records contribute zero in the displayed audit.  Hence the
probability of a qualified successful record
is at least
\((p_{n,\eta}(\mathcal D)-\zeta_{n,m,\eta,\tau})_+\).  Averaging the two
pointwise bounds proves
\cref{eq-lpn-decoder-success-charged-to-geom,eq-lpn-decoder-success-charged-to-abs}.
The Geom term is exactly the indicator product in
\cref{eq-lpn-decoder-qualified-geom-contribution} restricted to the verifier
good event.  Its expectation is
\cref{eq-lpn-decoder-qualified-family-geom-contribution}.  The bounds in
\cref{eq-lpn-decoder-qualified-family-geom-interval} follow from
\(\Pr(A\cap B)\ge(\Pr(A)-\Pr(B^c))_+\) and
\(\Pr(A\cap B)\le\Pr(A)\).

\end{proof}

\Cref{fig-lpn-decoder-postoutput-accounting} distinguishes the charged
post-output constructions from the prospective computation that produced the
decoder output.

\begin{figure}[tbp]
\centering
\resizebox{\linewidth}{!}{%
\begin{tikzpicture}[fw diagram,node distance=8mm and 10mm]
  \node[fw source,text width=2.4cm] (decoder) {represented decoder\\cost \(B_{\mathcal D}\)};
  \node[fw provenance,text width=2.5cm,right=of decoder] (store) {stored output\\\(\widehat s\) and seed};
  \node[fw use,text width=2.7cm,right=of store] (gate)
    {target\\admissibility\\\(D_{\widehat s}(s)=0\iff\widehat s=s\)};
  \node[fw geom,text width=2.6cm,above right=8mm and 12mm of gate] (geom)
    {ambient Geom branch\\visibility \(n/(n+1)\)};
  \node[fw box,text width=2.6cm,below right=8mm and 12mm of gate] (quot)
    {level quotient branch\\arrival at least \((1-\delta)/(eH)\)};
  \node[fw terminal,text width=2.8cm,right=22mm of gate] (audit)
    {retrospective accounting\\weighted by decoder success};
  \draw[fw arrow] (decoder)--(store);
  \draw[fw arrow] (store)--(gate);
  \draw[fw arrow] (gate)--(geom);
  \draw[fw arrow] (gate)--(quot);
  \draw[fw arrow] (geom)--(audit);
  \draw[fw arrow] (quot)--(audit);
  \draw[fw dashed arrow] (decoder) to[bend right=24]
    node[below,font=\scriptsize] {ancestor cost retained} (audit);
\end{tikzpicture}%
}
\caption{A decoder can cheaply build a target-centered carrier only after it
has produced \(\widehat s\).  The target-admissibility gate activates both
the ambient geometry and exact level quotient precisely on decoder success.
Their visibility therefore audits that success retrospectively and cannot be
used as a prospective source for the earlier decoding transitions.}
\label{fig-lpn-decoder-postoutput-accounting}
\end{figure}

\begin{corollary}[Quantitative post-output decoder audit]
\label{cor-lpn-decoder-success-profile-charge}

Under the hypotheses and notation of
\cref{thm-lpn-charged-decoder-geometry-level-quotient}, every represented
decoder satisfies

\begin{align}
\label{eq-lpn-decoder-success-upper-by-geom}
p_{n,\eta}(\mathcal D)
&\le
\zeta_{n,m,\eta,\tau}
+\Pr\!\left\{V_I^{-1}(1)=\{s\},\ \widehat s=s\right\},\\
\label{eq-lpn-decoder-success-upper-by-abs}
p_{n,\eta}(\mathcal D)
&\le
\zeta_{n,m,\eta,\tau}
+\frac{eH}{1-\delta}
\mathbb E\!\left[
\mathbf1_{\{V_I^{-1}(1)=\{s\},\ \widehat s=s\}}
V_I(q_{\widehat s})
\right].
\end{align}

Both use the single complete cost
\eqref{eq-lpn-decoder-geometry-complete-cost} and the same authenticated
decoder ancestor.  They are post-output accounting identities, not bounds by
a prospective source profile.

\end{corollary}

\begin{proof}

Use
\(p_{n,\eta}(\mathcal D)\le
\zeta_{n,m,\eta,\tau}+
\Pr\{V_I^{-1}(1)=\{s\},\widehat s=s\}\), then apply
\cref{eq-lpn-decoder-success-charged-to-abs} for the second inequality.

\end{proof}

\subsection{Selected-Current Hodge Classification}
\label{subsec-lpn-selected-current-hodge}

\begin{theorem}[Selected-current Hodge classification in LPN]
\label{thm-lpn-off-target-gradual-source-residual-dichotomy}

Work on the LPN verifier good event and fix a saturated budget \(B\).  Let
\[
\mathfrak g=(\mathfrak N,\Phi,D,J,\mathfrak K)
\in\mathcal C_{X,b}^{\mathrm{LPN}}(B)
\]
be a temporally admissible identity-supported record.  Apply the channel
projection to its actual selected current after the terminal-contact
coordinate has been accounted by
\cref{thm-lpn-no-free-direct-target-perturbation}.  In the represented
weighted edge space, the selected current has the orthogonal decomposition
\begin{equation}
\label{eq-lpn-off-target-current-caus-residual-split}
J_{\mathfrak K}
=
\Pi_{\operatorname{im}\nabla_\Phi}J_{\mathfrak K}
+
\bigl(I-\Pi_{\operatorname{im}\nabla_\Phi}\bigr)
J_{\mathfrak K}.
\end{equation}

This decomposition is internal to the admitted Geom record.  Neither
summand supplies a Geom source: the gradient component is derived \(\Caus\),
and the orthogonal component is \(\Res\).  In the absence of the underlying
target-qualified Geom record, both have zero prospective-source value.
The first summand is the \(\Caus\) gradient projection of the selected
operator current.  It contributes prospectively precisely through a
represented target-normalized potential, the same-generator authority
record, and the complete ambient visibility
\begin{equation}
\label{eq-lpn-off-target-actionable-direction-visibility}
M_\Phi\,
\chi_\Phi^{\mathrm{auth}}(D)\,
\operatorname{Align}_{\Caus}^{\mathrm{auth}}(\mathfrak K,D)\,
R_{\{s\}}(D)\,
\frac1{1+\rho_\Phi(D)}.
\end{equation}
This contribution is charged to the existing
\(G_{X,b,\mathfrak M,n}^{\mathrm{src,sat}}(B)\) restriction of the ambient
coordinate.  The directed projection
creates no additional source coordinate.

The second summand in
\cref{eq-lpn-off-target-current-caus-residual-split} is orthogonal to every
represented gradient in the budget-\(B\) observable or lifted algebra.
Consequently, for every represented target-normalized potential \(D'\),
\begin{equation}
\label{eq-lpn-off-target-residual-zero-gradient-pairing}
\left\langle
\bigl(I-\Pi_{\operatorname{im}\nabla_\Phi}\bigr)
J_{\mathfrak K},
-\nabla_\Phi D'
\right\rangle_E
=0.
\end{equation}
It therefore has zero authorized Caus alignment and belongs to the
post-projection \(\Res\) component.  If a family-uniform affordable
potential, current comparison, quotient, lift, or terminal readout later
gives this component a channel-qualified target-aligned representative, the
least-cost saturated reclassification moves that representative into the
corresponding exposed channel before the post-saturation residual is formed.

Quantitatively, if the selected-current gradient contribution of a record exceeds
\(\varepsilon\in(0,1)\), then
\begin{gather}
\label{eq-lpn-off-target-gradual-effective-gates}
M_\Phi=1,\qquad
\chi_\Phi^{\mathrm{auth}}(D)=1,\qquad
\operatorname{Align}_{\Caus}^{\mathrm{auth}}(\mathfrak K,D)>\varepsilon,
\qquad
R_{\{s\}}(D)>\varepsilon,\qquad
\frac1{1+\rho_\Phi(D)}>\varepsilon,\\
\label{eq-lpn-off-target-gradual-effective-pairing}
\left\langle J_{\mathfrak K},-\nabla_\Phi D\right\rangle_E
>
\sqrt{\varepsilon}\,
\|J_{\mathfrak K}\|_E\|\nabla_\Phi D\|_E.
\end{gather}
Thus an inverse-polynomial gradual contribution is already an
inverse-polynomial Caus-qualified Geom record in the single saturated ambient
profile.  A direction specified only by the semantic statement that it
points toward \(s\), without the represented potential values, local
differences, selected current, and authority comparison available before the
transition, activates neither
\(\chi_\Phi^{\mathrm{auth}}\) nor
\(\operatorname{Align}_{\Caus}^{\mathrm{auth}}\) and has zero structural
visibility at budget \(B\).

Together with
\cref{thm-lpn-no-free-direct-target-perturbation}, this exhausts the
identity-supported transition record:
direct target contact is neutral or charged through an existing source;
the selected-current gradient projection is the existing Caus-qualified Geom
coordinate; and its gradient-orthogonal complement is \(\Res\).

\end{theorem}

\begin{proof}

\Cref{thm-channel-exhaustiveness} applies the represented Hodge projection to
the actual selected edge current.  Its gradient subspace is
\(\operatorname{im}\nabla_\Phi\); orthogonal projection gives
\cref{eq-lpn-off-target-current-caus-residual-split} and
\cref{eq-lpn-off-target-residual-zero-gradient-pairing}.  By
\cref{def-authorized-caus-use-current,prop-derived-channels-create-no-structure},
the gradient summand is the derived \(\Caus\) use of a represented potential
and conductance.  Its dynamic use is admitted only by the same-generator
record of \cref{def-authority-compatible-caus-alignment}.  Substitution in
\cref{def-geom-extraction} gives
\cref{eq-lpn-off-target-actionable-direction-visibility}, so its complete
cost and target visibility are already included in
\(G_{X,b,\mathfrak M,n}^{\mathrm{src,sat}}(B)\).

The orthogonal summand has zero inner product with every gradient in the
represented subspace.  Its numerator in
\cref{def-quantitative-caus-alignment} is therefore zero.  The channel
decomposition assigns it to \(\Res\), and
\cref{prop-family-uniform-residual-reclassification,thm-no-uniform-residual-correlation,thm-irreducible-band-closure}
gives the stated least-cost reclassification property and excludes a
separate affordable target-aligned residual source.

If the visibility in
\cref{eq-lpn-off-target-actionable-direction-visibility} exceeds
\(\varepsilon\), all five factors lie in \([0,1]\), with the first two
binary.  Each quantitative factor therefore exceeds \(\varepsilon\), proving
\cref{eq-lpn-off-target-gradual-effective-gates}.  The alignment formula in
\cref{def-quantitative-caus-alignment} then gives
\cref{eq-lpn-off-target-gradual-effective-pairing}.  When the potential,
current identity, polarity, or authority comparison is absent,
\cref{def-authority-compatible-caus-alignment} sets the corresponding
authorized factor to zero.  Finally,
\cref{thm-lpn-no-free-direct-target-perturbation} exhausts the removed
terminal-contact part.

\end{proof}

\subsection{Geometric-Source Exclusion}
\label{subsec-lpn-geometric-source-exclusion}

\begin{lemma}[Absence of an independent geometric component in
LPN]\label{lem-lpn-no-independent-geom-source}
\label{thm-lpn-no-independent-geom-final}

For random binary LPN, no primitive local \(\Geom\) source has
inverse-polynomial target visibility. A generated construction whose
derivation uses the public residual map retains that map as an authenticated
ancestor. Its fiber map is a quotient candidate; it contributes a qualified
\(\Abs\) structural source only when its complete record satisfies the
framework's target-utility, compatibility, affordability, and temporal-source
conditions. Fiber formation alone supplies no \(\Abs\) source classification.
Every selected current in the ambient Geom mode is therefore either the derived
\(\Caus\) gradient projection of that same represented \(\Geom\) record
or a gradient-orthogonal residual term with zero authorized target alignment.
There is no additional source represented only by the semantic assertion
that a transition points toward the target.

\end{lemma}

\begin{proof}\phantomsection\label{proof-lpn-no-independent-geom-final}

Apply the
\hyperref[thm-valid-geom-source-characterization]{Characterization of
transfer-class geometric certificates theorem}. The
\hyperref[def-lpn-task-object]{LPN inversion-task definition} leaves its
relation slot undeveloped and supplies no primitive tuple consisting of a
target-oriented pointwise potential, exposed symmetric conductance,
positive killed gap, and target reach.  The complete
residual-independent ambient branch has target visibility at most \(2^{-n}\), and
its prospective selector advantage is zero, by
\cref{lem-lpn-syndrome-free-ambient-geom-bound}.  The coefficient-complete
exact decomposition
\cref{cor-lpn-coefficient-complete-exact-geom-decomposition} assigns every
remaining record one positive support--signature normal form.  Its complete
ambient value is evaluated, without replacing the selected record by an
ancestral restriction, in
\cref{thm-lpn-complete-geom-structure-exhaustion,cor-lpn-complete-ambient-geom-source-seal}.
A construction using target-oriented information from
\(b\) obtains it through a represented public derivation having the candidate
residual map among its authenticated ancestors. The
\hyperref[def-budgeted-derivability-closure]{Budgeted derivability closure}
retains that ancestry and its complete cost.  At a declared polynomial
saturated ceiling, the construction contributes only when its family-uniform
composite record satisfies the cost bound in
\cref{eq-affordable-composition-cost} and the resulting geometric record
passes the admission, temporal-source, and target-visibility conditions in
\cref{thm-represented-current-admission-no-semantic-support,def-pre-hit-temporally-admissible-source,def-geom-extraction}.
The possible information-theoretic dependence of the joint public input
\((A,b)\) on \(s\) supplies none of these conditions: it is an analytic
property of the input law, not a represented polynomial-cost readout,
potential, kernel, current, or selector.  In particular, a target projection
of the joint input enters the saturated source profile only through an
explicit represented postprocessing whose full ancestry and evaluation cost
fit that ceiling.  Conditional on such an admitted record, the
\hyperref[thm-affordable-composition-no-creation]{Affordable composition and
target-projection monotonicity theorem} assigns a deterministic postprocessing
no target projection beyond its represented joint input.  Without such a
record, the qualified-record convention assigns zero source contribution;
information-theoretic correlation alone is not an affordable structural
source.

When the construction exposes a fiber map \(q\), the
\hyperref[thm-abs-fiber-algebra]{Fiber algebra of a represented quotient
theorem} describes its quotient-measurable algebra. Structural-source status
is then decided by the existing gates: \(q\) contributes to the qualified
\(\Abs\) profile precisely when it satisfies
\cref{def-paper1-abs-usefulness-gate,def-abs-extraction} and, when invoked as a
prospective source, \cref{def-pre-hit-temporally-admissible-source}. A fiber
map failing those gates remains a represented postprocessing and is classified
only through the channel predicates it satisfies under
\cref{thm-channel-exhaustiveness}; it is not reclassified as an \(\Abs\)
source. Hence every
\(b\)-dependent construction retains its charged derivation provenance, while
none becomes an independent primitive \(\Geom\) source or a qualified
\(\Abs\) source solely by factoring through the residual map.
Its complete ambient geometric record is evaluated, without an
additional source type, by
\cref{lem-lpn-exhaustive-identity-perturbation-transport,thm-lpn-identity-geom-gate-evaluation}.
When such a derivation first produces a represented estimate
\(\widehat s\) and uses it to center a local target-oriented geometry,
\cref{thm-lpn-charged-decoder-geometry-level-quotient} retains the complete
decoder cost and evaluates every ambient gate in the post-output record.  The
same potential exposes the exact nonidentity level quotient
\(q_{\widehat s}(x)=d_H(x,\widehat s)\).  The corresponding calculations in
\cref{eq-lpn-decoder-success-charged-to-geom,%
eq-lpn-decoder-success-charged-to-abs} are accounting audits and supply no
prospective-source credit for the decoding transitions that produced
\(\widehat s\).
The reversible conductance of this or any other LPN carrier supplies no
\(\mathsf{Sym}\) tag.  On the random-code-rigidity event,
\cref{cor-lpn-symmetry-provenance-geom-closure} removes every independent
authenticated \(\mathsf{Sym}\) source and assigns each \(b\)-dependent
symmetry-breaking construction to its actual charged provenance cells.
When such a construction is learned, its target-qualified current is bounded
jointly by its generated target projection, authorized gradient component,
and decoder qualification probability through
\cref{cor-qualified-learned-current-accountability}.
The transition record supplies no semantic-direction exception.
\Cref{thm-lpn-no-free-direct-target-perturbation} exhausts its
terminal-contact part.  For the remaining selected current,
\cref{thm-lpn-off-target-gradual-source-residual-dichotomy} identifies every
represented actionable direction with the already charged Caus projection of
the complete Geom record; the orthogonal current has zero pairing with every
represented gradient and remains \(\Res\).  Since \(\Caus\) is derived by
\cref{prop-derived-channels-create-no-structure}, it contributes no
independent source beyond that Geom record.  A direction lacking the
represented potential, local differences, selected current, and
same-generator authority comparison has zero authorized alignment by
\cref{def-authority-compatible-caus-alignment}.

\end{proof}

\begin{lemma}[Absence of compensated quotient--geometric sources in LPN]
\label{lem-lpn-compensated-quotient-geom-seal}
\label{thm-lpn-coupled-source-seal-final}

Work on the full-rank and random-code-rigidity event of
\cref{thm-lpn-random-code-rigidity-final}.  At every budget in the declared
polynomial saturated ceiling,
\begin{equation}
\label{eq-lpn-compensated-source-vanishes}
G_{Q,\mathrm{rev},n}^{\mathrm{src,sat}}(B)
=
G_{Q,\mathrm{nr},n}^{\mathrm{src,sat}}(B)
=
G_{Q,\mathrm{ql},n}^{\mathrm{src,sat}}(B)
=0.
\end{equation}

\end{lemma}

\begin{proof}\phantomsection\label{proof-lpn-coupled-source-seal-final}

By
\cref{thm-complete-generic-source-theory,%
def-master-saturated-quotient-geometry-profile}, a record in either
compensated mode contains a represented nonidentity quotient \(q\), its
complete five-family derivation, a positive target-qualification factor, and
a reversible or general-resolvent stability factor.  The stability factor
compares the quotient and ambient dynamics; it creates neither quotient
fibers nor target alignment.

Apply the exhaustive five-family charge decomposition
\cref{cor-five-family-abs-envelope} to the represented quotient map and its
qualification record.  The \(\mathsf{Mult}\) and \(\mathsf{Ord}\) charges
vanish by
\cref{thm-lpn-no-multiplicative-source,thm-lpn-order-provenance}.
Filtration and valid-Lift vertices add zero incremental charge and return
inherited charge to an earlier represented ancestor by
\cref{thm-lpn-filtration-provenance,lem-lpn-exact-lift-composition}.
On the rigidity event, authenticated \(\mathsf{Sym}\) ancestry is
identity-equivalent by
\cref{thm-lpn-random-code-rigidity-final,cor-lpn-symmetry-provenance-geom-closure}.
The remaining \(\mathsf{Add}\) alternative is identity-equivalent by
\cref{thm-lpn-additive-quotient-identity}.  Mixed derivations retain these
same contextual charges and introduce no sixth alternative.

Thus every target-qualified quotient support is identity-equivalent or has
zero prospective source charge.  Identity support belongs to the ambient
mode \(\mathsf X\), whereas both compensated modes require
\(q\ne\mathrm{id}\).  Hence each compensated source family is empty.
Taking the two suprema and their maximum proves
\cref{eq-lpn-compensated-source-vanishes}.

\end{proof}
\section{Five-Family Provenance and Algebraic Source Classes}
\label{sec-lpn-five-family-algebraic-source-classes}
\label{subsec-lpn-five-family-abstraction-provenance}

The universal five-family normal form now specializes the exact abstraction
channel to the public LPN presentation.

\begin{corollary}[Five-family Abs provenance for LPN]
\label{cor-lpn-five-family-abs-provenance}

Every exact quotient evaluator admitted to the budgeted LPN quotient ledger
has additive, multiplicative, order, symmetry, or
filtration provenance in the sense of
\hyperref[def-five-family-abs-provenance]{five-family Abs evaluator
provenance}. Mixed derivations retain all applicable provenance tags and
their contextual joint charge. The complete LPN quotient profile obeys
the aggregate and individual bounds of
\hyperref[cor-five-family-abs-envelope]{the exhaustive five-family Abs
envelope}.

This is the LPN instance of the framework decomposition. The five
families are the provenance families of the universal normal form, not
an LPN-specific catalogue of algorithms or constructions.

\end{corollary}

\begin{proof}\phantomsection\label{proof-lpn-five-family-abs-provenance}

Apply
\hyperref[thm-five-family-abs-provenance]{the universal five-family Abs
provenance normal form} and
\hyperref[cor-five-family-abs-envelope]{its saturated-profile envelope}
to the budget-complete LPN presentation.

\end{proof}


\subsection{Additive and Symmetry Sources}
\label{subsec-lpn-additive-symmetry-sources}

Target alignment reduces both represented additive and authenticated symmetry
quotients to identity support.

\begin{lemma}[Additive LPN quotients are identity-equivalent]
\label{lem-lpn-additive-quotient-identity}
\label{thm-lpn-additive-quotient-identity}

Let \(q:X_I\to Q\) be a budget-admissible quotient map occurring either in
an exact \(\Abs\) record or as the support quotient of a compensated Geom
record.  Suppose its terminal quotient map has purely \(\mathsf{Add}\)
state-dependent provenance in the
\hyperref[thm-five-family-abs-provenance]{universal five-family Abs
provenance normal form}. Thus, after an affordable relabeling of its
image, its terminal quotient map is affine in the candidate variable:

\[
q(x)=L_Ix+c_I,
\qquad x\in X_I=\mathbb F_2^n.
\]

Every such \(q\) satisfying the
\hyperref[def-paper1-abs-usefulness-gate]{quotient-utility condition}
for the LPN target \(T_I=\{s\}\) obeys

\[
\ker L_I=\{0\}.
\]

Consequently \(q\) is injective and its nonempty-fiber partition is the
identity partition up to the represented affine relabeling. Hence the LPN
presentation has no proper target-qualified quotient support in the
\(\mathsf{Add}\) channel, whether exact or compensated.

\end{lemma}

\begin{proof}\phantomsection\label{proof-lpn-additive-quotient-identity}

The terminal-event clause of the
\hyperref[def-paper1-abs-usefulness-gate]{quotient-utility condition}
requires \(T_I\) to be a union of quotient fibers. Since
\(T_I=\{s\}\), the nonempty fiber containing \(s\) must equal
\(\{s\}\). Affinity gives

\[
q^{-1}(q(s))=s+\ker L_I.
\]

Therefore \(s+\ker L_I=\{s\}\), so \(\ker L_I=\{0\}\). The map
\(q\) is injective, and every nonempty fiber is a singleton. Its
partition is therefore the identity partition. Any proper additive
quotient has a nonzero kernel, so its target-containing fiber strictly
contains \(s\) and fails the terminal-event clause.  The proof uses the
terminal-event gate and the affine fiber identity, not exact intertwining.
It therefore applies to the exact and compensated quotient-support modes
selected by the framework decomposition.

\end{proof}

\begin{corollary}[Coefficient-complete exact Geom decomposition for LPN]
\label{cor-lpn-coefficient-complete-exact-geom-decomposition}

Every temporally admissible dynamically qualified LPN Geom source belongs to
exactly one of the 124 cells
\[
(\sigma,S)
\in
\mathfrak S_{\Geom}^{4}\times\mathfrak P_{\Geom,5}
\]
of \cref{thm-canonical-exact-signature-geom-decomposition}.  A represented
calculation from \((A,b,\eta)\) that uses a product, equality or fiber
operation, comparison or selection, authenticated group action, or stateful
iteration occurs instead in an exact signature containing respectively
\(\mathsf{Mult}\), \(\mathsf{Ord}\), \(\mathsf{Sym}\), or
\(\mathsf{Filt}\).  Its complete coefficient derivation remains in that
cell's active frontier and saturated cost.

The provenance tag \(\mathsf{Add}\) records ancestry within the existing
ambient Geom mode.  Every record using an additional constructor retains
that constructor in its exact provenance signature.  The complete ambient
record is evaluated by
\cref{thm-lpn-complete-geom-structure-exhaustion,cor-lpn-complete-ambient-geom-source-seal}.
No Hodge-orthogonal current occurs in an ambient Geom cell.

\end{corollary}

\begin{proof}

Apply \cref{thm-canonical-exact-signature-geom-decomposition} to the complete
LPN active record.  Evaluator normalization assigns every listed non-Add
operation its corresponding tag, including when the operation computes an
instance-dependent coefficient.  The complete-record evaluation is
\cref{thm-lpn-complete-geom-structure-exhaustion}; the current assignment is
the one in \cref{thm-canonical-exact-signature-geom-decomposition}.

\end{proof}

\begin{lemma}[Authenticated LPN symmetries induce code automorphisms]
\label{lem-lpn-symmetry-induces-code-automorphism}

Fix \(0<\eta<1/2\).  Let \(G\) be a represented group action whose
authenticated LPN action record occurs either in a budget-admissible exact
quotient or in the complete active ancestral set of a dynamically qualified Geom
record.  For every \(g\in G\), the authenticated action has the form

\[
g(x)=M_gx+h_g,
\qquad
P_g\in\operatorname{Sym}(m),
\]

and residual equivariance is equivalent to

\begin{equation}
\label{eq-lpn-symmetry-equivariance}
AM_g=P_gA,
\qquad
Ah_g+b=P_gb.
\end{equation}

Consequently

\[
P_g\mathcal C_A=\mathcal C_A,
\]

so \(P_g\) is a permutation automorphism of the public code.
In particular, this conclusion applies to every exact orbit quotient
\(q_G:X_I\to X_I/G\) with purely \(\mathsf{Sym}\) terminal provenance in
the universal five-family normal form.

\end{lemma}

\begin{proof}\phantomsection
\label{proof-lpn-symmetry-induces-code-automorphism}

The \(\mathsf{Sym}\) authentication clause requires the represented
action to preserve the active public LPN interface. Preservation of the
candidate affine interface gives \(g(x)=M_gx+h_g\) with \(M_g\)
invertible. Preservation of the Bernoulli product and Hamming sample
interface gives a coordinate permutation \(P_g\): Hamming isometries
may also flip coordinates, but such flips do not preserve
\(\operatorname{Bernoulli}(\eta)^{\otimes m}\) when
\(0<\eta<1/2\). Equivariance of the public residual primitive is

\[
r_I(gx)=P_gr_I(x)
\qquad(x\in\mathbb F_2^n).
\]

Substituting \(r_I(x)=Ax+b\) and comparing the linear and constant
terms gives \eqref{eq-lpn-symmetry-equivariance}. Since \(M_g\) is
invertible,

\[
P_g\mathcal C_A
=P_g\operatorname{Im}(A)
=\operatorname{Im}(AM_g)
=\operatorname{Im}(A),
\]

which proves the claim.

\end{proof}

\begin{lemma}[Symmetry-derived LPN quotients are identity-equivalent]
\label{lem-lpn-random-code-rigidity}
\label{thm-lpn-random-code-rigidity-final}

For uniform \(A\in\mathbb F_2^{m\times n}\) at fixed rate
\(n/m\to R\in(0,1)\), with probability \(1-o(1)\), every
budget-admissible quotient support occurring in an exact or compensated LPN
record with purely \(\mathsf{Sym}\) terminal provenance and satisfying the
\hyperref[def-paper1-abs-usefulness-gate]{quotient-utility condition}
is identity-equivalent. Hence the random LPN presentation has no proper
target-qualified exact or compensated quotient support in the
\(\mathsf{Sym}\) channel with high probability.

\end{lemma}

\begin{proof}\phantomsection\label{proof-lpn-random-code-rigidity-final}

With probability \(1-o(1)\), the matrix \(A\) has full column rank and
the public code \(\mathcal C_A\) has trivial permutation automorphism group:
the rank-failure probability is at most \(2^{n-m}\), and random-code rigidity
gives the second event at fixed rate \(R\in(0,1)\)
\citep{LefmannPhelpsRodl1993}. Work on their intersection and let \(q_G\) be
such a quotient. By
\hyperref[lem-lpn-symmetry-induces-code-automorphism]{the authenticated
LPN symmetry lemma}, every \(g\in G\) supplies a permutation
automorphism \(P_g\) of \(\mathcal C_A\). Triviality of the permutation
automorphism group gives \(P_g=I\). Equation
\eqref{eq-lpn-symmetry-equivariance} then becomes

\[
A(M_g-I)=0,
\qquad
Ah_g=0.
\]

Full column rank gives \(M_g=I\) and \(h_g=0\). Thus every element of
\(G\) acts trivially on \(X_I\). In the
\hyperref[thm-five-family-abs-provenance]{five-family normal form}, a
purely \(\mathsf{Sym}\) terminal map can obtain state dependence only
from its authenticated action constructors. With every such action
trivial, orbit maps and averaging operators are identity-equivalent;
an isotypic output is either the trivial component, hence
identity-equivalent, or state-independent, in which case it fails the
\hyperref[def-paper1-abs-usefulness-gate]{quotient-utility condition}
for the singleton target. Equivariant relabeling does not change
fibers. Therefore every admissible pure \(\mathsf{Sym}\) quotient has
the identity partition.

The intersection event has probability \(1-o(1)\), which proves the stated
random-presentation conclusion.

\end{proof}

For later family-functional bookkeeping, define the rigidity failure
probability
\begin{equation}
\label{eq-lpn-rigidity-failure-probability}
\varepsilon_{\mathrm{rig}}(n,m)
:=
\Pr_A\!\left\{
\operatorname{rank}(A)<n
\ \text{or}\
\operatorname{Aut}_{\mathrm{perm}}(\mathcal C_A)\ne\{I\}
\right\}.
\end{equation}
The proof of \cref{thm-lpn-random-code-rigidity-final} gives
\begin{equation}
\label{eq-lpn-rigidity-failure-rate}
\varepsilon_{\mathrm{rig}}(n,m)
\le 2^{n-m}+\varepsilon_{\mathrm{aut}}(n,m),
\qquad
\varepsilon_{\mathrm{aut}}(n,m)=o(1)
\end{equation}
at fixed rate \(n/m\to R\in(0,1)\).  Here
\(\varepsilon_{\mathrm{aut}}\) is the proportion of full-rank binary
\([m,n]\) codes having a nontrivial permutation automorphism, whose
vanishing is the random-code rigidity theorem
\citep{LefmannPhelpsRodl1993}.  This notation records the exceptional
presentation probability; it is not a new structural coordinate.

\begin{corollary}[Closure of symmetry provenance in LPN geometry]
\label{cor-lpn-symmetry-provenance-geom-closure}

On the full-rank and random-code-rigidity event of
\cref{thm-lpn-random-code-rigidity-final}, every authenticated
\(\mathsf{Sym}\) action in the complete active ancestral set of an LPN geometric
record acts trivially on the candidate and residual interfaces.  It therefore
supplies neither an independent \(\mathsf{Sym}\)-derived Geom source nor a
nonidentity \(\mathsf{Sym}\)-derived quotient.

The edge-reversal identity \(c_\Phi(x,z)=c_\Phi(z,x)\) remains the reversible
carrier condition and supplies no \(\mathsf{Sym}\) provenance.  A represented
\(b\)-dependent center, potential, conductance, ordering, or current that
lacks an authenticated equivariant LPN action retains its complete
derivation cost and its actual
\(\mathsf{Add}\), \(\mathsf{Mult}\), \(\mathsf{Ord}\), or
\(\mathsf{Filt}\) tags.  Its target contribution is evaluated in those
existing provenance cells by the full Geom gate product.  When the record is
learned and target qualification entails correct recovery, its admitted
current contribution also obeys
\cref{eq-qualified-learned-current-pointwise-ceiling,eq-qualified-learned-current-family-ceiling}.

\end{corollary}

\begin{proof}

By \cref{lem-lpn-symmetry-induces-code-automorphism}, every authenticated
LPN action induces a permutation automorphism of the public code.  Random-code
rigidity makes that automorphism the identity, and full column rank then gives
\(M_g=I\) and \(h_g=0\) exactly as in the
\hyperref[proof-lpn-random-code-rigidity-final]{proof of the LPN rigidity
theorem}.  Apply
\cref{thm-reversible-carrier-authenticated-symmetry-separation}: trivial
authenticated actions have singleton orbits and contribute no independent
\(\mathsf{Sym}\) source data, while symmetric conductance alone receives no
\(\mathsf{Sym}\) tag.

The exhaustive complete-ancestry classification
\cref{thm-exhaustive-five-family-geom-source-classification} retains every
other constructor tag and the complete represented cost.  Its ambient target
contribution is the product in \cref{def-geom-extraction}.  For a learned
record whose qualification is correct target recovery,
\cref{cor-qualified-learned-current-accountability} gives the two displayed
current ceilings.

\end{proof}

\subsection{Multiplicative and Order Sources}
\label{subsec-lpn-multiplicative-order-sources}

The multiplicative and order families are evaluated through their contextual
charge on the same public algebra.

\begin{lemma}[Absence of an independent multiplicative source in LPN]
\label{lem-lpn-no-multiplicative-source}
\label{thm-lpn-no-multiplicative-source}

For every binary LPN instance and every admissible budget \(B\), the
multiplicative contextual-charge envelope of the five-family frontier
vanishes:

\[
\mathcal C_{I,\mathfrak F_5}^{(\mathsf{Mult})}(B)=0.
\]

Thus the LPN presentation has no independent \(\mathsf{Mult}\) source,
including in mixed-provenance quotient derivations.

\end{lemma}

\begin{proof}\phantomsection\label{proof-lpn-no-multiplicative-source}

The state-dependent primitives of the
\hyperref[def-lpn-task-object]{LPN task presentation} are the candidate
coordinates and the affine residual map \(r_I(x)=Ax+b\); both have
\(\mathsf{Add}\) provenance. The public data \(A\), \(b\), and
\(\eta\) are state-independent coefficients. Derived residual weights
and comparisons have \(\mathsf{Ord}\) provenance, authenticated
actions have \(\mathsf{Sym}\) provenance, and clocked updates have
\(\mathsf{Filt}\) provenance. Hence the LPN primitive signature
contains no state-dependent primitive with \(\mathsf{Mult}\)
provenance.

Every \(\mathsf{Mult}\)-tagged vertex in the
\hyperref[cor-five-family-abs-envelope]{five-family frontier} is
therefore a derived multiplicative constructor. By
\hyperref[lem-derived-constructor-zero-contextual-charge]{the derived
constructor zero-charge lemma}, each such contextual extension has
charge zero. This remains true when the vertex carries additional
provenance tags, because the charge is attached to the same conditional
algebra extension. Summing over the multiplicative extensions and
taking the supremum in the definition of the aggregate envelope gives
the displayed identity for every \(B\).

\end{proof}

\begin{lemma}[Absence of an independent order source in LPN]
\label{lem-lpn-order-provenance}
\label{thm-lpn-order-provenance}

For every binary LPN instance and every admissible budget \(B\), the
order contextual-charge envelope of the five-family frontier vanishes:

\[
\mathcal C_{I,\mathfrak F_5}^{(\mathsf{Ord})}(B)=0.
\]

Thus the LPN presentation has no independent \(\mathsf{Ord}\) source,
including in mixed-provenance quotient derivations.

\end{lemma}

\begin{proof}\phantomsection\label{proof-lpn-order-provenance}

By the \hyperref[def-lpn-task-object]{LPN task presentation}, the only
state-dependent generators are the candidate coordinates and the
affine residual map \(r_I(x)=Ax+b\), which have \(\mathsf{Add}\)
provenance. The public data \(A\), \(b\), and \(\eta\) are
state-independent coefficients. Every weight, comparison, branch,
selection, metric or quadratic comparison, and ordered reduction in
the five-family frontier is therefore a derived \(\mathsf{Ord}\)
constructor.

The
\hyperref[lem-derived-constructor-zero-contextual-charge]{derived
constructor zero-charge lemma} gives zero contextual charge for every
such extension. This includes the affordable target-aligned Hamming
statistic \(x\mapsto |Ax+b|_1\): it is a represented postprocessing of
the earlier residual observable and does not enlarge its conditional
algebra. The same argument applies when an order vertex carries
additional provenance tags. Summing over the order extensions and
taking the supremum in the aggregate envelope gives the displayed
identity for every \(B\).

\end{proof}

\section{Filtration, Quotient, and Reversible-Coarsening Closure}
\label{sec-lpn-filtration-quotient-coarsening-closure}
\label{subsec-lpn-filtration-source-exhaustion}

Filtration constructors preserve the charged ancestry of their generated
flags, histories, and first-hit records.

\begin{lemma}[Absence of an independent filtration source in LPN]
\label{lem-lpn-filtration-provenance}
\label{thm-lpn-filtration-provenance}

For every binary LPN instance and every admissible budget, filtration
provenance supplies no independent generative source. Under the
\hyperref[thm-filtration-provenance-charging]{Filtration provenance
charging theorem}, a
\(\mathsf{Filt}\)-tagged contextual extension has zero incremental
charge or leaves the charge on an earlier represented ancestor in the
completed accounting DAG.  That ancestor is a prospective source only at a
use satisfying temporal source admissibility and the represented comparison
of \cref{cor-contextual-structural-charge-transfer}.  This
classification includes clocks, memory states, recurrences, outcome and
terminal transcripts, stored verifier flags, and exact first-hit records.

\end{lemma}

\begin{proof}\phantomsection\label{proof-lpn-filtration-provenance}

The \hyperref[def-lpn-task-object]{LPN inversion-task definition}
contains no primitive state-dependent filtration source. Every clock,
memory state, history, outcome record, and stored flag is generated by
a represented rule. The
\hyperref[def-budgeted-saturated-observable-closure]{budgeted saturated
represented-closure definition} retains that rule's complete update,
verifier, state-dependent arguments, transcript, and readout at their
accumulated cost.

Apply the
\hyperref[thm-filtration-provenance-charging]{Filtration provenance
charging theorem}. Its deterministic-constructor case is implemented
by the
\hyperref[lem-derived-constructor-zero-contextual-charge]{Derived
evaluator constructors carry zero contextual charge lemma}; hence a
represented clock, stored copy, memory update, recurrence step, or
history coordinate has zero incremental contextual charge. Its
valid-Lift case is implemented by the
\hyperref[lem-lpn-exact-lift-composition]{Part II exact Lift identity
specialized to LPN}; hence an auxiliary filtration-lift coordinate has
zero incremental charge and its target contribution remains on the
represented base record.

For an operational outcome or terminal transcript,
\cref{thm-five-family-abs-provenance} retains the generating update,
independent outcome coordinates, public verifier, stored flags, and terminal
map in one charged topological derivation. Independent outcome coordinates
are valid-Lift auxiliaries and add zero target charge by
\cref{lem-lpn-exact-lift-composition}. Once those coordinates and the earlier
state are exposed, every update, stored flag, and terminal record is a
represented deterministic postprocessing and adds zero contextual charge by
\cref{lem-derived-constructor-zero-contextual-charge}. The exact contextual
charge identity of \cref{cor-five-family-abs-envelope} therefore leaves every
nonzero charge on an earlier authenticated vertex of the completed
derivation.  It transfers that charge to a prospective source only under the
separate hypotheses of
\cref{cor-contextual-structural-charge-transfer,def-pre-hit-temporally-admissible-source}.

The public LPN verifier and the exact first-hit map are deterministic
postprocessings of the candidate record and completed flags. They add no
source by
\cref{thm-affordable-composition-no-creation,thm-filtration-provenance-charging}.
The flag-based first-hit representation is retrospective by
\cref{cor-first-hit-accounting-not-source}; an independent pre-hit
representation is instead charged and classified by
\cref{cor-no-retrospective-first-hit-source-loophole}. Hence no
\(\mathsf{Filt}\) opcode supplies independent LPN source credit.

\end{proof}

\begin{corollary}[Exhaustion of independent exact \(\Abs\) sources in LPN]
\label{cor-lpn-abs-source-exhaustion}

Fix \(0<\eta<1/2\) and a sample schedule with \(n/m\to R\in(0,1)\).
On the full-rank and random-code-rigidity event, every temporally admissible,
target-qualified, budget-admissible exact \(\Abs\) record is covered by the
following exhaustive alternatives:

\begin{enumerate}[label=\textup{(\roman*)}]
\item its \(\mathsf{Add}\) source is identity-equivalent;
\item its \(\mathsf{Mult}\) contextual charge is zero;
\item its \(\mathsf{Ord}\) contextual charge is zero;
\item its authenticated \(\mathsf{Sym}\) source is identity-equivalent; or
\item its \(\mathsf{Filt}\) vertex has zero incremental charge or leaves
      that charge on an earlier represented source ancestor.
\end{enumerate}

Mixed-provenance records retain every applicable alternative in the same
contextual charge decomposition.  Identity-equivalent records belong to the
ambient-support case of
\cref{thm-complete-generic-source-theory}; they supply no proper exact
\(\Abs\) source.  Consequently the LPN presentation supplies no independent
proper exact \(\Abs\) source at the declared ceiling.  A completed first-hit
quotient remains in the separate accounting profile of
\cref{def-core-refined-first-hit-accounting-profile} and receives no
prospective-source credit by
\cref{cor-first-hit-accounting-not-source}.
Equivalently, on the stated event and at every budget in the declared
ceiling,
\begin{equation}
\label{eq-lpn-exact-abs-source-vanishes}
m_{I_n}^{\mathrm{src,sat}}(B)=0.
\end{equation}

\end{corollary}

\begin{proof}

The exhaustive five-family envelope
\cref{cor-lpn-five-family-abs-provenance,cor-five-family-abs-envelope}
assigns every contextual charge to at least one of
\(\mathsf{Add}\), \(\mathsf{Mult}\), \(\mathsf{Ord}\),
\(\mathsf{Sym}\), and \(\mathsf{Filt}\).  Apply
\cref{thm-lpn-additive-quotient-identity,thm-lpn-no-multiplicative-source,%
thm-lpn-order-provenance,thm-lpn-random-code-rigidity-final,%
thm-lpn-filtration-provenance} in that order.

The \(\mathsf{Mult}\) and \(\mathsf{Ord}\) envelopes vanish exactly.
The \(\mathsf{Filt}\) result returns every nonzero inherited charge to an
earlier represented ancestor; iteration terminates because the retained
derivation is a finite directed acyclic graph.  A valid-Lift auxiliary
coordinate has zero incremental charge by
\cref{lem-lpn-exact-lift-composition}.  Thus a generated carrier creates no
additional exact-\(\Abs\) alternative.  The remaining authenticated
\(\mathsf{Add}\) and \(\mathsf{Sym}\) alternatives are
identity-equivalent by the two cited LPN seals.  The master normal form
\cref{thm-affordable-geom-abs-normal-form} assigns identity support to the
ambient Geom mode rather than to the nonidentity exact-\(\Abs\) mode.
These conclusions exhaust pure and mixed provenance.  Temporal source
admissibility and
\cref{cor-first-hit-accounting-not-source} give the final accounting
statement.

\end{proof}
\begin{corollary}[Provenance restriction of quotient-supported LPN geometry]
\label{cor-lpn-quotient-supported-geom-closure}

Adopt the fixed-positive-noise, fixed-rate setting of
\cref{cor-lpn-abs-source-exhaustion}, work on its good-presentation event,
fix the presented instance \(I_n\) used there, and let
\(B_{\mathrm{poly}}^\sharp\) be the common charged normalization ceiling.
Then the pointwise qualified source profiles satisfy
\begin{align}
\label{eq-lpn-exact-quotient-geom-closed}
G_{Q,\mathrm{ex},n}^{\mathrm{src,sat}}
  \bigl(b_{\mathrm{poly}}(n)\bigr)
&\le
m_{I_n}^{\mathrm{src,sat}}
  \bigl(b_{\mathrm{poly}}(n)\bigr)
 =0,\\
\label{eq-lpn-compensated-quotient-geom-closed}
\max\!\left\{
G_{Q,\mathrm{rev},n}^{\mathrm{src,sat}}
  \bigl(b_{\mathrm{poly}}(n)\bigr),
G_{Q,\mathrm{nr},n}^{\mathrm{src,sat}}
  \bigl(b_{\mathrm{poly}}(n)\bigr)
\right\}
&=0.
\end{align}

Thus every nonidentity quotient-supported Geom mode vanishes.  Identity
support belongs to the ambient mode \(\mathsf X\).

\end{corollary}

\begin{proof}

The prospective pointwise exact-support comparison
\cref{thm-exact-quotient-geometry-visibility-dominated} gives the inequality
in \cref{eq-lpn-exact-quotient-geom-closed};
\cref{eq-lpn-exact-abs-source-vanishes} gives its zero endpoint.

The compensated-source seal
\cref{eq-lpn-compensated-source-vanishes} gives
\cref{eq-lpn-compensated-quotient-geom-closed}.  Valid-Lift and generated
carriers are already included in that seal through their retained ancestry
and zero-incremental-charge identities.

\end{proof}
\begin{corollary}[Sharp reversible geometric-coarsening bounds for LPN]
\label{cor-lpn-sharp-reversible-geometric-coarsening-bound}

Adopt the fixed-positive-noise, fixed-rate setting and good-presentation event
of \cref{cor-lpn-quotient-supported-geom-closure}, and evaluate this
corollary at one presented instance \(I\).  At a saturated budget \(B\), put
\begin{equation}
\label{eq-lpn-reversible-nonidentity-envelope}
\alpha_{\mathrm{LPN},n}^{\mathrm{rev}}(B)
=
\max\!\left\{
G_{\mathrm{LPN},Q,\mathrm{ex},n}^{\mathrm{src,sat}}(B),
G_{\mathrm{LPN},Q,\mathrm{rev},n}^{\mathrm{src,sat}}(B)
\right\}.
\end{equation}
At \(B=b_{\mathrm{poly}}(n)\), it is bounded by
\[
\alpha_{\mathrm{LPN},n}^{\mathrm{rev}}
 \bigl(b_{\mathrm{poly}}(n)\bigr)
=0
\]
by
\cref{eq-lpn-exact-quotient-geom-closed,%
eq-lpn-compensated-quotient-geom-closed}.

Let \(q\ne\mathrm{id}\) be any represented terminal-compatible coarsening of
a reversible LPN geometric record whose complete master-certificate data have
cost at most \(B\).  Let \(K=-L_N\) be its killed normalized operator, let
\(\gamma>0\) be a represented lower spectral bound for \(K\), and retain the
represented discount \(\lambda>0\).  Then \(K=I-P_N\) and self-adjoint
sub-Markov contractivity give
\begin{equation}
\label{eq-lpn-normalized-reversible-operator-window}
\gamma I\preceq K\preceq2I.
\end{equation}
Consequently,
\begin{equation}
\label{eq-lpn-sharp-coarsened-geom-product-bound}
\overline A_q\,
M_{\mathfrak G_q}\,
C_{\mathfrak G_q}\,
\frac{R_{\mathfrak G_q}}{1+\rho_{\mathfrak G_q}}
\le
\alpha_{\mathrm{LPN},n}^{\mathrm{rev}}(B)
\frac{2+\lambda}{\gamma+\lambda}.
\end{equation}

Define the clipped obstruction scale
\begin{equation}
\label{eq-lpn-sharp-coarsening-obstruction-scale}
\beta_{\mathrm{LPN},n}^{\mathrm{rev}}
(B;\gamma,\lambda)
=
\min\!\left\{
1,
\alpha_{\mathrm{LPN},n}^{\mathrm{rev}}(B)
\frac{2+\lambda}{\gamma+\lambda}
\right\}.
\end{equation}
For every \(a,c,r,g\in(0,1]\) with
\begin{equation}
\label{eq-lpn-asymmetric-coarsening-threshold}
acrg>
\beta_{\mathrm{LPN},n}^{\mathrm{rev}}(B;\gamma,\lambda),
\end{equation}
the coarsening cannot satisfy simultaneously
\begin{equation}
\label{eq-lpn-asymmetric-coarsening-gates}
M_{\mathfrak G_q}=1,
\quad
\overline A_q>a,
\quad
C_{\mathfrak G_q}>c,
\quad
R_{\mathfrak G_q}>r,
\quad
\frac1{1+\rho_{\mathfrak G_q}}>g.
\end{equation}
The symmetric specialization is
\begin{equation}
\label{eq-lpn-symmetric-coarsening-obstruction}
M_{\mathfrak G_q}=0
\quad\text{or}\quad
\min\!\left\{
\overline A_q,
C_{\mathfrak G_q},
R_{\mathfrak G_q},
\frac1{1+\rho_{\mathfrak G_q}}
\right\}
\le
\left(
\beta_{\mathrm{LPN},n}^{\mathrm{rev}}(B;\gamma,\lambda)
\right)^{1/4}.
\end{equation}

These bounds combine with the existing geometric evaluations as follows.
Whenever \(C_{\mathfrak G_q}>c\), the local-consistency gate equals one and
the authorized current satisfies
\begin{equation}
\label{eq-lpn-coarsened-current-pairing-threshold}
\left\langle
J_{\mathfrak K_q},-\nabla_{\Phi_Q}D_Q
\right\rangle_E
>
\sqrt c\,
\|J_{\mathfrak K_q}\|_E
\|\nabla_{\Phi_Q}D_Q\|_E.
\end{equation}
The reach and regularity gates obey
\begin{align}
\label{eq-lpn-coarsened-reach-threshold}
R_{\mathfrak G_q}>r
&\Longrightarrow
\operatorname{ProjMass}_{\sigma(D_Q)}(T_Q)>r,\\
\label{eq-lpn-coarsened-regularity-threshold}
\frac1{1+\rho_{\mathfrak G_q}}>g
&\Longrightarrow
\rho_{\mathfrak G_q}<g^{-1}-1.
\end{align}
At the ambient level, an already admitted verifier-potential record satisfies
the exact formula in
\cref{eq-lpn-parameter-explicit-verifier-visibility}; its \(\Caus\) factor is
derived from that same Geom record by
\cref{prop-derived-channels-create-no-structure}.
Thus a coarsening that retains the same inverse-polynomial source scale must
obey
\cref{eq-lpn-coarsened-current-pairing-threshold,eq-lpn-coarsened-reach-threshold,eq-lpn-coarsened-regularity-threshold}
while
\cref{eq-lpn-sharp-coarsened-geom-product-bound} forbids their simultaneous
retention above the obstruction scale.

At the polynomial saturated ceiling,
\begin{equation}
\label{eq-lpn-polynomial-reversible-coarsening-obstruction-bound}
\beta_{\mathrm{LPN},n}^{\mathrm{rev}}
\bigl(b_{\mathrm{poly}}(n);\gamma,\lambda\bigr)
=0,
\end{equation}
and its fourth root is the corresponding symmetric obstruction in
\cref{eq-lpn-symmetric-coarsening-obstruction}.

\end{corollary}

\begin{proof}

For a self-adjoint sub-Markov contraction \(P_N\), its spectrum is contained
in \([-1,1]\).  Hence the spectrum of \(K=I-P_N\) is contained in
\([0,2]\), and the represented killed-gap comparison supplies the lower
bound \(\gamma I\preceq K\).  This proves
\cref{eq-lpn-normalized-reversible-operator-window}.

Apply
\cref{eq-strict-ambient-coarsening-visibility-bound} with
\(\Lambda=2\) and
\(\alpha_n=\alpha_{\mathrm{LPN},n}^{\mathrm{rev}}(B)\).  This proves
\cref{eq-lpn-sharp-coarsened-geom-product-bound}.  The asymmetric exclusion
and symmetric fourth-root obstruction are
\cref{eq-coercive-asymmetric-gate-exclusion,eq-coercive-symmetric-gate-obstruction}
under the same substitution.

Since
\(C_{\mathfrak G_q}
=\chi_{\Phi_Q}^{\mathrm{auth}}(D_Q)
\operatorname{Align}_{\Caus}^{\mathrm{auth}}(\mathfrak K_q,D_Q)\),
the inequality \(C_{\mathfrak G_q}>c\) forces the binary consistency gate
to equal one and the alignment factor to exceed \(c\).
\Cref{def-quantitative-caus-alignment} then gives
\cref{eq-lpn-coarsened-current-pairing-threshold}.  Projection monotonicity
in \cref{thm-generated-algebra-projection-ceiling} gives
\cref{eq-lpn-coarsened-reach-threshold}; the regularity implication is
algebraic.  The ambient verifier assertions are
\cref{cor-lpn-parameter-explicit-verifier-geometry}.

At the polynomial ceiling,
\cref{eq-lpn-exact-quotient-geom-closed,%
eq-lpn-compensated-quotient-geom-closed} sets the pointwise
quotient-supported envelope to zero at the same fixed presentation.  Moreover,
\(0<\gamma\le2\) and \(\lambda>0\) give
\[
\frac{2+\lambda}{\gamma+\lambda}
\le \frac2\gamma,
\]
so the transfer factor is polynomial whenever \(\gamma^{-1}\) is.
Substitution in
\cref{eq-lpn-sharp-coarsening-obstruction-scale} proves
\cref{eq-lpn-polynomial-reversible-coarsening-obstruction-bound}; taking the fourth
root gives the symmetric statement.

\end{proof}

\section{Residual Ambient Geometry and Five-Family Affordability}
\label{sec-lpn-residual-ambient-five-family-affordability}

\begin{corollary}[Exhaustive reversible-carrier classification in LPN]
\label{cor-lpn-no-nonreversible-geometric-escape}

On the good-presentation event, every temporally admissible dynamically
qualified LPN geometric source has a reversible \(\Geom\) carrier.  For an
identity-supported residual-ancestry record, every directed contribution to
target progress is either the authorized \(\Caus\) projection appearing in
the existing ambient visibility
\[
M_\Phi\,
\chi_\Phi^{\mathrm{auth}}(D)\,
\operatorname{Align}_{\Caus}^{\mathrm{auth}}(\mathfrak K,D)\,
R_{\{s\}}(D)\,
\frac1{1+\rho_\Phi(D)},
\]
or the gradient-orthogonal \(\Res\) component with zero authorized alignment.
There is no third, nonreversible ambient-geometry branch.

For a nonidentity quotient, use of a general-resolvent factor
\(S_{q,T}^{\mathrm{nr}}(\lambda)\) remains a quotient-supported
\(\Geom/\Caus\) transfer record and belongs to the general-resolvent family
already bounded in
\cref{eq-lpn-compensated-quotient-geom-closed}.  It supplies no independent
ambient source.  Consequently, after quotient-supported exhaustion, the
identity-supported residual-ancestry profile in
\cref{eq-lpn-geom-reduction-to-b-dependent-identity} is the displayed ambient
LPN coordinate, with all directed target alignment already included in its
\(\Caus\) gate.  That coordinate is charged directly to the saturated ambient
\(\Geom\) profile; absence of an effective nonidentity coarsening does not move
it to \(\Res\) or create an uncharged case.

\end{corollary}

\begin{proof}

Apply
\cref{thm-reversible-geom-exhaustive-directed-classification} to the selected
LPN generator.  Its pairing identities specialize to
\cref{eq-lpn-off-target-current-caus-residual-split,eq-lpn-off-target-residual-zero-gradient-pairing},
and
\cref{thm-lpn-off-target-gradual-source-residual-dichotomy} identifies the
first summand with the displayed ambient visibility and the second with
\(\Res\).  For \(q\ne\mathrm{id}\), the master normal form places a
general-resolvent comparison in the quotient-supported family, and
\cref{cor-lpn-quotient-supported-geom-closure} gives its stated bound.  The
three ambient alternatives are therefore exhausted by the cited channel and
quotient classifications.

\end{proof}

\begin{proposition}[Exhaustive reduction to classified ambient provenance]
\label{prop-lpn-remaining-syndrome-dependent-geom-branch}

Fix a presented instance \(I_n\) on the good-presentation event.  Its
temporally admissible, dynamically qualified geometric source profile is the
disjoint union of the following branches:

\begin{enumerate}[label=\textup{(\roman*)}]
\item identity-supported residual-independent ambient geometry, belonging to
      \(\mathcal C_{X,0}^{\mathrm{LPN}}\);
\item nonidentity quotient-supported geometry, belonging to
      \(\mathcal C_{Q,\mathrm{ex}}\),
      \(\mathcal C_{Q,\mathrm{rev}}\), or
      \(\mathcal C_{Q,\mathrm{nr}}\);
\item identity-supported residual-ancestry ambient geometry, belonging to
      \(\mathcal C_{X,b}^{\mathrm{LPN}}\), with its complete source ancestry
      classified by
      \cref{cor-lpn-residual-ancestry-five-family-geom-classification}.
\end{enumerate}

The family-average evaluation of the first branch satisfies
\eqref{eq-lpn-syndrome-free-ambient-geom-bound}.  The pointwise second branch
satisfies
\cref{eq-lpn-exact-quotient-geom-closed,eq-lpn-compensated-quotient-geom-closed}.
At the polynomial ceiling the complete pointwise geometric profile satisfies

\begin{equation}
\label{eq-lpn-geom-reduction-to-b-dependent-identity}
G_n^{\mathrm{src,sat}}
  \bigl(b_{\mathrm{poly}}(n)\bigr)
=
G_{X,n}^{\mathrm{src,sat}}
  \bigl(b_{\mathrm{poly}}(n)\bigr).
\end{equation}

Branch~\textup{(iii)} has the exact public-ancestry local form of
\cref{thm-lpn-public-ancestry-local-geometry-normal-form}.  At each charged
transition, that theorem writes
its selected sampler and potential difference exactly as

\[
P_I\bigl((x,\xi),(x+u,\xi')\bigr)
=\pi_I(u,\xi'\mid A,x,r_I(x),\xi)
\]

and

\[
\widehat D_I(A,x+u,r_I(x)+Au,\xi')
-\widehat D_I(A,x,r_I(x),\xi).
\]

These are reparameterizations of the same represented record.  The public
presentation supplies \(A\), \(x\), and \(r_I(x)\); the target-cancelling
displacement \(u=x+s\) is absent from its primitive signature.  A positive
mass assigned to that displacement is therefore counted only through the
represented pre-hit sampler and its existing target-alignment comparison,
as required by
\cref{thm-lpn-no-free-direct-target-perturbation,cor-lpn-target-edge-predicate-no-sampler}.

The residual-weight likelihood factors through the terminal-compatible
\((V_I,W_I)\) quotient and is exhausted by the exact, compensated, and
identity-support alternatives of
\cref{thm-lpn-residual-likelihood-quotient-classification}.  Its ambient
identity-support refinement obeys the simultaneous off-target bound
\cref{eq-lpn-uniform-off-target-perturbation-bound}; its only macroscopic
descent is the target-cancelling displacement in
\cref{eq-lpn-uniform-target-perturbation-margin}.  A different represented
potential receives precisely the complete ambient gate product of
\cref{eq-lpn-identity-geom-complete-gate-product}, evaluated on the same
sampler, local differences, current, authority record, reach comparison,
regularity comparison, and complete cost.  Operational products,
comparisons, mixtures, clocks, valid lifts, and derived \(\Caus\) records
retain their signatures and costs while contributing zero independent
target charge by
\cref{lem-derived-constructor-zero-contextual-charge,prop-derived-channels-create-no-structure}.

For a nonlinear full-carrier use of this record,
\cref{cor-nonlinear-nonquotient-dynamical-envelope} evaluates the selected
kernel before the source ancestry is charged backward.  In the
uniform-invariant native case,
\cref{thm-lpn-controlled-target-capacity-drift} gives its exact singleton
capacity, finite-horizon entrance bound, invariant drift ceiling, and the
noise-dependent conversion from residual descent to target contact.  The
canonical public carrier is evaluated in every functional-inequality slot by
\cref{thm-lpn-hypercube-functional-profile}; a learned carrier retains the
operator and population-transfer correction of
\cref{thm-lpn-learned-operator-capacity-transfer}.  These are numerical
outputs of the same ambient \(\Geom/\Caus\) record and preserve its complete
ancestry and cost.

A target-qualified nonidentity spectral or level-set coarsening lies in the
existing quotient-supported modes and satisfies
\cref{eq-lpn-sharp-coarsened-geom-product-bound,eq-lpn-exact-quotient-geom-closed,eq-lpn-compensated-quotient-geom-closed}.
With identity support, the authorized \(\Caus\) projection is the complete
directed part of the ambient record, and its gradient-orthogonal complement
is \(\Res\), by
\cref{cor-lpn-no-nonreversible-geometric-escape,cor-no-post-saturation-residual-source-alternative}.
Thus \cref{eq-lpn-geom-reduction-to-b-dependent-identity} is exactly the
maximum of the displayed budget-admissible, target-qualified structural
records.

\end{proposition}

\begin{proof}

By \cref{thm-affordable-geom-abs-normal-form,thm-exhaustive-geom-normal-form},
every dynamically qualified geometric record has identity support, exact
nonidentity quotient support, reversible compensated support, or general
resolvent-compensated support.  The last three alternatives give branch
\textup{(ii)}.  By
\cref{def-lpn-ambient-geom-ancestry-branches}, the identity-supported family
is the disjoint union of branches \textup{(i)} and \textup{(iii)}.  This
proves exhaustiveness and disjointness.

For branch~\textup{(iii)},
\cref{cor-lpn-residual-ancestry-five-family-geom-classification} applies the
framework source classification to the complete active tuple.  Thus the
binary residual-ancestry partition has no role as an additional structural
classification: every record in that branch already lies in the ambient row
of the exhaustive five-family provenance cover.  The exact substitution and
displacement pushforward of
\cref{thm-lpn-public-ancestry-local-geometry-normal-form} express its local
data using only the declared LPN presentation and preserve the same
qualification gates and complete cost.

The family-average estimate for branch~\textup{(i)} is
\cref{lem-lpn-syndrome-free-ambient-geom-bound}; the pointwise quotient
estimate for branch~\textup{(ii)} is
\cref{cor-lpn-quotient-supported-geom-closure}.  All three branches are
already among the five master modes of the saturated geometric profile.
The quotient-supported modes vanish by
\cref{eq-lpn-exact-quotient-geom-closed,eq-lpn-compensated-source-vanishes};
the remaining identity-supported modes comprise the ambient coordinate.
This proves
\cref{eq-lpn-geom-reduction-to-b-dependent-identity}.

\end{proof}

\begin{theorem}[Quantitative five-family affordability reduction for LPN
geometry]
\label{thm-lpn-quantitative-five-family-geom-affordability}

Fix the LPN family functional \(\mathfrak M_{n,m,\eta}\).  For
\(\sigma\in\mathfrak S_{\Geom}^{4}\) and
\(c\in\mathfrak J_{\Abs}^{5}\), write
\[
G_{\mathrm{LPN},\sigma,c,\mathfrak M,n}^{\mathrm{src,sat}}(B)
\]
for the LPN specialization of the provenance restriction in
\cref{def-active-geom-source-frontier}.  This is a restriction of the
existing prospective Geom profile, evaluated by
\(\mathfrak M_{n,m,\eta}\).  Put
\(\widehat B=\Psi_{\Geom,5}(B)\).  The exhaustive family-functional bound is

\begin{equation}
\label{eq-lpn-quantitative-geom-five-family-reduction}
\begin{aligned}
G_{\mathrm{LPN},\mathfrak M,n}^{\mathrm{src,sat}}(B)
\le \max\Biggl\{&2^{-n},
\max_{c\in\mathfrak J_{\Abs}^{5}}
G_{\mathrm{LPN},\mathsf X,c,\mathfrak M,n}^{b,\mathrm{src,sat}}
  (\widehat B),\\
&G_{\mathrm{LPN},Q,\mathrm{ex},\mathfrak M,n}^{\mathrm{src,sat}}(\widehat B),
G_{\mathrm{LPN},Q,\mathrm{rev},\mathfrak M,n}^{\mathrm{src,sat}}(\widehat B),
G_{\mathrm{LPN},Q,\mathrm{nr},\mathfrak M,n}^{\mathrm{src,sat}}(\widehat B)
\Biggr\}.
\end{aligned}
\end{equation}

Here the superscript \(b\) restricts the ambient support row to
\(\mathcal C_{X,b}^{\mathrm{LPN}}\).  For each provenance tag, this remaining
quantity is exactly

\begin{equation}
\label{eq-lpn-ambient-provenance-cell-exact-supremum}
\begin{aligned}
&G_{\mathrm{LPN},\mathsf X,c,\mathfrak M,n}^{b,\mathrm{src,sat}}(B)\\
&\quad=
\sup_{\substack{
\mathfrak g=(\mathfrak N,\Phi,D,J,\mathfrak K)
\in\mathcal C_{X,b}^{\mathrm{LPN}}(B)\\
c\in\operatorname{Prov}_{\Geom,5}(\mathfrak g)}}
\mathfrak M_{n,m,\eta}\!\left(
I\longmapsto
M_{\Phi}
\chi_{\Phi}^{\mathrm{auth}}(D)
\operatorname{Align}_{\Caus}^{\mathrm{auth}}(\mathfrak K,D)
R_{\{s\}}(D)
\frac{1}{1+\rho_{\Phi}(D)}
\right),
\end{aligned}
\end{equation}

with value zero when the restricted class is empty.  The derivation charged
in this supremum includes the complete canonical source ancestry, exposed
relation and conductance, normalized kernel, potential, current, authority
comparison, gap comparison, reach record, regularity comparison, and every
constructor used to evaluate them.

Let \(0<\varepsilon<1\).  If the left side of
\cref{eq-lpn-quantitative-geom-five-family-reduction} is larger than each of
the three quotient-supported entries and larger than
\(\max\{2^{-n},\varepsilon\}\),
then one tag \(c\in\mathfrak J_{\Abs}^{5}\) and one family-uniform record
\(\mathfrak g\) of complete cost at most \(\widehat B\) have joint gate
product \(Z_{\mathfrak g}\), equal to the product displayed in
\cref{eq-lpn-ambient-provenance-cell-exact-supremum}, such that

\begin{equation}
\label{eq-lpn-provenance-cell-positive-mass}
\begin{aligned}
\mathfrak M_{n,m,\eta}(Z_{\mathfrak g})&>\varepsilon,\\
\Pr_{I\sim\mathfrak I_{n,m,\eta}}
 \left\{Z_{\mathfrak g}(I)>\frac{\varepsilon}{2}\right\}
&>
\frac{\varepsilon}{2-\varepsilon}.
\end{aligned}
\end{equation}

On the event in
\cref{eq-lpn-provenance-cell-positive-mass}, the same represented record
satisfies simultaneously

\begin{equation}
\label{eq-lpn-provenance-cell-polynomial-geom-witness}
\begin{gathered}
c\in\operatorname{Prov}_{\Geom,5}(\mathfrak g),
\qquad
M_{\Phi}=1,
\qquad
\chi_{\Phi}^{\mathrm{auth}}(D)=1,\\
\operatorname{Align}_{\Caus}^{\mathrm{auth}}(\mathfrak K,D)>\frac{\varepsilon}{2},
\quad
R_{\{s\}}(D)>\frac{\varepsilon}{2},
\quad
\frac{1}{1+\rho_{\Phi}(D)}>\frac{\varepsilon}{2},
\end{gathered}
\end{equation}

Conversely, if a family-uniform record of complete cost at most \(B\)
satisfies all five gates in
\cref{eq-lpn-provenance-cell-polynomial-geom-witness}, with the three
quantitative gates exceeding \(a\in(0,1)\), on an instance event of
probability at least \(p\), then its family visibility is larger than
\(pa^3\).  Thus, at a polynomial ceiling, inverse-polynomial family
visibility is equivalent, up to explicit polynomial losses, to constructing
one complete record in one of these five ambient provenance restrictions
whose full gate product is inverse-polynomial on an inverse-polynomial-mass
set of LPN instances.  The criterion is a single saturated profile
evaluation; it contains no per-algorithm quantifier.

Separately, if \(0<\eta<1/2\), \(n/m\to R\in(0,1)\), and \(I_n\) is a
fixed presentation on the full-rank and random-code-rigidity event, then the
pointwise profile at the declared ceiling satisfies

\begin{equation}
\label{eq-lpn-polynomial-geom-affordability-reduction}
G_n^{\mathrm{src,sat}}
  \bigl(b_{\mathrm{poly}}(n)\bigr)
=
G_{X,n}^{\mathrm{src,sat}}
  \bigl(b_{\mathrm{poly}}(n)\bigr).
\end{equation}

This is the direct master-profile reduction obtained from
\cref{eq-lpn-exact-abs-source-vanishes,eq-lpn-exact-quotient-geom-closed,%
eq-lpn-compensated-source-vanishes}.  The ambient profile retains the
temporal-admissibility, target-qualification, and complete-cost gates.

\end{theorem}

\begin{proof}

Apply
\cref{thm-exhaustive-five-family-geom-source-classification} to the complete
LPN prospective Geom profile.  Its support coordinate is one of
\(\mathsf X,\mathsf Q_{\mathrm{ex}},\mathsf Q_{\mathrm{rev}},
\mathsf Q_{\mathrm{nr}}\), and its nonempty source signature contains at
least one tag in \(\mathfrak J_{\Abs}^{5}\).  The factorization and complete
cost are
\cref{eq-geom-active-frontier-factorization,eq-geom-active-frontier-cost}.
Taking the finite maximum gives the LPN specialization of
\cref{eq-geom-five-family-envelope-cover}.

For identity support,
\cref{def-lpn-ambient-geom-ancestry-branches} gives the disjoint
residual-independent and residual-ancestry partition.  The first part is at
most \(2^{-n}\) by
\cref{eq-lpn-syndrome-free-ambient-geom-bound}.  The second part retains its
five-family source signature by
\cref{cor-lpn-residual-ancestry-five-family-geom-classification}.  The three
nonidentity support tags are bounded by their unrestricted exact, reversible,
and general-resolvent quotient-supported Geom coordinates.  These
observations prove
\cref{eq-lpn-quantitative-geom-five-family-reduction}.

The ambient visibility identity
\cref{eq-lpn-identity-geom-complete-gate-product}, restricted to records whose
canonical signature contains \(c\), is exactly
\cref{eq-lpn-ambient-provenance-cell-exact-supremum}.  Every factor lies in
\([0,1]\), and the first two factors are binary.  Because the outer profile
value is larger than \(\varepsilon\), the definition of supremum supplies a
single family-uniform record \(\mathfrak g\) in one ambient provenance cell
with \(\mathbb E Z_{\mathfrak g}>\varepsilon\).  For
\(A=\{Z_{\mathfrak g}>\varepsilon/2\}\),

\[
\mathbb E Z_{\mathfrak g}
\le
\Pr(A)+\frac{\varepsilon}{2}\bigl(1-\Pr(A)\bigr),
\]

so \(\Pr(A)>\varepsilon/(2-\varepsilon)\).  On \(A\), both binary
gates equal one and every other factor exceeds \(\varepsilon/2\).  This
proves
\cref{eq-lpn-provenance-cell-positive-mass,eq-lpn-provenance-cell-polynomial-geom-witness}.
Conversely, on an event of probability at least \(p\), two active binary
gates and three factors larger than \(a\) give
\(Z_{\mathfrak g}>a^3\); hence
\(\mathbb E Z_{\mathfrak g}>pa^3\).
This is precisely the LPN specialization of the saturated ambient-profile
identity, perturbation obstruction cover, and inverse-polynomial witness
criterion in
\cref{thm-saturated-ambient-geom-accountability,thm-ambient-perturbation-characterization-obstruction,cor-saturated-ambient-geom-negligibility}.

Finally, on the fixed rigidity-event presentation,
\cref{eq-lpn-exact-abs-source-vanishes,eq-lpn-exact-quotient-geom-closed,eq-lpn-compensated-source-vanishes}
remove every quotient-supported master mode at the declared ceiling.  The
ambient mode retains the same public constructors, qualification gates, and
complete source records.  The master decomposition therefore proves
\cref{eq-lpn-polynomial-geom-affordability-reduction}.

\end{proof}

\section{First-Hit Accountability and Complete Source Exhaustion}
\label{sec-lpn-first-hit-complete-source-exhaustion}

\subsection{First-hit accounting and source modes}
\label{subsec-lpn-first-hit-accounting-source-modes}

\begin{corollary}[Core-refined first-hit accountability for LPN]
\label{cor-lpn-core-refined-first-hit-accountability}

Adopt the fixed-positive-noise and fixed-rate setting of
\cref{cor-lpn-abs-source-exhaustion}, and work on the intersection with the
good-presentation event of \cref{thm-lpn-public-verifier}.  Fix the declared
polynomial budget envelope \(\mathcal B_C\), put
\(\widehat B_C=\Psi(\mathcal B_C,n)\), and let \(B_C^\dagger\) be the common
charged compilation ceiling for the first-hit, five-family, and core
representations of the same admissible records.  Then

\begin{equation}
\label{eq-lpn-polynomial-selector-accountability}
\operatorname{Sel}_{I,n}^{\mathrm{sat}}(\widehat B_C)
\le
eH_{\widehat B_C}
\operatorname{HitCore}_{I,n}^{\mathrm{sat}}
 (B_C^\dagger).
\end{equation}

while the pointwise saturated success profile satisfies

\begin{equation}
\label{eq-lpn-pointwise-polynomial-success-accountability}
\operatorname{Succ}_{I,n,\eta}(\mathcal B_C)
\le
eH_C
\operatorname{HitCore}_{I,n}^{\mathrm{sat}}(B_C^\dagger).
\end{equation}

These are pointwise accounting bounds in the profile of
\cref{def-core-refined-first-hit-accounting-profile}.  Completed verifier
flags and exact first-hit quotients remain retrospective accounting records;
no displayed inequality promotes them to prospective sources.

\end{corollary}

\begin{proof}

Apply the pointwise form of
\cref{thm-core-refined-prospective-accountability} at
\(\widehat B_C\) for the selector inequality and at \(\mathcal B_C\) for
the success inequality.  This gives the displayed bounds directly.
The direct-contact and temporal classifications remain
\cref{thm-lpn-no-free-direct-target-perturbation,%
def-pre-hit-temporally-admissible-source,%
cor-first-hit-accounting-not-source}; they introduce no term outside the
core-refined first-hit profile.

\end{proof}

\begin{proposition}[Five master-source modes and the separate LPN accounting ledger]
\label{prop-lpn-five-mode-source-accounting-ledger}

Fix the hypotheses and declared ceiling of
\cref{cor-lpn-abs-source-exhaustion}.  The first four
rows below cover the five disjoint master source modes of
\cref{thm-complete-generic-source-theory}, with the two
nonzero-defect modes combined in row~\textup{(iv)}.  Row~\textup{(v)} is the
separate retrospective accounting profile.

\begin{enumerate}[label=\textup{(\roman*)}]

\item In the null-geometry exact-\(\Abs\) mode \(\mathsf A\), the qualified
record has the exact value
\begin{equation}
\label{eq-lpn-five-mode-ledger-pure-abs-value}
Q_I V_{\mathscr C,I}
=
Q_I A_{q,I}
\frac{\|P_q^{\mathrm{amb}}y\|^2}{\|y\|^2},
\qquad 0\le Q_I V_{\mathscr C,I}\le1.
\end{equation}
The five-family exhaustion
\cref{cor-lpn-abs-source-exhaustion} removes this independent proper
exact-\(\Abs\) source mode.  On a generated carrier, filtration and Lift
vertices add zero incremental charge and return every inherited charge to
the earlier represented source ancestor.  Identity-equivalent support is
assigned to the ambient mode \(\mathsf X\).

\item In the identity-supported ambient mode \(\mathsf X\), the
complete record has the exact same-record value
\begin{equation}
\label{eq-lpn-five-mode-ledger-ambient-value}
Q_I M_\Phi\chi_\Phi^{\mathrm{auth}}(D)
\operatorname{Align}_{\Caus}^{\mathrm{auth}}(\mathfrak K,D)
R_{\{s\}}(D)\frac{1}{1+\rho_\Phi(D)},
\end{equation}
and the exact factorwise ceilings are those of
\cref{thm-lpn-identity-geom-gate-evaluation,%
cor-lpn-complete-quantitative-ambient-geom-ledger}; no provenance reindexing
changes this product.  Its complete saturated value is bounded by
\cref{cor-lpn-complete-ambient-geom-source-seal} using
\cref{thm-lpn-residual-ancestry-attention-bound}.

\item In mode \(\mathsf Q_{\mathrm{ex}}\), exact quotient-supported geometry
is bounded by the same exact-\(\Abs\) coordinate:
\begin{equation}
\label{eq-lpn-five-mode-ledger-exact-quotient-value}
G_{Q,\mathrm{ex},n}^{\mathrm{src,sat}}(B)
\le m_{I_n}^{\mathrm{src,sat}}(B).
\end{equation}
Thus the exact support value is controlled by the same \(\mathsf A\)
profile through the displayed comparison and therefore vanishes by
\cref{eq-lpn-exact-abs-source-vanishes}.

\item In modes \(\mathsf Q_{\mathrm{rev}}\) and \(\mathsf Q_{\mathrm{nr}}\),
the master value is respectively
\begin{align}
\label{eq-lpn-five-mode-ledger-reversible-value}
Q_I\overline A_q(1-\vartheta_\lambda^T(q))
M_{\mathfrak G_q}C_{\mathfrak G_q}
\frac{R_{\mathfrak G_q}}{1+\rho_{\mathfrak G_q}},\\
\label{eq-lpn-five-mode-ledger-resolvent-value}
Q_I\overline A_q\max\{0,1-\Delta_\lambda(q)\}
M_{\mathfrak G_q}C_{\mathfrak G_q}
\frac{R_{\mathfrak G_q}}{1+\rho_{\mathfrak G_q}}.
\end{align}
Both nonzero-defect modes vanish by
\cref{eq-lpn-compensated-quotient-geom-closed}.  Its proof retains valid-Lift
and generated-carrier ancestry and cost.

\item The core-refined first-hit coordinate is a separate retrospective
accounting profile.  It has the exact finite maximum
\cref{eq-core-refined-first-hit-maximum}, every cell has the charge identity
\cref{eq-core-refined-first-hit-charge-identity}, and
\begin{equation}
\label{eq-lpn-five-mode-ledger-success-value}
\operatorname{Succ}_{I_n,n,\eta}(\mathcal B_C)
\le eH_C\operatorname{HitCore}_{I_n,n}^{\mathrm{sat}}
(B_C^\dagger).
\end{equation}
By \cref{cor-first-hit-accounting-not-source}, none of these accounting
identities transfers to a prospective source coordinate without the
represented comparison required by
\cref{cor-contextual-structural-charge-transfer}.

\end{enumerate}

Consequently, the exact-\(\Abs\) mode and the exact-quotient-supported Geom
mode vanish on the stated event, as do both compensated Geom modes.  Hence
the prospective master profile retains only the ambient Geom coordinate.
Its structure is exhausted by
\cref{thm-lpn-complete-geom-structure-exhaustion}; by
\cref{thm-lpn-residual-ancestry-attention-bound}, its complete saturated value
is bounded by \(2^{-n}+(1-\eta)^n\) through
\cref{cor-lpn-complete-ambient-geom-source-seal}.  The first-hit profile remains a
separate accounting maximum.  The
operational/core and state--coefficient decompositions are exact reindexings
and introduce no additional numerical factor.

\end{proposition}

\begin{proof}

Item~\textup{(i)} is
\cref{thm-lpn-additive-quotient-identity,%
cor-lpn-abs-source-exhaustion}.  Item~\textup{(ii)} is
\cref{thm-lpn-identity-geom-gate-evaluation,%
thm-lpn-complete-geom-structure-exhaustion,%
cor-lpn-complete-ambient-geom-source-seal}.  Item~\textup{(iii)}
is \cref{thm-exact-quotient-geometry-visibility-dominated,%
eq-lpn-exact-quotient-geom-closed}.  For
item~\textup{(iv)}, the two displayed products are the reversible and
general-resolvent specializations of
\cref{def-master-compensated-quotient-geometry-certificate}; the LPN
vanishing result is
\cref{thm-lpn-coupled-source-seal-final}.  Item~\textup{(v)} is
\cref{thm-core-refined-prospective-accountability}.  Finally,
\cref{thm-operational-core-provenance-separation} preserves the same
\(Q_I V_{\mathscr C,I}\) and complete cost under core reindexing.

\end{proof}

\subsection{Complete prospective value routing}
\label{subsec-lpn-complete-prospective-value-routing}

\begin{corollary}[No uncharged prospective LPN value]
\label{cor-lpn-no-uncharged-prospective-value}

Fix an LPN budget \(B\).  Every family-uniform, temporally admissible
pre-hit selector, controlled value, or valid-Lift arrival record is governed
by \cref{thm-no-uncharged-prospective-value-amplification}.  Hence a
prospective structural bound is obtained either from a represented
same-normalization comparison with an earlier qualified LPN source, or from
the complete record regarded as a directly qualified source in one of the
existing LPN master-profile cells.  In both cases every selected kernel,
potential, current, authority comparison, target interface, clock, and Lift
map remains in the charged record at the common class-preserving ceiling.

There is no separate LPN value-contraction coordinate.  Ancestry equality,
zero incremental \(\mathsf{Mult}\), \(\mathsf{Ord}\), \(\mathsf{Sym}\), or
\(\mathsf{Filt}\) charge, and valid-Lift inheritance do not by themselves
transfer a numerical ancestor bound to a derived output.  A derived output
without an inherited comparison is evaluated, if qualified, as the complete
ambient or quotient-supported source record that produces it.

\end{corollary}

\begin{proof}

Apply \cref{thm-no-uncharged-prospective-value-amplification} to the complete
LPN pre-hit record.  The master-cell assignment is the LPN specialization of
\cref{thm-affordable-geom-abs-normal-form,thm-no-unclassified-geom-source-after-quotient-assignment}.
Source inheritance under Caus, Lift, and boundary processing is
\cref{prop-derived-channels-create-no-structure}; it preserves provenance
and cost but supplies no unrepresented numerical comparison.

\end{proof}

\Cref{fig-no-uncharged-prospective-value} gives the framework-level routing
used here; no LPN-specific replacement coordinate is introduced.
\clearpage

\subsection{Complete Geom exhaustion and saturated evaluation}
\label{subsec-lpn-provenance-exhaustion-evaluated-ancestors}

\begin{theorem}[Complete structural exhaustion of LPN Geom]
\label{thm-lpn-complete-geom-structure-exhaustion}

Fix \(0<\eta<1/2\), work on the full-rank random-code-rigidity event
\(\mathcal G_{\rm rig}\) of
\cref{def-lpn-named-events-restricted-functionals}, and use the declared LPN
family functional restricted to that event.  Let \(B\) lie in the declared
polynomial saturated ceiling.  Every temporally admissible, dynamically
qualified prospective LPN Geom record is classified by a finite application
of the following structural exhaustion.  The primary support tags are
disjoint.  The ancestry rows may overlap and, for a mixed exact provenance
signature, are traversed whenever applicable before the routing terminates at
the earliest target-bearing source.

\begin{center}
\small
\begin{tabular}{@{}
  >{\raggedright\arraybackslash}p{0.24\linewidth}
  >{\raggedright\arraybackslash}p{0.31\linewidth}
  >{\raggedright\arraybackslash}p{0.37\linewidth}@{}}
\toprule
Support or source ancestry & Structural assignment & Evaluating result \\
\midrule
Exact nonidentity support
&
Exact quotient-supported \(\Geom/\Abs\)
&
\cref{eq-lpn-exact-quotient-geom-closed}
\\
Compensated nonidentity support
&
Reversible or general-resolvent quotient-supported \(\Geom\)
&
\cref{eq-lpn-compensated-quotient-geom-closed}
\\
Identity support without residual ancestry
&
Residual-independent ambient \(\Geom\)
&
\cref{lem-lpn-syndrome-free-ambient-geom-bound}
\\
Identity support with residual ancestry
&
Complete selected ambient \(\Geom/\Caus\) record, with every constructor,
comparison, transition, and cost retained
&
\cref{thm-no-unclassified-geom-source-after-quotient-assignment,%
cor-nonlinear-nonquotient-dynamical-envelope,%
thm-lpn-controlled-target-capacity-drift,%
thm-lpn-caus-action-residual-aiming,%
thm-lpn-hypercube-functional-profile,%
thm-lpn-learned-operator-capacity-transfer,%
thm-lpn-complete-controlled-learning-specialization}
\\
\(\mathsf{Mult}\) or \(\mathsf{Ord}\) ancestry
&
No independent source; inherited charge returns to its earlier source
ancestor
&
\cref{thm-lpn-no-multiplicative-source,thm-lpn-order-provenance}
\\
Authenticated \(\mathsf{Sym}\) ancestry
&
Identity-equivalent; no independent source
&
\cref{thm-lpn-random-code-rigidity-final,%
cor-lpn-symmetry-provenance-geom-closure}
\\
\(\mathsf{Filt}\), valid \(\Lift\), \(\Bdry\), or derived \(\Caus\)
&
Same source cell as the earlier represented ancestor
&
\cref{thm-lpn-filtration-provenance,lem-lpn-exact-lift-composition,%
thm-lpn-valid-lift-quantitative-transfer,%
cor-lpn-noise-indexed-caus-lift-interfaces,%
prop-derived-channels-create-no-structure}
\\
Spectral representation
&
Charged location, an existing quotient or ambient cell, terminal readout,
or superpolynomial explicit expansion
&
\cref{thm-lpn-no-independent-diffuse-spectral-source}
\\
Extremal readout without authorized intermediate transport
&
Terminal verifier; zero prospective source activity
&
\cref{cor-extremal-readout-requires-intermediate-transport}
\\
Gradient-orthogonal component
&
\(\Res\); zero authorized target alignment
&
\cref{cor-lpn-no-nonreversible-geometric-escape}
\\
\bottomrule
\end{tabular}
\end{center}

The table is exhaustive, including mixed exact provenance signatures.
Backward source charging on the finite coefficient-complete ancestral DAG
removes every zero-incremental or identity-equivalent decorator while
retaining the complete selected Geom record.  The classification is
preserved after every class-preserving polynomial enlargement of the
ceiling.

\end{theorem}

\begin{proof}

\Cref{thm-no-unclassified-geom-source-after-quotient-assignment} assigns
every qualified prospective Geom record one support tag and one exact,
coefficient-complete five-family signature.  The LPN specialization
\cref{prop-lpn-remaining-syndrome-dependent-geom-branch} refines the ambient
tag into the residual-independent and residual-ancestry branches.  The first
two rows of the table close the three quotient-supported tags.  The third
row is the residual-independent ambient estimate.

Consider a residual-ancestry ambient record.  Apply the
coefficient-complete five-family backward cut of
\cref{cor-lpn-residual-ancestry-five-family-geom-classification}.  The
\(\mathsf{Mult}\), \(\mathsf{Ord}\), authenticated \(\mathsf{Sym}\), and
\(\mathsf{Filt}\) rows of the table contribute no independent target-bearing
source.  Valid lifts, boundary readouts, and derived currents preserve the
same assignment.  Backward iteration terminates because the represented
ancestral DAG is finite.  The complete selected record remains in the
ambient master cell at its full saturated cost.

Before this source assignment is used, every nonlinear full-carrier
descendant is evaluated as its complete selected dynamical record by
\cref{cor-nonlinear-nonquotient-dynamical-envelope}.  Its target-boundary
capacity and global drift obey
\cref{thm-lpn-controlled-target-capacity-drift}; its authorized current and
noise-corrected aiming obey
\cref{thm-lpn-caus-action-residual-aiming,cor-lpn-noise-indexed-caus-lift-interfaces};
its public-carrier functional constants obey
\cref{thm-lpn-hypercube-functional-profile}; its valid lifted forms,
clocked laws, and target gates obey
\cref{thm-lpn-valid-lift-quantitative-transfer}; and its learned operator,
value, policy, and safe-spectrum coordinates obey
\cref{thm-lpn-learned-operator-capacity-transfer,thm-lpn-complete-controlled-learning-specialization}.
These are typed numerical outputs of the same complete record.  They remain
attached to that record while the backward cut identifies its structural
source.  Their explicit finite-parameter interfaces are
\cref{eq-lpn-dynamical-good-event-probability,%
eq-lpn-caus-parameter-event-probability,%
eq-lpn-action-safe-parity-krr-bound,%
eq-lpn-complete-dynamical-clean-risk}.

The complete ambient record retains the capacity, Caus, functional-profile,
Lift, and learning comparisons listed above.  Spectral
notation cannot introduce another terminal source by
\cref{thm-lpn-no-independent-diffuse-spectral-source}; an isolated extremum
is terminal verification by
\cref{cor-extremal-readout-requires-intermediate-transport}; and the
gradient-orthogonal component is \(\Res\) by
\cref{cor-lpn-no-nonreversible-geometric-escape}.  Thus every exact signature
has one source assignment, and every derived row retains the complete record
and the same numerical normalization during that assignment.  This proves
the asserted structural exhaustion.

\end{proof}

\begin{assumption}[Residual-ancestry compact ambient bound for LPN]
\label{thm-lpn-residual-ancestry-attention-bound}

Fix \(0<\eta<1/2\), the declared LPN family functional restricted to the
full-rank and random-code-rigidity event, and a budget \(B\) in the declared
polynomial saturated ceiling or a class-preserving polynomial enlargement.
For the residual-ancestry source family of
\cref{def-lpn-ambient-geom-ancestry-branches}, assume
\begin{equation}
\label{eq-lpn-residual-ancestry-attention-bound}
G_{X,b,\mathfrak M,n}^{\mathrm{src,sat}}(B)
\le (1-\eta)^n.
\end{equation}
The estimate applies to the complete record: score construction,
normalization, active acquisition, posterior update, Caus current, valid
Lift, and deployed transition are all charged at the same ceiling.
By \cref{thm-lpn-exact-constant-noise-spectral-perimeter}, the premise is
already discharged for explicitly materialized located spectra and for every
residual-parity fiber map that fails the quotient gate or is
identity-equivalent.  Its remaining content is the numerical evaluation of
compact, dynamically qualified, non-lumpable residual-ancestry ambient
records, including any polynomial evaluator with an unmaterialized diffuse
Walsh expansion.  The proof-grade empirical discharge in
\cref{thm-lpn-proof-grade-empirical-residual-discharge} lists the source
cover, capacity, attenuation, error, scale, crossover, and finite-prefix
certificates that together prove this premise.

\end{assumption}

This is the only statement in the LPN part declared as an assumption.  The
next theorem gives a finite proof-grade route for discharging each of its
fields; a confidence-only empirical estimate does not replace that route.
\Cref{fig-lpn-proof-grade-empirical-discharge} displays the resulting
finite-prefix and propagated-tail argument.

\begin{theorem}[Proof-grade empirical discharge of the residual-ancestry
bound]
\label{thm-lpn-proof-grade-empirical-residual-discharge}

Fix \(0<\eta<1/2\), a rate schedule \(m=m(n)\), a polynomial ceiling
\(B_n\), and a base size \(n_0\).  After the source reductions in
\cref{thm-lpn-exact-constant-noise-spectral-perimeter,cor-lpn-residual-ancestry-five-family-geom-classification},
let the single remaining source cell be the compact, dynamically qualified,
non-lumpable residual-ancestry ambient cell.  Suppose a proof-grade
scale-certified audit in
\cref{def-scale-certified-attenuating-source-audit} verifies for this cell
\begin{align}
\label{eq-lpn-empirical-discharge-capacity}
D_{b,n}
&\le \overline D_{b,0}\left(\frac n{n_0}\right)^d,\\
\label{eq-lpn-empirical-discharge-attenuation}
a_{b,n}
&\le \overline a_{b,0}e^{-\gamma(n-n_0)},\\
\label{eq-lpn-empirical-discharge-error}
e_{b,n}&\le\overline e_{b,n}
\end{align}
for represented constants \(d\ge0\), \(\gamma>0\), and every
\(n\ge n_0\).  These inequalities must cover the complete qualified source
cell, not only the trained record.  They imply the explicit cell bound
\begin{equation}
\label{eq-lpn-empirical-discharge-cell-bound}
G_{X,b,\mathfrak M,n}^{\mathrm{src,sat}}(B_n)
\le
\sqrt{\overline D_{b,0}}\,\overline a_{b,0}
\left(\frac n{n_0}\right)^{d/2}
e^{-\gamma(n-n_0)}
+\overline e_{b,n}.
\end{equation}

The attenuation check is explicitly noise-indexed.  In particular, if the
verified factorization gives
\begin{equation}
\label{eq-lpn-empirical-discharge-noise-attenuation}
a_{b,n}\le(1-2\eta)^{q_n},
\qquad
q_n\ge q_{n_0}+\alpha(n-n_0),
\end{equation}
then one may take
\begin{equation}
\label{eq-lpn-empirical-discharge-gamma}
\overline a_{b,0}=(1-2\eta)^{q_{n_0}},
\qquad
\gamma=\alpha\log\frac1{1-2\eta}.
\end{equation}
For a certified weight law \(q_n\ge\delta m(n)\) with
\(m(n)=n/R+O(1)\), any certified lower affine envelope for
\(q_n-q_{n_0}\) supplies the corresponding
\(\alpha=\delta/R+o(1)\); the verifier retains the floor and \(O(1)\)
terms rather than discarding them.

Put
\begin{equation}
\label{eq-lpn-empirical-discharge-lambda}
\lambda_\eta=\log\frac1{1-\eta}.
\end{equation}
Assume \(\gamma>\lambda_\eta\).  Let \(N_*\ge n_0\) be a represented
integer satisfying
\begin{align}
\label{eq-lpn-empirical-discharge-crossover-monotonicity}
N_*&\ge\frac{d}{2(\gamma-\lambda_\eta)},\\
\label{eq-lpn-empirical-discharge-crossover-leading}
\log\!\left(2\sqrt{\overline D_{b,0}}\,
                 \overline a_{b,0}\right)
+\frac d2\log\frac{N_*}{n_0}+\gamma n_0
&\le(\gamma-\lambda_\eta)N_*,\\
\label{eq-lpn-empirical-discharge-crossover-error}
\overline e_{b,n}&\le\frac12e^{-\lambda_\eta n}
\qquad(n\ge N_*).
\end{align}
If outward-rounded direct verification also proves
\begin{equation}
\label{eq-lpn-empirical-discharge-finite-prefix}
G_{X,b,\mathfrak M,n}^{\mathrm{src,sat}}(B_n)
\le(1-\eta)^n
\qquad(n_0\le n<N_*),
\end{equation}
then
\begin{equation}
\label{eq-lpn-empirical-discharge-final}
G_{X,b,\mathfrak M,n}^{\mathrm{src,sat}}(B_n)
\le(1-\eta)^n
\qquad(n\ge n_0).
\end{equation}
Consequently the exported certificate discharges
\cref{thm-lpn-residual-ancestry-attention-bound} at the audited ceiling.

Every premise above has a distinct verification field: source-cell coverage
comes from the cited LPN exhaustion theorems; the factorization, capacity,
attenuation, and error bounds come from the scale audit; the polynomial
ceiling comes from complete cost accounting; the crossover inequalities are
scalar interval checks; and the finite prefix is checked directly.  The
fields are not all finite computations: the crossover and finite-prefix
checks are finite, but the scale-law bounds
\labelcref{eq-lpn-empirical-discharge-capacity,eq-lpn-empirical-discharge-attenuation,eq-lpn-empirical-discharge-error}
are uniform-in-\(n\) proof obligations over the complete qualified cell and
must be discharged symbolically, so supplying them is equivalent in
strength to proving the residual-ancestry bound on that cell.  A neural
or Bayesian learner may propose all these records using
\cref{thm-empirical-saturated-bellman-upper-bound,cor-finite-network-simultaneous-audit-radius,thm-parameter-source-resolved-empirical-target-gain},
but only their
deterministic verification is used in
\cref{eq-lpn-empirical-discharge-final}.

\end{theorem}

\begin{proof}
The LPN source-exhaustion results cited in the statement verify
\cref{eq-scale-audit-source-cover} with the displayed residual-ancestry cell
as the sole unevaluated member.  Apply
\cref{eq-scale-audit-factorization,eq-scale-audit-capacity-attenuation} and
substitute
\cref{eq-lpn-empirical-discharge-capacity,eq-lpn-empirical-discharge-attenuation,eq-lpn-empirical-discharge-error};
this gives \cref{eq-lpn-empirical-discharge-cell-bound}.

Under \cref{eq-lpn-empirical-discharge-noise-attenuation},

\[
(1-2\eta)^{q_n}
\le
(1-2\eta)^{q_{n_0}}
\exp\!\left[-\alpha(n-n_0)
                 \log\frac1{1-2\eta}\right],
\]
which proves \cref{eq-lpn-empirical-discharge-gamma}.  The logarithm of the
leading term in \cref{eq-lpn-empirical-discharge-cell-bound}, divided by
\(e^{-\lambda_\eta n}\), is at most
\[
\log\!\left(\sqrt{\overline D_{b,0}}\,\overline a_{b,0}\right)
+\frac d2\log\frac n{n_0}+\gamma n_0
-(\gamma-\lambda_\eta)n.
\]
Its derivative is
\(d/(2n)-(\gamma-\lambda_\eta)\), which is nonpositive for
\(n\ge N_*\) by
\cref{eq-lpn-empirical-discharge-crossover-monotonicity}.  The check at
\(N_*\) in \cref{eq-lpn-empirical-discharge-crossover-leading} therefore
bounds the leading term by
\(\tfrac12e^{-\lambda_\eta n}\) for every \(n\ge N_*\).
Together with \cref{eq-lpn-empirical-discharge-crossover-error}, this gives
\cref{eq-lpn-empirical-discharge-final} on the propagated tail.  The direct
check \cref{eq-lpn-empirical-discharge-finite-prefix} supplies the remaining
finite sizes.  Since \(e^{-\lambda_\eta n}=(1-\eta)^n\), the conclusion is
exactly \cref{eq-lpn-residual-ancestry-attention-bound}.
\end{proof}

\begin{figure}[tbp]
\centering
\resizebox{.98\linewidth}{!}{%
\begin{tikzpicture}[fw diagram,node distance=7mm and 9mm]
  \node[fw source,text width=2.35cm] (public)
    {fixed-noise LPN\\\((n,m,\eta,R)\)};
  \node[fw provenance,text width=2.65cm,right=of public] (exhaust)
    {exact source exhaustion\\one remaining cell};
  \node[fw use,text width=2.4cm,above right=7mm and 10mm of exhaust] (capacity)
    {capacity audit\\\(D_{b,n}\le\overline D_{b,0}(n/n_0)^d\)};
  \node[fw geom,text width=2.55cm,right=of exhaust] (noise)
    {noise transfer\\\(a_{b,n}\le(1-2\eta)^{q_n}\)};
  \node[fw residual,text width=2.35cm,below right=7mm and 10mm of exhaust] (error)
    {error audit\\\(e_{b,n}\le\overline e_{b,n}\)};
  \node[fw box,text width=2.65cm,right=11mm of noise] (rates)
    {rate comparison\\\(\gamma>\lambda_\eta\)};
  \node[fw use,text width=2.4cm,above right=7mm and 11mm of rates] (prefix)
    {direct finite check\\\(n_0\le n<N_*\)};
  \node[fw geom,text width=2.4cm,below right=7mm and 11mm of rates] (tail)
    {verified propagation\\\(n\ge N_*\)};
  \node[fw terminal,text width=2.55cm,right=14mm of rates] (close)
    {\(G_{X,b}^{\mathrm{src,sat}}\le(1-\eta)^n\)};
  \draw[fw arrow] (public)--(exhaust);
  \draw[fw arrow] (exhaust)--(capacity);
  \draw[fw arrow] (exhaust)--(noise);
  \draw[fw arrow] (exhaust)--(error);
  \draw[fw arrow] (capacity)--(rates);
  \draw[fw arrow] (noise)--(rates);
  \draw[fw arrow] (error)--(rates);
  \draw[fw arrow] (rates)--(prefix);
  \draw[fw arrow] (rates)--(tail);
  \draw[fw arrow] (prefix)--(close);
  \draw[fw arrow] (tail)--(close);
\end{tikzpicture}%
}
\caption{Proof-grade LPN discharge.  Structural exhaustion fixes the sole
remaining source cell.  Capacity, noise transfer, and error certificates
establish the exponent comparison; direct verification covers the finite
prefix and scale propagation covers the tail.}
\label{fig-lpn-proof-grade-empirical-discharge}
\end{figure}

\begin{remark}[What attention compilation proves]
\label{rem-lpn-attention-compilation-scope}
Compile any selected transition into a complete attention record by
\cref{cor-attention-active-sampling-compilation}.  Then
\cref{thm-attention-structural-source-compilation,thm-attention-target-gain-envelope}
retain every score, normalization, value path, posterior update, and deployed
transition in the existing source-resolved ledger.  This proves that attention
adds no uncharged source coordinate.  It does not identify the complete public
experiment \((A,b)\) with its pivot-row statistic.  Such an identification
would require the represented pre-hit sufficiency condition of
\cref{cor-controlled-represented-prehit-sufficiency}; provenance
classification alone does not supply that condition.
\end{remark}

\begin{figure}[tbp]
\centering
\begin{tikzpicture}[fw diagram,node distance=6mm and 6mm]
\node[fw source] (pub) {LPN public record \((A,b)\)};
\node[fw use,right=of pub] (att) {complete normalized selector};
\node[fw provenance,above right=of att] (q) {qualified fibers\\closed by \(\Abs\)};
\node[fw geom,right=of att] (x) {full-carrier source};
\node[fw residual,below right=of att] (r) {terminal / \(\Res\)\\zero prospective gain};
\node[fw terminal,right=of x] (v) {\(\le(1-\eta)^n\)};
\draw[fw arrow] (pub)--(att);
\draw[fw arrow] (att)--(q);
\draw[fw arrow] (att)--(x);
\draw[fw arrow] (att)--(r);
\draw[fw arrow] (x)--(v);
\end{tikzpicture}
\caption{LPN specialization of the complete selector compilation.  Quotient,
terminal, residual, and materialized located-spectral rows are closed by
\cref{thm-lpn-exact-constant-noise-spectral-perimeter}.  The remaining compact
non-lumpable full-carrier row is evaluated by
\cref{thm-lpn-residual-ancestry-attention-bound};
\cref{thm-lpn-proof-grade-empirical-residual-discharge} gives its explicit
computer-assisted discharge protocol.}
\label{fig-lpn-attention-source-closure}
\end{figure}

\begin{corollary}[Complete ambient LPN source bound under the residual-ancestry bound]
\label{cor-lpn-complete-ambient-geom-source-seal}

Fix \(0<\eta<1/2\), and let \(\mathfrak M_{n,m,\eta}\) be the declared LPN
expectation functional on the full-rank and random-code-rigidity event.  At
every budget \(B\) in the declared polynomial saturated ceiling, assume
\cref{thm-lpn-residual-ancestry-attention-bound}.  Then
\begin{equation}
\label{eq-lpn-complete-ambient-geom-source-bound}
G_{X,\mathfrak M,n}^{\mathrm{src,sat}}(B)
\le 2^{-n}+(1-\eta)^n.
\end{equation}
The same bound holds after every class-preserving polynomial enlargement of
that ceiling.

\end{corollary}

\begin{proof}

Split the ambient source family by
\cref{def-lpn-ambient-geom-ancestry-branches}.  The residual-independent
branch is at most \(2^{-n}\) by
\cref{lem-lpn-syndrome-free-ambient-geom-bound}.

The residual-ancestry branch has the exhaustive five-family source
classification of
\cref{cor-lpn-residual-ancestry-five-family-geom-classification}; its
numerical value is bounded by
\cref{thm-lpn-residual-ancestry-attention-bound}.  Since the two ancestry
families are disjoint and exhaustive,
\[
G_{X,\mathfrak M,n}^{\mathrm{src,sat}}(B)
=\max\left\{
G_{X,0,\mathfrak M,n}^{\mathrm{src,sat}}(B),
G_{X,b,\mathfrak M,n}^{\mathrm{src,sat}}(B)
\right\}
\le\max\{2^{-n},(1-\eta)^n\}
\le2^{-n}+(1-\eta)^n.
\]
The residual-ancestry compact-ambient and residual-independent estimates are
stated at every class-preserving enlarged ceiling, so the same conclusion
holds there.

\end{proof}

\begin{corollary}[LPN specialization of the ambient obstruction criterion]
\label{cor-lpn-complete-ambient-obstruction-criterion}

Fix the LPN presentation, family functional, and saturated ceiling of
\cref{cor-lpn-complete-ambient-geom-source-seal}.  Suppose every complete
temporally admissible ambient record at that ceiling satisfies at least one
of
\[
M_\Phi=0,
\qquad
\chi_\Phi^{\rm auth}(D)=0,
\qquad
\operatorname{Align}_{\Caus}^{\rm auth}(\mathfrak K,D)\le\delta,
\qquad
R_{\{s\}}(D)\le\delta,
\qquad
\frac1{1+\rho_\Phi(D)}\le\delta.
\]
Then
\begin{equation}
\label{eq-lpn-complete-source-profile-obstruction-bound}
G_{X,\mathfrak M,n}^{\rm src,sat}(B)\le\delta,
\qquad
I_{\mathfrak M,n}^{\rm src,sat}(B)\le\delta,
\end{equation}
on the LPN good event.  For a polynomial ceiling, the complete prospective
source profile is negligible when the displayed cover holds with
\(\delta=n^{-k}\) for every fixed \(k\) and all sufficiently large \(n\).

\end{corollary}

\begin{proof}

The first conclusion is the direct LPN substitution in
\cref{thm-ambient-perturbation-characterization-obstruction}.  The exact and
compensated quotient-supported coordinates vanish by
\cref{eq-lpn-exact-abs-source-vanishes,eq-lpn-exact-quotient-geom-closed,eq-lpn-compensated-source-vanishes}.
The complete master decomposition
\cref{thm-complete-generic-source-theory,thm-exhaustive-five-family-geom-source-classification},
together with \cref{def-extraction-functional}, therefore gives
\(I_{\mathfrak M,n}^{\rm src,sat}=G_{X,\mathfrak M,n}^{\rm src,sat}\),
which proves the second conclusion without invoking the later LPN floor
theorem.  The polynomial statement is
\cref{cor-polynomial-ambient-perturbation-obstruction}.

\end{proof}

\subsection{Selector and valid-Lift consequences}
\label{subsec-lpn-selector-valid-lift-consequences}

\begin{corollary}[Prospective selector source routing]
\label{cor-lpn-prospective-selector-source-seal}

Under the hypotheses of
\cref{cor-lpn-complete-ambient-geom-source-seal}, let \(B\) be a declared
execution ceiling and let \(\widehat B=\Psi(B,n)\) be the prospective-selector
ceiling of \cref{thm-prospective-selector-accountability}.  Every positive
selector value is routed into the complete prospective profile.  Consequently,
\begin{equation}
\label{eq-lpn-prospective-selector-source-bound}
\operatorname{Sel}_{\mathfrak M,n,\eta}^{\mathrm{sat}}(\widehat B)
\le G_{X,\mathfrak M,n}^{\mathrm{src,sat}}(\widehat B)
\le 2^{-n}+(1-\eta)^n.
\end{equation}

\end{corollary}

\begin{proof}

A positive selector advantage is represented before the corresponding test
and is target-aligned by
\cref{def-represented-pre-hit-selector,thm-prospective-selector-accountability}.
Channel exhaustiveness and source no-creation place its complete qualified
record in the prospective \(\Geom/\Abs\) provenance profile at
\(\widehat B\); see
\cref{thm-channel-exhaustiveness,%
prop-structural-accounting-is-universal,%
cor-contextual-structural-charge-transfer}.  The exact-\(\Abs\) and all
quotient-supported Geom coordinates vanish by
\cref{eq-lpn-exact-abs-source-vanishes,%
eq-lpn-exact-quotient-geom-closed,%
eq-lpn-compensated-source-vanishes}.  The complete ambient coordinate is
bounded by \cref{eq-lpn-complete-ambient-geom-source-bound}.  A directly
qualified selector remains in that coordinate at complete cost by
\cref{thm-no-uncharged-prospective-value-amplification}; hence it is included
in, rather than additional to, the displayed bound.

\end{proof}

\begin{corollary}[Valid-Lift access preserves complete LPN value routing]
\label{cor-lpn-valid-lift-selector-access-seal}

Under the hypotheses of
\cref{cor-lpn-prospective-selector-source-seal}, let \(\mathscr L\) be any
represented valid-Lift class evaluated with the same family functional,
horizon \(H\), and enlarged ceiling \(\widehat B\).  Use the declared uniform
full-carrier neutral baseline \(\nu\); its singleton contact obeys
\begin{equation}
\label{eq-lpn-valid-lift-neutral-baseline}
\mathfrak M_{n,m,\eta}
\bigl(\operatorname{Enum}_{\nu}(\mathcal A,H)\bigr)
\le H2^{-n}.
\end{equation}
Then

\begin{equation}
\label{eq-lpn-valid-lift-selector-access-seal}
\operatorname{LiftSel}_{\mathscr L}^{\rm sat}(\widehat B,H)
\le
\operatorname{Sel}_{\mathfrak M,n,\eta}^{\rm sat}(\widehat B).
\end{equation}

For a clocked member of this class, let
\(E_t=P_t^\Omega U_{t+1}-U_tP_t^X\) be its represented intertwining defects,
and retain the initialization and gate-interface errors
\(\varepsilon_{\rm init},\varepsilon_{\rm gate}\) of
\cref{thm-controlled-clocked-lift-transfer}.  With the public LPN verifier
gate at threshold \(\tau\), its family-functional arrival probability obeys

\begin{align}
\label{eq-lpn-clocked-lift-sealed-arrival}
\mathfrak M_{n,m,\eta}
 \bigl(p_{\widehat T_I,H}^{\Omega}\bigr)
\le{}&
H\bigl(2^{-n}
       +\operatorname{Sel}_{\mathfrak M,n,\eta}^{\rm sat}(\widehat B)\bigr)\notag\\
&+\mathfrak M_{n,m,\eta}\!\left(
 \varepsilon_{\rm init}
 +\sum_{t=0}^{H-1}\|E_t\|
 +\varepsilon_{\rm gate}\right)\notag\\
&+\exp[-mD_{\mathrm{Ber}}(\tau\Vert\eta)]
 +(2^n-1)2^{-m(1-H_2(\tau))}.
\end{align}

Thus an exact valid lift contributes no interface defect and preserves the
complete prospective-value routing of its base selector.  A clocked lift
additionally contributes its defined interface defects, and the public
verifier contributes one explicit
\((n,m,\eta,\tau)\)-dependent exceptional-presentation term.

\end{corollary}

\begin{proof}

The inequality in \cref{eq-lpn-valid-lift-selector-access-seal} is
\cref{eq-controlled-liftsel-subprofile}.  Apply
\(\mathfrak M_{n,m,\eta}\) to the hazard identity in
\cref{thm-prospective-selector-accountability}, then use sublinearity,
\cref{eq-lpn-valid-lift-neutral-baseline}, and the saturated selector
supremum.  The base first-arrival probability is therefore at most
\[
\mathfrak M_{n,m,\eta}
 \bigl(p_{T_I,H}^{X}\bigr)
\le H\bigl(2^{-n}
 +\operatorname{Sel}_{\mathfrak M,n,\eta}^{\rm sat}(\widehat B)\bigr).
\]
Apply the clocked pullback estimate in
\cref{thm-controlled-clocked-lift-transfer}, followed by the public-gate
family comparison
\cref{eq-lpn-clocked-lift-public-gate-family-transfer}, to obtain
\cref{eq-lpn-clocked-lift-sealed-arrival}.

\end{proof}

\section{Exact Reduction of the Saturated LPN Profile}
\label{sec-lpn-exact-saturated-profile-reduction}
\label{subsec-lpn-exact-saturated-profile-reduction}

\Cref{cor-lpn-abs-source-exhaustion,%
cor-lpn-quotient-supported-geom-closure,%
thm-lpn-coupled-source-seal-final} remove the exact-\(\Abs\) mode and every
quotient-supported Geom mode from the prospective source ledger.  By the
master decomposition, the only remaining prospective coordinate is ambient
Geom.  The separate first-hit accounting ledger is retained by
\cref{def-core-refined-first-hit-accounting-profile,%
thm-core-refined-prospective-accountability}; temporal source admissibility
forbids an automatic transfer between the two ledgers.

\begin{theorem}[Exact saturated-profile ledger for LPN under the residual-ancestry bound]
\label{thm-lpn-saturated-floor}

Fix \(0<\eta<1/2\) and a sample schedule with \(n/m\to R\in(0,1)\).
On the full-rank and random-code-rigidity event, fix a presented instance
\(I_n\), and let
\(D_{\mathrm{poly}}^{\mathrm{hit}}\) be the first-hit normalization ceiling
supplied by \cref{thm-core-refined-prospective-accountability}.  Apply
\cref{thm-lpn-residual-ancestry-attention-bound} at
\(b_{\mathrm{poly}}(n)\).  Then

\begin{equation}
\label{eq-lpn-saturated-floor-abs-closed}
m_{I_n}^{\mathrm{src,sat}}
 \bigl(b_{\mathrm{poly}}(n)\bigr)
=
G_{Q,\mathrm{ex},n}^{\mathrm{src,sat}}
 \bigl(b_{\mathrm{poly}}(n)\bigr)
=
G_{Q,\mathrm{rev},n}^{\mathrm{src,sat}}
 \bigl(b_{\mathrm{poly}}(n)\bigr)
=
G_{Q,\mathrm{nr},n}^{\mathrm{src,sat}}
 \bigl(b_{\mathrm{poly}}(n)\bigr)
=0,
\end{equation}

\begin{align}
\label{eq-lpn-exact-saturated-source-obstruction}
I_n^{\mathrm{src,sat}}
 \bigl(b_{\mathrm{poly}}(n)\bigr)
&=
G_{X,n}^{\mathrm{src,sat}}
 \bigl(b_{\mathrm{poly}}(n)\bigr),\\
\label{eq-lpn-family-saturated-source-obstruction}
I_{\mathfrak M,n}^{\mathrm{src,sat}}
 \bigl(b_{\mathrm{poly}}(n)\bigr)
&=
G_{X,\mathfrak M,n}^{\mathrm{src,sat}}
 \bigl(b_{\mathrm{poly}}(n)\bigr),\\
\label{eq-lpn-family-saturated-ancestor-obstruction}
G_{X,\mathfrak M,n}^{\mathrm{src,sat}}
 \bigl(b_{\mathrm{poly}}(n)\bigr)
&\le 2^{-n}+(1-\eta)^n,\\
\label{eq-lpn-exact-first-hit-obstruction}
\operatorname{HitCore}_{\mathfrak M,n}^{\mathrm{sat}}
 (D_{\mathrm{poly}}^{\mathrm{hit}})
&=
\max_{\substack{K_{\mathrm{st}},K_{\mathrm{cf}}\\
S_{\mathrm{st}}\supseteq K_{\mathrm{st}}\\
S_{\mathrm{cf}}\supseteq K_{\mathrm{cf}}}}
\operatorname{HitCore}_{
 \mathfrak M,K_{\mathrm{st}},K_{\mathrm{cf}}
 \mid S_{\mathrm{st}},S_{\mathrm{cf}},n}^{\mathrm{sat},=}
 (D_{\mathrm{poly}}^{\mathrm{hit}}).
\end{align}

For the fixed declared polynomial execution envelope \(\mathcal B_C\), at
its common charged normalization ceiling \(B_C^\dagger\), the corresponding
pointwise success bound on the public-verifier good event is
\begin{equation}
\label{eq-lpn-exact-success-obstruction}
\operatorname{Succ}_{I_n,n,\eta}(\mathcal B_C)
\le
eH_C\,
\operatorname{HitCore}_{I_n,n}^{\mathrm{sat}}
 (B_C^\dagger).
\end{equation}

These identities and bounds are preserved by conservative extensions through
\hyperref[def-lpn-admissible-oracle]{presentation-admissible LPN
oracles}.  The first displays separate the exact pointwise provenance
reduction, the complete ambient source coordinate, and its saturated
evaluation.  Complete prospective values are routed by
\cref{cor-lpn-no-uncharged-prospective-value}; no separate derived-value
coordinate is inserted.  The next
display is the exact finite decomposition of the retrospective accounting
profile, and the pointwise accounting bound uses only that profile.  No
comparison between the two profile types is asserted.

\end{theorem}

\begin{proof}\phantomsection\label{proof-lpn-saturated-floor}

\Cref{eq-lpn-exact-abs-source-vanishes,eq-lpn-exact-quotient-geom-closed,%
eq-lpn-compensated-source-vanishes}
gives \cref{eq-lpn-saturated-floor-abs-closed}.
The exhaustive master decomposition
\cref{thm-complete-generic-source-theory,%
thm-exhaustive-five-family-geom-source-classification} expresses
\(G_n^{\mathrm{src,sat}}\) as the maximum of the ambient, exact-quotient,
reversible-compensated, and resolvent-compensated coordinates.  The preceding
vanishing identities leave exactly \(G_{X,n}^{\mathrm{src,sat}}\).  Since
\(I_n^{\mathrm{src,sat}}=m_n^{\mathrm{src,sat}}+G_n^{\mathrm{src,sat}}\)
by \cref{def-extraction-functional}, this proves
the equality in \cref{eq-lpn-exact-saturated-source-obstruction}.  Apply the
declared LPN family functional to the same exact decomposition to obtain
\cref{eq-lpn-family-saturated-source-obstruction}; then apply the complete
ambient seal \cref{eq-lpn-complete-ambient-geom-source-bound} to prove
\cref{eq-lpn-family-saturated-ancestor-obstruction}.

\Cref{eq-core-refined-first-hit-maximum} gives
\cref{eq-lpn-exact-first-hit-obstruction}.  Apply
the pointwise form of \cref{thm-core-refined-prospective-accountability} to obtain
\cref{eq-lpn-exact-success-obstruction}.  The LPN representation lemma
\cref{lem-lpn-no-creation-final} and oracle admissibility theorem
\cref{thm-oracle-admissibility-presentation-separation} retain the same
represented ancestry and class-preserving costs under admissible oracle
interfaces.

The completed first-hit record is accounting-admissible but not a
prospective source by
\cref{cor-first-hit-accounting-not-source}.  Its exact visibility is
nevertheless retained in HitCore and charged by
\cref{eq-core-refined-first-hit-charge-identity}; hence it is neither
discarded nor used as a retrospective structural source.

\end{proof}
\Cref{fig-lpn-closure} collects the exact source and accounting reduction in
\Cref{thm-lpn-saturated-floor}.

\begin{figure}[tbp]
\centering
\begin{tikzpicture}[fw diagram]
  \node[fw source,minimum width=3.5cm,text width=3.2cm] (master) at (-4.5,2.3)
    {prospective master profile\\\(\Abs\oplus\Geom\)};
  \node[fw residual,minimum width=4.4cm,text width=4.1cm] (closed) at (-4.5,0.4)
    {vanishing modes:\\
     \(m=G_{Q,\mathrm{ex}}=0\),\\
     \(G_{Q,\mathrm{rev}}=G_{Q,\mathrm{nr}}=0\)};
  \node[fw geom,minimum width=4.0cm,text width=3.7cm] (ambient) at (0.4,2.3)
    {by \cref{thm-lpn-residual-ancestry-attention-bound}:\\
     \(I^{\mathrm{src,sat}}\le2^{-n}+(1-\eta)^n\)};
  \node[fw terminal,minimum width=4.0cm] (identity) at (0.4,0.4)
    {\(\operatorname{Sel}^{\rm sat}\le I^{\rm src,sat}\)};
  \node[fw source,minimum width=3.3cm] (selector) at (-2.0,-1.5)
    {represented pre-hit selector};
  \node[fw box,minimum width=4.3cm] (success) at (3.0,-1.5)
    {\(\operatorname{Succ}\le H(2^{1-n}+(1-\eta)^n)\)};
  \node[fw terminal,minimum width=3.1cm] (hit) at (3.0,-2.8)
    {\(\operatorname{HitCore}^{\rm sat}\)};
  \draw[fw arrow] (master) -- (ambient);
  \draw[fw dashed arrow] (master) -- (closed);
  \draw[fw arrow] (ambient) -- (identity);
  \draw[fw arrow] (identity) -- (selector);
  \draw[fw arrow] (selector) -- (success);
  \draw[fw dashed arrow] (success) -- (hit);
\end{tikzpicture}
\caption{The exact LPN reduction.  Exact abstraction and all three
quotient-supported geometric modes vanish, and the complete saturated source
profile is bounded at the common charged ceiling by
\cref{thm-lpn-residual-ancestry-attention-bound}.  Prospective selector
accountability then gives the displayed success inequality.  The separate
\(\operatorname{HitCore}\) coordinate
retains completed first-hit accounting and is not a hardness hypothesis.
The source and success bounds are
\cref{thm-lpn-saturated-floor,thm-lpn-one-way}.}
\label{fig-lpn-closure}
\end{figure}

\section{Hardness from the Saturated Prospective Profile}
\label{sec-lpn-saturated-accounting-criterion}

The final LPN theorem applies prospective selector accountability at the
same represented budget.  The noise schedules
\(\eta=\eta(n)\), including the zero-noise and polynomial-enumeration scales,
are localized separately in
\cref{cor-lpn-quantitative-noise-phase-localization}.
\begin{theorem}[Admissible-oracle hardness of search-LPN under the residual-ancestry bound]
\label{thm-lpn-one-way}

Fix \(0<\eta<1/2\), choose \(0<\delta<1/2-\eta\), and set
\(\tau=\eta+\delta\). Let \(n/m\to R\in(0,1)\), with

\[
m\bigl(1-H_2(\tau)\bigr)
\ge n+\omega(\log n).
\]

Fix the declared polynomial execution envelope \(\mathcal B_C\), with
horizon \(H_C\), and put
\begin{equation}
\label{eq-lpn-universal-bayesian-compiled-ceiling}
B_C^{\rm uBayes}
=
\Psi_{\rm uBayes}
\bigl(\Psi_{\rm act}(\mathcal B_C,H_C,n,\mathbf 0),H_C,n\bigr),
\qquad
B_C^\sharp=\Psi(B_C^{\rm uBayes},n),
\qquad
H_C^\sharp=H_{B_C^\sharp}(n).
\end{equation}
For the declared LPN family functional restricted to the intersection of the
public-verifier good event of \cref{thm-lpn-public-verifier} and the
full-rank random-code-rigidity event of
\cref{thm-lpn-random-code-rigidity-final},
write
\(\mathfrak M^*=\mathfrak M^{\mathcal G_*(\tau)}_{n,m,\eta}\) in the
sense of \cref{def-lpn-named-events-restricted-functionals}.  The complete
saturated LPN source-profile bound of \cref{thm-lpn-saturated-floor} applies
at the class-preserving common ceiling \(B_C^\sharp\).  Apply
\cref{thm-lpn-residual-ancestry-attention-bound} at that ceiling.  Then
\begin{equation}
\label{eq-lpn-one-way-obstruction-bound}
\operatorname{Succ}_{\mathfrak M^*,n,\eta}(\mathcal B_C)
\le
H_C\bigl(2^{1-n}+(1-\eta)^n\bigr).
\end{equation}
Taking the declared LPN family functional and bounding success by one off
that event gives the full-family distributional inequality
\begin{equation}
\label{eq-lpn-one-way-family-obstruction-bound}
\begin{aligned}
\operatorname{Succ}_{\mathfrak M,n,\eta}(\mathcal B_C)
\le{}&
H_C\bigl(2^{1-n}+(1-\eta)^n\bigr)\\
&+\varepsilon_{\mathrm{rig}}(n,m)\\
&+\exp[-mD_{\mathrm{Ber}}(\tau\Vert\eta)]
+(2^n-1)2^{-m(1-H_2(\tau))}.
\end{aligned}
\end{equation}
The bound includes computations expressed through presentation-admissible LPN
oracle interfaces.  For fixed \(0<\eta<1/2\), polynomial \(H_C,H_C^\sharp\),
and the displayed sampling condition, every term on the right of
\cref{eq-lpn-one-way-family-obstruction-bound} is negligible.

\end{theorem}

\begin{proof}\phantomsection\label{proof-lpn-one-way}

On the good-presentation event of
\cref{thm-lpn-public-verifier}, the represented verifier accepts exactly
\(T_I=\{s\}\).  The LPN representation lemma
\cref{lem-lpn-no-creation-final} places the declared execution envelope,
including presentation-admissible oracle expansions and a final verifier
call, in the saturated represented computation class.  Apply the universal
compiler of
\cref{cor-universal-represented-computation-bayesian-ceiling} at
\cref{eq-lpn-universal-bayesian-compiled-ceiling}.  It preserves the complete
execution and first-hit laws.  The source-resolved LPN Bayesian theorem
\cref{thm-lpn-posterior-concentration-five-mode-exhaustion} retains the
recordwise framework candidates and the complete ambient record.
The compact-state compiler and
\cref{thm-lpn-optimal-latent-useful-belief-update} apply that same source
evaluation to every latent representation, optimal acquisition rule,
Bayes--Fisher update, learned decoder, and attention selector in the compiled
polynomial saturated class.  Their maximum useful belief update is the
profile in \cref{eq-lpn-optimal-latent-useful-gain-polynomial-ceiling} and is
bounded by \(2^{-n}+(1-\eta)^n\) at this same ceiling.

Since \(\mathcal G_*(\tau)\subseteq\mathcal G_{\rm ver}(\tau)\), the public
candidate verifier again accepts exactly \(\{s\}\) on \(\mathcal G_*(\tau)\);
the proof of
\cref{thm-lpn-prospective-learning-selector-accountability} therefore applies
verbatim with the restricting event \(\mathcal G_{\rm ver}(\tau)\) replaced by
\(\mathcal G_*(\tau)\) and the functional \(\mathfrak M^{\rm ver}_{n,m,\eta}\)
by the further-restricted functional \(\mathfrak M^*\).  Apply the resulting
LPN prospective selector bound
\cref{eq-lpn-prospective-neutral-selector-bound} to the compiled record.  The
uniform neutral baseline on the full LPN carrier contributes at most
\(H_C2^{-n}\), so
\[
\operatorname{Succ}_{\mathfrak M^*,n,\eta}(\mathcal B_C)
\le
H_C2^{-n}
+H_C\operatorname{Sel}_{\mathfrak M^*,n,\eta}^{\mathrm{sat}}(B_C^\sharp).
\]
Framework source routing and the complete saturated LPN profile evaluation
in \cref{thm-lpn-saturated-floor} give, at the same ceiling,
\[
\operatorname{Sel}_{\mathfrak M^*,n,\eta}^{\mathrm{sat}}(B_C^\sharp)
\le
I_{\mathfrak M^*,n}^{\mathrm{src,sat}}(B_C^\sharp)
=
G_{X,\mathfrak M^*,n}^{\mathrm{src,sat}}(B_C^\sharp)
\le 2^{-n}+(1-\eta)^n.
\]
Substitution gives
\begin{equation}
\label{eq-lpn-one-way-proof-selector-accounting}
\operatorname{Succ}_{\mathfrak M^*,n,\eta}(\mathcal B_C)
\le
H_C\bigl(2^{1-n}+(1-\eta)^n\bigr).
\end{equation}

Equation \cref{eq-lpn-one-way-proof-selector-accounting} is
\cref{eq-lpn-one-way-obstruction-bound}.  Valid Lift, Caus, boundary,
filtering, likelihood, normalization, and belief updates retain their
framework source cell and complete cost.

The source deduction uses only prospective records.  The completed first-hit
record belongs to the retrospective audit profile by
\cref{cor-first-hit-accounting-not-source}; neither it nor
\(\operatorname{HitCore}^{\mathrm{sat}}\) is a premise of
\cref{eq-lpn-one-way-proof-selector-accounting}.

Extend the restricted family expectation to the complete LPN family.  On the
complement of the rigidity event or the verifier-good event, use the unit
bound on success.  The union bound, together with
\cref{eq-lpn-rigidity-failure-probability,eq-lpn-verifier-good-event-probability},
gives the three exceptional-presentation terms in
\cref{eq-lpn-one-way-family-obstruction-bound}.

The complement of the good event is included explicitly in
\cref{eq-lpn-one-way-family-obstruction-bound}; its probability is an
evaluated error term, not an additional structural assumption.

\end{proof}

\begin{corollary}[Structural meta-complexity lower region for LPN]
\label{cor-lpn-structural-metacomplexity-lower-region}

Adopt the parameters, restricted functional \(\mathfrak M^*\), and common
charged ceiling \(B_C^\sharp\) of \cref{thm-lpn-one-way}, and put
\begin{equation}
\label{eq-lpn-metacomplexity-threshold}
\theta_{n,\eta}^{\mathrm{meta}}
=2^{-n}+(1-\eta)^n.
\end{equation}
For every \(n\) for which
\(\theta_{n,\eta}^{\mathrm{meta}}<1\), and therefore for every sufficiently
large \(n\) at fixed \(\eta>0\), the complete LPN pre-hit selector datum
satisfies
\begin{equation}
\label{eq-lpn-selector-meta-frontier-exclusion}
B_C^\sharp
\notin
\mathfrak P_{\mathfrak M^*,n,m,\eta,H_C}^{\mathrm{sel},>}
\bigl(\theta_{n,\eta}^{\mathrm{meta}}\bigr).
\end{equation}
For a scalar total-cost order whose sublevel \(b\) is represented by the
ceiling \(B_C^\sharp\), the Pareto exclusion implies
\begin{equation}
\label{eq-lpn-selector-scalar-metacomplexity-lower-bound}
\operatorname{SMC}_{\mathfrak M^*,n,m,\eta,H_C,\kappa}^{
\mathrm{sel},>}
\bigl(\theta_{n,\eta}^{\mathrm{meta}}\bigr)
\ge \kappa(B_C^\sharp),
\end{equation}
and the inequality is strict under the locally finite discrete scalar-cost
convention.  The exact boundary statement is the Pareto exclusion
\cref{eq-lpn-selector-meta-frontier-exclusion}; no injectivity of the
scalarization is required there.
Prospective meta-accountability gives
\begin{equation}
\label{eq-lpn-metacomplexity-success-bound}
\operatorname{Succ}_{\mathfrak M^*,n,\eta}(\mathcal B_C)
\le
H_C\bigl(2^{-n}+\theta_{n,\eta}^{\mathrm{meta}}\bigr)
=H_C\bigl(2^{1-n}+(1-\eta)^n\bigr).
\end{equation}

\end{corollary}

\begin{proof}
The proof of \cref{thm-lpn-one-way} establishes at the same ceiling
\begin{equation*}
\operatorname{Sel}_{\mathfrak M^*,n,\eta}^{\mathrm{sat}}(B_C^\sharp)
\le2^{-n}+(1-\eta)^n
=\theta_{n,\eta}^{\mathrm{meta}}.
\end{equation*}
By \cref{def-family-functional-selector-meta-datum}, this is exactly the
saturated profile of
\(\mathfrak D_{\mathfrak M^*,n,m,\eta,H_C}^{\mathrm{sel}}\) at that
ceiling.
Apply the exact generalized inverse
\cref{eq-exact-strict-profile-frontier-inverse} to obtain
\cref{eq-lpn-selector-meta-frontier-exclusion}.  The scalar statement is
\cref{eq-scalar-meta-inverse-implications}.  Finally, the uniform neutral
baseline contributes at most \(H_C2^{-n}\); substitute the frontier bound in
\cref{thm-structural-meta-accountability} to obtain
\cref{eq-lpn-metacomplexity-success-bound}.
\end{proof}

\begin{corollary}[Quantitative localization of the LPN noise phases]
\label{cor-lpn-quantitative-noise-phase-localization}

Let \(n/m\to R\in(0,1)\).  In this corollary the noise may be zero, may vary
with \(n\), or may be a fixed element of \((0,1/2)\), as specified in each
item.

\begin{enumerate}[label=\textup{(\roman*)}]
\item At \(\eta=0\), Gaussian elimination recovers \(s\) in polynomial time
with probability at least \(1-2^{n-m}\).

\item Under the polynomial enumeration hypotheses of
\cref{cor-lpn-polynomial-enumeration-window}, every schedule satisfying
\(\eta m\le c\log n\), for fixed \(c>0\), has a represented polynomial-time
recovery procedure with success at least
\(n^{-c+o(1)}-\operatorname{negl}(n)\).  For every fixed affordable radius
\(k_C\) and target probability \(0<\alpha<1\), the constant-success boundary
is \(\eta=\lambda_{k_C,\alpha}/m+o(m^{-1})\), where
\(\lambda_{k_C,\alpha}\) is determined by
\cref{eq-lpn-poisson-threshold-equation}.

\item By \cref{thm-lpn-residual-ancestry-attention-bound}, at fixed
\(0<\eta<\eta_{\mathrm{IT}}(R)=H_2^{-1}(1-R)\), the declared polynomial
execution envelope satisfies
\cref{eq-lpn-one-way-family-obstruction-bound}.  Its saturated structural
term and the rigidity and verifier-event errors are negligible under the
stated sampling condition.

\item For every fixed \(\eta>\eta_{\mathrm{IT}}(R)\), every decoder has
exponentially small exact-recovery probability.
\end{enumerate}

Thus the low-noise construction and the information-theoretic converse are
unconditional in their stated regimes.  The fixed-positive computational
bound uses
\cref{thm-lpn-residual-ancestry-attention-bound}.  For a finite scalar
budget and target success
\(\alpha\), the supremal affordable noise is
\(\eta_{\mathrm{comp}}^{\mathrm{sat}}\) from
\cref{eq-noise-indexed-computational-threshold}.  It is the endpoint of the
affordable interval whenever the budget absorbs represented flip
postprocessing.  Its explicit enumeration
lower bound, structural hard-side certificate, statistical-resolution
certificate, and information upper bound are
\cref{eq-lpn-enumeration-computational-threshold-lower-bound,eq-lpn-saturated-hard-side-threshold,eq-lpn-statistical-computational-threshold-lower-bound,eq-lpn-information-computational-threshold-upper-bound}.

\end{corollary}

\begin{proof}

At zero noise, \(b=As\).  The rank estimate in the proof of
\cref{prop-lpn-error-enumeration-success} gives
\(\Pr\{\operatorname{rank}(A)<n\}<2^{n-m}\); on its complement, row reduction
returns the unique secret.  This proves (i).

For (ii), the weight-zero branch of the represented enumerator succeeds on
\(E=0\).  Under \(\eta m\le c\log n\) and the stated polynomial sample
scaling,

\[
\Pr\{E=0\}=(1-\eta)^m\ge n^{-c+o(1)}.
\]

Subtracting the negligible algebraic failure in
\cref{eq-lpn-error-enumeration-algebraic-failure} gives the claim.  The
constant-success assertion is
\cref{eq-lpn-constant-success-enumeration-scale}.

For (iii), choose a constant \(\tau\) with
\(\eta<\tau<\eta_{\mathrm{IT}}(R)\).  Since
\(R<1-H_2(\tau)\),

\[
m(1-H_2(\tau))-n=\Theta(m),
\]

so the verification hypotheses of \cref{thm-lpn-one-way} hold; its
\cref{eq-lpn-one-way-family-obstruction-bound} gives the bound in (iii).
Part (iv) is the strong converse
\cref{eq-lpn-information-strong-converse}, with a fixed
\[
0<\epsilon<
\min\left\{
\eta,\frac{R-1+H_2(\eta)}{L_\eta}
\right\}.
\]

\end{proof}

\Cref{fig-lpn-three-noise-phases} collects the two asymptotic scales and the
finite-budget threshold relations.

\begin{figure}[tbp]
\centering
\begin{tikzpicture}[fw diagram,x=1cm,y=1cm]
  \node[font=\small\bfseries,anchor=east] at (-.25,2.2) {fixed \(\eta\)};
  \draw[draw=fwInk,line width=.8pt] (0,2.2)--(11.2,2.2);
  \fill[fwTeal] (0,2.2) circle[radius=2.2pt];
  \node[font=\scriptsize,anchor=south west] at (0,2.3)
    {\(\eta=0\): linear algebra};
  \path[fill=fwAmber!16,draw=fwAmber,rounded corners=1pt]
    (.25,1.42) rectangle (6.65,2.12);
  \node[font=\scriptsize,align=center,text width=5.9cm] at (3.45,1.77)
    {fixed \(0<\eta<\eta_{\rm IT}\): saturated bound\\[-1pt]
     \(\operatorname{Succ}\le
       H_C(2^{1-n}+(1-\eta)^n)+\varepsilon_{\rm good}\)};
  \draw[draw=fwBlue,line width=1.1pt] (6.8,1.32)--(6.8,2.42);
  \node[font=\scriptsize,fwBlue,anchor=south] at (6.8,2.43)
    {\(\eta_{\rm IT}=H_2^{-1}(1-R)\)};
  \path[fill=fwCoral!14,draw=fwCoral,rounded corners=1pt]
    (6.95,1.42) rectangle (11.2,2.12);
  \node[font=\scriptsize,align=center,text width=3.8cm] at (9.08,1.77)
    {information limit\\exponentially small success};
  \node[font=\scriptsize,anchor=north east] at (11.2,1.42) {\(1/2\)};

  \node[font=\small\bfseries,anchor=east] at (-.25,.25) {budget zoom};
  \draw[draw=fwInk,line width=.8pt] (0,.25)--(11.2,.25);
  \path[fill=fwTeal!16,draw=fwTeal,rounded corners=1pt]
    (.55,-.22) rectangle (3.0,.17);
  \node[font=\scriptsize,align=center,text width=2.2cm] at (1.78,0)
    {enumeration\\easy side};
  \draw[draw=fwTeal,line width=1.1pt] (3.15,-.3)--(3.15,.55);
  \node[font=\scriptsize,fwTeal,anchor=south] at (3.15,.57)
    {\(\eta_{\rm EE}\)};
  \draw[draw=fwViolet,line width=1.1pt] (5.6,-.3)--(5.6,.55);
  \node[font=\scriptsize,fwViolet,anchor=south] at (5.6,.57)
    {\(\eta_{\rm comp}^{\rm sat}\)};
  \draw[draw=fwBlue,line width=1.1pt] (9.0,-.3)--(9.0,.55);
  \node[font=\scriptsize,fwBlue,anchor=south] at (9.0,.57)
    {\(\eta_{\rm IT}^{\rm con}\)};
  \draw[fw arrow] (3.3,-.55)--node[below,font=\scriptsize]
    {supremal affordable noise lies in the certified bracket} (8.85,-.55);

  \node[fw provenance,text width=9.5cm,font=\scriptsize,align=center]
    (stat) at (5.6,-1.55)
    {\(\eta_{\rm Rad}\) is the separate sample-resolution endpoint;\\
      it is an easy-side computational bound when the empirical minimizer
      fits the same budget};
\end{tikzpicture}
\caption{The LPN noise regimes.  The upper row combines the
fixed-noise saturated-profile bound and the
unconditional information threshold from
\cref{thm-lpn-one-way,thm-lpn-information-threshold,cor-lpn-quantitative-noise-phase-localization}.
At a common high-probability target level \(\alpha=1-\delta\), whenever the
endpoints are defined and under the stated enumeration hypotheses, the lower
row gives the finite-budget relations
\(\eta_{\rm EE}\le\eta_{\rm comp}^{\rm sat}\le
\eta_{\rm IT}^{\rm con}\); the Rademacher endpoint joins the easy side only
when its represented empirical minimizer is affordable.  The diagram is
schematic; horizontal positions in both rows are not to scale.}
\label{fig-lpn-three-noise-phases}
\end{figure}

\begin{corollary}[Saturated statistical denoising obstruction for LPN]
\label{cor-lpn-universal-statistical-denoising-bound}

Adopt the hypotheses of \cref{thm-lpn-one-way}.  Fix a represented learning
ceiling \(B\) and an inverse-polynomial threshold
\(\varepsilon\in(0,1]\).  Let \(B_\varepsilon^{\mathrm{GL}}\) include the
represented bounded-score evaluator, Goldreich--Levin list decoder at
threshold \(\varepsilon\), and public verification pass, and let
\(p_{\mathrm{GL}}>0\) be its fixed amplified conditional success
probability.  At the common normalized ceiling
\((B_\varepsilon^{\mathrm{GL}})^\dagger\), the saturated clean-alignment
profile obeys

\begin{equation}
\label{eq-lpn-universal-clean-alignment-ceiling}
\begin{aligned}
\operatorname{Align}_{\mathfrak M,n,\eta}^{\mathrm{sat,stat}}(B,m)
\le{}&
\varepsilon
+\frac{H_{B_\varepsilon^{\mathrm{GL}}}}{p_{\mathrm{GL}}}
 \left(2^{-n}
 +\operatorname{Sel}_{\mathfrak M,n,\eta}^{\mathrm{sat}}
 \bigl((B_\varepsilon^{\mathrm{GL}})^\dagger\bigr)\right)
+\varepsilon_{\mathrm{rig}}(n,m)
+\exp[-mD_{\mathrm{Ber}}(\tau\Vert\eta)]\\
&+(2^n-1)2^{-m(1-H_2(\tau))}.
\end{aligned}
\end{equation}

For a sign-valued predictor in the same saturated class, put
\begin{align*}
\Delta_\varepsilon
={}&
\varepsilon
+\frac{H_{B_\varepsilon^{\mathrm{GL}}}}{p_{\mathrm{GL}}}
 \left(2^{-n}
 +\operatorname{Sel}_{\mathfrak M,n,\eta}^{\mathrm{sat}}
 \bigl((B_\varepsilon^{\mathrm{GL}})^\dagger\bigr)\right)
+\varepsilon_{\mathrm{rig}}(n,m)
+\exp[-mD_{\mathrm{Ber}}(\tau\Vert\eta)]\\
&+(2^n-1)2^{-m(1-H_2(\tau))}.
\end{align*}
Its family-average clean parity-prediction and fresh-row
noise-reconstruction risks satisfy

\begin{equation}
\label{eq-lpn-universal-denoising-risk-floor}
\mathbb E_{I,\omega} R_0(h)
\ge\frac{1-\Delta_\varepsilon}{2},
\qquad
\Pr_{I,\omega,a,e}\{\widehat e\ne e\}
\ge\frac{1-\Delta_\varepsilon}{2}.
\end{equation}

For every \(0<\delta_{\mathrm{stat}}<1\), on the simultaneous event of
\cref{thm-lpn-saturated-statistical-alignment}, the empirical-alignment
profile also obeys

\begin{equation}
\label{eq-lpn-universal-empirical-alignment-ceiling}
\sup_{\substack{\mathcal A\ \mathrm{represented}\\
\operatorname{Cost}(\mathcal A)\preceq B}}
\mathfrak M_n\!\left(
 I\longmapsto
 \mathbb E_{S,\omega}
 \left|\widehat c_{\eta,S}
 (h_{\mathcal A,I,S,\omega})\right|
\right)
\le
2\Gamma_{\mathrm{LPN},n,B,m,\delta_{\mathrm{stat}}}^{\mathrm{sat,stat}}
+(1-2\eta)\Delta_\varepsilon.
\end{equation}

Thus the statistical term controls adaptive sample fitting, while the
computational contribution is localized in the same saturated prospective
selector profile.  A negligible denoising obstruction follows when that
profile is negligible at the compiled Goldreich--Levin ceiling.  The
complete learner, optimizer, locator, retained
state, and downstream use obey the closed budget accounting of
\cref{cor-lpn-closed-budget-learning-accounting}.

\end{corollary}

\begin{proof}

For a represented evaluator \(h\) in the budget-\(B\) class, the
Goldreich--Levin heavy-Fourier algorithm returns, conditional on
\(\lvert\widehat h(s)\rvert\ge\varepsilon\), a polynomial-size list
containing \(s\) with probability at least \(p_{\mathrm{GL}}\)
\citep{GoldreichLevin1989}.  The standard bounded-score form follows by using
the represented internal randomness to sample a sign with conditional mean
\(h(a)\).  Evaluate the public predicate \(V_I\) of
\cref{thm-lpn-public-verifier} on every returned candidate.  On its
high-probability good-instance event, exactly one candidate is accepted and it
is \(s\).

The learner, evaluator, list decoder, and verification pass form one
family-uniform represented record at \(B_\varepsilon^{\mathrm{GL}}\).  All
intermediate records and admissible oracle calls retain their costs by
\cref{def-budgeted-derivability-closure,lem-lpn-no-creation-final}.  On the
verifier-good event, the family-averaged success of this one record is at
least \(p_{\mathrm{GL}}\) times the probability of
\(\lvert\widehat h(s)\rvert\ge\varepsilon\) within that event.  Apply the
family-functional prospective selector bound from
\cref{thm-prospective-selector-accountability}, and bound the complementary event by
\cref{eq-lpn-rigidity-failure-probability,eq-lpn-verifier-good-event-probability}.
Therefore
\[
\begin{aligned}
\Pr_{I,\omega}\!\left\{
\lvert\widehat h(s)\rvert\ge\varepsilon
\right\}
\le{}&
\frac{H_{B_\varepsilon^{\mathrm{GL}}}}{p_{\mathrm{GL}}}
 \left(2^{-n}
 +\operatorname{Sel}_{\mathfrak M,n,\eta}^{\mathrm{sat}}
 \bigl((B_\varepsilon^{\mathrm{GL}})^\dagger\bigr)\right)
+\varepsilon_{\mathrm{rig}}(n,m)
+\exp[-mD_{\mathrm{Ber}}(\tau\Vert\eta)]\\
&
+(2^n-1)2^{-m(1-H_2(\tau))}.
\end{aligned}
\]
Since \(\lvert\widehat h(s)\rvert\le1\), splitting its expectation at
\(\varepsilon\), applying the displayed probability bound, and taking the
saturated supremum proves
\cref{eq-lpn-universal-clean-alignment-ceiling}.

For a sign-valued predictor,

\[
R_0(h)=\frac{1-\mathbb E_a[h(a)f_s(a)]}{2}.
\]

Taking expectation over the represented randomness, the first inequality in
\eqref{eq-lpn-universal-denoising-risk-floor} follows from
\eqref{eq-lpn-universal-clean-alignment-ceiling}; the second is the exact
noise-reconstruction identity
\eqref{eq-lpn-noise-estimation-from-prediction}, averaged over the same seed.
Finally, apply the population--empirical transport inequality
\eqref{eq-lpn-uniform-statistical-alignment}, then take the saturated
supremum to obtain
\eqref{eq-lpn-universal-empirical-alignment-ceiling}.

\end{proof}

\begin{corollary}[Fixed-noise search-LPN hardness implies \(P\ne NP\)]
\label{cor-lpn-p-not-equal-np}

Use the standard finite binary encoding of the LPN presentation in
\cref{def-lpn-task-object}.  Fix rational constants
\(0<\eta<\tau<1/2\) and a sample schedule satisfying
\begin{equation}
\label{eq-lpn-p-np-public-verifier-schedule}
m\bigl(1-H_2(\tau)\bigr)\ge n+\omega(\log n),
\end{equation}
as in \cref{thm-lpn-public-verifier}.  Suppose fixed-noise search-LPN is hard for this
distribution in the following sense: every declared uniform polynomial
execution envelope \(\mathcal B_C\) satisfies
\begin{equation}
\label{eq-lpn-fixed-noise-search-hardness-premise}
\operatorname{Succ}_{\mathfrak M,n,\eta}(\mathcal B_C)
=\operatorname{negl}(n)
\end{equation}
at that fixed noise rate.  Then \(P\ne NP\).

In particular, any unconditional derivation of
\cref{eq-lpn-one-way-family-obstruction-bound} proves \(P\ne NP\).  Under the
present residual-ancestry premise, the derivation is supplied by
\cref{thm-lpn-residual-ancestry-attention-bound,thm-lpn-one-way}.
The implication from fixed-noise search-LPN hardness to \(P\ne NP\) itself
uses no additional complexity-barrier premise.

\end{corollary}

\begin{proof}

Suppose for contradiction that \(P=NP\).  For the threshold \(\tau\) fixed in
\cref{eq-lpn-p-np-public-verifier-schedule},
consider the prefix language

\[
L_{\mathrm{pref}}
=
\left\{
 ((A,b),u):
 \exists v\in\mathbb F_2^{\,n-|u|}
 \text{ such that }
 \lvert A(u\mathbin\Vert v)+b\rvert_1\le \tau m
\right\}.
\]
Here \(u\mathbin\Vert v\) denotes bit-string concatenation.

Membership in \(L_{\mathrm{pref}}\) is in \(NP\): the suffix \(v\) is
a polynomial-length witness, and the displayed predicate is the
polynomial-time predicate of
\hyperref[thm-lpn-public-verifier]{the public LPN verification
theorem}. Hence \(P=NP\) gives a uniform deterministic
polynomial-time decider for \(L_{\mathrm{pref}}\).

Starting with the empty prefix, query whether the current prefix
extended by \(0\) belongs to \(L_{\mathrm{pref}}\); retain \(0\) when
the answer is yes and retain \(1\) otherwise. After \(n\) queries this
procedure outputs a word accepted by the public predicate whenever one
exists. On the high-probability event in
\hyperref[thm-lpn-public-verifier]{the public LPN verification
theorem}, that word exists and is uniquely \(s\). Thus \(P=NP\) would place a
presentation-relative uniform deterministic prefix-search record with
success \(1-o(1)\) in one subenvelope \(\mathcal B_C\) of the declared
polynomial closure.  This contradicts
\cref{eq-lpn-fixed-noise-search-hardness-premise}.  Therefore \(P\ne NP\).  The final
specialization follows by applying
\cref{thm-lpn-one-way} at every polynomial ceiling covered by
\cref{thm-lpn-residual-ancestry-attention-bound}.

\end{proof}

\section{Complexity-Barrier Interpretation for LPN}
\label{sec-lpn-complexity-barriers}

The LPN barrier discussion specializes
\cref{prop-barrier-compliant-routes,rem-barrier-scope}.  Its three classical
components are indexed separately in
\cref{subsec-lpn-relativization,subsec-lpn-natural-proofs,subsec-lpn-algebrization};
the classification consequence is indexed in
\cref{subsec-lpn-oracle-blind-classification}.
These are presentation-extension tests, not exceptions to
\cref{cor-lpn-p-not-equal-np}: once fixed-noise search-LPN hardness is
established for the declared presentation, \(P\ne NP\) follows.

\subsection{Relativization}
\label{subsec-lpn-relativization}

The LPN partition gives exact obstruction criteria for its declared public
presentation and its presentation-admissible oracle interfaces.  A
conservative oracle preserves those bounds at class-preserving overhead.
The corresponding framework results are
\cref{thm-oracle-admissibility-presentation-separation,thm-framework-relativizes,cor-bgs-test};
the LPN specialization is \cref{def-lpn-admissible-oracle,thm-lpn-one-way}.

\subsection{Natural proofs}
\label{subsec-lpn-natural-proofs}

The natural-proofs test concerns an additional truth-table classifier defined
by constructivity and nonuniform usefulness. The framework classification is
\cref{thm-natural-proofs-position}, and the LPN profile and obstruction
results are \cref{thm-lpn-saturated-floor,thm-lpn-one-way}.

\subsection{Algebrization}
\label{subsec-lpn-algebrization}

A conservative algebraic interface preserves the LPN statement at
class-preserving overhead.  An interface that supplies new target-aligned
information defines an enriched LPN family, and its saturated profile records
the added source.  The corresponding framework results are
\cref{thm-oracle-admissibility-presentation-separation,thm-framework-algebrizes,prop-barrier-compliant-routes}.

Every conservative extension remains inside the same structural accounting.
An extension carrying new target-aligned information defines a different
presented family with its own saturated profile.

\subsection{Oracle-Blind Family Classification}
\label{subsec-lpn-oracle-blind-classification}

The presentation distinction in the preceding barrier tests yields the
following consequence for classifications that identify a problem with an
inadmissible solver-oracle enrichment.

\begin{corollary}[LPN witness against oracle-blind family classification]
\label{cor-oracle-blind-classification-unsound}

Adopt the hypotheses and conclusion of
\cref{cor-lpn-p-not-equal-np}.  Let
\(\mathfrak C\) be a classification of presented problem families
with the following two properties:

\begin{enumerate}
\item \(\mathfrak C\) classifies every presentation containing an
explicit uniform polynomial-time solver as polynomial-time solvable;
\item \(\mathfrak C\) assigns the same classification to a presentation
and to every enrichment of it by an inadmissible solver oracle.
\end{enumerate}

Then \(\mathfrak C\) is not sound for polynomial-time solvability.

\end{corollary}

\begin{proof}

Let \(\mathcal T_{\mathrm{LPN}}\) be the presentation of
\cref{def-lpn-task-object}, and adjoin a solver oracle
\(O_{\mathrm{LPN}}\) that returns its secret target \(s\).  The enriched
presentation \(\mathcal T_{\mathrm{LPN}}^{O_{\mathrm{LPN}}}\) has a uniform
one-query polynomial-time solver, so the first property makes
\(\mathfrak C\) classify it as polynomial-time solvable.  The oracle is
inadmissible for \(\mathcal T_{\mathrm{LPN}}\): otherwise
\cref{thm-oracle-admissibility-presentation-separation} would compile it into
the original saturated polynomial closure, contradicting
\cref{eq-lpn-fixed-noise-search-hardness-premise}.  The second property therefore assigns the same
polynomial-time label to \(\mathcal T_{\mathrm{LPN}}\), contradicting
\cref{eq-lpn-fixed-noise-search-hardness-premise}.  This is the LPN instance of
\cref{cor-presentation-coarse-no-p-np-classification}.

\end{proof}

\section{Conclusion}\label{sec-lpn-conclusion}
\label{summary-1}

The LPN specialization begins with the represented channel partition in
\cref{thm-lpn-five-term-partition-final}; the specialization is introduced
as a continuous running instance in \cref{ex-lpn-running-part5}.

The public data \((A,b,\eta)\) instantiate the
\hyperref[part:task-spaces]{Part~I task presentation},
\hyperref[part:dynamics]{Part~II component decomposition}, and
\hyperref[part:spectral-complexity]{Part~III saturated profile}; every uniform
polynomial computation using
\hyperref[def-lpn-admissible-oracle]{presentation-admissible oracle
interfaces} lies in its saturated represented closure.
The
\hyperref[thm-saturated-statistical-attack-alignment]{uniform saturated
statistical-attack theorem} specializes in
\cref{thm-lpn-saturated-statistical-alignment} to control every represented
LPN predictor on one simultaneous sample event.  Candidate recovery is
identified with clean parity alignment and transported through the noisy
sample in
\cref{cor-lpn-candidate-recovery-noise-accountability}.  The
exact execution-law identity
\eqref{eq-lpn-execution-law-recovery-alignment} identifies recovery with
average population alignment, and
\cref{thm-lpn-prospective-learning-selector-accountability} identifies
nonbaseline use of that alignment with the represented pre-hit selector and
retains its completed five-family accounting ancestry.  A prospective-source
comparison uses \cref{cor-contextual-structural-charge-transfer}; the
pointwise and family accounting bounds are
\cref{cor-lpn-core-refined-first-hit-accountability}.  The
\hyperref[cor-lpn-universal-statistical-denoising-bound]{LPN denoising
corollary} combines this statistical control with the saturated prospective
selector bound and the represented
heavy-Fourier decoder; it does not use the retrospective
\(\operatorname{HitCore}\) coordinate as a hardness premise.

The controlled-dynamics specialization evaluates the same qualified LPN
record through target capacity and drift, Caus action--aiming, the complete
hypercube functional profile, valid-Lift transfer, and the saturated
dynamical-learning envelope in
\cref{thm-lpn-controlled-target-capacity-drift,%
thm-lpn-caus-action-residual-aiming,%
thm-lpn-hypercube-functional-profile,%
thm-lpn-valid-lift-quantitative-transfer,%
cor-lpn-noise-indexed-caus-lift-interfaces,%
thm-lpn-complete-controlled-learning-specialization}.  The exact subset,
located-spectral, and threshold-guided checks are
\cref{cor-lpn-three-dynamical-stress-tests}; their dependency diagram is
\cref{fig-lpn-controlled-nonlinear-ambient-dependency}.  The resource-matched
BKW refinement is
\cref{thm-lpn-dependence-aware-bkw-recovery,thm-lpn-complete-bkw-structural-charging,cor-lpn-quantitative-improvement-over-bkw};
its numerical grid and noise frontiers are
\cref{tab-lpn-certified-bkw-grid,fig-lpn-bkw-noise-resource-frontiers}.
Valid lifted
selectors inherit the complete prospective routing by
\cref{cor-lpn-valid-lift-selector-access-seal}.
Optimal active learning is additionally represented by the stopped
occupation-flow/Bellman duality of
\cref{thm-exact-affordable-occupation-flow-primal,thm-bellman-barrier-duality-affordable-arrival}.
Its LPN specialization
\cref{thm-lpn-occupation-flow-bellman-support-certificate} gives a
modewise, parameter-explicit support bound for the best affordable
standard-presentation controller.  The learned-controller loss is then the
separate finite-sample term in
\cref{eq-lpn-complete-learning-control-lower,eq-lpn-complete-learning-control-radius}.
The structural Bayesian specialization fixes the uniform secret prior,
places \(b=As+e\) in the first target-dependent reveal slot, and charges
every likelihood, evidence normalization, belief update, and acquisition
rule.  Its target-contact and concentration bounds are
\cref{thm-lpn-quantitative-structural-bayesian-envelope}; the prospective
concentration readout retains the existing LPN five-mode decomposition, and
its deployed target gain has the explicit Add/Mult/Ord/Sym/Filt table and
noise-contracted same-record minimum in
\cref{thm-lpn-posterior-concentration-five-mode-exhaustion}.  The resulting
arrival bound is
\cref{eq-lpn-source-resolved-structural-bayesian-arrival}.  The learned
structural Bayesian loss is
\cref{cor-learned-lpn-structural-bayesian-policy}.

All LPN learning routes use the common source-resolved ledger of
\cref{thm-source-resolved-saturated-learning-compilation}: each reveal,
feature, posterior update, acquisition rule, controlled transition, and
deployment map retains its Add/Mult/Ord/Sym/Filt charge, binary contraction
\((1-2\eta)^2\), sample size \(m\), confidence level, horizon, and complete
cost.  The passive and sequential specialization is
\cref{eq-lpn-saturated-certificate-source-resolved-learning-candidate}, the
sparse-reward observation specialization is
\cref{thm-lpn-complete-sparse-reward-observation-accounting}, and the active
learning specialization is
\cref{thm-lpn-complete-master-active-sampling-envelope}.  Each contributes a
same-record target-contact candidate; none replaces the complete prospective
source evaluation.

The noise--resource comparison in \cref{sec-lpn-noise-resource-phase} gives
the exact-binomial error-enumeration lower-bound curve, the Rademacher
sample-resolution endpoint, the saturated computational obstruction, and
the sharp information threshold.  Their finite-budget relations are
collected in \cref{thm-lpn-three-threshold-phase-comparison};
\cref{cor-lpn-quantitative-noise-phase-localization,fig-lpn-three-noise-phases}
locate the vanishing-noise nonnegligible-success window, the fixed-positive
hardness region, and the information-impossible region.

The \hyperref[thm-channel-exhaustiveness]{generic component decomposition},
\hyperref[thm-lpn-no-independent-geom-final]{LPN geometry theorem},
\hyperref[thm-lpn-coupled-source-seal-final]{coupled-source theorem}, and
\hyperref[cor-lpn-abs-source-exhaustion]{LPN exact-\(\Abs\) source-exhaustion
corollary} reduce the prospective \(\Geom/\Abs\) source profile to the
ambient Geom coordinate.  Before the remaining compact ambient premise is
invoked, the constant-noise spectral perimeter
\cref{thm-lpn-exact-constant-noise-spectral-perimeter} exponentially closes
materialized located spectra and applies the exact terminal-rank gate to
residual-parity fiber maps.  The exact-\(\Abs\) and all quotient-supported
Geom coordinates vanish, and
\hyperref[cor-lpn-complete-ambient-geom-source-seal]{the complete ambient
source seal}, using
\cref{thm-lpn-residual-ancestry-attention-bound}, bounds the remaining coordinate by
\(2^{-n}+(1-\eta)^n\).  The only formal premise in this chain is therefore
\cref{thm-lpn-residual-ancestry-attention-bound}.  It is discharged by a
computer-assisted record exactly when the finite verification fields of
\cref{thm-lpn-proof-grade-empirical-residual-discharge} pass; no sampled
failure to find a source is substituted for saturated source coverage.  The
present manuscript proves this verification implication and specifies its
complete record, but it does not supply an accepted instantiated certificate
for the residual-ancestry cell.  Thus the fixed-noise hardness and separation
consequences remain conditional in this version on that internal certificate
obligation.  Supplying and deterministically verifying the record would
discharge the sole remaining premise without importing an LPN or
complexity-theoretic hardness assumption.  The
\hyperref[cor-no-retrospective-first-hit-source-loophole]{no retrospective
first-hit source loophole corollary} excludes a completed first-hit label as
an uncharged source and classifies every independent pre-hit representation
inside that same closure. These results prove the exact source and first-hit
bounds as separate ledgers in
\hyperref[thm-lpn-saturated-floor]{the saturated LPN ledger theorem}.

By \cref{thm-lpn-residual-ancestry-attention-bound}, the
\hyperref[thm-lpn-one-way]{search-LPN hardness theorem}
applies the complete saturated-profile bound and prospective selector accountability to
obtain the explicit negligible fixed-noise success bound, including
presentation-admissible LPN oracle interfaces.  This yields the
\hyperref[cor-lpn-p-not-equal-np]{separation consequence for
\(P\ne NP\)} and
\hyperref[cor-oracle-blind-classification-unsound]{oracle-blind
classification corollary}.  The roles of relativization, natural proofs,
and algebrization in these presentation-relative statements are indexed in
\cref{sec-lpn-complexity-barriers}.

\appendix
\section{Index of Notation}\label{sec-index-of-notation}
\label{index-of-notation}

The index records the load-bearing notation and the result that fixes
each symbol. The final column points to its first substantive use.

\begin{longtable}{@{}
  >{\raggedright\arraybackslash}p{0.15\textwidth}
  >{\raggedright\arraybackslash}p{0.34\textwidth}
  >{\raggedright\arraybackslash}p{0.20\textwidth}
  >{\raggedright\arraybackslash}p{0.20\textwidth}@{}}
\toprule
Symbol & Meaning & Defined in & First used in \\
\midrule
\endfirsthead
\toprule
Symbol & Meaning & Defined in & First used in \\
\midrule
\endhead
\multicolumn{4}{@{}l}{\emph{Spaces, algebras, and presentations}}\\
\addlinespace[2pt]
\(\mathcal T_I\) & Presented task on instance \(I\) &
  \cref{def-task-space-signature} & \cref{part:task-spaces} \\
\(X_I\) & Complete candidate-state space &
  \cref{def-task-space-signature} & \cref{the-task-space-signature} \\
\(\operatorname{Obs}_I\) & Declared primitive observable interface &
  \cref{def-task-space-signature} & \cref{the-task-space-signature} \\
\(\operatorname{Dyn}_I\) & Declared static relation slot &
  \cref{def-task-space-signature} & \cref{the-task-space-signature} \\
\(\operatorname{End}_I\) & Positive, negative, and run-dependent terminal data &
  \cref{def-task-space-signature} & \cref{the-task-space-signature} \\
\(\obsalg_I\) & Observable event algebra generated by the public interface &
  \cref{def-observable-algebra} & \cref{the-observable-algebra} \\
\(\mathfrak I_n\) & Size-\(n\) slice of a task family &
  \cref{def-task-family} & \cref{families-and-instances} \\
\(T_I\) & Positive terminal sector or target set &
  \cref{def-task-space-signature} & \cref{def-usefulness-environment} \\
\(y_I\) & Centered target contrast &
  \cref{def-generated-algebra-target-projection} &
  \cref{def-generated-algebra-target-projection} \\
\(\sem{\rho}_I\) & Denotation of represented derivation \(\rho\) on \(I\) &
  \cref{def-budgeted-derivability-closure} &
  \cref{def-budgeted-derivability-closure} \\
\addlinespace[5pt]
\multicolumn{4}{@{}l}{\emph{Dynamics and structural channels}}\\
\addlinespace[2pt]
\(L\) & Positive search generator & \cref{thm-finite-generator-form} &
  \cref{part:dynamics} \\
\(\mu_I\) & Declared reference measure for state and edge Hilbert spaces &
  \cref{def-reference-measure-legality} &
  \cref{def-reference-measure-legality} \\
\(H\) & Charged stopping horizon for a finite transcript carrier &
  \cref{def-resource-envelope} & \cref{def-full-configuration-lift} \\
\(\Geom\) & Reversible local geometric component &
  \cref{thm-channel-exhaustiveness} & \cref{part:dynamics} \\
\(\Caus\) & Authorized directed or drift component &
  \cref{thm-channel-exhaustiveness} & \cref{def-caus-authority-envelope} \\
\(\Abs\) & Qualified target-bearing quotient/lumping source; algebraic
  quotient provenance is recorded before qualification &
  \cref{def-paper1-abs-usefulness-gate,def-abs-extraction} &
  \cref{def-abs-lift-split,def-source-accounting-profile-separation} \\
\(\Lift\) & Valid represented lifted-state component &
  \cref{thm-channel-exhaustiveness} & \cref{def-valid-lift} \\
\(\Bdry\) & Boundary, killing, and terminal-readout component &
  \cref{thm-channel-exhaustiveness} & \cref{thm-channel-exhaustiveness} \\
\(\Res\) & Residual complement of the declared channel projections &
  \cref{thm-channel-exhaustiveness} &
  \cref{thm-exhaustive-residual-normal-form} \\
\(\Delta_e^{(m)}\) & Weighted single-edge geometric Laplacian &
  \cref{thm-exhaustive-geom-normal-form} &
  \cref{thm-exhaustive-geom-normal-form} \\
\(\kappa_x,P_\kappa\) & Valid-Lift fiber law and its averaging projection &
  \cref{thm-exhaustive-lift-normal-form} &
  \cref{thm-exhaustive-lift-normal-form} \\
\(L_{\mathrm{kill}}^{\Bdry},L_{\mathrm{val}}^{\Bdry}\) &
  Homogeneous loss and represented value blocks of the boundary channel &
  \cref{thm-exhaustive-bdry-normal-form} &
  \cref{thm-exhaustive-bdry-normal-form} \\
\(L_{\mathrm{dir}}^{\Res}\) & Hodge-orthogonal directed residual coordinate &
  \cref{thm-exhaustive-residual-normal-form} &
  \cref{cor-no-residual-provenance-outside-normal-form} \\
\(L_{\mathrm{nl}}^{\Res}\) & Quotient/lift-orthogonal nonlocal residual coordinate &
  \cref{thm-exhaustive-residual-normal-form} &
  \cref{cor-no-residual-provenance-outside-normal-form} \\
\(L^{\Res,B}\) & Budget-filtered residual ledger at budget \(B\) &
  \cref{thm-channel-exhaustiveness} &
  \cref{thm-exhaustive-residual-normal-form} \\
\(U,P,L_Q\) & Quotient lift, projection, and induced generator &
  \cref{def-quotient-lifts} & \cref{def-lumpability-defect} \\
\(E_q\) & Lumpability defect \(LU-UL_Q\) &
  \cref{def-lumpability-defect} & \cref{def-lumpability-defect} \\
\(\Phi\) & Represented Dirichlet geometry &
  \cref{def-static-dirichlet-geometry} & \cref{def-geom-extraction} \\
\(D\) & Target-normalized progress observable &
  \cref{def-geom-extraction} & \cref{def-geom-extraction} \\
\(\mathfrak A\) & Declared causal-authority envelope &
  \cref{def-caus-authority-envelope} &
  \cref{thm-caus-authority-no-substitution} \\
\addlinespace[5pt]
\multicolumn{4}{@{}l}{\emph{Costs, budgets, and saturation}}\\
\addlinespace[2pt]
\(\mathcal B\) & Resource structure, cost order, ceiling, and transfer class &
  \cref{def-budget-structure} & \cref{def-budget-structure} \\
\(b(n)\) & Declared resource ceiling at size \(n\) &
  \cref{def-budget-structure} & \cref{part:dynamics} \\
\(\mathcal R_{n,\le B}^{\mathrm{sat}}\) & Saturated represented objects of cost at most \(B\) &
  \cref{def-saturated-representation-degree} &
  \cref{def-saturated-representation-degree} \\
\(\mathcal D_{n,\le B}\) & Budgeted derivability closure &
  \cref{def-budgeted-derivability-closure} &
  \cref{def-budgeted-derivability-closure} \\
\(\satcost{R}\) & Least scalar saturated representation cost of \(R\) &
  \cref{def-saturated-representation-degree} &
  \cref{def-saturated-representation-degree} \\
\(\Psi\) & Class-preserving accounting enlargement &
  \cref{thm-universal-saturated-profile-accountability} &
  \cref{thm-dynamic-compensation-first-hit-accountability} \\
\(\lambda\) & Discount parameter for killed target arrival &
  \cref{def-stopwatch-score} & \cref{def-stopwatch-score} \\
\(B^*(n),d^*(n)\) & Saturated threshold budget and its log-cost coordinate &
  \cref{def-degree-invariant} & \cref{def-degree-invariant} \\
\(\mathfrak P_{\mathfrak D,n}^{>}(\theta)\) & Strict Pareto region of budgets
  supporting a complete record with gain greater than \(\theta\) &
  \cref{def-strict-structural-meta-frontier} &
  \cref{thm-strict-meta-profile-generalized-inverse} \\
\(\operatorname{SMC}_{\mathfrak D,n,\kappa}^{>}(\theta)\) & Scalar
  structural meta-complexity at strict gain threshold \(\theta\) &
  \cref{def-strict-structural-meta-frontier} &
  \cref{thm-structural-meta-accountability} \\
\(B_{\eta}^{*,\mathrm{src}}(n)\) & Prospective source-profile threshold cost &
  \cref{eq-prospective-source-threshold-cost} &
  \cref{cor-existing-saturated-threshold-as-metacomplexity} \\
\(\operatorname{RepMC}\),\newline\(\operatorname{ExtMC}\) & Exact-denotation
  representation cost and complete extraction cost &
  \cref{def-exact-denotation-meta-frontiers} &
  \cref{prop-meta-extraction-output-projection} \\
\(\operatorname{CC}_{\mathrm{meta}}\) & Minimum gate-count coordinate of
  an exact circuit representation &
  \cref{def-truth-table-circuit-presentation} &
  \cref{thm-circuit-minimization-specialization} \\
\(K_{\mathsf U}^{c,\mathrm{mode}}\),\newline
\(\operatorname{MinK}_{\mathsf U}^{c,\mathrm{mode}}\) & Universal-program
  representation cost and its threshold language &
  \cref{def-declared-universal-program-metacomplexity} &
  \cref{thm-program-minimization-specialization} \\
\(\operatorname{GapSMC}\) & Resource-and-gain promise version of a
  structural meta-problem &
  \cref{def-gap-structural-meta-problem} &
  \cref{def-gap-structural-meta-problem} \\
\(B^*_{\vartheta}\)\newline\((n;\mathrm{LPN}_\eta)\) & Visibility-indexed
  saturated threshold cost for the LPN family &
  \cref{eq-lpn-visibility-threshold-cost} &
  \cref{thm-lpn-three-threshold-phase-comparison} \\
\(B_{\mathrm{rec}}^*\)\newline\((n,m,\eta;\alpha)\) & Infimum scalarized saturated cost
  reaching recovery probability \(\alpha\) &
  \cref{def-noise-indexed-saturated-recovery-cost} &
  \cref{thm-three-noise-threshold-comparison} \\
\(B_{\mathrm{align}}^*\)\newline\((n,m,\eta;\theta)\) & Infimum scalarized
  cost reaching average population alignment \(\theta\) &
  \cref{eq-noise-indexed-population-alignment-cost} &
  \cref{cor-noise-indexed-learning-threshold-criterion} \\
\(\eta_{\mathrm{comp}}^{\mathrm{sat}}\) & Supremal affordable noise at a
  fixed scalar budget and success level; an interval endpoint under flip
closure &
  \cref{eq-noise-indexed-computational-threshold} &
  \cref{thm-lpn-three-threshold-phase-comparison} \\
\addlinespace[5pt]
\multicolumn{4}{@{}l}{\emph{Profiles and extraction coordinates}}\\
\addlinespace[2pt]
\(V_I(q)\) & Target-conditioned visibility of exact quotient \(q\) &
  \cref{def-abs-extraction} & \cref{def-abs-extraction} \\
\(A_T(\lambda)\) & Discounted target-arrival functional &
  \cref{def-arrival-functional,def-stopwatch-score} &
  \cref{def-stopwatch-score} \\
\(m_I^{\mathrm{sat}}(B)\) & Saturated exact-quotient visibility profile &
  \cref{def-abs-extraction} & \cref{def-abs-extraction} \\
\(G_{X,n}^{\mathrm{sat}}(B)\) & Ambient Geom/Caus extraction coordinate &
  \cref{def-geom-extraction} & \cref{def-geom-extraction} \\
\(G_n^{\mathrm{sat}}(B)\) & Active-geometry master subprofile &
  \cref{def-master-saturated-quotient-geometry-profile} &
  \cref{def-master-saturated-quotient-geometry-profile} \\
\(I_n^{\mathrm{sat}}(B)\) & Full two-coordinate saturated structural profile &
  \cref{def-extraction-functional} &
  \cref{thm-full-saturated-profile-decomposition} \\
\(\mathfrak U_n^{\mathrm{sat}}(B)\) & Master compensated quotient--geometry profile &
  \cref{def-master-saturated-quotient-geometry-profile} &
  \cref{thm-affordable-geom-abs-normal-form} \\
\(\mathscr C,V_{\mathscr C}\) & Master compensated certificate and its visibility &
  \cref{def-master-compensated-quotient-geometry-certificate} &
  \cref{def-master-compensated-quotient-geometry-certificate} \\
\(R_T(D)\) & Fraction of target contrast reachable from monotone functions of \(D\) &
  \cref{def-geom-extraction} & \cref{def-geom-extraction} \\
\(\rho_\Phi(D)\) & Dirichlet regularity of the progress observable &
  \cref{def-geom-extraction} & \cref{def-geom-extraction} \\
\(C_{\Phi,\mathfrak K}^{\mathrm{auth}}(D)\) & Authority-compatible causal alignment factor &
  \cref{def-geom-extraction} & \cref{def-authority-compatible-caus-alignment} \\
\(\operatorname{Pert}(\mathfrak g)\) & Authorized ambient perturbation data retained by an ambient source &
  \cref{def-authorized-ambient-perturbation-witness} &
  \cref{thm-ambient-perturbation-characterization-obstruction} \\
\(P_\Phi\) & Edge-restricted sample-and-hold kernel of an ambient perturbation record &
  \cref{eq-authorized-ambient-perturbation-kernel} &
  \cref{def-authorized-ambient-perturbation-witness} \\
\(E_\downarrow(\mathfrak g)\) & Exposed edges on which the authorized current descends the represented potential &
  \cref{def-authorized-ambient-perturbation-witness} &
  \cref{thm-ambient-perturbation-characterization-obstruction} \\
\(\operatorname{ch}_I(E)\) & Contextual target-aligned charge of an algebra extension &
  \cref{def-affordable-conditional-algebra-frontier} &
  \cref{thm-exact-affordable-algebra-structural-charging} \\
\(\operatorname{CoreSig}(\mathscr C)\) & Charge-bearing source signature of a
  complete normalized record &
  \cref{def-charge-bearing-core-signature} &
  \cref{thm-operational-core-provenance-separation} \\
\(\operatorname{CoreAct}^{\mathrm{sat}}(B)\) & Qualified source visibility in
  one operational-signature and charge-bearing-core cell &
  \cref{eq-core-activity-refined-cell} &
  \cref{thm-operational-core-provenance-separation} \\
\(\operatorname{HitCore}^{\mathrm{sat}}(B)\) & Core-refined exact first-hit
  accounting profile; not a prospective source profile &
  \cref{def-core-refined-first-hit-accounting-profile} &
  \cref{thm-core-refined-prospective-accountability} \\
\(\operatorname{Succ}_{\mathfrak M,n}(B)\) & Budget-\(B\) success profile under family functional \(\mathfrak M_n\) &
  \cref{def-uniform-saturated-success-profile} &
  \cref{thm-universal-saturated-profile-accountability} \\
\(\operatorname{Succ}^{\kappa}_{\mathfrak M,n}(s)\) & Scalar-sublevel
  saturated success profile &
  \cref{eq-scalarized-affordable-success-profile} &
  \cref{def-noise-indexed-saturated-recovery-cost} \\
\(\Gamma,J_{\Gamma,\ell}\) & Affordable frontier and its successive joint interfaces &
  \cref{def-affordable-conditional-algebra-frontier} &
  \cref{def-affordable-conditional-algebra-frontier} \\
\addlinespace[5pt]
\multicolumn{4}{@{}l}{\emph{Learning and kernel quantities}}\\
\addlinespace[2pt]
\(k_I,\mathcal H_{k_I}\) & Represented kernel and its RKHS &
  \cref{def-feature-rkhs-observable} & \cref{def-feature-rkhs-observable} \\
\(N_{\mathrm{eff}}(\lambda)\) & Effective dimension at ridge scale \(\lambda\) &
  \cref{def-effective-dimension} & \cref{def-effective-dimension} \\
\(D_{c,n},a_{c,n},e_{c,n}\) & Certified source-cell capacity,
  target-transfer attenuation, and factorization/tail/numerical error &
  \cref{def-scale-certified-attenuating-source-audit} &
  \cref{thm-empirical-asymptotic-saturated-source-closure} \\
\(\beta_c,\gamma_c,p_c\) & Capacity exponent, attenuation exponent, and
  subexponential scale factor in a verified source audit &
  \cref{thm-empirical-asymptotic-saturated-source-closure} &
  \cref{eq-empirical-asymptotic-source-profile-bound} \\
\(\lambda_\eta,N_*\) & Fixed-noise LPN target exponent and certified
  finite-to-asymptotic crossover size &
  \cref{eq-lpn-empirical-discharge-lambda,eq-lpn-empirical-discharge-crossover-monotonicity} &
  \cref{thm-lpn-proof-grade-empirical-residual-discharge} \\
\(A_m\) & Empirical kernel--target alignment &
  \cref{def-kernel-target-alignment} & \cref{def-kernel-target-alignment} \\
\(\mathcal H_{n,B,m}^{\mathrm{sat,stat}}\) & Sample-independent saturated statistical readout class &
  \cref{def-budget-saturated-statistical-readout-class} &
  \cref{thm-saturated-statistical-attack-alignment} \\
\(M_{I,n}^{\mathrm{sat,stat}}\)\newline\((B,m)\) & Saturated learned target-projection envelope &
  \cref{def-budget-saturated-statistical-readout-class} &
  \cref{thm-saturated-statistical-attack-alignment} \\
\(\operatorname{Align}_{\mathfrak M,n}^{\mathrm{sat,stat}}\) & Average-sample,
  average-seed population-alignment profile &
  \cref{def-average-seed-population-alignment-profile} &
  \cref{thm-population-alignment-generalization-memorization} \\
\(\operatorname{GenProj}_{\mathfrak M,n}^{\mathrm{sat,stat}}\) & Average
  generalizing target-projection profile &
  \cref{def-average-seed-population-alignment-profile} &
  \cref{thm-population-alignment-generalization-memorization} \\
\(\operatorname{Enum}_{\mathfrak M,n,\eta}^{\mathrm{sat}}\) & Saturated neutral-contact
  profile for the declared baseline &
  \cref{eq-noise-indexed-neutral-contact-profile} &
  \cref{thm-noise-indexed-learning-accountability} \\
\(\operatorname{Sel}_{\mathfrak M,n,\eta}^{\mathrm{sat}}\) & Saturated
  family-uniform prospective-selector profile &
  \cref{eq-noise-indexed-family-selector-profile} &
  \cref{thm-noise-indexed-learning-accountability} \\
\(\Gamma_{n,B,m,\delta}^{\mathrm{sat,stat}}\) & Uniform selection radius of the saturated readout class &
  \cref{def-budget-saturated-statistical-readout-class} &
  \cref{thm-saturated-statistical-attack-alignment} \\
\(\begin{gathered}\mathcal R^{\rm att,sat},\\
\operatorname{App}^{\rm sat}\end{gathered}\) &
  Saturated attainable clean risk and structural approximation floor &
  \cref{def-saturated-attainable-risk-approximation-floor} &
  \cref{thm-complete-quantitative-learning-decomposition} \\
\(N_{n,B,m}^{\mathrm{out}}\) & Number of complete saturated output records
  at a locally finite scalar budget &
  \cref{cor-finite-description-saturated-learning-radius} &
  \cref{cor-finite-description-saturated-learning-radius} \\
\(H_2(u)\) & Binary entropy in bits &
  \cref{eq-binary-entropy-definition} &
  \cref{thm-three-noise-threshold-comparison} \\
\(D_{\mathrm{Ber}}(u\Vert v)\) & Bernoulli relative entropy in nats &
  \cref{eq-bernoulli-divergence-definition} &
  \cref{thm-lpn-public-verifier} \\
\(d_{\mathrm{eff}}(\varepsilon)\) & Spectral prefix size capturing target mass \(1-\varepsilon\) &
  \cref{def-label-spectral-profile} & \cref{def-label-spectral-profile} \\
\(\mathrm{KL}^{\mathsf{str}}_{E,H}\) & Defect-restricted structural PAC--Bayes divergence &
  \cref{def-structural-pac-bayes-kl} & \cref{def-structural-pac-bayes-kl} \\
\(\beta(k),\mu_a,\delta_a\) & Absolute-regularity coefficient, retained
  block count, and post-coupling confidence level &
  \cref{def-dynamical-sample-certificate} &
  \cref{thm-beta-mixing-structural-blocking} \\
\(\mathfrak R_T^{\mathrm{seq}},\mathfrak q_T(\mathbf A)\) & Sequential
  structural Rademacher complexity and path-operator inverse-precision radius &
  \cref{def-sequential-structural-rademacher} &
  \cref{lem-sequential-energy-support} \\
\(\Gamma_{\rm dyn}^{\rm sat}\),
\(\operatorname{Exc}_{\rm corr}^{\rm sat}\)\newline
\(U_{\rm noise}^{\rm clean}\) & Uniform affordable dynamical deviation,
  corrupted excess, and clean binary-risk envelopes in \((n,b(n),m,\eta)\) &
  \cref{def-saturated-dynamical-noise-profiles} &
  \cref{thm-uniform-dynamical-label-noise-bound} \\
\(m_{\rm noise}^{\rm sat}\)\newline
\(\eta_{\rm stat,dyn}^{\rm sat}\),
\(\eta_{\rm comp,dyn}^{\rm sat}\) & Dynamical sample, statistical-noise,
  and computational-noise thresholds &
  \cref{def-dynamical-noise-thresholds} &
  \cref{thm-dynamical-noise-three-thresholds} \\
\(\operatorname{Err}_{q,\mathsf N}^{\rm sat}\) & Affordable reward,
  execution, observation, filter, or composite-channel error in its declared
  norm, indexed by \((n,b(n),m,\delta)\) &
  \cref{def-saturated-dynamical-channel-error-profiles} &
  \cref{thm-dynamical-affine-reward-channel-transport} \\
\(\varepsilon_{\rm obs}^{\rm rec},\tau_m\) & Observation reconstruction
  error and finite-row channel-estimation radius &
  \cref{def-dynamical-observation-reconstruction,thm-finite-alphabet-channel-estimation-entropy} &
  \cref{eq-observation-reconstruction-function-bound,eq-finite-channel-row-tv-estimation} \\
\(R_{\rm BB},B_{\rm BB},Q\) & Complete trained-model black-box audit,
  its charged cost, and its total oracle-query count &
  \cref{def-complete-black-box-predictor-audit} &
  \cref{thm-finite-query-simultaneous-black-box-audit} \\
\(\mathfrak C_{\rm BB}\) & Simultaneous black-box risk, noise, structural,
  geometric, quotient, and residual signature &
  \cref{def-complete-quantitative-black-box-signature} &
  \cref{thm-end-to-end-black-box-structural-certificate} \\
\(D_a(f),A_a(f)\) & Behavioral defect and target response of source
  intervention \(a\) &
  \cref{def-black-box-structural-intervention-ledger} &
  \cref{thm-simultaneous-black-box-structural-signature} \\
\(\mathcal E_{P,f},\mathcal V_f,\rho_{P,f}\) & Black-box output
  Dirichlet energy, variance, and Rayleigh ratio &
  \cref{def-black-box-perturbation-geometry} &
  \cref{thm-black-box-geometric-confidence-region} \\
\(q_f,\delta_{q_f}(P)\) & Output discretizer quotient and its perturbation-row
  lumpability defect &
  \cref{def-black-box-output-quotient} &
  \cref{thm-black-box-lumpability-compilation} \\
\addlinespace[5pt]
\multicolumn{4}{@{}l}{\emph{Controlled dynamical systems}}\\
\addlinespace[2pt]
\(P^a,P^\pi,L^\pi\) & Action kernel, selected policy kernel, and normalized generator &
  \cref{def-controlled-dynamical-presentation} &
  \cref{thm-controlled-policy-channel-decomposition} \\
\(\begin{gathered}K^A,P^{\pi,K},\\ r^{\pi,K}\end{gathered}\) & Execution channel and actual selected
  transition and reward &
  \cref{def-controlled-action-execution-channel} &
  \cref{thm-controlled-execution-channel-comparison} \\
\(\begin{gathered}\Omega_{\rm POMDP}^{[H]},b_t,\\ \widehat b_t,e_t\end{gathered}\) & Clocked stochastic
  POMDP carrier, analytic belief, represented filter, and filter error &
  \cref{def-represented-stochastic-pomdp-record,def-represented-pomdp-filter-comparison} &
  \cref{thm-controlled-pomdp-clocked-embedding} \\
\(\begin{gathered}\phi,q_\varepsilon,\rho_\phi(P),\\
M_\varepsilon(\phi),J_{R,t}^{\rm fib}\end{gathered}\) & Represented
  sparse-reward observation, resolution, regularity, cell count, and
  conditional within-cell target information &
  \cref{def-complete-sparse-reward-observation-record} &
  \cref{thm-observation-regularity-sparse-reward-exploration} \\
\(\begin{gathered}B_{\rm act},\\
\operatorname{ActArr}^{\rm sat},\\
\operatorname{ActSel}^{\rm sat},\\
B_{\rm act}^*,m_{\rm act}^*,\eta_{\rm act}^{\rm sat}\end{gathered}\) & Complete master
  active-sampling cost, optimal arrival and selector profiles, and their
  computational, sample, and noise frontiers &
  \cref{def-complete-master-active-sampling-record,def-family-uniform-active-saturated-profile,def-active-construction-thresholds} &
  \cref{thm-active-optimality-saturated-construction,thm-quantitative-active-optimum-monotonicity} \\
\(\begin{gathered}\mathfrak O_{I,H}^{\rm sat}(B),\\
\mathfrak D_{I,t}^{\rm sat}(B;v)\end{gathered}\) &
  Affordable stopped occupation-flow set and saturated Bellman support &
  \cref{def-stopped-occupation-flow-complete-active-record,def-affordable-bellman-support} &
  \cref{thm-exact-affordable-occupation-flow-primal,thm-bellman-barrier-duality-affordable-arrival} \\
\(\begin{gathered}\Pi_t,\widehat\Pi_t,\\
\operatorname{PostConc}^{\rm sat},\\
\operatorname{PostSrc}^{\rm sat},\\
\operatorname{BayesGain}^{\rm sat},\\
\operatorname{BayesArr}^{\rm sat}\end{gathered}\) &
  Analytic task posterior, represented structural belief, affordable
  concentration, prospective concentration-source, actionable gain, and
  Bayesian arrival profiles &
  \cref{def-complete-structural-bayesian-record,def-affordable-structural-posterior-concentration,def-source-resolved-actionable-bayesian-gain} &
  \cref{thm-structural-bayesian-target-contact-accountability,thm-exact-source-resolved-structural-bayesian-gain,thm-optimal-structural-bayesian-acquisition} \\
\scalebox{.82}{\(\begin{gathered}Z_t,b_{R,t}^{\rm lat},I_{R,t}^{\rm lat},\\
\Delta_{R,t}^{\rm suf},e_{R,t}^{\rm dec},\\
\operatorname{LatBayesArr}^{\rm sat},\\
\operatorname{LatBayesGain}^{\rm sat},\\
\operatorname{LatAttArr}^{\rm sat}\end{gathered}\)} &
  Compact latent state, code capacity, retained target information,
  sufficiency and decoder defects, useful belief gain, and latent
  Bayesian/attention arrival profiles &
  \cref{def-complete-compact-latent-belief-record,def-saturated-compact-latent-belief-profile} &
  \cref{thm-compact-latent-information-sufficiency,thm-optimal-compact-latent-acquisition};
  \cref{thm-joint-compact-latent-attention-acquisition} \\
\(\begin{gathered}C_{\mathcal Q}(K),\\
\vartheta_{\rm KL}(K;\mathcal Q)\end{gathered}\) & Conditional channel capacity
  in bits and KL strong data-processing coefficient &
  \cref{def-controlled-channel-information-ledger} &
  \cref{thm-controlled-multichannel-information-bound} \\
\(\begin{gathered}I_0^{(m)},C_{\rm pub}^{[H]},\\ D_{\min}(C)\end{gathered}\) & Training-information bound,
  deployment information budget, and inverse rate--distortion floor &
  \cref{eq-controlled-training-information-budget,eq-controlled-adaptive-multichannel-information} &
  \cref{eq-controlled-rate-distortion-inverse-floor} \\
\(\begin{gathered}\mathsf H_{\rm init}^{[H]},\mathsf H_{\rm ch}^{[H]},\\
\mathsf H_{\rm env}^{[H]},\\
\mathsf H_R^{[H]},H_\pi^{[H]},\\
\mathbf H_{\rm ctl}\end{gathered}\) & Initial retained-state,
  physical-channel, stochastic-transition, clean-reward, and policy entropy,
  and the typed,
  noninterchangeable controlled-entropy vector &
  \cref{def-controlled-channel-information-ledger,def-controlled-typed-entropy-outputs} &
  \cref{thm-controlled-belief-entropy-balance} \\
\(u_h^\pi,u_\beta^\pi\) & Finite-horizon and discounted reach--avoid committors &
  \cref{def-controlled-reach-avoid-committor} &
  \cref{thm-controlled-committor-bellman} \\
\(\operatorname{Cap}_\Phi(S,T)\) & Represented condenser capacity &
  \cref{def-controlled-condenser-capacity} &
  \cref{thm-controlled-reach-avoid-dirichlet-drift} \\
\(\gamma_\Phi,N_{\rm eff,\Phi}(\lambda)\) & Certified gap and effective dimension of a dynamical carrier &
  \cref{def-certified-dynamical-geometry} &
  \cref{thm-controlled-certified-mixing} \\
\(\mathfrak F_\Phi,\Lambda_\Phi^{\rm mix}\) & Audited functional-inequality
  profile and its spectral-profile coordinate &
  \cref{def-controlled-functional-inequality-profile} &
  \cref{thm-controlled-functional-inequality-consequences} \\
\(\mathcal A_{\Caus}\)\newline
\(\mathcal E_{\rm aim}\)\newline
\(\mathcal P_{\Caus}\) &
  Caus action, target-aiming energy, and authorized target power &
  \cref{def-controlled-complete-caus-flow-record} &
  \cref{thm-controlled-caus-action-aiming} \\
\(\operatorname{Act}^{\rm sat}_{\Caus}\)\newline
\(\operatorname{Aim}^{\rm sat}_{\Caus}\)\newline
\(\operatorname{Flow}^{\rm sat}_{\Caus}\) & Saturated affordable Caus-flow profiles &
  \cref{def-controlled-saturated-caus-action-profiles} &
  \cref{cor-controlled-uniform-caus-flow-bound} \\
\(U_{\Caus}^{\rm noise}\) & Clean-target Caus contact bound learned through
  corrupted binary supervision &
  \cref{def-controlled-clean-aiming-comparison} &
  \cref{thm-controlled-noise-corrected-caus-aiming} \\
\(Q_\pi,H_{F,0},\rho_\pi\) & Orthogonal fiber projection,
  boundary-zero fiber form, and lift leakage ratio &
  \cref{def-controlled-complete-form-lift-record} &
  \cref{thm-controlled-lift-capacity-schur} \\
\(\operatorname{LiftAmp}^{\rm sat}\)\newline
\(\operatorname{LiftSel}^{\rm sat}\)\newline
\(\operatorname{LiftAcc}^{\rm sat}\) & Saturated capacity-amplification,
  selector, and discounted-access profiles for valid lifts &
  \cref{def-controlled-saturated-lift-capacity-amplification,def-controlled-saturated-lift-access-envelope} &
  \cref{thm-controlled-valid-lift-block-resolvent,thm-controlled-clocked-lift-transfer} \\
\(\varepsilon_{T,F,H}^{\rm dep}\),
\(\varepsilon_{T,F,H}^{\rm flip}\)\newline
\(\varepsilon_T^\mu,\varepsilon_F^\mu\),
\(\omega_{\Phi,T,F}^{\rm Cap}\) & Deployment, direct-flip, reference-law,
  and capacity-stability errors for learned target and failure gates &
  \cref{def-controlled-noisy-gate-profiles} &
  \cref{thm-controlled-noisy-gate-lift-transfer} \\
\(N_{\rm eff}^{\mathcal A,\mathrm{safe}},
  \mathbf M_{D,H}^{\rm dyn}\) & Action-safe effective dimension and
  defect/harmonic control precision &
  \cref{def-action-safe-defect-precision} &
  \cref{thm-constrained-safe-spectral-capacity} \\
\(\mathbf U_{\rm dyn}\) & Coordinatewise dynamical-system certificate &
  \cref{def-coordinatewise-dynamical-certificate} &
  \cref{thm-universal-dynamical-certificate} \\
\(\mathbf U_{\rm learn}^{\rm dyn}\) & Typed learned-control certificate:
  model, geometry, value, policy, shift, binary and general channel noise,
  POMDP, information, entropy, safety, arrival, noisy capacity, return, and
  regret &
  \cref{def-learned-control-certificate-coordinates} &
  \cref{thm-master-structural-dynamical-learning-certificate} \\
\(\mathfrak G_{\rm dyn}\),
\(\mathfrak G_{\rm dyn}^{\rm clean}\)\newline
\(V_B^*,\Delta_B^{\rm ctrl}\)\newline
\(\Delta_B^{\rm POMDP}\)\newline
\(\overline\Delta_B^{\rm ctrl}\)\newline
\(\overline\Delta_B^{\rm POMDP}\) & Dynamical deviation
  and clean binary-risk envelopes, budget-reachable value, computational
  approximation gap, partial-observation restriction gap, and their
  represented operational upper certificates &
  \cref{def-dynamical-generalization-envelope} &
  \cref{thm-master-structural-dynamical-learning-certificate} \\
\(E_U=LU-U\widehat L\) & Represented generator-intertwining defect &
  \cref{eq-controlled-generator-intertwining-defect} &
  \cref{thm-controlled-continuum-lift-transfer} \\
\addlinespace[5pt]
\multicolumn{4}{@{}l}{\emph{Optimal structural extraction}}\\
\addlinespace[2pt]
\(U_{R,t}^{\rm opt}\) & Same-record minimum of information, source, Caus,
  Lift, boundary, and Bellman target-gain bounds &
  \cref{def-optimal-structural-extraction-certificate} &
  \cref{thm-optimal-structural-extraction-saturated-ceiling} \\
\(J_{R,t}^{\rm Bayes,src}\) & Source-resolved ordered-reveal information in
  bits & \cref{eq-source-resolved-bayesian-information-ledger} &
  \cref{cor-noise-sample-target-size-optimal-extraction} \\
\addlinespace[5pt]
\multicolumn{4}{@{}l}{\emph{Attention compilation}}\\
\addlinespace[2pt]
\(\Theta_{\rm att}\) & Complete depth, head, width, temperature, precision,
  and parameter record & \cref{def-complete-attention-record} &
  \cref{thm-attention-structural-source-compilation} \\
\(P_s\) & Represented masked Gibbs attention row &
  \cref{def-complete-attention-record} &
  \cref{thm-attention-gibbs-gradient-flow} \\
\(U_{R,t}^{\rm att}\) & Source-resolved attention target-gain envelope &
  \cref{def-source-resolved-attention-ledger} &
  \cref{thm-attention-target-gain-envelope} \\
\addlinespace[5pt]
\multicolumn{4}{@{}l}{\emph{LPN specialization}}\\
\addlinespace[2pt]
\((A,b,\eta)\) & Public binary LPN instance and noise rate &
  \cref{def-lpn-task-object} & \cref{part:lpn} \\
\(s\) & Latent secret and singleton target &
  \cref{def-lpn-task-object} & \cref{def-lpn-task-object} \\
\(r_I(x)=Ax+b\) & Public residual map of candidate \(x\) &
  \cref{def-lpn-task-object} & \cref{def-lpn-task-object} \\
\(\rho_\eta=1-2\eta\) & Binary Fourier attenuation at fixed noise &
  \cref{thm-lpn-complete-readout-coset-energy-envelope} &
  \cref{eq-lpn-complete-readout-coset-energy-pointwise} \\
\(\begin{gathered}\Gamma(\lambda)=A^{\mathsf T}\lambda,\\
\operatorname{Prov}(\lambda),\ w(\lambda)\end{gathered}\) &
  Semantic label, noisy-row provenance support, and provenance weight of a
  charged LPN combination &
  \cref{def-charged-binary-noisy-combination-record} &
  \cref{thm-lpn-rate-uniform-semantic-cancellation} \\
\(\mathcal E_{q_0,k}\) & Short-provenance combination with nonzero semantic
  support contained in at most \(k\) secret coordinates &
  \cref{eq-lpn-short-provenance-small-semantic-event} &
  \cref{eq-lpn-short-provenance-small-semantic-probability} \\
\(q=\lfloor\delta m\rfloor,\ L_B(n)\) & Short aligned-coset cutoff and
  charged ceiling on explicitly materialized spectral indices &
  \cref{thm-lpn-constant-noise-materialized-spectral-closure} &
  \cref{eq-lpn-constant-noise-materialized-family-bound} \\
\(\begin{gathered}q_{L,I},M_{L,A},\\c_{L,b}\end{gathered}\) &
  Residual-parity fiber map and its affine data &
  \cref{eq-lpn-residual-parity-spectral-map} &
  \cref{eq-lpn-residual-parity-target-fiber} \\
\(G_{X,b,\mathfrak M,n}^{\mathrm{src,sat}}(B)\) & Compact,
  dynamically qualified residual-ancestry ambient source coordinate &
  \cref{def-lpn-ambient-geom-ancestry-branches} &
  \cref{thm-lpn-exact-constant-noise-spectral-perimeter} \\
\(\begin{gathered}J_{I,t}^{\rm cur},\beta_{I,t}^{\rm cur},\\
B_{\rm cur}^{\rm LPN}\end{gathered}\) & Conditional within-cell target
  information, decoupled contact, and complete sparse-reward observation
  cost &
  \cref{eq-lpn-curiosity-within-cell-information,eq-lpn-curiosity-decoupled-contact,eq-lpn-complete-curiosity-cost} &
  \cref{thm-lpn-complete-sparse-reward-observation-accounting} \\
\(\begin{gathered}B_{\rm LPN}^{\rm act},q_B,\\
J_{n,m,\eta,q_B,t}^{\rm act},\\
B_{{\rm act},{\rm LPN}}^*\end{gathered}\) & LPN active-sampling ceiling,
  budgeted request count, parameter-explicit information budget, and optimal
  active recovery cost &
  \cref{def-lpn-master-active-sampling-ceiling,eq-lpn-budgeted-active-sample-count,eq-lpn-active-total-information-n-m-eta} &
  \hyperref[thm-lpn-complete-master-active-sampling-envelope]{master envelope},
  \hyperref[thm-lpn-optimal-active-construction-equivalence]{optimality} \\
\(V_I(x)\) & Public threshold verifier &
  \cref{thm-lpn-public-verifier} & \cref{the-public-residual-landscape} \\
\(N_m(k)\)\newline\(F_{m,\eta}(k)\)\newline\(k_B\) & Hamming-ball count, its binomial mass, and
  the exact affordable enumeration radius &
  \cref{eq-binary-hamming-ball-count,eq-lpn-error-ball-mass,%
eq-lpn-affordable-error-radius} &
  \hyperref[prop-lpn-error-enumeration-success]{error-enumeration result} \\
\(\eta_{\mathrm{EE}},\eta_{\mathrm{Rad}},\eta_{\mathrm{IT}}\) &
  Enumeration, statistical-resolution, and information endpoints &
  \cref{eq-lpn-exact-enumeration-noise-boundary,eq-lpn-explicit-rademacher-eta-threshold,eq-lpn-asymptotic-information-threshold} &
  \cref{cor-lpn-polynomial-enumeration-window} \\
\(\eta_{\mathrm{IT}}^{\mathrm{ach}}\)\newline
\(\eta_{\mathrm{IT}}^{\mathrm{con}}\) & Finite-sample information
  achievability and converse endpoints &
  \cref{eq-lpn-finite-information-achievability-endpoint,eq-lpn-finite-information-converse-endpoint} &
  \cref{thm-lpn-information-threshold} \\
\(\eta_{\mathrm X}(R)\) & Error-enumeration versus exhaustive-secret-search
  crossover & \cref{eq-lpn-exponential-attack-switch} &
  \cref{cor-lpn-rate-asymptotic-thresholds} \\
\(I_{\mathrm{LPN},n}^{\mathrm{sat}}\) & Full saturated LPN extraction profile &
  \cref{eq-lpn-saturated-extraction-functional} & \cref{part:lpn} \\
\(I_{\mathrm{LPN},n}^{\mathrm{src,sat}}\) & Saturated prospective LPN source
  profile &
  \cref{def-extraction-functional} &
  \cref{thm-lpn-saturated-floor} \\
\bottomrule
\end{longtable}

\subsection{Superscript and subscript conventions}
\label{subsec-index-superscript-subscript-conventions}

\begin{tabular}{@{}p{0.25\textwidth}p{0.65\textwidth}@{}}
\toprule
Decoration & Reading \\
\midrule
\(\cdot^{\mathrm{sat}}\) & Resource-saturated. \\
\(\cdot^{\mathrm{auth}}\) & Relative to the declared causal-authority envelope. \\
\(\cdot^{\mathrm{amb}}\) & Ambient, before quotient compression. \\
\(\cdot^{\mathrm{nf}}\) & Expanded derivation normal form. \\
\(\cdot^{\mathrm{env}}\) & Resource-filtered envelope. \\
\(\cdot^{\mathrm{vis}}\) & Visibility contribution. \\
\(\cdot^{\mathrm{imp}}\) & Imposed dynamics. \\
\(\cdot^{\mathrm{ctrl}}\) & Affordable controlled dynamics. \\
\(\cdot^{\mathrm{rev}}\) & Reversible compensation record. \\
\(\cdot^{\mathrm{nr}}\) & General nonreversible resolvent record. \\
\(\cdot^{\mathrm{ql}}\) & Quasi-lumpable compensated quotient family. \\
\bottomrule
\end{tabular}

\begingroup
\raggedright
\bibliographystyle{plainnat}
\bibliography{task_space_structural_complexity}
\endgroup

\end{document}
